% Auto-generated from data/segments.json + layout.json (+ transforms) — do not edit by hand.
% volume: 40Abhi12
% mode: printing
% segments: 5179
% transform_rules: 8
% toc_marks: 16

% Volume layout defaults (layout.json ``layout``).
\csromanlayoutapply{1}{1.5}{21.6pt}{8pt}{8pt}{65pt}{2.5em}%

\pitaka{อภิธมฺม\-ปิฏก}
\csromanheader{2}{\textbf{ธมฺมานุโลม}}
\csromanmatikaanchor{mtk.0}
\csromantocmarknopage{chapter}{ปฏฺฐานปาฬิ}{mtk.0}
\bookmark[dest={mtk.0},level=0]{ปฏฺฐานปาฬิ}
\gambhira{ติกติก\-ปฏฺฐาน\-ปาฬิ}
\namakkaram{นโม ตสฺส ภควโต อรหโต สมฺมา\-สมฺพุทฺธสฺส.}
\csromanmarkh{1. กุสลตฺติก}\csromanmarkhii{1. เวทนาตฺติก}\csromanheadernomark{2}{1. \textbf{กุสลตฺติก 1. เวทนาตฺติก}}
\prose{\csromanfolio{1}1. สุขาย\-เวทนาย\-สมฺปยุตฺตปท 1-7. ปฏิจฺจาทิวาร}

\prose{ปจฺจยจตุกฺกเหตุ}

\proseitem{1}{กุสลํ สุขาย เวทนาย สมฺปยุตฺตํ ธมฺมํ ปฏิจฺจ กุสโล สุขาย เวทนาย สมฺปยุตฺโต ธมฺโม อุปฺปชฺชติ เหตุปจฺจยา. (1)}

\prose{อกุสลํ สุขาย เวทนาย สมฺปยุตฺตํ ธมฺมํ ปฏิจฺจ อกุสโล สุขาย เวทนาย สมฺปยุตฺโต ธมฺโม อุปฺปชฺชติ เหตุปจฺจยา. (1)}

\prose{อพฺยากตํ สุขาย เวทนาย สมฺปยุตฺตํ ธมฺมํ ปฏิจฺจ อพฺยากโต สุขาย เวทนาย สมฺปยุตฺโต ธมฺโม อุปฺปชฺชติ เหตุปจฺจยา. (1)}

\proseitem{2}{เหตุยา ตีณิ, อารมฺมเณ ตีณิ ฯเปฯ อวิคเต ตีณิ.}

\vspace{\tipitakaparskip}
\csromancenter{(สํขิตฺตํ. ฯเปฯ สหชาตวารมฺปิ ฯเปฯ สมฺปยุตฺตวารมฺปิ ปฏิจฺจ\-วาร\-สทิสํ.)}
\proseitem{3}{กุสโล สุขาย เวทนาย สมฺปยุตฺโต ธมฺโม กุสลสฺส สุขาย เวทนาย สมฺปยุตฺตสฺส ธมฺมสฺส เหตุปจฺจเยน ปจฺจโย. (1)}

\prose{อกุสโล สุขาย เวทนาย สมฺปยุตฺโต ธมฺโม อกุสลสฺส สุขาย เวทนาย สมฺปยุตฺตสฺส ธมฺมสฺส เหตุปจฺจเยน ปจฺจโย. (1)}

\prose{\csromanfolio{2}อพฺยากโต สุขาย เวทนาย สมฺปยุตฺโต ธมฺโม อพฺยากตสฺส สุขาย เวทนาย สมฺปยุตฺตสฺส ธมฺมสฺส เหตุปจฺจเยน ปจฺจโย. (1)}

\prose{กุสโล สุขาย เวทนาย สมฺปยุตฺโต ธมฺโม กุสลสฺส สุขาย เวทนาย สมฺปยุตฺตสฺส ธมฺมสฺส อารมฺมณ\-ปจฺจเยน ปจฺจโย.}

\proseitem{4}{เหตุยา ตีณิ, อารมฺมเณ นว. (สํขิตฺตํ.) (ยถา กุสลตฺติเก ปญฺหาวารํ, เอวํ วิตฺถาเรตพฺพํ.)}

\prose{2. ทุกฺขาย\-เวทนาย\-สมฺปยุตฺตปท 1-7. ปฏิจฺจาทิวาร}

\prose{ปจฺจยจตุกฺกเหตุอารมฺมณ}

\proseitem{5}{อกุสลํ ทุกฺขาย เวทนาย สมฺปยุตฺตํ ธมฺมํ ปฏิจฺจ อกุสโล ทุกฺขาย เวทนาย สมฺปยุตฺโต ธมฺโม อุปฺปชฺชติ เหตุปจฺจยา. (1)}

\prose{อกุสลํ ทุกฺขาย เวทนาย สมฺปยุตฺตํ ธมฺมํ ปฏิจฺจ อกุสโล ทุกฺขาย เวทนาย สมฺปยุตฺโต ธมฺโม อุปฺปชฺชติ อารมฺมณ\-ปจฺจยา. (1)}

\prose{อพฺยากตํ ทุกฺขาย เวทนาย สมฺปยุตฺตํ ธมฺมํ ปฏิจฺจ อพฺยากโต ทุกฺขาย เวทนาย สมฺปยุตฺโต ธมฺโม อุปฺปชฺชติ อารมฺมณ\-ปจฺจยา. (1) (สํขิตฺตํ.)}

\proseitem{6}{เหตุยา เอกํ, อารมฺมเณ เทฺว ฯเปฯ อวิคเต เทฺว. (สํขิตฺตํ.) (สหชาตวารมฺปิ ฯเปฯ สมฺปยุตฺตวารมฺปิ ปฏิจฺจ\-วาร\-สทิสํ.)}

\proseitem{7}{อกุสโล ทุกฺขาย เวทนาย สมฺปยุตฺโต ธมฺโม อกุสลสฺส ทุกฺขาย เวทนาย สมฺปยุตฺตสฺส ธมฺมสฺส เหตุปจฺจเยน ปจฺจโย. (1)}

\prose{อกุสโล ทุกฺขาย เวทนาย สมฺปยุตฺโต ธมฺโม อกุสลสฺส ทุกฺขาย เวทนาย สมฺปยุตฺตสฺส ธมฺมสฺส อารมฺมณ\-ปจฺจเยน ปจฺจโย. (1)}

\prose{อพฺยากโต ทุกฺขาย เวทนาย สมฺปยุตฺโต ธมฺโม อกุสลสฺส ทุกฺขาย เวทนาย สมฺปยุตฺตสฺส ธมฺมสฺส อารมฺมณ\-ปจฺจเยน ปจฺจโย. (1) (สํขิตฺตํ.)}

\csromanmarkh{1. กุสลตฺติก}\csromanmarkhii{1. เวทนาตฺติก}\csromanheadernomark{2}{1. กุสลตฺติก 1. เวทนาตฺติก}
\proseitem{8}{\csromanfolio{3}เหตุยา เอกํ, อารมฺมเณ เทฺว. (สํขิตฺตํ. ยถา กุสลตฺติเก ปญฺหาวารํ, เอวํ วิตฺถาเรตพฺพํ.)}

\prose{3. อทุกฺขมสุขเวทนาย\-สมฺปยุตฺตปท 1-7. ปฏิจฺจาทิวาร}

\prose{ปจฺจยจตุกฺกเหตุ}

\proseitem{9}{กุสลํ อทุกฺขม\-สุขาย เวทนาย สมฺปยุตฺตํ ธมฺมํ ปฏิจฺจ กุสโล อทุกฺขม\-สุขาย เวทนาย สมฺปยุตฺโต ธมฺโม อุปฺปชฺชติ}

\prose{อกุสลํ อทุกฺขม\-สุขาย เวทนาย สมฺปยุตฺตํ ธมฺมํ ปฏิจฺจ อกุสโล อทุกฺขม\-สุขาย เวทนาย สมฺปยุตฺโต ธมฺโม อุปฺปชฺชติ}

\prose{อพฺยากตํ อทุกฺขม\-สุขาย เวทนาย สมฺปยุตฺตํ ธมฺมํ ปฏิจฺจ อพฺยากโต อทุกฺขม\-สุขาย เวทนาย สมฺปยุตฺโต ธมฺโม อุปฺปชฺชติ เหตุปจฺจยา. (1)}

\proseitem{10}{เหตุยา ตีณิ, อารมฺมเณ ตีณิ ฯเปฯ อวิคเต ตีณิ. (สํขิตฺตํ. สหชาตวารมฺปิ ฯเปฯ สมฺปยุตฺตวารมฺปิ ปฏิจฺจ\-วาร\-สทิสํ.)}

\proseitem{11}{กุสโล อทุกฺขม\-สุขาย เวทนาย สมฺปยุตฺโต ธมฺโม กุสลสฺส อทุกฺขม\-สุขาย เวทนาย สมฺปยุตฺตสฺส ธมฺมสฺส}

\prose{อกุสโล อทุกฺขม\-สุขาย เวทนาย สมฺปยุตฺโต ธมฺโม อกุสลสฺส อทุกฺขม\-สุขาย เวทนาย สมฺปยุตฺตสฺส ธมฺมสฺส เหตุปจฺจเยน ปจฺจโย. (1)}

\prose{อพฺยากโต อทุกฺขม\-สุขาย เวทนาย สมฺปยุตฺโต ธมฺโม อพฺยากตสฺส อทุกฺขม\-สุขาย เวทนาย สมฺปยุตฺตสฺส ธมฺมสฺส เหตุปจฺจเยน ปจฺจโย. (1) (สํขิตฺตํ.)}

\proseitem{12}{เหตุยา ตีณิ, อารมฺมเณ นว, อธิปติยา สตฺต, อนนฺตเร สตฺต ฯเปฯ อุปนิสฺสเย นว, อวิคเต ตีณิ. (สํขิตฺตํ. ยถา กุสลตฺติเก ปญฺหาวารํ, เอวํ วิตฺถาเรตพฺพํ.)}

\csromanmarkh{1. กุสลตฺติก}\csromanmarkhii{2. วิปากตฺติก}\csromanheadernomark{2}{1. \textbf{กุสลตฺติก 2. วิปากตฺติก}}
\proseitem{13}{\csromanfolio{4}อพฺยากตํ วิปากํ ธมฺมํ ปฏิจฺจ อพฺยากโต วิปาโก ธมฺโม อุปฺปชฺชติ เหตุปจฺจยา. (สํขิตฺตํ.)}

\prose{เหตุยา เอกํ, อารมฺมเณ เอกํ ฯเปฯ อวิคเต เอกํ. (สํขิตฺตํ.) (สหชาตวาเรปิ ฯเปฯ ปญฺหาวาเรปิ สพฺพตฺถ เอกํ.)}

\proseitem{14}{กุสลํ วิปาก\-ธมฺม\-ธมฺมํ ปฏิจฺจ กุสโล วิปาก\-ธมฺม\-ธมฺโม อุปฺปชฺชติ เหตุปจฺจยา. (1)}

\prose{อกุสลํ วิปาก\-ธมฺม\-ธมฺมํ ปฏิจฺจ อกุสโล วิปาก\-ธมฺม\-ธมฺโม อุปฺปชฺชติ เหตุปจฺจยา. (1) (สํขิตฺตํ.)}

\prose{เหตุยา เทฺว, อารมฺมเณ เทฺว ฯเปฯ อวิคเต เทฺว. (สํขิตฺตํ.)}

\prose{(สหชาตวารมฺปิ ฯเปฯ สมฺปยุตฺตวารมฺปิ ปฏิจฺจ\-วาร\-สทิสํ.)}

\proseitem{15}{กุสโล วิปาก\-ธมฺม\-ธมฺโม กุสลสฺส วิปาก\-ธมฺม\-ธมฺมสฺส เหตุปจฺจเยน ปจฺจโย. (1)}

\prose{อกุสโล วิปาก\-ธมฺม\-ธมฺโม อกุสลสฺส วิปาก\-ธมฺม\-ธมฺมสฺส}

\prose{กุสโล วิปาก\-ธมฺม\-ธมฺโม กุสลสฺส วิปาก\-ธมฺม\-ธมฺมสฺส อารมฺมณ\-ปจฺจเยน ปจฺจโย.  กุสโล วิปาก\-ธมฺม\-ธมฺโม อกุสลสฺส วิปาก\-ธมฺม\-ธมฺมสฺส อารมฺมณ\-ปจฺจเยน ปจฺจโย. (2)}

\prose{อกุสโล วิปาก\-ธมฺม\-ธมฺโม อกุสลสฺส วิปาก\-ธมฺม\-ธมฺมสฺส อารมฺมณ\-ปจฺจเยน ปจฺจโย.  อกุสโล วิปาก\-ธมฺม\-ธมฺโม กุสลสฺส วิปาก\-ธมฺม\-ธมฺมสฺส อารมฺมณ\-ปจฺจเยน ปจฺจโย. (2) (สํขิตฺตํ.)}

\proseitem{16}{เหตุยา เทฺว, อารมฺมเณ จตฺตาริ, อธิปติยา ตีณิ, อนนฺตเร เทฺว ฯเปฯ สหชาเต เทฺว, อุปนิสฺสเย จตฺตาริ ฯเปฯ อวิคเต เทฺว. (สํขิตฺตํ. ยถา กุสลตฺติเก ปญฺหาวารํ, เอวํ วิตฺถาเรตพฺพํ.)}

\csromanmarkh{1. กุสลตฺติก}\csromanmarkhii{3. อุปาทินฺนตฺติก}\csromanheadernomark{2}{1. กุสลตฺติก 3. อุปาทินฺนตฺติก}
\proseitem{17}{\csromanfolio{5}อพฺยากตํ เนว\-วิปาก\-น\-วิปากธมฺม\-ธมฺมํ ปฏิจฺจ อพฺยากโต เนว วิปาก\-น\-วิปาก\-ธมฺม\-ธมฺโม อุปฺปชฺชติ เหตุปจฺจยา. (สํขิตฺตํ.)}

\prose{เหตุยา เอกํ, อารมฺมเณ เอกํ ฯเปฯ อวิคเต เอกํ. (สํขิตฺตํ.)}

\vspace{\tipitakaparskip}
\csromancenter{(สหชาตวาเรปิ ฯเปฯ ปญฺหาวาเรปิ สพฺพตฺถ เอกํ.)}
\csromanmarkh{1. กุสลตฺติก}\csromanmarkhii{3. อุปาทินฺนตฺติก}\csromanheadernomark{2}{1. กุสลตฺติก 3. อุปาทินฺนตฺติก}
\proseitem{18}{อพฺยากตํ อุปาทินฺนุ\-ปาทานิยํ ธมฺมํ ปฏิจฺจ อพฺยากโต อุปาทินฺนุ\-ปาทานิโย ธมฺโม อุปฺปชฺชติ เหตุปจฺจยา. (สํขิตฺตํ.)}

\prose{เหตุยา เอกํ ฯเปฯ อวิคเต เอกํ. (สํขิตฺตํ.)}

\vspace{\tipitakaparskip}
\csromancenter{(สหชาตวาเรปิ ฯเปฯ ปญฺหาวาเรปิ สพฺพตฺถ เอกํ.)}
\proseitem{19}{กุสลํ อนุปาทินฺนุ\-ปาทานิยํ ธมฺมํ ปฏิจฺจ กุสโล อนุปาทินฺนุ\-ปาทานิโย ธมฺโม อุปฺปชฺชติ เหตุปจฺจยา. ตีณิ.}

\prose{อกุสลํ อนุปาทินฺนุ\-ปาทานิยํ ธมฺมํ ปฏิจฺจ อกุสโล อนุปาทินฺนุ\-ปาทานิโย ธมฺโม อุปฺปชฺชติ เหตุปจฺจยา. ตีณิ.}

\prose{อพฺยากตํ อนุปาทินฺนุ\-ปาทานิยํ ธมฺมํ ปฏิจฺจ อพฺยากโต อนุปาทินฺนุ\-ปาทานิโย ธมฺโม อุปฺปชฺชติ เหตุปจฺจยา. (1)}

\prose{กุสลํ อนุปาทินฺนุ\-ปาทานิยญฺจ อพฺยากตํ อนุปาทินฺนุ\-ปาทานิยญฺจ ธมฺมํ ปฏิจฺจ อพฺยากโต อนุปาทินฺนุ\-ปาทานิโย ธมฺโม อุปฺปชฺชติ}

\prose{อกุสลํ อนุปาทินฺนุ\-ปาทานิยญฺจ อพฺยากตํ อนุปาทินฺนุ\-ปาทานิยญฺจ ธมฺมํ ปฏิจฺจ อพฺยากโต อนุปาทินฺนุ\-ปาทานิโย ธมฺโม อุปฺปชฺชติ เหตุปจฺจยา. (1) (สํขิตฺตํ.)}

\proseitem{20}{เหตุยา นว, อวิคเต นว. (สํขิตฺตํ. สหชาตวารมฺปิ ฯเปฯ สมฺปยุตฺตวารมฺปิ ปฏิจฺจ\-วาร\-สทิสํ.)}

\proseitem{21}{\csromanfolio{6}กุสโล อนุปาทินฺนุ\-ปาทานิโย ธมฺโม กุสลสฺส อนุปาทินฺนุ\-ปาทานิยสฺส ธมฺมสฺส เหตุปจฺจเยน ปจฺจโย. ตีณิ.}

\prose{อกุสโล อนุปาทินฺนุ\-ปาทานิโย ธมฺโม อกุสลสฺส อนุปาทินฺนุ\-ปาทานิยสฺส ธมฺมสฺส เหตุปจฺจเยน ปจฺจโย. ตีณิ.}

\prose{อพฺยากโต อนุปาทินฺนุ\-ปาทานิโย ธมฺโม อพฺยากตสฺส อนุปาทินฺนุ\-ปาทานิยสฺส ธมฺมสฺส เหตุปจฺจเยน ปจฺจโย. (1)}

\proseitem{22}{กุสโล อนุปาทินฺนุ\-ปาทานิโย ธมฺโม กุสลสฺส อนุปาทินฺนุ\-ปาทานิยสฺส ธมฺมสฺส อารมฺมณ\-ปจฺจเยน ปจฺจโย. ตีณิ.}

\prose{อกุสโล อนุปาทินฺนุ\-ปาทานิโย ธมฺโม อกุสลสฺส อนุปาทินฺนุ\-ปาทานิยสฺส ธมฺมสฺส อารมฺมณ\-ปจฺจเยน ปจฺจโย. ตีณิ.}

\prose{อพฺยากโต อนุปาทินฺนุ\-ปาทานิโย ธมฺโม อพฺยากตสฺส อนุปาทินฺนุ\-ปาทานิยสฺส ธมฺมสฺส อารมฺมณ\-ปจฺจเยน ปจฺจโย. ตีณิ.}

\proseitem{23}{เหตุยา สตฺต, อารมฺมเณ นว, อวิคเต เอกาทส. (สํขิตฺตํ.)}

\prose{(ยถา กุสลตฺติเก ปญฺหาวารํ, เอวํ วิตฺถาเรตพฺพํ.)}

\proseitem{24}{กุสลํ อนุปาทินฺน\-อนุปาทานิยํ ธมฺมํ ปฏิจฺจ กุสโล อนุปาทินฺนุ\-ปาทานิโย ธมฺโม อุปฺปชฺชติ เหตุปจฺจยา. (1)}

\prose{อพฺยากตํ อนุปาทินฺน\-อนุปาทานิยํ ธมฺมํ ปฏิจฺจ อพฺยากโต อนุปาทินฺน\-อนุปาทานิโย ธมฺโม อุปฺปชฺชติ เหตุปจฺจยา. (1) (สํขิตฺตํ.)}

\prose{เหตุยา เทฺว, อวิคเต เทฺว. (สํขิตฺตํ. สหชาตวารมฺปิ ฯเปฯ ปญฺหาวารมฺปิ สพฺพตฺถ วิตฺถาโร.)}

\csromanmarkh{1. กุสลตฺติก}\csromanmarkhii{4. สํกิลิฏฺฐ\-ตฺติก}\csromanheadernomark{2}{1. กุสลตฺติก 4. สํกิลิฏฺฐ\-ตฺติก}
\proseitem{25}{อกุสลํ สํกิลิฏฺฐ\-สํกิเลสิกํ ธมฺมํ ปฏิจฺจ อกุสโล สํกิลิฏฺฐ\-สํกิเลสิโก ธมฺโม อุปฺปชฺชติ เหตุปจฺจยา. (สํขิตฺตํ.)}

\csromanmarkh{1. กุสลตฺติก}\csromanmarkhii{5. วิตกฺกตฺติก}\csromanheadernomark{2}{1. กุสลตฺติก 5. วิตกฺกตฺติก}
\prose{\csromanfolio{7}เหตุยา เอกํ ฯเปฯ อวิคเต เอกํ. (สํขิตฺตํ.)}

\prose{(สหชาตวาเรปิ ฯเปฯ ปญฺหาวาเรปิ สพฺพตฺถ เอกํ.)}

\proseitem{26}{กุสลํ อสํกิลิฏฺฐ\-สํกิเลสิกํ ธมฺมํ ปฏิจฺจ กุสโล อสํกิลิฏฺฐ\-สํกิเลสิโก ธมฺโม อุปฺปชฺชติ เหตุปจฺจยา. ตีณิ.}

\prose{อพฺยากตํ อสํกิลิฏฺฐ\-สํกิเลสิกํ ธมฺมํ ปฏิจฺจ อพฺยากโต อสํกิลิฏฺฐ\-สํกิเลสิโก ธมฺโม อุปฺปชฺชติ เหตุปจฺจยา. (1)}

\prose{กุสลํ อสํกิลิฏฺฐ\-สํกิเลสิกญฺจ อพฺยากตํ อสํกิลิฏฺฐ\-สํกิเลสิกญฺจ ธมฺมํ ปฏิจฺจ อพฺยากโต อสํกิลิฏฺฐ\-สํกิเลสิโก ธมฺโม อุปฺปชฺชติ เหตุปจฺจยา. (1)}

\proseitem{27}{เหตุยา ปญฺจ, อวิคเต ปญฺจ. (สํขิตฺตํ. สหชาตวาเรปิ ฯเปฯ ปญฺหาวาเรปิ สพฺพตฺถ วิตฺถาโร.)}

\proseitem{28}{กุสลํ อสํกิลิฏฺฐ\-อสํกิเลสิกํ ธมฺมํ ปฏิจฺจ กุสโล อสํกิลิฏฺฐ\-อสํกิเลสิโก ธมฺโม อุปฺปชฺชติ เหตุปจฺจยา. (1)}

\prose{อพฺยากตํ อสํกิลิฏฺฐ\-อสํกิเลสิกํ ธมฺมํ ปฏิจฺจ อพฺยากโต อสํกิลิฏฺฐ\-อสํกิเลสิโก ธมฺโม อุปฺปชฺชติ เหตุปจฺจยา. (1) (สํขิตฺตํ.)}

\prose{เหตุยา เทฺว, อวิคเต เทฺว. (สํขิตฺตํ. สหชาตวาเรปิ ฯเปฯ ปญฺหาวาเรปิ สพฺพตฺถ วิตฺถาเรตพฺพํ.)}

\csromanmarkh{1. กุสลตฺติก}\csromanmarkhii{5. วิตกฺกตฺติก}\csromanheadernomark{2}{1. กุสลตฺติก 5. วิตกฺกตฺติก}
\csromanheader{2}{1-7. ปฏิจฺจาทิวาร ปจฺจยจตุกฺกเหตุ}
\proseitem{29}{กุสลํ สวิตกฺก\-สวิจารํ ธมฺมํ ปฏิจฺจ กุสโล สวิตกฺก\-สวิจาโร ธมฺโม อุปฺปชฺชติ เหตุปจฺจยา. (1)}

\prose{อกุสลํ สวิตกฺก\-สวิจารํ ธมฺมํ ปฏิจฺจ อกุสโล สวิตกฺก\-สวิจาโร ธมฺโม อุปฺปชฺชติ เหตุปจฺจยา. (1)}

\prose{\csromanfolio{8}อพฺยากตํ สวิตกฺก\-สวิจารํ ธมฺมํ ปฏิจฺจ อพฺยากโต สวิตกฺก\-สวิจาโร ธมฺโม อุปฺปชฺชติ เหตุปจฺจยา. (1) (สํขิตฺตํ.)}

\prose{เหตุยา ตีณิ, อารมฺมเณ ตีณิ, อวิคเต ตีณิ. (สํขิตฺตํ. สหชาตวารมฺปิ ฯเปฯ สมฺปยุตฺตวารมฺปิ ปฏิจฺจ\-วาร\-สทิสํ.)}

\proseitem{30}{กุสโล สวิตกฺก\-สวิจาโร ธมฺโม กุสลสฺส สวิตกฺก\-สวิจารสฺส ธมฺมสฺส เหตุปจฺจเยน ปจฺจโย. (1)}

\prose{อกุสโล สวิตกฺก\-สวิจาโร ธมฺโม อกุสลสฺส สวิตกฺก\-สวิจารสฺส ธมฺมสฺส เหตุปจฺจเยน ปจฺจโย. (1)}

\prose{อพฺยากโต สวิตกฺก\-สวิจาโร ธมฺโม อพฺยากตสฺส สวิตกฺก\-สวิจารสฺส ธมฺมสฺส เหตุปจฺจเยน ปจฺจโย. (1) (สํขิตฺตํ.)}

\proseitem{31}{เหตุยา ตีณิ, อารมฺมเณ นว, อวิคเต ตีณิ. (สํขิตฺตํ.) (ยถา กุสลตฺติเก ปญฺหาวารํ, เอวํ วิตฺถาเรตพฺพํ.)}

\proseitem{32}{กุสลํ อวิตกฺก\-วิจารมตฺตํ ธมฺมํ ปฏิจฺจ กุสโล อวิตกฺก\-วิจารมตฺโต ธมฺโม อุปฺปชฺชติ เหตุปจฺจยา. (1)}

\prose{อพฺยากตํ อวิตกฺก\-วิจารมตฺตํ ธมฺมํ ปฏิจฺจ อพฺยากโต อวิตกฺก\-วิจารมตฺโต ธมฺโม อุปฺปชฺชติ เหตุปจฺจยา. (1) (สํขิตฺตํ.)}

\prose{เหตุยา เทฺว, อวิคเต เทฺว. (สํขิตฺตํ.)}

\prose{(สหชาตวาเรปิ ฯเปฯ สมฺปยุตฺตวาเรปิ ปญฺหาวาเรปิ สพฺพตฺถ วิตฺถาโร.)}

\proseitem{33}{กุสลํ อวิตกฺก\-อวิจารํ ธมฺมํ ปฏิจฺจ กุสโล อวิตกฺก\-อวิจาโร ธมฺโม อุปฺปชฺชติ เหตุปจฺจยา. ตีณิ.}

\prose{อพฺยากตํ อวิตกฺก\-อวิจารํ ธมฺมํ ปฏิจฺจ อพฺยากโต อวิตกฺก\-อวิจาโร ธมฺโม อุปฺปชฺชติ เหตุปจฺจยา. (1)}

\prose{กุสลํ อวิตกฺกอวิจารญฺจ อพฺยากตํ อวิตกฺกอวิจารญฺจ ธมฺมํ ปฏิจฺจ อพฺยากโต อวิตกฺก\-อวิจาโร ธมฺโม อุปฺปชฺชติ เหตุปจฺจยา. (1) (สํขิตฺตํ.)}

\csromanmarkh{1. กุสลตฺติก}\csromanmarkhii{6. ปีติตฺติก}\csromanheadernomark{2}{1. กุสลตฺติก 6. ปีติตฺติก}
\prose{\csromanfolio{9}เหตุยา ปญฺจ, อารมฺมเณ เทฺว, อวิคเต ปญฺจ. (สํขิตฺตํ. สหชาตวารมฺปิ ฯเปฯ ปญฺหาวารมฺปิ วิตฺถาเรตพฺพํ.)}

\csromanmarkh{1. กุสลตฺติก}\csromanmarkhii{6. ปีติตฺติก}\csromanheadernomark{2}{1. กุสลตฺติก 6. ปีติตฺติก}
\csromanheader{2}{1-7. ปฏิจฺจาทิวาร ปจฺจยจตุกฺกเหตุ}
\prose{34. กุสลํ ปีติสหคตํ ธมฺมํ ปฏิจฺจ กุสโล ปีติสหคโต}

\prose{อกุสลํ ปีติสหคตํ ธมฺมํ ปฏิจฺจ อกุสโล ปีติสหคโต}

\prose{อพฺยากตํ ปีติสหคตํ ธมฺมํ ปฏิจฺจ อพฺยากโต ปีติสหคโต ธมฺโม อุปฺปชฺชติ เหตุปจฺจยา. (1) (สํขิตฺตํ.)}

\prose{เหตุยา ตีณิ, อารมฺมเณ ตีณิ, อวิคเต ตีณิ. (สํขิตฺตํ. สหชาตวารมฺปิ ฯเปฯ สมฺปยุตฺตวารมฺปิ ปฏิจฺจ\-วาร\-สทิสํ.)}

\proseitem{35}{กุสโล ปีติสหคโต ธมฺโม กุสลสฺส ปีติสหคตสฺส ธมฺมสฺส}

\prose{อกุสโล ปีติสหคโต ธมฺโม อกุสลสฺส ปีติสหคตสฺส ธมฺมสฺส}

\prose{อพฺยากโต ปีติสหคโต ธมฺโม อพฺยากตสฺส ปีติสหคตสฺส ธมฺมสฺส เหตุปจฺจเยน ปจฺจโย. (1) (สํขิตฺตํ.)}

\prose{เหตุยา ตีณิ, อารมฺมเณ นว, อธิปติยา สตฺต, อนนฺตเร ปญฺจ, อวิคเต ตีณิ. (สํขิตฺตํ.) (ยถา กุสลตฺติเก ปญฺหาวารํ, เอวํ วิตฺถาเรตพฺพํ.)}

\prose{36. กุสลํ สุขสหคตํ ธมฺมํ ปฏิจฺจ กุสโล สุขสหคโต}

\prose{อกุสลํ สุขสหคตํ ธมฺมํ ปฏิจฺจ อกุสโล สุขสหคโต}

\prose{\csromanfolio{10}อพฺยากตํ สุขสหคตํ ธมฺมํ ปฏิจฺจ อพฺยากโต สุขสหคโต ธมฺโม อุปฺปชฺชติ เหตุปจฺจยา. (1) (สํขิตฺตํ.)}

\prose{เหตุยา ตีณิ, อารมฺมเณ ตีณิ, อวิคเต ตีณิ. (สํขิตฺตํ. สหชาตวารมฺปิ ฯเปฯ ปญฺหาวารมฺปิ วิตฺถาเรตพฺพํ.)}

\proseitem{37}{กุสลํ อุเปกฺขา\-สหคตํ ธมฺมํ ปฏิจฺจ กุสโล อุเปกฺขา\-สหคโต ธมฺโม อุปฺปชฺชติ เหตุปจฺจยา. (1)}

\prose{อกุสลํ อุเปกฺขา\-สหคตํ ธมฺมํ ปฏิจฺจ อกุสโล อุเปกฺขา\-สหคโต ธมฺโม อุปฺปชฺชติ เหตุปจฺจยา. (1)}

\prose{อพฺยากตํ อุเปกฺขา\-สหคตํ ธมฺมํ ปฏิจฺจ อพฺยากโต อุเปกฺขา\-สหคโต ธมฺโม อุปฺปชฺชติ เหตุปจฺจยา. (1) (สํขิตฺตํ.)}

\prose{เหตุยา ตีณิ, อวิคเต ตีณิ. (สํขิตฺตํ.)}

\vspace{\tipitakaparskip}
\csromancenter{(สหชาตวารมฺปิ ฯเปฯ ปญฺหาวารมฺปิ วิตฺถาเรตพฺพํ.)}
\csromanmarkh{1. กุสลตฺติก}\csromanmarkhii{7. ทสฺสนตฺติก}\csromanheadernomark{2}{1. \textbf{กุสลตฺติก 7. ทสฺสนตฺติก}}
\proseitem{38}{อกุสลํ ทสฺสเนน ปหาตพฺพํ ธมฺมํ ปฏิจฺจ อกุสโล ทสฺสเนน ปหาตพฺโพ ธมฺโม อุปฺปชฺชติ เหตุปจฺจยา. (สํขิตฺตํ.)}

\prose{เหตุยา เอกํ ฯเปฯ อวิคเต เอกํ. (สํขิตฺตํ.) (สหชาตวาเรปิ ฯเปฯ ปญฺหาวาเรปิ สพฺพตฺถ เอกํ.)}

\proseitem{39}{อกุสลํ ภาวนาย ปหาตพฺพํ ธมฺมํ ปฏิจฺจ อกุสโล ภาวนาย ปหาตพฺโพ ธมฺโม อุปฺปชฺชติ เหตุปจฺจยา. (1) (สํขิตฺตํ.)}

\prose{เหตุยา เอกํ ฯเปฯ อวิคเต เอกํ. (สํขิตฺตํ.)}

\proseitem{40}{กุสลํ เนวทสฺสเนน นภาวนาย ปหาตพฺพํ ธมฺมํ ปฏิจฺจ กุสโล เนวทสฺสเนน นภาวนาย ปหาตพฺโพ ธมฺโม อุปฺปชฺชติ เหตุปจฺจยา. ตีณิ.}

\csromanmarkh{1. กุสลตฺติก}\csromanmarkhii{8. ทสฺสน\-เหตุ\-กตฺติก}\csromanheadernomark{2}{1. กุสลตฺติก 8. ทสฺสน\-เหตุ\-กตฺติก}
\prose{\csromanfolio{11}อพฺยากตํ เนวทสฺสเนน นภาวนาย ปหาตพฺพํ ธมฺมํ ปฏิจฺจ อพฺยากโต เนวทสฺสเนน นภาวนาย ปหาตพฺโพ ธมฺโม อุปฺปชฺชติ}

\prose{กุสลํ เนวทสฺสเนน นภาวนาย ปหาตพฺพญฺจ อพฺยากตํ เนวทสฺสเนน นภาวนาย ปหาตพฺพญฺจ ธมฺมํ ปฏิจฺจ อพฺยากโต เนวทสฺสเนน นภาวนาย ปหาตพฺโพ ธมฺโม อุปฺปชฺชติ เหตุปจฺจยา. (1)}

\prose{เหตุยา ปญฺจ, อวิคเต ปญฺจ. (สํขิตฺตํ. สหชาตวาเรปิ ฯเปฯ ปญฺหาวาเรปิ สพฺพตฺถ วิตฺถาโร.)}

\csromanmarkh{1. กุสลตฺติก}\csromanmarkhii{8. ทสฺสน\-เหตุ\-กตฺติก}\csromanheadernomark{2}{1. กุสลตฺติก 8. ทสฺสน\-เหตุ\-กตฺติก}
\proseitem{41}{อกุสลํ ทสฺสเนน ปหาตพฺพ\-เหตุกํ ธมฺมํ ปฏิจฺจ อกุสโล ทสฺสเนน ปหาตพฺพ\-เหตุโก ธมฺโม อุปฺปชฺชติ เหตุปจฺจยา.}

\prose{เหตุยา เอกํ ฯเปฯ อวิคเต เอกํ. (สํขิตฺตํ.)}

\proseitem{42}{อกุสลํ ภาวนาย ปหาตพฺพ\-เหตุกํ ธมฺมํ ปฏิจฺจ อกุสโล ภาวนาย ปหาตพฺพ\-เหตุโก ธมฺโม อุปฺปชฺชติ เหตุปจฺจยา.}

\prose{เหตุยา เอกํ ฯเปฯ อวิคเต เอกํ. (สํขิตฺตํ.)}

\proseitem{43}{กุสลํ เนวทสฺสเนน นภาวนาย ปหาตพฺพ\-เหตุกํ ธมฺมํ ปฏิจฺจ กุสโล เนวทสฺสเนน นภาวนาย ปหาตพฺพ\-เหตุโก ธมฺโม อุปฺปชฺชติ เหตุปจฺจยา. (สํขิตฺตํ.)}

\prose{เหตุยา สตฺต, อารมฺมเณ เทฺว, อวิคเต สตฺต. (สํขิตฺตํ. สหชาตวาเรปิ ฯเปฯ ปญฺหาวาเรปิ สพฺพตฺถ วิตฺถาโร.)}

\csromanmarkh{1. กุสลตฺติก}\csromanmarkhii{9. อาจยคามิตฺติก}\csromanheadernomark{2}{1. \textbf{กุสลตฺติก 9. อาจยคามิตฺติก}}
\prose{\csromanfolio{12}44. กุสลํ อาจยคามิํ ธมฺมํ ปฏิจฺจ กุสโล อาจยคามี}

\prose{อกุสลํ อาจยคามิํ ธมฺม ปฏิจฺจ อกุสโล อาจยคามี ธมฺโม อุปฺปชฺชติ เหตุปจฺจยา. (1) (สํขิตฺตํ.)}

\prose{เหตุยา เทฺว ฯเปฯ อวิคเต เทฺว. (สํขิตฺตํ.) (สหชาตวาเรปิ ฯเปฯ ปญฺหาวาเรปิ สพฺพตฺถ วิตฺถาโร.)}

\proseitem{45}{กุสลํ อปจยคามิํ ธมฺมํ ปฏิจฺจ กุสโล อปจยคามี ธมฺโม อุปฺปชฺชติ เหตุปจฺจยา. (สํขิตฺตํ. ปฏิจฺจวาเรปิ ฯเปฯ ปญฺหาวาเรปิ สพฺพตฺถ เอกํ.)}

\proseitem{46}{อพฺยากตํ เนวา\-จยคามิ\-นา\-ปจยคามิํ ธมฺมํ ปฏิจฺจ อพฺยากโต เนวา\-จยคามิ\-นา\-ปจยคามี ธมฺโม อุปฺปชฺชติ เหตุปจฺจยา.}

\prose{เหตุยา เอกํ ฯเปฯ อวิคเต เอกํ. (สํขิตฺตํ.) (สหชาตวารมฺปิ ฯเปฯ ปญฺหาวารมฺปิ วิตฺถาเรตพฺพํ.)}

\csromanmarkh{1. กุสลตฺติก}\csromanmarkhii{10. เสกฺขตฺติก}\csromanheadernomark{2}{1. \textbf{กุสลตฺติก} 10. เสกฺขตฺติก}
\proseitem{47}{กุสลํ เสกฺขํ ธมฺมํ ปฏิจฺจ กุสโล เสกฺโข ธมฺโม อุปฺปชฺชติ เหตุปจฺจยา. (1)}

\prose{อพฺยากตํ เสกฺขํ ธมฺมํ ปฏิจฺจ อพฺยากโต เสกฺโข ธมฺโม อุปฺปชฺชติ เหตุปจฺจยา. (1) (สํขิตฺตํ.)}

\prose{เหตุยา เทฺว, อวิคเต เทฺว. (สํขิตฺตํ. สหชาตวาเรปิ ฯเปฯ ปญฺหาวาเรปิ สพฺพตฺถ วิตฺถาโร.)}

\csromanmarkh{1. กุสลตฺติก}\csromanmarkhii{11. ปริตฺตตฺติก}\csromanheadernomark{2}{1. กุสลตฺติก 11. ปริตฺตตฺติก}
\proseitem{48}{\csromanfolio{13}อพฺยากตํ อเสกฺขํ ธมฺมํ ปฏิจฺจ อพฺยากโต อเสกฺโข ธมฺโม อุปฺปชฺชติ เหตุปจฺจยา. (สํขิตฺตํ.)}

\prose{เหตุยา เอกํ ฯเปฯ อวิคเต เอกํ. (สํขิตฺตํ.) (สหชาตวาเรปิ ฯเปฯ ปญฺหาวาเรปิ สพฺพตฺถ เอกํ.)}

\proseitem{49}{กุสลํ เนว\-เสกฺข\-นาเสกฺขํ ธมฺมํ ปฏิจฺจ กุสโล เนว\-เสกฺข\-นา\-เสกฺโข ธมฺโม อุปฺปชฺชติ เหตุปจฺจยา. ตีณิ.}

\prose{อกุสลํ เนว\-เสกฺข\-นาเสกฺขํ ธมฺมํ ปฏิจฺจ อกุสโล เนว\-เสกฺข\-นา\-เสกฺโข ธมฺโม อุปฺปชฺชติ เหตุปจฺจยา. ตีณิ.}

\prose{อพฺยากตํ เนว\-เสกฺข\-นาเสกฺขํ ธมฺมํ ปฏิจฺจ อพฺยากโต เนว\-เสกฺข\-นา\-เสกฺโข ธมฺโม อุปฺปชฺชติ เหตุปจฺจยา. ตีณิ.}

\prose{เหตุยา นว, อวิคเต นว. (สํขิตฺตํ. สหชาตวารมฺปิ ฯเปฯ สมฺปยุตฺตวารมฺปิ ปฏิจฺจ\-วาร\-สทิสํ.)}

\proseitem{50}{กุสโล เนว\-เสกฺข\-นา\-เสกฺโข ธมฺโม กุสลสฺส เนว\-เสกฺข\-นาเสกฺขสฺส ธมฺมสฺส เหตุปจฺจเยน ปจฺจโย. (สํขิตฺตํ.)}

\prose{เหตุยา สตฺต, อารมฺมเณ นว, อวิคเต เตรส. (สํขิตฺตํ. ยถา กุสลตฺติเก ปญฺหาวารํ, เอวํ วิตฺถาเรตพฺพํ.)}

\csromanmarkh{1. กุสลตฺติก}\csromanmarkhii{11. ปริตฺตตฺติก}\csromanheadernomark{2}{1. กุสลตฺติก 11. ปริตฺตตฺติก}
\proseitem{51}{กุสลํ ปริตฺตํ ธมฺมํ ปฏิจฺจ กุสโล ปริตฺโต ธมฺโม อุปฺปชฺชติ เหตุปจฺจยา.  กุสลํ ปริตฺตํ ธมฺมํ ปฏิจฺจ อพฺยากโต ปริตฺโต ธมฺโม อุปฺปชฺชติ เหตุปจฺจยา.  กุสลํ ปริตฺตํ ธมฺมํ ปฏิจฺจ กุสโล ปริตฺโต จ อพฺยากโต ปริตฺโต จ ธมฺมา อุปฺปชฺชนฺติ เหตุปจฺจยา. (3)}

\prose{อกุสลํ ปริตฺตํ ธมฺมํ ปฏิจฺจ อกุสโล ปริตฺโต ธมฺโม อุปฺปชฺชติ เหตุปจฺจยา. ตีณิ.}

\prose{อพฺยากตํ ปริตฺตํ ธมฺมํ ปฏิจฺจ อพฺยากโต ปริตฺโต ธมฺโม อุปฺปชฺชติ เหตุปจฺจยา. (1)}

\prose{\csromanfolio{14}กุสลํ ปริตฺตญฺจ อพฺยากตํ ปริตฺตญฺจ ธมฺมํ ปฏิจฺจ อพฺยากโต ปริตฺโต ธมฺโม อุปฺปชฺชติ เหตุปจฺจยา. (1)}

\prose{อกุสลํ ปริตฺตญฺจ อพฺยากตํ ปริตฺตญฺจ ธมฺมํ ปฏิจฺจ อพฺยากโต ปริตฺโต ธมฺโม อุปฺปชฺชติ เหตุปจฺจยา. (1) (สํขิตฺตํ.)}

\prose{เหตุยา นว, อารมฺมเณ ตีณิ, อวิคเต นว. (สํขิตฺตํ. สหชาตวารมฺปิ ฯเปฯ สมฺปยุตฺตวารมฺปิ ปฏิจฺจ\-วาร\-สทิสํ.)}

\proseitem{52}{กุสโล ปริตฺโต ธมฺโม กุสลสฺส ปริตฺตสฺส ธมฺมสฺส เหตุปจฺจเยน ปจฺจโย. ตีณิ.}

\prose{อกุสโล ปริตฺโต ธมฺโม อกุสลสฺส ปริตฺตสฺส ธมฺมสฺส เหตุปจฺจเยน ปจฺจโย. ตีณิ.}

\prose{อพฺยากโต ปริตฺโต ธมฺโม อพฺยากตสฺส ปริตฺตสฺส ธมฺมสฺส}

\prose{กุสโล ปริตฺโต ธมฺโม กุสลสฺส ปริตฺตสฺส ธมฺมสฺส อารมฺมณ\-ปจฺจเยน ปจฺจโย. (สํขิตฺตํ.)}

\prose{เหตุยา สตฺต, อารมฺมเณ นว, อวิคเต เตรส. (สํขิตฺตํ. ยถา กุสลตฺติเก ปญฺหาวารํ, เอวํ วิตฺถาเรตพฺพํ.)}

\prose{มหคฺคตาทิปทเหตุ}

\prose{53. กุสลํ มหคฺคตํ ธมฺมํ ปฏิจฺจ กุสโล มหคฺคโต}

\prose{อพฺยากตํ มหคฺคตํ ธมฺมํ ปฏิจฺจ อพฺยากโต มหคฺคโต ธมฺโม อุปฺปชฺชติ เหตุปจฺจยา. (1) (สํขิตฺตํ.)}

\prose{เหตุยา เทฺว, อารมฺมเณ เทฺว, อวิคเต เทฺว. (สํขิตฺตํ. สหชาตวาเรปิ ฯเปฯ ปญฺหาวาเรปิ สพฺพตฺถ วิตฺถาโร.)}

\proseitem{54}{กุสลํ อปฺปมาณํ ธมฺมํ ปฏิจฺจ กุสโล อปฺปมาโณ ธมฺโม อุปฺปชฺชติ เหตุปจฺจยา. (1)}

\prose{อพฺยากตํ อปฺปมาณํ ธมฺมํ ปฏิจฺจ อพฺยากโต อปฺปมาโณ ธมฺโม อุปฺปชฺชติ เหตุปจฺจยา. (1) (สํขิตฺตํ.)}

\csromanmarkh{1. กุสลตฺติก}\csromanmarkhii{12. ปริตฺตา\-รมฺมณ\-ตฺติก}\csromanheadernomark{2}{1. กุสลตฺติก 12. ปริตฺตา\-รมฺมณ\-ตฺติก}
\prose{\csromanfolio{15}เหตุยา เทฺว, อวิคเต เทฺว. (สํขิตฺตํ. สหชาตวาเรปิ ฯเปฯ ปญฺหาวาเรปิ สพฺพตฺถ วิตฺถาโร.)}

\csromanmarkh{1. กุสลตฺติก}\csromanmarkhii{12. ปริตฺตา\-รมฺมณ\-ตฺติก}\csromanheadernomark{2}{1. กุสลตฺติก 12. ปริตฺตา\-รมฺมณ\-ตฺติก}
\proseitem{55}{กุสลํ ปริตฺตา\-รมฺมณํ ธมฺมํ ปฏิจฺจ กุสโล ปริตฺตารมฺมโณ ธมฺโม อุปฺปชฺชติ เหตุปจฺจยา. (1)}

\prose{อกุสลํ ปริตฺตา\-รมฺมณํ ธมฺมํ ปฏิจฺจ อกุสโล ปริตฺตารมฺมโณ}

\prose{อพฺยากตํ ปริตฺตา\-รมฺมณํ ธมฺมํ ปฏิจฺจ อพฺยากโต ปริตฺตารมฺมโณ ธมฺโม อุปฺปชฺชติ เหตุปจฺจยา. (1) (สํขิตฺตํ.)}

\prose{เหตุยา ตีณิ, อวิคเต ตีณิ. (สํขิตฺตํ.) (สหชาตวารมฺปิ ฯเปฯ ปญฺหาวารมฺปิ วิตฺถาเรตพฺพํ.)}

\proseitem{56}{กุสลํ มหคฺคตา\-รมฺมณํ ธมฺมํ ปฏิจฺจ กุสโล มหคฺคตา\-รมฺมโณ ธมฺโม อุปฺปชฺชติ เหตุปจฺจยา. (1)}

\prose{อกุสลํ มหคฺคตา\-รมฺมณํ ธมฺมํ ปฏิจฺจ อกุสโล มหคฺคตา\-รมฺมโณ ธมฺโม อุปฺปชฺชติ เหตุปจฺจยา. (1)}

\prose{อพฺยากตํ มหคฺคตา\-รมฺมณํ ธมฺมํ ปฏิจฺจ อพฺยากโต มหคฺคตา\-รมฺมโณ ธมฺโม อุปฺปชฺชติ เหตุปจฺจยา. (1) (สํขิตฺตํ.)}

\prose{เหตุยา ตีณิ, อวิคเต ตีณิ. (สํขิตฺตํ.) (สหชาตวาเรปิ ฯเปฯ ปญฺหาวาเรปิ สพฺพตฺถ วิตฺถาโร.)}

\proseitem{57}{กุสลํ อปฺปมาณา\-รมฺมณํ ธมฺมํ ปฏิจฺจ กุสโล อปฺปมาณารมฺมโณ ธมฺโม อุปฺปชฺชติ เหตุปจฺจยา. (1)}

\prose{อพฺยากตํ อปฺปมาณา\-รมฺมณํ ธมฺมํ ปฏิจฺจ อพฺยากโต อปฺปมาณารมฺมโณ ธมฺโม อุปฺปชฺชติ เหตุปจฺจยา. (1) (สํขิตฺตํ.)}

\prose{เหตุยา เทฺว, อวิคเต เทฺว. (สํขิตฺตํ.) (สหชาตวาเรปิ ฯเปฯ ปญฺหาวาเรปิ สพฺพตฺถ วิตฺถาโร.)}

\csromanmarkh{1. กุสลตฺติก}\csromanmarkhii{13. หีนตฺติก}\csromanheadernomark{2}{1. \textbf{กุสลตฺติก} 13. หีนตฺติก}
\proseitem{58}{\csromanfolio{16}อกุสลํ หีนํ ธมฺมํ ปฏิจฺจ อกุสโล หีโน ธมฺโม อุปฺปชฺชติ เหตุปจฺจยา. (สํขิตฺตํ.)}

\prose{เหตุยา เอกํ ฯเปฯ อวิคเต เอกํ. (สํขิตฺตํ. สพฺพตฺถ วิตฺถาโร.)}

\proseitem{59}{กุสลํ มชฺฌิมํ ธมฺมํ ปฏิจฺจ กุสโล มชฺฌิโม ธมฺโม อุปฺปชฺชติ เหตุปจฺจยา.  กุสลํ มชฺฌิมํ ธมฺมํ ปฏิจฺจ อพฺยากโต มชฺฌิโม ธมฺโม อุปฺปชฺชติ เหตุปจฺจยา.  กุสลํ มชฺฌิมํ ธมฺมํ ปฏิจฺจ กุสโล มชฺฌิโม จ อพฺยากโต มชฺฌิโม จ ธมฺมา อุปฺปชฺชนฺติ เหตุปจฺจยา. (3)}

\prose{อพฺยากตํ มชฺฌิมํ ธมฺมํ ปฏิจฺจ อพฺยากโต มชฺฌิโม ธมฺโม อุปฺปชฺชติ เหตุปจฺจยา. (1)}

\prose{กุสลํ มชฺฌิมญฺจ อพฺยากตํ มชฺฌิมญฺจ ธมฺมํ ปฏิจฺจ อพฺยากโต มชฺฌิโม ธมฺโม อุปฺปชฺชติ เหตุปจฺจยา. (1)}

\prose{กุสลํ มชฺฌิมํ ธมฺมํ ปฏิจฺจ กุสโล มชฺฌิโม ธมฺโม อุปฺปชฺชติ อารมฺมณ\-ปจฺจยา. (1)}

\prose{อพฺยากตํ มชฺฌิมํ ธมฺโม ปฏิจฺจ อพฺยากโต มชฺฌิโม ธมฺโม อุปฺปชฺชติ อารมฺมณ\-ปจฺจยา. (1) (สํขิตฺตํ.)}

\prose{เหตุยา ปญฺจ, อารมฺมเณ เทฺว, อวิคเต ปญฺจ. (สํขิตฺตํ. สหชาตวารมฺปิ ฯเปฯ สมฺปยุตฺตวารมฺปิ ปฏิจฺจ\-วาร\-สทิสํ.)}

\proseitem{60}{กุสโล มชฺฌิโม ธมฺโม กุสลสฺส มชฺฌิมสฺส ธมฺมสฺส เหตุปจฺจเยน ปจฺจโย. ตีณิ.}

\prose{อพฺยากโต มชฺฌิโม ธมฺโม อพฺยากตสฺส มชฺฌิมสฺส ธมฺมสฺส}

\prose{กุสโล มชฺฌิโม ธมฺโม กุสลสฺส มชฺฌิมสฺส ธมฺมสฺส อารมฺมณ\-ปจฺจเยน ปจฺจโย.  กุสโล มชฺฌิโม ธมฺโม อพฺยากตสฺส มชฺฌิมสฺส ธมฺมสฺส อารมฺมณ\-ปจฺจเยน ปจฺจโย. (2)}

\vspace{\tipitakaparskip}
\csromancenter{14. มิจฺฉตฺตนิยต\-ตฺติก อพฺยากโต มชฺฌิโม ธมฺโม อพฺยากตสฺส มชฺฌิมสฺส ธมฺมสฺส อารมฺมณ\-ปจฺจเยน ปจฺจโย.  อพฺยากโต มชฺฌิโม ธมฺโม กุสลสฺส มชฺฌิมสฺส ธมฺมสฺส อารมฺมณ\-ปจฺจเยน ปจฺจโย. (2) (สํขิตฺตํ.)}
\proseitem{61}{\csromanfolio{17}เหตุยา จตฺตาริ, อารมฺมเณ จตฺตาริ, อวิคเต สตฺต. (สํขิตฺตํ. ยถา กุสลตฺติเก ปญฺหาวารํ, เอวํ วิตฺถาเรตพฺพํ.)}

\prose{ปณีตปทเหตุ}

\proseitem{62}{กุสลํ ปณีตํ ธมฺมํ ปฏิจฺจ กุสโล ปณีโต ธมฺโม อุปฺปชฺชติ เหตุปจฺจยา. (1)}

\prose{อพฺยากตํ ปณีตํ ธมฺมํ ปฏิจฺจ อพฺยากโต ปณีโต ธมฺโม อุปฺปชฺชติ เหตุปจฺจยา. (1) (สํขิตฺตํ.)}

\prose{เหตุยา เทฺว, อวิคเต เทฺว. (สํขิตฺตํ. สหชาตวาเรปิ ฯเปฯ ปญฺหาวาเรปิ สพฺพตฺถ วิตฺถาโร.)}

\csromanmarkh{1. กุสลตฺติก}\csromanmarkhii{14. มิจฺฉตฺตนิยต\-ตฺติก}\csromanheadernomark{2}{1. กุสลตฺติก 14. มิจฺฉตฺตนิยต\-ตฺติก}
\proseitem{63}{อกุสลํ มิจฺฉตฺต\-นิยตํ ธมฺมํ ปฏิจฺจ อกุสโล มิจฺฉตฺต\-นิยโต ธมฺโม อุปฺปชฺชติ เหตุปจฺจยา. (สํขิตฺตํ.)}

\prose{เหตุยา เอกํ ฯเปฯ อวิคเต เอกํ. (สํขิตฺตํ. สพฺพตฺถ วิตฺถาโร.)}

\proseitem{64}{กุสลํ สมฺมตฺต\-นิยตํ ธมฺมํ ปฏิจฺจ กุสโล สมฺมตฺตนิยโต ธมฺโม อุปฺปชฺชติ เหตุปจฺจยา. (สํขิตฺตํ.)}

\prose{เหตุยา เอกํ ฯเปฯ อวิคเต เอกํ. (สํขิตฺตํ. สหชาตวาเรปิ ฯเปฯ ปญฺหาวาเรปิ สพฺพตฺถ เอกํ.)}

\proseitem{65}{กุสลํ อนิยตํ ธมฺมํ ปฏิจฺจ กุสโล อนิยโต ธมฺโม อุปฺปชฺชติ เหตุปจฺจยา.  กุสลํ อนิยตํ ธมฺมํ ปฏิจฺจ อพฺยากโต อนิยโต ธมฺโม อุปฺปชฺชติ เหตุปจฺจยา.  กุสลํ อนิยตํ ธมฺมํ \csromanfolio{18}ปฏิจฺจ กุสโล อนิยโต จ อพฺยากโต อนิยโต จ ธมฺมา อุปฺปชฺชนฺติ เหตุปจฺจยา. (3)}

\prose{อกุสลํ อนิยตํ ธมฺมํ ปฏิจฺจ อกุสโล อนิยโต ธมฺโม อุปฺปชฺชติ เหตุปจฺจยา. ตีณิ.}

\prose{อพฺยากตํ อนิยตํ ธมฺมํ ปฏิจฺจ อพฺยากโต อนิยโต ธมฺโม อุปฺปชฺชติ เหตุปจฺจยา. (1)}

\prose{กุสลํ อนิยตญฺจ อพฺยากตํ อนิยตญฺจ ธมฺมํ ปฏิจฺจ อพฺยากโต อนิยโต ธมฺโม อุปฺปชฺชติ เหตุปจฺจยา. (1)}

\prose{อกุสลํ อนิยตญฺจ อพฺยากตํ อนิยตญฺจ ธมฺมํ ปฏิจฺจ อพฺยากโต อนิยโต ธมฺโม อุปฺปชฺชติ เหตุปจฺจยา. (1) (สํขิตฺตํ.)}

\prose{เหตุยา นว, อารมฺมเณ ตีณิ, อวิคเต นว. (สํขิตฺตํ. สหชาตวารมฺปิ ฯเปฯ สมฺปยุตฺตวารมฺปิ ปฏิจฺจ\-วาร\-สทิสํ.)}

\proseitem{66}{กุสโล อนิยโต ธมฺโม กุสลสฺส อนิยตสฺส ธมฺมสฺส เหตุปจฺจเยน ปจฺจโย. ตีณิ.}

\prose{อกุสโล อนิยโต ธมฺโม อกุสลสฺส อนิยตสฺส ธมฺมสฺส เหตุปจฺจเยน ปจฺจโย. ตีณิ.}

\prose{อพฺยากโต อนิยโต ธมฺโม อพฺยากตสฺส อนิยตสฺส ธมฺมสฺส เหตุปจฺจเยน ปจฺจโย. (1) (สํขิตฺตํ.)}

\prose{เหตุยา สตฺต, อารมฺมเณ นว, อวิคเต เตรส. (สํขิตฺตํ. ยถา กุสลตฺติเก ปญฺหาวารํ, เอวํ วิตฺถาเรตพฺพํ.)}

\csromanmarkh{1. กุสลตฺติก}\csromanmarkhii{15. มคฺคา\-รมฺมณ\-ตฺติก}\csromanheadernomark{2}{1. กุสลตฺติก 15. มคฺคา\-รมฺมณ\-ตฺติก}
\proseitem{67}{กุสลํ มคฺคารมฺมณํ ธมฺมํ ปฏิจฺจ กุสโล มคฺคารมฺมโณ ธมฺโม อุปฺปชฺชติ เหตุปจฺจยา. (1)}

\csromanmarkh{1. กุสลตฺติก}\csromanmarkhii{17. อตีตตฺติก}\csromanheadernomark{2}{1. กุสลตฺติก 17. อตีตตฺติก}
\prose{\csromanfolio{19}อพฺยากตํ มคฺคารมฺมณํ ธมฺมํ ปฏิจฺจ อพฺยากโต มคฺคารมฺมโณ ธมฺโม อุปฺปชฺชติ เหตุปจฺจยา. (1) (สํขิตฺตํ.)}

\prose{เหตุยา เทฺว ฯเปฯ อวิคเต เทฺว. (สํขิตฺตํ. สพฺพตฺถ วิตฺถาโร.)}

\proseitem{68}{กุสลํ มคฺคเหตุกํ ธมฺมํ ปฏิจฺจ กุสโล มคฺคเหตุโก ธมฺโม อุปฺปชฺชติ เหตุปจฺจยา. (สํขิตฺตํ.)}

\prose{เหตุยา เอกํ ฯเปฯ อวิคเต เอกํ. (สพฺพตฺถ เอกํ. สํขิตฺตํ.)}

\prose{69. กุสลํ มคฺคาธิปติํ ธมฺมํ ปฏิจฺจ กุสโล มคฺคาธิปติ}

\prose{อพฺยากตํ มคฺคาธิปติํ ธมฺมํ ปฏิจฺจ อพฺยากโต มคฺคาธิปติ ธมฺโม อุปฺปชฺชติ เหตุปจฺจยา. (1) (สํขิตฺตํ.)}

\prose{เหตุยา เทฺว ฯเปฯ อวิคเต เทฺว. (สํขิตฺตํ. สพฺพตฺถ วิตฺถาโร.)}

\csromanmarkh{1. กุสลตฺติก}\csromanmarkhii{16. อุปฺปนฺนตฺติก}\csromanheadernomark{2}{1. \textbf{กุสลตฺติก} 16. อุปฺปนฺนตฺติก}
\csromanheader{2}{7. \textbf{ปญฺหาวาร ปจฺจยจตุกฺก}}
\proseitem{70}{กุสโล อุปฺปนฺโน ธมฺโม กุสลสฺส อุปฺปนฺนสฺส ธมฺมสฺส เหตุปจฺจเยน ปจฺจโย. (สํขิตฺตํ.)}

\prose{เหตุยา สตฺต. (สํขิตฺตํ.)}

\csromanmarkh{1. กุสลตฺติก}\csromanmarkhii{17. อตีตตฺติก}\csromanheadernomark{2}{1. กุสลตฺติก 17. อตีตตฺติก}
\csromanheader{2}{7. \textbf{ปญฺหาวาร ปจฺจยจตุกฺก}}
\proseitem{71}{กุสโล ปจฺจุปฺปนฺโน ธมฺโม กุสลสฺส ปจฺจุปฺปนฺนสฺส ธมฺมสฺส เหตุปจฺจเยน ปจฺจโย. (สํขิตฺตํ.)}

\prose{เหตุยา สตฺต. (สํขิตฺตํ.)}

\csromanmarkh{1. กุสลตฺติก}\csromanmarkhii{18. อตีตา\-รมฺมณ\-ตฺติก}\csromanheadernomark{2}{1. กุสลตฺติก 18. อตีตา\-รมฺมณ\-ตฺติก}
\prose{\csromanfolio{20}72. กุสลํ อตีตารมฺมณํ ธมฺมํ ปฏิจฺจ กุสโล อตีตารมฺมโณ}

\prose{อกุสลํ อตีตารมฺมณํ ธมฺมํ ปฏิจฺจ อกุสโล อตีตารมฺมโณ}

\prose{อพฺยากตํ อตีตารมฺมณํ ธมฺมํ ปฏิจฺจ อพฺยากโต อตีตารมฺมโณ ธมฺโม อุปฺปชฺชติ เหตุปจฺจยา. (1) (สํขิตฺตํ.)}

\prose{เหตุยา ตีณิ, อวิคเต ตีณิ. (สํขิตฺตํ.) (สหชาตวารมฺปิ ฯเปฯ สมฺปยุตฺตวารมฺปิ ปฏิจฺจ\-วาร\-สทิสํ.)}

\proseitem{73}{กุสโล อตีตารมฺมโณ ธมฺโม กุสลสฺส อตีตา\-รมฺมณสฺส ธมฺมสฺส เหตุปจฺจเยน ปจฺจโย. (1)}

\prose{อกุสโล อตีตารมฺมโณ ธมฺโม อกุสลสฺส อตีตา\-รมฺมณสฺส ธมฺมสฺส เหตุปจฺจเยน ปจฺจโย. (1)}

\prose{อพฺยากโต อตีตารมฺมโณ ธมฺโม อพฺยากตสฺส อตีตา\-รมฺมณสฺส ธมฺมสฺส เหตุปจฺจเยน ปจฺจโย. (1) (สํขิตฺตํ.)}

\prose{เหตุยา ตีณิ, อารมฺมเณ นว, อวิคเต ตีณิ. (สํขิตฺตํ. ยถา กุสลตฺติเก ปญฺหาวารํ, เอวํ วิตฺถาเรตพฺพํ.)}

\prose{อนาคตารมฺมณปทเหตุ}

\proseitem{74}{กุสลํ อนาคตา\-รมฺมณํ ธมฺมํ ปฏิจฺจ กุสโล อนาคตารมฺมโณ ธมฺโม อุปฺปชฺชติ เหตุปจฺจยา. (สํขิตฺตํ.)}

\prose{เหตุยา ตีณิ, อารมฺมเณ ตีณิ, อวิคเต ตีณิ. (สํขิตฺตํ. สหชาตวารมฺปิ ฯเปฯ สมฺปยุตฺตวารมฺปิ ปฏิจฺจ\-วาร\-สทิสํ.)}

\csromanmarkh{1. กุสลตฺติก}\csromanmarkhii{19. อชฺฌตฺตตฺติก}\csromanheadernomark{2}{1. กุสลตฺติก 19. อชฺฌตฺตตฺติก}
\proseitem{75}{\csromanfolio{21}กุสโล อนาคตารมฺมโณ ธมฺโม กุสลสฺส อนาคตา\-รมฺมณสฺส ธมฺมสฺส เหตุปจฺจเยน ปจฺจโย. (สํขิตฺตํ.)}

\prose{เหตุยา ตีณิ, อารมฺมเณ นว, อวิคเต ตีณิ. (สํขิตฺตํ. ยถา กุสลตฺติเก ปญฺหาวารํ, เอวํ วิตฺถาเรตพฺพํ.)}

\prose{ปจฺจุปฺปนฺนารมฺมณปทเหตุ}

\proseitem{76}{กุสลํ ปจฺจุปฺปนฺนา\-รมฺมณํ ธมฺมํ ปฏิจฺจ กุสโล ปจฺจุปฺปนฺนา\-รมฺมโณ ธมฺโม อุปฺปชฺชติ เหตุปจฺจยา. (สํขิตฺตํ.)}

\prose{เหตุยา ตีณิ, อารมฺมเณ ตีณิ, อวิคเต ตีณิ. (สํขิตฺตํ. สหชาตวารมฺปิ ฯเปฯ สมฺปยุตฺตวารมฺปิ ปฏิจฺจ\-วาร\-สทิสํ.)}

\proseitem{77}{กุสโล ปจฺจุปฺปนฺนา\-รมฺมโณ ธมฺโม กุสลสฺส ปจฺจุปฺปนฺนา\-รมฺมณสฺส ธมฺมสฺส เหตุปจฺจเยน ปจฺจโย. (1)}

\prose{อกุสโล ปจฺจุปฺปนฺนา\-รมฺมโณ ธมฺโม อกุสลสฺส ปจฺจุปฺปนฺนา\-รมฺมณสฺส ธมฺมสฺส เหตุปจฺจเยน ปจฺจโย. (1)}

\prose{อพฺยากโต ปจฺจุปฺปนฺนา\-รมฺมโณ ธมฺโม อพฺยากตสฺส ปจฺจุปฺปนฺนา\-รมฺมณสฺส ธมฺมสฺส เหตุปจฺจเยน ปจฺจโย. (1)}

\prose{เหตุยา ตีณิ, อารมฺมเณ ฉ, อวิคเต ตีณิ. (สํขิตฺตํ. ยถา กุสลตฺติเก ปญฺหาวารํ, เอวํ วิตฺถาเรตพฺพํ.)}

\csromanmarkh{1. กุสลตฺติก}\csromanmarkhii{19. อชฺฌตฺตตฺติก}\csromanheadernomark{2}{1. กุสลตฺติก 19. อชฺฌตฺตตฺติก}
\proseitem{78}{กุสลํ อชฺฌตฺตํ ธมฺมํ ปฏิจฺจ กุสโล อชฺฌตฺโต ธมฺโม อุปฺปชฺชติ เหตุปจฺจยา. ตีณิ.}

\prose{อกุสลํ อชฺฌตฺตํ ธมฺมํ ปฏิจฺจ อกุสโล อชฺฌตฺโต ธมฺโม อุปฺปชฺชติ เหตุปจฺจยา. ตีณิ.}

\prose{อพฺยากตํ อชฺฌตฺตํ ธมฺมํ ปฏิจฺจ อพฺยากโต อชฺฌตฺโต ธมฺโม อุปฺปชฺชติ เหตุปจฺจยา. ตีณิ.}

\prose{\csromanfolio{22}เหตุยา นว, อวิคเต นว. (สํขิตฺตํ.) (สหชาตวารมฺปิ ฯเปฯ สมฺปยุตฺตวารมฺปิ ปฏิจฺจ\-วาร\-สทิสํ, เอวํ วิตฺถาเรตพฺพํ.)}

\proseitem{79}{กุสโล อชฺฌตฺโต ธมฺโม กุสลสฺส อชฺฌตฺตสฺส ธมฺมสฺส เหตุปจฺจเยน ปจฺจโย. ตีณิ.}

\prose{อกุสโล อชฺฌตฺโต ธมฺโม อกุสลสฺส อชฺฌตฺตสฺส ธมฺมสฺส เหตุปจฺจเยน ปจฺจโย. ตีณิ.}

\prose{อพฺยากโต อชฺฌตฺโต ธมฺโม อพฺยากตสฺส อชฺฌตฺตสฺส ธมฺมสฺส เหตุปจฺจเยน ปจฺจโย. (1) (สํขิตฺตํ.)}

\prose{เหตุยา สตฺต, อารมฺมเณ นว, อวิคเต เตรส. (สํขิตฺตํ. ยถา กุสลตฺติเก ปญฺหาวารํ, เอวํ วิตฺถาเรตพฺพํ.)}

\prose{พหิทฺธาปทเหตุ}

\proseitem{80}{กุสลํ พหิทฺธา ธมฺมํ ปฏิจฺจ กุสโล พหิทฺธา ธมฺโม อุปฺปชฺชติ เหตุปจฺจยา. ตีณิ.}

\prose{อกุสลํ พหิทฺธา ธมฺมํ ปฏิจฺจ อกุสโล พหิทฺธา ธมฺโม อุปฺปชฺชติ เหตุปจฺจยา. ตีณิ.}

\prose{อพฺยากตํ พหิทฺธา ธมฺมํ ปฏิจฺจ อพฺยากโต พหิทฺธา ธมฺโม อุปฺปชฺชติ เหตุปจฺจยา. (1)}

\prose{กุสลํ พหิทฺธา จ อพฺยากตํ พหิทฺธา จ ธมฺมํ ปฏิจฺจ อพฺยากโต พหิทฺธา ธมฺโม อุปฺปชฺชติ เหตุปจฺจยา. (1)}

\prose{อกุสลํ พหิทฺธา จ อพฺยากตํ พหิทฺธา จ ธมฺมํ ปฏิจฺจ อพฺยากโต พหิทฺธา ธมฺโม อุปฺปชฺชติ เหตุปจฺจยา. (1) (สํขิตฺตํ.)}

\prose{เหตุยา นว ฯเปฯ วิปาเก เอกํ ฯเปฯ อวิคเต นว. (สํขิตฺตํ.) (สหชาตวารมฺปิ ฯเปฯ สมฺปยุตฺตวารมฺปิ ปฏิจฺจ\-วาร\-สทิสํ.)}

\proseitem{81}{กุสโล พหิทฺธา ธมฺโม กุสลสฺส พหิทฺธา ธมฺมสฺส เหตุปจฺจเยน ปจฺจโย. ตีณิ.}

\csromanmarkh{1. กุสลตฺติก}\csromanmarkhii{21. นิทสฺสนตฺติก}\csromanheadernomark{2}{1. กุสลตฺติก 21. นิทสฺสนตฺติก}
\prose{\csromanfolio{23}อกุสโล พหิทฺธา ธมฺโม อกุสลสฺส พหิทฺธา ธมฺมสฺส เหตุปจฺจเยน ปจฺจโย. ตีณิ.}

\prose{อพฺยากโต พหิทฺธา ธมฺโม อพฺยากตสฺส พหิทฺธา ธมฺมสฺส เหตุปจฺจเยน ปจฺจโย. (1) (สํขิตฺตํ.)}

\prose{เหตุยา สตฺต, อารมฺมเณ นว, อธิปติยา ทส ฯเปฯ อวิคเต เตรส.}

\prose{(สํขิตฺตํ. ยถา กุสลตฺติเก ปญฺหาวารํ, เอวํ วิตฺถาเรตพฺพํ.)}

\csromanmarkh{1. กุสลตฺติก}\csromanmarkhii{20. อชฺฌตฺตา\-รมฺมณ\-ตฺติก}\csromanheadernomark{2}{1. กุสลตฺติก 20. อชฺฌตฺตา\-รมฺมณ\-ตฺติก}
\proseitem{82}{กุสลํ อชฺฌตฺตา\-รมฺมณํ ธมฺมํ ปฏิจฺจ กุสโล อชฺฌตฺตา\-รมฺมโณ ธมฺโม อุปฺปชฺชติ เหตุปจฺจยา. (สํขิตฺตํ.)}

\prose{เหตุยา ตีณิ ฯเปฯ วิปาเก เอกํ ฯเปฯ อวิคเต ตีณิ. (สํขิตฺตํ.)}

\proseitem{83}{กุสลํ พหิทฺธา\-รมฺมณํ ธมฺมํ ปฏิจฺจ กุสโล พหิทฺธา\-รมฺมโณ ธมฺโม อุปฺปชฺชติ เหตุปจฺจยา. (สํขิตฺตํ.)}

\prose{เหตุยา ตีณิ ฯเปฯ วิปาเก เอกํ ฯเปฯ อวิคเต ตีณิ. (สํขิตฺตํ.)}

\csromanmarkh{1. กุสลตฺติก}\csromanmarkhii{21. สนิทสฺสน\-ตฺติก}\csromanheadernomark{2}{1. กุสลตฺติก 21. สนิทสฺสน\-ตฺติก}
\proseitem{84}{อพฺยากตํ อนิทสฺสน\-สปฺปฏิฆํ ธมฺมํ ปฏิจฺจ อพฺยากโต อนิทสฺสน\-สปฺปฏิโฆ ธมฺโม อุปฺปชฺชติ เหตุปจฺจยา. (สํขิตฺตํ.)}

\prose{เหตุยา เอกํ ฯเปฯ อวิคเต เอกํ. (สํขิตฺตํ.) (สหชาตวารมฺปิ ฯเปฯ ปญฺหาวารมฺปิ วิตฺถาเรตพฺพํ.)}

\proseitem{85}{กุสลํ อนิทสฺสน\-อปฺปฏิฆํ ธมฺมํ ปฏิจฺจ กุสโล อนิทสฺสน\-อปฺปฏิโฆ ธมฺโม อุปฺปชฺชติ เหตุปจฺจยา.  กุสลํ อนิทสฺสน\-อปฺปฏิฆํ ธมฺมํ \csromanfolio{24}ปฏิจฺจ อพฺยากโต อนิทสฺสน\-อปฺปฏิโฆ ธมฺโม อุปฺปชฺชติ เหตุปจฺจยา.  กุสลํ อนิทสฺสน\-อปฺปฏิฆํ ธมฺมํ ปฏิจฺจ กุสโล อนิทสฺสน\-อปฺปฏิโฆ จ อพฺยากโต อนิทสฺสน\-อปฺปฏิโฆ จ ธมฺมา อุปฺปชฺชนฺติ เหตุปจฺจยา. (3)}

\prose{อกุสลํ อนิทสฺสน\-อปฺปฏิฆํ ธมฺมํ ปฏิจฺจ อกุสโล อนิทสฺสน\-อปฺปฏิโฆ ธมฺโม อุปฺปชฺชติ เหตุปจฺจยา. ตีณิ.}

\prose{อพฺยากตํ อนิทสฺสน\-อปฺปฏิฆํ ธมฺมํ ปฏิจฺจ อพฺยากโต อนิทสฺสน\-อปฺปฏิโฆ ธมฺโม อุปฺปชฺชติ เหตุปจฺจยา. (1) กุสลํ อนิทสฺสนอปฺปฏิฆญฺจ อพฺยากตํ อนิทสฺสนอปฺปฏิฆญฺจ ธมฺมํ ปฏิจฺจ อพฺยากโต อนิทสฺสน\-อปฺปฏิโฆ ธมฺโม อุปฺปชฺชติ}

\prose{อกุสลํ อนิทสฺสนอปฺปฏิฆญฺจ อพฺยากตํ อนิทสฺสนอปฺปฏิฆญฺจ ธมฺมํ ปฏิจฺจ อพฺยากโต อนิทสฺสน\-อปฺปฏิโฆ ธมฺโม อุปฺปชฺชติ}

\prose{เหตุยา นว, อารมฺมเณ ตีณิ ฯเปฯ วิปาเก เอกํ ฯเปฯ อวิคเต นว. (สํขิตฺตํ. สหชาตวารมฺปิ ฯเปฯ สมฺปยุตฺตวารมฺปิ ปญฺหาวารมฺปิ วิตฺถาเรตพฺพํ.) กุสล\-ตฺติก\-สนิทสฺสน\-ตฺติกํ นิฏฺฐิตํ.}

\csromanmarkh{2. เวทนาตฺติก}\csromanmarkhii{1. กุสลตฺติก}\csromanheadernomark{2}{2. \textbf{เวทนาตฺติก 1. กุสลตฺติก}}
\proseitem{86}{สุขาย เวทนาย สมฺปยุตฺตํ กุสลํ ธมฺมํ ปฏิจฺจ สุขาย เวทนาย สมฺปยุตฺโต กุสโล ธมฺโม อุปฺปชฺชติ เหตุปจฺจยา. (1)}

\prose{อทุกฺขม\-สุขาย เวทนาย สมฺปยุตฺตํ กุสลํ ธมฺมํ ปฏิจฺจ อทุกฺขม\-สุขาย เวทนาย สมฺปยุตฺโต กุสโล ธมฺโม อุปฺปชฺชติ}

\prose{เหตุยา เทฺว ฯเปฯ อวิคเต เทฺว. (สํขิตฺตํ.)}

\prose{(สหชาตวารมฺปิ ฯเปฯ สมฺปยุตฺตวารมฺปิ ปฏิจฺจ\-วาร\-สทิสํ.)}

\proseitem{87}{สุขาย เวทนาย สมฺปยุตฺโต กุสโล ธมฺโม สุขาย เวทนาย สมฺปยุตฺตสฺส กุสลสฺส ธมฺมสฺส เหตุปจฺจเยน ปจฺจโย. (1)}

\csromanmarkh{2. เวทนาตฺติก}\csromanmarkhii{1. กุสลตฺติก}\csromanheadernomark{2}{2. เวทนาตฺติก 1. กุสลตฺติก}
\prose{\csromanfolio{25}อทุกฺขม\-สุขาย เวทนาย สมฺปยุตฺโต กุสโล ธมฺโม อทุกฺขม\-สุขาย เวทนาย สมฺปยุตฺตสฺส กุสลสฺส ธมฺมสฺส เหตุปจฺจเยน ปจฺจโย. (1) (สํขิตฺตํ.)}

\prose{เหตุยา เทฺว, อวิคเต เทฺว. (สํขิตฺตํ. สพฺพตฺถ วิตฺถาโร.)}

\proseitem{88}{สุขาย เวทนาย สมฺปยุตฺตํ อกุสลํ ธมฺมํ ปฏิจฺจ สุขาย เวทนาย สมฺปยุตฺโต อกุสโล ธมฺโม อุปฺปชฺชติ เหตุปจฺจยา. (1)}

\prose{ทุกฺขาย เวทนาย สมฺปยุตฺตํ อกุสลํ ธมฺมํ ปฏิจฺจ ทุกฺขาย เวทนาย สมฺปยุตฺโต อกุสโล ธมฺโม อุปฺปชฺชติ เหตุปจฺจยา. (1)}

\prose{อทุกฺขม\-สุขาย เวทนาย สมฺปยุตฺตํ อกุสลํ ธมฺมํ ปฏิจฺจ อทุกฺขม\-สุขาย เวทนาย สมฺปยุตฺโต อกุสโล ธมฺโม อุปฺปชฺชติ เหตุปจฺจยา. (1) (สํขิตฺตํ.)}

\prose{เหตุยา ตีณิ ฯเปฯ อวิคเต ตีณิ. (สํขิตฺตํ. สพฺพตฺถ วิตฺถาโร.)}

\proseitem{89}{สุขาย เวทนาย สมฺปยุตฺตํ อพฺยากตํ ธมฺมํ ปฏิจฺจ สุขาย เวทนาย สมฺปยุตฺโต อพฺยากโต ธมฺโม อุปฺปชฺชติ เหตุปจฺจยา. (1)}

\prose{อทุกฺขม\-สุขาย เวทนาย สมฺปยุตฺตํ อพฺยากตํ ธมฺมํ ปฏิจฺจ อทุกฺขม\-สุขาย เวทนาย สมฺปยุตฺโต อพฺยากโต ธมฺโม อุปฺปชฺชติ เหตุปจฺจยา. (1) (สํขิตฺตํ.)}

\prose{เหตุยา เทฺว, อวิคเต ตีณิ. (สํขิตฺตํ.)}

\csromanmarkh{3. วิปากตฺติก}\csromanmarkhii{1. กุสลตฺติก}\csromanheadernomark{2}{3. \textbf{วิปากตฺติก 1. กุสลตฺติก}}
\prose{1. ปฏิจฺจวารเหตุ}

\proseitem{90}{วิปาก\-ธมฺม\-ธมฺมํ กุสลํ ธมฺมํ ปฏิจฺจ วิปาก\-ธมฺม\-ธมฺโม กุสโล ธมฺโม อุปฺปชฺชติ เหตุปจฺจยา.}

\prose{เหตุยา เอกํ ฯเปฯ อวิคเต เอกํ. (สํขิตฺตํ.)}

\proseitem{91}{\csromanfolio{26}วิปาก\-ธมฺม\-ธมฺมํ อกุสลํ ธมฺมํ ปฏิจฺจ วิปาก\-ธมฺม\-ธมฺโม อกุสโล ธมฺโม อุปฺปชฺชติ เหตุปจฺจยา.}

\prose{เหตุยา เอกํ. (สพฺพตฺถ เอกํ.)}

\proseitem{92}{วิปากํ อพฺยากตํ ธมฺมํ ปฏิจฺจ วิปาโก อพฺยากโต ธมฺโม อุปฺปชฺชติ เหตุปจฺจยา. ตีณิ.}

\prose{เนว\-วิปาก\-น\-วิปากธมฺม\-ธมฺมํ อพฺยากตํ ธมฺมํ ปฏิจฺจ เนว\-วิปาก\-น\-วิปาก\-ธมฺม\-ธมฺโม อพฺยากโต ธมฺโม อุปฺปชฺชติ เหตุปจฺจยา. ตีณิ.}

\prose{วิปากํ อพฺยากตญฺจ เนว\-วิปาก\-น\-วิปากธมฺม\-ธมฺมํ อพฺยากตญฺจ ธมฺมํ ปฏิจฺจ วิปาโก อพฺยากโต ธมฺโม อุปฺปชฺชติ เหตุปจฺจยา. ตีณิ. (สํขิตฺตํ.)}

\prose{เหตุยา นว, อวิคเต นว. (สํขิตฺตํ. สพฺพตฺถ วิตฺถาโร.)}

\csromanmarkh{4. อุปาทินฺนตฺติก}\csromanmarkhii{1. กุสลตฺติก}\csromanheadernomark{2}{4. \textbf{อุปาทินฺนตฺติก 1. กุสลตฺติก}}
\prose{1. ปฏิจฺจวารเหตุ}

\proseitem{93}{อนุปาทินฺนุ\-ปาทานิยํ กุสลํ ธมฺมํ ปฏิจฺจ อนุปาทินฺนุ\-ปาทานิโย กุสโล ธมฺโม อุปฺปชฺชติ เหตุปจฺจยา. (1)}

\prose{อนุปาทินฺน\-อนุปาทานิยํ กุสลํ ธมฺมํ ปฏิจฺจ อนุปาทินฺน\-อนุปาทานิโย กุสโล ธมฺโม อุปฺปชฺชติ เหตุปจฺจยา. (1)}

\prose{เหตุยา เทฺว. (สพฺพตฺถ วิตฺถาโร.)}

\prose{อนุปาทินฺนุ\-ปาทานิยํ อกุสลํ ธมฺมํ ปฏิจฺจ อนุปาทินฺนุ\-ปาทานิโย อกุสโล ธมฺโม อุปฺปชฺชติ เหตุปจฺจยา. (สพฺพตฺถ เอกํ. สพฺพตฺถ วิตฺถาโร.)}

\proseitem{94}{อุปาทินฺนุ\-ปาทานิยํ อพฺยากตํ ธมฺมํ ปฏิจฺจ อุปาทินฺนุ\-ปาทานิโย อพฺยากโต ธมฺโม อุปฺปชฺชติ เหตุปจฺจยา. ตีณิ. (สพฺพตฺถ วิตฺถาโร.)}

\prose{อนุปาทินฺนุ\-ปาทานิยํ อพฺยากตํ ธมฺมํ ปฏิจฺจ อนุปาทินฺนุ\-ปาทานิโย อพฺยากโต ธมฺโม อุปฺปชฺชติ เหตุปจฺจยา. ตีณิ.}

\csromanmarkh{6. วิตกฺกตฺติก}\csromanmarkhii{1. กุสลตฺติก}\csromanheadernomark{2}{6. วิตกฺกตฺติก 1. กุสลตฺติก}
\prose{\csromanfolio{27}อนุปาทินฺน\-อนุปาทานิยํ อพฺยากตํ ธมฺมํ ปฏิจฺจ อนุปาทินฺน\-อนุปาทานิโย อพฺยากโต ธมฺโม อุปฺปชฺชติ เหตุปจฺจยา. (1)}

\prose{อนุปาทินฺนุ\-ปาทานิยํ อพฺยากตญฺจ อนุปาทินฺน\-อนุปาทานิยํ อพฺยากตญฺจ ธมฺมํ ปฏิจฺจ อนุปาทินฺนุ\-ปาทานิโย อพฺยากโต ธมฺโม อุปฺปชฺชติ เหตุปจฺจยา. (1)}

\prose{อุปาทินฺนุ\-ปาทานิยํ อพฺยากตญฺจ อนุปาทินฺนุ\-ปาทานิยํ อพฺยากตญฺจ ธมฺมํ ปฏิจฺจ อนุปาทินฺนุ\-ปาทานิโย อพฺยากโต ธมฺโม อุปฺปชฺชติ เหตุปจฺจยา. (1) (สํขิตฺตํ.)}

\prose{เหตุยา นว. (สพฺพตฺถ วิตฺถาโร.)}

\csromanmarkh{5. สํกิลิฏฺฐ\-ตฺติก}\csromanmarkhii{1. กุสลตฺติก}\csromanheadernomark{2}{5. สํกิลิฏฺฐ\-ตฺติก 1. กุสลตฺติก}
\prose{1. ปฏิจฺจวารเหตุ}

\proseitem{95}{อสํกิลิฏฺฐ\-สํกิเลสิกํ กุสลํ ธมฺมํ ปฏิจฺจ อสํกิลิฏฺฐ\-สํกิเลสิโก กุสโล ธมฺโม อุปฺปชฺชติ เหตุปจฺจยา.}

\prose{อสํกิลิฏฺฐ\-อสํกิเลสิกํ กุสลํ ธมฺมํ ปฏิจฺจ อสํกิลิฏฺฐ\-อสํกิเลสิโก กุสโล ธมฺโม อุปฺปชฺชติ เหตุปจฺจยา. (สพฺพตฺถ เทฺว.)}

\prose{สํกิลิฏฺฐ\-สํกิเลสิกํ อกุสลํ ธมฺมํ ปฏิจฺจ สํกิลิฏฺฐ\-สํกิเลสิโก อกุสโล ธมฺโม อุปฺปชฺชติ เหตุปจฺจยา. (สพฺพตฺถ เอกํ.)}

\prose{อสํกิลิฏฺฐ\-สํกิเลสิกํ อพฺยากตํ ธมฺมํ ปฏิจฺจ อสํกิลิฏฺฐ\-สํกิเลสิโก อพฺยากโต ธมฺโม อุปฺปชฺชติ เหตุปจฺจยา.}

\prose{เหตุยา ปญฺจ.}

\csromanmarkh{6. วิตกฺกตฺติก}\csromanmarkhii{1. กุสลตฺติก}\csromanheadernomark{2}{6. วิตกฺกตฺติก 1. กุสลตฺติก}
\prose{1. ปฏิจฺจวารเหตุ}

\proseitem{96}{สวิตกฺก\-สวิจารํ กุสลํ ธมฺมํ ปฏิจฺจ สวิตกฺก\-สวิจาโร กุสโล ธมฺโม อุปฺปชฺชติ เหตุปจฺจยา. ตีณิ.}

\prose{\csromanfolio{28}อวิตกฺก\-วิจารมตฺตํ กุสลํ ธมฺมํ ปฏิจฺจ อวิตกฺก\-วิจารมตฺโต กุสโล ธมฺโม อุปฺปชฺชติ เหตุปจฺจยา. จตฺตาริ.}

\prose{อวิตกฺก\-อวิจารํ กุสลํ ธมฺมํ ปฏิจฺจ อวิตกฺก\-อวิจาโร กุสโล}

\prose{อวิตกฺก\-อวิจารํ กุสลํ ธมฺมํ ปฏิจฺจ อวิตกฺก\-วิจารมตฺโต กุสโล ธมฺโม อุปฺปชฺชติ เหตุปจฺจยา. (1)}

\prose{อวิตกฺก\-วิจารมตฺตํ กุสลญฺจ อวิตกฺก\-อวิจารํ กุสลญฺจ ธมฺมํ ฯเปฯ. สวิตกฺก\-สวิจารํ กุสลญฺจ อวิตกฺก\-วิจารมตฺตํ กุสลญฺจ ธมฺมํ ปฏิจฺจ สวิตกฺก\-สวิจาโร กุสโล ธมฺโม อุปฺปชฺชติ}

\prose{เหตุยา เอกาทส.}

\proseitem{97}{สวิตกฺก\-สวิจารํ อกุสลํ ธมฺมํ ปฏิจฺจ สวิตกฺก\-สวิจาโร อกุสโล ธมฺโม อุปฺปชฺชติ เหตุปจฺจยา. ตีณิ.}

\prose{อวิตกฺก\-วิจารมตฺตํ อกุสลํ ธมฺมํ ปฏิจฺจ สวิตกฺก\-สวิจาโร อกุสโล ธมฺโม อุปฺปชฺชติ เหตุปจฺจยา. (1)}

\prose{สวิตกฺก\-สวิจารํ อกุสลญฺจ อวิตกฺก\-วิจารมตฺตํ อกุสลญฺจ ธมฺมํ ปฏิจฺจ สวิตกฺก\-สวิจาโร อกุสโล ธมฺโม อุปฺปชฺชติ}

\prose{เหตุยา ปญฺจ.}

\proseitem{98}{สวิตกฺก\-สวิจารํ อพฺยากตํ ธมฺมํ ปฏิจฺจ สวิตกฺก\-สวิจาโร อพฺยากโต ธมฺโม อุปฺปชฺชติ เหตุปจฺจยา.}

\prose{เหตุยา สตฺตติํส.}

\csromanmarkh{7. ปีติตฺติก}\csromanmarkhii{1. กุสลตฺติก}\csromanheadernomark{2}{7. \textbf{ปีติตฺติก 1. กุสลตฺติก}}
\prose{1. ปฏิจฺจวารเหตุ}

\proseitem{99}{ปีติสหคตํ กุสลํ ธมฺมํ ปฏิจฺจ ปีติสหคโต กุสโล ธมฺโม อุปฺปชฺชติ เหตุปจฺจยา. ตีณิ.}

\csromanmarkh{7. ปีติตฺติก}\csromanmarkhii{1. กุสลตฺติก}\csromanheadernomark{2}{7. ปีติตฺติก 1. กุสลตฺติก}
\prose{\csromanfolio{29}สุขสหคตํ กุสลํ ธมฺมํ ปฏิจฺจ สุขสหคโต กุสโล ธมฺโม อุปฺปชฺชติ เหตุปจฺจยา. ตีณิ.}

\prose{อุเปกฺขา\-สหคตํ กุสลํ ธมฺมํ ปฏิจฺจ อุเปกฺขา\-สหคโต กุสโล}

\prose{ปีติสหคตํ กุสลญฺจ สุขสหคตํ กุสลญฺจ ธมฺมํ ปฏิจฺจ ปีติสหคโต กุสโล ธมฺโม อุปฺปชฺชติ เหตุปจฺจยา. ตีณิ. (สพฺพตฺถ ทส. สพฺพตฺถ วิตฺถาโร.)}

\proseitem{100}{ปีติสหคตํ อกุสลํ ธมฺมํ ปฏิจฺจ ปีติสหคโต อกุสโล ธมฺโม อุปฺปชฺชติ เหตุปจฺจยา. ตีณิ.}

\prose{สุขสหคตํ อกุสลํ ธมฺมํ ปฏิจฺจ สุขสหคโต อกุสโล ธมฺโม อุปฺปชฺชติ เหตุปจฺจยา. ตีณิ.}

\prose{อุเปกฺขา\-สหคตํ อกุสลํ ธมฺมํ ปฏิจฺจ อุเปกฺขา\-สหคโต อกุสโล ธมฺโม อุปฺปชฺชติ เหตุปจฺจยา. (1)}

\prose{ปีติสหคตํ อกุสลญฺจ สุขสหคตํ อกุสลญฺจ ธมฺมํ ปฏิจฺจ ปีติสหคโต อกุสโล ธมฺโม อุปฺปชฺชติ เหตุปจฺจยา. ตีณิ. (สพฺพตฺถ ทส. สพฺพตฺถ วิตฺถาโร.)}

\proseitem{101}{ปีติสหคตํ อพฺยากตํ ธมฺมํ ปฏิจฺจ ปีติสหคโต อพฺยากโต ธมฺโม อุปฺปชฺชติ เหตุปจฺจยา. ตีณิ.}

\prose{สุขสหคตํ อพฺยากตํ ธมฺมํ ปฏิจฺจ สุขสหคโต อพฺยากโต ธมฺโม อุปฺปชฺชติ เหตุปจฺจยา. ตีณิ.}

\prose{อุเปกฺขา\-สหคตํ อพฺยากตํ ธมฺมํ ปฏิจฺจ อุเปกฺขา\-สหคโต อพฺยากโต ธมฺโม อุปฺปชฺชติ เหตุปจฺจยา. (1)}

\prose{ปีติสหคตํ อพฺยากตญฺจ สุขสหคตํ อพฺยากตญฺจ ธมฺมํ ปฏิจฺจ ปีติสหคโต อพฺยากโต ธมฺโม อุปฺปชฺชติ เหตุปจฺจยา. ตีณิ. (สพฺพตฺถ ทส. สพฺพตฺถ วิตฺถาโร.)}

\csromanmarkh{8. ทสฺสนตฺติก}\csromanmarkhii{1. กุสลตฺติก}\csromanheadernomark{2}{8. \textbf{ทสฺสนตฺติก 1. กุสลตฺติก}}
\prose{\csromanfolio{30}1. ปฏิจฺจวารเหตุ}

\proseitem{102}{เนวทสฺสเนน นภาวนาย ปหาตพฺพํ กุสลํ ธมฺมํ ปฏิจฺจ เนวทสฺสเนน นภาวนาย ปหาตพฺโพ กุสโล ธมฺโม อุปฺปชฺชติ เหตุปจฺจยา. (สพฺพตฺถ เอกํ.)}

\proseitem{103}{ทสฺสเนน ปหาตพฺพํ อกุสลํ ธมฺมํ ปฏิจฺจ ทสฺสเนน ปหาตพฺโพ อกุสโล ธมฺโม อุปฺปชฺชติ เหตุปจฺจยา. (1)}

\prose{ภาวนาย ปหาตพฺพํ อกุสลํ ธมฺมํ ปฏิจฺจ ภาวนาย ปหาตพฺโพ อกุสโล ธมฺโม อุปฺปชฺชติ เหตุปจฺจยา. (1) (สพฺพตฺถ เทฺว. สพฺพตฺถ วิตฺถาโร.)}

\proseitem{104}{เนวทสฺสเนน นภาวนาย ปหาตพฺพํ อพฺยากตํ ธมฺมํ ปฏิจฺจ เนวทสฺสเนน นภาวนาย ปหาตพฺโพ อพฺยากโต ธมฺโม อุปฺปชฺชติ เหตุปจฺจยา. (สพฺพตฺถ เอกํ.)}

\csromanmarkh{9. ทสฺสน\-เหตุ\-ตฺติก}\csromanmarkhii{1. กุสลตฺติก}\csromanheadernomark{2}{9. ทสฺสน\-เหตุ\-ตฺติก 1. กุสลตฺติก}
\prose{1. ปฏิจฺจวารเหตุ}

\proseitem{105}{เนวทสฺสเนน นภาวนาย ปหาตพฺพ\-เหตุกํ กุสลํ ธมฺมํ ปฏิจฺจ เนวทสฺสเนน นภาวนาย ปหาตพฺพ\-เหตุโก กุสโล ธมฺโม อุปฺปชฺชติ เหตุปจฺจยา. (สพฺพตฺถ เอกํ. สพฺพตฺถ วิตฺถาโร.)}

\proseitem{106}{ทสฺสเนน ปหาตพฺพ\-เหตุกํ อกุสลํ ธมฺมํ ปฏิจฺจ ทสฺสเนน ปหาตพฺพ\-เหตุโก อกุสโล ธมฺโม อุปฺปชฺชติ เหตุปจฺจยา.}

\prose{ภาวนาย ปหาตพฺพ\-เหตุกํ อกุสลํ ธมฺมํ ปฏิจฺจ ภาวนาย ปหาตพฺพ\-เหตุโก อกุสโล ธมฺโม อุปฺปชฺชติ เหตุปจฺจยา. (สํขิตฺตํ.)}

\prose{เหตุยา ฉ, อารมฺมเณ ทส, อธิปติยา เทฺว ฯเปฯ อวิคเต ทส. (สพฺพตฺถ วิตฺถาโร.)}

\csromanmarkh{11. เสกฺขตฺติก}\csromanmarkhii{1. กุสลตฺติก}\csromanheadernomark{2}{11. เสกฺขตฺติก 1. กุสลตฺติก}
\proseitem{107}{\csromanfolio{31}เนวทสฺสเนน นภาวนาย ปหาตพฺพ\-เหตุกํ อพฺยากตํ ธมฺมํ ปฏิจฺจ เนวทสฺสเนน นภาวนาย ปหาตพฺพ\-เหตุโก อพฺยากโต ธมฺโม อุปฺปชฺชติ เหตุปจฺจยา. (สพฺพตฺถ เอกํ.)}

\csromanmarkh{10. อาจยคามิตฺติก}\csromanmarkhii{1. กุสลตฺติก}\csromanheadernomark{2}{10. \textbf{อาจยคามิตฺติก 1. กุสลตฺติก}}
\prose{1. ปฏิจฺจวารเหตุ}

\prose{108. อาจยคามิํ กุสลํ ธมฺมํ ปฏิจฺจ อาจยคามี กุสโล}

\prose{อปจยคามิํ กุสลํ ธมฺมํ ปฏิจฺจ อปจยคามี กุสโล ธมฺโม อุปฺปชฺชติ เหตุปจฺจยา. (1) (สพฺพตฺถ เทฺว. สพฺพตฺถ วิตฺถาโร.)}

\prose{อาจยคามิํ อกุสลํ ธมฺมํ ปฏิจฺจ อาจยคามี อกุสโล ธมฺโม อุปฺปชฺชติ เหตุปจฺจยา. (สพฺพตฺถ เอกํ.)}

\prose{เนวา\-จยคามิ\-นา\-ปจยคามิํ อพฺยากตํ ธมฺมํ ปฏิจฺจ เนวา\-จยคามิ\-นา\-ปจยคามี อพฺยากโต ธมฺโม อุปฺปชฺชติ เหตุปจฺจยา. (สพฺพตฺถ เอกํ.)}

\csromanmarkh{11. เสกฺขตฺติก}\csromanmarkhii{1. กุสลตฺติก}\csromanheadernomark{2}{11. เสกฺขตฺติก 1. กุสลตฺติก}
\prose{1. ปฏิจฺจวารเหตุ}

\proseitem{109}{เสกฺขํ กุสลํ ธมฺมํ ปฏิจฺจ เสกฺโข กุสโล ธมฺโม อุปฺปชฺชติ เหตุปจฺจยา. (1)}

\prose{เนว\-เสกฺข\-นาเสกฺขํ กุสลํ ธมฺมํ ปฏิจฺจ เนว\-เสกฺข\-นา\-เสกฺโข กุสโล ธมฺโม อุปฺปชฺชติ เหตุปจฺจยา. (1) (สพฺพตฺถ เทฺว. สพฺพตฺถ วิตฺถาโร.)}

\prose{เนว\-เสกฺข\-นาเสกฺขํ อกุสลํ ธมฺมํ ปฏิจฺจ เนว\-เสกฺข\-นา\-เสกฺโข อกุสโล ธมฺโม อุปฺปชฺชติ เหตุปจฺจยา. (สพฺพตฺถ เอกํ.)}

\prose{เสกฺขํ อพฺยากตํ ธมฺมํ ปฏิจฺจ เสกฺโข อพฺยากโต ธมฺโม อุปฺปชฺชติ เหตุปจฺจยา. (สํขิตฺตํ.)}

\prose{เหตุยา นว. (สพฺพตฺถ วิตฺถาเรตพฺพํ.)}

\csromanmarkh{12. ปริตฺตติก}\csromanmarkhii{1. กุสลตฺติก}\csromanheadernomark{2}{12. \textbf{ปริตฺตติก 1. กุสลตฺติก}}
\prose{\csromanfolio{32}1. ปฏิจฺจวารเหตุ}

\proseitem{110}{ปริตฺตํ กุสลํ ธมฺมํ ปฏิจฺจ ปริตฺโต กุสโล ธมฺโม อุปฺปชฺชติ เหตุปจฺจยา. (1)}

\prose{มหคฺคตํ กุสลํ ธมฺมํ ปฏิจฺจ มหคฺคโต กุสโล ธมฺโม อุปฺปชฺชติ เหตุปจฺจยา. (1)}

\prose{(สพฺพตฺถ ตีณิ. สพฺพตฺถ วิตฺถาโร.)}

\proseitem{111}{ปริตฺตํ อกุสลํ ธมฺมํ ปฏิจฺจ ปริตฺโต อกุสโล ธมฺโม อุปฺปชฺชติ เหตุปจฺจยา. (สพฺพตฺถ เอกํ.)}

\prose{ปริตฺตํ อพฺยากตํ ธมฺมํ ปฏิจฺจ ปริตฺโต อพฺยากโต ธมฺโม อุปฺปชฺชติ เหตุปจฺจยา. ตีณิ.}

\prose{มหคฺคตํ อพฺยากตํ ธมฺมํ ปฏิจฺจ มหคฺคโต อพฺยากโต ธมฺโม อุปฺปชฺชติ เหตุปจฺจยา. ตีณิ.}

\prose{อปฺปมาณํ อพฺยากตํ ธมฺมํ ปฏิจฺจ อปฺปมาโณ อพฺยากโต ธมฺโม อุปฺปชฺชติ เหตุปจฺจยา. (สํขิตฺตํ.)}

\prose{เหตุยา เตรส.}

\csromanmarkh{13. ปริตฺตา\-รมฺมณ\-ตฺติก}\csromanmarkhii{1. กุสลตฺติก}\csromanheadernomark{2}{13. ปริตฺตา\-รมฺมณ\-ตฺติก 1. กุสลตฺติก}
\prose{1. ปฏิจฺจวารเหตุ}

\proseitem{112}{ปริตฺตา\-รมฺมณํ กุสลํ ธมฺมํ ปฏิจฺจ ปริตฺตารมฺมโณ กุสโล ธมฺโม อุปฺปชฺชติ เหตุปจฺจยา. (1)}

\prose{มหคฺคตา\-รมฺมณํ กุสลํ ธมฺมํ ปฏิจฺจ มหคฺคตา\-รมฺมโณ กุสโล ธมฺโม อุปฺปชฺชติ เหตุปจฺจยา. (1)}

\csromanmarkh{14. หีนตฺติก}\csromanmarkhii{1. กุสลตฺติก}\csromanheadernomark{2}{14. หีนตฺติก 1. กุสลตฺติก}
\prose{\csromanfolio{33}อปฺปมาณา\-รมฺมณํ กุสลํ ธมฺมํ ปฏิจฺจ อปฺปมาณารมฺมโณ กุสโล ธมฺโม อุปฺปชฺชติ เหตุปจฺจยา. (1) (สพฺพตฺถ ตีณิ. สพฺพตฺถ วิตฺถาโร.)}

\proseitem{113}{ปริตฺตา\-รมฺมณํ อกุสลํ ธมฺมํ ปฏิจฺจ ปริตฺตารมฺมโณ อกุสโล ธมฺโม อุปฺปชฺชติ เหตุปจฺจยา. (1)}

\prose{มหคฺคตา\-รมฺมณํ อกุสลํ ธมฺมํ ปฏิจฺจ มหคฺคตา\-รมฺมโณ อกุสโล ธมฺโม อุปฺปชฺชติ เหตุปจฺจยา. (1) (สพฺพตฺถ เทฺว. สพฺพตฺถ วิตฺถาโร.)}

\proseitem{114}{ปริตฺตา\-รมฺมณํ อพฺยากตํ ธมฺมํ ปฏิจฺจ ปริตฺตารมฺมโณ อพฺยากโต ธมฺโม อุปฺปชฺชติ เหตุปจฺจยา. (1)}

\prose{มหคฺคตา\-รมฺมณํ อพฺยากตํ ธมฺมํ ปฏิจฺจ มหคฺคตา\-รมฺมโณ อพฺยากโต ธมฺโม อุปฺปชฺชติ เหตุปจฺจยา. (1)}

\prose{อปฺปมาณา\-รมฺมณํ อพฺยากตํ ธมฺมํ ปฏิจฺจ อปฺปมาณารมฺมโณ อพฺยากโต ธมฺโม อุปฺปชฺชติ เหตุปจฺจยา. (1) (สพฺพตฺถ ตีณิ. สพฺพตฺถ วิตฺถาโร.)}

\csromanmarkh{14. หีนตฺติก}\csromanmarkhii{1. กุสลตฺติก}\csromanheadernomark{2}{14. หีนตฺติก 1. กุสลตฺติก}
\prose{1. ปฏิจฺจวารเหตุ}

\proseitem{115}{มชฺฌิมํ กุสลํ ธมฺมํ ปฏิจฺจ มชฺฌิโม กุสโล ธมฺโม อุปฺปชฺชติ เหตุปจฺจยา. (1)}

\prose{(สพฺพตฺถ เทฺว. สพฺพตฺถ วิตฺถาโร.)}

\proseitem{116}{หีนํ อกุสลํ ธมฺมํ ปฏิจฺจ หีโน อกุสโล ธมฺโม อุปฺปชฺชติ เหตุปจฺจยา. (สพฺพตฺถ เอกํ.)}

\prose{117. มชฺฌิมํ อพฺยากตํ ธมฺมํ ปฏิจฺจ มชฺฌิโม อพฺยากโต}

\prose{ปณีตํ อพฺยากตํ ธมฺมํ ปฏิจฺจ ปณีโต อพฺยากโต ธมฺโม อุปฺปชฺชติ เหตุปจฺจยา. ตีณิ.}

\prose{\csromanfolio{34}มชฺฌิมํ อพฺยากตญฺจ ปณีตํ อพฺยากตญฺจ ธมฺมํ ปฏิจฺจ มชฺฌิโม อพฺยากโต ธมฺโม อุปฺปชฺชติ เหตุปจฺจยา. (1) (สํขิตฺตํ.)}

\prose{เหตุยา ปญฺจ. (สพฺพตฺถ วิตฺถาโร.)}

\csromanmarkh{15. มิจฺฉตฺตตฺติก}\csromanmarkhii{1. กุสลตฺติก}\csromanheadernomark{2}{15. \textbf{มิจฺฉตฺตตฺติก 1. กุสลตฺติก}}
\prose{1. ปฏิจฺจวารเหตุ}

\proseitem{118}{สมฺมตฺต\-นิยตํ กุสลํ ธมฺมํ ปฏิจฺจ สมฺมตฺตนิยโต กุสโล ธมฺโม อุปฺปชฺชติ เหตุปจฺจยา. (1)}

\prose{(สพฺพตฺถ เทฺว. สพฺพตฺถ วิตฺถาโร.)}

\proseitem{119}{มิจฺฉตฺต\-นิยตํ อกุสลํ ธมฺมํ ปฏิจฺจ มิจฺฉตฺต\-นิยโต อกุสโล ธมฺโม อุปฺปชฺชติ เหตุปจฺจยา. (1)}

\prose{(สพฺพตฺถ เทฺว. สพฺพตฺถ วิตฺถาโร.)}

\prose{อนิยตํ อพฺยากตํ ธมฺมํ ปฏิจฺจ อนิยโต อพฺยากโต ธมฺโม อุปฺปชฺชติ เหตุปจฺจยา. (สพฺพตฺถ เอกํ.)}

\csromanmarkh{16. มคฺคา\-รมฺมณ\-ตฺติก}\csromanmarkhii{1. กุสลตฺติก}\csromanheadernomark{2}{16. มคฺคา\-รมฺมณ\-ตฺติก 1. กุสลตฺติก}
\prose{1. ปฏิจฺจวารเหตุ}

\proseitem{120}{มคฺคารมฺมณํ กุสลํ ธมฺมํ ปฏิจฺจ มคฺคารมฺมโณ กุสโล ธมฺโม อุปฺปชฺชติ เหตุปจฺจยา. ตีณิ.}

\prose{มคฺคเหตุกํ กุสลํ ธมฺมํ ปฏิจฺจ มคฺคเหตุโก กุสโล ธมฺโม อุปฺปชฺชติ เหตุปจฺจยา. ตีณิ.}

\prose{มคฺคาธิปติํ กุสลํ ธมฺมํ ปฏิจฺจ มคฺคาธิปติ กุสโล ธมฺโม อุปฺปชฺชติ เหตุปจฺจยา. ปญฺจ.}

\csromanmarkh{18. อตีตตฺติก}\csromanmarkhii{1. กุสลตฺติก}\csromanheadernomark{2}{18. อตีตตฺติก 1. กุสลตฺติก}
\prose{\csromanfolio{35}มคฺคารมฺมณํ กุสลญฺจ มคฺคาธิปติํ กุสลญฺจ ธมฺมํ ปฏิจฺจ มคฺคารมฺมโณ กุสโล ธมฺโม อุปฺปชฺชติ เหตุปจฺจยา. ตีณิ.}

\prose{มคฺคเหตุกํ กุสลญฺจ มคฺคาธิปติํ กุสลญฺจ ธมฺมํ ปฏิจฺจ มคฺคเหตุโก กุสโล ธมฺโม อุปฺปชฺชติ เหตุปจฺจยา. ตีณิ. (สํขิตฺตํ.)}

\prose{เหตุยา สตฺตรส ฯเปฯ อวิคเต สตฺตรส. (สํขิตฺตํ.)}

\proseitem{121}{มคฺคารมฺมณํ อพฺยากตํ ธมฺมํ ปฏิจฺจ มคฺคารมฺมโณ อพฺยากโต ธมฺโม อุปฺปชฺชติ เหตุปจฺจยา. ตีณิ.}

\prose{มคฺคาธิปติํ อพฺยากตํ ธมฺมํ ปฏิจฺจ มคฺคาธิปติ อพฺยากโต ธมฺโม อุปฺปชฺชติ เหตุปจฺจยา. ตีณิ.}

\prose{(สพฺพตฺถ วิตฺถาโร.)}

\csromanmarkh{17. อุปฺปนฺนตฺติก}\csromanmarkhii{1. กุสลตฺติก}\csromanheadernomark{2}{17. \textbf{อุปฺปนฺนตฺติก 1. กุสลตฺติก}}
\prose{7. ปญฺหาวารเหตุ}

\proseitem{122}{อุปฺปนฺโน กุสโล ธมฺโม อุปฺปนฺนสฺส กุสลสฺส ธมฺมสฺส เหตุปจฺจเยน ปจฺจโย. (สพฺพตฺถ เอกํ. สพฺพตฺถ วิตฺถาโร.)}

\prose{อุปฺปนฺโน อกุสโล ธมฺโม อุปฺปนฺนสฺส อกุสลสฺส ธมฺมสฺส เหตุปจฺจเยน ปจฺจโย. (สพฺพตฺถ เอกํ. สพฺพตฺถ วิตฺถาโร.)}

\prose{อุปฺปนฺโน อพฺยากโต ธมฺโม อุปฺปนฺนสฺส อพฺยากตสฺส ธมฺมสฺส เหตุปจฺจเยน ปจฺจโย. (สํขิตฺตํ.)}

\prose{เหตุยา เอกํ, อารมฺมเณ ตีณิ ฯเปฯ อุปนิสฺสเย ตีณิ ฯเปฯ อวิคเต เอกํ. (สพฺพตฺถ วิตฺถาโร.)}

\csromanmarkh{18. อตีตตฺติก}\csromanmarkhii{1. กุสลตฺติก}\csromanheadernomark{2}{18. อตีตตฺติก 1. กุสลตฺติก}
\prose{7. ปญฺหาวารเหตุ}

\proseitem{123}{ปจฺจุปฺปนฺโน กุสโล ธมฺโม ปจฺจุปฺปนฺนสฺส กุสลสฺส ธมฺมสฺส เหตุปจฺจเยน ปจฺจโย. (สพฺพตฺถ เอกํ.)}

\prose{\csromanfolio{36}ปจฺจุปฺปนฺโน อกุสโล ธมฺโม ปจฺจุปฺปนฺนสฺส อกุสลสฺส ธมฺมสฺส เหตุปจฺจเยน ปจฺจโย. (สพฺพตฺถ เอกํ. สพฺพตฺถ วิตฺถาโร.)}

\prose{ปจฺจุปฺปนฺโน อพฺยากโต ธมฺโม ปจฺจุปฺปนฺนสฺส อพฺยากตสฺส ธมฺมสฺส เหตุปจฺจเยน ปจฺจโย.}

\prose{เหตุยา เอกํ. (สพฺพตฺถ วิตฺถาโร.)}

\csromanmarkh{19. อตีตา\-รมฺมณ\-ตฺติก}\csromanmarkhii{1. กุสลตฺติก}\csromanheadernomark{2}{19. อตีตา\-รมฺมณ\-ตฺติก 1. กุสลตฺติก}
\prose{1. ปฏิจฺจวารเหตุ}

\prose{124. อตีตารมฺมณํ กุสลํ ธมฺมํ ปฏิจฺจ อตีตารมฺมโณ กุสโล}

\prose{อนาคตา\-รมฺมณํ กุสลํ ธมฺมํ ปฏิจฺจ อนาคตารมฺมโณ กุสโล}

\prose{ปจฺจุปฺปนฺนา\-รมฺมณํ กุสลํ ธมฺมํ ปฏิจฺจ ปจฺจุปฺปนฺนา\-รมฺมโณ กุสโล ธมฺโม อุปฺปชฺชติ เหตุปจฺจยา. (1) (สพฺพตฺถ ตีณิ. สพฺพตฺถ วิตฺถาโร.)}

\proseitem{125}{อตีตารมฺมณํ อกุสลํ ธมฺมํ ปฏิจฺจ อตีตารมฺมโณ อกุสโล ธมฺโม อุปฺปชฺชติ เหตุปจฺจยา. (1)}

\prose{อนาคตา\-รมฺมณํ อกุสลํ ธมฺมํ ปฏิจฺจ อนาคตารมฺมโณ อกุสโล ธมฺโม อุปฺปชฺชติ เหตุปจฺจยา. (1)}

\prose{ปจฺจุปฺปนฺนา\-รมฺมณํ อกุสลํ ธมฺมํ ปฏิจฺจ ปจฺจุปฺปนฺนา\-รมฺมโณ อกุสโล ธมฺโม อุปฺปชฺชติ เหตุปจฺจยา. (1) (สพฺพตฺถ ตีณิ. สพฺพตฺถ วิตฺถาโร.)}

\proseitem{126}{อตีตารมฺมณํ อพฺยากตํ ธมฺมํ ปฏิจฺจ อตีตารมฺมโณ อพฺยากโต ธมฺโม อุปฺปชฺชติ เหตุปจฺจยา. (1)}

\prose{อนาคตา\-รมฺมณํ อพฺยากตํ ธมฺมํ ปฏิจฺจ อนาคตารมฺมโณ อพฺยากโต ธมฺโม อุปฺปชฺชติ เหตุปจฺจยา. (1)}

\prose{ปจฺจุปฺปนฺนา\-รมฺมณํ อพฺยากตํ ธมฺมํ ปฏิจฺจ ปจฺจุปฺปนฺนา\-รมฺมโณ อพฺยากโต ธมฺโม อุปฺปชฺชติ เหตุปจฺจยา. (1) (สพฺพตฺถ ตีณิ. สพฺพตฺถ วิตฺถาโร.)}

\vspace{\tipitakaparskip}
\csromancenter{สนิทสฺสน\-ตฺติก 1. กุสลตฺติก 20. อชฺฌตฺตตฺติก 1. กุสลตฺติก}
\prose{\csromanfolio{37}1. ปฏิจฺจวารเหตุ}

\proseitem{127}{อชฺฌตฺตํ กุสลํ ธมฺมํ ปฏิจฺจ อชฺฌตฺโต กุสโล ธมฺโม อุปฺปชฺชติ เหตุปจฺจยา. (1)}

\prose{(สพฺพตฺถ เทฺว. สพฺพตฺถ วิตฺถาโร.)}

\proseitem{128}{อชฺฌตฺตํ อกุสลํ ธมฺมํ ปฏิจฺจ อชฺฌตฺโต อกุสโล ธมฺโม อุปฺปชฺชติ เหตุปจฺจยา. (1)}

\prose{(สพฺพตฺถ เทฺว. สพฺพตฺถ วิตฺถาโร.)}

\proseitem{129}{อชฺฌตฺตํ อพฺยากตํ ธมฺมํ ปฏิจฺจ อชฺฌตฺโต อพฺยากโต ธมฺโม อุปฺปชฺชติ เหตุปจฺจยา. (1)}

\prose{พหิทฺธา อพฺยากตํ ธมฺมํ ปฏิจฺจ พหิทฺธา อพฺยากโต ธมฺโม อุปฺปชฺชติ เหตุปจฺจยา. (1) (สพฺพตฺถ เทฺว. สพฺพตฺถ วิตฺถาโร.)}

\csromanmarkh{21. อชฺฌตฺตา\-รมฺมณ\-ตฺติก}\csromanmarkhii{1. กุสลตฺติก}\csromanheadernomark{2}{21. อชฺฌตฺตา\-รมฺมณ\-ตฺติก 1. กุสลตฺติก}
\prose{1. ปฏิจฺจวารเหตุ}

\proseitem{130}{อชฺฌตฺตา\-รมฺมณํ กุสลํ ธมฺมํ ปฏิจฺจ อชฺฌตฺตา\-รมฺมโณ กุสโล ธมฺโม อุปฺปชฺชติ เหตุปจฺจยา. (สพฺพตฺถ เทฺว. สพฺพตฺถ วิตฺถาโร.)}

\prose{อชฺฌตฺตา\-รมฺมณํ อกุสลํ ธมฺมํ ปฏิจฺจ อชฺฌตฺตา\-รมฺมโณ อกุสโล ธมฺโม อุปฺปชฺชติ เหตุปจฺจยา. (สพฺพตฺถ เทฺว. สพฺพตฺถ วิตฺถาโร.)}

\prose{อชฺฌตฺตา\-รมฺมณํ อพฺยากตํ ธมฺมํ ปฏิจฺจ อชฺฌตฺตา\-รมฺมโณ อพฺยากโต ธมฺโม อุปฺปชฺชติ เหตุปจฺจยา. (สพฺพตฺถ เทฺว. สพฺพตฺถ วิตฺถาโร.)}

\csromanmarkh{22. สนิทสฺสน\-ตฺติก}\csromanmarkhii{1. กุสลตฺติก}\csromanheadernomark{2}{22. สนิทสฺสน\-ตฺติก 1. กุสลตฺติก}
\prose{1. ปฏิจฺจวารเหตุ}

\proseitem{131}{อนิทสฺสน\-อปฺปฏิฆํ กุสลํ ธมฺมํ ปฏิจฺจ อนิทสฺสน\-อปฺปฏิโฆ กุสโล ธมฺโม อุปฺปชฺชติ เหตุปจฺจยา. (สพฺพตฺถ เอกํ.)}

\prose{\csromanfolio{38}อนิทสฺสน\-อปฺปฏิฆํ อกุสลํ ธมฺมํ ปฏิจฺจ อนิทสฺสน\-อปฺปฏิโฆ อกุสโล ธมฺโม อุปฺปชฺชติ เหตุปจฺจยา. (สพฺพตฺถ เอกํ.)}

\proseitem{132}{อนิทสฺสน\-สปฺปฏิฆํ อพฺยากตํ ธมฺมํ ปฏิจฺจ อนิทสฺสน\-สปฺปฏิโฆ อพฺยากโต ธมฺโม อุปฺปชฺชติ เหตุปจฺจยา. (เอกํ.) . อนิทสฺสน\-สปฺปฏิฆํ อพฺยากตํ ธมฺมํ ปฏิจฺจ สนิทสฺสน\-สปฺปฏิโฆ อพฺยากโต ธมฺโม อุปฺปชฺชติ เหตุปจฺจยา. (เทฺว.) . อนิทสฺสน\-สปฺปฏิฆํ อพฺยากตํ ธมฺมํ ปฏิจฺจ อนิทสฺสน\-อปฺปฏิโฆ อพฺยากโต ธมฺโม อุปฺปชฺชติ เหตุปจฺจยา. (ตีณิ.) . อนิทสฺสน\-สปฺปฏิฆํ อพฺยากตํ ธมฺมํ ปฏิจฺจ สนิทสฺสน\-สปฺปฏิโฆ อพฺยากโต จ อนิทสฺสน\-อปฺปฏิโฆ อพฺยากโต จ ธมฺมา อุปฺปชฺชนฺติ เหตุปจฺจยา. (จตฺตาริ.) . อนิทสฺสน\-สปฺปฏิฆํ อพฺยากตํ ธมฺมํ ปฏิจฺจ อนิทสฺสน\-สปฺปฏิโฆ อพฺยากโต จ อนิทสฺสน\-อปฺปฏิโฆ อพฺยากโต จ ธมฺมา อุปฺปชฺชนฺติ เหตุปจฺจยา. (ปญฺจ.) . อนิทสฺสน\-สปฺปฏิฆํ อพฺยากตํ ธมฺมํ ปฏิจฺจ สนิทสฺสน\-สปฺปฏิโฆ อพฺยากโต จ อนิทสฺสน\-สปฺปฏิโฆ อพฺยากโต จ ธมฺมา อุปฺปชฺชนฺติ เหตุปจฺจยา. (ฉ.) . อนิทสฺสน\-สปฺปฏิฆํ อพฺยากตํ ธมฺมํ ปฏิจฺจ สนิทสฺสน\-สปฺปฏิโฆ อพฺยากโต จ อนิทสฺสน\-สปฺปฏิโฆ อพฺยากโต จ อนิทสฺสน\-อปฺปฏิโฆ อพฺยากโต จ ธมฺมา อุปฺปชฺชนฺติ เหตุปจฺจยา. (สตฺต.)}

\proseitem{133}{อนิทสฺสน\-อปฺปฏิฆํ อพฺยากตํ ธมฺมํ ปฏิจฺจ อนิทสฺสน\-อปฺปฏิโฆ อพฺยากโต ธมฺโม อุปฺปชฺชติ เหตุปจฺจยา. (สตฺต.) . อนิทสฺสน\-อปฺปฏิฆํ อพฺยากตญฺจ อนิทสฺสน\-สปฺปฏิฆํ อพฺยากตญฺจ ธมฺมํ ปฏิจฺจ สนิทสฺสน\-สปฺปฏิโฆ อพฺยากโต ธมฺโม อุปฺปชฺชติ เหตุปจฺจยา. (สตฺต.) (สํขิตฺตํ.)}

\proseitemwithcloser{134}{เหตุยา เอกวีส, อวิคเต เอกวีส. (สํขิตฺตํ.) (สหชาตวารมฺปิ ฯเปฯ ปญฺหาวารมฺปิ วิตฺถาเรตพฺพํ.)}{ธมฺมานุโลเม ติกติก\-ปฏฺฐานํ นิฏฺฐิตํ.}
\csromanpalirecto
\pitaka{อภิธมฺม\-ปิฏก}
\csromanheader{2}{\textbf{ธมฺมานุโลม}}
\csromanmatikaanchor{mtk.1}
\csromanmatikaanchor{mtk.3}
\csromantocmarknopage{section}{ปญฺจมภาค}{mtk.1}
\bookmark[dest={mtk.1},level=1]{ปญฺจมภาค}
\gambhira{ทุกทุก\-ปฏฺฐาน\-ปาฬิ}
\namakkaram{นโม ตสฺส ภควโต อรหโต สมฺมา\-สมฺพุทฺธสฺส.}
\csromanmarkh{1. เหตุทุก}\csromanmarkhii{1. สเหตุกทุก}\csromanheadernomark{2}{1. \textbf{เหตุทุก 1. สเหตุกทุก}}
\csromanheader{2}{\textbf{สเหตุกปท 1-7. ปฏิจฺจาทิวาร}}
\prose{\csromanfolio{39}ปจฺจยจตุกฺกเหตุ}

\proseitem{1}{เหตุํ สเหตุกํ ธมฺมํ ปฏิจฺจ เหตุ สเหตุโก ธมฺโม อุปฺปชฺชติ เหตุปจฺจยา.  เหตุํ สเหตุกํ ธมฺมํ ปฏิจฺจ นเหตุ สเหตุโก ธมฺโม อุปฺปชฺชติ เหตุปจฺจยา.  เหตุํ สเหตุกํ ธมฺมํ ปฏิจฺจ เหตุ สเหตุโก จ นเหตุ สเหตุโก จ ธมฺมา อุปฺปชฺชนฺติ เหตุปจฺจยา. (3)}

\proseitem{2}{นเหตุํ สเหตุกํ ธมฺมํ ปฏิจฺจ นเหตุ สเหตุโก ธมฺโม อุปฺปชฺชติ เหตุปจฺจยา.  นเหตุํ สเหตุกํ ธมฺมํ ปฏิจฺจ เหตุ สเหตุโก ธมฺโม อุปฺปชฺชติ เหตุปจฺจยา.  นเหตุํ สเหตุกํ ธมฺมํ ปฏิจฺจ เหตุ สเหตุโก จ นเหตุ สเหตุโก จ ธมฺมา อุปฺปชฺชนฺติ เหตุปจฺจยา. (3)}

\proseitem{3}{เหตุํ สเหตุกญฺจ นเหตุํ สเหตุกญฺจ ธมฺมํ ปฏิจฺจ เหตุ สเหตุโก ธมฺโม อุปฺปชฺชติ เหตุปจฺจยา.  เหตุํ สเหตุกญฺจ นเหตุํ สเหตุกญฺจ ธมฺมํ ปฏิจฺจ นเหตุ สเหตุโก ธมฺโม อุปฺปชฺชติ \csromanfolio{40}เหตุปจฺจยา.  เหตุํ สเหตุกญฺจ นเหตุํ สเหตุกญฺจ ธมฺมํ ปฏิจฺจ เหตุ สเหตุโก จ นเหตุ สเหตุโก จ ธมฺมา อุปฺปชฺชนฺติ เหตุปจฺจยา. (3)}

\proseitem{4}{เหตุยา นว, อารมฺมเณ นว ฯเปฯ อวิคเต นว. (สํขิตฺตํ.)}

\prose{ปจฺจนียนอธิปติ}

\proseitem{5}{เหตุํ สเหตุกํ ธมฺมํ ปฏิจฺจ เหตุ สเหตุโก ธมฺโม อุปฺปชฺชติ นอธิปติ\-ปจฺจยา. (สํขิตฺตํ.)}

\prose{นอธิปติยา นว, นปุเรชาเต นว, นปจฺฉาชาเต นว, นกมฺเม ตีณิ, นวิปาเก นว, นวิปฺปยุตฺเต นว. (สํขิตฺตํ.)}

\prose{เหตุปจฺจยา นอธิปติยา นว. (สํขิตฺตํ.)}

\prose{นอธิปติ\-ปจฺจยา เหตุยา นว. (สํขิตฺตํ.)}

\prose{(สหชาตวารมฺปิ ปจฺจยวารมฺปิ นิสฺสยวารมฺปิ สํสฏฺฐ\-วารมฺปิ สมฺปยุตฺตวารมฺปิ ปฏิจฺจ\-วาร\-สทิสํ.)}

\csromanheader{2}{\textbf{เหตุ อารมฺมณ}}
\proseitem{6}{เหตุ สเหตุโก ธมฺโม เหตุสฺส สเหตุกสฺส ธมฺมสฺส เหตุปจฺจเยน ปจฺจโย.  เหตุ สเหตุโก ธมฺโม นเหตุสฺส สเหตุกสฺส ธมฺมสฺส เหตุปจฺจเยน ปจฺจโย.  เหตุ สเหตุโก ธมฺโม เหตุสฺส สเหตุกสฺส จ นเหตุสฺส สเหตุกสฺส จ ธมฺมสฺส เหตุปจฺจเยน ปจฺจโย. (3)}

\proseitem{7}{เหตุ สเหตุโก ธมฺโม เหตุสฺส สเหตุกสฺส ธมฺมสฺส อารมฺมณ\-ปจฺจเยน ปจฺจโย.  เหตุ สเหตุโก ธมฺโม นเหตุสฺส สเหตุกสฺส ธมฺมสฺส อารมฺมณ\-ปจฺจเยน ปจฺจโย.  เหตุ สเหตุโก ธมฺโม เหตุสฺส สเหตุกสฺส จ นเหตุสฺส สเหตุกสฺส จ ธมฺมสฺส อารมฺมณ\-ปจฺจเยน ปจฺจโย. (3)}

\prose{นเหตุ สเหตุโก ธมฺโม นเหตุสฺส สเหตุกสฺส ธมฺมสฺส อารมฺมณ\-ปจฺจเยน ปจฺจโย.  นเหตุ สเหตุโก ธมฺโม เหตุสฺส สเหตุกสฺส ธมฺมสฺส อารมฺมณ\-ปจฺจเยน ปจฺจโย.  นเหตุ สเหตุโก ธมฺโม เหตุสฺส สเหตุกสฺส จ นเหตุสฺส สเหตุกสฺส จ ธมฺมสฺส อารมฺมณ\-ปจฺจเยน ปจฺจโย. (3)}

\csromanmarkh{1. เหตุทุก}\csromanmarkhii{1. สเหตุกทุก}\csromanheadernomark{2}{1. เหตุทุก 1. สเหตุกทุก}
\prose{\csromanfolio{41}เหตุ สเหตุโก จ นเหตุ สเหตุโก จ ธมฺมา เหตุสฺส สเหตุกสฺส ธมฺมสฺส อารมฺมณ\-ปจฺจเยน ปจฺจโย.  เหตุ สเหตุโก จ นเหตุ สเหตุโก จ ธมฺมา นเหตุสฺส สเหตุกสฺส ธมฺมสฺส อารมฺมณ\-ปจฺจเยน ปจฺจโย.  เหตุ สเหตุโก จ นเหตุ สเหตุโก จ ธมฺมา เหตุสฺส สเหตุกสฺส จ นเหตุสฺส สเหตุกสฺส จ ธมฺมสฺส อารมฺมณ\-ปจฺจเยน ปจฺจโย. (3) (สํขิตฺตํ.)}

\proseitem{8}{เหตุยา ตีณิ, อารมฺมเณ นว ฯเปฯ อุปนิสฺสเย นว, อวิคเต นว.}

\vspace{\tipitakaparskip}
\csromancenter{(ยถา กุสลตฺติเก ปญฺหาวารํ, เอวํ วิตฺถาเรตพฺพํ.)}
\csromanheader{2}{\textbf{อเหตุกปท ปจฺจยจตุกฺก}}
\proseitem{9}{เหตุํ อเหตุกํ ธมฺมํ ปฏิจฺจ นเหตุ อเหตุโก ธมฺโม อุปฺปชฺชติ เหตุปจฺจยา. (1)}

\prose{นเหตุํ อเหตุกํ ธมฺมํ ปฏิจฺจ นเหตุ อเหตุโก ธมฺโม อุปฺปชฺชติ เหตุปจฺจยา. (1)}

\prose{เหตุํ อเหตุกญฺจ นเหตุํ อเหตุกญฺจ ธมฺมํ ปฏิจฺจ นเหตุ อเหตุโก ธมฺโม อุปฺปชฺชติ เหตุปจฺจยา. (1) (สํขิตฺตํ.)}

\prose{เหตุยา ตีณิ, อารมฺมเณ เอกํ, อวิคเต ตีณิ. (สํขิตฺตํ. สหชาตวารมฺปิ ฯเปฯ สมฺปยุตฺตวารมฺปิ ปฏิจฺจ\-วาร\-สทิสํ.)}

\prose{10. เหตุ อเหตุโก ธมฺโม นเหตุสฺส อเหตุกสฺส ธมฺมสฺส}

\prose{เหตุ อเหตุโก ธมฺโม เหตุสฺส อเหตุกสฺส ธมฺมสฺส อารมฺมณ\-ปจฺจเยน ปจฺจโย.  เหตุ อเหตุโก ธมฺโม นเหตุสฺส อเหตุกสฺส ธมฺมสฺส อารมฺมณ\-ปจฺจเยน ปจฺจโย. (2)}

\prose{นเหตุ อเหตุโก ธมฺโม นเหตุสฺส อเหตุกสฺส ธมฺมสฺส อารมฺมณ\-ปจฺจเยน ปจฺจโย.  นเหตุ อเหตุโก ธมฺโม เหตุสฺส อเหตุกสฺส ธมฺมสฺส อารมฺมณ\-ปจฺจเยน ปจฺจโย. (2) (สํขิตฺตํ.)}

\proseitem{11}{\csromanfolio{42}เหตุยา เอกํ, อารมฺมเณ จตฺตาริ, อวิคเต จตฺตาริ. (สํขิตฺตํ.) (ยถา กุสลตฺติเก ปญฺหาวารํ, เอวํ วิตฺถาเรตพฺพํ.)}

\csromanmarkh{1. เหตุทุก}\csromanmarkhii{2. เหตุ\-สมฺปยุตฺตทุก}\csromanheadernomark{2}{1. เหตุทุก 2. เหตุ\-สมฺปยุตฺตทุก}
\csromanheader{2}{1-7. ปฏิจฺจาทิวาร ปจฺจยจตุกฺกเหตุ}
\proseitem{12}{เหตุํ เหตุ\-สมฺปยุตฺตํ ธมฺมํ ปฏิจฺจ เหตุ เหตุสมฺปยุตฺโต ธมฺโม อุปฺปชฺชติ เหตุปจฺจยา.  เหตุํ เหตุ\-สมฺปยุตฺตํ ธมฺมํ ปฏิจฺจ นเหตุ เหตุสมฺปยุตฺโต ธมฺโม อุปฺปชฺชติ เหตุปจฺจยา.  เหตุํ เหตุ\-สมฺปยุตฺตํ ธมฺมํ ปฏิจฺจ เหตุ เหตุสมฺปยุตฺโต จ นเหตุ เหตุสมฺปยุตฺโต จ ธมฺมา อุปฺปชฺชนฺติ เหตุปจฺจยา. (3)}

\prose{นเหตุํ เหตุ\-สมฺปยุตฺตํ ธมฺมํ ปฏิจฺจ นเหตุ เหตุสมฺปยุตฺโต ธมฺโม อุปฺปชฺชติ เหตุปจฺจยา. ตีณิ.}

\prose{เหตุํ เหตุ\-สมฺปยุตฺตญฺจ นเหตุํ เหตุ\-สมฺปยุตฺตญฺจ ธมฺมํ ปฏิจฺจ เหตุ เหตุสมฺปยุตฺโต ธมฺโม อุปฺปชฺชติ เหตุปจฺจยา. ตีณิ.}

\proseitem{13}{เหตุยา นว, อารมฺมเณ นว ฯเปฯ อวิคเต นว. (สํขิตฺตํ.)}

\proseitem{14}{เหตุ เหตุสมฺปยุตฺโต ธมฺโม เหตุสฺส เหตุ\-สมฺปยุตฺตสฺส ธมฺมสฺส เหตุปจฺจเยน ปจฺจโย. ตีณิ. (สํขิตฺตํ.)}

\prose{เหตุยา ตีณิ ฯเปฯ อวิคเต นว. (สํขิตฺตํ.)}

\prose{(ยถา กุสลตฺติเก ปญฺหาวารํ, เอวํ วิตฺถาเรตพฺพํ.)}

\prose{15. เหตุํ เหตุ\-วิปฺปยุตฺตํ ธมฺมํ ปฏิจฺจ นเหตุ เหตุวิปฺปยุตฺโต}

\prose{นเหตุํ เหตุ\-วิปฺปยุตฺตํ ธมฺมํ ปฏิจฺจ นเหตุ เหตุวิปฺปยุตฺโต}

\csromanmarkh{1. เหตุทุก}\csromanmarkhii{5. นเหตุ\-สเหตุกทุก}\csromanheadernomark{2}{1. เหตุทุก 5. นเหตุ\-สเหตุกทุก}
\prose{\csromanfolio{43}เหตุํ เหตุ\-วิปฺปยุตฺตญฺจ นเหตุํ เหตุ\-วิปฺปยุตฺตญฺจ ธมฺมํ ปฏิจฺจ นเหตุ เหตุวิปฺปยุตฺโต ธมฺโม อุปฺปชฺชติ เหตุปจฺจยา. (1)}

\prose{เหตุยา ตีณิ, อารมฺมเณ เอกํ ฯเปฯ อวิคเต ตีณิ.}

\prose{(สหชาตวารมฺปิ ฯเปฯ สมฺปยุตฺตวารมฺปิ ปฏิจฺจ\-วาร\-สทิสํ วิตฺถาเรตพฺพํ.)}

\proseitem{16}{เหตุ เหตุวิปฺปยุตฺโต ธมฺโม นเหตุสฺส เหตุ\-วิปฺปยุตฺตสฺส ธมฺมสฺส เหตุปจฺจเยน ปจฺจโย. (สํขิตฺตํ.)}

\prose{เหตุยา เอกํ, อารมฺมเณ จตฺตาริ ฯเปฯ อวิคเต จตฺตาริ. (สํขิตฺตํ.) (ยถา กุสลตฺติเก ปญฺหาวารํ, เอวํ วิตฺถาเรตพฺพํ.)}

\csromanmarkh{1. เหตุทุก}\csromanmarkhii{3. เหตุ\-สเหตุกทุก}\csromanheadernomark{2}{1. เหตุทุก 3. เหตุ\-สเหตุกทุก}
\prose{1. ปฏิจฺจวาร ปจฺจยจตุกฺกเหตุ}

\proseitem{17}{เหตุํ เหตุญฺเจว สเหตุกญฺจ ธมฺมํ ปฏิจฺจ เหตุ เหตุ เจว สเหตุโก ธมฺโม อุปฺปชฺชติ เหตุปจฺจยา. (สพฺพตฺถ เอกํ.)}

\prose{นเหตุํ สเหตุกญฺเจว น จ เหตุํ ธมฺมํ ปฏิจฺจ นเหตุ สเหตุโก เจว น จ เหตุ ธมฺโม อุปฺปชฺชติ เหตุปจฺจยา. (สพฺพตฺถ เอกํ.)}

\csromanmarkh{1. เหตุทุก}\csromanmarkhii{4. เหตุ\-เหตุ\-สมฺปยุตฺตทุก}\csromanheadernomark{2}{1. เหตุทุก 4. เหตุ\-เหตุ\-สมฺปยุตฺตทุก}
\prose{1. ปฏิจฺจวาร ปจฺจยจตุกฺกเหตุ}

\proseitem{18}{เหตุํ เหตุญฺเจว เหตุ\-สมฺปยุตฺตญฺจ ธมฺมํ ปฏิจฺจ เหตุ เหตุ เจว เหตุสมฺปยุตฺโต จ ธมฺโม อุปฺปชฺชติ เหตุปจฺจยา. (สพฺพตฺถ เอกํ.)}

\prose{นเหตุํ เหตุ\-สมฺปยุตฺต\-ญฺเจว น จ เหตุํ ธมฺมํ ปฏิจฺจ นเหตุ เหตุสมฺปยุตฺโต เจว น จ เหตุ ธมฺโม อุปฺปชฺชติ เหตุปจฺจยา. (สพฺพตฺถ เอกํ.)}

\csromanmarkh{1. เหตุทุก}\csromanmarkhii{5. นเหตุ\-สเหตุกทุก}\csromanheadernomark{2}{1. เหตุทุก 5. นเหตุ\-สเหตุกทุก}
\proseitem{19}{นเหตุํ สเหตุกํ ธมฺมํ ปฏิจฺจ นเหตุ สเหตุโก ธมฺโม อุปฺปชฺชติ เหตุปจฺจยา. (สพฺพตฺถ เอกํ.)}

\prosewithcloser{\csromanfolio{44}นเหตุํ อเหตุกํ ธมฺมํ ปฏิจฺจ นเหตุ อเหตุโก ธมฺโม อุปฺปชฺชติ เหตุปจฺจยา. (สพฺพตฺถ เอกํ.)}{เหตุโคจฺฉกํ นิฏฺฐิตํ.}
\csromanmarkh{1. เหตุทุก}\csromanmarkhii{6. สปฺปจฺจยทุก}\csromanheadernomark{2}{1. \textbf{เหตุทุก 6. สปฺปจฺจยทุก}}
\proseitem{20}{เหตุํ สปฺปจฺจยํ ธมฺมํ ปฏิจฺจ เหตุ สปฺปจฺจโย ธมฺโม อุปฺปชฺชติ เหตุปจฺจยา. ตีณิ.}

\prose{นเหตุํ สปฺปจฺจยํ ธมฺมํ ปฏิจฺจ นเหตุ สปฺปจฺจโย ธมฺโม อุปฺปชฺชติ เหตุปจฺจยา. ตีณิ.}

\prose{เหตุํ สปฺปจฺจยญฺจ นเหตุํ สปฺปจฺจยญฺจ ธมฺมํ ปฏิจฺจ เหตุ สปฺปจฺจโย ธมฺโม อุปฺปชฺชติ เหตุปจฺจยา. ตีณิ. (สํขิตฺตํ.)}

\prose{เหตุยา นว, อารมฺมเณ นว ฯเปฯ อวิคเต นว. (สํขิตฺตํ.) (สหชาตวารมฺปิ ฯเปฯ สมฺปยุตฺตวารมฺปิ ปฏิจฺจ\-วาร\-สทิสํ.)}

\proseitem{21}{เหตุ สปฺปจฺจโย ธมฺโม เหตุสฺส สปฺปจฺจยสฺส ธมฺมสฺส เหตุปจฺจเยน ปจฺจโย. (สํขิตฺตํ.)}

\prose{เหตุยา ตีณิ.}

\vspace{\tipitakaparskip}
\csromancenter{(สํขิตฺตํ. ปญฺหาวารมฺปิ เอวํ วิตฺถาเรตพฺพํ.)}
\csromanmarkh{1. เหตุทุก}\csromanmarkhii{7. สงฺขตทุก}\csromanheadernomark{2}{1. \textbf{เหตุทุก 7. สงฺขตทุก}}
\proseitem{22}{เหตุํ สงฺขตํ ธมฺมํ ปฏิจฺจ เหตุ สงฺขโต ธมฺโม อุปฺปชฺชติ เหตุปจฺจยา.}

\prose{เหตุยา นว ฯเปฯ อวิคเต นว. (สปฺปจฺจย\-ทุก\-สทิสํ.)}

\csromanmarkh{1. เหตุทุก}\csromanmarkhii{8. สนิทสฺสนทุก}\csromanheadernomark{2}{1. \textbf{เหตุทุก 8. สนิทสฺสนทุก}}
\proseitem{23}{เหตุํ อนิทสฺสนํ ธมฺมํ ปฏิจฺจ เหตุ อนิทสฺสโน ธมฺโม อุปฺปชฺชติ เหตุปจฺจยา. ตีณิ.}

\csromanmarkh{1. เหตุทุก}\csromanmarkhii{10. รูปีทุก}\csromanheadernomark{2}{1. เหตุทุก 10. รูปีทุก}
\prose{\csromanfolio{45}นเหตุํ อนิทสฺสนํ ธมฺมํ ปฏิจฺจ นเหตุ อนิทสฺสโน ธมฺโม อุปฺปชฺชติ เหตุปจฺจยา. ตีณิ.}

\prose{เหตุํ อนิทสฺสนญฺจ นเหตุํ อนิทสฺสนญฺจ ธมฺมํ ปฏิจฺจ เหตุ อนิทสฺสโน ธมฺโม อุปฺปชฺชติ เหตุปจฺจยา. ตีณิ. (สํขิตฺตํ.)}

\prose{เหตุยา นว ฯเปฯ อวิคเต นว. (สํขิตฺตํ.) (สหชาตวารมฺปิ ฯเปฯ ปญฺหาวารมฺปิ วิตฺถาเรตพฺพํ.)}

\csromanmarkh{1. เหตุทุก}\csromanmarkhii{9. สปฺปฏิฆทุก}\csromanheadernomark{2}{1. \textbf{เหตุทุก 9. สปฺปฏิฆทุก}}
\proseitem{24}{นเหตุํ สปฺปฏิฆํ ธมฺมํ ปฏิจฺจ นเหตุ สปฺปฏิโฆ ธมฺโม อุปฺปชฺชติ เหตุปจฺจยา. (สพฺพตฺถ เอกํ.)}

\proseitem{25}{เหตุํ อปฺปฏิฆํ ธมฺมํ ปฏิจฺจ เหตุ อปฺปฏิโฆ ธมฺโม อุปฺปชฺชติ เหตุปจฺจยา.  เหตุํ อปฺปฏิฆํ ธมฺมํ ปฏิจฺจ นเหตุ อปฺปฏิโฆ ธมฺโม อุปฺปชฺชติ เหตุปจฺจยา.  เหตุํ อปฺปฏิฆํ ธมฺมํ ปฏิจฺจ เหตุ อปฺปฏิโฆ จ นเหตุ อปฺปฏิโฆ จ ธมฺมา อุปฺปชฺชนฺติ เหตุปจฺจยา. (3)}

\prose{นเหตุํ อปฺปฏิฆํ ธมฺมํ ปฏิจฺจ นเหตุ อปฺปฏิโฆ ธมฺโม อุปฺปชฺชติ เหตุปจฺจยา. ตีณิ.}

\prose{เหตุํ อปฺปฏิฆญฺจ นเหตุํ อปฺปฏิฆญฺจ ธมฺมํ ปฏิจฺจ เหตุ อปฺปฏิโฆ ธมฺโม อุปฺปชฺชติ เหตุปจฺจยา. ตีณิ. (สํขิตฺตํ.)}

\prose{เหตุยา นว, อารมฺมเณ นว ฯเปฯ อวิคเต นว. (สํขิตฺตํ.) (สหชาตวารมฺปิ ฯเปฯ ปญฺหาวารมฺปิ วิตฺถาเรตพฺพํ.)}

\csromanmarkh{1. เหตุทุก}\csromanmarkhii{10. รูปีทุก}\csromanheadernomark{2}{1. เหตุทุก 10. รูปีทุก}
\proseitem{26}{นเหตุํ รูปิํ ธมฺมํ ปฏิจฺจ นเหตุ รูปี ธมฺโม อุปฺปชฺชติ เหตุปจฺจยา. (สพฺพตฺถ เอกํ.)}

\proseitem{27}{เหตุํ อรูปิํ ธมฺมํ ปฏิจฺจ เหตุ อรูปี ธมฺโม อุปฺปชฺชติ เหตุปจฺจยา. ตีณิ.}

\prose{นเหตุํ อรูปิํ ธมฺมํ ปฏิจฺจ นเหตุ อรูปี ธมฺโม อุปฺปชฺชติ เหตุปจฺจยา. ตีณิ.}

\prose{\csromanfolio{46}เหตุํ อรูปิญฺจ นเหตุํ อรูปิญฺจ ธมฺมํ ปฏิจฺจ เหตุ อรูปี ธมฺโม อุปฺปชฺชติ เหตุปจฺจยา. ตีณิ. (สํขิตฺตํ.)}

\prose{เหตุยา นว, อารมฺมเณ นว ฯเปฯ อวิคเต นว. (สํขิตฺตํ.) (สหชาตวารมฺปิ ฯเปฯ ปญฺหาวารมฺปิ วิตฺถาเรตพฺพํ.)}

\csromanmarkh{1. เหตุทุก}\csromanmarkhii{11. โลกิยทุก}\csromanheadernomark{2}{1. เหตุทุก 11. โลกิยทุก}
\proseitem{28}{เหตุํ โลกิยํ ธมฺมํ ปฏิจฺจ เหตุ โลกิโย ธมฺโม อุปฺปชฺชติ เหตุปจฺจยา. ตีณิ.}

\prose{นเหตุํ โลกิยํ ธมฺมํ ปฏิจฺจ นเหตุ โลกิโย ธมฺโม อุปฺปชฺชติ เหตุปจฺจยา. ตีณิ.}

\prose{เหตุํ โลกิยญฺจ นเหตุํ โลกิยญฺจ ธมฺมํ ปฏิจฺจ เหตุ โลกิโย ธมฺโม อุปฺปชฺชติ เหตุปจฺจยา. ตีณิ. (สํขิตฺตํ.)}

\prose{เหตุยา นว, อารมฺมเณ นว ฯเปฯ อวิคเต นว. (สํขิตฺตํ.) (สหชาตวารมฺปิ ฯเปฯ ปญฺหาวารมฺปิ เอวํ วิตฺถาเรตพฺพํ.)}

\proseitem{29}{เหตุํ โลกุตฺตรํ ธมฺมํ ปฏิจฺจ เหตุ โลกุตฺตโร ธมฺโม อุปฺปชฺชติ เหตุปจฺจยา. ตีณิ.}

\prose{นเหตุํ โลกุตฺตรํ ธมฺมํ ปฏิจฺจ นเหตุ โลกุตฺตโร ธมฺโม อุปฺปชฺชติ เหตุปจฺจยา. ตีณิ.}

\prose{เหตุํ โลกุตฺตรญฺจ นเหตุํ โลกุตฺตรญฺจ ธมฺมํ ปฏิจฺจ เหตุ โลกุตฺตโร ธมฺโม อุปฺปชฺชติ เหตุปจฺจยา. ตีณิ. (สํขิตฺตํ.)}

\prose{เหตุยา นว, อารมฺมเณ นว ฯเปฯ อวิคเต นว. (สํขิตฺตํ.)}

\prose{(สหชาตวารมฺปิ ฯเปฯ สมฺปยุตฺตวารมฺปิ ปญฺหาวารมฺปิ ปฏิจฺจ\-วาร\-สทิสํ วิตฺถาเรตพฺพํ.)}

\csromanmarkh{1. เหตุทุก}\csromanmarkhii{12. เกนจิ\-วิญฺเญยฺยทุก}\csromanheadernomark{2}{1. เหตุทุก 12. เกนจิ\-วิญฺเญยฺยทุก}
\proseitem{30}{เหตุํ เกนจิ วิญฺเญยฺยํ ธมฺมํ ปฏิจฺจ เหตุ เกนจิ วิญฺเญยฺโย ธมฺโม อุปฺปชฺชติ เหตุปจฺจยา. (สํขิตฺตํ.)}

\prose{เหตุยา นว ฯเปฯ อวิคเต นว. (สํขิตฺตํ.)}

\vspace{\tipitakaparskip}
\csromancenter{(สหชาตวารมฺปิ ฯเปฯ ปญฺหาวารมฺปิ วิตฺถาเรตพฺพํ.)}
\csromanmarkh{1. เหตุทุก}\csromanmarkhii{14. สาสวทุก}\csromanheadernomark{2}{1. เหตุทุก 14. สาสวทุก}
\proseitem{31}{\csromanfolio{47}เหตุํ เกนจิ นวิญฺเญยฺยํ ธมฺมํ ปฏิจฺจ เหตุ เกนจิ นวิญฺเญยฺโย ธมฺโม อุปฺปชฺชติ เหตุปจฺจยา. (สํขิตฺตํ.)}

\prosewithcloser{เหตุยา นว ฯเปฯ อวิคเต นว. (สพฺพตฺถ วิตฺถาโร.)}{จูฬนฺตรทุกํ นิฏฺฐิตํ.}
\csromanmarkh{1. เหตุทุก}\csromanmarkhii{13. อาสวทุก}\csromanheadernomark{2}{1. เหตุทุก 13. อาสวทุก}
\proseitem{32}{เหตุํ อาสวํ ธมฺมํ ปฏิจฺจ เหตุ อาสโว ธมฺโม อุปฺปชฺชติ เหตุปจฺจยา.  เหตุํ อาสวํ ธมฺมํ ปฏิจฺจ นเหตุ อาสโว ธมฺโม อุปฺปชฺชติ เหตุปจฺจยา.  เหตุํ อาสวํ ธมฺมํ ปฏิจฺจ เหตุ อาสโว จ นเหตุ อาสโว จ ธมฺมา อุปฺปชฺชนฺติ เหตุปจฺจยา. (3)}

\prose{นเหตุํ อาสวํ ธมฺมํ ปฏิจฺจ นเหตุ อาสโว ธมฺโม อุปฺปชฺชติ}

\prose{เหตุํ อาสวญฺจ นเหตุํ อาสวญฺจ ธมฺมํ ปฏิจฺจ เหตุ อาสโว ธมฺโม อุปฺปชฺชติ เหตุปจฺจยา. (1) (สํขิตฺตํ.)}

\prose{เหตุยา ปญฺจ, อารมฺมเณ ปญฺจ ฯเปฯ อวิคเต ปญฺจ. (สพฺพตฺถ วิตฺถาโร.)}

\proseitem{33}{เหตุํ โนอาสวํ ธมฺมํ ปฏิจฺจ เหตุ โนอาสโว ธมฺโม อุปฺปชฺชติ เหตุปจฺจยา. ตีณิ.}

\prose{นเหตุํ โนอาสวํ ธมฺมํ ปฏิจฺจ นเหตุ โนอาสโว ธมฺโม อุปฺปชฺชติ เหตุปจฺจยา. ตีณิ.}

\prose{เหตุํ โนอาสวญฺจ นเหตุํ โนอาสวญฺจ ธมฺมํ ปฏิจฺจ เหตุ โนอาสโว ธมฺโม อุปฺปชฺชติ เหตุปจฺจยา. ตีณิ. (สํขิตฺตํ.)}

\prose{เหตุยา นว ฯเปฯ อวิคเต นว. (สพฺพตฺถ วิตฺถาโร.)}

\csromanmarkh{1. เหตุทุก}\csromanmarkhii{14. สาสวทุก}\csromanheadernomark{2}{1. เหตุทุก 14. สาสวทุก}
\proseitem{34}{เหตุํ สาสวํ ธมฺมํ ปฏิจฺจ เหตุ สาสโว ธมฺโม อุปฺปชฺชติ เหตุปจฺจยา. ตีณิ.}

\prose{นเหตุํ สาสวํ ธมฺมํ ปฏิจฺจ นเหตุ สาสโว ธมฺโม อุปฺปชฺชติ เหตุปจฺจยา. ตีณิ.}

\prose{\csromanfolio{48}เหตุํ สาสวญฺจ นเหตุํ สาสวญฺจ ธมฺมํ ปฏิจฺจ เหตุ สาสโว ธมฺโม อุปฺปชฺชติ เหตุปจฺจยา. ตีณิ. (สํขิตฺตํ.)}

\prose{เหตุยา นว ฯเปฯ อวิคเต นว. (สพฺพตฺถ วิตฺถาโร.)}

\proseitem{35}{เหตุํ อนาสวํ ธมฺมํ ปฏิจฺจ เหตุ อนาสโว ธมฺโม อุปฺปชฺชติ เหตุปจฺจยา. ตีณิ.}

\prose{นเหตุํ อนาสวํ ธมฺมํ ปฏิจฺจ นเหตุ อนาสโว ธมฺโม อุปฺปชฺชติ เหตุปจฺจยา. ตีณิ.}

\prose{เหตุํ อนาสวญฺจ นเหตุํ อนาสวญฺจ ธมฺมํ ปฏิจฺจ เหตุ อนาสโว ธมฺโม อุปฺปชฺชติ เหตุปจฺจยา. ตีณิ. (สํขิตฺตํ.)}

\prose{เหตุยา นว ฯเปฯ อวิคเต นว. (สพฺพตฺถ วิตฺถาโร.)}

\csromanmarkh{1. เหตุทุก}\csromanmarkhii{15. อาสว\-สมฺปยุตฺตทุก}\csromanheadernomark{2}{1. เหตุทุก 15. อาสว\-สมฺปยุตฺตทุก}
\proseitem{36}{เหตุํ อาสว\-สมฺปยุตฺตํ ธมฺมํ ปฏิจฺจ เหตุ อาสว\-สมฺปยุตฺโต ธมฺโม อุปฺปชฺชติ เหตุปจฺจยา. (สํขิตฺตํ.)}

\prose{เหตุยา นว ฯเปฯ อวิคเต นว. (สพฺพตฺถ วิตฺถาโร.)}

\proseitem{37}{เหตุํ อาสว\-วิปฺปยุตฺตํ ธมฺมํ ปฏิจฺจ เหตุ อาสว\-วิปฺปยุตฺโต ธมฺโม อุปฺปชฺชติ เหตุปจฺจยา. (สํขิตฺตํ.)}

\prose{เหตุยา นว ฯเปฯ อวิคเต นว. (สพฺพตฺถ วิตฺถาโร.)}

\csromanmarkh{1. เหตุทุก}\csromanmarkhii{16. อาสว\-สาสวทุก}\csromanheadernomark{2}{1. เหตุทุก 16. อาสว\-สาสวทุก}
\proseitem{38}{เหตุํ อาสวญฺเจว สาสวญฺจ ธมฺมํ ปฏิจฺจ เหตุ อาสโว เจว สาสโว จ ธมฺโม อุปฺปชฺชติ เหตุปจฺจยา. (สํขิตฺตํ.)}

\prose{เหตุยา ปญฺจ ฯเปฯ อวิคเต ปญฺจ. (สพฺพตฺถ วิตฺถาโร.)}

\proseitem{39}{เหตุํ สาสวญฺเจว โน จ อาสวํ ธมฺมํ ปฏิจฺจ เหตุ สาสโว เจว โน จ อาสโว ธมฺโม อุปฺปชฺชติ เหตุปจฺจยา. (สํขิตฺตํ.)}

\prose{เหตุยา นว ฯเปฯ อวิคเต นว. (สพฺพตฺถ วิตฺถาโร.)}

\vspace{\tipitakaparskip}
\csromancenter{เหตุทุก 15-53. สญฺโญชนา\-ทิทุก 1. เหตุทุก 17. อาสว\-อาสว\-สมฺปยุตฺตทุก}
\proseitem{40}{\csromanfolio{49}เหตุํ อาสวญฺเจว อาสว\-สมฺปยุตฺตญฺจ ธมฺมํ ปฏิจฺจ เหตุ อาสโว เจว อาสว\-สมฺปยุตฺโต จ ธมฺโม อุปฺปชฺชติ เหตุปจฺจยา.}

\prose{เหตุยา ปญฺจ ฯเปฯ อวิคเต ปญฺจ. (สพฺพตฺถ วิตฺถาโร.)}

\proseitem{41}{เหตุํ อาสว\-สมฺปยุตฺต\-ญฺเจว โน จ อาสวํ ธมฺมํ ปฏิจฺจ เหตุ อาสว\-สมฺปยุตฺโต เจว โน จ อาสโว ธมฺโม อุปฺปชฺชติ เหตุปจฺจยา.}

\prose{เหตุยา ปญฺจ ฯเปฯ อวิคเต ปญฺจ. (สพฺพตฺถ วิตฺถาโร.)}

\csromanmarkh{1. เหตุทุก}\csromanmarkhii{18. อาสว\-วิปฺปยุตฺต\-สาสวทุก}\csromanheadernomark{2}{1. เหตุทุก 18. อาสว\-วิปฺปยุตฺต\-สาสวทุก}
\proseitem{42}{เหตุํ อาสว\-วิปฺปยุตฺตํ สาสวํ ธมฺมํ ปฏิจฺจ เหตุ อาสว\-วิปฺปยุตฺโต สาสโว ธมฺโม อุปฺปชฺชติ เหตุปจฺจยา. (สํขิตฺตํ.)}

\prose{เหตุยา นว ฯเปฯ อวิคเต นว. (สพฺพตฺถ วิตฺถาโร.)}

\prose{เหตุํ อาสว\-วิปฺปยุตฺตํ อนาสวํ ธมฺมํ ปฏิจฺจ เหตุ อาสว\-วิปฺปยุตฺโต อนาสโว ธมฺโม อุปฺปชฺชติ เหตุปจฺจยา. (สํขิตฺตํ.)}

\prosewithcloser{เหตุยา นว ฯเปฯ อวิคเต นว. (สพฺพตฺถ วิตฺถาโร.)}{เหตุ\-ทุก\-อาสว\-โคจฺฉกํ นิฏฺฐิตํ.}
\csromanheader{2}{1. เหตุทุก 19-53 สญฺโญชนา\-ทิทุก}
\proseitem{43}{เหตุํ สญฺโญชนํ ธมฺมํ ปฏิจฺจ ฯเปฯ. เหตุํ คนฺถํ ธมฺมํ ปฏิจฺจ ฯเปฯ. เหตุํ โอฆํ ธมฺมํ ปฏิจฺจ ฯเปฯ. เหตุํ โยคํ ธมฺมํ ปฏิจฺจ ฯเปฯ. เหตุํ นีวรณํ ธมฺมํ ปฏิจฺจ ฯเปฯ. เหตุํ โนปรามาสํ ธมฺมํ ปฏิจฺจ เหตุ โน ปรามาโส ธมฺโม อุปฺปชฺชติ เหตุปจฺจยา. (สํขิตฺตํ.)}

\prosewithcloser{เหตุยา นว ฯเปฯ อวิคเต นว. (สพฺพตฺถ โคจฺฉกํ วิตฺถาเรตพฺพํ.)}{เหตุ\-ทุก\-ปรามาส\-โคจฺฉกํ นิฏฺฐิตํ.}
\csromanmarkh{1. เหตุทุก}\csromanmarkhii{54. สารมฺมณทุก}\csromanheadernomark{2}{1. เหตุทุก 54. สารมฺมณทุก}
\proseitem{44}{\csromanfolio{50}เหตุํ สารมฺมณํ ธมฺมํ ปฏิจฺจ เหตุ สารมฺมโณ ธมฺโม อุปฺปชฺชติ เหตุปจฺจยา. ตีณิ.}

\prose{นเหตุํ สารมฺมณํ ธมฺมํ ปฏิจฺจ นเหตุ สารมฺมโณ ธมฺโม อุปฺปชฺชติ เหตุปจฺจยา. ตีณิ.}

\prose{เหตุํ สารมฺมณญฺจ นเหตุํ สารมฺมณญฺจ ธมฺมํ ปฏิจฺจ เหตุ สารมฺมโณ ธมฺโม อุปฺปชฺชติ เหตุปจฺจยา. ตีณิ. (สํขิตฺตํ.)}

\prose{เหตุยา นว ฯเปฯ อวิคเต นว. (สพฺพตฺถ วิตฺถาโร.)}

\prose{นเหตุํ อนารมฺมณํ ธมฺมํ ปฏิจฺจ นเหตุ อนารมฺมโณ ธมฺโม อุปฺปชฺชติ เหตุปจฺจยา. (สพฺพตฺถ เอกํ.)}

\csromanmarkh{1. เหตุทุก}\csromanmarkhii{55. จิตฺตทุก}\csromanheadernomark{2}{1. เหตุทุก 55. จิตฺตทุก}
\proseitem{45}{เหตุํ โนจิตฺตํ ธมฺมํ ปฏิจฺจ เหตุ โนจิตฺโต ธมฺโม อุปฺปชฺชติ เหตุปจฺจยา. (สํขิตฺตํ.)}

\prose{เหตุยา นว ฯเปฯ อวิคเต นว. (สพฺพตฺถ วิตฺถาโร.)}

\csromanmarkh{1. เหตุทุก}\csromanmarkhii{56. เจตสิกทุก}\csromanheadernomark{2}{1. เหตุทุก 56. เจตสิกทุก}
\proseitem{46}{เหตุํ เจตสิกํ ธมฺมํ ปฏิจฺจ เหตุ เจตสิโก ธมฺโม อุปฺปชฺชติ เหตุปจฺจยา. (สํขิตฺตํ.)}

\prose{เหตุยา นว ฯเปฯ อวิคเต นว. (สพฺพตฺถ วิตฺถาโร.)}

\prose{นเหตุํ อเจตสิกํ ธมฺมํ ปฏิจฺจ นเหตุ อเจตสิโก ธมฺโม อุปฺปชฺชติ เหตุปจฺจยา. (สพฺพตฺถ เอกํ.)}

\csromanmarkh{1. เหตุทุก}\csromanmarkhii{57. จิตฺต\-สมฺปยุตฺตทุก}\csromanheadernomark{2}{1. เหตุทุก 57. จิตฺต\-สมฺปยุตฺตทุก}
\proseitem{47}{เหตุํ จิตฺต\-สมฺปยุตฺตํ ธมฺมํ ปฏิจฺจ เหตุ จิตฺต\-สมฺปยุตฺโต ธมฺโม อุปฺปชฺชติ เหตุปจฺจยา. (สํขิตฺตํ.)}

\prose{เหตุยา นว ฯเปฯ อวิคเต นว. (สพฺพตฺถ วิตฺถาโร.)}

\prose{นเหตุํ จิตฺต\-วิปฺปยุตฺตํ ธมฺมํ ปฏิจฺจ นเหตุ จิตฺต\-วิปฺปยุตฺโต ธมฺโม อุปฺปชฺชติ เหตุปจฺจยา. (สพฺพตฺถ เอกํ.)}

\vspace{\tipitakaparskip}
\csromancenter{เหตุทุก 62. จิตฺต\-สํสฏฺฐ\-สมุฏฺฐานทุก 1. เหตุทุก 58. จิตฺต\-สํสฏฺฐทุก}
\proseitem{48}{\csromanfolio{51}เหตุํ จิตฺต\-สํสฏฺฐํ ธมฺมํ ปฏิจฺจ เหตุ จิตฺตสํสฏฺโฐ ธมฺโม อุปฺปชฺชติ เหตุปจฺจยา. (สํขิตฺตํ.)}

\prose{เหตุยา นว ฯเปฯ อวิคเต นว. (สพฺพตฺถ วิตฺถาโร.)}

\prose{นเหตุํ จิตฺต\-วิสํสฏฺฐํ ธมฺมํ ปฏิจฺจ นเหตุ จิตฺต\-วิสํสฏฺโฐ ธมฺโม อุปฺปชฺชติ เหตุปจฺจยา. (สพฺพตฺถ เอกํ.)}

\csromanmarkh{1. เหตุทุก}\csromanmarkhii{59. จิตฺต\-สมุฏฺฐานทุก}\csromanheadernomark{2}{1. เหตุทุก 59. จิตฺต\-สมุฏฺฐานทุก}
\proseitem{49}{เหตุํ จิตฺต\-สมุฏฺฐานํ ธมฺมํ ปฏิจฺจ เหตุ จิตฺต\-สมุฏฺฐาโน ธมฺโม อุปฺปชฺชติ เหตุปจฺจยา. (สํขิตฺตํ.)}

\prose{เหตุยา นว ฯเปฯ อวิคเต นว. (สพฺพตฺถ วิตฺถาโร.)}

\prose{นเหตุํ โนจิตฺต\-สมุฏฺฐานํ ธมฺมํ ปฏิจฺจ นเหตุ โนจิตฺต\-สมุฏฺฐาโน ธมฺโม อุปฺปชฺชติ เหตุปจฺจยา. (สพฺพตฺถ เอกํ.)}

\csromanmarkh{1. เหตุทุก}\csromanmarkhii{60. จิตฺต\-สหภูทุก}\csromanheadernomark{2}{1. เหตุทุก 60. จิตฺต\-สหภูทุก}
\proseitem{50}{เหตุํ จิตฺตสหภุํ ธมฺมํ ปฏิจฺจ เหตุ จิตฺตสหภู ธมฺโม อุปฺปชฺชติ เหตุปจฺจยา. (สพฺพตฺถ นว.)}

\prose{นเหตุํ โนจิตฺต\-สหภุํ ธมฺมํ ปฏิจฺจ นเหตุ โนจิตฺตสหภู ธมฺโม อุปฺปชฺชติ เหตุปจฺจยา. (สพฺพตฺถ เอกํ. สพฺพตฺถ วิตฺถาโร.)}

\csromanmarkh{1. เหตุทุก}\csromanmarkhii{61. จิตฺตา\-นุปริวตฺติทุก}\csromanheadernomark{2}{1. เหตุทุก 61. จิตฺตา\-นุปริวตฺติทุก}
\proseitem{51}{เหตุํ จิตฺตา\-นุปริวตฺติํ ธมฺมํ ปฏิจฺจ เหตุ จิตฺตา\-นุปริวตฺตี ธมฺโม อุปฺปชฺชติ เหตุปจฺจยา. (สพฺพตฺถ นว.)}

\prose{นเหตุํ โนจิตฺตา\-นุปริวตฺติํ ธมฺมํ ปฏิจฺจ นเหตุ โนจิตฺตา\-นุปริวตฺตี ธมฺโม อุปฺปชฺชติ เหตุปจฺจยา. (สพฺพตฺถ เอกํ. สพฺพตฺถ วิตฺถาโร.)}

\csromanmarkh{1. เหตุทุก}\csromanmarkhii{62. จิตฺต\-สํสฏฺฐ\-สมุฏฺฐานทุก}\csromanheadernomark{2}{1. เหตุทุก 62. จิตฺต\-สํสฏฺฐ\-สมุฏฺฐานทุก}
\proseitem{52}{เหตุํ จิตฺต\-สํสฏฺฐ\-สมุฏฺฐานํ ธมฺมํ ปฏิจฺจ เหตุ จิตฺต\-สํสฏฺฐ\-สมุฏฺฐาโน ธมฺโม อุปฺปชฺชติ เหตุปจฺจยา. (สพฺพตฺถ นว.)}

\prose{นเหตุํ โนจิตฺต\-สํสฏฺฐ\-สมุฏฺฐานํ ธมฺมํ ปฏิจฺจ นเหตุ โนจิตฺต\-สํสฏฺฐ\-สมุฏฺฐาโน ธมฺโม อุปฺปชฺชติ เหตุปจฺจยา. (สพฺพตฺถ เอกํ. สพฺพตฺถ วิตฺถาโร.)}

\prose{\csromanfolio{52}1. เหตุทุก 63. จิตฺต\-สํสฏฺฐ\-สมุฏฺฐาน\-สหภูทุก}

\proseitem{53}{เหตุํ จิตฺต\-สํสฏฺฐ\-สมุฏฺฐาน\-สหภุํ ธมฺมํ ปฏิจฺจ เหตุ จิตฺต\-สํสฏฺฐ\-สมุฏฺฐาน\-สหภู ธมฺโม อุปฺปชฺชติ เหตุปจฺจยา. (สพฺพตฺถ นว.)}

\prose{นเหตุํ โนจิตฺต\-สํสฏฺฐ\-สมุฏฺฐาน\-สหภุํ ธมฺมํ ปฏิจฺจ นเหตุ โนจิตฺต\-สํสฏฺฐ\-สมุฏฺฐาน\-สหภู ธมฺโม อุปฺปชฺชติ เหตุปจฺจยา. (สพฺพตฺถ เอกํ. สพฺพตฺถ วิตฺถาโร.)}

\prose{1. เหตุทุก 64. จิตฺต\-สํสฏฺฐ\-สมุฏฺฐานา\-นุปริวตฺติทุก}

\proseitem{54}{เหตุํ จิตฺต\-สํสฏฺฐ\-สมุฏฺฐานา\-นุปริวตฺติํ ธมฺมํ ปฏิจฺจ เหตุ จิตฺต\-สํสฏฺฐ\-สมุฏฺฐานา\-นุปริวตฺตี ธมฺโม อุปฺปชฺชติ เหตุปจฺจยา. (สพฺพตฺถ นว.)}

\prose{นเหตุํ โนจิตฺต\-สํสฏฺฐ\-สมุฏฺฐานา\-นุปริวตฺติํ ธมฺมํ ปฏิจฺจ นเหตุ โนจิตฺต\-สํสฏฺฐ\-สมุฏฺฐานา\-นุปริวตฺตี ธมฺโม อุปฺปชฺชติ เหตุปจฺจยา. (สพฺพตฺถ เอกํ. สพฺพตฺถ วิตฺถาโร.)}

\csromanmarkh{1. เหตุทุก}\csromanmarkhii{65. อชฺฌตฺติกทุก}\csromanheadernomark{2}{1. เหตุทุก 65. อชฺฌตฺติกทุก}
\proseitem{55}{นเหตุํ อชฺฌตฺติกํ ธมฺมํ ปฏิจฺจ นเหตุ อชฺฌตฺติโก ธมฺโม อุปฺปชฺชติ เหตุปจฺจยา. (สพฺพตฺถ เอกํ. สพฺพตฺถ วิตฺถาโร.)}

\prose{เหตุํ พาหิรํ ธมฺมํ ปฏิจฺจ เหตุ พาหิโร ธมฺโม อุปฺปชฺชติ เหตุปจฺจยา. (สพฺพตฺถ นว. สพฺพตฺถ วิตฺถาโร.)}

\csromanmarkh{1. เหตุทุก}\csromanmarkhii{66. อุปาทาทุก}\csromanheadernomark{2}{1. เหตุทุก 66. อุปาทาทุก}
\proseitem{56}{เหตุํ โนอุปาทา ธมฺมํ ปฏิจฺจ เหตุ โนอุปาทา ธมฺโม อุปฺปชฺชติ เหตุปจฺจยา. (สํขิตฺตํ.)}

\prose{เหตุยา นว ฯเปฯ อวิคเต นว. (สํขิตฺตํ. สพฺพตฺถ วิตฺถาโร.)}

\csromanmarkh{1. เหตุทุก}\csromanmarkhii{67. อุปาทินฺนทุก}\csromanheadernomark{2}{1. เหตุทุก 67. อุปาทินฺนทุก}
\proseitem{57}{เหตุํ อุปาทินฺนํ ธมฺมํ ปฏิจฺจ เหตุ อุปาทินฺโน ธมฺโม อุปฺปชฺชติ เหตุปจฺจยา. (สํขิตฺตํ.)}

\prose{เหตุยา นว ฯเปฯ อวิคเต นว. (สพฺพตฺถ วิตฺถาโร.)}

\csromanmarkh{1. เหตุทุก}\csromanmarkhii{82. ปิฏฺฐิทุก}\csromanheadernomark{2}{1. เหตุทุก 82. ปิฏฺฐิทุก}
\prose{\csromanfolio{53}เหตุํ อนุปาทินฺนํ ธมฺมํ ปฏิจฺจ เหตุ อนุปาทินฺโน ธมฺโม อุปฺปชฺชติ เหตุปจฺจยา. (สํขิตฺตํ.)}

\prosewithcloser{เหตุยา นว ฯเปฯ อวิคเต นว. (สพฺพตฺถ วิตฺถาโร.)}{เหตุ\-ทุก\-มหนฺตรทุกํ นิฏฺฐิตํ.}
\csromanmarkh{1. เหตุทุก}\csromanmarkhii{68. อุปาทาน\-โคจฺฉก}\csromanheadernomark{2}{1. เหตุทุก 68. อุปาทาน\-โคจฺฉก}
\proseitem{58}{เหตุํ อุปาทานํ ธมฺมํ ปฏิจฺจ นเหตุ อุปาทาโน ธมฺโม อุปฺปชฺชติ เหตุปจฺจยา.}

\prose{เหตุยา เทฺว ฯเปฯ อวิคเต เทฺว. (สํขิตฺตํ.)}

\csromanmarkh{1. เหตุทุก}\csromanmarkhii{74. กิเลสโคจฺฉก}\csromanheadernomark{2}{1. เหตุทุก 74. กิเลสโคจฺฉก}
\proseitem{59}{เหตุํ กิเลสํ ธมฺมํ ปฏิจฺจ เหตุ กิเลโส ธมฺโม อุปฺปชฺชติ เหตุปจฺจยา.}

\prose{เหตุยา นว ฯเปฯ อวิคเต นว. (สํขิตฺตํ.)}

\csromanmarkh{1. เหตุทุก}\csromanmarkhii{82. ปิฏฺฐิทุก}\csromanheadernomark{2}{1. เหตุทุก 82. ปิฏฺฐิทุก}
\proseitem{60}{เหตุํ ทสฺสเนน ปหาตพฺพํ ธมฺมํ ปฏิจฺจ เหตุ ทสฺสเนน ปหาตพฺโพ ธมฺโม อุปฺปชฺชติ เหตุปจฺจยา. (สพฺพตฺถ นว. วิปากํ นตฺถิ.)}

\prose{เหตุํ นทสฺสเนน ปหาตพฺพํ ธมฺมํ ปฏิจฺจ เหตุ นทสฺสเนน ปหาตพฺโพ ธมฺโม อุปฺปชฺชติ เหตุ ปจฺจยา. (สพฺพตฺถ นว. วิปากํ นตฺถิ.)}

\proseitem{61}{เหตุํ ภาวนาย ปหาตพฺพํ ธมฺมํ ปฏิจฺจ เหตุ ภาวนาย ปหาตพฺโพ ธมฺโม อุปฺปชฺชติ เหตุปจฺจยา. (สพฺพตฺถ นว. วิปากํ นตฺถิ.)}

\prose{เหตุํ นภาวนาย ปหาตพฺพํ ธมฺมํ ปฏิจฺจ เหตุ นภาวนาย ปหาตพฺโพ ธมฺโม อุปฺปชฺชติ เหตุปจฺจยา. (สพฺพตฺถ นว. วิปากํ นตฺถิ.)}

\proseitem{62}{เหตุํ ทสฺสเนน ปหาตพฺพ\-เหตุกํ ธมฺมํ ปฏิจฺจ เหตุ ทสฺสเนน ปหาตพฺพ\-เหตุโก ธมฺโม อุปฺปชฺชติ เหตุปจฺจยา. (สพฺพตฺถ นว.)}

\prose{เหตุํ นทสฺสเนน ปหาตพฺพ\-เหตุกํ ธมฺมํ ปฏิจฺจ เหตุ นทสฺสเนน ปหาตพฺพ\-เหตุโก ธมฺโม อุปฺปชฺชติ เหตุปจฺจยา. (สพฺพตฺถ นว.)}

\proseitem{63}{\csromanfolio{54}เหตุํ ภาวนาย ปหาตพฺพ\-เหตุกํ ธมฺมํ ปฏิจฺจ เหตุ ภาวนาย ปหาตพฺพ\-เหตุโก ธมฺโม อุปฺปชฺชติ เหตุปจฺจยา. (สพฺพตฺถ นว.)}

\prose{เหตุํ นภาวนาย ปหาตพฺพ\-เหตุกํ ธมฺมํ ปฏิจฺจ เหตุ นภาวนาย ปหาตพฺพ\-เหตุโก ธมฺโม อุปฺปชฺชติ เหตุปจฺจยา. (สพฺพตฺถ นว.)}

\proseitem{64}{เหตุํ สวิตกฺกํ ธมฺมํ ปฏิจฺจ เหตุ สวิตกฺโก ธมฺโม อุปฺปชฺชติ เหตุปจฺจยา. (สพฺพตฺถ นว.)}

\prose{เหตุํ อวิตกฺกํ ธมฺมํ ปฏิจฺจ เหตุ อวิตกฺโก ธมฺโม อุปฺปชฺชติ เหตุปจฺจยา. (สพฺพตฺถ นว.)}

\proseitem{65}{เหตุํ สวิจารํ ธมฺมํ ปฏิจฺจ เหตุ สวิจาโร ธมฺโม อุปฺปชฺชติ เหตุปจฺจยา. (สพฺพตฺถ นว.)}

\prose{เหตุํ อวิจารํ ธมฺมํ ปฏิจฺจ เหตุ อวิจาโร ธมฺโม อุปฺปชฺชติ เหตุปจฺจยา. (สพฺพตฺถ นว.)}

\proseitem{66}{เหตุํ สปฺปีติกํ ธมฺมํ ปฏิจฺจ ฯเปฯ. เหตุํ อปฺปีติกํ ธมฺมํ ปฏิจฺจ ฯเปฯ.}

\prose{เหตุํ ปีติสหคตํ ธมฺมํ ปฏิจฺจ ฯเปฯ. เหตุํ นปีติสหคตํ ธมฺมํ ปฏิจฺจ ฯเปฯ.}

\prose{เหตุํ สุขสหคตํ ธมฺมํ ปฏิจฺจ ฯเปฯ. เหตุํ นสุขสหคตํ ธมฺมํ ปฏิจฺจ ฯเปฯ.}

\prose{เหตุํ อุเปกฺขา\-สหคตํ ธมฺมํ ปฏิจฺจ ฯเปฯ. เหตุํ นอุเปกฺขา\-สหคตํ ธมฺมํ ปฏิจฺจ ฯเปฯ.}

\prose{เหตุํ กามาวจรํ ธมฺมํ ปฏิจฺจ ฯเปฯ. เหตุํ นกามาวจรํ ธมฺมํ ปฏิจฺจ ฯเปฯ.}

\prose{เหตุํ รูปาวจรํ ธมฺมํ ปฏิจฺจ ฯเปฯ. เหตุํ นรูปาวจรํ ธมฺมํ ปฏิจฺจ ฯเปฯ.}

\prose{เหตุํ อรูปาวจรํ ธมฺมํ ปฏิจฺจ ฯเปฯ. เหตุํ นอรูปาวจรํ ธมฺมํ ปฏิจฺจ ฯเปฯ.}

\prose{เหตุํ ปริยาปนฺนํ ธมฺมํ ปฏิจฺจ ฯเปฯ. เหตุํ อปริยาปนฺนํ ธมฺมํ ปฏิจฺจ ฯเปฯ.}

\prose{เหตุํ นิยฺยานิกํ ธมฺมํ ปฏิจฺจ ฯเปฯ. เหตุํ อนิยฺยานิกํ ธมฺมํ ปฏิจฺจ ฯเปฯ.}

\prose{เหตุํ นิยตํ ธมฺมํ ปฏิจฺจ ฯเปฯ. เหตุํ อนิยตํ ธมฺมํ ปฏิจฺจ ฯเปฯ.}

\prose{เหตุํ สอุตฺตรํ ธมฺมํ ปฏิจฺจ ฯเปฯ. เหตุํ อนุตฺตรํ ธมฺมํ ปฏิจฺจ ฯเปฯ.}

\csromanmarkh{5. เหตุ\-เหตุ\-สมฺปยุตฺตทุก}\csromanmarkhii{1. เหตุทุก}\csromanheadernomark{2}{5. เหตุ\-เหตุ\-สมฺปยุตฺตทุก 1. เหตุทุก}
\prose{\csromanfolio{55}เหตุํ สรณํ ธมฺมํ ปฏิจฺจ ฯเปฯ. เหตุํ อรณํ ธมฺมํ ปฏิจฺจ เหตุ อรโณ ธมฺโม อุปฺปชฺชติ เหตุปจฺจยา. (สํขิตฺตํ.)}

\prosewithcloser{เหตุยา นว ฯเปฯ อวิคเต นว. (สํขิตฺตํ. สพฺพตฺถ วิตฺถาโร.)}{เหตุ\-ทุก\-ปิฏฺฐิทุกํ นิฏฺฐิตํ.}
\csromanmarkh{2. สเหตุกทุก}\csromanmarkhii{1. เหตุทุก}\csromanheadernomark{2}{2. \textbf{สเหตุกทุก 1. เหตุทุก}}
\proseitem{67}{สเหตุกํ เหตุํ ธมฺมํ ปฏิจฺจ สเหตุโก เหตุ ธมฺโม อุปฺปชฺชติ เหตุปจฺจยา. (สพฺพตฺถ เอกํ.)}

\prose{สเหตุกํ นเหตุํ ธมฺมํ ปฏิจฺจ สเหตุโก นเหตุ ธมฺโม อุปฺปชฺชติ เหตุปจฺจยา.}

\prose{เหตุยา นว.}

\csromanmarkh{3. เหตุ\-สมฺปยุตฺตทุก}\csromanmarkhii{1. เหตุทุก}\csromanheadernomark{2}{3. เหตุ\-สมฺปยุตฺตทุก 1. เหตุทุก}
\proseitem{68}{เหตุ\-สมฺปยุตฺตํ เหตุํ ธมฺมํ ปฏิจฺจ เหตุสมฺปยุตฺโต เหตุ ธมฺโม อุปฺปชฺชติ เหตุปจฺจยา. (สพฺพตฺถ เอกํ.)}

\prose{เหตุ\-สมฺปยุตฺตํ นเหตุํ ธมฺมํ ปฏิจฺจ เหตุสมฺปยุตฺโต นเหตุ ธมฺโม อุปฺปชฺชติ เหตุปจฺจยา. (สํขิตฺตํ.)}

\prose{เหตุยา นว.}

\csromanmarkh{4. เหตุ\-สเหตุกทุก}\csromanmarkhii{1. เหตุทุก}\csromanheadernomark{2}{4. เหตุ\-สเหตุกทุก 1. เหตุทุก}
\proseitem{69}{เหตุญฺเจว สเหตุกญฺจ เหตุํ ธมฺมํ ปฏิจฺจ เหตุ เจว สเหตุโก จ เหตุ ธมฺโม อุปฺปชฺชติ เหตุปจฺจยา. (สพฺพตฺถ เอกํ.)}

\prose{สเหตุกญฺเจว น จ เหตุํ นเหตุํ ธมฺมํ ปฏิจฺจ สเหตุโก เจว น จ เหตุ นเหตุ ธมฺโม อุปฺปชฺชติ เหตุปจฺจยา. (สพฺพตฺถ เอกํ.)}

\csromanmarkh{5. เหตุ\-เหตุ\-สมฺปยุตฺตทุก}\csromanmarkhii{1. เหตุทุก}\csromanheadernomark{2}{5. เหตุ\-เหตุ\-สมฺปยุตฺตทุก 1. เหตุทุก}
\proseitem{70}{เหตุญฺเจว เหตุ\-สมฺปยุตฺตญฺจ เหตุํ ธมฺมํ ปฏิจฺจ เหตุ เจว เหตุสมฺปยุตฺโต จ เหตุ ธมฺโม อุปฺปชฺชติ เหตุปจฺจยา. (สพฺพตฺถ เอกํ.)}

\prose{\csromanfolio{56}เหตุ\-สมฺปยุตฺต\-ญฺเจว น จ เหตุํ นเหตุํ ธมฺมํ ปฏิจฺจ เหตุสมฺปยุตฺโต เจว น จ เหตุ นเหตุ ธมฺโม อุปฺปชฺชติ เหตุปจฺจยา. (สพฺพตฺถ เอกํ.)}

\csromanmarkh{6. นเหตุ\-สเหตุกทุก}\csromanmarkhii{1. เหตุทุก}\csromanheadernomark{2}{6. นเหตุ\-สเหตุกทุก 1. เหตุทุก}
\proseitem{71}{นเหตุ\-สเหตุกํ นเหตุํ ธมฺมํ ปฏิจฺจ นเหตุสเหตุโก นเหตุ ธมฺโม อุปฺปชฺชติ เหตุปจฺจยา.}

\prose{นเหตุํ อเหตุกํ นเหตุํ ธมฺมํ ปฏิจฺจ นเหตุ อเหตุโก นเหตุ ธมฺโม อุปฺปชฺชติ เหตุปจฺจยา.}

\csromanmarkh{7. จูฬนฺตรทุก}\csromanmarkhii{1. เหตุทุก}\csromanheadernomark{2}{7. \textbf{จูฬนฺตรทุก 1. เหตุทุก}}
\proseitem{72}{สปฺปจฺจยํ เหตุํ ธมฺมํ ปฏิจฺจ สปฺปจฺจโย เหตุ ธมฺโม อุปฺปชฺชติ เหตุปจฺจยา. (สพฺพตฺถ เอกํ.)}

\prose{สปฺปจฺจยํ นเหตุํ ธมฺมํ ปฏิจฺจ สปฺปจฺจโย นเหตุ ธมฺโม อุปฺปชฺชติ เหตุปจฺจยา. (สพฺพตฺถ เอกํ.)}

\proseitem{73}{สงฺขตํ เหตุํ ธมฺมํ ปฏิจฺจ สงฺขโต เหตุ ธมฺโม อุปฺปชฺชติ เหตุปจฺจยา. (สพฺพตฺถ เอกํ.)}

\prose{สงฺขตํ นเหตุํ ธมฺมํ ปฏิจฺจ สงฺขโต นเหตุ ธมฺโม อุปฺปชฺชติ เหตุปจฺจยา. (สพฺพตฺถ เอกํ.)}

\proseitem{74}{อนิทสฺสนํ เหตุํ ธมฺมํ ปฏิจฺจ อนิทสฺสโน เหตุ ธมฺโม อุปฺปชฺชติ เหตุปจฺจยา. (สพฺพตฺถ เอกํ.)}

\prose{อนิทสฺสนํ นเหตุํ ธมฺมํ ปฏิจฺจ อนิทสฺสโน นเหตุ ธมฺโม อุปฺปชฺชติ เหตุปจฺจยา. (สํขิตฺตํ.)}

\prose{เหตุยา ตีณิ, อารมฺมเณ เอกํ ฯเปฯ อวิคเต ตีณิ. (สํขิตฺตํ.)}

\proseitem{75}{อปฺปฏิฆํ เหตุํ ธมฺมํ ปฏิจฺจ สปฺปฏิโฆ เหตุ ธมฺโม อุปฺปชฺชติ เหตุปจฺจยา. (สพฺพตฺถ เอกํ.)}

\csromanmarkh{7. จูฬนฺตรทุก}\csromanmarkhii{1. เหตุทุก}\csromanheadernomark{2}{7. จูฬนฺตรทุก 1. เหตุทุก}
\prose{\csromanfolio{57}สปฺปฏิฆํ นเหตุํ ธมฺมํ ปฏิจฺจ สปฺปฏิโฆ นเหตุ ธมฺโม อุปฺปชฺชติ เหตุปจฺจยา. ตีณิ.}

\prose{อปฺปฏิฆํ นเหตุํ ธมฺมํ ปฏิจฺจ อปฺปฏิโฆ นเหตุ ธมฺโม อุปฺปชฺชติ เหตุปจฺจยา. ตีณิ.}

\prose{สปฺปฏิฆํ นเหตุญฺจ อปฺปฏิฆํ นเหตุญฺจ ธมฺมํ ปฏิจฺจ สปฺปฏิโฆ นเหตุ ธมฺโม อุปฺปชฺชติ เหตุปจฺจยา. ตีณิ. (สํขิตฺตํ.)}

\prose{เหตุยา นว, อารมฺมเณ เอกํ ฯเปฯ อญฺญมญฺเญ ฉ ฯเปฯ อวิคเต นว.}

\proseitem{76}{อรูปิํ เหตุํ ธมฺมํ ปฏิจฺจ อรูปี เหตุ ธมฺโม อุปฺปชฺชติ เหตุปจฺจยา. (สพฺพตฺถ เอกํ.)}

\prose{รูปิํ นเหตุํ ธมฺมํ ปฏิจฺจ รูปี นเหตุ ธมฺโม อุปฺปชฺชติ เหตุปจฺจยา. (สํขิตฺตํ.)}

\prose{เหตุยา นว.}

\prose{77. โลกิยํ เหตุํ ธมฺมํ ปฏิจฺจ โลกิโย เหตุ ธมฺโม อุปฺปชฺชติ}

\prose{โลกุตฺตรํ เหตุํ ธมฺมํ ปฏิจฺจ โลกุตฺตโร เหตุ ธมฺโม อุปฺปชฺชติ เหตุปจฺจยา. (1) (สพฺพตฺถ เทฺว.)}

\prose{โลกิยํ นเหตุํ ธมฺมํ ปฏิจฺจ โลกิโย นเหตุ ธมฺโม อุปฺปชฺชติ}

\prose{โลกุตฺตรํ นเหตุํ ธมฺมํ ปฏิจฺจ โลกุตฺตโร นเหตุ ธมฺโม อุปฺปชฺชติ เหตุปจฺจยา. ตีณิ.}

\prose{โลกิยํ นเหตุญฺจ โลกุตฺตรํ นเหตุญฺจ ธมฺมํ ปฏิจฺจ โลกิโย นเหตุ ธมฺโม อุปฺปชฺชติ เหตุปจฺจยา. (1) (สํขิตฺตํ.)}

\prose{เหตุยา ปญฺจ.}

\proseitem{78}{เกนจิ วิญฺเญยฺยํ เหตุํ ธมฺมํ ปฏิจฺจ ฯเปฯ. นเกนจิ วิญฺเญยฺยํ เหตุํ ธมฺมํ ปฏิจฺจ. (สพฺพตฺถ นว.)}

\prose{เกนจิ วิญฺเญยฺยํ นเหตุํ ธมฺมํ ปฏิจฺจ ฯเปฯ. นเกนจิ วิญฺเญยฺยํ นเหตุํ ธมฺมํ ปฏิจฺจ. (สพฺพตฺถ นว.)}

\csromanmarkh{14. อาสวโคจฺฉก}\csromanmarkhii{1. เหตุทุก}\csromanheadernomark{2}{14. \textbf{อาสวโคจฺฉก 1. เหตุทุก}}
\proseitem{79}{\csromanfolio{58}อาสวํ เหตุํ ธมฺมํ ปฏิจฺจ ฯเปฯ. โนอาสวํ เหตุํ ธมฺมํ ปฏิจฺจ. (สํขิตฺตํ.)}

\prose{อาสวํ นเหตุํ ธมฺมํ ปฏิจฺจ ฯเปฯ. โนอาสวํ นเหตุํ ธมฺมํ ปฏิจฺจ. (สํขิตฺตํ.)}

\proseitem{80}{สาสวํ เหตุํ ธมฺมํ ปฏิจฺจ ฯเปฯ. อนาสวํ เหตุํ ธมฺมํ ปฏิจฺจ. (สพฺพตฺถ เทฺว.)}

\prose{สาสวํ นเหตุํ ธมฺมํ ปฏิจฺจ ฯเปฯ. อนาสวํ นเหตุํ ธมฺมํ ปฏิจฺจ ฯเปฯ.}

\proseitem{81}{อาสว\-สมฺปยุตฺตํ เหตุํ ธมฺมํ ปฏิจฺจ ฯเปฯ. อาสว\-วิปฺปยุตฺตํ เหตุํ ธมฺมํ ปฏิจฺจ. (สํขิตฺตํ.)}

\prose{อาสว\-สมฺปยุตฺตํ นเหตุํ ธมฺมํ ปฏิจฺจ ฯเปฯ. อาสว\-วิปฺปยุตฺตํ นเหตุํ ธมฺมํ ปฏิจฺจ. (สํขิตฺตํ.)}

\proseitem{82}{อาสวญฺเจว สาสวญฺจ เหตุํ ธมฺมํ ปฏิจฺจ ฯเปฯ. สาสวญฺเจว โน จ อาสวํ เหตุํ ธมฺมํ ปฏิจฺจ.}

\prose{อาสวญฺเจว สาสวญฺจ นเหตุํ ธมฺมํ ปฏิจฺจ ฯเปฯ. สาสวญฺเจว โน จ อาสวํ นเหตุํ ธมฺมํ ปฏิจฺจ. (สํขิตฺตํ.)}

\proseitem{83}{อาสวญฺเจว อาสว\-สมฺปยุตฺตญฺจ เหตุํ ธมฺมํ ปฏิจฺจ ฯเปฯ. อาสว\-สมฺปยุตฺต\-ญฺเจว โน จ อาสวํ เหตุํ ธมฺมํ ปฏิจฺจ.}

\prose{อาสวญฺเจว อาสว\-สมฺปยุตฺตญฺจ นเหตุํ ธมฺมํ ปฏิจฺจ ฯเปฯ. อาสว\-สมฺปยุตฺต\-ญฺเจว โน จ อาสวํ นเหตุํ ธมฺมํ ปฏิจฺจ.}

\proseitem{84}{อาสว\-วิปฺปยุตฺตํ สาสวํ เหตุํ ธมฺมํ ปฏิจฺจ ฯเปฯ. อาสว\-วิปฺปยุตฺตํ อนาสวํ เหตุํ ธมฺมํ ปฏิจฺจ.}

\prose{อาสว\-วิปฺปยุตฺตํ สาสวํ นเหตุํ ธมฺมํ ปฏิจฺจ ฯเปฯ. อาสว\-วิปฺปยุตฺตํ อนาสวํ นเหตุํ ธมฺมํ ปฏิจฺจ.}

\csromanmarkh{55. มหนฺตรทุก}\csromanmarkhii{1. เหตุทุก 20. สญฺโญชนา\-ทิทุก 1. เหตุทุก}\csromanheadernomark{2}{55. มหนฺตรทุก 1. เหตุทุก 20. สญฺโญชนา\-ทิทุก 1. เหตุทุก}
\proseitem{85}{\csromanfolio{59}สญฺโญชนํ เหตุํ ธมฺมํ ปฏิจฺจ ฯเปฯ. คนฺถํ เหตุํ ธมฺมํ ปฏิจฺจ ฯเปฯ.}

\prose{โอฆํ เหตุํ ธมฺมํ ปฏิจฺจ ฯเปฯ. โยคํ เหตุํ ธมฺมํ ปฏิจฺจ ฯเปฯ.}

\prose{นีวรณํ เหตุํ ธมฺมํ ปฏิจฺจ. (สํขิตฺตํ.)}

\prose{โนปรามาสํ เหตุํ ธมฺมํ ปฏิจฺจ. (สพฺพตฺถ เอกํ. สํขิตฺตํ.)}

\csromanmarkh{55. มหนฺตรทุก}\csromanmarkhii{1. เหตุทุก}\csromanheadernomark{2}{55. มหนฺตรทุก 1. เหตุทุก}
\proseitem{86}{สารมฺมณํ เหตุํ ธมฺมํ ปฏิจฺจ. (สพฺพตฺถ เอกํ.) สารมฺมณํ นเหตุํ ธมฺมํ ปฏิจฺจ. (สํขิตฺตํ.)}

\proseitem{87}{โนจิตฺตํ เหตุํ ธมฺมํ ปฏิจฺจ.}

\prose{เจตสิกํ เหตุํ ธมฺมํ ปฏิจฺจ.}

\prose{จิตฺต\-สมฺปยุตฺตํ เหตุํ ธมฺมํ ปฏิจฺจ.}

\prose{จิตฺต\-สํสฏฺฐํ เหตุํ ธมฺมํ ปฏิจฺจ.}

\prose{จิตฺต\-สมุฏฺฐานํ เหตุํ ธมฺมํ ปฏิจฺจ.}

\prose{จิตฺตสหภุํ เหตุํ ธมฺมํ ปฏิจฺจ.}

\prose{จิตฺตา\-นุปริวตฺติํ เหตุํ ธมฺมํ ปฏิจฺจ.}

\prose{จิตฺต\-สํสฏฺฐ\-สมุฏฺฐานํ เหตุํ ธมฺมํ ปฏิจฺจ.}

\prose{จิตฺต\-สํสฏฺฐ\-สมุฏฺฐาน\-สหภุํ เหตุํ ธมฺมํ ปฏิจฺจ.}

\prose{จิตฺต\-สํสฏฺฐ\-สมุฏฺฐานา\-นุปริวตฺติํ เหตุํ ธมฺมํ ปฏิจฺจ.}

\prose{พาหิรํ เหตุํ ธมฺมํ ปฏิจฺจ.}

\prose{โนอุปาทา เหตุํ ธมฺมํ ปฏิจฺจ.}

\prose{อุปาทินฺนํ เหตุํ ธมฺมํ ปฏิจฺจ ฯเปฯ. อนุปาทินฺนํ เหตุํ ธมฺมํ ปฏิจฺจ.}

\csromanmarkh{69. อุปาทาน\-โคจฺฉก}\csromanmarkhii{1. เหตุทุก}\csromanheadernomark{2}{69. อุปาทาน\-โคจฺฉก 1. เหตุทุก}
\vspace{\tipitakaparskip}
\csromancenter{อุปาทานํ เหตุํ ธมฺมํ ปฏิจฺจ ฯเปฯ.}
\csromanmarkh{75. กิเลสโคจฺฉก}\csromanmarkhii{1. เหตุทุก}\csromanheadernomark{2}{75. \textbf{กิเลสโคจฺฉก 1. เหตุทุก}}
\vspace{\tipitakaparskip}
\csromancenter{กิเลสํ เหตุํ ธมฺมํ ปฏิจฺจ ฯเปฯ.}
\csromanmarkh{83. ปิฏฺฐิทุก}\csromanmarkhii{1. เหตุทุก}\csromanheadernomark{2}{83. \textbf{ปิฏฺฐิทุก 1. เหตุทุก}}
\proseitem{90}{\csromanfolio{60}ทสฺสเนน ปหาตพฺพํ เหตุํ ธมฺมํ ปฏิจฺจ ฯเปฯ. นทสฺสเนน ปหาตพฺพํ เหตุํ ธมฺมํ ปฏิจฺจ.}

\proseitem{91}{ภาวนาย ปหาตพฺพํ เหตุํ ธมฺมํ ปฏิจฺจ ฯเปฯ. นภาวนาย ปหาตพฺพํ เหตุํ ธมฺมํ ปฏิจฺจ.}

\proseitem{92}{ทสฺสเนน ปหาตพฺพ\-เหตุกํ เหตุํ ธมฺมํ ปฏิจฺจ ฯเปฯ. นทสฺสเนน ปหาตพฺพ\-เหตุกํ เหตุํ ธมฺมํ ปฏิจฺจ.}

\proseitem{93}{ภาวนาย ปหาตพฺพ\-เหตุกํ เหตุํ ธมฺมํ ปฏิจฺจ ฯเปฯ. นภาวนาย ปหาตพฺพ\-เหตุกํ เหตุํ ธมฺมํ ปฏิจฺจ.}

\proseitem{94}{สวิตกฺกํ เหตุํ ธมฺมํ ปฏิจฺจ ฯเปฯ. อวิตกฺกํ เหตุํ ธมฺมํ ปฏิจฺจ ฯเปฯ.}

\prose{สวิจารํ เหตุํ ธมฺมํ ปฏิจฺจ ฯเปฯ. อวิจารํ เหตุํ ธมฺมํ ปฏิจฺจ ฯเปฯ.}

\prose{สปฺปีติกํ เหตุํ ธมฺมํ ปฏิจฺจ ฯเปฯ. อปฺปีติกํ เหตุํ ธมฺมํ ปฏิจฺจ. (สํขิตฺตํ.)}

\prose{ปีติสหคตํ เหตุํ ธมฺมํ ปฏิจฺจ ฯเปฯ. นปีติสหคตํ เหตุํ ธมฺมํ ปฏิจฺจ ฯเปฯ.}

\prose{สุขสหคตํ เหตุํ ธมฺมํ ปฏิจฺจ ฯเปฯ. นสุขสหคตํ เหตุํ ธมฺมํ ปฏิจฺจ ฯเปฯ.}

\prose{อุเปกฺขา\-สหคตํ เหตุํ ธมฺมํ ปฏิจฺจ ฯเปฯ. นอุเปกฺขา\-สหคตํ เหตุํ ธมฺมํ ปฏิจฺจ ฯเปฯ.}

\prose{กามาวจรํ เหตุํ ธมฺมํ ปฏิจฺจ ฯเปฯ. นกามาวจรํ เหตุํ ธมฺมํ ปฏิจฺจ ฯเปฯ.}

\csromanmarkh{83. ปิฏฺฐิทุก}\csromanmarkhii{1. เหตุทุก}\csromanheadernomark{2}{83. ปิฏฺฐิทุก 1. เหตุทุก}
\prose{\csromanfolio{61}รูปาวจรํ เหตุํ ธมฺมํ ปฏิจฺจ ฯเปฯ. นรูปาวจรํ เหตุํ ธมฺมํ ปฏิจฺจ ฯเปฯ.}

\prose{อรูปาวจรํ เหตุํ ธมฺมํ ปฏิจฺจ ฯเปฯ. นอรูปาวจรํ เหตุํ ธมฺมํ ปฏิจฺจ ฯเปฯ.}

\prose{ปริยาปนฺนํ เหตุํ ธมฺมํ ปฏิจฺจ ฯเปฯ. อปริยาปนฺนํ เหตุํ ธมฺมํ ปฏิจฺจ ฯเปฯ.}

\prose{นิยฺยานิกํ เหตุํ ธมฺมํ ปฏิจฺจ ฯเปฯ. อนิยฺยานิกํ เหตุํ ธมฺมํ ปฏิจฺจ ฯเปฯ.}

\prose{นิยตํ เหตุํ ธมฺมํ ปฏิจฺจ ฯเปฯ. อนิยตํ เหตุํ ธมฺมํ ปฏิจฺจ ฯเปฯ.}

\prose{สอุตฺตรํ เหตุํ ธมฺมํ ปฏิจฺจ ฯเปฯ. อนุตฺตรํ เหตุํ ธมฺมํ ปฏิจฺจ ฯเปฯ. (สพฺพตฺถ เทฺว.)}

\prose{สรณํ เหตุํ ธมฺมํ ปฏิจฺจ ฯเปฯ. อรณํ เหตุํ ธมฺมํ ปฏิจฺจ อรโณ เหตุ ธมฺโม อุปฺปชฺชติ เหตุปจฺจยา. (สํขิตฺตํ.)}

\prose{อรณํ นเหตุํ ธมฺมํ ปฏิจฺจ อรโณ นเหตุ ธมฺโม อุปฺปชฺชติ เหตุปจฺจยา.}

\prosewithcloser{(สพฺพตฺถ วิตฺถาโร.)}{ธมฺมานุโลเม ทุกทุก\-ปฏฺฐานํ นิฏฺฐิตํ.}
\nitthitam{อนุโลม\-ปฏฺฐานํ นิฏฺฐิตํ.}
\csromanpalirecto
\pitaka{อภิธมฺม\-ปิฏก}
\csromanheader{2}{ธมฺม\-ปจฺจนีย}
\csromanmatikaanchor{mtk.2}
\csromanmatikaanchor{mtk.6}
\csromantocmarknopage{chapter}{ติกติกปฏฺฐานปาฬิ}{mtk.2}
\bookmark[dest={mtk.2},level=0]{ติกติกปฏฺฐานปาฬิ}
\csromantocmarknopage{chapter}{ทุกทุกปฏฺฐานปาฬิ}{mtk.3}
\bookmark[dest={mtk.3},level=0]{ทุกทุกปฏฺฐานปาฬิ}
\gambhira{ติกปฏฺฐาน\-ปาฬิ}
\namakkaram{นโม ตสฺส ภควโต อรหโต สมฺมา\-สมฺพุทฺธสฺส.}
\csromanmarkh{1. กุสลตฺติก}\csromanmarkhii{1. ปฏิจฺจวาร}\csromanheadernomark{2}{1. \textbf{กุสลตฺติก 1. ปฏิจฺจวาร}}
\prose{\csromanfolio{63}ปจฺจยจตุกฺกเหตุอารมฺมณ}

\proseitem{1}{นกุสลํ ธมฺมํ ปฏิจฺจ นกุสโล ธมฺโม อุปฺปชฺชติ เหตุปจฺจยา. อกุสลํ อพฺยากตํ เอกํ ขนฺธํ ปฏิจฺจ อกุสลา อพฺยากตา ตโย ขนฺธา จิตฺต\-สมุฏฺฐานญฺจ รูปํ ฯเปฯ เทฺว ขนฺเธ ปฏิจฺจ เทฺว ขนฺธา จิตฺต\-สมุฏฺฐานญฺจ รูปํ.  นกุสลํ ธมฺมํ ปฏิจฺจ นอกุสโล ธมฺโม อุปฺปชฺชติ เหตุปจฺจยา.  นกุสลํ ธมฺมํ ปฏิจฺจ นอพฺยากโต ธมฺโม อุปฺปชฺชติ เหตุปจฺจยา.  นกุสลํ ธมฺมํ ปฏิจฺจ นกุสโล จ นอพฺยากโต จ ธมฺมา อุปฺปชฺชนฺติ เหตุปจฺจยา.  นกุสลํ ธมฺมํ ปฏิจฺจ นกุสโล จ นอกุสโล จ ธมฺมา อุปฺปชฺชนฺติ เหตุปจฺจยา. ปญฺจ.}

\proseitem{2}{นอกุสลํ ธมฺมํ ปฏิจฺจ นอกุสโล ธมฺโม อุปฺปชฺชติ เหตุปจฺจยา.  นอกุสลํ ธมฺมํ ปฏิจฺจ นกุสโล ธมฺโม อุปฺปชฺชติ เหตุปจฺจยา.  นอกุสลํ ธมฺมํ ปฏิจฺจ นอพฺยากโต ธมฺโม อุปฺปชฺชติ เหตุปจฺจยา.  นอกุสลํ ธมฺมํ ปฏิจฺจ นอกุสโล จ นอพฺยากโต จ ธมฺมา อุปฺปชฺชนฺติ เหตุปจฺจยา.  นอกุสลํ ธมฺมํ ปฏิจฺจ นกุสโล จ นอกุสโล จ ธมฺมา อุปฺปชฺชนฺติ เหตุปจฺจยา. ปญฺจ.}

\csromanpalirecto
\csromanmatikaanchor{mtk.4}
\csromantocmarknopage{chapter}{ติกปฏฺฐานปาฬิ}{mtk.4}
\bookmark[dest={mtk.4},level=0]{ติกปฏฺฐานปาฬิ}
\gambhira{ธมฺม\-ปจฺจนีย ติกปฏฺฐาน\-ปาฬิ}
\proseitem{3}{\csromanfolio{64}นอพฺยากตํ ธมฺมํ ปฏิจฺจ นอพฺยากโต ธมฺโม อุปฺปชฺชติ เหตุปจฺจยา.  นอพฺยากตํ ธมฺมํ ปฏิจฺจ นกุสโล ธมฺโม อุปฺปชฺชติ เหตุปจฺจยา.  นอพฺยากตํ ธมฺมํ ปฏิจฺจ นอกุสโล ธมฺโม อุปฺปชฺชติ เหตุปจฺจยา.  นอพฺยากตํ ธมฺมํ ปฏิจฺจ นกุสโล จ นอพฺยากโต จ ธมฺมา อุปฺปชฺชนฺติ เหตุปจฺจยา.  นอพฺยากตํ ธมฺมํ ปฏิจฺจ นอกุสโล จ นอพฺยากโต จ ธมฺมา อุปฺปชฺชนฺติ เหตุปจฺจยา.  นอพฺยากตํ ธมฺมํ ปฏิจฺจ นกุสโล จ นอกุสโล จ ธมฺมา อุปฺปชฺชนฺติ เหตุปจฺจยา. ฉ.}

\proseitem{4}{นกุสลญฺจ นอพฺยากตญฺจ ธมฺมํ ปฏิจฺจ นกุสโล ธมฺโม อุปฺปชฺชติ เหตุปจฺจยา.  นกุสลญฺจ นอพฺยากตญฺจ ธมฺมํ ปฏิจฺจ นอกุสโล ธมฺโม อุปฺปชฺชติ เหตุปจฺจยา.  นกุสลญฺจ นอพฺยากตญฺจ ธมฺมํ ปฏิจฺจ นอพฺยากโต ธมฺโม อุปฺปชฺชติ เหตุปจฺจยา.  นกุสลญฺจ นอพฺยากตญฺจ ธมฺมํ ปฏิจฺจ นกุสโล จ นอพฺยากโต จ ธมฺมา อุปฺปชฺชนฺติ เหตุปจฺจยา.  นกุสลญฺจ นอพฺยากตญฺจ ธมฺมํ ปฏิจฺจ นกุสโล จ นอกุสโล จ ธมฺมา อุปฺปชฺชนฺติ เหตุปจฺจยา. ปญฺจ.}

\proseitem{5}{นอกุสลญฺจ นอพฺยากตญฺจ ธมฺมํ ปฏิจฺจ นกุสโล ธมฺโม อุปฺปชฺชติ เหตุปจฺจยา.  นอกุสลญฺจ นอพฺยากตญฺจ ธมฺมํ ปฏิจฺจ นอกุสโล ธมฺโม อุปฺปชฺชติ เหตุปจฺจยา.  นอกุสลญฺจ นอพฺยากตญฺจ ธมฺมํ ปฏิจฺจ นอพฺยากโต ธมฺโม อุปฺปชฺชติ เหตุปจฺจยา.  นอกุสลญฺจ นอพฺยากตญฺจ ธมฺมํ ปฏิจฺจ นอกุสโล จ นอพฺยากโต จ ธมฺมา อุปฺปชฺชนฺติ เหตุปจฺจยา.  นอกุสลญฺจ นอพฺยากตญฺจ ธมฺมํ ปฏิจฺจ นกุสโล จ นอกุสโล จ ธมฺมา อุปฺปชฺชนฺติ เหตุปจฺจยา. ปญฺจ.}

\proseitem{6}{นกุสลญฺจ นอกุสลญฺจ ธมฺมํ ปฏิจฺจ นกุสโล ธมฺโม อุปฺปชฺชติ เหตุปจฺจยา.  นกุสลญฺจ นอกุสลญฺจ ธมฺมํ ปฏิจฺจ นอกุสโล ธมฺโม อุปฺปชฺชติ เหตุปจฺจยา.  นกุสลญฺจ นอกุสลญฺจ ธมฺมํ ปฏิจฺจ นกุสโล จ นอกุสโล จ ธมฺมา อุปฺปชฺชนฺติ เหตุปจฺจยา. ตีณิ.}

\prose{นกุสลํ ธมฺมํ ปฏิจฺจ นกุสโล ธมฺโม อุปฺปชฺชติ อารมฺมณ\-ปจฺจยา. (สํขิตฺตํ.)}

\proseitem{7}{เหตุยา เอกูนติํส, อารมฺมเณ จตุวีส ฯเปฯ อวิคเต เอกูนติํส.}

\prose{\csromanfolio{65}(สหชาตวารมฺปิ ปจฺจยวารมฺปิ นิสฺสยวารมฺปิ สํสฏฺฐ\-วารมฺปิ สมฺปยุตฺตวารมฺปิ ปฏิจฺจ\-วาร\-สทิสํ.)}

\csromanmarkh{1. กุสลตฺติก}\csromanmarkhii{7. ปญฺหาวาร}\csromanheadernomark{2}{1. กุสลตฺติก 7. ปญฺหาวาร}
\prose{ปจฺจยจตุกฺกเหตุอารมฺมณาทิ}

\proseitem{8}{นกุสโล ธมฺโม นกุสลสฺส ธมฺมสฺส เหตุปจฺจเยน ปจฺจโย ฯเปฯ. นกุสโล ธมฺโม นกุสลสฺส ธมฺมสฺส อารมฺมณ\-ปจฺจเยน ปจฺจโย.  นกุสโล ธมฺโม นอกุสลสฺส ธมฺมสฺส อารมฺมณ\-ปจฺจเยน ปจฺจโย.  นกุสโล ธมฺโม นอพฺยากตสฺส ธมฺมสฺส อารมฺมณ\-ปจฺจเยน ปจฺจโย.  นกุสโล ธมฺโม นกุสลสฺส จ นอพฺยากตสฺส จ ธมฺมสฺส อารมฺมณ\-ปจฺจเยน ปจฺจโย.  นกุสโล ธมฺโม นอกุสลสฺส จ นอพฺยากตสฺส จ ธมฺมสฺส อารมฺมณ\-ปจฺจเยน ปจฺจโย.  นกุสโล ธมฺโม นกุสลสฺส จ นอกุสลสฺส จ ธมฺมสฺส อารมฺมณ\-ปจฺจเยน ปจฺจโย. ฉ.}

\prose{นอกุสโล ธมฺโม นอกุสลสฺส ธมฺมสฺส อารมฺมณ\-ปจฺจเยน ปจฺจโย. ฉ.}

\prose{นอพฺยากโต ธมฺโม นอพฺยากตสฺส ธมฺมสฺส อารมฺมณ\-ปจฺจเยน ปจฺจโย. ฉ.}

\prose{นกุสโล จ นอพฺยากโต จ ธมฺมา นกุสลสฺส ธมฺมสฺส อารมฺมณ\-ปจฺจเยน ปจฺจโย. ฉ.}

\prose{นอกุสโล จ นอพฺยากโต จ ธมฺมา นกุสลสฺส ธมฺมสฺส อารมฺมณ\-ปจฺจเยน ปจฺจโย. ฉ.}

\prose{นกุสโล จ นอกุสโล จ ธมฺมา นกุสลสฺส ธมฺมสฺส อารมฺมณ\-ปจฺจเยน ปจฺจโย. ฉ.}

\proseitem{9}{นกุสโล ธมฺโม นกุสลสฺส ธมฺมสฺส อธิปติ\-ปจฺจเยน ปจฺจโย. อนนฺตร\-ปจฺจเยน ปจฺจโย. สมนนฺตร\-ปจฺจเยน ปจฺจโย. สหชาต\-ปจฺจเยน ปจฺจโย. อญฺญมญฺญ\-ปจฺจเยน ปจฺจโย. นิสฺสย\-ปจฺจเยน ปจฺจโย. อุปนิสฺสย\-ปจฺจเยน ปจฺจโย. ปุเรชาต\-ปจฺจเยน ปจฺจโย. นกุสโล ธมฺโม นอกุสลสฺส ธมฺมสฺส ปุเรชาต\-ปจฺจเยน ปจฺจโย ฯเปฯ.  นกุสโล ธมฺโม นกุสลสฺส จ นอกุสลสฺส จ ธมฺมสฺส ปุเรชาต\-ปจฺจเยน ปจฺจโย. ฉ.}

\prose{\csromanfolio{66}นอกุสโล ธมฺโม นอกุสลสฺส ธมฺมสฺส ปุเรชาต\-ปจฺจเยน ปจฺจโย.  นอกุสโล ธมฺโม นกุสลสฺส ธมฺมสฺส ปุเรชาต\-ปจฺจเยน ปจฺจโย ฯเปฯ.  นอกุสโล ธมฺโม นกุสลสฺส จ นอกุสลสฺส จ ธมฺมสฺส ปุเรชาต\-ปจฺจเยน ปจฺจโย. ฉ.}

\prose{นกุสโล จ นอกุสโล จ ธมฺมา นกุสลสฺส ธมฺมสฺส ปุเรชาต\-ปจฺจเยน ปจฺจโย.  นกุสโล จ นอกุสโล จ ธมฺมา นอกุสลสฺส ธมฺมสฺส ปุเรชาต\-ปจฺจเยน ปจฺจโย ฯเปฯ.  นกุสโล จ นอกุสโล จ ธมฺมา นกุสลสฺส จ นอกุสลสฺส จ ธมฺมสฺส ปุเรชาต\-ปจฺจเยน ปจฺจโย. ฉ. (สํขิตฺตํ.)}

\proseitem{10}{เหตุยา เอกูนติํส, อารมฺมเณ ฉตฺติํส, อธิปติยา ปญฺจติํส, อนนฺตเร จตุตฺติํส, สมนนฺตเร จตุตฺติํส, สหชาเต เอกูนติํส, อญฺญมญฺเญ จตุวีส, นิสฺสเย จตุตฺติํส, อุปนิสฺสเย ฉตฺติํส, ปุเรชาเต อฏฺฐารส, ปจฺฉาชาเต อฏฺฐารส, อาเสวเน จตุวีส, กมฺเม เอกูนติํส, วิปาเก นว, อาหาเร เอกูนติํส ฯเปฯ อวิคเต จตุตฺติํส.}

\prose{(ยถา กุสลตฺติเก ปญฺหาวารสฺส อนุโลมมฺปิ ปจฺจนียมฺปิ อนุโลมปจฺจนียมฺปิ ปจฺจนียานุโลมมฺปิ คณิตํ, เอวํ คเณตพฺพํ.)}

\csromanmarkh{2. เวทนาตฺติก}\csromanmarkhii{1. ปฏิจฺจวาร}\csromanheadernomark{2}{2. \textbf{เวทนาตฺติก 1. ปฏิจฺจวาร}}
\proseitem{11}{นสุขาย เวทนาย สมฺปยุตฺตํ ธมฺมํ ปฏิจฺจ นสุขาย เวทนาย สมฺปยุตฺโต ธมฺโม อุปฺปชฺชติ เหตุปจฺจยา.  นสุขาย เวทนาย สมฺปยุตฺตํ ธมฺมํ ปฏิจฺจ นทุกฺขาย เวทนาย สมฺปยุตฺโต ธมฺโม อุปฺปชฺชติ เหตุปจฺจยา.  นสุขาย เวทนาย สมฺปยุตฺตํ ธมฺมํ ปฏิจฺจ นอทุกฺขม\-สุขาย เวทนาย สมฺปยุตฺโต ธมฺโม อุปฺปชฺชติ เหตุปจฺจยา.  นสุขาย เวทนาย สมฺปยุตฺตํ ธมฺมํ ปฏิจฺจ นสุขาย เวทนาย สมฺปยุตฺโต จ นอทุกฺขม\-สุขาย เวทนาย สมฺปยุตฺโต จ ธมฺมา อุปฺปชฺชนฺติ เหตุปจฺจยา.  นสุขาย เวทนาย สมฺปยุตฺตํ ธมฺมํ ปฏิจฺจ นทุกฺขาย เวทนาย สมฺปยุตฺโต จ นอทุกฺขม\-สุขาย เวทนาย สมฺปยุตฺโต จ ธมฺมา อุปฺปชฺชนฺติ เหตุปจฺจยา.  นสุขาย เวทนาย สมฺปยุตฺตํ ธมฺมํ ปฏิจฺจ นสุขาย เวทนาย สมฺปยุตฺโต จ นทุกฺขาย เวทนาย สมฺปยุตฺโต จ ธมฺมา อุปฺปชฺชนฺติ เหตุปจฺจยา.  }

\vspace{\tipitakaparskip}
\csromancenter{สํกิลิฏฺฐ\-ตฺติก นสุขาย เวทนาย สมฺปยุตฺตํ ธมฺมํ ปฏิจฺจ นสุขาย เวทนาย สมฺปยุตฺโต จ นทุกฺขาย เวทนาย สมฺปยุตฺโต จ นอทุกฺขม\-สุขาย เวทนาย สมฺปยุตฺโต จ ธมฺมา อุปฺปชฺชนฺติ เหตุปจฺจยา. (7)}
\prose{\csromanfolio{67}นทุกฺขาย เวทนาย สมฺปยุตฺตํ ธมฺมํ ปฏิจฺจ นทุกฺขาย เวทนาย สมฺปยุตฺโต ธมฺโม อุปฺปชฺชติ เหตุปจฺจยา. สตฺต.}

\prose{นอทุกฺขม\-สุขาย เวทนาย สมฺปยุตฺตํ ธมฺมํ ปฏิจฺจ นอทุกฺขม\-สุขาย เวทนาย สมฺปยุตฺโต ธมฺโม อุปฺปชฺชติ เหตุปจฺจยา. สตฺต. (สํขิตฺตํ.)}

\csromanmarkh{3. วิปากตฺติก}\csromanmarkhii{1. ปฏิจฺจวาร}\csromanheadernomark{2}{3. \textbf{วิปากตฺติก 1. ปฏิจฺจวาร}}
\proseitem{12}{นวิปากํ ธมฺมํ ปฏิจฺจ นวิปาโก ธมฺโม อุปฺปชฺชติ เหตุปจฺจยา ฯเปฯ. นวิปาก\-ธมฺม\-ธมฺมํ ปฏิจฺจ นวิปาก\-ธมฺม\-ธมฺโม อุปฺปชฺชติ เหตุปจฺจยา ฯเปฯ.}

\prose{นเน\-ว\-วิปาก\-น\-วิปากธมฺม\-ธมฺมํ ปฏิจฺจ นเน\-ว\-วิปาก\-น\-วิปาก\-ธมฺม\-ธมฺโม อุปฺปชฺชติ เหตุปจฺจยา.}

\csromanmarkh{4. อุปาทินฺนตฺติก}\csromanmarkhii{1. ปฏิจฺจวาร}\csromanheadernomark{2}{4. \textbf{อุปาทินฺนตฺติก 1. ปฏิจฺจวาร}}
\proseitem{13}{นอุปาทินฺนุปาทานิยํ ธมฺมํ ปฏิจฺจ นอุปาทินฺนุปาทานิโย ธมฺโม อุปฺปชฺชติ เหตุปจฺจยา ฯเปฯ.}

\prose{นอนุปาทินฺนุ\-ปาทานิยํ ธมฺมํ ปฏิจฺจ นอนุปาทินฺนุ\-ปาทานิโย ธมฺโม อุปฺปชฺชติ เหตุปจฺจยา ฯเปฯ.}

\prose{นอนุปาทินฺน\-อนุปาทานิยํ ธมฺมํ ปฏิจฺจ นอนุปาทินฺน\-อนุปาทานิโย ธมฺโม อุปฺปชฺชติ เหตุปจฺจยา. (สํขิตฺตํ.)}

\csromanmarkh{5. สํกิลิฏฺฐ\-ตฺติก}\csromanmarkhii{1. ปฏิจฺจวาร}\csromanheadernomark{2}{5. สํกิลิฏฺฐ\-ตฺติก 1. ปฏิจฺจวาร}
\proseitem{14}{นสํกิลิฏฺฐ\-สํกิเลสิกํ ธมฺมํ ปฏิจฺจ นสํกิลิฏฺฐ\-สํกิเลสิโก ธมฺโม อุปฺปชฺชติ เหตุปจฺจยา ฯเปฯ.}

\prose{นอสํกิลิฏฺฐ\-สํกิเลสิกํ ธมฺมํ ปฏิจฺจ นอสํกิลิฏฺฐ\-สํกิเลสิโก ธมฺโม อุปฺปชฺชติ เหตุปจฺจยา ฯเปฯ.}

\prose{นอสํกิลิฏฺฐ\-อสํกิเลสิกํ ธมฺมํ ปฏิจฺจ นอสํกิลิฏฺฐ\-อสํกิเลสิโก ธมฺโม อุปฺปชฺชติ เหตุปจฺจยา. (สํขิตฺตํ.)}

\csromanmarkh{ธมฺม\-ปจฺจนีย ติกปฏฺฐาน\-ปาฬิ}\csromanmarkhii{6. วิตกฺกตฺติก 1. ปฏิจฺจวาร}\csromanheadernomark{2}{ธมฺม\-ปจฺจนีย ติกปฏฺฐาน\-ปาฬิ 6. วิตกฺกตฺติก 1. ปฏิจฺจวาร}
\proseitem{15}{\csromanfolio{68}นสวิตกฺก\-สวิจารํ ธมฺมํ ปฏิจฺจ นสวิตกฺก\-สวิจาโร ธมฺโม อุปฺปชฺชติ เหตุปจฺจยา ฯเปฯ.}

\prose{นอวิตกฺก\-วิจารมตฺตํ ธมฺมํ ปฏิจฺจ นอวิตกฺก\-วิจารมตฺโต ธมฺโม อุปฺปชฺชติ เหตุปจฺจยา ฯเปฯ.}

\prose{นอวิตกฺก\-อวิจารํ ธมฺมํ ปฏิจฺจ นอวิตกฺก\-อวิจาโร ธมฺโม อุปฺปชฺชติ เหตุปจฺจยา. (สํขิตฺตํ.) 7. ปีติตฺติก 1. ปฏิจฺจวาร}

\proseitem{16}{นปีติสหคตํ ธมฺมํ ปฏิจฺจ นปีติสหคโต ธมฺโม อุปฺปชฺชติ เหตุปจฺจยา ฯเปฯ.}

\prose{นสุขสหคตํ ธมฺมํ ปฏิจฺจ ฯเปฯ.  นฺเอาเปกฺขา\-สหคตํ ธมฺมํ ปฏิจฺจ. (สํขิตฺตํ.) 8. ทสฺสนตฺติก 1. ปฏิจฺจวาร}

\proseitem{17}{นทสฺสเนน ปหาตพฺพํ ธมฺมํ ปฏิจฺจ ฯเปฯ.  นภาวนาย ปหาตพฺพํ ธมฺมํ ปฏิจฺจ ฯเปฯ.  นเน\-ว\-ทสฺสเนน นภาวนาย ปหาตพฺพํ ธมฺมํ ปฏิจฺจ. (สํขิตฺตํ.) 9. ทสฺสน\-เหตุ\-ตฺติก 1. ปฏิจฺจวาร}

\proseitem{18}{นทสฺสเนน ปหาตพฺพ\-เหตุกํ ธมฺมํ ปฏิจฺจ ฯเปฯ.  นภาวนาย ปหาตพฺพ\-เหตุกํ ธมฺมํ ปฏิจฺจ ฯเปฯ.  นเน\-ว\-ทสฺสเนน นภาวนาย ปหาตพฺพ\-เหตุกํ ธมฺมํ ปฏิจฺจ. 10. อาจยคามิตฺติก 1. ปฏิจฺจวาร}

\proseitem{19}{นอาจยคามิํ ธมฺมํ ปฏิจฺจ ฯเปฯ. นอปจยคามิํ ธมฺมํ ปฏิจฺจ ฯเปฯ.  นเน\-วา\-จยคามิ\-นา\-ปจยคามิํ ธมฺมํ ปฏิจฺจ. (สํขิตฺตํ.) 11. เสกฺขตฺติก 1. ปฏิจฺจวาร}

\proseitem{20}{นเสกฺขํ ธมฺมํ ปฏิจฺจ ฯเปฯ.  นอเสกฺขํ ธมฺมํ ปฏิจฺจ ฯเปฯ.  นเน\-ว\-เสกฺข\-นาเสกฺขํ ธมฺมํ ปฏิจฺจ. (สํขิตฺตํ.)}

\csromanmarkh{17. อุปฺปนฺนตฺติก}\csromanmarkhii{12. ปริตฺตตฺติก 1. ปฏิจฺจวาร}\csromanheadernomark{2}{17. อุปฺปนฺนตฺติก \textbf{12. ปริตฺตตฺติก 1. ปฏิจฺจวาร}}
\proseitem{21}{\csromanfolio{69}นปริตฺตํ ธมฺมํ ปฏิจฺจ ฯเปฯ.  นมหคฺคตํ ธมฺมํ ปฏิจฺจ ฯเปฯ.  นอปฺปมาณํ ธมฺมํ ปฏิจฺจ. (สํขิตฺตํ.) 13. ปริตฺตา\-รมฺมณ\-ตฺติก 1. ปฏิจฺจวาร}

\proseitem{22}{นปริตฺตา\-รมฺมณํ ธมฺมํ ปฏิจฺจ ฯเปฯ.  นมหคฺคตา\-รมฺมณํ ธมฺมํ ปฏิจฺจ ฯเปฯ.  นอปฺปมาณา\-รมฺมณํ ธมฺมํ ปฏิจฺจ. (สํขิตฺตํ.) 14. หีนตฺติก 1. ปฏิจฺจวาร}

\proseitem{23}{นหีนํ ธมฺมํ ปฏิจฺจ ฯเปฯ.  นมชฺฌิมํ ธมฺมํ ปฏิจฺจ ฯเปฯ.  นปณีตํ ธมฺมํ ปฏิจฺจ. (สํขิตฺตํ.) 15. มิจฺฉตฺตตฺติก 1. ปฏิจฺจวาร}

\proseitem{24}{นมิจฺฉตฺต\-นิยตํ ธมฺมํ ปฏิจฺจ ฯเปฯ.  นสมฺมตฺต\-นิยตํ ธมฺมํ ปฏิจฺจ ฯเปฯ.  นอนิยตํ ธมฺมํ ปฏิจฺจ. (สํขิตฺตํ.) 16. มคฺคา\-รมฺมณ\-ตฺติก 1. ปฏิจฺจวาร}

\proseitem{25}{นมคฺคา\-รมฺมณํ ธมฺมํ ปฏิจฺจ ฯเปฯ.  นมคฺคเหตุกํ ธมฺมํ ปฏิจฺจ ฯเปฯ.  นมคฺคา\-ธิปติํ ธมฺมํ ปฏิจฺจ. 17. อุปฺปนฺนตฺติก 1. ปฏิจฺจวาร}

\proseitem{26}{นอนุปฺปนฺนํ ธมฺมํ ปฏิจฺจ นอนุปฺปนฺโน ธมฺโม อุปฺปชฺชติ เหตุปจฺจยา.  นอนุปฺปนฺนํ ธมฺมํ ปฏิจฺจ นอุปฺปาที ธมฺโม อุปฺปชฺชติ เหตุปจฺจยา.  นอนุปฺปนฺนํ ธมฺมํ ปฏิจฺจ นอนุปฺปนฺโน จ นอุปฺปาที จ ธมฺมา อุปฺปชฺชนฺติ เหตุปจฺจยา. (3)}

\prose{นอุปฺปาทิํ ธมฺมํ ปฏิจฺจ นอุปฺปาที ธมฺโม อุปฺปชฺชติ เหตุปจฺจยา.  นอุปฺปาทิํ ธมฺมํ ปฏิจฺจ นอนุปฺปนฺโน ธมฺโม อุปฺปชฺชติ เหตุปจฺจยา.  นอุปฺปาทิํ ธมฺมํ ปฏิจฺจ นอนุปฺปนฺโน จ นอุปฺปาที จ ธมฺมา อุปฺปชฺชนฺติ เหตุปจฺจยา. (3)}

\prose{นอนุปฺปนฺนญฺจ นอุปฺปาทิญฺจ ธมฺมํ ปฏิจฺจ นอนุปฺปนฺโน ธมฺโม อุปฺปชฺชติ เหตุปจฺจยา.  นอนุปฺปนฺนญฺจ นอุปฺปาทิญฺจ ธมฺมํ ปฏิจฺจ นอุปฺปาที ธมฺโม อุปฺปชฺชติ เหตุปจฺจยา.  นอนุปฺปนฺนญฺจ นอุปฺปาทิญฺจ ธมฺมํ ปฏิจฺจ นอนุปฺปนฺโน จ นอุปฺปาที จ ธมฺมา อุปฺปชฺชนฺติ เหตุปจฺจยา. (3) (สํขิตฺตํ.)}

\csromanmarkh{ธมฺม\-ปจฺจนีย ติกปฏฺฐาน\-ปาฬิ}\csromanmarkhii{18. อตีตตฺติก 1-7. ปฏิจฺจาทิวาร}\csromanheadernomark{2}{ธมฺม\-ปจฺจนีย ติกปฏฺฐาน\-ปาฬิ 18. อตีตตฺติก 1-7. ปฏิจฺจาทิวาร}
\proseitem{27}{\csromanfolio{70}นอตีตํ ธมฺมํ ปฏิจฺจ นอตีโต ธมฺโม อุปฺปชฺชติ เหตุปจฺจยา. (สํขิตฺตํ.)}

\prose{นอตีโต ธมฺโม นอตีตสฺส ธมฺมสฺส เหตุปจฺจเยน ปจฺจโย.  นอตีโต ธมฺโม นอนาคตสฺส ธมฺมสฺส เหตุปจฺจเยน ปจฺจโย.  นอตีโต ธมฺโม นอตีตสฺส จ นอนาคตสฺส จ ธมฺมสฺส เหตุปจฺจเยน ปจฺจโย. (3)}

\prose{นอนาคโต ธมฺโม นอนาคตสฺส ธมฺมสฺส เหตุปจฺจเยน ปจฺจโย.  นอนาคโต ธมฺโม นอตีตสฺส ธมฺมสฺส เหตุปจฺจเยน ปจฺจโย.  นอนาคโต ธมฺโม นอตีตสฺส จ นอนาคตสฺส จ ธมฺมสฺส เหตุปจฺจเยน ปจฺจโย. (3)}

\prose{นอตีโต จ นอนาคโต จ ธมฺมา นอตีตสฺส ธมฺมสฺส เหตุปจฺจเยน ปจฺจโย.  นอตีโต จ นอนาคโต จ ธมฺมา นอนาคตสฺส ธมฺมสฺส เหตุปจฺจเยน ปจฺจโย.  นอตีโต จ นอนาคโต จ ธมฺมา นอตีตสฺส จ นอนาคตสฺส จ ธมฺมสฺส เหตุปจฺจเยน ปจฺจโย. (3)}

\csromanmarkh{19. อตีตา\-รมฺมณ\-ตฺติก}\csromanmarkhii{1. ปฏิจฺจวาร}\csromanheadernomark{2}{19. อตีตา\-รมฺมณ\-ตฺติก 1. ปฏิจฺจวาร}
\proseitem{28}{นอตีตา\-รมฺมณํ ธมฺมํ ปฏิจฺจ ฯเปฯ. นอนาคตา\-รมฺมณํ ธมฺมํ ปฏิจฺจ ฯเปฯ. นปจฺจุปฺปนฺนา\-รมฺมณํ ธมฺมํ ปฏิจฺจ.}

\csromanmarkh{20. อชฺฌตฺตตฺติก}\csromanmarkhii{1. ปฏิจฺจวาร}\csromanheadernomark{2}{20. \textbf{อชฺฌตฺตตฺติก 1. ปฏิจฺจวาร}}
\proseitem{29}{นอชฺฌตฺตํ ธมฺมํ ปฏิจฺจ นอชฺฌตฺโต ธมฺโม อุปฺปชฺชติ เหตุปจฺจยา. นพหิทฺธา ธมฺมํ ปฏิจฺจ นพหิทฺธา ธมฺโม อุปฺปชฺชติ เหตุปจฺจยา. (สํขิตฺตํ.)}

\csromanmarkh{21. อชฺฌตฺตา\-รมฺมณ\-ตฺติก}\csromanmarkhii{1. ปฏิจฺจวาร}\csromanheadernomark{2}{21. อชฺฌตฺตา\-รมฺมณ\-ตฺติก 1. ปฏิจฺจวาร}
\proseitem{30}{นอชฺฌตฺตา\-รมฺมณํ ธมฺมํ ปฏิจฺจ นอชฺฌตฺตา\-รมฺมโณ ธมฺโม อุปฺปชฺชติ เหตุปจฺจยา.  นอชฺฌตฺตา\-รมฺมณํ ธมฺมํ ปฏิจฺจ นพหิทฺธา\-รมฺมโณ ธมฺโม อุปฺปชฺชติ เหตุปจฺจยา. (2)}

\csromanheader{2}{22. สนิทสฺสน\-ตฺติก}
\prose{\csromanfolio{71}นพหิทฺธา\-รมฺมณํ ธมฺมํ ปฏิจฺจ นพหิทฺธา\-รมฺมโณ ธมฺโม อุปฺปชฺชติ เหตุปจฺจยา.  นพหิทฺธา\-รมฺมณํ ธมฺมํ ปฏิจฺจ นอชฺฌตฺตา\-รมฺมโณ ธมฺโม อุปฺปชฺชติ เหตุปจฺจยา. (2) (สํขิตฺตํ.)}

\csromanmarkh{22. สนิทสฺสน\-ตฺติก}\csromanmarkhii{1. ปฏิจฺจวาร}\csromanheadernomark{2}{22. สนิทสฺสน\-ตฺติก 1. ปฏิจฺจวาร}
\prosecont{นสนิทสฺสน\-สปฺปฏิฆํ ธมฺมํ ปฏิจฺจ นสนิทสฺสน\-สปฺปฏิโฆ ธมฺโม อุปฺปชฺชติ เหตุปจฺจยา.  นสนิทสฺสน\-สปฺปฏิฆํ ธมฺมํ ปฏิจฺจ นอนิทสฺสน\-สปฺปฏิโฆ ธมฺโม อุปฺปชฺชติ เหตุปจฺจยา.  นสนิทสฺสน\-สปฺปฏิฆํ ธมฺมํ ปฏิจฺจ นอนิทสฺสน\-อปฺปฏิโฆ ธมฺโม อุปฺปชฺชติ เหตุปจฺจยา.  นสนิทสฺสน\-สปฺปฏิฆํ ธมฺมํ ปฏิจฺจ นสนิทสฺสน\-สปฺปฏิโฆ จ นอนิทสฺสน\-อปฺปฏิโฆ จ ธมฺมา อุปฺปชฺชนฺติ เหตุปจฺจยา.  นสนิทสฺสน\-สปฺปฏิฆํ ธมฺมํ ปฏิจฺจ นอนิทสฺสน\-สปฺปฏิโฆ จ นอนิทสฺสน\-อปฺปฏิโฆ จ ธมฺมา อุปฺปชฺชนฺติ เหตุปจฺจยา.  นสนิทสฺสน\-สปฺปฏิฆํ ธมฺมํ ปฏิจฺจ นสนิทสฺสน\-สปฺปฏิโฆ จ นอนิทสฺสน\-สปฺปฏิโฆ จ ธมฺมา อุปฺปชฺชนฺติ เหตุปจฺจยา. (6)}

\proseitem{32}{นอนิทสฺสน\-สปฺปฏิฆํ ธมฺมํ ปฏิจฺจ นอนิทสฺสน\-สปฺปฏิโฆ ธมฺโม อุปฺปชฺชติ เหตุปจฺจยา.  นอนิทสฺสน\-สปฺปฏิฆํ ธมฺมํ ปฏิจฺจ นสนิทสฺสน\-สปฺปฏิโฆ ธมฺโม อุปฺปชฺชติ เหตุปจฺจยา.  นอนิทสฺสน\-สปฺปฏิฆํ ธมฺมํ ปฏิจฺจ นอนิทสฺสน\-อปฺปฏิโฆ ธมฺโม อุปฺปชฺชติ เหตุปจฺจยา.  นอนิทสฺสน\-สปฺปฏิฆํ ธมฺมํ ปฏิจฺจ นสนิทสฺสน\-สปฺปฏิโฆ จ นอนิทสฺสน\-อปฺปฏิโฆ จ ธมฺมา อุปฺปชฺชนฺติ เหตุปจฺจยา.  นอนิทสฺสน\-สปฺปฏิฆํ ธมฺมํ ปฏิจฺจ นอนิทสฺสน\-สปฺปฏิโฆ จ นอนิทสฺสน\-อปฺปฏิโฆ จ ธมฺมา อุปฺปชฺชนฺติ เหตุปจฺจยา.  นอนิทสฺสน\-สปฺปฏิฆํ ธมฺมํ ปฏิจฺจ นสนิทสฺสน\-สปฺปฏิโฆ จ นอนิทสฺสน\-สปฺปฏิโฆ จ ธมฺมา อุปฺปชฺชนฺติ เหตุปจฺจยา. (6)}

\proseitem{33}{นอนิทสฺสน\-อปฺปฏิฆํ ธมฺมํ ปฏิจฺจ นอนิทสฺสน\-อปฺปฏิโฆ ธมฺโม อุปฺปชฺชติ เหตุปจฺจยา.  นอนิทสฺสน\-อปฺปฏิฆํ ธมฺมํ ปฏิจฺจ นสนิทสฺสน\-สปฺปฏิโฆ ธมฺโม อุปฺปชฺชติ เหตุปจฺจยา.  นอนิทสฺสน\-อปฺปฏิฆํ ธมฺมํ ปฏิจฺจ นอนิทสฺสน\-สปฺปฏิโฆ ธมฺโม อุปฺปชฺชติ เหตุปจฺจยา. ฉ.}

\prose{นสนิทสฺสนสปฺปฏิฆญฺจ นอนิทสฺสนอปฺปฏิฆญฺจ ธมฺมํ ปฏิจฺจ นสนิทสฺสน\-สปฺปฏิโฆ ธมฺโม อุปฺปชฺชติ เหตุปจฺจยา. ฉ.}

\prose{\csromanfolio{72}นสนิทสฺสนสปฺปฏิฆญฺจ นอนิทสฺสนสปฺปฏิฆญฺจ ธมฺมํ ปฏิจฺจ นสนิทสฺสน\-สปฺปฏิโฆ ธมฺโม อุปฺปชฺชติ เหตุปจฺจยา. ฉ.}

\proseitem{34}{เหตุยา ติํส ฯเปฯ อวิคเต ติํส.}

\prosewithcloser{(ยถา กุสลตฺติเก สหชาตวารมฺปิ ปจฺจยวารมฺปิ นิสฺสยวารมฺปิ สํสฏฺฐ\-วารมฺปิ สมฺปยุตฺตวารมฺปิ ปญฺหาวารมฺปิ, คณิตํ, เอวํ คเณตพฺพํ.)}{ธมฺม\-ปจฺจนีเย ติกปฏฺฐานํ นิฏฺฐิตํ.}
\csromanpalirecto
\pitaka{อภิธมฺม\-ปิฏก}
\csromanmatikaanchor{mtk.5}
\csromantocmarknopage{chapter}{ทุกปฏฺฐานปาฬิ}{mtk.5}
\bookmark[dest={mtk.5},level=0]{ทุกปฏฺฐานปาฬิ}
\csromantocmarknopage{chapter}{ทุกติกปฏฺฐานปาฬิ}{mtk.6}
\bookmark[dest={mtk.6},level=0]{ทุกติกปฏฺฐานปาฬิ}
\gambhira{ทุกปฏฺฐาน\-ปาฬิ}
\namakkaram{นโม ตสฺส ภควโต อรหโต สมฺมา\-สมฺพุทฺธสฺส.}
\csromanheader{2}{1. \textbf{เหตุทุก 1-7. ปฏิจฺจาทิวาร}}
\prose{\csromanfolio{73}ปจฺจยจตุกฺกเหตุ}

\proseitem{1}{นเหตุํ ธมฺมํ ปฏิจฺจ นเหตุ ธมฺโม อุปฺปชฺชติ เหตุปจฺจยา. นเหตุํ เอกํ ขนฺธํ ปฏิจฺจ ตโย ขนฺธา จิตฺต\-สมุฏฺฐานญฺจ รูปํ ฯเปฯ เทฺว ขนฺเธ ปฏิจฺจ เทฺว ขนฺธา จิตฺต\-สมุฏฺฐานญฺจ รูปํ ปฏิสนฺธิ\-กฺขเณ ฯเปฯ. (ยาว มหาภูตา.) . นเหตุํ ธมฺมํ ปฏิจฺจ น นเหตุ ธมฺโม อุปฺปชฺชติ เหตุปจฺจยา. นเหตู ขนฺเธ ปฏิจฺจ เหตู. ปฏิสนฺธิ\-กฺขเณ ฯเปฯ.  นเหตุํ ธมฺมํ ปฏิจฺจ นเหตุ จ น นเหตุ จ ธมฺมา อุปฺปชฺชนฺติ เหตุปจฺจยา. (3)}

\prose{น นเหตุํ ธมฺมํ ปฏิจฺจ น นเหตุ ธมฺโม อุปฺปชฺชติ เหตุปจฺจยา. น นเหตุํ ธมฺมํ ปฏิจฺจ นเหตุ ธมฺโม อุปฺปชฺชติ เหตุปจฺจยา.  น นเหตุํ ธมฺมํ ปฏิจฺจ นเหตุ จ น นเหตุ จ ธมฺมา อุปฺปชฺชนฺติ เหตุปจฺจยา. (3)}

\prose{นเหตุญฺจ น นเหตุญฺจ ธมฺมํ ปฏิจฺจ นเหตุ ธมฺโม อุปฺปชฺชติ เหตุปจฺจยา.  นเหตุญฺจ น นเหตุญฺจ ธมฺมํ ปฏิจฺจ น นเหตุ ธมฺโม อุปฺปชฺชติ เหตุปจฺจยา.  นเหตุญฺจ น นเหตุญฺจ ธมฺมํ ปฏิจฺจ นเหตุ จ น นเหตุ จ ธมฺมา อุปฺปชฺชนฺติ เหตุปจฺจยา. (3)}

\proseitem{2}{เหตุยา นว, อารมฺมเณ นว ฯเปฯ อวิคเต นว.}

\prose{(สหชาตวารมฺปิ ฯเปฯ สมฺปยุตฺตวารมฺปิ ปฏิจฺจ\-วาร\-สทิสํ, เอวํ วิตฺถาเรตพฺพํ.)}

\csromanheader{2}{ธมฺม\-ปจฺจนีย ทุกปฏฺฐาน\-ปาฬิ เหตุอารมฺมณ}
\proseitem{3}{\csromanfolio{74}น นเหตุ ธมฺโม น นเหตุสฺส ธมฺมสฺส เหตุปจฺจเยน ปจฺจโย.  น นเหตุ ธมฺโม นเหตุสฺส ธมฺมสฺส เหตุปจฺจเยน ปจฺจโย.  น นเหตุ ธมฺโม นเหตุสฺส จ น นเหตุสฺส จ ธมฺมสฺส เหตุปจฺจเยน ปจฺจโย. (3)}

\prose{นเหตุ ธมฺโม นเหตุสฺส ธมฺมสฺส อารมฺมณ\-ปจฺจเยน ปจฺจโย. นเหตุ ธมฺโม น นเหตุสฺส ธมฺมสฺส อารมฺมณ\-ปจฺจเยน ปจฺจโย.  นเหตุ ธมฺโม นเหตุสฺส จ น นเหตุสฺส จ ธมฺมสฺส อารมฺมณ\-ปจฺจเยน ปจฺจโย. (3)}

\prose{น นเหตุ ธมฺโม น นเหตุสฺส ธมฺมสฺส อารมฺมณ\-ปจฺจเยน ปจฺจโย.  น นเหตุ ธมฺโม นเหตุสฺส ธมฺมสฺส อารมฺมณ\-ปจฺจเยน ปจฺจโย.  น นเหตุ ธมฺโม นเหตุสฺส จ น นเหตุสฺส จ ธมฺมสฺส อารมฺมณ\-ปจฺจเยน ปจฺจโย.  (3)}

\prose{นเหตุ จ น นเหตุ จ ธมฺมา นเหตุสฺส ธมฺมสฺส อารมฺมณ\-ปจฺจเยน ปจฺจโย.  นเหตุ จ น นเหตุ จ ธมฺมา น นเหตุสฺส ธมฺมสฺส อารมฺมณ\-ปจฺจเยน ปจฺจโย.  นเหตุ จ น นเหตุ จ ธมฺมา นเหตุสฺส จ น นเหตุสฺส จ ธมฺมสฺส อารมฺมณ\-ปจฺจเยน ปจฺจโย. (3) (สํขิตฺตํ.)}

\proseitem{4}{เหตุยา ตีณิ, อารมฺมเณ นว ฯเปฯ อวิคเต นว.}

\vspace{\tipitakaparskip}
\csromancenter{(ยถา กุสลตฺติเก ปญฺหาวารํ, เอวํ วิตฺถาเรตพฺพํ.)}
\csromanheader{2}{2. \textbf{สเหตุกทุก 1-7. ปฏิจฺจาทิวาร}}
\prose{ปจฺจยจตุกฺกเหตุ}

\proseitem{5}{นสเหตุกํ ธมฺมํ ปฏิจฺจ น สเหตุโก ธมฺโม อุปฺปชฺชติ เหตุปจฺจยา. วิจิกิจฺฉา\-สหคตํ อุทฺธจฺจ\-สหคตํ โมหํ ปฏิจฺจ จิตฺต\-สมุฏฺฐานํ รูปํ. เอกํ มหาภูตํ ปฏิจฺจ ตโย มหาภูตา ฯเปฯ เทฺว มหาภูเต ปฏิจฺจ เทฺว มหาภูตา. มหาภูเต ปฏิจฺจ จิตฺต\-สมุฏฺฐานํ รูปํ, กฏตฺตารูปํ, อุปาทารูปํ.  นสเหตุกํ ธมฺมํ ปฏิจฺจ นอเหตุโก ธมฺโม อุปฺปชฺชติ เหตุปจฺจยา. วิจิกิจฺฉา\-สหคตํ อุทฺธจฺจ\-สหคตํ โมหํ ปฏิจฺจ สมฺปยุตฺตกา ขนฺธา. ปฏิสนฺธิ\-กฺขเณ วตฺถุํ ปฏิจฺจ สเหตุกา ขนฺธา.  นสเหตุกํ ธมฺมํ ปฏิจฺจ นสเหตุโก จ นอเหตุโก จ ธมฺมา อุปฺปชฺชนฺติ เหตุปจฺจยา. (3)}

\csromanheader{2}{3. เหตุ\-สมฺปยุตฺตทุก}
\prose{\csromanfolio{75}นอเหตุกํ ธมฺมํ ปฏิจฺจ นอเหตุโก ธมฺโม อุปฺปชฺชติ เหตุปจฺจยา. ตีณิ.}

\prose{นสเหตุกญฺจ นอเหตุกญฺจ ธมฺมํ ปฏิจฺจ นสเหตุโก ธมฺโม อุปฺปชฺชติ เหตุปจฺจยา.  นสเหตุกญฺจ นอเหตุกญฺจ ธมฺมํ ปฏิจฺจ นอเหตุโก ธมฺโม อุปฺปชฺชติ เหตุปจฺจยา.  นสเหตุกญฺจ นอเหตุกญฺจ ธมฺมํ ปฏิจฺจ นสเหตุโก จ นอเหตุโก จ ธมฺมา อุปฺปชฺชนฺติ เหตุปจฺจยา. (3) (สํขิตฺตํ.)}

\proseitem{6}{เหตุยา นว, อารมฺมเณ ฉ ฯเปฯ อวิคเต นว.}

\vspace{\tipitakaparskip}
\csromancenter{(สหชาตวารมฺปิ ฯเปฯ สมฺปยุตฺตวารมฺปิ ปฏิจฺจ\-วาร\-สทิสํ, เอวํ วิตฺถาเรตพฺพํ.)}
\prose{เหตุอารมฺมณ}

\proseitem{7}{นสเหตุโก ธมฺโม นสเหตุกสฺส ธมฺมสฺส เหตุปจฺจเยน ปจฺจโย.  นสเหตุโก ธมฺโม นอเหตุกสฺส ธมฺมสฺส เหตุปจฺจเยน ปจฺจโย.  นสเหตุโก ธมฺโม นสเหตุกสฺส จ นอเหตุกสฺส จ ธมฺมสฺส เหตุปจฺจเยน ปจฺจโย. (3)}

\prose{นอเหตุโก ธมฺโม นอเหตุกสฺส ธมฺมสฺส เหตุปจฺจเยน ปจฺจโย. ตีณิ.}

\prose{นสเหตุโก ธมฺโม นสเหตุกสฺส ธมฺมสฺส อารมฺมณ\-ปจฺจเยน ปจฺจโย. (สํขิตฺตํ.)}

\proseitem{8}{เหตุยา ฉ, อารมฺมเณ นว ฯเปฯ อวิคเต นว.}

\prose{(ยถา กุสลตฺติเก ปญฺหาวารํ, เอวํ วิตฺถาเรตพฺพํ.)}

\csromanmarkh{3. เหตุ\-สมฺปยุตฺตทุก}\csromanmarkhii{1. ปฏิจฺจวาร}\csromanheadernomark{2}{3. เหตุ\-สมฺปยุตฺตทุก 1. ปฏิจฺจวาร}
\proseitem{9}{นเหตุ\-สมฺปยุตฺตํ ธมฺมํ ปฏิจฺจ นเหตุ\-สมฺปยุตฺโต ธมฺโม อุปฺปชฺชติ เหตุปจฺจยา.  นเหตุ\-สมฺปยุตฺตํ ธมฺมํ ปฏิจฺจ นเหตุ\-วิปฺปยุตฺโต ธมฺโม อุปฺปชฺชติ เหตุปจฺจยา.  นเหตุ\-สมฺปยุตฺตํ ธมฺมํ ปฏิจฺจ นเหตุ\-สมฺปยุตฺโต จ นเหตุ\-วิปฺปยุตฺโต จ ธมฺมา อุปฺปชฺชนฺติ เหตุปจฺจยา. (3)}

\prose{นเหตุ\-วิปฺปยุตฺตํ ธมฺมํ ปฏิจฺจ นเหตุ\-วิปฺปยุตฺโต ธมฺโม อุปฺปชฺชติ เหตุปจฺจยา.  นเหตุ\-วิปฺปยุตฺตํ ธมฺมํ ปฏิจฺจ นเหตุ\-สมฺปยุตฺโต ธมฺโม อุปฺปชฺชติ}

\prose{\csromanfolio{76}เหตุปจฺจยา.  นเหตุ\-วิปฺปยุตฺตํ ธมฺมํ ปฏิจฺจ นเหตุ\-สมฺปยุตฺโต จ นเหตุ\-วิปฺปยุตฺโต จ ธมฺมา อุปฺปชฺชนฺติ เหตุปจฺจยา. (3)}

\prose{นเหตุ\-สมฺปยุตฺตญฺจ นเหตุ\-วิปฺปยุตฺตญฺจ ธมฺมํ ปฏิจฺจ นเหตุ\-สมฺปยุตฺโต ธมฺโม อุปฺปชฺชติ เหตุปจฺจยา. ตีณิ. (สํขิตฺตํ.)}

\proseitem{10}{เหตุยา นว, อารมฺมเณ ฉ ฯเปฯ อวิคเต นว.}

\vspace{\tipitakaparskip}
\csromancenter{(สหชาตวาเรปิ ฯเปฯ ปญฺหาวาเรปิ สพฺพตฺถ วิตฺถาเรตพฺพํ.)}
\csromanmarkh{4. เหตุ\-สเหตุกทุก}\csromanmarkhii{1. ปฏิจฺจวาร}\csromanheadernomark{2}{4. เหตุ\-สเหตุกทุก 1. ปฏิจฺจวาร}
\proseitem{11}{นเหตุญฺเจว นอเหตุกญฺจ ธมฺมํ ปฏิจฺจ นเหตุ เจว นอเหตุโก จ ธมฺโม อุปฺปชฺชติ เหตุปจฺจยา.  นเหตุญฺเจว นอเหตุกญฺจ ธมฺมํ ปฏิจฺจ นอเหตุโก เจว น นเหตุ จ ธมฺโม อุปฺปชฺชติ เหตุปจฺจยา.  นเหตุญฺเจว นอเหตุกญฺจ ธมฺมํ ปฏิจฺจ นเหตุ เจว นอเหตุโก จ นอเหตุโก เจว น นเหตุ จ ธมฺมา อุปฺปชฺชนฺติ เหตุปจฺจยา. (3)}

\prose{นอเหตุกญฺเจว น นเหตุญฺจ ธมฺมํ ปฏิจฺจ นอเหตุโก เจว น นเหตุ จ ธมฺโม อุปฺปชฺชติ เหตุปจฺจยา.  นอเหตุกญฺเจว น นเหตุญฺจ ธมฺมํ ปฏิจฺจ นเหตุ เจว นอเหตุโก จ ธมฺโม อุปฺปชฺชติ เหตุปจฺจยา.  นอเหตุกญฺเจว น นเหตุญฺจ ธมฺมํ ปฏิจฺจ นเหตุ เจว นอเหตุโก จ นอเหตุโก เจว น นเหตุ จ ธมฺมา อุปฺปชฺชนฺติ เหตุปจฺจยา. (3)}

\prose{นเหตุญฺเจว นอเหตุกญฺจ นอเหตุกญฺเจว น นเหตุญฺจ ธมฺมํ ปฏิจฺจ นเหตุ เจว นอเหตุโก จ ธมฺโม อุปฺปชฺชติ เหตุปจฺจยา. ตีณิ.}

\prose{เหตุยา นว, อารมฺมเณ นว ฯเปฯ อวิคเต นว.}

\vspace{\tipitakaparskip}
\csromancenter{(สหชาตวาเรปิ ฯเปฯ ปญฺหาวาเรปิ สพฺพตฺถ วิตฺถาโร.)}
\csromanmarkh{5. เหตุ\-เหตุ\-สมฺปยุตฺตทุก}\csromanmarkhii{1. ปฏิจฺจวาร}\csromanheadernomark{2}{5. เหตุ\-เหตุ\-สมฺปยุตฺตทุก 1. ปฏิจฺจวาร}
\proseitem{12}{นเหตุญฺเจว นเหตุ\-วิปฺปยุตฺตญฺจ ธมฺมํ ปฏิจฺจ นเหตุ เจว นเหตุ\-วิปฺปยุตฺโต จ ธมฺโม อุปฺปชฺชติ เหตุปจฺจยา.  นเหตุญฺเจว นเหตุ\-วิปฺปยุตฺตญฺจ ธมฺมํ ปฏิจฺจ นเหตุ\-วิปฺปยุตฺโต เจว น นเหตุ จ ธมฺโม อุปฺปชฺชติ}

\vspace{\tipitakaparskip}
\csromancenter{นเหตุ\-สเหตุกทุก เหตุปจฺจยา.  นเหตุญฺเจว นเหตุ\-วิปฺปยุตฺตญฺจ ธมฺมํ ปฏิจฺจ นเหตุ เจว นเหตุ\-วิปฺปยุตฺโต จ นเหตุ\-วิปฺปยุตฺโต เจว น นเหตุ จ ธมฺมา อุปฺปชฺชนฺติ เหตุปจฺจยา. (3)}
\prosecont{\csromanfolio{77}นเหตุ\-วิปฺปยุตฺต\-ญฺเจว น นเหตุญฺจ ธมฺมํ ปฏิจฺจ นเหตุ\-วิปฺปยุตฺโต เจว น นเหตุ จ ธมฺโม อุปฺปชฺชติ เหตุปจฺจยา.  นเหตุ\-วิปฺปยุตฺต\-ญฺเจว น นเหตุญฺจ ธมฺมํ ปฏิจฺจ นเหตุ เจว นเหตุ\-วิปฺปยุตฺโต จ ธมฺโม อุปฺปชฺชติ เหตุปจฺจยา.  นเหตุ\-วิปฺปยุตฺต\-ญฺเจว น นเหตุญฺจ ธมฺมํ ปฏิจฺจ นเหตุ เจว นเหตุ\-วิปฺปยุตฺโต จ นเหตุ\-วิปฺปยุตฺโต เจว น นเหตุ จ ธมฺมา อุปฺปชฺชนฺติ เหตุปจฺจยา. (3)}

\prose{นเหตุญฺเจว นเหตุ\-วิปฺปยุตฺตญฺจ นเหตุ\-วิปฺปยุตฺต\-ญฺเจว น นเหตุญฺจ ธมฺมํ ปฏิจฺจ นเหตุ เจว นเหตุ\-วิปฺปยุตฺโต จ ธมฺโม อุปฺปชฺชติ เหตุปจฺจยา. ตีณิ. (สํขิตฺตํ.)}

\prose{เหตุยา นว ฯเปฯ อวิคเต นว. (สพฺพตฺถ วิตฺถาโร.)}

\csromanheader{2}{6. นเหตุ\-สเหตุกทุก 1-7. ปฏิจฺจาทิวาร}
\proseitem{13}{นเหตุํ นสเหตุกํ ธมฺมํ ปฏิจฺจ นเหตุ นสเหตุโก ธมฺโม อุปฺปชฺชติ เหตุปจฺจยา. เอกํ มหาภูตํ ปฏิจฺจ ตโย มหาภูตา ฯเปฯ. นเหตุํ นสเหตุกํ ธมฺมํ ปฏิจฺจ นเหตุ นอเหตุโก ธมฺโม อุปฺปชฺชติ เหตุปจฺจยา. ปฏิสนฺธิ\-กฺขเณ วตฺถุํ ปฏิจฺจ นเหตุสเหตุกา ขนฺธา.  นเหตุํ นสเหตุกํ ธมฺมํ ปฏิจฺจ นเหตุ นสเหตุโก จ นเหตุ นอเหตุโก จ ธมฺมา อุปฺปชฺชนฺติ เหตุปจฺจยา. (3)}

\prose{นเหตุํ นอเหตุกํ ธมฺมํ ปฏิจฺจ นเหตุ นอเหตุโก ธมฺโม อุปฺปชฺชติ เหตุปจฺจยา.  นเหตุํ นอเหตุกํ ธมฺมํ ปฏิจฺจ นเหตุ นสเหตุโก ธมฺโม อุปฺปชฺชติ เหตุปจฺจยา.  นเหตุํ นอเหตุกํ ธมฺมํ ปฏิจฺจ นเหตุ นสเหตุโก จ นเหตุ นอเหตุโก จ ธมฺมา อุปฺปชฺชนฺติ เหตุปจฺจยา. (3)}

\prose{นเหตุํ นสเหตุกญฺจ นเหตุํ นอเหตุกญฺจ ธมฺมํ ปฏิจฺจ นเหตุ นสเหตุโก ธมฺโม อุปฺปชฺชติ เหตุปจฺจยา.  นเหตุํ นสเหตุกญฺจ นเหตุํ นอเหตุกญฺจ ธมฺมํ ปฏิจฺจ นเหตุ นอเหตุโก ธมฺโม อุปฺปชฺชติ เหตุปจฺจยา.  นเหตุํ นสเหตุกญฺจ นเหตุํ นอเหตุกญฺจ ธมฺมํ \csromanfolio{78}ปฏิจฺจ นเหตุ นสเหตุโก จ นเหตุ นอเหตุโก จ ธมฺมา อุปฺปชฺชนฺติ เหตุปจฺจยา. (3) (สํขิตฺตํ.)}

\proseitem{14}{เหตุยา นว, อารมฺมเณ จตฺตาริ ฯเปฯ อวิคเต นว.}

\prose{(สหชาตวารมฺปิ ฯเปฯ สมฺปยุตฺตวารมฺปิ ปฏิจฺจ\-วาร\-สทิสํ.)}

\csromanheader{2}{\textbf{อารมฺมณ}}
\proseitem{15}{นเหตุ นสเหตุโก ธมฺโม นเหตุสฺส นสเหตุกสฺส ธมฺมสฺส อารมฺมณ\-ปจฺจเยน ปจฺจโย.  นเหตุ นสเหตุโก ธมฺโม นเหตุสฺส นอเหตุกสฺส ธมฺมสฺส อารมฺมณ\-ปจฺจเยน ปจฺจโย. (2)}

\prosewithcloser{นเหตุ นอเหตุโก ธมฺโม นเหตุสฺส นอเหตุกสฺส ธมฺมสฺส อารมฺมณ\-ปจฺจเยน ปจฺจโย.  นเหตุ นอเหตุโก ธมฺโม นเหตุสฺส นสเหตุกสฺส ธมฺมสฺส อารมฺมณ\-ปจฺจเยน ปจฺจโย. (2) (สํขิตฺตํ.)}{16. อารมฺมเณ จตฺตาริ ฯเปฯ อวิคเต สตฺต. เหตุโคจฺฉกํ นิฏฺฐิตํ.}
\csromanheader{2}{7. \textbf{สปฺปจฺจยทุก 1-7. ปฏิจฺจาทิวาร}}
\proseitem{17}{นอปฺปจฺจยํ ธมฺมํ ปฏิจฺจ นอปฺปจฺจโย ธมฺโม อุปฺปชฺชติ เหตุปจฺจยา.}

\prose{เหตุยา เอกํ. (สพฺพตฺถ วิตฺถาโร.)}

\proseitem{18}{นอปฺปจฺจโย ธมฺโม นอปฺปจฺจยสฺส ธมฺมสฺส เหตุปจฺจเยน ปจฺจโย.}

\prose{นสปฺปจฺจโย ธมฺโม นอปฺปจฺจยสฺส ธมฺมสฺส อารมฺมณ\-ปจฺจเยน ปจฺจโย. (สํขิตฺตํ.)}

\prose{เหตุยา เอกํ, อารมฺมเณ เทฺว. (สพฺพตฺถ วิตฺถาโร.)}

\csromanmarkh{8. สงฺขตทุก}\csromanmarkhii{1. ปฏิจฺจวาร}\csromanheadernomark{2}{8. \textbf{สงฺขตทุก 1. ปฏิจฺจวาร}}
\proseitem{19}{นอสงฺขตํ ธมฺมํ ปฏิจฺจ นอสงฺขโต ธมฺโม อุปฺปชฺชติ เหตุปจฺจยา. เหตุยา เอกํ. (สปฺปจฺจย\-ทุก\-สทิสํ.)}

\csromanmarkh{10. สปฺปฏิฆทุก}\csromanmarkhii{9. สนิทสฺสนทุก 1. ปฏิจฺจวาร}\csromanheadernomark{2}{10. สปฺปฏิฆทุก \textbf{9. สนิทสฺสนทุก 1. ปฏิจฺจวาร}}
\proseitem{20}{\csromanfolio{79}นสนิทสฺสนํ ธมฺมํ ปฏิจฺจ นสนิทสฺสโน ธมฺโม อุปฺปชฺชติ เหตุปจฺจยา.  นสนิทสฺสนํ ธมฺมํ ปฏิจฺจ นอนิทสฺสโน ธมฺโม อุปฺปชฺชติ เหตุปจฺจยา.  นสนิทสฺสนํ ธมฺมํ ปฏิจฺจ นสนิทสฺสโน จ นอนิทสฺสโน จ ธมฺมา อุปฺปชฺชนฺติ เหตุปจฺจยา. (3)}

\prose{เหตุยา ตีณิ.}

\csromanheader{2}{10. สปฺปฏิฆทุก 1-7. ปฏิจฺจาทิวาร}
\proseitem{21}{นสปฺปฏิฆํ ธมฺมํ ปฏิจฺจ นสปฺปฏิโฆ ธมฺโม อุปฺปชฺชติ เหตุปจฺจยา.  นสปฺปฏิฆํ ธมฺมํ ปฏิจฺจ นอปฺปฏิโฆ ธมฺโม อุปฺปชฺชติ เหตุปจฺจยา.  นสปฺปฏิฆํ ธมฺมํ ปฏิจฺจ นสปฺปฏิโฆ จ นอปฺปฏิโฆ จ ธมฺมา อุปฺปชฺชนฺติ เหตุปจฺจยา. (3)}

\prose{นอปฺปฏิฆํ ธมฺมํ ปฏิจฺจ นอปฺปฏิโฆ ธมฺโม อุปฺปชฺชติ เหตุปจฺจยา. ตีณิ.}

\prose{นสปฺปฏิฆญฺจ นอปฺปฏิฆญฺจ ธมฺมํ ปฏิจฺจ นสปฺปฏิโฆ ธมฺโม อุปฺปชฺชติ เหตุปจฺจยา. ตีณิ. (สํขิตฺตํ.)}

\prose{เหตุยา นว, อารมฺมเณ เอกํ.}

\vspace{\tipitakaparskip}
\csromancenter{(สหชาตวารมฺปิ ฯเปฯ สมฺปยุตฺตวารมฺปิ ปฏิจฺจ\-วาร\-สทิสํ.)}
\proseitem{22}{นสปฺปฏิโฆ ธมฺโม นสปฺปฏิฆสฺส ธมฺมสฺส เหตุปจฺจเยน ปจฺจโย.  นสปฺปฏิโฆ ธมฺโม นอปฺปฏิฆสฺส ธมฺมสฺส เหตุปจฺจเยน ปจฺจโย.  นสปฺปฏิโฆ ธมฺโม นสปฺปฏิฆสฺส จ นอปฺปฏิฆสฺส จ ธมฺมสฺส เหตุปจฺจเยน ปจฺจโย. (3)}

\prose{นสปฺปฏิโฆ ธมฺโม นสปฺปฏิฆสฺส ธมฺมสฺส อารมฺมณ\-ปจฺจเยน ปจฺจโย.  นอปฺปฏิโฆ ธมฺโม นสปฺปฏิฆสฺส ธมฺมสฺส อารมฺมณ\-ปจฺจเยน ปจฺจโย. (2) (สํขิตฺตํ.)}

\proseitem{23}{เหตุยา ตีณิ, อารมฺมเณ เทฺว, อธิปติยา จตฺตาริ ฯเปฯ อวิคเต นว.}

\csromanmarkh{ธมฺม\-ปจฺจนีย ทุกปฏฺฐาน\-ปาฬิ}\csromanmarkhii{11. รูปีทุก 1. ปฏิจฺจวาร}\csromanheadernomark{2}{ธมฺม\-ปจฺจนีย ทุกปฏฺฐาน\-ปาฬิ 11. รูปีทุก 1. ปฏิจฺจวาร}
\proseitem{24}{\csromanfolio{80}นรูปิํ ธมฺมํ ปฏิจฺจ นรูปี ธมฺโม อุปฺปชฺชติ เหตุปจฺจยา.  นรูปิํ ธมฺมํ ปฏิจฺจ นอรูปี ธมฺโม อุปฺปชฺชติ เหตุปจฺจยา.  นรูปิํ ธมฺมํ ปฏิจฺจ นรูปี จ นอรูปี จ ธมฺมา อุปฺปชฺชนฺติ เหตุปจฺจยา. (3)}

\prose{นอรูปิํ ธมฺมํ ปฏิจฺจ นอรูปี ธมฺโม อุปฺปชฺชติ เหตุปจฺจยา.  นอรูปิํ ธมฺมํ ปฏิจฺจ นรูปี ธมฺโม อุปฺปชฺชติ เหตุปจฺจยา.  นอรูปิํ ธมฺมํ ปฏิจฺจ นรูปี จ นอรูปี จ ธมฺมา อุปฺปชฺชนฺติ เหตุปจฺจยา. (3)}

\prose{นรูปิญฺจ นอรูปิญฺจ ธมฺมํ ปฏิจฺจ นรูปี ธมฺโม อุปฺปชฺชติ เหตุปจฺจยา.  นรูปิญฺจ นอรูปิญฺจ ธมฺมํ ปฏิจฺจ นอรูปี ธมฺโม อุปฺปชฺชติ เหตุปจฺจยา.  นรูปิญฺจ นอรูปิญฺจ ธมฺมํ ปฏิจฺจ นรูปี จ นอรูปี จ ธมฺมา อุปฺปชฺชนฺติ เหตุปจฺจยา. (3) (สํขิตฺตํ.)}

\proseitem{25}{เหตุยา นว, อารมฺมเณ ตีณิ ฯเปฯ อวิคเต นว.}

\vspace{\tipitakaparskip}
\csromancenter{(สหชาตวาเรปิ ฯเปฯ ปญฺหาวาเรปิ สพฺพตฺถ วิตฺถาโร.)}
\csromanmarkh{12. โลกิยทุก}\csromanmarkhii{1. ปฏิจฺจวาร}\csromanheadernomark{2}{12. \textbf{โลกิยทุก 1. ปฏิจฺจวาร}}
\proseitem{26}{นโลกิยํ ธมฺมํ ปฏิจฺจ นโลกิโย ธมฺโม อุปฺปชฺชติ เหตุปจฺจยา.  นโลกิยํ ธมฺมํ ปฏิจฺจ นโลกุตฺตโร ธมฺโม อุปฺปชฺชติ เหตุปจฺจยา.  นโลกิยํ ธมฺมํ ปฏิจฺจ นโลกิโย จ นโลกุตฺตโร จ ธมฺมา อุปฺปชฺชนฺติ เหตุปจฺจยา. (3)}

\prose{นโลกุตฺตรํ ธมฺมํ ปฏิจฺจ นโลกุตฺตโร ธมฺโม อุปฺปชฺชติ เหตุปจฺจยา.}

\prose{นโลกิยญฺจ นโลกุตฺตรญฺจ ธมฺมํ ปฏิจฺจ นโลกุตฺตโร ธมฺโม อุปฺปชฺชติ เหตุปจฺจยา. (2) (สํขิตฺตํ.)}

\prose{เหตุยา ปญฺจ, อารมฺมเณ เทฺว ฯเปฯ อวิคเต ปญฺจ. (สพฺพตฺถ วิตฺถาโร.)}

\csromanmarkh{13. เกนจิ\-วิญฺเญยฺยทุก}\csromanmarkhii{1. ปฏิจฺจวาร}\csromanheadernomark{2}{13. เกนจิ\-วิญฺเญยฺยทุก 1. ปฏิจฺจวาร}
\proseitem{27}{นเกนจิ วิญฺเญยฺยํ ธมฺมํ ปฏิจฺจ นเกนจิ วิญฺเญยฺโย ธมฺโม อุปฺปชฺชติ เหตุปจฺจยา.  นเกนจิ วิญฺเญยฺยํ ธมฺมํ ปฏิจฺจ นนเกนจิ วิญฺเญยฺโย ธมฺโม อุปฺปชฺชติ เหตุปจฺจยา.  นเกนจิ วิญฺเญยฺยํ ธมฺมํ ปฏิจฺจ นเกนจิ วิญฺเญยฺโย จ นนเกนจิ วิญฺเญยฺโย จ ธมฺมา อุปฺปชฺชนฺติ เหตุปจฺจยา. (3)}

\csromanheader{2}{15. สาสวทุก}
\prose{\csromanfolio{81}นนเกนจิ วิญฺเญยฺยํ ธมฺมํ ปฏิจฺจ นนเกนจิ วิญฺเญยฺโย ธมฺโม อุปฺปชฺชติ เหตุปจฺจยา. ตีณิ.}

\prose{นเกนจิ วิญฺเญยฺยญฺจ นนเกนจิ วิญฺเญยฺยญฺจ ธมฺมํ ปฏิจฺจ นเกนจิ วิญฺเญยฺโย ธมฺโม อุปฺปชฺชติ เหตุปจฺจยา. ตีณิ. (สํขิตฺตํ.)}

\prosewithcloser{เหตุยา นว ฯเปฯ อวิคเต นว. (สพฺพตฺถ วิตฺถาโร.)}{จูฬนฺตรทุกํ นิฏฺฐิตํ.}
\csromanmarkh{14. อาสวทุก}\csromanmarkhii{1. ปฏิจฺจวาร}\csromanheadernomark{2}{14. \textbf{อาสวทุก 1. ปฏิจฺจวาร}}
\proseitem{28}{โนอาสวํ ธมฺมํ ปฏิจฺจ โนอาสโว ธมฺโม อุปฺปชฺชติ เหตุปจฺจยา.  โนอาสวํ ธมฺมํ ปฏิจฺจ นโนอาสโว ธมฺโม อุปฺปชฺชติ เหตุปจฺจยา.  โนอาสวํ ธมฺมํ ปฏิจฺจ โนอาสโว จ นโนอาสโว จ ธมฺมา อุปฺปชฺชนฺติ เหตุปจฺจยา. (3)}

\prose{นโนอาสวํ ธมฺมํ ปฏิจฺจ นโนอาสโว ธมฺโม อุปฺปชฺชติ เหตุปจฺจยา. ตีณิ.}

\prose{โนอาสวญฺจ นโนอาสวญฺจ ธมฺมํ ปฏิจฺจ โนอาสโว ธมฺโม อุปฺปชฺชติ เหตุปจฺจยา.  โนอาสวญฺจ นโนอาสวญฺจ ธมฺมํ ปฏิจฺจ นโนอาสโว ธมฺโม อุปฺปชฺชติ เหตุปจฺจยา. โนอาสวญฺจ นโนอาสวญฺจ ธมฺมํ ปฏิจฺจ โนอาสโว จ นโนอาสโว จ ธมฺมา อุปฺปชฺชนฺติ เหตุปจฺจยา. (3) (สํขิตฺตํ.)}

\prose{เหตุยา นว, อารมฺมเณ นว ฯเปฯ อวิคเต นว. (สพฺพตฺถ วิตฺถาโร.)}

\csromanmarkh{15. สาสวทุก}\csromanmarkhii{1. ปฏิจฺจวาร}\csromanheadernomark{2}{15. สาสวทุก 1. ปฏิจฺจวาร}
\proseitem{29}{นสาสวํ ธมฺมํ ปฏิจฺจ นสาสโว ธมฺโม อุปฺปชฺชติ เหตุปจฺจยา.  นสาสวํ ธมฺมํ ปฏิจฺจ นอนาสโว ธมฺโม อุปฺปชฺชติ เหตุปจฺจยา.  นสาสวํ ธมฺมํ ปฏิจฺจ นสาสโว จ นอนาสโว จ ธมฺมา อุปฺปชฺชนฺติ เหตุปจฺจยา. (3)}

\prose{นอนาสวํ ธมฺมํ ปฏิจฺจ นอนาสโว ธมฺโม อุปฺปชฺชติ}

\csromanpalirecto
\gambhira{ธมฺม\-ปจฺจนีย ทุกปฏฺฐาน\-ปาฬิ}
\prose{\csromanfolio{82}นสาสวญฺจ นอนาสวญฺจ ธมฺมํ ปฏิจฺจ นอนาสโว ธมฺโม อุปฺปชฺชติ เหตุปจฺจยา. (1) (สํขิตฺตํ.)}

\prose{เหตุยา ปญฺจ, อารมฺมเณ เทฺว ฯเปฯ อวิคเต ปญฺจ. (สพฺพตฺถ วิตฺถาโร.)}

\csromanmarkh{16. อาสว\-สมฺปยุตฺตทุก}\csromanmarkhii{1. ปฏิจฺจวาร}\csromanheadernomark{2}{16. อาสว\-สมฺปยุตฺตทุก 1. ปฏิจฺจวาร}
\proseitem{30}{นอาสว\-สมฺปยุตฺตํ ธมฺมํ ปฏิจฺจ นอาสว\-สมฺปยุตฺโต ธมฺโม อุปฺปชฺชติ เหตุปจฺจยา.  นอาสว\-สมฺปยุตฺตํ ธมฺมํ ปฏิจฺจ นอาสว\-วิปฺปยุตฺโต ธมฺโม อุปฺปชฺชติ เหตุปจฺจยา.  นอาสว\-สมฺปยุตฺตํ ธมฺมํ ปฏิจฺจ นอาสว\-สมฺปยุตฺโต จ นอาสว\-วิปฺปยุตฺโต จ ธมฺมา อุปฺปชฺชนฺติ เหตุปจฺจยา. (3)}

\prose{นอาสว\-วิปฺปยุตฺตํ ธมฺมํ ปฏิจฺจ นอาสว\-วิปฺปยุตฺโต ธมฺโม อุปฺปชฺชติ เหตุปจฺจยา.  นอาสว\-วิปฺปยุตฺตํ ธมฺมํ ปฏิจฺจ นอาสว\-สมฺปยุตฺโต ธมฺโม อุปฺปชฺชติ เหตุปจฺจยา.  นอาสว\-วิปฺปยุตฺตํ ธมฺมํ ปฏิจฺจ นอาสว\-สมฺปยุตฺโต จ นอาสว\-วิปฺปยุตฺโต จ ธมฺมา อุปฺปชฺชนฺติ เหตุปจฺจยา. (3)}

\prose{นอาสว\-สมฺปยุตฺตญฺจ นอาสว\-วิปฺปยุตฺตญฺจ ธมฺมํ ปฏิจฺจ นอาสว\-สมฺปยุตฺโต ธมฺโม อุปฺปชฺชติ เหตุปจฺจยา.  นอาสว\-สมฺปยุตฺตญฺจ นอาสว\-วิปฺปยุตฺตญฺจ ธมฺมํ ปฏิจฺจ นอาสว\-วิปฺปยุตฺโต ธมฺโม อุปฺปชฺชติ เหตุปจฺจยา.  นอาสว\-สมฺปยุตฺตญฺจ นอาสว\-วิปฺปยุตฺตญฺจ ธมฺมํ ปฏิจฺจ นอาสว\-สมฺปยุตฺโต จ นอาสว\-วิปฺปยุตฺโต จ ธมฺมา อุปฺปชฺชนฺติ เหตุปจฺจยา. (3) (สํขิตฺตํ.)}

\prose{เหตุยา นว, อารมฺมเณ ฉ ฯเปฯ อวิคเต นว. (สพฺพตฺถ วิตฺถาโร.)}

\csromanmarkh{17. อาสว\-สาสวทุก}\csromanmarkhii{1. ปฏิจฺจวาร}\csromanheadernomark{2}{17. อาสว\-สาสวทุก 1. ปฏิจฺจวาร}
\proseitem{31}{นอาสวญฺเจว นอนาสวญฺจ ธมฺมํ ปฏิจฺจ นอาสโว เจว นอนาสโว จ ธมฺโม อุปฺปชฺชติ เหตุปจฺจยา.  นอาสวญฺเจว นอนาสวญฺจ ธมฺมํ ปฏิจฺจ นอนาสโว เจว นโน จ อาสโว ธมฺโม อุปฺปชฺชติ เหตุปจฺจยา.  นอาสวญฺเจว นอนาสวญฺจ ธมฺมํ ปฏิจฺจ โนอาสโว เจว นอนาสโว จ นอนาสโว เจว นโน จ อาสโว จ ธมฺมา อุปฺปชฺชนฺติ เหตุปจฺจยา. (3)}

\prose{นอนาสวญฺเจว นโน จ อาสวํ ธมฺมํ ปฏิจฺจ นอนาสโว เจว นโน จ อาสโว ธมฺโม อุปฺปชฺชติ เหตุปจฺจยา.  นอนาสวญฺเจว นโน จ อาสวํ ธมฺมํ ปฏิจฺจ นอาสโว เจว นอนาสโว จ ธมฺโม อุปฺปชฺชติ เหตุปจฺจยา.  นอนาสวญฺเจว นโน จ อาสวํ ธมฺมํ ปฏิจฺจ นอาสโว เจว}

\proseitem{18}{\csromanfolio{83}อาสว\-อาสว\-สมฺปยุตฺตทุก นอนาสโว จ นอนาสโว เจว นโน จ อาสโว จ ธมฺมา อุปฺปชฺชนฺติ เหตุปจฺจยา. (3)}

\prose{นอาสวญฺเจว น นอนาสวญฺจ นอนาสวญฺเจว นโนอาสวญฺจ ธมฺมํ ปฏิจฺจ นอาสโว เจว นอนาสโว จ ธมฺโม อุปฺปชฺชติ เหตุปจฺจยา.  นอาสวญฺเจว นอนาสวญฺจ นอนาสวญฺเจว นโน จ อาสวํ ธมฺมํ ปฏิจฺจ นอนาสโว เจว นโน จ อาสโว ธมฺโม อุปฺปชฺชติ เหตุปจฺจยา.  นอาสวญฺเจว นอนาสวญฺจ นอนาสวญฺเจว นโนอาสวญฺจ ธมฺมํ ปฏิจฺจ นอาสโว เจว นอนาสโว จ นอนาสโว เจว นโน จ อาสโว จ ธมฺมา อุปฺปชฺชนฺติ เหตุปจฺจยา. (3) (สํขิตฺตํ.)}

\prose{เหตุยา นว, อารมฺมเณ นว ฯเปฯ อวิคเต นว. (สพฺพตฺถ วิตฺถาโร.)}

\prose{18. อาสว\-อาสว\-สมฺปยุตฺตทุก 1. ปฏิจฺจวาร}

\proseitem{32}{นอาสวญฺเจว นอาสว\-วิปฺปยุตฺตญฺจ ธมฺมํ ปฏิจฺจ นอาสโว เจว นอาสว\-วิปฺปยุตฺโต จ ธมฺโม อุปฺปชฺชติ เหตุปจฺจยา.  นอาสวญฺเจว นอาสว\-วิปฺปยุตฺตญฺจ ธมฺมํ ปฏิจฺจ นอาสว\-วิปฺปยุตฺโต เจว นโน จ อาสโว ธมฺโม อุปฺปชฺชติ เหตุปจฺจยา.  นอาสวญฺเจว นอาสว\-วิปฺปยุตฺตญฺจ ธมฺมํ ปฏิจฺจ นอาสโว เจว นอาสว\-วิปฺปยุตฺโต จ นอาสว\-วิปฺปยุตฺโต เจว นโนอาสโว จ ธมฺมา อุปฺปชฺชนฺติ เหตุปจฺจยา. (3)}

\prose{นอาสว\-วิปฺปยุตฺต\-ญฺเจว นโนอาสวญฺจ ธมฺมํ ปฏิจฺจ นอาสว\-วิปฺปยุตฺโต เจว นโนอาสโว จ ธมฺโม อุปฺปชฺชติ เหตุปจฺจยา.  นอาสว\-วิปฺปยุตฺต\-ญฺเจว นโนอาสวญฺจ ธมฺมํ ปฏิจฺจ นอาสโว เจว นอาสว\-วิปฺปยุตฺโต จ ธมฺโม อุปฺปชฺชติ เหตุปจฺจยา.  นอาสว\-วิปฺปยุตฺต\-ญฺเจว นโน จ อาสวํ ธมฺมํ ปฏิจฺจ นอาสโว เจว นอาสว\-วิปฺปยุตฺโต จ นอาสว\-วิปฺปยุตฺโต เจว นโน จ อาสโว จ ธมฺมา อุปฺปชฺชนฺติ เหตุปจฺจยา. (3)}

\prose{นอาสวญฺเจว นอาสว\-วิปฺปยุตฺตญฺจ นอาสว\-วิปฺปยุตฺต\-ญฺเจว นโนอาสวญฺจ ธมฺมํ ปฏิจฺจ นอาสโว เจว นอาสว\-วิปฺปยุตฺโต จ ธมฺโม อุปฺปชฺชติ เหตุปจฺจยา.  นอาสวญฺเจว นอาสว\-วิปฺปยุตฺตญฺจ นอาสว\-วิปฺปยุตฺต\-ญฺเจว นโน จ อาสวญฺจ ธมฺมํ ปฏิจฺจ นอาสว\-วิปฺปยุตฺโต เจว นโน จ อาสโว ธมฺโม อุปฺปชฺชติ เหตุปจฺจยา.  นอาสวญฺเจว นอาสว\-วิปฺปยุตฺตญฺจ นอาสว\-วิปฺปยุตฺต\-ญฺเจว นโน จ อาสวญฺจ ธมฺมํ ปฏิจฺจ}

\prose{\csromanfolio{84}นอาสโว เจว นอาสว\-วิปฺปยุตฺโต จ นอาสว\-วิปฺปยุตฺโต เจว นโน จ อาสโว จ ธมฺมา อุปฺปชฺชนฺติ เหตุปจฺจยา. (3) (สํขิตฺตํ.)}

\prose{เหตุยา นว, อารมฺมเณ นว ฯเปฯ อวิคเต นว. (สพฺพตฺถ วิตฺถาโร.)}

\csromanmarkh{19. อาสว\-วิปฺปยุตฺต\-สาสวทุก}\csromanmarkhii{1. ปฏิจฺจวาร}\csromanheadernomark{2}{19. อาสว\-วิปฺปยุตฺต\-สาสวทุก 1. ปฏิจฺจวาร}
\proseitem{33}{อาสว\-วิปฺปยุตฺตํ นสาสวํ ธมฺมํ ปฏิจฺจ อาสว\-วิปฺปยุตฺโต นสาสโว ธมฺโม อุปฺปชฺชติ เหตุปจฺจยา.  อาสว\-วิปฺปยุตฺตํ นสาสวํ ธมฺมํ ปฏิจฺจ อาสว\-วิปฺปยุตฺโต นอนาสโว ธมฺโม อุปฺปชฺชติ เหตุปจฺจยา.  อาสว\-วิปฺปยุตฺตํ นสาสวํ ธมฺมํ ปฏิจฺจ อาสว\-วิปฺปยุตฺโต นสาสโว จ อาสว\-วิปฺปยุตฺโต นอนาสโว จ ธมฺมา อุปฺปชฺชนฺติ เหตุปจฺจยา. (3)}

\prose{อาสว\-วิปฺปยุตฺตํ นอนาสวํ ธมฺมํ ปฏิจฺจ อาสว\-วิปฺปยุตฺโต นอนาสโว ธมฺโม อุปฺปชฺชติ เหตุปจฺจยา. (1)}

\prose{อาสว\-วิปฺปยุตฺตํ นสาสวญฺจ อาสว\-วิปฺปยุตฺตํ นอนาสวญฺจ ธมฺมํ ปฏิจฺจ อาสว\-วิปฺปยุตฺโต นอนาสโว ธมฺโม อุปฺปชฺชติ เหตุปจฺจยา. (สํขิตฺตํ.)}

\prosewithcloser{เหตุยา ปญฺจ, อารมฺมเณ เทฺว ฯเปฯ อวิคเต ปญฺจ. (สพฺพตฺถ วิตฺถาโร.)}{อาสวโคจฺฉกํ นิฏฺฐิตํ.}
\csromanheader{2}{20. \textbf{สญฺโญชนทุก}}
\proseitem{34}{โนสญฺโญชนํ ธมฺมํ ปฏิจฺจ โนสญฺโญชโน ธมฺโม อุปฺปชฺชติ เหตุปจฺจยา.  โนสญฺโญชนํ ธมฺมํ ปฏิจฺจ นโนสญฺโญชโน ธมฺโม อุปฺปชฺชติ เหตุปจฺจยา.  โนสญฺโญชนํ ธมฺมํ ปฏิจฺจ โนสญฺโญชโน จ นโนสญฺโญชโน จ ธมฺมา อุปฺปชฺชนฺติ เหตุปจฺจยา. (3)}

\prose{นโนสญฺโญชนํ ธมฺมํ ปฏิจฺจ นโนสญฺโญชโน ธมฺโม อุปฺปชฺชติ เหตุปจฺจยา.  นโนสญฺโญชนํ ธมฺมํ ปฏิจฺจ โนสญฺโญชโน ธมฺโม อุปฺปชฺชติ เหตุปจฺจยา.  นโนสญฺโญชนํ ธมฺมํ ปฏิจฺจ โนสญฺโญชโน จ นโนสญฺโญชโน จ ธมฺมา อุปฺปชฺชนฺติ เหตุปจฺจยา. (3)}

\prose{โนสญฺโญชนญฺจ นโน\-สญฺโญชนญฺจ ธมฺมํ ปฏิจฺจ โนสญฺโญชโน ธมฺโม อุปฺปชฺชติ เหตุปจฺจยา.  โนสญฺโญชนญฺจ นโน\-สญฺโญชนญฺจ}

\vspace{\tipitakaparskip}
\csromancenter{สญฺโญชน\-สมฺปยุตฺตทุก ธมฺมํ ปฏิจฺจ นโนสญฺโญชโน ธมฺโม อุปฺปชฺชติ เหตุปจฺจยา.  โนสญฺโญชนญฺจ นโน\-สญฺโญชนญฺจ ธมฺมํ ปฏิจฺจ โนสญฺโญชโน จ นโนสญฺโญชโน จ ธมฺมา อุปฺปชฺชนฺติ เหตุปจฺจยา. (3) (สํขิตฺตํ.)}
\prose{\csromanfolio{85}เหตุยา นว, อารมฺมเณ นว ฯเปฯ อวิคเต นว. (สพฺพตฺถ วิตฺถาโร.) 21. \textbf{ สญฺโญชนิยทุก}}

\proseitem{35}{นสญฺโญชนิยํ ธมฺมํ ปฏิจฺจ นสญฺโญชนิโย ธมฺโม อุปฺปชฺชติ เหตุปจฺจยา.  นสญฺโญชนิยํ ธมฺมํ ปฏิจฺจ นอสญฺโญชนิโย ธมฺโม อุปฺปชฺชติ เหตุปจฺจยา.  นสญฺโญชนิยํ ธมฺมํ ปฏิจฺจ นสญฺโญชนิโย จ นอสญฺโญชนิโย จ ธมฺมา อุปฺปชฺชนฺติ เหตุปจฺจยา. (3)}

\prose{นอสญฺโญ\-ชนิยํ ธมฺมํ ปฏิจฺจ นอสญฺโญชนิโย ธมฺโม อุปฺปชฺชติ เหตุปจฺจยา. (1)}

\prose{นสญฺโญชนิยญฺจ นอสญฺโญ\-ชนิยญฺจ ธมฺมํ ปฏิจฺจ นอสญฺโญชนิโย ธมฺโม อุปฺปชฺชติ เหตุปจฺจยา. (1) (สํขิตฺตํ.)}

\prose{เหตุยา ปญฺจ, อารมฺมเณ เทฺว ฯเปฯ อวิคเต ปญฺจ. (สพฺพตฺถ วิตฺถาโร.) 22. สญฺโญชน\-สมฺปยุตฺตทุก นสญฺโญชน\-สมฺปยุตฺตํ ธมฺมํ ปฏิจฺจ นสญฺโญชน\-สมฺปยุตฺโต ธมฺโม อุปฺปชฺชติ เหตุปจฺจยา.  นสญฺโญชน\-สมฺปยุตฺตํ ธมฺมํ ปฏิจฺจ นสญฺโญชน\-วิปฺปยุตฺโต ธมฺโม อุปฺปชฺชติ เหตุปจฺจยา.  นสญฺโญชน\-สมฺปยุตฺตํ ธมฺมํ ปฏิจฺจ นสญฺโญชน\-สมฺปยุตฺโต จ นสญฺโญชน\-วิปฺปยุตฺโต จ ธมฺมา อุปฺปชฺชนฺติ เหตุปจฺจยา. (3)}

\proseitem{36}{นสญฺโญชน\-วิปฺปยุตฺตํ ธมฺมํ ปฏิจฺจ นสญฺโญชน\-วิปฺปยุตฺโต ธมฺโม อุปฺปชฺชติ เหตุปจฺจยา.  นสญฺโญชน\-วิปฺปยุตฺตํ ธมฺมํ ปฏิจฺจ นสญฺโญชน\-สมฺปยุตฺโต ธมฺโม อุปฺปชฺชติ เหตุปจฺจยา.  นสญฺโญชน\-วิปฺปยุตฺตํ ธมฺมํ ปฏิจฺจ นสญฺโญชน\-สมฺปยุตฺโต จ นสญฺโญชน\-วิปฺปยุตฺโต จ ธมฺมา อุปฺปชฺชนฺติ เหตุปจฺจยา. (3) นสญฺโญชน\-สมฺปยุตฺตญฺจ นสญฺโญชน\-วิปฺปยุตฺตญฺจ ธมฺมํ ปฏิจฺจ นสญฺโญชน\-สมฺปยุตฺโต ธมฺโม อุปฺปชฺชติ เหตุปจฺจยา.  นสญฺโญชน\-สมฺปยุตฺตญฺจ นสญฺโญชน\-วิปฺปยุตฺตญฺจ ธมฺมํ ปฏิจฺจ นสญฺโญชน\-วิปฺปยุตฺโต ธมฺโม อุปฺปชฺชติ}

\prose{\csromanfolio{86}เหตุปจฺจยา.  นสญฺโญชน\-สมฺปยุตฺตญฺจ นสญฺโญชน\-วิปฺปยุตฺตญฺจ ธมฺมํ ปฏิจฺจ นสญฺโญชน\-สมฺปยุตฺโต จ นสญฺโญชน\-วิปฺปยุตฺโต จ ธมฺมา อุปฺปชฺชนฺติ เหตุปจฺจยา. (3) (สํขิตฺตํ.)}

\prose{เหตุยา นว, อารมฺมเณ ฉ ฯเปฯ อวิคเต นว. (สพฺพตฺถ วิตฺถาโร.)}

\csromanheader{2}{23. สญฺโญชน\-สญฺโญชนิยทุก}
\proseitem{37}{นสญฺโญชน\-ญฺเจว นอสญฺโญ\-ชนิยญฺจ ธมฺมํ ปฏิจฺจ นสญฺโญชโน เจว นอสญฺโญชนิโย จ ธมฺโม อุปฺปชฺชติ เหตุปจฺจยา.  นสญฺโญชน\-ญฺเจว นอสญฺโญ\-ชนิยญฺจ ธมฺมํ ปฏิจฺจ นอสญฺโญชนิโย เจว นโนอสญฺโญชโน จ ธมฺโม อุปฺปชฺชติ เหตุปจฺจยา.  นสญฺโญชน\-ญฺเจว นอสญฺโญ\-ชนิยญฺจ ธมฺมํ ปฏิจฺจ นสญฺโญชโน เจว นอสญฺโญชนิโย จ นอสญฺโญชนิโย เจว นโน จ สญฺโญชโน จ ธมฺมา อุปฺปชฺชนฺติ เหตุปจฺจยา. (3)}

\prose{นอสญฺโญ\-ชนิย\-ญฺเจว นโน จ สญฺโญชนํ ธมฺมํ ปฏิจฺจ นอสญฺโญชนิโย เจว นโน จ สญฺโญชโน ธมฺโม อุปฺปชฺชติ เหตุปจฺจยา. ตีณิ.}

\prose{นสญฺโญชน\-ญฺเจว นอสญฺโญ\-ชนิยญฺจ นอสญฺโญ\-ชนิย\-ญฺเจว นโน\-สญฺโญชนญฺจ ธมฺมํ ปฏิจฺจ นสญฺโญชโน เจว นอสญฺโญชนิโย จ ธมฺโม อุปฺปชฺชติ เหตุปจฺจยา. ตีณิ. (สํขิตฺตํ.)}

\prose{เหตุยา นว, อารมฺมเณ นว ฯเปฯ อวิคเต นว. (สพฺพตฺถ วิตฺถาโร.)}

\csromanheader{2}{24. สญฺโญชน\-สญฺโญชน\-สมฺปยุตฺตทุก}
\proseitem{38}{นสญฺโญชน\-ญฺเจว นสญฺโญชน\-วิปฺปยุตฺตญฺจ ธมฺมํ ปฏิจฺจ นสญฺโญชโน เจว นสญฺโญชน\-วิปฺปยุตฺโต จ ธมฺโม อุปฺปชฺชติ เหตุปจฺจยา.  นสญฺโญชน\-ญฺเจว นสญฺโญชน\-วิปฺปยุตฺตญฺจ ธมฺมํ ปฏิจฺจ นสญฺโญชน\-วิปฺปยุตฺโต เจว นโนสญฺโญชโน จ ธมฺโม อุปฺปชฺชติ เหตุปจฺจยา.  นสญฺโญชน\-ญฺเจว นสญฺโญชน\-วิปฺปยุตฺตญฺจ ธมฺมํ ปฏิจฺจ นสญฺโญชโน เจว นสญฺโญชน\-วิปฺปยุตฺโต จ นสญฺโญชน\-วิปฺปยุตฺโต เจว นโน จ สญฺโญชโน จ ธมฺมา อุปฺปชฺชนฺติ เหตุปจฺจยา. (3)}

\csromanheader{2}{25. สญฺโญชน\-วิปฺปยุตฺต\-สญฺโญชนิยทุก}
\prose{\csromanfolio{87}นสญฺโญชน\-วิปฺปยุตฺต\-ญฺเจว นโน จ สญฺโญชนํ ธมฺมํ ปฏิจฺจ นสญฺโญชน\-วิปฺปยุตฺโต เจว นโน จ สญฺโญชโน ธมฺโม อุปฺปชฺชติ เหตุปจฺจยา.  นสญฺโญชน\-วิปฺปยุตฺต\-ญฺเจว นโน จ สญฺโญชนํ ธมฺมํ ปฏิจฺจ นสญฺโญชโน เจว นสญฺโญชน\-วิปฺปยุตฺโต จ ธมฺโม อุปฺปชฺชติ เหตุปจฺจยา.  นสญฺโญชน\-วิปฺปยุตฺต\-ญฺเจว นโน จ สญฺโญชนํ ธมฺมํ ปฏิจฺจ นสญฺโญชโน เจว นสญฺโญชน\-วิปฺปยุตฺโต จ นสญฺโญชน\-วิปฺปยุตฺโต เจว นโน จ สญฺโญชโน จ ธมฺมา อุปฺปชฺชนฺติ เหตุปจฺจยา. (3)}

\prose{นสญฺโญชน\-ญฺเจว นสญฺโญชน\-วิปฺปยุตฺตญฺจ นสญฺโญชน\-วิปฺปยุตฺต\-ญฺเจว นโน จ สญฺโญชนญฺจ ธมฺมํ ปฏิจฺจ นสญฺโญชโน เจว นสญฺโญชน\-วิปฺปยุตฺโต จ ธมฺโม อุปฺปชฺชติ เหตุปจฺจยา. ตีณิ. (สํขิตฺตํ.)}

\prose{เหตุยา นว ฯเปฯ อวิคเต นว. (สพฺพตฺถ วิตฺถาโร.)}

\csromanheader{2}{25. สญฺโญชน\-วิปฺปยุตฺต\-สญฺโญชนิยทุก}
\proseitem{39}{สญฺโญชน\-วิปฺปยุตฺตํ นสญฺโญชนิยํ ธมฺมํ ปฏิจฺจ สญฺโญชน\-วิปฺปยุตฺโต นสญฺโญชนิโย ธมฺโม อุปฺปชฺชติ เหตุปจฺจยา.  สญฺโญชน\-วิปฺปยุตฺตํ นสญฺโญชนิยํ ธมฺมํ ปฏิจฺจ สญฺโญชน\-วิปฺปยุตฺโต นอสญฺโญชนิโย ธมฺโม อุปฺปชฺชติ เหตุปจฺจยา.  สญฺโญชน\-วิปฺปยุตฺตํ นสญฺโญชนิยํ ธมฺมํ ปฏิจฺจ สญฺโญชน\-วิปฺปยุตฺโต นสญฺโญชนิโย จ สญฺโญชน\-วิปฺปยุตฺโต นอสญฺโญชนิโย จ ธมฺมา อุปฺปชฺชนฺติ เหตุปจฺจยา. (3)}

\prose{สญฺโญชน\-วิปฺปยุตฺตํ นอสญฺโญ\-ชนิยํ ธมฺมํ ปฏิจฺจ สญฺโญชน\-วิปฺปยุตฺโต นอสญฺโญชนิโย ธมฺโม อุปฺปชฺชติ เหตุปจฺจยา. (1)}

\prose{สญฺโญชน\-วิปฺปยุตฺตํ นสญฺโญชนิยญฺจ สญฺโญชน\-วิปฺปยุตฺตํ นอสญฺโญ\-ชนิยญฺจ ธมฺมํ ปฏิจฺจ สญฺโญชน\-วิปฺปยุตฺโต นอสญฺโญชนิโย ธมฺโม อุปฺปชฺชติ เหตุปจฺจยา. (1) (สํขิตฺตํ.)}

\prosewithcloser{เหตุยา ปญฺจ ฯเปฯ อวิคเต ปญฺจ. (สพฺพตฺถ วิตฺถาโร.)}{สญฺโญชน\-โคจฺฉกํ นิฏฺฐิตํ.}
\csromanmarkh{ธมฺม\-ปจฺจนีย ทุกปฏฺฐาน\-ปาฬิ}\csromanmarkhii{26. คนฺถทุก}\csromanheadernomark{2}{ธมฺม\-ปจฺจนีย ทุกปฏฺฐาน\-ปาฬิ 26. คนฺถทุก}
\proseitem{40}{\csromanfolio{88}โนคนฺถํ ธมฺมํ ปฏิจฺจ โนคนฺโถ ธมฺโม อุปฺปชฺชติ เหตุปจฺจยา.  โนคนฺถํ ธมฺมํ ปฏิจฺจ นโนคนฺโถ ธมฺโม อุปฺปชฺชติ เหตุปจฺจยา.  โนคนฺถํ ธมฺมํ ปฏิจฺจ โนคนฺโถ จ นโนคนฺโถ จ ธมฺมา อุปฺปชฺชนฺติ เหตุปจฺจยา. (3)}

\prose{นโนคนฺถํ ธมฺมํ ปฏิจฺจ นโนคนฺโถ ธมฺโม อุปฺปชฺชติ เหตุปจฺจยา.  นโนคนฺถํ ธมฺมํ ปฏิจฺจ โนคนฺโถ ธมฺโม อุปฺปชฺชติ เหตุปจฺจยา.  นโนคนฺถํ ธมฺมํ ปฏิจฺจ โนคนฺโถ จ นโนคนฺโถ จ ธมฺมา อุปฺปชฺชนฺติ เหตุปจฺจยา. (3)}

\prose{โนคนฺถญฺจ นโนคนฺถญฺจ ธมฺมํ ปฏิจฺจ โนคนฺโถ ธมฺโม อุปฺปชฺชติ เหตุปจฺจยา.  โนคนฺถญฺจ นโนคนฺถญฺจ ธมฺมํ ปฏิจฺจ นโนคนฺโถ ธมฺโม อุปฺปชฺชติ เหตุปจฺจยา.  โนคนฺถญฺจ นโนคนฺถญฺจ ธมฺมํ ปฏิจฺจ โนคนฺโถ จ นโนคนฺโถ จ ธมฺมา อุปฺปชฺชนฺติ เหตุปจฺจยา. (3) (สํขิตฺตํ.)}

\prose{เหตุยา นว ฯเปฯ อวิคเต นว. (สพฺพตฺถ นว.)}

\csromanheader{2}{27. \textbf{คนฺถนิยทุก}}
\proseitem{41}{นคนฺถนิยํ ธมฺมํ ปฏิจฺจ นคนฺถนิโย ธมฺโม อุปฺปชฺชติ เหตุปจฺจยา.  นคนฺถนิยํ ธมฺมํ ปฏิจฺจ นอคนฺถนิโย ธมฺโม อุปฺปชฺชติ เหตุปจฺจยา.  นคนฺถนิยํ ธมฺมํ ปฏิจฺจ นคนฺถนิโย จ นอคนฺถนิโย จ ธมฺมา อุปฺปชฺชนฺติ เหตุปจฺจยา. (3)}

\prose{นอคนฺถนิยํ ธมฺมํ ปฏิจฺจ นอคนฺถนิโย ธมฺโม อุปฺปชฺชติ}

\prose{นคนฺถนิยญฺจ นอคนฺถนิยญฺจ ธมฺมํ ปฏิจฺจ นอคนฺถนิโย ธมฺโม อุปฺปชฺชติ เหตุปจฺจยา. (1) (สํขิตฺตํ.)}

\prose{เหตุยา ปญฺจ ฯเปฯ อวิคเต ปญฺจ. (สพฺพตฺถ ปญฺจ.)}

\csromanheader{2}{28. คนฺถ\-สมฺปยุตฺตทุก}
\prosecont{นคนฺถ\-สมฺปยุตฺตํ ธมฺมํ ปฏิจฺจ นคนฺถ\-สมฺปยุตฺโต ธมฺโม อุปฺปชฺชติ เหตุปจฺจยา.  นคนฺถ\-สมฺปยุตฺตํ ธมฺมํ ปฏิจฺจ นคนฺถ\-วิปฺปยุตฺโต ธมฺโม อุปฺปชฺชติ เหตุปจฺจยา.  นคนฺถ\-สมฺปยุตฺตํ ธมฺมํ ปฏิจฺจ นคนฺถ\-สมฺปยุตฺโต จ นคนฺถ\-วิปฺปยุตฺโต จ ธมฺมา อุปฺปชฺชนฺติ เหตุปจฺจยา. (3)}

\csromanheader{2}{29. คนฺถ\-คนฺถนิยทุก}
\prose{\csromanfolio{89}นคนฺถ\-วิปฺปยุตฺตํ ธมฺมํ ปฏิจฺจ นคนฺถ\-วิปฺปยุตฺโต ธมฺโม อุปฺปชฺชติ เหตุปจฺจยา.  นคนฺถ\-วิปฺปยุตฺตํ ธมฺมํ ปฏิจฺจ นคนฺถ\-สมฺปยุตฺโต ธมฺโม อุปฺปชฺชติ เหตุปจฺจยา.  นคนฺถ\-วิปฺปยุตฺตํ ธมฺมํ ปฏิจฺจ นคนฺถ\-สมฺปยุตฺโต จ นคนฺถ\-วิปฺปยุตฺโต จ ธมฺมา อุปฺปชฺชนฺติ เหตุปจฺจยา. (3)}

\prose{นคนฺถ\-สมฺปยุตฺตญฺจ นคนฺถ\-วิปฺปยุตฺตญฺจ ธมฺมํ ปฏิจฺจ นคนฺถ\-สมฺปยุตฺโต ธมฺโม อุปฺปชฺชติ เหตุปจฺจยา.  นคนฺถ\-สมฺปยุตฺตญฺจ นคนฺถ\-วิปฺปยุตฺตญฺจ ธมฺมํ ปฏิจฺจ นคนฺถ\-วิปฺปยุตฺโต ธมฺโม อุปฺปชฺชติ เหตุปจฺจยา.  นคนฺถ\-สมฺปยุตฺตญฺจ นคนฺถ\-วิปฺปยุตฺตญฺจ ธมฺมํ ปฏิจฺจ นคนฺถ\-สมฺปยุตฺโต จ นคนฺถ\-วิปฺปยุตฺโต จ ธมฺมา อุปฺปชฺชนฺติ เหตุปจฺจยา. (3) (สํขิตฺตํ.)}

\prose{เหตุยา นว, อารมฺมเณ ฉ, อธิปติยา นว ฯเปฯ อวิคเต นว. (สพฺพตฺถ วิตฺถาโร.)}

\csromanheader{2}{29. คนฺถ\-คนฺถนิยทุก}
\proseitem{43}{นคนฺถญฺเจว นอคนฺถนิยญฺจ ธมฺมํ ปฏิจฺจ นคนฺโถ เจว นอคนฺถนิโย จ ธมฺโม อุปฺปชฺชติ เหตุปจฺจยา.  นคนฺถญฺเจว นอคนฺถนิยญฺจ ธมฺมํ ปฏิจฺจ นอคนฺถนิโย เจว นโน จ คนฺโถ ธมฺโม อุปฺปชฺชติ เหตุปจฺจยา.  นคนฺถญฺเจว นอคนฺถนิยญฺจ ธมฺมํ ปฏิจฺจ นคนฺโถ เจว นอคนฺถนิโย จ นอคนฺถนิโย เจว นโน จ คนฺโถ จ ธมฺมา อุปฺปชฺชนฺติ เหตุปจฺจยา. (3)}

\prose{นอคนฺถนิย\-ญฺเจว นโน จ คนฺถํ ธมฺมํ ปฏิจฺจ นอคนฺถนิโย เจว นโนคนฺโถ จ ธมฺโม อุปฺปชฺชติ เหตุปจฺจยา.  นอคนฺถนิย\-ญฺเจว นโน จ คนฺถํ ธมฺมํ ปฏิจฺจ นคนฺโถ เจว นอคนฺถนิโย จ ธมฺโม อุปฺปชฺชติ เหตุปจฺจยา.  นอคนฺถนิย\-ญฺเจว นโน จ คนฺถํ ธมฺมํ ปฏิจฺจ นคนฺโถ เจว นอคนฺถนิโย จ นอคนฺถนิโย เจว นโน จ คนฺโถ จ ธมฺมา อุปฺปชฺชนฺติ เหตุปจฺจยา. (3)}

\prose{นคนฺถญฺเจว นอคนฺถนิยญฺจ นอคนฺถนิย\-ญฺเจว นโน จ คนฺถญฺจ ธมฺมํ ปฏิจฺจ นคนฺโถ เจว นอคนฺถนิโย จ ธมฺโม อุปฺปชฺชติ เหตุปจฺจยา. ตีณิ. (สํขิตฺตํ.)}

\prose{เหตุยา นว ฯเปฯ อวิคเต นว. (สพฺพตฺถ วิตฺถาโร.)}

\csromanmarkh{ธมฺม\-ปจฺจนีย ทุกปฏฺฐาน\-ปาฬิ}\csromanmarkhii{30. คนฺถ\-คนฺถ\-สมฺปยุตฺตทุก}\csromanheadernomark{2}{ธมฺม\-ปจฺจนีย ทุกปฏฺฐาน\-ปาฬิ 30. คนฺถ\-คนฺถ\-สมฺปยุตฺตทุก}
\prosecont{\csromanfolio{90}นคนฺถญฺเจว นคนฺถ\-วิปฺปยุตฺตญฺจ ธมฺมํ ปฏิจฺจ นคนฺโถ เจว นคนฺถ\-วิปฺปยุตฺโต จ ธมฺโม อุปฺปชฺชติ เหตุปจฺจยา.  นคนฺถญฺเจว นคนฺถ\-วิปฺปยุตฺตญฺจ ธมฺมํ ปฏิจฺจ นคนฺถ\-วิปฺปยุตฺโต เจว นโน จ คนฺโถ ธมฺโม อุปฺปชฺชติ เหตุปจฺจยา.  นคนฺถญฺเจว นคนฺถ\-วิปฺปยุตฺตญฺจ ธมฺมํ ปฏิจฺจ นคนฺโถ เจว นคนฺถ\-วิปฺปยุตฺโต จ นคนฺถ\-วิปฺปยุตฺโต เจว นโน จ คนฺโถ จ ธมฺมา อุปฺปชฺชนฺติ เหตุปจฺจยา. (3) นคนฺถ\-วิปฺปยุตฺต\-ญฺเจว นโน จ คนฺถํ ธมฺมํ ปฏิจฺจ นคนฺถ\-วิปฺปยุตฺโต เจว นโน จ คนฺโถ ธมฺโม อุปฺปชฺชติ เหตุปจฺจยา.  นคนฺถ\-วิปฺปยุตฺต\-ญฺเจว นโน จ คนฺถํ ธมฺมํ ปฏิจฺจ นคนฺโถ เจว นคนฺถ\-วิปฺปยุตฺโต จ ธมฺโม อุปฺปชฺชติ เหตุปจฺจยา.  นคนฺถ\-วิปฺปยุตฺต\-ญฺเจว นโน จ คนฺถํ ธมฺมํ ปฏิจฺจ นคนฺโถ เจว นคนฺถ\-วิปฺปยุตฺโต จ นคนฺถ\-วิปฺปยุตฺโต เจว นโน จ คนฺโถ จ ธมฺมา อุปฺปชฺชนฺติ เหตุปจฺจยา. (3) นคนฺถญฺเจว นคนฺถ\-วิปฺปยุตฺตญฺจ นคนฺถ\-วิปฺปยุตฺต\-ญฺเจว นโน จ คนฺถญฺจ ธมฺมํ ปฏิจฺจ นคนฺโถ เจว นคนฺถ\-วิปฺปยุตฺโต จ ธมฺโม อุปฺปชฺชติ เหตุปจฺจยา.  นคนฺถญฺเจว นคนฺถ\-วิปฺปยุตฺตญฺจ นคนฺถ\-วิปฺปยุตฺต\-ญฺเจว นโน จ คนฺถญฺจ ธมฺมํ ปฏิจฺจ นคนฺถ\-วิปฺปยุตฺโต เจว นโน จ คนฺโถ ธมฺโม อุปฺปชฺชติ เหตุปจฺจยา.  นคนฺถญฺเจว นคนฺถ\-วิปฺปยุตฺตญฺจ นคนฺถ\-วิปฺปยุตฺต\-ญฺเจว นโน จ คนฺถญฺจ ธมฺมํ ปฏิจฺจ นคนฺโถ เจว นคนฺถ\-วิปฺปยุตฺโต จ นคนฺถ\-วิปฺปยุตฺโต เจว นโน จ คนฺโถ จ ธมฺมา อุปฺปชฺชนฺติ เหตุปจฺจยา. (3) (สํขิตฺตํ.)}

\prose{เหตุยา นว ฯเปฯ อวิคเต นว. (สพฺพตฺถ วิตฺถาโร.)}

\csromanheader{2}{31. คนฺถ\-วิปฺปยุตฺต\-คนฺถนิยทุก}
\prosecont{คนฺถ\-วิปฺปยุตฺตํ นคนฺถนิยํ ธมฺมํ ปฏิจฺจ คนฺถ\-วิปฺปยุตฺโต นคนฺถนิโย ธมฺโม อุปฺปชฺชติ เหตุปจฺจยา.  คนฺถ\-วิปฺปยุตฺตํ นคนฺถนิยํ ธมฺมํ ปฏิจฺจ คนฺถ\-วิปฺปยุตฺโต นอคนฺถนิโย ธมฺโม อุปฺปชฺชติ เหตุปจฺจยา.  คนฺถ\-วิปฺปยุตฺตํ นคนฺถนิยํ ธมฺมํ ปฏิจฺจ คนฺถ\-วิปฺปยุตฺโต นคนฺถนิโย จ คนฺถ\-วิปฺปยุตฺโต นอคนฺถนิโย จ ธมฺมา อุปฺปชฺชนฺติ เหตุปจฺจยา. (3)}

\proseitem{45}{คนฺถ\-วิปฺปยุตฺตํ นอคนฺถนิยํ ธมฺมํ ปฏิจฺจ คนฺถ\-วิปฺปยุตฺโต นอคนฺถนิโย ธมฺโม อุปฺปชฺชติ เหตุปจฺจยา. (1)}

\csromanheader{2}{50-54. ปรามาสทุกาทิ}
\prose{\csromanfolio{91}คนฺถ\-วิปฺปยุตฺตํ นคนฺถนิยญฺจ คนฺถ\-วิปฺปยุตฺตํ นอคนฺถนิยญฺจ ธมฺมํ ปฏิจฺจ คนฺถ\-วิปฺปยุตฺโต นอคนฺถนิโย ธมฺโม อุปฺปชฺชติ เหตุปจฺจยา. (1) (สํขิตฺตํ.)}

\prosewithcloser{เหตุยา ปญฺจ, อารมฺมเณ เทฺว ฯเปฯ อวิคเต ปญฺจ. (สพฺพตฺถ วิตฺถาโร.)}{คนฺถ\-โคจฺฉกํ นิฏฺฐิตํ.}
\csromanheader{2}{32-49. โอฆาทิทุก}
\proseitem{46}{โนโอฆํ ธมฺมํ ปฏิจฺจ ฯเปฯ. โนโยคํ ธมฺมํ ปฏิจฺจ ฯเปฯ.}

\vspace{\tipitakaparskip}
\csromancenter{(โอฆ\-โคจฺฉก\-โยค\-โคจฺฉเกสุ อามสนํ อาสวโคจฺฉเก อามสนสทิสํ.)}
\proseitem{47}{โนนีวรณํ ธมฺมํ ปฏิจฺจ ฯเปฯ.}

\vspace{\tipitakaparskip}
\csromancenter{(นีวรณ\-โคจฺฉเก อามสนํ สญฺโญชน\-โคจฺฉเก อามสนสทิสํ.)}
\nitthitam{นีวรณ\-โคจฺฉกํ นิฏฺฐิตํ.}
\csromanheader{2}{50-54. ปรามาสทุกาทิ}
\proseitem{48}{โนปรามาสํ ธมฺมํ ปฏิจฺจ โนปรามาโส ธมฺโม อุปฺปชฺชติ เหตุปจฺจยา. (สํขิตฺตํ.)}

\proseitem{49}{นปรามฏฺฐํ ธมฺมํ ปฏิจฺจ นปรามฏฺโฐ ธมฺโม อุปฺปชฺชติ เหตุปจฺจยา. (สํขิตฺตํ.)}

\proseitem{50}{นปรามาส\-สมฺปยุตฺตํ ธมฺมํ ปฏิจฺจ นปรามาส\-สมฺปยุตฺโต ธมฺโม อุปฺปชฺชติ เหตุปจฺจยา.}

\proseitem{51}{นปรามาส\-ญฺเจว นอปรามฏฺฐญฺจ ธมฺมํ นปรามาโส เจว นอปรามฏฺโฐ จ ธมฺโม อุปฺปชฺชติ เหตุปจฺจยา.  นปรามาส\-ญฺเจว นอปรามฏฺฐญฺจ ธมฺมํ ปฏิจฺจ นอปรามฏฺโฐ เจว นโน จ ปรามาโส ธมฺโม อุปฺปชฺชติ เหตุปจฺจยา.  นปรามาส\-ญฺเจว นอปรามฏฺฐญฺจ ธมฺมํ ปฏิจฺจ}

\prose{\csromanfolio{92}นปรามาโส เจว นอปรามฏฺโฐ จ นอปรามฏฺโฐ เจว นโนปรามาโส จ ธมฺมา อุปฺปชฺชนฺติ เหตุปจฺจยา. (3)}

\prose{นอปรามฏฺฐ\-ญฺเจว นโน จ ปรามาสํ ธมฺมํ ปฏิจฺจ นปรามาโส เจว นอปรามฏฺโฐ จ ธมฺโม อุปฺปชฺชติ เหตุปจฺจยา. (1)}

\prose{นปรามาส\-ญฺเจว นอปรามฏฺฐญฺจ นอปรามฏฺฐ\-ญฺเจว นโน จ ปรามาสญฺจ ธมฺมํ ปฏิจฺจ นปรามาโส เจว นอปรามฏฺโฐ จ ธมฺโม อุปฺปชฺชติ เหตุปจฺจยา. (1) (สํขิตฺตํ.)}

\proseitem{52}{ปรามาส\-วิปฺปยุตฺตํ นปรามฏฺฐํ ธมฺมํ ปฏิจฺจ ปรามาส\-วิปฺปยุตฺโต นปรามฏฺโฐ ธมฺโม อุปฺปชฺชติ เหตุปจฺจยา.  ปรามาส\-วิปฺปยุตฺตํ นปรามฏฺฐํ ธมฺมํ ปฏิจฺจ ปรามาส\-วิปฺปยุตฺโต นอปรามฏฺโฐ ธมฺโม อุปฺปชฺชติ เหตุปจฺจยา.  ปรามาส\-วิปฺปยุตฺตํ นปรามฏฺฐํ ธมฺมํ ปฏิจฺจ ปรามาส\-วิปฺปยุตฺโต นปรามฏฺโฐ จ ปรามาส\-วิปฺปยุตฺโต นอปรามฏฺโฐ จ ธมฺมา อุปฺปชฺชนฺติ เหตุปจฺจยา. (3)}

\prose{ปรามาส\-วิปฺปยุตฺตํ นอปรามฏฺฐํ ธมฺมํ ปฏิจฺจ ปรามาส\-วิปฺปยุตฺโต นอปรามฏฺโฐ ธมฺโม อุปฺปชฺชติ เหตุปจฺจยา. (1)}

\prose{ปรามาส\-วิปฺปยุตฺตํ นปรามฏฺฐญฺจ ปรามาส\-วิปฺปยุตฺตํ นอปรามฏฺฐญฺจ ธมฺมํ ปฏิจฺจ ปรามาส\-วิปฺปยุตฺโต นอปรามฏฺโฐ}

\prosewithcloser{เหตุยา ปญฺจ, อารมฺมเณ เทฺว ฯเปฯ อวิคเต ปญฺจ.}{ปรามาส\-โคจฺฉกํ นิฏฺฐิตํ.}
\csromanheader{2}{55. \textbf{สารมฺมณทุก}}
\proseitem{53}{นสารมฺมณํ ธมฺมํ ปฏิจฺจ นสารมฺมโณ ธมฺโม อุปฺปชฺชติ เหตุปจฺจยา.  นสารมฺมณํ ธมฺมํ ปฏิจฺจ นอนารมฺมโณ ธมฺโม อุปฺปชฺชติ เหตุปจฺจยา.  นสารมฺมณํ ธมฺมํ ปฏิจฺจ นสารมฺมโณ จ นอนารมฺมโณ จ ธมฺมา อุปฺปชฺชนฺติ เหตุปจฺจยา. (3)}

\prose{นอนารมฺมณํ ธมฺมํ ปฏิจฺจ นอนารมฺมโณ ธมฺโม อุปฺปชฺชติ เหตุปจฺจยา.  นอนารมฺมณํ ธมฺมํ ปฏิจฺจ นสารมฺมโณ ธมฺโม อุปฺปชฺชติ}

\vspace{\tipitakaparskip}
\csromancenter{เจตสิกทุก เหตุปจฺจยา.  นอนารมฺมณํ ธมฺมํ ปฏิจฺจ นสารมฺมโณ จ นอนารมฺมโณ จ ธมฺมา อุปฺปชฺชนฺติ เหตุปจฺจยา. (3)}
\prose{\csromanfolio{93}นสารมฺมณญฺจ นอนารมฺมณญฺจ ธมฺมํ ปฏิจฺจ นสารมฺมโณ ธมฺโม อุปฺปชฺชติ เหตุปจฺจยา.  นสารมฺมณญฺจ นอนารมฺมณญฺจ ธมฺมํ ปฏิจฺจ นอนารมฺมโณ ธมฺโม อุปฺปชฺชติ เหตุปจฺจยา.  นสารมฺมณญฺจ นอนารมฺมณญฺจ ธมฺมํ ปฏิจฺจ นสารมฺมโณ จ นอนารมฺมโณ จ ธมฺมา อุปฺปชฺชนฺติ เหตุปจฺจยา. (3) (สํขิตฺตํ.)}

\prose{เหตุยา นว ฯเปฯ อวิคเต นว. (สพฺพตฺถ วิตฺถาโร.)}

\csromanheader{2}{56. \textbf{จิตฺตทุก}}
\proseitem{54}{นจิตฺตํ ธมฺมํ ปฏิจฺจ นจิตฺโต ธมฺโม อุปฺปชฺชติ เหตุปจฺจยา.  นจิตฺตํ ธมฺมํ ปฏิจฺจ นโนจิตฺโต ธมฺโม อุปฺปชฺชติ เหตุปจฺจยา.  นจิตฺตํ ธมฺมํ ปฏิจฺจ นจิตฺโต จ นโนจิตฺโต จ ธมฺมา อุปฺปชฺชนฺติ เหตุปจฺจยา. (3)}

\prose{นโนจิตฺตํ ธมฺมํ ปฏิจฺจ นจิตฺโต ธมฺโม อุปฺปชฺชติ เหตุปจฺจยา. (1)}

\prose{นจิตฺตญฺจ นโนจิตฺตญฺจ ธมฺมํ ปฏิจฺจ นจิตฺโต ธมฺโม อุปฺปชฺชติ เหตุปจฺจยา. (1) (สํขิตฺตํ.)}

\prose{เหตุยา ปญฺจ. (สพฺพตฺถ ปญฺจ.)}

\csromanheader{2}{57. เจตสิกทุก}
\proseitem{55}{นเจตสิกํ ธมฺมํ ปฏิจฺจ นเจตสิโก ธมฺโม อุปฺปชฺชติ เหตุปจฺจยา.  นเจตสิกํ ธมฺมํ ปฏิจฺจ นอเจตสิโก ธมฺโม อุปฺปชฺชติ เหตุปจฺจยา.  นเจตสิกํ ธมฺมํ ปฏิจฺจ นเจตสิโก จ นอเจตสิโก จ ธมฺมา อุปฺปชฺชนฺติ เหตุปจฺจยา. (3)}

\prose{นอเจตสิกํ ธมฺมํ ปฏิจฺจ นอเจตสิโก ธมฺโม อุปฺปชฺชติ เหตุปจฺจยา.  นอเจตสิกํ ธมฺมํ ปฏิจฺจ นเจตสิโก ธมฺโม อุปฺปชฺชติ เหตุปจฺจยา.  นอเจตสิกํ ธมฺมํ ปฏิจฺจ นเจตสิโก จ นอเจตสิโก จ ธมฺมา อุปฺปชฺชนฺติ เหตุปจฺจยา. (3)}

\prose{นเจตสิกญฺจ นอเจตสิกญฺจ ธมฺมํ ปฏิจฺจ นเจตสิโก ธมฺโม อุปฺปชฺชติ เหตุปจฺจยา.  นเจตสิกญฺจ นอเจตสิกญฺจ ธมฺมํ ปฏิจฺจ นอเจตสิโก \csromanfolio{94}ธมฺโม อุปฺปชฺชติ เหตุปจฺจยา.  นเจตสิกญฺจ นอเจตสิกญฺจ ธมฺมํ ปฏิจฺจ นเจตสิโก จ นอเจตสิโก จ ธมฺมา อุปฺปชฺชนฺติ เหตุปจฺจยา. (3)}

\prose{เหตุยา นว ฯเปฯ อวิคเต นว. (สพฺพตฺถ วิตฺถาโร.)}

\csromanheader{2}{58. จิตฺต\-สมฺปยุตฺตทุก}
\proseitem{56}{นจิตฺต\-สมฺปยุตฺตํ ธมฺมํ ปฏิจฺจ นจิตฺต\-สมฺปยุตฺโต ธมฺโม อุปฺปชฺชติ เหตุปจฺจยา. จิตฺตํ ปฏิจฺจ จิตฺต\-สมุฏฺฐานํ รูปํ.  นจิตฺต\-สมฺปยุตฺตํ ธมฺมํ ปฏิจฺจ นจิตฺต\-วิปฺปยุตฺโต ธมฺโม อุปฺปชฺชติ เหตุปจฺจยา. จิตฺตํ ปฏิจฺจ จิตฺต\-สมฺปยุตฺตกา ขนฺธา.}

\prose{(ปจฺจนีย\-ทุก\-มตฺตานิ น วตฺตพฺพํ ธมฺมํ ปูเรตุํ นว นว ปญฺหานิ กโรตุ.)}

\csromanheader{2}{59. จิตฺต\-สํสฏฺฐทุก}
\proseitem{55}{นจิตฺต\-สํสฏฺฐํ ธมฺมํ ปฏิจฺจ นจิตฺต\-สํสฏฺโฐ ธมฺโม อุปฺปชฺชติ เหตุปจฺจยา. เหตุยา นว.}

\csromanheader{2}{60. จิตฺต\-สมุฏฺฐานทุก}
\proseitem{58}{นจิตฺต\-สมุฏฺฐานํ ธมฺมํ ปฏิจฺจ นจิตฺต\-สมุฏฺฐาโน ธมฺโม อุปฺปชฺชติ เหตุปจฺจยา.}

\prose{เหตุยา นว.}

\csromanheader{2}{61. จิตฺต\-สหภูทุก}
\proseitem{59}{นจิตฺต\-สหภุํ ธมฺมํ ปฏิจฺจ นจิตฺตสหภู ธมฺโม อุปฺปชฺชติ เหตุปจฺจยา.}

\prose{เหตุยา นว.}

\csromanheader{2}{62. จิตฺตา\-นุปริวตฺติทุก}
\proseitem{60}{นจิตฺตา\-นุปริวตฺติํ ธมฺมํ ปฏิจฺจ นจิตฺตา\-นุปริวตฺตี ธมฺโม อุปฺปชฺชติ เหตุปจฺจยา. เหตุยา นว.}

\csromanheader{2}{63. จิตฺต\-สํสฏฺฐ\-สมุฏฺฐานทุก}
\proseitem{61}{นจิตฺต\-สํสฏฺฐ\-สมุฏฺฐานํ ธมฺมํ ปฏิจฺจ นจิตฺต\-สํสฏฺฐ\-สมุฏฺฐาโน ธมฺโม อุปฺปชฺชติ เหตุปจฺจยา.  เหตุยา นว.}

\csromanmarkh{68. อุปาทินฺนทุก}\csromanmarkhii{64. จิตฺต\-สํสฏฺฐ\-สมุฏฺฐาน\-สหภูทุก}\csromanheadernomark{2}{68. อุปาทินฺนทุก 64. จิตฺต\-สํสฏฺฐ\-สมุฏฺฐาน\-สหภูทุก}
\proseitem{62}{\csromanfolio{95}นจิตฺต\-สํสฏฺฐ\-สมุฏฺฐาน\-สหภุํ ธมฺมํ ปฏิจฺจ นจิตฺต\-สํสฏฺฐ\-สมุฏฺฐาน\-สหภู ธมฺโม อุปฺปชฺชติ เหตุปจฺจยา.  เหตุยา นว.}

\csromanheader{2}{65. จิตฺต\-สํสฏฺฐ\-สมุฏฺฐานา\-นุปริวตฺติทุก}
\proseitem{63}{นจิตฺต\-สํสฏฺฐ\-สมุฏฺฐานา\-นุ\-ปริวตฺติํ ธมฺมํ ปฏิจฺจ นจิตฺต\-สํสฏฺฐ\-สมุฏฺฐานา\-นุ\-ปริวตฺตี ธมฺโม อุปฺปชฺชติ เหตุปจฺจยา.}

\prose{เหตุยา นว ฯเปฯ อวิคเต นว. (สพฺพตฺถ วิตฺถาโร.)}

\csromanheader{2}{66. \textbf{อชฺฌตฺติกทุก}}
\proseitem{64}{นอชฺฌตฺติกํ ธมฺมํ ปฏิจฺจ นอชฺฌตฺติโก ธมฺโม อุปฺปชฺชติ เหตุปจฺจยา. ตีณิ.}

\prose{นพาหิรํ ธมฺมํ ปฏิจฺจ นพาหิโร ธมฺโม อุปฺปชฺชติ เหตุปจฺจยา. ปฏิสนฺธิ\-กฺขเณ จิตฺตํ ปฏิจฺจ อชฺฌตฺติกํ กฏตฺตารูปํ. (สํขิตฺตํ.)}

\prose{เหตุยา นว ฯเปฯ อวิคเต นว. (สพฺพตฺถ วิตฺถาโร.)}

\csromanheader{2}{67. \textbf{อุปาทาทุก}}
\proseitem{65}{นอุปาทา ธมฺมํ ปฏิจฺจ นอุปาทา ธมฺโม อุปฺปชฺชติ เหตุปจฺจยา.  นอุปาทา ธมฺมํ ปฏิจฺจ นโนอุปาทา ธมฺโม อุปฺปชฺชติ เหตุปจฺจยา.  นอุปาทา ธมฺมํ ปฏิจฺจ นอุปาทา จ นโนอุปาทา จ ธมฺมา อุปฺปชฺชนฺติ เหตุปจฺจยา. (3)}

\prose{นโนอุปาทา ธมฺมํ ปฏิจฺจ นอุปาทา ธมฺโม อุปฺปชฺชติ}

\prose{นอุปาทา จ นโนอุปาทา จ ธมฺมํ ปฏิจฺจ นอุปาทา ธมฺโม อุปฺปชฺชติ เหตุปจฺจยา. (1) (สํขิตฺตํ.)}

\prose{เหตุยา ปญฺจ.}

\csromanheader{2}{68. อุปาทินฺนทุก}
\prose{66. นอุปาทินฺนํ ธมฺมํ ปฏิจฺจ นอุปาทินฺโน ธมฺโม อุปฺปชฺชติ}

\prose{นอนุปาทินฺนํ ธมฺมํ ปฏิจฺจ นอนุปาทินฺโน ธมฺโม อุปฺปชฺชติ เหตุปจฺจยา.  นอนุปาทินฺนํ ธมฺมํ ปฏิจฺจ นอุปาทินฺโน ธมฺโม อุปฺปชฺชติ เหตุปจฺจยา.  }

\prose{\csromanfolio{96}นอนุปาทินฺนํ ธมฺมํ ปฏิจฺจ นฺเอาปาทินฺโน จ นอนุปาทินฺโน จ ธมฺมา อุปฺปชฺชนฺติ เหตุปจฺจยา. (3)}

\prosewithcloser{นอุปาทินฺนญฺจ นอนุปาทินฺนญฺจ ธมฺมํ ปฏิจฺจ นอุปาทินฺโน ธมฺโม อุปฺปชฺชติ เหตุปจฺจยา. (1) (สํขิตฺตํ.)}{เหตุยา ปญฺจ. (สพฺพตฺถ วิตฺถาโร.) มหนฺตรทุกํ นิฏฺฐิตํ.}
\csromanheader{2}{69. อุปาทาน\-โคจฺฉก}
\proseitem{67}{นอุปาทานํ ธมฺมํ ปฏิจฺจ นอุปาทาโน ธมฺโม อุปฺปชฺชติ เหตุปจฺจยา. (สํขิตฺตํ.) 75. กิเลสทุก}

\proseitem{68}{โนกิเลสํ ธมฺมํ ปฏิจฺจ โนกิเลโส ธมฺโม อุปฺปชฺชติ เหตุปจฺจยา.  โนกิเลสํ ธมฺมํ ปฏิจฺจ นโนกิเลโส ธมฺโม อุปฺปชฺชติ เหตุปจฺจยา.  โนกิเลสํ ธมฺมํ ปฏิจฺจ โนกิเลโส จ นโนกิเลโส จ ธมฺมา อุปฺปชฺชนฺติ เหตุปจฺจยา. (3)}

\prose{นโนกิเลสํ ธมฺมํ ปฏิจฺจ นโนกิเลโส ธมฺโม อุปฺปชฺชติ เหตุปจฺจยา. ตีณิ.}

\prose{โนกิเลสญฺจ นโนกิเลสญฺจ ธมฺมํ ปฏิจฺจ โนกิเลโส ธมฺโม อุปฺปชฺชติ เหตุปจฺจยา. ตีณิ. (สํขิตฺตํ.)}

\prose{เหตุยา นว. (สพฺพตฺถ นว.) \textbf{76. สํกิเลสิกทุก}}

\proseitem{69}{นสํกิเลสิกํ ธมฺมํ ปฏิจฺจ นสํกิเลสิโก ธมฺโม อุปฺปชฺชติ เหตุปจฺจยา.  นสํกิเลสิกํ ธมฺมํ ปฏิจฺจ นอสํกิเลสิโก ธมฺโม อุปฺปชฺชติ เหตุปจฺจยา.  นสํกิเลสิกํ ธมฺมํ ปฏิจฺจ นสํกิเลสิโก จ นอสํกิเลสิโก จ ธมฺมา อุปฺปชฺชนฺติ เหตุปจฺจยา. (3)}

\prose{นอสํกิเลสิกํ ธมฺมํ ปฏิจฺจ นอสํกิเลสิโก ธมฺโม อุปฺปชฺชติ เหตุปจฺจยา. (1)}

\csromanheader{2}{79. กิเลสสํกิเลสิกทุก}
\prose{\csromanfolio{97}นสํกิเลสิกญฺจ นอสํกิเลสิกญฺจ ธมฺมํ ปฏิจฺจ นอสํกิเลสิโก ธมฺโม อุปฺปชฺชติ เหตุปจฺจยา. (1) (สํขิตฺตํ.)}

\prose{เหตุยา ปญฺจ ฯเปฯ อวิคเต ปญฺจ. (สพฺพตฺถ วิตฺถาโร.)}

\csromanheader{2}{77. \textbf{สํกิลิฏฺฐทุก}}
\proseitem{70}{นสํกิลิฏฺฐํ ธมฺมํ ปฏิจฺจ นสํกิลิฏฺโฐ ธมฺโม อุปฺปชฺชติ}

\prose{นอสํกิลิฏฺฐํ ธมฺมํ ปฏิจฺจ นอสํกิลิฏฺโฐ ธมฺโม อุปฺปชฺชติ เหตุปจฺจยา.  นอสํกิลิฏฺฐํ ธมฺมํ ปฏิจฺจ นสํกิลิฏฺโฐ ธมฺโม อุปฺปชฺชติ เหตุปจฺจยา.  นอสํกิลิฏฺฐํ ธมฺมํ ปฏิจฺจ นสํกิลิฏฺโฐ จ นอสํกิลิฏฺโฐ จ ธมฺมา อุปฺปชฺชนฺติ เหตุปจฺจยา. (3)}

\prose{นสํกิลิฏฺฐญฺจ นอสํกิลิฏฺฐญฺจ ธมฺมํ ปฏิจฺจ นสํกิลิฏฺโฐ ธมฺโม อุปฺปชฺชติ เหตุปจฺจยา. (1) (สํขิตฺตํ.)}

\prose{เหตุยา ปญฺจ ฯเปฯ อวิคเต ปญฺจ. (สพฺพตฺถ วิตฺถาโร.)}

\csromanheader{2}{78. กิเลส\-สมฺปยุตฺตทุก}
\prose{71. นกิเลส\-สมฺปยุตฺตํ ธมฺมํ ปฏิจฺจ นกิเลส\-สมฺปยุตฺโต}

\prose{นกิเลส\-วิปฺปยุตฺตํ ธมฺมํ ปฏิจฺจ นกิเลส\-วิปฺปยุตฺโต ธมฺโม อุปฺปชฺชติ เหตุปจฺจยา. ตีณิ.}

\prose{นกิเลส\-สมฺปยุตฺตญฺจ นกิเลส\-วิปฺปยุตฺตญฺจ ธมฺมํ ปฏิจฺจ นกิเลส\-สมฺปยุตฺโต ธมฺโม อุปฺปชฺชติ เหตุปจฺจยา. (1) (สํขิตฺตํ.)}

\prose{เหตุยา ปญฺจ, อวิคเต ปญฺจ. (สพฺพตฺถ วิตฺถาโร.)}

\csromanheader{2}{79. กิเลสสํกิเลสิกทุก}
\proseitem{72}{นกิเลสญฺเจว นอสํกิเลสิกญฺจ ธมฺมํ ปฏิจฺจ นกิเลโส เจว นอสํกิเลสิโก จ ธมฺโม อุปฺปชฺชติ เหตุปจฺจยา. ตีณิ.}

\prose{นอสํกิเลสิก\-ญฺเจว นโน จ กิเลสํ ธมฺมํ ปฏิจฺจ นอสํกิเลสิโก เจว นโน กิเลโส ธมฺโม อุปฺปชฺชติ เหตุปจฺจยา. ตีณิ.}

\prose{\csromanfolio{98}นกิเลสญฺเจว นอสํกิเลสิกญฺจ นอสํกิเลสิก\-ญฺเจว นโน จ กิเลสญฺจ ธมฺมํ ปฏิจฺจ นกิเลโส เจว นอสํกิเลสิโก จ ธมฺโม อุปฺปชฺชติ เหตุปจฺจยา. ตีณิ. (สํขิตฺตํ.)}

\prose{เหตุยา นว.}

\csromanheader{2}{80. กิเลส\-สํกิลิฏฺฐทุก}
\proseitem{73}{นกิเลสญฺเจว นอสํกิลิฏฺฐญฺจ ธมฺมํ ปฏิจฺจ นกิเลโส เจว นอสํกิลิฏฺโฐ จ ธมฺโม อุปฺปชฺชติ เหตุปจฺจยา. ตีณิ.}

\prose{นอสํกิลิฏฺฐ\-ญฺเจว นโน จ กิเลสํ ธมฺมํ ปฏิจฺจ นอสํกิลิฏฺโฐ เจว นโน จ กิเลโส ธมฺโม อุปฺปชฺชติ เหตุปจฺจยา. ตีณิ.}

\prose{นกิเลสญฺเจว นอสํกิลิฏฺฐญฺจ นอสํกิลิฏฺฐ\-ญฺเจว นโน จ กิเลสญฺจ ธมฺมํ ปฏิจฺจ นกิเลโส เจว นอสํกิลิฏฺโฐ จ ธมฺโม อุปฺปชฺชติ เหตุปจฺจยา. ตีณิ. (สํขิตฺตํ.)}

\prose{เหตุยา นว.}

\csromanheader{2}{81. กิเลส\-กิเลส\-สมฺปยุตฺตทุก}
\proseitem{74}{นกิเลสญฺเจว นกิเลส\-วิปฺปยุตฺตญฺจ ธมฺมํ ปฏิจฺจ นกิเลโส เจว นกิเลส\-วิปฺปยุตฺโต จ ธมฺโม อุปฺปชฺชติ เหตุปจฺจยา. ตีณิ.}

\prose{นกิเลส\-วิปฺปยุตฺต\-ญฺเจว นโน จ กิเลสํ ธมฺมํ ปฏิจฺจ นกิเลส\-วิปฺปยุตฺโต เจว นโน จ กิเลโส ธมฺโม อุปฺปชฺชติ เหตุปจฺจยา. ตีณิ.}

\prose{นกิเลสญฺเจว นกิเลส\-วิปฺปยุตฺตญฺจ นกิเลส\-วิปฺปยุตฺต\-ญฺเจว นโน จ กิเลสญฺจ ธมฺมํ ปฏิจฺจ นกิเลโส เจว นกิเลส\-วิปฺปยุตฺโต จ ธมฺโม อุปฺปชฺชติ เหตุปจฺจยา. ตีณิ. (สํขิตฺตํ.)}

\prose{เหตุยา นว.}

\csromanheader{2}{82. กิเลส\-วิปฺปยุตฺต\-สํกิเลสิกทุก}
\prosecont{กิเลส\-วิปฺปยุตฺตํ นสํกิเลสิกํ ธมฺมํ ปฏิจฺจ กิเลส\-วิปฺปยุตฺโต นสํกิเลสิโก ธมฺโม อุปฺปชฺชติ เหตุปจฺจยา.  กิเลส\-วิปฺปยุตฺตํ นสํกิเลสิกํ ธมฺมํ ปฏิจฺจ กิเลส\-วิปฺปยุตฺโต นอสํกิเลสิโก ธมฺโม อุปฺปชฺชติ เหตุปจฺจยา.  กิเลส\-วิปฺปยุตฺตํ นสํกิเลสิกํ ธมฺมํ ปฏิจฺจ}

\vspace{\tipitakaparskip}
\csromancenter{ทสฺสเนน\-ปหาตพฺพ\-เหตุกทุก กิเลส\-วิปฺปยุตฺโต นสํกิเลสิโก จ กิเลส\-วิปฺปยุตฺโต นอสํกิเลสิโก จ ธมฺมา อุปฺปชฺชนฺติ เหตุปจฺจยา. (3)}
\prose{\csromanfolio{99}กิเลส\-วิปฺปยุตฺตํ นอสํกิเลสิกํ ธมฺมํ ปฏิจฺจ กิเลส\-วิปฺปยุตฺโต นอสํกิเลสิโก ธมฺโม อุปฺปชฺชติ เหตุปจฺจยา. (1)}

\prose{กิเลส\-วิปฺปยุตฺตํ นสํกิเลสิกญฺจ กิเลส\-วิปฺปยุตฺตํ นอสํกิเลสิกญฺจ ธมฺมํ ปฏิจฺจ กิเลส\-วิปฺปยุตฺโต นอสํกิเลสิโก ธมฺโม อุปฺปชฺชติ เหตุปจฺจยา. (1) (สํขิตฺตํ.)}

\prose{เหตุยา ปญฺจ. (สพฺพตฺถ วิตฺถาโร.)}

\csromanheader{2}{83. ทสฺสเนน\-ปหาตพฺพทุก}
\prose{76. นทสฺสเนน ปหาตพฺพํ ธมฺมํ ปฏิจฺจ นทสฺสเนน ปหาตพฺโพ}

\prose{นนทสฺสเนน ปหาตพฺพํ ธมฺมํ ปฏิจฺจ นนทสฺสเนน ปหาตพฺโพ ธมฺโม อุปฺปชฺชติ เหตุปจฺจยา. ตีณิ.}

\prose{นทสฺสเนน ปหาตพฺพญฺจ นนทสฺสเนน ปหาตพฺพญฺจ ธมฺมํ ปฏิจฺจ นทสฺสเนน ปหาตพฺโพ ธมฺโม อุปฺปชฺชติ เหตุปจฺจยา. (1)}

\prose{เหตุยา ปญฺจ, อารมฺมเณ เทฺว.}

\csromanheader{2}{84. ภาวนาย\-ปหาตพฺพทุก}
\proseitem{77}{นภาวนาย ปหาตพฺพํ ธมฺมํ ปฏิจฺจ นภาวนาย ปหาตพฺโพ ธมฺโม อุปฺปชฺชติ เหตุปจฺจยา. (1)}

\prose{นนภาวนาย ปหาตพฺพํ ธมฺมํ ปฏิจฺจ นนภาวนาย ปหาตพฺโพ ธมฺโม อุปฺปชฺชติ เหตุปจฺจยา. ตีณิ.}

\prose{นภาวนาย ปหาตพฺพญฺจ นนภาวนาย ปหาตพฺพญฺจ ธมฺมํ ปฏิจฺจ นภาวนาย ปหาตพฺโพ ธมฺโม อุปฺปชฺชติ เหตุปจฺจยา. (1)}

\prose{เหตุยา ปญฺจ.}

\csromanheader{2}{85. ทสฺสเนน\-ปหาตพฺพ\-เหตุกทุก}
\proseitem{78}{นทสฺสเนน ปหาตพฺพ\-เหตุกํ ธมฺมํ ปฏิจฺจ นทสฺสเนน ปหาตพฺพ\-เหตุโก ธมฺโม อุปฺปชฺชติ เหตุปจฺจยา. ตีณิ.}

\prose{\csromanfolio{100}นนทสฺสเนน ปหาตพฺพ\-เหตุกํ ธมฺมํ ปฏิจฺจ นนทสฺสเนน ปหาตพฺพ\-เหตุโก ธมฺโม อุปฺปชฺชติ เหตุปจฺจยา. ตีณิ.}

\prose{นทสฺสเนน ปหาตพฺพ\-เหตุ\-กญฺจ นนทสฺสเนน ปหาตพฺพ\-เหตุ\-กญฺจ ธมฺมํ ปฏิจฺจ นทสฺสเนน ปหาตพฺพ\-เหตุโก ธมฺโม อุปฺปชฺชติ เหตุปจฺจยา. ตีณิ. (สํขิตฺตํ.)}

\prose{เหตุยา นว.}

\csromanheader{2}{86. ภาวนาย\-ปหาตพฺพ\-เหตุกทุก}
\proseitem{79}{นภาวนาย ปหาตพฺพ\-เหตุกํ ธมฺมํ ปฏิจฺจ นภาวนาย ปหาตพฺพ\-เหตุโก ธมฺโม อุปฺปชฺชติ เหตุปจฺจยา. ตีณิ.}

\prose{นนภาวนาย ปหาตพฺพ\-เหตุกํ ธมฺมํ ปฏิจฺจ นนภาวนาย ปหาตพฺพ\-เหตุโก ธมฺโม อุปฺปชฺชติ เหตุปจฺจยา. ตีณิ.}

\prose{นภาวนาย ปหาตพฺพ\-เหตุ\-กญฺจ นนภาวนาย ปหาตพฺพ\-เหตุ\-กญฺจ ธมฺมํ ปฏิจฺจ นภาวนาย ปหาตพฺพ\-เหตุโก ธมฺโม อุปฺปชฺชติ เหตุปจฺจยา. ตีณิ. (สํขิตฺตํ.)}

\prose{เหตุยา นว.}

\csromanheader{2}{87. \textbf{สวิตกฺกทุก}}
\proseitem{80}{นสวิตกฺกํ ธมฺมํ ปฏิจฺจ นสวิตกฺโก ธมฺโม อุปฺปชฺชติ เหตุปจฺจยา. ตีณิ.}

\prose{นอวิตกฺกํ ธมฺมํ ปฏิจฺจ นอวิตกฺโก ธมฺโม อุปฺปชฺชติ เหตุปจฺจยา. ตีณิ.}

\prose{นสวิตกฺกญฺจ นอวิตกฺกญฺจ ธมฺมํ ปฏิจฺจ นสวิตกฺโก ธมฺโม อุปฺปชฺชติ เหตุปจฺจยา. ตีณิ. (สํขิตฺตํ.)}

\prose{เหตุยา นว.}

\csromanheader{2}{88. \textbf{สวิจารทุก}}
\proseitem{81}{นสวิจารํ ธมฺมํ ปฏิจฺจ นสวิจาโร ธมฺโม อุปฺปชฺชติ เหตุปจฺจยา. ตีณิ.}

\prose{นอวิจารํ ธมฺมํ ปฏิจฺจ นอวิจาโร ธมฺโม อุปฺปชฺชติ เหตุปจฺจยา. ตีณิ.}

\csromanheader{2}{91. สุขสหคตทุก}
\prose{\csromanfolio{101}นสวิจารญฺจ นอวิจารญฺจ ธมฺมํ ปฏิจฺจ นสวิจาโร ธมฺโม อุปฺปชฺชติ เหตุปจฺจยา. ตีณิ. (สํขิตฺตํ.)}

\prose{เหตุยา นว.}

\csromanheader{2}{89. \textbf{สปฺปีติกทุก}}
\proseitem{82}{นสปฺปีติกํ ธมฺมํ ปฏิจฺจ นสปฺปีติโก ธมฺโม อุปฺปชฺชติ เหตุปจฺจยา. ตีณิ.}

\prose{นอปฺปีติกํ ธมฺมํ ปฏิจฺจ นอปฺปีติโก ธมฺโม อุปฺปชฺชติ เหตุปจฺจยา. ตีณิ.}

\prose{นสปฺปีติกญฺจ นอปฺปีติกญฺจ ธมฺมํ ปฏิจฺจ นสปฺปีติโก ธมฺโม อุปฺปชฺชติ เหตุปจฺจยา. ตีณิ. (สํขิตฺตํ.)}

\prose{เหตุยา นว.}

\csromanheader{2}{90. \textbf{ปีติสหคตทุก}}
\proseitem{83}{นปีติสหคตํ ธมฺมํ ปฏิจฺจ นปีติสหคโต ธมฺโม อุปฺปชฺชติ เหตุปจฺจยา. ตีณิ.}

\prose{นนปี\-ติ\-สหคตํ ธมฺมํ ปฏิจฺจ นนปี\-ติ\-สหคโต ธมฺโม อุปฺปชฺชติ เหตุปจฺจยา. ตีณิ.}

\prose{นปีติสหคตญฺจ นนปีติสหคตญฺจ ธมฺมํ ปฏิจฺจ นปีติสหคโต ธมฺโม อุปฺปชฺชติ เหตุปจฺจยา. ตีณิ. (สํขิตฺตํ.)}

\prose{เหตุยา นว.}

\csromanheader{2}{91. สุขสหคตทุก}
\proseitem{84}{นสุขสหคตํ ธมฺมํ ปฏิจฺจ นสุขสหคโต ธมฺโม อุปฺปชฺชติ เหตุปจฺจยา. ตีณิ.}

\prose{นนสุข\-สหคตํ ธมฺมํ ปฏิจฺจ นนสุข\-สหคโต ธมฺโม อุปฺปชฺชติ เหตุปจฺจยา. ตีณิ.}

\prose{นสุขสหคตญฺจ นนสุขสหคตญฺจ ธมฺมํ ปฏิจฺจ นสุขสหคโต ธมฺโม อุปฺปชฺชติ เหตุปจฺจยา. ตีณิ. (สํขิตฺตํ.)}

\prose{เหตุยา นว.}

\csromanmarkh{ธมฺม\-ปจฺจนีย ทุกปฏฺฐาน\-ปาฬิ}\csromanmarkhii{92. อุเปกฺขาสหคตทุก}\csromanheadernomark{2}{ธมฺม\-ปจฺจนีย ทุกปฏฺฐาน\-ปาฬิ 92. อุเปกฺขาสหคตทุก}
\proseitem{85}{\csromanfolio{102}นอุเปกฺขาสหคตํ ธมฺมํ ปฏิจฺจ นอุเปกฺขาสหคโต ธมฺโม อุปฺปชฺชติ เหตุปจฺจยา. ตีณิ.}

\prose{นนอุเปกฺขาสหคตํ ธมฺมํ ปฏิจฺจ นนอุเปกฺขาสหคโต ธมฺโม อุปฺปชฺชติ เหตุปจฺจยา. ตีณิ.}

\prose{นอุเปกฺขาสหคตญฺจ นนอุเปกฺขาสหคตญฺจ ธมฺมํ ปฏิจฺจ นอุเปกฺขาสหคโต ธมฺโม อุปฺปชฺชติ เหตุปจฺจยา. ตีณิ. (สํขิตฺตํ.)}

\prose{เหตุยา นว.}

\csromanheader{2}{93. \textbf{กามาวจรทุก}}
\proseitem{86}{นกามาวจรํ ธมฺมํ ปฏิจฺจ นกามาวจโร ธมฺโม อุปฺปชฺชติ เหตุปจฺจยา. ตีณิ.}

\prose{นนกามา\-วจรํ ธมฺมํ ปฏิจฺจ นนกามาวจโร ธมฺโม อุปฺปชฺชติ เหตุปจฺจยา. ตีณิ.}

\prose{นกามาวจรญฺจ นนกามาวจรญฺจ ธมฺมํ ปฏิจฺจ นกามาวจโร ธมฺโม อุปฺปชฺชติ เหตุปจฺจยา. ตีณิ. (สํขิตฺตํ.)}

\prose{เหตุยา นว.}

\csromanheader{2}{94. \textbf{รูปาวจรทุก}}
\proseitem{87}{นรูปาวจรํ ธมฺมํ ปฏิจฺจ นรูปาวจโร ธมฺโม อุปฺปชฺชติ เหตุปจฺจยา. ตีณิ.}

\prose{นนรูปาวจรํ ธมฺมํ ปฏิจฺจ นนรูปาวจโร ธมฺโม อุปฺปชฺชติ เหตุปจฺจยา. ตีณิ.}

\prose{นรูปาวจรญฺจ นนรูปาวจรญฺจ ธมฺมํ ปฏิจฺจ นรูปาวจโร ธมฺโม อุปฺปชฺชติ เหตุปจฺจยา. ตีณิ. (สํขิตฺตํ.)}

\prose{เหตุยา นว.}

\csromanheader{2}{95. \textbf{อรูปาวจรทุก}}
\vspace{\tipitakaparskip}
\csromancenter{นอรูปาวจรํ ธมฺมํ ปฏิจฺจ นอรูปาวจโร ธมฺโม อุปฺปชฺชติ เหตุปจฺจยา. เอกํ.}
\csromanheader{2}{98. นิยตทุก}
\prose{\csromanfolio{103}นนอรูปาวจรํ ธมฺมํ ปฏิจฺจ นนอรูปาวจโร ธมฺโม อุปฺปชฺชติ เหตุปจฺจยา. ตีณิ.}

\prose{นอรูปาวจรญฺจ นนอรูปาวจรญฺจ ธมฺมํ ปฏิจฺจ นอรูปาวจโร ธมฺโม อุปฺปชฺชติ เหตุปจฺจยา. เอกํ. (สํขิตฺตํ.)}

\prose{เหตุยา ปญฺจ.}

\csromanheader{2}{96. \textbf{ปริยาปนฺนทุก}}
\proseitem{89}{นปริยาปนฺนํ ธมฺมํ ปฏิจฺจ นปริยาปนฺโน ธมฺโม อุปฺปชฺชติ เหตุปจฺจยา. ตีณิ.}

\prose{นอปริยาปนฺนํ ธมฺมํ ปฏิจฺจ นอปริยาปนฺโน ธมฺโม อุปฺปชฺชติ เหตุปจฺจยา. เอกํ.}

\prose{นปริยาปนฺนญฺจ นอปริยาปนฺนญฺจ ธมฺมํ ปฏิจฺจ นอปริยาปนฺโน ธมฺโม อุปฺปชฺชติ เหตุปจฺจยา. เอกํ. (สํขิตฺตํ.)}

\prose{เหตุยา ปญฺจ.}

\csromanheader{2}{97. \textbf{นิยฺยานิกทุก}}
\proseitem{90}{นนิยฺยานิกํ ธมฺมํ ปฏิจฺจ นนิยฺยานิโก ธมฺโม อุปฺปชฺชติ เหตุปจฺจยา. เอกํ.}

\prose{นอนิยฺยานิกํ ธมฺมํ ปฏิจฺจ นอนิยฺยานิโก ธมฺโม อุปฺปชฺชติ เหตุปจฺจยา. ตีณิ.}

\prose{นนิยฺยานิกญฺจ นอนิยฺยานิกญฺจ ธมฺมํ ปฏิจฺจ นนิยฺยานิโก ธมฺโม อุปฺปชฺชติ เหตุปจฺจยา. เอกํ. (สํขิตฺตํ.)}

\prose{เหตุยา ปญฺจ.}

\csromanheader{2}{98. นิยตทุก}
\proseitem{91}{นนิยตํ ธมฺมํ ปฏิจฺจ นนิยโต ธมฺโม อุปฺปชฺชติ เหตุปจฺจยา. เอกํ.}

\prose{นอนิยตํ ธมฺมํ ปฏิจฺจ นอนิยโต ธมฺโม อุปฺปชฺชติ เหตุปจฺจยา. ตีณิ.}

\prose{\csromanfolio{104}นนิยตญฺจ นอนิยตญฺจ ธมฺมํ ปฏิจฺจ นนิยโต ธมฺโม อุปฺปชฺชติ เหตุปจฺจยา. เอกํ. (สํขิตฺตํ.)}

\prose{เหตุยา ปญฺจ.}

\prose{99. สอุตฺตรทุก}

\proseitem{92}{นสอุตฺตรํ ธมฺมํ ปฏิจฺจ นสอุตฺตโร ธมฺโม อุปฺปชฺชติ เหตุปจฺจยา. ตีณิ.}

\prose{นอนุตฺตรํ ธมฺมํ ปฏิจฺจ นอนุตฺตโร ธมฺโม อุปฺปชฺชติ เหตุปจฺจยา. เอกํ.}

\prose{นสอุตฺตรญฺจ นอนุตฺตรญฺจ ธมฺมํ ปฏิจฺจ นอนุตฺตโร ธมฺโม อุปฺปชฺชติ เหตุปจฺจยา. เอกํ. (สํขิตฺตํ.)}

\prose{เหตุยา ปญฺจ.}

\csromanheader{2}{100. \textbf{สรณทุก}}
\proseitem{93}{นสรณํ ธมฺมํ ปฏิจฺจ นสรโณ ธมฺโม อุปฺปชฺชติ เหตุปจฺจยา. เอกํ.}

\prose{นอรณํ ธมฺมํ ปฏิจฺจ นอรโณ ธมฺโม อุปฺปชฺชติ เหตุปจฺจยา. ตีณิ.}

\prose{นสรณญฺจ นอรณญฺจ ธมฺมํ ปฏิจฺจ นสรโณ ธมฺโม อุปฺปชฺชติ เหตุปจฺจยา. เอกํ. (สํขิตฺตํ.)}

\prosewithcloser{เหตุยา ปญฺจ.}{ปิฏฺฐิทุกํ นิฏฺฐิตํ.}
\nitthitam{ธมฺม\-ปจฺจนีเย ทุกปฏฺฐานํ นิฏฺฐิตํ.}
\csromanpalirecto
\pitaka{อภิธมฺม\-ปิฏก}
\gambhira{ทุกติก\-ปฏฺฐาน\-ปาฬิ}
\namakkaram{นโม ตสฺส ภควโต อรหโต สมฺมา\-สมฺพุทฺธสฺส.}
\csromanmarkh{1. เหตุทุก}\csromanmarkhii{1. กุสลตฺติก}\csromanheadernomark{2}{1. \textbf{เหตุทุก 1. กุสลตฺติก}}
\csromanheader{2}{1-7. ปฏิจฺจาทิวาร ปจฺจยจตุกฺกเหตุ}
\proseitem{1}{\csromanfolio{105}นเหตุํ นกุสลํ ธมฺมํ ปฏิจฺจ นเหตุ นกุสโล ธมฺโม อุปฺปชฺชติ เหตุปจฺจยา. นเหตุํ อกุสลํ อพฺยากตํ เอกํ ขนฺธํ ปฏิจฺจ ตโย ขนฺธา จิตฺต\-สมุฏฺฐานญฺจ รูปํ.  นเหตุํ นกุสลํ ธมฺมํ ปฏิจฺจ นนเหตุํ นกุสโล ธมฺโม อุปฺปชฺชติ เหตุปจฺจยา. นเหตู อกุสเล อพฺยากเต ขนฺเธ ปฏิจฺจ เหตู.  นเหตุํ นกุสลํ ธมฺมํ ปฏิจฺจ นเหตุ นกุสโล จ นนเหตุ นกุสโล จ ธมฺมา อุปฺปชฺชนฺติ เหตุปจฺจยา. (3)}

\prose{นนเหตุํ นกุสลํ ธมฺมํ ปฏิจฺจ นนเหตุ นกุสโล ธมฺโม อุปฺปชฺชติ เหตุปจฺจยา.  นนเหตุํ นกุสลํ ธมฺมํ ปฏิจฺจ นเหตุ นกุสโล ธมฺโม อุปฺปชฺชติ เหตุปจฺจยา.  นนเหตุํ นกุสลํ ธมฺมํ ปฏิจฺจ นเหตุ นกุสโล จ นนเหตุ นกุสโล จ ธมฺมา อุปฺปชฺชนฺติ เหตุปจฺจยา. (3)}

\prose{นเหตุํ นกุสลญฺจ นนเหตุํ นกุสลญฺจ ธมฺมํ ปฏิจฺจ นเหตุ นกุสโล ธมฺโม อุปฺปชฺชติ เหตุปจฺจยา.  นเหตุํ นกุสลญฺจ นนเหตุํ นกุสลญฺจ ธมฺมํ ปฏิจฺจ นนเหตุ นกุสโล ธมฺโม อุปฺปชฺชติ เหตุปจฺจยา.  นเหตุํ นกุสลญฺจ นนเหตุํ นกุสลญฺจ ธมฺมํ ปฏิจฺจ นเหตุ นกุสโล จ นนเหตุ นกุสโล จ ธมฺมา อุปฺปชฺชนฺติ เหตุปจฺจยา. (3) (สํขิตฺตํ.)}

\csromanpalirecto
\gambhira{ธมฺม\-ปจฺจนีย ทุกติก\-ปฏฺฐาน\-ปาฬิ}
\proseitem{2}{\csromanfolio{106}เหตุยา นว ฯเปฯ อวิคเต นว.}

\prose{นเหตุยา เทฺว.}

\prose{(สหชาตวาเรปิ ฯเปฯ สมฺปยุตฺตวาเรปิ สพฺพตฺถ วิตฺถาโร.)}

\csromanheader{2}{\textbf{เหตุ อารมฺมณ}}
\proseitem{3}{นนเหตุ นกุสโล ธมฺโม นนเหตุสฺส นกุสลสฺส ธมฺมสฺส เหตุปจฺจเยน ปจฺจโย.  นนเหตุ นกุสโล ธมฺโม นเหตุสฺส นกุสลสฺส ธมฺมสฺส เหตุปจฺจเยน ปจฺจโย.  นนเหตุ นกุสโล ธมฺโม นเหตุสฺส นกุสลสฺส จ นนเหตุสฺส นกุสลสฺส จ ธมฺมสฺส เหตุปจฺจเยน ปจฺจโย. (3)}

\prose{นเหตุ นกุสโล ธมฺโม นเหตุสฺส นกุสลสฺส ธมฺมสฺส อารมฺมณ\-ปจฺจเยน ปจฺจโย. (สํขิตฺตํ.)}

\proseitem{4}{เหตุยา ตีณิ, อารมฺมเณ นว ฯเปฯ อวิคเต นว. (ปญฺหาวารํ เอวํ วิตฺถาเรตพฺพํ.)}

\proseitem{5}{นเหตุํ นอกุสลํ ธมฺมํ ปฏิจฺจ นเหตุ นอกุสโล ธมฺโม อุปฺปชฺชติ เหตุปจฺจยา. นเหตุํ กุสลํ อพฺยากตํ เอกํ ขนฺธํ ปฏิจฺจ ตโย ขนฺธา จิตฺต\-สมุฏฺฐานญฺจ รูปํ ฯเปฯ เทฺว ขนฺเธ ฯเปฯ.  นเหตุํ นอกุสลํ ธมฺมํ ปฏิจฺจ นนเหตุ นอกุสโล ธมฺโม อุปฺปชฺชติ เหตุปจฺจยา. นเหตู กุสเล อพฺยากเต ขนฺเธ ปฏิจฺจ เหตู.  นเหตุํ นอกุสลํ ธมฺมํ ปฏิจฺจ นเหตุ นอกุสโล จ นนเหตุ นอกุสโล จ ธมฺมา อุปฺปชฺชนฺติ เหตุปจฺจยา. (3)}

\prose{นนเหตุํ นอกุสลํ ธมฺมํ ปฏิจฺจ นนเหตุ นอกุสโล ธมฺโม อุปฺปชฺชติ เหตุปจฺจยา.  นนเหตุํ นอกุสลํ ธมฺมํ ปฏิจฺจ นเหตุ นอกุสโล ธมฺโม อุปฺปชฺชติ เหตุปจฺจยา.  นนเหตุํ นอกุสลํ ธมฺมํ ปฏิจฺจ นเหตุ นอกุสโล จ นนเหตุ นอกุสโล จ ธมฺมา อุปฺปชฺชนฺติ เหตุปจฺจยา. (3)}

\prose{นเหตุํ นอกุสลญฺจ นนเหตุํ นอกุสลญฺจ ธมฺมํ ปฏิจฺจ นเหตุ นอกุสโล ธมฺโม อุปฺปชฺชติ เหตุปจฺจยา.  นเหตุํ นอกุสลญฺจ นนเหตุํ นอกุสลญฺจ ธมฺมํ ปฏิจฺจ นนเหตุ นอกุสโล ธมฺโม อุปฺปชฺชติ}

\vspace{\tipitakaparskip}
\csromancenter{เหตุทุก 2. เวทนาตฺติก เหตุปจฺจยา.  นเหตุํ นอกุสลญฺจ นนเหตุํ นอกุสลญฺจ ธมฺมํ ปฏิจฺจ นเหตุ นอกุสโล จ นนเหตุ นอกุสโล จ ธมฺมา อุปฺปชฺชนฺติ เหตุปจฺจยา. (3) (สํขิตฺตํ.)}
\prose{\csromanfolio{107}เหตุยา นว ฯเปฯ อวิคเต นว. (สพฺพตฺถ นว.)}

\proseitem{6}{นเหตุํ นอพฺยากตํ ธมฺมํ ปฏิจฺจ นเหตุ นอพฺยากโต ธมฺโม อุปฺปชฺชติ เหตุปจฺจยา. กุสลากุสลํ นเหตุํ เอกํ ขนฺธํ ปฏิจฺจ ตโย ขนฺธา. ตโย ขนฺเธ ปฏิจฺจ เอโก ขนฺธา ฯเปฯ.  นเหตุํ นอพฺยากตํ ธมฺมํ ปฏิจฺจ นนเหตุ นอพฺยากโต ธมฺโม อุปฺปชฺชติ เหตุปจฺจยา.  นเหตุํ นอพฺยากตํ ธมฺมํ ปฏิจฺจ นเหตุ นอพฺยากโต จ นนเหตุ นอพฺยากโต จ ธมฺมา อุปฺปชฺชนฺติ เหตุปจฺจยา. (3)}

\prose{นนเหตุํ นอพฺยากตํ ธมฺมํ ปฏิจฺจ นนเหตุ นอพฺยากโต ธมฺโม อุปฺปชฺชติ เหตุปจฺจยา. ตีณิ.}

\prose{นเหตุํ นอพฺยากตญฺจ นนเหตุํ นอพฺยากตญฺจ ธมฺมํ ปฏิจฺจ นเหตุ นอพฺยากโต ธมฺโม อุปฺปชฺชติ เหตุปจฺจยา.  นเหตุํ นอพฺยากตญฺจ นนเหตุํ นอพฺยากตญฺจ ธมฺมํ ปฏิจฺจ นนเหตุ นอพฺยากโต ธมฺโม อุปฺปชฺชติ เหตุปจฺจยา.  นเหตุํ นอพฺยากตญฺจ นนเหตุํ นอพฺยากตญฺจ ธมฺมํ ปฏิจฺจ นเหตุ นอพฺยากโต จ นนเหตุ นอพฺยากโต จ ธมฺมา อุปฺปชฺชนฺติ เหตุปจฺจยา. (3) (สํขิตฺตํ.)}

\prose{เหตุยา นว ฯเปฯ อวิคเต นว. (สพฺพตฺถ นว.)}

\csromanmarkh{1. เหตุทุก}\csromanmarkhii{2. เวทนาตฺติก}\csromanheadernomark{2}{1. เหตุทุก 2. เวทนาตฺติก}
\prosecont{นเหตุํ นสุขาย เวทนาย สมฺปยุตฺตํ ธมฺมํ ปฏิจฺจ นเหตุ นสุขาย เวทนาย สมฺปยุตฺโต ธมฺโม อุปฺปชฺชติ เหตุปจฺจยา.  นเหตุํ นสุขาย เวทนาย สมฺปยุตฺตํ ธมฺมํ ปฏิจฺจ นนเหตุ นสุขาย เวทนาย สมฺปยุตฺโต ธมฺโม อุปฺปชฺชติ เหตุปจฺจยา.}

\prose{เหตุยา นว. (สพฺพตฺถ นว.)}

\proseitem{8}{นเหตุํ นทุกฺขาย เวทนาย สมฺปยุตฺตํ ธมฺมํ ปฏิจฺจ นเหตุ นทุกฺขาย เวทนาย สมฺปยุตฺโต ธมฺโม อุปฺปชฺชติ เหตุปจฺจยา.}

\prose{เหตุยา นว. (สพฺพตฺถ นว.)}

\proseitem{9}{\csromanfolio{108}นเหตุํ นอทุกฺขม\-สุขาย เวทนาย สมฺปยุตฺตํ ธมฺมํ ปฏิจฺจ นเหตุ นอทุกฺขม\-สุขาย เวทนาย สมฺปยุตฺโต ธมฺโม อุปฺปชฺชติ เหตุปจฺจยา. (สํขิตฺตํ.)}

\prose{เหตุยา นว. (สพฺพตฺถ นว.)}

\csromanheader{2}{1. \textbf{เหตุทุก 3-9. วิปากตฺติกาทิ}}
\proseitem{10}{นเหตุํ นวิปากํ ธมฺมํ ปฏิจฺจ ฯเปฯ.}

\prose{นเหตุํ นวิปาก\-ธมฺม\-ธมฺมํ ปฏิจฺจ ฯเปฯ.}

\prose{นเหตุํ นเน\-ว\-วิปาก\-น\-วิปากธมฺม\-ธมฺมํ ปฏิจฺจ ฯเปฯ.}

\proseitem{11}{นเหตุํ นอุปาทินฺนุปาทานิยํ ธมฺมํ ปฏิจฺจ ฯเปฯ.}

\prose{นเหตุํ นอนุปาทินฺนุ\-ปาทานิยํ ธมฺมํ ปฏิจฺจ ฯเปฯ.}

\prose{นเหตุํ นอนุปาทินฺน\-อนุปาทานิยํ ธมฺมํ ปฏิจฺจ ฯเปฯ.}

\proseitem{12}{นเหตุํ นสํกิลิฏฺฐ\-สํกิเลสิกํ ธมฺมํ ปฏิจฺจ ฯเปฯ.}

\prose{นเหตุํ นอสํกิลิฏฺฐ\-สํกิเลสิกํ ธมฺมํ ปฏิจฺจ ฯเปฯ.}

\prose{นเหตุํ นอสํกิลิฏฺฐ\-อสํกิเลสิกํ ธมฺมํ ปฏิจฺจ ฯเปฯ.}

\proseitem{13}{นเหตุํ นสวิตกฺก\-สวิจารํ ธมฺมํ ปฏิจฺจ ฯเปฯ.}

\prose{นเหตุํ นอวิตกฺก\-วิจารมตฺตํ ธมฺมํ ปฏิจฺจ ฯเปฯ.}

\prose{นเหตุํ นอวิตกฺก\-อวิจารํ ธมฺมํ ปฏิจฺจ ฯเปฯ.}

\proseitem{14}{นเหตุํ นปีติสหคตํ ธมฺมํ ปฏิจฺจ ฯเปฯ.}

\prose{นเหตุํ นสุขสหคตํ ธมฺมํ ปฏิจฺจ ฯเปฯ.}

\prose{นเหตุํ นอุเปกฺขาสหคตํ ธมฺมํ ปฏิจฺจ ฯเปฯ.}

\proseitem{15}{นเหตุํ นทสฺสเนน ปหาตพฺพํ ธมฺมํ ปฏิจฺจ ฯเปฯ.}

\prose{นเหตุํ นภาวนาย ปหาตพฺพํ ธมฺมํ ปฏิจฺจ ฯเปฯ.}

\prose{นเหตุํ นเน\-ว\-ทสฺสเนน นภาวนาย ปหาตพฺพํ ธมฺมํ ปฏิจฺจ ฯเปฯ.}

\csromanheader{2}{1. เหตุทุก 10-21. อาจยคามิตฺติกาทิ}
\proseitem{16}{\csromanfolio{109}นเหตุํ นทสฺสเนน ปหาตพฺพ\-เหตุกํ ธมฺมํ ปฏิจฺจ ฯเปฯ.}

\prose{นเหตุํ นภาวนาย ปหาตพฺพ\-เหตุกํ ธมฺมํ ปฏิจฺจ ฯเปฯ.}

\prose{นเหตุํ นเน\-ว\-ทสฺสเนน นภาวนาย ปหาตพฺพ\-เหตุกํ ธมฺมํ ปฏิจฺจ ฯเปฯ.}

\csromanheader{2}{1. เหตุทุก 10-21. อาจยคามิตฺติกาทิ}
\proseitem{17}{นเหตุํ นอาจยคามิํ ธมฺมํ ปฏิจฺจ ฯเปฯ.}

\prose{นเหตุํ นอปจยคามิํ ธมฺมํ ปฏิจฺจ ฯเปฯ.}

\prose{นเหตุํ นเนวอาจยคามิํ นาปจยคามิํ ธมฺมํ ปฏิจฺจ ฯเปฯ.}

\proseitem{18}{นเหตุํ นเสกฺขํ ธมฺมํ ปฏิจฺจ ฯเปฯ.}

\prose{นเหตุํ นอเสกฺขํ ธมฺมํ ปฏิจฺจ ฯเปฯ.}

\prose{นเหตุํ นเน\-ว\-เสกฺข\-นาเสกฺขํ ธมฺมํ ปฏิจฺจ ฯเปฯ.}

\proseitem{19}{นเหตุํ นปริตฺตํ ธมฺมํ ปฏิจฺจ ฯเปฯ.}

\prose{นเหตุํ นมหคฺคตํ ธมฺมํ ปฏิจฺจ ฯเปฯ.}

\prose{นเหตุํ นอปฺปมาณํ ธมฺมํ ปฏิจฺจ ฯเปฯ.}

\proseitem{20}{นเหตุํ นปริตฺตา\-รมฺมณํ ธมฺมํ ปฏิจฺจ ฯเปฯ.}

\prose{นเหตุํ นมหคฺคตา\-รมฺมณํ ธมฺมํ ปฏิจฺจ ฯเปฯ.}

\prose{นเหตุํ นอปฺปมาณา\-รมฺมณํ ธมฺมํ ปฏิจฺจ ฯเปฯ.}

\proseitem{21}{นเหตุํ นหีนํ ธมฺมํ ปฏิจฺจ ฯเปฯ.}

\prose{นเหตุํ นมชฺฌิมํ ธมฺมํ ปฏิจฺจ ฯเปฯ.}

\prose{นเหตุํ นปณีตํ ธมฺมํ ปฏิจฺจ ฯเปฯ.}

\proseitem{22}{นเหตุํ นมิจฺฉตฺต\-นิยตํ ธมฺมํ ปฏิจฺจ ฯเปฯ.}

\prose{นเหตุํ นสมฺมตฺต\-นิยตํ ธมฺมํ ปฏิจฺจ ฯเปฯ.}

\prose{นเหตุํ นอนิยตํ ธมฺมํ ปฏิจฺจ ฯเปฯ.}

\proseitem{23}{\csromanfolio{110}นเหตุํ นมคฺคา\-รมฺมณํ ธมฺมํ ปฏิจฺจ ฯเปฯ.}

\prose{นเหตุํ นมคฺคเหตุกํ ธมฺมํ ปฏิจฺจ ฯเปฯ.}

\prose{นเหตุํ นมคฺคา\-ธิปติํ ธมฺมํ ปฏิจฺจ ฯเปฯ.}

\proseitem{24}{นเหตุํ นอนุปฺปนฺนํ ธมฺมํ ปฏิจฺจ ฯเปฯ.}

\prose{นเหตุํ นอุปฺปาทิํ ธมฺมํ ปฏิจฺจ ฯเปฯ.}

\proseitem{25}{นเหตุํ นอตีตํ ธมฺมํ ปฏิจฺจ ฯเปฯ.}

\prose{นเหตุํ นอนาคตํ ธมฺมํ ปฏิจฺจ ฯเปฯ.}

\proseitem{26}{นเหตุํ นอตีตา\-รมฺมณํ ธมฺมํ ปฏิจฺจ ฯเปฯ.}

\prose{นเหตุํ นอนาคตา\-รมฺมณํ ธมฺมํ ปฏิจฺจ ฯเปฯ.}

\prose{นเหตุํ นปจฺจุปฺปนฺนา\-รมฺมณํ ธมฺมํ ปฏิจฺจ ฯเปฯ.}

\proseitem{27}{นเหตุํ นอชฺฌตฺตํ ธมฺมํ ปฏิจฺจ ฯเปฯ.}

\prose{นเหตุํ นพหิทฺธา ธมฺมํ ปฏิจฺจ ฯเปฯ.}

\proseitem{28}{นเหตุํ นอชฺฌตฺตา\-รมฺมณํ ธมฺมํ ปฏิจฺจ ฯเปฯ.}

\prose{นเหตุํ นพหิทฺธา\-รมฺมณํ ธมฺมํ ปฏิจฺจ ฯเปฯ.}

\csromanmarkh{1. เหตุทุก}\csromanmarkhii{22. สนิทสฺสน\-ตฺติก}\csromanheadernomark{2}{1. เหตุทุก 22. สนิทสฺสน\-ตฺติก}
\csromanheader{2}{1-7. ปฏิจฺจาทิวาร ปจฺจยจตุกฺกเหตุ}
\proseitem{29}{นเหตุํ นสนิทสฺสน\-สปฺปฏิฆํ ธมฺมํ ปฏิจฺจ นเหตุ นสนิทสฺสน\-สปฺปฏิโฆ ธมฺโม อุปฺปชฺชติ เหตุปจฺจยา.  นเหตุํ นสนิทสฺสน\-สปฺปฏิฆํ ธมฺมํ ปฏิจฺจ นนเหตุ นสนิทสฺสน\-สปฺปฏิโฆ ธมฺโม อุปฺปชฺชติ เหตุปจฺจยา.  นเหตุํ นสนิทสฺสน\-สปฺปฏิฆํ ธมฺมํ ปฏิจฺจ นเหตุ นสนิทสฺสน\-สปฺปฏิโฆ จ นนเหตุ นสนิทสฺสน\-สปฺปฏิโฆ จ ธมฺมา อุปฺปชฺชนฺติ เหตุปจฺจยา. (3)}

\prose{นนเหตุํ นสนิทสฺสน\-สปฺปฏิฆํ ธมฺมํ ปฏิจฺจ นนเหตุ นสนิทสฺสน\-สปฺปฏิโฆ ธมฺโม อุปฺปชฺชติ เหตุปจฺจยา.  นนเหตุํ นสนิทสฺสน\-สปฺปฏิฆํ ธมฺมํ}

\vspace{\tipitakaparskip}
\csromancenter{เหตุทุก 22. สนิทสฺสน\-ตฺติก ปฏิจฺจ นเหตุ นสนิทสฺสน\-สปฺปฏิโฆ ธมฺโม อุปฺปชฺชติ เหตุปจฺจยา.  นนเหตุํ นสนิทสฺสน\-สปฺปฏิฆํ ธมฺมํ ปฏิจฺจ นเหตุ นสนิทสฺสน\-สปฺปฏิโฆ จ นนเหตุ นสนิทสฺสน\-สปฺปฏิโฆ จ ธมฺมา อุปฺปชฺชนฺติ เหตุปจฺจยา. (3)}
\prose{\csromanfolio{111}นเหตุํ นสนิทสฺสนสปฺปฏิฆญฺจ นนเหตุํ นสนิทสฺสนสปฺปฏิฆญฺจ ธมฺมํ ปฏิจฺจ นเหตุ นสนิทสฺสน\-สปฺปฏิโฆ ธมฺโม อุปฺปชฺชติ เหตุปจฺจยา.  นเหตุํ นสนิทสฺสนสปฺปฏิฆญฺจ นนเหตุํ นสนิทสฺสนสปฺปฏิฆญฺจ ธมฺมํ ปฏิจฺจ นนเหตุ นสนิทสฺสน\-สปฺปฏิโฆ ธมฺโม อุปฺปชฺชติ เหตุปจฺจยา.  นเหตุํ นสนิทสฺสนสปฺปฏิฆญฺจ นนเหตุํ นสนิทสฺสนสปฺปฏิฆญฺจ ธมฺมํ ปฏิจฺจ นเหตุ นสนิทสฺสน\-สปฺปฏิโฆ จ นนเหตุ นสนิทสฺสน\-สปฺปฏิโฆ จ ธมฺมา อุปฺปชฺชนฺติ เหตุปจฺจยา. (3) (สํขิตฺตํ.)}

\prose{เหตุยา นว ฯเปฯ อวิคเต นว. (สพฺพตฺถ นว.)}

\proseitem{30}{นเหตุํ นอนิทสฺสน\-สปฺปฏิฆํ ธมฺมํ ปฏิจฺจ นเหตุ นอนิทสฺสน\-สปฺปฏิโฆ ธมฺโม อุปฺปชฺชติ เหตุปจฺจยา.  นเหตุํ นอนิทสฺสน\-สปฺปฏิฆํ ธมฺมํ ปฏิจฺจ นนเหตุ นอนิทสฺสน\-สปฺปฏิโฆ ธมฺโม อุปฺปชฺชติ เหตุปจฺจยา.  นเหตุํ นอนิทสฺสน\-สปฺปฏิฆํ ธมฺมํ ปฏิจฺจ นเหตุ นอนิทสฺสน\-สปฺปฏิโฆ จ นนเหตุ นอนิทสฺสน\-สปฺปฏิโฆ จ ธมฺมา อุปฺปชฺชนฺติ เหตุปจฺจยา. (3)}

\prose{นนเหตุํ นอนิทสฺสน\-สปฺปฏิฆํ ธมฺมํ ปฏิจฺจ นนเหตุ นอนิทสฺสน\-สปฺปฏิโฆ ธมฺโม อุปฺปชฺชติ เหตุปจฺจยา.  นนเหตุํ นอนิทสฺสน\-สปฺปฏิฆํ ธมฺมํ ปฏิจฺจ นเหตุ นอนิทสฺสน\-สปฺปฏิโฆ ธมฺโม อุปฺปชฺชติ เหตุปจฺจยา.  นนเหตุํ นอนิทสฺสน\-สปฺปฏิฆํ ธมฺมํ ปฏิจฺจ นเหตุ นอนิทสฺสน\-สปฺปฏิโฆ จ นนเหตุ นอนิทสฺสน\-สปฺปฏิโฆ จ ธมฺมา อุปฺปชฺชนฺติ เหตุปจฺจยา. (3)}

\prose{นเหตุํ นอนิทสฺสนสปฺปฏิฆญฺจ นนเหตุํ นอนิทสฺสนสปฺปฏิฆญฺจ ธมฺมํ ปฏิจฺจ นเหตุ นอนิทสฺสน\-สปฺปฏิโฆ ธมฺโม อุปฺปชฺชติ เหตุปจฺจยา.  นเหตุํ นอนิทสฺสนสปฺปฏิฆญฺจ นนเหตุํ นอนิทสฺสนสปฺปฏิฆญฺจ ธมฺมํ ปฏิจฺจ นนเหตุ นอนิทสฺสน\-สปฺปฏิโฆ ธมฺโม อุปฺปชฺชติ เหตุปจฺจยา.  นเหตุํ นอนิทสฺสนสปฺปฏิฆญฺจ นนเหตุํ นอนิทสฺสนสปฺปฏิฆญฺจ ธมฺมํ ปฏิจฺจ นเหตุ นอนิทสฺสน\-สปฺปฏิโฆ จ นนเหตุ นอนิทสฺสน\-สปฺปฏิโฆ จ ธมฺมา อุปฺปชฺชนฺติ เหตุปจฺจยา. (3) (สํขิตฺตํ.)}

\prose{เหตุยา นว ฯเปฯ อวิคเต นว. (สหชาตวาเรปิ. สมฺปยุตฺตวาเรปิ สพฺพตฺถ วิตฺถาโร.)}

\csromanheader{2}{ธมฺม\-ปจฺจนีย ทุกติก\-ปฏฺฐาน\-ปาฬิ เหตุอารมฺมณ}
\proseitem{31}{\csromanfolio{112}นนเหตุ นอนิทสฺสน\-สปฺปฏิโฆ ธมฺโม นนเหตุสฺส นอนิทสฺสน\-สปฺปฏิฆสฺส ธมฺมสฺส เหตุปจฺจเยน ปจฺจโย.  นนเหตุ นอนิทสฺสน\-สปฺปฏิโฆ ธมฺโม นเหตุสฺส นอนิทสฺสน\-สปฺปฏิฆสฺส ธมฺมสฺส เหตุปจฺจเยน ปจฺจโย.  นนเหตุ นอนิทสฺสน\-สปฺปฏิโฆ ธมฺโม นเหตุสฺส นอนิทสฺสน\-สปฺปฏิฆสฺส จ นนเหตุสฺส นอนิทสฺสน\-สปฺปฏิฆสฺส จ ธมฺมสฺส เหตุปจฺจเยน ปจฺจโย. (3)}

\prose{นเหตุ นอนิทสฺสน\-สปฺปฏิโฆ ธมฺโม นเหตุสฺส นอนิทสฺสน\-สปฺปฏิฆสฺส ธมฺมสฺส อารมฺมณ\-ปจฺจเยน ปจฺจโย.}

\prose{เหตุยา ตีณิ, อารมฺมเณ นว ฯเปฯ อุปนิสฺสเย นว, ปุเรชาเต ปจฺฉาชาเต ตีณิ, อาเสวเน นว, กมฺเม ตีณิ, วิปาเก นว, อาหาเร ตีณิ, อินฺทฺริเย นว, ฌาเน ตีณิ, มคฺเค สมฺปยุตฺเต นว, วิปฺปยุตฺเต ปญฺจ ฯเปฯ อวิคเต นว. (ปญฺหาวารํ เอวํ วิตฺถาเรตพฺพํ.)}

\proseitem{32}{นเหตุํ นอนิทสฺสน\-อปฺปฏิฆํ ธมฺมํ ปฏิจฺจ นเหตุ นอนิทสฺสน\-อปฺปฏิโฆ ธมฺโม อุปฺปชฺชติ เหตุปจฺจยา. (1) (สํขิตฺตํ.)}

\prose{เหตุยา เอกํ. (สพฺพตฺถ วิตฺถาโร.)}

\csromanmarkh{2. สเหตุกทุก}\csromanmarkhii{1. กุสลตฺติก}\csromanheadernomark{2}{2. \textbf{สเหตุกทุก 1. กุสลตฺติก}}
\prosecont{นสเหตุกํ นกุสลํ ธมฺมํ ปฏิจฺจ นสเหตุโก นกุสโล ธมฺโม อุปฺปชฺชติ เหตุปจฺจยา. วิจิกิจฺฉา\-สหคตํ อุทฺธจฺจ\-สหคตํ โมหํ ปฏิจฺจ จิตฺต\-สมุฏฺฐานํ รูปํ. เอกํ มหาภูตํ ปฏิจฺจ ตโย มหาภูตา ฯเปฯ เทฺว มหาภูเต ฯเปฯ.  นสเหตุกํ นกุสลํ ธมฺมํ ปฏิจฺจ นอเหตุโก นกุสโล ธมฺโม อุปฺปชฺชติ เหตุปจฺจยา. วิจิกิจฺฉา\-สหคตํ อุทฺธจฺจ\-สหคตํ โมหํ ปฏิจฺจ สมฺปยุตฺตกา ขนฺธา. ปฏิสนฺธิ\-กฺขเณ วตฺถุํ ปฏิจฺจ สเหตุกา ขนฺธา.  นสเหตุกํ นกุสลํ ธมฺมํ ปฏิจฺจ นสเหตุโก นกุสโล จ นอเหตุโก นกุสโล จ ธมฺมา อุปฺปชฺชนฺติ เหตุปจฺจยา. (3)}

\proseitem{33}{นอเหตุกํ นกุสลํ ธมฺมํ นอเหตุโก นกุสโล ธมฺโม อุปฺปชฺชติ เหตุปจฺจยา. ตีณิ.}

\csromanmarkh{4. เหตุ\-สเหตุกทุก}\csromanmarkhii{1. กุสลตฺติก}\csromanheadernomark{2}{4. เหตุ\-สเหตุกทุก 1. กุสลตฺติก}
\prose{\csromanfolio{113}นสเหตุกํ นกุสลญฺจ นอเหตุกํ นกุสลญฺจ ธมฺมํ ปฏิจฺจ นสเหตุโก นกุสโล ธมฺโม อุปฺปชฺชติ เหตุปจฺจยา. ตีณิ. (สํขิตฺตํ.)}

\prose{เหตุยา นว.}

\proseitem{34}{นสเหตุกํ นอกุสลํ ธมฺมํ ปฏิจฺจ นสเหตุโก นอกุสโล ธมฺโม อุปฺปชฺชติ เหตุปจฺจยา. ตีณิ.}

\prose{นอเหตุกํ นอกุสลํ ธมฺมํ ปฏิจฺจ นอเหตุโก นอกุสโล ธมฺโม อุปฺปชฺชติ เหตุปจฺจยา. ตีณิ.}

\prose{นสเหตุกํ นอกุสลญฺจ นอเหตุกํ นอกุสลญฺจ ธมฺมํ ปฏิจฺจ นสเหตุโก นอกุสโล ธมฺโม อุปฺปชฺชติ เหตุปจฺจยา. ตีณิ.}

\prose{เหตุยา นว.}

\proseitem{35}{นสเหตุกํ นอพฺยากตํ ธมฺมํ ปฏิจฺจ นอเหตุโก นอพฺยากโต ธมฺโม อุปฺปชฺชติ เหตุปจฺจยา. (1)}

\prose{นอเหตุกํ นอพฺยากตํ ธมฺมํ ปฏิจฺจ นอเหตุโก นอพฺยากโต}

\prose{นสเหตุกํ นอพฺยากตญฺจ นอเหตุกํ นอพฺยากตญฺจ ธมฺมํ ปฏิจฺจ นอเหตุโก นอพฺยากโต ธมฺโม อุปฺปชฺชติ เหตุปจฺจยา. (1)}

\prose{เหตุยา ตีณิ. (เหตุสมฺปยุตฺตทุกํ สเหตุกทุกสทิสํ.)}

\vspace{\tipitakaparskip}
\csromancenter{(เหตุ\-สมฺปยุตฺตทุกํ สเหตุก\-ทุก\-สทิสํ.)}
\csromanmarkh{4. เหตุ\-สเหตุกทุก}\csromanmarkhii{1. กุสลตฺติก}\csromanheadernomark{2}{4. เหตุ\-สเหตุกทุก 1. กุสลตฺติก}
\proseitem{36}{นเหตุญฺเจว นอเหตุกญฺจ นกุสลํ ธมฺมํ ปฏิจฺจ นเหตุ เจว นอเหตุโก จ นกุสโล ธมฺโม อุปฺปชฺชติ เหตุปจฺจยา.  นเหตุญฺเจว นอเหตุกญฺจ นกุสลํ ธมฺมํ ปฏิจฺจ นอเหตุโก เจว นน จ เหตุ นกุสโล ธมฺโม อุปฺปชฺชติ เหตุปจฺจยา.  นเหตุญฺเจว นอเหตุกญฺจ นกุสลํ ธมฺมํ ปฏิจฺจ นเหตุ เจว นอเหตุโก จ นกุสโล จ นอเหตุโก เจว นน จ เหตุ นกุสโล จ ธมฺมา อุปฺปชฺชนฺติ เหตุปจฺจยา. (3)}

\prose{\csromanfolio{114}นอเหตุกญฺเจว นน จ เหตุํ นกุสลํ ธมฺมํ ปฏิจฺจ นอเหตุโก เจว นน จ เหตุ นกุสโล ธมฺโม อุปฺปชฺชติ เหตุปจฺจยา. ตีณิ.}

\prose{นเหตุญฺเจว นอเหตุกญฺจ นกุสลญฺจ นอเหตุกญฺเจว นน จ เหตุํ นกุสลญฺจ ธมฺมํ ปฏิจฺจ นเหตุ เจว นอเหตุโก จ นกุสโล ธมฺโม อุปฺปชฺชติ เหตุปจฺจยา. ตีณิ. (สํขิตฺตํ.)}

\prose{เหตุยา นว.}

\proseitem{37}{นเหตุญฺเจว นอเหตุกญฺจ นอกุสลํ ธมฺมํ ปฏิจฺจ นเหตุ เจว นอเหตุโก จ นอกุสโล ธมฺโม อุปฺปชฺชติ เหตุปจฺจยา.}

\prose{เหตุยา นว.}

\proseitem{38}{นเหตุญฺเจว นอเหตุกญฺจ นอพฺยากตํ ธมฺมํ ปฏิจฺจ นเหตุ เจว นอเหตุโก จ นอพฺยากโต ธมฺโม อุปฺปชฺชติ เหตุปจฺจยา.}

\prose{เหตุยา นว.}

\vspace{\tipitakaparskip}
\csromancenter{(เหตุ\-เหตุ\-สมฺปยุตฺตทุกํ เหตุ\-สเหตุก\-ทุก\-สทิสํ.)}
\csromanmarkh{6. นเหตุ\-สเหตุกทุก}\csromanmarkhii{1. กุสลตฺติก}\csromanheadernomark{2}{6. นเหตุ\-สเหตุกทุก 1. กุสลตฺติก}
\proseitem{39}{นเหตุํ นสเหตุกํ นกุสลํ ธมฺมํ ปฏิจฺจ นเหตุ นสเหตุโก นกุสโล ธมฺโม อุปฺปชฺชติ เหตุปจฺจยา. (สํขิตฺตํ.)}

\prose{เหตุยา นว.}

\proseitem{40}{นเหตุํ นสเหตุกํ นอกุสลํ ธมฺมํ ปฏิจฺจ นเหตุ นสเหตุโก นอกุสโล ธมฺโม อุปฺปชฺชติ เหตุปจฺจยา. (สํขิตฺตํ.)}

\prose{เหตุยา นว.}

\proseitem{41}{นเหตุํ นอเหตุกํ นอพฺยากตํ ธมฺมํ ปฏิจฺจ นเหตุ นอเหตุโก นอพฺยากโต ธมฺโม อุปฺปชฺชติ เหตุปจฺจยา. (เอกปญฺหเมว. เอเตน อุปาเยน กุสลากุสลทุกํ กุสล\-ตฺติกเมว เอตฺถ วฏฺฏติ.)}

\csromanmarkh{14. อาสวโคจฺฉก}\csromanmarkhii{1. กุสลตฺติก 7. จูฬนฺตรทุก 1. กุสลตฺติก}\csromanheadernomark{2}{14. อาสวโคจฺฉก 1. กุสลตฺติก \textbf{7. จูฬนฺตรทุก 1. กุสลตฺติก}}
\proseitem{42}{\csromanfolio{115}นอปจฺจยํ นกุสลํ ธมฺมํ ปฏิจฺจ ฯเปฯ. (สพฺพตฺถ เอกํ.) นอสงฺขตํ นกุสลํ ธมฺมํ ปฏิจฺจ ฯเปฯ.}

\proseitem{43}{นสนิทสฺสนํ นกุสลํ ธมฺมํ ปฏิจฺจ.}

\proseitem{44}{นสปฺปฏิฆํ นกุสลํ ธมฺมํ ปฏิจฺจ ฯเปฯ.  นอปฺปฏิฆํ นกุสลํ ธมฺมํ ปฏิจฺจ.}

\proseitem{45}{นรูปิํ นกุสลํ ธมฺมํ ปฏิจฺจ ฯเปฯ.  นอรูปิํ นกุสลํ ธมฺมํ ปฏิจฺจ.}

\proseitem{46}{นโลกิยํ นกุสลํ ธมฺมํ ปฏิจฺจ ฯเปฯ.  นโลกุตฺตรํ นกุสลํ ธมฺมํ ปฏิจฺจ.}

\proseitem{47}{นเกนจิ วิญฺเญยฺยํ นกุสลํ ธมฺมํ ปฏิจฺจ ฯเปฯ.  นนเกนจิ วิญฺเญยฺยํ นกุสลํ ธมฺมํ ปฏิจฺจ.}

\csromanmarkh{14. อาสวโคจฺฉก}\csromanmarkhii{1. กุสลตฺติก}\csromanheadernomark{2}{14. อาสวโคจฺฉก 1. กุสลตฺติก}
\proseitem{48}{นอาสวํ นกุสลํ ธมฺมํ ปฏิจฺจ ฯเปฯ.  นโนอาสวํ นกุสลํ ธมฺมํ ปฏิจฺจ.}

\proseitem{49}{นสาสวํ นกุสลํ ธมฺมํ ปฏิจฺจ ฯเปฯ.  นอนาสวํ กุสลํ ธมฺมํ ปฏิจฺจ.}

\proseitem{50}{นอาสว\-สมฺปยุตฺตํ นกุสลํ ธมฺมํ ปฏิจฺจ ฯเปฯ.  นอาสว\-วิปฺปยุตฺตํ นกุสลํ ธมฺมํ ปฏิจฺจ.}

\proseitem{51}{นอาสวญฺเจว นอนาสวญฺจ นกุสลํ ธมฺมํ ปฏิจฺจ ฯเปฯ.  นอนาสวญฺเจว นโน จ อาสวํ นกุสลํ ธมฺมํ ปฏิจฺจ.}

\proseitem{52}{นอาสวญฺเจว นอาสว\-วิปฺปยุตฺตญฺจ นกุสลํ ธมฺมํ ปฏิจฺจ ฯเปฯ.  นอาสว\-วิปฺปยุตฺต\-ญฺเจว นโน จ อาสวํ นกุสลํ ธมฺมํ ปฏิจฺจ.}

\proseitem{53}{อาสว\-วิปฺปยุตฺตํ นสาสวํ นกุสลํ ธมฺมํ ปฏิจฺจ ฯเปฯ.  อาสว\-วิปฺปยุตฺตํ นอนาสวํ นกุสลํ ธมฺมํ ปฏิจฺจ.}

\csromanmarkh{ธมฺม\-ปจฺจนีย ทุกติก\-ปฏฺฐาน\-ปาฬิ 20-54. ฉโคจฺฉก}\csromanmarkhii{1. กุสลตฺติก}\csromanheadernomark{2}{ธมฺม\-ปจฺจนีย ทุกติก\-ปฏฺฐาน\-ปาฬิ 20-54. ฉโคจฺฉก 1. กุสลตฺติก}
\proseitem{54}{\csromanfolio{116}โนสญฺโญชนํ นกุสลํ ธมฺมํ ปฏิจฺจ ฯเปฯ.  นโนสญฺโญชนํ นกุสลํ ธมฺมํ ปฏิจฺจ.}

\proseitem{55}{โนคนฺถํ นกุสลํ ธมฺมํ ปฏิจฺจ ฯเปฯ.  โนโอฆํ ฯเปฯ.  โนโยคํ ฯเปฯ.  โนนีวรณํ.}

\prose{โนปรามาสํ นกุสลํ ธมฺมํ ปฏิจฺจ.}

\csromanmarkh{55. มหนฺตรทุก}\csromanmarkhii{1. กุสลตฺติก}\csromanheadernomark{2}{55. \textbf{มหนฺตรทุก 1. กุสลตฺติก}}
\proseitem{56}{นสารมฺมณํ นกุสลํ ธมฺมํ ปฏิจฺจ ฯเปฯ.  นอนารมฺมณํ นกุสลํ ธมฺมํ ปฏิจฺจ.}

\proseitem{57}{นจิตฺตํ นกุสลํ ธมฺมํ ปฏิจฺจ ฯเปฯ.  นโนจิตฺตํ นกุสลํ ธมฺมํ ปฏิจฺจ.}

\prose{58. นเจตสิกํ นกุสลํ ธมฺมํ ปฏิจฺจ. (สํขิตฺตํ.)}

\prose{59. นจิตฺต\-สมฺปยุตฺตํ นกุสลํ ธมฺมํ ปฏิจฺจ. (สํขิตฺตํ.)}

\proseitem{60}{นจิตฺต\-สํสฏฺฐํ นกุสลํ ธมฺมํ ปฏิจฺจ.}

\proseitem{61}{โนจิตฺต\-สมุฏฺฐานํ นกุสลํ ธมฺมํ ปฏิจฺจ.}

\proseitem{62}{โนจิตฺต\-สหภุํ นกุสลํ ธมฺมํ ปฏิจฺจ.}

\proseitem{63}{โนจิตฺตา\-นุปริวตฺติํ นกุสลํ ธมฺมํ ปฏิจฺจ ฯเปฯ.  โนจิตฺต\-สํสฏฺฐ\-สมุฏฺฐานํ นกุสลํ ธมฺมํ ปฏิจฺจ ฯเปฯ.  โนจิตฺต\-สํสฏฺฐ\-สมุฏฺฐาน\-สหภุํ นกุสลํ ธมฺมํ ปฏิจฺจ ฯเปฯ.  โนจิตฺต\-สํสฏฺฐ\-สมุฏฺฐานา\-นุปริวตฺติํ นกุสลํ ธมฺมํ ปฏิจฺจ.}

\proseitem{64}{นอชฺฌตฺติกํ นกุสลํ ธมฺมํ ปฏิจฺจ ฯเปฯ.  นพาหิรํ นกุสลํ ธมฺมํ ปฏิจฺจ.}

\proseitem{65}{นอุปาทา นกุสลํ ธมฺมํ ปฏิจฺจ.}

\proseitem{66}{นอุปาทินฺนํ นกุสลํ ธมฺมํ ปฏิจฺจ.}

\csromanmarkh{83. ปิฏฺฐิทุก}\csromanmarkhii{1. กุสลตฺติก 69. อุปาทาน\-โคจฺฉก 1. กุสลตฺติก}\csromanheadernomark{2}{83. ปิฏฺฐิทุก 1. กุสลตฺติก 69. อุปาทาน\-โคจฺฉก 1. กุสลตฺติก}
\vspace{\tipitakaparskip}
\csromancenter{โนอุปาทานํ นกุสลํ ธมฺมํ ปฏิจฺจ.}
\csromanmarkh{75. กิเลสโคจฺฉก}\csromanmarkhii{1. กุสลตฺติก}\csromanheadernomark{2}{75. \textbf{กิเลสโคจฺฉก 1. กุสลตฺติก}}
\proseitem{68}{\csromanfolio{117}โนกิเลสํ นกุสลํ ธมฺมํ ปฏิจฺจ ฯเปฯ.  นโนกิเลสํ นกุสลํ ธมฺมํ ปฏิจฺจ.}

\proseitem{69}{นสํกิเลสิกํ นกุสลํ ธมฺมํ ปฏิจฺจ ฯเปฯ.  นอสํกิเลสิกํ นกุสลํ ธมฺมํ ปฏิจฺจ.}

\proseitem{70}{นสํกิลิฏฺฐํ นกุสลํ ธมฺมํ ปฏิจฺจ ฯเปฯ.  นอสํกิลิฏฺฐํ นกุสลํ ธมฺมํ ปฏิจฺจ.}

\proseitem{71}{นกิเลส\-สมฺปยุตฺตํ นกุสลํ ธมฺมํ ปฏิจฺจ ฯเปฯ.  นกิเลส\-วิปฺปยุตฺตํ นกุสลํ ธมฺมํ ปฏิจฺจ.}

\proseitem{72}{นกิเลสญฺเจว นอสํกิเลสิกญฺจ นกุสลํ ธมฺมํ ปฏิจฺจ ฯเปฯ.  นอสํกิเลสิก\-ญฺเจว นโน จ กิเลสํ นกุสลํ ธมฺมํ ปฏิจฺจ.}

\proseitem{73}{นกิเลสญฺเจว นอสํกิลิฏฺฐญฺจ นกุสลํ ธมฺมํ ปฏิจฺจ ฯเปฯ. นอสํกิลิฏฺฐ\-ญฺเจว นโน จ กิเลสํ นกุสลํ ธมฺมํ ปฏิจฺจ ฯเปฯ. นกิเลสญฺเจว นกิเลส\-วิปฺปยุตฺตญฺจ นกุสลํ ธมฺมํ ปฏิจฺจ ฯเปฯ. นกิเลส\-วิปฺปยุตฺต\-ญฺเจว นโน จ กิเลสํ นกุสลํ ธมฺมํ ปฏิจฺจ ฯเปฯ. (สพฺพตฺถ นว. วิปากํ นตฺถิ.)}

\prose{กิเลส\-วิปฺปยุตฺตํ นสํกิเลสิกํ นกุสลํ ธมฺมํ ปฏิจฺจ ฯเปฯ. กิเลส\-วิปฺปยุตฺตํ นอสํกิเลสิกํ นกุสลํ ธมฺมํ ปฏิจฺจ.}

\csromanmarkh{83. ปิฏฺฐิทุก}\csromanmarkhii{1. กุสลตฺติก}\csromanheadernomark{2}{83. ปิฏฺฐิทุก 1. กุสลตฺติก}
\proseitem{74}{นทสฺสเนน ปหาตพฺพํ นกุสลํ ธมฺมํ ปฏิจฺจ ฯเปฯ.  นนทสฺสเนน ปหาตพฺพํ นกุสลํ ธมฺมํ ปฏิจฺจ.}

\proseitem{75}{นภาวนาย ปหาตพฺพํ นกุสลํ ธมฺมํ ปฏิจฺจ ฯเปฯ.  นนภาวนาย ปหาตพฺพํ นกุสลํ ธมฺมํ ปฏิจฺจ.}

\proseitem{76}{\csromanfolio{118}นทสฺสเนน ปหาตพฺพ\-เหตุกํ นกุสลํ ธมฺมํ ปฏิจฺจ ฯเปฯ. นนทสฺสเนน ปหาตพฺพ\-เหตุกํ นกุสลํ ธมฺมํ ปฏิจฺจ ฯเปฯ. นภาวนาย ปหาตพฺพ\-เหตุกํ นกุสลํ ธมฺมํ ปฏิจฺจ ฯเปฯ. นนภาวนาย ปหาตพฺพ\-เหตุกํ นกุสลํ ธมฺมํ ปฏิจฺจ.}

\proseitem{77}{นสวิตกฺกํ นกุสลํ ธมฺมํ ปฏิจฺจ ฯเปฯ. นอวิตกฺกํ นกุสลํ ธมฺมํ ปฏิจฺจ ฯเปฯ. นสวิจารํ นกุสลํ ธมฺมํ ปฏิจฺจ ฯเปฯ. นอวิจารํ นกุสลํ ธมฺมํ ปฏิจฺจ ฯเปฯ. นสปฺปีติกํ นกุสลํ ธมฺมํ ปฏิจฺจ ฯเปฯ. นอปฺปีติกํ นกุสลํ ธมฺมํ ปฏิจฺจ ฯเปฯ. นปีติสหคตํ นกุสลํ ธมฺมํ ปฏิจฺจ ฯเปฯ. นนปี\-ติ\-สหคตํ นกุสลํ ธมฺมํ ปฏิจฺจ ฯเปฯ. นสุขสหคตํ นกุสลํ ธมฺมํ ปฏิจฺจ ฯเปฯ. นนสุข\-สหคตํ นกุสลํ ธมฺมํ ปฏิจฺจ ฯเปฯ. นอุเปกฺขาสหคตํ นกุสลํ ธมฺมํ ปฏิจฺจ ฯเปฯ. นนอุเปกฺขาสหคตํ นกุสลํ ธมฺมํ ปฏิจฺจ.}

\proseitem{78}{นกามาวจรํ นกุสลํ ธมฺมํ ปฏิจฺจ ฯเปฯ. นนกามา\-วจรํ นกุสลํ ธมฺมํ ปฏิจฺจ ฯเปฯ. นรูปาวจรํ นกุสลํ ธมฺมํ ปฏิจฺจ ฯเปฯ. นนรูปาวจรํ นกุสลํ ธมฺมํ ปฏิจฺจ.}

\proseitem{79}{นอรูปาวจรํ นกุสลํ ธมฺมํ ปฏิจฺจ ฯเปฯ. นนอรูปาวจรํ นกุสลํ ธมฺมํ ปฏิจฺจ.}

\proseitem{80}{นปริยาปนฺนํ นกุสลํ ธมฺมํ ปฏิจฺจ ฯเปฯ. นอปริยาปนฺนํ นกุสลํ ธมฺมํ ปฏิจฺจ.}

\proseitem{81}{นนิยฺยานิกํ นกุสลํ ธมฺมํ ปฏิจฺจ ฯเปฯ. นอนิยฺยานิกํ นกุสลํ ธมฺมํ ปฏิจฺจ. (สํขิตฺตํ. สพฺพตฺถ เอกํ.)}

\prose{นนิยตํ นกุสลํ ธมฺมํ ปฏิจฺจ ฯเปฯ. นอนิยตํ นกุสลํ ธมฺมํ ปฏิจฺจ. (สํขิตฺตํ.)}

\proseitem{82}{นสอุตฺตรํ นกุสลํ ธมฺมํ ปฏิจฺจ ฯเปฯ. นอนุตฺตรํ นกุสลํ ธมฺมํ ปฏิจฺจ. (สํขิตฺตํ.)}

\prose{83. นสรณํ นกุสลํ ธมฺมํ ปฏิจฺจ นสรโณ นกุสโล}

\csromanmarkh{100. สรณทุก}\csromanmarkhii{2. เวทนาตฺติกาทิ}\csromanheadernomark{2}{100. สรณทุก 2. เวทนาตฺติกาทิ}
\prose{\csromanfolio{119}นอรณํ นกุสลํ ธมฺมํ ปฏิจฺจ นอรโณ นกุสโล ธมฺโม อุปฺปชฺชติ เหตุปจฺจยา.  นอรณํ นกุสลํ ธมฺมํ ปฏิจฺจ นสรโณ นกุสโล ธมฺโม อุปฺปชฺชติ เหตุปจฺจยา.  นอรณํ นกุสลํ ธมฺมํ ปฏิจฺจ นสรโณ นกุสโล จ นอรโณ นกุสโล จ ธมฺมา อุปฺปชฺชนฺติ เหตุปจฺจยา. (3)}

\prose{นสรณํ นกุสลญฺจ นอรณํ นกุสลญฺจ ธมฺมํ ปฏิจฺจ นสรโณ นกุสโล ธมฺโม อุปฺปชฺชติ เหตุปจฺจยา. (1) (สํขิตฺตํ.) เหตุยา ปญฺจ ฯเปฯ อวิคเต ปญฺจ.}

\proseitem{84}{นสรณํ นอกุสลํ ธมฺมํ ปฏิจฺจ นสรโณ นอกุสโล ธมฺโม อุปฺปชฺชติ เหตุปจฺจยา. (สํขิตฺตํ.) เหตุยา เอกํ. (สพฺพตฺถ เอกํ.)}

\prose{85. นสรณํ นอพฺยากตํ ธมฺมํ ปฏิจฺจ นสรโณ นอพฺยากโต}

\prose{นอรณํ นอพฺยากตํ ธมฺมํ ปฏิจฺจ นอรโณ นอพฺยากโต ธมฺโม อุปฺปชฺชติ เหตุปจฺจยา. (1) (สํขิตฺตํ.)}

\prose{เหตุยา เทฺว, อารมฺมเณ เทฺว ฯเปฯ อวิคเต เทฺว.}

\vspace{\tipitakaparskip}
\csromancenter{(สหชาตวาเรปิ ฯเปฯ ปญฺหาวาเรปิ สพฺพตฺถ วิตฺถาโร.)}
\csromanmarkh{100. สรณทุก}\csromanmarkhii{2. เวทนาตฺติกาทิ}\csromanheadernomark{2}{100. สรณทุก 2. เวทนาตฺติกาทิ}
\prose{ปจฺจยจตุกฺกเหตุ}

\proseitem{86}{นสรณํ นสุขาย เวทนาย สมฺปยุตฺตํ ธมฺมํ ปฏิจฺจ ฯเปฯ. นสรณํ นทุกฺขาย เวทนาย สมฺปยุตฺตํ ธมฺมํ ปฏิจฺจ ฯเปฯ. นสรณํ นอทุกฺขม\-สุขาย เวทนาย สมฺปยุตฺตํ ธมฺมํ ปฏิจฺจ.}

\proseitem{87}{นสรณํ นวิปากํ ธมฺมํ ปฏิจฺจ.}

\prose{นสรณํ นวิปาก\-ธมฺม\-ธมฺมํ ปฏิจฺจ.}

\prose{นสรณํ นเน\-ว\-วิปาก\-น\-วิปากธมฺม\-ธมฺมํ ปฏิจฺจ.}

\proseitem{88}{\csromanfolio{120}นสรณํ นอุปาทินฺนุปาทานิยํ ธมฺมํ ปฏิจฺจ.}

\prose{นสรณํ นอนุปาทินฺนุ\-ปาทานิยํ ธมฺมํ ปฏิจฺจ.}

\prose{นสรณํ นอนุปาทินฺน\-อนุปาทานิยํ ธมฺมํ ปฏิจฺจ.}

\csromanmarkh{100. สรณทุก}\csromanmarkhii{22. สนิทสฺสน\-ตฺติก}\csromanheadernomark{2}{100. สรณทุก 22. สนิทสฺสน\-ตฺติก}
\proseitem{89}{นสรณํ นสนิทสฺสน\-สปฺปฏิฆํ ธมฺมํ ปฏิจฺจ นสรโณ นสนิทสฺสน\-สปฺปฏิโฆ ธมฺโม อุปฺปชฺชติ เหตุปจฺจยา. (1)}

\prose{นอรณํ นสนิทสฺสน\-สปฺปฏิฆํ ธมฺมํ ปฏิจฺจ นอรโณ นสนิทสฺสน\-สปฺปฏิโฆ ธมฺโม อุปฺปชฺชติ เหตุปจฺจยา.  นอรณํ นสนิทสฺสน\-สปฺปฏิฆํ ธมฺมํ ปฏิจฺจ นสรโณ นสนิทสฺสน\-สปฺปฏิโฆ ธมฺโม อุปฺปชฺชติ เหตุปจฺจยา.  นอรณํ นสนิทสฺสน\-สปฺปฏิฆํ ธมฺมํ ปฏิจฺจ นสรโณ นสนิทสฺสน\-สปฺปฏิโฆ จ นอรโณ นสนิทสฺสน\-สปฺปฏิโฆ จ ธมฺมา อุปฺปชฺชนฺติ เหตุปจฺจยา. (3)}

\prose{นสรณํ นสนิทสฺสนสปฺปฏิฆญฺจ นอรณํ นสนิทสฺสนสปฺปฏิฆญฺจ ธมฺมํ ปฏิจฺจ นสรโณ นสนิทสฺสน\-สปฺปฏิโฆ ธมฺโม อุปฺปชฺชติ เหตุปจฺจยา. (1) (สํขิตฺตํ.)}

\prose{เหตุยา ปญฺจ.}

\proseitem{90}{นสรณํ นอนิทสฺสน\-สปฺปฏิฆํ ธมฺมํ ปฏิจฺจ นสรโณ นอนิทสฺสน\-สปฺปฏิโฆ ธมฺโม อุปฺปชฺชติ เหตุปจฺจยา. (1)}

\prose{นอรณํ นอนิทสฺสน\-สปฺปฏิฆํ ธมฺมํ ปฏิจฺจ นสรโณ นอนิทสฺสน\-สปฺปฏิโฆ ธมฺโม อุปฺปชฺชติ เหตุปจฺจยา.  นอรณํ นอนิทสฺสน\-สปฺปฏิฆํ ธมฺมํ ปฏิจฺจ นสรโณ นอนิทสฺสน\-สปฺปฏิโฆ จ นอรโณ นอนิทสฺสน\-สปฺปฏิโฆ จ ธมฺมา อุปฺปชฺชนฺติ เหตุปจฺจยา. (3)}

\csromanmarkh{100. สรณทุก}\csromanmarkhii{22. สนิทสฺสน\-ตฺติก}\csromanheadernomark{2}{100. สรณทุก 22. สนิทสฺสน\-ตฺติก}
\prose{\csromanfolio{121}นสรณํ นอนิทสฺสนสปฺปฏิฆญฺจ นอรณํ นอนิทสฺสนสปฺปฏิฆญฺจ ธมฺมํ ปฏิจฺจ นสรโณ นอนิทสฺสน\-สปฺปฏิโฆ ธมฺโม อุปฺปชฺชติ เหตุปจฺจยา. (1) (สํขิตฺตํ.)}

\prose{เหตุยา ปญฺจ, อารมฺมเณ เทฺว ฯเปฯ อวิคเต ปญฺจ.}

\vspace{\tipitakaparskip}
\csromancenter{(สหชาตวาเรปิ ฯเปฯ ปญฺหาวาเรปิ สพฺพตฺถ วิตฺถาเรตพฺพํ.)}
\proseitem{91}{นสรณํ นอนิทสฺสน\-อปฺปฏิฆํ ธมฺมํ ปฏิจฺจ นสรโณ นอนิทสฺสน\-อปฺปฏิโฆ ธมฺโม อุปฺปชฺชติ เหตุปจฺจยา. (สํขิตฺตํ.)}

\prosewithcloser{เหตุยา เอกํ. (สพฺพตฺถ วิตฺถาโร.)}{ธมฺม\-ปจฺจนีเย ทุกติก\-ปฏฺฐานํ นิฏฺฐิตํ.}
\csromanpalirecto
\pitaka{อภิธมฺม\-ปิฏก}
\csromanmatikaanchor{mtk.7}
\csromantocmarknopage{chapter}{ติกทุกปฏฺฐานปาฬิ}{mtk.7}
\bookmark[dest={mtk.7},level=0]{ติกทุกปฏฺฐานปาฬิ}
\gambhira{ติกทุก\-ปฏฺฐาน\-ปาฬิ}
\namakkaram{นโม ตสฺส ภควโต อรหโต สมฺมา\-สมฺพุทฺธสฺส.}
\csromanmarkh{1. กุสลตฺติก}\csromanmarkhii{1. เหตุทุก}\csromanheadernomark{2}{1. \textbf{กุสลตฺติก 1. เหตุทุก}}
\prose{\csromanfolio{123}1. ปฏิจฺจวารเหตุ}

\proseitem{1}{นกุสลํ นเหตุํ ธมฺมํ ปฏิจฺจ นกุสโล นเหตุ ธมฺโม อุปฺปชฺชติ เหตุปจฺจยา.  นกุสลํ นเหตุํ ธมฺมํ ปฏิจฺจ นอกุสโล นเหตุ ธมฺโม อุปฺปชฺชติ เหตุปจฺจยา.  นกุสลํ นเหตุํ ธมฺมํ ปฏิจฺจ นอพฺยากโต นเหตุ ธมฺโม อุปฺปชฺชติ เหตุปจฺจยา.  นกุสลํ นเหตุํ ธมฺมํ ปฏิจฺจ นกุสโล นเหตุ จ นอพฺยากโต นเหตุ จ ธมฺมา อุปฺปชฺชนฺติ เหตุปจฺจยา.  นกุสลํ นเหตุํ ธมฺมํ ปฏิจฺจ นกุสโล นเหตุ จ นอกุสโล นเหตุ จ ธมฺมา อุปฺปชฺชนฺติ เหตุปจฺจยา. ปญฺจ.}

\proseitem{2}{นอกุสลํ นเหตุํ ธมฺมํ ปฏิจฺจ นอกุสโล นเหตุ ธมฺโม อุปฺปชฺชติ เหตุปจฺจยา.  นอกุสลํ นเหตุํ ธมฺมํ ปฏิจฺจ นกุสโล นเหตุ ธมฺโม อุปฺปชฺชติ เหตุปจฺจยา.  นอกุสลํ นเหตุํ ธมฺมํ ปฏิจฺจ นอพฺยากโต นเหตุ ธมฺโม อุปฺปชฺชติ เหตุปจฺจยา.  นอกุสลํ นเหตุํ ธมฺมํ ปฏิจฺจ นอกุสโล นเหตุ จ นอพฺยากโต นเหตุ จ ธมฺมา อุปฺปชฺชนฺติ เหตุปจฺจยา.  นอกุสลํ นเหตุํ ธมฺมํ ปฏิจฺจ นกุสโล นเหตุ จ นอกุสโล นเหตุ จ ธมฺมา อุปฺปชฺชนฺติ เหตุปจฺจยา. ปญฺจ.}

\csromanpalirecto
\gambhira{ธมฺม\-ปจฺจนีย ติกทุก\-ปฏฺฐาน\-ปาฬิ}
\proseitem{3}{\csromanfolio{124}นอพฺยากตํ นเหตุํ ธมฺมํ ปฏิจฺจ นอพฺยากโต นเหตุ ธมฺโม อุปฺปชฺชติ เหตุปจฺจยา.  นอพฺยากตํ นเหตุํ ธมฺมํ ปฏิจฺจ นกุสโล นเหตุ ธมฺโม อุปฺปชฺชติ เหตุปจฺจยา.  นอพฺยากตํ นเหตุํ ธมฺมํ ปฏิจฺจ นอกุสโล นเหตุ ธมฺโม อุปฺปชฺชติ เหตุปจฺจยา.  นอพฺยากตํ นเหตุํ ธมฺมํ ปฏิจฺจ นกุสโล นเหตุ จ นอพฺยากโต นเหตุ จ ธมฺมา อุปฺปชฺชนฺติ เหตุปจฺจยา.  นอพฺยากตํ นเหตุํ ธมฺมํ ปฏิจฺจ นอกุสโล นเหตุ จ นอพฺยากโต นเหตุ จ ธมฺมา อุปฺปชฺชนฺติ เหตุปจฺจยา.  นอพฺยากตํ นเหตุํ ธมฺมํ ปฏิจฺจ นกุสโล นเหตุ จ นอกุสโล นเหตุ จ ธมฺมา อุปฺปชฺชนฺติ เหตุปจฺจยา. ฉ.}

\proseitem{4}{นกุสลํ นเหตุญฺจ นอพฺยากตํ นเหตุญฺจ ธมฺมํ ปฏิจฺจ นกุสโล นเหตุ ธมฺโม อุปฺปชฺชติ เหตุปจฺจยา.  นกุสลํ นเหตุญฺจ นอพฺยากตํ นเหตุญฺจ ธมฺมํ ปฏิจฺจ นอกุสโล นเหตุ ธมฺโม อุปฺปชฺชติ เหตุปจฺจยา.  นกุสลํ นเหตุญฺจ นอพฺยากตํ นเหตุญฺจ ธมฺมํ ปฏิจฺจ นอพฺยากโต นเหตุ ธมฺโม อุปฺปชฺชติ เหตุปจฺจยา.  นกุสลํ นเหตุญฺจ นอพฺยากตํ นเหตุญฺจ ธมฺมํ ปฏิจฺจ นกุสโล นเหตุ จ นอพฺยากโต นเหตุ จ ธมฺมา อุปฺปชฺชนฺติ เหตุปจฺจยา.  นกุสลํ นเหตุญฺจ นอพฺยากตํ นเหตุญฺจ ธมฺมํ ปฏิจฺจ นกุสโล นเหตุ จ นอกุสโล นเหตุ จ ธมฺมา อุปฺปชฺชนฺติ เหตุปจฺจยา. ปญฺจ.}

\proseitem{5}{นอกุสลํ นเหตุญฺจ นอพฺยากตํ นเหตุญฺจ ธมฺมํ ปฏิจฺจ นกุสโล นเหตุ ธมฺโม อุปฺปชฺชติ เหตุปจฺจยา.  นอกุสลํ นเหตุญฺจ นอพฺยากตํ นเหตุญฺจ ธมฺมํ ปฏิจฺจ นอกุสโล นเหตุ ธมฺโม อุปฺปชฺชติ เหตุปจฺจยา.  นอกุสลํ นเหตุญฺจ นอพฺยากตํ นเหตุญฺจ ธมฺมํ ปฏิจฺจ นอพฺยากโต นเหตุ ธมฺโม อุปฺปชฺชติ เหตุปจฺจยา.  นอกุสลํ นเหตุญฺจ นอพฺยากตํ นเหตุญฺจ ธมฺมํ ปฏิจฺจ นอกุสโล นเหตุ จ นอพฺยากโต นเหตุ จ ธมฺมา อุปฺปชฺชนฺติ เหตุปจฺจยา.  นอกุสลํ นเหตุญฺจ นอพฺยากตํ นเหตุญฺจ ธมฺมํ ปฏิจฺจ นกุสโล นเหตุ จ นอกุสโล นเหตุ จ ธมฺมา อุปฺปชฺชนฺติ เหตุปจฺจยา. ปญฺจ.}

\proseitem{6}{นกุสลํ นเหตุญฺจ นอกุสลํ นเหตุญฺจ ธมฺมํ ปฏิจฺจ นกุสโล นเหตุ ธมฺโม อุปฺปชฺชติ เหตุปจฺจยา.  นกุสลํ นเหตุญฺจ นอกุสลํ นเหตุญฺจ ธมฺมํ ปฏิจฺจ นอกุสโล นเหตุ ธมฺโม อุปฺปชฺชติ เหตุปจฺจยา.  นกุสลํ นเหตุญฺจ นอกุสลํ นเหตุญฺจ ธมฺมํ ปฏิจฺจ}

\vspace{\tipitakaparskip}
\csromancenter{2. สเหตุกทุก นกุสโล นเหตุ จ นอกุสโล นเหตุ จ ธมฺมา อุปฺปชฺชนฺติ เหตุปจฺจยา. ตีณิ. (สํขิตฺตํ.)}
\prose{\csromanfolio{125}เหตุยา เอกูนติํส, อารมฺมเณ จตุวีส ฯเปฯ อวิคเต เอกูนติํส. (สพฺพตฺถ วิตฺถาโร.)}

\proseitem{7}{นกุสลํ นนเหตุํ ธมฺมํ ปฏิจฺจ นกุสโล นนเหตุ ธมฺโม อุปฺปชฺชติ เหตุปจฺจยา.  นกุสลํ นนเหตุํ ธมฺมํ ปฏิจฺจ นอกุสโล นนเหตุ ธมฺโม อุปฺปชฺชติ เหตุปจฺจยา.  นกุสลํ นนเหตุํ ธมฺมํ ปฏิจฺจ นอพฺยากโต นนเหตุ ธมฺโม อุปฺปชฺชติ เหตุปจฺจยา.  นกุสลํ นนเหตุํ ธมฺมํ ปฏิจฺจ นกุสโล นนเหตุ จ นอพฺยากโต นนเหตุ จ ธมฺมา อุปฺปชฺชนฺติ เหตุปจฺจยา.  นกุสลํ นนเหตุํ ธมฺมํ ปฏิจฺจ นกุสโล นนเหตุ จ นอกุสโล นนเหตุ จ ธมฺมา อุปฺปชฺชนฺติ เหตุปจฺจยา. (5)}

\prose{นอกุสลํ นนเหตุํ ธมฺมํ ปฏิจฺจ นอกุสโล นนเหตุ ธมฺโม อุปฺปชฺชติ เหตุปจฺจยา. ปญฺจ.}

\prose{นอพฺยากตํ นนเหตุํ ธมฺมํ ปฏิจฺจ นอพฺยากโต นนเหตุ ธมฺโม อุปฺปชฺชติ เหตุปจฺจยา. ปญฺจ.}

\prose{นกุสลํ นนเหตุญฺจ นอพฺยากตํ นนเหตุญฺจ ธมฺมํ ปฏิจฺจ นกุสโล นนเหตุ ธมฺโม อุปฺปชฺชติ เหตุปจฺจยา. ตีณิ.}

\prose{นอกุสลํ นนเหตุญฺจ นอพฺยากตํ นนเหตุญฺจ ธมฺมํ ปฏิจฺจ นอกุสโล นนเหตุ ธมฺโม อุปฺปชฺชติ เหตุปจฺจยา. ตีณิ.}

\prose{นกุสลํ นนเหตุญฺจ นอกุสลํ นนเหตุญฺจ ธมฺมํ ปฏิจฺจ นกุสโล นนเหตุ ธมฺโม อุปฺปชฺชติ เหตุปจฺจยา. ตีณิ. (สํขิตฺตํ.)}

\prose{เหตุยา จตุวีส, อารมฺมเณ จตุวีส ฯเปฯ อวิคเต จตุวีส. (สพฺพตฺถ วิตฺถาโร.)}

\csromanmarkh{1. กุสลตฺติก}\csromanmarkhii{2. สเหตุกทุก}\csromanheadernomark{2}{1. กุสลตฺติก 2. สเหตุกทุก}
\proseitem{8}{นกุสลํ นสเหตุกํ ธมฺมํ ปฏิจฺจ นกุสโล นสเหตุโก ธมฺโม อุปฺปชฺชติ เหตุปจฺจยา.  นกุสลํ นสเหตุกํ ธมฺมํ ปฏิจฺจ นอกุสโล นสเหตุโก ธมฺโม อุปฺปชฺชติ เหตุปจฺจยา.  นกุสลํ นสเหตุกํ \csromanfolio{126}ธมฺมํ ปฏิจฺจ นกุสโล นสเหตุโก จ นอกุสโล นสเหตุโก จ ธมฺมา อุปฺปชฺชนฺติ เหตุปจฺจยา. (3)}

\prose{นอกุสลํ นสเหตุกํ ธมฺมํ ปฏิจฺจ นอกุสโล นสเหตุโก ธมฺโม อุปฺปชฺชติ เหตุปจฺจยา.  นอกุสลํ นสเหตุกํ ธมฺมํ ปฏิจฺจ นกุสโล นสเหตุโก ธมฺโม อุปฺปชฺชติ เหตุปจฺจยา.  นอกุสลํ นสเหตุกํ ธมฺมํ ปฏิจฺจ นกุสโล นสเหตุโก จ นอกุสโล นสเหตุโก จ ธมฺมา อุปฺปชฺชนฺติ เหตุปจฺจยา. (3)}

\prose{นอพฺยากตํ นสเหตุกํ ธมฺมํ ปฏิจฺจ นกุสโล นสเหตุโก ธมฺโม อุปฺปชฺชติ เหตุปจฺจยา. ตีณิ.}

\prose{นกุสลํ นสเหตุกญฺจ นอพฺยากตํ นสเหตุกญฺจ ธมฺมํ ปฏิจฺจ นกุสโล นสเหตุโก ธมฺโม อุปฺปชฺชติ เหตุปจฺจยา. ตีณิ.}

\prose{นกุสลํ นสเหตุกญฺจ นอกุสลํ นสเหตุกญฺจ ธมฺมํ ปฏิจฺจ นกุสโล นสเหตุโก ธมฺโม อุปฺปชฺชติ เหตุปจฺจยา. ตีณิ.}

\prose{เหตุยา ปนฺนรส, อารมฺมเณ นว ฯเปฯ อธิปติยา นว ฯเปฯ วิปาเก นว ฯเปฯ อวิคเต ปนฺนรส. (สพฺพตฺถ วิตฺถาโร.)}

\proseitem{9}{นกุสลํ นอเหตุกํ ธมฺมํ ปฏิจฺจ นกุสโล นอเหตุโก ธมฺโม อุปฺปชฺชติ เหตุปจฺจยา. ปญฺจ.}

\prose{นอกุสลํ นอเหตุกํ ธมฺมํ ปฏิจฺจ นอกุสโล นอเหตุโก ธมฺโม อุปฺปชฺชติ เหตุปจฺจยา. ปญฺจ.}

\prose{นอพฺยากตํ นอเหตุกํ ธมฺมํ ปฏิจฺจ นอพฺยากโต นอเหตุโก ธมฺโม อุปฺปชฺชติ เหตุปจฺจยา. ปญฺจ.}

\prose{นกุสลํ นอเหตุกญฺจ นอพฺยากตํ นอเหตุกญฺจ ธมฺมํ ปฏิจฺจ นกุสโล นอเหตุโก ธมฺโม อุปฺปชฺชติ เหตุปจฺจยา. ตีณิ.}

\prose{นอกุสลํ นอเหตุกญฺจ นอพฺยากตํ นอเหตุกญฺจ ธมฺมํ ปฏิจฺจ นอกุสโล นอเหตุโก ธมฺโม อุปฺปชฺชติ เหตุปจฺจยา. ตีณิ.}

\prose{นกุสลํ นอเหตุกญฺจ นอกุสลํ นอเหตุกญฺจ ธมฺมํ ปฏิจฺจ นกุสโล นอเหตุโก ธมฺโม อุปฺปชฺชติ เหตุปจฺจยา. ตีณิ.}

\prose{เหตุยา จตุวีส ฯเปฯ อวิคเต จตุวีส. (สพฺพตฺถ วิตฺถาโร.)}

\vspace{\tipitakaparskip}
\csromancenter{14. อาสวโคจฺฉก 1. กุสลตฺติก 3. เหตุ\-สมฺปยุตฺต\-ทุกาทิ}
\proseitem{10}{\csromanfolio{127}นกุสลํ นเหตุ\-สมฺปยุตฺตํ ธมฺมํ ปฏิจฺจ.}

\proseitem{11}{นกุสลํ นเหตุ\-วิปฺปยุตฺตํ ธมฺมํ ปฏิจฺจ.}

\proseitem{12}{นกุสลํ นเหตุญฺเจว นอเหตุกญฺจ ธมฺมํ ปฏิจฺจ ฯเปฯ. นกุสลํ นอเหตุกญฺเจว นนเหตุญฺจ ธมฺมํ ปฏิจฺจ ฯเปฯ. นกุสลํ นเหตุญฺเจว นเหตุ\-วิปฺปยุตฺตญฺจ ธมฺมํ ปฏิจฺจ ฯเปฯ. นกุสลํ นเหตุ\-วิปฺปยุตฺต\-ญฺเจว นนเหตุญฺจ ธมฺมํ ปฏิจฺจ.}

\proseitem{13}{นกุสลํ นเหตุํ นสเหตุกํ ธมฺมํ ปฏิจฺจ.}

\prose{นกุสลํ นเหตุํ นอเหตุกํ ธมฺมํ ปฏิจฺจ.}

\csromanmarkh{1. กุสลตฺติก}\csromanmarkhii{7. จูฬนฺตรทุก}\csromanheadernomark{2}{1. \textbf{กุสลตฺติก 7. จูฬนฺตรทุก}}
\proseitem{14}{นกุสลํ นอปฺปจฺจยํ ธมฺมํ ปฏิจฺจ ฯเปฯ. นกุสลํ นอสงฺขตํ ธมฺมํ ปฏิจฺจ.}

\proseitem{15}{นกุสลํ นสนิทสฺสนํ ธมฺมํ ปฏิจฺจ.}

\proseitem{16}{นกุสลํ นสปฺปฏิฆํ ธมฺมํ ปฏิจฺจ.}

\proseitem{17}{นกุสลํ นอปฺปฏิฆํ ธมฺมํ ปฏิจฺจ.}

\proseitem{18}{นกุสลํ นรูปิํ ธมฺมํ ปฏิจฺจ.}

\proseitem{19}{นกุสลํ นอรูปิํ ธมฺมํ ปฏิจฺจ.}

\proseitem{20}{นกุสลํ นโลกิยํ ธมฺมํ ปฏิจฺจ.}

\proseitem{21}{นกุสลํ นโลกุตฺตรํ ธมฺมํ ปฏิจฺจ.}

\proseitem{22}{นกุสลํ นเกนจิ วิญฺเญยฺยํ ธมฺมํ ปฏิจฺจ ฯเปฯ. นกุสลํ นนเกนจิ วิญฺเญยฺยํ ธมฺมํ ปฏิจฺจ.}

\csromanmarkh{1. กุสลตฺติก}\csromanmarkhii{14. อาสวโคจฺฉก}\csromanheadernomark{2}{1. กุสลตฺติก 14. อาสวโคจฺฉก}
\proseitem{23}{นกุสลํ โนอาสวํ ธมฺมํ ปฏิจฺจ.}

\prose{นกุสลํ นโนอาสวํ ธมฺมํ ปฏิจฺจ.}

\proseitem{24}{\csromanfolio{128}นกุสลํ นสาสวํ ธมฺมํ ปฏิจฺจ.}

\prose{นกุสลํ นอนาสวํ ธมฺมํ ปฏิจฺจ.}

\proseitem{25}{นกุสลํ นอาสว\-สมฺปยุตฺตํ ธมฺมํ ปฏิจฺจ.}

\prose{นกุสลํ นอาสว\-วิปฺปยุตฺตํ ธมฺมํ ปฏิจฺจ.}

\proseitem{26}{นกุสลํ โนอาสวญฺเจว นอนาสวญฺจ ธมฺมํ ปฏิจฺจ.}

\prose{นกุสลํ นอนาสวญฺเจว นโน จ อาสวํ ธมฺมํ ปฏิจฺจ.}

\proseitem{27}{นกุสลํ นอนาสวญฺเจว นอาสว\-วิปฺปยุตฺตญฺจ ธมฺมํ ปฏิจฺจ.}

\prose{นกุสลํ นอาสว\-วิปฺปยุตฺต\-ญฺเจว นโน จ อาสวํ ธมฺมํ ปฏิจฺจ.}

\proseitem{28}{นกุสลํ อาสว\-วิปฺปยุตฺตํ นสาสวํ ธมฺมํ ปฏิจฺจ.}

\prose{นกุสลํ อาสว\-วิปฺปยุตฺตํ นอนาสวํ ธมฺมํ ปฏิจฺจ.}

\csromanmarkh{1. กุสลตฺติก}\csromanmarkhii{20. สญฺโญชน\-โคจฺฉกาทิ}\csromanheadernomark{2}{1. กุสลตฺติก 20. สญฺโญชน\-โคจฺฉกาทิ}
\proseitem{29}{นกุสลํ โนสญฺโญชนํ ธมฺมํ ปฏิจฺจ.}

\proseitem{30}{นกุสลํ โนคนฺถํ ธมฺมํ ปฏิจฺจ ฯเปฯ. นกุสลํ โนโอฆํ ธมฺมํ ปฏิจฺจ ฯเปฯ. นกุสลํ โนโยคํ ธมฺมํ ปฏิจฺจ ฯเปฯ. นกุสลํ โนนีวรณํ ธมฺมํ ปฏิจฺจ ฯเปฯ. นกุสลํ โนปรามาสํ ธมฺมํ ปฏิจฺจ.}

\csromanmarkh{1. กุสลตฺติก}\csromanmarkhii{55. มหนฺตรทุก}\csromanheadernomark{2}{1. \textbf{กุสลตฺติก} 55. มหนฺตรทุก}
\proseitem{31}{นกุสลํ นสารมฺมณํ ธมฺมํ ปฏิจฺจ.}

\prose{นกุสลํ นอนารมฺมณํ ธมฺมํ ปฏิจฺจ. (สํขิตฺตํ.)}

\proseitem{32}{นกุสลํ นจิตฺตํ ธมฺมํ ปฏิจฺจ. (สํขิตฺตํ. นโนจิตฺตปทํ น ลพฺภติ.)}

\proseitem{33}{นกุสลํ นเจตสิกํ ธมฺมํ ปฏิจฺจ.}

\proseitem{34}{นกุสลํ นจิตฺต\-สมฺปยุตฺตํ ธมฺมํ ปฏิจฺจ ฯเปฯ. นกุสลํ นจิตฺต\-สํสฏฺฐํ ธมฺมํ ปฏิจฺจ ฯเปฯ. นกุสลํ นจิตฺต\-สมุฏฺฐานํ ธมฺมํ ปฏิจฺจ.}

\csromanmarkh{1. กุสลตฺติก}\csromanmarkhii{83. ปิฏฺฐิทุก}\csromanheadernomark{2}{1. กุสลตฺติก 83. ปิฏฺฐิทุก}
\proseitem{35}{\csromanfolio{129}นกุสลํ นจิตฺต\-สหภุํ ธมฺมํ ปฏิจฺจ.}

\proseitem{36}{นกุสลํ นจิตฺตา\-นุปริวตฺติํ ธมฺมํ ปฏิจฺจ ฯเปฯ. นกุสลํ นจิตฺต\-สํสฏฺฐ\-สมุฏฺฐานํ ธมฺมํ ปฏิจฺจ ฯเปฯ. นกุสลํ นจิตฺต\-สํสฏฺฐ\-สมุฏฺฐาน\-สหภุํ ธมฺมํ ปฏิจฺจ ฯเปฯ. นกุสลํ นจิตฺต\-สํสฏฺฐ\-สมุฏฺฐานา\-นุ\-ปริวตฺติํ ธมฺมํ ปฏิจฺจ.}

\proseitem{37}{นกุสลํ นอชฺฌตฺติกํ ธมฺมํ ปฏิจฺจ.}

\prose{นกุสลํ นพาหิรํ ธมฺมํ ปฏิจฺจ.}

\proseitem{38}{นกุสลํ นอุปาทา ธมฺมํ ปฏิจฺจ.}

\proseitem{39}{นกุสลํ นอุปาทินฺนํ ธมฺมํ ปฏิจฺจ.}

\prose{นกุสลํ นอนุปาทินฺนํ ธมฺมํ ปฏิจฺจ.}

\csromanmarkh{1. กุสลตฺติก}\csromanmarkhii{69. อุปาทานทุกาทิ}\csromanheadernomark{2}{1. \textbf{กุสลตฺติก} 69. อุปาทานทุกาทิ}
\proseitem{40}{นกุสลํ โนอุปาทานํ ธมฺมํ ปฏิจฺจ.}

\prose{นกุสลํ นโนอุปาทานํ ธมฺมํ ปฏิจฺจ.}

\proseitem{41}{นกุสลํ โนกิเลสํ ธมฺมํ ปฏิจฺจ.}

\prose{นกุสลํ นโนกิเลสํ ธมฺมํ ปฏิจฺจ.}

\csromanmarkh{1. กุสลตฺติก}\csromanmarkhii{83. ปิฏฺฐิทุก}\csromanheadernomark{2}{1. กุสลตฺติก 83. ปิฏฺฐิทุก}
\proseitem{42}{นกุสลํ นทสฺสเนน ปหาตพฺพํ ธมฺมํ ปฏิจฺจ ฯเปฯ. นกุสลํ นนทสฺสเนน ปหาตพฺพํ ธมฺมํ ปฏิจฺจ ฯเปฯ. นกุสลํ นภาวนาย ปหาตพฺพํ ธมฺมํ ปฏิจฺจ ฯเปฯ. นกุสลํ นภาวนาย ปหาตพฺพํ ธมฺมํ ปฏิจฺจ ฯเปฯ. นกุสลํ นทสฺสเนน ปหาตพฺพ\-เหตุกํ ธมฺมํ ปฏิจฺจ ฯเปฯ. นกุสลํ นนทสฺสเนน ปหาตพฺพ\-เหตุกํ ธมฺมํ ปฏิจฺจ ฯเปฯ. นกุสลํ นภาวนาย ปหาตพฺพ\-เหตุกํ ธมฺมํ ปฏิจฺจ.}

\prose{นกุสลํ นนภาวนาย ปหาตพฺพ\-เหตุกํ ธมฺมํ ปฏิจฺจ.}

\proseitem{43}{นกุสลํ นสวิตกฺกํ ธมฺมํ ปฏิจฺจ ฯเปฯ. นกุสลํ นอวิตกฺกํ ธมฺมํ ปฏิจฺจ ฯเปฯ. นกุสลํ นสวิจารํ ธมฺมํ ปฏิจฺจ.}

\prose{นกุสลํ นอวิจารํ ธมฺมํ ปฏิจฺจ.}

\proseitem{44}{\csromanfolio{130}นกุสลํ นสปฺปีติกํ ธมฺมํ ปฏิจฺจ ฯเปฯ. นกุสลํ นอปฺปีติกํ ธมฺมํ ปฏิจฺจ ฯเปฯ. นกุสลํ นปีติสหคตํ ธมฺมํ ปฏิจฺจ ฯเปฯ. นกุสลํ นนปี\-ติ\-สหคตํ ธมฺมํ ปฏิจฺจ ฯเปฯ. นกุสลํ นสุขสหคตํ ธมฺมํ ปฏิจฺจ ฯเปฯ. นกุสลํ นนสุข\-สหคตํ ธมฺมํ ปฏิจฺจ ฯเปฯ. นกุสลํ นอุเปกฺขาสหคตํ ธมฺมํ ปฏิจฺจ.}

\prose{นกุสลํ นนอุเปกฺขาสหคตํ ธมฺมํ ปฏิจฺจ.}

\proseitem{45}{นกุสลํ นกามาวจรํ ธมฺมํ ปฏิจฺจ.}

\prose{นกุสลํ นนกามา\-วจรํ ธมฺมํ ปฏิจฺจ.}

\proseitem{46}{นกุสลํ นรูปาวจรํ ธมฺมํ ปฏิจฺจ ฯเปฯ. นกุสลํ นนรูปาวจรํ ธมฺมํ ปฏิจฺจ ฯเปฯ. นกุสลํ นอรูปาวจรํ ธมฺมํ ปฏิจฺจ.}

\prose{นกุสลํ นนอรูปาวจรํ ธมฺมํ ปฏิจฺจ.}

\proseitem{47}{นกุสลํ นปริยาปนฺนํ ธมฺมํ ปฏิจฺจ.}

\prose{นกุสลํ นอปริยาปนฺนํ ธมฺมํ ปฏิจฺจ.}

\proseitem{48}{นกุสลํ นนิยฺยานิกํ ธมฺมํ ปฏิจฺจ.}

\prose{นกุสลํ นอนิยฺยานิกํ ธมฺมํ ปฏิจฺจ.}

\proseitem{49}{นกุสลํ นนิยตํ ธมฺมํ ปฏิจฺจ.}

\prose{นกุสลํ นอนิยตํ ธมฺมํ ปฏิจฺจ.}

\proseitem{50}{นกุสลํ นสอุตฺตรํ ธมฺมํ ปฏิจฺจ.}

\prose{นกุสลํ นอนุตฺตรํ ธมฺมํ ปฏิจฺจ.}

\proseitem{51}{นกุสลํ นสรณํ ธมฺมํ ปฏิจฺจ นกุสโล นสรโณ ธมฺโม อุปฺปชฺชติ เหตุปจฺจยา.  นกุสลํ นสรณํ ธมฺมํ ปฏิจฺจ นอกุสโล นสรโณ ธมฺโม อุปฺปชฺชติ เหตุปจฺจยา.  นกุสลํ นสรณํ ธมฺมํ ปฏิจฺจ นกุสโล นสรโณ จ นอกุสโล นสรโณ จ ธมฺมา อุปฺปชฺชนฺติ เหตุปจฺจยา. (3)}

\prose{นอกุสลํ นสรณํ ธมฺมํ ปฏิจฺจ นอกุสโล นสรโณ ธมฺโม อุปฺปชฺชติ เหตุปจฺจยา. ปญฺจ.}

\csromanmarkh{2. เวทนาตฺติก}\csromanmarkhii{1. เหตุทุก}\csromanheadernomark{2}{2. เวทนาตฺติก 1. เหตุทุก}
\prose{\csromanfolio{131}นอพฺยากตํ นสรณํ ธมฺมํ ปฏิจฺจ นอพฺยากโต นสรโณ ธมฺโม อุปฺปชฺชติ เหตุปจฺจยา. ปญฺจ.}

\prose{นอกุสลํ นสรณญฺจ นอพฺยากตํ นสรณญฺจ ธมฺมํ ปฏิจฺจ นกุสโล นสรโณ ธมฺโม อุปฺปชฺชติ เหตุปจฺจยา. ปญฺจ.}

\prose{นกุสลํ นสรณญฺจ นอกุสลํ นสรณญฺจ ธมฺมํ ปฏิจฺจ นกุสโล นสรโณ ธมฺโม อุปฺปชฺชติ เหตุปจฺจยา. ตีณิ. (สํขิตฺตํ.)}

\prose{เหตุยา เอกวีส, อารมฺมเณ สตฺตรส ฯเปฯ อวิคเต เอกวีส. (สพฺพตฺถ วิตฺถาโร.)}

\prosewithcloser{นกุสลํ นอรณํ ธมฺมํ ปฏิจฺจ นกุสโล นอรโณ ธมฺโม อุปฺปชฺชติ เหตุปจฺจยา. (สํขิตฺตํ.) เหตุยา นว.}{กุสล\-ตฺติก\-ปิฏฺฐิทุกํ นิฏฺฐิตํ.}
\csromanmarkh{2. เวทนาตฺติก}\csromanmarkhii{1. เหตุทุก}\csromanheadernomark{2}{2. เวทนาตฺติก 1. เหตุทุก}
\prosecont{นสุขาย เวทนาย สมฺปยุตฺตํ นเหตุํ ธมฺมํ ปฏิจฺจ นสุขาย เวทนาย สมฺปยุตฺโต นเหตุ ธมฺโม อุปฺปชฺชติ เหตุปจฺจยา.  นสุขาย เวทนาย สมฺปยุตฺตํ นเหตุํ ธมฺมํ ปฏิจฺจ นทุกฺขาย เวทนาย สมฺปยุตฺโต นเหตุ ธมฺโม อุปฺปชฺชติ เหตุปจฺจยา.  นสุขาย เวทนาย สมฺปยุตฺตํ นเหตุํ ธมฺมํ ปฏิจฺจ นอทุกฺขม\-สุขาย เวทนาย สมฺปยุตฺโต นเหตุ ธมฺโม อุปฺปชฺชติ เหตุปจฺจยา ฯเปฯ. สตฺต.}

\proseitem{52}{นทุกฺขาย เวทนาย สมฺปยุตฺตํ นเหตุํ ธมฺมํ ปฏิจฺจ นทุกฺขาย เวทนาย สมฺปยุตฺโต นเหตุ ธมฺโม อุปฺปชฺชติ เหตุปจฺจยา. สตฺต. (สํขิตฺตํ.)}

\prose{นอทุกฺขม\-สุขาย เวทนาย สมฺปยุตฺตํ นเหตุํ ธมฺมํ ปฏิจฺจ. สตฺต. (สํขิตฺตํ.)}

\prose{นสุขาย เวทนาย สมฺปยุตฺตํ นนเหตุํ ธมฺมํ ปฏิจฺจ นสุขาย เวทนาย สมฺปยุตฺโต นนเหตุ ธมฺโม อุปฺปชฺชติ เหตุปจฺจยา. ปญฺจ.}

\csromanmarkh{ธมฺม\-ปจฺจนีย ติกทุก\-ปฏฺฐาน\-ปาฬิ}\csromanmarkhii{3. วิปากตฺติกาทิ 1. เหตุทุก}\csromanheadernomark{2}{ธมฺม\-ปจฺจนีย ติกทุก\-ปฏฺฐาน\-ปาฬิ 3. วิปากตฺติกาทิ 1. เหตุทุก}
\proseitem{53}{\csromanfolio{132}นวิปากํ นเหตุํ ธมฺมํ ปฏิจฺจ ฯเปฯ. นวิปาก\-ธมฺม\-ธมฺมํ นเหตุํ ธมฺมํ ปฏิจฺจ ฯเปฯ. นเน\-ว\-วิปาก\-น\-วิปากธมฺม\-ธมฺมํ นเหตุํ ธมฺมํ ปฏิจฺจ.}

\proseitem{54}{นอุปาทินฺนุปาทานิยํ นเหตุํ ธมฺมํ ปฏิจฺจ ฯเปฯ. นอนุปาทินฺนุ\-ปาทานิยํ นเหตุํ ธมฺมํ ปฏิจฺจ ฯเปฯ. นอนุปาทินฺน\-อนุปาทานิยํ นเหตุํ ธมฺมํ ปฏิจฺจ.}

\proseitem{55}{นสํกิลิฏฺฐ\-สํกิเลสิกํ นเหตุํ ธมฺมํ ปฏิจฺจ ฯเปฯ. นอสํกิลิฏฺฐ\-สํกิเลสิกํ นเหตุํ ธมฺมํ ปฏิจฺจ ฯเปฯ. นอสํกิลิฏฺฐ\-อสํกิเลสิกํ นเหตุํ ธมฺมํ ปฏิจฺจ.}

\proseitem{56}{นสวิตกฺก\-สวิจารํ นเหตุํ ธมฺมํ ปฏิจฺจ ฯเปฯ. นอวิตกฺก\-วิจารมตฺตํ นเหตุํ ธมฺมํ ปฏิจฺจ ฯเปฯ. นอวิตกฺก\-อวิจารํ นเหตุํ ธมฺมํ ปฏิจฺจ.}

\proseitem{57}{นปีติสหคตํ นเหตุํ ธมฺมํ ปฏิจฺจ ฯเปฯ. นสุขสหคตํ นเหตุํ ธมฺมํ ปฏิจฺจ ฯเปฯ. นอุเปกฺขาสหคตํ นเหตุํ ธมฺมํ ปฏิจฺจ.}

\proseitem{58}{นทสฺสเนน ปหาตพฺพํ นเหตุํ ธมฺมํ ปฏิจฺจ ฯเปฯ. นภาวนาย ปหาตพฺพํ นเหตุํ ธมฺมํ ปฏิจฺจ ฯเปฯ. นเน\-ว\-ทสฺสเนน นภาวนาย ปหาตพฺพํ นเหตุํ ธมฺมํ ปฏิจฺจ ฯเปฯ. นทสฺสเนน ปหาตพฺพ\-เหตุกํ นเหตุํ ธมฺมํ ปฏิจฺจ ฯเปฯ. นภาวนาย ปหาตพฺพ\-เหตุกํ นเหตุํ ธมฺมํ ปฏิจฺจ ฯเปฯ. นเน\-ว\-ทสฺสเนน นภาวนาย ปหาตพฺพ\-เหตุกํ นเหตุํ ธมฺมํ ปฏิจฺจ.}

\proseitem{59}{นอาจยคามิํ นเหตุํ ธมฺมํ ปฏิจฺจ ฯเปฯ. นอปจยคามิํ นเหตุํ ธมฺมํ ปฏิจฺจ ฯเปฯ. นเน\-วา\-จยคามิ\-นา\-ปจยคามิํ นเหตุํ ธมฺมํ ปฏิจฺจ.}

\prose{นเสกฺขํ นเหตุํ ธมฺมํ ปฏิจฺจ ฯเปฯ. นอเสกฺขํ นเหตุํ ธมฺมํ ปฏิจฺจ ฯเปฯ. นเน\-ว\-เสกฺข\-นาเสกฺขํ นเหตุํ ธมฺมํ ปฏิจฺจ.}

\proseitem{60}{นปริตฺตํ นเหตุํ ธมฺมํ ปฏิจฺจ ฯเปฯ. นมหคฺคตํ นเหตุํ ธมฺมํ ปฏิจฺจ ฯเปฯ. นอปฺปมาณํ นเหตุํ ธมฺมํ ปฏิจฺจ.}

\proseitem{61}{นปริตฺตา\-รมฺมณํ นเหตุํ ธมฺมํ ปฏิจฺจ ฯเปฯ. นมหคฺคตา\-รมฺมณํ นเหตุํ ธมฺมํ ปฏิจฺจ ฯเปฯ. นอปฺปมาณา\-รมฺมณํ นเหตุํ ธมฺมํ ปฏิจฺจ.}

\csromanmarkh{22. สนิทสฺสน\-ตฺติก}\csromanmarkhii{1. เหตุทุกาทิ}\csromanheadernomark{2}{22. สนิทสฺสน\-ตฺติก 1. เหตุทุกาทิ}
\proseitem{62}{\csromanfolio{133}นหีนํ นเหตุํ ธมฺมํ ปฏิจฺจ ฯเปฯ. นมชฺฌิมํ นเหตุํ ธมฺมํ ปฏิจฺจ ฯเปฯ. นปณีตํ นเหตุํ ธมฺมํ ปฏิจฺจ.}

\proseitem{63}{นมิจฺฉตฺต\-นิยตํ นเหตุํ ธมฺมํ ปฏิจฺจ ฯเปฯ. นสมฺมตฺต\-นิยตํ นเหตุํ ธมฺมํ ปฏิจฺจ ฯเปฯ. นอนิยตํ นเหตุํ ธมฺมํ ปฏิจฺจ.}

\proseitem{64}{นมคฺคา\-รมฺมณํ นเหตุํ ธมฺมํ ปฏิจฺจ ฯเปฯ. นมคฺคเหตุกํ นเหตุํ ธมฺมํ ปฏิจฺจ ฯเปฯ. นมคฺคา\-ธิปติํ นเหตุํ ธมฺมํ ปฏิจฺจ.}

\proseitem{65}{นอนุปฺปนฺนํ นเหตุํ ธมฺมํ ปฏิจฺจ ฯเปฯ. นอุปฺปาทิํ นเหตุํ ธมฺมํ ปฏิจฺจ.}

\proseitem{66}{นอตีตํ นเหตุํ ธมฺมํ ปฏิจฺจ ฯเปฯ. นอนาคตํ นเหตุํ ธมฺมํ ปฏิจฺจ.}

\proseitem{67}{นอตีตา\-รมฺมณํ นเหตุํ ธมฺมํ ปฏิจฺจ ฯเปฯ. นอนาคตา\-รมฺมณํ นเหตุํ ธมฺมํ ปฏิจฺจ ฯเปฯ. นปจฺจุปฺปนฺนา\-รมฺมณํ นเหตุํ ธมฺมํ ปฏิจฺจ.}

\proseitem{68}{นอชฺฌตฺตํ นเหตุํ ธมฺมํ ปฏิจฺจ ฯเปฯ. นพหิทฺธา นเหตุํ ธมฺมํ ปฏิจฺจ.}

\prose{นอชฺฌตฺตา\-รมฺมณํ นเหตุํ ธมฺมํ ปฏิจฺจ ฯเปฯ. นพหิทฺธา\-รมฺมณํ นเหตุํ ธมฺมํ ปฏิจฺจ.}

\csromanmarkh{22. สนิทสฺสน\-ตฺติก}\csromanmarkhii{1. เหตุทุกาทิ}\csromanheadernomark{2}{22. สนิทสฺสน\-ตฺติก 1. เหตุทุกาทิ}
\proseitem{69}{นสนิทสฺสน\-สปฺปฏิฆํ นเหตุํ ธมฺมํ ปฏิจฺจ นสนิทสฺสน\-สปฺปฏิโฆ นเหตุ ธมฺโม อุปฺปชฺชติ เหตุปจฺจยา.  นสนิทสฺสน\-สปฺปฏิฆํ นเหตุํ ธมฺมํ ปฏิจฺจ นอนิทสฺสน\-สปฺปฏิโฆ นเหตุ ธมฺโม อุปฺปชฺชติ เหตุปจฺจยา.  นสนิทสฺสน\-สปฺปฏิฆํ นเหตุํ ธมฺมํ ปฏิจฺจ นอนิทสฺสน\-อปฺปฏิโฆ นเหตุ ธมฺโม อุปฺปชฺชติ เหตุปจฺจยา.  นสนิทสฺสน\-สปฺปฏิฆํ นเหตุํ ธมฺมํ ปฏิจฺจ นสนิทสฺสน\-สปฺปฏิโฆ นเหตุ จ นอนิทสฺสน\-อปฺปฏิโฆ นเหตุ จ ธมฺมา อุปฺปชฺชนฺติ เหตุปจฺจยา.  นสนิทสฺสน\-สปฺปฏิฆํ นเหตุํ ธมฺมํ ปฏิจฺจ นอนิทสฺสน\-สปฺปฏิโฆ นเหตุ จ นอนิทสฺสน\-อปฺปฏิโฆ นเหตุ จ ธมฺมา อุปฺปชฺชนฺติ เหตุปจฺจยา.  นสนิทสฺสน\-สปฺปฏิฆํ นเหตุํ ธมฺมํ ปฏิจฺจ นสนิทสฺสน\-สปฺปฏิโฆ นเหตุ จ นอนิทสฺสน\-สปฺปฏิโฆ นเหตุ จ ธมฺมา อุปฺปชฺชนฺติ เหตุปจฺจยา. (6)}

\prose{นอนิทสฺสน\-สปฺปฏิฆํ นเหตุํ ธมฺมํ ปฏิจฺจ นอนิทสฺสน\-สปฺปฏิโฆ นเหตุ ธมฺโม อุปฺปชฺชติ เหตุปจฺจยา.  นอนิทสฺสน\-สปฺปฏิฆํ นเหตุํ ธมฺมํ ปฏิจฺจ \csromanfolio{134}นสนิทสฺสน\-สปฺปฏิโฆ นเหตุ ธมฺโม อุปฺปชฺชติ เหตุปจฺจยา.  นอนิทสฺสน\-สปฺปฏิฆํ นเหตุํ ธมฺมํ ปฏิจฺจ นอนิทสฺสน\-อปฺปฏิโฆ นเหตุ ธมฺโม อุปฺปชฺชติ เหตุปจฺจยา.  นอนิทสฺสน\-สปฺปฏิฆํ นเหตุํ ธมฺมํ ปฏิจฺจ นสนิทสฺสน\-สปฺปฏิโฆ นเหตุ จ นอนิทสฺสน\-อปฺปฏิโฆ นเหตุ จ ธมฺมา อุปฺปชฺชนฺติ เหตุปจฺจยา.  นอนิทสฺสน\-สปฺปฏิฆํ นเหตุํ ธมฺมํ ปฏิจฺจ นอนิทสฺสน\-สปฺปฏิโฆ นเหตุ จ นอนิทสฺสน\-อปฺปฏิโฆ นเหตุ จ ธมฺมา อุปฺปชฺชนฺติ เหตุปจฺจยา.  นอนิทสฺสน\-อปฺปฏิฆํ นเหตุํ ธมฺมํ ปฏิจฺจ นสนิทสฺสน\-สปฺปฏิโฆ นเหตุ จ นอนิทสฺสน\-สปฺปฏิโฆ นเหตุ จ ธมฺมา อุปฺปชฺชนฺติ เหตุปจฺจยา. (6) นอนิทสฺสน\-อปฺปฏิฆํ นเหตุํ ธมฺมํ ปฏิจฺจ นอนิทสฺสน\-อปฺปฏิโฆ นเหตุ ธมฺโม อุปฺปชฺชติ เหตุปจฺจยา. ฉ.}

\prose{นสนิทสฺสน\-สปฺปฏิฆํ นเหตุญฺจ นอนิทสฺสน\-อปฺปฏิฆํ นเหตุญฺจ ธมฺมํ ปฏิจฺจ นสนิทสฺสน\-สปฺปฏิโฆ นเหตุ ธมฺโม อุปฺปชฺชติ เหตุปจฺจยา. ฉ.}

\prose{นสนิทสฺสน\-สปฺปฏิฆํ นเหตุญฺจ นอนิทสฺสน\-สปฺปฏิฆํ นเหตุญฺจ ธมฺมํ ปฏิจฺจ นสนิทสฺสน\-สปฺปฏิโฆ นเหตุ ธมฺโม อุปฺปชฺชติ เหตุปจฺจยา. ฉ. (สํขิตฺตํ.)}

\prose{เหตุยา ติํส, อารมฺมเณ นว. (สพฺพตฺถ วิตฺถาโร.)}

\proseitem{70}{นสนิทสฺสน\-สปฺปฏิฆํ นนเหตุํ ธมฺมํ ปฏิจฺจ นสนิทสฺสน\-สปฺปฏิโฆ นนเหตุ ธมฺโม อุปฺปชฺชติ เหตุปจฺจยา. (นว.)}

\prose{นสนิทสฺสน\-สปฺปฏิฆํ นสเหตุกํ ธมฺมํ ปฏิจฺจ.}

\prose{นสนิทสฺสน\-สปฺปฏิฆํ นอเหตุกํ ธมฺมํ ปฏิจฺจ.}

\proseitem{71}{นสนิทสฺสน\-สปฺปฏิฆํ นเหตุ\-สมฺปยุตฺตํ ธมฺมํ ปฏิจฺจ.}

\prose{นสนิทสฺสน\-สปฺปฏิฆํ นเหตุ\-วิปฺปยุตฺตํ ธมฺมํ ปฏิจฺจ.}

\proseitem{72}{นสนิทสฺสน\-สปฺปฏิฆํ นเหตุญฺเจว นอเหตุกญฺจ ธมฺมํ ปฏิจฺจ ฯเปฯ. นสนิทสฺสน\-สปฺปฏิฆํ นอเหตุกญฺเจว นนเหตุญฺจ ธมฺมํ ปฏิจฺจ ฯเปฯ. นสนิทสฺสน\-สปฺปฏิฆํ นเหตุญฺเจว นเหตุ\-วิปฺปยุตฺตญฺจ ธมฺมํ ปฏิจฺจ.}

\prose{นสนิทสฺสน\-สปฺปฏิฆํ นเหตุ\-วิปฺปยุตฺต\-ญฺเจว นนเหตุญฺจ ธมฺมํ ปฏิจฺจ.}

\csromanmarkh{22. สนิทสฺสน\-ตฺติก}\csromanmarkhii{55. มหนฺตรทุก}\csromanheadernomark{2}{22. สนิทสฺสน\-ตฺติก 55. มหนฺตรทุก}
\prose{\csromanfolio{135}นสนิทสฺสน\-สปฺปฏิฆํ นเหตุํ นสเหตุกํ ธมฺมํ ปฏิจฺจ.}

\prose{นสนิทสฺสน\-สปฺปฏิฆํ นเหตุํ นอเหตุกํ ธมฺมํ ปฏิจฺจ.}

\csromanmarkh{22. สนิทสฺสน\-ตฺติก}\csromanmarkhii{7. จูฬนฺตรทุก}\csromanheadernomark{2}{22. สนิทสฺสน\-ตฺติก 7. จูฬนฺตรทุก}
\proseitem{73}{นสนิทสฺสน\-สปฺปฏิฆํ นอปฺปจฺจยํ ธมฺมํ ปฏิจฺจ ฯเปฯ. นสนิทสฺสน\-สปฺปฏิฆํ นอสงฺขตํ ธมฺมํ ปฏิจฺจ.}

\proseitem{74}{นสนิทสฺสน\-สปฺปฏิฆํ นสนิทสฺสนํ ธมฺมํ ปฏิจฺจ.}

\proseitem{75}{นสนิทสฺสน\-สปฺปฏิฆํ นสปฺปฏิฆํ ธมฺมํ ปฏิจฺจ.}

\prose{นสนิทสฺสน\-สปฺปฏิฆํ นอปฺปฏิฆํ ธมฺมํ ปฏิจฺจ.}

\proseitem{76}{นสนิทสฺสน\-สปฺปฏิฆํ นรูปิํ ธมฺมํ ปฏิจฺจ. นสนิทสฺสน\-สปฺปฏิฆํ นอรูปิํ ธมฺมํ ปฏิจฺจ.}

\proseitem{77}{นสนิทสนสปฺปฏิฆํ นโลกิยํ ธมฺมํ ปฏิจฺจ.}

\prose{นสนิทสฺสน\-สปฺปฏิฆํ นโลกุตฺตรํ ธมฺมํ ปฏิจฺจ.}

\proseitem{78}{นสนิทสฺสน\-สปฺปฏิฆํ นเกนจิ วิญฺเญยฺยํ ธมฺมํ ปฏิจฺจ ฯเปฯ. นสนิทสฺสน\-สปฺปฏิฆํ นนเกนจิ วิญฺเญยฺยํ ธมฺมํ ปฏิจฺจ.}

\csromanmarkh{22. สนิทสฺสน\-ตฺติก}\csromanmarkhii{14. อาสว\-โคจฺฉกาทิ}\csromanheadernomark{2}{22. สนิทสฺสน\-ตฺติก 14. อาสว\-โคจฺฉกาทิ}
\proseitem{79}{นสนิทสฺสน\-สปฺปฏิฆํ โนอาสวํ ธมฺมํ ปฏิจฺจ.}

\proseitem{80}{นสนิทสฺสน\-สปฺปฏิฆํ โนสญฺโญชนํ ธมฺมํ ปฏิจฺจ ฯเปฯ. นสนิทสฺสน\-สปฺปฏิฆํ โนคนฺถํ ธมฺมํ ปฏิจฺจ ฯเปฯ. นสนิทสฺสน\-สปฺปฏิฆํ โนโอฆํ ธมฺมํ ปฏิจฺจ ฯเปฯ. นสนิทสฺสน\-สปฺปฏิฆํ โนโยคํ ธมฺมํ ปฏิจฺจ ฯเปฯ. นสนิทสฺสน\-สปฺปฏิฆํ โนนีวรณํ ธมฺมํ ปฏิจฺจ ฯเปฯ. นสนิทสฺสน\-สปฺปฏิฆํ โนปรามาสํ ธมฺมํ ปฏิจฺจ.}

\csromanmarkh{22. สนิทสฺสน\-ตฺติก}\csromanmarkhii{55. มหนฺตรทุก}\csromanheadernomark{2}{22. สนิทสฺสน\-ตฺติก 55. มหนฺตรทุก}
\proseitem{81}{นสนิทสฺสน\-สปฺปฏิฆํ นสารมฺมณํ ธมฺมํ ปฏิจฺจ.}

\proseitem{82}{\csromanfolio{136}นสนิทสฺสน\-สปฺปฏิฆํ นจิตฺตํ ธมฺมํ ปฏิจฺจ.}

\prose{(นโนจิตฺตปทํ น ลพฺภติ.)}

\proseitem{83}{นสนิทสฺสน\-สปฺปฏิฆํ นเจตสิกํ ธมฺมํ ปฏิจฺจ.}

\proseitem{84}{นสนิทสฺสน\-สปฺปฏิฆํ นจิตฺต\-สมฺปยุตฺตํ ธมฺมํ ปฏิจฺจ ฯเปฯ. นสนิทสฺสน\-สปฺปฏิฆํ นจิตฺต\-สํสฏฺฐํ ธมฺมํ ปฏิจฺจ.}

\proseitem{85}{นสนิทสฺสน\-สปฺปฏิฆํ นจิตฺต\-สมุฏฺฐานํ ธมฺมํ ปฏิจฺจ ฯเปฯ. นสนิทสฺสน\-สปฺปฏิฆํ นจิตฺต\-สหภุํ ธมฺมํ ปฏิจฺจ ฯเปฯ. นสนิทสฺสน\-สปฺปฏิฆํ นจิตฺตา\-นุปริวตฺติํ ธมฺมํ ปฏิจฺจ ฯเปฯ. นสนิทสฺสน\-สปฺปฏิฆํ นจิตฺต\-สํสฏฺฐ\-สมุฏฺฐานํ ธมฺมํ ปฏิจฺจ ฯเปฯ. นสนิทสฺสน\-สปฺปฏิฆํ นจิตฺต\-สํสฏฺฐ\-สมุฏฺฐาน\-สหภุํ ธมฺมํ ปฏิจฺจ ฯเปฯ. นสนิทสฺสน\-สปฺปฏิฆํ นจิตฺต\-สํสฏฺฐ\-สมุฏฺฐานา\-นุ\-ปริวตฺติํ ธมฺมํ ปฏิจฺจ.}

\proseitem{86}{นสนิทสฺสน\-สปฺปฏิฆํ นอชฺฌตฺติกํ ธมฺมํ ปฏิจฺจ.}

\prose{นสนิทสฺสน\-สปฺปฏิฆํ นพาหิรํ ธมฺมํ ปฏิจฺจ.}

\proseitem{87}{นสนิทสฺสน\-สปฺปฏิฆํ นอุปาทา ธมฺมํ ปฏิจฺจ.}

\proseitem{88}{นสนิทสฺสน\-สปฺปฏิฆํ นอุปาทินฺนํ ธมฺมํ ปฏิจฺจ ฯเปฯ. นสนิทสฺสน\-สปฺปฏิฆํ นอนุปาทินฺนํ ธมฺมํ ปฏิจฺจ.}

\csromanmarkh{22. สนิทสฺสนตฺตฺอิก}\csromanmarkhii{69. อุปาทานาทิทุก}\csromanheadernomark{2}{22. \textbf{สนิทสฺสนตฺตฺ}อิก 69. อุปาทานาทิทุก}
\proseitem{89}{นสนิทสฺสน\-สปฺปฏิฆํ โนอุปาทานํ ธมฺมํ ปฏิจฺจ ฯเปฯ. นสนิทสฺสน\-สปฺปฏิฆํ โนกิเลสํ ธมฺมํ ปฏิจฺจ.}

\csromanmarkh{22. สนิทสฺสน\-ตฺติก}\csromanmarkhii{83. ปิฏฺฐิทุก}\csromanheadernomark{2}{22. สนิทสฺสน\-ตฺติก 83. ปิฏฺฐิทุก}
\proseitem{90}{นสนิทสฺสน\-สปฺปฏิฆํ นทสฺสเนน ปหาตพฺพํ ธมฺมํ ปฏิจฺจ.}

\prose{เหตุยา ติํส, อารมฺมเณ นว ฯเปฯ อญฺญมญฺเญ ปญฺจวีส ฯเปฯ อวิคเต ติํส. (ปิฏฺฐิทุกํ วิตฺถาเรตพฺพํ.)}

\proseitem{91}{นสนิทสฺสน\-สปฺปฏิฆํ นสรณํ ธมฺมํ ปฏิจฺจ นสนิทสฺสน\-สปฺปฏิโฆ นสรโณ ธมฺโม อุปฺปชฺชติ เหตุปจฺจยา.  นสนิทสฺสน\-สปฺปฏิฆํ นสรณํ}

\vspace{\tipitakaparskip}
\csromancenter{สนิทสฺสน\-ตฺติก 83. ปิฏฺฐิทุก ธมฺมํ ปฏิจฺจ นอนิทสฺสน\-สปฺปฏิโฆ นสรโณ ธมฺโม อุปฺปชฺชติ เหตุปจฺจยา.  นสนิทสฺสน\-สปฺปฏิฆํ นสรณํ ธมฺมํ ปฏิจฺจ นอนิทสฺสน\-อปฺปฏิโฆ นสรโณ ธมฺโม อุปฺปชฺชติ เหตุปจฺจยา ฯเปฯ. ฉ.}
\prose{\csromanfolio{137}นอนิทสฺสน\-สปฺปฏิฆํ นสรณํ ธมฺมํ ปฏิจฺจ นอนิทสฺสน\-สปฺปฏิโฆ นสรโณ ธมฺโม อุปฺปชฺชติ เหตุปจฺจยา.  นอนิทสฺสน\-สปฺปฏิฆํ นสรณํ ธมฺมํ ปฏิจฺจ นสนิทสฺสน\-สปฺปฏิโฆ นสรโณ ธมฺโม อุปฺปชฺชติ เหตุปจฺจยา ฯเปฯ. ฉ.}

\prose{นอนิทสฺสน\-สปฺปฏิฆํ นสรณํ ธมฺมํ ปฏิจฺจ นอนิทสฺสน\-อปฺปฏิโฆ นสรโณ ธมฺโม อุปฺปชฺชติ เหตุปจฺจยา. ฉ.}

\prose{นสนิทสฺสน\-สปฺปฏิฆํ นสรณญฺจ นอนิทสฺสน\-อปฺปฏิฆํ นสรณญฺจ ธมฺมํ ปฏิจฺจ นสนิทสฺสน\-สปฺปฏิโฆ นสรโณ ธมฺโม อุปฺปชฺชติ เหตุปจฺจยา. ฉ.}

\prose{นสนิทสฺสน\-สปฺปฏิฆํ นสรณญฺจ นอนิทสฺสน\-สปฺปฏิฆํ นสรณญฺจ ธมฺมํ ปฏิจฺจ นสนิทสฺสน\-สปฺปฏิโฆ นสรโณ ธมฺโม อุปฺปชฺชติ เหตุปจฺจยา. ฉ.}

\prose{เหตุยา ติํส, อารมฺมเณ นว ฯเปฯ อวิคเต ติํส. (สพฺพตฺถ วิตฺถาโร.)}

\proseitem{92}{นสนิทสฺสน\-สปฺปฏิฆํ นอรณํ ธมฺมํ ปฏิจฺจ นสนิทสฺสน\-สปฺปฏิโฆ นอรโณ ธมฺโม อุปฺปชฺชติ เหตุปจฺจยา.  นสนิทสฺสน\-สปฺปฏิฆํ นอรณํ ธมฺมํ ปฏิจฺจ นอนิทสฺสน\-สปฺปฏิโฆ นอรโณ ธมฺโม อุปฺปชฺชติ เหตุปจฺจยา.  นสนิทสฺสน\-สปฺปฏิฆํ นอรณํ ธมฺมํ ปฏิจฺจ นสนิทสฺสน\-สปฺปฏิโฆ นอรโณ จ นอนิทสฺสน\-สปฺปฏิโฆ นอรโณ จ ธมฺมา อุปฺปชฺชนฺติ เหตุปจฺจยา. (3)}

\proseitem{93}{นอนิทสฺสน\-สปฺปฏิฆํ นอรณํ ธมฺมํ ปฏิจฺจ นอนิทสฺสน\-สปฺปฏิโฆ นอรโณ ธมฺโม อุปฺปชฺชติ เหตุปจฺจยา.  นอนิทสฺสน\-สปฺปฏิฆํ นอรณํ ธมฺมํ ปฏิจฺจ นสนิทสฺสน\-สปฺปฏิโฆ นอรโณ ธมฺโม อุปฺปชฺชติ เหตุปจฺจยา.  นอนิทสฺสน\-สปฺปฏิฆํ นอรณํ ธมฺมํ ปฏิจฺจ นสนิทสฺสน\-สปฺปฏิโฆ นอรโณ จ นอนิทสฺสน\-สปฺปฏิโฆ นอรโณ จ ธมฺมา อุปฺปชฺชนฺติ เหตุปจฺจยา. (3)}

\proseitem{94}{\csromanfolio{138}นสนิทสฺสน\-สปฺปฏิฆํ นอรณญฺจ นอนิทสฺสน\-สปฺปฏิฆํ นอรณญฺจ ธมฺมํ ปฏิจฺจ นสนิทสฺสน\-สปฺปฏิโฆ นอรโณ ธมฺโม อุปฺปชฺชติ เหตุปจฺจยา.  นสนิทสฺสน\-สปฺปฏิฆํ นอรณญฺจ นอนิทสฺสน\-สปฺปฏิฆํ นอรณญฺจ ธมฺมํ ปฏิจฺจ นอนิทสฺสน\-สปฺปฏิโฆ นอรโณ ธมฺโม อุปฺปชฺชติ เหตุปจฺจยา.  นสนิทสฺสน\-สปฺปฏิฆํ นอรณญฺจ นอนิทสฺสน\-สปฺปฏิฆํ นอรณญฺจ ธมฺมํ ปฏิจฺจ นสนิทสฺสน\-สปฺปฏิโฆ นอรโณ จ นอนิทสฺสน\-สปฺปฏิโฆ นอรโณ จ ธมฺมา อุปฺปชฺชนฺติ เหตุปจฺจยา. (3)}

\prose{เหตุยา นว, อารมฺมเณ นว ฯเปฯ อวิคเต นว.}

\vspace{\tipitakaparskip}
\csromancenter{(สหชาตวาเรปิ ปจฺจยวาเรปิ นิสฺสยวาเรปิ สํสฏฺฐ\-วาเรปิ สมฺปยุตฺตวาเรปิ สพฺพตฺถ นว.)}
\nitthitam{ธมฺม\-ปจฺจนีเย ติกทุก\-ปฏฺฐานํ นิฏฺฐิตํ.}
\pitaka{อภิธมฺม\-ปิฏก}
\namakkaram{นโม ตสฺส ภควโต อรหโต สมฺมา\-สมฺพุทฺธสฺส.}
\csromanmarkh{1. กุสลตฺติก}\csromanmarkhii{1. เวทนาตฺติก}\csromanheadernomark{2}{1. \textbf{กุสลตฺติก 1. เวทนาตฺติก}}
\prose{\csromanfolio{139}1. ปฏิจฺจวารเหตุ นกุสลํ นสุขาย เวทนาย สมฺปยุตฺตํ ธมฺมํ ปฏิจฺจ นกุสโล นสุขาย เวทนาย สมฺปยุตฺโต ธมฺโม อุปฺปชฺชติ เหตุปจฺจยา.  นกุสลํ นสุขาย เวทนาย สมฺปยุตฺตํ ธมฺมํ ปฏิจฺจ นอกุสโล นสุขาย เวทนาย สมฺปยุตฺโต ธมฺโม อุปฺปชฺชติ เหตุปจฺจยา.  นกุสลํ นสุขาย เวทนาย สมฺปยุตฺตํ ธมฺมํ ปฏิจฺจ นอพฺยากโต นสุขาย เวทนาย สมฺปยุตฺโต ธมฺโม อุปฺปชฺชติ เหตุปจฺจยา.  นกุสลํ นสุขาย เวทนาย สมฺปยุตฺตํ ธมฺมํ ปฏิจฺจ นกุสโล นสุขาย เวทนาย สมฺปยุตฺโต จ นอพฺยากโต นสุขาย เวทนาย สมฺปยุตฺโต จ ธมฺมา อุปฺปชฺชนฺติ เหตุปจฺจยา.  นกุสลํ นสุขาย เวทนาย สมฺปยุตฺตํ ธมฺมํ ปฏิจฺจ นกุสโล นสุขาย เวทนาย สมฺปยุตฺโต จ นอกุสโล นสุขาย เวทนาย สมฺปยุตฺโต จ ธมฺมา อุปฺปชฺชนฺติ เหตุปจฺจยา. (5)}

\proseitem{2}{นอกุสลํ นสุขาย เวทนาย สมฺปยุตฺตํ ธมฺมํ ปฏิจฺจ นอกุสโล นสุขาย เวทนาย สมฺปยุตฺโต ธมฺโม อุปฺปชฺชติ เหตุปจฺจยา.  นอกุสลํ นสุขาย เวทนาย สมฺปยุตฺตํ ธมฺมํ ปฏิจฺจ นกุสโล นสุขาย เวทนาย สมฺปยุตฺโต ธมฺโม อุปฺปชฺชติ เหตุปจฺจยา.  นอกุสลํ นสุขาย เวทนาย สมฺปยุตฺตํ ธมฺมํ ปฏิจฺจ นอพฺยากโต นสุขาย เวทนาย \csromanfolio{140}สมฺปยุตฺโต ธมฺโม อุปฺปชฺชติ เหตุปจฺจยา.  นอกุสลํ นสุขาย เวทนาย สมฺปยุตฺตํ ธมฺมํ ปฏิจฺจ นอกุสโล นสุขาย เวทนาย สมฺปยุตฺโต จ นอพฺยากโต นสุขาย เวทนาย สมฺปยุตฺโต จ ธมฺมา อุปฺปชฺชนฺติ เหตุปจฺจยา.  นอกุสลํ นสุขาย เวทนาย สมฺปยุตฺตํ ธมฺมํ ปฏิจฺจ นกุสโล นสุขาย เวทนาย สมฺปยุตฺโต จ นอกุสโล นสุขาย เวทนาย สมฺปยุตฺโต จ ธมฺมา อุปฺปชฺชนฺติ เหตุปจฺจยา. (5)}

\prose{นอพฺยากตํ นสุขาย เวทนาย สมฺปยุตฺตํ ธมฺมํ ปฏิจฺจ นอพฺยากโต นสุขาย เวทนาย สมฺปยุตฺโต ธมฺโม อุปฺปชฺชติ เหตุปจฺจยา.  นอพฺยากตํ นสุขาย เวทนาย สมฺปยุตฺตํ ธมฺมํ ปฏิจฺจ นกุสโล นสุขาย เวทนาย สมฺปยุตฺโต ธมฺโม อุปฺปชฺชติ เหตุปจฺจยา.  นอพฺยากตํ นสุขาย เวทนาย สมฺปยุตฺตํ ธมฺมํ ปฏิจฺจ นอกุสโล นสุขาย เวทนาย สมฺปยุตฺโต ธมฺโม อุปฺปชฺชติ เหตุปจฺจยา.  นอพฺยากตํ นสุขาย เวทนาย สมฺปยุตฺตํ ธมฺมํ ปฏิจฺจ นกุสโล นสุขาย เวทนาย สมฺปยุตฺโต จ นอพฺยากโต นสุขาย เวทนาย สมฺปยุตฺโต จ ธมฺมา อุปฺปชฺชนฺติ เหตุปจฺจยา.  นอพฺยากตํ นสุขาย เวทนาย สมฺปยุตฺตํ ธมฺมํ ปฏิจฺจ นอกุสโล นสุขาย เวทนาย สมฺปยุตฺโต จ นอพฺยากโต นสุขาย เวทนาย สมฺปยุตฺโต จ ธมฺมา อุปฺปชฺชนฺติ เหตุปจฺจยา.  นอพฺยากตํ นสุขาย เวทนาย สมฺปยุตฺตํ ธมฺมํ ปฏิจฺจ นกุสโล นสุขาย เวทนาย สมฺปยุตฺโต จ นอกุสโล นสุขาย เวทนาย สมฺปยุตฺโต จ ธมฺมา อุปฺปชฺชนฺติ เหตุปจฺจยา. (6)}

\proseitem{4}{นกุสลํ นสุขาย เวทนาย สมฺปยุตฺตญฺจ นอพฺยากตํ นสุขาย เวทนาย สมฺปยุตฺตญฺจ ธมฺมํ ปฏิจฺจ นกุสโล นสุขาย เวทนาย สมฺปยุตฺโต ธมฺโม อุปฺปชฺชติ เหตุปจฺจยา.  นกุสลํ นสุขาย เวทนาย สมฺปยุตฺตญฺจ นอพฺยากตํ นสุขาย เวทนาย สมฺปยุตฺตญฺจ ธมฺมํ ปฏิจฺจ นอกุสโล นสุขาย เวทนาย สมฺปยุตฺโต ธมฺโม อุปฺปชฺชติ เหตุปจฺจยา.  นกุสลํ นสุขาย เวทนาย สมฺปยุตฺตญฺจ นอพฺยากตํ นสุขาย เวทนาย สมฺปยุตฺตญฺจ ธมฺมํ ปฏิจฺจ นอพฺยากโต นสุขาย เวทนาย สมฺปยุตฺโต ธมฺโม อุปฺปชฺชติ เหตุปจฺจยา.  นกุสลํ นสุขาย เวทนาย สมฺปยุตฺตญฺจ นอพฺยากตํ นสุขาย เวทนาย สมฺปยุตฺตญฺจ ธมฺมํ ปฏิจฺจ นกุสโล นสุขาย เวทนาย สมฺปยุตฺโต จ นอพฺยากโต นสุขาย เวทนาย สมฺปยุตฺโต จ ธมฺมา อุปฺปชฺชนฺติ}

\vspace{\tipitakaparskip}
\csromancenter{1. เวทนาตฺติก เหตุปจฺจยา.  นกุสลํ นสุขาย เวทนาย สมฺปยุตฺตญฺจ นอพฺยากตํ นสุขาย เวทนาย สมฺปยุตฺตญฺจ ธมฺมํ ปฏิจฺจ นกุสโล นสุขาย เวทนาย สมฺปยุตฺโต จ นอกุสโล นสุขาย เวทนาย สมฺปยุตฺโต จ ธมฺมา อุปฺปชฺชนฺติ เหตุปจฺจยา. (5)}
\proseitem{5}{\csromanfolio{141}นอกุสลํ นสุขาย เวทนาย สมฺปยุตฺตญฺจ นอพฺยากตํ นสุขาย เวทนาย สมฺปยุตฺตญฺจ ธมฺมํ ปฏิจฺจ นกุสโล นสุขาย เวทนาย สมฺปยุตฺโต ธมฺโม อุปฺปชฺชติ เหตุปจฺจยา.  นอกุสลํ นสุขาย เวทนาย สมฺปยุตฺตญฺจ นอพฺยากตํ นสุขาย เวทนาย สมฺปยุตฺตญฺจ ธมฺมํ ปฏิจฺจ นอกุสโล นสุขาย เวทนาย สมฺปยุตฺโต ธมฺโม อุปฺปชฺชติ เหตุปจฺจยา.  นอกุสลํ นสุขาย เวทนาย สมฺปยุตฺตญฺจ นอพฺยากตํ นสุขาย เวทนาย สมฺปยุตฺตญฺจ ธมฺมํ ปฏิจฺจ นอพฺยากโต นสุขาย เวทนาย สมฺปยุตฺโต ธมฺโม อุปฺปชฺชติ เหตุปจฺจยา.  นอกุสลํ นสุขาย เวทนาย สมฺปยุตฺตญฺจ นอพฺยากตํ นสุขาย เวทนาย สมฺปยุตฺตญฺจ ธมฺมํ ปฏิจฺจ นอกุสโล นสุขาย เวทนาย สมฺปยุตฺโต จ นอพฺยากโต นสุขาย เวทนาย สมฺปยุตฺโต จ ธมฺมา อุปฺปชฺชนฺติ เหตุปจฺจยา.  นอกุสลํ นสุขาย เวทนาย สมฺปยุตฺตญฺจ นอพฺยากตํ นสุขาย เวทนาย สมฺปยุตฺตญฺจ ธมฺมํ ปฏิจฺจ นกุสโล นสุขาย เวทนาย สมฺปยุตฺโต จ นอกุสโล นสุขาย เวทนาย สมฺปยุตฺโต จ ธมฺมา อุปฺปชฺชนฺติ เหตุปจฺจยา. (5)}

\proseitem{6}{นกุสลํ นสุขาย เวทนาย สมฺปยุตฺตญฺจ นอกุสลํ นสุขาย เวทนาย สมฺปยุตฺตญฺจ ธมฺมํ ปฏิจฺจ นกุสโล นสุขาย เวทนาย สมฺปยุตฺโต ธมฺโม อุปฺปชฺชติ เหตุปจฺจยา.  นกุสลํ นสุขาย เวทนาย สมฺปยุตฺตญฺจ นอกุสลํ นสุขาย เวทนาย สมฺปยุตฺตญฺจ ธมฺมํ ปฏิจฺจ นอกุสโล นสุขาย เวทนาย สมฺปยุตฺโต ธมฺโม อุปฺปชฺชติ เหตุปจฺจยา.  นกุสลํ นสุขาย เวทนาย สมฺปยุตฺตญฺจ นอกุสลํ นสุขาย เวทนาย สมฺปยุตฺตญฺจ ธมฺมํ ปฏิจฺจ นกุสโล นสุขาย เวทนาย สมฺปยุตฺโต จ นอกุสโล นสุขาย เวทนาย สมฺปยุตฺโต จ ธมฺมา อุปฺปชฺชนฺติ เหตุปจฺจยา. (3)}

\prose{เหตุยา เอกูนติํส, อารมฺมเณ จตุวีส ฯเปฯ อวิคเต เอกูนติํส. (สหชาตวาเรปิ ฯเปฯ ปญฺหาวาเรปิ วิตฺถาโร.)}

\vspace{\tipitakaparskip}
\csromancenter{(สหชาตวาเรปิ ฯเปฯ ปญฺหาวาเรปิ วิตฺถาโร.)}
\csromanpalirecto
\gambhira{ธมฺม\-ปจฺจนีย ติกติก\-ปฏฺฐาน\-ปาฬิ ทุกฺขายเวทนายสมฺปยุตฺตปทเหตุ}
\prosecont{\csromanfolio{142}นกุสลํ นทุกฺขาย เวทนาย สมฺปยุตฺตํ ธมฺมํ ปฏิจฺจ นกุสโล นทุกฺขาย เวทนาย สมฺปยุตฺโต ธมฺโม อุปฺปชฺชติ เหตุปจฺจยา.  นกุสลํ นทุกฺขาย เวทนาย สมฺปยุตฺตํ ธมฺมํ ปฏิจฺจ นอกุสโล นทุกฺขาย เวทนาย สมฺปยุตฺโต ธมฺโม อุปฺปชฺชติ เหตุปจฺจยา.  นกุสลํ นทุกฺขาย เวทนาย สมฺปยุตฺตํ ธมฺมํ ปฏิจฺจ นอพฺยากโต นทุกฺขาย เวทนาย สมฺปยุตฺโต ธมฺโม อุปฺปชฺชติ เหตุปจฺจยา.  นกุสลํ นทุกฺขาย เวทนาย สมฺปยุตฺตํ ธมฺมํ ปฏิจฺจ นกุสโล นทุกฺขาย เวทนาย สมฺปยุตฺโต จ นอพฺยากโต นทุกฺขาย เวทนาย สมฺปยุตฺโต จ ธมฺมา อุปฺปชฺชนฺติ เหตุปจฺจยา.  นกุสลํ นทุกฺขาย เวทนาย สมฺปยุตฺตํ ธมฺมํ ปฏิจฺจ นกุสโล นทุกฺขาย เวทนาย สมฺปยุตฺโต จ นอกุสโล นทุกฺขาย เวทนาย สมฺปยุตฺโต จ ธมฺมา อุปฺปชฺชนฺติ เหตุปจฺจยา. (5)}

\proseitem{7}{นอกุสลํ นทุกฺขาย เวทนาย สมฺปยุตฺตํ ธมฺมํ ปฏิจฺจ นอกุสโล นทุกฺขาย เวทนาย สมฺปยุตฺโต ธมฺโม อุปฺปชฺชติ เหตุปจฺจยา. ปญฺจ.}

\prose{นอพฺยากตํ นทุกฺขาย เวทนาย สมฺปยุตฺตํ ธมฺมํ ปฏิจฺจ นอพฺยากโต นทุกฺขาย เวทนาย สมฺปยุตฺโต ธมฺโม อุปฺปชฺชติ เหตุปจฺจยา. ฉ.}

\prose{นกุสลํ นทุกฺขาย เวทนาย สมฺปยุตฺตญฺจ นอพฺยากตํ นทุกฺขาย เวทนาย สมฺปยุตฺตญฺจ ธมฺมํ ปฏิจฺจ นกุสโล นทุกฺขาย เวทนาย สมฺปยุตฺโต ธมฺโม อุปฺปชฺชติ เหตุปจฺจยา. ปญฺจ.}

\prose{นอกุสลํ นทุกฺขาย เวทนาย สมฺปยุตฺตญฺจ นอพฺยากตํ นทุกฺขาย เวทนาย สมฺปยุตฺตญฺจ ธมฺมํ ปฏิจฺจ นกุสโล นทุกฺขาย เวทนาย สมฺปยุตฺโต ธมฺโม อุปฺปชฺชติ เหตุปจฺจยา. ปญฺจ.}

\prose{นกุสลํ นทุกฺขาย เวทนาย สมฺปยุตฺตญฺจ นอกุสลํ นทุกฺขาย เวทนาย สมฺปยุตฺตญฺจ ธมฺมํ ปฏิจฺจ นกุสโล นทุกฺขาย เวทนาย สมฺปยุตฺโต ธมฺโม อุปฺปชฺชติ เหตุปจฺจยา. ตีณิ.}

\prose{เหตุยา เอกูนติํส ฯเปฯ อวิคเต เอกูนติํส. (สพฺพตฺถ วิตฺถาเรตพฺพํ.)}

\csromanmarkh{1. กุสลตฺติก}\csromanmarkhii{2. วิปากตฺติก ตติยปทเหตุ}\csromanheadernomark{2}{1. กุสลตฺติก 2. วิปากตฺติก ตติยปทเหตุ}
\proseitem{8}{\csromanfolio{143}นกุสลํ นอทุกฺขม\-สุขาย เวทนาย สมฺปยุตฺตํ ธมฺมํ ปฏิจฺจ นกุสโล นอทุกฺขม\-สุขาย เวทนาย สมฺปยุตฺโต ธมฺโม อุปฺปชฺชติ เหตุปจฺจยา. ปญฺจ.}

\prose{นอกุสลํ นอทุกฺขม\-สุขาย เวทนาย สมฺปยุตฺตํ ธมฺมํ ปฏิจฺจ นอกุสโล นอทุกฺขม\-สุขาย เวทนาย สมฺปยุตฺโต ธมฺโม อุปฺปชฺชติ เหตุปจฺจยา. ปญฺจ.}

\prose{นอพฺยากตํ นอทุกฺขม\-สุขาย เวทนาย สมฺปยุตฺตํ ธมฺมํ ปฏิจฺจ นอพฺยากโต นอทุกฺขม\-สุขาย เวทนาย สมฺปยุตฺโต ธมฺโม อุปฺปชฺชติ เหตุปจฺจยา. ฉ.}

\prose{นกุสลํ นอทุกฺขม\-สุขาย เวทนาย สมฺปยุตฺตญฺจ นอพฺยากตํ นอทุกฺขม\-สุขาย เวทนาย สมฺปยุตฺตญฺจ ธมฺมํ ปฏิจฺจ นกุสโล นอทุกฺขม\-สุขาย เวทนาย สมฺปยุตฺโต ธมฺโม อุปฺปชฺชติ เหตุปจฺจยา. ปญฺจ.}

\prose{นอกุสลํ นอทุกฺขม\-สุขาย เวทนาย สมฺปยุตฺตญฺจ นอพฺยากตํ นอทุกฺขม\-สุขาย เวทนาย สมฺปยุตฺตญฺจ ธมฺมํ ปฏิจฺจ นกุสโล นอทุกฺขม\-สุขาย เวทนาย สมฺปยุตฺโต ธมฺโม อุปฺปชฺชติ เหตุปจฺจยา. ปญฺจ.}

\prose{นกุสลํ นอทุกฺขม\-สุขาย เวทนาย สมฺปยุตฺตญฺจ นอกุสลํ นอทุกฺขม\-สุขาย เวทนาย สมฺปยุตฺตญฺจ ธมฺมํ ปฏิจฺจ นกุสโล นอทุกฺขม\-สุขาย เวทนาย สมฺปยุตฺโต ธมฺโม อุปฺปชฺชติ เหตุปจฺจยา. ตีณิ. (สํขิตฺตํ.)}

\prose{เหตุยา เอกูนติํส ฯเปฯ อวิคเต เอกูนติํส. (สพฺพตฺถ วิตฺถาโร.)}

\csromanmarkh{1. กุสลตฺติก}\csromanmarkhii{2. วิปากตฺติก}\csromanheadernomark{2}{1. กุสลตฺติก 2. วิปากตฺติก}
\proseitem{9}{นกุสลํ นวิปากํ ธมฺมํ ปฏิจฺจ นกุสโล นวิปาโก ธมฺโม อุปฺปชฺชติ เหตุปจฺจยา.  นกุสลํ นวิปากํ ธมฺมํ ปฏิจฺจ นอกุสโล นวิปาโก ธมฺโม อุปฺปชฺชติ เหตุปจฺจยา.  นกุสลํ นวิปากํ ธมฺมํ ปฏิจฺจ นอพฺยากโต นวิปาโก ธมฺโม อุปฺปชฺชติ เหตุปจฺจยา.  นกุสลํ นวิปากํ ธมฺมํ ปฏิจฺจ นกุสโล นวิปาโก จ นอพฺยากโต นวิปาโก จ ธมฺมา อุปฺปชฺชนฺติ}

\prose{\csromanfolio{144}เหตุปจฺจยา.  นกุสลํ นวิปากํ ธมฺมํ ปฏิจฺจ นกุสโล นวิปาโก จ นอกุสโล นวิปาโก จ ธมฺมา อุปฺปชฺชนฺติ เหตุปจฺจยา. (5)}

\prose{นอกุสลํ นวิปากํ ธมฺมํ ปฏิจฺจ นอกุสโล นวิปาโก ธมฺโม อุปฺปชฺชติ เหตุปจฺจยา. ปญฺจ.}

\prose{นอพฺยากตํ นวิปากํ ธมฺมํ ปฏิจฺจ นอพฺยากโต นวิปาโก ธมฺโม อุปฺปชฺชติ เหตุปจฺจยา. ฉ.}

\prose{นกุสลํ นวิปากญฺจ นอพฺยากตํ นวิปากญฺจ ธมฺมํ ปฏิจฺจ นกุสโล นวิปาโก ธมฺโม อุปฺปชฺชติ เหตุปจฺจยา. ปญฺจ.}

\prose{นอกุสลํ นวิปากญฺจ นอพฺยากตํ นวิปากญฺจ ธมฺมํ ปฏิจฺจ นกุสโล นวิปาโก ธมฺโม อุปฺปชฺชติ เหตุปจฺจยา. ปญฺจ.}

\prose{นกุสลํ นวิปากญฺจ นอกุสลํ นวิปากญฺจ ธมฺมํ ปฏิจฺจ นกุสโล นวิปาโก ธมฺโม อุปฺปชฺชติ เหตุปจฺจยา. ตีณิ. (สํขิตฺตํ.)}

\prose{เหตุยา เอกูนติํส ฯเปฯ อวิคเต เอกูนติํส. (สพฺพตฺถ วิตฺถาโร.)}

\proseitem{10}{นกุสลํ นวิปาก\-ธมฺม\-ธมฺมํ ปฏิจฺจ นกุสโล นวิปาก\-ธมฺม\-ธมฺโม อุปฺปชฺชติ เหตุปจฺจยา.  นกุสลํ นวิปาก\-ธมฺม\-ธมฺมํ ปฏิจฺจ นอกุสโล นวิปาก\-ธมฺม\-ธมฺโม อุปฺปชฺชติ เหตุปจฺจยา.  นกุสลํ นวิปาก\-ธมฺม\-ธมฺมํ ปฏิจฺจ นกุสโล นวิปาก\-ธมฺม\-ธมฺโม จ นอกุสโล นวิปาก\-ธมฺม\-ธมฺโม จ ธมฺมา อุปฺปชฺชนฺติ เหตุปจฺจยา. (3)}

\prose{นอกุสลํ นวิปาก\-ธมฺม\-ธมฺมํ ปฏิจฺจ นอกุสโล นวิปาก\-ธมฺม\-ธมฺโม อุปฺปชฺชติ เหตุปจฺจยา.  นอกุสลํ นวิปาก\-ธมฺม\-ธมฺมํ ปฏิจฺจ นกุสโล นวิปาก\-ธมฺม\-ธมฺโม อุปฺปชฺชติ เหตุปจฺจยา.  นอกุสลํ นวิปาก\-ธมฺม\-ธมฺมํ ปฏิจฺจ นกุสโล นวิปาก\-ธมฺม\-ธมฺโม จ นอกุสโล นวิปาก\-ธมฺม\-ธมฺโม จ ธมฺมา อุปฺปชฺชนฺติ เหตุปจฺจยา. (3)}

\prose{นกุสลํ นวิปาก\-ธมฺม\-ธมฺมญฺจ นอกุสลํ นวิปาก\-ธมฺม\-ธมฺมญฺจ ธมฺมํ ปฏิจฺจ นกุสโล นวิปาก\-ธมฺม\-ธมฺโม อุปฺปชฺชติ เหตุปจฺจยา.  นกุสลํ นวิปาก\-ธมฺม\-ธมฺมญฺจ นอกุสลํ นวิปาก\-ธมฺม\-ธมฺมญฺจ ธมฺมํ ปฏิจฺจ นอกุสโล นวิปาก\-ธมฺม\-ธมฺโม อุปฺปชฺชติ เหตุปจฺจยา.  นกุสลํ นวิปาก\-ธมฺม\-ธมฺมญฺจ}

\vspace{\tipitakaparskip}
\csromancenter{6. ปีติตฺติก นอกุสลํ นวิปาก\-ธมฺม\-ธมฺมญฺจ ธมฺมํ ปฏิจฺจ นกุสโล นวิปาก\-ธมฺม\-ธมฺโม จ นอกุสโล นวิปาก\-ธมฺม\-ธมฺโม จ ธมฺมา อุปฺปชฺชนฺติ เหตุปจฺจยา. (3) (สํขิตฺตํ.)}
\prose{\csromanfolio{145}เหตุยา นว ฯเปฯ อวิคเต นว. (สพฺพตฺถ นว.)}

\proseitem{11}{นกุสลํ นเน\-ว\-วิปาก\-น\-วิปากธมฺม\-ธมฺมํ ปฏิจฺจ นกุสโล นเน\-ว\-วิปาก\-น\-วิปาก\-ธมฺม\-ธมฺโม อุปฺปชฺชติ เหตุปจฺจยา.}

\csromanmarkh{1. กุสลตฺติก}\csromanmarkhii{3. อุปาทินฺนตฺติก}\csromanheadernomark{2}{1. \textbf{กุสลตฺติก 3. อุปาทินฺนตฺติก}}
\proseitem{12}{นกุสลํ นอุปาทินฺนุปาทานิยํ ธมฺมํ ปฏิจฺจ.}

\proseitem{13}{นกุสลํ นอนุปาทินฺนุ\-ปาทานิยํ ธมฺมํ ปฏิจฺจ.}

\proseitem{14}{นกุสลํ นอนุปาทินฺน\-อนุปาทานิยํ ธมฺมํ ปฏิจฺจ.}

\csromanmarkh{1. กุสลตฺติก}\csromanmarkhii{4. สํกิลิฏฺฐ\-ตฺติก}\csromanheadernomark{2}{1. กุสลตฺติก 4. สํกิลิฏฺฐ\-ตฺติก}
\proseitem{15}{นกุสลํ นสํกิลิฏฺฐ\-สํกิเลสิกํ ธมฺมํ ปฏิจฺจ.}

\proseitem{16}{นกุสลํ นอสํกิลิฏฺฐ\-สํกิเลสิกํ ธมฺมํ ปฏิจฺจ.}

\proseitem{17}{นกุสลํ นอสํกิลิฏฺฐ\-อสํกิเลสิกํ ธมฺมํ ปฏิจฺจ.}

\csromanmarkh{1. กุสลตฺติก}\csromanmarkhii{5. วิตกฺกตฺติก}\csromanheadernomark{2}{1. \textbf{กุสลตฺติก 5. วิตกฺกตฺติก}}
\proseitem{18}{นกุสลํ นสวิตกฺก\-สวิจารํ ธมฺมํ ปฏิจฺจ.}

\proseitem{19}{นกุสลํ นอวิตกฺก\-วิจารมตฺตํ ธมฺมํ ปฏิจฺจ.}

\proseitem{20}{นกุสลํ นอวิตกฺก\-อวิจารํ ธมฺมํ ปฏิจฺจ.}

\csromanmarkh{1. กุสลตฺติก}\csromanmarkhii{6. ปีติตฺติก}\csromanheadernomark{2}{1. กุสลตฺติก 6. ปีติตฺติก}
\proseitem{21}{นกุสลํ นปีติสหคตํ ธมฺมํ ปฏิจฺจ ฯเปฯ. นกุสลํ นสุขสหคตํ ธมฺมํ ปฏิจฺจ ฯเปฯ. นกุสลํ นอุเปกฺขาสหคตํ ธมฺมํ ปฏิจฺจ.}

\csromanmarkh{ธมฺม\-ปจฺจนีย ติกติก\-ปฏฺฐาน\-ปาฬิ}\csromanmarkhii{1. กุสลตฺติก 7-12. ทสฺสนตฺติกาทิ}\csromanheadernomark{2}{ธมฺม\-ปจฺจนีย ติกติก\-ปฏฺฐาน\-ปาฬิ 1. กุสลตฺติก 7-12. ทสฺสนตฺติกาทิ}
\proseitem{22}{\csromanfolio{146}นกุสลํ นทสฺสเนน ปหาตพฺพํ ธมฺมํ ปฏิจฺจ ฯเปฯ. นกุสลํ นภาวนาย ปหาตพฺพํ ธมฺมํ ปฏิจฺจ.}

\prose{นกุสลํ นเน\-ว\-ทสฺสเนน นภาวนาย ปหาตพฺพํ ธมฺมํ ปฏิจฺจ.}

\proseitem{23}{นกุสลํ นทสฺสเนน ปหาตพฺพ\-เหตุกํ ธมฺมํ ปฏิจฺจ ฯเปฯ. นกุสลํ นภาวนาย ปหาตพฺพ\-เหตุกํ ธมฺมํ ปฏิจฺจ.}

\prose{นกุสลํ นเน\-ว\-ทสฺสเนน นภาวนาย ปหาตพฺพ\-เหตุกํ ธมฺมํ ปฏิจฺจ.}

\proseitem{24}{นกุสลํ นอาจยคามิํ ธมฺมํ ปฏิจฺจ.}

\prose{นกุสลํ นอปจยคามิํ ธมฺมํ ปฏิจฺจ.}

\prose{นกุสลํ นเน\-วา\-จยคามิ\-นา\-ปจยคามิํ ธมฺมํ ปฏิจฺจ.}

\proseitem{25}{นกุสลํ นเสกฺขํ ธมฺมํ ปฏิจฺจ ฯเปฯ. นกุสลํ นอเสกฺขํ ธมฺมํ ปฏิจฺจ. นกุสลํ นเน\-ว\-เสกฺข\-นาเสกฺขํ ธมฺมํ ปฏิจฺจ.}

\proseitem{26}{นกุสลํ นปริตฺตํ ธมฺมํ ปฏิจฺจ.}

\prose{นกุสลํ นมหคฺคตํ ธมฺมํ ปฏิจฺจ ฯเปฯ. นกุสลํ นอปฺปมาณํ ธมฺมํ ปฏิจฺจ.}

\proseitem{27}{นกุสลํ นปริตฺตา\-รมฺมณํ ธมฺมํ ปฏิจฺจ ฯเปฯ. นกุสลํ นมหคฺคตา\-รมฺมณํ ธมฺมํ ปฏิจฺจ ฯเปฯ. นกุสลํ นอปฺปมาณา\-รมฺมณํ ธมฺมํ ปฏิจฺจ.}

\csromanheader{2}{1. \textbf{กุสลตฺติก} 13-20. หีนตฺติกาทิ}
\proseitem{28}{นกุสลํ นหีนํ ธมฺมํ ปฏิจฺจ.}

\prose{นกุสลํ นมชฺฌิมํ ธมฺมํ ปฏิจฺจ.}

\prose{นกุสลํ นปณีตํ ธมฺมํ ปฏิจฺจ.}

\proseitem{29}{นกุสลํ นมิจฺฉตฺต\-นิยตํ ธมฺมํ ปฏิจฺจ ฯเปฯ. นกุสลํ นสมฺมตฺต\-นิยตํ ธมฺมํ ปฏิจฺจ.}

\prose{นกุสลํ นอนิยตํ ธมฺมํ ปฏิจฺจ.}

\csromanmarkh{1. กุสลตฺติก}\csromanmarkhii{21. สนิทสฺสน\-ตฺติก}\csromanheadernomark{2}{1. กุสลตฺติก 21. สนิทสฺสน\-ตฺติก}
\proseitem{30}{\csromanfolio{147}นกุสลํ นมคฺคา\-รมฺมณํ ธมฺมํ ปฏิจฺจ ฯเปฯ. นกุสลํ นมคฺคเหตุกํ ธมฺมํ ปฏิจฺจ ฯเปฯ. นกุสลํ นมคฺคา\-ธิปติํ ธมฺมํ ปฏิจฺจ.}

\proseitem{31}{นกุสลํ นอนุปฺปนฺนํ ธมฺมํ ปฏิจฺจ ฯเปฯ. นกุสลํ นอุปฺปาทิํ ธมฺมํ ปฏิจฺจ.}

\proseitem{32}{นกุสลํ นอตีตํ ธมฺมํ ปฏิจฺจ ฯเปฯ. นกุสลํ นอนาคตํ ธมฺมํ ปฏิจฺจ.}

\proseitem{33}{นกุสลํ นอตีตา\-รมฺมณํ ธมฺมํ ปฏิจฺจ ฯเปฯ. นกุสลํ นอนาคตา\-รมฺมณํ ธมฺมํ ปฏิจฺจ ฯเปฯ. นกุสลํ นปจฺจุนฺนารมฺมณํ ธมฺมํ ปฏิจฺจ.}

\proseitem{34}{นกุสลํ นอชฺฌตฺตํ ธมฺมํ ปฏิจฺจ ฯเปฯ. นกุสลํ นพหิทฺธา ธมฺมํ ปฏิจฺจ.}

\proseitem{35}{นกุสลํ นอชฺฌตฺตา\-รมฺมณํ ธมฺมํ ปฏิจฺจ ฯเปฯ. นกุสลํ นพหิทฺธา\-รมฺมณํ ธมฺมํ ปฏิจฺจ.}

\csromanmarkh{1. กุสลตฺติก}\csromanmarkhii{21. สนิทสฺสน\-ตฺติก}\csromanheadernomark{2}{1. กุสลตฺติก 21. สนิทสฺสน\-ตฺติก}
\proseitem{36}{นกุสลํ นสนิทสฺสน\-สปฺปฏิฆํ ธมฺมํ ปฏิจฺจ นกุสโล นสนิทสฺสน\-สปฺปฏิโฆ ธมฺโม อุปฺปชฺชติ เหตุปจฺจยา.  นกุสลํ นสนิทสฺสน\-สปฺปฏิฆํ ธมฺมํ ปฏิจฺจ นอกุสโล นสนิทสฺสน\-สปฺปฏิโฆ ธมฺโม อุปฺปชฺชติ เหตุปจฺจยา.  นกุสลํ นสนิทสฺสน\-สปฺปฏิฆํ ธมฺมํ ปฏิจฺจ นอพฺยากโต นสนิทสฺสน\-สปฺปฏิโฆ ธมฺโม อุปฺปชฺชติ เหตุปจฺจยา.  นกุสลํ นสนิทสฺสน\-สปฺปฏิฆํ ธมฺมํ ปฏิจฺจ นกุสโล นสนิทสฺสน\-สปฺปฏิโฆ จ นอพฺยากโต นสนิทสฺสน\-สปฺปฏิโฆ จ ธมฺมา อุปฺปชฺชนฺติ เหตุปจฺจยา.  นกุสลํ นสนิทสฺสน\-สปฺปฏิฆํ ธมฺมํ ปฏิจฺจ นกุสโล นสนิทสฺสน\-สปฺปฏิโฆ จ นอกุสโล นสนิทสฺสน\-สปฺปฏิโฆ จ ธมฺมา อุปฺปชฺชนฺติ เหตุปจฺจยา. (5)}

\prose{นอกุสลํ นสนิทสฺสน\-สปฺปฏิฆํ ธมฺมํ ปฏิจฺจ นอกุสโล นสนิทสฺสน\-สปฺปฏิโฆ ธมฺโม อุปฺปชฺชติ เหตุปจฺจยา.  นอกุสลํ นสนิทสฺสน\-สปฺปฏิฆํ ธมฺมํ ปฏิจฺจ นกุสโล นสนิทสฺสน\-สปฺปฏิโฆ ธมฺโม อุปฺปชฺชติ เหตุปจฺจยา.  นอกุสลํ นสนิทสฺสน\-สปฺปฏิฆํ ธมฺมํ ปฏิจฺจ นอพฺยากโต นสนิทสฺสน\-สปฺปฏิโฆ ธมฺโม อุปฺปชฺชติ เหตุปจฺจยา.  นอกุสลํ นสนิทสฺสน\-สปฺปฏิฆํ ธมฺมํ ปฏิจฺจ นอกุสโล นสนิทสฺสน\-สปฺปฏิโฆ จ นอพฺยากโต นสนิทสฺสน\-สปฺปฏิโฆ จ ธมฺมา อุปฺปชฺชนฺติ เหตุปจฺจยา.  }

\prose{\csromanfolio{148}นอกุสลํ นสนิทสฺสน\-สปฺปฏิฆํ ธมฺมํ ปฏิจฺจ นกุสโล นสนิทสฺสน\-สปฺปฏิโฆ จ นอกุสโล นสนิทสฺสน\-สปฺปฏิโฆ จ ธมฺมา อุปฺปชฺชนฺติ เหตุปจฺจยา. (5)}

\prose{นอพฺยากตํ นสนิทสฺสน\-สปฺปฏิฆํ ธมฺมํ ปฏิจฺจ นอพฺยากโต นสนิทสฺสน\-สปฺปฏิโฆ ธมฺโม อุปฺปชฺชติ เหตุปจฺจยา.  นอพฺยากตํ นสนิทสฺสน\-สปฺปฏิฆํ ธมฺมํ ปฏิจฺจ นกุสโล นสนิทสฺสน\-สปฺปฏิโฆ ธมฺโม อุปฺปชฺชติ เหตุปจฺจยา.  นอพฺยากตํ นสนิทสฺสน\-สปฺปฏิฆํ ธมฺมํ ปฏิจฺจ นอกุสโล นสนิทสฺสน\-สปฺปฏิโฆ ธมฺโม อุปฺปชฺชติ เหตุปจฺจยา.  นอพฺยากตํ นสนิทสฺสน\-สปฺปฏิฆํ ธมฺมํ ปฏิจฺจ นกุสโล นสนิทสฺสน\-สปฺปฏิโฆ จ นอพฺยากโต นสนิทสฺสน\-สปฺปฏิโฆ จ ธมฺมา อุปฺปชฺชนฺติ เหตุปจฺจยา.  นอพฺยากตํ นสนิทสฺสน\-สปฺปฏิฆํ ธมฺมํ ปฏิจฺจ นอกุสโล นสนิทสฺสน\-สปฺปฏิโฆ จ นอพฺยากโต นสนิทสฺสน\-สปฺปฏิโฆ จ ธมฺมา อุปฺปชฺชนฺติ เหตุปจฺจยา.  นอพฺยากตํ นสนิทสฺสน\-สปฺปฏิฆํ ธมฺมํ ปฏิจฺจ นกุสโล นสนิทสฺสน\-สปฺปฏิโฆ จ นอกุสโล นสนิทสฺสน\-สปฺปฏิโฆ จ ธมฺมา อุปฺปชฺชนฺติ เหตุปจฺจยา. (6)}

\prose{นกุสลํ นสนิทสฺสนสปฺปฏิฆญฺจ นอพฺยากตํ นสนิทสฺสนสปฺปฏิฆญฺจ ธมฺมํ ปฏิจฺจ นกุสโล นสนิทสฺสน\-สปฺปฏิโฆ ธมฺโม อุปฺปชฺชติ เหตุปจฺจยา. ปญฺจ. นอกุสลํ นสนิทสฺสนสปฺปฏิฆญฺจ นอพฺยากตํ นสนิทสฺสนสปฺปฏิฆญฺจ ธมฺมํ ปฏิจฺจ นกุสโล นสนิทสฺสน\-สปฺปฏิโฆ ธมฺโม อุปฺปชฺชติ เหตุปจฺจยา. ปญฺจ.}

\prose{นกุสลํ นสนิทสฺสนสปฺปฏิฆญฺจ นอกุสลํ นสนิทสฺสนสปฺปฏิฆญฺจ ธมฺมํ ปฏิจฺจ นกุสโล นสนิทสฺสน\-สปฺปฏิโฆ ธมฺโม อุปฺปชฺชติ เหตุปจฺจยา.  นกุสลํ นสนิทสฺสนสปฺปฏิฆญฺจ นอกุสลํ นสนิทสฺสนสปฺปฏิฆญฺจ ธมฺมํ ปฏิจฺจ นอกุสโล นสนิทสฺสน\-สปฺปฏิโฆ ธมฺโม อุปฺปชฺชติ เหตุปจฺจยา.  นกุสลํ นสนิทสฺสนสปฺปฏิฆญฺจ นอกุสลํ นสนิทสฺสนสปฺปฏิฆญฺจ ธมฺมํ ปฏิจฺจ นกุสโล นสนิทสฺสน\-สปฺปฏิโฆ จ นอกุสโล นสนิทสฺสน\-สปฺปฏิโฆ จ ธมฺมา อุปฺปชฺชนฺติ เหตุปจฺจยา. (3) (สํขิตฺตํ.)}

\prose{เหตุยา เอกูนติํส ฯเปฯ อวิคเต เอกูนติํส. (สพฺพตฺถ วิตฺถาโร.)}

\csromanmarkh{2. เวทนาตฺติก}\csromanmarkhii{1. กุสลตฺติก}\csromanheadernomark{2}{2. เวทนาตฺติก 1. กุสลตฺติก}
\proseitem{37}{\csromanfolio{149}นกุสลํ นอนิทสฺสน\-สปฺปฏิฆํ ธมฺมํ ปฏิจฺจ นกุสโล นอนิทสฺสน\-สปฺปฏิโฆ ธมฺโม อุปฺปชฺชติ เหตุปจฺจยา. (สํขิตฺตํ.)}

\prose{เหตุยา เอกูนติํส ฯเปฯ อวิคเต เอกูนติํส. (สพฺพตฺถ วิตฺถาโร.)}

\proseitem{38}{นกุสลํ นอนิทสฺสน\-อปฺปฏิฆํ ธมฺมํ ปฏิจฺจ นกุสโล นอนิทสฺสน\-อปฺปฏิโฆ ธมฺโม อุปฺปชฺชติ เหตุปจฺจยา. (สํขิตฺตํ.)}

\prose{เหตุยา นว. (สพฺพตฺถ วิตฺถาโร.)}

\csromanmarkh{2. เวทนาตฺติก}\csromanmarkhii{1. กุสลตฺติก}\csromanheadernomark{2}{2. เวทนาตฺติก 1. กุสลตฺติก}
\proseitem{39}{นสุขาย เวทนาย สมฺปยุตฺตํ นกุสลํ ธมฺมํ ปฏิจฺจ นสุขาย เวทนาย สมฺปยุตฺโต นกุสโล ธมฺโม อุปฺปชฺชติ เหตุปจฺจยา.  นสุขาย เวทนาย สมฺปยุตฺตํ นกุสลํ ธมฺมํ ปฏิจฺจ นทุกฺขาย เวทนาย สมฺปยุตฺโต นกุสโล ธมฺโม อุปฺปชฺชติ เหตุปจฺจยา.  นสุขาย เวทนาย สมฺปยุตฺตํ นกุสลํ ธมฺมํ ปฏิจฺจ นอทุกฺขม\-สุขาย เวทนาย สมฺปยุตฺโต นกุสโล ธมฺโม อุปฺปชฺชติ เหตุปจฺจยา.  นสุขาย เวทนาย สมฺปยุตฺตํ นกุสลํ ธมฺมํ ปฏิจฺจ นสุขาย เวทนาย สมฺปยุตฺโต นกุสโล จ นอทุกฺขม\-สุขาย เวทนาย สมฺปยุตฺโต นกุสโล จ ธมฺมา อุปฺปชฺชนฺติ เหตุปจฺจยา ฯเปฯ. สตฺต.}

\prose{นทุกฺขาย เวทนาย สมฺปยุตฺตํ นกุสลํ ธมฺมํ ปฏิจฺจ นทุกฺขาย เวทนาย สมฺปยุตฺโต นกุสโล ธมฺโม อุปฺปชฺชติ เหตุปจฺจยา.  นทุกฺขาย เวทนาย สมฺปยุตฺตํ นกุสลํ ธมฺมํ ปฏิจฺจ นสุขาย เวทนาย สมฺปยุตฺโต นกุสโล ธมฺโม อุปฺปชฺชติ เหตุปจฺจยา.  นทุกฺขาย เวทนาย สมฺปยุตฺตํ นกุสลํ ธมฺมํ ปฏิจฺจ นอทุกฺขม\-สุขาย เวทนาย สมฺปยุตฺโต นกุสโล ธมฺโม อุปฺปชฺชติ เหตุปจฺจยา.  นทุกฺขาย เวทนาย สมฺปยุตฺตํ นกุสลํ ธมฺมํ ปฏิจฺจ นสุขาย เวทนาย สมฺปยุตฺโต นกุสโล จ นอทุกฺขม\-สุขาย เวทนาย สมฺปยุตฺโต นกุสโล จ ธมฺมา อุปฺปชฺชนฺติ เหตุปจฺจยา ฯเปฯ. สตฺต.}

\prose{นอทุกฺขม\-สุขาย เวทนาย สมฺปยุตฺตํ นกุสลํ ธมฺมํ ปฏิจฺจ นอทุกฺขม\-สุขาย เวทนาย สมฺปยุตฺโต นกุสโล ธมฺโม อุปฺปชฺชติ เหตุปจฺจยา.  นอทุกฺขม\-สุขาย เวทนาย สมฺปยุตฺตํ นกุสลํ ธมฺมํ ปฏิจฺจ \csromanfolio{150}นสุขาย เวทนาย สมฺปยุตฺโต นกุสโล ธมฺโม อุปฺปชฺชติ เหตุปจฺจยา.  นอทุกฺขม\-สุขาย เวทนาย สมฺปยุตฺตํ นกุสลํ ธมฺมํ ปฏิจฺจ นทุกฺขาย เวทนาย สมฺปยุตฺโต นกุสโล ธมฺโม อุปฺปชฺชติ เหตุปจฺจยา.  นอทุกฺขม\-สุขาย เวทนาย สมฺปยุตฺตํ นกุสลํ ธมฺมํ ปฏิจฺจ นสุขาย เวทนาย สมฺปยุตฺโต นกุสโล จ นอทุกฺขม\-สุขาย เวทนาย สมฺปยุตฺโต นกุสโล จ ธมฺมา อุปฺปชฺชนฺติ เหตุปจฺจยา ฯเปฯ. สตฺต.}

\proseitem{40}{นสุขาย เวทนาย สมฺปยุตฺตํ นกุสลญฺจ นอทุกฺขม\-สุขาย เวทนาย สมฺปยุตฺตํ นกุสลญฺจ ธมฺมํ ปฏิจฺจ นสุขาย เวทนาย สมฺปยุตฺโต นกุสโล ธมฺโม อุปฺปชฺชติ เหตุปจฺจยา.  นสุขาย เวทนาย สมฺปยุตฺตํ นกุสลญฺจ น- อทุกฺขม\-สุขาย เวทนาย สมฺปยุตฺตํ นกุสลญฺจ ธมฺมํ ปฏิจฺจ นทุกฺขาย เวทนาย สมฺปยุตฺโต นกุสโล ธมฺโม อุปฺปชฺชติ เหตุปจฺจยา.  นสุขาย เวทนาย สมฺปยุตฺตํ นกุสลญฺจ นอทุกฺขม\-สุขาย เวทนาย สมฺปยุตฺตํ นกุสลญฺจ ธมฺมํ ปฏิจฺจ นอทุกฺขม\-สุขาย เวทนาย สมฺปยุตฺโต นกุสโล ธมฺโม อุปฺปชฺชติ เหตุปจฺจยา.  นสุขาย เวทนาย สมฺปยุตฺตํ นกุสลญฺจ นอทุกฺขม\-สุขาย เวทนาย สมฺปยุตฺตํ นกุสลญฺจ ธมฺมํ ปฏิจฺจ นสุขาย เวทนาย สมฺปยุตฺโต นกุสโล จ นอทุกฺขม\-สุขาย เวทนาย สมฺปยุตฺโต นกุสโล จ ธมฺมา อุปฺปชฺชนฺติ เหตุปจฺจยา ฯเปฯ. สตฺต.}

\prose{นทุกฺขาย เวทนาย สมฺปยุตฺตํ นกุสลญฺจ นอทุกฺขม\-สุขาย เวทนาย สมฺปยุตฺตํ นกุสลญฺจ ธมฺมํ ปฏิจฺจ นสุขาย เวทนาย สมฺปยุตฺโต นกุสโล ธมฺโม อุปฺปชฺชติ เหตุปจฺจยา.  นทุกฺขาย เวทนาย สมฺปยุตฺตํ นกุสลญฺจ น- อทุกฺขม\-สุขาย เวทนาย สมฺปยุตฺตํ นกุสลญฺจ ธมฺมํ ปฏิจฺจ นทุกฺขาย เวทนาย สมฺปยุตฺโต นกุสโล ธมฺโม อุปฺปชฺชติ เหตุปจฺจยา.  นทุกฺขาย เวทนาย สมฺปยุตฺตํ นกุสลญฺจ นอทุกฺขม\-สุขาย เวทนาย สมฺปยุตฺตํ นกุสลญฺจ ธมฺมํ ปฏิจฺจ นอทุกฺขม\-สุขาย เวทนาย สมฺปยุตฺโต นกุสโล ธมฺโม อุปฺปชฺชติ เหตุปจฺจยา ฯเปฯ. สตฺต. (สํขิตฺตํ.)}

\prose{เหตุยา เอกูนปญฺญาส, อารมฺมเณ เอกูนปญฺญาส ฯเปฯ อวิคเต เอกูนปญฺญาส.}

\csromanmarkh{6-21. วิตกฺกตฺติกาทิ}\csromanmarkhii{1. กุสลตฺติก}\csromanheadernomark{2}{6-21. วิตกฺกตฺติกาทิ 1. กุสลตฺติก}
\proseitem{41}{\csromanfolio{151}นสุขาย เวทนาย สมฺปยุตฺตํ นอกุสลํ ธมฺมํ ปฏิจฺจ นสุขาย เวทนาย สมฺปยุตฺโต นอกุสโล ธมฺโม อุปฺปชฺชติ เหตุปจฺจยา. (สํขิตฺตํ.) . เหตุยา เอกูนปญฺญาส.}

\proseitem{42}{นสุขาย เวทนาย สมฺปยุตฺตํ นอพฺยากตํ ธมฺมํ ปฏิจฺจ นสุขาย เวทนาย สมฺปยุตฺโต นอพฺยากโต ธมฺโม อุปฺปชฺชติ เหตุปจฺจยา. (สํขิตฺตํ.) . เหตุยา อฏฺฐจตฺตาลีส. (สพฺพตฺถ วิตฺถาโร.)}

\csromanmarkh{2. เวทนาตฺติก}\csromanmarkhii{2. วิปากตฺติก}\csromanheadernomark{2}{2. \textbf{เวทนาตฺติก 2. วิปากตฺติก}}
\proseitem{43}{นสุขาย เวทนาย สมฺปยุตฺตํ นวิปากํ ธมฺมํ ปฏิจฺจ นสุขาย เวทนาย สมฺปยุตฺโต นวิปาโก ธมฺโม อุปฺปชฺชติ เหตุปจฺจยา.}

\prose{เหตุยา เอกูนปญฺญาส ฯเปฯ วิคเต เอกูนปญฺญาส. (สํขิตฺตํ.)}

\csromanmarkh{3. วิปากตฺติก}\csromanmarkhii{1. กุสลตฺติก}\csromanheadernomark{2}{3. วิปากตฺติก 1. กุสลตฺติก}
\proseitem{44}{นวิปากํ นกุสลํ ธมฺมํ ปฏิจฺจ ฯเปฯ. นวิปาก\-ธมฺม\-ธมฺมํ นอกุสลํ ธมฺมํ ปฏิจฺจ.}

\prose{นเน\-ว\-วิปาก\-น\-วิปากธมฺม\-ธมฺมํ นอพฺยากตํ ธมฺมํ ปฏิจฺจ.}

\csromanmarkh{4. อุปาทินฺนตฺติก}\csromanmarkhii{1. กุสลตฺติก}\csromanheadernomark{2}{4. \textbf{อุปาทินฺนตฺติก 1. กุสลตฺติก}}
\proseitem{45}{นอุปาทินฺนุปาทานิยํ นกุสลํ\csromansharedfootnote{f22a090cb8b33468}{อิโต ปฏฺฐาย ยาว อชฺฌตฺตารมฺมณตฺติกา สพฺพโปตฺถเกสุ อกุสลาพฺยากตปทานิ น ทิสฺสนฺติ, เวทนาตฺติก สนิทสฺสนตฺติเกสุ วิย นนุ เตหิปิ ภวิตพฺพํ.} ธมฺมํ ปฏิจฺจ ฯเปฯ. นอนุปาทินฺนุ\-ปาทานิยํ นกุสลํ ธมฺมํ ปฏิจฺจ ฯเปฯ. นอนุปาทินฺน\-อนุปาทานิยํ นกุสลํ ธมฺมํ ปฏิจฺจ.}

\csromanmarkh{5. สํกิลิฏฺฐ\-ตฺติก}\csromanmarkhii{1. กุสลตฺติก}\csromanheadernomark{2}{5. สํกิลิฏฺฐ\-ตฺติก 1. กุสลตฺติก}
\proseitem{46}{นสํกิลิฏฺฐ\-สํกิเลสิกํ นกุสลํ ธมฺมํ ปฏิจฺจ ฯเปฯ. นอสํกิลิฏฺฐ\-สํกิเลสิกํ นกุสลํ ธมฺมํ ปฏิจฺจ ฯเปฯ. นอสํกิลิฏฺฐ\-อสํกิเลสิกํ นกุสลํ ธมฺมํ ปฏิจฺจ.}

\csromanmarkh{6-21. วิตกฺกตฺติกาทิ}\csromanmarkhii{1. กุสลตฺติก}\csromanheadernomark{2}{6-21. วิตกฺกตฺติกาทิ 1. กุสลตฺติก}
\proseitem{47}{นสวิตกฺก\-สวิจารํ นกุสลํ ธมฺมํ ปฏิจฺจ ฯเปฯ. นอวิตกฺก\-วิจารมตฺตํ นกุสลํ ธมฺมํ ปฏิจฺจ ฯเปฯ. นอวิตกฺก\-อวิจารํ นกุสลํ ธมฺมํ ปฏิจฺจ.}

\csromanpalirecto
\gambhira{ธมฺม\-ปจฺจนีย ติกติก\-ปฏฺฐาน\-ปาฬิ}
\proseitem{48}{\csromanfolio{152}นปีติสหคตํ นกุสลํ ธมฺมํ ปฏิจฺจ ฯเปฯ. นสุขสหคตํ นกุสลํ ธมฺมํ ปฏิจฺจ ฯเปฯ. นอุเปกฺขาสหคตํ นกุสลํ ธมฺมํ ปฏิจฺจ.}

\proseitem{49}{นทสฺสเนน ปหาตพฺพํ นกุสลํ ธมฺมํ ปฏิจฺจ ฯเปฯ. นภาวนาย ปหาตพฺพํ นกุสลํ ธมฺมํ ปฏิจฺจ ฯเปฯ. นเน\-ว\-ทสฺสเนน นภาวนาย ปหาตพฺพํ นกุสลํ ธมฺมํ ปฏิจฺจ.}

\proseitem{50}{นทสฺสเนน ปหาตพฺพ\-เหตุกํ นกุสลํ ธมฺมํ ปฏิจฺจ ฯเปฯ. นภาวนาย ปหาตพฺพ\-เหตุกํ นกุสลํ ธมฺมํ ปฏิจฺจ ฯเปฯ. นเน\-ว\-ทสฺสเนน นภาวนาย ปหาตพฺพ\-เหตุกํ นกุสลํ ธมฺมํ ปฏิจฺจ.}

\proseitem{51}{นอาจยคามิํ นกุสลํ ธมฺมํ ปฏิจฺจ ฯเปฯ. นอปจยคามิํ นกุสลํ ธมฺมํ ปฏิจฺจ ฯเปฯ. นเน\-วา\-จยคามิ\-นา\-ปจยคามิํ นกุสลํ ธมฺมํ ปฏิจฺจ.}

\proseitem{52}{นเสกฺขํ นกุสลํ ธมฺมํ ปฏิจฺจ ฯเปฯ. นอเสกฺขํ นกุสลํ ธมฺมํ ปฏิจฺจ ฯเปฯ. นเนวาเสกฺขนาเสกฺขํ นกุสลํ ธมฺมํ ปฏิจฺจ.}

\proseitem{53}{นปริตฺตํ นกุสลํ ธมฺมํ ปฏิจฺจ ฯเปฯ. นมหคฺคตํ นกุสลํ ธมฺมํ ปฏิจฺจ ฯเปฯ. นอปฺปมาณํ นกุสลํ ธมฺมํ ปฏิจฺจ.}

\proseitem{54}{นปริตฺตา\-รมฺมณํ นกุสลํ ธมฺมํ ปฏิจฺจ ฯเปฯ. นมหคฺคตา\-รมฺมณํ นกุสลํ ธมฺมํ ปฏิจฺจ ฯเปฯ. นอปฺปมาณา\-รมฺมณํ นกุสลํ ธมฺมํ ปฏิจฺจ.}

\proseitem{55}{นหีนํ นกุสลํ ธมฺมํ ปฏิจฺจ ฯเปฯ. นมชฺฌิมํ นกุสลํ ธมฺมํ ปฏิจฺจ ฯเปฯ. นปณีตํ นกุสลํ ธมฺมํ ปฏิจฺจ.}

\proseitem{56}{นมิจฺฉตฺต\-นิยตํ นกุสลํ ธมฺมํ ปฏิจฺจ ฯเปฯ. นสมฺมตฺต\-นิยตํ นกุสลํ ธมฺมํ ปฏิจฺจ ฯเปฯ. นอนิยตํ นกุสลํ ธมฺมํ ปฏิจฺจ.}

\proseitem{57}{นมคฺคา\-รมฺมณํ นกุสลํ ธมฺมํ ปฏิจฺจ ฯเปฯ. นมคฺคเหตุกํ นกุสลํ ธมฺมํ ปฏิจฺจ ฯเปฯ. นมคฺคา\-ธิปติํ นกุสลํ ธมฺมํ ปฏิจฺจ.}

\proseitem{58}{นอนุปฺปนฺนํ นกุสลํ ธมฺมํ ปฏิจฺจ ฯเปฯ. นอุปฺปาทิํ นกุสลํ ธมฺมํ ปฏิจฺจ.}

\proseitem{59}{นอตีตํ นกุสลํ ธมฺมํ ปฏิจฺจ ฯเปฯ. นอนาคตํ นกุสลํ ธมฺมํ ปฏิจฺจ.}

\proseitem{60}{นอตีตา\-รมฺมณํ นกุสลํ ธมฺมํ ปฏิจฺจ ฯเปฯ. นอนาคตา\-รมฺมณํ นกุสลํ ธมฺมํ ปฏิจฺจ ฯเปฯ. นปจฺจุปฺปนฺนา\-รมฺมณํ นกุสลํ ธมฺมํ ปฏิจฺจ.}

\csromanmarkh{22. สนิทสฺสน\-ตฺติก}\csromanmarkhii{1. กุสลตฺติก}\csromanheadernomark{2}{22. สนิทสฺสน\-ตฺติก 1. กุสลตฺติก}
\proseitem{61}{\csromanfolio{153}นอชฺฌตฺตํ นอกุสลํ ธมฺมํ ปฏิจฺจ ฯเปฯ. นพหิทฺธา นกุสลํ ธมฺมํ ปฏิจฺจ.}

\proseitem{62}{นอชฺฌตฺตา\-รมฺมณํ นกุสลํ ธมฺมํ ปฏิจฺจ ฯเปฯ. นพหิทฺธา\-รมฺมณํ นกุสลํ ธมฺมํ ปฏิจฺจ.}

\csromanmarkh{22. สนิทสฺสน\-ตฺติก}\csromanmarkhii{1. กุสลตฺติก}\csromanheadernomark{2}{22. สนิทสฺสน\-ตฺติก 1. กุสลตฺติก}
\proseitem{63}{นสนิทสฺสน\-สปฺปฏิฆํ นกุสลํ ธมฺมํ ปฏิจฺจ นสนิทสฺสน\-สปฺปฏิโฆ นกุสโล ธมฺโม อุปฺปชฺชติ เหตุปจฺจยา.  นสนิทสฺสน\-สปฺปฏิฆํ นกุสลํ ธมฺมํ ปฏิจฺจ นอนิทสฺสน\-สปฺปฏิโฆ นกุสโล ธมฺโม อุปฺปชฺชติ เหตุปจฺจยา.  นสนิทสฺสน\-สปฺปฏิฆํ นกุสลํ ธมฺมํ ปฏิจฺจ นอนิทสฺสน\-อปฺปฏิโฆ นกุสโล ธมฺโม อุปฺปชฺชติ เหตุปจฺจยา ฯเปฯ. ฉ.}

\prose{นอนิทสฺสน\-สปฺปฏิฆํ นกุสลํ ธมฺมํ ปฏิจฺจ นอนิทสฺสน\-สปฺปฏิโฆ นกุสโล ธมฺโม อุปฺปชฺชติ เหตุปจฺจยา.  นอนิทสฺสน\-สปฺปฏิฆํ นกุสลํ ธมฺมํ ปฏิจฺจ นสนิทสฺสน\-สปฺปฏิโฆ นกุสโล ธมฺโม อุปฺปชฺชติ เหตุปจฺจยา.  น- อนิทสฺสน\-สปฺปฏิฆํ นกุสลํ ธมฺมํ ปฏิจฺจ นอนิทสฺสน\-อปฺปฏิโฆ นกุสโล ธมฺโม อุปฺปชฺชติ เหตุปจฺจยา ฯเปฯ. ฉ.}

\prose{นอนิทสฺสน\-อปฺปฏิฆํ นกุสลํ ธมฺมํ ปฏิจฺจ นอนิทสฺสน\-อปฺปฏิโฆ นกุสโล ธมฺโม อุปฺปชฺชติ เหตุปจฺจยา. ฉ.}

\prose{นสนิทสฺสน\-สปฺปฏิฆํ นกุสลญฺจ นอนิทสฺสน\-อปฺปฏิฆํ นกุสลญฺจ ธมฺมํ ปฏิจฺจ นสนิทสฺสน\-สปฺปฏิโฆ นกุสโล ธมฺโม อุปฺปชฺชติ เหตุปจฺจยา. ฉ.}

\prose{นอนิทสฺสน\-สปฺปฏิฆํ นกุสลญฺจ นอนิทสฺสน\-อปฺปฏิฆํ นกุสลญฺจ ธมฺมํ ปฏิจฺจ นสนิทสฺสน\-สปฺปฏิโฆ นกุสโล ธมฺโม อุปฺปชฺชติ เหตุปจฺจยา. ฉ. (สํขิตฺตํ.) . เหตุยา ติํส, อารมฺมเณ นว ฯเปฯ อวิคเต ติํส. (สพฺพตฺถ วิตฺถาโร.)}

\proseitem{64}{นสนิทสฺสน\-สปฺปฏิฆํ นอกุสลํ ธมฺมํ ปฏิจฺจ นสนิทสฺสน\-สปฺปฏิโฆ นอกุสโล ธมฺโม อุปฺปชฺชติ เหตุปจฺจยา.  เหตุยา ติํส, อารมฺมเณ นว ฯเปฯ อวิคเต ติํส. (สพฺพตฺถ วิตฺถาโร.)}

\proseitem{65}{นสนิทสฺสน\-สปฺปฏิฆํ นอพฺยากตํ ธมฺมํ ปฏิจฺจ นสนิทสฺสน\-สปฺปฏิโฆ นอพฺยากโต ธมฺโม อุปฺปชฺชติ เหตุปจฺจยา.}

\prose{เหตุยา นว. (สพฺพตฺถ นว.)}

\csromanmarkh{ธมฺม\-ปจฺจนีย ติกติก\-ปฏฺฐาน\-ปาฬิ}\csromanmarkhii{22. สนิทสฺสน\-ตฺติก 2-21. เวทนาตฺติกาทิ}\csromanheadernomark{2}{ธมฺม\-ปจฺจนีย ติกติก\-ปฏฺฐาน\-ปาฬิ 22. สนิทสฺสน\-ตฺติก 2-21. เวทนาตฺติกาทิ}
\proseitem{66}{\csromanfolio{154}นสนิทสฺสน\-สปฺปฏิฆํ นสุขาย เวทนาย สมฺปยุตฺตํ ธมฺมํ ปฏิจฺจ. เหตุยา ติํส.}

\prose{นสนิทสฺสน\-สปฺปฏิฆํ นทุกฺขาย เวทนาย สมฺปยุตฺตํ ธมฺมํ ปฏิจฺจ. เหตุยา ติํส.}

\prose{นสนิทสฺสน\-สปฺปฏิฆํ นอทุกฺขม\-สุขาย เวทนาย สมฺปยุตฺตํ ธมฺมํ ปฏิจฺจ. เหตุยา ติํส.}

\proseitem{67}{นสนิทสฺสน\-สปฺปฏิฆํ นวิปากํ ธมฺมํ ปฏิจฺจ ฯเปฯ.}

\prose{นสนิทสฺสน\-สปฺปฏิฆํ นวิปาก\-ธมฺม\-ธมฺมํ ปฏิจฺจ. เหตุยา ติํส.}

\prose{นสนิทสฺสน\-สปฺปฏิฆํ นเน\-ว\-วิปาก\-น\-วิปากธมฺม\-ธมฺมํ ปฏิจฺจ. เหตุยา นว.}

\proseitem{68}{นสนิทสฺสน\-สปฺปฏิฆํ นอุปาทินฺนุปาทานิยํ ธมฺมํ ปฏิจฺจ ฯเปฯ.}

\prose{นสนิทสฺสน\-สปฺปฏิฆํ นอนุปาทินฺนุ\-ปาทานิยํ ธมฺมํ ปฏิจฺจ ฯเปฯ.}

\prose{นสนิทสฺสน\-สปฺปฏิฆํ นอนุปาทินฺน\-อนุปาทานิยํ ธมฺมํ ปฏิจฺจ.}

\prose{เหตุยา ติํส.}

\proseitem{69}{นสนิทสฺสน\-สปฺปฏิฆํ นสํกิลิฏฺฐ\-สํกิเลสิกํ ธมฺมํ ปฏิจฺจ.}

\prose{เหตุยา ติํส.}

\prose{นสนิทสฺสน\-สปฺปฏิฆํ นอสํกิลิฏฺฐ\-สํกิเลสิกํ ธมฺมํ ปฏิจฺจ ฯเปฯ.}

\prose{นสนิทสฺสน\-สปฺปฏิฆํ นอสํกิลิฏฺฐ\-อสํกิเลสิกํ ธมฺมํ ปฏิจฺจ.}

\prose{เหตุยา ติํส.}

\proseitem{70}{นสนิทสฺสน\-สปฺปฏิฆํ นสวิตกฺก\-สวิจารํ ธมฺมํ ปฏิจฺจ ฯเปฯ.}

\prose{นสนิทสฺสน\-สปฺปฏิฆํ นอวิตกฺก\-วิจารมตฺตํ ธมฺมํ ปฏิจฺจ.}

\prose{นสนิทสฺสน\-สปฺปฏิฆํ นอวิตกฺก\-อวิจารํ ธมฺมํ ปฏิจฺจ ฯเปฯ.}

\proseitem{71}{นสนิทสฺสน\-สปฺปฏิฆํ นปีติสหคตํ ธมฺมํ ปฏิจฺจ ฯเปฯ. นสนิทสฺสน\-สปฺปฏิฆํ นสุขสหคตํ ธมฺมํ ปฏิจฺจ ฯเปฯ. นสนิทสฺสน\-สปฺปฏิฆํ นอุเปกฺขาสหคตํ ธมฺมํ ปฏิจฺจ.}

\csromanheader{2}{22. สนิทสฺสน\-ตฺติก 2-21. เวทนาตฺติกาทิ}
\proseitem{72}{\csromanfolio{155}นสนิทสฺสน\-สปฺปฏิฆํ นทสฺสเนน ปหาตพฺพํ ธมฺมํ ปฏิจฺจ ฯเปฯ. นสนิทสฺสน\-สปฺปฏิฆํ นภาวนาย ปหาตพฺพํ ธมฺมํ ปฏิจฺจ.}

\prose{นสนิทสฺสน\-สปฺปฏิฆํ นเน\-ว\-ทสฺสเนน นภาวนาย ปหาตพฺพํ ธมฺมํ ปฏิจฺจ.}

\proseitem{73}{นสนิทสฺสน\-สปฺปฏิฆํ นทสฺสเนน ปหาตพฺพ\-เหตุกํ ธมฺมํ ปฏิจฺจ ฯเปฯ. นสนิทสฺสน\-สปฺปฏิฆํ นภาวนาย ปหาตพฺพ\-เหตุกํ ธมฺมํ ปฏิจฺจ.}

\prose{นสนิทสฺสน\-สปฺปฏิฆํ นเน\-ว\-ทสฺสเนน นภาวนาย ปหาตพฺพ\-เหตุกํ ธมฺมํ ปฏิจฺจ.}

\proseitem{74}{นสนิทสฺสน\-สปฺปฏิฆํ นอาจยคามิํ ธมฺมํ ปฏิจฺจ ฯเปฯ. นสนิทสฺสน\-สปฺปฏิฆํ นอปจยคามิํ ธมฺมํ ปฏิจฺจ.}

\prose{นสนิทสฺสน\-สปฺปฏิฆํ นเน\-วา\-จยคามิ\-นา\-ปจยคามิํ ธมฺมํ ปฏิจฺจ.}

\proseitem{75}{นสนิทสฺสน\-สปฺปฏิฆํ นเสกฺขํ ธมฺมํ ปฏิจฺจ ฯเปฯ. นสนิทสฺสน\-สปฺปฏิฆํ นอเสกฺขํ ธมฺมํ ปฏิจฺจ.}

\prose{นสนิทสฺสน\-สปฺปฏิฆํ นเนวาเสกฺขนาเสกฺขํ ธมฺมํ ปฏิจฺจ ฯเปฯ.}

\proseitem{76}{นสนิทสฺสน\-สปฺปฏิฆํ นปริตฺตํ ธมฺมํ ปฏิจฺจ.}

\prose{นสนิทสฺสน\-สปฺปฏิฆํ นมหคฺคตํ ธมฺมํ ปฏิจฺจ ฯเปฯ. นสนิทสฺสน\-สปฺปฏิฆํ นอปฺปมาณํ ธมฺมํ ปฏิจฺจ.}

\proseitem{77}{นสนิทสฺสน\-สปฺปฏิฆํ นปริตฺตา\-รมฺมณํ ธมฺมํ ปฏิจฺจ ฯเปฯ. นสนิทสฺสน\-สปฺปฏิฆํ นมหคฺคตา\-รมฺมณํ ธมฺมํ ปฏิจฺจ ฯเปฯ. นสนิทสฺสน\-สปฺปฏิฆํ นอปฺปมาณา\-รมฺมณํ ธมฺมํ ปฏิจฺจ.}

\proseitem{78}{นสนิทสฺสน\-สปฺปฏิฆํ นหีนํ ธมฺมํ ปฏิจฺจ.}

\prose{นสนิทสฺสน\-สปฺปฏิฆํ นมชฺฌิมํ ธมฺมํ ปฏิจฺจ.}

\prose{นสนิทสฺสน\-สปฺปฏิฆํ นปณีตํ ธมฺมํ ปฏิจฺจ.}

\proseitem{79}{นสนิทสฺสน\-สปฺปฏิฆํ นมิจฺฉตฺต\-นิยตํ ธมฺมํ ปฏิจฺจ ฯเปฯ. นสนิทสฺสน\-สปฺปฏิฆํ นสมฺมตฺต\-นิยตํ ธมฺมํ ปฏิจฺจ.}

\prose{นสนิทสฺสน\-สปฺปฏิฆํ นอนิยตํ ธมฺมํ ปฏิจฺจ ฯเปฯ.}

\proseitem{80}{\csromanfolio{156}นสนิทสฺสน\-สปฺปฏิฆํ นมคฺคา\-รมฺมณํ ธมฺมํ ปฏิจฺจ ฯเปฯ. นสนิทสฺสน\-สปฺปฏิฆํ นมคฺคเหตุกํ ธมฺมํ ปฏิจฺจ ฯเปฯ. นสนิทสฺสน\-สปฺปฏิฆํ นมคฺคา\-ธิปติํ ธมฺมํ ปฏิจฺจ.}

\proseitem{81}{นสนิทสฺสน\-สปฺปฏิฆํ นอนุปฺปนฺนํ ธมฺมํ ปฏิจฺจ ฯเปฯ. นสนิทสฺสน\-สปฺปฏิฆํ นอุปฺปาทิํ ธมฺมํ ปฏิจฺจ.}

\proseitem{82}{นสนิทสฺสน\-สปฺปฏิฆํ นอตีตํ ธมฺมํ ปฏิจฺจ ฯเปฯ. นสนิทสฺสน\-สปฺปฏิฆํ นอนาคตํ ธมฺมํ ปฏิจฺจ.}

\proseitem{83}{นสนิทสฺสน\-สปฺปฏิฆํ นอตีตา\-รมฺมณํ ธมฺมํ ปฏิจฺจ ฯเปฯ. นสนิทสฺสน\-สปฺปฏิฆํ นอนาคตา\-รมฺมณํ ธมฺมํ ปฏิจฺจ ฯเปฯ. นสนิทสฺสน\-สปฺปฏิฆํ นปจฺจุปฺปนฺนา\-รมฺมณํ ธมฺมํ ปฏิจฺจ.}

\proseitem{84}{นสนิทสฺสน\-สปฺปฏิฆํ นอชฺฌตฺตํ ธมฺมํ ปฏิจฺจ ฯเปฯ. นสนิทสฺสน\-สปฺปฏิฆํ นพหิทฺธา ธมฺมํ ปฏิจฺจ.}

\proseitem{85}{นสนิทสฺสน\-สปฺปฏิฆํ นอชฺฌตฺตา\-รมฺมณํ ธมฺมํ ปฏิจฺจ นสนิทสฺสน\-สปฺปฏิโฆ นอชฺฌตฺตา\-รมฺมโณ ธมฺโม อุปฺปชฺชติ เหตุปจฺจยา.  นสนิทสฺสน\-สปฺปฏิฆํ นอชฺฌตฺตา\-รมฺมณํ ธมฺมํ ปฏิจฺจ นอนิทสฺสน\-สปฺปฏิโฆ นอชฺฌตฺตา\-รมฺมโณ ธมฺโม อุปฺปชฺชติ เหตุปจฺจยา.}

\prose{เหตุยา ติํส ฯเปฯ อวิคเต ติํส. (สพฺพตฺถ วิตฺถาโร.)}

\proseitem{86}{นสนิทสฺสน\-สปฺปฏิฆํ นพหิทฺธา\-รมฺมณํ ธมฺมํ ปฏิจฺจ นสนิทสฺสน\-สปฺปฏิโฆ นพหิทฺธา\-รมฺมโณ ธมฺโม อุปฺปชฺชติ เหตุปจฺจยา.  นสนิทสฺสน\-สปฺปฏิฆํ นพหิทฺธา\-รมฺมณํ ธมฺมํ ปฏิจฺจ นอนิทสฺสน\-สปฺปฏิโฆ นพหิทฺธา\-รมฺมโณ ธมฺโม อุปฺปชฺชติ เหตุปจฺจยา.  นสนิทสฺสน\-สปฺปฏิฆํ นพหิทฺธา\-รมฺมณํ ธมฺมํ ปฏิจฺจ นอนิทสฺสน\-อปฺปฏิโฆ นพหิทฺธา\-รมฺมโณ ธมฺโม อุปฺปชฺชติ เหตุปจฺจยา ฯเปฯ. ฉ.}

\prose{นอนิทสฺสน\-สปฺปฏิฆํ นพหิทฺธา\-รมฺมณํ ธมฺมํ ปฏิจฺจ นอนิทสฺสน\-สปฺปฏิโฆ นพหิทฺธา\-รมฺมโณ ธมฺโม อุปฺปชฺชติ เหตุปจฺจยา. ฉ.}

\prose{นอนิทสฺสน\-อปฺปฏิฆํ นพหิทฺธา\-รมฺมณํ ธมฺมํ ปฏิจฺจ นอนิทสฺสน\-อปฺปฏิโฆ นพหิทฺธา\-รมฺมโณ ธมฺโม อุปฺปชฺชติ เหตุปจฺจยา. ฉ.}

\csromanheader{2}{22. สนิทสฺสน\-ตฺติก 2-21. เวทนาตฺติกาทิ}
\prose{\csromanfolio{157}นสนิทสฺสน\-สปฺปฏิฆํ นพหิทฺธา\-รมฺมณญฺจ นอนิทสฺสน\-อปฺปฏิฆํ นพหิทฺธา\-รมฺมณญฺจ ธมฺมํ ปฏิจฺจ นสนิทสฺสน\-สปฺปฏิโฆ นพหิทฺธา\-รมฺมโณ ธมฺโม อุปฺปชฺชติ เหตุปจฺจยา. ฉ.}

\prose{นอนิทสฺสน\-สปฺปฏิฆํ นพหิทฺธา\-รมฺมณญฺจ นอนิทสฺสน\-อปฺปฏิฆํ นพหิทฺธา\-รมฺมณญฺจ ธมฺมํ ปฏิจฺจ นสนิทสฺสน\-สปฺปฏิโฆ นพหิทฺธา\-รมฺมโณ ธมฺโม อุปฺปชฺชติ เหตุปจฺจยา. ฉ. (สํขิตฺตํ.)}

\prose{เหตุยา ติํส, อารมฺมเณ นว ฯเปฯ อวิคเต ติํส.}

\prosewithcloser{(สหชาตวาเรปิ ปจฺจยวาเรปิ นิสฺสยวาเรปิ สํสฏฺฐ\-วาเรปิ สมฺปยุตฺตวาเรปิ ปญฺหาวาเรปิ วิตฺถาโร.)}{ธมฺม\-ปจฺจนีเย ติกติก\-ปฏฺฐานํ นิฏฺฐิตํ.}
\pitaka{อภิธมฺม\-ปิฏก}
\namakkaram{นโม ตสฺส ภควโต อรหโต สมฺมา\-สมฺพุทฺธสฺส.}
\csromanheader{2}{1. \textbf{เหตุทุก 1-5. สเหตุกทุกาทิ}}
\proseitem{1}{\csromanfolio{159}นเหตุํ นสเหตุกํ ธมฺมํ ปฏิจฺจ นเหตุ นสเหตุโก ธมฺโม อุปฺปชฺชติ เหตุปจฺจยา. (สํขิตฺตํ.)}

\prose{เหตุยา ตีณิ, อารมฺมเณ เอกํ ฯเปฯ อวิคเต ตีณิ.}

\proseitem{2}{นเหตุํ นสเหตุกํ ธมฺมํ ปฏิจฺจ นเหตุ นสเหตุโก ธมฺโม อุปฺปชฺชติ นเหตุปจฺจยา. (1)}

\prose{นเหตุํ นสเหตุกํ ธมฺมํ ปฏิจฺจ นเหตุ นสเหตุโก ธมฺโม อุปฺปชฺชติ นอารมฺมณ\-ปจฺจยา. (สํขิตฺตํ.)}

\prose{นเหตุยา เอกํ, นอารมฺมเณ ตีณิ, นอธิปติยา ตีณิ.}

\prose{เหตุอารมฺมณ}

\prose{3. นนเหตุ นสเหตุโก ธมฺโม นเหตุสฺส นสเหตุกสฺส ธมฺมสฺส}

\prose{นเหตุ นสเหตุโก ธมฺโม นเหตุสฺส นสเหตุกสฺส ธมฺมสฺส อารมฺมณ\-ปจฺจเยน ปจฺจโย.  นเหตุ นสเหตุโก ธมฺโม นนเหตุสฺส นสเหตุกสฺส ธมฺมสฺส อารมฺมณ\-ปจฺจเยน ปจฺจโย.}

\prose{เหตุยา เอกํ, อารมฺมเณ จตฺตาริ ฯเปฯ อวิคเต จตฺตาริ.}

\csromanpalirecto
\gambhira{ธมฺม\-ปจฺจนีย ทุกทุก\-ปฏฺฐาน\-ปาฬิ}
\proseitem{4}{\csromanfolio{160}นเหตุํ นอเหตุกํ ธมฺมํ ปฏิจฺจ นเหตุ นอเหตุโก ธมฺโม อุปฺปชฺชติ เหตุปจฺจยา.  นเหตุํ นอเหตุกํ ธมฺมํ ปฏิจฺจ นนเหตุ นอเหตุโก ธมฺโม อุปฺปชฺชติ เหตุปจฺจยา.  นเหตุํ นอเหตุกํ ธมฺมํ ปฏิจฺจ นเหตุ นอเหตุโก จ นนเหตุ นอเหตุโก จ ธมฺมา อุปฺปชฺชนฺติ เหตุปจฺจยา. (3)}

\prose{นนเหตุํ นอเหตุกํ ธมฺมํ ปฏิจฺจ นนเหตุ นอเหตุโก ธมฺโม อุปฺปชฺชติ เหตุปจฺจยา.  นนเหตุํ นอเหตุกํ ธมฺมํ ปฏิจฺจ นเหตุ นอเหตุโก ธมฺโม อุปฺปชฺชติ เหตุปจฺจยา.  นนเหตุํ นอเหตุกํ ธมฺมํ ปฏิจฺจ นเหตุ นอเหตุโก จ นนเหตุ นอเหตุโก จ ธมฺมา อุปฺปชฺชนฺติ เหตุปจฺจยา. (3)}

\prose{นเหตุํ นอเหตุกญฺจ นนเหตุํ นอเหตุกญฺจ ธมฺมํ ปฏิจฺจ นเหตุ นอเหตุโก ธมฺโม อุปฺปชฺชติ เหตุปจฺจยา.  นเหตุํ นอเหตุกญฺจ นนเหตุํ นอเหตุกญฺจ ธมฺมํ ปฏิจฺจ นนเหตุ นอเหตุโก ธมฺโม อุปฺปชฺชติ เหตุปจฺจยา.  นเหตุํ นอเหตุกญฺจ นนเหตุํ นอเหตุกญฺจ ธมฺมํ ปฏิจฺจ นเหตุ นอเหตุโก จ นนเหตุ นอเหตุโก จ ธมฺมา อุปฺปชฺชนฺติ เหตุปจฺจยา. (3) (สํขิตฺตํ.)}

\prose{เหตุยา นว, อารมฺมเณ นว ฯเปฯ อวิคเต นว. (สพฺพตฺถ นว.)}

\proseitem{5}{นเหตุํ นเหตุ\-สมฺปยุตฺตํ ธมฺมํ ปฏิจฺจ นเหตุ นเหตุ\-สมฺปยุตฺโต ธมฺโม อุปฺปชฺชติ เหตุปจฺจยา. ตีณิ.}

\proseitem{6}{นเหตุํ นเหตุ\-วิปฺปยุตฺตํ ธมฺมํ ปฏิจฺจ นเหตุ นเหตุ\-วิปฺปยุตฺโต ธมฺโม อุปฺปชฺชติ เหตุปจฺจยา. ตีณิ.}

\prose{นนเหตุํ นเหตุ\-วิปฺปยุตฺตํ ธมฺมํ ปฏิจฺจ นนเหตุ นเหตุ\-วิปฺปยุตฺโต ธมฺโม อุปฺปชฺชติ เหตุปจฺจยา. ตีณิ.}

\prose{นเหตุํ นเหตุ\-วิปฺปยุตฺตญฺจ นนเหตุํ นเหตุ\-วิปฺปยุตฺตญฺจ ธมฺมํ ปฏิจฺจ นเหตุ นเหตุ\-วิปฺปยุตฺโต ธมฺโม อุปฺปชฺชติ เหตุปจฺจยา. ตีณิ.}

\prose{เหตุยา นว. (สพฺพตฺถ วิตฺถาโร.)}

\csromanmarkh{1. เหตุทุก}\csromanmarkhii{6. จูฬนฺตรทุก}\csromanheadernomark{2}{1. เหตุทุก 6. จูฬนฺตรทุก}
\proseitem{7}{\csromanfolio{161}นเหตุํ นเหตุญฺเจว นอเหตุกญฺจ ธมฺมํ ปฏิจฺจ นเหตุ นเหตุ เจว นอเหตุโก จ ธมฺโม อุปฺปชฺชติ เหตุปจฺจยา. (ยาว ปญฺหาวาเรปิ เอกํ.)}

\prose{นนเหตุํ นอเหตุกญฺเจว นน จ เหตุํ ธมฺมํ ปฏิจฺจ นนเหตุ นอเหตุโก เจว นน จ เหตุ ธมฺโม อุปฺปชฺชติ เหตุปจฺจยา.}

\prose{นเหตุํ นเหตุญฺเจว นเหตุ\-วิปฺปยุตฺตญฺจ ธมฺมํ ปฏิจฺจ นเหตุ นเหตุ เจว นเหตุ\-วิปฺปยุตฺโต จ ธมฺโม อุปฺปชฺชติ เหตุปจฺจยา.}

\prose{นนเหตุํ นเหตุ\-วิปฺปยุตฺต\-ญฺเจว นน จ เหตุํ ธมฺมํ ปฏิจฺจ นนเหตุ นเหตุ\-วิปฺปยุตฺโต เจว นน จ เหตุ ธมฺโม อุปฺปชฺชติ เหตุปจฺจยา.}

\prose{นเหตุํ นสเหตุกํ ธมฺมํ ปฏิจฺจ นเหตุ นสเหตุโก ธมฺโม อุปฺปชฺชติ เหตุปจฺจยา.}

\prose{นเหตุํ นอเหตุกํ ธมฺมํ ปฏิจฺจ นเหตุ นอเหตุโก ธมฺโม อุปฺปชฺชติ เหตุปจฺจยา. (สพฺพตฺถ เอกํ.)}

\csromanmarkh{1. เหตุทุก}\csromanmarkhii{6. จูฬนฺตรทุก}\csromanheadernomark{2}{1. เหตุทุก 6. จูฬนฺตรทุก}
\proseitem{8}{นเหตุํ นอปฺปจฺจยํ ธมฺมํ ปฏิจฺจ นเหตุ นอปฺปจฺจโย ธมฺโม อุปฺปชฺชติ เหตุปจฺจยา.  นเหตุํ นอปฺปจฺจยํ ธมฺมํ ปฏิจฺจ นนเหตุ นอปฺปจฺจโย ธมฺโม อุปฺปชฺชติ เหตุปจฺจยา.  นเหตุํ นอปฺปจฺจยํ ธมฺมํ ปฏิจฺจ นเหตุ นอปฺปจฺจโย จ นนเหตุ นอปฺปจฺจโย จ ธมฺมา อุปฺปชฺชนฺติ เหตุปจฺจยา. (3)}

\prose{นนเหตุํ นอปฺปจฺจยํ ธมฺมํ ปฏิจฺจ นนเหตุ นอปฺปจฺจโย ธมฺโม อุปฺปชฺชติ เหตุปจฺจยา.  นนเหตุํ นอปฺปจฺจยํ ธมฺมํ ปฏิจฺจ นเหตุ นอปฺปจฺจโย ธมฺโม อุปฺปชฺชติ เหตุปจฺจยา.  นนเหตุํ นอปฺปจฺจยํ ธมฺมํ ปฏิจฺจ นเหตุ นอปฺปจฺจโย จ นนเหตุ นอปฺปจฺจโย จ ธมฺมา อุปฺปชฺชนฺติ เหตุปจฺจยา. (3)}

\prose{นเหตุํ นอปฺปจฺจยญฺจ นนเหตุํ นอปฺปจฺจยญฺจ ธมฺมํ ปฏิจฺจ นเหตุ นอปฺปจฺจโย ธมฺโม อุปฺปชฺชติ เหตุปจฺจยา.  นเหตุํ นอปฺปจฺจยญฺจ นนเหตุํ นอปฺปจฺจยญฺจ ธมฺมํ ปฏิจฺจ นนเหตุ นอปฺปจฺจโย ธมฺโม อุปฺปชฺชติ เหตุปจฺจยา.  นเหตุํ นอปฺปจฺจยญฺจ นนเหตุํ นอปฺปจฺจยญฺจ ธมฺมํ ปฏิจฺจ นเหตุ นอปฺปจฺจโย จ นนเหตุ นอปฺปจฺจโย จ ธมฺมา อุปฺปชฺชนฺติ เหตุปจฺจยา. (3)}

\prose{\csromanfolio{162}เหตุยา นว.}

\prose{นเหตุํ นอสงฺขตํ ธมฺมํ ปฏิจฺจ.}

\prose{นเหตุํ นสนิทสฺสนํ ธมฺมํ ปฏิจฺจ.}

\prose{นเหตุํ นสปฺปฏิฆํ ธมฺมํ ปฏิจฺจ.}

\prose{นเหตุํ นอปฺปฏิฆํ ธมฺมํ ปฏิจฺจ.}

\prose{นเหตุํ นรูปิํ ธมฺมํ ปฏิจฺจ.}

\prose{นเหตุํ นอรูปิํ ธมฺมํ ปฏิจฺจ.}

\prose{นเหตุํ นโลกิยํ ธมฺมํ ปฏิจฺจ ฯเปฯ. นเหตุํ นโลกุตฺตรํ ธมฺมํ ปฏิจฺจ.}

\prose{นเหตุํ นเกนจิ วิญฺเญยฺยํ ธมฺมํ ปฏิจฺจ ฯเปฯ. นเหตุํ นนเกนจิ วิญฺเญยฺยํ ธมฺมํ ปฏิจฺจ.}

\csromanheader{2}{1. \textbf{เหตุทุก 13-}18. อาสวโคจฺฉก}
\proseitem{9}{นเหตุํ โนอาสวํ ธมฺมํ ปฏิจฺจ.}

\prose{นเหตุํ นโนอาสวํ ธมฺมํ ปฏิจฺจ.}

\prose{นเหตุํ นสาสวํ ธมฺมํ ปฏิจฺจ ฯเปฯ. นเหตุํ นอนาสวํ ธมฺมํ ปฏิจฺจ.}

\prose{นเหตุํ นอาสว\-สมฺปยุตฺตํ ธมฺมํ ปฏิจฺจ.}

\prose{นเหตุํ นอาสว\-วิปฺปยุตฺตํ ธมฺมํ ปฏิจฺจ.}

\prose{นเหตุํ โนอาสวญฺเจว นอนาสวญฺจ ธมฺมํ ปฏิจฺจ.}

\prose{นเหตุํ นอนาสวญฺเจว นโน จ อาสวํ ธมฺมํ ปฏิจฺจ.}

\prose{นเหตุํ นอาสวญฺเจว นอาสว\-วิปฺปยุตฺตญฺจ ธมฺมํ ปฏิจฺจ.}

\prose{นเหตุํ นอาสว\-วิปฺปยุตฺต\-ญฺเจว นโน จ อาสวํ ธมฺมํ ปฏิจฺจ.}

\prose{นเหตุํ อาสว\-วิปฺปยุตฺตํ นสาสวํ ธมฺมํ ปฏิจฺจ ฯเปฯ. นเหตุํ อาสว\-วิปฺปยุตฺตํ นอนาสวํ ธมฺมํ ปฏิจฺจ.}

\csromanheader{2}{1. เหตุทุก 19-53. สญฺโญชนา\-ทิทุก}
\proseitem{10}{นเหตุํ โนสญฺโญชนํ ธมฺมํ ปฏิจฺจ ฯเปฯ. นเหตุํ โนคนฺถํ ธมฺมํ ปฏิจฺจ. นเหตุํ โนโอฆํ ธมฺมํ ปฏิจฺจ ฯเปฯ. นเหตุํ โนโยคํ ธมฺมํ ปฏิจฺจ.}

\csromanmarkh{1. เหตุทุก}\csromanmarkhii{82. ปิฏฺฐิทุก}\csromanheadernomark{2}{1. เหตุทุก 82. ปิฏฺฐิทุก}
\prose{\csromanfolio{163}นเหตุํ โนนีวรณํ ธมฺมํ ปฏิจฺจ.}

\prose{นเหตุํ โนปรามาสํ ธมฺมํ ปฏิจฺจ.}

\csromanmarkh{1. เหตุทุก}\csromanmarkhii{54. มหนฺตรทุก}\csromanheadernomark{2}{1. เหตุทุก 54. มหนฺตรทุก}
\prose{11. นเหตุํ นสารมฺมณํ ธมฺมํ ปฏิจฺจ. (สํขิตฺตํ.)}

\prose{นเหตุํ นจิตฺตํ ธมฺมํ ปฏิจฺจ.}

\prose{นเหตุํ นเจตสิกํ ธมฺมํ ปฏิจฺจ.}

\prose{นเหตุํ นจิตฺต\-สมฺปยุตฺตํ ธมฺมํ ปฏิจฺจ ฯเปฯ. นเหตุํ นจิตฺต\-สํสฏฺฐํ ธมฺมํ ปฏิจฺจ ฯเปฯ.}

\prose{นเหตุํ นจิตฺต\-สมุฏฺฐานํ ธมฺมํ ปฏิจฺจ.}

\prose{นเหตุํ นจิตฺต\-สหภุํ ธมฺมํ ปฏิจฺจ ฯเปฯ. นเหตุํ นจิตฺตา\-นุปริวตฺติํ ธมฺมํ ปฏิจฺจ.}

\prose{นเหตุํ นจิตฺต\-สํสฏฺฐ\-สมุฏฺฐานํ ธมฺมํ ปฏิจฺจ ฯเปฯ. นเหตุํ นจิตฺต\-สํสฏฺฐ\-สมุฏฺฐาน\-สหภุํ ธมฺมํ ปฏิจฺจ ฯเปฯ. นเหตุํ นจิตฺต\-สํสฏฺฐ\-สมุฏฺฐานา\-นุ\-ปริวตฺติํ ธมฺมํ ปฏิจฺจ.}

\prose{นเหตุํ นอชฺฌตฺติกํ ธมฺมํ ปฏิจฺจ.}

\prose{นเหตุํ นพาหิรํ ธมฺมํ ปฏิจฺจ.}

\prose{นเหตุํ นอุปาทาธมฺมํ ปฏิจฺจ.}

\prose{นเหตุํ นอุปาทินฺนํ ธมฺมํ ปฏิจฺจ.}

\prose{นเหตุํ นอนุปาทินฺนํ ธมฺมํ ปฏิจฺจ.}

\prose{นเหตุํ โนอุปาทานํ ธมฺมํ ปฏิจฺจ ฯเปฯ. นเหตุํ โนกิเลสํ ธมฺมํ ปฏิจฺจ ฯเปฯ.}

\csromanmarkh{1. เหตุทุก}\csromanmarkhii{82. ปิฏฺฐิทุก}\csromanheadernomark{2}{1. เหตุทุก 82. ปิฏฺฐิทุก}
\proseitem{12}{นเหตุํ นทสฺสเนน ปหาตพฺพํ ธมฺมํ ปฏิจฺจ.}

\prose{นเหตุํ นนทสฺสเนน ปหาตพฺพํ ธมฺมํ ปฏิจฺจ.}

\proseitem{13}{นเหตุํ นภาวนาย ปหาตพฺพํ ธมฺมํ ปฏิจฺจ ฯเปฯ. นเหตุํ นนภาวนาย ปหาตพฺพํ ธมฺมํ ปฏิจฺจ ฯเปฯ. นเหตุํ นทสฺสเนน ปหาตพฺพ\-เหตุกํ ธมฺมํ ปฏิจฺจ ฯเปฯ. นเหตุํ นนทสฺสเนน ปหาตพฺพ\-เหตุกํ ธมฺมํ ปฏิจฺจ ฯเปฯ. นเหตุํ นภาวนาย ปหาตพฺพ\-เหตุกํ ธมฺมํ ปฏิจฺจ.}

\prose{นเหตุํ นนภาวนาย ปหาตพฺพ\-เหตุกํ ธมฺมํ ปฏิจฺจ.}

\proseitem{14}{\csromanfolio{164}นเหตุํ นสวิตกฺกํ ธมฺมํ ปฏิจฺจ ฯเปฯ. นเหตุํ นอวิตกฺกํ ธมฺมํ ปฏิจฺจ ฯเปฯ. นเหตุํ นสวิจารํ ธมฺมํ ปฏิจฺจ ฯเปฯ. นเหตุํ นอวิจารํ ธมฺมํ ปฏิจฺจ ฯเปฯ. นเหตุํ นสปฺปีติกํ ธมฺมํ ปฏิจฺจ ฯเปฯ. นเหตุํ นอปฺปีติกํ ธมฺมํ ปฏิจฺจ ฯเปฯ. นเหตุํ นปีติสหคตํ ธมฺมํ ปฏิจฺจ ฯเปฯ. นเหตุํ นนปี\-ติ\-สหคตํ ธมฺมํ ปฏิจฺจ ฯเปฯ. นเหตุํ นสุขสหคตํ ธมฺมํ ปฏิจฺจ ฯเปฯ. นเหตุํ นนสุข\-สหคตํ ธมฺมํ ปฏิจฺจ ฯเปฯ. นเหตุํ นอุเปกฺขาสหคตํ ธมฺมํ ปฏิจฺจ ฯเปฯ. นเหตุํ นนอุเปกฺขาสหคตํ ธมฺมํ ปฏิจฺจ ฯเปฯ.}

\proseitem{15}{นเหตุํ นกามาวจรํ ธมฺมํ ปฏิจฺจ ฯเปฯ. นเหตุํ นนกามา\-วจรํ ธมฺมํ ปฏิจฺจ ฯเปฯ. นเหตุํ นรูปาวจรํ ธมฺมํ ปฏิจฺจ ฯเปฯ. นเหตุํ นนรูปาวจรํ ธมฺมํ ปฏิจฺจ ฯเปฯ. นเหตุํ นอรูปาวจรํ ธมฺมํ ปฏิจฺจ ฯเปฯ. นเหตุํ นนอรูปาวจรํ ธมฺมํ ปฏิจฺจ ฯเปฯ.}

\proseitem{16}{นเหตุํ นปริยาปนฺนํ ธมฺมํ ปฏิจฺจ ฯเปฯ. นเหตุํ นอปริยาปนฺนํ ธมฺมํ ปฏิจฺจ ฯเปฯ. นเหตุํ นนิยฺยานิกํ ธมฺมํ ปฏิจฺจ ฯเปฯ. นเหตุํ นอนิยฺยานิกํ ธมฺมํ ปฏิจฺจ ฯเปฯ. นเหตุํ นนิยตํ ธมฺมํ ปฏิจฺจ ฯเปฯ. นเหตุํ นอนิยตํ ธมฺมํ ปฏิจฺจ ฯเปฯ.}

\prose{นเหตุํ นสอุตฺตรํ ธมฺมํ ปฏิจฺจ ฯเปฯ. นเหตุํ นอนุตฺตรํ ธมฺมํ ปฏิจฺจ ฯเปฯ.}

\proseitem{17}{นเหตุํ นสรณํ ธมฺมํ ปฏิจฺจ นเหตุ นสรโณ ธมฺโม อุปฺปชฺชติ เหตุปจฺจยา. (นว.)}

\prose{นเหตุํ นอรณํ ธมฺมํ ปฏิจฺจ นเหตุ นอรโณ ธมฺโม อุปฺปชฺชติ เหตุปจฺจยา.  นเหตุํ นอรณํ ธมฺมํ ปฏิจฺจ นนเหตุ นอรโณ ธมฺโม อุปฺปชฺชติ เหตุปจฺจยา.  นเหตุํ นอรณํ ธมฺมํ ปฏิจฺจ นเหตุ นอรโณ จ นนเหตุ นอรโณ จ ธมฺมา อุปฺปชฺชนฺติ เหตุปจฺจยา. (3)}

\prose{นนเหตุํ นอรณํ ธมฺมํ ปฏิจฺจ นนเหตุ นอรโณ ธมฺโม อุปฺปชฺชติ เหตุปจฺจยา.  นนเหตุํ นอรณํ ธมฺมํ ปฏิจฺจ นเหตุ นอรโณ ธมฺโม อุปฺปชฺชติ เหตุปจฺจยา.  นนเหตุํ นอรณํ ธมฺมํ ปฏิจฺจ นเหตุ นอรโณ จ นนเหตุ นอรโณ จ ธมฺมา อุปฺปชฺชนฺติ เหตุปจฺจยา. (3)}

\prose{นเหตุํ นอรณญฺจ นนเหตุํ นอรณญฺจ ธมฺมํ ปฏิจฺจ นเหตุ นอรโณ ธมฺโม อุปฺปชฺชติ เหตุปจฺจยา.  นเหตุํ นอรณญฺจ นนเหตุํ นอรณญฺจ ธมฺมํ ปฏิจฺจ นนเหตุ นอรโณ ธมฺโม อุปฺปชฺชติ เหตุปจฺจยา.  นเหตุํ นอรณญฺจ นนเหตุํ นอรณญฺจ ธมฺมํ ปฏิจฺจ นเหตุ นอรโณ จ นนเหตุ}

\vspace{\tipitakaparskip}
\csromancenter{(สํขิตฺตํ.) เหตุยา นว. (สพฺพตฺถ นว.)}
\csromanmarkh{2. สเหตุกาทิทุก}\csromanmarkhii{1. เหตุทุก}\csromanheadernomark{2}{2. สเหตุกาทิทุก 1. เหตุทุก}
\proseitem{18}{\csromanfolio{165}นสเหตุกํ นเหตุํ ธมฺมํ ปฏิจฺจ นสเหตุโก นเหตุ ธมฺโม อุปฺปชฺชติ เหตุปจฺจยา.  นสเหตุกํ นเหตุํ ธมฺมํ ปฏิจฺจ นอเหตุโก นเหตุ ธมฺโม อุปฺปชฺชติ เหตุปจฺจยา.  นสเหตุกํ นเหตุํ ธมฺมํ ปฏิจฺจ นสเหตุโก นเหตุ จ นอเหตุโก นเหตุ จ ธมฺมา อุปฺปชฺชนฺติ เหตุปจฺจยา. (3)}

\prose{นอเหตุกํ นเหตุํ ธมฺมํ ปฏิจฺจ นอเหตุโก นเหตุ ธมฺโม อุปฺปชฺชติ เหตุปจฺจยา.  นอเหตุกํ นเหตุํ ธมฺมํ ปฏิจฺจ นสเหตุโก นเหตุ ธมฺโม อุปฺปชฺชติ เหตุปจฺจยา.  นอเหตุกํ นเหตุํ ธมฺมํ ปฏิจฺจ นสเหตุโก นเหตุ จ นอเหตุโก นเหตุ จ ธมฺมา อุปฺปชฺชนฺติ เหตุปจฺจยา. (3)}

\prose{นสเหตุกํ นเหตุญฺจ นอเหตุกํ นเหตุญฺจ ธมฺมํ ปฏิจฺจ นสเหตุโก นเหตุ ธมฺโม อุปฺปชฺชติ เหตุปจฺจยา.  นสเหตุกํ นเหตุญฺจ นอเหตุกํ นเหตุญฺจ ธมฺมํ ปฏิจฺจ นอเหตุโก นเหตุ ธมฺโม อุปฺปชฺชติ เหตุปจฺจยา.  นสเหตุกํ นเหตุญฺจ นอเหตุกํ นเหตุญฺจ ธมฺมํ ปฏิจฺจ นสเหตุโก นเหตุ จ นอเหตุโก นเหตุ จ ธมฺมา อุปฺปชฺชนฺติ เหตุปจฺจยา. (3) (สํขิตฺตํ.)}

\prose{เหตุยา นว. (สพฺพตฺถ วิตฺถาโร.)}

\proseitem{19}{นอเหตุกํ นนเหตุํ ธมฺมํ ปฏิจฺจ นอเหตุโก นนเหตุ ธมฺโม อุปฺปชฺชติ เหตุปจฺจยา. (สํขิตฺตํ.)}

\prose{เหตุยา เอกํ ฯเปฯ อวิคเต เอกํ. (สพฺพตฺถ เอกํ.)}

\proseitem{20}{นเหตุ\-สมฺปยุตฺตํ นเหตุํ ธมฺมํ ปฏิจฺจ.}

\prose{นเหตุ\-วิปฺปยุตฺตํ นนเหตุํ ธมฺมํ ปฏิจฺจ.}

\proseitem{21}{นเหตุญฺเจว นอเหตุกญฺจ นเหตุํ ธมฺมํ ปฏิจฺจ ฯเปฯ. นอเหตุกญฺเจว นน จ เหตุํ ธมฺมํ ปฏิจฺจ ฯเปฯ. นเหตุญฺเจว นเหตุ\-วิปฺปยุตฺตญฺจ นเหตุํ ธมฺมํ ปฏิจฺจ ฯเปฯ. นเหตุ\-วิปฺปยุตฺต\-ญฺเจว นน จ เหตุํ ธมฺมํ ปฏิจฺจ.}

\proseitem{22}{\csromanfolio{166}นเหตุํ นสเหตุกํ นเหตุํ ธมฺมํ ปฏิจฺจ ฯเปฯ. นเหตุํ นอเหตุกํ นเหตุํ ธมฺมํ ปฏิจฺจ.}

\csromanmarkh{7. จูฬนฺตรทุก}\csromanmarkhii{1. เหตุทุก}\csromanheadernomark{2}{7. \textbf{จูฬนฺตรทุก 1. เหตุทุก}}
\proseitem{23}{นอปฺปจฺจยํ นเหตุํ ธมฺมํ ปฏิจฺจ ฯเปฯ. นอสงฺขตํ นเหตุํ ธมฺมํ ปฏิจฺจ.}

\prose{นสนิทสฺสนํ นเหตุํ ธมฺมํ ปฏิจฺจ.}

\prose{นสปฺปฏิฆํ นเหตุํ ธมฺมํ ปฏิจฺจ ฯเปฯ. นอปฺปฏิฆํ นเหตุํ ธมฺมํ ปฏิจฺจ.}

\prose{นรูปิํ นเหตุํ ธมฺมํ ปฏิจฺจ ฯเปฯ. นอรูปิํ นเหตุํ ธมฺมํ ปฏิจฺจ.}

\prose{นโลกิยํ นเหตุํ ธมฺมํ ปฏิจฺจ ฯเปฯ. นโลกุตฺตรํ นเหตุํ ธมฺมํ ปฏิจฺจ.}

\prose{นเกนจิ วิญฺเญยฺยํ นเหตุํ ธมฺมํ ปฏิจฺจ ฯเปฯ. นนเกนจิ วิญฺเญยฺยํ นเหตุํ ธมฺมํ ปฏิจฺจ.}

\csromanmarkh{14. อาสวโคจฺฉก}\csromanmarkhii{1. เหตุทุก}\csromanheadernomark{2}{14. \textbf{อาสวโคจฺฉก 1. เหตุทุก}}
\proseitem{24}{โนอาสวํ นเหตุํ ธมฺมํ ปฏิจฺจ ฯเปฯ. นโนอาสวํ นเหตุํ ธมฺมํ ปฏิจฺจ.}

\prose{นสาสวํ นเหตุํ ธมฺมํ ปฏิจฺจ ฯเปฯ. นอนาสวํ นเหตุํ ธมฺมํ ปฏิจฺจ.}

\prose{นอาสว\-สมฺปยุตฺตํ นเหตุํ ธมฺมํ ปฏิจฺจ ฯเปฯ. นอาสว\-วิปฺปยุตฺตํ นเหตุํ ธมฺมํ ปฏิจฺจ.}

\prose{นอาสวญฺเจว นอนาสวํ นเหตุํ ธมฺมํ ปฏิจฺจ ฯเปฯ. นอนาสวญฺเจว นโน จ อาสวํ นเหตุํ ธมฺมํ ปฏิจฺจ.}

\prose{นอาสวญฺเจว นอาสว\-วิปฺปยุตฺตญฺจ นเหตุํ ธมฺมํ ปฏิจฺจ.}

\prose{นอาสว\-วิปฺปยุตฺต\-ญฺเจว นโน จ อาสวํ นเหตุํ ธมฺมํ ปฏิจฺจ ฯเปฯ. อาสว\-วิปฺปยุตฺตํ นสาสวํ นเหตุํ ธมฺมํ ปฏิจฺจ ฯเปฯ. อาสว\-วิปฺปยุตฺตํ นอนาสวํ นเหตุํ ธมฺมํ ปฏิจฺจ.}

\csromanmarkh{20. สญฺโญชนา\-ทิทุก}\csromanmarkhii{1. เหตุทุก}\csromanheadernomark{2}{20. สญฺโญชนา\-ทิทุก 1. เหตุทุก}
\proseitem{25}{โนสญฺโญชนํ นเหตุํ ธมฺมํ ปฏิจฺจ.}

\proseitem{26}{โนคนฺถํ นเหตุํ ธมฺมํ ปฏิจฺจ.}

\csromanmarkh{83. ปิฏฺฐิทุก}\csromanmarkhii{1. เหตุทุก}\csromanheadernomark{2}{83. ปิฏฺฐิทุก 1. เหตุทุก}
\proseitem{27}{\csromanfolio{167}โนโอฆํ นเหตุํ ธมฺมํ ปฏิจฺจ ฯเปฯ. โนโยคํ นเหตุํ ธมฺมํ ปฏิจฺจ.}

\proseitem{28}{โนนีวรณํ นเหตุํ ธมฺมํ ปฏิจฺจ.}

\proseitem{29}{โนปรามาสํ นเหตุํ ธมฺมํ ปฏิจฺจ.}

\csromanmarkh{55. มหนฺตรทุก}\csromanmarkhii{1. เหตุทุก}\csromanheadernomark{2}{55. \textbf{มหนฺตรทุก 1. เหตุทุก}}
\prose{30. นสารมฺมณํ นเหตุํ ธมฺมํ ปฏิจฺจ. (สํขิตฺตํ.)}

\proseitem{31}{นจิตฺตํ นเหตุํ ธมฺมํ ปฏิจฺจ.}

\prose{นเจตสิกํ นเหตุํ ธมฺมํ ปฏิจฺจ.}

\prose{นจิตฺต\-สมฺปยุตฺตํ นเหตุํ ธมฺมํ ปฏิจฺจ ฯเปฯ. นจิตฺต\-สํสฏฺฐํ นเหตุํ ธมฺมํ ปฏิจฺจ.}

\prose{นจิตฺต\-สมุฏฺฐานํ นเหตุํ ธมฺมํ ปฏิจฺจ.}

\proseitem{32}{นจิตฺต\-สหภุํ นเหตุํ ธมฺมํ ปฏิจฺจ ฯเปฯ. นจิตฺตา\-นุปริวตฺติํ นเหตุํ ธมฺมํ ปฏิจฺจ ฯเปฯ. นจิตฺต\-สํสฏฺฐ\-สมุฏฺฐานํ นเหตุํ ธมฺมํ ปฏิจฺจ ฯเปฯ. นจิตฺต\-สํสฏฺฐ\-สมุฏฺฐาน\-สหภุํ นเหตุํ ธมฺมํ ปฏิจฺจ ฯเปฯ. นจิตฺต\-สํสฏฺฐ\-สมุฏฺฐานา\-นุ\-ปริวตฺติํ นเหตุํ ธมฺมํ ปฏิจฺจ.}

\proseitem{33}{นอชฺฌตฺติกํ นเหตุํ ธมฺมํ ปฏิจฺจ ฯเปฯ. นพาหิรํ นเหตุํ ธมฺมํ ปฏิจฺจ.}

\prose{โนอุปาทา นเหตุํ ธมฺมํ ปฏิจฺจ.}

\proseitem{34}{นอุปาทินฺนํ นเหตุํ ธมฺมํ ปฏิจฺจ ฯเปฯ. นอนุปาทินฺนํ นเหตุํ ธมฺมํ ปฏิจฺจ.}

\csromanmarkh{69. อุปาทานทุกาทิ}\csromanmarkhii{1. เหตุทุก}\csromanheadernomark{2}{69. \textbf{อุปาทานทุกาทิ 1. เหตุทุก}}
\proseitem{35}{โนอุปาทานํ นเหตุํ ธมฺมํ ปฏิจฺจ.}

\proseitem{36}{โนกิเลสํ นเหตุํ ธมฺมํ ปฏิจฺจ ฯเปฯ. นโนกิเลสํ นเหตุํ ธมฺมํ ปฏิจฺจ.}

\csromanmarkh{83. ปิฏฺฐิทุก}\csromanmarkhii{1. เหตุทุก}\csromanheadernomark{2}{83. ปิฏฺฐิทุก 1. เหตุทุก}
\proseitem{37}{นทสฺสเนน ปหาตพฺพํ นเหตุํ ธมฺมํ ปฏิจฺจ ฯเปฯ. นนทสฺสเนน ปหาตพฺพํ นเหตุํ ธมฺมํ ปฏิจฺจ ฯเปฯ. นภาวนาย ปหาตพฺพํ นเหตุํ ธมฺมํ \csromanfolio{168}ปฏิจฺจ ฯเปฯ. นนภาวนาย ปหาตพฺพํ นเหตุํ ธมฺมํ ปฏิจฺจ ฯเปฯ. นทสฺสเนน ปหาตพฺพ\-เหตุกํ นเหตุํ ธมฺมํ ปฏิจฺจ ฯเปฯ. นนทสฺสเนน ปหาตพฺพ\-เหตุกํ นเหตุํ ธมฺมํ ปฏิจฺจ ฯเปฯ. นภาวนาย ปหาตพฺพ\-เหตุกํ นเหตุํ ธมฺมํ ปฏิจฺจ ฯเปฯ. นนภาวนาย ปหาตพฺพ\-เหตุกํ นเหตุํ ธมฺมํ ปฏิจฺจ.}

\proseitem{38}{นสวิตกฺกํ นเหตุํ ธมฺมํ ปฏิจฺจ ฯเปฯ. นอวิตกฺกํ นเหตุํ ธมฺมํ ปฏิจฺจ ฯเปฯ. นสวิจารํ นเหตุํ ธมฺมํ ปฏิจฺจ ฯเปฯ. นอวิจารํ นเหตุํ ธมฺมํ ปฏิจฺจ ฯเปฯ. นสปฺปีติกํ นเหตุํ ธมฺมํ ปฏิจฺจ ฯเปฯ. นอปฺปีติกํ นเหตุํ ธมฺมํ ปฏิจฺจ ฯเปฯ. นปีติสหคตํ นเหตุํ ธมฺมํ ปฏิจฺจ ฯเปฯ. นนปี\-ติ\-สหคตํ นเหตุํ ธมฺมํ ปฏิจฺจ ฯเปฯ. นสุขสหคตํ นเหตุํ ธมฺมํ ปฏิจฺจ ฯเปฯ. นนสุข\-สหคตํ นเหตุํ ธมฺมํ ปฏิจฺจ ฯเปฯ. นอุเปกฺขาสหคตํ นเหตุํ ธมฺมํ ปฏิจฺจ ฯเปฯ. นนอุเปกฺขาสหคตํ นเหตุํ ธมฺมํ ปฏิจฺจ.}

\proseitem{39}{นกามาวจรํ นเหตุํ ธมฺมํ ปฏิจฺจ ฯเปฯ. นนกามา\-วจรํ นเหตุํ ธมฺมํ ปฏิจฺจ ฯเปฯ. นรูปาวจรํ นเหตุํ ธมฺมํ ปฏิจฺจ ฯเปฯ. นนรูปาวจรํ นเหตุํ ธมฺมํ ปฏิจฺจ.}

\proseitem{40}{นอรูปาวจรํ นเหตุํ ธมฺมํ ปฏิจฺจ ฯเปฯ. นนอรูปาวจรํ นเหตุํ ธมฺมํ ปฏิจฺจ.}

\prose{นปริยาปนฺนํ นเหตุํ ธมฺมํ ปฏิจฺจ ฯเปฯ. นอปริยาปนฺนํ นเหตุํ ธมฺมํ ปฏิจฺจ.}

\proseitem{41}{นนิยฺยานิกํ นเหตุํ ธมฺมํ ปฏิจฺจ ฯเปฯ. นอนิยฺยานิกํ นเหตุํ ธมฺมํ ปฏิจฺจ ฯเปฯ. นนิยตํ นเหตุํ ธมฺมํ ปฏิจฺจ ฯเปฯ. นอนิยตํ นเหตุํ ธมฺมํ ปฏิจฺจ.}

\prose{นสอุตฺตรํ นเหตุํ ธมฺมํ ปฏิจฺจ ฯเปฯ. นอนุตฺตรํ นเหตุํ ธมฺมํ ปฏิจฺจ.}

\proseitem{42}{นสรณํ นเหตุํ ธมฺมํ ปฏิจฺจ นสรโณ นเหตุ ธมฺโม อุปฺปชฺชติ เหตุปจฺจยา. (1)}

\prose{นอรณํ นเหตุํ ธมฺมํ ปฏิจฺจ นอรโณ นเหตุ ธมฺโม อุปฺปชฺชติ เหตุปจฺจยา.  นอรณํ นเหตุํ ธมฺมํ ปฏิจฺจ นสรโณ นเหตุ ธมฺโม อุปฺปชฺชติ เหตุปจฺจยา.  นอรณํ นเหตุํ ธมฺมํ ปฏิจฺจ นสรโณ นเหตุ จ นอรโณ นเหตุ จ ธมฺมา อุปฺปชฺชนฺติ เหตุปจฺจยา. (3)}

\prose{นสรณํ นเหตุญฺจ นอรณํ นเหตุญฺจ ธมฺมํ ปฏิจฺจ นสรโณ นเหตุ ธมฺโม อุปฺปชฺชติ เหตุปจฺจยา. (1)}

\prose{เหตุยา ปญฺจ ฯเปฯ อวิคเต ปญฺจ. (สพฺพตฺถ วิตฺถาโร.)}

\csromanmarkh{100. สรณทุก}\csromanmarkhii{7. จูฬนฺตรทุก}\csromanheadernomark{2}{100. สรณทุก 7. จูฬนฺตรทุก}
\prose{\csromanfolio{169}43. นสรณํ นนเหตุํ ธมฺมํ ปฏิจฺจ นสรโณ นนเหตุ}

\prose{นอรณํ นนเหตุํ ธมฺมํ ปฏิจฺจ นอรโณ นนเหตุ ธมฺโม อุปฺปชฺชติ เหตุปจฺจยา. (1)}

\prose{เหตุยา เทฺว ฯเปฯ อวิคเต เทฺว. (สพฺพตฺถ วิตฺถาโร.)}

\csromanmarkh{100. สรณทุก}\csromanmarkhii{2. สเหตุกทุกาทิ}\csromanheadernomark{2}{100. \textbf{สรณทุก 2. สเหตุกทุกาทิ}}
\proseitem{44}{นสรณํ นสเหตุกํ ธมฺมํ ปฏิจฺจ ฯเปฯ. นสรณํ นอเหตุกํ ธมฺมํ ปฏิจฺจ ฯเปฯ. นสรณํ นเหตุ\-สมฺปยุตฺตํ ธมฺมํ ปฏิจฺจ.}

\prose{นสรณํ นเหตุ\-วิปฺปยุตฺตํ ธมฺมํ ปฏิจฺจ.}

\proseitem{45}{นสรณํ นเหตุญฺเจว นอเหตุกญฺจ ธมฺมํ ปฏิจฺจ ฯเปฯ. นสรณํ นอเหตุกญฺเจว นน จ เหตุํ ธมฺมํ ปฏิจฺจ ฯเปฯ. นสรณํ นเหตุญฺเจว นเหตุ\-วิปฺปยุตฺตญฺจ ธมฺมํ ปฏิจฺจ ฯเปฯ. นสรณํ นเหตุ\-วิปฺปยุตฺต\-ญฺเจว นน จ เหตุํ ธมฺมํ ปฏิจฺจ.}

\prose{นสรณํ นเหตุํ นสเหตุกํ ธมฺมํ ปฏิจฺจ.}

\prose{นสรณํ นเหตุํ นอเหตุกํ ธมฺมํ ปฏิจฺจ.}

\csromanmarkh{100. สรณทุก}\csromanmarkhii{7. จูฬนฺตรทุก}\csromanheadernomark{2}{100. สรณทุก 7. จูฬนฺตรทุก}
\proseitem{46}{นสรณํ นอปฺปจฺจยํ ธมฺมํ ปฏิจฺจ ฯเปฯ. นสรณํ นอสงฺขตํ ธมฺมํ ปฏิจฺจ.}

\prose{นสรณํ นสนิทสฺสนํ ธมฺมํ ปฏิจฺจ.}

\prose{นสรณํ นสปฺปฏิฆํ ธมฺมํ ปฏิจฺจ.}

\prose{นสรณํ นอปฺปฏิฆํ ธมฺมํ ปฏิจฺจ.}

\prose{นสรณํ นรูปิํ ธมฺมํ ปฏิจฺจ.}

\prose{นสรณํ นอรูปิํ ธมฺมํ ปฏิจฺจ.}

\prose{นสรณํ นโลกุตฺตรํ ธมฺมํ ปฏิจฺจ.}

\prose{นสรณํ นเกนจิ วิญฺเญยฺยํ ธมฺมํ ปฏิจฺจ ฯเปฯ. นสรณํ นนเกนจิ วิญฺเญยฺยํ ธมฺมํ ปฏิจฺจ.}

\csromanmarkh{ธมฺม\-ปจฺจนีย ทุกทุก\-ปฏฺฐาน\-ปาฬิ}\csromanmarkhii{100. สรณทุก 14. อาสวโคจฺฉก}\csromanheadernomark{2}{ธมฺม\-ปจฺจนีย ทุกทุก\-ปฏฺฐาน\-ปาฬิ 100. สรณทุก 14. อาสวโคจฺฉก}
\proseitem{47}{\csromanfolio{170}นสรณํ โนอาสวํ ธมฺมํ ปฏิจฺจ.}

\prose{นสรณํ นโนอาสวํ ธมฺมํ ปฏิจฺจ.}

\prose{นสรณํ นสาสวํ ธมฺมํ ปฏิจฺจ.}

\prose{นสรณํ นอนาสวํ ธมฺมํ ปฏิจฺจ.}

\prose{นสรณํ นอาสว\-สมฺปยุตฺตํ ธมฺมํ ปฏิจฺจ.}

\prose{นสรณํ นอาสว\-วิปฺปยุตฺตํ ธมฺมํ ปฏิจฺจ.}

\prose{นสรณํ นอาสวญฺเจว นอนาสวญฺจ ธมฺมํ ปฏิจฺจ.}

\prose{นสรณํ นอนาสวญฺเจว นโน จ อาสวํ ธมฺมํ ปฏิจฺจ.}

\prose{นสรณํ นอาสวญฺเจว นอาสว\-วิปฺปยุตฺตญฺจ ธมฺมํ ปฏิจฺจ ฯเปฯ. นสรณํ นอาสว\-วิปฺปยุตฺต\-ญฺเจว นโน จ อาสวํ ธมฺมํ ปฏิจฺจ.}

\prose{นสรณํ นอาสว\-วิปฺปยุตฺตํ นสาสวํ ธมฺมํ ปฏิจฺจ.}

\prose{นสรณํ นอาสว\-วิปฺปยุตฺตํ นอนาสวํ ธมฺมํ ปฏิจฺจ.}

\csromanmarkh{100. สรณทุก}\csromanmarkhii{20. สญฺโญชน\-ทุกาทิ}\csromanheadernomark{2}{100. สรณทุก 20. สญฺโญชน\-ทุกาทิ}
\proseitem{48}{นสรณํ โนสญฺโญชนํ ธมฺมํ ปฏิจฺจ.}

\proseitem{49}{นสรณํ โนคนฺถํ ธมฺมํ ปฏิจฺจ ฯเปฯ. นสรณํ โนโอฆํ ธมฺมํ ปฏิจฺจ ฯเปฯ. นสรณํ โนโยคํ ธมฺมํ ปฏิจฺจ ฯเปฯ. นสรณํ โนนีวรณํ ธมฺมํ ปฏิจฺจ ฯเปฯ. นสรณํ โนปรามาสํ ธมฺมํ ปฏิจฺจ.}

\csromanmarkh{100. สรณทุก}\csromanmarkhii{55. มหนฺตรทุก}\csromanheadernomark{2}{100. \textbf{สรณทุก 5}5. มหนฺตรทุก}
\prose{50. นสรณํ นสารมฺมณํ ธมฺมํ ปฏิจฺจ. (สํขิตฺตํ.)}

\prose{นสรณํ นจิตฺตํ ธมฺมํ ปฏิจฺจ.}

\prose{นสรณํ นเจตสิกํ ธมฺมํ ปฏิจฺจ.}

\prose{นสรณํ นจิตฺต\-สมฺปยุตฺตํ ธมฺมํ ปฏิจฺจ ฯเปฯ. นสรณํ นจิตฺต\-สํสฏฺฐํ ธมฺมํ ปฏิจฺจ.}

\csromanmarkh{100. สรณทุก}\csromanmarkhii{83. ทสฺสนทุกาทิ}\csromanheadernomark{2}{100. สรณทุก 83. ทสฺสนทุกาทิ}
\prose{\csromanfolio{171}นสรณํ นจิตฺต\-สมุฏฺฐานํ ธมฺมํ ปฏิจฺจ.}

\prose{นสรณํ นจิตฺต\-สหภุํ ธมฺมํ ปฏิจฺจ ฯเปฯ. นสรณํ นจิตฺตา\-นุปริวตฺติํ ธมฺมํ ปฏิจฺจ.}

\proseitem{51}{นสรณํ นจิตฺต\-สํสฏฺฐ\-สมุฏฺฐานํ ธมฺมํ ปฏิจฺจ ฯเปฯ. นสรณํ นจิตฺต\-สํสฏฺฐ\-สมุฏฺฐาน\-สหภุํ ธมฺมํ ปฏิจฺจ ฯเปฯ. นสรณํ นจิตฺต\-สํสฏฺฐ\-สมุฏฺฐานา\-นุ\-ปริวตฺติํ ธมฺมํ ปฏิจฺจ.}

\prose{นสรณํ นอชฺฌตฺติกํ ธมฺมํ ปฏิจฺจ.}

\prose{นสรณํ นพาหิรํ ธมฺมํ ปฏิจฺจ.}

\prose{นสรณํ นอุปาทา ธมฺมํ ปฏิจฺจ.}

\prose{นสรณํ นอุปาทินฺนํ ธมฺมํ ปฏิจฺจ.}

\prose{นสรณํ นอนุปาทินฺนํ ธมฺมํ ปฏิจฺจ.}

\csromanmarkh{100. สรณทุก}\csromanmarkhii{69. อุปาทานาทิทุก}\csromanheadernomark{2}{100. \textbf{สรณทุก 6}9. อุปาทานาทิทุก}
\proseitem{52}{นสรณํ นอุปาทานํ ธมฺมํ ปฏิจฺจ.}

\proseitem{53}{นสรณํ โนกิเลสํ ธมฺมํ ปฏิจฺจ.}

\csromanmarkh{100. สรณทุก}\csromanmarkhii{83. ทสฺสนทุกาทิ}\csromanheadernomark{2}{100. สรณทุก 83. ทสฺสนทุกาทิ}
\proseitem{54}{นสรณํ นทสฺสเนน ปหาตพฺพํ ธมฺมํ ปฏิจฺจ ฯเปฯ. นสรณํ นนทสฺสเนน ปหาตพฺพํ ธมฺมํ ปฏิจฺจ ฯเปฯ. นสรณํ นภาวนาย ปหาตพฺพํ ธมฺมํ ปฏิจฺจ ฯเปฯ. นสรณํ นนภาวนาย ปหาตพฺพํ ธมฺมํ ปฏิจฺจ ฯเปฯ. นสรณํ นทสฺสเนน ปหาตพฺพ\-เหตุกํ ธมฺมํ ปฏิจฺจ ฯเปฯ. นสรณํ นนทสฺสเนน ปหาตพฺพ\-เหตุกํ ธมฺมํ ปฏิจฺจ ฯเปฯ. นสรณํ นภาวนาย ปหาตพฺพ\-เหตุกํ ธมฺมํ ปฏิจฺจ.}

\prose{นสรณํ นนภาวนาย ปหาตพฺพ\-เหตุกํ ธมฺมํ ปฏิจฺจ.}

\proseitem{55}{นสรณํ นสวิตกฺกํ ธมฺมํ ปฏิจฺจ ฯเปฯ. นสรณํ นอวิตกฺกํ ธมฺมํ ปฏิจฺจ ฯเปฯ. นสรณํ นสวิจารํ ธมฺมํ ปฏิจฺจ.}

\prose{นสรณํ นอวิจารํ ธมฺมํ ปฏิจฺจ.}

\proseitem{56}{นสรณํ นสปฺปีติกํ ธมฺมํ ปฏิจฺจ ฯเปฯ. นสรณํ นอปฺปีติกํ ธมฺมํ ปฏิจฺจ ฯเปฯ. นสรณํ นปีติสหคตํ ธมฺมํ ปฏิจฺจ ฯเปฯ. นสรณํ นนปี\-ติ\-สหคตํ ธมฺมํ ปฏิจฺจ ฯเปฯ. นสรณํ นสุขสหคตํ ธมฺมํ ปฏิจฺจ ฯเปฯ. นสรณํ นนสุข\-สหคตํ ธมฺมํ ปฏิจฺจ ฯเปฯ. นสรณํ น- อุเปกฺขา\-สหคตํ ธมฺมํ ปฏิจฺจ.}

\proseitem{57}{\csromanfolio{172}นสรณํ นนอุเปกฺขาสหคตํ ธมฺมํ ปฏิจฺจ.}

\prose{นสรณํ นกามาวจรํ ธมฺมํ ปฏิจฺจ.}

\prose{นสรณํ นนกามา\-วจรํ ธมฺมํ ปฏิจฺจ.}

\proseitem{58}{นสรณํ นรูปาวจรํ ธมฺมํ ปฏิจฺจ ฯเปฯ. นสรณํ นนรูปาวจรํ ธมฺมํ ปฏิจฺจ ฯเปฯ. นสรณํ นอรูปาวจรํ ธมฺมํ ปฏิจฺจ.}

\prose{นสรณํ นนอรูปาวจรํ ธมฺมํ ปฏิจฺจ.}

\prose{นสรณํ นปริยาปนฺนํ ธมฺมํ ปฏิจฺจ.}

\prose{นสรณํ นอปริยาปนฺนํ ธมฺมํ ปฏิจฺจ.}

\proseitem{59}{นสรณํ นนิยฺยานิกํ ธมฺมํ ปฏิจฺจ.}

\prose{นสรณํ นอนิยฺยานิกํ ธมฺมํ ปฏิจฺจ.}

\prose{นสรณํ นนิยตํ ธมฺมํ ปฏิจฺจ.}

\prose{นสรณํ นอนิยตํ ธมฺมํ ปฏิจฺจ.}

\proseitem{60}{นสรณํ นสอุตฺตรํ ธมฺมํ ปฏิจฺจ นสรโณ นสอุตฺตโร ธมฺโม อุปฺปชฺชติ เหตุปจฺจยา.}

\prose{นสรณํ นอนุตฺตรํ ธมฺมํ ปฏิจฺจ นสรโณ นอนุตฺตโร}

\prose{นอรณํ นอนุตฺตรํ ธมฺมํ ปฏิจฺจ นอรโณ นอนุตฺตโร ธมฺโม อุปฺปชฺชติ เหตุปจฺจยา.  นอรณํ นอนุตฺตรํ ธมฺมํ ปฏิจฺจ นสรโณ นอนุตฺตโร ธมฺโม อุปฺปชฺชติ เหตุปจฺจยา.  นอรณํ นอนุตฺตรํ ธมฺมํ ปฏิจฺจ นสรโณ นอนุตฺตโร จ นอรโณ นอนุตฺตโร จ ธมฺมา อุปฺปชฺชนฺติ เหตุปจฺจยา. (3)}

\prose{นสรณํ นอนุตฺตรญฺจ นอรณํ นอนุตฺตรญฺจ ธมฺมํ ปฏิจฺจ นสรโณ นอนุตฺตโร ธมฺโม อุปฺปชฺชติ เหตุปจฺจยา. (1) (สํขิตฺตํ.)}

\prosewithcloser{เหตุยา ปญฺจ. (สพฺพตฺถ วิตฺตาโร. สหชาตวารมฺปิ ปจฺจยวารมฺปิ นิสฺสยวารมฺปิ สํสฏฺฐ\-วารมฺปิ สมฺปยุตฺตวารมฺปิ ปญฺหาวารมฺปิ วิตฺถาเรตพฺพํ.)}{ธมฺม\-ปจฺจนีเย ทุกทุก\-ปฏฺฐานํ นิฏฺฐิตํ.}
\nitthitam{ปจฺจนีย\-ปฏฺฐานํ นิฏฺฐิตํ.}
\csromanheader{2}{\textbf{อภิธมฺมปิฏิก}}
\csromanheader{2}{ธมฺมา\-นุโลม\-ปจฺจนีย}
\namakkaram{นโม ตสฺส ภควโต อรหโต สมฺมา\-สมฺพุทฺธสฺส.}
\csromanmarkh{1. กุสลตฺติก}\csromanmarkhii{1. ปฏิจฺจวาร}\csromanheadernomark{2}{1. \textbf{กุสลตฺติก 1. ปฏิจฺจวาร}}
\prose{\csromanfolio{173}ปจฺจยจตุกฺกเหตุ}

\proseitem{1}{กุสลํ ธมฺมํ ปฏิจฺจ นกุสโล ธมฺโม อุปฺปชฺชติ เหตุปจฺจยา. กุสเล ขนฺเธ ปฏิจฺจ จิตฺต\-สมุฏฺฐานํ รูปํ.  กุสลํ ธมฺมํ ปฏิจฺจ นอกุสโล ธมฺโม อุปฺปชฺชติ เหตุปจฺจยา. กุสลํ เอกํ ขนฺธํ ปฏิจฺจ ตโย ขนฺธา จิตฺต\-สมุฏฺฐานญฺจ รูปํ ฯเปฯ.  กุสลํ ธมฺมํ ปฏิจฺจ นอพฺยากโต ธมฺโม อุปฺปชฺชติ เหตุปจฺจยา. กุสลํ เอกํ ขนฺธํ ปฏิจฺจ ตโย ขนฺธา ฯเปฯ.  กุสลํ ธมฺมํ ปฏิจฺจ นอกุสโล จ นอพฺยากโต จ ธมฺมา อุปฺปชฺชนฺติ เหตุปจฺจยา. กุสลํ เอกํ ขนฺธํ ปฏิจฺจ ตโย ขนฺธา ฯเปฯ.  กุสลํ ธมฺมํ ปฏิจฺจ นกุสโล จ นอกุสโล จ ธมฺมา อุปฺปชฺชนฺติ เหตุปจฺจยา. กุสเล ขนฺเธ ปฏิจฺจ จิตฺต\-สมุฏฺฐานํ รูปํ. (5)}

\proseitem{2}{อกุสลํ ธมฺมํ ปฏิจฺจ นอกุสโล ธมฺโม อุปฺปชฺชติ เหตุปจฺจยา. อกุสเล ขนฺเธ ปฏิจฺจ จิตฺต\-สมุฏฺฐานํ รูปํ.  อกุสลํ ธมฺมํ ปฏิจฺจ นกุสโล ธมฺโม อุปฺปชฺชติ เหตุปจฺจยา. อกุสลํ เอกํ ขนฺธํ ปฏิจฺจ ตโย ขนฺธา จิตฺต\-สมุฏฺฐานญฺจ รูปํ ฯเปฯ.  อกุสลํ ธมฺมํ ปฏิจฺจ นอพฺยากโต ธมฺโม อุปฺปชฺชติ เหตุปจฺจยา.  อกุสลํ ธมฺมํ ปฏิจฺจ นกุสโล จ นอพฺยากโต จ ธมฺมา อุปฺปชฺชนฺติ เหตุปจฺจยา.  อกุสลํ ธมฺมํ ปฏิจฺจ นกุสโล จ นอกุสโล จ ธมฺมา อุปฺปชฺชนฺติ เหตุปจฺจยา. (5)}

\csromanpalirecto
\gambhira{ธมฺมา\-นุโลม\-ปจฺจนีย ติกปฏฺฐาน\-ปาฬิ}
\proseitem{3}{\csromanfolio{174}อพฺยากตํ ธมฺมํ ปฏิจฺจ นกุสโล ธมฺโม อุปฺปชฺชติ เหตุปจฺจยา. วิปากาพฺยากตํ กิริยาพฺยากตํ เอกํ ขนฺธํ ปฏิจฺจ ตโย ขนฺธา จิตฺต\-สมุฏฺฐานญฺจ รูปํ ฯเปฯ. ปฏิสนฺธิ\-กฺขเณ ฯเปฯ เอกํ มหาภูตํ ปฏิจฺจ ตโย มหาภูตา ฯเปฯ เทฺว มหาภูเต ปฏิจฺจ เทฺว มหาภูตา, มหาภูเต ปฏิจฺจ จิตฺต\-สมุฏฺฐานํ รูปํ.  อพฺยากตํ ธมฺมํ ปฏิจฺจ นอกุสโล ธมฺโม อุปฺปชฺชติ เหตุปจฺจยา. วิปากาพฺยากตํ ฯเปฯ. ปฏิสนฺธิ\-กฺขเณ ฯเปฯ มหาภูตํ ฯเปฯ.  อพฺยากตํ ธมฺมํ ปฏิจฺจ นกุสโล จ นอกุสโล จ ธมฺมา อุปฺปชฺชนฺติ เหตุปจฺจยา. วิปากาพฺยากตํ กิริยาพฺยากตํ ฯเปฯ. (3) (สํขิตฺตํ.)}

\proseitem{4}{กุสลญฺจ อพฺยากตญฺจ ธมฺมํ ปฏิจฺจ นกุสโล ธมฺโม อุปฺปชฺชติ เหตุปจฺจยา. กุสเล ขนฺเธ จ มหาภูเต จ ปฏิจฺจ จิตฺต\-สมุฏฺฐานํ รูปํ.  กุสลญฺจ อพฺยากตญฺจ ธมฺมํ ปฏิจฺจ นอกุสโล ธมฺโม อุปฺปชฺชติ เหตุปจฺจยา.  กุสลญฺจ อพฺยากตญฺจ ธมฺมํ ปฏิจฺจ นกุสโล จ นอกุสโล จ ธมฺมา อุปฺปชฺชนฺติ เหตุปจฺจยา. (3)}

\prose{อกุสลญฺจ อพฺยากตญฺจ ธมฺมํ ปฏิจฺจ นกุสโล ธมฺโม อุปฺปชฺชติ เหตุปจฺจยา.  อกุสลญฺจ อพฺยากตญฺจ ธมฺมํ ปฏิจฺจ นอกุสโล ธมฺโม อุปฺปชฺชติ เหตุปจฺจยา.  อกุสลญฺจ อพฺยากตญฺจ ธมฺมํ ปฏิจฺจ นกุสโล จ นอกุสโล จ ธมฺมา อุปฺปชฺชนฺติ เหตุปจฺจยา. (3) (จิตฺตสมุฏฺฐาน\-รูปเมว เอตฺถ วตฺตติ. เอกูนวีสติ ปญฺหา กาตพฺพา.)}

\csromanheader{2}{\textbf{อารมฺมณ}}
\proseitem{5}{กุสลํ ธมฺมํ ปฏิจฺจ นอกุสโล ธมฺโม อุปฺปชฺชติ อารมฺมณ\-ปจฺจยา.  กุสลํ ธมฺมํ ปฏิจฺจ นอพฺยากโต ธมฺโม อุปฺปชฺชติ อารมฺมณ\-ปจฺจยา.  กุสลํ ธมฺมํ ปฏิจฺจ นอกุสโล จ นอพฺยากโต จ ธมฺมา อุปฺปชฺชนฺติ อารมฺมณ\-ปจฺจยา. (3)}

\prose{อกุสลํ ธมฺมํ ปฏิจฺจ นกุสโล ธมฺโม อุปฺปชฺชติ อารมฺมณ\-ปจฺจยา.  อกุสลํ ธมฺมํ ปฏิจฺจ นอพฺยากโต ธมฺโม อุปฺปชฺชติ อารมฺมณ\-ปจฺจยา.  อกุสลํ ธมฺมํ ปฏิจฺจ นกุสโล จ นอพฺยากโต จ ธมฺมา อุปฺปชฺชนฺติ อารมฺมณ\-ปจฺจยา. (3)}

\prose{อพฺยากตํ ธมฺมํ ปฏิจฺจ นกุสโล ธมฺโม อุปฺปชฺชติ อารมฺมณ\-ปจฺจยา.  อพฺยากตํ ธมฺมํ ปฏิจฺจ นอกุสโล ธมฺโม อุปฺปชฺชติ อารมฺมณ\-ปจฺจยา.   \csromanfolio{175}อพฺยากตํ ธมฺมํ ปฏิจฺจ นกุสโล จ นอกุสโล จ ธมฺมา อุปฺปชฺชนฺติ อารมฺมณ\-ปจฺจยา. (3) (สํขิตฺตํ.)}

\prose{เหตุยา เอกูนวีส, อารมฺมเณ นว, อธิปติยา เอกูนวีส, อนนฺตเร นว, สมนนฺตเร นว, สหชาเต เอกูนวีส ฯเปฯ อวิคเต เอกูนวีส.}

\prose{ปจฺจนียนเหตุนอารมฺมณ}

\proseitem{6}{อกุสลํ ธมฺมํ ปฏิจฺจ นกุสโล ธมฺโม อุปฺปชฺชติ นเหตุปจฺจยา.  อกุสลํ ธมฺมํ ปฏิจฺจ นอพฺยากโต ธมฺโม อุปฺปชฺชติ นเหตุปจฺจยา.  อกุสลํ ธมฺมํ ปฏิจฺจ นกุสโล จ นอพฺยากโต จ ธมฺมา อุปฺปชฺชนฺติ นเหตุปจฺจยา. (3)}

\prose{อพฺยากตํ ธมฺมํ ปฏิจฺจ นกุสโล ธมฺโม อุปฺปชฺชติ นเหตุปจฺจยา.  อพฺยากตํ ธมฺมํ ปฏิจฺจ นอกุสโล ธมฺโม อุปฺปชฺชติ นเหตุปจฺจยา.  อพฺยากตํ ธมฺมํ ปฏิจฺจ นกุสโล จ นอกุสโล จ ธมฺมา อุปฺปชฺชนฺติ นเหตุปจฺจยา. (3)}

\proseitem{7}{กุสลํ ธมฺมํ ปฏิจฺจ นกุสโล ธมฺโม อุปฺปชฺชติ นอารมฺมณ\-ปจฺจยา.  กุสลํ ธมฺมํ ปฏิจฺจ นอกุสโล ธมฺโม อุปฺปชฺชติ นอารมฺมณ\-ปจฺจยา.  กุสลํ ธมฺมํ ปฏิจฺจ นกุสโล จ นอกุสโล จ ธมฺมา อุปฺปชฺชนฺติ นอารมฺมณ\-ปจฺจยา. (3) อกุสลํ ธมฺมํ ปฏิจฺจ นอกุสโล ธมฺโม อุปฺปชฺชติ นอารมฺมณ\-ปจฺจยา.  อกุสลํ ธมฺมํ ปฏิจฺจ นกุสโล ธมฺโม อุปฺปชฺชติ นอารมฺมณ\-ปจฺจยา.  อกุสลํ ธมฺมํ ปฏิจฺจ นกุสโล จ นอกุสโล จ ธมฺมา อุปฺปชฺชนฺติ นอารมฺมณ\-ปจฺจยา. (3) (สํขิตฺตํ.)}

\prose{นเหตุยา ฉ, นอารมฺมเณ ปนฺนรส, นอธิปติยา เอกูนวีส ฯเปฯ โนวิคเต ปนฺนรส.}

\prose{(ปจฺจนียํ วิตฺถาเรตพฺพํ. สหชาตวารมฺปิ ปจฺจยวารมฺปิ วิตฺถาเรตพฺพํ. ปจฺจยวาเรปิ เหตุยา ฉพฺพีส, อารมฺมเณ อฏฺฐารส ฯเปฯ อวิคเต ฉพฺพีส. นิสฺสยวารมฺปิ สํสฏฺฐ\-วารมฺปิ สมฺปยุตฺตวารมฺปิ ปฏิจฺจ\-วาร\-สทิสํ.)}

\csromanmarkh{ธมฺมา\-นุโลม\-ปจฺจนีย ติกปฏฺฐาน\-ปาฬิ}\csromanmarkhii{1. กุสลตฺติก 7. ปญฺหาวาร}\csromanheadernomark{2}{ธมฺมา\-นุโลม\-ปจฺจนีย ติกปฏฺฐาน\-ปาฬิ 1. กุสลตฺติก 7. ปญฺหาวาร}
\prose{\csromanfolio{176}ปจฺจยจตุกฺกเหตุอารมฺมณ}

\proseitem{8}{กุสโล ธมฺโม นกุสลสฺส ธมฺมสฺส เหตุปจฺจเยน ปจฺจโย.  กุสโล ธมฺโม นอกุสลสฺส ธมฺมสฺส เหตุปจฺจเยน ปจฺจโย.  กุสโล ธมฺโม นอพฺยากตสฺส ธมฺมสฺส เหตุปจฺจเยน ปจฺจโย.  กุสโล ธมฺโม นอกุสลสฺส จ นอพฺยากตสฺส จ ธมฺมสฺส เหตุปจฺจเยน ปจฺจโย.  กุสโล ธมฺโม นกุสลสฺส จ นอกุสลสฺส จ ธมฺมสฺส เหตุปจฺจเยน ปจฺจโย. (5)}

\proseitem{9}{อกุสโล ธมฺโม นอกุสลสฺส ธมฺมสฺส เหตุปจฺจเยน ปจฺจโย.  อกุสโล ธมฺโม นกุสลสฺส ธมฺมสฺส เหตุปจฺจเยน ปจฺจโย.  อกุสโล ธมฺโม นอพฺยากตสฺส ธมฺมสฺส เหตุปจฺจเยน ปจฺจโย.  อกุสโล ธมฺโม นกุสลสฺส จ นอพฺยากตสฺส จ ธมฺมสฺส เหตุปจฺจเยน ปจฺจโย.  อกุสโล ธมฺโม นกุสลสฺส จ นอกุสลสฺส จ ธมฺมสฺส เหตุปจฺจเยน ปจฺจโย. (5)}

\prose{อพฺยากโต ธมฺโม นกุสลสฺส ธมฺมสฺส เหตุปจฺจเยน ปจฺจโย.  อพฺยากโต ธมฺโม นอกุสลสฺส ธมฺมสฺส เหตุปจฺจเยน ปจฺจโย.  อพฺยากโต ธมฺโม นกุสลสฺส จ นอกุสลสฺส จ ธมฺมสฺส เหตุปจฺจเยน ปจฺจโย. (3)}

\proseitem{10}{กุสโล ธมฺโม นกุสลสฺส ธมฺมสฺส อารมฺมณ\-ปจฺจเยน ปจฺจโย. (ฉ ปญฺหา.)}

\prose{อกุสโล ธมฺโม นอกุสลสฺส ธมฺมสฺส อารมฺมณ\-ปจฺจเยน ปจฺจโย. (ฉ ปญฺหา.)}

\prose{อพฺยากโต ธมฺโม นอพฺยากตสฺส ธมฺมสฺส อารมฺมณ\-ปจฺจเยน ปจฺจโย. (ฉ ปญฺหา. สํขิตฺตํ.)}

\proseitem{11}{เหตุยา เตรส, อารมฺมเณ อฏฺฐารส, อธิปติยา สตฺตรส, อนนฺตเร โสฬส, สมนนฺตเร โสฬส, สหชาเต เอกูนวีส, อญฺญมญฺเญ นว, นิสฺสเย ฉพฺพีส, อุปนิสฺสเย อฏฺฐารส, ปุเรชาเต ฉ, ปจฺฉาชาเต นว, อาเสวเน นว, กมฺเม เตรส, วิปาเก ตีณิ, อาหาเร เตรส ฯเปฯ มคฺเค เตรส, สมฺปยุตฺเต นว, วิปฺปยุตฺเต ทฺวาทส ฯเปฯ อวิคเต ฉพฺพีส. (ปญฺหาวารํ วิตฺถาเรตพฺพํ.)}

\csromanmarkh{2. เวทนาตฺติก}\csromanmarkhii{2. เวทนาตฺติก 1. ปฏิจฺจวาร}\csromanheadernomark{2}{2. เวทนาตฺติก \textbf{2. เวทนาตฺติก 1. ปฏิจฺจวาร}}
\prose{\csromanfolio{177}ปจฺจยจตุกฺกเหตุ สุขาย เวทนาย สมฺปยุตฺตํ ธมฺมํ ปฏิจฺจ นสุขาย เวทนาย สมฺปยุตฺโต ธมฺโม อุปฺปชฺชติ เหตุปจฺจยา. สุขาย เวทนาย สมฺปยุตฺเต ขนฺเธ ปฏิจฺจ สุขเวทนา จิตฺต\-สมุฏฺฐานญฺจ รูปํ. ปฏิสนฺธิ\-กฺขเณ ฯเปฯ. (มหาภูตา นตฺถิ.) . สุขาย เวทนาย สมฺปยุตฺตํ ธมฺมํ ปฏิจฺจ นทุกฺขาย เวทนาย สมฺปยุตฺโต ธมฺโม อุปฺปชฺชติ เหตุปจฺจยา. สุขาย เวทนาย สมฺปยุตฺตํ เอกํ ขนฺธํ ปฏิจฺจ ตโย ขนฺธา จิตฺต\-สมุฏฺฐานญฺจ รูปํ ฯเปฯ. สุขาย เวทนาย สมฺปยุตฺตํ ธมฺมํ ปฏิจฺจ นอทุกฺขม\-สุขาย เวทนาย สมฺปยุตฺโต ธมฺโม อุปฺปชฺชติ เหตุปจฺจยา.  สุขาย เวทนาย สมฺปยุตฺตํ ธมฺมํ ปฏิจฺจ นสุขาย เวทนาย สมฺปยุตฺโต จ นอทุกฺขม\-สุขาย เวทนาย สมฺปยุตฺโต จ ธมฺมา อุปฺปชฺชนฺติ เหตุปจฺจยา.  สุขาย เวทนาย สมฺปยุตฺตํ ธมฺมํ ปฏิจฺจ นทุกฺขาย เวทนาย สมฺปยุตฺโต จ นอทุกฺขม\-สุขาย เวทนาย สมฺปยุตฺโต จ ธมฺมา อุปฺปชฺชนฺติ เหตุปจฺจยา.  สุขาย เวทนาย สมฺปยุตฺตํ ธมฺมํ ปฏิจฺจ นสุขาย เวทนาย สมฺปยุตฺโต จ นทุกฺขาย เวทนาย สมฺปยุตฺโต จ ธมฺมา อุปฺปชฺชนฺติ เหตุปจฺจยา.  สุขาย เวทนาย สมฺปยุตฺตํ ธมฺมํ ปฏิจฺจ นสุขาย เวทนาย สมฺปยุตฺโต จ นทุกฺขาย เวทนาย สมฺปยุตฺโต จ นอทุกฺขม\-สุขาย เวทนาย สมฺปยุตฺโต จ ธมฺมา อุปฺปชฺชนฺติ เหตุปจฺจยา. (7)}

\proseitem{13}{ทุกฺขาย เวทนาย สมฺปยุตฺตํ ธมฺมํ ปฏิจฺจ นทุกฺขาย เวทนาย สมฺปยุตฺโต ธมฺโม อุปฺปชฺชติ เหตุปจฺจยา.  ทุกฺขาย เวทนาย สมฺปยุตฺตํ ธมฺมํ ปฏิจฺจ นสุขาย เวทนาย สมฺปยุตฺโต ธมฺโม อุปฺปชฺชติ เหตุปจฺจยา.  ทุกฺขาย เวทนาย สมฺปยุตฺตํ ธมฺมํ ปฏิจฺจ นอทุกฺขม\-สุขาย เวทนาย สมฺปยุตฺโต ธมฺโม อุปฺปชฺชติ เหตุปจฺจยา.  ทุกฺขาย เวทนาย สมฺปยุตฺตํ ธมฺมํ ปฏิจฺจ นสุขาย เวทนาย สมฺปยุตฺโต จ นอทุกฺขม\-สุขาย เวทนาย สมฺปยุตฺโต จ ธมฺมา อุปฺปชฺชนฺติ เหตุปจฺจยา.  ทุกฺขาย เวทนาย สมฺปยุตฺตํ ธมฺมํ ปฏิจฺจ นทุกฺขาย เวทนาย สมฺปยุตฺโต จ นอทุกฺขม\-สุขาย เวทนาย สมฺปยุตฺโต จ ธมฺมา อุปฺปชฺชนฺติ เหตุปจฺจยา.  ทุกฺขาย เวทนาย สมฺปยุตฺตํ ธมฺมํ ปฏิจฺจ นสุขาย เวทนาย สมฺปยุตฺโต จ นทุกฺขาย เวทนาย สมฺปยุตฺโต จ ธมฺมา อุปฺปชฺชนฺติ เหตุปจฺจยา.  ทุกฺขาย เวทนาย สมฺปยุตฺตํ ธมฺมํ ปฏิจฺจ นสุขาย เวทนาย สมฺปยุตฺโต จ นทุกฺขาย เวทนาย}

\prose{\csromanfolio{178}สมฺปยุตฺโต จ นอทุกฺขม\-สุขาย เวทนาย สมฺปยุตฺโต จ ธมฺมา อุปฺปชฺชนฺติ เหตุปจฺจยา. (7) อทุกฺขม\-สุขาย เวทนาย สมฺปยุตฺตํ ธมฺมํ ปฏิจฺจ นอทุกฺขม\-สุขาย เวทนาย สมฺปยุตฺโต ธมฺโม อุปฺปชฺชติ เหตุปจฺจยา.  อทุกฺขม\-สุขาย เวทนาย สมฺปยุตฺตํ ธมฺมํ ปฏิจฺจ นสุขาย เวทนาย สมฺปยุตฺโต ธมฺโม อุปฺปชฺชติ เหตุปจฺจยา.  อทุกฺขม\-สุขาย เวทนาย สมฺปยุตฺตํ ธมฺมํ ปฏิจฺจ นทุกฺขาย เวทนาย สมฺปยุตฺโต ธมฺโม อุปฺปชฺชติ เหตุปจฺจยา.  อทุกฺขม\-สุขาย เวทนาย สมฺปยุตฺตํ ธมฺมํ ปฏิจฺจ นสุขาย เวทนาย สมฺปยุตฺโต จ นอทุกฺขม\-สุขาย เวทนาย สมฺปยุตฺโต จ ธมฺมา อุปฺปชฺชนฺติ เหตุปจฺจยา.  อทุกฺขม\-สุขาย เวทนาย สมฺปยุตฺตํ ธมฺมํ ปฏิจฺจ นทุกฺขาย เวทนาย สมฺปยุตฺโต จ นอทุกฺขม\-สุขาย เวทนาย สมฺปยุตฺโต จ ธมฺมา อุปฺปชฺชนฺติ เหตุปจฺจยา.  อทุกฺขม\-สุขาย เวทนาย สมฺปยุตฺตํ ธมฺมํ ปฏิจฺจ นสุขาย เวทนาย สมฺปยุตฺโต จ นทุกฺขาย เวทนาย สมฺปยุตฺโต จ ธมฺมา อุปฺปชฺชนฺติ เหตุปจฺจยา.  อทุกฺขม\-สุขาย เวทนาย สมฺปยุตฺตํ ธมฺมํ ปฏิจฺจ นสุขาย เวทนาย สมฺปยุตฺโต จ นทุกฺขาย เวทนาย สมฺปยุตฺโต จ นอทุกฺขม\-สุขาย เวทนาย สมฺปยุตฺโต จ ธมฺมา อุปฺปชฺชนฺติ เหตุปจฺจยา. (7) (สํขิตฺตํ.)}

\proseitem{14}{เหตุยา เอกวีส, อารมฺมเณ เอกวีส ฯเปฯ อวิคเต เอกวีส.}

\prose{ปจฺจนีย\-น\-เหตุ}

\proseitem{15}{สุขาย เวทนาย สมฺปยุตฺตํ ธมฺมํ ปฏิจฺจ นสุขาย เวทนาย สมฺปยุตฺโต ธมฺโม อุปฺปชฺชติ นเหตุปจฺจยา. (สํขิตฺตํ.)}

\prose{นเหตุยา เอกวีส, นอารมฺมเณ เอกวีส ฯเปฯ นวิปฺปยุตฺโต จุทฺทส ฯเปฯ โนวิคเต เอกวีส.}

\prose{(สหชาตวารมฺปิ ปจฺจยวารมฺปิ นิสฺสยวารมฺปิ สํสฏฺฐ\-วารมฺปิ สมฺปยุตฺตวารมฺปิ ปฏิจฺจ\-วาร\-สทิสํ.)}

\csromanmarkh{2. เวทนาตฺติก}\csromanmarkhii{7. ปญฺหาวาร}\csromanheadernomark{2}{2. \textbf{เวทนาตฺติก 7. ปญฺหาวาร}}
\prose{ปจฺจยจตุกฺกเหตุ}

\proseitem{16}{สุขาย เวทนาย สมฺปยุตฺโต ธมฺโม นสุขาย เวทนาย สมฺปยุตฺตสฺส ธมฺมสฺส เหตุปจฺจเยน ปจฺจโย.  สุขาย เวทนาย \csromanfolio{179}สมฺปยุตฺโต ธมฺโม นทุกฺขาย เวทนาย สมฺปยุตฺตสฺส ธมฺมสฺส เหตุปจฺจเยน ปจฺจโย.  สุขาย เวทนาย สมฺปยุตฺโต ธมฺโม นอทุกฺขม\-สุขาย เวทนาย สมฺปยุตฺตสฺส ธมฺมสฺส เหตุปจฺจเยน ปจฺจโย.  สุขาย เวทนาย สมฺปยุตฺโต ธมฺโม นสุขาย เวทนาย สมฺปยุตฺตสฺส จ นอทุกฺขม\-สุขาย เวทนาย สมฺปยุตฺตสฺส จ ธมฺมสฺส เหตุปจฺจเยน ปจฺจโย.  สุขาย เวทนาย สมฺปยุตฺโต ธมฺโม นทุกฺขาย เวทนาย สมฺปยุตฺตสฺส จ นอทุกฺขม\-สุขาย เวทนาย สมฺปยุตฺตสฺส จ ธมฺมสฺส เหตุปจฺจเยน ปจฺจโย.  สุขาย เวทนาย สมฺปยุตฺโต ธมฺโม นสุขาย เวทนาย สมฺปยุตฺตสฺส จ นทุกฺขาย เวทนาย สมฺปยุตฺตสฺส จ ธมฺมสฺส เหตุปจฺจเยน ปจฺจโย.  สุขาย เวทนาย สมฺปยุตฺโต ธมฺโม นสุขาย เวทนาย สมฺปยุตฺตสฺส จ นทุกฺขาย เวทนาย สมฺปยุตฺตสฺส จ นอทุกฺขม\-สุขาย เวทนาย สมฺปยุตฺตสฺส จ ธมฺมสฺส เหตุปจฺจเยน ปจฺจโย. (7) (สํขิตฺตํ.)}

\prose{เหตุยา เอกวีส, อารมฺมเณ เอกวีส ฯเปฯ อวิคเต เอกวีส. (ปญฺหาวารํ วิตฺถาเรตพฺพํ.)}

\csromanmarkh{3. วิปากตฺติก}\csromanmarkhii{1. ปฏิจฺจวาร}\csromanheadernomark{2}{3. วิปากตฺติก 1. ปฏิจฺจวาร}
\prose{ปจฺจยจตุกฺกเหตุอารมฺมณ}

\proseitem{17}{วิปากํ ธมฺมํ ปฏิจฺจ นวิปาโก ธมฺโม อุปฺปชฺชติ เหตุปจฺจยา.  วิปากํ ธมฺมํ ปฏิจฺจ นวิปาก\-ธมฺม\-ธมฺโม อุปฺปชฺชติ เหตุปจฺจยา.  วิปากํ ธมฺมํ ปฏิจฺจ นเน\-ว\-วิปาก\-น\-วิปาก\-ธมฺม\-ธมฺโม อุปฺปชฺชติ เหตุปจฺจยา.  วิปากํ ธมฺมํ ปฏิจฺจ นวิปาก\-ธมฺม\-ธมฺโม จ นเน\-ว\-วิปาก\-น\-วิปาก\-ธมฺม\-ธมฺโม จ ธมฺมา อุปฺปชฺชนฺติ เหตุปจฺจยา.  วิปากํ ธมฺมํ ปฏิจฺจ นวิปาโก จ นวิปาก\-ธมฺม\-ธมฺโม จ ธมฺมา อุปฺปชฺชนฺติ เหตุปจฺจยา. (ปญฺจ ปญฺหา.)}

\proseitem{18}{วิปาก\-ธมฺม\-ธมฺมํ ปฏิจฺจ นวิปาก\-ธมฺม\-ธมฺโม อุปฺปชฺชติ เหตุปจฺจยา.  วิปาก\-ธมฺม\-ธมฺมํ ปฏิจฺจ นวิปาโก ธมฺโม อุปฺปชฺชติ เหตุปจฺจยา.  วิปาก\-ธมฺม\-ธมฺมํ ปฏิจฺจ นเน\-ว\-วิปาก\-น\-วิปาก\-ธมฺม\-ธมฺโม อุปฺปชฺชติ เหตุปจฺจยา.  วิปาก\-ธมฺม\-ธมฺมํ ปฏิจฺจ นวิปาโก จ นเน\-ว\-วิปาก\-น\-วิปาก\-ธมฺม\-ธมฺโม จ ธมฺมา อุปฺปชฺชนฺติ เหตุปจฺจยา.  วิปาก\-ธมฺม\-ธมฺมํ ปฏิจฺจ นวิปาโก จ นวิปาก\-ธมฺม\-ธมฺโม จ ธมฺมา อุปฺปชฺชนฺติ เหตุปจฺจยา. (ปญฺจ ปญฺหา.) \csromanfolio{180}เนว\-วิปาก\-น\-วิปากธมฺม\-ธมฺมํ ปฏิจฺจ นเน\-ว\-วิปาก\-น\-วิปาก\-ธมฺม\-ธมฺโม อุปฺปชฺชติ เหตุปจฺจยา.  เนว\-วิปาก\-น\-วิปากธมฺม\-ธมฺมํ ปฏิจฺจ นวิปาโก ธมฺโม อุปฺปชฺชติ เหตุปจฺจยา.  เนว\-วิปาก\-น\-วิปากธมฺม\-ธมฺมํ ปฏิจฺจ นวิปาก\-ธมฺม\-ธมฺโม อุปฺปชฺชติ เหตุปจฺจยา.  เนว\-วิปาก\-น\-วิปากธมฺม\-ธมฺมํ ปฏิจฺจ นวิปาก\-ธมฺม\-ธมฺโม จ นเน\-ว\-วิปาก\-น\-วิปาก\-ธมฺม\-ธมฺโม จ ธมฺมา อุปฺปชฺชนฺติ เหตุปจฺจยา.  เนว\-วิปาก\-น\-วิปากธมฺม\-ธมฺมํ ปฏิจฺจ นวิปาโก จ นวิปาก\-ธมฺม\-ธมฺโม จ ธมฺมา อุปฺปชฺชนฺติ เหตุปจฺจยา. (ปญฺจ ปญฺหา.)}

\proseitem{20}{วิปากญฺจ เนว\-วิปาก\-น\-วิปากธมฺม\-ธมฺมญฺจ ธมฺมํ ปฏิจฺจ นวิปาโก ธมฺโม อุปฺปชฺชติ เหตุปจฺจยา.  วิปากญฺจ เนว\-วิปาก\-น\-วิปากธมฺม\-ธมฺมญฺจ ธมฺมํ ปฏิจฺจ นวิปาก\-ธมฺม\-ธมฺโม อุปฺปชฺชติ เหตุปจฺจยา.  วิปากญฺจ เนว\-วิปาก\-น\-วิปากธมฺม\-ธมฺมญฺจ ธมฺมํ ปฏิจฺจ นเน\-ว\-วิปาก\-น\-วิปาก\-ธมฺม\-ธมฺโม อุปฺปชฺชติ เหตุปจฺจยา.  วิปากญฺจ เนว\-วิปาก\-น\-วิปากธมฺม\-ธมฺมญฺจ ธมฺมํ ปฏิจฺจ นวิปาก\-ธมฺม\-ธมฺโม จ นเน\-ว\-วิปาก\-น\-วิปาก\-ธมฺม\-ธมฺโม จ ธมฺมา อุปฺปชฺชนฺติ เหตุปจฺจยา.  วิปากญฺจ เนว\-วิปาก\-น\-วิปากธมฺม\-ธมฺมญฺจ ธมฺมํ ปฏิจฺจ นวิปาโก จ นวิปาก\-ธมฺม\-ธมฺโม จ ธมฺมา อุปฺปชฺชนฺติ เหตุปจฺจยา. (5)}

\prose{วิปาก\-ธมฺม\-ธมฺมญฺจ เนว\-วิปาก\-น\-วิปากธมฺม\-ธมฺมญฺจ ธมฺมํ ปฏิจฺจ นวิปาโก ธมฺโม อุปฺปชฺชติ เหตุปจฺจยา.  วิปาก\-ธมฺม\-ธมฺมญฺจ เนว\-วิปาก\-น\-วิปากธมฺม\-ธมฺมญฺจ ธมฺมํ ปฏิจฺจ นวิปาก\-ธมฺม\-ธมฺโม อุปฺปชฺชติ เหตุปจฺจยา.  วิปาก\-ธมฺม\-ธมฺมญฺจ เนว\-วิปาก\-น\-วิปากธมฺม\-ธมฺมญฺจ ธมฺมํ ปฏิจฺจ นวิปาโก จ นวิปาก\-ธมฺม\-ธมฺโม จ ธมฺมา อุปฺปชฺชนฺติ เหตุปจฺจยา. (3)}

\proseitem{21}{วิปากํ ธมฺมํ ปฏิจฺจ นวิปาก\-ธมฺม\-ธมฺโม อุปฺปชฺชติ อารมฺมณ\-ปจฺจยา. ตีณิ.}

\prose{วิปาก\-ธมฺม\-ธมฺมํ ปฏิจฺจ นวิปาโก ธมฺโม อุปฺปชฺชติ อารมฺมณ\-ปจฺจยา. ตีณิ. เนว\-วิปาก\-น\-วิปากธมฺม\-ธมฺมํ ปฏิจฺจ นเน\-ว\-วิปาก\-น\-วิปาก\-ธมฺม\-ธมฺโม อุปฺปชฺชติ อารมฺมณ\-ปจฺจยา. (ปญฺจ ปญฺหา.)}

\prose{วิปากญฺจ เนว\-วิปาก\-น\-วิปากธมฺม\-ธมฺมญฺจ ธมฺมํ ปฏิจฺจ นวิปาก\-ธมฺม\-ธมฺโม อุปฺปชฺชติ อารมฺมณ\-ปจฺจยา. ตีณิ. (สํขิตฺตํ.)}

\proseitem{22}{เหตุยา เตวีส, อารมฺมเณ จุทฺทส ฯเปฯ อวิคเต เตวีส. (สํขิตฺตํ.)}

\prose{นเหตุยา อฏฺฐารส, นอารมฺมเณ ปนฺนรส.}

\prose{\csromanfolio{181}(สหชาตวารมฺปิ ปจฺจยวารมฺปิ นิสฺสยวารมฺปิ สํสฏฺฐ\-วารมฺปิ สมฺปยุตฺตวารมฺปิ วิตฺถาเรตพฺพํ.)}

\csromanmarkh{3. วิปากตฺติก}\csromanmarkhii{7. ปญฺหาวาร}\csromanheadernomark{2}{3. วิปากตฺติก 7. ปญฺหาวาร}
\prose{ปจฺจยจตุกฺกเหตุอารมฺมณ}

\proseitem{23}{วิปาโก ธมฺโม นวิปากสฺส ธมฺมสฺส เหตุปจฺจเยน ปจฺจโย.  วิปาโก ธมฺโม นวิปาก\-ธมฺม\-ธมฺมสฺส เหตุปจฺจเยน ปจฺจโย.  วิปาโก ธมฺโม นเน\-ว\-วิปาก\-น\-วิปากธมฺม\-ธมฺมสฺส เหตุปจฺจเยน ปจฺจโย.  วิปาโก ธมฺโม นวิปาก\-ธมฺม\-ธมฺมสฺส จ นเน\-ว\-วิปาก\-น\-วิปากธมฺม\-ธมฺมสฺส จ ธมฺมสฺส เหตุปจฺจเยน ปจฺจโย.  วิปาโก ธมฺโม นวิปากสฺส จ นวิปาก\-ธมฺม\-ธมฺมสฺส จ ธมฺมสฺส เหตุปจฺจเยน ปจฺจโย. (5)}

\proseitem{24}{วิปาโก ธมฺโม นวิปากสฺส ธมฺมสฺส อารมฺมณ\-ปจฺจเยน ปจฺจโย.  วิปาโก ธมฺโม นวิปาก\-ธมฺม\-ธมฺมสฺส อารมฺมณ\-ปจฺจเยน ปจฺจโย.  วิปาโก ธมฺโม นเน\-ว\-วิปาก\-น\-วิปากธมฺม\-ธมฺมสฺส อารมฺมณ\-ปจฺจเยน ปจฺจโย.  วิปาโก ธมฺโม นวิปากสฺส จ นเน\-ว\-วิปาก\-น\-วิปากธมฺม\-ธมฺมสฺส จ ธมฺมสฺส อารมฺมณ\-ปจฺจเยน ปจฺจโย.  วิปาโก ธมฺโม นวิปาก\-ธมฺม\-ธมฺมสฺส จ นเน\-ว\-วิปาก\-น\-วิปากธมฺม\-ธมฺมสฺส จ ธมฺมสฺส อารมฺมณ\-ปจฺจเยน ปจฺจโย.  วิปาโก ธมฺโม นวิปากสฺส จ นวิปาก\-ธมฺม\-ธมฺมสฺส จ ธมฺมสฺส อารมฺมณ\-ปจฺจเยน ปจฺจโย. (6)}

\prose{วิปาก\-ธมฺม\-ธมฺโม นวิปาก\-ธมฺม\-ธมฺมสฺส อารมฺมณ\-ปจฺจเยน ปจฺจโย.  วิปาก\-ธมฺม\-ธมฺโม นวิปากสฺส ธมฺมสฺส อารมฺมณ\-ปจฺจเยน ปจฺจโย.  วิปาก\-ธมฺม\-ธมฺโม นเน\-ว\-วิปาก\-น\-วิปากธมฺม\-ธมฺมสฺส อารมฺมณ\-ปจฺจเยน ปจฺจโย.  วิปาก\-ธมฺม\-ธมฺโม นวิปากสฺส จ นเน\-ว\-วิปาก\-น\-วิปากธมฺม\-ธมฺมสฺส จ ธมฺมสฺส อารมฺมณ\-ปจฺจเยน ปจฺจโย.  วิปาก\-ธมฺม\-ธมฺโม นวิปาก\-ธมฺม\-ธมฺมสฺส จ นเน\-ว\-วิปาก\-น\-วิปากธมฺม\-ธมฺมสฺส จ ธมฺมสฺส อารมฺมณ\-ปจฺจเยน ปจฺจโย.  วิปาก\-ธมฺม\-ธมฺโม นวิปากสฺส จ นวิปาก\-ธมฺม\-ธมฺมสฺส จ ธมฺมสฺส อารมฺมณ\-ปจฺจเยน ปจฺจโย. (6)}

\prose{เนว\-วิปาก\-น\-วิปาก\-ธมฺม\-ธมฺโม นเน\-ว\-วิปาก\-น\-วิปากธมฺม\-ธมฺมสฺส อารมฺมณ\-ปจฺจเยน ปจฺจโย ฯเปฯ. ฉ. (สํขิตฺตํ.)}

\proseitem{25}{\csromanfolio{182}เหตุยา เตรส, อารมฺมเณ อฏฺฐารส, อธิปติยา สตฺตรส, อนนฺตเร โสฬส ฯเปฯ ปุเรชาเต ฉ, ปจฺฉาชาเต นว, อาเสวเน ฉ, กมฺเม จุทฺทส, วิปาเก ปญฺจ, อินฺทฺริเย อฏฺฐารส ฯเปฯ วิปฺปยุตฺเต ทฺวาทส ฯเปฯ อวิคเต ฉพฺพีส. (ปญฺหาวารํ วิตฺถาเรตพฺพํ.)}

\csromanmarkh{4. อุปาทินฺนตฺติก}\csromanmarkhii{1. ปฏิจฺจวาร}\csromanheadernomark{2}{4. \textbf{อุปาทินฺนตฺติก 1. ปฏิจฺจวาร}}
\prose{ปจฺจยจตุกฺกเหตุ}

\proseitem{26}{อุปาทินฺนุ\-ปาทานิยํ ธมฺมํ ปฏิจฺจ นอุปาทินฺนุปาทานิโย ธมฺโม อุปฺปชฺชติ เหตุปจฺจยา.  อุปาทินฺนุ\-ปาทานิยํ ธมฺมํ ปฏิจฺจ นอนุปาทินฺนุ\-ปาทานิโย ธมฺโม อุปฺปชฺชติ เหตุปจฺจยา.  อุปาทินฺนุ\-ปาทานิยํ ธมฺมํ ปฏิจฺจ นอนุปาทินฺน\-อนุปาทานิโย ธมฺโม อุปฺปชฺชติ เหตุปจฺจยา.  อุปาทินฺนุ\-ปาทานิยํ ธมฺมํ ปฏิจฺจ นอุปาทินฺนุปาทานิโย จ นอนุปาทินฺน\-อนุปาทานิโย จ ธมฺมา อุปฺปชฺชนฺติ เหตุปจฺจยา.  อุปาทินฺนุ\-ปาทานิยํ ธมฺมํ ปฏิจฺจ นอนุปาทินฺนุ\-ปาทานิโย จ นอนุปาทินฺน\-อนุปาทานิโย จ ธมฺมา อุปฺปชฺชนฺติ เหตุปจฺจยา. (5)}

\prose{อนุปาทินฺนุ\-ปาทานิยํ ธมฺมํ ปฏิจฺจ นอุปาทินฺนุปาทานิโย ธมฺโม อุปฺปชฺชติ เหตุปจฺจยา.  อนุปาทินฺนุ\-ปาทานิยํ ธมฺมํ ปฏิจฺจ นอนุปาทินฺน\-อนุปาทานิโย ธมฺโม อุปฺปชฺชติ เหตุปจฺจยา.  อนุปาทินฺนุ\-ปาทานิยํ ธมฺมํ ปฏิจฺจ นอุปาทินฺนุปาทานิโย จ นอนุปาทินฺน\-อนุปาทานิโย จ ธมฺมา อุปฺปชฺชนฺติ เหตุปจฺจยา. (3)}

\prose{อนุปาทินฺน\-อนุปาทานิยํ ธมฺมํ ปฏิจฺจ นอนุปาทินฺน\-อนุปาทานิโย ธมฺโม อุปฺปชฺชติ เหตุปจฺจยา.  อนุปาทินฺน\-อนุปาทานิยํ ธมฺมํ ปฏิจฺจ นอุปาทินฺนุปาทานิโย ธมฺโม อุปฺปชฺชติ เหตุปจฺจยา.  อนุปาทินฺน\-อนุปาทานิยํ ธมฺมํ ปฏิจฺจ นอนุปาทินฺนุ\-ปาทานิโย ธมฺโม อุปฺปชฺชติ เหตุปจฺจยา.  อนุปาทินฺน\-อนุปาทานิยํ ธมฺมํ ปฏิจฺจ นอุปาทินฺนุปาทานิโย จ นอนุปาทินฺน\-อนุปาทานิโย จ ธมฺมา อุปฺปชฺชนฺติ เหตุปจฺจยา.  อนุปาทินฺน\-อนุปาทานิยํ ธมฺมํ ปฏิจฺจ นอุปาทินฺนุปาทานิโย จ นอนุปาทินฺนุ\-ปาทานิโย จ ธมฺมา อุปฺปชฺชนฺติ เหตุปจฺจยา. (5)}

\prose{อนุปาทินฺนุ\-ปาทานิยญฺจ อนุปาทินฺน\-อนุปาทานิยญฺจ ธมฺมํ ปฏิจฺจ นอุปาทินฺนุปาทานิโย ธมฺโม อุปฺปชฺชติ เหตุปจฺจยา.  อนุปาทินฺนุ\-ปาทานิยญฺจ อนุปาทินฺน\-อนุปาทานิยญฺจ ธมฺมํ ปฏิจฺจ นอนุปาทินฺน\-อนุปาทานิโย ธมฺโม อุปฺปชฺชติ เหตุปจฺจยา.  อนุปาทินฺนุ\-ปาทานิยญฺจ อนุปาทินฺน\-อนุปาทานิยญฺจ ธมฺมํ ปฏิจฺจ}

\vspace{\tipitakaparskip}
\csromancenter{อุปาทินฺนตฺติก นอุปาทินฺนุปาทานิโย จ นอนุปาทินฺน\-อนุปาทานิโย จ ธมฺมา อุปฺปชฺชนฺติ เหตุปจฺจยา. (3)}
\prose{\csromanfolio{183}อุปาทินฺนุ\-ปาทานิยญฺจ อนุปาทินฺนุ\-ปาทานิยญฺจ ธมฺมํ ปฏิจฺจ นอุปาทินฺนุปาทานิโย ธมฺโม อุปฺปชฺชติ เหตุปจฺจยา.  อุปาทินฺนุ\-ปาทานิยญฺจ อนุปาทินฺนุ\-ปาทานิยญฺจ ธมฺมํ ปฏิจฺจ นอนุปาทินฺน\-อนุปาทานิโย ธมฺโม อุปฺปชฺชติ เหตุปจฺจยา.  อุปาทินฺนุ\-ปาทานิยญฺจ อนุปาทินฺนุ\-ปาทานิยญฺจ ธมฺมํ ปฏิจฺจ นอุปาทินฺนุปาทานิโย จ นอนุปาทินฺน\-อนุปาทานิโย จ ธมฺมา อุปฺปชฺชนฺติ เหตุปจฺจยา. (3) (สํขิตฺตํ.)}

\proseitem{27}{เหตุยา เอกูนวีส, อารมฺมเณ นว, อธิปติยา เอกาทส ฯเปฯ สหชาเต เอกูนวีส.}

\vspace{\tipitakaparskip}
\csromancenter{(สหชาตวารมฺปิ ปจฺจยวารมฺปิ นิสฺสยวารมฺปิ สํสฏฺฐ\-วารมฺปิ สมฺปยุตฺตวารมฺปิ วิตฺถาเรตพฺพํ.)}
\csromanmarkh{4. อุปาทินฺนตฺติก}\csromanmarkhii{7. ปญฺหาวาร}\csromanheadernomark{2}{4. อุปาทินฺนตฺติก 7. ปญฺหาวาร}
\prose{ปจฺจยจตุกฺกเหตุอารมฺมณ}

\proseitem{28}{อุปาทินฺนุ\-ปาทานิโย ธมฺโม นอุปาทินฺนุปาทานิยสฺส ธมฺมสฺส เหตุปจฺจเยน ปจฺจโย.  อุปาทินฺนุ\-ปาทานิโย ธมฺโม นอนุปาทินฺนุ\-ปาทานิยสฺส ธมฺมสฺส เหตุปจฺจเยน ปจฺจโย ฯเปฯ.}

\proseitem{29}{อุปาทินฺนุ\-ปาทานิโย ธมฺโม นอุปาทินฺนุปาทานิยสฺส ธมฺมสฺส อารมฺมณ\-ปจฺจเยน ปจฺจโย.  อุปาทินฺนุ\-ปาทานิโย ธมฺโม นอนุปาทินฺนุ\-ปาทานิยสฺส ธมฺมสฺส อารมฺมณ\-ปจฺจเยน ปจฺจโย.  อุปาทินฺนุ\-ปาทานิโย ธมฺโม นอนุปาทินฺน\-อนุปาทานิยสฺส ธมฺมสฺส อารมฺมณ\-ปจฺจเยน ปจฺจโย.  อุปาทินฺนุ\-ปาทานิโย ธมฺโม นอุปาทินฺนุปาทานิยสฺส จ นอนุปาทินฺน\-อนุปาทานิยสฺส จ ธมฺมสฺส อารมฺมณ\-ปจฺจเยน ปจฺจโย.  อุปาทินฺนุ\-ปาทานิโย ธมฺโม นอนุปาทินฺนุ\-ปาทานิยสฺส จ นอนุปาทินฺน\-อนุปาทานิยสฺส จ ธมฺมสฺส อารมฺมณ\-ปจฺจเยน ปจฺจโย. (5)}

\prose{อนุปาทินฺนุ\-ปาทานิโย ธมฺโม นอนุปาทินฺนุ\-ปาทานิยสฺส ธมฺมสฺส อารมฺมณ\-ปจฺจเยน ปจฺจโย.  อนุปาทินฺนุ\-ปาทานิโย ธมฺโม นอุปาทินฺนุปาทานิยสฺส ธมฺมสฺส อารมฺมณ\-ปจฺจเยน ปจฺจโย.  อนุปาทินฺนุ\-ปาทานิโย ธมฺโม นอนุปาทินฺน\-อนุปาทานิยสฺส ธมฺมสฺส อารมฺมณ\-ปจฺจเยน ปจฺจโย.   \csromanfolio{184}อนุปาทินฺนุ\-ปาทานิโย ธมฺโม นฺเอาปาทินฺนุ\-ปาทานิยสฺส จ นอนุปาทินฺน\-อนุปาทานิยสฺส จ ธมฺมสฺส อารมฺมณ\-ปจฺจเยน ปจฺจโย.  อนุปาทินฺนุ\-ปาทานิโย ธมฺโม นอนุปาทินฺนุ\-ปาทานิยสฺส จ นอนุปาทินฺน\-อนุปาทานิยสฺส จ ธมฺมสฺส อารมฺมณ\-ปจฺจเยน ปจฺจโย. (5)}

\prose{อนุปาทินฺน\-อนุปาทานิโย ธมฺโม นอนุปาทินฺน\-อนุปาทานิยสฺส ธมฺมสฺส อารมฺมณ\-ปจฺจเยน ปจฺจโย ฯเปฯ. ปญฺจ. (สํขิตฺตํ.)}

\proseitem{30}{เหตุยา เตรส, อารมฺมเณ ปนฺนรส, อธิปติยา เอกาทส. (ปญฺหาวารมฺปิ วิตฺถาเรตพฺพํ.)}

\csromanmarkh{5. สํกิลิฏฺฐ\-ตฺติก}\csromanmarkhii{1. ปฏิจฺจวาร}\csromanheadernomark{2}{5. สํกิลิฏฺฐ\-ตฺติก 1. ปฏิจฺจวาร}
\prose{1. ปจฺจยานุโลมเหตุ สํกิลิฏฺฐ\-สํกิเลสิกํ ธมฺมํ ปฏิจฺจ นสํกิลิฏฺฐ\-สํกิเลสิโก ธมฺโม อุปฺปชฺชติ เหตุปจฺจยา.  สํกิลิฏฺฐ\-สํกิเลสิกํ ธมฺมํ ปฏิจฺจ นอสํกิลิฏฺฐ\-สํกิเลสิโก ธมฺโม อุปฺปชฺชติ เหตุปจฺจยา.  สํกิลิฏฺฐ\-สํกิเลสิกํ ธมฺมํ ปฏิจฺจ นอสํกิลิฏฺฐ\-อสํกิเลสิโก ธมฺโม อุปฺปชฺชติ เหตุปจฺจยา.  สํกิลิฏฺฐ\-สํกิเลสิกํ ธมฺมํ ปฏิจฺจ นสํกิลิฏฺฐ\-สํกิเลสิโก จ นอสํกิลิฏฺฐ\-อสํกิเลสิโก จ ธมฺมา อุปฺปชฺชนฺติ เหตุปจฺจยา.  สํกิลิฏฺฐ\-สํกิเลสิกํ ธมฺมํ ปฏิจฺจ นอสํกิลิฏฺฐ\-สํกิเลสิโก จ นอสํกิลิฏฺฐ\-อสํกิเลสิโก จ ธมฺมา อุปฺปชฺชนฺติ เหตุปจฺจยา. (5) อสํกิลิฏฺฐ\-สํกิเลสิกํ ธมฺมํ ปฏิจฺจ นสํกิลิฏฺฐ\-สํกิเลสิโก ธมฺโม อุปฺปชฺชติ เหตุปจฺจยา.  อสํกิลิฏฺฐ\-สํกิเลสิกํ ธมฺมํ ปฏิจฺจ นอสํกิลิฏฺฐ\-อสํกิเลสิโก ธมฺโม อุปฺปชฺชติ เหตุปจฺจยา.  อสํกิลิฏฺฐ\-สํกิเลสิกํ ธมฺมํ ปฏิจฺจ นสํกิลิฏฺฐ\-สํกิเลสิโก จ นอสํกิลิฏฺฐ\-อสํกิเลสิโก จ ธมฺมา อุปฺปชฺชนฺติ เหตุปจฺจยา. (3)}

\proseitem{31}{อสํกิลิฏฺฐ\-อสํกิเลสิกํ ธมฺมํ ปฏิจฺจ นอสํกิลิฏฺฐ\-อสํกิเลสิโก ธมฺโม อุปฺปชฺชติ เหตุปจฺจยา.  อสํกิลิฏฺฐ\-อสํกิเลสิกํ ธมฺมํ ปฏิจฺจ นสํกิลิฏฺฐ\-สํกิเลสิโก ธมฺโม อุปฺปชฺชติ เหตุปจฺจยา.  อสํกิลิฏฺฐ\-อสํกิเลสิกํ ธมฺมํ ปฏิจฺจ นอสํกิลิฏฺฐ\-สํกิเลสิโก ธมฺโม อุปฺปชฺชติ}

\vspace{\tipitakaparskip}
\csromancenter{วิตกฺกตฺติก เหตุปจฺจยา.  อสํกิลิฏฺฐ\-อสํกิเลสิกํ ธมฺมํ ปฏิจฺจ นสํกิลิฏฺฐ\-สํกิเลสิโก จ นอสํกิลิฏฺฐ\-อสํกิเลสิโก จ ธมฺมา อุปฺปชฺชนฺติ เหตุปจฺจยา.  อสํกิลิฏฺฐ\-อสํกิเลสิกํ ธมฺมํ ปฏิจฺจ นสํกิลิฏฺฐ\-สํกิเลสิโก จ นอสํกิลิฏฺฐ\-สํกิเลสิโก จ ธมฺมา อุปฺปชฺชนฺติ เหตุปจฺจยา. (5)}
\prose{\csromanfolio{185}อสํกิลิฏฺฐ\-สํกิเลสิกญฺจ อสํกิลิฏฺฐ\-อสํกิเลสิกญฺจ ธมฺมํ ปฏิจฺจ นสํกิลิฏฺฐ\-สํกิเลสิโก ธมฺโม อุปฺปชฺชติ เหตุปจฺจยา. ตีณิ.}

\prose{สํกิลิฏฺฐ\-สํกิเลสิกญฺจ อสํกิลิฏฺฐ\-สํกิเลสิกญฺจ ธมฺมํ ปฏิจฺจ นสํกิลิฏฺฐ\-สํกิเลสิโก ธมฺโม อุปฺปชฺชติ เหตุปจฺจยา ฯเปฯ. ตีณิ.}

\prose{เหตุยา เอกูนวีส, อารมฺมเณ นว ฯเปฯ อวิคเต เอกูนวีส. (สหชาตวารมฺปิ ปจฺจยวารมฺปิ นิสฺสยวารมฺปิ สํสฏฺฐ\-วารมฺปิ สมฺปยุตฺตวารมฺปิ วิตฺถาเรตพฺพํ. ปญฺหาวาเร น สทิสํ.)}

\proseitem{32}{เหตุยา เตรส, อารมฺมเณ ปนฺนรส, อธิปติยา ปนฺนรส, อนนฺตเร โสฬส ฯเปฯ ปุเรชาเต ฉ, ปจฺฉาชาเต นว, อาเสวเน อฏฺฐ, กมฺเม เตรส, วิปาเก อฏฺฐ, อาหาเร เตรส ฯเปฯ มคฺเค เตรส, วิปฺปยุตฺเต ทฺวาทส ฯเปฯ อวิคเต ฉพฺพีส.}

\csromanmarkh{6. วิตกฺกตฺติก}\csromanmarkhii{1. ปฏิจฺจวาร}\csromanheadernomark{2}{6. วิตกฺกตฺติก 1. ปฏิจฺจวาร}
\prose{1. ปจฺจยานุโลมเหตุ}

\proseitem{33}{สวิตกฺก\-สวิจารํ ธมฺมํ ปฏิจฺจ นสวิตกฺก\-สวิจาโร ธมฺโม อุปฺปชฺชติ เหตุปจฺจยา.  สวิตกฺก\-สวิจารํ ธมฺมํ ปฏิจฺจ นอวิตกฺก\-วิจารมตฺโต ธมฺโม อุปฺปชฺชติ เหตุปจฺจยา.  สวิตกฺก\-สวิจารํ ธมฺมํ ปฏิจฺจ นอวิตกฺก\-อวิจาโร ธมฺโม อุปฺปชฺชติ เหตุปจฺจยา.  สวิตกฺก\-สวิจารํ ธมฺมํ ปฏิจฺจ นสวิตกฺก\-สวิจาโร จ นอวิตกฺก\-อวิจาโร จ ธมฺมา อุปฺปชฺชนฺติ เหตุปจฺจยา.  สวิตกฺก\-สวิจารํ ธมฺมํ ปฏิจฺจ นอวิตกฺก\-วิจารมตฺโต จ นอวิตกฺก\-อวิจาโร จ ธมฺมา อุปฺปชฺชนฺติ เหตุปจฺจยา.  สวิตกฺก\-สวิจารํ ธมฺมํ ปฏิจฺจ นสวิตกฺก\-สวิจาโร จ นอวิตกฺก\-วิจารมตฺโต จ ธมฺมา อุปฺปชฺชนฺติ เหตุปจฺจยา.  สวิตกฺก\-สวิจารํ ธมฺมํ ปฏิจฺจ นสวิตกฺก\-สวิจาโร จ นอวิตกฺก\-วิจารมตฺโต จ นอวิตกฺก\-อวิจาโร จ ธมฺมา อุปฺปชฺชนฺติ เหตุปจฺจยา. (7)}

\proseitem{34}{\csromanfolio{186}อวิตกฺก\-วิจารมตฺตํ ธมฺมํ ปฏิจฺจ นอวิตกฺก\-วิจารมตฺโต ธมฺโม อุปฺปชฺชติ เหตุปจฺจยา.  อวิตกฺก\-วิจารมตฺตํ ธมฺมํ ปฏิจฺจ นสวิตกฺก\-สวิจาโร ธมฺโม อุปฺปชฺชติ เหตุปจฺจยา.  อวิตกฺก\-วิจารมตฺตํ ธมฺมํ ปฏิจฺจ นอวิตกฺก\-อวิจาโร ธมฺโม อุปฺปชฺชติ เหตุปจฺจยา.  อวิตกฺก\-วิจารมตฺตํ ธมฺมํ ปฏิจฺจ นสวิตกฺก\-สวิจาโร จ นอวิตกฺก\-อวิจาโร จ ธมฺมา อุปฺปชฺชนฺติ เหตุปจฺจยา.  อวิตกฺก\-วิจารมตฺตํ ธมฺมํ ปฏิจฺจ นอวิตกฺก\-วิจารมตฺโต จ นอวิตกฺก\-อวิจาโร จ ธมฺมา อุปฺปชฺชนฺติ เหตุปจฺจยา.  อวิตกฺก\-วิจารมตฺตํ ธมฺมํ ปฏิจฺจ นสวิตกฺก\-สวิจาโร จ นอวิตกฺก\-วิจารมตฺโต จ ธมฺมา อุปฺปชฺชนฺติ เหตุปจฺจยา.  อวิตกฺก\-วิจารมตฺตํ ธมฺมํ ปฏิจฺจ นสวิตกฺก\-สวิจาโร จ นอวิตกฺก\-วิจารมตฺโต จ นอวิตกฺก\-อวิจาโร จ ธมฺมา อุปฺปชฺชนฺติ เหตุปจฺจยา. (7)}

\prose{อวิตกฺก\-อวิจารํ ธมฺมํ ปฏิจฺจ นอวิตกฺก\-อวิจาโร ธมฺโม อุปฺปชฺชติ เหตุปจฺจยา ฯเปฯ. สตฺต.}

\prose{สวิตกฺกสวิจารญฺจ อวิตกฺกอวิจารญฺจ ธมฺมํ ปฏิจฺจ นสวิตกฺก\-สวิจาโร ธมฺโม อุปฺปชฺชติ เหตุปจฺจยา ฯเปฯ. สตฺต.}

\prose{อวิตกฺกวิจารมตฺตญฺจ อวิตกฺกอวิจารญฺจ ธมฺมํ ปฏิจฺจ นสวิตกฺก\-สวิจาโร ธมฺโม อุปฺปชฺชติ เหตุปจฺจยา ฯเปฯ. สตฺต.}

\prose{สวิตกฺกสวิจารญฺจ อวิตกฺกวิจารมตฺตญฺจ ธมฺมํ ปฏิจฺจ นสวิตกฺก\-สวิจาโร ธมฺโม อุปฺปชฺชติ เหตุปจฺจยา ฯเปฯ. สตฺต.}

\prose{สวิตกฺกสวิจารญฺจ อวิตกฺกวิจารมตฺตญฺจ อวิตกฺกอวิจารญฺจ ธมฺมํ ปฏิจฺจ นสวิตกฺก\-สวิจาโร ธมฺโม อุปฺปชฺชติ เหตุปจฺจยา ฯเปฯ. สตฺต. (สํขิตฺตํ.)}

\prose{เหตุยา เอกูนปญฺญาส, อารมฺมเณ เอกูนปญฺญาส ฯเปฯ อวิคเต เอกูนปญฺญาส.}

\vspace{\tipitakaparskip}
\csromancenter{(สหชาตวารมฺปิ ฯเปฯ ปญฺหาวารมฺปิ วิตฺถาเรตพฺพํ.)}
\csromanmarkh{7. ปีติตฺติก}\csromanmarkhii{1. ปฏิจฺจวาร}\csromanheadernomark{2}{7. \textbf{ปีติตฺติก 1. ปฏิจฺจวาร}}
\prose{1. ปจฺจยานุโลมเหตุ}

\proseitem{35}{ปีติสหคตํ ธมฺมํ ปฏิจฺจ นปีติสหคโต ธมฺโม อุปฺปชฺชติ เหตุปจฺจยา.  ปีติสหคตํ ธมฺมํ ปฏิจฺจ นสุขสหคโต ธมฺโม อุปฺปชฺชติ เหตุปจฺจยา.  ปีติสหคตํ ธมฺมํ ปฏิจฺจ นอุเปกฺขาสหคโต ธมฺโม}

\vspace{\tipitakaparskip}
\csromancenter{ทสฺสนตฺติก อุปฺปชฺชติ เหตุปจฺจยา.  ปีติสหคตํ ธมฺมํ ปฏิจฺจ นปีติสหคโต จ นอุเปกฺขาสหคโต จ ธมฺมา อุปฺปชฺชนฺติ เหตุปจฺจยา.  ปีติสหคตํ ธมฺมํ ปฏิจฺจ นสุขสหคโต จ นอุเปกฺขาสหคโต จ ธมฺมา อุปฺปชฺชนฺติ เหตุปจฺจยา.  ปีติสหคตํ ธมฺมํ ปฏิจฺจ นปีติสหคโต จ นสุขสหคโต จ ธมฺมา อุปฺปชฺชนฺติ เหตุปจฺจยา.  ปีติสหคตํ ธมฺมํ ปฏิจฺจ นปีติสหคโต จ นสุขสหคโต จ นอุเปกฺขาสหคโต จ ธมฺมา อุปฺปชฺชนฺติ เหตุปจฺจยา. (7)}
\prose{\csromanfolio{187}สุขสหคตํ ธมฺมํ ปฏิจฺจ นสุขสหคโต ธมฺโม อุปฺปชฺชติ เหตุปจฺจยา ฯเปฯ. สตฺต.}

\prose{อุเปกฺขา\-สหคตํ ธมฺมํ ปฏิจฺจ นอุเปกฺขาสหคโต ธมฺโม อุปฺปชฺชติ เหตุปจฺจยา ฯเปฯ. สตฺต.}

\prose{ปีติสหคตญฺจ สุขสหคตญฺจ ธมฺมํ ปฏิจฺจ นปีติสหคโต ธมฺโม อุปฺปชฺชติ เหตุปจฺจยา ฯเปฯ. สตฺต. (สํขิตฺตํ.)}

\prose{เหตุยา อฏฺฐวีส, อารมฺมเณ จตุวีส ฯเปฯ อวิคเต อฏฺฐวีส.}

\prose{(สหชาตวารมฺปิ ฯเปฯ สมฺปยุตฺตวารมฺปิ ปญฺหาวารมฺปิ สพฺพตฺถ วิตฺถาเรตพฺพํ.)}

\csromanmarkh{8. ทสฺสนตฺติก}\csromanmarkhii{1. ปฏิจฺจวาร}\csromanheadernomark{2}{8. ทสฺสนตฺติก 1. ปฏิจฺจวาร}
\prose{1. ปจฺจยานุโลมเหตุ}

\proseitem{36}{ทสฺสเนน ปหาตพฺพํ ธมฺมํ ปฏิจฺจ นทสฺสเนน ปหาตพฺโพ ธมฺโม อุปฺปชฺชติ เหตุปจฺจยา.  ทสฺสเนน ปหาตพฺพํ ธมฺมํ ปฏิจฺจ นภาวนาย ปหาตพฺโพ ธมฺโม อุปฺปชฺชติ เหตุปจฺจยา.  ทสฺสเนน ปหาตพฺพํ ธมฺมํ ปฏิจฺจ นเน\-ว\-ทสฺสเนน นภาวนาย ปหาตพฺโพ ธมฺโม อุปฺปชฺชติ เหตุปจฺจยา.  ทสฺสเนน ปหาตพฺพํ ธมฺมํ ปฏิจฺจ นภาวนาย ปหาตพฺโพ จ นเน\-ว\-ทสฺสเนน นภาวนาย ปหาตพฺโพ จ ธมฺมา อุปฺปชฺชนฺติ เหตุปจฺจยา.  ทสฺสเนน ปหาตพฺพํ ธมฺมํ ปฏิจฺจ นทสฺสเนน ปหาตพฺโพ จ นภาวนาย ปหาตพฺโพ จ ธมฺมา อุปฺปชฺชนฺติ เหตุปจฺจยา. (5)}

\prose{ภาวนาย ปหาตพฺพํ ธมฺมํ ปฏิจฺจ นภาวนาย ปหาตพฺโพ ธมฺโม อุปฺปชฺชติ เหตุปจฺจยา. ปญฺจ.}

\prose{\csromanfolio{188}เนวทสฺสเนน นภาวนาย ปหาตพฺพํ ธมฺมํ ปฏิจฺจ นทสฺสเนน ปหาตพฺโพ ธมฺโม อุปฺปชฺชติ เหตุปจฺจยา. ตีณิ.}

\prose{ทสฺสเนน ปหาตพฺพญฺจ เนวทสฺสเนน นภาวนาย ปหาตพฺพญฺจ ธมฺมํ ปฏิจฺจ นทสฺสเนน ปหาตพฺโพ ธมฺโม อุปฺปชฺชติ เหตุปจฺจยา. ตีณิ.}

\prose{ภาวนาย ปหาตพฺพญฺจ เนวทสฺสเนน นภาวนาย ปหาตพฺพญฺจ ธมฺมํ ปฏิจฺจ นทสฺสเนน ปหาตพฺโพ ธมฺโม อุปฺปชฺชติ เหตุปจฺจยา. ตีณิ. (สํขิตฺตํ.)}

\prose{เหตุยา เอกูนวีส ฯเปฯ อวิคเต เอกูนวีส.}

\vspace{\tipitakaparskip}
\csromancenter{(สหชาตวารมฺปิ ฯเปฯ สมฺปยุตฺตวารมฺปิ ปญฺหาวารมฺปิ วิตฺถาเรตพฺพํ.)}
\csromanmarkh{9. ทสฺสน\-เหตุ\-ตฺติก}\csromanmarkhii{1. ปฏิจฺจวาร}\csromanheadernomark{2}{9. ทสฺสน\-เหตุ\-ตฺติก 1. ปฏิจฺจวาร}
\proseitem{37}{ทสฺสเนน ปหาตพฺพ\-เหตุกํ ธมฺมํ ปฏิจฺจ นทสฺสเนน ปหาตพฺพ\-เหตุโก ธมฺโม อุปฺปชฺชติ เหตุปจฺจยา. (สํขิตฺตํ.)}

\prose{เหตุยา ฉพฺพีส ฯเปฯ อวิคเต ฉพฺพีส. (วิตฺถาเรตพฺพํ.)}

\csromanmarkh{10. อาจยคามิตฺติก}\csromanmarkhii{1. ปฏิจฺจวาร}\csromanheadernomark{2}{10. \textbf{อาจยคามิตฺติก 1. ปฏิจฺจวาร}}
\prose{1. ปจฺจยานุโลมเหตุ}

\proseitem{38}{อาจยคามิํ ธมฺมํ ปฏิจฺจ นอาจยคามี ธมฺโม อุปฺปชฺชติ เหตุปจฺจยา.  อาจยคามิํ ธมฺมํ ปฏิจฺจ นอปจยคามี ธมฺโม อุปฺปชฺชติ เหตุปจฺจยา.  อาจยคามิํ ธมฺมํ ปฏิจฺจ นเน\-วา\-จยคามิ\-นา\-ปจยคามี ธมฺโม อุปฺปชฺชติ เหตุปจฺจยา.  อาจยคามิํ ธมฺมํ ปฏิจฺจ นอปจยคามี จ นเน\-วา\-จยคามิ\-นา\-ปจยคามี จ ธมฺมา อุปฺปชฺชนฺติ เหตุปจฺจยา.  อาจยคามิํ ธมฺมํ ปฏิจฺจ นอาจยคามี จ นอปจยคามี จ ธมฺมา อุปฺปชฺชนฺติ เหตุปจฺจยา. (5)}

\prose{อปจยคามิํ ธมฺมํ ปฏิจฺจ นอปจยคามี ธมฺโม อุปฺปชฺชติ เหตุปจฺจยา.  อปจยคามิํ ธมฺมํ ปฏิจฺจ นอาจยคามี ธมฺโม อุปฺปชฺชติ เหตุปจฺจยา.  อปจยคามิํ ธมฺมํ ปฏิจฺจ นเน\-วา\-จยคามิ\-นา\-ปจยคามี ธมฺโม อุปฺปชฺชติ เหตุปจฺจยา.  อปจยคามิํ ธมฺมํ ปฏิจฺจ นอาจยคามี จ นเน\-วา\-จยคามิ\-นา\-ปจยคามี จ ธมฺมา อุปฺปชฺชนฺติ เหตุปจฺจยา.  อปจยคามิํ ธมฺมํ ปฏิจฺจ นอาจยคามี จ นอปจยคามี จ ธมฺมา อุปฺปชฺชนฺติ เหตุปจฺจยา. (5)}

\csromanheader{2}{11. เสกฺขตฺติก}
\prose{\csromanfolio{189}เนวา\-จยคามิ\-นา\-ปจยคามิํ ธมฺมํ ปฏิจฺจ นอาจยคามี ธมฺโม อุปฺปชฺชติ เหตุปจฺจยา.  เนวา\-จยคามิ\-นา\-ปจยคามิํ ธมฺมํ ปฏิจฺจ นอปจยคามี ธมฺโม อุปฺปชฺชติ เหตุปจฺจยา.  เนวา\-จยคามิ\-นา\-ปจยคามิํ ธมฺมํ ปฏิจฺจ นอาจยคามี จ นอปจยคามี จ ธมฺมา อุปฺปชฺชนฺติ เหตุปจฺจยา. (3)}

\prose{อาจยคามิญฺจ เนวา\-จยคามิ\-นา\-ปจยคามิญฺจ ธมฺมํ ปฏิจฺจ นอาจยคามี ธมฺโม อุปฺปชฺชติ เหตุปจฺจยา.  อาจยคามิญฺจ เนวา\-จยคามิ\-นา\-ปจยคามิญฺจ ธมฺมํ ปฏิจฺจ นอปจยคามี ธมฺโม อุปฺปชฺชติ เหตุปจฺจยา.  อาจยคามิญฺจ เนวา\-จยคามิ\-นา\-ปจยคามิญฺจ ธมฺมํ ปฏิจฺจ นอาจยคามี จ นอปจยคามี จ ธมฺมา อุปฺปชฺชนฺติ เหตุปจฺจยา. (3)}

\prose{อปจยคามิญฺจ เนวา\-จยคามิ\-นา\-ปจยคามิญฺจ ธมฺมํ ปฏิจฺจ นอาจยคามี ธมฺโม อุปฺปชฺชติ เหตุปจฺจยา.  อปจยคามิญฺจ เนวา\-จยคามิ\-นา\-ปจยคามิญฺจ ธมฺมํ ปฏิจฺจ นอปจยคามี ธมฺโม อุปฺปชฺชติ เหตุปจฺจยา.  อปจยคามิญฺจ เนวา\-จยคามิ\-นา\-ปจยคามิญฺจ ธมฺมํ ปฏิจฺจ นอาจยคามี จ นอปจยคามี จ ธมฺมา อุปฺปชฺชนฺติ เหตุปจฺจยา. (3) (สํขิตฺตํ.)}

\prose{เหตุยา เอกูนวีส. (สพฺพตฺถ วิตฺถาโร.)}

\csromanmarkh{11. เสกฺขตฺติก}\csromanmarkhii{1. ปฏิจฺจวาร}\csromanheadernomark{2}{11. เสกฺขตฺติก 1. ปฏิจฺจวาร}
\prose{1. ปจฺจยานุโลมเหตุ}

\proseitem{39}{เสกฺขํ ธมฺมํ ปฏิจฺจ นเสกฺโข ธมฺโม อุปฺปชฺชติ เหตุปจฺจยา.  เสกฺขํ ธมฺมํ ปฏิจฺจ นอเสกฺโข ธมฺโม อุปฺปชฺชติ เหตุปจฺจยา.  เสกฺขํ ธมฺมํ ปฏิจฺจ นเน\-ว\-เสกฺข\-นา\-เสกฺโข ธมฺโม อุปฺปชฺชติ เหตุปจฺจยา.  เสกฺขํ ธมฺมํ ปฏิจฺจ นอเสกฺโข จ นเน\-ว\-เสกฺข\-นา\-เสกฺโข จ ธมฺมา อุปฺปชฺชนฺติ เหตุปจฺจยา.  เสกฺขํ ธมฺมํ ปฏิจฺจ นเสกฺโข จ นอเสกฺโข จ ธมฺมา อุปฺปชฺชนฺติ เหตุปจฺจยา. (5)}

\prose{อเสกฺขํ ธมฺมํ ปฏิจฺจ นอเสกฺโข ธมฺโม อุปฺปชฺชติ เหตุปจฺจยา.  อเสกฺขํ ธมฺมํ ปฏิจฺจ นเสกฺโข ธมฺโม อุปฺปชฺชติ เหตุปจฺจยา.  อเสกฺขํ ธมฺมํ ปฏิจฺจ นเน\-ว\-เสกฺข\-นา\-เสกฺโข ธมฺโม อุปฺปชฺชติ เหตุปจฺจยา.  อเสกฺขํ ธมฺมํ ปฏิจฺจ นเสกฺโข จ นเน\-ว\-เสกฺข\-นา\-เสกฺโข จ ธมฺมา อุปฺปชฺชนฺติ เหตุปจฺจยา.  อเสกฺขํ ธมฺมํ ปฏิจฺจ นเสกฺโข จ นอเสกฺโข จ ธมฺมา อุปฺปชฺชนฺติ เหตุปจฺจยา. (5)}

\prose{\csromanfolio{190}เนว\-เสกฺข\-นาเสกฺขํ ธมฺมํ ปฏิจฺจ นเสกฺโข ธมฺโม อุปฺปชฺชติ เหตุปจฺจยา.  เนว\-เสกฺข\-นาเสกฺขํ ธมฺมํ ปฏิจฺจ นอเสกฺโข ธมฺโม อุปฺปชฺชติ เหตุปจฺจยา.  เนว\-เสกฺข\-นาเสกฺขํ ธมฺมํ ปฏิจฺจ นเสกฺโข จ นอเสกฺโข จ ธมฺมา อุปฺปชฺชนฺติ เหตุปจฺจยา. (3)}

\prose{เสกฺขญฺจ เนว\-เสกฺข\-นาเสกฺขญฺจ ธมฺมํ ปฏิจฺจ นเสกฺโข ธมฺโม อุปฺปชฺชติ เหตุปจฺจยา.  เสกฺขญฺจ เนว\-เสกฺข\-นาเสกฺขญฺจ ธมฺมํ ปฏิจฺจ นอเสกฺโข ธมฺโม อุปฺปชฺชติ เหตุปจฺจยา.  เสกฺขญฺจ เนว\-เสกฺข\-นาเสกฺขญฺจ ธมฺมํ ปฏิจฺจ นเสกฺโข จ นอเสกฺโข จ ธมฺมา อุปฺปชฺชนฺติ เหตุปจฺจยา. (3)}

\prose{อเสกฺขญฺจ เนว\-เสกฺข\-นาเสกฺขญฺจ ธมฺมํ ปฏิจฺจ นเสกฺโข ธมฺโม อุปฺปชฺชติ เหตุปจฺจยา.  อเสกฺขญฺจ เนว\-เสกฺข\-นาเสกฺขญฺจ ธมฺมํ ปฏิจฺจ นอเสกฺโข ธมฺโม อุปฺปชฺชติ เหตุปจฺจยา.  อเสกฺขญฺจ เนว\-เสกฺข\-นาเสกฺขญฺจ ธมฺมํ ปฏิจฺจ นเสกฺโข จ นอเสกฺโข จ ธมฺมา อุปฺปชฺชนฺติ เหตุปจฺจยา. (3) (สํขิตฺตํ.)}

\prose{เหตุยา เอกูนวีส. (สพฺพตฺถ วิตฺถาโร.)}

\csromanmarkh{12. ปริตฺตตฺติก}\csromanmarkhii{1. ปฏิจฺจวาร}\csromanheadernomark{2}{12. \textbf{ปริตฺตตฺติก 1. ปฏิจฺจวาร}}
\prose{1. ปจฺจยานุโลมเหตุ}

\proseitem{40}{ปริตฺตํ ธมฺมํ ปฏิจฺจ นปริตฺโต ธมฺโม อุปฺปชฺชติ เหตุปจฺจยา.  ปริตฺตํ ธมฺมํ ปฏิจฺจ นมหคฺคโต ธมฺโม อุปฺปชฺชติ เหตุปจฺจยา.  ปริตฺตํ ธมฺมํ ปฏิจฺจ นอปฺปมาโณ ธมฺโม อุปฺปชฺชติ เหตุปจฺจยา.  ปริตฺตํ ธมฺมํ ปฏิจฺจ นปริตฺโต จ นอปฺปมาโณ จ ธมฺมา อุปฺปชฺชนฺติ เหตุปจฺจยา.  ปริตฺตํ ธมฺมํ ปฏิจฺจ นมหคฺคโต จ นอปฺปมาโณ จ ธมฺมา อุปฺปชฺชนฺติ เหตุปจฺจยา. (5)}

\prose{มหคฺคตํ ธมฺมํ ปฏิจฺจ นมหคฺคโต ธมฺโม อุปฺปชฺชติ เหตุปจฺจยา.  มหคฺคตํ ธมฺมํ ปฏิจฺจ นปริตฺโต ธมฺโม อุปฺปชฺชติ เหตุปจฺจยา.  มหคฺคตํ ธมฺมํ ปฏิจฺจ นอปฺปมาโณ ธมฺโม อุปฺปชฺชติ เหตุปจฺจยา.  มหคฺคตํ ธมฺมํ ปฏิจฺจ นปริตฺโต จ นอปฺปมาโณ จ ธมฺมา อุปฺปชฺชนฺติ เหตุปจฺจยา.  มหคฺคตํ ธมฺมํ ปฏิจฺจ นมหคฺคโต จ นอปฺปมาโณ จ ธมฺมา อุปฺปชฺชนฺติ เหตุปจฺจยา. (5)}

\prose{อปฺปมาณํ ธมฺมํ ปฏิจฺจ นอปฺปมาโณ ธมฺโม อุปฺปชฺชติ เหตุปจฺจยา.  อปฺปมาณํ ธมฺมํ ปฏิจฺจ นปริตฺโต ธมฺโม อุปฺปชฺชติ เหตุปจฺจยา.  อปฺปมาณํ ธมฺมํ}

\vspace{\tipitakaparskip}
\csromancenter{มิจฺฉตฺตตฺติก ปฏิจฺจ นมหคฺคโต ธมฺโม อุปฺปชฺชติ เหตุปจฺจยา.  อปฺปมาณํ ธมฺมํ ปฏิจฺจ นมหคฺคโต จ นอปฺปมาโณ จ ธมฺมา อุปฺปชฺชนฺติ เหตุปจฺจยา.  อปฺปมาณํ ธมฺมํ ปฏิจฺจ นปริตฺโต จ นมหคฺคโต จ ธมฺมา อุปฺปชฺชนฺติ เหตุปจฺจยา. (5)}
\prose{\csromanfolio{191}ปริตฺตญฺจ อปฺปมาณญฺจ ธมฺมํ ปฏิจฺจ นมหคฺคโต ธมฺโม อุปฺปชฺชติ เหตุปจฺจยา.  ปริตฺตญฺจ อปฺปมาณญฺจ ธมฺมํ ปฏิจฺจ นอปฺปมาโณ ธมฺโม อุปฺปชฺชติ เหตุปจฺจยา.  ปริตฺตญฺจ อปฺปมาณญฺจ ธมฺมํ ปฏิจฺจ นมหคฺคโต จ นอปฺปมาโณ จ ธมฺมา อุปฺปชฺชนฺติ เหตุปจฺจยา. (3)}

\prose{ปริตฺตญฺจ มหคฺคตญฺจ ธมฺมํ ปฏิจฺจ นปริตฺโต ธมฺโม อุปฺปชฺชติ เหตุปจฺจยา.  ปริตฺตญฺจ มหคฺคตญฺจ ธมฺมํ ปฏิจฺจ นมหคฺคโต ธมฺโม อุปฺปชฺชติ เหตุปจฺจยา.  ปริตฺตญฺจ มหคฺคตญฺจ ธมฺมํ ปฏิจฺจ นอปฺปมาโณ ธมฺโม อุปฺปชฺชติ เหตุปจฺจยา.  ปริตฺตญฺจ มหคฺคตญฺจ ธมฺมํ ปฏิจฺจ นปริตฺโต จ นอปฺปมาโณ จ ธมฺมา อุปฺปชฺชนฺติ เหตุปจฺจยา.  ปริตฺตญฺจ มหคฺคตญฺจ ธมฺมํ ปฏิจฺจ นมหคฺคโต จ นอปฺปมาโณ จ ธมฺมา อุปฺปชฺชนฺติ เหตุปจฺจยา. (5)}

\prose{เหตุยา เตวีส, อารมฺมเณ จุทฺทส. (สพฺพตฺถ วิตฺถาเรตพฺพํ.)}

\csromanmarkh{13. ปริตฺตา\-รมฺมณ\-ตฺติก}\csromanmarkhii{1. ปฏิจฺจวาร}\csromanheadernomark{2}{13. ปริตฺตา\-รมฺมณ\-ตฺติก 1. ปฏิจฺจวาร}
\proseitem{41}{ปริตฺตา\-รมฺมณํ ธมฺมํ ปฏิจฺจ นปริตฺตา\-รมฺมโณ ธมฺโม อุปฺปชฺชติ เหตุปจฺจยา.}

\prose{เหตุยา เอกวีส ฯเปฯ อวิคเต เอกวีส.}

\csromanmarkh{14. หีนตฺติก}\csromanmarkhii{1. ปฏิจฺจวาร}\csromanheadernomark{2}{14. \textbf{หีนตฺติก 1. ปฏิจฺจวาร}}
\proseitem{42}{หีนํ ธมฺมํ ปฏิจฺจ นหีโน ธมฺโม อุปฺปชฺชติ เหตุปจฺจยา. (สํกิลิฏฺฐ\-สํกิเลสิก\-ตฺติก\-สทิสํ.)}

\prose{เหตุยา เอกูนวีส ฯเปฯ อวิคเต เอกูนวีส.}

\csromanmarkh{15. มิจฺฉตฺตตฺติก}\csromanmarkhii{1. ปฏิจฺจวาร}\csromanheadernomark{2}{15. มิจฺฉตฺตตฺติก 1. ปฏิจฺจวาร}
\prose{1. ปจฺจยานุโลมเหตุ}

\proseitem{43}{มิจฺฉตฺต\-นิยตํ ธมฺมํ ปฏิจฺจ นมิจฺฉตฺต\-นิยโต ธมฺโม อุปฺปชฺชติ เหตุปจฺจยา.  มิจฺฉตฺต\-นิยตํ ธมฺมํ ปฏิจฺจ นสมฺมตฺตนิยโต ธมฺโม อุปฺปชฺชติ}

\prose{\csromanfolio{192}เหตุปจฺจยา.  มิจฺฉตฺต\-นิยตํ ธมฺมํ ปฏิจฺจ นอนิยโต ธมฺโม อุปฺปชฺชติ เหตุปจฺจยา.  มิจฺฉตฺต\-นิยตํ ธมฺมํ ปฏิจฺจ นสมฺมตฺตนิยโต จ นอนิยโต จ ธมฺมา อุปฺปชฺชนฺติ เหตุปจฺจยา.  มิจฺฉตฺต\-นิยตํ ธมฺมํ ปฏิจฺจ นมิจฺฉตฺต\-นิยโต จ นสมฺมตฺตนิยโต จ ธมฺมา อุปฺปชฺชนฺติ เหตุปจฺจยา. (5)}

\prose{สมฺมตฺต\-นิยตํ ธมฺมํ ปฏิจฺจ นสมฺมตฺตนิยโต ธมฺโม อุปฺปชฺชติ เหตุปจฺจยา.  สมฺมตฺต\-นิยตํ ธมฺมํ ปฏิจฺจ นมิจฺฉตฺต\-นิยโต ธมฺโม อุปฺปชฺชติ เหตุปจฺจยา.  สมฺมตฺต\-นิยตํ ธมฺมํ ปฏิจฺจ นอนิยโต ธมฺโม อุปฺปชฺชติ เหตุปจฺจยา.  สมฺมตฺต\-นิยตํ ธมฺมํ ปฏิจฺจ นมิจฺฉตฺต\-นิยโต จ นอนิยโต จ ธมฺมา อุปฺปชฺชนฺติ เหตุปจฺจยา.  สมฺมตฺต\-นิยตํ ธมฺมํ ปฏิจฺจ นมิจฺฉตฺต\-นิยโต จ นสมฺมตฺตนิยโต จ ธมฺมา อุปฺปชฺชนฺติ เหตุปจฺจยา. (5)}

\prose{อนิยตํ ธมฺมํ ปฏิจฺจ นมิจฺฉตฺต\-นิยโต ธมฺโม อุปฺปชฺชติ เหตุปจฺจยา. ตีณิ.}

\prose{มิจฺฉตฺตนิยตญฺจ อนิยตญฺจ ธมฺมํ ปฏิจฺจ นมิจฺฉตฺต\-นิยโต ธมฺโม อุปฺปชฺชติ เหตุปจฺจยา. ตีณิ.}

\prose{สมฺมตฺตนิยตญฺจ อนิยตญฺจ ธมฺมํ ปฏิจฺจ นมิจฺฉตฺต\-นิยโต ธมฺโม อุปฺปชฺชติ เหตุปจฺจยา. ตีณิ. (สํขิตฺตํ.)}

\prose{เหตุยา เอกูนวีส. (สพฺพตฺถ วิตฺถาเรตพฺพํ.)}

\csromanmarkh{16. มคฺคา\-รมฺมณ\-ตฺติก}\csromanmarkhii{1. ปฏิจฺจวาร}\csromanheadernomark{2}{16. มคฺคา\-รมฺมณ\-ตฺติก 1. ปฏิจฺจวาร}
\proseitem{44}{มคฺคารมฺมณํ ธมฺมํ ปฏิจฺจ นมคฺคารมฺมโณ ธมฺโม อุปฺปชฺชติ เหตุปจฺจยา.  มคฺคารมฺมณํ ธมฺมํ ปฏิจฺจ นมคฺคเหตุโก ธมฺโม อุปฺปชฺชติ เหตุปจฺจยา.  มคฺคารมฺมณํ ธมฺมํ ปฏิจฺจ นมคฺคาธิปติ ธมฺโม อุปฺปชฺชติ เหตุปจฺจยา.  มคฺคารมฺมณํ ธมฺมํ ปฏิจฺจ นมคฺคารมฺมโณ จ นมคฺคาธิปติ จ ธมฺมา อุปฺปชฺชนฺติ เหตุปจฺจยา.  มคฺคารมฺมณํ ธมฺมํ ปฏิจฺจ นมคฺคเหตุโก จ นมคฺคาธิปติ จ ธมฺมา อุปฺปชฺชนฺติ เหตุปจฺจยา.  มคฺคารมฺมณํ ธมฺมํ ปฏิจฺจ นมคฺคารมฺมโณ จ นมคฺคเหตุโก จ ธมฺมา อุปฺปชฺชนฺติ เหตุปจฺจยา.  มคฺคารมฺมณํ ธมฺมํ ปฏิจฺจ นมคฺคารมฺมโณ จ นมคฺคเหตุโก จ นมคฺคาธิปติ จ ธมฺมา อุปฺปชฺชนฺติ เหตุปจฺจยา. (7) (สํขิตฺตํ.)}

\prose{เหตุยา ปญฺจติํส ฯเปฯ อวิคเต ปญฺจติํส. (สพฺพตฺถ วิตฺถาเรตพฺพํ.)}

\csromanmarkh{20. อชฺฌตฺตตฺติก}\csromanmarkhii{17. อุปฺปนฺนตฺติก 7. ปญฺหาวาร}\csromanheadernomark{2}{20. อชฺฌตฺตตฺติก \textbf{17. อุปฺปนฺนตฺติก 7. ปญฺหาวาร}}
\proseitem{45}{\csromanfolio{193}อุปฺปนฺโน ธมฺโม นอนุปฺปนฺนสฺส ธมฺมสฺส เหตุปจฺจเยน ปจฺจโย. (สํขิตฺตํ.)}

\prose{เหตุยา ตีณิ, อารมฺมเณ นว.}

\csromanmarkh{18. อตีตตฺติก}\csromanmarkhii{7. ปญฺหาวาร}\csromanheadernomark{2}{18. \textbf{อตีตตฺติก 7. ปญฺหาวาร}}
\proseitem{46}{ปจฺจุปฺปนฺโน ธมฺโม นอตีตสฺส ธมฺมสฺส เหตุปจฺจเยน ปจฺจโย. (สํขิตฺตํ.)}

\prose{เหตุยา ตีณิ, อารมฺมเณ นว.}

\csromanmarkh{19. อตีตา\-รมฺมณ\-ตฺติก}\csromanmarkhii{1. ปฏิจฺจวาร}\csromanheadernomark{2}{19. อตีตา\-รมฺมณ\-ตฺติก 1. ปฏิจฺจวาร}
\proseitem{47}{อตีตารมฺมณํ ธมฺมํ ปฏิจฺจ นอตีตารมฺมโณ ธมฺโม อุปฺปชฺชติ เหตุปจฺจยา.  อตีตารมฺมณํ ธมฺมํ ปฏิจฺจ นอนาคตา\-รมฺมโณ ธมฺโม อุปฺปชฺชติ เหตุปจฺจยา.  อตีตารมฺมณํ ธมฺมํ ปฏิจฺจ นปจฺจุปฺปนฺนา\-รมฺมโณ ธมฺโม อุปฺปชฺชติ เหตุปจฺจยา.  อตีตารมฺมณํ ธมฺมํ ปฏิจฺจ นอตีตารมฺมโณ จ นปจฺจุปฺปนฺนา\-รมฺมโณ จ ธมฺมา อุปฺปชฺชนฺติ เหตุปจฺจยา.  อตีตารมฺมณํ ธมฺมํ ปฏิจฺจ นอนาคตา\-รมฺมโณ จ นปจฺจุปฺปนฺนา\-รมฺมโณ จ ธมฺมา อุปฺปชฺชนฺติ เหตุปจฺจยา.  อตีตารมฺมณํ ธมฺมํ ปฏิจฺจ นอตีตารมฺมโณ จ นอนาคตา\-รมฺมโณ จ ธมฺมา อุปฺปชฺชนฺติ เหตุปจฺจยา.  อตีตารมฺมณํ ธมฺมํ ปฏิจฺจ นอตีตารมฺมโณ จ นอนาคตา\-รมฺมโณ จ นปจฺจุปฺปนฺนา\-รมฺมโณ จ ธมฺมา อุปฺปชฺชนฺติ เหตุปจฺจยา. (7) (สํขิตฺตํ.)}

\prose{เหตุยา เอกวีส ฯเปฯ อวิคเต เอกวีส.}

\csromanmarkh{20. อชฺฌตฺตตฺติก}\csromanmarkhii{1. ปฏิจฺจวาร}\csromanheadernomark{2}{20. อชฺฌตฺตตฺติก 1. ปฏิจฺจวาร}
\prose{48. อชฺฌตฺตํ ธมฺมํ ปฏิจฺจ นพหิทฺธา ธมฺโม อุปฺปชฺชติ}

\prose{พหิทฺธา ธมฺมํ ปฏิจฺจ นอชฺฌตฺโต ธมฺโม อุปฺปชฺชติ เหตุปจฺจยา. (1) เหตุยา เทฺว. (สพฺพตฺถ วิตฺถาโร.)}

\csromanmarkh{ธมฺมา\-นุโลม\-ปจฺจนีย ติกปฏฺฐาน\-ปาฬิ}\csromanmarkhii{21. อชฺฌตฺตา\-รมฺมณ\-ตฺติก 1. ปฏิจฺจวาร}\csromanheadernomark{2}{ธมฺมา\-นุโลม\-ปจฺจนีย ติกปฏฺฐาน\-ปาฬิ 21. อชฺฌตฺตา\-รมฺมณ\-ตฺติก 1. ปฏิจฺจวาร}
\proseitem{49}{\csromanfolio{194}อชฺฌตฺตา\-รมฺมณํ ธมฺมํ ปฏิจฺจ นอชฺฌตฺตา\-รมฺมโณ ธมฺโม อุปฺปชฺชติ เหตุปจฺจยา.}

\prose{เหตุยา ฉ.}

\csromanheader{2}{22. สนิทสฺสน\-ตฺติก 1-7. ปฏิจฺจาทิวาร}
\prose{1. ปจฺจยานุโลมเหตุ}

\proseitem{50}{อนิทสฺสน\-สปฺปฏิฆํ ธมฺมํ ปฏิจฺจ นอนิทสฺสน\-สปฺปฏิโฆ ธมฺโม อุปฺปชฺชติ เหตุปจฺจยา.  อนิทสฺสน\-สปฺปฏิฆํ ธมฺมํ ปฏิจฺจ นสนิทสฺสน\-สปฺปฏิโฆ ธมฺโม อุปฺปชฺชติ เหตุปจฺจยา.  อนิทสฺสน\-สปฺปฏิฆํ ธมฺมํ ปฏิจฺจ นอนิทสฺสน\-อปฺปฏิโฆ ธมฺโม อุปฺปชฺชติ เหตุปจฺจยา.  อนิทสฺสน\-สปฺปฏิฆํ ธมฺมํ ปฏิจฺจ นสนิทสฺสน\-สปฺปฏิโฆ จ นอนิทสฺสน\-อปฺปฏิโฆ จ ธมฺมา อุปฺปชฺชนฺติ เหตุปจฺจยา.  อนิทสฺสน\-สปฺปฏิฆํ ธมฺมํ ปฏิจฺจ นอนิทสฺสน\-สปฺปฏิโฆ จ นอนิทสฺสน\-อปฺปฏิโฆ จ ธมฺมา อุปฺปชฺชนฺติ เหตุปจฺจยา.  อนิทสฺสน\-สปฺปฏิฆํ ธมฺมํ ปฏิจฺจ นสนิทสฺสน\-สปฺปฏิโฆ จ นอนิทสฺสน\-สปฺปฏิโฆ จ ธมฺมา อุปฺปชฺชนฺติ เหตุปจฺจยา. (6)}

\prose{อนิทสฺสน\-อปฺปฏิฆํ ธมฺมํ ปฏิจฺจ นอนิทสฺสน\-อปฺปฏิโฆ ธมฺโม อุปฺปชฺชติ เหตุปจฺจยา. ฉ.}

\prose{อนิทสฺสนสปฺปฏิฆญฺจ อนิทสฺสนอปฺปฏิฆญฺจ ธมฺมํ ปฏิจฺจ นสนิทสฺสน\-สปฺปฏิโฆ ธมฺโม อุปฺปชฺชติ เหตุปจฺจยา. ฉ.}

\prose{เหตุยา อฏฺฐารส, อารมฺมเณ ตีณิ ฯเปฯ อวิคเต อฏฺฐารส. (สพฺพตฺถ วิตฺถาโร. สหชาตวารมฺปิ ฯเปฯ สมฺปยุตฺตวารมฺปิ วิตฺถาเรตพฺพํ.)}

\prose{เหตุอารมฺมณ}

\proseitem{51}{อนิทสฺสน\-อปฺปฏิโฆ ธมฺโม นอนิทสฺสน\-อปฺปฏิฆสฺส ธมฺมสฺส เหตุปจฺจเยน ปจฺจโย.  อนิทสฺสน\-อปฺปฏิโฆ ธมฺโม นสนิทสฺสน\-สปฺปฏิฆสฺส ธมฺมสฺส เหตุปจฺจเยน ปจฺจโย.  อนิทสฺสน\-อปฺปฏิโฆ ธมฺโม นอนิทสฺสน\-สปฺปฏิฆสฺส ธมฺมสฺส เหตุปจฺจเยน ปจฺจโย.  }

\vspace{\tipitakaparskip}
\csromancenter{สนิทสฺสน\-ตฺติก อนิทสฺสน\-อปฺปฏิโฆ ธมฺโม นสนิทสฺสน\-สปฺปฏิฆสฺส จ นอนิทสฺสน\-อปฺปฏิฆสฺส จ ธมฺมสฺส เหตุปจฺจเยน ปจฺจโย.  อนิทสฺสน\-อปฺปฏิโฆ ธมฺโม นอนิทสฺสน\-สปฺปฏิฆสฺส จ นอนิทสฺสน\-อปฺปฏิฆสฺส จ ธมฺมสฺส เหตุปจฺจเยน ปจฺจโย.  อนิทสฺสน\-อปฺปฏิโฆ ธมฺโม นสนิทสฺสน\-สปฺปฏิฆสฺส จ นอนิทสฺสน\-สปฺปฏิฆสฺส จ ธมฺมสฺส เหตุปจฺจเยน ปจฺจโย. (6)}
\prose{\csromanfolio{195}สนิทสฺสน\-สปฺปฏิโฆ ธมฺโม นสนิทสฺสน\-สปฺปฏิฆสฺส ธมฺมสฺส อารมฺมณ\-ปจฺจเยน ปจฺจโย. ตีณิ. (สํขิตฺตํ.)}

\proseitemwithcloser{52}{เหตุยา ฉ, อารมฺมเณ นว. (ปญฺหาวารํ วิตฺถาเรตพฺพํ.)}{ธมฺมา\-นุโลม\-ปจฺจนีเย ติกปฏฺฐานํ นิฏฺฐิตํ.}
\pitaka{อภิธมฺม\-ปิฏก}
\namakkaram{นโม ตสฺส ภควโต อรหโต สมฺมา\-สมฺพุทฺธสฺส.}
\csromanmarkh{1. เหตุทุก}\csromanmarkhii{1. ปฏิจฺจาทิวาร}\csromanheadernomark{2}{1. \textbf{เหตุทุก 1. ปฏิจฺจาทิวาร}}
\prose{\csromanfolio{197}ปจฺจยานุโลมเหตุ}

\proseitem{1}{เหตุํ ธมฺมํ ปฏิจฺจ นเหตุ ธมฺโม อุปฺปชฺชติ เหตุปจฺจยา. เหตุํ ธมฺมํ ปฏิจฺจ สมฺปยุตฺตกา ขนฺธา จิตฺต\-สมุฏฺฐานญฺจ รูปํ. ปฏิสนฺธิ\-กฺขเณ ฯเปฯ. เหตุํ ธมฺมํ ปฏิจฺจ นนเหตุ ธมฺโม อุปฺปชฺชติ เหตุปจฺจยา.  เหตุํ ธมฺมํ ปฏิจฺจ นเหตุ จ นนเหตุ จ ธมฺมา อุปฺปชฺชนฺติ เหตุปจฺจยา. (3)}

\prose{นเหตุํ ธมฺมํ ปฏิจฺจ นนเหตุ ธมฺโม อุปฺปชฺชติ เหตุปจฺจยา.  นเหตุํ ธมฺมํ ปฏิจฺจ นเหตุ ธมฺโม อุปฺปชฺชติ เหตุปจฺจยา.  นเหตุํ ธมฺมํ ปฏิจฺจ นเหตุ จ นนเหตุ จ ธมฺมา อุปฺปชฺชนฺติ เหตุปจฺจยา. (3)}

\prose{เหตุญฺจ นเหตุญฺจ ธมฺมํ ปฏิจฺจ นเหตุ ธมฺโม อุปฺปชฺชติ เหตุปจฺจยา.  เหตุญฺจ นเหตุญฺจ ธมฺมํ ปฏิจฺจ นนเหตุ ธมฺโม อุปฺปชฺชติ เหตุปจฺจยา.  เหตุญฺจ นเหตุญฺจ ธมฺมํ ปฏิจฺจ นเหตุ จ นนเหตุ จ ธมฺมา อุปฺปชฺชนฺติ เหตุปจฺจยา. (3) (สํขิตฺตํ.)}

\prose{เหตุยา นว, อารมฺมเณ นว ฯเปฯ อวิคเต นว.}

\vspace{\tipitakaparskip}
\csromancenter{(สหชาตวารมฺปิ ฯเปฯ สมฺปยุตฺตวารมฺปิ ปฏิจฺจ\-วาร\-สทิสํ.)}
\csromanheader{2}{ธมฺมา\-นุโลม\-ปจฺจนีย ทุกปฏฺฐาน\-ปาฬิ เหตุอารมฺมณ}
\proseitem{2}{\csromanfolio{198}เหตุ ธมฺโม นเหตุสฺส ธมฺมสฺส เหตุปจฺจเยน ปจฺจโย. ตีณิ.}

\prose{เหตุ ธมฺโม นเหตุสฺส ธมฺมสฺส อารมฺมณ\-ปจฺจเยน ปจฺจโย. ตีณิ.}

\prose{นเหตุ ธมฺโม นนเหตุสฺส ธมฺมสฺส อารมฺมณ\-ปจฺจเยน ปจฺจโย. ตีณิ.}

\prose{เหตุ จ นเหตุ จ ธมฺมา นเหตุสฺส ธมฺมสฺส อารมฺมณ\-ปจฺจเยน ปจฺจโย. ตีณิ. (สํขิตฺตํ.)}

\proseitem{3}{เหตุยา ตีณิ, อารมฺมเณ นว ฯเปฯ อวิคเต นว. (ปญฺหาวารมฺปิ เอวํ วิตฺถาเรตพฺพํ.)}

\csromanmarkh{2. สเหตุกทุก}\csromanmarkhii{1. ปฏิจฺจวาร}\csromanheadernomark{2}{2. \textbf{สเหตุกทุก 1. ปฏิจฺจวาร}}
\proseitem{4}{สเหตุกํ ธมฺมํ ปฏิจฺจ นสเหตุโก ธมฺโม อุปฺปชฺชติ เหตุปจฺจยา.  สเหตุกํ ธมฺมํ ปฏิจฺจ นอเหตุโก ธมฺโม อุปฺปชฺชติ เหตุปจฺจยา.  สเหตุกํ ธมฺมํ ปฏิจฺจ นสเหตุโก จ นอเหตุโก จ ธมฺมา อุปฺปชฺชนฺติ เหตุปจฺจยา. (3)}

\prose{อเหตุกํ ธมฺมํ ปฏิจฺจ นอเหตุโก ธมฺโม อุปฺปชฺชติ เหตุปจฺจยา.  อเหตุกํ ธมฺมํ ปฏิจฺจ นสเหตุโก ธมฺโม อุปฺปชฺชติ เหตุปจฺจยา.  อเหตุกํ ธมฺมํ ปฏิจฺจ นสเหตุโก จ นอเหตุโก จ ธมฺมา อุปฺปชฺชนฺติ เหตุปจฺจยา. (3)}

\prose{สเหตุกญฺจ อเหตุกญฺจ ธมฺมํ ปฏิจฺจ นสเหตุโก ธมฺโม อุปฺปชฺชติ เหตุปจฺจยา.  สเหตุกญฺจ อเหตุกญฺจ ธมฺมํ ปฏิจฺจ นอเหตุโก ธมฺโม อุปฺปชฺชติ เหตุปจฺจยา.  สเหตุกญฺจ อเหตุกญฺจ ธมฺมํ ปฏิจฺจ นสเหตุโก จ นอเหตุโก จ ธมฺมา อุปฺปชฺชนฺติ เหตุปจฺจยา. (3) (สํขิตฺตํ.)}

\prose{เหตุยา นว, อารมฺมเณ ฉ ฯเปฯ อวิคเต นว.}

\vspace{\tipitakaparskip}
\csromancenter{(สหชาตวารมฺปิ ฯเปฯ ปญฺหาวารมฺปิ วิตฺถาเรตพฺพํ.)}
\csromanmarkh{4. เหตุ\-สเหตุกทุก}\csromanmarkhii{3. เหตุ\-สมฺปยุตฺตทุก 1. ปฏิจฺจวาร}\csromanheadernomark{2}{4. เหตุ\-สเหตุกทุก 3. เหตุ\-สมฺปยุตฺตทุก 1. ปฏิจฺจวาร}
\proseitem{5}{\csromanfolio{199}เหตุ\-สมฺปยุตฺตํ ธมฺมํ ปฏิจฺจ นเหตุ\-สมฺปยุตฺโต ธมฺโม อุปฺปชฺชติ เหตุปจฺจยา.  เหตุ\-สมฺปยุตฺตํ ธมฺมํ ปฏิจฺจ นเหตุ\-วิปฺปยุตฺโต ธมฺโม อุปฺปชฺชติ เหตุปจฺจยา.  เหตุ\-สมฺปยุตฺตํ ธมฺมํ ปฏิจฺจ นเหตุ\-สมฺปยุตฺโต จ นเหตุ\-วิปฺปยุตฺโต จ ธมฺมา อุปฺปชฺชนฺติ เหตุปจฺจยา. (3)}

\prose{เหตุ\-วิปฺปยุตฺตํ ธมฺมํ ปฏิจฺจ นเหตุ\-วิปฺปยุตฺโต ธมฺโม อุปฺปชฺชติ เหตุปจฺจยา.  เหตุ\-วิปฺปยุตฺตํ ธมฺมํ ปฏิจฺจ นเหตุ\-สมฺปยุตฺโต ธมฺโม อุปฺปชฺชติ เหตุปจฺจยา.  เหตุ\-วิปฺปยุตฺตํ ธมฺมํ ปฏิจฺจ นเหตุ\-สมฺปยุตฺโต จ นเหตุ\-วิปฺปยุตฺโต จ ธมฺมา อุปฺปชฺชนฺติ เหตุปจฺจยา. (3)}

\prose{เหตุ\-สมฺปยุตฺตญฺจ เหตุ\-วิปฺปยุตฺตญฺจ ธมฺมํ ปฏิจฺจ นเหตุ\-สมฺปยุตฺโต ธมฺโม อุปฺปชฺชติ เหตุปจฺจยา.  เหตุ\-สมฺปยุตฺตญฺจ เหตุ\-วิปฺปยุตฺตญฺจ ธมฺมํ ปฏิจฺจ นเหตุ\-วิปฺปยุตฺโต ธมฺโม อุปฺปชฺชติ เหตุปจฺจยา.  เหตุ\-สมฺปยุตฺตญฺจ เหตุ\-วิปฺปยุตฺตญฺจ ธมฺมํ ปฏิจฺจ นเหตุ\-สมฺปยุตฺโต จ นเหตุ\-วิปฺปยุตฺโต จ ธมฺมา อุปฺปชฺชนฺติ เหตุปจฺจยา. (3) (สํขิตฺตํ.)}

\prose{เหตุยา นว.}

\vspace{\tipitakaparskip}
\csromancenter{(สหชาตวารมฺปิ ฯเปฯ ปญฺหาวารมฺปิ วิตฺถาเรตพฺพํ.)}
\csromanmarkh{4. เหตุ\-สเหตุกทุก}\csromanmarkhii{1. ปฏิจฺจวาร}\csromanheadernomark{2}{4. เหตุ\-สเหตุกทุก 1. ปฏิจฺจวาร}
\proseitem{6}{เหตุญฺเจว สเหตุกญฺจ ธมฺมํ ปฏิจฺจ นเหตุ เจว นอเหตุโก จ ธมฺโม อุปฺปชฺชติ เหตุปจฺจยา.  เหตุญฺเจว สเหตุกญฺจ ธมฺมํ ปฏิจฺจ นอเหตุโก เจว นนเหตุ จ ธมฺโม อุปฺปชฺชติ เหตุปจฺจยา.  เหตุญฺเจว สเหตุกญฺจ ธมฺมํ ปฏิจฺจ นเหตุ เจว นอเหตุโก จ นอเหตุโก เจว นนเหตุ จ ธมฺมา อุปฺปชฺชนฺติ เหตุปจฺจยา. (3)}

\prose{สเหตุกญฺเจว น จ เหตุํ ธมฺมํ ปฏิจฺจ นอเหตุโก เจว นนเหตุ จ ธมฺโม อุปฺปชฺชติ เหตุปจฺจยา.  สเหตุกญฺเจว น จ เหตุํ ธมฺมํ ปฏิจฺจ นเหตุ เจว นอเหตุโก จ ธมฺโม อุปฺปชฺชติ เหตุปจฺจยา.  สเหตุกญฺเจว น จ เหตุํ ธมฺมํ ปฏิจฺจ นเหตุ เจว นอเหตุโก จ นอเหตุโก เจว นนเหตุ จ ธมฺมา อุปฺปชฺชนฺติ เหตุปจฺจยา. (3)}

\prose{เหตุญฺเจว สเหตุกญฺจ สเหตุกญฺเจว น จ เหตุํ ธมฺมํ ปฏิจฺจ นเหตุ เจว นอเหตุโก จ ธมฺโม อุปฺปชฺชติ เหตุปจฺจยา.  เหตุญฺเจว สเหตุกญฺจ สเหตุกญฺเจว น จ เหตุญฺจ ธมฺมํ ปฏิจฺจ นอเหตุโก \csromanfolio{200}เจว นนเหตุ จ ธมฺโม อุปฺปชฺชติ เหตุปจฺจยา.  เหตุญฺเจว สเหตุกญฺจ สเหตุกญฺเจว น จ เหตุํ ธมฺมํ ปฏิจฺจ นเหตุ เจว นอเหตุโก จ นอเหตุโก เจว นนเหตุ จ ธมฺมา อุปฺปชฺชนฺติ เหตุปจฺจยา. (3)}

\prose{เหตุยา นว.}

\vspace{\tipitakaparskip}
\csromancenter{(สหชาตวารมฺปิ ฯเปฯ ปญฺหาวารมฺปิ วิตฺถาเรตพฺพํ.)}
\csromanmarkh{5. เหตุ\-เหตุ\-สมฺปยุตฺตทุก}\csromanmarkhii{1. ปฏิจฺจวาร}\csromanheadernomark{2}{5. เหตุ\-เหตุ\-สมฺปยุตฺตทุก 1. ปฏิจฺจวาร}
\proseitem{7}{เหตุญฺเจว เหตุ\-สมฺปยุตฺตญฺจ ธมฺมํ ปฏิจฺจ นเหตุ เจว นเหตุ\-วิปฺปยุตฺโต จ ธมฺโม อุปฺปชฺชติ เหตุปจฺจยา.  เหตุญฺเจว เหตุ\-สมฺปยุตฺตญฺจ ธมฺมํ ปฏิจฺจ นเหตุ\-วิปฺปยุตฺโต เจว นนเหตุ จ ธมฺโม อุปฺปชฺชติ เหตุปจฺจยา.  เหตุญฺเจว เหตุ\-สมฺปยุตฺตญฺจ ธมฺมํ ปฏิจฺจ นเหตุ เจว นเหตุ\-วิปฺปยุตฺโต จ นเหตุ\-วิปฺปยุตฺโต เจว นนเหตุ จ ธมฺมา อุปฺปชฺชนฺติ เหตุปจฺจยา. (3)}

\prose{เหตุ\-สมฺปยุตฺต\-ญฺเจว น จ เหตุํ ธมฺมํ ปฏิจฺจ นเหตุ\-วิปฺปยุตฺโต เจว นนเหตุ จ ธมฺโม อุปฺปชฺชติ เหตุปจฺจยา.  เหตุ\-สมฺปยุตฺต\-ญฺเจว น จ เหตุํ ธมฺมํ ปฏิจฺจ นเหตุ เจว นเหตุ\-วิปฺปยุตฺโต จ ธมฺโม อุปฺปชฺชติ เหตุปจฺจยา.  เหตุ\-สมฺปยุตฺต\-ญฺเจว น จ เหตุํ ธมฺมํ ปฏิจฺจ นเหตุ เจว นเหตุ\-วิปฺปยุตฺโต จ นเหตุ\-วิปฺปยุตฺโต เจว นนเหตุ จ ธมฺมา อุปฺปชฺชนฺติ เหตุปจฺจยา. (3)}

\prose{เหตุญฺเจว เหตุ\-สมฺปยุตฺตญฺจ เหตุ\-สมฺปยุตฺต\-ญฺเจว น จ เหตุญฺจ ธมฺมํ ปฏิจฺจ นเหตุ เจว นเหตุ\-วิปฺปยุตฺโต จ ธมฺโม อุปฺปชฺชติ เหตุปจฺจยา.  เหตุญฺเจว เหตุ\-สมฺปยุตฺตญฺจ เหตุ\-สมฺปยุตฺต\-ญฺเจว น จ เหตุญฺจ ธมฺมํ ปฏิจฺจ นเหตุ\-วิปฺปยุตฺโต เจว นนเหตุ ธมฺโม จ อุปฺปชฺชติ เหตุปจฺจยา.  เหตุญฺเจว เหตุ\-สมฺปยุตฺตญฺจ เหตุ\-สมฺปยุตฺต\-ญฺเจว น จ เหตุญฺจ ธมฺมํ ปฏิจฺจ นเหตุ เจว นเหตุ\-วิปฺปยุตฺโต จ นเหตุ\-วิปฺปยุตฺโต เจว นนเหตุ จ ธมฺมา อุปฺปชฺชนฺติ เหตุปจฺจยา. (3) (สํขิตฺตํ.)}

\prose{เหตุยา นว.}

\vspace{\tipitakaparskip}
\csromancenter{(สหชาตวารมฺปิ ฯเปฯ ปญฺหาวารมฺปิ วิตฺถาเรตพฺพํ.)}
\csromanmarkh{10. สปฺปฏิฆทุก}\csromanmarkhii{6. นเหตุ\-สเหตุกทุก 1. ปฏิจฺจวาร}\csromanheadernomark{2}{10. สปฺปฏิฆทุก 6. นเหตุ\-สเหตุกทุก 1. ปฏิจฺจวาร}
\proseitem{8}{\csromanfolio{201}นเหตุํ สเหตุกํ ธมฺมํ ปฏิจฺจ นเหตุ นสเหตุโก ธมฺโม อุปฺปชฺชติ เหตุปจฺจยา. (สํขิตฺตํ.)}

\prosewithcloser{เหตุยา นว.}{(สหชาตวารมฺปิ ฯเปฯ ปญฺหาวารมฺปิ วิตฺถาเรตพฺพํ.) เหตุโคจฺฉกํ นิฏฺฐิตํ.}
\csromanmarkh{7. สปฺปจฺจยทุก}\csromanmarkhii{1. ปฏิจฺจาทิวาร}\csromanheadernomark{2}{7. \textbf{สปฺปจฺจยทุก 1. ปฏิจฺจาทิวาร}}
\proseitem{9}{สปฺปจฺจยํ ธมฺมํ ปฏิจฺจ นอปฺปจฺจโย ธมฺโม อุปฺปชฺชติ เหตุปจฺจยา. (สํขิตฺตํ.) . เหตุยา เอกํ.}

\proseitem{10}{สปฺปจฺจโย ธมฺโม นอปฺปจฺจยสฺส ธมฺมสฺส อารมฺมณ\-ปจฺจเยน ปจฺจโย.}

\prose{อปฺปจฺจโย ธมฺโม นอปฺปจฺจยสฺส ธมฺมสฺส อารมฺมณ\-ปจฺจเยน ปจฺจโย. (สงฺขตํ สปฺปจฺจย\-สทิสํ.)}

\csromanmarkh{9. สนิทสฺสนทุก}\csromanmarkhii{1. ปฏิจฺจวาร}\csromanheadernomark{2}{9. \textbf{สนิทสฺสนทุก 1. ปฏิจฺจวาร}}
\proseitem{11}{อนิทสฺสนํ ธมฺมํ ปฏิจฺจ นอนิทสฺสโน ธมฺโม อุปฺปชฺชติ เหตุปจฺจยา.  อนิทสฺสนํ ธมฺมํ ปฏิจฺจ นสนิทสฺสโน ธมฺโม อุปฺปชฺชติ เหตุปจฺจยา.  อนิทสฺสนํ ธมฺมํ ปฏิจฺจ นสนิทสฺสโน จ นอนิทสฺสโน จ ธมฺมา อุปฺปชฺชนฺติ เหตุปจฺจยา. (3) (สํขิตฺตํ.)}

\prose{เหตุยา ตีณิ. (สพฺพตฺถ วิตฺถาโร.)}

\csromanmarkh{10. สปฺปฏิฆทุก}\csromanmarkhii{1. ปฏิจฺจวาร}\csromanheadernomark{2}{10. สปฺปฏิฆทุก 1. ปฏิจฺจวาร}
\proseitem{12}{สปฺปฏิฆํ ธมฺมํ ปฏิจฺจ นสปฺปฏิโฆ ธมฺโม อุปฺปชฺชติ เหตุปจฺจยา.  สปฺปฏิฆํ ธมฺมํ ปฏิจฺจ นอปฺปฏิโฆ ธมฺโม อุปฺปชฺชติ เหตุปจฺจยา.  สปฺปฏิฆํ ธมฺมํ ปฏิจฺจ นสปฺปฏิโฆ จ นอปฺปฏิโฆ จ ธมฺมา อุปฺปชฺชนฺติ เหตุปจฺจยา. (3)}

\prose{อปฺปฏิฆํ ธมฺมํ ปฏิจฺจ นอปฺปฏิโฆ ธมฺโม อุปฺปชฺชติ เหตุปจฺจยา.  อปฺปฏิฆํ ธมฺมํ ปฏิจฺจ นสปฺปฏิโฆ ธมฺโม อุปฺปชฺชติ เหตุปจฺจยา.  อปฺปฏิฆํ ธมฺมํ ปฏิจฺจ นสปฺปฏิโฆ จ นอปฺปฏิโฆ จ ธมฺมา อุปฺปชฺชนฺติ เหตุปจฺจยา. (3)}

\csromanpalirecto
\gambhira{ธมฺมา\-นุโลม\-ปจฺจนีย ทุกปฏฺฐาน\-ปาฬิ}
\prose{\csromanfolio{202}สปฺปฏิฆญฺจ อปฺปฏิฆญฺจ ธมฺมํ ปฏิจฺจ นสปฺปฏิโฆ ธมฺโม อุปฺปชฺชติ เหตุปจฺจยา.  สปฺปฏิฆญฺจ อปฺปฏิฆญฺจ ธมฺมํ ปฏิจฺจ นอปฺปฏิโฆ ธมฺโม อุปฺปชฺชติ เหตุปจฺจยา.  สปฺปฏิฆญฺจ อปฺปฏิฆญฺจ ธมฺมํ ปฏิจฺจ นสปฺปฏิโฆ จ นอปฺปฏิโฆ จ ธมฺมา อุปฺปชฺชนฺติ เหตุปจฺจยา. (3) (\csromansymbolfoottext{( )}{* อยํ สงฺขฺยา วิจาเรตพฺพา, นโลกิยนโลกุตฺตรธมฺโม นาม นตฺถิ.}สํขิตฺตํ.)}

\prose{เหตุยา นว. (สพฺพตฺถ วิตฺถาโร.)}

\csromanmarkh{11. รูปีทุก}\csromanmarkhii{1. ปฏิจฺจวาร}\csromanheadernomark{2}{11. \textbf{รูปีทุก 1. ปฏิจฺจวาร}}
\proseitem{13}{รูปิํ ธมฺมํ ปฏิจฺจ นรูปี ธมฺโม อุปฺปชฺชติ เหตุปจฺจยา.}

\prose{เหตุยา นว. (สพฺพตฺถ วิตฺถาโร.)}

\csromanmarkh{12. โลกิยทุก}\csromanmarkhii{1. ปฏิจฺจวาร}\csromanheadernomark{2}{12. \textbf{โลกิยทุก 1. ปฏิจฺจวาร}}
\prose{14. โลกิยํ ธมฺมํ ปฏิจฺจ นโลกุตฺตโร ธมฺโม อุปฺปชฺชติ}

\prose{โลกุตฺตรํ ธมฺมํ ปฏิจฺจ นโลกิโย ธมฺโม อุปฺปชฺชติ เหตุปจฺจยา.  โลกุตฺตรํ ธมฺมํ ปฏิจฺจ นโลกุตฺตโร ธมฺโม อุปฺปชฺชติ เหตุปจฺจยา.  (โลกุตฺตรํ ธมฺมํ ปฏิจฺจ นโลกิโย จ นโลกุตฺตโร จ ธมฺมา อุปฺปชฺชนฺติ เหตุปจฺจยา.)* (3)}

\prose{โลกิยญฺจ โลกุตฺตรญฺจ ธมฺมํ ปฏิจฺจ นโลกุตฺตโร ธมฺโม อุปฺปชฺชติ เหตุปจฺจยา. (1) (สํขิตฺตํ.)}

\prose{เหตุยา ปญฺจ. (สพฺพตฺถ ปญฺจ.)}

\csromanmarkh{13. เกนจิ\-วิญฺเญยฺยทุก}\csromanmarkhii{1. ปฏิจฺจวาร}\csromanheadernomark{2}{13. เกนจิ\-วิญฺเญยฺยทุก 1. ปฏิจฺจวาร}
\proseitem{15}{เกนจิ วิญฺเญยฺยํ ธมฺมํ ปฏิจฺจ นเกนจิ วิญฺเญยฺโย ธมฺโม อุปฺปชฺชติ เหตุปจฺจยา.  เกนจิ วิญฺเญยฺยํ ธมฺมํ ปฏิจฺจ นนเกนจิ วิญฺเญยฺโย ธมฺโม อุปฺปชฺชติ เหตุปจฺจยา.  เกนจิ วิญฺเญยฺยํ ธมฺมํ ปฏิจฺจ นเกนจิ วิญฺเญยฺโย จ นนเกนจิ วิญฺเญยฺโย จ ธมฺมา อุปฺปชฺชนฺติ เหตุปจฺจยา. (3)}

\prose{นเกนจิ วิญฺเญยฺยํ ธมฺมํ ปฏิจฺจ นนเกนจิ วิญฺเญยฺโย ธมฺโม อุปฺปชฺชติ เหตุปจฺจยา. ตีณิ.}

\csromanheader{2}{15. สาสวทุก}
\prosewithcloser{\csromanfolio{203}เกนจิ วิญฺเญยฺยญฺจ นเกนจิ วิญฺเญยฺยญฺจ ธมฺมํ ปฏิจฺจ นเกนจิ วิญฺเญยฺโย ธมฺโม อุปฺปชฺชติ เหตุปจฺจยา. ตีณิ. (สํขิตฺตํ.) เหตุยา นว. (สพฺพตฺถ นว.)}{จูฬนฺตรทุกํ นิฏฺฐิตํ.}
\csromanmarkh{14. อาสวทุก}\csromanmarkhii{1. ปฏิจฺจวาร}\csromanheadernomark{2}{14. \textbf{อาสวทุก 1. ปฏิจฺจวาร}}
\proseitem{16}{อาสวํ ธมฺมํ ปฏิจฺจ โนอาสโว ธมฺโม อุปฺปชฺชติ เหตุปจฺจยา.  อาสวํ ธมฺมํ ปฏิจฺจ นโนอาสโว ธมฺโม อุปฺปชฺชติ เหตุปจฺจยา.  อาสวํ ธมฺมํ ปฏิจฺจ โนอาสโว จ นโนอาสโว จ ธมฺมา อุปฺปชฺชนฺติ เหตุปจฺจยา. (3)}

\prose{โนอาสวํ ธมฺมํ ปฏิจฺจ นโนอาสโว ธมฺโม อุปฺปชฺชติ เหตุปจฺจยา.  โนอาสวํ ธมฺมํ ปฏิจฺจ โนอาสโว ธมฺโม อุปฺปชฺชติ เหตุปจฺจยา.  โนอาสวํ ธมฺมํ ปฏิจฺจ โนอาสโว จ นโนอาสโว จ ธมฺมา อุปฺปชฺชนฺติ เหตุปจฺจยา. (3)}

\prose{อาสวญฺจ โนอาสวญฺจ ธมฺมํ ปฏิจฺจ โนอาสโว ธมฺโม อุปฺปชฺชติ เหตุปจฺจยา.  อาสวญฺจ โนอาสวญฺจ ธมฺมํ ปฏิจฺจ นโนอาสโว ธมฺโม อุปฺปชฺชติ เหตุปจฺจยา.  อาสวญฺจ โนอาสวญฺจ ธมฺมํ ปฏิจฺจ โนอาสโว จ นโนอาสโว จ ธมฺมา อุปฺปชฺชนฺติ เหตุปจฺจยา. (3) (สํขิตฺตํ.) เหตุยา นว. (สพฺพตฺถ นว.)}

\csromanmarkh{15. สาสวทุก}\csromanmarkhii{1. ปฏิจฺจวาร}\csromanheadernomark{2}{15. สาสวทุก 1. ปฏิจฺจวาร}
\prose{17. สาสวํ ธมฺมํ ปฏิจฺจ นอนาสโว ธมฺโม อุปฺปชฺชติ}

\prose{อนาสวํ ธมฺมํ ปฏิจฺจ นอนาสโว ธมฺโม อุปฺปชฺชติ เหตุปจฺจยา.  อนาสวํ ธมฺมํ ปฏิจฺจ นสาสโว ธมฺโม อุปฺปชฺชติ เหตุปจฺจยา.  อนาสวํ ธมฺมํ ปฏิจฺจ นสาสโว จ นอนาสโว จ ธมฺมา อุปฺปชฺชนฺติ เหตุปจฺจยา. (3)}

\prose{สาสวญฺจ อนาสวญฺจ ธมฺมํ ปฏิจฺจ นอนาสโว ธมฺโม อุปฺปชฺชติ เหตุปจฺจยา. (1) (สํขิตฺตํ.)}

\prose{เหตุยา ปญฺจ. (สพฺพตฺถ วิตฺตาโร.)}

\csromanmarkh{ธมฺมา\-นุโลม\-ปจฺจนีย ทุกปฏฺฐาน\-ปาฬิ}\csromanmarkhii{16. อาสว\-สมฺปยุตฺตทุก 1. ปฏิจฺจวาร}\csromanheadernomark{2}{ธมฺมา\-นุโลม\-ปจฺจนีย ทุกปฏฺฐาน\-ปาฬิ 16. อาสว\-สมฺปยุตฺตทุก 1. ปฏิจฺจวาร}
\proseitem{18}{\csromanfolio{204}อาสว\-สมฺปยุตฺตํ ธมฺมํ ปฏิจฺจ นอาสว\-สมฺปยุตฺโต ธมฺโม อุปฺปชฺชติ เหตุปจฺจยา.  อาสว\-สมฺปยุตฺตํ ธมฺมํ ปฏิจฺจ นอาสว\-วิปฺปยุตฺโต ธมฺโม อุปฺปชฺชติ เหตุปจฺจยา.  อาสว\-สมฺปยุตฺตํ ธมฺมํ ปฏิจฺจ นอาสว\-สมฺปยุตฺโต จ นอาสว\-วิปฺปยุตฺโต จ ธมฺมา อุปฺปชฺชนฺติ เหตุปจฺจยา. (3)}

\prose{อาสว\-วิปฺปยุตฺตํ ธมฺมํ ปฏิจฺจ นอาสว\-วิปฺปยุตฺโต ธมฺโม อุปฺปชฺชติ เหตุปจฺจยา. อาสว\-วิปฺปยุตฺตํ ธมฺมํ ปฏิจฺจ นอาสว\-สมฺปยุตฺโต ธมฺโม อุปฺปชฺชติ เหตุปจฺจยา.  อาสว\-วิปฺปยุตฺตํ ธมฺมํ ปฏิจฺจ นอาสว\-สมฺปยุตฺโต จ นอาสว\-วิปฺปยุตฺโต จ ธมฺมา อุปฺปชฺชนฺติ เหตุปจฺจยา. (3)}

\prose{อาสว\-สมฺปยุตฺตญฺจ อาสว\-วิปฺปยุตฺตญฺจ ธมฺมํ ปฏิจฺจ นอาสว\-สมฺปยุตฺโต ธมฺโม อุปฺปชฺชติ เหตุปจฺจยา.  อาสว\-สมฺปยุตฺตญฺจ อาสว\-วิปฺปยุตฺตญฺจ ธมฺมํ ปฏิจฺจ นอาสว\-วิปฺปยุตฺโต ธมฺโม อุปฺปชฺชติ เหตุปจฺจยา.  อาสว\-สมฺปยุตฺตญฺจ อาสว\-วิปฺปยุตฺตญฺจ ธมฺมํ ปฏิจฺจ นอาสว\-สมฺปยุตฺโต จ นอาสว\-วิปฺปยุตฺโต จ ธมฺมา อุปฺปชฺชนฺติ เหตุปจฺจยา. (3) (สํขิตฺตํ.)}

\prose{เหตุยา นว. (สพฺพตฺถ นว.)}

\csromanmarkh{17. อาสว\-สาสวทุก}\csromanmarkhii{1. ปฏิจฺจวาร}\csromanheadernomark{2}{17. อาสว\-สาสวทุก 1. ปฏิจฺจวาร}
\proseitem{19}{อาสวญฺเจว สาสวญฺจ ธมฺมํ ปฏิจฺจ นอาสโว เจว นอนาสโว จ ธมฺโม อุปฺปชฺชติ เหตุปจฺจยา.  อาสวญฺเจว สาสวญฺจ ธมฺมํ ปฏิจฺจ นอนาสโว เจว นโน จ อาสโว ธมฺโม อุปฺปชฺชติ เหตุปจฺจยา.  อาสวญฺเจว สาสวญฺจ ธมฺมํ ปฏิจฺจ นอาสโว เจว นอนาสโว จ นอนาสโว เจว นโน จ อาสโว จ ธมฺมา อุปฺปชฺชนฺติ เหตุปจฺจยา. (3)}

\prose{สาสวญฺเจว โน จ อาสวํ ธมฺมํ ปฏิจฺจ นอนาสโว เจว นโน จ อาสโว ธมฺโม อุปฺปชฺชติ เหตุปจฺจยา. ตีณิ.}

\prose{อาสวญฺเจว สาสวญฺจ สาสวญฺเจว โน จ อาสวญฺจ ธมฺมํ ปฏิจฺจ นอาสโว เจว นอนาสโว จ ธมฺโม อุปฺปชฺชติ เหตุปจฺจยา. ตีณิ. เหตุยา นว. (สพฺพตฺถ นว.)}

\csromanmarkh{50. ปรามาสทุก}\csromanmarkhii{18. อาสว\-อาสว\-สมฺปยุตฺตทุก 1. ปฏิจฺจวาร}\csromanheadernomark{2}{50. ปรามาสทุก 18. อาสว\-อาสว\-สมฺปยุตฺตทุก 1. ปฏิจฺจวาร}
\proseitem{20}{\csromanfolio{205}อาสวญฺเจว อาสว\-สมฺปยุตฺตญฺจ ธมฺมํ ปฏิจฺจ นอาสโว เจว นอาสว\-วิปฺปยุตฺโต จ ธมฺโม อุปฺปชฺชติ เหตุปจฺจยา. ตีณิ.}

\prose{อาสว\-สมฺปยุตฺต\-ญฺเจว โน จ อาสวํ ธมฺมํ ปฏิจฺจ นอาสว\-วิปฺปยุตฺโต เจว นโน จ อาสโว ธมฺโม อุปฺปชฺชติ เหตุปจฺจยา. ตีณิ.}

\prose{อาสวญฺเจว อาสว\-สมฺปยุตฺตญฺจ อาสว\-วิปฺปยุตฺต\-ญฺเจว โน จ อาสวญฺจ ธมฺมํ ปฏิจฺจ โนอาสโว เจว นอาสว\-วิปฺปยุตฺโต จ ธมฺโม อุปฺปชฺชติ เหตุปจฺจยา. ตีณิ. (สํขิตฺตํ.) เหตุยา นว. (สพฺพตฺถ นว.)}

\csromanmarkh{19. อาสว\-วิปฺปยุตฺต\-สาสวทุก}\csromanmarkhii{1. ปฏิจฺจวาร}\csromanheadernomark{2}{19. อาสว\-วิปฺปยุตฺต\-สาสวทุก 1. ปฏิจฺจวาร}
\proseitem{21}{อาสว\-วิปฺปยุตฺตํ สาสวํ ธมฺมํ ปฏิจฺจ อาสว\-วิปฺปยุตฺโต นอนาสโว ธมฺโม อุปฺปชฺชติ เหตุปจฺจยา. (1)}

\prose{อาสว\-วิปฺปยุตฺตํ อนาสวํ ธมฺมํ ปฏิจฺจ อาสว\-วิปฺปยุตฺโต นอนาสโว ธมฺโม อุปฺปชฺชติ เหตุปจฺจยา. ตีณิ.}

\prose{อาสว\-วิปฺปยุตฺตํ สาสวญฺจ อาสว\-วิปฺปยุตฺตํ อนาสวญฺจ ธมฺมํ ปฏิจฺจ อาสว\-วิปฺปยุตฺโต นโนอนาสโว ธมฺโม อุปฺปชฺชติ เหตุปจฺจยา. (1)}

\prosewithcloser{เหตุยา ปญฺจ ฯเปฯ อวิคเต ปญฺจ. (สพฺพตฺถ วิตฺถาโร.)}{อาสวโคจฺฉกํ นิฏฺฐิตํ.}
\csromanmarkh{20. สญฺโญชน\-ทุกาทิ}\csromanmarkhii{1. ปฏิจฺจวาร}\csromanheadernomark{2}{20. สญฺโญชน\-ทุกาทิ 1. ปฏิจฺจวาร}
\proseitem{22}{สญฺโญชนํ ธมฺมํ ปฏิจฺจ โนสญฺโญชโน ธมฺโม อุปฺปชฺชติ เหตุปจฺจยา ฯเปฯ. คนฺถํ ธมฺมํ ปฏิจฺจ โนคนฺโถ ธมฺโม อุปฺปชฺชติ เหตุปจฺจยา ฯเปฯ. โอฆํ ธมฺมํ ปฏิจฺจ โนโอโฆ ธมฺโม อุปฺปชฺชติ เหตุปจฺจยา ฯเปฯ. โยคํ ธมฺมํ ปฏิจฺจ โนโยโค ธมฺโม อุปฺปชฺชติ เหตุปจฺจยา ฯเปฯ. นีวรณํ ธมฺมํ ปฏิจฺจ โนนีวรโณ ธมฺโม อุปฺปชฺชติ เหตุปจฺจยา.}

\csromanmarkh{50. ปรามาสทุก}\csromanmarkhii{1. ปฏิจฺจวาร}\csromanheadernomark{2}{50. ปรามาสทุก 1. ปฏิจฺจวาร}
\proseitem{23}{ปรามาสํ ธมฺมํ ปฏิจฺจ โนปรามาโส ธมฺโม อุปฺปชฺชติ เหตุปจฺจยา. (อาสว\-โคจฺฉก\-สทิสํ.)}

\csromanmarkh{ธมฺมา\-นุโลม\-ปจฺจนีย ทุกปฏฺฐาน\-ปาฬิ}\csromanmarkhii{55. สารมฺมณทุก 1. ปฏิจฺจวาร}\csromanheadernomark{2}{ธมฺมา\-นุโลม\-ปจฺจนีย ทุกปฏฺฐาน\-ปาฬิ 55. สารมฺมณทุก 1. ปฏิจฺจวาร}
\proseitem{24}{\csromanfolio{206}สารมฺมณํ ธมฺมํ ปฏิจฺจ นสารมฺมโณ ธมฺโม อุปฺปชฺชติ เหตุปจฺจยา. ตีณิ.}

\prose{อนารมฺมณํ ธมฺมํ ปฏิจฺจ นอนารมฺมโณ ธมฺโม อุปฺปชฺชติ เหตุปจฺจยา. ตีณิ.}

\prose{สารมฺมณญฺจ อนารมฺมณญฺจ ธมฺมํ ปฏิจฺจ นสารมฺมโณ ธมฺโม อุปฺปชฺชติ เหตุปจฺจยา. ตีณิ. (สํขิตฺตํ.)}

\prose{เหตุยา นว ฯเปฯ อวิคเต นว. (สพฺพตฺถ วิตฺถาโร.)}

\csromanmarkh{56. จิตฺตทุก}\csromanmarkhii{1. ปฏิจฺจวาร}\csromanheadernomark{2}{56. \textbf{จิตฺตทุก 1. ปฏิจฺจวาร}}
\proseitem{25}{จิตฺตํ ธมฺมํ ปฏิจฺจ โนจิตฺโต ธมฺโม อุปฺปชฺชติ เหตุปจฺจยา. เอกํ.}

\prose{โนจิตฺตํ ธมฺมํ ปฏิจฺจ นโนจิตฺโต ธมฺโม อุปฺปชฺชติ เหตุปจฺจยา. ตีณิ.}

\prose{จิตฺตญฺจ โนจิตฺตญฺจ ธมฺมํ ปฏิจฺจ โนจิตฺโต ธมฺโม อุปฺปชฺชติ เหตุปจฺจยา. (1) (สํขิตฺตํ. ปญฺจ.)}

\csromanmarkh{57. เจตสิกทุกาทิ}\csromanmarkhii{1. ปฏิจฺจวาร}\csromanheadernomark{2}{57. \textbf{เจตสิกทุกาทิ 1. ปฏิจฺจวาร}}
\proseitem{26}{เจตสิกํ ธมฺมํ ปฏิจฺจ นเจตสิโก ธมฺโม อุปฺปชฺชติ เหตุปจฺจยา.}

\proseitem{27}{จิตฺต\-สมฺปยุตฺตํ ธมฺมํ ปฏิจฺจ นจิตฺต\-สมฺปยุตฺโต ธมฺโม อุปฺปชฺชติ เหตุปจฺจยา ฯเปฯ.}

\prose{จิตฺต\-สํสฏฺฐํ ธมฺมํ ปฏิจฺจ นจิตฺต\-สํสฏฺโฐ ธมฺโม อุปฺปชฺชติ เหตุปจฺจยา.}

\proseitem{28}{จิตฺต\-สมุฏฺฐานํ ธมฺมํ ปฏิจฺจ นจิตฺต\-สมุฏฺฐาโน ธมฺโม อุปฺปชฺชติ เหตุปจฺจยา.}

\proseitem{29}{จิตฺตสหภุํ ธมฺมํ ปฏิจฺจ นจิตฺตสหภู ธมฺโม อุปฺปชฺชติ เหตุปจฺจยา ฯเปฯ. จิตฺตา\-นุปริวตฺติํ ธมฺมํ ปฏิจฺจ นจิตฺตา\-นุปริวตฺตี ธมฺโม อุปฺปชฺชติ เหตุปจฺจยา ฯเปฯ. จิตฺต\-สํสฏฺฐ\-สมุฏฺฐานํ ธมฺมํ ปฏิจฺจ นจิตฺต\-สํสฏฺฐ\-สมุฏฺฐาโน ธมฺโม อุปฺปชฺชติ เหตุปจฺจยา ฯเปฯ. จิตฺต\-สํสฏฺฐ\-สมุฏฺฐาน\-สหภุํ ธมฺมํ ปฏิจฺจ นจิตฺต\-สํสฏฺฐ\-สมุฏฺฐาน\-สหภู ธมฺโม อุปฺปชฺชติ เหตุปจฺจยา ฯเปฯ. จิตฺต\-สํสฏฺฐ\-สมุฏฺฐานา\-นุปริวตฺติํ ธมฺมํ ปฏิจฺจ นจิตฺต\-สํสฏฺฐ\-สมุฏฺฐานา\-นุ\-ปริวตฺตี ธมฺโม อุปฺปชฺชติ เหตุปจฺจยา.}

\csromanheader{2}{83. ทสฺสเนน\-ปหาตพฺพทุก}
\proseitem{30}{\csromanfolio{207}อชฺฌตฺติกํ ธมฺมํ ปฏิจฺจ นอชฺฌตฺติโก ธมฺโม อุปฺปชฺชติ เหตุปจฺจยา. ตีณิ.}

\prose{พาหิรํ ธมฺมํ ปฏิจฺจ นพาหิโร ธมฺโม อุปฺปชฺชติ เหตุปจฺจยา. ตีณิ.}

\prose{อชฺฌตฺติกญฺจ พาหิรญฺจ ธมฺมํ ปฏิจฺจ นอชฺฌตฺติโก ธมฺโม อุปฺปชฺชติ เหตุปจฺจยา. ตีณิ.}

\proseitem{31}{อุปาทา ธมฺมํ ปฏิจฺจ โนอุปาทา ธมฺโม อุปฺปชฺชติ เหตุปจฺจยา. (1)}

\prose{โนอุปาทา ธมฺมํ ปฏิจฺจ นโนอุปาทา ธมฺโม อุปฺปชฺชติ เหตุปจฺจยา. ตีณิ.}

\prose{อุปาทา จ โนอุปาทา จ ธมฺมํ ปฏิจฺจ โนอุปาทา ธมฺโม อุปฺปชฺชติ}

\proseitem{32}{อุปาทินฺนํ ธมฺมํ ปฏิจฺจ นอุปาทินฺโน ธมฺโม อุปฺปชฺชติ เหตุปจฺจยา. ตีณิ.}

\prose{อนุปาทินฺนํ ธมฺมํ ปฏิจฺจ นอนุปาทินฺโน ธมฺโม อุปฺปชฺชติ}

\prose{อุปาทินฺนญฺจ อนุปาทินฺนญฺจ ธมฺมํ ปฏิจฺจ นอุปาทินฺโน ธมฺโม อุปฺปชฺชติ เหตุปจฺจยา. (1) (สํขิตฺตํ.)}

\prosewithcloser{เหตุยา ปญฺจ ฯเปฯ อวิคเต ปญฺจ. (สพฺพตฺถ วิตฺถาโร.)}{มหนฺตรทุกํ นิฏฺฐิตํ.}
\csromanheader{2}{69. อุปาทาน\-โคจฺฉก}
\vspace{\tipitakaparskip}
\csromancenter{อุปาทานํ ธมฺมํ ปฏิจฺจ โนอุปาทาโน ธมฺโม อุปฺปชฺชติ เหตุปจฺจยา. เหตุยา นว.}
\csromanheader{2}{75. \textbf{กิเลสโคจฺฉก}}
\vspace{\tipitakaparskip}
\csromancenter{กิเลสํ ธมฺมํ ปฏิจฺจ โนกิเลโส ธมฺโม อุปฺปชฺชติ เหตุปจฺจยา. เหตุยา นว.}
\csromanheader{2}{83. ทสฺสเนน\-ปหาตพฺพทุก}
\proseitem{35}{ทสฺสเนน ปหาตพฺพํ ธมฺมํ ปฏิจฺจ นทสฺสเนน ปหาตพฺโพ ธมฺโม อุปฺปชฺชติ เหตุปจฺจยา.  ทสฺสเนน ปหาตพฺพํ ธมฺมํ ปฏิจฺจ นนทสฺสเนน \csromanfolio{208}ปหาตพฺโพ ธมฺโม อุปฺปชฺชติ เหตุปจฺจยา.  ทสฺสเนน ปหาตพฺพํ ธมฺมํ ปฏิจฺจ นทสฺสเนน ปหาตพฺโพ จ นนทสฺสเนน ปหาตพฺโพ จ ธมฺมา อุปฺปชฺชนฺติ เหตุปจฺจยา. (3)}

\prose{นทสฺสเนน ปหาตพฺพํ ธมฺมํ ปฏิจฺจ นทสฺสเนน ปหาตพฺโพ}

\prose{ทสฺสเนน ปหาตพฺพญฺจ นทสฺสเนน ปหาตพฺพญฺจ ธมฺมํ ปฏิจฺจ นทสฺสเนน ปหาตพฺโพ ธมฺโม อุปฺปชฺชติ เหตุปจฺจยา. (1) (สํขิตฺตํ.)}

\prose{เหตุยา ปญฺจ ฯเปฯ อวิคเต ปญฺจ.}

\csromanheader{2}{84. ภาวนาย\-ปหาตพฺพทุก}
\proseitem{36}{ภาวนาย ปหาตพฺพํ ธมฺมํ ปฏิจฺจ นภาวนาย ปหาตพฺโพ ธมฺโม อุปฺปชฺชติ เหตุปจฺจยา.  ภาวนาย ปหาตพฺพํ ธมฺมํ ปฏิจฺจ นนภาวนาย ปหาตพฺโพ ธมฺโม อุปฺปชฺชติ เหตุปจฺจยา.  ภาวนาย ปหาตพฺพํ ธมฺมํ ปฏิจฺจ นภาวนาย ปหาตพฺโพ จ นนภาวนาย ปหาตพฺโพ จ ธมฺมา อุปฺปชฺชนฺติ เหตุปจฺจยา. (3)}

\prose{นภาวนาย ปหาตพฺพํ ธมฺมํ ปฏิจฺจ นภาวนาย ปหาตพฺโพ}

\prose{ภาวนาย ปหาตพฺพญฺจ นภาวนาย ปหาตพฺพญฺจ ธมฺมํ ปฏิจฺจ นภาวนาย ปหาตพฺโพ ธมฺโม อุปฺปชฺชติ เหตุปจฺจยา. (1)}

\prose{เหตุยา ปญฺจ.}

\csromanheader{2}{85. ทสฺสเนน\-ปหาตพฺพ\-เหตุกทุก}
\proseitem{37}{ทสฺสเนน ปหาตพฺพ\-เหตุกํ ธมฺมํ ปฏิจฺจ นทสฺสเนน ปหาตพฺพ\-เหตุโก ธมฺโม อุปฺปชฺชติ เหตุปจฺจยา.  ทสฺสเนน ปหาตพฺพ\-เหตุกํ ธมฺมํ ปฏิจฺจ นนทสฺสเนน ปหาตพฺพ\-เหตุโก ธมฺโม อุปฺปชฺชติ เหตุปจฺจยา. (สํขิตฺตํ.)}

\prose{เหตุยา นว.}

\csromanmarkh{89-92. สปฺปีติกาทิทุก}\csromanmarkhii{86. ภาวนาย\-ปหาตพฺพ\-เหตุกทุก}\csromanheadernomark{2}{89-92. สปฺปีติกาทิทุก 86. ภาวนาย\-ปหาตพฺพ\-เหตุกทุก}
\proseitem{38}{\csromanfolio{209}ภาวนาย ปหาตพฺพ\-เหตุกํ ธมฺมํ ปฏิจฺจ นภาวนาย ปหาตพฺพ\-เหตุโก ธมฺโม อุปฺปชฺชติ เหตุปจฺจยา.  ภาวนาย ปหาตพฺพ\-เหตุกํ ธมฺมํ ปฏิจฺจ นนภาวนาย ปหาตพฺพ\-เหตุโก ธมฺโม อุปฺปชฺชติ เหตุปจฺจยา.  เหตุยา นว.}

\csromanheader{2}{87-88. สวิตกฺกาทิทุก}
\proseitem{39}{สวิตกฺกํ ธมฺมํ ปฏิจฺจ นสวิตกฺโก ธมฺโม อุปฺปชฺชติ เหตุปจฺจยา. ตีณิ.}

\prose{อวิตกฺกํ ธมฺมํ ปฏิจฺจ นอวิตกฺโก ธมฺโม อุปฺปชฺชติ เหตุปจฺจยา. (สํขิตฺตํ.) เหตุยา นว.}

\proseitem{40}{สวิจารํ ธมฺมํ ปฏิจฺจ นสวิจาโร ธมฺโม อุปฺปชฺชติ เหตุปจฺจยา. ตีณิ.}

\prose{อวิจารํ ธมฺมํ ปฏิจฺจ นอวิจาโร ธมฺโม อุปฺปชฺชติ เหตุปจฺจยา. (สํขิตฺตํ.) เหตุยา นว.}

\csromanheader{2}{89-92. สปฺปีติกาทิทุก}
\proseitem{41}{สปฺปีติกํ ธมฺมํ ปฏิจฺจ นสปฺปีติโก ธมฺโม อุปฺปชฺชติ เหตุปจฺจยา. ตีณิ.}

\prose{อปฺปีติกํ ธมฺมํ ปฏิจฺจ นอปฺปีติโก ธมฺโม อุปฺปชฺชติ เหตุปจฺจยา.  เหตุยา นว.}

\proseitem{42}{ปีติสหคตํ ธมฺมํ ปฏิจฺจ นปีติสหคโต ธมฺโม อุปฺปชฺชติ เหตุปจฺจยา. ตีณิ.}

\prose{นปีติสหคตํ ธมฺมํ ปฏิจฺจ นนปี\-ติ\-สหคโต ธมฺโม อุปฺปชฺชติ เหตุปจฺจยา.  เหตุยา นว.}

\proseitem{43}{\csromanfolio{210}สุขสหคตํ ธมฺมํ ปฏิจฺจ นสุขสหคโต ธมฺโม อุปฺปชฺชติ เหตุปจฺจยา. ตีณิ.}

\prose{นสุขสหคตํ ธมฺมํ ปฏิจฺจ นนสุข\-สหคโต ธมฺโม อุปฺปชฺชติ เหตุปจฺจยา.  เหตุยา นว.}

\proseitem{44}{อุเปกฺขา\-สหคตํ ธมฺมํ ปฏิจฺจ นอุเปกฺขาสหคโต ธมฺโม อุปฺปชฺชติ เหตุปจฺจยา. ตีณิ.}

\prose{นอุเปกฺขาสหคตํ ธมฺมํ ปฏิจฺจ นนอุเปกฺขาสหคโต ธมฺโม อุปฺปชฺชติ เหตุปจฺจยา.  เหตุยา นว.}

\csromanheader{2}{93-95. กามาวจราทิทุก}
\proseitem{45}{กามาวจรํ ธมฺมํ ปฏิจฺจ นกามาวจโร ธมฺโม อุปฺปชฺชติ เหตุปจฺจยา. ตีณิ.}

\prose{นกามาวจรํ ธมฺมํ ปฏิจฺจ นนกามาวจโร ธมฺโม อุปฺปชฺชติ เหตุปจฺจยา. เหตุยา นว.}

\proseitem{46}{รูปาวจรํ ธมฺมํ ปฏิจฺจ นรูปาวจโร ธมฺโม อุปฺปชฺชติ เหตุปจฺจยา. ตีณิ.}

\prose{นรูปาวจรํ ธมฺมํ ปฏิจฺจ นนรูปาวจโร ธมฺโม อุปฺปชฺชติ เหตุปจฺจยา. เหตุยา นว.}

\proseitem{47}{อรูปาวจรํ ธมฺมํ ปฏิจฺจ นอรูปาวจโร ธมฺโม อุปฺปชฺชติ เหตุปจฺจยา. ตีณิ.}

\prose{นอรูปาวจรํ ธมฺมํ ปฏิจฺจ นนอรูปาวจโร ธมฺโม อุปฺปชฺชติ เหตุปจฺจยา. (1)}

\prose{อรูปาวจรญฺจ นอรูปาวจรญฺจ ธมฺมํ ปฏิจฺจ นอรูปาวจโร ธมฺโม อุปฺปชฺชติ เหตุปจฺจยา. (1). เหตุยา ปญฺจ.}

\prose{\csromanfolio{211}99. สฺเอาตฺตรทุก \textbf{96. ปริยาปนฺนทุก}}

\proseitem{48}{ปริยาปนฺนํ ธมฺมํ ปฏิจฺจ นอปริยาปนฺโน ธมฺโม อุปฺปชฺชติ เหตุปจฺจยา. (1)}

\prose{อปริยาปนฺนํ ธมฺมํ ปฏิจฺจ นอปริยาปนฺโน ธมฺโม อุปฺปชฺชติ เหตุปจฺจยา.  อปริยาปนฺนํ ธมฺมํ ปฏิจฺจ นปริยาปนฺโน ธมฺโม อุปฺปชฺชติ เหตุปจฺจยา.  อปริยาปนฺนํ ธมฺมํ ปฏิจฺจ นปริยาปนฺโน จ นอปริยาปนฺโน จ ธมฺมา อุปฺปชฺชนฺติ เหตุปจฺจยา. (3)}

\prose{ปริยาปนฺนญฺจ อปริยาปนฺนญฺจ ธมฺมํ ปฏิจฺจ นอปริยาปนฺโน ธมฺโม อุปฺปชฺชติ เหตุปจฺจยา. (1) (สํขิตฺตํ.) . เหตุยา ปญฺจ.}

\csromanheader{2}{97. \textbf{นิยฺยานิกทุก}}
\proseitem{49}{นิยฺยานิกํ ธมฺมํ ปฏิจฺจ นนิยฺยานิโก ธมฺโม อุปฺปชฺชติ เหตุปจฺจยา. ตีณิ.}

\prose{อนิยฺยานิกํ ธมฺมํ ปฏิจฺจ นนิยฺยานิโก ธมฺโม อุปฺปชฺชติ}

\prose{นิยฺยานิกญฺจ อนิยฺยานิกญฺจ ธมฺมํ ปฏิจฺจ นนิยฺยานิโก ธมฺโม อุปฺปชฺชติ เหตุปจฺจยา. (1). เหตุยา ปญฺจ.}

\csromanheader{2}{98. \textbf{นิยตทุก}}
\proseitem{50}{นิยตํ ธมฺมํ ปฏิจฺจ นนิยโต ธมฺโม อุปฺปชฺชติ เหตุปจฺจยา. ตีณิ.}

\prose{อนิยตํ ธมฺมํ ปฏิจฺจ นนิยโต ธมฺโม อุปฺปชฺชติ เหตุปจฺจยา. (1)}

\prose{นิยตญฺจ อนิยตญฺจ ธมฺมํ ปฏิจฺจ นนิยโต ธมฺโม อุปฺปชฺชติ เหตุปจฺจยา. (1). เหตุยา ปญฺจ.}

\prose{99. สอุตฺตรทุก}

\prose{51. สอุตฺตรํ ธมฺมํ ปฏิจฺจ นอนุตฺตโร ธมฺโม อุปฺปชฺชติ}

\prose{อนุตฺตรํ ธมฺมํ ปฏิจฺจ นอนุตฺตโร ธมฺโม อุปฺปชฺชติ เหตุปจฺจยา.  อนุตฺตรํ ธมฺมํ ปฏิจฺจ นสอุตฺตโร ธมฺโม อุปฺปชฺชติ เหตุปจฺจยา.  อนุตฺตรํ ธมฺมํ ปฏิจฺจ นสอุตฺตโร จ นอนุตฺตโร จ ธมฺมา อุปฺปชฺชนฺติ เหตุปจฺจยา. (3)}

\prose{\csromanfolio{212}สอุตฺตรญฺจ อนุตฺตรญฺจ ธมฺมํ ปฏิจฺจ นอนุตฺตโร ธมฺโม อุปฺปชฺชติ เหตุปจฺจยา. (1). เหตุยา ปญฺจ.}

\csromanmarkh{100. สรณทุก}\csromanmarkhii{1. ปฏิจฺจวาร}\csromanheadernomark{2}{100. \textbf{สรณทุก 1. ปฏิจฺจวาร}}
\proseitem{52}{สรณํ ธมฺมํ ปฏิจฺจ นสรโณ ธมฺโม อุปฺปชฺชติ เหตุปจฺจยา.  สรณํ ธมฺมํ ปฏิจฺจ นอรโณ ธมฺโม อุปฺปชฺชติ เหตุปจฺจยา.  สรณํ ธมฺมํ ปฏิจฺจ นสรโณ จ นอรโณ จ ธมฺมา อุปฺปชฺชนฺติ เหตุปจฺจยา. (3)}

\prose{อรณํ ธมฺมํ ปฏิจฺจ นอรโณ ธมฺโม อุปฺปชฺชติ เหตุปจฺจยา. (1)}

\prose{สรณญฺจ อรณญฺจ ธมฺมํ ปฏิจฺจ นสรโณ ธมฺโม อุปฺปชฺชติ เหตุปจฺจยา. (1). เหตุยา ปญฺจ, อารมฺมเณ เทฺว ฯเปฯ อวิคเต ปญฺจ.}

\csromanheader{2}{\textbf{ปจฺจนีย นเหตุ}}
\prosecont{สรณํ ธมฺมํ ปฏิจฺจ นอรโณ ธมฺโม อุปฺปชฺชติ นเหตุปจฺจยา. (1)}

\proseitem{53}{อรณํ ธมฺมํ ปฏิจฺจ นนอรโณ ธมฺโม อุปฺปชฺชติ นเหตุปจฺจยา. (สํขิตฺตํ.)}

\prose{นเหตุยา เทฺว, นอารมฺมเณ ตีณิ ฯเปฯ โนวิคเต ตีณิ.}

\prose{(สหชาตวารมฺปิ ปจฺจยวารมฺปิ นิสฺสยวารมฺปิ สํสฏฺฐ\-วารมฺปิ สมฺปยุตฺตวารมฺปิ ปฏิจฺจ\-วาร\-สทิสํ วิตฺถาเรตพฺพํ.)}

\csromanmarkh{100. สรณทุก}\csromanmarkhii{7. ปญฺหาวาร}\csromanheadernomark{2}{100. \textbf{สรณทุก 7. ปญฺหาวาร}}
\proseitem{54}{สรโณ ธมฺโม นสรณสฺส ธมฺมสฺส เหตุปจฺจเยน ปจฺจโย.  สรโณ ธมฺโม นอรณสฺส ธมฺมสฺส เหตุปจฺจเยน ปจฺจโย.  สรโณ ธมฺโม นสรณสฺส จ นอรณสฺส จ ธมฺมสฺส เหตุปจฺจเยน ปจฺจโย. (3)}

\prose{อรโณ ธมฺโม นสรณสฺส ธมฺมสฺส เหตุปจฺจเยน ปจฺจโย. (1)}

\proseitem{55}{สรโณ ธมฺโม นสรณสฺส ธมฺมสฺส อารมฺมณ\-ปจฺจเยน ปจฺจโย.  สรโณ ธมฺโม นอรณสฺส ธมฺมสฺส อารมฺมณ\-ปจฺจเยน ปจฺจโย. (สํขิตฺตํ.)}

\prose{เหตุยา จตฺตาริ, อารมฺมเณ จตฺตาริ, อธิปติยา ปญฺจ, อนนฺตเร จตฺตาริ ฯเปฯ อวิคเต สตฺต.}

\csromanheader{2}{ปจฺจนียุ\-ทฺธาร ปจฺจนียุ\-ทฺธาร}
\proseitem{56}{\csromanfolio{213}สรโณ ธมฺโม นสรณสฺส ธมฺมสฺส อารมฺมณ\-ปจฺจเยน ปจฺจโย. สหชาต\-ปจฺจเยน ปจฺจโย. อุปนิสฺสย\-ปจฺจเยน ปจฺจโย. ปจฺฉาชาต\-ปจฺจเยน ปจฺจโย. กมฺมปจฺจเยน ปจฺจโย.}

\prose{สรโณ ธมฺโม นอรณสฺส ธมฺมสฺส อารมฺมณ\-ปจฺจเยน ปจฺจโย. สหชาต\-ปจฺจเยน ปจฺจโย. อุปนิสฺสย\-ปจฺจเยน ปจฺจโย. (สํขิตฺตํ.)}

\proseitem{57}{นเหตุยา สตฺต, นอารมฺมเณ สตฺต, นอธิปติยา สตฺต, นอนนฺตเร สตฺต ฯเปฯ โนอวิคเต จตฺตาริ.}

\prose{เหตุปจฺจยา นอารมฺมเณ จตฺตาริ. (สํขิตฺตํ.)}

\prose{นเหตุปจฺจยา อารมฺมเณ จตฺตาริ. (สํขิตฺตํ.)}

\prosewithcloser{(ยถา กุสลตฺติเก ปญฺหาวารํ, เอวํ วิตฺถาเรตพฺพํ.)}{ธมฺมา\-นุโลม\-ปจฺจนีเย ทุกปฏฺฐานํ นิฏฺฐิตํ.}
\pitaka{อภิธมฺม\-ปิฏก}
\namakkaram{นโม ตสฺส ภควโต อรหโต สมฺมา\-สมฺพุทฺธสฺส.}
\csromanmarkh{1. เหตุทุก}\csromanmarkhii{1. กุสลตฺติก}\csromanheadernomark{2}{1. \textbf{เหตุทุก 1. กุสลตฺติก}}
\csromanheader{2}{1. \textbf{กุสลปท 1-7. ปฏิจฺจาทิวาร}}
\proseitem{1}{\csromanfolio{215}เหตุํ กุสลํ ธมฺมํ ปฏิจฺจ นเหตุ นกุสโล ธมฺโม อุปฺปชฺชติ เหตุปจฺจยา.  นเหตุํ กุสลํ ธมฺมํ ปฏิจฺจ นนเหตุ นกุสโล ธมฺโม อุปฺปชฺชติ เหตุปจฺจยา.  เหตุํ กุสลญฺจ นเหตุํ กุสลญฺจ ธมฺมํ ปฏิจฺจ นเหตุ นกุสโล ธมฺโม อุปฺปชฺชติ เหตุปจฺจยา. ตีณิ. (จิตฺตสมุฏฺฐานเมว, อารมฺมณํ นตฺถิ.)}

\prose{เหตุยา ตีณิ, อธิปติยา ตีณิ ฯเปฯ อวิคเต ตีณิ.}

\vspace{\tipitakaparskip}
\csromancenter{(สหชาตวารมฺปิ ฯเปฯ นิสฺสยวารมฺปิ ปฏิจฺจ\-วาร\-สทิสํ.)}
\prose{2. เหตุ กุสโล ธมฺโม นเหตุสฺส นกุสลสฺส ธมฺมสฺส}

\prose{เหตุ กุสโล ธมฺโม นเหตุสฺส นกุสลสฺส ธมฺมสฺส อารมฺมณ\-ปจฺจเยน ปจฺจโย. ตีณิ.}

\prose{นเหตุ กุสโล ธมฺโม นนเหตุสฺส นกุสลสฺส ธมฺมสฺส อารมฺมณ\-ปจฺจเยน ปจฺจโย. ตีณิ.}

\prose{เหตุ กุสโล จ นเหตุ กุสโล จ ธมฺมา นเหตุสฺส นกุสลสฺส ธมฺมสฺส อารมฺมณ\-ปจฺจเยน ปจฺจโย. ตีณิ. (สํขิตฺตํ.)}

\csromanheader{2}{ธมฺมา\-นุโลม\-ปจฺจนีย ทุกติก\-ปฏฺฐาน\-ปาฬิ}
\proseitem{3}{\csromanfolio{216}เหตุยา เอกํ, อารมฺมเณ นว, อธิปติยา นว ฯเปฯ อวิคเต ตีณิ. (ปญฺหาวารมฺปิ วิตฺถาเรตพฺพํ.)}

\csromanheader{2}{2. \textbf{อกุสลปท 1-7. ปฏิจฺจาทิวาร}}
\proseitem{4}{เหตุํ อกุสลํ ธมฺมํ ปฏิจฺจ นเหตุ นอกุสโล ธมฺโม อุปฺปชฺชติ เหตุปจฺจยา. (1)}

\prose{นเหตุํ อกุสลํ ธมฺมํ ปฏิจฺจ นเหตุ นอกุสโล ธมฺโม อุปฺปชฺชติ เหตุปจฺจยา. (1)}

\prose{เหตุํ อกุสลญฺจ นเหตุํ อกุสลญฺจ ธมฺมํ ปฏิจฺจ นเหตุ นอกุสโล ธมฺโม อุปฺปชฺชติ เหตุปจฺจยา. (1) (สํขิตฺตํ.) . เหตุยา ตีณิ, อธิปติยา ตีณิ ฯเปฯ อวิคเต ตีณิ.}

\vspace{\tipitakaparskip}
\csromancenter{(สหชาตวารมฺปิ ฯเปฯ นิสฺสยวารมฺปิ ปฏิจฺจ\-วาร\-สทิสํ.)}
\prose{5. เหตุ อกุสโล ธมฺโม นเหตุสฺส นอกุสลสฺส ธมฺมสฺส}

\prose{เหตุ อกุสโล ธมฺโม นเหตุสฺส นอกุสลสฺส ธมฺมสฺส อารมฺมณ\-ปจฺจเยน ปจฺจโย. ตีณิ.}

\prose{นเหตุ อกุสโล ธมฺโม นนเหตุสฺส นอกุสลสฺส ธมฺมสฺส อารมฺมณ\-ปจฺจเยน ปจฺจโย. ตีณิ.}

\prose{เหตุ อกุสโล จ นเหตุ อกุสโล จ ธมฺมา นเหตุสฺส นอกุสลสฺส ธมฺมสฺส อารมฺมณ\-ปจฺจเยน ปจฺจโย. ตีณิ. (สํขิตฺตํ.)}

\proseitem{6}{เหตุยา เอกํ, อารมฺมเณ นว, อธิปติยา เอกํ ฯเปฯ อวิคเต ตีณิ.}

\csromanmarkh{3. อพฺยากตปท}\csromanmarkhii{3. ปจฺจยวาร}\csromanheadernomark{2}{3. \textbf{อพฺยากตปท 3. ปจฺจยวาร}}
\proseitem{7}{นเหตุํ อพฺยากตํ ธมฺมํ ปจฺจยา นนเหตุ นอพฺยากโต ธมฺโม อุปฺปชฺชติ เหตุปจฺจยา.  เหตุยา ตีณิ.}

\vspace{\tipitakaparskip}
\csromancenter{(นิสฺสยวารมฺปิ ฯเปฯ ปญฺหาวารมฺปิ วิตฺถาเรตพฺพํ.)}
\csromanmarkh{1. เหตุทุก}\csromanmarkhii{4. อุปาทินฺนตฺติก 1. เหตุทุก 2. เวทนาตฺติก}\csromanheadernomark{2}{1. เหตุทุก 4. อุปาทินฺนตฺติก \textbf{1. เหตุทุก 2. เวทนาตฺติก}}
\proseitem{8}{\csromanfolio{217}เหตุํ สุขาย เวทนาย สมฺปยุตฺตํ ธมฺมํ ปฏิจฺจ นเหตุ นสุขาย เวทนาย สมฺปยุตฺโต ธมฺโม อุปฺปชฺชติ เหตุปจฺจยา. (1)}

\prose{นเหตุํ สุขาย เวทนาย สมฺปยุตฺตํ ธมฺมํ ปฏิจฺจ นเหตุ นสุขาย เวทนาย สมฺปยุตฺโต ธมฺโม อุปฺปชฺชติ เหตุปจฺจยา. (1)}

\prose{เหตุํ สุขาย เวทนาย สมฺปยุตฺตญฺจ นเหตุํ สุขาย เวทนาย สมฺปยุตฺตญฺจ ธมฺมํ ปฏิจฺจ นเหตุ นสุขาย เวทนาย สมฺปยุตฺโต ธมฺโม อุปฺปชฺชติ เหตุปจฺจยา. (1) (สํขิตฺตํ.) . เหตุยา ตีณิ, อารมฺมเณ ตีณิ ฯเปฯ อวิคเต ตีณิ.}

\vspace{\tipitakaparskip}
\csromancenter{(สหชาตวารมฺปิ ฯเปฯ ปญฺหาวารมฺปิ วิตฺถาเรตพฺพํ.)}
\proseitem{9}{เหตุํ ทุกฺขาย เวทนาย สมฺปยุตฺตํ ธมฺมํ ปฏิจฺจ นเหตุ นทุกฺขาย เวทนาย สมฺปยุตฺโต ธมฺโม อุปฺปชฺชติ เหตุปจฺจยา. (สํขิตฺตํ.) . เหตุยา ตีณิ. (สพฺพตฺถ วิตฺถาโร.)}

\proseitem{10}{เหตุํ อทุกฺขม\-สุขาย เวทนาย สมฺปยุตฺตํ ธมฺมํ ปฏิจฺจ นเหตุ นอทุกฺขม\-สุขาย เวทนาย สมฺปยุตฺโต ธมฺโม อุปฺปชฺชติ เหตุปจฺจยา. (สํขิตฺตํ.) . เหตุยา ตีณิ. (สพฺพตฺถ วิตฺถาโร.)}

\csromanmarkh{1. เหตุทุก}\csromanmarkhii{3. วิปากตฺติก}\csromanheadernomark{2}{1. \textbf{เหตุทุก 3. วิปากตฺติก}}
\proseitem{11}{เหตุํ วิปากํ ธมฺมํ ปฏิจฺจ นเหตุ นวิปาโก ธมฺโม อุปฺปชฺชติ เหตุปจฺจยา.  เหตุยา ตีณิ.}

\prose{เหตุํ วิปาก\-ธมฺม\-ธมฺมํ ปฏิจฺจ นเหตุ นวิปาก\-ธมฺม\-ธมฺโม อุปฺปชฺชติ เหตุปจฺจยา.  เหตุยา ตีณิ.}

\prose{นเหตุํ เนว\-วิปาก\-น\-วิปากธมฺม\-ธมฺมํ ปฏิจฺจ นนเหตุ นเน\-ว\-วิปาก\-น\-วิปาก\-ธมฺม\-ธมฺโม อุปฺปชฺชติ เหตุปจฺจยา.  เหตุยา ตีณิ.}

\csromanmarkh{1. เหตุทุก}\csromanmarkhii{4. อุปาทินฺนตฺติก}\csromanheadernomark{2}{1. เหตุทุก 4. อุปาทินฺนตฺติก}
\proseitem{12}{เหตุํ อุปาทินฺนุ\-ปาทานิยํ ธมฺมํ ปฏิจฺจ นเหตุ นอุปาทินฺนุปาทานิโย ธมฺโม อุปฺปชฺชติ เหตุปจฺจยา.  เหตุยา ตีณิ.}

\csromanheader{2}{ธมฺมา\-นุโลม\-ปจฺจนีย ทุกติก\-ปฏฺฐาน\-ปาฬิ}
\prose{\csromanfolio{218}นเหตุํ อนุปาทินฺนุ\-ปาทานิยํ ธมฺมํ ปฏิจฺจ นนเหตุ อนุปาทินฺนุ\-ปาทานิโย ธมฺโม อุปฺปชฺชติ เหตุปจฺจยา.}

\prose{เหตุํ อนุปาทินฺน\-อนุปาทานิยํ ธมฺมํ ปฏิจฺจ นเหตุ นอนุปาทินฺน\-อนุปาทานิโย ธมฺโม อุปฺปชฺชติ เหตุปจฺจยา.  เหตุยา ตีณิ.}

\csromanmarkh{1. เหตุทุก}\csromanmarkhii{5. สํกิลิฏฺฐ\-ตฺติก}\csromanheadernomark{2}{1. เหตุทุก 5. สํกิลิฏฺฐ\-ตฺติก}
\proseitem{13}{เหตุํ สํกิลิฏฺฐ\-สํกิเลสิกํ ธมฺมํ ปฏิจฺจ นเหตุ นสํกิลิฏฺฐ\-สํกิเลสิโก ธมฺโม อุปฺปชฺชติ เหตุปจฺจยา.  เหตุยา ตีณิ.}

\prose{นเหตุํ อสํกิลิฏฺฐ\-สํกิเลสิกํ ธมฺมํ ปจฺจยา นนเหตุ นอสํกิลิฏฺฐ\-สํกิเลสิโก ธมฺโม อุปฺปชฺชติ เหตุปจฺจยา.  เหตุยา ตีณิ.}

\prose{เหตุํ อสํกิลิฏฺฐ\-อสํกิเลสิกํ ธมฺมํ ปฏิจฺจ นเหตุ นอสํกิลิฏฺฐ\-อสํกิเลสิโก ธมฺโม อุปฺปชฺชติ เหตุปจฺจยา.  เหตุยา ตีณิ.}

\csromanmarkh{1. เหตุทุก}\csromanmarkhii{6. วิตกฺกตฺติก}\csromanheadernomark{2}{1. \textbf{เหตุทุก 6. วิตกฺกตฺติก}}
\proseitem{14}{เหตุํ สวิตกฺก\-สวิจารํ ธมฺมํ ปฏิจฺจ นเหตุ นสวิตกฺก\-สวิจาโร ธมฺโม อุปฺปชฺชติ เหตุปจฺจยา.  เหตุยา ตีณิ.}

\prose{เหตุํ อวิตกฺก\-วิจารมตฺตํ ธมฺมํ ปฏิจฺจ นเหตุ นอวิตกฺก\-วิจารมตฺโต ธมฺโม อุปฺปชฺชติ เหตุปจฺจยา.  เหตุยา ตีณิ.}

\prose{นเหตุํ อวิตกฺก\-อวิจารํ ธมฺมํ ปฏิจฺจ นนเหตุ นอวิตกฺก\-อวิจาโร ธมฺโม อุปฺปชฺชติ เหตุปจฺจยา.  เหตุยา ตีณิ.}

\csromanmarkh{1. เหตุทุก}\csromanmarkhii{7. ปีติตฺติก}\csromanheadernomark{2}{1. \textbf{เหตุทุก 7. ปีติตฺติก}}
\proseitem{15}{เหตุํ ปีติสหคตํ ธมฺมํ ปฏิจฺจ นเหตุ นปีติสหคโต ธมฺโม อุปฺปชฺชติ เหตุปจฺจยา.  เหตุยา ตีณิ.}

\prose{เหตุํ สุขสหคตํ ธมฺมํ ปฏิจฺจ นเหตุ นสุขสหคโต ธมฺโม อุปฺปชฺชติ เหตุปจฺจยา.  เหตุยา ตีณิ.}

\prose{เหตุํ อุเปกฺขา\-สหคตํ ธมฺมํ ปฏิจฺจ นเหตุ นอุเปกฺขาสหคโต ธมฺโม อุปฺปชฺชติ เหตุปจฺจยา.  เหตุยา ตีณิ.}

\csromanmarkh{1. เหตุทุก}\csromanmarkhii{11. เสกฺขตฺติก 1. เหตุทุก 8. ทสฺสนตฺติก}\csromanheadernomark{2}{1. เหตุทุก 11. เสกฺขตฺติก \textbf{1. เหตุทุก 8. ทสฺสนตฺติก}}
\proseitem{16}{\csromanfolio{219}เหตุํ ทสฺสเนน ปหาตพฺพํ ธมฺมํ ปฏิจฺจ นเหตุ นทสฺสเนน ปหาตพฺโพ ธมฺโม อุปฺปชฺชติ เหตุปจฺจยา.  เหตุยา ตีณิ.}

\prose{เหตุํ ภาวนาย ปหาตพฺพํ ธมฺมํ ปฏิจฺจ นเหตุ นภาวนาย ปหาตพฺโพ ธมฺโม อุปฺปชฺชติ เหตุปจฺจยา.  เหตุยา ตีณิ.}

\prose{นเหตุํ เนวทสฺสเนน นภาวนาย ปหาตพฺพํ ธมฺมํ ปจฺจยา นนเหตุ นเน\-ว\-ทสฺสเนน นภาวนาย ปหาตพฺโพ ธมฺโม อุปฺปชฺชติ เหตุปจฺจยา.  เหตุยา ตีณิ.}

\csromanmarkh{1. เหตุทุก}\csromanmarkhii{9. ทสฺสน\-เหตุ\-ตฺติก}\csromanheadernomark{2}{1. เหตุทุก 9. ทสฺสน\-เหตุ\-ตฺติก}
\proseitem{17}{เหตุํ ทสฺสเนน ปหาตพฺพ\-เหตุกํ ธมฺมํ ปฏิจฺจ นเหตุ นทสฺสเนน ปหาตพฺพ\-เหตุโก ธมฺโม อุปฺปชฺชติ เหตุปจฺจยา.  เหตุยา ตีณิ.}

\prose{เหตุํ ภาวนาย ปหาตพฺพ\-เหตุกํ ธมฺมํ ปฏิจฺจ นเหตุ นภาวนาย ปหาตพฺพ\-เหตุโก ธมฺโม อุปฺปชฺชติ เหตุปจฺจยา.  เหตุยา ตีณิ.}

\prose{เหตุํ เนวทสฺสเนน นภาวนาย ปหาตพฺพ\-เหตุกํ ธมฺมํ ปฏิจฺจ นเหตุ นเน\-ว\-ทสฺสเนน นภาวนาย ปหาตพฺพ\-เหตุโก ธมฺโม อุปฺปชฺชติ เหตุปจฺจยา.  เหตุยา เอกํ.}

\csromanmarkh{1. เหตุทุก}\csromanmarkhii{10. อาจยคามิตฺติก}\csromanheadernomark{2}{1. เหตุทุก 10. อาจยคามิตฺติก}
\proseitem{18}{เหตุํ อาจยคามิํ ธมฺมํ ปฏิจฺจ นเหตุ นอาจยคามี ธมฺโม อุปฺปชฺชติ เหตุปจฺจยา.  เหตุยา ตีณิ.}

\prose{เหตุํ อปจยคามิํ ธมฺมํ ปฏิจฺจ นเหตุ นอปจยคามี ธมฺโม อุปฺปชฺชติ เหตุปจฺจยา.  เหตุยา ตีณิ.}

\prose{นเหตุํ เนวา\-จยคามิ\-นา\-ปจยคามิํ ธมฺมํ ปจฺจยา นนเหตุ นเน\-วา\-จยคามิ\-นา\-ปจยคามี ธมฺโม อุปฺปชฺชติ เหตุปจฺจยา.  เหตุยา ตีณิ.}

\csromanmarkh{1. เหตุทุก}\csromanmarkhii{11. เสกฺขตฺติก}\csromanheadernomark{2}{1. เหตุทุก 11. เสกฺขตฺติก}
\proseitem{19}{เหตุํ เสกฺขํ ธมฺมํ ปฏิจฺจ นเหตุ นเสกฺโข ธมฺโม อุปฺปชฺชติ เหตุปจฺจยา.  เหตุยา ตีณิ.}

\csromanheader{2}{ธมฺมา\-นุโลม\-ปจฺจนีย ทุกติก\-ปฏฺฐาน\-ปาฬิ}
\prose{\csromanfolio{220}เหตุํ อเสกฺขํ ธมฺมํ ปฏิจฺจ นเหตุ นอเสกฺโข ธมฺโม อุปฺปชฺชติ เหตุปจฺจยา.  เหตุยา ตีณิ.}

\prose{นเหตุํ เนว\-เสกฺข\-นาเสกฺขํ ธมฺมํ ปจฺจยา นนเหตุ นเน\-ว\-เสกฺข\-นา\-เสกฺโข ธมฺโม อุปฺปชฺชติ เหตุปจฺจยา.  เหตุยา ตีณิ.}

\csromanmarkh{1. เหตุทุก}\csromanmarkhii{12. ปริตฺตตฺติก}\csromanheadernomark{2}{1. เหตุทุก 12. ปริตฺตตฺติก}
\proseitem{20}{นเหตุํ ปริตฺตํ ธมฺมํ ปฏิจฺจ นนเหตุ นปริตฺโต ธมฺโม อุปฺปชฺชติ เหตุปจฺจยา. (ตีณิ.)}

\prose{เหตุํ มหคฺคตํ ธมฺมํ ปฏิจฺจ นเหตุ นมหคฺคโต ธมฺโม อุปฺปชฺชติ เหตุปจฺจยา.  เหตุยา ตีณิ.}

\prose{เหตุํ อปฺปมาณํ ธมฺมํ ปฏิจฺจ นเหตุ นอปฺปมาโณ ธมฺโม อุปฺปชฺชติ เหตุปจฺจยา.  เหตุยา ตีณิ.}

\csromanmarkh{1. เหตุทุก}\csromanmarkhii{13. ปริตฺตา\-รมฺมณ\-ตฺติก}\csromanheadernomark{2}{1. เหตุทุก 13. ปริตฺตา\-รมฺมณ\-ตฺติก}
\proseitem{21}{เหตุํ ปริตฺตา\-รมฺมณํ ธมฺมํ ปฏิจฺจ นเหตุ นปริตฺตา\-รมฺมโณ ธมฺโม อุปฺปชฺชติ เหตุปจฺจยา. (ตีณิ.)}

\prose{เหตุํ มหคฺคตา\-รมฺมณํ ธมฺมํ ปฏิจฺจ นเหตุ นมหคฺคตา\-รมฺมโณ ธมฺโม อุปฺปชฺชติ เหตุปจฺจยา. ตีณิ.}

\prose{เหตุํ อปฺปมาณา\-รมฺมณํ ธมฺมํ ปฏิจฺจ นเหตุ นอปฺปมาณารมฺมโณ ธมฺโม อุปฺปชฺชติ เหตุปจฺจยา.  เหตุยา ตีณิ.}

\csromanmarkh{1. เหตุทุก}\csromanmarkhii{14. หีนตฺติก}\csromanheadernomark{2}{1. เหตุทุก 14. หีนตฺติก}
\proseitem{22}{เหตุํ หีนํ ธมฺมํ ปฏิจฺจ นเหตุ นหีโน ธมฺโม อุปฺปชฺชติ เหตุปจฺจยา.  เหตุยา ตีณิ.}

\prose{นเหตุํ มชฺฌิมํ ธมฺมํ ปจฺจยา นนเหตุ นมชฺฌิโม ธมฺโม อุปฺปชฺชติ เหตุปจฺจยา.  เหตุยา ตีณิ.}

\prose{เหตุํ ปณีตํ ธมฺมํ ปฏิจฺจ นเหตุ นปณีโต ธมฺโม อุปฺปชฺชติ เหตุปจฺจยา.  เหตุยา ตีณิ.}

\csromanmarkh{1. เหตุทุก}\csromanmarkhii{18. อตีตตฺติก 1. เหตุทุก 15. มิจฺฉตฺตนิยต\-ตฺติก}\csromanheadernomark{2}{1. เหตุทุก 18. อตีตตฺติก 1. เหตุทุก 15. มิจฺฉตฺตนิยต\-ตฺติก}
\proseitem{23}{\csromanfolio{221}เหตุํ มิจฺฉตฺต\-นิยตํ ธมฺมํ ปฏิจฺจ นเหตุ นมิจฺฉตฺต\-นิยโต ธมฺโม อุปฺปชฺชติ เหตุปจฺจยา. (ตีณิ.)}

\prose{เหตุํ สมฺมตฺต\-นิยตํ ธมฺมํ ปฏิจฺจ นเหตุ นสมฺมตฺตนิยโต ธมฺโม อุปฺปชฺชติ เหตุปจฺจยา. ตีณิ.}

\prose{นเหตุํ อนิยตํ ธมฺมํ ปจฺจยา นนเหตุ นอนิยโต ธมฺโม อุปฺปชฺชติ เหตุปจฺจยา.  เหตุยา ตีณิ.}

\csromanmarkh{1. เหตุทุก}\csromanmarkhii{16. มคฺคา\-รมฺมณ\-ตฺติก}\csromanheadernomark{2}{1. เหตุทุก 16. มคฺคา\-รมฺมณ\-ตฺติก}
\proseitem{24}{เหตุํ มคฺคารมฺมณํ ธมฺมํ ปฏิจฺจ นเหตุ นมคฺคารมฺมโณ ธมฺโม อุปฺปชฺชติ เหตุปจฺจยา. ตีณิ.}

\prose{เหตุํ มคฺคเหตุกํ ธมฺมํ ปฏิจฺจ นเหตุ นมคฺคเหตุโก ธมฺโม อุปฺปชฺชติ เหตุปจฺจยา. ตีณิ.}

\prose{เหตุํ มคฺคาธิปติํ ธมฺมํ ปฏิจฺจ นเหตุ นมคฺคาธิปติ ธมฺโม อุปฺปชฺชติ เหตุปจฺจยา.  เหตุยา ตีณิ.}

\csromanmarkh{1. เหตุทุก}\csromanmarkhii{17. อุปฺปนฺนตฺติก}\csromanheadernomark{2}{1. เหตุทุก 17. อุปฺปนฺนตฺติก}
\proseitem{25}{เหตุ อนุปฺปนฺโน ธมฺโม นเหตุสฺส นอนุปฺปนฺนสฺส ธมฺมสฺส อารมฺมณ\-ปจฺจเยน ปจฺจโย.  อารมฺมเณ นว, อธิปติยา อุปนิสฺสเย นว.}

\prose{เหตุ อุปฺปาที ธมฺโม นเหตุสฺส นอุปฺปาทิสฺส ธมฺมสฺส อารมฺมณ\-ปจฺจเยน ปจฺจโย.  อารมฺมเณ นว.}

\csromanmarkh{1. เหตุทุก}\csromanmarkhii{18. อตีตตฺติก}\csromanheadernomark{2}{1. เหตุทุก 18. อตีตตฺติก}
\proseitem{26}{เหตุ อตีโต ธมฺโม นเหตุสฺส นอตีตสฺส ธมฺมสฺส อารมฺมณ\-ปจฺจเยน ปจฺจโย.  อารมฺมเณ นว.}

\prose{เหตุ อนาคโต ธมฺโม นเหตุสฺส นอนาคตสฺส ธมฺมสฺส อารมฺมณ\-ปจฺจเยน ปจฺจโย.  อารมฺมเณ นว.}

\csromanmarkh{ธมฺมา\-นุโลม\-ปจฺจนีย ทุกติก\-ปฏฺฐาน\-ปาฬิ}\csromanmarkhii{1. เหตุทุก 19. อตีตา\-รมฺมณ\-ตฺติก}\csromanheadernomark{2}{ธมฺมา\-นุโลม\-ปจฺจนีย ทุกติก\-ปฏฺฐาน\-ปาฬิ 1. เหตุทุก 19. อตีตา\-รมฺมณ\-ตฺติก}
\proseitem{27}{\csromanfolio{222}เหตุํ อตีตารมฺมณํ ธมฺมํ ปฏิจฺจ นเหตุ นอตีตารมฺมโณ ธมฺโม อุปฺปชฺชติ เหตุปจฺจยา.  เหตุยา ตีณิ.}

\prose{เหตุํ อนาคตา\-รมฺมณํ ธมฺมํ ปฏิจฺจ นเหตุ นอนาคตา\-รมฺมโณ ธมฺโม อุปฺปชฺชติ เหตุปจฺจยา.  เหตุยา ตีณิ.}

\prose{เหตุํ ปจฺจุปฺปนฺนา\-รมฺมณํ ธมฺมํ ปฏิจฺจ นเหตุ นปจฺจุปฺปนฺนา\-รมฺมโณ ธมฺโม อุปฺปชฺชติ เหตุปจฺจยา.  เหตุยา ตีณิ.}

\csromanmarkh{1. เหตุทุก}\csromanmarkhii{20. อชฺฌตฺตา\-รมฺมณ\-ตฺติก}\csromanheadernomark{2}{1. เหตุทุก 20. อชฺฌตฺตา\-รมฺมณ\-ตฺติก}
\proseitem{28}{เหตุํ อชฺฌตฺตา\-รมฺมณํ ธมฺมํ ปฏิจฺจ นเหตุ นอชฺฌตฺตา\-รมฺมโณ ธมฺโม อุปฺปชฺชติ เหตุปจฺจยา. (ตีณิ.)}

\prose{เหตุํ พหิทฺธา\-รมฺมณํ ธมฺมํ ปฏิจฺจ นเหตุ นพหิทฺธา\-รมฺมโณ ธมฺโม อุปฺปชฺชติ เหตุปจฺจยา.  เหตุยา ตีณิ.}

\csromanmarkh{1. เหตุทุก}\csromanmarkhii{21. สนิทสฺสน\-ตฺติก}\csromanheadernomark{2}{1. เหตุทุก 21. สนิทสฺสน\-ตฺติก}
\proseitem{29}{นเหตุ สนิทสฺสน\-สปฺปฏิโฆ ธมฺโม นนเหตุสฺส นสนิทสฺสน\-สปฺปฏิฆสฺส ธมฺมสฺส อารมฺมณ\-ปจฺจเยน ปจฺจโย ฯเปฯ. นเหตุ สนิทสฺสน\-สปฺปฏิโฆ ธมฺโม นเหตุสฺส นสนิทสฺสน\-สปฺปฏิฆสฺส จ นนเหตุสฺส นสนิทสฺสน\-สปฺปฏิฆสฺส จ ธมฺมสฺส อารมฺมณ\-ปจฺจเยน ปจฺจโย.  อารมฺมเณ ตีณิ, อธิปติยา อุปนิสฺสเย ปุเรชาเต อตฺถิยา อวิคเต ตีณิ.}

\prose{นเหตุํ อนิทสฺสน\-สปฺปฏิฆํ ธมฺมํ ปฏิจฺจ นเหตุ นอนิทสฺสน\-สปฺปฏิโฆ ธมฺโม อุปฺปชฺชติ เหตุปจฺจยา.  เหตุยา เอกํ ฯเปฯ อวิคเต เอกํ.}

\prose{เหตุํ อนิทสฺสน\-อปฺปฏิฆํ ธมฺมํ ปฏิจฺจ นเหตุ นอนิทสฺสน\-อปฺปฏิโฆ ธมฺโม อุปฺปชฺชติ เหตุปจฺจยา.  นเหตุํ อนิทสฺสน\-อปฺปฏิฆํ ธมฺมํ ปฏิจฺจ นเหตุ นอนิทสฺสน\-อปฺปฏิโฆ ธมฺโม อุปฺปชฺชติ เหตุปจฺจยา.  เหตุํ อนิทสฺสนอปฺปฏิฆญฺจ นเหตุํ อนิทสฺสนอปฺปฏิฆญฺจ ธมฺมํ ปฏิจฺจ นเหตุ นอนิทสฺสน\-อปฺปฏิโฆ ธมฺโม อุปฺปชฺชติ เหตุปจฺจยา. (3)}

\prose{เหตุยา ตีณิ.}

\vspace{\tipitakaparskip}
\csromancenter{เหตุ\-สเหตุกทุก 1. กุสลตฺติก 2. สเหตุกทุก 1. กุสลตฺติก}
\proseitem{30}{\csromanfolio{223}สเหตุกํ กุสลํ ธมฺมํ ปฏิจฺจ นสเหตุโก นกุสโล ธมฺโม อุปฺปชฺชติ เหตุปจฺจยา.  เหตุยา เอกํ.}

\prose{สเหตุกํ อกุสลํ ธมฺมํ ปฏิจฺจ นสเหตุโก นอกุสโล ธมฺโม อุปฺปชฺชติ เหตุปจฺจยา.  เหตุยา ตีณิ.}

\prose{อเหตุกํ อพฺยากตํ ธมฺมํ ปจฺจยา นอเหตุโก นอพฺยากโต ธมฺโม อุปฺปชฺชติ เหตุปจฺจยา.  เหตุยา เอกํ.}

\csromanmarkh{3. เหตุ\-สมฺปยุตฺตทุก}\csromanmarkhii{1. กุสลตฺติก}\csromanheadernomark{2}{3. เหตุ\-สมฺปยุตฺตทุก 1. กุสลตฺติก}
\proseitem{31}{เหตุ\-สมฺปยุตฺตํ กุสลํ ธมฺมํ ปฏิจฺจ นเหตุ\-สมฺปยุตฺโต นกุสโล ธมฺโม อุปฺปชฺชติ เหตุปจฺจยา.  เหตุยา เอกํ.}

\prose{เหตุ\-สมฺปยุตฺตํ อกุสลํ ธมฺมํ ปฏิจฺจ นเหตุ\-สมฺปยุตฺโต นอกุสโล ธมฺโม อุปฺปชฺชติ เหตุปจฺจยา.  เหตุยา ตีณิ.}

\prose{เหตุ\-วิปฺปยุตฺตํ อพฺยากตํ ธมฺมํ ปจฺจยา นเหตุ\-วิปฺปยุตฺโต นอพฺยากโต ธมฺโม อุปฺปชฺชติ เหตุปจฺจยา.  เหตุยา เอกํ.}

\csromanmarkh{4. เหตุ\-สเหตุกทุก}\csromanmarkhii{1. กุสลตฺติก}\csromanheadernomark{2}{4. เหตุ\-สเหตุกทุก 1. กุสลตฺติก}
\proseitem{32}{เหตุ เจว สเหตุโก จ กุสโล ธมฺโม นเหตุสฺส เจว นอเหตุกสฺส จ นกุสลสฺส ธมฺมสฺส อารมฺมณ\-ปจฺจเยน ปจฺจโย.  เหตุ เจว สเหตุโก จ กุสโล ธมฺโม นอเหตุกสฺส เจว นน จ เหตุสฺส นกุสลสฺส ธมฺมสฺส อารมฺมณ\-ปจฺจเยน ปจฺจโย.  เหตุ เจว สเหตุโก จ กุสโล ธมฺโม นเหตุสฺส เจว นอเหตุกสฺส จ นกุสลสฺส จ นอเหตุกสฺส เจว นน จ เหตุสฺส นกุสลสฺส จ ธมฺมสฺส อารมฺมณ\-ปจฺจเยน ปจฺจโย. (3)}

\prose{สเหตุโก เจว น จ เหตุ กุสโล ธมฺโม นอเหตุกสฺส เจว นน จ เหตุสฺส นกุสลสฺส ธมฺมสฺส อารมฺมณ\-ปจฺจเยน ปจฺจโย.  สเหตุโก เจว น จ เหตุ กุสโล ธมฺโม นเหตุสฺส เจว นอเหตุกสฺส จ นกุสลสฺส ธมฺมสฺส อารมฺมณ\-ปจฺจเยน ปจฺจโย.  }

\vspace{\tipitakaparskip}
\csromancenter{ทุกติก\-ปฏฺฐาน\-ปาฬิ สเหตุโก เจว น จ เหตุ กุสโล ธมฺโม นเหตุสฺส เจว นอเหตุกสฺส จ นกุสลสฺส จ นอเหตุกสฺส เจว นน จ เหตุสฺส นกุสลสฺส จ ธมฺมสฺส อารมฺมณ\-ปจฺจเยน ปจฺจโย. (3)}
\prose{\csromanfolio{224}เหตุ เจว สเหตุโก กุสโล จ สเหตุโก เจว น จ เหตุ กุสโล จ ธมฺมา นเหตุสฺส เจว นอเหตุกสฺส นกุสลสฺส ธมฺมสฺส อารมฺมณ\-ปจฺจเยน ปจฺจโย.  เหตุ เจว สเหตุโก กุสโล จ สเหตุโก เจว น จ เหตุ กุสโล จ ธมฺมา นอเหตุกสฺส เจว นน จ เหตุสฺส นกุสลสฺส ธมฺมสฺส อารมฺมณ\-ปจฺจเยน ปจฺจโย.  เหตุ เจว สเหตุโก กุสโล จ สเหตุโก เจว น จ เหตุ กุสโล จ ธมฺมา นเหตุสฺส เจว นอเหตุกสฺส นกุสลสฺส จ นอเหตุกสฺส เจว นน เหตุสฺส นกุสลสฺส จ ธมฺมสฺส อารมฺมณ\-ปจฺจเยน ปจฺจโย. (3)}

\prose{อารมฺมเณ นว.}

\proseitem{33}{เหตุ เจว สเหตุโก จ อกุสโล ธมฺโม นเหตุสฺส เจว นอเหตุกสฺส นอกุสลสฺส จ ธมฺมสฺส อารมฺมณ\-ปจฺจเยน ปจฺจโย. (เอเตน อุปาเยน นว ปญฺหา กาตพฺพา.)}

\proseitem{34}{เหตุ เจว สเหตุโก จ อพฺยากโต ธมฺโม นเหตุสฺส เจว นอเหตุกสฺส จ นอพฺยากตสฺส ธมฺมสฺส อารมฺมณ\-ปจฺจเยน ปจฺจโย. (นว ปญฺหา กาตพฺพา.) (สํขิตฺตํ.) (เหตุ\-เหตุ\-สมฺปยุตฺตทุกํ เหตุ\-สเหตุก\-ทุก\-สทิสํ. สํขิตฺตํ. นว ปญฺหา.)}

\csromanmarkh{6. นเหตุ\-สเหตุกทุก}\csromanmarkhii{1. กุสลตฺติก}\csromanheadernomark{2}{6. นเหตุ\-สเหตุกทุก 1. กุสลตฺติก}
\proseitem{35}{นเหตุํ สเหตุกํ กุสลํ ธมฺมํ ปฏิจฺจ นเหตุ นสเหตุโก นกุสโล ธมฺโม อุปฺปชฺชติ เหตุปจฺจยา.  เหตุยา เอกํ.}

\prose{นเหตุํ สเหตุกํ อกุสลํ ธมฺมํ ปฏิจฺจ นเหตุ นสเหตุโก นอกุสโล ธมฺโม อุปฺปชฺชติ เหตุปจฺจยา.  เหตุยา เอกํ.}

\prosewithcloser{นเหตุํ อเหตุกํ อพฺยากตํ ธมฺมํ ปจฺจยา นเหตุ นอเหตุโก นอพฺยากโต ธมฺโม อุปฺปชฺชติ เหตุปจฺจยา.  เหตุยา เอกํ.}{เหตุโคจฺฉกํ นิฏฺฐิตํ.}
\csromanmarkh{12. โลกิยทุก}\csromanmarkhii{1. กุสลตฺติก 7. สปฺปจฺจยทุก 1. กุสลตฺติก}\csromanheadernomark{2}{12. โลกิยทุก 1. กุสลตฺติก \textbf{7. สปฺปจฺจยทุก 1. กุสลตฺติก}}
\proseitem{36}{\csromanfolio{225}สปฺปจฺจยํ กุสลํ ธมฺมํ ปฏิจฺจ นอปฺปจฺจโย นกุสโล ธมฺโม อุปฺปชฺชติ เหตุปจฺจยา.  เหตุยา เอกํ.}

\prose{สปฺปจฺจยํ อกุสลํ ธมฺมํ ปฏิจฺจ นอปฺปจฺจโย นอกุสโล ธมฺโม อุปฺปชฺชติ เหตุปจฺจยา.  เหตุยา เอกํ.}

\prose{สปฺปจฺจยํ อพฺยากตํ ธมฺมํ ปจฺจยา นอปฺปจฺจโย นอพฺยากโต ธมฺโม อุปฺปชฺชติ เหตุปจฺจยา.  เหตุยา เอกํ ฯเปฯ. (สํขิตฺตํ. สปฺปจฺจย\-สทิสํ.)}

\csromanmarkh{9. สนิทสฺสนทุก}\csromanmarkhii{1. กุสลตฺติก}\csromanheadernomark{2}{9. \textbf{สนิทสฺสนทุก 1. กุสลตฺติก}}
\proseitem{37}{อนิทสฺสนํ กุสลํ ธมฺมํ ปฏิจฺจ นอนิทสฺสโน นกุสโล ธมฺโม อุปฺปชฺชติ เหตุปจฺจยา.  เหตุยา ตีณิ. (อกุสลํ กุสลสทิสํ.)}

\prose{อนิทสฺสนํ อพฺยากตํ ธมฺมํ ปจฺจยา นสนิทสฺสโน นอพฺยากโต ธมฺโม อุปฺปชฺชติ เหตุปจฺจยา.  เหตุยา เอกํ.}

\prose{อปฺปฏิฆํ กุสลํ ธมฺมํ ปฏิจฺจ นอปฺปฏิโฆ นกุสโล ธมฺโม อุปฺปชฺชติ เหตุปจฺจยา.  เหตุยา ตีณิ.}

\prose{อปฺปฏิฆํ อกุสลํ ธมฺมํ ปฏิจฺจ นอปฺปฏิโฆ นอกุสโล ธมฺโม อุปฺปชฺชติ เหตุปจฺจยา.  เหตุยา ตีณิ.}

\prose{อพฺยากเต เอกํ.}

\csromanmarkh{11. รูปีทุก}\csromanmarkhii{1. กุสลตฺติก}\csromanheadernomark{2}{11. \textbf{รูปีทุก 1. กุสลตฺติก}}
\proseitem{38}{อรูปิํ กุสลํ ธมฺมํ ปฏิจฺจ นอรูปี นกุสโล ธมฺโม อุปฺปชฺชติ เหตุปจฺจยา.  เหตุยา เอกํ.}

\prose{อรูปิํ อกุสลํ ธมฺมํ ปฏิจฺจ นอรูปี นอกุสโล ธมฺโม อุปฺปชฺชติ เหตุปจฺจยา.  เหตุยา เอกํ.}

\prose{รูปิํ อพฺยากตํ ธมฺมํ ปจฺจยา นรูปี นอพฺยากโต ธมฺโม อุปฺปชฺชติ เหตุปจฺจยา.  เหตุยา เอกํ.}

\csromanmarkh{12. โลกิยทุก}\csromanmarkhii{1. กุสลตฺติก}\csromanheadernomark{2}{12. โลกิยทุก 1. กุสลตฺติก}
\proseitem{39}{โลกิยํ กุสลํ ธมฺมํ ปฏิจฺจ นโลกุตฺตโร นกุสโล ธมฺโม อุปฺปชฺชติ เหตุปจฺจยา.  เหตุยา เทฺว.}

\csromanheader{2}{ธมฺมา\-นุโลม\-ปจฺจนีย ทุกติก\-ปฏฺฐาน\-ปาฬิ}
\prose{\csromanfolio{226}โลกิยํ อกุสลํ ธมฺมํ ปฏิจฺจ นโลกุตฺตโร นอกุสโล ธมฺโม อุปฺปชฺชติ เหตุปจฺจยา.  เหตุยา เอกํ.}

\prose{โลกิยํ อพฺยากตํ ธมฺมํ ปจฺจยา นโลกิโย นอพฺยากโต ธมฺโม อุปฺปชฺชติ เหตุปจฺจยา.  เหตุยา เทฺว.}

\csromanmarkh{13. เกนจิ\-วิญฺเญยฺยทุก}\csromanmarkhii{1. กุสลตฺติก}\csromanheadernomark{2}{13. เกนจิ\-วิญฺเญยฺยทุก 1. กุสลตฺติก}
\proseitem{40}{เกนจิ วิญฺเญยฺยํ กุสลํ ธมฺมํ ปฏิจฺจ นเกนจิ วิญฺเญยฺโย นกุสโล ธมฺโม อุปฺปชฺชติ เหตุปจฺจยา.  เหตุยา นว.}

\prose{เกนจิ วิญฺเญยฺยํ อกุสลํ ธมฺมํ ปฏิจฺจ นเกนจิ วิญฺเญยฺโย นอกุสโล ธมฺโม อุปฺปชฺชติ เหตุปจฺจยา.  เหตุยา นว.}

\prosewithcloser{เกนจิ วิญฺเญยฺยํ อพฺยากตํ ธมฺมํ ปจฺจยา นเกนจิ วิญฺเญยฺโย นอพฺยากโต ธมฺโม อุปฺปชฺชติ เหตุปจฺจยา.  เหตุยา นว.}{จูฬนฺตรทุกํ นิฏฺฐิตํ.}
\csromanmarkh{14. อาสวทุก}\csromanmarkhii{1. กุสลตฺติก}\csromanheadernomark{2}{14. \textbf{อาสวทุก 1. กุสลตฺติก}}
\proseitem{41}{โนอาสวํ กุสลํ ธมฺมํ ปฏิจฺจ โนอาสโว นกุสโล ธมฺโม อุปฺปชฺชติ เหตุปจฺจยา.  เหตุยา เอกํ.}

\prose{อาสวํ อกุสลํ ธมฺมํ ปฏิจฺจ โนอาสโว นอกุสโล ธมฺโม อุปฺปชฺชติ เหตุปจฺจยา.  เหตุยา ตีณิ.}

\prose{โนอาสวํ อพฺยากตํ ธมฺมํ ปจฺจยา นโนอาสโว นอพฺยากโต ธมฺโม อุปฺปชฺชติ เหตุปจฺจยา.  เหตุยา ตีณิ.}

\csromanmarkh{15. สาสวทุก}\csromanmarkhii{1. กุสลตฺติก}\csromanheadernomark{2}{15. \textbf{สาสวทุก 1. กุสลตฺติก}}
\proseitem{42}{สาสวํ กุสลํ ธมฺมํ ปฏิจฺจ นอนาสโว นกุสโล ธมฺโม อุปฺปชฺชติ เหตุปจฺจยา.  อนาสวํ กุสลํ ธมฺมํ ปฏิจฺจ นอนาสโว นกุสโล ธมฺโม อุปฺปชฺชติ เหตุปจฺจยา.  เหตุยา เทฺว.}

\prose{สาสวํ อกุสลํ ธมฺมํ ปฏิจฺจ นอนาสโว นอกุสโล ธมฺโม อุปฺปชฺชติ เหตุปจฺจยา.  เหตุยา เอกํ.}

\csromanmarkh{19. อาสว\-วิปฺปยุตฺต\-สาสวทุก}\csromanmarkhii{1. กุสลตฺติก}\csromanheadernomark{2}{19. อาสว\-วิปฺปยุตฺต\-สาสวทุก 1. กุสลตฺติก}
\prose{\csromanfolio{227}สาสวํ อพฺยากตํ ธมฺมํ ปจฺจยา นสาสโว นอพฺยากโต ธมฺโม อุปฺปชฺชติ เหตุปจฺจยา.  เหตุยา เทฺว.}

\csromanmarkh{16. อาสว\-สมฺปยุตฺตทุก}\csromanmarkhii{1. กุสลตฺติก}\csromanheadernomark{2}{16. อาสว\-สมฺปยุตฺตทุก 1. กุสลตฺติก}
\proseitem{43}{อาสว\-วิปฺปยุตฺตํ กุสลํ ธมฺมํ ปฏิจฺจ นอาสว\-สมฺปยุตฺโต นกุสโล ธมฺโม อุปฺปชฺชติ เหตุปจฺจยา.  เหตุยา เอกํ.}

\prose{อาสว\-สมฺปยุตฺตํ อกุสลํ ธมฺมํ ปฏิจฺจ นอาสว\-สมฺปยุตฺโต นอกุสโล ธมฺโม อุปฺปชฺชติ เหตุปจฺจยา.  เหตุยา ตีณิ.}

\prose{อาสว\-วิปฺปยุตฺตํ อพฺยากตํ ธมฺมํ ปจฺจยา นอาสว\-สมฺปยุตฺโต นอพฺยากโต ธมฺโม อุปฺปชฺชติ เหตุปจฺจยา.  เหตุยา ตีณิ.}

\csromanmarkh{17. อาสว\-สาสวทุก}\csromanmarkhii{1. กุสลตฺติก}\csromanheadernomark{2}{17. อาสว\-สาสวทุก 1. กุสลตฺติก}
\proseitem{44}{สาสวญฺเจว โน จ อาสวํ กุสลํ ธมฺมํ ปฏิจฺจ นอาสโว เจว นอนาสโว จ นกุสโล ธมฺโม อุปฺปชฺชติ เหตุปจฺจยา.  เหตุยา เอกํ.}

\prose{อาสวญฺเจว สาสวญฺจ อกุสลํ ธมฺมํ ปฏิจฺจ โนอาสโว เจว นอนาสโว จ นอกุสโล ธมฺโม อุปฺปชฺชติ เหตุปจฺจยา.  เหตุยา ตีณิ.}

\prose{สาสวญฺเจว โน จ อาสวํ อพฺยากตํ ธมฺมํ ปจฺจยา นอนาสโว เจว นโน จ อาสโว นอพฺยากโต ธมฺโม อุปฺปชฺชติ เหตุปจฺจยา.  เหตุยา ตีณิ. (อาสวอาสวสมฺปยุตฺตทุกํ อาสวสาสวทุกสทิสํ.)}

\vspace{\tipitakaparskip}
\csromancenter{(อาสว\-อาสว\-สมฺปยุตฺตทุกํ อาสว\-สาสว\-ทุก\-สทิสํ.)}
\csromanmarkh{19. อาสว\-วิปฺปยุตฺต\-สาสวทุก}\csromanmarkhii{1. กุสลตฺติก}\csromanheadernomark{2}{19. อาสว\-วิปฺปยุตฺต\-สาสวทุก 1. กุสลตฺติก}
\proseitem{45}{อาสว\-วิปฺปยุตฺตํ สาสวํ กุสลํ ธมฺมํ ปฏิจฺจ อาสว\-วิปฺปยุตฺโต นอนาสโว นกุสโล ธมฺโม อุปฺปชฺชติ เหตุปจฺจยา. เอกํ.}

\prose{อาสว\-วิปฺปยุตฺตํ อนาสวํ กุสลํ ธมฺมํ ปฏิจฺจ อาสว\-วิปฺปยุตฺโต นอนาสโว นกุสโล ธมฺโม อุปฺปชฺชติ เหตุปจฺจยา. เอกํ.}

\csromanheader{2}{ธมฺมา\-นุโลม\-ปจฺจนีย ทุกติก\-ปฏฺฐาน\-ปาฬิ}
\prose{\csromanfolio{228}อาสว\-วิปฺปยุตฺตํ สาสวํ อกุสลํ ธมฺมํ ปฏิจฺจ อาสว\-วิปฺปยุตฺโต นอนาสโว นอกุสโล ธมฺโม อุปฺปชฺชติ เหตุปจฺจยา.  เหตุยา เอกํ.}

\prosewithcloser{อาสว\-วิปฺปยุตฺตํ สาสวํ อพฺยากตํ ธมฺมํ ปจฺจยา อาสว\-วิปฺปยุตฺโต นสาสโว นอพฺยากโต ธมฺโม อุปฺปชฺชติ เหตุปจฺจยา.  อาสว\-วิปฺปยุตฺตํ สาสวํ อพฺยากตํ ธมฺมํ ปจฺจยา อาสว\-วิปฺปยุตฺโต นอนาสโว นอพฺยากโต ธมฺโม อุปฺปชฺชติ เหตุปจฺจยา.  เหตุยา เทฺว.}{อาสวโคจฺฉกํ นิฏฺฐิตํ.}
\csromanmarkh{20. สญฺโญชนา\-ทิทุก}\csromanmarkhii{1. กุสลตฺติก}\csromanheadernomark{2}{20. สญฺโญชนา\-ทิทุก 1. กุสลตฺติก}
\proseitem{46}{โนสญฺโญชนํ กุสลํ ธมฺมํ ปฏิจฺจ โนสญฺโญชโน นกุสโล ธมฺโม อุปฺปชฺชติ เหตุปจฺจยา ฯเปฯ. โนคนฺถํ กุสลํ ธมฺมํ ปฏิจฺจ โนคนฺโถ นกุสโล ธมฺโม อุปฺปชฺชติ เหตุปจฺจยา ฯเปฯ. โนโอฆํ กุสลํ ธมฺมํ ปฏิจฺจ โนโอโฆ นกุสโล ธมฺโม อุปฺปชฺชติ เหตุปจฺจยา ฯเปฯ. โนโยคํ กุสลํ ธมฺมํ ปฏิจฺจ โนโยโค นกุสโล ธมฺโม อุปฺปชฺชติ เหตุปจฺจยา ฯเปฯ. โนนีวรณํ กุสลํ ธมฺมํ ปฏิจฺจ โนนีวรโณ นกุสโล ธมฺโม อุปฺปชฺชติ เหตุปจฺจยา ฯเปฯ. โนปรามาสํ กุสลํ ธมฺมํ ปฏิจฺจ โนปรามาโส นกุสโล ธมฺโม อุปฺปชฺชติ เหตุปจฺจยา.}

\csromanmarkh{55. สารมฺมณทุก}\csromanmarkhii{1. กุสลตฺติก}\csromanheadernomark{2}{55. \textbf{สารมฺมณทุก 1. กุสลตฺติก}}
\proseitem{47}{สารมฺมณํ กุสลํ ธมฺมํ ปฏิจฺจ นสารมฺมโณ นกุสโล ธมฺโม อุปฺปชฺชติ เหตุปจฺจยา.  เหตุยา เอกํ.}

\prose{สารมฺมณํ อกุสลํ ธมฺมํ ปฏิจฺจ นสารมฺมโณ นอกุสโล ธมฺโม อุปฺปชฺชติ เหตุปจฺจยา.  เหตุยา เอกํ.}

\prose{อนารมฺมณํ อพฺยากตํ ธมฺมํ ปจฺจยา นอนารมฺมโณ นอพฺยากโต ธมฺโม อุปฺปชฺชติ เหตุปจฺจยา.  เหตุยา เอกํ.}

\csromanmarkh{56. จิตฺตทุก}\csromanmarkhii{1. กุสลตฺติก}\csromanheadernomark{2}{56. \textbf{จิตฺตทุก 1. กุสลตฺติก}}
\proseitem{48}{จิตฺตํ กุสลํ ธมฺมํ ปฏิจฺจ นจิตฺโต นกุสโล ธมฺโม อุปฺปชฺชติ เหตุปจฺจยา.  เหตุยา ตีณิ.}

\csromanmarkh{60. จิตฺต\-สมุฏฺฐานทุก}\csromanmarkhii{1. กุสลตฺติก}\csromanheadernomark{2}{60. จิตฺต\-สมุฏฺฐานทุก 1. กุสลตฺติก}
\prose{\csromanfolio{229}จิตฺตํ อกุสลํ ธมฺมํ ปฏิจฺจ นจิตฺโต นอกุสโล ธมฺโม อุปฺปชฺชติ เหตุปจฺจยา.  เหตุยา ตีณิ.}

\prose{โนจิตฺตํ อพฺยากตํ ธมฺมํ ปจฺจยา นโนจิตฺโต นอพฺยากโต ธมฺโม อุปฺปชฺชติ เหตุปจฺจยา.  เหตุยา ตีณิ.}

\csromanmarkh{57. เจตสิกทุก}\csromanmarkhii{1. กุสลตฺติก}\csromanheadernomark{2}{57. \textbf{เจตสิกทุก 1. กุสลตฺติก}}
\proseitem{49}{เจตสิกํ กุสลํ ธมฺมํ ปฏิจฺจ นเจตสิโก นกุสโล ธมฺโม อุปฺปชฺชติ เหตุปจฺจยา.  เหตุยา ตีณิ.}

\prose{เจตสิกํ อกุสลํ ธมฺมํ ปฏิจฺจ นเจตสิโก นอกุสโล ธมฺโม อุปฺปชฺชติ เหตุปจฺจยา.  เหตุยา ตีณิ.}

\prose{อเจตสิกํ อพฺยากตํ ธมฺมํ ปจฺจยา นอเจตสิโก นอพฺยากโต ธมฺโม อุปฺปชฺชติ เหตุปจฺจยา.  เหตุยา ตีณิ.}

\csromanmarkh{58. จิตฺต\-สมฺปยุตฺตทุก}\csromanmarkhii{1. กุสลตฺติก}\csromanheadernomark{2}{58. จิตฺต\-สมฺปยุตฺตทุก 1. กุสลตฺติก}
\proseitem{50}{จิตฺต\-สมฺปยุตฺตํ กุสลํ ธมฺมํ ปฏิจฺจ นจิตฺต\-สมฺปยุตฺโต นกุสโล ธมฺโม อุปฺปชฺชติ เหตุปจฺจยา.  เหตุยา เอกํ.}

\prose{จิตฺต\-สมฺปยุตฺตํ อกุสลํ ธมฺมํ ปฏิจฺจ นจิตฺต\-สมฺปยุตฺโต นอกุสโล ธมฺโม อุปฺปชฺชติ เหตุปจฺจยา.  เหตุยา เอกํ. (อพฺยากตมูลํ ตีณิเยว ปจฺจยวเสน.)}

\csromanmarkh{59. จิตฺต\-สํสฏฺฐทุก}\csromanmarkhii{1. กุสลตฺติก}\csromanheadernomark{2}{59. จิตฺต\-สํสฏฺฐทุก 1. กุสลตฺติก}
\proseitem{51}{จิตฺต\-สํสฏฺฐํ กุสลํ ธมฺมํ ปฏิจฺจ นจิตฺต\-สํสฏฺโฐ นกุสโล ธมฺโม อุปฺปชฺชติ เหตุปจฺจยา.  เหตุยา เอกํ.}

\prose{จิตฺต\-สํสฏฺฐํ อกุสลํ ธมฺมํ ปฏิจฺจ นจิตฺต\-สํสฏฺโฐ นอกุสโล ธมฺโม อุปฺปชฺชติ เหตุปจฺจยา.  เหตุยา เอกํ. (อพฺยากตมูลํ ตีณิเยว ปจฺจยวเสน.)}

\csromanmarkh{60. จิตฺต\-สมุฏฺฐานทุก}\csromanmarkhii{1. กุสลตฺติก}\csromanheadernomark{2}{60. จิตฺต\-สมุฏฺฐานทุก 1. กุสลตฺติก}
\proseitem{52}{จิตฺต\-สมุฏฺฐานํ กุสลํ ธมฺมํ ปฏิจฺจ นโน\-จิตฺตสมุฏฺฐาโน นกุสโล ธมฺโม อุปฺปชฺชติ เหตุปจฺจยา.  เหตุยา ตีณิ.}

\csromanheader{2}{ธมฺมา\-นุโลม\-ปจฺจนีย ทุกติก\-ปฏฺฐาน\-ปาฬิ}
\prose{\csromanfolio{230}จิตฺต\-สมุฏฺฐานํ อกุสลํ ธมฺมํ ปฏิจฺจ นโน\-จิตฺตสมุฏฺฐาโน นอกุสโล ธมฺโม อุปฺปชฺชติ เหตุปจฺจยา.  เหตุยา ตีณิ. (อพฺยากตมูลํ ตีณิเยว ปจฺจยวเสน.)}

\csromanmarkh{61. จิตฺต\-สหภูทุก}\csromanmarkhii{1. กุสลตฺติก}\csromanheadernomark{2}{61. จิตฺต\-สหภูทุก 1. กุสลตฺติก}
\proseitem{53}{จิตฺตสหภุํ กุสลํ ธมฺมํ ปฏิจฺจ โนจิตฺตสหภู นกุสโล ธมฺโม อุปฺปชฺชติ เหตุปจฺจยา.  เหตุยา นว.}

\prose{จิตฺตสหภุํ อกุสลํ ธมฺมํ ปฏิจฺจ โนจิตฺตสหภู นอกุสโล ธมฺโม อุปฺปชฺชติ เหตุปจฺจยา.  เหตุยา นว. (อพฺยากตมูลํ ตีณิเยว ปจฺจยวเสน.)}

\csromanmarkh{62. จิตฺตา\-นุปริวตฺติทุก}\csromanmarkhii{1. กุสลตฺติก}\csromanheadernomark{2}{62. จิตฺตา\-นุปริวตฺติทุก 1. กุสลตฺติก}
\proseitem{54}{จิตฺตา\-นุปริวตฺติํ กุสลํ ธมฺมํ ปฏิจฺจ โนจิตฺตา\-นุปริวตฺตี นกุสโล ธมฺโม อุปฺปชฺชติ เหตุปจฺจยา.  เหตุยา นว.}

\prose{จิตฺตา\-นุปริวตฺติํ อกุสลํ ธมฺมํ ปฏิจฺจ โนจิตฺตา\-นุปริวตฺตี นอกุสโล ธมฺโม อุปฺปชฺชติ เหตุปจฺจยา.  เหตุยา นว. (อพฺยากตมูลํ ตีณิเยว ปจฺจยวเสน.)}

\csromanmarkh{63. จิตฺต\-สํสฏฺฐ\-สมุฏฺฐานทุก}\csromanmarkhii{1. กุสลตฺติก}\csromanheadernomark{2}{63. จิตฺต\-สํสฏฺฐ\-สมุฏฺฐานทุก 1. กุสลตฺติก}
\proseitem{55}{จิตฺต\-สํสฏฺฐ\-สมุฏฺฐานํ กุสลํ ธมฺมํ ปฏิจฺจ โนจิตฺต\-สํสฏฺฐ\-สมุฏฺฐาโน นกุสโล ธมฺโม อุปฺปชฺชติ เหตุปจฺจยา.  เหตุยา ตีณิ.}

\prose{จิตฺต\-สํสฏฺฐ\-สมุฏฺฐานํ อกุสลํ ธมฺมํ ปฏิจฺจ โนจิตฺต\-สํสฏฺฐ\-สมุฏฺฐาโน นอกุสโล ธมฺโม อุปฺปชฺชติ เหตุปจฺจยา.  เหตุยา ตีณิ. (อพฺยากตมูลํ ตีณิเยว ปจฺจยวเสน.)}

\prose{64. จิตฺต\-สํสฏฺฐ\-สมุฏฺฐาน\-สหภูทุก 1. กุสลตฺติก}

\proseitem{56}{จิตฺตสํสฏฺฐสมุฏฺฏฺหานสหภุํ กุสลํ ธมฺมํ ปฏิจฺจ โนจิตฺต\-สํสฏฺฐ\-สมุฏฺฐาน\-สหภู นกุสโล ธมฺโม อุปฺปชฺชติ เหตุปจฺจยา.  เหตุยา ตีณิ.}

\prose{จิตฺตสํสฏฺฐสมุฏฺฏฺหานสหภุํ อกุสลํ ธมฺมํ ปฏิจฺจ โนจิตฺต\-สํสฏฺฐ\-สมุฏฺฐาน\-สหภู นอกุสโล ธมฺโม อุปฺปชฺชติ เหตุปจฺจยา.  เหตุยา ตีณิ. (อพฺยากตมูลํ ตีณิเยว ปจฺจยวเสน.)}

\vspace{\tipitakaparskip}
\csromancenter{อุปาทินฺนทุก 1. กุสลตฺติก 65. จิตฺต\-สํสฏฺฐ\-สมุฏฺฐานา\-นุปริวตฺติทุก 1. กุสลตฺติก}
\proseitem{57}{\csromanfolio{231}จิตฺตสํสฏฺฐสมุฏฺฐฺอานานุปริวตฺติํ กุสลํ ธมฺมํ ปฏิจฺจ โนจิตฺต\-สํสฏฺฐ\-สมุฏฺฐานา\-นุปริวตฺตี นกุสโล ธมฺโม อุปฺปชฺชติ เหตุปจฺจยา.  เหตุยา ตีณิ.}

\prose{จิตฺตสํสฏฺฐสมุฏฺฐฺอานานุปริวตฺติํ อกุสลํ ธมฺมํ ปฏิจฺจ โนจิตฺต\-สํสฏฺฐ\-สมุฏฺฐานา\-นุปริวตฺตี นอกุสโล ธมฺโม อุปฺปชฺชติ เหตุปจฺจยา.  เหตุยา ตีณิ. (อพฺยากตมูลํ ตีณิเยว.)}

\csromanmarkh{66. อชฺฌตฺติกทุก}\csromanmarkhii{1. กุสลตฺติก}\csromanheadernomark{2}{66. \textbf{อชฺฌตฺติกทุก 1. กุสลตฺติก}}
\proseitem{58}{อชฺฌตฺติกํ กุสลํ ธมฺมํ ปฏิจฺจ นอชฺฌตฺติโก นกุสโล ธมฺโม อุปฺปชฺชติ เหตุปจฺจยา.  พาหิรํ กุสลํ ธมฺมํ ปฏิจฺจ นอชฺฌตฺติโก นกุสโล ธมฺโม อุปฺปชฺชติ เหตุปจฺจยา.  อชฺฌตฺติกํ กุสลญฺจ พาหิรํ กุสลญฺจ ธมฺมํ ปฏิจฺจ นอชฺฌตฺติโก นกุสโล ธมฺโม อุปฺปชฺชติ เหตุปจฺจยา.  เหตุยา ตีณิ.}

\prose{อชฺฌตฺติกํ อกุสลํ ธมฺมํ ปฏิจฺจ นอชฺฌตฺติโก นอกุสโล ธมฺโม อุปฺปชฺชติ เหตุปจฺจยา.  พาหิรํ อกุสลํ ธมฺมํ ปฏิจฺจ นอชฺฌตฺติโก นอกุสโล ธมฺโม อุปฺปชฺชติ เหตุปจฺจยา.  อชฺฌตฺติกํ อกุสลญฺจ พาหิรํ อกุสลญฺจ ธมฺมํ ปฏิจฺจ นอชฺฌตฺติโก นอกุสโล ธมฺโม อุปฺปชฺชติ เหตุปจฺจยา.  เหตุยา ตีณิ. (อพฺยากตมูลํ ตีณิเยว.)}

\csromanmarkh{67. อุปาทาทุก}\csromanmarkhii{1. กุสลตฺติก}\csromanheadernomark{2}{67. \textbf{อุปาทาทุก 1. กุสลตฺติก}}
\proseitem{59}{โนอุปาทา กุสลํ ธมฺมํ ปฏิจฺจ นโนอุปาทา นกุสโล ธมฺโม อุปฺปชฺชติ เหตุปจฺจยา.  เหตุยา ตีณิ.}

\prose{โนอุปาทา อกุสลํ ธมฺมํ ปฏิจฺจ นโนอุปาทา นอกุสโล ธมฺโม อุปฺปชฺชติ เหตุปจฺจยา.  เหตุยา ตีณิ. (อพฺยากตมูลํ เอกํเยว.)}

\csromanmarkh{68. อุปาทินฺนทุก}\csromanmarkhii{1. กุสลตฺติก}\csromanheadernomark{2}{68. อุปาทินฺนทุก 1. กุสลตฺติก}
\proseitem{60}{อนุปาทินฺนํ กุสลํ ธมฺมํ ปฏิจฺจ นอุปาทินฺโน นกุสโล ธมฺโม อุปฺปชฺชติ เหตุปจฺจยา.  เหตุยา เอกํ.}

\csromanheader{2}{ธมฺมา\-นุโลม\-ปจฺจนีย ทุกติก\-ปฏฺฐาน\-ปาฬิ}
\prosewithcloser{\csromanfolio{232}อนุปาทินฺนํ อกุสลํ ธมฺมํ ปฏิจฺจ นอุปาทินฺโน นอกุสโล ธมฺโม อุปฺปชฺชติ เหตุปจฺจยา.  เหตุยา เอกํ. (อพฺยากตมูลํ เอกํเยว.)}{มหนฺตรทุกํ นิฏฺฐิตํ.}
\csromanmarkh{69. ทฺวิโคจฺฉก}\csromanmarkhii{1. กุสลตฺติก}\csromanheadernomark{2}{69. \textbf{ทฺวิโคจฺฉก 1. กุสลตฺติก}}
\proseitem{61}{โนอุปาทานํ กุสลํ ธมฺมํ ปฏิจฺจ โนอุปาทาโน นกุสโล ธมฺโม อุปฺปชฺชติ เหตุปจฺจยา.  เหตุยา เอกํ.}

\prose{อุปาทานํ อกุสลํ ธมฺมํ ปฏิจฺจ โนอุปาทาโน นอกุสโล ธมฺโม อุปฺปชฺชติ เหตุปจฺจยา.  เหตุยา ตีณิ.}

\prose{โนกิเลสํ กุสลํ ธมฺมํ ปฏิจฺจ โนกิเลโส นกุสโล ธมฺโม อุปฺปชฺชติ เหตุปจฺจยา.  เหตุยา เอกํ.}

\csromanmarkh{83. ทสฺสเนน\-ปหาตพฺพทุก}\csromanmarkhii{1. กุสลตฺติก}\csromanheadernomark{2}{83. ทสฺสเนน\-ปหาตพฺพทุก 1. กุสลตฺติก}
\proseitem{62}{นทสฺสเนน ปหาตพฺพํ กุสลํ ธมฺมํ ปฏิจฺจ นทสฺสเนน ปหาตพฺโพ นกุสโล ธมฺโม อุปฺปชฺชติ เหตุปจฺจยา.  เหตุยา เอกํ.}

\prose{ทสฺสเนน ปหาตพฺพํ อกุสลํ ธมฺมํ ปฏิจฺจ นทสฺสเนน ปหาตพฺโพ นอกุสโล ธมฺโม อุปฺปชฺชติ เหตุปจฺจยา.  นทสฺสเนน ปหาตพฺพํ อกุสลํ ธมฺมํ ปฏิจฺจ นทสฺสเนน ปหาตพฺโพ นอกุสโล ธมฺโม อุปฺปชฺชติ เหตุปจฺจยา.  เหตุยา เทฺว.}

\prose{นทสฺสเนน ปหาตพฺพํ อพฺยากตํ ธมฺมํ ปจฺจยา นนทสฺสเนน ปหาตพฺโพ นอพฺยากโต ธมฺโม อุปฺปชฺชติ เหตุปจฺจยา.  เหตุยา เทฺว.}

\csromanmarkh{84. ภาวนาย\-ปหาตพฺพทุก}\csromanmarkhii{1. กุสลตฺติก}\csromanheadernomark{2}{84. ภาวนาย\-ปหาตพฺพทุก 1. กุสลตฺติก}
\proseitem{63}{นภาวนาย ปหาตพฺพํ กุสลํ ธมฺมํ ปฏิจฺจ นภาวนาย ปหาตพฺโพ นกุสโล ธมฺโม อุปฺปชฺชติ เหตุปจฺจยา.  เหตุยา เอกํ.}

\prose{ภาวนาย ปหาตพฺพํ อกุสลํ ธมฺมํ ปฏิจฺจ นภาวนาย ปหาตพฺโพ นอกุสโล ธมฺโม อุปฺปชฺชติ เหตุปจฺจยา.  เหตุยา เทฺว.}

\prose{นภาวนาย ปหาตพฺพํ อพฺยากตํ ธมฺมํ ปจฺจยา นนภาวนาย ปหาตพฺโพ นอพฺยากโต ธมฺโม อุปฺปชฺชติ เหตุปจฺจยา.  เหตุยา เทฺว.}

\vspace{\tipitakaparskip}
\csromancenter{สวิจารทุก 1. กุสลตฺติก 85. ทสฺสเนน\-ปหาตพฺพ\-เหตุกทุก 1. กุสลตฺติก}
\proseitem{64}{\csromanfolio{233}นทสฺสเนน ปหาตพฺพ\-เหตุกํ กุสลํ ธมฺมํ ปฏิจฺจ นทสฺสเนน ปหาตพฺพ\-เหตุโก นกุสโล ธมฺโม อุปฺปชฺชติ เหตุปจฺจยา.  เหตุยา เอกํ.}

\prose{ทสฺสเนน ปหาตพฺพ\-เหตุกํ อกุสลํ ธมฺมํ ปฏิจฺจ นทสฺสเนน ปหาตพฺพ\-เหตุโก นอกุสโล ธมฺโม อุปฺปชฺชติ เหตุปจฺจยา.  เหตุยา ตีณิ.}

\prose{นทสฺสเนน ปหาตพฺพ\-เหตุกํ อพฺยากตํ ธมฺมํ ปจฺจยา นนทสฺสเนน ปหาตพฺพ\-เหตุโก นอพฺยากโต ธมฺโม อุปฺปชฺชติ เหตุปจฺจยา.  เหตุยา เทฺว. (อพฺยากตวาเร สพฺพตฺถ ปจฺจยวเสน คเณตพฺพํ.)}

\vspace{\tipitakaparskip}
\csromancenter{(อพฺยากตวาเร สพฺพตฺถ ปจฺจยวเสน คเณตพฺพํ.)}
\csromanmarkh{86. ภาวนาย\-ปหาตพฺพ\-เหตุกทุก}\csromanmarkhii{1. กุสลตฺติก}\csromanheadernomark{2}{86. ภาวนาย\-ปหาตพฺพ\-เหตุกทุก 1. กุสลตฺติก}
\proseitem{65}{นภาวนาย ปหาตพฺพ\-เหตุกํ กุสลํ ธมฺมํ ปฏิจฺจ นภาวนาย ปหาตพฺพ\-เหตุโก นกุสโล ธมฺโม อุปฺปชฺชติ เหตุปจฺจยา.  เหตุยา เอกํ.}

\prose{ภาวนาย ปหาตพฺพ\-เหตุกํ อกุสลํ ธมฺมํ ปฏิจฺจ นภาวนาย ปหาตพฺพ\-เหตุโก นอกุสโล ธมฺโม อุปฺปชฺชติ เหตุปจฺจยา.  เหตุยา ตีณิ. (อพฺยากตมูลํ เทฺว.)}

\csromanmarkh{87. สวิตกฺกทุก}\csromanmarkhii{1. กุสลตฺติก}\csromanheadernomark{2}{87. \textbf{สวิตกฺกทุก 1. กุสลตฺติก}}
\proseitem{66}{สวิตกฺกํ กุสลํ ธมฺมํ ปฏิจฺจ นสวิตกฺโก นกุสโล ธมฺโม อุปฺปชฺชติ เหตุปจฺจยา.  เหตุยา ตีณิ.}

\prose{สวิตกฺกํ อกุสลํ ธมฺมํ ปฏิจฺจ นสวิตกฺโก นอกุสโล ธมฺโม อุปฺปชฺชติ เหตุปจฺจยา.  เหตุยา ตีณิ. (อพฺยากตมูลํ ตีณิเยว.)}

\csromanmarkh{88. สวิจารทุก}\csromanmarkhii{1. กุสลตฺติก}\csromanheadernomark{2}{88. สวิจารทุก 1. กุสลตฺติก}
\proseitem{67}{สวิจารํ กุสลํ ธมฺมํ ปฏิจฺจ นสวิจาโร นกุสโล ธมฺโม อุปฺปชฺชติ เหตุปจฺจยา.  เหตุยา ตีณิ.}

\prose{สวิจารํ อกุสลํ ธมฺมํ ปฏิจฺจ นสวิจาโร นอกุสโล ธมฺโม อุปฺปชฺชติ เหตุปจฺจยา.  เหตุยา ตีณิ. (อพฺยากตมูลํ ตีณิเยว.)}

\csromanmarkh{ธมฺมา\-นุโลม\-ปจฺจนีย ทุกติก\-ปฏฺฐาน\-ปาฬิ}\csromanmarkhii{89. สปฺปีติกทุก 1. กุสลตฺติก}\csromanheadernomark{2}{ธมฺมา\-นุโลม\-ปจฺจนีย ทุกติก\-ปฏฺฐาน\-ปาฬิ 89. สปฺปีติกทุก 1. กุสลตฺติก}
\proseitem{68}{\csromanfolio{234}สปฺปีติกํ กุสลํ ธมฺมํ ปฏิจฺจ นสปฺปีติโก นกุสโล ธมฺโม อุปฺปชฺชติ เหตุปจฺจยา.  อปฺปีติกํ กุสลํ ธมฺมํ ปฏิจฺจ นสปฺปีติโก นกุสโล ธมฺโม อุปฺปชฺชติ เหตุปจฺจยา.  สปฺปีติกํ กุสลญฺจ อปฺปีติกํ กุสลญฺจ ธมฺมํ ปฏิจฺจ นสปฺปีติโก นกุสโล ธมฺโม อุปฺปชฺชติ เหตุปจฺจยา.  เหตุยา ตีณิ.}

\prose{สปฺปีติกํ อกุสลํ ธมฺมํ ปฏิจฺจ นสปฺปีติโก นอกุสโล ธมฺโม อุปฺปชฺชติ เหตุปจฺจยา.  อปฺปีติกํ อกุสลํ ธมฺมํ ปฏิจฺจ นสปฺปีติโก นอกุสโล ธมฺโม อุปฺปชฺชติ เหตุปจฺจยา.  สปฺปีติกํ อกุสลญฺจ อปฺปีติกํ อกุสลญฺจ ธมฺมํ ปฏิจฺจ นสปฺปีติโก อกุสโล ธมฺโม อุปฺปชฺชติ เหตุปจฺจยา.  เหตุยา ตีณิ. (อพฺยากตมูลํ ตีณิเยว.)}

\csromanmarkh{90. ปีติสหคตทุก}\csromanmarkhii{1. กุสลตฺติก}\csromanheadernomark{2}{90. \textbf{ปีติสหคตทุก 1. กุสลตฺติก}}
\proseitem{69}{ปีติสหคตํ กุสลํ ธมฺมํ ปฏิจฺจ นปีติสหคโต นกุสโล ธมฺโม อุปฺปชฺชติ เหตุปจฺจยา.  เหตุยา ตีณิ.}

\prose{ปีติสหคตํ อกุสลํ ธมฺมํ ปฏิจฺจ นปีติสหคโต นอกุสโล ธมฺโม อุปฺปชฺชติ เหตุปจฺจยา.  เหตุยา ตีณิ. (อพฺยากตมูลํ ตีณิเยว.)}

\csromanmarkh{91. สุขสหคตทุก}\csromanmarkhii{1. กุสลตฺติก}\csromanheadernomark{2}{91. \textbf{สุขสหคตทุก 1. กุสลตฺติก}}
\proseitem{70}{สุขสหคตํ กุสลํ ธมฺมํ ปฏิจฺจ นสุขสหคโต นกุสโล ธมฺโม อุปฺปชฺชติ เหตุปจฺจยา.  นสุขสหคตํ กุสลํ ธมฺมํ ปฏิจฺจ นสุขสหคโต นกุสโล ธมฺโม อุปฺปชฺชติ เหตุปจฺจยา.  สุขสหคตํ กุสลญฺจ นสุขสหคตํ กุสลญฺจ ธมฺมํ ปฏิจฺจ นสุขสหคโต นกุสโล ธมฺโม อุปฺปชฺชติ เหตุปจฺจยา.  เหตุยา ตีณิ.}

\prose{สุขสหคตํ อกุสลํ ธมฺมํ ปฏิจฺจ นสุขสหคโต นอกุสโล ธมฺโม อุปฺปชฺชติ เหตุปจฺจยา.  นสุขสหคตํ อกุสลํ ธมฺมํ ปฏิจฺจ นสุขสหคโต นอกุสโล ธมฺโม อุปฺปชฺชติ เหตุปจฺจยา.  สุขสหคตํ อกุสลญฺจ นสุขสหคตํ อกุสลญฺจ ธมฺมํ ปฏิจฺจ นสุขสหคโต นอกุสโล ธมฺโม อุปฺปชฺชติ เหตุปจฺจยา.  เหตุยา ตีณิ. (อพฺยากตมูลํ ตีณิเยว.)}

\vspace{\tipitakaparskip}
\csromancenter{อรูปาวจรทุก 1. กุสลตฺติก 92. \textbf{ อุเปกฺขาสหคตทุก 1.} กุสลตฺติก}
\proseitem{71}{\csromanfolio{235}อุเปกฺขา\-สหคตํ กุสลํ ธมฺมํ ปฏิจฺจ นอุเปกฺขาสหคโต นกุสโล ธมฺโม อุปฺปชฺชติ เหตุปจฺจยา.  นอุเปกฺขา\-สหคตํ กุสลํ ธมฺมํ ปฏิจฺจ นอุเปกฺขาสหคโต นกุสโล ธมฺโม อุปฺปชฺชติ เหตุปจฺจยา.  อุเปกฺขา\-สหคตํ กุสลญฺจ นอุเปกฺขา\-สหคตํ กุสลญฺจ ธมฺมํ ปฏิจฺจ นอุเปกฺขาสหคโต นกุสโล ธมฺโม อุปฺปชฺชติ เหตุปจฺจยา.  เหตุยา ตีณิ.}

\prose{อุเปกฺขา\-สหคตํ อกุสลํ ธมฺมํ ปฏิจฺจ นอุเปกฺขาสหคโต นอกุสโล ธมฺโม อุปฺปชฺชติ เหตุปจฺจยา.  นอุเปกฺขา\-สหคตํ อกุสลํ ธมฺมํ ปฏิจฺจ นอุเปกฺขาสหคโต นอกุสโล ธมฺโม อุปฺปชฺชติ เหตุปจฺจยา.  อุเปกฺขา\-สหคตํ อกุสลญฺจ นอุเปกฺขา\-สหคตํ อกุสลญฺจ ธมฺมํ ปฏิจฺจ นอุเปกฺขาสหคโต นอกุสโล ธมฺโม อุปฺปชฺชติ เหตุปจฺจยา.  เหตุยา ตีณิ. (อพฺยากตมูลํ ตีณิเยว.)}

\csromanmarkh{93. กามาวจรทุก}\csromanmarkhii{1. กุสลตฺติก}\csromanheadernomark{2}{93. \textbf{กามาวจรทุก 1. กุสลตฺติก}}
\proseitem{72}{กามาวจรํ กุสลํ ธมฺมํ ปฏิจฺจ นนกามาวจโร นกุสโล ธมฺโม อุปฺปชฺชติ เหตุปจฺจยา.  นกามาวจรํ กุสลํ ธมฺมํ ปฏิจฺจ นนกามาวจโร นกุสโล ธมฺโม อุปฺปชฺชติ เหตุปจฺจยา.  เหตุยา เทฺว.}

\prose{กามาวจรํ อกุสลํ ธมฺมํ ปฏิจฺจ นนกามาวจโร นอกุสโล ธมฺโม อุปฺปชฺชติ เหตุปจฺจยา.  เหตุยา เอกํ. (อพฺยากตมูลํ เทฺว.)}

\csromanmarkh{94. รูปาวจรทุก}\csromanmarkhii{1. กุสลตฺติก}\csromanheadernomark{2}{94. \textbf{รูปาวจรทุก 1. กุสลตฺติก}}
\proseitem{73}{รูปาวจรํ กุสลํ ธมฺมํ ปฏิจฺจ นรูปาวจโร นกุสโล ธมฺโม อุปฺปชฺชติ เหตุปจฺจยา.  นรูปาวจรํ กุสลํ ธมฺมํ ปฏิจฺจ นรูปาวจโร นกุสโล ธมฺโม อุปฺปชฺชติ เหตุปจฺจยา.  เหตุยา เทฺว.}

\prose{นรูปาวจรํ อกุสลํ ธมฺมํ ปฏิจฺจ นรูปาวจโร นอกุสโล ธมฺโม อุปฺปชฺชติ เหตุปจฺจยา.  เหตุยา เอกํ. (อพฺยากตมูลํ เทฺว.)}

\csromanmarkh{95. อรูปาวจรทุก}\csromanmarkhii{1. กุสลตฺติก}\csromanheadernomark{2}{95. อรูปาวจรทุก 1. กุสลตฺติก}
\proseitem{74}{อรูปาวจรํ กุสลํ ธมฺมํ ปฏิจฺจ นอรูปาวจโร นกุสโล ธมฺโม อุปฺปชฺชติ เหตุปจฺจยา.  นอรูปาวจรํ กุสลํ ธมฺมํ ปฏิจฺจ นอรูปาวจโร นกุสโล ธมฺโม อุปฺปชฺชติ เหตุปจฺจยา.  เหตุยา เทฺว.}

\csromanheader{2}{ธมฺมา\-นุโลม\-ปจฺจนีย ทุกติก\-ปฏฺฐาน\-ปาฬิ}
\prose{\csromanfolio{236}นอรูปาวจรํ อกุสลํ ธมฺมํ ปฏิจฺจ นอรูปาวจโร นอกุสโล ธมฺโม อุปฺปชฺชติ เหตุปจฺจยา. เหตุยา เอกํ. (อพฺยากตมูลํ เทฺว.)}

\csromanmarkh{96. ปริยาปนฺนทุก}\csromanmarkhii{1. กุสลตฺติก}\csromanheadernomark{2}{96. \textbf{ปริยาปนฺนทุก 1. กุสลตฺติก}}
\proseitem{75}{ปริยาปนฺนํ กุสลํ ธมฺมํ ปฏิจฺจ นอปริยาปนฺโน นกุสโล ธมฺโม อุปฺปชฺชติ เหตุปจฺจยา.  อปริยาปนฺนํ กุสลํ ธมฺมํ ปฏิจฺจ นอปริยาปนฺโน นกุสโล ธมฺโม อุปฺปชฺชติ เหตุปจฺจยา.  เหตุยา เทฺว.}

\prose{ปริยาปนฺนํ อกุสลํ ธมฺมํ ปฏิจฺจ นอปริยาปนฺโน นอกุสโล ธมฺโม อุปฺปชฺชติ เหตุปจฺจยา.  เหตุยา เอกํ. (อพฺยากตมูลํ เทฺว.)}

\csromanmarkh{97. นิยฺยานิกทุก}\csromanmarkhii{1. กุสลตฺติก}\csromanheadernomark{2}{97. \textbf{นิยฺยานิกทุก 1. กุสลตฺติก}}
\proseitem{76}{นิยฺยานิกํ กุสลํ ธมฺมํ ปฏิจฺจ นนิยฺยานิโก นกุสโล ธมฺโม อุปฺปชฺชติ เหตุปจฺจยา.  อนิยฺยานิกํ กุสลํ ธมฺมํ ปฏิจฺจ นนิยฺยานิโก นกุสโล ธมฺโม อุปฺปชฺชติ เหตุปจฺจยา.  เหตุยา เทฺว.}

\prose{อนิยฺยานิกํ อกุสลํ ธมฺมํ ปฏิจฺจ นนิยฺยานิโก นอกุสโล ธมฺโม อุปฺปชฺชติ เหตุปจฺจยา.  เหตุยา เอกํ. (อพฺยากตมูลํ เทฺว.)}

\csromanmarkh{98. นิยตทุก}\csromanmarkhii{1. กุสลตฺติก}\csromanheadernomark{2}{98. \textbf{นิยตทุก 1. กุสลตฺติก}}
\proseitem{77}{นิยตํ กุสลํ ธมฺมํ ปฏิจฺจ นนิยโต นกุสโล ธมฺโม อุปฺปชฺชติ เหตุปจฺจยา.  อนิยตํ กุสลํ ธมฺมํ ปฏิจฺจ นนิยโต นกุสโล ธมฺโม อุปฺปชฺชติ เหตุปจฺจยา.  เหตุยา เทฺว.}

\prose{นิยตํ อกุสลํ ธมฺมํ ปฏิจฺจ นนิยโต นอกุสโล ธมฺโม อุปฺปชฺชติ เหตุปจฺจยา.  อนิยตํ อกุสลํ ธมฺมํ ปฏิจฺจ นนิยโต นอกุสโล ธมฺโม อุปฺปชฺชติ เหตุปจฺจยา.  เหตุยา เทฺว. (อพฺยากตมูเล เทฺว.)}

\prose{99. สอุตฺตรทุก 1. กุสลตฺติก}

\proseitem{78}{สอุตฺตรํ กุสลํ ธมฺมํ ปฏิจฺจ นอนุตฺตโร นกุสโล ธมฺโม อุปฺปชฺชติ เหตุปจฺจยา.  อนุตฺตรํ กุสลํ ธมฺมํ ปฏิจฺจ นอนุตฺตโร นกุสโล ธมฺโม อุปฺปชฺชติ เหตุปจฺจยา.  เหตุยา เทฺว.}

\prose{สอุตฺตรํ อกุสลํ ธมฺมํ ปฏิจฺจ นอนุตฺตโร นอกุสโล ธมฺโม อุปฺปชฺชติ เหตุปจฺจยา.  เหตุยา เอกํ. (อพฺยากตมูเล เทฺว.)}

\csromanmarkh{100. สรณทุก}\csromanmarkhii{2. เวทนาตฺติก 100. สรณทุก 1. กุสลตฺติก}\csromanheadernomark{2}{100. สรณทุก 2. เวทนาตฺติก \textbf{100. สรณทุก 1. กุสลตฺติก}}
\proseitem{79}{\csromanfolio{237}อรณํ กุสลํ ธมฺมํ ปฏิจฺจ นสรโณ นกุสโล ธมฺโม อุปฺปชฺชติ เหตุปจฺจยา.  เหตุยา เอกํ.}

\prose{สรณํ อกุสลํ ธมฺมํ ปฏิจฺจ นสรโณ นอกุสโล ธมฺโม อุปฺปชฺชติ เหตุปจฺจยา.  เหตุยา เอกํ.}

\prose{อรณํ อพฺยากตํ ธมฺมํ ปจฺจยา นอรโณ นอพฺยากโต ธมฺโม อุปฺปชฺชติ เหตุปจฺจยา.  อรณํ อพฺยากตํ ธมฺมํ ปจฺจยา นสรโณ นอพฺยากโต ธมฺโม อุปฺปชฺชติ เหตุปจฺจยา.}

\prose{เหตุยา เทฺว.}

\vspace{\tipitakaparskip}
\csromancenter{(สหชาตวารมฺปิ ฯเปฯ สมฺปยุตฺตวารมฺปิ ปฏิจฺจ\-วาร\-สทิสํ.)}
\proseitem{80}{อรโณ อพฺยากโต ธมฺโม นอรณสฺส นอพฺยากตสฺส ธมฺมสฺส อารมฺมณ\-ปจฺจเยน ปจฺจโย.  อรโณ อพฺยากโต ธมฺโม นสรณสฺส นอพฺยากตสฺส ธมฺมสฺส อารมฺมณ\-ปจฺจเยน ปจฺจโย. (2)}

\prose{อารมฺมเณ เทฺว ฯเปฯ อวิคเต เทฺว.}

\csromanmarkh{100. สรณทุก}\csromanmarkhii{2. เวทนาตฺติก}\csromanheadernomark{2}{100. สรณทุก 2. เวทนาตฺติก}
\prosecont{สรณํ สุขาย เวทนาย สมฺปยุตฺตํ ธมฺมํ ปฏิจฺจ นสรโณ นสุขาย เวทนาย สมฺปยุตฺโต ธมฺโม อุปฺปชฺชติ เหตุปจฺจยา.  สรณํ สุขาย เวทนาย สมฺปยุตฺตํ ธมฺมํ ปฏิจฺจ นอรโณ นสุขาย เวทนาย สมฺปยุตฺโต ธมฺโม อุปฺปชฺชติ เหตุปจฺจยา.  สรณํ สุขาย เวทนาย สมฺปยุตฺตํ ธมฺมํ ปฏิจฺจ นสรโณ นสุขาย เวทนาย สมฺปยุตฺโต จ นอรโณ นสุขาย เวทนาย สมฺปยุตฺโต จ ธมฺมา อุปฺปชฺชนฺติ เหตุปจฺจยา. (3)}

\proseitem{81}{อรณํ สุขาย เวทนาย สมฺปยุตฺตํ ธมฺมํ ปฏิจฺจ นสรโณ นสุขาย เวทนาย สมฺปยุตฺโต ธมฺโม อุปฺปชฺชติ เหตุปจฺจยา. (1)}

\prose{เหตุยา จตฺตาริ.}

\proseitem{82}{สรณํ ทุกฺขาย เวทนาย สมฺปยุตฺตํ ธมฺมํ ปฏิจฺจ นสรโณ นทุกฺขาย เวทนาย สมฺปยุตฺโต ธมฺโม อุปฺปชฺชติ เหตุปจฺจยา.}

\prose{เหตุยา ตีณิ.}

\csromanheader{2}{ธมฺมา\-นุโลม\-ปจฺจนีย ทุกติก\-ปฏฺฐาน\-ปาฬิ}
\prose{\csromanfolio{238}สรณํ อทุกฺขม\-สุขาย เวทนาย สมฺปยุตฺตํ ธมฺมํ ปฏิจฺจ นสรโณ นอทุกฺขม\-สุขาย เวทนาย สมฺปยุตฺโต ธมฺโม อุปฺปชฺชติ เหตุปจฺจยา.}

\prose{เหตุยา จตฺตาริ.}

\csromanmarkh{100. สรณทุก}\csromanmarkhii{3. วิปากตฺติก}\csromanheadernomark{2}{100. \textbf{สรณทุก 3. วิปากตฺติก}}
\proseitem{83}{อรณํ วิปากํ ธมฺมํ ปฏิจฺจ นสรโณ นวิปาโก ธมฺโม อุปฺปชฺชติ เหตุปจฺจยา.  เหตุยา เอกํ.}

\prose{สรณํ วิปาก\-ธมฺม\-ธมฺมํ ปฏิจฺจ นสรโณ นวิปาก\-ธมฺม\-ธมฺโม อุปฺปชฺชติ เหตุปจฺจยา.  อรณํ วิปาก\-ธมฺม\-ธมฺมํ ปฏิจฺจ นสรโณ นวิปาก\-ธมฺม\-ธมฺโม อุปฺปชฺชติ เหตุปจฺจยา.  เหตุยา เทฺว.}

\prose{อรณํ เนว\-วิปาก\-น\-วิปากธมฺม\-ธมฺมํ ปฏิจฺจ นสรโณ นเน\-ว\-วิปาก\-น\-วิปาก\-ธมฺม\-ธมฺโม อุปฺปชฺชติ เหตุปจฺจยา.  เหตุยา เอกํ.}

\csromanmarkh{100. สรณทุก}\csromanmarkhii{4. อุปาทินฺนตฺติก}\csromanheadernomark{2}{100. \textbf{สรณทุก 4. อุปาทินฺนตฺติก}}
\proseitem{84}{อรณํ อุปาทินฺนุ\-ปาทานิยํ ธมฺมํ ปฏิจฺจ นสรโณ นอุปาทินฺนุปาทานิโย ธมฺโม อุปฺปชฺชติ เหตุปจฺจยา.  เหตุยา เอกํ.}

\prose{อรณํ อนุปาทินฺน\-อนุปาทานิยํ ธมฺมํ ปฏิจฺจ นสรโณ นอนุปาทินฺน\-อนุปาทานิโย ธมฺโม อุปฺปชฺชติ เหตุปจฺจยา.  เหตุยา เอกํ.}

\csromanmarkh{100. สรณทุก}\csromanmarkhii{5. สํกิลิฏฺฐ\-ตฺติก}\csromanheadernomark{2}{100. สรณทุก 5. สํกิลิฏฺฐ\-ตฺติก}
\proseitem{85}{สรณํ สํกิลิฏฺฐ\-สํกิเลสิกํ ธมฺมํ ปฏิจฺจ นสรโณ นสํกิลิฏฺฐ\-สํกิเลสิโก ธมฺโม อุปฺปชฺชติ เหตุปจฺจยา.  เหตุยา เอกํ.}

\prose{อรณํ อสํกิลิฏฺฐ\-สํกิเลสิกํ ธมฺมํ ปจฺจยา นอรโณ นอสํกิลิฏฺฐ\-สํกิเลสิโก ธมฺโม อุปฺปชฺชติ เหตุปจฺจยา.  อรณํ อสํกิลิฏฺฐ\-สํกิเลสิกํ ธมฺมํ ปจฺจยา นสรโณ นอสํกิลิฏฺฐ\-สํกิเลสิโก ธมฺโม อุปฺปชฺชติ เหตุปจฺจยา.  เหตุยา เทฺว.}

\prose{อรณํ อสํกิลิฏฺฐ\-อสํกิเลสิกํ ธมฺมํ ปฏิจฺจ นสรโณ นอสํกิลิฏฺฐ\-อสํกิเลสิโก ธมฺโม อุปฺปชฺชติ เหตุปจฺจยา.  เหตุยา เอกํ.}

\csromanmarkh{100. สรณทุก}\csromanmarkhii{8. ทสฺสนตฺติก 100. สรณทุก 6. วิตกฺกตฺติก}\csromanheadernomark{2}{100. สรณทุก 8. ทสฺสนตฺติก \textbf{100. สรณทุก 6. วิตกฺกตฺติก}}
\proseitem{86}{\csromanfolio{239}สรณํ สวิตกฺก\-สวิจารํ ธมฺมํ ปฏิจฺจ นสรโณ นสวิตกฺก\-สวิจาโร ธมฺโม อุปฺปชฺชติ เหตุปจฺจยา. ตีณิ.}

\prose{อรณํ สวิตกฺก\-สวิจารํ ธมฺมํ ปฏิจฺจ นอรโณ นสวิตกฺก\-สวิจาโร ธมฺโม อุปฺปชฺชติ เหตุปจฺจยา. เอกํ.}

\proseitem{87}{สรณํ อวิตกฺก\-วิจารมตฺตํ ธมฺมํ ปฏิจฺจ นสรโณ นอวิตกฺก\-วิจารมตฺโต ธมฺโม อุปฺปชฺชติ เหตุปจฺจยา.  เหตุยา จตฺตาริ.}

\prose{อรณํ อวิตกฺก\-อวิจารํ ธมฺมํ ปฏิจฺจ นสรโณ นอวิตกฺก\-อวิจาโร ธมฺโม อุปฺปชฺชติ เหตุปจฺจยา.  เหตุยา เอกํ.}

\csromanmarkh{100. สรณทุก}\csromanmarkhii{7. ปีติตฺติก}\csromanheadernomark{2}{100. \textbf{สรณทุก 7. ปีติตฺติก}}
\proseitem{88}{สรณํ ปีติสหคตํ ธมฺมํ ปฏิจฺจ นสรโณ นปีติสหคโต ธมฺโม อุปฺปชฺชติ เหตุปจฺจยา. ตีณิ.}

\prose{อรณํ ปีติสหคตํ ธมฺมํ ปฏิจฺจ นสรโณ นปีติสหคโต ธมฺโม อุปฺปชฺชติ เหตุปจฺจยา. เอกํ.}

\prose{สรณํ สุขสหคตํ ธมฺมํ ปฏิจฺจ นสรโณ นสุขสหคโต ธมฺโม อุปฺปชฺชติ เหตุปจฺจยา. ตีณิ.}

\prose{อรณํ สุขสหคตํ ธมฺมํ ปฏิจฺจ นสรโณ นสุขสหคโต ธมฺโม อุปฺปชฺชติ เหตุปจฺจยา. เอกํ.}

\prose{สรณํ อุเปกฺขา\-สหคตํ ธมฺมํ ปฏิจฺจ นสรโณ นอุเปกฺขาสหคโต ธมฺโม อุปฺปชฺชติ เหตุปจฺจยา.  เหตุยา จตฺตาริ.}

\csromanmarkh{100. สรณทุก}\csromanmarkhii{8. ทสฺสนตฺติก}\csromanheadernomark{2}{100. สรณทุก 8. ทสฺสนตฺติก}
\proseitem{89}{สรณํ ทสฺสเนน ปหาตพฺพํ ธมฺมํ ปฏิจฺจ นสรโณ นทสฺสเนน ปหาตพฺโพ ธมฺโม อุปฺปชฺชติ เหตุปจฺจยา.  เหตุยา เอกํ.}

\prose{สรณํ ภาวนาย ปหาตพฺพํ ธมฺมํ ปฏิจฺจ นสรโณ นภาวนาย ปหาตพฺโพ ธมฺโม อุปฺปชฺชติ เหตุปจฺจยา.  เหตุยา เอกํ.}

\csromanheader{2}{ธมฺมา\-นุโลม\-ปจฺจนีย ทุกติก\-ปฏฺฐาน\-ปาฬิ}
\prose{\csromanfolio{240}อรณํ เนวทสฺสเนน นภาวนาย ปหาตพฺพํ ธมฺมํ ปจฺจยา นอรโณ นเน\-ว\-ทสฺสเนน นภาวนาย ปหาตพฺโพ ธมฺโม อุปฺปชฺชติ เหตุปจฺจยา.  เหตุยา เอกํ.}

\csromanmarkh{100. สรณทุก}\csromanmarkhii{9. ทสฺสน\-เหตุ\-ตฺติก}\csromanheadernomark{2}{100. สรณทุก 9. ทสฺสน\-เหตุ\-ตฺติก}
\proseitem{90}{สรณํ ทสฺสเนน ปหาตพฺพ\-เหตุกํ ธมฺมํ ปจฺจยา นสรโณ นทสฺสเนน ปหาตพฺพ\-เหตุโก ธมฺโม อุปฺปชฺชติ เหตุปจฺจยา.  เหตุยา เอกํ.}

\prose{สรณํ ภาวนาย ปหาตพฺพ\-เหตุกํ ธมฺมํ ปฏิจฺจ นสรโณ นภาวนาย ปหาตพฺพ\-เหตุโก ธมฺโม อุปฺปชฺชติ เหตุปจฺจยา.  เหตุยา เอกํ.}

\prose{อรณํ เนวทสฺสเนน นภาวนาย ปหาตพฺพ\-เหตุกํ ธมฺมํ ปฏิจฺจ นอรโณ นเน\-ว\-ทสฺสเนน นภาวนาย ปหาตพฺพ\-เหตุโก ธมฺโม อุปฺปชฺชติ เหตุปจฺจยา.  เหตุยา เอกํ.}

\csromanmarkh{100. สรณทุก}\csromanmarkhii{10. อาจยคามิตฺติก}\csromanheadernomark{2}{100. \textbf{สรณทุก 1}0. อาจยคามิตฺติก}
\proseitem{91}{สรณํ อาจยคามิํ ธมฺมํ ปฏิจฺจ นสรโณ นอาจยคามี ธมฺโม อุปฺปชฺชติ เหตุปจฺจยา.  เหตุยา เทฺว.}

\prose{อรณํ อปจยคามิํ ธมฺมํ ปฏิจฺจ นสรโณ นอปจยคามี ธมฺโม อุปฺปชฺชติ เหตุปจฺจยา.  เหตุยา เอกํ.}

\prose{อรณํ เนวา\-จยคามิ\-นา\-ปจยคามิํ ธมฺมํ ปจฺจยา นอรโณ นเน\-วา\-จยคามิ\-นา\-ปจยคามี ธมฺโม อุปฺปชฺชติ เหตุปจฺจยา.}

\csromanmarkh{100. สรณทุก}\csromanmarkhii{11. เสกฺขตฺติก}\csromanheadernomark{2}{100. \textbf{สรณทุก 1}1. เสกฺขตฺติก}
\proseitem{92}{อรณํ เสกฺขํ ธมฺมํ ปฏิจฺจ นสรโณ นเสกฺโข ธมฺโม อุปฺปชฺชติ เหตุปจฺจยา.  เหตุยา เอกํ.}

\prose{อรณํ อเสกฺขํ ธมฺมํ ปฏิจฺจ นสรโณ นอเสกฺโข ธมฺโม อุปฺปชฺชติ เหตุปจฺจยา.  เหตุยา เอกํ.}

\csromanmarkh{100. สรณทุก}\csromanmarkhii{14. หีนตฺติก}\csromanheadernomark{2}{100. สรณทุก 14. หีนตฺติก}
\prose{\csromanfolio{241}อรณํ เนว\-เสกฺข\-นาเสกฺขํ ธมฺมํ ปจฺจยา นสรโณ นเน\-ว\-เสกฺข\-นา\-เสกฺโข ธมฺโม อุปฺปชฺชติ เหตุปจฺจยา.  เหตุยา เอกํ.}

\csromanmarkh{100. สรณทุก}\csromanmarkhii{12. ปริตฺตตฺติก}\csromanheadernomark{2}{100. \textbf{สรณทุก 1}2. ปริตฺตตฺติก}
\proseitem{93}{อรณํ ปริตฺตํ ธมฺมํ ปฏิจฺจ นสรโณ นปริตฺโต ธมฺโม อุปฺปชฺชติ เหตุปจฺจยา.  เหตุยา เอกํ.}

\prose{อรณํ มหคฺคตํ ธมฺมํ ปฏิจฺจ นสรโณ นมหคฺคโต ธมฺโม อุปฺปชฺชติ เหตุปจฺจยา.  เหตุยา เอกํ.}

\prose{อรณํ อปฺปมาณํ ธมฺมํ ปฏิจฺจ นสรโณ นอปฺปมาโณ ธมฺโม อุปฺปชฺชติ เหตุปจฺจยา.  เหตุยา เอกํ.}

\csromanmarkh{100. สรณทุก}\csromanmarkhii{13. ปริตฺตา\-รมฺมณ\-ตฺติก}\csromanheadernomark{2}{100. สรณทุก 13. ปริตฺตา\-รมฺมณ\-ตฺติก}
\proseitem{94}{สรณํ ปริตฺตา\-รมฺมณํ ธมฺมํ ปฏิจฺจ นสรโณ นปริตฺตา\-รมฺมโณ ธมฺโม อุปฺปชฺชติ เหตุปจฺจยา.  อรณํ ปริตฺตา\-รมฺมณํ ธมฺมํ ปฏิจฺจ นสรโณ นปริตฺตา\-รมฺมโณ ธมฺโม อุปฺปชฺชติ เหตุปจฺจยา.  เหตุยา เทฺว.}

\prose{สรณํ มหคฺคตา\-รมฺมณํ ธมฺมํ ปฏิจฺจ นสรโณ นมหคฺคตา\-รมฺมโณ ธมฺโม อุปฺปชฺชติ เหตุปจฺจยา.  อรณํ มหคฺคตา\-รมฺมณํ ธมฺมํ ปฏิจฺจ นสรโณ นมหคฺคตา\-รมฺมโณ ธมฺโม อุปฺปชฺชติ เหตุปจฺจยา.  เหตุยา เทฺว.}

\prose{อรณํ อปฺปมาณา\-รมฺมณํ ธมฺมํ ปฏิจฺจ นสรโณ นอปฺปมาณารมฺมโณ ธมฺโม อุปฺปชฺชติ เหตุปจฺจยา.  เหตุยา เอกํ.}

\csromanmarkh{100. สรณทุก}\csromanmarkhii{14. หีนตฺติก}\csromanheadernomark{2}{100. สรณทุก 14. หีนตฺติก}
\proseitem{95}{สรณํ หีนํ ธมฺมํ ปฏิจฺจ นสรโณ นหีโน ธมฺโม อุปฺปชฺชติ เหตุปจฺจยา.  เหตุยา เอกํ.}

\prose{อรณํ มชฺฌิมํ ธมฺมํ ปจฺจยา นอรโณ นมชฺฌิโม ธมฺโม อุปฺปชฺชติ เหตุปจฺจยา.  อรณํ มชฺฌิมํ ธมฺมํ ปจฺจยา นสรโณ นมชฺฌิโม ธมฺโม อุปฺปชฺชติ เหตุปจฺจยา.  เหตุยา เทฺว.}

\prose{อรณํ ปณีตํ ธมฺมํ ปฏิจฺจ นสรโณ นปณีโต ธมฺโม อุปฺปชฺชติ เหตุปจฺจยา.  เหตุยา เอกํ.}

\vspace{\tipitakaparskip}
\csromancenter{ทุกติก\-ปฏฺฐาน\-ปาฬิ 100. สรณทุก 15. มิจฺฉตฺตนิยต\-ตฺติก}
\proseitem{96}{\csromanfolio{242}สรณํ มิจฺฉตฺต\-นิยตํ ธมฺมํ ปฏิจฺจ นสรโณ นมิจฺฉตฺต\-นิยโต ธมฺโม อุปฺปชฺชติ เหตุปจฺจยา.  เหตุยา เอกํ.}

\prose{อรณํ สมฺมตฺต\-นิยตํ ธมฺมํ ปฏิจฺจ นสรโณ นสมฺมตฺตนิยโต ธมฺโม อุปฺปชฺชติ เหตุปจฺจยา.  เหตุยา เอกํ.}

\prose{อรณํ อนิยตํ ธมฺมํ ปจฺจยา นอรโณ นอนิยโต ธมฺโม อุปฺปชฺชติ เหตุปจฺจยา.  อรณํ อนิยตํ ธมฺมํ ปจฺจยา นสรโณ นอนิยโต ธมฺโม อุปฺปชฺชติ เหตุปจฺจยา.  เหตุยา เทฺว.}

\csromanmarkh{100. สรณทุก}\csromanmarkhii{16. มคฺคา\-รมฺมณ\-ตฺติก}\csromanheadernomark{2}{100. สรณทุก 16. มคฺคา\-รมฺมณ\-ตฺติก}
\proseitem{97}{อรณํ มคฺคารมฺมณํ ธมฺมํ ปฏิจฺจ นสรโณ นมคฺคารมฺมโณ ธมฺโม อุปฺปชฺชติ เหตุปจฺจยา.  เหตุยา เอกํ.}

\prose{อรณํ มคฺคเหตุกํ ธมฺมํ ปฏิจฺจ นสรโณ นมคฺคเหตุโก ธมฺโม อุปฺปชฺชติ เหตุปจฺจยา.  เหตุยา เอกํ.}

\prose{อรณํ มคฺคาธิปติํ ธมฺมํ ปฏิจฺจ นสรโณ นมคฺคาธิปติ ธมฺโม อุปฺปชฺชติ เหตุปจฺจยา.  เหตุยา เอกํ.}

\csromanmarkh{100. สรณทุก}\csromanmarkhii{17. อุปฺปนฺนตฺติก}\csromanheadernomark{2}{100. \textbf{สรณทุก 1}7. อุปฺปนฺนตฺติก}
\proseitem{98}{สรโณ อนุปฺปนฺโน ธมฺโม นสรณสฺส นอนุปฺปนฺนสฺส ธมฺมสฺส อารมฺมณ\-ปจฺจเยน ปจฺจโย.  อารมฺมเณ จตฺตาริ.}

\prose{อรโณ อุปฺปาที ธมฺโม นอรณสฺส นอุปฺปาทิสฺส ธมฺมสฺส อารมฺมณ\-ปจฺจเยน ปจฺจโย.  อรโณ อุปฺปาที ธมฺโม นสรณสฺส นอุปฺปาทิสฺส ธมฺมสฺส อารมฺมณ\-ปจฺจเยน ปจฺจโย.  อารมฺมเณ เทฺว.}

\csromanmarkh{100. สรณทุก}\csromanmarkhii{18. อตีตตฺติก}\csromanheadernomark{2}{100. \textbf{สรณทุก 1}8. อตีตตฺติก}
\proseitem{99}{สรโณ อตีโต ธมฺโม นสรณสฺส นอตีตสฺส ธมฺมสฺส อารมฺมณ\-ปจฺจเยน ปจฺจโย.  อารมฺมเณ จตฺตาริ. (อนาคโต อตีตสทิโส.)}

\vspace{\tipitakaparskip}
\csromancenter{สรณทุก 22. สนิทสฺสน\-ตฺติก 100. สรณทุก 19. อตีตา\-รมฺมณ\-ตฺติก}
\prose{\csromanfolio{243}สรณํ อตีตารมฺมณํ ธมฺมํ ปฏิจฺจ นสรโณ นอตีตารมฺมโณ ธมฺโม อุปฺปชฺชติ เหตุปจฺจยา.  อรณํ อตีตารมฺมณํ ธมฺมํ ปฏิจฺจ นสรโณ นอตีตารมฺมโณ ธมฺโม อุปฺปชฺชติ เหตุปจฺจยา.  เหตุยา เทฺว.}

\prose{สรณํ อนาคตา\-รมฺมณํ ธมฺมํ ปฏิจฺจ นสรโณ นอนาคตา\-รมฺมโณ ธมฺโม อุปฺปชฺชติ เหตุปจฺจยา.  อรณํ อนาคตา\-รมฺมณํ ธมฺมํ ปฏิจฺจ นสรโณ นอนาคตา\-รมฺมโณ ธมฺโม อุปฺปชฺชติ เหตุปจฺจยา.  เหตุยา เทฺว.}

\prose{สรณํ ปจฺจุปฺปนฺนา\-รมฺมณํ ธมฺมํ ปฏิจฺจ นสรโณ นปจฺจุปฺปนฺนา\-รมฺมโณ ธมฺโม อุปฺปชฺชติ เหตุปจฺจยา.  อรณํ ปจฺจุปฺปนฺนา\-รมฺมณํ ธมฺมํ ปฏิจฺจ นสรโณ นปจฺจุปฺปนฺนา\-รมฺมโณ ธมฺโม อุปฺปชฺชติ เหตุปจฺจยา.  เหตุยา เทฺว.}

\csromanmarkh{100. สรณทุก}\csromanmarkhii{21. อชฺฌตฺตา\-รมฺมณ\-ตฺติก}\csromanheadernomark{2}{100. สรณทุก 21. อชฺฌตฺตา\-รมฺมณ\-ตฺติก}
\proseitem{101}{สรณํ อชฺฌตฺตา\-รมฺมณํ ธมฺมํ ปฏิจฺจ นสรโณ นอชฺฌตฺตา\-รมฺมโณ ธมฺโม อุปฺปชฺชติ เหตุปจฺจยา.  อรณํ อชฺฌตฺตา\-รมฺมณํ ธมฺมํ ปฏิจฺจ นสรโณ นอชฺฌตฺตา\-รมฺมโณ ธมฺโม อุปฺปชฺชติ เหตุปจฺจยา.  เหตุยา เทฺว.}

\prose{สรณํ พหิทฺธา\-รมฺมณํ ธมฺมํ ปฏิจฺจ นสรโณ นพหิทฺธา\-รมฺมโณ ธมฺโม อุปฺปชฺชติ เหตุปจฺจยา.  อรณํ พหิทฺธา\-รมฺมณํ ธมฺมํ ปฏิจฺจ นสรโณ นพหิทฺธา\-รมฺมโณ ธมฺโม อุปฺปชฺชติ เหตุปจฺจยา.  เหตุยา เทฺว.}

\csromanmarkh{100. สรณทุก}\csromanmarkhii{22. สนิทสฺสน\-ตฺติก}\csromanheadernomark{2}{100. สรณทุก 22. สนิทสฺสน\-ตฺติก}
\proseitem{102}{อรณํ อนิทสฺสน\-สปฺปฏิฆํ ธมฺมํ ปฏิจฺจ นสรโณ นอนิทสฺสน\-สปฺปฏิโฆ ธมฺโม อุปฺปชฺชติ เหตุปจฺจยา.  เอกํ.}

\prose{สรณํ อนิทสฺสน\-อปฺปฏิฆํ ธมฺมํ ปฏิจฺจ นสรโณ นอนิทสฺสน\-อปฺปฏิโฆ ธมฺโม อุปฺปชฺชติ เหตุปจฺจยา.  อรณํ อนิทสฺสน\-อปฺปฏิฆํ ธมฺมํ ปฏิจฺจ นสรโณ นอนิทสฺสน\-อปฺปฏิโฆ ธมฺโม อุปฺปชฺชติ เหตุปจฺจยา.  สรณํ อนิทสฺสนอปฺปฏิฆญฺจ อรณํ อนิทสฺสนอปฺปฏิฆญฺจ ธมฺมํ ปฏิจฺจ นสรโณ นอนิทสฺสน\-อปฺปฏิโฆ ธมฺโม อุปฺปชฺชติ เหตุปจฺจยา. (3)}

\prose{เหตุยา ตีณิ ฯเปฯ อวิคเต ตีณิ.}

\csromanheader{2}{ธมฺมา\-นุโลม\-ปจฺจนีย ทุกติก\-ปฏฺฐาน\-ปาฬิ นเหตุ นอารมฺมณ}
\proseitem{103}{\csromanfolio{244}อรณํ อนิทสฺสน\-อปฺปฏิฆํ ธมฺมํ ปฏิจฺจ นสรโณ นอนิทสฺสน\-อปฺปฏิโฆ ธมฺโม อุปฺปชฺชติ นเหตุปจฺจยา. (1)}

\prose{สรณํ อนิทสฺสน\-อปฺปฏิฆํ ธมฺมํ ปฏิจฺจ นสรโณ นอนิทสฺสน\-อปฺปฏิโฆ ธมฺโม อุปฺปชฺชติ นอารมฺมณ\-ปจฺจยา.}

\prose{นเหตุยา เอกํ, นอารมฺมเณ ตีณิ ฯเปฯ โนวิคเต ตีณิ.}

\prose{(สหชาตวารมฺปิ ฯเปฯ สมฺปยุตฺตวารมฺปิ ปฏิจฺจ\-วาร\-สทิสํ.)}

\proseitem{104}{สรโณ อนิทสฺสน\-อปฺปฏิโฆ ธมฺโม นสรณสฺส นอนิทสฺสน\-อปฺปฏิฆสฺส ธมฺมสฺส เหตุปจฺจเยน ปจฺจโย. (1)}

\prose{อรโณ อนิทสฺสน\-อปฺปฏิโฆ ธมฺโม นสรณสฺส นอนิทสฺสน\-อปฺปฏิฆสฺส ธมฺมสฺส เหตุปจฺจเยน ปจฺจโย. (1)}

\prose{เหตุยา เทฺว, อธิปติยา เทฺว ฯเปฯ อวิคเต ตีณิ.}

\csromanheader{2}{ปจฺจนียุ\-ทฺธาร}
\proseitem{105}{สรโณ อนิทสฺสน\-อปฺปฏิโฆ ธมฺโม นสรณสฺส นอนิทสฺสน\-อปฺปฏิฆสฺส ธมฺมสฺส สหชาต\-ปจฺจเยน ปจฺจโย. ปจฺฉาชาต\-ปจฺจเยน ปจฺจโย. กมฺมปจฺจเยน ปจฺจโย. (สํขิตฺตํ.)}

\prose{นเหตุยา ตีณิ, นอารมฺมเณ ตีณิ ฯเปฯ โนอวิคเต ตีณิ.}

\prosewithcloser{(ยถา กุสลตฺติเก ปญฺหาวารํ, เอวํ วิตฺถาเรตพฺพํ.)}{ธมฺมา\-นุโลม\-ปจฺจนีเย ทุกติก\-ปฏฺฐานํ นิฏฺฐิตํ.}
\pitaka{อภิธมฺม\-ปิฏก}
\namakkaram{นโม ตสฺส ภควโต อรหโต สมฺมา\-สมฺพุทฺธสฺส.}
\csromanmarkh{1. กุสลตฺติก}\csromanmarkhii{1. เหตุทุก}\csromanheadernomark{2}{1. \textbf{กุสลตฺติก 1. เหตุทุก}}
\csromanheader{2}{1-7. ปฏิจฺจาทิวาร ปจฺจยจตุกฺกเหตุ}
\proseitem{1}{\csromanfolio{245}กุสลํ เหตุํ ธมฺมํ ปฏิจฺจ นกุสโล นเหตุ ธมฺโม อุปฺปชฺชติ เหตุปจฺจยา.  กุสลํ เหตุํ ธมฺมํ ปฏิจฺจ นอกุสโล นเหตุ ธมฺโม อุปฺปชฺชติ เหตุปจฺจยา.  กุสลํ เหตุํ ธมฺมํ ปฏิจฺจ นอพฺยากโต นเหตุ ธมฺโม อุปฺปชฺชติ เหตุปจฺจยา.  กุสลํ เหตุํ ธมฺมํ ปฏิจฺจ นอกุสโล นเหตุ จ นอพฺยากโต นเหตุ จ ธมฺมา อุปฺปชฺชนฺติ เหตุปจฺจยา.  กุสลํ เหตุํ ธมฺมํ ปฏิจฺจ นกุสโล นเหตุ จ นอกุสโล นเหตุ จ ธมฺมา อุปฺปชฺชนฺติ เหตุปจฺจยา. (5)}

\prose{อกุสลํ เหตุํ ธมฺมํ ปฏิจฺจ นอกุสโล นเหตุ ธมฺโม อุปฺปชฺชติ เหตุปจฺจยา.  อกุสลํ เหตุํ ธมฺมํ ปฏิจฺจ นกุสโล นเหตุ ธมฺโม อุปฺปชฺชติ เหตุปจฺจยา.  อกุสลํ เหตุํ ธมฺมํ ปฏิจฺจ นอพฺยากโต นเหตุ ธมฺโม อุปฺปชฺชติ เหตุปจฺจยา.  อกุสลํ เหตุํ ธมฺมํ ปฏิจฺจ นกุสโล นเหตุ จ นอพฺยากโต นเหตุ จ ธมฺมา อุปฺปชฺชนฺติ เหตุปจฺจยา.  อกุสลํ เหตุํ ธมฺมํ ปฏิจฺจ นกุสโล นเหตุ จ นอกุสโล นเหตุ จ ธมฺมา อุปฺปชฺชนฺติ เหตุปจฺจยา. (5)}

\csromanheader{2}{ธมฺมา\-นุโลม\-ปจฺจนีย ติกทุก\-ปฏฺฐาน\-ปาฬิ}
\prose{\csromanfolio{246}อพฺยากตํ เหตุํ ธมฺมํ ปฏิจฺจ นกุสโล นเหตุ ธมฺโม อุปฺปชฺชติ เหตุปจฺจยา.  อพฺยากตํ เหตุํ ธมฺมํ ปฏิจฺจ นอกุสโล นเหตุ ธมฺโม อุปฺปชฺชติ เหตุปจฺจยา.  อพฺยากตํ เหตุํ ธมฺมํ ปฏิจฺจ นกุสโล นเหตุ จ นอกุสโล นเหตุ จ ธมฺมา อุปฺปชฺชนฺติ เหตุปจฺจยา. (3) (สํขิตฺตํ.)}

\prose{เหตุยา เตรส, อารมฺมเณ นว ฯเปฯ อวิคเต เตรส.}

\prose{ปจฺจนียนอารมฺมณ}

\proseitem{2}{กุสลํ เหตุํ ธมฺมํ ปฏิจฺจ นกุสโล นเหตุ ธมฺโม อุปฺปชฺชติ นอารมฺมณ\-ปจฺจยา. (สํขิตฺตํ.)}

\prose{นอารมฺมเณ นว, นอธิปติยา เตรส ฯเปฯ นวิปฺปยุตฺเต นว.}

\prose{(สหชาตวารมฺปิ ฯเปฯ สมฺปยุตฺตวารมฺปิ ปฏิจฺจ\-วาร\-สทิสํ วิตฺถาเรตพฺพํ.)}

\csromanheader{2}{\textbf{เหตุ อารมฺมณ}}
\proseitem{3}{กุสโล เหตุ ธมฺโม นกุสลสฺส นเหตุสฺส ธมฺมสฺส เหตุปจฺจเยน ปจฺจโย.  กุสโล เหตุ ธมฺโม นอกุสลสฺส นเหตุสฺส ธมฺมสฺส เหตุปจฺจเยน ปจฺจโย.  กุสโล เหตุ ธมฺโม นอพฺยากตสฺส นเหตุสฺส ธมฺมสฺส เหตุปจฺจเยน ปจฺจโย. กุสโล เหตุ ธมฺโม นอกุสลสฺส นเหตุสฺส จ นอพฺยากตสฺส นเหตุสฺส จ ธมฺมสฺส เหตุปจฺจเยน ปจฺจโย.  กุสโล เหตุ ธมฺโม นกุสลสฺส นเหตุสฺส จ นอกุสลสฺส นเหตุสฺส จ ธมฺมสฺส เหตุปจฺจเยน ปจฺจโย. (5)}

\prose{อกุสโล เหตุ ธมฺโม นอกุสลสฺส นเหตุสฺส ธมฺมสฺส เหตุปจฺจเยน ปจฺจโย. ปญฺจ.}

\prose{อพฺยากโต เหตุ ธมฺโม นกุสลสฺส นเหตุสฺส ธมฺมสฺส เหตุปจฺจเยน ปจฺจโย. ตีณิ.}

\prose{กุสโล เหตุ ธมฺโม นกุสลสฺส นเหตุสฺส ธมฺมสฺส อารมฺมณ\-ปจฺจเยน ปจฺจโย. (สํขิตฺตํ.)}

\proseitem{4}{เหตุยา เตรส, อารมฺมเณ อฏฺฐารส ฯเปฯ อวิคเต เตรส. (ปญฺหาวารํ วิตฺถาเรตพฺพํ.)}

\csromanmarkh{1. กุสลตฺติก}\csromanmarkhii{2. สเหตุกทุก}\csromanheadernomark{2}{1. กุสลตฺติก 2. สเหตุกทุก}
\proseitem{5}{\csromanfolio{247}กุสลํ นเหตุํ ธมฺมํ ปฏิจฺจ นอกุสโล นนเหตุ ธมฺโม อุปฺปชฺชติ เหตุปจฺจยา.  กุสลํ นเหตุํ ธมฺมํ ปฏิจฺจ นอพฺยากโต นนเหตุ ธมฺโม อุปฺปชฺชติ เหตุปจฺจยา. (สํขิตฺตํ.)}

\prose{เหตุยา นว, อารมฺมเณ นว ฯเปฯ อวิคเต นว. (สพฺพตฺถ นว.)}

\csromanmarkh{1. กุสลตฺติก}\csromanmarkhii{2. สเหตุกทุก}\csromanheadernomark{2}{1. กุสลตฺติก 2. สเหตุกทุก}
\csromanheader{2}{1-7. ปฏิจฺจาทิวาร ปจฺจยจตุกฺกเหตุ}
\proseitem{6}{กุสลํ สเหตุกํ ธมฺมํ ปฏิจฺจ นกุสโล นสเหตุโก ธมฺโม อุปฺปชฺชติ เหตุปจฺจยา.  กุสลํ สเหตุกํ ธมฺมํ ปฏิจฺจ นอกุสโล นสเหตุโก ธมฺโม อุปฺปชฺชติ เหตุปจฺจยา.  กุสลํ สเหตุกํ ธมฺมํ ปฏิจฺจ นกุสโล นสเหตุโก จ นอกุสโล นสเหตุโก จ ธมฺมา อุปฺปชฺชนฺติ เหตุปจฺจยา. ตีณิ.}

\prose{อกุสลํ สเหตุกํ ธมฺมํ ปฏิจฺจ นอกุสโล นสเหตุโก ธมฺโม อุปฺปชฺชติ เหตุปจฺจยา.  อกุสลํ สเหตุกํ ธมฺมํ ปฏิจฺจ นกุสโล นสเหตุโก ธมฺโม อุปฺปชฺชติ เหตุปจฺจยา.  อกุสลํ สเหตุกํ ธมฺมํ ปฏิจฺจ นกุสโล นสเหตุโก จ นอกุสโล นสเหตุโก จ ธมฺมา อุปฺปชฺชนฺติ เหตุปจฺจยา. ตีณิ.}

\prose{อพฺยากตํ สเหตุกํ ธมฺมํ ปฏิจฺจ นกุสโล นสเหตุโก ธมฺโม อุปฺปชฺชติ เหตุปจฺจยา.  อพฺยากตํ สเหตุกํ ธมฺมํ ปฏิจฺจ นอกุสโล นสเหตุโก ธมฺโม อุปฺปชฺชติ เหตุปจฺจยา.  อพฺยากตํ สเหตุกํ ธมฺมํ ปฏิจฺจ นกุสโล นสเหตุโก จ นอกุสโล นสเหตุโก จ ธมฺมา อุปฺปชฺชนฺติ เหตุปจฺจยา. ตีณิ.}

\prose{เหตุยา นว, อารมฺมเณ ตีณิ ฯเปฯ อวิคเต เอกาทส.}

\vspace{\tipitakaparskip}
\csromancenter{(สหชาตวารมฺปิ ฯเปฯ สมฺปยุตฺตวารมฺปิ ปฏิจฺจ\-วาร\-สทิสํ.)}
\proseitem{7}{กุสโล สเหตุโก ธมฺโม นกุสลสฺส นสเหตุกสฺส ธมฺมสฺส เหตุปจฺจเยน ปจฺจโย.  กุสโล สเหตุโก ธมฺโม นอกุสลสฺส นสเหตุกสฺส ธมฺมสฺส เหตุปจฺจเยน ปจฺจโย.  กุสโล สเหตุโก ธมฺโม นกุสลสฺส นสเหตุกสฺส จ นอกุสลสฺส นสเหตุกสฺส จ ธมฺมสฺส เหตุปจฺจเยน ปจฺจโย. ตีณิ.}

\csromanheader{2}{ธมฺมา\-นุโลม\-ปจฺจนีย ติกทุก\-ปฏฺฐาน\-ปาฬิ}
\prose{\csromanfolio{248}อกุสโล สเหตุโก ธมฺโม นอกุสลสฺส นสเหตุกสฺส ธมฺมสฺส เหตุปจฺจเยน ปจฺจโย.  อกุสโล สเหตุโก ธมฺโม นกุสลสฺส นสเหตุกสฺส ธมฺมสฺส เหตุปจฺจเยน ปจฺจโย.  อกุสโล สเหตุโก ธมฺโม นกุสลสฺส นสเหตุกสฺส จ นอกุสลสฺส นสเหตุกสฺส จ ธมฺมสฺส เหตุปจฺจเยน ปจฺจโย. ตีณิ.}

\prose{อพฺยากโต สเหตุโก ธมฺโม นกุสลสฺส นสเหตุกสฺส ธมฺมสฺส เหตุปจฺจเยน ปจฺจโย.  อพฺยากโต สเหตุโก ธมฺโม นอกุสลสฺส นสเหตุกสฺส ธมฺมสฺส เหตุปจฺจเยน ปจฺจโย.  อพฺยากโต สเหตุโก ธมฺโม นกุสลสฺส นสเหตุกสฺส จ นอกุสลสฺส นสเหตุกสฺส จ ธมฺมสฺส เหตุปจฺจเยน ปจฺจโย. ตีณิ. (สํขิตฺตํ.)}

\proseitem{8}{เหตุยา นว, อารมฺมเณ ปนฺนรส ฯเปฯ อวิคเต เอกาทส.}

\proseitem{9}{อกุสลํ อเหตุกํ ธมฺมํ ปฏิจฺจ นกุสโล นอเหตุโก ธมฺโม อุปฺปชฺชติ เหตุปจฺจยา.  อกุสลํ อเหตุกํ ธมฺมํ ปฏิจฺจ นอพฺยากโต นอเหตุโก ธมฺโม อุปฺปชฺชติ เหตุปจฺจยา.  อกุสลํ อเหตุกํ ธมฺมํ ปฏิจฺจ นกุสโล นอเหตุโก จ นอพฺยากโต นอเหตุโก จ ธมฺมา อุปฺปชฺชนฺติ เหตุปจฺจยา. ตีณิ.}

\prose{อพฺยากตํ อเหตุกํ ธมฺมํ ปฏิจฺจ นกุสโล นอเหตุโก ธมฺโม อุปฺปชฺชติ เหตุปจฺจยา.  อพฺยากตํ อเหตุกํ ธมฺมํ ปฏิจฺจ นอกุสโล นอเหตุโก ธมฺโม อุปฺปชฺชติ เหตุปจฺจยา.  อพฺยากตํ อเหตุกํ ธมฺมํ ปฏิจฺจ นกุสโล นอเหตุโก จ นอกุสโล นอเหตุโก จ ธมฺมา อุปฺปชฺชนฺติ เหตุปจฺจยา. ตีณิ.}

\prose{เหตุยา ฉ, อารมฺมเณ ฉ ฯเปฯ อวิคเต ฉ.}

\csromanmarkh{1. กุสลตฺติก}\csromanmarkhii{3. เหตุ\-สมฺปยุตฺตทุก}\csromanheadernomark{2}{1. กุสลตฺติก 3. เหตุ\-สมฺปยุตฺตทุก}
\proseitem{10}{กุสลํ เหตุ\-สมฺปยุตฺตํ ธมฺมํ ปฏิจฺจ นกุสโล นเหตุ\-สมฺปยุตฺโต ธมฺโม อุปฺปชฺชติ เหตุปจฺจยา. (สเหตุก\-ทุก\-สทิสํ.)}

\csromanmarkh{1. กุสลตฺติก}\csromanmarkhii{4. เหตุ\-สเหตุกทุก}\csromanheadernomark{2}{1. กุสลตฺติก 4. เหตุ\-สเหตุกทุก}
\proseitem{11}{กุสลํ เหตุญฺเจว สเหตุกญฺจ ธมฺมํ ปฏิจฺจ นอกุสโล นเหตุ เจว นอเหตุโก จ ธมฺโม อุปฺปชฺชติ เหตุปจฺจยา.  กุสลํ}

\vspace{\tipitakaparskip}
\csromancenter{4. เหตุ\-สเหตุกทุก เหตุญฺเจว สเหตุกญฺจ ธมฺมํ ปฏิจฺจ นอพฺยากโต นเหตุ เจว นอเหตุโก จ ธมฺโม อุปฺปชฺชติ เหตุปจฺจยา.  กุสลํ เหตุญฺเจว สเหตุกญฺจ ธมฺมํ ปฏิจฺจ นอกุสโล นเหตุ เจว นอเหตุโก จ นอพฺยากโต นเหตุ เจว นอเหตุโก จ ธมฺมา อุปฺปชฺชนฺติ เหตุปจฺจยา. (3)}
\prose{\csromanfolio{249}อกุสลํ เหตุญฺเจว สเหตุญฺจ ธมฺมํ ปฏิจฺจ นกุสโล นเหตุ เจว นอเหตุโก จ ธมฺโม อุปฺปชฺชติ เหตุปจฺจยา. ตีณิ.}

\prose{อพฺยากตํ เหตุญฺเจว สเหตุกญฺจ ธมฺมํ ปฏิจฺจ นกุสโล นเหตุ เจว นอเหตุโก จ ธมฺโม อุปฺปชฺชติ เหตุปจฺจยา.  อพฺยากตํ เหตุญฺเจว สเหตุกญฺจ ธมฺมํ ปฏิจฺจ นอกุสโล นเหตุ เจว นอเหตุโก จ ธมฺโม อุปฺปชฺชติ เหตุปจฺจยา.  อพฺยากตํ เหตุญฺเจว สเหตุกญฺจ ธมฺมํ ปฏิจฺจ นกุสโล นเหตุ เจว นอเหตุโก จ นอกุสโล นเหตุ เจว นอเหตุโก จ ธมฺมา อุปฺปชฺชนฺติ เหตุปจฺจยา. (3) (สํขิตฺตํ.)}

\prose{เหตุยา นว.}

\proseitem{12}{กุสลํ สเหตุกญฺเจว น จ เหตุํ ธมฺมํ ปฏิจฺจ นอกุสโล นอเหตุโก เจว นนเหตุ จ ธมฺโม อุปฺปชฺชติ เหตุปจฺจยา.  กุสลํ สเหตุกญฺเจว น จ เหตุํ ธมฺมํ ปฏิจฺจ นอพฺยากโต นอเหตุโก เจว นนเหตุ จ ธมฺโม อุปฺปชฺชติ เหตุปจฺจยา.  กุสลํ สเหตุกญฺเจว น จ เหตุํ ธมฺมํ ปฏิจฺจ นอกุสโล นอเหตุโก เจว นนเหตุ จ นอพฺยากโต นอเหตุโก เจว นนเหตุ จ ธมฺมา อุปฺปชฺชนฺติ เหตุปจฺจยา. ตีณิ.}

\prose{อกุสลํ สเหตุกญฺเจว น จ เหตุํ ธมฺมํ ปฏิจฺจ นกุสโล นอเหตุโก เจว นนเหตุ จ ธมฺโม อุปฺปชฺชติ เหตุปจฺจยา.  อกุสลํ สเหตุกญฺเจว น จ เหตุํ ธมฺมํ ปฏิจฺจ นอพฺยากโต นอเหตุโก เจว นนเหตุ จ ธมฺโม อุปฺปชฺชติ เหตุปจฺจยา.  อกุสลํ สเหตุกญฺเจว น จ เหตุํ ธมฺมํ ปฏิจฺจ นอกุสโล นอเหตุโก เจว นนเหตุ จ นอพฺยากโต นอเหตุโก เจว นนเหตุ จ ธมฺมา อุปฺปชฺชนฺติ เหตุปจฺจยา. ตีณิ.}

\prose{อพฺยากตํ สเหตุกญฺเจว น จ เหตุํ ธมฺมํ ปฏิจฺจ นกุสโล นอเหตุโก เจว นนเหตุ จ ธมฺโม อุปฺปชฺชติ เหตุปจฺจยา.  }

\vspace{\tipitakaparskip}
\csromancenter{ติกทุก\-ปฏฺฐาน\-ปาฬิ อพฺยากตํ สเหตุกญฺเจว น จ เหตุํ ธมฺมํ ปฏิจฺจ นอกุสโล นอเหตุโก เจว นนเหตุ จ ธมฺโม อุปฺปชฺชติ เหตุปจฺจยา.  อพฺยากตํ สเหตุกญฺเจว น จ เหตุํ ธมฺมํ ปฏิจฺจ นกุสโล นอเหตุโก เจว นนเหตุ จ นอกุสโล นอเหตุโก เจว นนเหตุ จ ธมฺมา อุปฺปชฺชนฺติ เหตุปจฺจยา. ตีณิ.}
\csromanmarkh{1. กุสลตฺติก}\csromanmarkhii{5. เหตุ\-เหตุ\-สมฺปยุตฺตทุก}\csromanheadernomark{2}{1. กุสลตฺติก 5. เหตุ\-เหตุ\-สมฺปยุตฺตทุก}
\proseitem{13}{\csromanfolio{250}กุสลํ เหตุญฺเจว เหตุ\-สมฺปยุตฺตญฺจ ธมฺมํ ปฏิจฺจ นอกุสโล นเหตุ เจว นเหตุ\-วิปฺปยุตฺโต จ ธมฺโม อุปฺปชฺชติ เหตุปจฺจยา.  กุสลํ เหตุญฺเจว เหตุ\-สมฺปยุตฺตญฺจ ธมฺมํ ปฏิจฺจ นอพฺยากโต นเหตุ เจว นเหตุ\-วิปฺปยุตฺโต จ ธมฺโม อุปฺปชฺชติ เหตุปจฺจยา.  กุสลํ เหตุญฺเจว เหตุ\-สมฺปยุตฺตญฺจ ธมฺมํ ปฏิจฺจ นอกุสโล นเหตุ เจว นเหตุ\-วิปฺปยุตฺโต จ นอพฺยากโต นเหตุ เจว นเหตุ\-วิปฺปยุตฺโต จ ธมฺมา อุปฺปชฺชนฺติ เหตุปจฺจยา. ตีณิ.}

\prose{(อกุสลํ ตีณิ กาตพฺพํ. อพฺยากตํ ตีณิ กาตพฺพํ.)}

\prose{เหตุยา นว.}

\proseitem{14}{กุสลํ เหตุ\-สมฺปยุตฺต\-ญฺเจว น จ เหตุํ ธมฺมํ ปฏิจฺจ นอกุสโล นเหตุ\-วิปฺปยุตฺโต เจว นนเหตุ จ ธมฺโม อุปฺปชฺชติ เหตุปจฺจยา.  กุสลํ เหตุ\-สมฺปยุตฺต\-ญฺเจว น จ เหตุํ ธมฺมํ ปฏิจฺจ นอพฺยากโต นเหตุ\-วิปฺปยุตฺโต เจว นนเหตุ จ ธมฺโม อุปฺปชฺชติ เหตุปจฺจยา.  กุสลํ เหตุ\-สมฺปยุตฺต\-ญฺเจว น จ เหตุํ ธมฺมํ ปฏิจฺจ นอกุสโล นเหตุ\-วิปฺปยุตฺโต เจว นนเหตุ จ นอพฺยากโต นเหตุ\-วิปฺปยุตฺโต เจว นนเหตุ จ ธมฺมา อุปฺปชฺชนฺติ เหตุปจฺจยา. ตีณิ.}

\prose{อกุสลํ เหตุ\-สมฺปยุตฺต\-ญฺเจว น จ เหตุํ ธมฺมํ ปฏิจฺจ. ตีณิ. อพฺยากตํ ตีณิ.  เหตุยา นว.}

\csromanmarkh{1. กุสลตฺติก}\csromanmarkhii{6. เหตุ\-สเหตุกทุก}\csromanheadernomark{2}{1. กุสลตฺติก 6. เหตุ\-สเหตุกทุก}
\proseitem{15}{กุสลํ นเหตุํ สเหตุกํ ธมฺมํ ปฏิจฺจ นกุสโล นเหตุ นสเหตุโก ธมฺโม อุปฺปชฺชติ เหตุปจฺจยา.  กุสลํ นเหตุํ สเหตุกํ ธมฺมํ ปฏิจฺจ นอกุสโล นเหตุ นสเหตุโก ธมฺโม}

\vspace{\tipitakaparskip}
\csromancenter{7. สปฺปจฺจยทุก อุปฺปชฺชติ เหตุปจฺจยา.  กุสลํ นเหตุํ สเหตุกํ ธมฺมํ ปฏิจฺจ นกุสโล นเหตุ นสเหตุโก จ นอกุสโล นเหตุ นสเหตุโก จ ธมฺมา อุปฺปชฺชนฺติ เหตุปจฺจยา. ตีณิ.}
\prose{\csromanfolio{251}อกุสลํ นเหตุํ สเหตุกํ. ตีณิ.}

\prose{อพฺยากตํ นเหตุํ สเหตุกํ ธมฺมํ ปฏิจฺจ นกุสโล ฯเปฯ. นอกุสโล ฯเปฯ. นกุสโล นเหตุ นสเหตุโก จ นอกุสโล นเหตุ นสเหตุโก จ ธมฺมา อุปฺปชฺชนฺติ เหตุปจฺจยา. ตีณิ.}

\proseitemwithcloser{16}{อพฺยากตํ นเหตุํ อเหตุกํ ธมฺมํ ปฏิจฺจ นกุสโล นเหตุ นอเหตุโก ธมฺโม อุปฺปชฺชติ เหตุปจฺจยา.  อพฺยากตํ นเหตุํ อเหตุกํ ธมฺมํ ปฏิจฺจ นอกุสโล นเหตุ นอเหตุโก ธมฺโม อุปฺปชฺชติ เหตุปจฺจยา.  อพฺยากตํ นเหตุํ อเหตุกํ ธมฺมํ ปฏิจฺจ นกุสโล นเหตุ นอเหตุโก จ นอกุสโล นเหตุ นอเหตุโก จ ธมฺมา อุปฺปชฺชนฺติ เหตุปจฺจยา. (3)}{เหตุโคจฺฉกํ นิฏฺฐิตํ.}
\csromanmarkh{1. กุสลตฺติก}\csromanmarkhii{7. สปฺปจฺจยทุก}\csromanheadernomark{2}{1. กุสลตฺติก 7. สปฺปจฺจยทุก}
\proseitem{17}{อพฺยากโต อปฺปจฺจโย ธมฺโม นอพฺยากตสฺส นอปฺปจฺจยสฺส ธมฺมสฺส อารมฺมณ\-ปจฺจเยน ปจฺจโย.  อพฺยากโต อปฺปจฺจโย ธมฺโม นกุสลสฺส นอปฺปจฺจยสฺส ธมฺมสฺส อารมฺมณ\-ปจฺจเยน ปจฺจโย.  อพฺยากโต อปฺปจฺจโย ธมฺโม นอกุสลสฺส นอปฺปจฺจยสฺส ธมฺมสฺส อารมฺมณ\-ปจฺจเยน ปจฺจโย.  อพฺยากโต อปฺปจฺจโย ธมฺโม นอกุสลสฺส น อปฺปจฺจยสฺส จ นอพฺยากตสฺส นอปฺปจฺจยสฺส จ ธมฺมสฺส อารมฺมณ\-ปจฺจเยน ปจฺจโย.  อพฺยากโต อปฺปจฺจโย ธมฺโม นกุสลสฺส นอปฺปจฺจยสฺส จ นอกุสลสฺส นอปฺปจฺจยสฺส จ ธมฺมสฺส อารมฺมณ\-ปจฺจเยน ปจฺจโย. (5)}

\prose{อารมฺมเณ ปญฺจ. (อสงฺขตํ อปฺปจฺจย\-สทิสํ.)}

\csromanmarkh{ธมฺมา\-นุโลม\-ปจฺจนีย ติกทุก\-ปฏฺฐาน\-ปาฬิ}\csromanmarkhii{1. กุสลตฺติก 9. สนิทสฺสนทุก}\csromanheadernomark{2}{ธมฺมา\-นุโลม\-ปจฺจนีย ติกทุก\-ปฏฺฐาน\-ปาฬิ 1. กุสลตฺติก 9. สนิทสฺสนทุก}
\proseitem{18}{\csromanfolio{252}อพฺยากโต สนิทสฺสโน ธมฺโม นอพฺยากตสฺส นสนิทสฺสนสฺส ธมฺมสฺส อารมฺมณ\-ปจฺจเยน ปจฺจโย.  อพฺยากโต สนิทสฺสโน ธมฺโม นกุสลสฺส นสนิทสฺสนสฺส ฯเปฯ. (ฉ ปญฺหา กาตพฺพา.)}

\prose{กุสลํ อนิทสฺสนํ ธมฺมํ ปฏิจฺจ นกุสโล นอนิทสฺสโน ธมฺโม อุปฺปชฺชติ เหตุปจฺจยา.  อกุสเลน ตีณิเยว. อพฺยากเตน ตีณิเยว.  เหตุยา นว.}

\csromanmarkh{1. กุสลตฺติก}\csromanmarkhii{10. สปฺปฏิฆทุก}\csromanheadernomark{2}{1. \textbf{กุสลตฺติก} 10. สปฺปฏิฆทุก}
\prosecont{อพฺยากตํ สปฺปฏิฆํ ธมฺมํ ปฏิจฺจ นกุสโล นสปฺปฏิโฆ ธมฺโม อุปฺปชฺชติ เหตุปจฺจยา.  อพฺยากตํ สปฺปฏิฆํ ธมฺมํ ปฏิจฺจ นอกุสโล นสปฺปฏิโฆ ธมฺโม อุปฺปชฺชติ เหตุปจฺจยา.  อพฺยากตํ สปฺปฏิฆํ ธมฺมํ ปฏิจฺจ นกุสโล นสปฺปฏิโฆ จ นอกุสโล นสปฺปฏิโฆ จ ธมฺมา อุปฺปชฺชนฺติ เหตุปจฺจยา.  เหตุยา ตีณิ.}

\proseitem{19}{กุสลํ อปฺปฏิเฆน ตีณิ. อกุสลํ อปฺปฏิเฆน ตีณิ. อพฺยากตํ อปฺปฏิเฆน ตีณิ. กุสลํ อปฺปฏิฆญฺจ อพฺยากตํ อปฺปฏิฆญฺจ ตีณิ. อกุสลํ อปฺปฏิฆญฺจ อพฺยากตํ อปฺปฏิฆญฺจ ตีณิ.  เหตุยา ปนฺนรส.}

\csromanmarkh{1. กุสลตฺติก}\csromanmarkhii{11. รูปีทุก}\csromanheadernomark{2}{1. \textbf{กุสลตฺติก} 11. รูปีทุก}
\proseitem{20}{อพฺยากตํ รูปิํ ธมฺมํ ปฏิจฺจ นกุสโล นรูปี ธมฺโม อุปฺปชฺชติ เหตุปจฺจยา.  อพฺยากตํ รูปิํ ธมฺมํ ปฏิจฺจ นอกุสโล นรูปี ธมฺโม อุปฺปชฺชติ เหตุปจฺจยา.  อพฺยากตํ รูปิํ ธมฺมํ ปฏิจฺจ นกุสโล นรูปี จ นอกุสโล น รูปี จ ธมฺมา อุปฺปชฺชนฺติ เหตุปจฺจยา.  เหตุยา ตีณิ.}

\prose{กุสลํ อรูปิํ ธมฺมํ ปฏิจฺจ ตีณิ. อกุสลํ อรูปิํ ธมฺมํ ปฏิจฺจ ตีณิ. อพฺยากตํ อรูปิํ ธมฺมํ ปฏิจฺจ ตีณิ.  เหตุยา นว.}

\csromanmarkh{1. กุสลตฺติก}\csromanmarkhii{12. โลกิยทุก}\csromanheadernomark{2}{1. \textbf{กุสลตฺติก} 12. โลกิยทุก}
\proseitem{21}{อพฺยากตํ โลกิยํ ธมฺมํ ปจฺจยา นอพฺยากโต นโลกิโย ธมฺโม อุปฺปชฺชติ เหตุปจฺจยา.  อพฺยากตํ โลกิยํ ธมฺมํ ปจฺจยา นกุสโล นโลกิโย ธมฺโม อุปฺปชฺชติ เหตุปจฺจยา.  อพฺยากตํ}

\vspace{\tipitakaparskip}
\csromancenter{13. เกนจิ\-วิญฺเญยฺยทุก โลกิยํ ธมฺมํ ปจฺจยา นอกุสโล นโลกิโย ธมฺโม อุปฺปชฺชติ เหตุปจฺจยา.  อพฺยากตํ โลกิยํ ธมฺมํ ปจฺจยา นอกุสโล นโลกิโย จ นอพฺยากโต นโลกิโย จ ธมฺมา อุปฺปชฺชนฺติ เหตุปจฺจยา.  อพฺยากตํ โลกิยํ ธมฺมํ ปจฺจยา นกุสโล นโลกิโย จ นอกุสโล นโลกิโย จ ธมฺมา อุปฺปชฺชนฺติ เหตุปจฺจยา.  เหตุยา ปญฺจ.}
\prose{\csromanfolio{253}กุสลํ โลกุตฺตรํ ธมฺมํ ปฏิจฺจ นกุสโล นโลกุตฺตโร ธมฺโม อุปฺปชฺชติ เหตุปจฺจยา.  กุสลํ โลกุตฺตรํ ธมฺมํ ปฏิจฺจ นอกุสโล นโลกุตฺตโร ธมฺโม อุปฺปชฺชติ เหตุปจฺจยา.  กุสลํ โลกุตฺตรํ ธมฺมํ ปฏิจฺจ นกุสโล นโลกุตฺตโร จ นอกุสโล นโลกุตฺตโร จ ธมฺมา อุปฺปชฺชนฺติ เหตุปจฺจยา. ตีณิ.}

\prose{อพฺยากตํ โลกุตฺตรํ ธมฺมํ ปฏิจฺจ นกุสโล นโลกุตฺตโร ธมฺโม อุปฺปชฺชติ เหตุปจฺจยา.  อพฺยากตํ โลกุตฺตรํ ธมฺมํ ปฏิจฺจ นอกุสโล นโลกุตฺตโร ธมฺโม อุปฺปชฺชติ เหตุปจฺจยา.  อพฺยากตํ โลกุตฺตรํ ธมฺมํ ปฏิจฺจ นกุสโล นโลกุตฺตโร จ นอกุสโล นโลกุตฺตโร จ ธมฺมา อุปฺปชฺชนฺติ เหตุปจฺจยา. ตีณิ.}

\csromanmarkh{1. กุสลตฺติก}\csromanmarkhii{13. เกนจิ\-วิญฺเญยฺยทุก}\csromanheadernomark{2}{1. กุสลตฺติก 13. เกนจิ\-วิญฺเญยฺยทุก}
\proseitem{22}{กุสลํ เกนจิ วิญฺเญยฺยํ ธมฺมํ ปฏิจฺจ นกุสโล นเกนจิ วิญฺเญยฺโย ธมฺโม อุปฺปชฺชติ เหตุปจฺจยา.  กุสลํ เกนจิ วิญฺเญยฺยํ ธมฺมํ ปฏิจฺจ นอกุสโล นเกนจิ วิญฺเญยฺโย ธมฺโม อุปฺปชฺชติ เหตุปจฺจยา.  กุสลํ เกนจิ วิญฺเญยฺยํ ธมฺมํ ปฏิจฺจ นอพฺยากโต นเกนจิ วิญฺเญยฺโย ธมฺโม อุปฺปชฺชติ เหตุปจฺจยา. (สํขิตฺตํ.) . เหตุยา เอกูนวีสติ.}

\proseitem{23}{กุสลํ เกนจิ นวิญฺเญยฺยํ ธมฺมํ ปฏิจฺจ นกุสโล นเกนจิ นวิญฺเญยฺโย ธมฺโม อุปฺปชฺชติ เหตุปจฺจยา.  กุสลํ เกนจิ นวิญฺเญยฺยํ ธมฺมํ ปฏิจฺจ นอกุสโล นเกนจิ นวิญฺเญยฺโย ธมฺโม อุปฺปชฺชติ เหตุปจฺจยา.  กุสลํ เกนจิ นวิญฺเญยฺยํ ธมฺมํ ปฏิจฺจ นอพฺยากโต นเกนจิ นวิญฺเญยฺโย ธมฺโม อุปฺปชฺชติ เหตุปจฺจยา.}

\prose{(เอเตน อุปาเยน เกนจิ นวิญฺเญยฺเย เอกูนวีสติ ปญฺหา กาตพฺพา.)}

\csromanmarkh{ธมฺมา\-นุโลม\-ปจฺจนีย ติกทุก\-ปฏฺฐาน\-ปาฬิ}\csromanmarkhii{1. กุสลตฺติก 14. อาสวโคจฺฉก}\csromanheadernomark{2}{ธมฺมา\-นุโลม\-ปจฺจนีย ติกทุก\-ปฏฺฐาน\-ปาฬิ 1. กุสลตฺติก 14. อาสวโคจฺฉก}
\proseitem{24}{\csromanfolio{254}อกุสลํ อาสวํ ธมฺมํ ปฏิจฺจ นอกุสโล นอาสโว ธมฺโม ฯเปฯ.  นกุสโล นอาสโว ธมฺโม ฯเปฯ.  นอพฺยากโต นอาสโว ธมฺโม ฯเปฯ.  นกุสโล นอาสโว จ นอพฺยากโต นอาสโว จ ธมฺมา ฯเปฯ.  นกุสโล นอาสโว จ นอกุสโล นอาสโว จ ธมฺมา อุปฺปชฺชนฺติ เหตุปจฺจยา.  เหตุยา ปญฺจ, อารมฺมเณ ตีณิ ฯเปฯ อวิคเต ปญฺจ.}

\prose{อกุสลํ โนอาสวํ ธมฺมํ ปฏิจฺจ นกุสโล นโนอาสโว ธมฺโม อุปฺปชฺชติ เหตุปจฺจยา.  อกุสลํ โนอาสวํ ธมฺมํ ปฏิจฺจ นอพฺยากโต นโนอาสโว ธมฺโม ฯเปฯ.  นกุสโล นโนอาสโว จ นอพฺยากโต นโนอาสโว จ ธมฺมา อุปฺปชฺชนฺติ เหตุปจฺจยา.  เหตุยา ตีณิ ฯเปฯ อวิคเต ตีณิ.}

\prose{25. อพฺยากตํ สาสวํ ธมฺมํ ปฏิจฺจ ฯเปฯ. (โลกิย\-ทุก\-สทิสํ.)}

\proseitem{26}{อกุสลํ อาสว\-สมฺปยุตฺตํ ธมฺมํ ปฏิจฺจ นอกุสโล นอาสว\-สมฺปยุตฺโต ธมฺโม ฯเปฯ.  นกุสโล นอาสว\-สมฺปยุตฺโต ธมฺโม ฯเปฯ.  นอพฺยากโต นอาสว\-สมฺปยุตฺโต ธมฺโม ฯเปฯ.  นกุสโล นอาสว\-สมฺปยุตฺโต จ นอพฺยากโต นอาสว\-สมฺปยุตฺโต จ ธมฺมา ฯเปฯ.  นกุสโล นอาสว\-สมฺปยุตฺโต จ นอกุสโล นอาสว\-สมฺปยุตฺโต จ ธมฺมา อุปฺปชฺชนฺติ เหตุปจฺจยา.  เหตุยา ปญฺจ.}

\prose{อกุสลํ อาสว\-วิปฺปยุตฺตํ ธมฺมํ ปฏิจฺจ นกุสโล นอาสว\-วิปฺปยุตฺโต ธมฺโม ฯเปฯ.  นอพฺยากโต นอาสว\-วิปฺปยุตฺโต ธมฺโม ฯเปฯ.  นกุสโล นอาสว\-วิปฺปยุตฺโต จ นอพฺยากโต นอาสว\-วิปฺปยุตฺโต จ ธมฺมา อุปฺปชฺชนฺติ เหตุปจฺจยา.  เหตุยา ตีณิ.}

\proseitem{27}{อกุสลํ อาสวญฺเจว สาสวญฺจ ธมฺมํ ปฏิจฺจ นอกุสโล นอาสโว เจว นอนาสโว จ ธมฺโม ฯเปฯ.  นกุสโล นอาสโว เจว นอนาสโว จ ธมฺโม ฯเปฯ.  นอพฺยากโต นอาสโว เจว นอนาสโว จ ธมฺโม ฯเปฯ.  นกุสโล นอาสโว เจว นอนาสโว จ นอพฺยากโต นอาสโว เจว นอนาสโว จ ธมฺมา ฯเปฯ.  นกุสโล}

\vspace{\tipitakaparskip}
\csromancenter{55. สารมฺมณทุก นอาสโว เจว นอนาสโว จ นอกุสโล นอาสโว เจว นอนาสโว จ ธมฺมา อุปฺปชฺชนฺติ เหตุปจฺจยา.  เหตุยา ปญฺจ.}
\prose{\csromanfolio{255}อกุสลํ สาสวญฺเจว โน จ อาสวํ ธมฺมํ ปฏิจฺจ นกุสโล นอนาสโว เจว นโน จ อาสโว ธมฺโม ฯเปฯ.  นอพฺยากโต นอนาสโว เจว นโน จ อาสโว ธมฺโม ฯเปฯ.  นกุสโล นอนาสโว เจว นโน จ อาสโว นอพฺยากโต นอนาสโว เจว นโนอาสโว จ ธมฺมา อุปฺปชฺชนฺติ เหตุปจฺจยา.  เหตุยา ตีณิ.}

\proseitem{28}{อกุสลํ อาสวญฺเจว อาสว\-สมฺปยุตฺตญฺจ ธมฺมํ ปฏิจฺจ นกุสโล นอาสโว เจว นอาสว\-วิปฺปยุตฺโต จ ธมฺโม ฯเปฯ.  นอพฺยากโต นอาสโว เจว นอาสว\-วิปฺปยุตฺโต จ ธมฺโม ฯเปฯ.  นกุสโล นอาสโว เจว นอาสว\-วิปฺปยุตฺโต จ นอพฺยากโต นอาสโว เจว นอาสว\-วิปฺปยุตฺโต จ ธมฺมา อุปฺปชฺชนฺติ เหตุปจฺจยา.  เหตุยา ตีณิ.}

\prose{อกุสลํ อาสว\-สมฺปยุตฺต\-ญฺเจว โน จ อาสวํ ธมฺมํ ปฏิจฺจ นกุสโล นอาสว\-วิปฺปยุตฺโต เจว นโน จ อาสโว ธมฺโม ฯเปฯ. นอพฺยากโต นอาสว\-วิปฺปยุตฺโต เจว นโน จ อาสโว ธมฺโม ฯเปฯ. นกุสโล นอาสว\-วิปฺปยุตฺโต เจว นโน อาสโว จ นอพฺยากโต นอาสว\-วิปฺปยุตฺโต เจว นโน จ อาสโว จ ธมฺมา อุปฺปชฺชนฺติ เหตุปจฺจยา.  เหตุยา ตีณิ.}

\proseitem{29}{อพฺยากตํ อาสว\-วิปฺปยุตฺตํ สาสวํ ธมฺมํ ปฏิจฺจ ฯเปฯ. (โลกิย\-ทุก\-สทิสํ.)}

\csromanmarkh{1. กุสลตฺติก}\csromanmarkhii{20. ฉโคจฺฉกทุก}\csromanheadernomark{2}{1. \textbf{กุสลตฺติก} 20. ฉโคจฺฉกทุก}
\proseitem{30}{อกุสลํ สญฺโญชนํ ธมฺมํ ปฏิจฺจ ฯเปฯ. คนฺถํ ฯเปฯ. โอฆํ ฯเปฯ. โยคํ ฯเปฯ. นีวรณํ ฯเปฯ. ปรามาสํ. (สํขิตฺตํ.)}

\csromanmarkh{1. กุสลตฺติก}\csromanmarkhii{55. สารมฺมณทุก}\csromanheadernomark{2}{1. กุสลตฺติก 55. สารมฺมณทุก}
\proseitem{31}{กุสลํ สารมฺมณํ ธมฺมํ ปฏิจฺจ นกุสโล นสารมฺมโณ ธมฺโม อุปฺปชฺชติ เหตุปจฺจยา.  กุสลํ สารมฺมณํ ธมฺมํ ปฏิจฺจ นอกุสโล นสารมฺมโณ ธมฺโม อุปฺปชฺชติ เหตุปจฺจยา.  กุสลํ สารมฺมณํ ธมฺมํ}

\vspace{\tipitakaparskip}
\csromancenter{ติกทุก\-ปฏฺฐาน\-ปาฬิ ปฏิจฺจ นกุสโล นสารมฺมโณ จ นอกุสโล นสารมฺมโณ จ ธมฺมา อุปฺปชฺชนฺติ เหตุปจฺจยา. ตีณิ.}
\prose{\csromanfolio{256}อกุสลํ สารมฺมณํ ธมฺมํ ปฏิจฺจ นอกุสโล นสารมฺมโณ ธมฺโม อุปฺปชฺชติ เหตุปจฺจยา.  อกุสลํ สารมฺมณํ ธมฺมํ ปฏิจฺจ นกุสโล นสารมฺมโณ ธมฺโม อุปฺปชฺชติ เหตุปจฺจยา.  อกุสลํ สารมฺมณํ ธมฺมํ ปฏิจฺจ นกุสโล นสารมฺมโณ จ นอกุสโล นสารมฺมโณ จ ธมฺมา อุปฺปชฺชนฺติ เหตุปจฺจยา. ตีณิ.}

\prose{อพฺยากตํ สารมฺมณํ ธมฺมํ ปฏิจฺจ นกุสโล นสารมฺมโณ ธมฺโม อุปฺปชฺชติ เหตุปจฺจยา. ตีณิ.}

\prose{เหตุยา นว.}

\proseitem{32}{อพฺยากตํ อนารมฺมณํ ธมฺมํ ปฏิจฺจ นกุสโล นอนารมฺมโณ ธมฺโม อุปฺปชฺชติ เหตุปจฺจยา.  อพฺยากตํ อนารมฺมณํ ธมฺมํ ปฏิจฺจ นอกุสโล นอนารมฺมโณ ธมฺโม อุปฺปชฺชติ เหตุปจฺจยา.  อพฺยากตํ อนารมฺมณํ ธมฺมํ ปฏิจฺจ นกุสโล นอนารมฺมโณ จ นอกุสโล นอนารมฺมโณ จ ธมฺมา อุปฺปชฺชนฺติ เหตุปจฺจยา. (สํขิตฺตํ.) . เหตุยา ตีณิ.}

\csromanmarkh{1. กุสลตฺติก}\csromanmarkhii{56. จิตฺตทุก}\csromanheadernomark{2}{1. \textbf{กุสลตฺติก} 56. จิตฺตทุก}
\proseitem{33}{กุสลํ จิตฺตํ ธมฺมํ ปฏิจฺจ นกุสโล นจิตฺโต ธมฺโม อุปฺปชฺชติ เหตุปจฺจยา.  กุสลํ จิตฺตํ ธมฺมํ ปฏิจฺจ นอกุสโล นจิตฺโต ธมฺโม อุปฺปชฺชติ เหตุปจฺจยา.  กุสลํ จิตฺตํ ธมฺมํ ปฏิจฺจ นอพฺยากโต นจิตฺโต ธมฺโม อุปฺปชฺชติ เหตุปจฺจยา.  กุสลํ จิตฺตํ ธมฺมํ ปฏิจฺจ นอกุสโล นจิตฺโต จ นอพฺยากโต นจิตฺโต จ ธมฺมา อุปฺปชฺชนฺติ เหตุปจฺจยา.  กุสลํ จิตฺตํ ธมฺมํ ปฏิจฺจ นกุสโล นจิตฺโต จ นอกุสโล นจิตฺโต จ ธมฺมา อุปฺปชฺชนฺติ เหตุปจฺจยา. ปญฺจ.}

\prose{อกุสลํ จิตฺตํ ธมฺมํ ปฏิจฺจ นอกุสโล นจิตฺโต ธมฺโม อุปฺปชฺชติ เหตุปจฺจยา.  อกุสลํ จิตฺตํ ธมฺมํ ปฏิจฺจ นกุสโล นจิตฺโต ธมฺโม อุปฺปชฺชติ เหตุปจฺจยา.  อกุสลํ จิตฺตํ ธมฺมํ ปฏิจฺจ นอพฺยากโต นจิตฺโต ธมฺโม อุปฺปชฺชติ เหตุปจฺจยา.  อกุสลํ จิตฺตํ ธมฺมํ ปฏิจฺจ นกุสโล นจิตฺโต จ นอพฺยากโต นจิตฺโต จ ธมฺมา อุปฺปชฺชนฺติ เหตุปจฺจยา.  }

\vspace{\tipitakaparskip}
\csromancenter{57. เจตสิกทุก อกุสลํ จิตฺตํ ธมฺมํ ปฏิจฺจ นกุสโล นจิตฺโต จ นอกุสโล นจิตฺโต จ ธมฺมา อุปฺปชฺชนฺติ เหตุปจฺจยา. ปญฺจ.}
\prose{\csromanfolio{257}อพฺยากตํ จิตฺตํ ธมฺมํ ปฏิจฺจ นกุสโล นจิตฺโต ธมฺโม อุปฺปชฺชติ เหตุปจฺจยา.  อพฺยากตํ จิตฺตํ ธมฺมํ ปฏิจฺจ นอกุสโล นจิตฺโต ธมฺโม อุปฺปชฺชติ เหตุปจฺจยา.  อพฺยากตํ จิตฺตํ ธมฺมํ ปฏิจฺจ นกุสโล นจิตฺโต จ นอกุสโล นจิตฺโต จ ธมฺมา อุปฺปชฺชนฺติ เหตุปจฺจยา. ตีณิ. เหตุยา เตรส.}

\proseitem{34}{กุสลํ โนจิตฺตํ ธมฺมํ ปฏิจฺจ นอกุสโล นโนจิตฺโต ธมฺโม อุปฺปชฺชติ เหตุปจฺจยา.  กุสลํ โนจิตฺตํ ธมฺมํ ปฏิจฺจ นอพฺยากโต นโนจิตฺโต ธมฺโม อุปฺปชฺชติ เหตุปจฺจยา.  กุสลํ โนจิตฺตํ ธมฺมํ ปฏิจฺจ นอกุสโล นโนจิตฺโต จ นอพฺยากโต นโนจิตฺโต จ ธมฺมา อุปฺปชฺชนฺติ เหตุปจฺจยา. ตีณิ.}

\prose{อกุสลํ โนจิตฺตํ ธมฺมํ ปฏิจฺจ นกุสโล นโนจิตฺโต ธมฺโม อุปฺปชฺชติ เหตุปจฺจยา.  อกุสลํ โนจิตฺตํ ธมฺมํ ปฏิจฺจ นอพฺยากโต นโนจิตฺโต ธมฺโม อุปฺปชฺชติ เหตุปจฺจยา.  อกุสลํ โนจิตฺตํ ธมฺมํ ปฏิจฺจ นกุสโล นโนจิตฺโต จ นอพฺยากโต นโนจิตฺโต จ ธมฺมา อุปฺปชฺชนฺติ เหตุปจฺจยา. ตีณิ.}

\prose{อพฺยากตานิ ตีณิ.  เหตุยา นว.}

\csromanmarkh{1. กุสลตฺติก}\csromanmarkhii{57. เจตสิกทุก}\csromanheadernomark{2}{1. กุสลตฺติก 57. เจตสิกทุก}
\proseitem{35}{กุสลํ เจตสิกํ ธมฺมํ ปฏิจฺจ นกุสโล นเจตสิโก ธมฺโม อุปฺปชฺชติ เหตุปจฺจยา.  กุสลํ เจตสิกํ ธมฺมํ ปฏิจฺจ นอกุสโล นเจตสิโก ธมฺโม อุปฺปชฺชติ เหตุปจฺจยา.  กุสลํ เจตสิกํ ธมฺมํ ปฏิจฺจ นอพฺยากโต นเจตสิโก ธมฺโม อุปฺปชฺชติ เหตุปจฺจยา.  กุสลํ เจตสิกํ ธมฺมํ ปฏิจฺจ นอกุสโล นเจตสิโก จ นอพฺยากโต นเจตสิโก จ ธมฺมา อุปฺปชฺชนฺติ เหตุปจฺจยา.  กุสลํ เจตสิกํ ธมฺมํ ปฏิจฺจ นกุสโล นเจตสิโก จ นอกุสโล นเจตสิโก จ ธมฺมา อุปฺปชฺชนฺติ เหตุปจฺจยา. ปญฺจ.}

\prose{อกุสลํ เจตสิกํ ธมฺมํ ปฏิจฺจ นอกุสโล นเจตสิโก ธมฺโม อุปฺปชฺชติ เหตุปจฺจยา.  อกุสลํ ธมฺมํ ปฏิจฺจ นกุสโล}

\vspace{\tipitakaparskip}
\csromancenter{ติกทุก\-ปฏฺฐาน\-ปาฬิ นเจตสิโก ธมฺโม อุปฺปชฺชติ เหตุปจฺจยา.  อกุสลํ เจตสิกํ ธมฺมํ ปฏิจฺจ นอพฺยากโต นเจตสิโก ธมฺโม อุปฺปชฺชติ เหตุปจฺจยา.  อกุสลํ เจตสิกํ ธมฺมํ ปฏิจฺจ นกุสโล นเจตสิโก จ นอพฺยากโต นเจตสิโก จ ธมฺมา อุปฺปชฺชนฺติ เหตุปจฺจยา.  อกุสลํ เจตสิกํ ธมฺมํ ปฏิจฺจ นกุสโล นเจตสิโก จ นอกุสโล นเจตสิโก จ ธมฺมา อุปฺปชฺชนฺติ เหตุปจฺจยา. ปญฺจ.}
\prose{\csromanfolio{258}อพฺยากตํ เจตสิกํ ธมฺมํ ปฏิจฺจ นกุสโล นเจตสิโก ธมฺโม อุปฺปชฺชติ เหตุปจฺจยา.  อพฺยากตํ เจตสิกํ ธมฺมํ ปฏิจฺจ นอกุสโล นเจตสิโก ธมฺโม อุปฺปชฺชติ เหตุปจฺจยา.  อพฺยากตํ เจตสิกํ ธมฺมํ ปฏิจฺจ นกุสโล นเจตสิโก จ นอกุสโล นเจตสิโก จ ธมฺมา อุปฺปชฺชนฺติ เหตุปจฺจยา. ตีณิ.}

\prose{เหตุยา เตรส.  อเจตสิกานิ นว.}

\csromanmarkh{1. กุสลตฺติก}\csromanmarkhii{58. จิตฺต\-สมฺปยุตฺตทุก}\csromanheadernomark{2}{1. กุสลตฺติก 58. จิตฺต\-สมฺปยุตฺตทุก}
\prosecont{กุสลํ จิตฺต\-สมฺปยุตฺตํ ธมฺมํ ปฏิจฺจ นกุสโล นจิตฺต\-สมฺปยุตฺโต ธมฺโม อุปฺปชฺชติ เหตุปจฺจยา.  กุสลํ จิตฺต\-สมฺปยุตฺตํ ธมฺมํ ปฏิจฺจ นอกุสโล นจิตฺต\-สมฺปยุตฺโต ธมฺโม อุปฺปชฺชติ เหตุปจฺจยา.  กุสลํ จิตฺต\-สมฺปยุตฺตํ ธมฺมํ ปฏิจฺจ นอพฺยากโต นจิตฺต\-สมฺปยุตฺโต ธมฺโม อุปฺปชฺชติ เหตุปจฺจยา.  กุสลํ จิตฺต\-สมฺปยุตฺตํ ธมฺมํ ปฏิจฺจ นอกุสโล นจิตฺต\-สมฺปยุตฺโต จ นอพฺยากโต นจิตฺต\-สมฺปยุตฺโต จ ธมฺมา อุปฺปชฺชนฺติ เหตุปจฺจยา.  กุสลํ จิตฺต\-สมฺปยุตฺตํ ธมฺมํ ปฏิจฺจ นกุสโล นจิตฺต\-สมฺปยุตฺโต จ นอกุสโล นจิตฺต\-สมฺปยุตฺโต จ ธมฺมา อุปฺปชฺชนฺติ เหตุปจฺจยา. ปญฺจ.}

\proseitem{36}{อกุสลํ จิตฺต\-สมฺปยุตฺตํ ธมฺมํ ปฏิจฺจ นอกุสโล นจิตฺต\-สมฺปยุตฺโต ธมฺโม อุปฺปชฺชติ เหตุปจฺจยา. ปญฺจ.}

\prose{อพฺยากตํ จิตฺต\-สมฺปยุตฺตํ ธมฺมํ ปฏิจฺจ นกุสโล นจิตฺต\-สมฺปยุตฺโต ธมฺโม อุปฺปชฺชติ เหตุปจฺจยา.  อพฺยากตํ จิตฺต\-สมฺปยุตฺตํ ธมฺมํ ปฏิจฺจ นอกุสโล นจิตฺต\-สมฺปยุตฺโต ธมฺโม อุปฺปชฺชติ เหตุปจฺจยา.  อพฺยากตํ จิตฺต\-สมฺปยุตฺตํ ธมฺมํ ปฏิจฺจ นกุสโล นจิตฺต\-สมฺปยุตฺโต จ นอกุสโล นจิตฺต\-สมฺปยุตฺโต จ ธมฺมา อุปฺปชฺชนฺติ เหตุปจฺจยา. ตีณิ. (สํขิตฺตํ.) . เหตุยา เตรส.  จิตฺต\-วิปฺปยุตฺเต ตีณิ.}

\vspace{\tipitakaparskip}
\csromancenter{60. จิตฺต\-สมุฏฺฐานทุก 1. กุสลตฺติก 59. จิตฺต\-สํสฏฺฐทุก}
\prosecont{\csromanfolio{259}กุสลํ จิตฺต\-สํสฏฺฐํ ธมฺมํ ปฏิจฺจ นกุสโล นจิตฺต\-สํสฏฺโฐ ธมฺโม อุปฺปชฺชติ เหตุปจฺจยา.  กุสลํ จิตฺต\-สํสฏฺฐํ ธมฺมํ ปฏิจฺจ นอกุสโล นจิตฺต\-สํสฏฺโฐ ธมฺโม อุปฺปชฺชติ เหตุปจฺจยา.  กุสลํ จิตฺต\-สํสฏฺฐํ ธมฺมํ ปฏิจฺจ นอพฺยากโต นจิตฺต\-สํสฏฺโฐ ธมฺโม อุปฺปชฺชติ เหตุปจฺจยา.  กุสลํ จิตฺต\-สํสฏฺฐํ ธมฺมํ ปฏิจฺจ นอกุสโล นจิตฺต\-สํสฏฺโฐ จ นอพฺยากโต นจิตฺต\-สํสฏฺโฐ จ ธมฺมา อุปฺปชฺชนฺติ เหตุปจฺจยา.  กุสลํ จิตฺต\-สํสฏฺฐํ ธมฺมํ ปฏิจฺจ นกุสโล นจิตฺต\-สํสฏฺโฐ จ นอกุสโล นจิตฺต\-สํสฏฺโฐ จ ธมฺมา อุปฺปชฺชนฺติ เหตุปจฺจยา. ปญฺจ.}

\proseitem{37}{อกุสลํ จิตฺต\-สํสฏฺฐํ ธมฺมํ ปฏิจฺจ นอกุสโล นจิตฺต\-สํสฏฺโฐ ธมฺโม อุปฺปชฺชติ เหตุปจฺจยา.  อกุสลํ จิตฺต\-สํสฏฺฐํ ธมฺมํ ปฏิจฺจ นกุสโล นจิตฺต\-สํสฏฺโฐ ธมฺโม อุปฺปชฺชติ เหตุปจฺจยา.  อกุสลํ จิตฺต\-สํสฏฺฐํ ธมฺมํ ปฏิจฺจ นอพฺยากโต นจิตฺต\-สํสฏฺโฐ ธมฺโม อุปฺปชฺชติ เหตุปจฺจยา.  อกุสลํ จิตฺต\-สํสฏฺฐํ ธมฺมํ ปฏิจฺจ นกุสโล นจิตฺต\-สํสฏฺโฐ จ นอพฺยากโต นจิตฺต\-สํสฏฺโฐ จ ธมฺมา อุปฺปชฺชนฺติ เหตุปจฺจยา.  อกุสลํ จิตฺต\-สํสฏฺฐํ ธมฺมํ ปฏิจฺจ นกุสโล นจิตฺต\-สํสฏฺโฐ จ นอกุสโล นจิตฺต\-สํสฏฺโฐ จ ธมฺมา อุปฺปชฺชนฺติ เหตุปจฺจยา. ปญฺจ.}

\prose{อพฺยากตํ จิตฺต\-สํสฏฺฐํ ธมฺมํ ปฏิจฺจ นกุสโล นจิตฺต\-สํสฏฺโฐ ธมฺโม อุปฺปชฺชติ เหตุปจฺจยา.  อพฺยากตํ จิตฺต\-สํสฏฺฐํ ธมฺมํ ปฏิจฺจ นอกุสโล นจิตฺต\-สํสฏฺโฐ ธมฺโม อุปฺปชฺชติ เหตุปจฺจยา.  อพฺยากตํ จิตฺต\-สํสฏฺฐํ ธมฺมํ ปฏิจฺจ นกุสโล นจิตฺต\-สํสฏฺโฐ จ นอกุสโล นจิตฺต\-สํสฏฺโฐ จ ธมฺมา อุปฺปชฺชนฺติ เหตุปจฺจยา. ตีณิ. (สํขิตฺตํ.) . เหตุยา เตรส.}

\csromanmarkh{1. กุสลตฺติก}\csromanmarkhii{60. จิตฺต\-สมุฏฺฐานทุก}\csromanheadernomark{2}{1. กุสลตฺติก 60. จิตฺต\-สมุฏฺฐานทุก}
\proseitem{38}{กุสลํ จิตฺต\-สมุฏฺฐานํ ธมฺมํ ปฏิจฺจ นอกุสโล นจิตฺต\-สมุฏฺฐาโน ธมฺโม อุปฺปชฺชติ เหตุปจฺจยา.  กุสลํ จิตฺต\-สมุฏฺฐานํ ธมฺมํ ปฏิจฺจ นอพฺยากโต นจิตฺต\-สมุฏฺฐาโน ธมฺโม อุปฺปชฺชติ เหตุปจฺจยา.  กุสลํ จิตฺต\-สมุฏฺฐานํ ธมฺมํ ปฏิจฺจ นอกุสโล นจิตฺต\-สมุฏฺฐาโน จ นอพฺยากโต นจิตฺต\-สมุฏฺฐาโน จ ธมฺมา อุปฺปชฺชนฺติ เหตุปจฺจยา. ตีณิ.}

\csromanheader{2}{ธมฺมา\-นุโลม\-ปจฺจนีย ติกทุก\-ปฏฺฐาน\-ปาฬิ}
\prose{\csromanfolio{260}อกุสลํ จิตฺต\-สมุฏฺฐานํ ธมฺมํ ปฏิจฺจ นกุสโล นจิตฺต\-สมุฏฺฐาโน ธมฺโม อุปฺปชฺชติ เหตุปจฺจยา.  อกุสลํ จิตฺต\-สมุฏฺฐานํ ธมฺมํ ปฏิจฺจ นอพฺยากโต นจิตฺต\-สมุฏฺฐาโน ธมฺโม อุปฺปชฺชติ เหตุปจฺจยา.  อกุสลํ จิตฺต\-สมุฏฺฐานํ ธมฺมํ ปฏิจฺจ นกุสโล นจิตฺต\-สมุฏฺฐาโน จ นอพฺยากโต นจิตฺต\-สมุฏฺฐาโน จ ธมฺมา อุปฺปชฺชนฺติ เหตุปจฺจยา. ตีณิ.}

\prose{อพฺยากตํ จิตฺต\-สมุฏฺฐานํ ธมฺมํ ปฏิจฺจ นกุสโล นจิตฺต\-สมุฏฺฐาโน ธมฺโม อุปฺปชฺชติ เหตุปจฺจยา. ตีณิ.  เหตุยา นว ฯเปฯ อวิคเต นว.}

\proseitem{39}{กุสลํ โนจิตฺต\-สมุฏฺฐานํ ธมฺมํ ปฏิจฺจ นกุสโล นโน\-จิตฺตสมุฏฺฐาโน ธมฺโม อุปฺปชฺชติ เหตุปจฺจยา. ปญฺจ.}

\prose{อกุสลํ โนจิตฺต\-สมุฏฺฐานํ ธมฺมํ ปฏิจฺจ นอกุสโล นโน\-จิตฺตสมุฏฺฐาโน ธมฺโม อุปฺปชฺชติ เหตุปจฺจยา. ปญฺจ.}

\prose{อพฺยากตํ โนจิตฺต\-สมุฏฺฐานํ ธมฺมํ ปฏิจฺจ นกุสโล นโน\-จิตฺตสมุฏฺฐาโน ธมฺโม อุปฺปชฺชติ เหตุปจฺจยา. ตีณิ.  เหตุยา เตรส ฯเปฯ อวิคเต เตรส.}

\csromanmarkh{1. กุสลตฺติก}\csromanmarkhii{61. จิตฺต\-สหภูทุก}\csromanheadernomark{2}{1. กุสลตฺติก 61. จิตฺต\-สหภูทุก}
\proseitem{40}{กุสลํ จิตฺตสหภุํ ธมฺมํ ปฏิจฺจ นกุสโล นจิตฺตสหภู ธมฺโม อุปฺปชฺชติ เหตุปจฺจยา. ปญฺจ.}

\prose{อกุสลํ จิตฺตสหภุํ ธมฺมํ ปฏิจฺจ นอกุสโล นจิตฺตสหภู ธมฺโม อุปฺปชฺชติ เหตุปจฺจยา. ปญฺจ.}

\prose{อพฺยากตํ จิตฺตสหภุํ ธมฺมํ ปฏิจฺจ นกุสโล นจิตฺตสหภู ธมฺโม อุปฺปชฺชติ เหตุปจฺจยา. ตีณิ.  เหตุยา เตรส ฯเปฯ อวิคเต เตรส.}

\proseitem{41}{กุสลํ โนจิตฺต\-สหภุํ ธมฺมํ ปฏิจฺจ นกุสโล นโนจิตฺตสหภู ธมฺโม อุปฺปชฺชติ เหตุปจฺจยา. ปญฺจ.}

\prose{อกุสลํ โนจิตฺต\-สหภุํ ธมฺมํ ปฏิจฺจ นอกุสโล นโนจิตฺตสหภู ธมฺโม อุปฺปชฺชติ เหตุปจฺจยา. ปญฺจ.}

\csromanmarkh{1. กุสลตฺติก}\csromanmarkhii{69. อุปาทานทุกาทิ}\csromanheadernomark{2}{1. กุสลตฺติก 69. อุปาทานทุกาทิ}
\prose{\csromanfolio{261}อพฺยากตํ โนจิตฺต\-สหภุํ ธมฺมํ ปฏิจฺจ นกุสโล นโนจิตฺตสหภู ธมฺโม อุปฺปชฺชติ เหตุปจฺจยา. ตีณิ. (สํขิตฺตํ.) . เหตุยา เตรส ฯเปฯ อวิคเต เตรส.}

\csromanmarkh{1. กุสลตฺติก}\csromanmarkhii{62. จิตฺตา\-นุปริวตฺติทุก}\csromanheadernomark{2}{1. กุสลตฺติก 62. จิตฺตา\-นุปริวตฺติทุก}
\proseitem{42}{กุสลํ จิตฺตา\-นุปริวตฺติํ ธมฺมํ ปฏิจฺจ. (สํขิตฺตํ.) . เหตุยา เตรส.}

\prose{กุสลํ โนจิตฺตา\-นุปริวตฺติํ ธมฺมํ ปฏิจฺจ. (สํขิตฺตํ.) . เหตุยา เตรส. (เอเต สํขิตฺตา. ทุกตฺตยํ จิตฺต\-ทุก\-สทิสํ.)}

\csromanmarkh{1. กุสลตฺติก}\csromanmarkhii{66. อชฺฌตฺติก\-ทุกาทิ}\csromanheadernomark{2}{1. กุสลตฺติก 66. อชฺฌตฺติก\-ทุกาทิ}
\proseitem{43}{กุสลํ อชฺฌตฺติกํ ธมฺมํ ปฏิจฺจ. (สํขิตฺตํ. จิตฺต\-ทุก\-สทิสํ.)}

\prose{กุสลํ พาหิรํ ธมฺมํ ปฏิจฺจ นอกุสโล นพาหิโร ธมฺโม อุปฺปชฺชติ เหตุปจฺจยา. ตีณิ.}

\proseitem{44}{อพฺยากตํ อุปาทา ธมฺมํ ปฏิจฺจ นกุสโล นอุปาทา ธมฺโม อุปฺปชฺชติ เหตุปจฺจยา.  เหตุยา ตีณิ.}

\prose{กุสลํ โนอุปาทา ธมฺมํ ปฏิจฺจ นกุสโล นโนอุปาทา ธมฺโม อุปฺปชฺชติ เหตุปจฺจยา. ตีณิ.  นว.}

\proseitem{45}{อพฺยากตํ อุปาทินฺนํ ธมฺมํ ปฏิจฺจ นกุสโล นอุปาทินฺโน ธมฺโม อุปฺปชฺชติ เหตุปจฺจยา.  อพฺยากตํ อุปาทินฺนํ ธมฺมํ ปฏิจฺจ นอกุสโล นอุปาทินฺโน ธมฺโม อุปฺปชฺชติ เหตุปจฺจยา.  อพฺยากตํ อุปาทินฺนํ ธมฺมํ ปฏิจฺจ นกุสโล นอุปาทินฺโน จ นอกุสโล นอุปาทินฺโน จ ธมฺมา อุปฺปชฺชนฺติ เหตุปจฺจยา.  เหตุยา ตีณิ.}

\csromanmarkh{1. กุสลตฺติก}\csromanmarkhii{69. อุปาทานทุกาทิ}\csromanheadernomark{2}{1. กุสลตฺติก 69. อุปาทานทุกาทิ}
\proseitem{46}{อกุสลํ อุปาทานํ ธมฺมํ ปฏิจฺจ นอกุสโล โนอุปาทาโน ธมฺโม อุปฺปชฺชติ เหตุปจฺจยา.}

\prose{อกุสลํ กิเลสํ ธมฺมํ ปฏิจฺจ นอกุสโล โนกิเลโส ธมฺโม อุปฺปชฺชติ เหตุปจฺจยา.}

\csromanmarkh{ธมฺมา\-นุโลม\-ปจฺจนีย ติกทุก\-ปฏฺฐาน\-ปาฬิ}\csromanmarkhii{1. กุสลตฺติก 83. ปิฏฺฐิทุก}\csromanheadernomark{2}{ธมฺมา\-นุโลม\-ปจฺจนีย ติกทุก\-ปฏฺฐาน\-ปาฬิ 1. กุสลตฺติก 83. ปิฏฺฐิทุก}
\proseitem{47}{\csromanfolio{262}อกุสลํ ทสฺสเนน ปหาตพฺพํ ธมฺมํ ปฏิจฺจ นอกุสโล นทสฺสเนน ปหาตพฺโพ ธมฺโม อุปฺปชฺชติ เหตุปจฺจยา.  อกุสลํ ทสฺสเนน ปหาตพฺพํ ธมฺมํ ปฏิจฺจ นกุสโล นทสฺสเนน ปหาตพฺโพ ธมฺโม อุปฺปชฺชติ เหตุปจฺจยา.  อกุสลํ ทสฺสเนน ปหาตพฺพํ ธมฺมํ ปฏิจฺจ นกุสโล นทสฺสเนน ปหาตพฺโพ จ นอกุสโล นทสฺสเนน ปหาตพฺโพ จ ธมฺมา อุปฺปชฺชนฺติ เหตุปจฺจยา.  เหตุยา ตีณิ.}

\prose{อพฺยากตํ นทสฺสเนน ปหาตพฺพํ ธมฺมํ ปจฺจยา นอพฺยากโต นนทสฺสเนน ปหาตพฺโพ ธมฺโม อุปฺปชฺชติ เหตุปจฺจยา.  อพฺยากตํ นทสฺสเนน ปหาตพฺพํ ธมฺมํ ปจฺจยา นกุสโล นนทสฺสเนน ปหาตพฺโพ ธมฺโม อุปฺปชฺชติ เหตุปจฺจยา.  อพฺยากตํ นทสฺสเนน ปหาตพฺพํ ธมฺมํ ปจฺจยา นกุสโล นนทสฺสเนน ปหาตพฺโพ จ นอพฺยากโต นนทสฺสเนน ปหาตพฺโพ จ ธมฺมา อุปฺปชฺชนฺติ เหตุปจฺจยา.  เหตุยา ตีณิ.}

\proseitem{48}{อกุสลํ ภาวนาย ปหาตพฺพํ ธมฺมํ ปฏิจฺจ นอกุสโล นภาวนาย ปหาตพฺโพ ธมฺโม อุปฺปชฺชติ เหตุปจฺจยา.  อกุสลํ ภาวนาย ปหาตพฺพํ ธมฺมํ ปฏิจฺจ นกุสโล นภาวนาย ปหาตพฺโพ ธมฺโม อุปฺปชฺชติ เหตุปจฺจยา.  อกุสลํ ภาวนาย ปหาตพฺพํ ธมฺมํ ปฏิจฺจ นกุสโล นภาวนาย ปหาตพฺโพ จ นอกุสโล นภาวนาย ปหาตพฺโพ จ ธมฺมา อุปฺปชฺชนฺติ เหตุปจฺจยา.  เหตุยา ตีณิ.}

\prose{อพฺยากตํ นภาวนาย ปหาตพฺพํ ธมฺมํ ปจฺจยา นอพฺยากโต นนภาวนาย ปหาตพฺโพ ธมฺโม อุปฺปชฺชติ เหตุปจฺจยา.  เหตุยา ตีณิ.}

\proseitem{49}{อกุสลํ ทสฺสเนน ปหาตพฺพ\-เหตุกํ ธมฺมํ ปฏิจฺจ นอกุสโล นทสฺสเนน ปหาตพฺพ\-เหตุโก ธมฺโม อุปฺปชฺชติ เหตุปจฺจยา.  เหตุยา ตีณิ.}

\prose{อกุสลํ นทสฺสเนน ปหาตพฺพ\-เหตุกํ ธมฺมํ ปฏิจฺจ นกุสโล นนทสฺสเนน ปหาตพฺพ\-เหตุโก ธมฺโม อุปฺปชฺชติ เหตุปจฺจยา.  เหตุยา ตีณิ.}

\csromanmarkh{1. กุสลตฺติก}\csromanmarkhii{83. ปิฏฺฐิทุก}\csromanheadernomark{2}{1. กุสลตฺติก 83. ปิฏฺฐิทุก}
\prose{\csromanfolio{263}อพฺยากตํ นทสฺสเนน ปหาตพฺพ\-เหตุกํ ธมฺมํ ปจฺจยา นอพฺยากโต นนทสฺสเนน ปหาตพฺพ\-เหตุโก ธมฺโม อุปฺปชฺชติ เหตุปจฺจยา.  เหตุยา ตีณิ.}

\proseitem{50}{อกุสลํ ภาวนาย ปหาตพฺพ\-เหตุกํ ธมฺมํ ปฏิจฺจ นอกุสโล นภาวนาย ปหาตพฺพ\-เหตุโก ธมฺโม อุปฺปชฺชติ เหตุปจฺจยา.  เหตุยา ตีณิ.}

\prose{อกุสลํ นภาวนาย ปหาตพฺพ\-เหตุกํ ธมฺมํ ปฏิจฺจ นกุสโล นนภาวนาย ปหาตพฺพ\-เหตุโก ธมฺโม อุปฺปชฺชติ เหตุปจฺจยา.  เหตุยา ตีณิ.}

\proseitem{51}{กุสลํ สวิตกฺกํ ธมฺมํ ปฏิจฺจ นกุสโล นสวิตกฺโก ธมฺโม อุปฺปชฺชติ เหตุปจฺจยา.  เหตุยา เตรส.}

\prose{กุสลํ อวิตกฺกํ ธมฺมํ ปฏิจฺจ นอกุสโล นอวิตกฺโก ธมฺโม อุปฺปชฺชติ เหตุปจฺจยา.  เหตุยา นว.}

\proseitem{52}{กุสลํ สวิจารํ ธมฺมํ ปฏิจฺจ นกุสโล นสวิจาโร ธมฺโม อุปฺปชฺชติ เหตุปจฺจยา.  เหตุยา เตรส.}

\prose{กุสลํ อวิจารํ ธมฺมํ ปฏิจฺจ นอกุสโล นอวิจาโร ธมฺโม อุปฺปชฺชติ เหตุปจฺจยา.  เหตุยา นว.}

\proseitem{53}{กุสลํ สปฺปีติกํ ธมฺมํ ปฏิจฺจ นกุสโล นสปฺปีติโก ธมฺโม อุปฺปชฺชติ เหตุปจฺจยา.  เหตุยา เตรส.}

\prose{กุสลํ อปฺปีติกํ ธมฺมํ ปฏิจฺจ นอกุสโล นอปฺปีติโก ธมฺโม อุปฺปชฺชติ เหตุปจฺจยา.  เหตุยา นว.}

\proseitem{54}{กุสลํ ปีติสหคตํ ธมฺมํ ปฏิจฺจ นกุสโล นปีติสหคโต ธมฺโม อุปฺปชฺชติ เหตุปจฺจยา.  เหตุยา เตรส.}

\prose{กุสลํ นปีติสหคตํ ธมฺมํ ปฏิจฺจ นอกุสโล นนปี\-ติ\-สหคโต ธมฺโม อุปฺปชฺชติ เหตุปจฺจยา.  เหตุยา นว.}

\proseitem{55}{กุสลํ สุขสหคตํ ธมฺมํ ปฏิจฺจ นกุสโล นสุขสหคโต ธมฺโม อุปฺปชฺชติ เหตุปจฺจยา.  เหตุยา เตรส.}

\csromanheader{2}{ธมฺมา\-นุโลม\-ปจฺจนีย ติกทุก\-ปฏฺฐาน\-ปาฬิ}
\prose{\csromanfolio{264}กุสลํ นสุขสหคตํ ธมฺมํ ปฏิจฺจ นอกุสโล นนสุข\-สหคโต ธมฺโม อุปฺปชฺชติ เหตุปจฺจยา.  เหตุยา นว.}

\proseitem{56}{กุสลํ อุเปกฺขา\-สหคตํ ธมฺมํ ปฏิจฺจ นกุสโล นอุเปกฺขาสหคโต ธมฺโม อุปฺปชฺชติ เหตุปจฺจยา.  เหตุยา เตรส.}

\prose{กุสลํ นอุเปกฺขาสหคตํ ธมฺมํ ปฏิจฺจ นกุสโล นนอุเปกฺขาสหคโต ธมฺโม อุปฺปชฺชติ เหตุปจฺจยา.  เหตุยา นว.}

\proseitem{57}{อพฺยากตํ กามาวจรํ ธมฺมํ ปฏิจฺจ นกุสโล นกามาวจโร ธมฺโม อุปฺปชฺชติ เหตุปจฺจยา.  เหตุยา ตีณิ.}

\prose{กุสลํ นกามาวจรํ ธมฺมํ ปฏิจฺจ นกุสโล นนกามาวจโร ธมฺโม อุปฺปชฺชติ เหตุปจฺจยา.  เหตุยา ฉ.}

\proseitem{58}{กุสลํ รูปาวจรํ ธมฺมํ ปฏิจฺจ นกุสโล นรูปาวจโร ธมฺโม อุปฺปชฺชติ เหตุปจฺจยา.  เหตุยา ฉ.}

\prose{อพฺยากตํ นรูปาวจรํ ธมฺมํ ปฏิจฺจ นกุสโล นนรูปาวจโร ธมฺโม อุปฺปชฺชติ เหตุปจฺจยา.  เหตุยา ตีณิ.}

\proseitem{59}{กุสลํ อรูปาวจรํ ธมฺมํ ปฏิจฺจ นกุสโล นอรูปาวจโร ธมฺโม อุปฺปชฺชติ เหตุปจฺจยา.  เหตุยา ฉ.}

\prose{อพฺยากตํ นอรูปาวจรํ ธมฺมํ ปจฺจยา นอพฺยากโต นนอรูปาวจโร ธมฺโม อุปฺปชฺชติ เหตุปจฺจยา.  เหตุยา ปญฺจ.}

\proseitem{60}{อพฺยากตํ ปริยาปนฺนํ ธมฺมํ ปจฺจยา นอพฺยากโต นปริยาปนฺโน ธมฺโม อุปฺปชฺชติ เหตุปจฺจยา.  เหตุยา ปญฺจ.}

\prose{กุสลํ อปริยาปนฺนํ ธมฺมํ ปฏิจฺจ นกุสโล นอปริยาปนฺโน ธมฺโม อุปฺปชฺชติ เหตุปจฺจยา.  เหตุยา ฉ.}

\proseitem{61}{กุสลํ นิยฺยานิกํ ธมฺมํ ปฏิจฺจ นกุสโล นนิยฺยานิโก ธมฺโม อุปฺปชฺชติ เหตุปจฺจยา.  เหตุยา ตีณิ.}

\prose{อพฺยากตํ อนิยฺยานิกํ ธมฺมํ ปจฺจยา นอพฺยากโต นอนิยฺยานิโก ธมฺโม อุปฺปชฺชติ เหตุปจฺจยา.  เหตุยา ตีณิ.}

\csromanmarkh{2. เวทนาตฺติก}\csromanmarkhii{1. เหตุทุก}\csromanheadernomark{2}{2. เวทนาตฺติก 1. เหตุทุก}
\proseitem{62}{\csromanfolio{265}กุสลํ นิยตํ ธมฺมํ ปฏิจฺจ นกุสโล นนิยโต ธมฺโม อุปฺปชฺชติ เหตุปจฺจยา.  เหตุยา ฉ.}

\prose{อพฺยากตํ อนิยตํ ธมฺมํ ปจฺจยา นอพฺยากโต นอนิยโต ธมฺโม อุปฺปชฺชติ เหตุปจฺจยา.  เหตุยา ปญฺจ.}

\proseitem{63}{อพฺยากตํ สอุตฺตรํ ธมฺมํ ปจฺจยา นอพฺยากโต นสอุตฺตโร ธมฺโม อุปฺปชฺชติ เหตุปจฺจยา.  เหตุยา ปญฺจ.}

\prose{กุสลํ อนุตฺตรํ ธมฺมํ ปฏิจฺจ นกุสโล นอนุตฺตโร ธมฺโม อุปฺปชฺชติ เหตุปจฺจยา.  เหตุยา ฉ.}

\proseitem{64}{อกุสลํ สรณํ ธมฺมํ ปฏิจฺจ นอกุสโล นสรโณ ธมฺโม อุปฺปชฺชติ เหตุปจฺจยา.  อกุสลํ สรณํ ธมฺมํ ปฏิจฺจ นกุสโล นสรโณ ธมฺโม อุปฺปชฺชติ เหตุปจฺจยา.  อกุสลํ สรณํ ธมฺมํ ปฏิจฺจ นกุสโล นสรโณ จ นอกุสโล นสรโณ จ ธมฺมา อุปฺปชฺชนฺติ เหตุปจฺจยา.  เหตุยา ตีณิ.}

\prose{อพฺยากตํ อรณํ ธมฺมํ ปจฺจยา นอพฺยากโต นอรโณ ธมฺโม อุปฺปชฺชติ เหตุปจฺจยา.  อพฺยากตํ อรณํ ธมฺมํ ปจฺจยา นกุสโล นอรโณ ธมฺโม อุปฺปชฺชติ เหตุปจฺจยา.  อพฺยากตํ อรณํ ธมฺมํ ปจฺจยา นกุสโล นอรโณ จ นอพฺยากโต นอรโณ จ ธมฺมา อุปฺปชฺชนฺติ เหตุปจฺจยา.  เหตุยา ตีณิ.}

\csromanmarkh{2. เวทนาตฺติก}\csromanmarkhii{1. เหตุทุก}\csromanheadernomark{2}{2. เวทนาตฺติก 1. เหตุทุก}
\prosecont{สุขาย เวทนาย สมฺปยุตฺตํ เหตุํ ธมฺมํ ปฏิจฺจ นสุขาย เวทนาย สมฺปยุตฺโต นเหตุ ธมฺโม อุปฺปชฺชติ เหตุปจฺจยา.  สุขาย เวทนาย สมฺปยุตฺตํ เหตุํ ธมฺมํ ปฏิจฺจ นทุกฺขาย เวทนาย สมฺปยุตฺโต นเหตุ ธมฺโม อุปฺปชฺชติ เหตุปจฺจยา.  สุขาย เวทนาย สมฺปยุตฺตํ เหตุํ ธมฺมํ ปฏิจฺจ นอทุกฺขม\-สุขาย เวทนาย สมฺปยุตฺโต นเหตุ ธมฺโม อุปฺปชฺชติ เหตุปจฺจยา ฯเปฯ. สตฺต ปญฺหา.}

\proseitem{65}{ทุกฺขาย เวทนาย สมฺปยุตฺตํ เหตุํ ธมฺมํ ปฏิจฺจ นทุกฺขาย เวทนาย สมฺปยุตฺโต นเหตุ ธมฺโม อุปฺปชฺชติ เหตุปจฺจยา. สตฺต ปญฺหา.}

\csromanheader{2}{ธมฺมา\-นุโลม\-ปจฺจนีย ติกทุก\-ปฏฺฐาน\-ปาฬิ}
\prose{\csromanfolio{266}อทุกฺขม\-สุขาย เวทนาย สมฺปยุตฺตํ เหตุํ ธมฺมํ ปฏิจฺจ นอทุกฺขม\-สุขาย เวทนาย สมฺปยุตฺโต นเหตุ ธมฺโม อุปฺปชฺชติ เหตุปจฺจยา. สตฺต ปญฺหา.  เหตุยา เอกวีส ฯเปฯ อวิคเต เอกวีส. (สพฺพตฺถ เอกวีส.)}

\proseitem{66}{สุขาย เวทนาย สมฺปยุตฺตํ นเหตุํ ธมฺมํ ปฏิจฺจ นทุกฺขาย เวทนาย สมฺปยุตฺโต นนเหตุ ธมฺโม อุปฺปชฺชติ เหตุปจฺจยา.  สุขาย เวทนาย สมฺปยุตฺตํ นเหตุํ ธมฺมํ ปฏิจฺจ นอทุกฺขม\-สุขาย เวทนาย สมฺปยุตฺโต นนเหตุ ธมฺโม อุปฺปชฺชติ เหตุปจฺจยา.  สุขาย เวทนาย สมฺปยุตฺตํ นเหตุํ ธมฺมํ ปฏิจฺจ นทุกฺขาย เวทนาย สมฺปยุตฺโต นนเหตุ จ นอทุกฺขม\-สุขาย เวทนาย สมฺปยุตฺโต นนเหตุ จ ธมฺมา อุปฺปชฺชนฺติ เหตุปจฺจยา. ตีณิ.}

\prose{ทุกฺขาย เวทนาย สมฺปยุตฺตํ นเหตุํ ธมฺมํ ปฏิจฺจ นสุขาย เวทนาย สมฺปยุตฺโต นนเหตุ ธมฺโม อุปฺปชฺชติ เหตุปจฺจยา. ตีณิ.}

\prose{อทุกฺขม\-สุขาย เวทนาย สมฺปยุตฺตํ นเหตุํ ธมฺมํ ปฏิจฺจ นสุขาย เวทนาย สมฺปยุตฺโต นนเหตุ ธมฺโม อุปฺปชฺชติ เหตุปจฺจยา. ตีณิ.  เหตุยา นว.}

\csromanmarkh{3. วิปากตฺติก}\csromanmarkhii{1. เหตุทุก}\csromanheadernomark{2}{3. \textbf{วิปากตฺติก 1. เหตุทุก}}
\proseitem{67}{วิปากํ เหตุํ ธมฺมํ ปฏิจฺจ นวิปาโก นเหตุ ธมฺโม อุปฺปชฺชติ เหตุปจฺจยา.  วิปากํ เหตุํ ธมฺมํ ปฏิจฺจ นวิปาก\-ธมฺม\-ธมฺโม นเหตุ ธมฺโม อุปฺปชฺชติ เหตุปจฺจยา.  วิปากํ เหตุํ ธมฺมํ ปฏิจฺจ นเน\-ว\-วิปาก\-น\-วิปาก\-ธมฺม\-ธมฺโม นเหตุ ธมฺโม อุปฺปชฺชติ เหตุปจฺจยา.  วิปากํ เหตุํ ธมฺมํ ปฏิจฺจ นวิปาก\-ธมฺม\-ธมฺโม นเหตุ จ นเน\-ว\-วิปาก\-น\-วิปาก\-ธมฺม\-ธมฺโม นเหตุ จ ธมฺมา อุปฺปชฺชนฺติ เหตุปจฺจยา.  วิปากํ เหตุํ ธมฺมํ ปฏิจฺจ นวิปาโก นเหตุ จ นวิปาก\-ธมฺม\-ธมฺโม นเหตุ จ ธมฺมา อุปฺปชฺชนฺติ เหตุปจฺจยา. ปญฺจ. วิปาก\-ธมฺม\-ธมฺมํ เหตุํ ธมฺมํ ปฏิจฺจ นวิปาก\-ธมฺม\-ธมฺโม นเหตุ ธมฺโม อุปฺปชฺชติ เหตุปจฺจยา.  วิปาก\-ธมฺม\-ธมฺมํ เหตุํ ธมฺมํ ปฏิจฺจ นวิปาโก นเหตุ ธมฺโม อุปฺปชฺชติ เหตุปจฺจยา.  วิปาก\-ธมฺม\-ธมฺมํ เหตุํ ธมฺมํ ปฏิจฺจ นเน\-ว\-วิปาก\-น\-วิปาก\-ธมฺม\-ธมฺโม นเหตุ ธมฺโม อุปฺปชฺชติ เหตุปจฺจยา.  วิปาก\-ธมฺม\-ธมฺมํ เหตุํ ธมฺมํ ปฏิจฺจ นวิปาโก นเหตุ จ นเน\-ว\-วิปาก\-น\-วิปาก\-ธมฺม\-ธมฺโม นเหตุ จ ธมฺมา อุปฺปชฺชนฺติ เหตุปจฺจยา.  วิปาก\-ธมฺม\-ธมฺมํ เหตุํ ธมฺมํ ปฏิจฺจ}

\vspace{\tipitakaparskip}
\csromancenter{อุปาทินฺนตฺติก 1. เหตุทุก นวิปาโก นเหตุ จ นวิปาก\-ธมฺม\-ธมฺโม นเหตุ จ ธมฺมา อุปฺปชฺชนฺติ เหตุปจฺจยา. ปญฺจ.}
\prosecont{\csromanfolio{267}เนว\-วิปาก\-น\-วิปากธมฺม\-ธมฺมํ เหตุํ ธมฺมํ ปฏิจฺจ นวิปาโก นเหตุ ธมฺโม อุปฺปชฺชติ เหตุปจฺจยา.  เนว\-วิปาก\-น\-วิปากธมฺม\-ธมฺมํ เหตุํ ธมฺมํ ปฏิจฺจ นวิปาก\-ธมฺม\-ธมฺโม นเหตุ ธมฺโม อุปฺปชฺชติ เหตุปจฺจยา.  เนว\-วิปาก\-น\-วิปากธมฺม\-ธมฺมํ เหตุํ ธมฺมํ ปฏิจฺจ นวิปาโก นเหตุ จ นวิปาก\-ธมฺม\-ธมฺโม นเหตุ จ ธมฺมา อุปฺปชฺชนฺติ เหตุปจฺจยา. ตีณิ.  เหตุยา เตรส.}

\proseitem{68}{วิปากํ นเหตุํ ธมฺมํ ปฏิจฺจ นวิปาก\-ธมฺม\-ธมฺโม นนเหตุ ธมฺโม อุปฺปชฺชติ เหตุปจฺจยา.  วิปากํ นเหตุํ ธมฺมํ ปฏิจฺจ นเน\-ว\-วิปาก\-น\-วิปาก\-ธมฺม\-ธมฺโม นนเหตุ ธมฺโม อุปฺปชฺชติ เหตุปจฺจยา.  วิปากํ นเหตุํ ธมฺมํ ปฏิจฺจ นวิปาก\-ธมฺม\-ธมฺโม นนเหตุ จ นเน\-ว\-วิปาก\-น\-วิปาก\-ธมฺม\-ธมฺโม นนเหตุ จ ธมฺมา อุปฺปชฺชนฺติ เหตุปจฺจยา. ตีณิ.}

\prose{วิปาก\-ธมฺม\-ธมฺมํ นเหตุํ ธมฺมํ ปฏิจฺจ นวิปาโก นนเหตุ ธมฺโม อุปฺปชฺชติ เหตุปจฺจยา.  วิปาก\-ธมฺม\-ธมฺมํ นเหตุํ ธมฺมํ ปฏิจฺจ นเน\-ว\-วิปาก\-น\-วิปาก\-ธมฺม\-ธมฺโม นนเหตุ ธมฺโม อุปฺปชฺชติ เหตุปจฺจยา.  วิปาก\-ธมฺม\-ธมฺมํ นเหตุํ ธมฺมํ ปฏิจฺจ นวิปาโก นนเหตุ จ นเน\-ว\-วิปาก\-น\-วิปาก\-ธมฺม\-ธมฺโม นนเหตุ จ ธมฺมา อุปฺปชฺชนฺติ เหตุปจฺจยา. ตีณิ.}

\prose{เนว\-วิปาก\-น\-วิปากธมฺม\-ธมฺมํ นเหตุํ ธมฺมํ ปฏิจฺจ นเน\-ว\-วิปาก\-น\-วิปาก\-ธมฺม\-ธมฺโม นนเหตุ ธมฺโม อุปฺปชฺชติ เหตุปจฺจยา.  เนว\-วิปาก\-น\-วิปากธมฺม\-ธมฺมํ นเหตุํ ธมฺมํ ปฏิจฺจ นวิปาโก นนเหตุ ธมฺโม อุปฺปชฺชติ เหตุปจฺจยา ฯเปฯ. เนว\-วิปาก\-น\-วิปากธมฺม\-ธมฺมํ นเหตุํ ธมฺมํ ปฏิจฺจ นวิปาโก นนเหตุ จ นวิปาก\-ธมฺม\-ธมฺโม นนเหตุ จ ธมฺมา อุปฺปชฺชนฺติ เหตุปจฺจยา. ปญฺจ.}

\prose{วิปากํ นเหตุญฺจ เนว\-วิปาก\-น\-วิปากธมฺม\-ธมฺมํ นเหตุญฺจ ธมฺมํ ปฏิจฺจ นวิปาก\-ธมฺม\-ธมฺโม นนเหตุ ธมฺโม อุปฺปชฺชติ เหตุปจฺจยา. ตีณิ.  เหตุยา จุทฺทส.}

\csromanmarkh{4. อุปาทินฺนตฺติก}\csromanmarkhii{1. เหตุทุก}\csromanheadernomark{2}{4. อุปาทินฺนตฺติก 1. เหตุทุก}
\proseitem{69}{อุปาทินฺนุ\-ปาทานิยํ เหตุํ ธมฺมํ ปฏิจฺจ นอุปาทินฺนุปาทานิโย นเหตุ ธมฺโม อุปฺปชฺชติ เหตุปจฺจยา.  อุปาทินฺนุ\-ปาทานิยํ เหตุํ ธมฺมํ ปฏิจฺจ นอนุปาทินฺนุ\-ปาทานิโย นเหตุ ธมฺโม อุปฺปชฺชติ เหตุปจฺจยา.  }

\vspace{\tipitakaparskip}
\csromancenter{ติกทุก\-ปฏฺฐาน\-ปาฬิ อุปาทินฺนุ\-ปาทานิยํ เหตุํ ธมฺมํ ปฏิจฺจ นอนุปาทินฺน\-อนุปาทานิโย นเหตุ ธมฺโม อุปฺปชฺชติ เหตุปจฺจยา.  อุปาทินฺนุ\-ปาทานิยํ เหตุํ ธมฺมํ ปฏิจฺจ นฺเอาปาทินฺนุ\-ปาทานิโย นเหตุ จ นอนุปาทินฺน\-อนุปาทานิโย นเหตุ จ ธมฺมา อุปฺปชฺชนฺติ เหตุปจฺจยา.  อุปาทินฺนุ\-ปาทานิยํ เหตุํ ธมฺมํ ปฏิจฺจ นอนุปาทินฺนุ\-ปาทานิโย นเหตุ จ นอนุปาทินฺน\-อนุปาทานิโย นเหตุ จ ธมฺมา อุปฺปชฺชนฺติ เหตุปจฺจยา. ปญฺจ.}
\prose{\csromanfolio{268}อนุปาทินฺนุ\-ปาทานิยํ เหตุํ ธมฺมํ ปฏิจฺจ นอุปาทินฺนุปาทานิโย นเหตุ ธมฺโม อุปฺปชฺชติ เหตุปจฺจยา. ตีณิ.}

\prose{อนุปาทินฺน\-อนุปาทานิยํ เหตุํ ธมฺมํ ปฏิจฺจ นอนุปาทินฺน\-อนุปาทานิโย นเหตุ ธมฺโม อุปฺปชฺชติ เหตุปจฺจยา. ปญฺจ ปญฺหา.  เหตุยา เตรส.}

\proseitem{70}{อุปาทินฺนุ\-ปาทานิยํ นเหตุํ ธมฺมํ ปฏิจฺจ นอนุปาทินฺนุ\-ปาทานิโย นนเหตุ ธมฺโม อุปฺปชฺชติ เหตุปจฺจยา. ตีณิ.  เหตุยา นว.}

\csromanmarkh{5. สํกิลิฏฺฐ\-ตฺติก}\csromanmarkhii{1. เหตุทุก}\csromanheadernomark{2}{5. สํกิลิฏฺฐ\-ตฺติก 1. เหตุทุก}
\proseitem{71}{สํกิลิฏฺฐ\-สํกิเลสิกํ เหตุํ ธมฺมํ ปฏิจฺจ นสํกิลิฏฺฐ\-สํกิเลสิโก นเหตุ ธมฺโม อุปฺปชฺชติ เหตุปจฺจยา.  เหตุยา เตรส.}

\prose{สํกิลิฏฺฐ\-สํกิเลสิกํ นเหตุํ ธมฺมํ ปฏิจฺจ นอสํกิลิฏฺฐ\-สํกิเลสิโก นนเหตุ ธมฺโม อุปฺปชฺชติ เหตุปจฺจยา. ตีณิ.  เหตุยา นว.}

\csromanmarkh{6. วิตกฺกตฺติก}\csromanmarkhii{1. เหตุทุก}\csromanheadernomark{2}{6. \textbf{วิตกฺกตฺติก 1. เหตุทุก}}
\proseitem{72}{สวิตกฺก\-สวิจารํ เหตุํ ธมฺมํ ปฏิจฺจ นสวิตกฺก\-สวิจาโร นเหตุ ธมฺโม อุปฺปชฺชติ เหตุปจฺจยา.  เหตุยา ปนฺนรส.}

\prose{สวิตกฺก\-สวิจารํ นเหตุํ ธมฺมํ ปฏิจฺจ นอวิตกฺก\-วิจารมตฺโต นนเหตุ ธมฺโม อุปฺปชฺชติ เหตุปจฺจยา.  เหตุยา อฏฺฐวีส.}

\csromanmarkh{7. ปีติตฺติก}\csromanmarkhii{1. เหตุทุก}\csromanheadernomark{2}{7. \textbf{ปีติตฺติก 1. เหตุทุก}}
\proseitem{73}{ปีติสหคตํ เหตุํ ธมฺมํ ปฏิจฺจ นปีติสหคโต นเหตุ ธมฺโม อุปฺปชฺชติ เหตุปจฺจยา.  เหตุยา อฏฺฐวีส.}

\csromanmarkh{11. เสกฺขตฺติก}\csromanmarkhii{1. เหตุทุก}\csromanheadernomark{2}{11. เสกฺขตฺติก 1. เหตุทุก}
\prose{\csromanfolio{269}ปีติสหคตํ นเหตุํ ธมฺมํ ปฏิจฺจ นอุเปกฺขาสหคโต นนเหตุ ธมฺโม อุปฺปชฺชติ เหตุปจฺจยา.  เหตุยา อฏฺฐารส\csromansharedfootnote{e0ee4924e42dd2cb}{เหตุยา อฏฺฐ?}.}

\csromanmarkh{8. ทสฺสนตฺติก}\csromanmarkhii{1. เหตุทุก}\csromanheadernomark{2}{8. \textbf{ทสฺสนตฺติก 1. เหตุทุก}}
\proseitem{74}{ทสฺสเนน ปหาตพฺพํ เหตุํ ธมฺมํ ปฏิจฺจ นทสฺสเนน ปหาตพฺโพ นเหตุ ธมฺโม อุปฺปชฺชติ เหตุปจฺจยา.  เหตุยา เตรส.}

\prose{ทสฺสเนน ปหาตพฺพํ นเหตุํ ธมฺมํ ปฏิจฺจ นภาวนาย ปหาตพฺโพ นนเหตุ ธมฺโม อุปฺปชฺชติ เหตุปจฺจยา.  เหตุยา นว.}

\csromanmarkh{9. ทสฺสน\-เหตุ\-ตฺติก}\csromanmarkhii{1. เหตุทุก}\csromanheadernomark{2}{9. ทสฺสน\-เหตุ\-ตฺติก 1. เหตุทุก}
\proseitem{75}{ทสฺสเนน ปหาตพฺพ\-เหตุกํ เหตุํ ธมฺมํ ปฏิจฺจ นทสฺสเนน ปหาตพฺพ\-เหตุโก นเหตุ ธมฺโม อุปฺปชฺชติ เหตุปจฺจยา.  เหตุยา โสฬส.}

\prose{ทสฺสเนน ปหาตพฺพ\-เหตุกํ นเหตุํ ธมฺมํ ปฏิจฺจ นภาวนาย ปหาตพฺพ\-เหตุโก นนเหตุ ธมฺโม อุปฺปชฺชติ เหตุปจฺจยา.  เหตุยา นว.}

\csromanmarkh{10. อาจยคามิตฺติก}\csromanmarkhii{1. เหตุทุก}\csromanheadernomark{2}{10. \textbf{อาจยคามิตฺติก 1. เหตุทุก}}
\proseitem{76}{อาจยคามิํ เหตุํ ธมฺมํ ปฏิจฺจ นอาจยคามี นเหตุ ธมฺโม อุปฺปชฺชติ เหตุปจฺจยา.  เหตุยา เตรส.}

\prose{อาจยคามิํ นเหตุํ ธมฺมํ ปฏิจฺจ นอปจยคามี นนเหตุ ธมฺโม อุปฺปชฺชติ เหตุปจฺจยา.  เหตุยา นว.}

\csromanmarkh{11. เสกฺขตฺติก}\csromanmarkhii{1. เหตุทุก}\csromanheadernomark{2}{11. เสกฺขตฺติก 1. เหตุทุก}
\proseitem{77}{เสกฺขํ เหตุํ ธมฺมํ ปฏิจฺจ นเสกฺโข นเหตุ ธมฺโม อุปฺปชฺชติ เหตุปจฺจยา.  เหตุยา เตรส.}

\prose{เสกฺขํ นเหตุํ ธมฺมํ ปฏิจฺจ นอเสกฺโข นนเหตุ ธมฺโม อุปฺปชฺชติ เหตุปจฺจยา.  เหตุยา นว.}

\csromanmarkh{ธมฺมา\-นุโลม\-ปจฺจนีย ติกทุก\-ปฏฺฐาน\-ปาฬิ}\csromanmarkhii{12. ปริตฺตตฺติก 1. เหตุทุก}\csromanheadernomark{2}{ธมฺมา\-นุโลม\-ปจฺจนีย ติกทุก\-ปฏฺฐาน\-ปาฬิ 12. ปริตฺตตฺติก 1. เหตุทุก}
\proseitem{78}{\csromanfolio{270}ปริตฺตํ เหตุํ ธมฺมํ ปฏิจฺจ นมหคฺคโต นเหตุ ธมฺโม อุปฺปชฺชติ เหตุปจฺจยา.  เหตุยา เตรส.}

\prose{ปริตฺตํ นเหตุํ ธมฺมํ ปฏิจฺจ นปริตฺโต นนเหตุ ธมฺโม อุปฺปชฺชติ เหตุปจฺจยา.  เหตุยา จุทฺทส.}

\csromanmarkh{13. ปริตฺตา\-รมฺมณ\-ตฺติก}\csromanmarkhii{1. เหตุทุก}\csromanheadernomark{2}{13. ปริตฺตา\-รมฺมณ\-ตฺติก 1. เหตุทุก}
\proseitem{79}{ปริตฺตา\-รมฺมณํ เหตุํ ธมฺมํ ปฏิจฺจ นปริตฺตา\-รมฺมโณ นเหตุ ธมฺโม อุปฺปชฺชติ เหตุปจฺจยา.  เหตุยา เอกวีส.}

\prose{ปริตฺตา\-รมฺมณํ นเหตุํ ธมฺมํ ปฏิจฺจ นมหคฺคตา\-รมฺมโณ นนเหตุ ธมฺโม อุปฺปชฺชติ เหตุปจฺจยา.  เหตุยา นว.}

\csromanmarkh{14. หีนตฺติก}\csromanmarkhii{1. เหตุทุก}\csromanheadernomark{2}{14. \textbf{หีนตฺติก 1. เหตุทุก}}
\proseitem{80}{หีนํ เหตุํ ธมฺมํ ปฏิจฺจ นหีโน นเหตุ ธมฺโม อุปฺปชฺชติ เหตุปจฺจยา.  เหตุยา เตรส.}

\prose{หีนํ นเหตุํ ธมฺมํ ปฏิจฺจ นมชฺฌิโม นนเหตุ ธมฺโม อุปฺปชฺชติ เหตุปจฺจยา.  เหตุยา นว.}

\csromanmarkh{15. มิจฺฉตฺตนิยต\-ตฺติก}\csromanmarkhii{1. เหตุทุก}\csromanheadernomark{2}{15. มิจฺฉตฺตนิยต\-ตฺติก 1. เหตุทุก}
\proseitem{81}{มิจฺฉตฺต\-นิยตํ เหตุํ ธมฺมํ ปฏิจฺจ นมิจฺฉตฺต\-นิยโต นเหตุ ธมฺโม อุปฺปชฺชติ เหตุปจฺจยา.  เหตุยา เตรส.}

\prose{มิจฺฉตฺต\-นิยตํ นเหตุํ ธมฺมํ ปฏิจฺจ นสมฺมตฺตนิยโต นนเหตุ ธมฺโม อุปฺปชฺชติ เหตุปจฺจยา.  เหตุยา นว.}

\csromanmarkh{16. มคฺคา\-รมฺมณ\-ตฺติก}\csromanmarkhii{1. เหตุทุก}\csromanheadernomark{2}{16. มคฺคา\-รมฺมณ\-ตฺติก 1. เหตุทุก}
\proseitem{82}{มคฺคารมฺมณํ เหตุํ ธมฺมํ ปฏิจฺจ นมคฺคารมฺมโณ นเหตุ ธมฺโม อุปฺปชฺชติ เหตุปจฺจยา.  เหตุยา ปญฺจวีส.}

\prose{มคฺคารมฺมณํ นเหตุํ ธมฺมํ ปฏิจฺจ นมคฺคเหตุโก นนเหตุ ธมฺโม อุปฺปชฺชติ เหตุปจฺจยา.  เหตุยา เตรส\csromansharedfootnote{e0ee4924e42dd2cb}{เหตุยา อฏฺฐ?}.}

\vspace{\tipitakaparskip}
\csromancenter{อชฺฌตฺตา\-รมฺมณ\-ตฺติก 1. เหตุทุก 17. อุปฺปนฺนตฺติก 1. เหตุทุก}
\proseitem{83}{\csromanfolio{271}อุปฺปนฺโน เหตุ ธมฺโม นอนุปฺปนฺนสฺส นเหตุสฺส ธมฺมสฺส เหตุปจฺจเยน ปจฺจโย.  อุปฺปนฺโน เหตุ ธมฺโม นอุปฺปาทิสฺส นเหตุสฺส ธมฺมสฺส เหตุปจฺจเยน ปจฺจโย.  อุปฺปนฺโน เหตุ ธมฺโม นอนุปฺปนฺนสฺส นเหตุสฺส จ นอุปฺปาทิสฺส นเหตุสฺส จ ธมฺมสฺส เหตุปจฺจเยน ปจฺจโย.  เหตุยา ตีณิ.}

\csromanmarkh{18. อตีตตฺติก}\csromanmarkhii{1. เหตุทุก}\csromanheadernomark{2}{18. \textbf{อตีตตฺติก 1. เหตุทุก}}
\proseitem{84}{ปจฺจุปฺปนฺโน เหตุ ธมฺโม นอตีตสฺส นเหตุสฺส ธมฺมสฺส เหตุปจฺจเยน ปจฺจโย.  ปจฺจุปฺปนฺโน เหตุ ธมฺโม นอนาคตสฺส นเหตุสฺส ธมฺมสฺส เหตุปจฺจเยน ปจฺจโย.  ปจฺจุปฺปนฺโน เหตุ ธมฺโม นอตีตสฺส นเหตุสฺส จ นอนาคตสฺส นเหตุสฺส จ ธมฺมสฺส เหตุปจฺจเยน ปจฺจโย.  เหตุยา ตีณิ.}

\csromanmarkh{19. อตีตา\-รมฺมณ\-ตฺติก}\csromanmarkhii{1. เหตุทุก}\csromanheadernomark{2}{19. อตีตา\-รมฺมณ\-ตฺติก 1. เหตุทุก}
\proseitem{85}{อตีตารมฺมณํ เหตุํ ธมฺมํ ปฏิจฺจ นอตีตารมฺมโณ นเหตุ ธมฺโม อุปฺปชฺชติ เหตุปจฺจยา.  เหตุยา เอกวีส.}

\prose{อตีตารมฺมณํ นเหตุํ ธมฺมํ ปฏิจฺจ นอนาคตา\-รมฺมโณ นนเหตุ ธมฺโม อุปฺปชฺชติ เหตุปจฺจยา.  เหตุยา นว.}

\csromanmarkh{20. อชฺฌตฺตตฺติก}\csromanmarkhii{1. เหตุทุก}\csromanheadernomark{2}{20. \textbf{อชฺฌตฺตตฺติก 1. เหตุทุก}}
\proseitem{86}{อชฺฌตฺตํ เหตุํ ธมฺมํ ปฏิจฺจ นพหิทฺธา นเหตุ ธมฺโม อุปฺปชฺชติ เหตุปจฺจยา.  พหิทฺธา เหตุํ ธมฺมํ ปฏิจฺจ นอชฺฌตฺโต นเหตุ ธมฺโม อุปฺปชฺชติ เหตุปจฺจยา.  เหตุยา เทฺว.}

\prose{อชฺฌตฺตํ นเหตุํ ธมฺมํ ปฏิจฺจ นพหิทฺธา นนเหตุ ธมฺโม อุปฺปชฺชติ เหตุปจฺจยา.  พหิทฺธา นเหตุํ ธมฺมํ ปฏิจฺจ นอชฺฌตฺโต นนเหตุ ธมฺโม อุปฺปชฺชติ เหตุปจฺจยา.  เหตุยา เทฺว.}

\csromanmarkh{21. อชฺฌตฺตา\-รมฺมณ\-ตฺติก}\csromanmarkhii{1. เหตุทุก}\csromanheadernomark{2}{21. อชฺฌตฺตา\-รมฺมณ\-ตฺติก 1. เหตุทุก}
\proseitem{87}{อชฺฌตฺตา\-รมฺมณํ เหตุํ ธมฺมํ ปฏิจฺจ นอชฺฌตฺตา\-รมฺมโณ นเหตุ ธมฺโม อุปฺปชฺชติ เหตุปจฺจยา.  เหตุยา ฉ.}

\csromanheader{2}{ธมฺมา\-นุโลม\-ปจฺจนีย ติกทุก\-ปฏฺฐาน\-ปาฬิ}
\prose{\csromanfolio{272}อชฺฌตฺตา\-รมฺมณํ นเหตุํ ธมฺมํ ปฏิจฺจ นพหิทฺธา\-รมฺมโณ นนเหตุ ธมฺโม อุปฺปชฺชติ เหตุปจฺจยา.  เหตุยา เทฺว. 22. สนิทสฺสน\-ตฺติก 1. เหตุทุก}

\proseitem{88}{อนิทสฺสน\-อปฺปฏิฆํ เหตุํ ธมฺมํ ปฏิจฺจ นอนิทสฺสน\-อปฺปฏิโฆ นเหตุ ธมฺโม อุปฺปชฺชติ เหตุปจฺจยา.  อนิทสฺสน\-อปฺปฏิฆํ เหตุํ ธมฺมํ ปฏิจฺจ นสนิทสฺสน\-สปฺปฏิโฆ นเหตุ ธมฺโม อุปฺปชฺชติ เหตุปจฺจยา ฯเปฯ.  อนิทสฺสน\-อปฺปฏิฆํ เหตุํ ธมฺมํ ปฏิจฺจ นสนิทสฺสน\-สปฺปฏิโฆ นเหตุ จ นอนิทสฺสน\-สปฺปฏิโฆ นเหตุ จ ธมฺมา อุปฺปชฺชนฺติ เหตุปจฺจยา.  เหตุยา ฉ.}

\proseitem{89}{อนิทสฺสน\-อปฺปฏิฆํ นเหตุํ ธมฺมํ ปฏิจฺจ นสนิทสฺสน\-สปฺปฏิโฆ นนเหตุ ธมฺโม อุปฺปชฺชติ เหตุปจฺจยา.  อนิทสฺสน\-อปฺปฏิฆํ นเหตุํ ธมฺมํ ปฏิจฺจ นอนิทสฺสน\-สปฺปฏิโฆ นนเหตุ ธมฺโม อุปฺปชฺชติ เหตุปจฺจยา.  อนิทสฺสน\-อปฺปฏิฆํ นเหตุํ ธมฺมํ ปฏิจฺจ นสนิทสฺสน\-สปฺปฏิโฆ นนเหตุ จ นอนิทสฺสน\-สปฺปฏิโฆ นนเหตุ จ ธมฺมา อุปฺปชฺชนฺติ เหตุปจฺจยา.  เหตุยา ตีณิ. 22. สนิทสฺสน\-ตฺติก 2. สเหตุทุก}

\proseitem{90}{อนิทสฺสน\-อปฺปฏิฆํ สเหตุกํ ธมฺมํ ปฏิจฺจ นอนิทสฺสน\-อปฺปฏิโฆ นสเหตุโก ธมฺโม อุปฺปชฺชติ เหตุปจฺจยา.  อนิทสฺสน\-อปฺปฏิฆํ สเหตุกํ ธมฺมํ ปฏิจฺจ นสนิทสฺสน\-สปฺปฏิโฆ นสเหตุโก ธมฺโม อุปฺปชฺชติ เหตุปจฺจยา ฯเปฯ.  อนิทสฺสน\-อปฺปฏิฆํ สเหตุกํ ธมฺมํ ปฏิจฺจ นสนิทสฺสน\-สปฺปฏิโฆ นสเหตุโก จ นอนิทสฺสน\-สปฺปฏิโฆ นสเหตุโก จ ธมฺมา อุปฺปชฺชนฺติ เหตุปจฺจยา.  เหตุยา ฉ.}

\prose{อนิทสฺสน\-อปฺปฏิฆํ อเหตุกํ ธมฺมํ ปฏิจฺจ นสนิทสฺสน\-สปฺปฏิโฆ นอเหตุโก ธมฺโม อุปฺปชฺชติ เหตุปจฺจยา. (สํขิตฺตํ.) . เหตุยา ตีณิ.}

\csromanmarkh{22. สนิทสฺสน\-ตฺติก}\csromanmarkhii{3. เหตุ\-สมฺปยุตฺตทุก}\csromanheadernomark{2}{22. สนิทสฺสน\-ตฺติก 3. เหตุ\-สมฺปยุตฺตทุก}
\proseitem{91}{อนิทสฺสน\-อปฺปฏิฆํ เหตุ\-สมฺปยุตฺตํ ธมฺมํ ปฏิจฺจ นอนิทสฺสน\-อปฺปฏิโฆ นเหตุ\-สมฺปยุตฺโต ธมฺโม อุปฺปชฺชติ เหตุปจฺจยา.  }

\vspace{\tipitakaparskip}
\csromancenter{สนิทสฺสน\-ตฺติก 6. นเหตุ\-สเหตุกทุก อนิทสฺสน\-อปฺปฏิฆํ เหตุ\-สมฺปยุตฺตํ ธมฺมํ ปฏิจฺจ นสนิทสฺสน\-สปฺปฏิโฆ นเหตุ\-สมฺปยุตฺโต ธมฺโม อุปฺปชฺชติ เหตุปจฺจยา ฯเปฯ.  อนิทสฺสน\-อปฺปฏิฆํ เหตุ\-สมฺปยุตฺตํ ธมฺมํ ปฏิจฺจ นสนิทสฺสน\-สปฺปฏิโฆ นเหตุ\-สมฺปยุตฺโต จ นอนิทสฺสน\-สปฺปฏิโฆ นเหตุ\-สมฺปยุตฺโต จ ธมฺมา อุปฺปชฺชนฺติ เหตุปจฺจยา.  เหตุยา ฉ.}
\prose{\csromanfolio{273}อนิทสฺสน\-อปฺปฏิฆํ เหตุ\-วิปฺปยุตฺตํ ธมฺมํ ปฏิจฺจ นสนิทสฺสน\-สปฺปฏิโฆ นเหตุ\-วิปฺปยุตฺโต ธมฺโม อุปฺปชฺชติ เหตุปจฺจยา.  เหตุยา ตีณิ.}

\csromanmarkh{22. สนิทสฺสน\-ตฺติก}\csromanmarkhii{4. เหตุ\-สเหตุกทุก}\csromanheadernomark{2}{22. สนิทสฺสน\-ตฺติก 4. เหตุ\-สเหตุกทุก}
\proseitem{92}{อนิทสฺสน\-อปฺปฏิฆํ เหตุญฺเจว สเหตุกญฺจ ธมฺมํ ปฏิจฺจ นสนิทสฺสน\-สปฺปฏิโฆ นเหตุ เจว นอเหตุโก จ ธมฺโม อุปฺปชฺชติ เหตุปจฺจยา.  เหตุยา ตีณิ.}

\prose{อนิทสฺสน\-อปฺปฏิฆํ สเหตุกญฺเจว น จ เหตุํ ธมฺมํ ปฏิจฺจ นสนิทสฺสน\-สปฺปฏิโฆ นอเหตุโก เจว นนเหตุ จ ธมฺโม อุปฺปชฺชติ เหตุปจฺจยา.  เหตุยา ตีณิ.}

\csromanmarkh{22. สนิทสฺสน\-ตฺติก}\csromanmarkhii{5. เหตุ\-เหตุ\-สมฺปยุตฺตทุก}\csromanheadernomark{2}{22. สนิทสฺสน\-ตฺติก 5. เหตุ\-เหตุ\-สมฺปยุตฺตทุก}
\proseitem{93}{อนิทสฺสน\-อปฺปฏิฆํ เหตุญฺเจว เหตุ\-สมฺปยุตฺตญฺจ ธมฺมํ ปฏิจฺจ นสนิทสฺสน\-สปฺปฏิโฆ นเหตุ เจว นเหตุ\-วิปฺปยุตฺโต จ ธมฺโม อุปฺปชฺชติ เหตุปจฺจยา.  เหตุยา ตีณิ. อนิทสฺสน\-อปฺปฏิฆํ เหตุ\-สมฺปยุตฺต\-ญฺเจว น จ เหตุํ ธมฺมํ ปฏิจฺจ นสนิทสฺสน\-สปฺปฏิโฆ นเหตุ\-วิปฺปยุตฺโต เจว นนเหตุ จ ธมฺโม อุปฺปชฺชติ เหตุปจฺจยา.  เหตุยา ตีณิ.}

\csromanmarkh{22. สนิทสฺสน\-ตฺติก}\csromanmarkhii{6. นเหตุ\-สเหตุกทุก}\csromanheadernomark{2}{22. สนิทสฺสน\-ตฺติก 6. นเหตุ\-สเหตุกทุก}
\proseitem{94}{อนิทสฺสน\-อปฺปฏิฆํ นเหตุํ สเหตุกํ ธมฺมํ ปฏิจฺจ นอนิทสฺสน\-อปฺปฏิโฆ นเหตุ นสเหตุโก ธมฺโม อุปฺปชฺชติ เหตุปจฺจยา.  เหตุยา ฉ.}

\csromanheader{2}{ธมฺมา\-นุโลม\-ปจฺจนีย ติกทุก\-ปฏฺฐาน\-ปาฬิ}
\prose{\csromanfolio{274}อนิทสฺสน\-อปฺปฏิฆํ นเหตุํ อเหตุกํ ธมฺมํ ปฏิจฺจ นสนิทสฺสน\-สปฺปฏิโฆ นเหตุ นอเหตุโก ธมฺโม อุปฺปชฺชติ เหตุปจฺจยา.  เหตุยา ตีณิ.}

\csromanmarkh{22. สนิทสฺสน\-ตฺติก}\csromanmarkhii{7. สปฺปจฺจยทุก}\csromanheadernomark{2}{22. สนิทสฺสน\-ตฺติก 7. สปฺปจฺจยทุก}
\proseitem{95}{อนิทสฺสน\-อปฺปฏิโฆ อปฺปจฺจโย ธมฺโม นสนิทสฺสน\-สปฺปฏิฆสฺส นอปฺปจฺจยสฺส ธมฺมสฺส อารมฺมณ\-ปจฺจเยน ปจฺจโย.}

\prose{อารมฺมเณ ตีณิ, อธิปติยา ตีณิ, อุปนิสฺสเย ตีณิ. (สํขิตฺตํ.)}

\csromanmarkh{22. สนิทสฺสน\-ตฺติก}\csromanmarkhii{9. สนิทสฺสนทุก}\csromanheadernomark{2}{22. สนิทสฺสน\-ตฺติก 9. สนิทสฺสนทุก}
\proseitem{96}{สนิทสฺสน\-สปฺปฏิโฆ สนิทสฺสโน ธมฺโม นสนิทสฺสน\-สปฺปฏิฆสฺส นสนิทสฺสนสฺส ธมฺมสฺส อารมฺมณ\-ปจฺจเยน ปจฺจโย. (สํขิตฺตํ.) . อารมฺมเณ ตีณิ, อธิปติยา ตีณิ, อุปนิสฺสเย ปุเรชาเต อตฺถิยา อวิคเต ตีณิ.}

\prose{อนิทสฺสน\-สปฺปฏิฆํ อนิทสฺสนํ ธมฺมํ ปฏิจฺจ นอนิทสฺสน\-สปฺปฏิโฆ นอนิทสฺสโน ธมฺโม อุปฺปชฺชติ เหตุปจฺจยา.  เหตุยา นว.}

\csromanmarkh{22. สนิทสฺสนตฺตฺอิก}\csromanmarkhii{10. สปฺปฏิฆทุก}\csromanheadernomark{2}{22. \textbf{สนิทสฺสนตฺตฺ}อิก 10. สปฺปฏิฆทุก}
\proseitem{97}{อนิทสฺสน\-สปฺปฏิฆํ สปฺปฏิฆํ ธมฺมํ ปฏิจฺจ นอนิทสฺสน\-สปฺปฏิโฆ นสปฺปฏิโฆ ธมฺโม อุปฺปชฺชติ เหตุปจฺจยา.  อนิทสฺสน\-สปฺปฏิฆํ สปฺปฏิฆํ ธมฺมํ ปฏิจฺจ นสนิทสฺสน\-สปฺปฏิโฆ นสปฺปฏิโฆ ธมฺโม อุปฺปชฺชติ เหตุปจฺจยา.  อนิทสฺสน\-สปฺปฏิฆํ สปฺปฏิฆํ ธมฺมํ ปฏิจฺจ นสนิทสฺสน\-สปฺปฏิโฆ นสปฺปฏิโฆ จ นอนิทสฺสน\-สปฺปฏิโฆ นสปฺปฏิโฆ จ ธมฺมา อุปฺปชฺชนฺติ เหตุปจฺจยา.  เหตุยา ตีณิ.}

\prose{อนิทสฺสน\-อปฺปฏิฆํ อปฺปฏิฆํ ธมฺมํ ปฏิจฺจ นอนิทสฺสน\-อปฺปฏิโฆ นอปฺปฏิโฆ ธมฺโม อุปฺปชฺชติ เหตุปจฺจยา.  อนิทสฺสน\-อปฺปฏิฆํ อปฺปฏิฆํ ธมฺมํ ปฏิจฺจ นสนิทสฺสน\-สปฺปฏิโฆ นอปฺปฏิโฆ ธมฺโม อุปฺปชฺชติ เหตุปจฺจยา ฯเปฯ. อนิทสฺสน\-อปฺปฏิฆํ อปฺปฏิฆํ ธมฺมํ ปฏิจฺจ นอนิทสฺสน\-สปฺปฏิโฆ นอปฺปฏิโฆ จ นอนิทสฺสน\-อปฺปฏิโฆ นอปฺปฏิโฆ จ ธมฺมา อุปฺปชฺชนฺติ เหตุปจฺจยา.  เหตุยา ปญฺจ.}

\vspace{\tipitakaparskip}
\csromancenter{สนิทสฺสน\-ตฺติก 13. เกนจิ\-วิญฺเญยฺยทุก 22. สนิทสฺสน\-ตฺติก 11. รูปีทุก}
\proseitem{98}{\csromanfolio{275}อนิทสฺสน\-อปฺปฏิฆํ รูปิํ ธมฺมํ ปฏิจฺจ นสนิทสฺสน\-สปฺปฏิโฆ นรูปี ธมฺโม อุปฺปชฺชติ เหตุปจฺจยา.  เหตุยา ตีณิ.}

\prose{อนิทสฺสน\-อปฺปฏิฆํ อรูปิํ ธมฺมํ ปฏิจฺจ นอนิทสฺสน\-อปฺปฏิโฆ นอรูปี ธมฺโม อุปฺปชฺชติ เหตุปจฺจยา.  อนิทสฺสน\-อปฺปฏิฆํ อรูปิํ ธมฺมํ ปฏิจฺจ นสนิทสฺสน\-สปฺปฏิโฆ นอรูปี ธมฺโม อุปฺปชฺชติ เหตุปจฺจยา ฯเปฯ.  อนิทสฺสน\-อปฺปฏิฆํ อรูปิํ ธมฺมํ ปฏิจฺจ นสนิทสฺสน\-สปฺปฏิโฆ นอรูปี จ นอนิทสฺสน\-สปฺปฏิโฆ นอรูปี จ ธมฺมา อุปฺปชฺชนฺติ เหตุปจฺจยา.  เหตุยา ฉ.}

\csromanmarkh{22. สนิทสฺสน\-ตฺติก}\csromanmarkhii{12. โลกิยทุก}\csromanheadernomark{2}{22. สนิทสฺสน\-ตฺติก 12. โลกิยทุก}
\proseitem{99}{อนิทสฺสน\-อปฺปฏิฆํ โลกิยํ ธมฺมํ ปจฺจยา นสนิทสฺสน\-สปฺปฏิโฆ นโลกิโย ธมฺโม อุปฺปชฺชติ เหตุปจฺจยา. (สํขิตฺตํ.) . เหตุยา ตีณิ ฯเปฯ อวิคเต ตีณิ.}

\prose{อนิทสฺสน\-อปฺปฏิฆํ โลกุตฺตรํ ธมฺมํ ปฏิจฺจ นอนิทสฺสน\-อปฺปฏิโฆ นโลกุตฺตโร ธมฺโม อุปฺปชฺชติ เหตุปจฺจยา.  เหตุยา ฉ.}

\csromanmarkh{22. สนิทสฺสน\-ตฺติก}\csromanmarkhii{13. เกนจิ\-วิญฺเญยฺยทุก}\csromanheadernomark{2}{22. สนิทสฺสน\-ตฺติก 13. เกนจิ\-วิญฺเญยฺยทุก}
\proseitem{100}{อนิทสฺสน\-สปฺปฏิฆํ เกนจิ วิญฺเญยฺยํ ธมฺมํ ปฏิจฺจ นอนิทสฺสน\-สปฺปฏิโฆ นเกนจิ วิญฺเญยฺโย ธมฺโม อุปฺปชฺชติ เหตุปจฺจยา.  อนิทสฺสน\-สปฺปฏิฆํ เกนจิ วิญฺเญยฺยํ ธมฺมํ ปฏิจฺจ นสนิทสฺสน\-สปฺปฏิโฆ นเกนจิ วิญฺเญยฺโย ธมฺโม อุปฺปชฺชติ เหตุปจฺจยา.  อนิทสฺสน\-สปฺปฏิฆํ เกนจิ วิญฺเญยฺยํ ธมฺมํ ปฏิจฺจ นอนิทสฺสน\-อปฺปฏิโฆ นเกนจิ วิญฺเญยฺโย ธมฺโม อุปฺปชฺชติ เหตุปจฺจยา ฯเปฯ.  อนิทสฺสน\-สปฺปฏิฆํ เกนจิ วิญฺเญยฺยํ ธมฺมํ ปฏิจฺจ นสนิทสฺสน\-สปฺปฏิโฆ นเกนจิ วิญฺเญยฺโย จ นอนิทสฺสน\-สปฺปฏิโฆ นเกนจิ วิญฺเญยฺโย จ ธมฺมา อุปฺปชฺชนฺติ เหตุปจฺจยา. ฉ.}

\prose{อนิทสฺสน\-อปฺปฏิฆํ เกนจิ วิญฺเญยฺยํ ธมฺมํ ปฏิจฺจ นอนิทสฺสน\-อปฺปฏิโฆ นเกนจิ วิญฺเญยฺโย ธมฺโม อุปฺปชฺชติ เหตุปจฺจยา. ฉ.}

\prose{อนิทสฺสน\-สปฺปฏิฆํ เกนจิ วิญฺเญยฺยญฺจ อนิทสฺสน\-อปฺปฏิฆํ เกนจิ วิญฺเญยฺยญฺจ ธมฺมํ ปฏิจฺจ นสนิทสฺสน\-สปฺปฏิโฆ นเกนจิ วิญฺเญยฺโย ธมฺโม อุปฺปชฺชติ เหตุปจฺจยา. ฉ.  เหตุยา อฏฺฐารส.}

\csromanheader{2}{ธมฺมา\-นุโลม\-ปจฺจนีย ติกทุก\-ปฏฺฐาน\-ปาฬิ}
\prose{\csromanfolio{276}อนิทสฺสน\-สปฺปฏิฆํ นเกนจิ วิญฺเญยฺยํ ธมฺมํ ปฏิจฺจ นสนิทสฺสน\-สปฺปฏิโฆ นนเกนจิ วิญฺเญยฺโย ธมฺโม อุปฺปชฺชติ เหตุปจฺจยา.  เหตุยา อฏฺฐารส.}

\csromanmarkh{22. สนิทสฺสน\-ตฺติก}\csromanmarkhii{14. อาสวทุก}\csromanheadernomark{2}{22. สนิทสฺสน\-ตฺติก 14. อาสวทุก}
\proseitem{101}{อนิทสฺสน\-อปฺปฏิฆํ อาสวํ ธมฺมํ ปฏิจฺจ นอนิทสฺสน\-อปฺปฏิโฆ โนอาสโว ธมฺโม อุปฺปชฺชติ เหตุปจฺจยา.  อนิทสฺสน\-อปฺปฏิฆํ อาสวํ ธมฺมํ ปฏิจฺจ นสนิทสฺสน\-สปฺปฏิโฆ โนอาสโว ธมฺโม อุปฺปชฺชติ เหตุปจฺจยา ฯเปฯ.  อนิทสฺสน\-อปฺปฏิฆํ อาสวํ ธมฺมํ ปฏิจฺจ นสนิทสฺสน\-สปฺปฏิโฆ โนอาสโว จ นอนิทสฺสน\-สปฺปฏิโฆ โนอาสโว จ ธมฺมา อุปฺปชฺชนฺติ เหตุปจฺจยา.  เหตุยา ฉ.}

\prose{อนิทสฺสน\-อปฺปฏิฆํ โนอาสวํ ธมฺมํ ปฏิจฺจ นสนิทสฺสน\-สปฺปฏิโฆ นโนอาสโว ธมฺโม อุปฺปชฺชติ เหตุปจฺจยา.  อนิทสฺสน\-อปฺปฏิฆํ โนอาสวํ ธมฺมํ ปฏิจฺจ นอนิทสฺสน\-สปฺปฏิโฆ นโนอาสโว ธมฺโม อุปฺปชฺชติ เหตุปจฺจยา.  อนิทสฺสน\-อปฺปฏิฆํ โนอาสวํ ธมฺมํ ปฏิจฺจ นสนิทสฺสน\-สปฺปฏิโฆ นโนอาสโว จ นอนิทสฺสน\-สปฺปฏิโฆ นโนอาสโว จ ธมฺมา อุปฺปชฺชนฺติ เหตุปจฺจยา.  เหตุยา ตีณิ.}

\csromanmarkh{22. สนิทสฺสน\-ตฺติก}\csromanmarkhii{15. สาสวทุก}\csromanheadernomark{2}{22. สนิทสฺสน\-ตฺติก 15. สาสวทุก}
\proseitem{102}{อนิทสฺสน\-อปฺปฏิฆํ สาสวํ ธมฺมํ ปจฺจยา นสนิทสฺสน\-สปฺปฏิโฆ นสาสโว ธมฺโม อุปฺปชฺชติ เหตุปจฺจยา.  อนิทสฺสน\-อปฺปฏิฆํ สาสวํ ธมฺมํ ปจฺจยา นอนิทสฺสน\-สปฺปฏิโฆ นสาสโว ธมฺโม อุปฺปชฺชติ เหตุปจฺจยา.  อนิทสฺสน\-อปฺปฏิฆํ สาสวํ ธมฺมํ ปจฺจยา นสนิทสฺสน\-สปฺปฏิโฆ นสาสโว จ นอนิทสฺสน\-สปฺปฏิโฆ นสาสโว จ ธมฺมา อุปฺปชฺชนฺติ เหตุปจฺจยา.  เหตุยา ตีณิ.}

\prose{อนิทสฺสน\-อปฺปฏิฆํ อนาสวํ ธมฺมํ ปฏิจฺจ นอนิทสฺสน\-อปฺปฏิโฆ นอนาสโว ธมฺโม อุปฺปชฺชติ เหตุปจฺจยา.  อนิทสฺสน\-อปฺปฏิฆํ อนาสวํ ธมฺมํ ปฏิจฺจ นสนิทสฺสน\-สปฺปฏิโฆ นอนาสโว ธมฺโม อุปฺปชฺชติ เหตุปจฺจยา ฯเปฯ.  อนิทสฺสน\-อปฺปฏิฆํ อนาสวํ ธมฺมํ ปฏิจฺจ นสนิทสฺสน\-สปฺปฏิโฆ นอนาสโว จ นอนิทสฺสน\-สปฺปฏิโฆ นอนาสโว จ ธมฺมา อุปฺปชฺชนฺติ เหตุปจฺจยา.  เหตุยา ฉ.}

\vspace{\tipitakaparskip}
\csromancenter{สนิทสฺสน\-ตฺติก 17. อาสว\-สาสวทุก 22. สนิทสฺสน\-ตฺติก 16. อาสว\-สมฺปยุตฺตทุก}
\proseitem{103}{\csromanfolio{277}อนิทสฺสน\-อปฺปฏิฆํ อาสว\-สมฺปยุตฺตํ ธมฺมํ ปฏิจฺจ นอนิทสฺสน\-อปฺปฏิโฆ นอาสว\-สมฺปยุตฺโต ธมฺโม อุปฺปชฺชติ เหตุปจฺจยา.  อนิทสฺสน\-อปฺปฏิฆํ อาสว\-สมฺปยุตฺตํ ธมฺมํ ปฏิจฺจ นสนิทสฺสน\-สปฺปฏิโฆ นอาสว\-สมฺปยุตฺโต ธมฺโม อุปฺปชฺชติ เหตุปจฺจยา ฯเปฯ.  อนิทสฺสน\-อปฺปฏิฆํ อาสว\-สมฺปยุตฺตํ ธมฺมํ ปฏิจฺจ นสนิทสฺสน\-สปฺปฏิโฆ นอาสว\-สมฺปยุตฺโต จ นอนิทสฺสน\-สปฺปฏิโฆ นอาสว\-สมฺปยุตฺโต จ ธมฺมา อุปฺปชฺชนฺติ เหตุปจฺจยา.  เหตุยา ฉ.}

\prose{อนิทสฺสน\-อปฺปฏิฆํ อาสว\-วิปฺปยุตฺตํ ธมฺมํ ปฏิจฺจ นสนิทสฺสน\-สปฺปฏิโฆ นอาสว\-วิปฺปยุตฺโต ธมฺโม อุปฺปชฺชติ เหตุปจฺจยา.  เหตุยา ตีณิ.}

\csromanmarkh{22. สนิทสฺสน\-ตฺติก}\csromanmarkhii{17. อาสว\-สาสวทุก}\csromanheadernomark{2}{22. สนิทสฺสน\-ตฺติก 17. อาสว\-สาสวทุก}
\proseitem{104}{อนิทสฺสน\-อปฺปฏิฆํ อาสวญฺเจว สาสวญฺจ ธมฺมํ ปฏิจฺจ นอนิทสฺสน\-อปฺปฏิโฆ นอาสโว เจว นอนาสโว จ ธมฺโม อุปฺปชฺชติ เหตุปจฺจยา.  อนิทสฺสน\-อปฺปฏิฆํ อาสวญฺเจว สาสวญฺจ ธมฺมํ ปฏิจฺจ นสนิทสฺสน\-สปฺปฏิโฆ นอาสโว เจว นอนาสโว จ ธมฺโม อุปฺปชฺชติ เหตุปจฺจยา ฯเปฯ.  อนิทสฺสน\-อปฺปฏิฆํ อาสวญฺเจว สาสวญฺจ ธมฺมํ ปฏิจฺจ นสนิทสฺสน\-สปฺปฏิโฆ นอาสโว เจว นอนาสโว จ นอนิทสฺสน\-สปฺปฏิโฆ นอาสโว เจว นอนาสโว จ ธมฺมา อุปฺปชฺชนฺติ เหตุปจฺจยา.  เหตุยา ฉ.}

\prose{อนิทสฺสน\-อปฺปฏิฆํ สาสวญฺเจว โน จ อาสวํ ธมฺมํ ปฏิจฺจ นสนิทสฺสน\-สปฺปฏิโฆ นอนาสโว เจว นโนอาสโว จ ธมฺโม อุปฺปชฺชติ เหตุปจฺจยา.  อนิทสฺสน\-อปฺปฏิฆํ สาสวญฺเจว โน จ อาสวํ ธมฺมํ ปฏิจฺจ นอนิทสฺสน\-สปฺปฏิโฆ นอนาสโว เจว นโน จ อาสโว ธมฺโม อุปฺปชฺชติ เหตุปจฺจยา.  อนิทสฺสน\-สปฺปฏิฆํ สาสวญฺเจว โน จ อาสวํ ธมฺมํ ปฏิจฺจ นสนิทสฺสน\-สปฺปฏิโฆ นอนาสโว เจว นโน จ อาสโว นอนิทสฺสน\-สปฺปฏิโฆ นอนาสโว เจว นโน อาสโว จ ธมฺมา อุปฺปชฺชนฺติ เหตุปจฺจยา.  เหตุยา ตีณิ.}

\vspace{\tipitakaparskip}
\csromancenter{ติกทุก\-ปฏฺฐาน\-ปาฬิ 22. สนิทสฺสน\-ตฺติก 18. อาสว\-อาสว\-สมฺปยุตฺตทุก}
\proseitem{105}{\csromanfolio{278}อนิทสฺสน\-อปฺปฏิฆํ อาสวญฺเจว อาสว\-สมฺปยุตฺตญฺจ ธมฺมํ ปฏิจฺจ นสนิทสฺสน\-สปฺปฏิโฆ นอาสโว เจว นอาสว\-วิปฺปยุตฺโต จ ธมฺโม อุปฺปชฺชติ เหตุปจฺจยา.  อนิทสฺสน\-อปฺปฏิฆํ อาสวญฺเจว อาสว\-สมฺปยุตฺตญฺจ ธมฺมํ ปฏิจฺจ นอนิทสฺสน\-สปฺปฏิโฆ นอาสโว เจว นอาสว\-วิปฺปยุตฺโต จ ธมฺโม อุปฺปชฺชติ เหตุปจฺจยา.  อนิทสฺสน\-อปฺปฏิฆํ อาสวญฺเจว อาสว\-สมฺปยุตฺตญฺจ ธมฺมํ ปฏิจฺจ นสนิทสฺสน\-สปฺปฏิโฆ นอาสโว เจว นอาสว\-วิปฺปยุตฺโต จ นอนิทสฺสน\-สปฺปฏิโฆ นอาสโว เจว นอาสว\-วิปฺปยุตฺโต จ ธมฺมา อุปฺปชฺชนฺติ เหตุปจฺจยา.  เหตุยา ตีณิ.}

\prose{อนิทสฺสน\-อปฺปฏิฆํ อาสว\-สมฺปยุตฺต\-ญฺเจว โน จ อาสวํ ธมฺมํ ปฏิจฺจ นสนิทสฺสน\-สปฺปฏิโฆ นอาสว\-วิปฺปยุตฺโต เจว นโน จ อาสโว ธมฺโม อุปฺปชฺชติ เหตุปจฺจยา.  เหตุยา ตีณิ.}

\prose{22. สนิทสฺสน\-ตฺติก 19. อาสว\-วิปฺปยุตฺต\-สาสวทุก}

\proseitem{106}{อนิทสฺสน\-อปฺปฏิฆํ อาสว\-วิปฺปยุตฺตํ สาสวํ ธมฺมํ ปจฺจยา นสนิทสฺสน\-สปฺปฏิโฆ อาสว\-วิปฺปยุตฺโต นสาสโว ธมฺโม อุปฺปชฺชติ เหตุปจฺจยา.  อนิทสฺสน\-อปฺปฏิฆํ อาสว\-วิปฺปยุตฺตํ สาสวํ ธมฺมํ ปจฺจยา นอนิทสฺสน\-สปฺปฏิโฆ อาสว\-วิปฺปยุตฺโต นสาสโว ธมฺโม อุปฺปชฺชติ เหตุปจฺจยา.  อนิทสฺสน\-อปฺปฏิฆํ อาสว\-วิปฺปยุตฺตํ สาสวํ ธมฺมํ ปจฺจยา นสนิทสฺสน\-สปฺปฏิโฆ อาสว\-วิปฺปยุตฺโต นสาสโว จ นอนิทสฺสน\-สปฺปฏิโฆ อาสว\-วิปฺปยุตฺโต นสาสโว จ ธมฺมา อุปฺปชฺชนฺติ เหตุปจฺจยา.  เหตุยา ตีณิ.}

\prose{อนิทสฺสน\-อปฺปฏิฆํ อาสว\-วิปฺปยุตฺตํ อนาสวํ ธมฺมํ ปฏิจฺจ นอนิทสฺสน\-อปฺปฏิโฆ อาสว\-วิปฺปยุตฺโต นอนาสโว ธมฺโม อุปฺปชฺชติ เหตุปจฺจยา.  อนิทสฺสน\-อปฺปฏิฆํ อาสว\-วิปฺปยุตฺตํ อนาสวํ ธมฺมํ ปฏิจฺจ นสนิทสฺสน\-สปฺปฏิโฆ อาสว\-วิปฺปยุตฺโต นอนาสโว ธมฺโม อุปฺปชฺชติ เหตุปจฺจยา ฯเปฯ.  อนิทสฺสน\-อปฺปฏิฆํ อาสว\-วิปฺปยุตฺตํ อนาสวํ ธมฺมํ ปฏิจฺจ นสนิทสฺสน\-สปฺปฏิโฆ อาสว\-วิปฺปยุตฺโต นอนาสโว จ นอนิทสฺสน\-สปฺปฏิโฆ อาสว\-วิปฺปยุตฺโต นอนาสโว จ ธมฺมา อุปฺปชฺชนฺติ เหตุปจฺจยา.  เหตุยา ฉ.}

\vspace{\tipitakaparskip}
\csromancenter{สนิทสฺสน\-ตฺติก 56. จิตฺตทุก 22. สนิทสฺสน\-ตฺติก 20-54. สญฺโญชนา\-ทิทุก}
\proseitem{107}{\csromanfolio{279}อนิทสฺสน\-อปฺปฏิฆํ สญฺโญชนํ ธมฺมํ ปฏิจฺจ นอนิทสฺสน\-อปฺปฏิโฆ โนสญฺโญชโน ธมฺโม อุปฺปชฺชติ เหตุปจฺจยา.}

\proseitem{108}{อนิทสฺสน\-อปฺปฏิฆํ คนฺถํ ธมฺมํ ปฏิจฺจ นอนิทสฺสน\-อปฺปฏิโฆ โนคนฺโถ ธมฺโม อุปฺปชฺชติ เหตุปจฺจยา ฯเปฯ.}

\prose{อนิทสฺสน\-อปฺปฏิฆํ โอฆํ ธมฺมํ ปฏิจฺจ นอนิทสฺสน\-อปฺปฏิโฆ โนโอโฆ ธมฺโม อุปฺปชฺชติ เหตุปจฺจยา ฯเปฯ.}

\prose{อนิทสฺสน\-อปฺปฏิฆํ โยคํ ธมฺมํ ปฏิจฺจ นอนิทสฺสน\-อปฺปฏิโฆ โนโยโค ธมฺโม อุปฺปชฺชติ เหตุปจฺจยา.}

\proseitem{109}{อนิทสฺสน\-อปฺปฏิฆํ นีวรณํ ธมฺมํ ปฏิจฺจ นอนิทสฺสน\-อปฺปฏิโฆ โนนีวรโณ ธมฺโม อุปฺปชฺชติ เหตุปจฺจยา.}

\proseitem{110}{อนิทสฺสน\-อปฺปฏิฆํ ปรามาสํ ธมฺมํ ปฏิจฺจ นอนิทสฺสน\-อปฺปฏิโฆ โนปรามาโส ธมฺโม อุปฺปชฺชติ เหตุปจฺจยา.}

\csromanmarkh{22. สนิทสฺสนตฺตฺอิก}\csromanmarkhii{55. สารมฺมณทุก}\csromanheadernomark{2}{22. \textbf{สนิทสฺสนตฺตฺ}อิก 55. สารมฺมณทุก}
\proseitem{111}{อนิทสฺสน\-อปฺปฏิฆํ สารมฺมณํ ธมฺมํ ปฏิจฺจ นอนิทสฺสน\-อปฺปฏิโฆ นสารมฺมโณ ธมฺโม อุปฺปชฺชติ เหตุปจฺจยา.  อนิทสฺสน\-อปฺปฏิฆํ สารมฺมณํ ธมฺมํ ปฏิจฺจ นสนิทสฺสน\-สปฺปฏิโฆ นสารมฺมโณ ธมฺโม อุปฺปชฺชติ เหตุปจฺจยา ฯเปฯ.  อนิทสฺสน\-อปฺปฏิฆํ สารมฺมณํ ธมฺมํ ปฏิจฺจ นสนิทสฺสน\-สปฺปฏิโฆ นสารมฺมโณ จ นอนิทสฺสน\-สปฺปฏิโฆ นสารมฺมโณ จ ธมฺมา อุปฺปชฺชนฺติ เหตุปจฺจยา.}

\prose{อนิทสฺสน\-อปฺปฏิฆํ อนารมฺมณํ ธมฺมํ ปฏิจฺจ นสนิทสฺสน\-สปฺปฏิโฆ นอนารมฺมโณ ธมฺโม อุปฺปชฺชติ เหตุปจฺจยา.  เหตุยา ตีณิ.}

\csromanmarkh{22. สนิทสฺสน\-ตฺติก}\csromanmarkhii{56. จิตฺตทุก}\csromanheadernomark{2}{22. สนิทสฺสน\-ตฺติก 56. จิตฺตทุก}
\proseitem{112}{อนิทสฺสน\-อปฺปฏิฆํ จิตฺตํ ธมฺมํ ปฏิจฺจ นอนิทสฺสน\-อปฺปฏิโฆ โนจิตฺโต ธมฺโม อุปฺปชฺชติ เหตุปจฺจยา.  อนิทสฺสน\-อปฺปฏิฆํ จิตฺตํ ธมฺมํ ปฏิจฺจ นสนิทสฺสน\-สปฺปฏิโฆ โนจิตฺโต ธมฺโม อุปฺปชฺชติ เหตุปจฺจยา ฯเปฯ.  }

\vspace{\tipitakaparskip}
\csromancenter{ติกทุก\-ปฏฺฐาน\-ปาฬิ อนิทสฺสน\-อปฺปฏิฆํ จิตฺตํ ธมฺมํ ปฏิจฺจ นสนิทสฺสน\-สปฺปฏิโฆ โนจิตฺโต จ นอนิทสฺสน\-สปฺปฏิโฆ โนจิตฺโต จ ธมฺมา อุปฺปชฺชนฺติ เหตุปจฺจยา.  เหตุยา ฉ.}
\prose{\csromanfolio{280}อนิทสฺสน\-อปฺปฏิฆํ โนจิตฺตํ ธมฺมํ ปฏิจฺจ นสนิทสฺสน\-สปฺปฏิโฆ นโนจิตฺโต ธมฺโม อุปฺปชฺชติ เหตุปจฺจยา.  เหตุยา ตีณิ.}

\csromanmarkh{22. สนิทสฺสน\-ตฺติก}\csromanmarkhii{57. เจตสิกทุก}\csromanheadernomark{2}{22. สนิทสฺสน\-ตฺติก 57. เจตสิกทุก}
\proseitem{113}{อนิทสฺสน\-อปฺปฏิฆํ เจตสิกํ ธมฺมํ ปฏิจฺจ นอนิทสฺสน\-อปฺปฏิโฆ นเจตสิโก ธมฺโม อุปฺปชฺชติ เหตุปจฺจยา.  อนิทสฺสน\-อปฺปฏิฆํ เจตสิกํ ธมฺมํ ปฏิจฺจ นสนิทสฺสน\-สปฺปฏิโฆ นเจตสิโก ธมฺโม อุปฺปชฺชติ เหตุปจฺจยา ฯเปฯ.  อนิทสฺสน\-อปฺปฏิฆํ เจตสิกํ ธมฺมํ ปฏิจฺจ นสนิทสฺสน\-สปฺปฏิโฆ นเจตสิโก จ นอนิทสฺสน\-สปฺปฏิโฆ นเจตสิโก จ ธมฺมา อุปฺปชฺชนฺติ เหตุปจฺจยา.  เหตุยา ฉ.}

\csromanmarkh{22. สนิทสฺสน\-ตฺติกฺอ}\csromanmarkhii{58. จิตฺต\-สมฺปยุตฺตทุก}\csromanheadernomark{2}{22. สนิทสฺสน\-ตฺติกฺอ 58. จิตฺต\-สมฺปยุตฺตทุก}
\proseitem{114}{อนิทสฺสน\-อปฺปฏิฆํ จิตฺต\-สมฺปยุตฺตํ ธมฺมํ ปฏิจฺจ นอนิทสฺสน\-อปฺปฏิโฆ นจิตฺต\-สมฺปยุตฺโต ธมฺโม อุปฺปชฺชติ เหตุปจฺจยา.  อนิทสฺสน\-อปฺปฏิฆํ จิตฺต\-สมฺปยุตฺตํ ธมฺมํ ปฏิจฺจ นสนิทสฺสน\-สปฺปฏิโฆ นจิตฺต\-สมฺปยุตฺโต ธมฺโม อุปฺปชฺชติ เหตุปจฺจยา ฯเปฯ.  อนิทสฺสน\-อปฺปฏิฆํ จิตฺต\-สมฺปยุตฺตํ ธมฺมํ ปฏิจฺจ นสนิทสฺสน\-สปฺปฏิโฆ นจิตฺต\-สมฺปยุตฺโต จ นอนิทสฺสน\-สปฺปฏิโฆ นจิตฺต\-สมฺปยุตฺโต จ ธมฺมา อุปฺปชฺชนฺติ เหตุปจฺจยา.  เหตุยา ฉ.}

\csromanmarkh{22. สนิทสฺสน\-ตฺติกฺอ}\csromanmarkhii{59. จิตฺต\-สํสฏฺฐทุก}\csromanheadernomark{2}{22. สนิทสฺสน\-ตฺติกฺอ 59. จิตฺต\-สํสฏฺฐทุก}
\proseitem{115}{อนิทสฺสน\-อปฺปฏิฆํ จิตฺต\-สํสฏฺฐํ ธมฺมํ ปฏิจฺจ นอนิทสฺสน\-อปฺปฏิโฆ นจิตฺต\-สํสฏฺโฐ ธมฺโม อุปฺปชฺชติ เหตุปจฺจยา.  อนิทสฺสน\-อปฺปฏิฆํ จิตฺต\-สํสฏฺฐํ ธมฺมํ ปฏิจฺจ นสนิทสฺสน\-สปฺปฏิโฆ นจิตฺต\-สํสฏฺโฐ ธมฺโม อุปฺปชฺชติ เหตุปจฺจยา ฯเปฯ.  อนิทสฺสน\-อปฺปฏิฆํ จิตฺต\-สํสฏฺฐํ ธมฺมํ ปฏิจฺจ นสนิทสฺสน\-สปฺปฏิโฆ นจิตฺต\-สํสฏฺโฐ จ นอนิทสฺสน\-สปฺปฏิโฆ นจิตฺต\-สํสฏฺโฐ จ ธมฺมา อุปฺปชฺชนฺติ เหตุปจฺจยา.  เหตุยา ฉ.}

\vspace{\tipitakaparskip}
\csromancenter{สนิทสฺสน\-ตฺติก 66. อชฺฌตฺตทุก 22. สนิทสฺสน\-ตฺติก 60. จิตฺต\-สมุฏฺฐานทุก}
\proseitem{116}{\csromanfolio{281}อนิทสฺสน\-อปฺปฏิฆํ จิตฺต\-สมุฏฺฐานํ ธมฺมํ ปฏิจฺจ นอนิทสฺสน\-อปฺปฏิโฆ โนจิตฺต\-สมุฏฺฐาโน ธมฺโม อุปฺปชฺชติ เหตุปจฺจยา.  เหตุยา ฉ.}

\csromanmarkh{22. สนิทสฺสน\-ตฺติก}\csromanmarkhii{61. จิตฺต\-สหภูทุก}\csromanheadernomark{2}{22. สนิทสฺสน\-ตฺติก 61. จิตฺต\-สหภูทุก}
\proseitem{117}{อนิทสฺสน\-อปฺปฏิฆํ จิตฺตสหภุํ ธมฺมํ ปฏิจฺจ นอนิทสฺสน\-อปฺปฏิโฆ โนจิตฺตสหภู ธมฺโม อุปฺปชฺชติ เหตุปจฺจยา.  เหตุยา ฉ.}

\csromanmarkh{22. สนิทสฺสน\-ตฺติก}\csromanmarkhii{62. จิตฺตา\-นุปริวตฺติทุก}\csromanheadernomark{2}{22. สนิทสฺสน\-ตฺติก 62. จิตฺตา\-นุปริวตฺติทุก}
\proseitem{118}{อนิทสฺสน\-อปฺปฏิฆํ จิตฺตา\-นุปริวตฺติํ ธมฺมํ ปฏิจฺจ นอนิทสฺสน\-อปฺปฏิโฆ โนจิตฺตา\-นุปริวตฺตี ธมฺโม อุปฺปชฺชติ เหตุปจฺจยา.  เหตุยา ฉ.}

\prose{22. สนิทสฺสน\-ตฺติก 63-65. จิตฺต\-สํสฏฺฐ\-สมุฏฺฐาน\-ทุกาทิ}

\proseitem{119}{อนิทสฺสน\-อปฺปฏิฆํ จิตฺต\-สํสฏฺฐ\-สมุฏฺฐานํ ธมฺมํ ปฏิจฺจ นอนิทสฺสน\-อปฺปฏิโฆ โนจิตฺต\-สํสฏฺฐ\-สมุฏฺฐาโน ธมฺโม อุปฺปชฺชติ เหตุปจฺจยา.  เหตุยา ฉ.}

\proseitem{120}{อนิทสฺสน\-อปฺปฏิฆํ จิตฺต\-สํสฏฺฐ\-สมุฏฺฐาน\-สหภุํ ธมฺมํ ปฏิจฺจ นอนิทสฺสน\-อปฺปฏิโฆ โนจิตฺต\-สํสฏฺฐ\-สมุฏฺฐาน\-สหภู ธมฺโม อุปฺปชฺชติ เหตุปจฺจยา. เหตุยา ฉ.}

\proseitem{121}{อนิทสฺสน\-อปฺปฏิฆํ จิตฺต\-สํสฏฺฐ\-สมุฏฺฐานา\-นุปริวตฺติํ ธมฺมํ ปฏิจฺจ นอนิทสฺสน\-อปฺปฏิโฆ โนจิตฺต\-สํสฏฺฐ\-สมุฏฺฐานา\-นุปริวตฺตี ธมฺโม อุปฺปชฺชติ เหตุปจฺจยา.  เหตุยา ฉ.}

\csromanmarkh{22. สนิทสฺสน\-ตฺติก}\csromanmarkhii{66. อชฺฌตฺตทุก}\csromanheadernomark{2}{22. สนิทสฺสน\-ตฺติก 66. อชฺฌตฺตทุก}
\proseitem{122}{อนิทสฺสน\-อปฺปฏิฆํ อชฺฌตฺติกํ ธมฺมํ ปฏิจฺจ นอนิทสฺสน\-อปฺปฏิโฆ นอชฺฌตฺติโก ธมฺโม อุปฺปชฺชติ เหตุปจฺจยา.  อนิทสฺสน\-อปฺปฏิฆํ อชฺฌตฺติกํ ธมฺมํ ปฏิจฺจ นสนิทสฺสน\-สปฺปฏิโฆ นอชฺฌตฺติโก ธมฺโม อุปฺปชฺชติ เหตุปจฺจยา ฯเปฯ.  อนิทสฺสน\-อปฺปฏิฆํ อชฺฌตฺติกํ ธมฺมํ ปฏิจฺจ นสนิทสฺสน\-สปฺปฏิโฆ นอชฺฌตฺติโก จ นอนิทสฺสน\-สปฺปฏิโฆ นอชฺฌตฺติโก จ ธมฺมา อุปฺปชฺชนฺติ เหตุปจฺจยา.  เหตุยา ฉ, อารมฺมเณ ตีณิ ฯเปฯ อวิคเต ฉ.}

\csromanheader{2}{ธมฺมา\-นุโลม\-ปจฺจนีย ติกทุก\-ปฏฺฐาน\-ปาฬิ}
\prose{\csromanfolio{282}อนิทสฺสน\-สปฺปฏิฆํ พาหิรํ ธมฺมํ ปฏิจฺจ นสนิทสฺสน\-สปฺปฏิโฆ นพาหิโร ธมฺโม อุปฺปชฺชติ เหตุปจฺจยา.  เหตุยา เอกาทส.}

\csromanmarkh{22. สนิทสฺสน\-ตฺติก}\csromanmarkhii{67. อุปาทาทุก}\csromanheadernomark{2}{22. สนิทสฺสน\-ตฺติก 67. อุปาทาทุก}
\proseitem{123}{อนิทสฺสน\-อปฺปฏิฆํ อุปาทา ธมฺมํ ปฏิจฺจ นสนิทสฺสน\-สปฺปฏิโฆ โนอุปาทา ธมฺโม อุปฺปชฺชติ เหตุปจฺจยา.  เหตุยา ตีณิ.}

\prose{อนิทสฺสน\-สปฺปฏิฆํ โนอุปาทา ธมฺมํ ปฏิจฺจ นอนิทสฺสน\-สปฺปฏิโฆ นโนอุปาทา ธมฺโม อุปฺปชฺชติ เหตุปจฺจยา. ฉ.}

\prose{อนิทสฺสน\-อปฺปฏิฆํ โนอุปาทา ธมฺมํ ปฏิจฺจ นอนิทสฺสน\-อปฺปฏิโฆ นโนอุปาทา ธมฺโม อุปฺปชฺชติ เหตุปจฺจยา. ฉ.}

\prose{อนิทสฺสน\-สปฺปฏิฆํ โนอุปาทา จ อนิทสฺสน\-อปฺปฏิฆํ โนอุปาทา จ ธมฺมํ ปฏิจฺจ นสนิทสฺสน\-สปฺปฏิโฆ นโนอุปาทา ธมฺโม อุปฺปชฺชติ เหตุปจฺจยา. ฉ. เหตุยา อฏฺฐารส.}

\csromanmarkh{22. สนิทสฺสน\-ตฺติก}\csromanmarkhii{68. อุปาทินฺนทุก}\csromanheadernomark{2}{22. สนิทสฺสน\-ตฺติก 68. อุปาทินฺนทุก}
\proseitem{124}{อนิทสฺสน\-อปฺปฏิฆํ อุปาทินฺนํ ธมฺมํ ปฏิจฺจ นอนิทสฺสน\-อปฺปฏิโฆ นอุปาทินฺโน ธมฺโม อุปฺปชฺชติ เหตุปจฺจยา.  อนิทสฺสน\-อปฺปฏิฆํ อุปาทินฺนํ ธมฺมํ ปฏิจฺจ นสนิทสฺสน\-สปฺปฏิโฆ นอุปาทินฺโน ธมฺโม อุปฺปชฺชติ เหตุปจฺจยา ฯเปฯ.  อนิทสฺสน\-อปฺปฏิฆํ อุปาทินฺนํ ธมฺมํ ปฏิจฺจ นสนิทสฺสน\-สปฺปฏิโฆ นอุปาทินฺโน จ นอนิทสฺสน\-สปฺปฏิโฆ นอุปาทินฺโน จ ธมฺมา อุปฺปชฺชนฺติ เหตุปจฺจยา.  เหตุยา ฉ.}

\proseitem{125}{สนิทสฺสน\-สปฺปฏิโฆ อนุปาทินฺโน ธมฺโม นสนิทสฺสน\-สปฺปฏิฆสฺส นอนุปาทินฺนสฺส ธมฺมสฺส อารมฺมณ\-ปจฺจเยน ปจฺจโย. (สํขิตฺตํ.) . อารมฺมเณ นว.}

\csromanheader{2}{22. สนิทสฺสน\-ตฺติก 69-82. อุปาทานทุกาทิ}
\proseitem{126}{อนิทสฺสน\-อปฺปฏิฆํ อุปาทานํ ธมฺมํ ปฏิจฺจ นอนิทสฺสน\-อปฺปฏิโฆ โนอุปาทาโน ธมฺโม อุปฺปชฺชติ เหตุปจฺจยา.  เหตุยา ฉ.}

\proseitem{127}{อนิทสฺสน\-อปฺปฏิฆํ กิเลสํ ธมฺมํ ปฏิจฺจ นอนิทสฺสน\-อปฺปฏิโฆ โนกิเลโส ธมฺโม อุปฺปชฺชติ เหตุปจฺจยา.  เหตุยา ฉ.}

\proseitem{22}{\csromanfolio{283}สนิทสฺสน\-ตฺติก 85. ทสฺสเนน\-ปหาตพฺพ\-เหตุกทุก 22. สนิทสฺสน\-ตฺติก 83. ทสฺสเนน\-ปหาตพฺพทุก}

\proseitem{128}{อนิทสฺสน\-อปฺปฏิฆํ ทสฺสเนน ปหาตพฺพํ ธมฺมํ ปฏิจฺจ นอนิทสฺสน\-อปฺปฏิโฆ นทสฺสเนน ปหาตพฺโพ ธมฺโม อุปฺปชฺชติ เหตุปจฺจยา.  เหตุยา ฉ.}

\prose{อนิทสฺสน\-อปฺปฏิฆํ นทสฺสเนน ปหาตพฺพํ ธมฺมํ ปจฺจยา นสนิทสฺสน\-สปฺปฏิโฆ นนทสฺสเนน ปหาตพฺโพ ธมฺโม อุปฺปชฺชติ เหตุปจฺจยา.  อนิทสฺสน\-อปฺปฏิฆํ นทสฺสเนน ปหาตพฺพํ ธมฺมํ ปจฺจยา นอนิทสฺสน\-สปฺปฏิโฆ นนทสฺสเนน ปหาตพฺโพ ธมฺโม อุปฺปชฺชติ เหตุปจฺจยา.  อนิทสฺสน\-อปฺปฏิฆํ นทสฺสเนน ปหาตพฺพํ ธมฺมํ ปจฺจยา นสนิทสฺสน\-สปฺปฏิโฆ นนทสฺสเนน ปหาตพฺโพ จ นอนิทสฺสน\-สปฺปฏิโฆ นนทสฺสเนน ปหาตพฺโพ จ ธมฺมา อุปฺปชฺชนฺติ เหตุปจฺจยา.  เหตุยา ตีณิ.}

\csromanmarkh{22. สนิทสฺสน\-ตฺติก}\csromanmarkhii{84. ภาวนาย\-ปหาตพฺพทุก}\csromanheadernomark{2}{22. สนิทสฺสน\-ตฺติก 84. ภาวนาย\-ปหาตพฺพทุก}
\proseitem{129}{อนิทสฺสน\-อปฺปฏิฆํ ภาวนาย ปหาตพฺพํ ธมฺมํ ปฏิจฺจ นอนิทสฺสน\-อปฺปฏิโฆ นภาวนาย ปหาตพฺโพ ธมฺโม อุปฺปชฺชติ เหตุปจฺจยา.  เหตุยา ฉ.}

\prose{อนิทสฺสน\-อปฺปฏิฆํ นภาวนาย ปหาตพฺพํ ธมฺมํ ปจฺจยา นสนิทสฺสน\-สปฺปฏิโฆ นนภาวนาย ปหาตพฺโพ ธมฺโม อุปฺปชฺชติ เหตุปจฺจยา.  เหตุยา ตีณิ.}

\prose{22. สนิทสฺสน\-ตฺติก 85. ทสฺสเนน\-ปหาตพฺพ\-เหตุกทุก}

\proseitem{130}{อนิทสฺสน\-อปฺปฏิฆํ ทสฺสเนน ปหาตพฺพ\-เหตุกํ ธมฺมํ ปฏิจฺจ นอนิทสฺสน\-อปฺปฏิโฆ นทสฺสเนน ปหาตพฺพ\-เหตุโก ธมฺโม อุปฺปชฺชติ เหตุปจฺจยา.  เหตุยา ฉ.}

\prose{อนิทสฺสน\-อปฺปฏิฆํ นทสฺสเนน ปหาตพฺพ\-เหตุกํ ธมฺมํ ปฏิจฺจ นสนิทสฺสน\-สปฺปฏิโฆ นนทสฺสเนน ปหาตพฺพ\-เหตุโก ธมฺโม อุปฺปชฺชติ เหตุปจฺจยา.  เหตุยา ตีณิ.}

\vspace{\tipitakaparskip}
\csromancenter{ติกทุก\-ปฏฺฐาน\-ปาฬิ 22. สนิทสฺสน\-ตฺติก 86. ภาวนาย\-ปหาตพฺพ\-เหตุกทุก}
\proseitem{131}{\csromanfolio{284}อนิทสฺสน\-อปฺปฏิฆํ ภาวนาย ปหาตพฺพ\-เหตุกํ ธมฺมํ ปฏิจฺจ นอนิทสฺสน\-อปฺปฏิโฆ นภาวนาย ปหาตพฺพ\-เหตุโก ธมฺโม อุปฺปชฺชติ เหตุปจฺจยา.  เหตุยา ฉ.}

\prose{อนิทสฺสน\-อปฺปฏิฆํ นภาวนาย ปหาตพฺพ\-เหตุกํ ธมฺมํ ปฏิจฺจ นสนิทสฺสน\-สปฺปฏิโฆ นนภาวนาย ปหาตพฺพ\-เหตุโก ธมฺโม อุปฺปชฺชติ เหตุปจฺจยา.  เหตุยา ตีณิ.}

\csromanmarkh{22. สนิทสฺสนตฺตฺอิก}\csromanmarkhii{87. สวิตกฺกทุก}\csromanheadernomark{2}{22. \textbf{สนิทสฺสนตฺตฺ}อิก 87. สวิตกฺกทุก}
\proseitem{132}{อนิทสฺสน\-อปฺปฏิฆํ สวิตกฺกํ ธมฺมํ ปฏิจฺจ นอนิทสฺสน\-อปฺปฏิโฆ นสวิตกฺโก ธมฺโม อุปฺปชฺชติ เหตุปจฺจยา.  อนิทสฺสน\-อปฺปฏิฆํ สวิตกฺกํ ธมฺมํ ปฏิจฺจ นสนิทสฺสน\-สปฺปฏิโฆ นสวิตกฺโก ธมฺโม อุปฺปชฺชติ เหตุปจฺจยา ฯเปฯ.  อนิทสฺสน\-อปฺปฏิฆํ สวิตกฺกํ ธมฺมํ ปฏิจฺจ นสนิทสฺสน\-สปฺปฏิโฆ นสวิตกฺโก จ นอนิทสฺสน\-สปฺปฏิโฆ นสวิตกฺโก จ ธมฺมา อุปฺปชฺชนฺติ เหตุปจฺจยา.  เหตุยา ฉ.}

\proseitem{133}{อนิทสฺสน\-อปฺปฏิฆํ อวิตกฺกํ ธมฺมํ ปฏิจฺจ นสนิทสฺสน\-สปฺปฏิโฆ นอวิตกฺโก ธมฺโม อุปฺปชฺชติ เหตุปจฺจยา.  อนิทสฺสน\-อปฺปฏิฆํ อวิตกฺกํ ธมฺมํ ปฏิจฺจ นอนิทสฺสน\-สปฺปฏิโฆ นอวิตกฺโก ธมฺโม อุปฺปชฺชติ เหตุปจฺจยา.  อนิทสฺสน\-อปฺปฏิฆํ อวิตกฺกํ ธมฺมํ ปฏิจฺจ นสนิทสฺสน\-สปฺปฏิโฆ นอวิตกฺโก จ นอนิทสฺสน\-สปฺปฏิโฆ นอวิตกฺโก จ ธมฺมา อุปฺปชฺชนฺติ เหตุปจฺจยา.  เหตุยา ตีณิ.}

\csromanmarkh{22. สนิทสฺสน\-ตฺติก}\csromanmarkhii{88. สวิจารทุก}\csromanheadernomark{2}{22. สนิทสฺสน\-ตฺติก 88. สวิจารทุก}
\proseitem{134}{อนิทสฺสน\-อปฺปฏิฆํ สวิจารํ ธมฺมํ ปฏิจฺจ นอนิทสฺสน\-อปฺปฏิโฆ นสวิจาโร ธมฺโม อุปฺปชฺชติ เหตุปจฺจยา.  เหตุยา ฉ.}

\prose{อนิทสฺสน\-อปฺปฏิฆํ อวิจารํ ธมฺมํ ปฏิจฺจ นสนิทสฺสน\-สปฺปฏิโฆ นอวิจาโร ธมฺโม อุปฺปชฺชติ เหตุปจฺจยา.  อนิทสฺสน\-อปฺปฏิฆํ อวิจารํ ธมฺมํ ปฏิจฺจ นอนิทสฺสน\-สปฺปฏิโฆ นอวิจาโร ธมฺโม อุปฺปชฺชติ เหตุปจฺจยา.  อนิทสฺสน\-อปฺปฏิฆํ อวิจารํ ธมฺมํ ปฏิจฺจ นสนิทสฺสน\-สปฺปฏิโฆ นอวิจาโร จ นอนิทสฺสน\-สปฺปฏิโฆ นอวิจาโร จ ธมฺมา อุปฺปชฺชนฺติ เหตุปจฺจยา.  เหตุยา ตีณิ.}

\vspace{\tipitakaparskip}
\csromancenter{สนิทสฺสน\-ตฺติก 91. สุขสหคตทุก 22. สนิทสฺสน\-ตฺติก 89. สปฺปีติกทุก}
\proseitem{135}{\csromanfolio{285}อนิทสฺสน\-อปฺปฏิฆํ สปฺปีติกํ ธมฺมํ ปฏิจฺจ นอนิทสฺสน\-อปฺปฏิโฆ นสปฺปีติโก ธมฺโม อุปฺปชฺชติ เหตุปจฺจยา.  อนิทสฺสน\-อปฺปฏิฆํ สปฺปีติกํ ธมฺมํ ปฏิจฺจ นสนิทสฺสน\-สปฺปฏิโฆ นสปฺปีติโก ธมฺโม อุปฺปชฺชติ เหตุปจฺจยา ฯเปฯ.  อนิทสฺสน\-อปฺปฏิฆํ สปฺปีติกํ ธมฺมํ ปฏิจฺจ นสนิทสฺสน\-สปฺปฏิโฆ นสปฺปีติโก จ นอนิทสฺสน\-สปฺปฏิโฆ นสปฺปีติโก จ ธมฺมา อุปฺปชฺชนฺติ เหตุปจฺจยา.  เหตุยา ฉ.}

\prose{อนิทสฺสน\-อปฺปฏิฆํ อปฺปีติกํ ธมฺมํ ปฏิจฺจ นสนิทสฺสน\-สปฺปฏิโฆ นอปฺปีติโก ธมฺโม อุปฺปชฺชติ เหตุปจฺจยา.  อนิทสฺสน\-อปฺปฏิฆํ อปฺปีติกํ ธมฺมํ ปฏิจฺจ นอนิทสฺสน\-สปฺปฏิโฆ นอปฺปีติโก ธมฺโม อุปฺปชฺชติ เหตุปจฺจยา.  อนิทสฺสน\-อปฺปฏิฆํ อปฺปีติกํ ธมฺมํ ปฏิจฺจ นสนิทสฺสน\-สปฺปฏิโฆ นอปฺปีติโก จ นอนิทสฺสน\-สปฺปฏิโฆ นอปฺปีติโก จ ธมฺมา อุปฺปชฺชนฺติ เหตุปจฺจยา.  เหตุยา ตีณิ.}

\csromanmarkh{22. สนิทสฺสน\-ตฺติก}\csromanmarkhii{90. ปีติสหคตทุก}\csromanheadernomark{2}{22. สนิทสฺสน\-ตฺติก 90. ปีติสหคตทุก}
\proseitem{136}{อนิทสฺสน\-อปฺปฏิฆํ ปีติสหคตํ ธมฺมํ ปฏิจฺจ นอนิทสฺสน\-อปฺปฏิโฆ นปีติสหคโต ธมฺโม อุปฺปชฺชติ เหตุปจฺจยา.  อนิทสฺสน\-อปฺปฏิฆํ ปีติสหคตํ ธมฺมํ ปฏิจฺจ นสนิทสฺสน\-สปฺปฏิโฆ นปีติสหคโต ธมฺโม อุปฺปชฺชติ เหตุปจฺจยา ฯเปฯ.  อนิทสฺสน\-อปฺปฏิฆํ ปีติสหคตํ ธมฺมํ ปฏิจฺจ นสนิทสฺสน\-สปฺปฏิโฆ นปีติสหคโต จ นอนิทสฺสน\-สปฺปฏิโฆ นปีติสหคโต จ ธมฺมา อุปฺปชฺชนฺติ เหตุปจฺจยา.  เหตุยา ฉ.}

\prose{อนิทสฺสน\-อปฺปฏิฆํ นปีติสหคตํ ธมฺมํ ปฏิจฺจ นสนิทสฺสน\-สปฺปฏิโฆ นนปี\-ติ\-สหคโต ธมฺโม อุปฺปชฺชติ เหตุปจฺจยา.  เหตุยา ตีณิ.}

\csromanmarkh{22. สนิทสฺสน\-ตฺติก}\csromanmarkhii{91. สุขสหคตทุก}\csromanheadernomark{2}{22. สนิทสฺสน\-ตฺติก 91. สุขสหคตทุก}
\proseitem{137}{อนิทสฺสน\-อปฺปฏิฆํ สุขสหคตํ ธมฺมํ ปฏิจฺจ นอนิทสฺสน\-อปฺปฏิโฆ นสุขสหคโต ธมฺโม อุปฺปชฺชติ เหตุปจฺจยา.  เหตุยา ฉ.}

\prose{อนิทสฺสน\-อปฺปฏิฆํ นสุขสหคตํ ธมฺมํ ปฏิจฺจ นสนิทสฺสน\-สปฺปฏิโฆ นนสุข\-สหคโต ธมฺโม อุปฺปชฺชติ เหตุปจฺจยา.  เหตุยา ตีณิ.}

\vspace{\tipitakaparskip}
\csromancenter{ติกทุก\-ปฏฺฐาน\-ปาฬิ 22. สนิทสฺสน\-ตฺติก 92. อุเปกฺขาสหคตทุก}
\proseitem{138}{\csromanfolio{286}อนิทสฺสน\-อปฺปฏิฆํ อุเปกฺขา\-สหคตํ ธมฺมํ ปฏิจฺจ นอนิทสฺสน\-อปฺปฏิโฆ นอุเปกฺขาสหคโต ธมฺโม อุปฺปชฺชติ เหตุปจฺจยา.  อนิทสฺสน\-อปฺปฏิฆํ อุเปกฺขา\-สหคตํ ธมฺมํ ปฏิจฺจ นสนิทสฺสน\-สปฺปฏิโฆ นอุเปกฺขาสหคโต ธมฺโม อุปฺปชฺชติ เหตุปจฺจยา ฯเปฯ.  อนิทสฺสน\-อปฺปฏิฆํ อุเปกฺขา\-สหคตํ ธมฺมํ ปฏิจฺจ นสนิทสฺสน\-สปฺปฏิโฆ นอุเปกฺขาสหคโต จ นอนิทสฺสน\-สปฺปฏิโฆ นอุเปกฺขาสหคโต จ ธมฺมา อุปฺปชฺชนฺติ เหตุปจฺจยา.  เหตุยา ฉ.}

\prose{อนิทสฺสน\-อปฺปฏิฆํ นอุเปกฺขาสหคตํ ธมฺมํ ปฏิจฺจ นสนิทสฺสน\-สปฺปฏิโฆ นนอุเปกฺขาสหคโต ธมฺโม อุปฺปชฺชติ เหตุปจฺจยา.  เหตุยา ตีณิ.}

\csromanmarkh{22. สนิทสฺสนตฺตฺอิก}\csromanmarkhii{93. กามาวจรทุก}\csromanheadernomark{2}{22. \textbf{สนิทสฺสนตฺตฺ}อิก 93. กามาวจรทุก}
\proseitem{139}{อนิทสฺสน\-อปฺปฏิฆํ กามาวจรํ ธมฺมํ ปฏิจฺจ นสนิทสฺสน\-สปฺปฏิโฆ นกามาวจโร ธมฺโม อุปฺปชฺชติ เหตุปจฺจยา.  เหตุยา ตีณิ.}

\prose{อนิทสฺสน\-อปฺปฏิฆํ นกามาวจรํ ธมฺมํ ปฏิจฺจ นอนิทสฺสน\-อปฺปฏิโฆ นนกามาวจโร ธมฺโม อุปฺปชฺชติ เหตุปจฺจยา ฯเปฯ.  อนิทสฺสน\-อปฺปฏิฆํ นกามาวจรํ ธมฺมํ ปฏิจฺจ นสนิทสฺสน\-สปฺปฏิโฆ นนกามาวจโร จ นอนิทสฺสน\-สปฺปฏิโฆ นนกามาวจโร จ ธมฺมา อุปฺปชฺชนฺติ เหตุปจฺจยา.  เหตุยา ฉ.}

\csromanmarkh{22. สนิทสฺสนตฺตฺอิก}\csromanmarkhii{94. รูปาวจรทุก}\csromanheadernomark{2}{22. \textbf{สนิทสฺสนตฺตฺ}อิก 94. รูปาวจรทุก}
\proseitem{140}{อนิทสฺสน\-อปฺปฏิฆํ รูปาวจรํ ธมฺมํ ปฏิจฺจ นอนิทสฺสน\-อปฺปฏิโฆ นรูปาวจโร ธมฺโม อุปฺปชฺชติ เหตุปจฺจยา.  เหตุยา ฉ.}

\prose{อนิทสฺสน\-อปฺปฏิฆํ นรูปาวจรํ ธมฺมํ ปฏิจฺจ นสนิทสฺสน\-สปฺปฏิโฆ นนรูปาวจโร ธมฺโม อุปฺปชฺชติ เหตุปจฺจยา.  เหตุยา ตีณิ.}

\csromanmarkh{22. สนิทสฺสน\-ตฺติก}\csromanmarkhii{95. อรูปาวจรทุก}\csromanheadernomark{2}{22. สนิทสฺสน\-ตฺติก 95. อรูปาวจรทุก}
\proseitem{141}{อนิทสฺสน\-อปฺปฏิฆํ อรูปาวจรํ ธมฺมํ ปฏิจฺจ นอนิทสฺสน\-อปฺปฏิโฆ นอรูปาวจโร ธมฺโม อุปฺปชฺชติ เหตุปจฺจยา.  เหตุยา ฉ.}

\prose{\csromanfolio{287}22. สนิทสฺสน\-ตฺติก 99. สอุตฺตรทุก}

\prose{อนิทสฺสน\-อปฺปฏิฆํ นอรูปาวจรํ ธมฺมํ ปจฺจยา นสนิทสฺสน\-สปฺปฏิโฆ นนอรูปาวจโร ธมฺโม อุปฺปชฺชติ เหตุปจฺจยา.  เหตุยา ตีณิ.}

\csromanmarkh{22. สนิทสฺสน\-ตฺติก}\csromanmarkhii{96. ปริยาปนฺนทุก}\csromanheadernomark{2}{22. สนิทสฺสน\-ตฺติก 96. ปริยาปนฺนทุก}
\proseitem{142}{อนิทสฺสน\-อปฺปฏิฆํ ปริยาปนฺนํ ธมฺมํ ปจฺจยา นสนิทสฺสน\-สปฺปฏิโฆ นปริยาปนฺโน ธมฺโม อุปฺปชฺชติ เหตุปจฺจยา.  เหตุยา ตีณิ.}

\prose{อนิทสฺสน\-อปฺปฏิฆํ อปริยาปนฺนํ ธมฺมํ ปฏิจฺจ นอนิทสฺสน\-อปฺปฏิโฆ นอปริยาปนฺโน ธมฺโม อุปฺปชฺชติ เหตุปจฺจยา.  อนิทสฺสน\-อปฺปฏิฆํ อปริยาปนฺนํ ธมฺมํ ปฏิจฺจ นสนิทสฺสน\-สปฺปฏิโฆ นอปริยาปนฺโน ธมฺโม อุปฺปชฺชติ เหตุปจฺจยา ฯเปฯ.  อนิทสฺสน\-อปฺปฏิฆํ อปริยาปนฺนํ ธมฺมํ ปฏิจฺจ นสนิทสฺสน\-สปฺปฏิโฆ นอปริยาปนฺโน จ นอนิทสฺสน\-สปฺปฏิโฆ นอปริยาปนฺโน จ ธมฺมา อุปฺปชฺชนฺติ เหตุปจฺจยา.  เหตุยา ฉ.}

\csromanmarkh{22. สนิทสฺสนตฺตฺอิก}\csromanmarkhii{97. นิยฺยานิกทุก}\csromanheadernomark{2}{22. \textbf{สนิทสฺสนตฺตฺ}อิก 97. นิยฺยานิกทุก}
\proseitem{143}{อนิทสฺสน\-อปฺปฏิฆํ นิยฺยานิกํ ธมฺมํ ปฏิจฺจ นอนิทสฺสน\-อปฺปฏิโฆ นนิยฺยานิโก ธมฺโม อุปฺปชฺชติ เหตุปจฺจยา.  เหตุยา ฉ.}

\prose{อนิทสฺสน\-อปฺปฏิฆํ อนิยฺยานิกํ ธมฺมํ ปจฺจยา นสนิทสฺสน\-สปฺปฏิโฆ นอนิยฺยานิโก ธมฺโม อุปฺปชฺชติ เหตุปจฺจยา.  เหตุยา ตีณิ.}

\csromanmarkh{22. สนิทสฺสน\-ตฺติก}\csromanmarkhii{98. นิยตทุก}\csromanheadernomark{2}{22. สนิทสฺสน\-ตฺติก 98. นิยตทุก}
\proseitem{144}{อนิทสฺสน\-อปฺปฏิฆํ นิยตํ ธมฺมํ ปฏิจฺจ นอนิทสฺสน\-อปฺปฏิโฆ นนิยโต ธมฺโม อุปฺปชฺชติ เหตุปจฺจยา.  เหตุยา ฉ.}

\prose{อนิทสฺสน\-อปฺปฏิฆํ อนิยตํ ธมฺมํ ปจฺจยา นสนิทสฺสน\-สปฺปฏิโฆ นอนิยโต ธมฺโม อุปฺปชฺชติ เหตุปจฺจยา.  เหตุยา ตีณิ.}

\prose{22. สนิทสฺสน\-ตฺติก 99. สอุตฺตรทุก}

\proseitem{145}{อนิทสฺสน\-อปฺปฏิฆํ สอุตฺตรํ ธมฺมํ ปจฺจยา นสนิทสฺสน\-สปฺปฏิโฆ นสอุตฺตโร ธมฺโม อุปฺปชฺชติ เหตุปจฺจยา.  เหตุยา ตีณิ.}

\prose{อนิทสฺสน\-อปฺปฏิฆํ อนุตฺตรํ ธมฺมํ ปฏิจฺจ นอนิทสฺสน\-อปฺปฏิโฆ นอนุตฺตโร ธมฺโม อุปฺปชฺชติ เหตุปจฺจยา.  เหตุยา ฉ.}

\csromanmarkh{ธมฺมา\-นุโลม\-ปจฺจนีย ติกทุก\-ปฏฺฐาน\-ปาฬิ}\csromanmarkhii{22. สนิทสฺสน\-ตฺติก 100. สรณทุก}\csromanheadernomark{2}{ธมฺมา\-นุโลม\-ปจฺจนีย ติกทุก\-ปฏฺฐาน\-ปาฬิ 22. สนิทสฺสน\-ตฺติก 100. สรณทุก}
\proseitem{146}{\csromanfolio{288}อนิทสฺสน\-อปฺปฏิฆํ สรณํ ธมฺมํ ปฏิจฺจ นอนิทสฺสน\-อปฺปฏิโฆ นสรโณ ธมฺโม อุปฺปชฺชติ เหตุปจฺจยา.  เหตุยา ฉ.}

\prose{อนิทสฺสน\-อปฺปฏิฆํ อรณํ ธมฺมํ ปจฺจยา นสนิทสฺสน\-สปฺปฏิโฆ นอรโณ ธมฺโม อุปฺปชฺชติ เหตุปจฺจยา.  เหตุยา ตีณิ.}

\vspace{\tipitakaparskip}
\csromancenter{(สหชาตวารมฺปิ ฯเปฯ สมฺปยุตฺตวารมฺปิ ปฏิจฺจ\-วาร\-สทิสํ.)}
\proseitem{147}{สนิทสฺสน\-สปฺปฏิโฆ อรโณ ธมฺโม นสนิทสฺสน\-สปฺปฏิฆสฺส นอรณสฺส ธมฺมสฺส อารมฺมณ\-ปจฺจเยน ปจฺจโย.}

\prosewithcloser{(ยถา กุสลตฺติเก ปญฺหาวารสฺส อนุโลมมฺปิ ปจฺจนียมฺปิ อนุโลมปจฺจนียมฺปิ ปจฺจนียานุโลมมฺปิ คณิตํ, เอวํ คเณตพฺพํ.)}{ธมฺมา\-นุโลม\-ปจฺจนีเย ติกทุก\-ปฏฺฐานํ นิฏฺฐิตํ.}
\pitaka{อภิธมฺม\-ปิฏก}
\namakkaram{นโม ตสฺส ภควโต อรหโต สมฺมา\-สมฺพุทฺธสฺส.}
\csromanmarkh{1. กุสลตฺติก}\csromanmarkhii{1. เวทนาตฺติก}\csromanheadernomark{2}{1. \textbf{กุสลตฺติก 1. เวทนาตฺติก}}
\prosecont{\csromanfolio{289}กุสลํ สุขาย เวทนาย สมฺปยุตฺตํ ธมฺมํ ปฏิจฺจ นกุสโล นสุขาย เวทนาย สมฺปยุตฺโต ธมฺโม อุปฺปชฺชติ เหตุปจฺจยา.  กุสลํ สุขาย เวทนาย สมฺปยุตฺตํ ธมฺมํ ปฏิจฺจ นอกุสโล นสุขาย เวทนาย สมฺปยุตฺโต ธมฺโม อุปฺปชฺชติ เหตุปจฺจยา ฯเปฯ.  กุสลํ สุขาย เวทนาย สมฺปยุตฺตํ ธมฺมํ ปฏิจฺจ นกุสโล นสุขาย เวทนาย สมฺปยุตฺโต จ นอกุสโล นสุขาย เวทนาย สมฺปยุตฺโต จ ธมฺมา อุปฺปชฺชนฺติ เหตุปจฺจยา. ปญฺจ. อกุสลํ สุขาย เวทนาย สมฺปยุตฺตํ ธมฺมํ ปฏิจฺจ นอกุสโล นสุขาย เวทนาย สมฺปยุตฺโต ธมฺโม อุปฺปชฺชติ เหตุปจฺจยา.  อกุสลํ สุขาย เวทนาย สมฺปยุตฺตํ ธมฺมํ ปฏิจฺจ นกุสโล นสุขาย เวทนาย สมฺปยุตฺโต ธมฺโม อุปฺปชฺชติ เหตุปจฺจยา ฯเปฯ.  อกุสลํ สุขาย เวทนาย สมฺปยุตฺตํ ธมฺมํ ปฏิจฺจ นกุสโล นสุขาย เวทนาย สมฺปยุตฺโต จ นอกุสโล นสุขาย เวทนาย สมฺปยุตฺโต จ ธมฺมา อุปฺปชฺชนฺติ เหตุปจฺจยา. ปญฺจ. อพฺยากตํ สุขาย เวทนาย สมฺปยุตฺตํ ธมฺมํ ปฏิจฺจ นกุสโล นสุขาย เวทนาย สมฺปยุตฺโต ธมฺโม อุปฺปชฺชติ เหตุปจฺจยา.  อพฺยากตํ สุขาย เวทนาย สมฺปยุตฺตํ ธมฺมํ ปฏิจฺจ นอกุสโล นสุขาย}

\vspace{\tipitakaparskip}
\csromancenter{ติกติก\-ปฏฺฐาน\-ปาฬิ เวทนาย สมฺปยุตฺโต ธมฺโม อุปฺปชฺชติ เหตุปจฺจยา.  อพฺยากตํ สุขาย เวทนาย สมฺปยุตฺตํ ธมฺมํ ปฏิจฺจ นกุสโล นสุขาย เวทนาย สมฺปยุตฺโต จ นอกุสโล นสุขาย เวทนาย สมฺปยุตฺโต จ ธมฺมา อุปฺปชฺชนฺติ เหตุปจฺจยา. ตีณิ. (สํขิตฺตํ.)}
\prose{\csromanfolio{290}เหตุยา เตรส, อารมฺมเณ นว ฯเปฯ อวิคเต เตรส.}

\prose{ปจฺจนีย\-น\-เหตุ นอารมฺมณ}

\proseitem{2}{อพฺยากตํ สุขาย เวทนาย สมฺปยุตฺตํ ธมฺมํ ปฏิจฺจ นกุสโล นสุขาย เวทนาย สมฺปยุตฺโต ธมฺโม อุปฺปชฺชติ นเหตุปจฺจยา. ตีณิ.}

\proseitem{3}{กุสลํ สุขาย เวทนาย สมฺปยุตฺตํ ธมฺมํ ปฏิจฺจ นกุสโล นสุขาย เวทนาย สมฺปยุตฺโต ธมฺโม อุปฺปชฺชติ นอารมฺมณ\-ปจฺจยา.  กุสลํ สุขาย เวทนาย สมฺปยุตฺตํ ธมฺมํ ปฏิจฺจ นอกุสโล นสุขาย เวทนาย สมฺปยุตฺโต ธมฺโม อุปฺปชฺชติ นอารมฺมณ\-ปจฺจยา.  กุสลํ สุขาย เวทนาย สมฺปยุตฺตํ ธมฺมํ ปฏิจฺจ นกุสโล นสุขาย เวทนาย สมฺปยุตฺโต จ นอกุสโล นสุขาย เวทนาย สมฺปยุตฺโต จ ธมฺมา อุปฺปชฺชนฺติ นอารมฺมณ\-ปจฺจยา. ตีณิ.}

\prose{อกุสลํ สุขาย เวทนาย สมฺปยุตฺตํ ธมฺมํ ปฏิจฺจ นอกุสโล นสุขาย เวทนาย สมฺปยุตฺโต ธมฺโม อุปฺปชฺชติ นอารมฺมณ\-ปจฺจยา. ตีณิ.}

\prose{อพฺยากตํ สุขาย เวทนาย สมฺปยุตฺตํ ธมฺมํ ปฏิจฺจ นกุสโล นสุขาย เวทนาย สมฺปยุตฺโต ธมฺโม อุปฺปชฺชติ นอารมฺมณ\-ปจฺจยา.  อพฺยากตํ สุขาย เวทนาย สมฺปยุตฺตํ ธมฺมํ ปฏิจฺจ นอกุสโล นสุขาย เวทนาย สมฺปยุตฺโต ธมฺโม อุปฺปชฺชติ นอารมฺมณ\-ปจฺจยา.  อพฺยากตํ สุขาย เวทนาย สมฺปยุตฺตํ ธมฺมํ ปฏิจฺจ นกุสโล นสุขาย เวทนาย สมฺปยุตฺโต จ นอกุสโล นสุขาย เวทนาย สมฺปยุตฺโต จ ธมฺมา อุปฺปชฺชนฺติ นอารมฺมณ\-ปจฺจยา. ตีณิ.}

\proseitem{4}{นเหตุยา ตีณิ, นอารมฺมเณ นว, นอธิปติยา เตรส ฯเปฯ โนวิคเต นว.}

\prose{เหตุปจฺจยา นอารมฺมเณ นว.}

\prose{นเหตุปจฺจยา อารมฺมเณ ตีณิ.}

\csromanmarkh{1. กุสลตฺติก}\csromanmarkhii{1. เวทนาตฺติก}\csromanheadernomark{2}{1. กุสลตฺติก 1. เวทนาตฺติก}
\prose{\csromanfolio{291}(สหชาตวารมฺปิ ปจฺจยวารมฺปิ นิสฺสยวารมฺปิ สํสฏฺฐ\-วารมฺปิ สมฺปยุตฺตวารมฺปิ ปฏิจฺจ\-วาร\-สทิสํ.)}

\proseitem{5}{กุสโล สุขาย เวทนาย สมฺปยุตฺโต ธมฺโม นกุสลสฺส นสุขาย เวทนาย สมฺปยุตฺตสฺส ธมฺมสฺส เหตุปจฺจเยน ปจฺจโย ฯเปฯ.}

\prose{อารมฺมเณ อฏฺฐารส. (ปญฺหาวารมฺปิ วิตฺถาเรตพฺพํ.)}

\prose{อกุสลปทเหตุ}

\proseitem{6}{อกุสลํ ทุกฺขาย เวทนาย สมฺปยุตฺตํ ธมฺมํ ปฏิจฺจ นอกุสโล นทุกฺขาย เวทนาย สมฺปยุตฺโต ธมฺโม อุปฺปชฺชติ เหตุปจฺจยา.  อกุสลํ ทุกฺขาย เวทนาย สมฺปยุตฺตํ ธมฺมํ ปฏิจฺจ นกุสโล นทุกฺขาย เวทนาย สมฺปยุตฺโต ธมฺโม อุปฺปชฺชติ เหตุปจฺจยา ฯเปฯ.  อกุสลํ ทุกฺขาย เวทนาย สมฺปยุตฺตํ ธมฺมํ ปฏิจฺจ นกุสโล นทุกฺขาย เวทนาย สมฺปยุตฺโต จ นอกุสโล นทุกฺขาย เวทนาย สมฺปยุตฺโต จ ธมฺมา อุปฺปชฺชนฺติ เหตุปจฺจยา.  เหตุยา ปญฺจ, อารมฺมเณ ฉ ฯเปฯ อวิคเต อฏฺฐ.}

\prose{อพฺยากตปทปจฺจนีย}

\proseitem{7}{อพฺยากตํ ทุกฺขาย เวทนาย สมฺปยุตฺตํ ธมฺมํ ปฏิจฺจ นกุสโล นทุกฺขาย เวทนาย สมฺปยุตฺโต ธมฺโม อุปฺปชฺชติ นเหตุปจฺจยา.  อพฺยากตํ ทุกฺขาย เวทนาย สมฺปยุตฺตํ ธมฺมํ ปฏิจฺจ นอกุสโล นทุกฺขาย เวทนาย สมฺปยุตฺโต ธมฺโม อุปฺปชฺชติ นเหตุปจฺจยา.  อพฺยากตํ ทุกฺขาย เวทนาย สมฺปยุตฺตํ ธมฺมํ ปฏิจฺจ นกุสโล นทุกฺขาย เวทนาย สมฺปยุตฺโต จ นอกุสโล นทุกฺขาย เวทนาย สมฺปยุตฺโต จ ธมฺมา อุปฺปชฺชนฺติ เหตุปจฺจยา.}

\prose{นเหตุยา ตีณิ.}

\vspace{\tipitakaparskip}
\csromancenter{(สหชาตวารมฺปิ ฯเปฯ ปญฺหาวารมฺปิ วิตฺถาเรตพฺพํ.)}
\proseitem{8}{กุสลํ อทุกฺขม\-สุขาย เวทนาย สมฺปยุตฺตํ ธมฺมํ ปฏิจฺจ นกุสโล นอทุกฺขม\-สุขาย เวทนาย สมฺปยุตฺโต ธมฺโม อุปฺปชฺชติ เหตุปจฺจยา. (สํขิตฺตํ.) . เหตุยา เตรส. (สพฺพตฺถ วิตฺถาโร.)}

\csromanmarkh{ธมฺมา\-นุโลม\-ปจฺจนีย ติกติก\-ปฏฺฐาน\-ปาฬิ}\csromanmarkhii{1. กุสลตฺติก 2. วิปากตฺติก}\csromanheadernomark{2}{ธมฺมา\-นุโลม\-ปจฺจนีย ติกติก\-ปฏฺฐาน\-ปาฬิ 1. กุสลตฺติก 2. วิปากตฺติก}
\proseitem{9}{\csromanfolio{292}อพฺยากตํ วิปากํ ธมฺมํ ปฏิจฺจ นกุสโล นวิปาโก ธมฺโม อุปฺปชฺชติ เหตุปจฺจยา.  เหตุยา ตีณิ.}

\proseitem{10}{กุสลํ วิปาก\-ธมฺม\-ธมฺมํ ปฏิจฺจ นกุสโล นวิปาก\-ธมฺม\-ธมฺโม อุปฺปชฺชติ เหตุปจฺจยา.  กุสลํ วิปาก\-ธมฺม\-ธมฺมํ ปฏิจฺจ นอกุสโล นวิปาก\-ธมฺม\-ธมฺโม อุปฺปชฺชติ เหตุปจฺจยา.  กุสลํ วิปาก\-ธมฺม\-ธมฺมํ ปฏิจฺจ นกุสโล นวิปาก\-ธมฺม\-ธมฺโม จ นอกุสโล นวิปาก\-ธมฺม\-ธมฺโม จ ธมฺมา อุปฺปชฺชนฺติ เหตุปจฺจยา. ตีณิ.}

\prose{อกุสลํ วิปาก\-ธมฺม\-ธมฺมํ ปฏิจฺจ นอกุสโล นวิปาก\-ธมฺม\-ธมฺโม อุปฺปชฺชติ เหตุปจฺจยา.  อกุสลํ วิปาก\-ธมฺม\-ธมฺมํ ปฏิจฺจ นกุสโล นวิปาก\-ธมฺม\-ธมฺโม อุปฺปชฺชติ เหตุปจฺจยา.  อกุสลํ วิปาก\-ธมฺม\-ธมฺมํ ปฏิจฺจ นกุสโล นวิปาก\-ธมฺม\-ธมฺโม จ นอกุสโล นวิปาก\-ธมฺม\-ธมฺโม จ ธมฺมา อุปฺปชฺชนฺติ เหตุปจฺจยา. (3). เหตุยา ฉ.}

\proseitem{11}{อพฺยากตํ เนว\-วิปาก\-น\-วิปากธมฺม\-ธมฺมํ ปฏิจฺจ นกุสโล นเน\-ว\-วิปาก\-น\-วิปาก\-ธมฺม\-ธมฺโม อุปฺปชฺชติ เหตุปจฺจยา.  อพฺยากตํ เนว\-วิปาก\-น\-วิปากธมฺม\-ธมฺมํ ปฏิจฺจ นอกุสโล นเน\-ว\-วิปาก\-น\-วิปาก\-ธมฺม\-ธมฺโม อุปฺปชฺชติ เหตุปจฺจยา.  อพฺยากตํ เนว\-วิปาก\-น\-วิปากธมฺม\-ธมฺมํ ปฏิจฺจ นกุสโล นเน\-ว\-วิปาก\-น\-วิปาก\-ธมฺม\-ธมฺโม จ นอกุสโล เนว\-วิปาก\-น\-วิปาก\-ธมฺม\-ธมฺโม จ ธมฺมา อุปฺปชฺชนฺติ เหตุปจฺจยา.  เหตุยา ตีณิ.}

\csromanmarkh{1. กุสลตฺติก}\csromanmarkhii{3. อุปาทินฺนตฺติก}\csromanheadernomark{2}{1. \textbf{กุสลตฺติก 3. อุปาทินฺนตฺติก}}
\proseitem{12}{อพฺยากตํ อุปาทินฺนุ\-ปาทานิยํ ธมฺมํ ปฏิจฺจ นกุสโล นอุปาทินฺนุปาทานิโย ธมฺโม อุปฺปชฺชติ เหตุปจฺจยา.  อพฺยากตํ อุปาทินฺนุ\-ปาทานิยํ ธมฺมํ ปฏิจฺจ นอกุสโล นอุปาทินฺนุปาทานิโย ธมฺโม อุปฺปชฺชติ เหตุปจฺจยา.  อพฺยากตํ อุปาทินฺนุ\-ปาทานิยํ ธมฺมํ ปฏิจฺจ นกุสโล นอุปาทินฺนุปาทานิโย จ นอกุสโล นอุปาทินฺนุปาทานิโย จ ธมฺมา อุปฺปชฺชนฺติ เหตุปจฺจยา.  เหตุยา ตีณิ.}

\proseitem{13}{กุสโล อนุปาทินฺนุ\-ปาทานิโย ธมฺโม นกุสลสฺส นอนุปาทินฺนุ\-ปาทานิยสฺส ธมฺมสฺส อารมฺมณ\-ปจฺจเยน ปจฺจโย.  กุสโล}

\vspace{\tipitakaparskip}
\csromancenter{4. สํกิลิฏฺฐ\-ตฺติก อนุปาทินฺนุ\-ปาทานิโย ธมฺโม นอกุสลสฺส นอนุปาทินฺนุ\-ปาทานิยสฺส ธมฺมสฺส อารมฺมณ\-ปจฺจเยน ปจฺจโย.  กุสโล อนุปาทินฺนุ\-ปาทานิโย ธมฺโม นกุสลสฺส นอนุปาทินฺนุ\-ปาทานิยสฺส จ นอกุสลสฺส นอนุปาทินฺนุ\-ปาทานิยสฺส จ ธมฺมสฺส อารมฺมณ\-ปจฺจเยน ปจฺจโย. ตีณิ.}
\prose{\csromanfolio{293}อกุสเล ตีณิ. อพฺยากตํ อนุปาทินฺนุ\-ปาทานิเย ตีณิเยว ฯเปฯ.}

\proseitem{14}{กุสลํ อนุปาทินฺน\-อนุปาทานิยํ ธมฺมํ ปฏิจฺจ นกุสโล นอนุปาทินฺน\-อนุปาทานิโย ธมฺโม อุปฺปชฺชติ เหตุปจฺจยา. ตีณิ.}

\prose{อพฺยากตํ อนุปาทินฺน\-อนุปาทานิยํ ธมฺมํ ปฏิจฺจ นกุสโล นอนุปาทินฺน\-อนุปาทานิโย ธมฺโม อุปฺปชฺชติ เหตุปจฺจยา. ตีณิ.  เหตุยา ฉ.}

\csromanmarkh{1. กุสลตฺติก}\csromanmarkhii{4. สํกิลิฏฺฐ\-ตฺติก}\csromanheadernomark{2}{1. กุสลตฺติก 4. สํกิลิฏฺฐ\-ตฺติก}
\proseitem{15}{อกุสลํ สํกิลิฏฺฐ\-สํกิเลสิกํ ธมฺมํ ปฏิจฺจ นอกุสโล นสํกิลิฏฺฐ\-สํกิเลสิโก ธมฺโม อุปฺปชฺชติ เหตุปจฺจยา.  เหตุยา ตีณิ.}

\prose{อพฺยากตํ อสํกิลิฏฺฐ\-สํกิเลสิกํ ธมฺมํ ปจฺจยา นอพฺยากโต นอสํกิลิฏฺฐ\-สํกิเลสิโก ธมฺโม อุปฺปชฺชติ เหตุปจฺจยา.  อพฺยากตํ อสํกิลิฏฺฐ\-สํกิเลสิกํ ธมฺมํ ปจฺจยา นกุสโล นอสํกิลิฏฺฐ\-สํกิเลสิโก ธมฺโม อุปฺปชฺชติ เหตุปจฺจยา.  อพฺยากตํ อสํกิลิฏฺฐ\-สํกิเลสิกํ ธมฺมํ ปจฺจยา นอกุสโล นอสํกิลิฏฺฐ\-สํกิเลสิโก ธมฺโม อุปฺปชฺชติ เหตุปจฺจยา.  อพฺยากตํ อสํกิลิฏฺฐ\-สํกิเลสิกํ ธมฺมํ ปจฺจยา นกุสโล นอสํกิลิฏฺฐ\-สํกิเลสิโก จ นอพฺยากโต นอสํกิลิฏฺฐ\-สํกิเลสิโก จ ธมฺมา อุปฺปชฺชนฺติ เหตุปจฺจยา.  อพฺยากตํ อสํกิลิฏฺฐ\-สํกิเลสิกํ ธมฺมํ ปจฺจยา นอกุสโล นอสํกิลิฏฺฐสํกิเลโก จ นอพฺยากโต นอสํกิลิฏฺฐ\-สํกิเลสิโก จ ธมฺมา อุปฺปชฺชนฺติ เหตุปจฺจยา.  อพฺยากตํ อสํกิลิฏฺฐ\-สํกิเลสิกํ ธมฺมํ ปจฺจยา นกุสโล นอสํกิลิฏฺฐ\-สํกิเลสิโก จ นอกุสโล นอสํกิลิฏฺฐ\-สํกิเลสิโก จ ธมฺมา อุปฺปชฺชนฺติ เหตุปจฺจยา. เหตุยา ฉ.}

\proseitem{16}{กุสลํ อสํกิลิฏฺฐ\-อสํกิเลสิกํ ธมฺมํ ปฏิจฺจ นกุสโล นอสํกิลิฏฺฐ\-อสํกิเลสิโก ธมฺโม อุปฺปชฺชติ เหตุปจฺจยา. ตีณิ.}

\csromanheader{2}{ธมฺมา\-นุโลม\-ปจฺจนีย ติกติก\-ปฏฺฐาน\-ปาฬิ}
\prose{\csromanfolio{294}อพฺยากตํ อสํกิลิฏฺฐ\-อสํกิเลสิกํ ธมฺมํ ปฏิจฺจ นกุสโล นอสํกิลิฏฺฐ\-อสํกิเลสิโก ธมฺโม อุปฺปชฺชติ เหตุปจฺจยา. ตีณิ.  เหตุยา ฉ.}

\csromanmarkh{1. กุสลตฺติก}\csromanmarkhii{5. วิตกฺกตฺติก}\csromanheadernomark{2}{1. \textbf{กุสลตฺติก 5. วิตกฺกตฺติก}}
\proseitem{17}{กุสลํ สวิตกฺก\-สวิจารํ ธมฺมํ ปฏิจฺจ นกุสโล นสวิตกฺก\-สวิจาโร ธมฺโม อุปฺปชฺชติ เหตุปจฺจยา.  เหตุยา เตรส. (กุสเล ปญฺจ. อกุสเล ปญฺจ. อพฺยากเต ตีณิ.)}

\prose{กุสลํ อวิตกฺก\-วิจารมตฺตํ ธมฺมํ ปฏิจฺจ นกุสโล นอวิตกฺก\-วิจารมตฺโต ธมฺโม อุปฺปชฺชติ เหตุปจฺจยา.  เหตุยา เตรส.}

\prose{กุสลํ อวิตกฺก\-อวิจารํ ธมฺมํ ปฏิจฺจ นกุสโล นอวิตกฺก\-อวิจาโร ธมฺโม อุปฺปชฺชติ เหตุปจฺจยา. ตีณิ.}

\prose{อพฺยากตํ อวิตกฺก\-อวิจารํ ธมฺมํ ปฏิจฺจ นกุสโล นอวิตกฺก\-อวิจาโร ธมฺโม อุปฺปชฺชติ เหตุปจฺจยา. ตีณิ.  เหตุยา ฉ.}

\csromanmarkh{1. กุสลตฺติก}\csromanmarkhii{6. ปีติตฺติก}\csromanheadernomark{2}{1. \textbf{กุสลตฺติก 6. ปีติตฺติก}}
\proseitem{18}{กุสลํ ปีติสหคตํ ธมฺมํ ปฏิจฺจ นกุสโล นปีติสหคโต ธมฺโม อุปฺปชฺชติ เหตุปจฺจยา. ปญฺจ.  เหตุยา เตรส.}

\prose{กุสลํ สุขสหคตํ ธมฺมํ ปฏิจฺจ นกุสโล นสุขสหคโต ธมฺโม อุปฺปชฺชติ เหตุปจฺจยา. ปญฺจ.  เหตุยา เตรส.}

\prose{กุสลํ อุเปกฺขา\-สหคตํ ธมฺมํ ปฏิจฺจ นกุสโล นอุเปกฺขาสหคโต ธมฺโม อุปฺปชฺชติ เหตุปจฺจยา. ปญฺจ.  เหตุยา เตรส.}

\csromanmarkh{1. กุสลตฺติก}\csromanmarkhii{7. ทสฺสนตฺติก}\csromanheadernomark{2}{1. \textbf{กุสลตฺติก 7. ทสฺสนตฺติก}}
\proseitem{19}{อกุสลํ ทสฺสเนน ปหาตพฺพํ ธมฺมํ ปฏิจฺจ นอกุสโล นทสฺสเนน ปหาตพฺโพ ธมฺโม อุปฺปชฺชติ เหตุปจฺจยา.  เหตุยา ตีณิ.}

\prose{อกุสลํ ภาวนาย ปหาตพฺพํ ธมฺมํ ปฏิจฺจ นอกุสโล นภาวนาย ปหาตพฺโพ ธมฺโม อุปฺปชฺชติ เหตุปจฺจยา.  เหตุยา ตีณิ.}

\csromanmarkh{1. กุสลตฺติก}\csromanmarkhii{10. เสกฺขตฺติก}\csromanheadernomark{2}{1. กุสลตฺติก 10. เสกฺขตฺติก}
\prose{\csromanfolio{295}อพฺยากตํ เนวทสฺสเนน นภาวนาย ปหาตพฺพํ ธมฺมํ ปจฺจยา นอพฺยากโต นเน\-ว\-ทสฺสเนน นภาวนาย ปหาตพฺโพ ธมฺโม อุปฺปชฺชติ เหตุปจฺจยา.  เหตุยา ตีณิ.}

\csromanmarkh{1. กุสลตฺติก}\csromanmarkhii{8. ทสฺสน\-เหตุ\-ตฺติก}\csromanheadernomark{2}{1. กุสลตฺติก 8. ทสฺสน\-เหตุ\-ตฺติก}
\proseitem{20}{อกุสลํ ทสฺสเนน ปหาตพฺพ\-เหตุกํ ธมฺมํ ปฏิจฺจ นอกุสโล นทสฺสเนน ปหาตพฺพ\-เหตุโก ธมฺโม อุปฺปชฺชติ เหตุปจฺจยา.  เหตุยา ตีณิ.}

\prose{อกุสลํ ภาวนาย ปหาตพฺพ\-เหตุกํ ธมฺมํ ปฏิจฺจ นอกุสโล นภาวนาย ปหาตพฺพ\-เหตุโก ธมฺโม อุปฺปชฺชติ เหตุปจฺจยา.  เหตุยา ตีณิ.}

\prose{อกุสลํ เนวทสฺสเนน นภาวนาย ปหาตพฺพ\-เหตุกํ ธมฺมํ ปฏิจฺจ นกุสโล นเน\-ว\-ทสฺสเนน นภาวนาย ปหาตพฺพ\-เหตุโก ธมฺโม อุปฺปชฺชติ เหตุปจฺจยา.  เหตุยา ตีณิ.}

\prose{อพฺยากตํ เนวทสฺสเนน นภาวนาย ปหาตพฺพ\-เหตุกํ ธมฺมํ ปจฺจยา นอพฺยากโต นเน\-ว\-ทสฺสเนน นภาวนาย ปหาตพฺพ\-เหตุโก ธมฺโม อุปฺปชฺชติ เหตุปจฺจยา.  เหตุยา ตีณิ.}

\csromanmarkh{1. กุสลตฺติก}\csromanmarkhii{9. อาจยคามิตฺติก}\csromanheadernomark{2}{1. \textbf{กุสลตฺติก 9. อาจยคามิตฺติก}}
\proseitem{21}{กุสลํ อาจยคามิํ ธมฺมํ ปฏิจฺจ นกุสโล นอาจยคามี ธมฺโม อุปฺปชฺชติ เหตุปจฺจยา.  เหตุยา ฉ.}

\prose{กุสลํ อปจยคามิํ ธมฺมํ ปฏิจฺจ นกุสโล นอปจยคามี ธมฺโม อุปฺปชฺชติ เหตุปจฺจยา.  เหตุยา ตีณิ.}

\prose{อพฺยากตํ เนวา\-จยคามิ\-นา\-ปจยคามิํ ธมฺมํ ปจฺจยา นอพฺยากโต นเน\-วา\-จยคามิ\-นา\-ปจยคามี ธมฺโม อุปฺปชฺชติ เหตุปจฺจยา.  เหตุยา ปญฺจ.}

\csromanmarkh{1. กุสลตฺติก}\csromanmarkhii{10. เสกฺขตฺติก}\csromanheadernomark{2}{1. กุสลตฺติก 10. เสกฺขตฺติก}
\proseitem{22}{กุสลํ เสกฺขํ ธมฺมํ ปฏิจฺจ นกุสโล นเสกฺโข ธมฺโม อุปฺปชฺชติ เหตุปจฺจยา. (เทฺว มูลานิ.) . เหตุยา ฉ.}

\csromanheader{2}{ธมฺมา\-นุโลม\-ปจฺจนีย ติกติก\-ปฏฺฐาน\-ปาฬิ}
\prose{\csromanfolio{296}อพฺยากตํ อเสกฺขํ ธมฺมํ ปฏิจฺจ นกุสโล นอเสกฺโข ธมฺโม อุปฺปชฺชติ เหตุปจฺจยา.  เหตุยา ตีณิ.}

\prose{อพฺยากตํ เนว\-เสกฺข\-นาเสกฺขํ ธมฺมํ ปจฺจยา นอพฺยากโต นเน\-ว\-เสกฺข\-นา\-เสกฺโข ธมฺโม อุปฺปชฺชติ เหตุปจฺจยา.  เหตุยา ปญฺจ.}

\csromanmarkh{1. กุสลตฺติก}\csromanmarkhii{11. ปริตฺตตฺติก}\csromanheadernomark{2}{1. \textbf{กุสลตฺติก} 11. ปริตฺตตฺติก}
\proseitem{23}{อพฺยากตํ ปริตฺตํ ธมฺมํ ปฏิจฺจ นกุสโล นปริตฺโต ธมฺโม อุปฺปชฺชติ เหตุปจฺจยา.  เหตุยา ตีณิ.}

\prose{กุสลํ มหคฺคตํ ธมฺมํ ปฏิจฺจ นกุสโล นมหคฺคโต ธมฺโม อุปฺปชฺชติ เหตุปจฺจยา.  เหตุยา ฉ.}

\prose{กุสลํ อปฺปมาณํ ธมฺมํ ปฏิจฺจ นกุสโล นอปฺปมาโณ ธมฺโม อุปฺปชฺชติ เหตุปจฺจยา.  เหตุยา ฉ.}

\csromanmarkh{1. กุสลตฺติก}\csromanmarkhii{12. ปริตฺตา\-รมฺมณ\-ตฺติก}\csromanheadernomark{2}{1. กุสลตฺติก 12. ปริตฺตา\-รมฺมณ\-ตฺติก}
\proseitem{24}{กุสลํ ปริตฺตา\-รมฺมณํ ธมฺมํ ปฏิจฺจ นกุสโล นปริตฺตา\-รมฺมโณ ธมฺโม อุปฺปชฺชติ เหตุปจฺจยา.  เหตุยา นว.}

\prose{กุสลํ มหคฺคตา\-รมฺมณํ ธมฺมํ ปฏิจฺจ นกุสโล นมหคฺคตา\-รมฺมโณ ธมฺโม อุปฺปชฺชติ เหตุปจฺจยา.  เหตุยา นว.}

\prose{กุสลํ อปฺปมาณา\-รมฺมณํ ธมฺมํ ปฏิจฺจ นกุสโล นอปฺปมาณารมฺมโณ ธมฺโม อุปฺปชฺชติ เหตุปจฺจยา.  เหตุยา ฉ.}

\csromanmarkh{1. กุสลตฺติก}\csromanmarkhii{13. หีนตฺติก}\csromanheadernomark{2}{1. \textbf{กุสลตฺติก} 13. หีนตฺติก}
\proseitem{25}{อกุสลํ หีนํ ธมฺมํ ปฏิจฺจ นอกุสโล นหีโน ธมฺโม อุปฺปชฺชติ เหตุปจฺจยา.  เหตุยา ตีณิ.}

\prose{อพฺยากตํ มชฺฌิมํ ธมฺมํ ปจฺจยา นอพฺยากโต นมชฺฌิโม ธมฺโม อุปฺปชฺชติ เหตุปจฺจยา.  เหตุยา ฉ.}

\prose{กุสลํ ปณีตํ ธมฺมํ ปฏิจฺจ นกุสโล นปณีโต ธมฺโม อุปฺปชฺชติ เหตุปจฺจยา.  เหตุยา ฉ.}

\vspace{\tipitakaparskip}
\csromancenter{18. อตีตา\-รมฺมณ\-ตฺติก 1. กุสลตฺติก 14. มิจฺฉตฺตนิยตฺตตฺติก}
\proseitem{26}{\csromanfolio{297}อกุสลํ มิจฺฉตฺต\-นิยตํ ธมฺมํ ปฏิจฺจ นอกุสโล นมิจฺฉตฺต\-นิยโต ธมฺโม อุปฺปชฺชติ เหตุปจฺจยา.  เหตุยา ตีณิ.}

\prose{กุสลํ สมฺมตฺต\-นิยตํ ธมฺมํ ปฏิจฺจ นกุสโล นสมฺมตฺตนิยโต ธมฺโม อุปฺปชฺชติ เหตุปจฺจยา.  เหตุยา ตีณิ.}

\prose{อพฺยากตํ อนิยตํ ธมฺมํ ปจฺจยา นอพฺยากโต นอนิยโต ธมฺโม อุปฺปชฺชติ เหตุปจฺจยา.  เหตุยา ปญฺจ.}

\csromanmarkh{1. กุสลตฺติก}\csromanmarkhii{15. มคฺคา\-รมฺมณ\-ตฺติก}\csromanheadernomark{2}{1. กุสลตฺติก 15. มคฺคา\-รมฺมณ\-ตฺติก}
\proseitem{27}{กุสลํ มคฺคารมฺมณํ ธมฺมํ ปฏิจฺจ นกุสโล นมคฺคารมฺมโณ ธมฺโม อุปฺปชฺชติ เหตุปจฺจยา.  เหตุยา ฉ.}

\prose{กุสลํ มคฺคเหตุกํ ธมฺมํ ปฏิจฺจ นกุสโล นมคฺคเหตุโก ธมฺโม อุปฺปชฺชติ เหตุปจฺจยา.  เหตุยา ตีณิ.}

\prose{กุสลํ มคฺคาธิปติํ ธมฺมํ ปฏิจฺจ นกุสโล นมคฺคาธิปติ ธมฺโม อุปฺปชฺชติ เหตุปจฺจยา.  เหตุยา ฉ.}

\csromanmarkh{1. กุสลตฺติก}\csromanmarkhii{16. อุปฺปนฺนตฺติก}\csromanheadernomark{2}{1. \textbf{กุสลตฺติก} 16. อุปฺปนฺนตฺติก}
\proseitem{28}{กุสโล อนุปฺปนฺโน ธมฺโม นกุสลสฺส นอนุปฺปนฺนสฺส ธมฺมสฺส อารมฺมณ\-ปจฺจเยน ปจฺจโย.  อารมฺมเณ อฏฺฐารส. (อุปฺปาที อนุปฺปนฺน\-สทิสํ.)}

\csromanmarkh{1. กุสลตฺติก}\csromanmarkhii{17. อตีตตฺติก}\csromanheadernomark{2}{1. \textbf{กุสลตฺติก} 17. อตีตตฺติก}
\proseitem{29}{กุสโล อตีโต ธมฺโม นกุสลสฺส นอตีตสฺส ธมฺมสฺส อารมฺมณ\-ปจฺจเยน ปจฺจโย.}

\prose{อารมฺมเณ อฏฺฐารส. (อนาคตํ อตีตสทิสํ.)}

\csromanmarkh{1. กุสลตฺติก}\csromanmarkhii{18. อตีตา\-รมฺมณ\-ตฺติก}\csromanheadernomark{2}{1. กุสลตฺติก 18. อตีตา\-รมฺมณ\-ตฺติก}
\proseitem{30}{กุสลํ อตีตารมฺมณํ ธมฺมํ ปฏิจฺจ นกุสโล นอตีตารมฺมโณ ธมฺโม อุปฺปชฺชติ เหตุปจฺจยา.  เหตุยา นว.}

\csromanheader{2}{ธมฺมา\-นุโลม\-ปจฺจนีย ติกติก\-ปฏฺฐาน\-ปาฬิ}
\prose{\csromanfolio{298}กุสลํ อนาคตา\-รมฺมณํ ธมฺมํ ปฏิจฺจ นกุสโล นอนาคตา\-รมฺมโณ ธมฺโม อุปฺปชฺชติ เหตุปจฺจยา.  เหตุยา นว.}

\prose{กุสลํ ปจฺจุปฺปนฺนา\-รมฺมณํ ธมฺมํ ปฏิจฺจ นกุสโล นปจฺจุปฺปนฺนา\-รมฺมโณ ธมฺโม อุปฺปชฺชติ เหตุปจฺจยา.  เหตุยา นว.}

\csromanheader{2}{1. กุสลตฺติก 20-19. อชฺฌตฺต\-ตฺติก\-ทฺวย}
\proseitem{31}{กุสโล อชฺฌตฺโต ธมฺโม นอชฺฌตฺตสฺส นกุสลสฺส ธมฺมสฺส อารมฺมณ\-ปจฺจเยน ปจฺจโย. (สํขิตฺตํ.) . อารมฺมเณ อฏฺฐารส.}

\prose{กุสโล พหิทฺธา ธมฺโม นพหิทฺธา นกุสลสฺส ธมฺมสฺส อารมฺมณ\-ปจฺจเยน ปจฺจโย. (สํขิตฺตํ.) . อารมฺมเณ อฏฺฐารส.}

\proseitem{32}{กุสลํ อชฺฌตฺตา\-รมฺมณํ ธมฺมํ ปฏิจฺจ นกุสโล นอชฺฌตฺตา\-รมฺมโณ ธมฺโม อุปฺปชฺชติ เหตุปจฺจยา.  เหตุยา นว.}

\prose{กุสลํ พหิทฺธา\-รมฺมณํ ธมฺมํ ปฏิจฺจ นกุสโล นพหิทฺธา\-รมฺมโณ ธมฺโม อุปฺปชฺชติ เหตุปจฺจยา.  เหตุยา นว.}

\csromanmarkh{1. กุสลตฺติก}\csromanmarkhii{21. สนิทสฺสน\-ตฺติก}\csromanheadernomark{2}{1. กุสลตฺติก 21. สนิทสฺสน\-ตฺติก}
\proseitem{33}{อพฺยากโต สนิทสฺสน\-สปฺปฏิโฆ ธมฺโม นอพฺยากตสฺส นสนิทสฺสน\-สปฺปฏิฆสฺส ธมฺมสฺส อารมฺมณ\-ปจฺจเยน ปจฺจโย.}

\prose{อารมฺมเณ ฉ. (ตีณิ เวทิตกํ กาตพฺพํ.)}

\proseitem{34}{อพฺยากตํ อนิทสฺสน\-สปฺปฏิฆํ ธมฺมํ ปฏิจฺจ กุสโล นอนิทสฺสน\-สปฺปฏิโฆ ธมฺโม อุปฺปชฺชติ เหตุปจฺจยา.  อพฺยากตํ อนิทสฺสน\-สปฺปฏิฆํ ธมฺมํ ปฏิจฺจ นอกุสโล นอนิทสฺสน\-สปฺปฏิโฆ ธมฺโม อุปฺปชฺชติ เหตุปจฺจยา.  อพฺยากตํ อนิทสฺสน\-สปฺปฏิฆํ ธมฺมํ ปฏิจฺจ นกุสโล นอนิทสฺสน\-สปฺปฏิโฆ จ นอกุสโล นอนิทสฺสน\-สปฺปฏิโฆ จ ธมฺมา อุปฺปชฺชนฺติ เหตุปจฺจยา.  เหตุยา ตีณิ.}

\proseitem{35}{กุสลํ อนิทสฺสน\-อปฺปฏิฆํ ธมฺมํ ปฏิจฺจ นกุสโล นอนิทสฺสน\-อปฺปฏิโฆ ธมฺโม อุปฺปชฺชติ เหตุปจฺจยา.  กุสลํ อนิทสฺสน\-อปฺปฏิฆํ ธมฺมํ ปฏิจฺจ นอกุสโล นอนิทสฺสน\-อปฺปฏิโฆ ธมฺโม อุปฺปชฺชติ}

\vspace{\tipitakaparskip}
\csromancenter{เวทนาตฺติก 1. กุสลตฺติก เหตุปจฺจยา.  กุสลํ อนิทสฺสน\-อปฺปฏิฆํ ธมฺมํ ปฏิจฺจ นกุสโล นอนิทสฺสน\-อปฺปฏิโฆ จ นอกุสโล นอนิทสฺสน\-อปฺปฏิโฆ จ ธมฺมา อุปฺปชฺชนฺติ เหตุปจฺจยา. ตีณิ.}
\prose{\csromanfolio{299}อกุสลํ อนิทสฺสน\-อปฺปฏิฆํ ธมฺมํ ปฏิจฺจ นอกุสโล นอนิทสฺสน\-อปฺปฏิโฆ ธมฺโม อุปฺปชฺชติ เหตุปจฺจยา.  อกุสลํ อนิทสฺสน\-อปฺปฏิฆํ ธมฺมํ ปฏิจฺจ นกุสโล นอนิทสฺสน\-อปฺปฏิโฆ ธมฺโม อุปฺปชฺชติ เหตุปจฺจยา.  อกุสลํ อนิทสฺสน\-อปฺปฏิฆํ ธมฺมํ ปฏิจฺจ นกุสโล นอนิทสฺสน\-อปฺปฏิโฆ จ นอกุสโล นอนิทสฺสน\-อปฺปฏิโฆ จ ธมฺมา อุปฺปชฺชนฺติ เหตุปจฺจยา. ตีณิ.}

\prose{อพฺยากตํ อนิทสฺสน\-อปฺปฏิฆํ ธมฺมํ ปฏิจฺจ นกุสโล นอนิทสฺสน\-อปฺปฏิโฆ ธมฺโม อุปฺปชฺชติ เหตุปจฺจยา.  อพฺยากตํ อนิทสฺสน\-อปฺปฏิฆํ ธมฺมํ ปฏิจฺจ นอกุสโล นอนิทสฺสน\-อปฺปฏิโฆ ธมฺโม อุปฺปชฺชติ เหตุปจฺจยา.  อพฺยากตํ อนิทสฺสน\-อปฺปฏิฆํ ธมฺมํ ปฏิจฺจ นกุสโล นอนิทสฺสน\-อปฺปฏิโฆ จ นอกุสโล นอนิทสฺสน\-อปฺปฏิโฆ จ ธมฺมา อุปฺปชฺชนฺติ เหตุปจฺจยา. ตีณิ. กุสลํ อนิทสฺสนอปฺปฏิฆญฺจ อพฺยากตํ อนิทสฺสนอปฺปฏิฆญฺจ ธมฺมํ ปฏิจฺจ นกุสโล นอนิทสฺสน\-อปฺปฏิโฆ ธมฺโม อุปฺปชฺชติ เหตุปจฺจยา. ตีณิ.}

\prose{อกุสลํ อนิทสฺสนอปฺปฏิฆญฺจ อพฺยากตํ อนิทสฺสนอปฺปฏิฆญฺจ ธมฺมํ ปฏิจฺจ นกุสโล นอนิทสฺสน\-อปฺปฏิโฆ ธมฺโม อุปฺปชฺชติ เหตุปจฺจยา. ตีณิ. เหตุยา ปนฺนรส.}

\csromanmarkh{2. เวทนาตฺติก}\csromanmarkhii{1. กุสลตฺติก}\csromanheadernomark{2}{2. เวทนาตฺติก 1. กุสลตฺติก}
\prosecont{สุขาย เวทนาย สมฺปยุตฺตํ กุสลํ ธมฺมํ ปฏิจฺจ นสุขาย เวทนาย สมฺปยุตฺโต นกุสโล ธมฺโม อุปฺปชฺชติ เหตุปจฺจยา.  สุขาย เวทนาย สมฺปยุตฺตํ กุสลํ ธมฺมํ ปฏิจฺจ นทุกฺขาย เวทนาย สมฺปยุตฺโต นกุสโล ธมฺโม อุปฺปชฺชติ เหตุปจฺจยา.  สุขาย เวทนาย สมฺปยุตฺตํ กุสลํ ธมฺมํ ปฏิจฺจ นอทุกฺขม\-สุขาย เวทนาย สมฺปยุตฺโต นกุสโล ธมฺโม อุปฺปชฺชติ เหตุปจฺจยา.  สุขาย เวทนาย สมฺปยุตฺตํ กุสลํ ธมฺมํ ปฏิจฺจ นสุขาย เวทนาย สมฺปยุตฺโต นกุสโล จ นอทุกฺขม\-สุขาย เวทนาย สมฺปยุตฺโต นกุสโล จ ธมฺมา อุปฺปชฺชนฺติ}

\vspace{\tipitakaparskip}
\csromancenter{ติกติก\-ปฏฺฐาน\-ปาฬิ เหตุปจฺจยา.  สุขาย เวทนาย สมฺปยุตฺตํ กุสลํ ธมฺมํ ปฏิจฺจ นทุกฺขาย เวทนาย สมฺปยุตฺโต นกุสโล จ นอทุกฺขม\-สุขาย เวทนาย สมฺปยุตฺโต นกุสโล จ ธมฺมา อุปฺปชฺชนฺติ เหตุปจฺจยา.  สุขาย เวทนาย สมฺปยุตฺตํ กุสลํ ธมฺมํ ปฏิจฺจ นสุขาย เวทนาย สมฺปยุตฺโต นกุสโล จ นทุกฺขาย เวทนาย สมฺปยุตฺโต นกุสโล จ ธมฺมา อุปฺปชฺชนฺติ เหตุปจฺจยา.  สุขาย เวทนาย สมฺปยุตฺตํ กุสลํ ธมฺมํ ปฏิจฺจ นสุขาย เวทนาย สมฺปยุตฺโต นกุสโล จ นทุกฺขาย เวทนาย สมฺปยุตฺโต นกุสโล จ นอทุกฺขม\-สุขาย เวทนาย สมฺปยุตฺโต นกุสโล จ ธมฺมา อุปฺปชฺชนฺติ เหตุปจฺจยา. สตฺต.}
\prose{\csromanfolio{300}อทุกฺขม\-สุขาย เวทนาย สมฺปยุตฺตํ กุสลํ ธมฺมํ ปฏิจฺจ นอทุกฺขม\-สุขาย เวทนาย สมฺปยุตฺโต นกุสโล ธมฺโม อุปฺปชฺชติ เหตุปจฺจยา. สตฺต. เหตุยา จุทฺทส.}

\proseitem{37}{สุขาย เวทนาย สมฺปยุตฺตํ อกุสลํ ธมฺมํ ปฏิจฺจ นสุขาย เวทนาย สมฺปยุตฺโต นอกุสโล ธมฺโม อุปฺปชฺชติ เหตุปจฺจยา. สตฺต.}

\prose{ทุกฺขาย เวทนาย สมฺปยุตฺตํ อกุสลํ ธมฺมํ ปฏิจฺจ นทุกฺขาย เวทนาย สมฺปยุตฺโต นอกุสโล ธมฺโม อุปฺปชฺชติ เหตุปจฺจยา. สตฺต.}

\prose{อทุกฺขม\-สุขาย เวทนาย สมฺปยุตฺตํ อกุสลํ ธมฺมํ ปฏิจฺจ นอทุกฺขม\-สุขาย เวทนาย สมฺปยุตฺโต นอกุสโล ธมฺโม อุปฺปชฺชติ เหตุปจฺจยา. สตฺต. (เอกวีสติ ปญฺหา.)}

\proseitem{38}{สุขาย เวทนาย สมฺปยุตฺโต อพฺยากโต ธมฺโม นสุขาย เวทนาย สมฺปยุตฺตสฺส นอพฺยากตสฺส ธมฺมสฺส อารมฺมณ\-ปจฺจเยน ปจฺจโย. (สํขิตฺตํ.)}

\prose{อารมฺมเณ เอกวีสติ.}

\vspace{\tipitakaparskip}
\csromancenter{(ทุกฺขาย เวทนาย สมฺปยุตฺต\-อพฺยากตมูลํ อทุกฺขม\-สุขาย เวทนาย สมฺปยุตฺตอพฺยากตมูลมฺปิ กาตพฺพํ.)}
\csromanmarkh{3. วิปากตฺติก}\csromanmarkhii{1. กุสลตฺติก}\csromanheadernomark{2}{3. \textbf{วิปากตฺติก 1. กุสลตฺติก}}
\proseitem{39}{วิปาก\-ธมฺม\-ธมฺมํ กุสลํ ธมฺมํ ปฏิจฺจ นวิปาก\-ธมฺม\-ธมฺโม นกุสโล ธมฺโม อุปฺปชฺชติ เหตุปจฺจยา.  เหตุยา ตีณิ.}

\csromanmarkh{6. วิตกฺกตฺติก}\csromanmarkhii{1. กุสลตฺติก}\csromanheadernomark{2}{6. วิตกฺกตฺติก 1. กุสลตฺติก}
\prose{\csromanfolio{301}วิปาก\-ธมฺม\-ธมฺมํ อกุสลํ ธมฺมํ ปฏิจฺจ นวิปาก\-ธมฺม\-ธมฺโม นอกุสโล ธมฺโม อุปฺปชฺชติ เหตุปจฺจยา.  เหตุยา ตีณิ.}

\prose{เนว\-วิปาก\-น\-วิปากธมฺม\-ธมฺมํ อพฺยากตํ ธมฺมํ ปจฺจยา นเน\-ว\-วิปาก\-น\-วิปาก\-ธมฺม\-ธมฺโม นอพฺยากโต ธมฺโม อุปฺปชฺชติ เหตุปจฺจยา.  เหตุยา ตีณิ.}

\csromanmarkh{4. อุปาทินฺนตฺติก}\csromanmarkhii{1. กุสลตฺติก}\csromanheadernomark{2}{4. \textbf{อุปาทินฺนตฺติก 1. กุสลตฺติก}}
\proseitem{40}{อนุปาทินฺนุ\-ปาทานิยํ กุสลํ ธมฺมํ ปฏิจฺจ นอุปาทินฺนุปาทานิโย นกุสโล ธมฺโม อุปฺปชฺชติ เหตุปจฺจยา. ตีณิ.}

\prose{อนุปาทินฺน\-อนุปาทานิยํ กุสลํ ธมฺมํ ปฏิจฺจ นอนุปาทินฺน\-อนุปาทานิโย นกุสโล ธมฺโม อุปฺปชฺชติ เหตุปจฺจยา. ตีณิ. เหตุยา ฉ.}

\prose{อนุปาทินฺนุ\-ปาทานิยํ อกุสลํ ธมฺมํ ปฏิจฺจ นอุปาทินฺนุปาทานิโย นอกุสโล ธมฺโม อุปฺปชฺชติ เหตุปจฺจยา.  เหตุยา ตีณิ.}

\prose{อุปาทินฺนุ\-ปาทานิยํ อพฺยากตํ ธมฺมํ ปจฺจยา นอุปาทินฺนุปาทานิโย นอพฺยากโต ธมฺโม อุปฺปชฺชติ เหตุปจฺจยา.  เหตุยา ปญฺจ.}

\csromanmarkh{5. สํกิลิฏฺฐ\-ตฺติก}\csromanmarkhii{1. กุสลตฺติก}\csromanheadernomark{2}{5. สํกิลิฏฺฐ\-ตฺติก 1. กุสลตฺติก}
\proseitem{41}{อสํกิลิฏฺฐ\-สํกิเลสิกํ กุสลํ ธมฺมํ ปฏิจฺจ นสํกิลิฏฺฐ\-สํกิเลสิโก นกุสโล ธมฺโม อุปฺปชฺชติ เหตุปจฺจยา. ตีณิ.}

\prose{อสํกิลิฏฺฐ\-อสํกิเลสิกํ กุสลํ ธมฺมํ ปฏิจฺจ นอสํกิลิฏฺฐ\-อสํกิเลสิโก นกุสโล ธมฺโม อุปฺปชฺชติ เหตุปจฺจยา. ตีณิ. เหตุยา ฉ.}

\prose{สํกิลิฏฺฐ\-สํกิเลสิกํ อกุสลํ ธมฺมํ ปฏิจฺจ นสํกิลิฏฺฐ\-สํกิเลสิโก นอกุสโล ธมฺโม อุปฺปชฺชติ เหตุปจฺจยา. ตีณิ.}

\prose{อสํกิลิฏฺฐ\-สํกิเลสิกํ อพฺยากตํ ธมฺมํ ปจฺจยา นอสํกิลิฏฺฐ\-สํกิเลสิโก นอพฺยากโต ธมฺโม อุปฺปชฺชติ เหตุปจฺจยา. (สํขิตฺตํ.) . เหตุยา ฉ ฯเปฯ อวิคเต ฉ.}

\csromanmarkh{6. วิตกฺกตฺติก}\csromanmarkhii{1. กุสลตฺติก}\csromanheadernomark{2}{6. วิตกฺกตฺติก 1. กุสลตฺติก}
\proseitem{42}{สวิตกฺก\-สวิจารํ กุสลํ ธมฺมํ ปฏิจฺจ นสวิตกฺก\-สวิจาโร นกุสโล ธมฺโม อุปฺปชฺชติ เหตุปจฺจยา. ตีณิ.}

\csromanheader{2}{ธมฺมา\-นุโลม\-ปจฺจนีย ติกติก\-ปฏฺฐาน\-ปาฬิ}
\prose{\csromanfolio{302}อวิตกฺก\-วิจารมตฺตํ กุสลํ ธมฺมํ ปฏิจฺจ นอวิตกฺก\-วิจารมตฺโต นกุสโล ธมฺโม อุปฺปชฺชติ เหตุปจฺจยา.  เหตุยา ปนฺนรส.}

\prose{สวิตกฺก\-สวิจารํ อกุสลํ ธมฺมํ ปฏิจฺจ นสวิตกฺก\-สวิจาโร นอกุสโล ธมฺโม อุปฺปชฺชติ เหตุปจฺจยา.  เหตุยา นว.}

\csromanmarkh{7. ปีติตฺติก}\csromanmarkhii{1. กุสลตฺติก}\csromanheadernomark{2}{7. \textbf{ปีติตฺติก 1. กุสลตฺติก}}
\proseitem{43}{ปีติสหคตํ กุสลํ ธมฺมํ ปฏิจฺจ นปีติสหคโต นกุสโล ธมฺโม อุปฺปชฺชติ เหตุปจฺจยา. สตฺต.}

\prose{สุขสหคตํ กุสลํ ธมฺมํ ปฏิจฺจ นสุขสหคโต นกุสโล ธมฺโม อุปฺปชฺชติ เหตุปจฺจยา. สตฺต.}

\prose{อุเปกฺขา\-สหคตํ กุสลํ ธมฺมํ ปฏิจฺจ นอุเปกฺขาสหคโต นกุสโล ธมฺโม อุปฺปชฺชติ เหตุปจฺจยา. สตฺต.  (เวทนา\-ตฺติก\-สทิสํ. สํขิตฺตํ.)}

\csromanmarkh{22. สนิทสฺสน\-ตฺติก}\csromanmarkhii{1. กุสลตฺติก}\csromanheadernomark{2}{22. สนิทสฺสน\-ตฺติก 1. กุสลตฺติก}
\proseitem{44}{อนิทสฺสน\-อปฺปฏิฆํ กุสลํ ธมฺมํ ปฏิจฺจ นอนิทสฺสน\-อปฺปฏิโฆ นกุสโล ธมฺโม อุปฺปชฺชติ เหตุปจฺจยา.  อนิทสฺสน\-อปฺปฏิฆํ กุสลํ ธมฺมํ ปฏิจฺจ นสนิทสฺสน\-สปฺปฏิโฆ นกุสโล ธมฺโม อุปฺปชฺชติ เหตุปจฺจยา ฯเปฯ.  อนิทสฺสน\-อปฺปฏิฆํ กุสลํ ธมฺมํ ปฏิจฺจ นสนิทสฺสน\-สปฺปฏิโฆ นกุสโล จ นอนิทสฺสน\-สปฺปฏิโฆ นกุสโล จ ธมฺมา อุปฺปชฺชนฺติ เหตุปจฺจยา.  เหตุยา ฉ.}

\prose{อนิทสฺสน\-อปฺปฏิฆํ อกุสลํ ธมฺมํ ปฏิจฺจ นอนิทสฺสน\-อปฺปฏิโฆ นอกุสโล ธมฺโม อุปฺปชฺชติ เหตุปจฺจยา.  อนิทสฺสน\-อปฺปฏิฆํ อกุสลํ ธมฺมํ ปฏิจฺจ นสนิทสฺสน\-สปฺปฏิโฆ นกุสโล ธมฺโม อุปฺปชฺชติ เหตุปจฺจยา ฯเปฯ.  อนิทสฺสน\-อปฺปฏิฆํ อกุสลํ ธมฺมํ ปฏิจฺจ นสนิทสฺสน\-สปฺปฏิโฆ นกุสโล จ นอนิทสฺสน\-สปฺปฏิโฆ นอกุสโล จ ธมฺมา อุปฺปชฺชนฺติ เหตุปจฺจยา. (สํขิตฺตํ.) . เหตุยา ฉ ฯเปฯ อวิคเต ฉ.}

\prose{อนิทสฺสน\-อปฺปฏิฆํ อพฺยากตํ ธมฺมํ ปจฺจยา นสนิทสฺสน\-สปฺปฏิโฆ นอพฺยากโต ธมฺโม อุปฺปชฺชติ เหตุปจฺจยา.  อนิทสฺสน\-อปฺปฏิฆํ อพฺยากตํ ธมฺมํ ปจฺจยา นอนิทสฺสน\-สปฺปฏิโฆ นอพฺยากโต ธมฺโม อุปฺปชฺชติ}

\vspace{\tipitakaparskip}
\csromancenter{สนิทสฺสน\-ตฺติก 4. อุปาทินฺนตฺติก เหตุปจฺจยา.  อนิทสฺสน\-อปฺปฏิฆํ อพฺยากตํ ธมฺมํ ปจฺจยา นสนิทสฺสน\-สปฺปฏิโฆ นอพฺยากโต จ นอนิทสฺสน\-สปฺปฏิโฆ นอพฺยากโต จ ธมฺมา อุปฺปชฺชนฺติ เหตุปจฺจยา.  เหตุยา ตีณิ ฯเปฯ อวิคเต ตีณิ.}
\csromanmarkh{22. สนิทสฺสน\-ตฺติก}\csromanmarkhii{2. เวทนาตฺติก}\csromanheadernomark{2}{22. สนิทสฺสน\-ตฺติก 2. เวทนาตฺติก}
\proseitem{45}{\csromanfolio{303}อนิทสฺสน\-อปฺปฏิฆํ สุขาย เวทนาย สมฺปยุตฺตํ ธมฺมํ ปฏิจฺจ นอนิทสฺสน\-อปฺปฏิโฆ นสุขาย เวทนาย สมฺปยุตฺโต ธมฺโม อุปฺปชฺชติ เหตุปจฺจยา.  เหตุยา ฉ.}

\prose{อนิทสฺสน\-อปฺปฏิฆํ ทุกฺขาย เวทนาย สมฺปยุตฺตํ ธมฺมํ ปฏิจฺจ นอนิทสฺสน\-อปฺปฏิโฆ นทุกฺขาย เวทนาย สมฺปยุตฺโต ธมฺโม อุปฺปชฺชติ เหตุปจฺจยา.  เหตุยา ฉ.}

\prose{อนิทสฺสน\-อปฺปฏิฆํ อทุกฺขม\-สุขาย เวทนาย สมฺปยุตฺตํ ธมฺมํ ปฏิจฺจ นอนิทสฺสน\-อปฺปฏิโฆ นอทุกฺขม\-สุขาย เวทนาย สมฺปยุตฺโต ธมฺโม อุปฺปชฺชติ เหตุปจฺจยา.  เหตุยา ฉ.}

\csromanmarkh{22. สนิทสฺสน\-ตฺติก}\csromanmarkhii{3. วิปากตฺติก}\csromanheadernomark{2}{22. สนิทสฺสน\-ตฺติก 3. วิปากตฺติก}
\proseitem{46}{อนิทสฺสน\-อปฺปฏิฆํ วิปากํ ธมฺมํ ปฏิจฺจ นอนิทสฺสน\-อปฺปฏิโฆ นวิปาโก ธมฺโม อุปฺปชฺชติ เหตุปจฺจยา.  เหตุยา ฉ.}

\prose{อนิทสฺสน\-อปฺปฏิฆํ วิปาก\-ธมฺม\-ธมฺมํ ปฏิจฺจ นอนิทสฺสน\-อปฺปฏิโฆ นวิปาก\-ธมฺม\-ธมฺโม อุปฺปชฺชติ เหตุปจฺจยา.  เหตุยา ฉ.}

\prose{อนิทสฺสน\-อปฺปฏิฆํ เนว\-วิปาก\-น\-วิปากธมฺม\-ธมฺมํ ปฏิจฺจ นสนิทสฺสน\-สปฺปฏิโฆ นเน\-ว\-วิปาก\-น\-วิปาก\-ธมฺม\-ธมฺโม อุปฺปชฺชติ เหตุปจฺจยา.  เหตุยา ตีณิ.}

\csromanmarkh{22. สนิทสฺสน\-ตฺติก}\csromanmarkhii{4. อุปาทินฺนตฺติก}\csromanheadernomark{2}{22. สนิทสฺสน\-ตฺติก 4. อุปาทินฺนตฺติก}
\proseitem{47}{อนิทสฺสน\-อปฺปฏิฆํ อุปาทินฺนุ\-ปาทานิยํ ธมฺมํ ปฏิจฺจ นอนิทสฺสน\-อปฺปฏิโฆ นอุปาทินฺนุปาทานิโย ธมฺโม อุปฺปชฺชติ เหตุปจฺจยา.  เหตุยา ฉ.}

\prose{สนิทสฺสน\-สปฺปฏิโฆ อนุปาทินฺนุ\-ปาทานิโย ธมฺโม นสนิทสฺสน\-สปฺปฏิฆสฺส นอนุปาทินฺนุ\-ปาทานิยสฺส ธมฺมสฺส อารมฺมณ\-ปจฺจเยน ปจฺจโย.}

\setlength{\csromangathapagewidth}{0pt}
\setlength{\csromangathawakwidth}{0pt}
\csromangathameasuregroup{}{ธมฺมานุโลมปจฺจนีย ติกติกปฏฺฐานปาฬิ (สํขิตฺตํ.) .}
\csromangathameasurewak{ธมฺมานุโลมปจฺจนีย ติกติกปฏฺฐานปาฬิ (สํขิตฺตํ.) .}
\setlength{\csromangathaleftcol}{0pt}
\csromangathagroup{\csromangathastanza{\csromangathawak{\csromanfolio{304}ธมฺมา\-นุโลม\-ปจฺจนีย ติกติก\-ปฏฺฐาน\-ปาฬิ (สํขิตฺตํ.) .}}}

\prose{อารมฺมเณ นว, อนนฺตเร ตีณิ ฯเปฯ อุปนิสฺสเย ปุเรชาเต นว ฯเปฯ อวิคเต นว.}

\proseitem{47}{อนิทสฺสน\-อปฺปฏิฆํ อนุปาทินฺน\-อนุปาทานิยํ ธมฺมํ ปฏิจฺจ นอนิทสฺสน\-อปฺปฏิโฆ นอนุปาทินฺน\-อนุปาทานิโย ธมฺโม อุปฺปชฺชติ เหตุปจฺจยา.  เหตุยา ฉ.}

\csromanmarkh{22. สนิทสฺสน\-ตฺติก}\csromanmarkhii{5. สํกิลิฏฺฐ\-ตฺติก}\csromanheadernomark{2}{22. สนิทสฺสน\-ตฺติก 5. สํกิลิฏฺฐ\-ตฺติก}
\proseitem{48}{อนิทสฺสน\-อปฺปฏิฆํ สํกิลิฏฺฐ\-สํกิเลสิกํ ธมฺมํ ปฏิจฺจ นอนิทสฺสน\-อปฺปฏิโฆ นสํกิลิฏฺฐ\-สํกิเลสิโก ธมฺโม อุปฺปชฺชติ เหตุปจฺจยา.  เหตุยา ฉ.}

\prose{อนิทสฺสน\-อปฺปฏิฆํ อสํกิลิฏฺฐ\-สํกิเลสิกํ ธมฺมํ ปจฺจยา นสนิทสฺสน\-สปฺปฏิโฆ นอสํกิลิฏฺฐ\-สํกิเลสิโก ธมฺโม อุปฺปชฺชติ เหตุปจฺจยา.  เหตุยา ตีณิ.}

\prose{อนิทสฺสน\-อปฺปฏิฆํ อสํกิลิฏฺฐ\-อสํกิเลสิกํ ธมฺมํ ปฏิจฺจ นอนิทสฺสน\-อปฺปฏิโฆ นอสํกิลิฏฺฐ\-อสํกิเลสิโก ธมฺโม อุปฺปชฺชติ เหตุปจฺจยา.  เหตุยา ฉ.}

\csromanmarkh{22. สนิทสฺสน\-ตฺติก}\csromanmarkhii{6. วิตกฺกตฺติก}\csromanheadernomark{2}{22. สนิทสฺสน\-ตฺติก 6. วิตกฺกตฺติก}
\proseitem{49}{อนิทสฺสน\-อปฺปฏิฆํ สวิตกฺก\-สวิจารํ ธมฺมํ ปฏิจฺจ นอนิทสฺสน\-อปฺปฏิโฆ นสวิตกฺก\-สวิจาโร ธมฺโม อุปฺปชฺชติ เหตุปจฺจยา.  เหตุยา ฉ.}

\prose{อนิทสฺสน\-อปฺปฏิฆํ อวิตกฺก\-วิจารมตฺตํ ธมฺมํ ปฏิจฺจ นอนิทสฺสน\-อปฺปฏิโฆ นอวิตกฺก\-วิจารมตฺโต ธมฺโม อุปฺปชฺชติ เหตุปจฺจยา.  เหตุยา ฉ.}

\prose{อนิทสฺสน\-อปฺปฏิฆํ อวิตกฺก\-อวิจารํ ธมฺมํ ปฏิจฺจ นสนิทสฺสน\-สปฺปฏิโฆ นอวิตกฺก\-อวิจาโร ธมฺโม อุปฺปชฺชติ เหตุปจฺจยา.  เหตุยา ตีณิ.}

\csromanmarkh{22. สนิทสฺสน\-ตฺติก}\csromanmarkhii{21. อชฺฌตฺตา\-รมฺมณ\-ตฺติก}\csromanheadernomark{2}{22. สนิทสฺสน\-ตฺติก 21. อชฺฌตฺตา\-รมฺมณ\-ตฺติก}
\proseitem{50}{อนิทสฺสน\-อปฺปฏิฆํ อชฺฌตฺตา\-รมฺมณํ ธมฺมํ ปฏิจฺจ นอนิทสฺสน\-อปฺปฏิโฆ นอชฺฌตฺตา\-รมฺมโณ ธมฺโม อุปฺปชฺชติ เหตุปจฺจยา.  เหตุยา ฉ.}

\csromanmarkh{22. สนิทสฺสน\-ตฺติก}\csromanmarkhii{21. อชฺฌตฺตา\-รมฺมณ\-ตฺติก}\csromanheadernomark{2}{22. สนิทสฺสน\-ตฺติก 21. อชฺฌตฺตา\-รมฺมณ\-ตฺติก}
\prosewithcloser{\csromanfolio{305}อนิทสฺสน\-อปฺปฏิฆํ พหิทฺธา\-รมฺมณํ ธมฺมํ ปฏิจฺจ นอนิทสฺสน\-อปฺปฏิโฆ นพหิทฺธา\-รมฺมโณ ธมฺโม อุปฺปชฺชติ เหตุปจฺจยา.  อนิทสฺสน\-อปฺปฏิฆํ พหิทฺธา\-รมฺมณํ ธมฺมํ ปฏิจฺจ นสนิทสฺสน\-สปฺปฏิโฆ นพหิทฺธา\-รมฺมโณ ธมฺโม อุปฺปชฺชติ เหตุปจฺจยา.  อนิทสฺสน\-อปฺปฏิฆํ พหิทฺธา\-รมฺมณํ ธมฺมํ ปฏิจฺจ นอนิทสฺสน\-สปฺปฏิโฆ นพหิทฺธา\-รมฺมโณ ธมฺโม อุปฺปชฺชติ เหตุปจฺจยา ฯเปฯ.  อนิทสฺสน\-อปฺปฏิฆํ พหิทฺธา\-รมฺมณํ ธมฺมํ ปฏิจฺจ นสนิทสฺสน\-สปฺปฏิโฆ นพหิทฺธา\-รมฺมโณ จ นอนิทสฺสน\-สปฺปฏิโฆ นพหิทฺธา\-รมฺมโณ จ ธมฺมา อุปฺปชฺชนฺติ เหตุปจฺจยา.  เหตุยา ฉ ฯเปฯ อวิคเต ฉ. (ปญฺหาวารํ วิตฺถาเรตพฺพํ.)}{ธมฺมา\-นุโลม\-ปจฺจนีเย ติกติก\-ปฏฺฐานํ นิฏฺฐิตํ.}
\pitaka{อภิธมฺม\-ปิฏก}
\namakkaram{นโม ตสฺส ภควโต อรหโต สมฺมา\-สมฺพุทฺธสฺส.}
\csromanmarkh{1. เหตุทุก}\csromanmarkhii{1. สเหตุกทุก}\csromanheadernomark{2}{1. \textbf{เหตุทุก 1. สเหตุกทุก}}
\proseitem{1}{\csromanfolio{307}เหตุํ สเหตุกํ ธมฺมํ ปฏิจฺจ นเหตุ นสเหตุโก ธมฺโม อุปฺปชฺชติ เหตุปจฺจยา.  นเหตุํ สเหตุกํ ธมฺมํ ปฏิจฺจ นเหตุ นสเหตุโก ธมฺโม อุปฺปชฺชติ เหตุปจฺจยา.  เหตุํ สเหตุกญฺจ นเหตุํ สเหตุกญฺจ ธมฺมํ ปฏิจฺจ นเหตุ นสเหตุโก ธมฺโม อุปฺปชฺชติ เหตุปจฺจยา. (สํขิตฺตํ.) . เหตุยา ตีณิ, อารมฺมเณ เอกํ ฯเปฯ อวิคเต ปญฺจ.}

\proseitem{2}{นเหตุํ สเหตุกํ ธมฺมํ ปฏิจฺจ นนเหตุ นสเหตุโก ธมฺโม อุปฺปชฺชติ นเหตุปจฺจยา. (สํขิตฺตํ.) . นเหตุยา เอกํ, นอารมฺมเณ ตีณิ ฯเปฯ โนวิคเต ตีณิ.}

\prose{(สหชาตวารมฺปิ ฯเปฯ นมฺปยุตฺตวารมฺปิ ปฏิจฺจ\-วาร\-สทิสํ.)}

\prose{เหตุอารมฺมณ}

\prose{3. เหตุ สเหตุโก ธมฺโม นเหตุสฺส นสเหตุกสฺส ธมฺมสฺส}

\prose{เหตุ สเหตุโก ธมฺโม นเหตุสฺส นสเหตุกสฺส ธมฺมสฺส อารมฺมณ\-ปจฺจเยน ปจฺจโย.  เหตุ สเหตุโก ธมฺโม นนเหตุสฺส}

\vspace{\tipitakaparskip}
\csromancenter{ทุกทุก\-ปฏฺฐาน\-ปาฬิ นสเหตุกสฺส ธมฺมสฺส อารมฺมณ\-ปจฺจเยน ปจฺจโย.  เหตุ สเหตุโก ธมฺโม นเหตุสฺส นสเหตุกสฺส จ นนเหตุสฺส นสเหตุกสฺส จ ธมฺมสฺส อารมฺมณ\-ปจฺจเยน ปจฺจโย. (3)}
\prose{\csromanfolio{308}นเหตุ สเหตุโก ธมฺโม นนเหตุสฺส นสเหตุกสฺส ธมฺมสฺส อารมฺมณ\-ปจฺจเยน ปจฺจโย.  นเหตุ สเหตุโก ธมฺโม นเหตุสฺส นสเหตุกสฺส ธมฺมสฺส อารมฺมณ\-ปจฺจเยน ปจฺจโย.  นเหตุ สเหตุโก ธมฺโม นเหตุสฺส นสเหตุกสฺส จ นนเหตุสฺส นสเหตุกสฺส จ ธมฺมสฺส อารมฺมณ\-ปจฺจเยน ปจฺจโย. (3) (สํขิตฺตํ.)}

\prose{เหตุยา เอกํ, อารมฺมเณ ฉ, อธิปติยา เทฺว ฯเปฯ อวิคเต ปญฺจ. (ปญฺหาวารํ วิตฺถาเรตพฺพํ.)}

\proseitem{4}{เหตุํ อเหตุกํ ธมฺมํ ปฏิจฺจ นเหตุ นอเหตุโก ธมฺโม อุปฺปชฺชติ เหตุปจฺจยา. (1)}

\prose{นเหตุํ อเหตุกํ ธมฺมํ ปฏิจฺจ นนเหตุ นอเหตุโก ธมฺโม อุปฺปชฺชติ เหตุปจฺจยา. ตีณิ. (สํขิตฺตํ.) . เหตุยา จตฺตาริ.}

\csromanmarkh{1. เหตุทุก}\csromanmarkhii{2. เหตุ\-สมฺปยุตฺตทุก}\csromanheadernomark{2}{1. เหตุทุก 2. เหตุ\-สมฺปยุตฺตทุก}
\proseitem{5}{เหตุํ เหตุ\-สมฺปยุตฺตํ ธมฺมํ ปฏิจฺจ นเหตุ นเหตุ\-สมฺปยุตฺโต ธมฺโม อุปฺปชฺชติ เหตุปจฺจยา. ตีณิ. (สเหตุก\-สทิสํ.)}

\prose{เหตุํ เหตุ\-วิปฺปยุตฺตํ ธมฺมํ ปฏิจฺจ นเหตุ นเหตุ\-วิปฺปยุตฺโต}

\prose{นเหตุํ เหตุ\-วิปฺปยุตฺตํ ธมฺมํ ปฏิจฺจ นนเหตุ นเหตุ\-วิปฺปยุตฺโต ธมฺโม อุปฺปชฺชติ เหตุปจฺจยา. ตีณิ.  เหตุยา จตฺตาริ.}

\csromanheader{2}{1. เหตุทุก 3-5. เหตุ\-สเหตุก\-ทุกาทิ}
\proseitem{6}{เหตุํ เหตุญฺเจว สเหตุกญฺจ ธมฺมํ ปฏิจฺจ นเหตุ นเหตุ เจว นอเหตุโก จ ธมฺโม อุปฺปชฺชติ เหตุปจฺจยา.  เหตุยา เอกํ.}

\prose{นเหตุํ สเหตุกญฺเจว น จ เหตุํ ธมฺมํ ปฏิจฺจ นนเหตุ นอเหตุโก เจว นนเหตุ จ ธมฺโม อุปฺปชฺชติ เหตุปจฺจยา.  เหตุยา เอกํ.}

\csromanheader{2}{1. เหตุทุก 6-12. จูฬนฺตรทุก}
\proseitem{7}{\csromanfolio{309}เหตุํ เหตุญฺเจว เหตุ\-สมฺปยุตฺตญฺจ ธมฺมํ ปฏิจฺจ นเหตุ นเหตุ เจว นเหตุ\-วิปฺปยุตฺโต จ ธมฺโม อุปฺปชฺชติ เหตุปจฺจยา.  เหตุยา เอกํ.}

\prose{นเหตุํ เหตุ\-สมฺปยุตฺต\-ญฺเจว น จ เหตุํ ธมฺมํ ปฏิจฺจ นนเหตุ นเหตุ วิปฺปยุตฺโต เจว นนเหตุ จ ธมฺโม อุปฺปชฺชติ เหตุปจฺจยา.  เหตุยา เอกํ.}

\proseitem{8}{นเหตุํ นเหตุํ สเหตุกํ ธมฺมํ ปฏิจฺจ นเหตุ นเหตุ นสเหตุโก ธมฺโม อุปฺปชฺชติ เหตุปจฺจยา.  เหตุยา อธิปติยา เอกํ.}

\prose{นเหตุํ นเหตุํ อเหตุกํ ธมฺมํ ปฏิจฺจ นเหตุ นเหตุ นอเหตุโก ธมฺโม อุปฺปชฺชติ เหตุปจฺจยา.  เหตุยา เอกํ.}

\csromanheader{2}{1. เหตุทุก 6-12. จูฬนฺตรทุก}
\proseitem{9}{นเหตุ อปฺปจฺจโย ธมฺโม นนเหตุสฺส นอปฺปจฺจยสฺส ธมฺมสฺส อารมฺมณ\-ปจฺจเยน ปจฺจโย ฯเปฯ.  นเหตุ อปฺปจฺจโย ธมฺโม นเหตุสฺส นอปฺปจฺจยสฺส ธมฺมสฺส อารมฺมณ\-ปจฺจเยน ปจฺจโย.  นเหตุ อปฺปจฺจโย ธมฺโม นเหตุสฺส นอปฺปจฺจยสฺส จ นนเหตุสฺส นอปฺปจฺจยสฺส จ ธมฺมสฺส อารมฺมณ\-ปจฺจเยน ปจฺจโย ฯเปฯ. (อสงฺขตํ อปฺปจฺจย\-สทิสํ.)}

\proseitem{10}{นเหตุ สนิทสฺสโน ธมฺโม นนเหตุสฺส นสนิทสฺสนสฺส ธมฺมสฺส อารมฺมณ\-ปจฺจเยน ปจฺจโย.  นเหตุ สนิทสฺสโน ธมฺโม นเหตุสฺส นสนิทสฺสนสฺส ธมฺมสฺส อารมฺมณ\-ปจฺจเยน ปจฺจโย.  นเหตุ สนิทสฺสโน ธมฺโม นเหตุสฺส นสนิทสฺสนสฺส จ นนเหตุสฺส นสนิทสฺสนสฺส จ ธมฺมสฺส อารมฺมณ\-ปจฺจเยน ปจฺจโย.  อารมฺมเณ ตีณิ.}

\prose{เหตุํ อนิทสฺสนํ ธมฺมํ ปฏิจฺจ นเหตุ นอนิทสฺสโน ธมฺโม อุปฺปชฺชติ เหตุปจฺจยา.  เหตุยา ตีณิ.}

\proseitem{11}{นเหตุํ สปฺปฏิฆํ ธมฺมํ ปฏิจฺจ นเหตุ นสปฺปฏิโฆ ธมฺโม อุปฺปชฺชติ เหตุปจฺจยา.  เหตุยา เอกํ.}

\prose{เหตุํ อปฺปฏิฆํ ธมฺมํ ปฏิจฺจ นเหตุ นอปฺปฏิโฆ ธมฺโม อุปฺปชฺชติ เหตุปจฺจยา.  เหตุยา ตีณิ.}

\csromanheader{2}{ธมฺมา\-นุโลม\-ปจฺจนีย ทุกทุก\-ปฏฺฐาน\-ปาฬิ}
\proseitem{12}{\csromanfolio{310}นเหตุํ รูปิํ ธมฺมํ ปฏิจฺจ นนเหตุ นรูปี ธมฺโม อุปฺปชฺชติ เหตุปจฺจยา.  เหตุยา ตีณิ.}

\prose{เหตุํ อรูปิํ ธมฺมํ ปฏิจฺจ นเหตุ นอรูปี ธมฺโม อุปฺปชฺชติ เหตุปจฺจยา.  เหตุยา ตีณิ.}

\proseitem{13}{นเหตุํ โลกิยํ ธมฺมํ ปจฺจยา นนเหตุ นโลกิโย ธมฺโม อุปฺปชฺชติ เหตุปจฺจยา. เหตุยา ตีณิ.}

\prose{เหตุํ โลกุตฺตรํ ธมฺมํ ปฏิจฺจ นเหตุ นโลกุตฺตโร ธมฺโม อุปฺปชฺชติ เหตุปจฺจยา.  นเหตุํ โลกุตฺตรํ ธมฺมํ ปฏิจฺจ นเหตุ นโลกุตฺตโร ธมฺโม อุปฺปชฺชติ เหตุปจฺจยา.  เหตุํ โลกุตฺตรญฺจ นเหตุํ โลกุตฺตรญฺจ ธมฺมํ ปฏิจฺจ นเหตุ นโลกุตฺตโร ธมฺโม อุปฺปชฺชติ เหตุปจฺจยา.  เหตุยา ตีณิ.}

\proseitem{14}{เหตุํ เกนจิ วิญฺเญยฺยํ ธมฺมํ ปฏิจฺจ นเหตุ นเกนจิ วิญฺเญยฺโย ธมฺโม อุปฺปชฺชติ เหตุปจฺจยา.  เหตุยา นว.}

\prose{เหตุํ นเกนจิ วิญฺเญยฺยํ ธมฺมํ ปฏิจฺจ นเหตุ นนเกนจิ วิญฺเญยฺโย ธมฺโม อุปฺปชฺชติ เหตุปจฺจยา.  เหตุยา นว.}

\csromanmarkh{1. เหตุทุก}\csromanmarkhii{13. อาสวโคจฺฉก}\csromanheadernomark{2}{1. เหตุทุก 13. อาสวโคจฺฉก}
\proseitem{15}{เหตุํ อาสวํ ธมฺมํ ปฏิจฺจ นเหตุ นอาสโว ธมฺโม อุปฺปชฺชติ เหตุปจฺจยา.  เหตุยา ปญฺจ.}

\prose{เหตุํ โนอาสวํ ธมฺมํ ปฏิจฺจ นนเหตุ นโนอาสโว ธมฺโม อุปฺปชฺชติ เหตุปจฺจยา.  เหตุยา ปญฺจ.}

\proseitem{16}{นเหตุํ สาสวํ ธมฺมํ ปจฺจยา นนเหตุ นสาสโว ธมฺโม อุปฺปชฺชติ เหตุปจฺจยา.  เหตุยา ตีณิ.}

\prose{เหตุํ อนาสวํ ธมฺมํ ปฏิจฺจ นเหตุ นอนาสโว ธมฺโม อุปฺปชฺชติ เหตุปจฺจยา.  เหตุยา ตีณิ.}

\proseitem{17}{เหตุํ อาสว\-สมฺปยุตฺตํ ธมฺมํ ปฏิจฺจ นเหตุ นอาสว\-สมฺปยุตฺโต ธมฺโม อุปฺปชฺชติ เหตุปจฺจยา.  เหตุยา นว.}

\csromanmarkh{1. เหตุทุก}\csromanmarkhii{19. สญฺโญชนาทิโคจฺฉก}\csromanheadernomark{2}{1. เหตุทุก 19. สญฺโญชนาทิโคจฺฉก}
\prose{\csromanfolio{311}เหตุํ อาสว\-วิปฺปยุตฺตํ ธมฺมํ ปฏิจฺจ นเหตุ นอาสว\-วิปฺปยุตฺโต ธมฺโม อุปฺปชฺชติ เหตุปจฺจยา.  เหตุยา ตีณิ.}

\proseitem{18}{เหตุํ อาสวญฺเจว สาสวญฺจ ธมฺมํ ปฏิจฺจ นเหตุ นอาสโว เจว นอนาสโว จ ธมฺโม อุปฺปชฺชติ เหตุปจฺจยา.  เหตุยา ปญฺจ.}

\prose{เหตุํ สาสวญฺเจว โน จ อาสวํ ธมฺมํ ปฏิจฺจ นนเหตุ นอนาสโว เจว นโน จ อาสโว ธมฺโม อุปฺปชฺชติ เหตุปจฺจยา.  เหตุยา ปญฺจ.}

\proseitem{19}{เหตุํ อาสวญฺเจว อาสว\-สมฺปยุตฺตญฺจ ธมฺมํ ปฏิจฺจ นเหตุ นอาสโว เจว นอาสว\-วิปฺปยุตฺโต จ ธมฺโม อุปฺปชฺชติ เหตุปจฺจยา.  เหตุยา ตีณิ.}

\prose{นเหตุํ อาสว\-สมฺปยุตฺต\-ญฺเจว โน จ อาสวํ ธมฺมํ ปฏิจฺจ นนเหตุ นอาสว\-วิปฺปยุตฺโต เจว นโนอาสโว จ ธมฺโม อุปฺปชฺชติ เหตุปจฺจยา.  เหตุยา ตีณิ.}

\proseitem{20}{นเหตุํ อาสว\-วิปฺปยุตฺตํ สาสวํ ธมฺมํ ปจฺจยา นนเหตุ อาสว\-วิปฺปยุตฺโต นสาสโว ธมฺโม อุปฺปชฺชติ เหตุปจฺจยา. (โลกิยสทิสํ.)}

\prose{เหตุํ อาสว\-วิปฺปยุตฺตํ อนาสวํ ธมฺมํ ปฏิจฺจ นเหตุ อาสว\-วิปฺปยุตฺโต นอนาสโว ธมฺโม อุปฺปชฺชติ เหตุปจฺจยา.  นเหตุํ อาสว\-วิปฺปยุตฺตํ อนาสวํ ธมฺมํ ปฏิจฺจ นเหตุ อาสว\-วิปฺปยุตฺโต นอนาสโว ธมฺโม อุปฺปชฺชติ เหตุปจฺจยา.  เหตุํ อาสว\-วิปฺปยุตฺตํ อนาสวญฺจ นเหตุํ อาสว\-วิปฺปยุตฺตํ อนาสวญฺจ ธมฺมํ ปฏิจฺจ นเหตุ อาสว\-วิปฺปยุตฺโต นอนาสโว ธมฺโม อุปฺปชฺชติ เหตุปจฺจยา.  เหตุยา ตีณิ.}

\csromanmarkh{1. เหตุทุก}\csromanmarkhii{19. สญฺโญชนาทิโคจฺฉก}\csromanheadernomark{2}{1. เหตุทุก 19. สญฺโญชนาทิโคจฺฉก}
\proseitem{21}{เหตุํ สญฺโญชนํ ธมฺมํ ปฏิจฺจ นเหตุ โนสญฺโญชโน ธมฺโม อุปฺปชฺชติ เหตุปจฺจยา.  เหตุยา ตีณิ.}

\proseitem{22}{เหตุํ คนฺถํ ธมฺมํ ปฏิจฺจ นเหตุ โนคนฺโถ ธมฺโม อุปฺปชฺชติ เหตุปจฺจยา.  เหตุยา นว.}

\csromanheader{2}{ธมฺมา\-นุโลม\-ปจฺจนีย ทุกทุก\-ปฏฺฐาน\-ปาฬิ}
\proseitem{23}{\csromanfolio{312}เหตุํ โอฆํ ธมฺมํ ปฏิจฺจ นเหตุ โนโอโฆ ธมฺโม อุปฺปชฺชติ เหตุปจฺจยา.  เหตุยา ปญฺจ.}

\proseitem{24}{เหตุํ โยคํ ธมฺมํ ปฏิจฺจ นเหตุ โนโยโค ธมฺโม อุปฺปชฺชติ เหตุปจฺจยา.  เหตุยา ปญฺจ.}

\proseitem{25}{เหตุํ นีวรณํ ธมฺมํ ปฏิจฺจ นเหตุ โนนีวรโณ ธมฺโม อุปฺปชฺชติ เหตุปจฺจยา.  เหตุยา ตีณิ.}

\proseitem{26}{นเหตุํ ปรามาสํ ธมฺมํ ปฏิจฺจ นนเหตุ โนปรามาโส ธมฺโม อุปฺปชฺชติ เหตุปจฺจยา.  เหตุยา ตีณิ.}

\csromanmarkh{1. เหตุทุก}\csromanmarkhii{54. มหนฺตรทุก}\csromanheadernomark{2}{1. เหตุทุก 54. มหนฺตรทุก}
\proseitem{27}{เหตุํ สารมฺมณํ ธมฺมํ ปฏิจฺจ นเหตุ นสารมฺมโณ ธมฺโม อุปฺปชฺชติ เหตุปจฺจยา.  เหตุยา ตีณิ.}

\prose{นเหตุํ อนารมฺมณํ ธมฺมํ ปฏิจฺจ นนเหตุ นอารมฺมโณ ธมฺโม อุปฺปชฺชติ เหตุปจฺจยา.  เหตุยา ตีณิ.}

\proseitem{28}{นเหตุํ จิตฺตํ ธมฺมํ ปฏิจฺจ นนเหตุ โนจิตฺโต ธมฺโม อุปฺปชฺชติ เหตุปจฺจยา.  เหตุยา ตีณิ.}

\prose{เหตุํ โนจิตฺตํ ธมฺมํ ปฏิจฺจ นเหตุ นโนจิตฺโต ธมฺโม อุปฺปชฺชติ เหตุปจฺจยา.  เหตุยา ตีณิ.}

\proseitem{29}{เหตุํ เจตสิกํ ธมฺมํ ปฏิจฺจ นเหตุ นเจตสิโก ธมฺโม อุปฺปชฺชติ เหตุปจฺจยา.  เหตุยา ตีณิ.}

\prose{นเหตุํ อเจตสิกํ ธมฺมํ ปฏิจฺจ นนเหตุ นอเจตสิโก ธมฺโม อุปฺปชฺชติ เหตุปจฺจยา.  เหตุยา ตีณิ.}

\proseitem{30}{เหตุํ จิตฺต\-สมฺปยุตฺตํ ธมฺมํ ปฏิจฺจ นเหตุ นจิตฺต\-สมฺปยุตฺโต ธมฺโม อุปฺปชฺชติ เหตุปจฺจยา.  เหตุยา ตีณิ.}

\prose{นเหตุํ จิตฺต\-วิปฺปยุตฺตํ ธมฺมํ ปฏิจฺจ นนเหตุ นจิตฺต\-วิปฺปยุตฺโต ธมฺโม อุปฺปชฺชติ เหตุปจฺจยา.  เหตุยา ตีณิ.}

\csromanmarkh{1. เหตุทุก}\csromanmarkhii{54. มหนฺตรทุก}\csromanheadernomark{2}{1. เหตุทุก 54. มหนฺตรทุก}
\proseitem{31}{\csromanfolio{313}เหตุํ จิตฺต\-สํสฏฺฐํ ธมฺมํ ปฏิจฺจ นเหตุ นจิตฺต\-สํสฏฺโฐ ธมฺโม อุปฺปชฺชติ เหตุปจฺจยา.  เหตุยา ตีณิ.}

\prose{นเหตุํ จิตฺต\-วิสํสฏฺฐํ ธมฺมํ ปฏิจฺจ นนเหตุ นจิตฺต\-วิสํสฏฺโฐ ธมฺโม อุปฺปชฺชติ เหตุปจฺจยา.  เหตุยา ตีณิ.}

\proseitem{32}{เหตุํ จิตฺต\-สมุฏฺฐานํ ธมฺมํ ปฏิจฺจ นเหตุ โนจิตฺต\-สมุฏฺฐาโน ธมฺโม อุปฺปชฺชติ เหตุปจฺจยา.  เหตุยา ตีณิ.}

\prose{นเหตุํ โนจิตฺต\-สมุฏฺฐานํ ธมฺมํ ปฏิจฺจ นนเหตุ นโน\-จิตฺตสมุฏฺฐาโน ธมฺโม อุปฺปชฺชติ เหตุปจฺจยา.  เหตุยา ตีณิ.}

\proseitem{33}{เหตุํ จิตฺตสหภุํ ธมฺมํ ปฏิจฺจ นเหตุ โนจิตฺตสหภู ธมฺโม อุปฺปชฺชติ เหตุปจฺจยา.  เหตุยา ตีณิ.}

\prose{นเหตุํ โนจิตฺต\-สหภุํ ธมฺมํ ปฏิจฺจ นนเหตุ นโนจิตฺตสหภู ธมฺโม อุปฺปชฺชติ เหตุปจฺจยา.  เหตุยา ตีณิ.}

\proseitem{34}{เหตุํ จิตฺตา\-นุปริวตฺติํ ธมฺมํ ปฏิจฺจ นเหตุ โนจิตฺตา\-นุปริวตฺตี ธมฺโม อุปฺปชฺชติ เหตุปจฺจยา.  เหตุยา ตีณิ.}

\prose{นเหตุํ โนจิตฺตา\-นุปริวตฺติํ ธมฺมํ ปฏิจฺจ นนเหตุ นโนจิตฺตานุปริวตฺตี ธมฺโม อุปฺปชฺชติ เหตุปจฺจยา. เหตุยา ตีณิ.}

\proseitem{35}{เหตุํ จิตฺต\-สํสฏฺฐ\-สมุฏฺฐานํ ธมฺมํ ปฏิจฺจ นเหตุ โนจิตฺต\-สํสฏฺฐ\-สมุฏฺฐาโน ธมฺโม อุปฺปชฺชติ เหตุปจฺจยา.  เหตุยา ตีณิ.}

\prose{นเหตุํ โนจิตฺต\-สํสฏฺฐ\-สมุฏฺฐานํ ธมฺมํ ปฏิจฺจ นนเหตุ นโน\-จิตฺต\-สํสฏฺฐ\-สมุฏฺฐาโน ธมฺโม อุปฺปชฺชติ เหตุปจฺจยา.  เหตุยา ตีณิ.}

\proseitem{36}{เหตุํ จิตฺต\-สํสฏฺฐ\-สมุฏฺฐาน\-สหภุํ ธมฺมํ ปฏิจฺจ นเหตุ โนจิตฺต\-สํสฏฺฐ\-สมุฏฺฐาน\-สหภู ธมฺโม อุปฺปชฺชติ เหตุปจฺจยา.  เหตุยา ตีณิ.}

\prose{นเหตุํ โนจิตฺต\-สํสฏฺฐ\-สมุฏฺฐาน\-สหภุํ ธมฺมํ ปฏิจฺจ นนเหตุ นโนจิตฺตสํสฏฺฐสมุฏฺฐานสหภู ธมฺโม อุปฺปชฺชติ เหตุปจฺจยา.  เหตุยา ตีณิ.}

\proseitem{37}{เหตุํ จิตฺต\-สํสฏฺฐ\-สมุฏฺฐานา\-นุปริวตฺติํ ธมฺมํ ปฏิจฺจ นเหตุ โนจิตฺต\-สํสฏฺฐ\-สมุฏฺฐานา\-นุปริวตฺตี ธมฺโม อุปฺปชฺชติ เหตุปจฺจยา.  เหตุยา ตีณิ.}

\csromanheader{2}{ธมฺมา\-นุโลม\-ปจฺจนีย ทุกทุก\-ปฏฺฐาน\-ปาฬิ}
\prose{\csromanfolio{314}นเหตุํ โนจิตฺต\-สํสฏฺฐ\-สมุฏฺฐานา\-นุปริวตฺติํ ธมฺมํ ปฏิจฺจ นนเหตุ นโนจิตฺตสํสฏฺฐสมุฏฺฐานานุปริวตฺตี ธมฺโม อุปฺปชฺชติ เหตุปจฺจยา.}

\prose{เหตุยา ตีณิ. (สํขิตฺตํ. นยวเสน วิตฺถาเรตพฺพํ.)}

\csromanmarkh{1. เหตุทุก}\csromanmarkhii{82. ทสฺสเนน\-ปหาตพฺพ\-ทุกาทิ}\csromanheadernomark{2}{1. เหตุทุก 82. ทสฺสเนน\-ปหาตพฺพ\-ทุกาทิ}
\proseitem{38}{เหตุํ ทสฺสเนน ปหาตพฺพํ ธมฺมํ ปฏิจฺจ นเหตุ นทสฺสเนน ปหาตพฺโพ ธมฺโม อุปฺปชฺชติ เหตุปจฺจยา.  เหตุยา ตีณิ.}

\prose{นเหตุํ นทสฺสเนน ปหาตพฺพํ ธมฺมํ ปจฺจยา นนเหตุ นนทสฺสเนน ปหาตพฺโพ ธมฺโม อุปฺปชฺชติ เหตุปจฺจยา. ตีณิเมว.}

\proseitem{39}{เหตุํ ทสฺสเนน ปหาตพฺพ\-เหตุกํ ธมฺมํ ปฏิจฺจ นเหตุ นทสฺสเนน ปหาตพฺพ\-เหตุโก ธมฺโม อุปฺปชฺชติ เหตุปจฺจยา.  เหตุยา ตีณิ.}

\prose{นเหตุํ นทสฺสเนน ปหาตพฺพ\-เหตุกํ ธมฺมํ ปฏิจฺจ นนเหตุ นนทสฺสเนน ปหาตพฺพ\-เหตุโก ธมฺโม อุปฺปชฺชติ เหตุปจฺจยา.  เหตุยา เอกํ.}

\csromanmarkh{1. เหตุทุก}\csromanmarkhii{99. สรณทุก}\csromanheadernomark{2}{1. เหตุทุก 99. สรณทุก}
\proseitem{40}{เหตุํ สรณํ ธมฺมํ ปฏิจฺจ นเหตุ นสรโณ ธมฺโม อุปฺปชฺชติ เหตุปจฺจยา.  นเหตุํ สรณํ ธมฺมํ ปฏิจฺจ นเหตุ นสรโณ ธมฺโม อุปฺปชฺชติ เหตุปจฺจยา.  เหตุํ สรณญฺจ นเหตุํ สรณญฺจ ธมฺมํ ปฏิจฺจ นเหตุ นสรโณ ธมฺโม อุปฺปชฺชติ เหตุปจฺจยา.  เหตุยา ตีณิ.}

\prose{นเหตุํ อรณํ ธมฺมํ ปจฺจยา นนเหตุ นอรโณ ธมฺโม อุปฺปชฺชติ เหตุปจฺจยา.  นเหตุํ อรณํ ธมฺมํ ปจฺจยา นเหตุ นอรโณ ธมฺโม อุปฺปชฺชติ เหตุปจฺจยา.  นเหตุํ อรณํ ธมฺมํ ปจฺจยา นเหตุ นอรโณ จ นนเหตุ นอรโณ จ ธมฺมา อุปฺปชฺชนฺติ เหตุปจฺจยา.  เหตุยา ตีณิ.}

\csromanmarkh{2. สเหตุกทุก}\csromanmarkhii{1. เหตุทุก}\csromanheadernomark{2}{2. \textbf{สเหตุกทุก 1. เหตุทุก}}
\proseitem{41}{สเหตุกํ เหตุํ ธมฺมํ ปฏิจฺจ นสเหตุโก นเหตุ ธมฺโม อุปฺปชฺชติ เหตุปจฺจยา. ตีณิ.}

\csromanmarkh{7-13. จูฬนฺตรทุก}\csromanmarkhii{1. เหตุทุก}\csromanheadernomark{2}{7-13. จูฬนฺตรทุก 1. เหตุทุก}
\prose{\csromanfolio{315}อเหตุกํ เหตุํ ธมฺมํ ปฏิจฺจ นอเหตุโก นเหตุ ธมฺโม อุปฺปชฺชติ เหตุปจฺจยา. ตีณิ.  เหตุยา ฉ.}

\prose{สเหตุกํ นเหตุํ ธมฺมํ ปฏิจฺจ นอเหตุโก นนเหตุ ธมฺโม อุปฺปชฺชติ เหตุปจฺจยา.  เหตุยา ตีณิ.}

\csromanmarkh{3. เหตุ\-สมฺปยุตฺตทุก}\csromanmarkhii{1. เหตุทุก}\csromanheadernomark{2}{3. เหตุ\-สมฺปยุตฺตทุก 1. เหตุทุก}
\proseitem{42}{เหตุ\-สมฺปยุตฺตํ เหตุํ ธมฺมํ ปฏิจฺจ นเหตุ\-สมฺปยุตฺโต นเหตุ ธมฺโม อุปฺปชฺชติ เหตุปจฺจยา. ตีณิ.}

\prose{เหตุ\-วิปฺปยุตฺตํ เหตุํ ธมฺมํ ปฏิจฺจ นเหตุ\-วิปฺปยุตฺโต นเหตุ ธมฺโม อุปฺปชฺชติ เหตุปจฺจยา. ตีณิ.  เหตุยา ฉ.}

\prose{เหตุ\-สมฺปยุตฺตํ นเหตุํ ธมฺมํ ปฏิจฺจ นเหตุ\-วิปฺปยุตฺโต นนเหตุ ธมฺโม อุปฺปชฺชติ เหตุปจฺจยา.  เหตุยา ตีณิ.}

\csromanmarkh{4. เหตุ\-สเหตุกทุก}\csromanmarkhii{1. เหตุทุก}\csromanheadernomark{2}{4. เหตุ\-สเหตุกทุก 1. เหตุทุก}
\proseitem{43}{เหตุญฺเจว สเหตุกญฺจ เหตุํ ธมฺมํ ปฏิจฺจ นเหตุ เจว นอเหตุโก จ นเหตุ ธมฺโม อุปฺปชฺชติ เหตุปจฺจยา.  เหตุยา เอกํ.}

\prose{สเหตุกญฺเจว น จ เหตุํ นเหตุํ ธมฺมํ ปฏิจฺจ นอเหตุโก เจว นน จ เหตุ นนเหตุ ธมฺโม อุปฺปชฺชติ เหตุปจฺจยา.  เหตุยา เอกํ.}

\csromanmarkh{5. เหตุ\-เหตุ\-สมฺปยุตฺตทุก}\csromanmarkhii{1. เหตุทุก}\csromanheadernomark{2}{5. เหตุ\-เหตุ\-สมฺปยุตฺตทุก 1. เหตุทุก}
\proseitem{44}{เหตุญฺเจว เหตุ\-สมฺปยุตฺตญฺจ เหตุํ ธมฺมํ ปฏิจฺจ นเหตุ เจว นเหตุ\-วิปฺปยุตฺโต จ นเหตุ ธมฺโม อุปฺปชฺชติ เหตุปจฺจยา.  เหตุยา เอกํ.}

\prose{เหตุ\-สมฺปยุตฺต\-ญฺเจว น จ เหตุํ นเหตุํ ธมฺมํ ปฏิจฺจ นเหตุ\-วิปฺปยุตฺโต เจว นน จ เหตุ นนเหตุ ธมฺโม อุปฺปชฺชติ เหตุปจฺจยา.  เหตุยา เอกํ.}

\csromanmarkh{7-13. จูฬนฺตรทุก}\csromanmarkhii{1. เหตุทุก}\csromanheadernomark{2}{7-13. จูฬนฺตรทุก 1. เหตุทุก}
\proseitem{45}{สปฺปจฺจยํ เหตุํ ธมฺมํ ปฏิจฺจ นอปฺปจฺจโย นเหตุ ธมฺโม อุปฺปชฺชติ เหตุปจฺจยา.  เหตุยา เอกํ.}

\csromanheader{2}{ธมฺมา\-นุโลม\-ปจฺจนีย ทุกทุก\-ปฏฺฐาน\-ปาฬิ}
\prose{\csromanfolio{316}สปฺปจฺจยํ นเหตุํ ธมฺมํ ปฏิจฺจ นอปฺปจฺจโย นนเหตุ ธมฺโม อุปฺปชฺชติ เหตุปจฺจยา.  เหตุยา เอกํ. (สงฺขตํ สปฺปจฺจย\-สทิสํ.)}

\proseitem{46}{อนิทสฺสนํ เหตุํ ธมฺมํ ปฏิจฺจ นอนิทสฺสโน นเหตุ ธมฺโม อุปฺปชฺชติ เหตุปจฺจยา.  เหตุยา ตีณิ.}

\prose{อนิทสฺสนํ นเหตุํ ธมฺมํ ปฏิจฺจ นสนิทสฺสโน นนเหตุ ธมฺโม อุปฺปชฺชติ เหตุปจฺจยา.  เหตุยา เอกํ.}

\proseitem{47}{อปฺปฏิฆํ เหตุํ ธมฺมํ ปฏิจฺจ นอปฺปฏิโฆ นเหตุ ธมฺโม อุปฺปชฺชติ เหตุปจฺจยา.  เหตุยา ตีณิ.}

\prose{อปฺปฏิฆํ นเหตุํ ธมฺมํ ปฏิจฺจ นสปฺปฏิโฆ นนเหตุ ธมฺโม อุปฺปชฺชติ เหตุปจฺจยา.  เหตุยา เอกํ.}

\proseitem{48}{อรูปิํ เหตุํ ธมฺมํ ปฏิจฺจ นอรูปี นเหตุ ธมฺโม อุปฺปชฺชติ เหตุปจฺจยา.  เหตุยา ตีณิ.}

\prose{รูปิํ นเหตุํ ธมฺมํ ปฏิจฺจ นรูปี นนเหตุ ธมฺโม อุปฺปชฺชติ เหตุปจฺจยา.}

\prose{อรูปิํ นเหตุํ ธมฺมํ ปฏิจฺจ นอรูปี นนเหตุ ธมฺโม อุปฺปชฺชติ เหตุปจฺจยา.  เหตุยา ตีณิ.}

\proseitem{49}{โลกิยํ เหตุํ ธมฺมํ ปฏิจฺจ นโลกุตฺตโร นเหตุ ธมฺโม อุปฺปชฺชติ เหตุปจฺจยา. เอกํ.}

\prose{โลกุตฺตรํ เหตุํ ธมฺมํ ปฏิจฺจ นโลกุตฺตโร นเหตุ ธมฺโม อุปฺปชฺชติ เหตุปจฺจยา. ตีณิ.  เหตุยา จตฺตาริ.}

\prose{โลกิยํ นเหตุํ ธมฺมํ ปฏิจฺจ นโลกุตฺตโร นนเหตุ ธมฺโม อุปฺปชฺชติ เหตุปจฺจยา. เอกํ.}

\prose{โลกุตฺตรํ นเหตุํ ธมฺมํ ปฏิจฺจ นโลกิโย นนเหตุ ธมฺโม อุปฺปชฺชติ เหตุปจฺจยา. เอกํ.  เหตุยา เทฺว.}

\proseitem{50}{เกนจิ วิญฺเญยฺยํ เหตุํ ธมฺมํ ปฏิจฺจ นเกนจิ วิญฺเญยฺโย นเหตุ ธมฺโม อุปฺปชฺชติ เหตุปจฺจยา. ตีณิ.}

\prose{นเกนจิ วิญฺเญยฺยํ เหตุํ ธมฺมํ ปฏิจฺจ นนเกนจิ วิญฺเญยฺโย นเหตุ ธมฺโม อุปฺปชฺชติ เหตุปจฺจยา. ตีณิ.}

\csromanmarkh{16. อาสว\-สมฺปยุตฺตทุก}\csromanmarkhii{1. เหตุทุก}\csromanheadernomark{2}{16. อาสว\-สมฺปยุตฺตทุก 1. เหตุทุก}
\prose{\csromanfolio{317}เกนจิ วิญฺเญยฺยํ เหตุญฺจ นเกนจิ วิญฺเญยฺยํ เหตุญฺจ ธมฺมํ ปฏิจฺจ นเกนจิ วิญฺเญยฺโย นเหตุ ธมฺโม อุปฺปชฺชติ เหตุปจฺจยา.  เหตุยา นว.}

\prose{เกนจิ วิญฺเญยฺยํ นเหตุํ ธมฺมํ ปฏิจฺจ นเกนจิ วิญฺเญยฺโย นนเหตุ ธมฺโม อุปฺปชฺชติ เหตุปจฺจยา.  เหตุยา นว.}

\csromanmarkh{14. อาสวทุก}\csromanmarkhii{1. เหตุทุก}\csromanheadernomark{2}{14. \textbf{อาสวทุก 1. เหตุทุก}}
\proseitem{51}{อาสวํ เหตุํ ธมฺมํ ปฏิจฺจ โนอาสโว นเหตุ ธมฺโม อุปฺปชฺชติ เหตุปจฺจยา. ตีณิ.}

\prose{โนอาสวํ เหตุํ ธมฺมํ ปฏิจฺจ นโนอาสโว นเหตุ ธมฺโม อุปฺปชฺชติ เหตุปจฺจยา. เอกํ.}

\prose{อาสวํ เหตุญฺจ โนอาสวํ เหตุญฺจ ธมฺมํ ปฏิจฺจ โนอาสโว นเหตุ ธมฺโม อุปฺปชฺชติ เหตุปจฺจยา. เอกํ. (สพฺพตฺถ ปญฺจ.)}

\prose{โนอาสวํ นเหตุํ ธมฺมํ ปฏิจฺจ นโนอาสโว นนเหตุ ธมฺโม อุปฺปชฺชติ เหตุปจฺจยา.  เหตุยา ปญฺจ.}

\csromanmarkh{15. สาสวทุก}\csromanmarkhii{1. เหตุทุก}\csromanheadernomark{2}{15. \textbf{สาสวทุก 1. เหตุทุก}}
\proseitem{52}{สาสวํ เหตุํ ธมฺมํ ปฏิจฺจ นอนาสโว นเหตุ ธมฺโม อุปฺปชฺชติ เหตุปจฺจยา. เอกํ.}

\prose{อนาสวํ เหตุํ ธมฺมํ ปฏิจฺจ นอนาสโว นเหตุ ธมฺโม อุปฺปชฺชติ เหตุปจฺจยา. ตีณิ.  เหตุยา จตฺตาริ. (โลกิย\-ทุก\-สทิสํ.)}

\csromanmarkh{16. อาสว\-สมฺปยุตฺตทุก}\csromanmarkhii{1. เหตุทุก}\csromanheadernomark{2}{16. อาสว\-สมฺปยุตฺตทุก 1. เหตุทุก}
\proseitem{53}{อาสว\-สมฺปยุตฺตํ เหตุํ ธมฺมํ ปฏิจฺจ นอาสว\-สมฺปยุตฺโต นเหตุ ธมฺโม อุปฺปชฺชติ เหตุปจฺจยา. ตีณิ.}

\prose{อาสว\-วิปฺปยุตฺตํ เหตุํ ธมฺมํ ปฏิจฺจ นอาสว\-วิปฺปยุตฺโต นเหตุ ธมฺโม อุปฺปชฺชติ เหตุปจฺจยา. ตีณิ.}

\prose{อาสว\-สมฺปยุตฺตํ เหตุญฺจ อาสว\-วิปฺปยุตฺตํ เหตุญฺจ ธมฺมํ ปฏิจฺจ นอาสว\-สมฺปยุตฺโต นเหตุ ธมฺโม อุปฺปชฺชติ เหตุปจฺจยา. ตีณิ.  เหตุยา นว.}

\prose{อาสว\-สมฺปยุตฺตํ นเหตุํ ธมฺมํ ปฏิจฺจ นอาสว\-สมฺปยุตฺโต นนเหตุ ธมฺโม อุปฺปชฺชติ เหตุปจฺจยา. ตีณิ.}

\csromanheader{2}{ธมฺมา\-นุโลม\-ปจฺจนีย ทุกทุก\-ปฏฺฐาน\-ปาฬิ}
\prose{\csromanfolio{318}อาสว\-วิปฺปยุตฺตํ นเหตุํ ธมฺมํ ปฏิจฺจ นอาสว\-สมฺปยุตฺโต นนเหตุ ธมฺโม อุปฺปชฺชติ เหตุปจฺจยา. เอกํ.  เหตุยา จตฺตาริ.}

\csromanmarkh{17. อาสว\-สาสวทุก}\csromanmarkhii{1. เหตุทุก}\csromanheadernomark{2}{17. อาสว\-สาสวทุก 1. เหตุทุก}
\proseitem{54}{อาสวญฺเจว สาสวญฺจ เหตุํ ธมฺมํ ปฏิจฺจ นอาสโว เจว นอนาสโว จ นเหตุ ธมฺโม อุปฺปชฺชติ เหตุปจฺจยา. ตีณิ.}

\prose{สาสวญฺเจว โน จ อาสวํ เหตุํ ธมฺมํ ปฏิจฺจ นอาสโว เจว นอนาสโว จ นเหตุ ธมฺโม อุปฺปชฺชติ เหตุปจฺจยา. เอกํ.}

\prose{อาสวญฺเจว สาสวํ เหตุญฺจ สาสวญฺเจว โน จ อาสวํ เหตุญฺจ ธมฺมํ ปฏิจฺจ นอาสโว เจว นอนาสโว จ นเหตุ ธมฺโม อุปฺปชฺชติ เหตุปจฺจยา. เอกํ.  เหตุยา ปญฺจ.}

\prose{อาสวญฺเจว สาสวญฺจ นเหตุํ ธมฺมํ ปฏิจฺจ นอนาสโว เจว นโนอาสโว จ นนเหตุ ธมฺโม อุปฺปชฺชติ เหตุปจฺจยา.  เหตุยา ปญฺจ.  }

\prose{18. อาสว\-อาสว\-สมฺปยุตฺตทุก 1. เหตุทุก}

\proseitem{55}{อาสวญฺเจว อาสว\-สมฺปยุตฺตญฺจ เหตุํ ธมฺมํ ปฏิจฺจ นอาสโว เจว นอาสว\-วิปฺปยุตฺโต จ นเหตุ ธมฺโม อุปฺปชฺชติ เหตุปจฺจยา. ตีณิ.}

\prose{อาสว\-สมฺปยุตฺต\-ญฺเจว โน จ อาสวํ เหตุํ ธมฺมํ ปฏิจฺจ นอาสโว เจว นอาสว\-วิปฺปยุตฺโต จ นเหตุ ธมฺโม อุปฺปชฺชติ เหตุปจฺจยา. เอกํ.  เหตุยา จตฺตาริ.}

\prose{อาสวญฺเจว อาสว\-สมฺปยุตฺตญฺจ นเหตุํ ธมฺมํ ปฏิจฺจ นอาสว\-วิปฺปยุตฺโต เจว นโน จ อาสโว นนเหตุ ธมฺโม อุปฺปชฺชติ เหตุปจฺจยา. เอกํ.}

\prose{อาสว\-สมฺปยุตฺต\-ญฺเจว โน จ อาสวํ นเหตุํ ธมฺมํ ปฏิจฺจ นอาสว\-วิปฺปยุตฺโต เจว นโน จ อาสโว นนเหตุ ธมฺโม อุปฺปชฺชติ เหตุปจฺจยา. ตีณิ.  เหตุยา จตฺตาริ.}

\csromanmarkh{19. อาสว\-วิปฺปยุตฺต\-สาสวทุก}\csromanmarkhii{1. เหตุทุก}\csromanheadernomark{2}{19. อาสว\-วิปฺปยุตฺต\-สาสวทุก 1. เหตุทุก}
\proseitem{56}{อาสว\-วิปฺปยุตฺตํ สาสวํ เหตุํ ธมฺมํ ปฏิจฺจ อาสว\-วิปฺปยุตฺโต นอนาสโว นเหตุ ธมฺโม อุปฺปชฺชติ เหตุปจฺจยา. เอกํ.}

\csromanmarkh{55. มหนฺตรทุก}\csromanmarkhii{1. เหตุทุก}\csromanheadernomark{2}{55. มหนฺตรทุก 1. เหตุทุก}
\prose{\csromanfolio{319}อาสว\-วิปฺปยุตฺตํ อนาสวํ เหตุํ ธมฺมํ ปฏิจฺจ อาสว\-วิปฺปยุตฺโต นอนาสโว นเหตุ ธมฺโม อุปฺปชฺชติ เหตุปจฺจยา. ตีณิ. เหตุยา จตฺตาริ.}

\prose{อาสว\-วิปฺปยุตฺตํ สาสวํ นเหตุํ ธมฺมํ ปฏิจฺจ อาสว\-วิปฺปยุตฺโต นอนาสโว นนเหตุ ธมฺโม อุปฺปชฺชติ เหตุปจฺจยา.  อาสว\-วิปฺปยุตฺตํ อนาสวํ นเหตุํ ธมฺมํ ปฏิจฺจ อาสว\-วิปฺปยุตฺโต นสาสโว นนเหตุ ธมฺโม อุปฺปชฺชติ เหตุปจฺจยา.  เหตุยา เทฺว.}

\csromanmarkh{20. สญฺโญชน\-โคจฺฉกาทิ}\csromanmarkhii{1. เหตุทุก}\csromanheadernomark{2}{20. สญฺโญชน\-โคจฺฉกาทิ 1. เหตุทุก}
\proseitem{57}{สญฺโญชนํ เหตุํ ธมฺมํ ปฏิจฺจ โนสญฺโญชโน นเหตุ ธมฺโม อุปฺปชฺชติ เหตุปจฺจยา.}

\proseitem{58}{คนฺถํ เหตุํ ธมฺมํ ปฏิจฺจ โนคนฺโถ นเหตุ ธมฺโม อุปฺปชฺชติ เหตุปจฺจยา.}

\proseitem{59}{โอฆํ เหตุํ ธมฺมํ ปฏิจฺจ โนโอโฆ นเหตุ ธมฺโม อุปฺปชฺชติ เหตุปจฺจยา.}

\proseitem{60}{โยคํ เหตุํ ธมฺมํ ปฏิจฺจ โนโยโค นเหตุ ธมฺโม อุปฺปชฺชติ เหตุปจฺจยา.}

\proseitem{61}{นีวรณํ เหตุํ ธมฺมํ ปฏิจฺจ โนนีวรโณ นเหตุ ธมฺโม อุปฺปชฺชติ เหตุปจฺจยา.}

\proseitem{62}{โนปรามาสํ เหตุํ ธมฺมํ ปฏิจฺจ นโนปรามาโส นเหตุ ธมฺโม อุปฺปชฺชติ เหตุปจฺจยา.  เหตุยา ตีณิ.}

\csromanmarkh{55. มหนฺตรทุก}\csromanmarkhii{1. เหตุทุก}\csromanheadernomark{2}{55. มหนฺตรทุก 1. เหตุทุก}
\proseitem{63}{สารมฺมณํ เหตุํ ธมฺมํ ปฏิจฺจ นสารมฺมโณ นเหตุ ธมฺโม อุปฺปชฺชติ เหตุปจฺจยา.  เหตุยา ตีณิ.}

\prose{สารมฺมณํ นเหตุํ ธมฺมํ ปฏิจฺจ นอนารมฺมโณ นนเหตุ ธมฺโม อุปฺปชฺชติ เหตุปจฺจยา.  อนารมฺมณํ นเหตุํ ธมฺมํ ปฏิจฺจ นอนารมฺมโณ นนเหตุ}

\vspace{\tipitakaparskip}
\csromancenter{ทุกทุก\-ปฏฺฐาน\-ปาฬิ ธมฺโม อุปฺปชฺชติ เหตุปจฺจยา.  สารมฺมณํ นเหตุญฺจ อนารมฺมณํ นเหตุญฺจ ธมฺมํ ปฏิจฺจ นอนารมฺมโณ นนเหตุ ธมฺโม อุปฺปชฺชติ เหตุปจฺจยา.  เหตุยา ตีณิ.}
\proseitem{64}{\csromanfolio{320}โนจิตฺตํ เหตุํ ธมฺมํ ปฏิจฺจ นโนจิตฺโต นเหตุ ธมฺโม อุปฺปชฺชติ เหตุปจฺจยา.  เหตุยา ตีณิ.}

\prose{จิตฺตํ นเหตุํ ธมฺมํ ปฏิจฺจ โนจิตฺโต นนเหตุ ธมฺโม อุปฺปชฺชติ เหตุปจฺจยา.  เหตุยา ตีณิ.}

\proseitem{65}{เจตสิกํ เหตุํ ธมฺมํ ปฏิจฺจ นเจตสิโก นเหตุ ธมฺโม อุปฺปชฺชติ เหตุปจฺจยา.  เหตุยา ตีณิ.}

\prose{เจตสิกํ นเหตุํ ธมฺมํ ปฏิจฺจ นอเจตสิโก นนเหตุ ธมฺโม อุปฺปชฺชติ เหตุปจฺจยา.  เหตุยา ตีณิ.}

\proseitem{66}{จิตฺต\-สมฺปยุตฺตํ เหตุํ ธมฺมํ ปฏิจฺจ นจิตฺต\-สมฺปยุตฺโต นเหตุ ธมฺโม อุปฺปชฺชติ เหตุปจฺจยา.  เหตุยา ตีณิ. (สํขิตฺตํ.)}

\proseitem{67}{จิตฺต\-สํสฏฺฐํ เหตุํ ธมฺมํ ปฏิจฺจ นจิตฺต\-สํสฏฺโฐ นเหตุ ธมฺโม อุปฺปชฺชติ เหตุปจฺจยา.  เหตุยา ตีณิ.}

\proseitem{68}{จิตฺต\-สมุฏฺฐานํ เหตุํ ธมฺมํ ปฏิจฺจ โนจิตฺต\-สมุฏฺฐาโน นเหตุ ธมฺโม อุปฺปชฺชติ เหตุปจฺจยา.  เหตุยา ตีณิ.}

\proseitem{69}{จิตฺตสหภุํ เหตุํ ธมฺมํ ปฏิจฺจ โนจิตฺตสหภู นเหตุ ธมฺโม อุปฺปชฺชติ เหตุปจฺจยา.  เหตุยา ตีณิ.}

\proseitem{70}{จิตฺตา\-นุปริวตฺติํ เหตุํ ธมฺมํ ปฏิจฺจ โนจิตฺตา\-นุปริวตฺตี นเหตุ ธมฺโม อุปฺปชฺชติ เหตุปจฺจยา.  เหตุยา ตีณิ.}

\proseitem{71}{จิตฺต\-สํสฏฺฐ\-สมุฏฺฐานํ เหตุํ ธมฺมํ ปฏิจฺจ โนจิตฺต\-สํสฏฺฐ\-สมุฏฺฐาโน นเหตุ ธมฺโม อุปฺปชฺชติ เหตุปจฺจยา.  เหตุยา ตีณิ.}

\proseitem{72}{จิตฺต\-สํสฏฺฐ\-สมุฏฺฐาน\-สหภุํ เหตุํ ธมฺมํ ปฏิจฺจ โนจิตฺต\-สํสฏฺฐ\-สมุฏฺฐาน\-สหภู นเหตุ ธมฺโม อุปฺปชฺชติ เหตุปจฺจยา.  เหตุยา ตีณิ.}

\csromanmarkh{69. อุปาทาน\-โคจฺฉกาทิ}\csromanmarkhii{1. เหตุทุก}\csromanheadernomark{2}{69. อุปาทาน\-โคจฺฉกาทิ 1. เหตุทุก}
\proseitem{73}{\csromanfolio{321}จิตฺต\-สํสฏฺฐ\-สมุฏฺฐานา\-นุปริวตฺติํ เหตุํ ธมฺมํ ปฏิจฺจ โนจิตฺต\-สํสฏฺฐ\-สมุฏฺฐานา\-นุปริวตฺตี นเหตุ ธมฺโม อุปฺปชฺชติ เหตุปจฺจยา.  เหตุยา ตีณิ.}

\proseitem{74}{พาหิรํ เหตุํ ธมฺมํ ปฏิจฺจ นพาหิโร นเหตุ ธมฺโม อุปฺปชฺชติ เหตุปจฺจยา.  เหตุยา ตีณิ.}

\prose{อชฺฌตฺติกํ นเหตุํ ธมฺมํ ปฏิจฺจ นอชฺฌตฺติโก นนเหตุ ธมฺโม อุปฺปชฺชติ เหตุปจฺจยา.  เหตุยา ตีณิ.}

\proseitem{75}{โนอุปาทา เหตุํ ธมฺมํ ปฏิจฺจ นโนอุปาทา นเหตุ ธมฺโม อุปฺปชฺชติ เหตุปจฺจยา.  เหตุยา ตีณิ.}

\prose{อุปาทา นเหตุํ ธมฺมํ ปฏิจฺจ โนอุปาทา นนเหตุ ธมฺโม อุปฺปชฺชติ เหตุปจฺจยา.  เหตุยา ตีณิ.}

\proseitem{76}{อุปาทินฺนํ เหตุํ ธมฺมํ ปฏิจฺจ นอุปาทินฺโน นเหตุ ธมฺโม อุปฺปชฺชติ เหตุปจฺจยา. ตีณิ.}

\prose{อนุปาทินฺนํ เหตุํ ธมฺมํ ปฏิจฺจ นอุปาทินฺโน นเหตุ ธมฺโม อุปฺปชฺชติ เหตุปจฺจยา เอกํ.  เหตุยา จตฺตาริ.}

\prose{อุปาทินฺนํ นเหตุํ ธมฺมํ ปฏิจฺจ นอนุปาทินฺโน นนเหตุ ธมฺโม อุปฺปชฺชติ เหตุปจฺจยา. เอกํ.}

\prose{อนุปาทินฺนํ นเหตุํ ธมฺมํ ปฏิจฺจ นอุปาทินฺโน นนเหตุ ธมฺโม อุปฺปชฺชติ เหตุปจฺจยา. เอกํ.  เหตุยา เทฺว.}

\csromanmarkh{69. อุปาทาน\-โคจฺฉกาทิ}\csromanmarkhii{1. เหตุทุก}\csromanheadernomark{2}{69. อุปาทาน\-โคจฺฉกาทิ 1. เหตุทุก}
\proseitem{77}{อุปาทานํ เหตุํ ธมฺมํ ปฏิจฺจ โนอุปาทาโน นเหตุ ธมฺโม อุปฺปชฺชติ เหตุปจฺจยา.}

\proseitem{78}{กิเลสํ เหตุํ ธมฺมํ ปฏิจฺจ โนกิเลโส นเหตุ ธมฺโม อุปฺปชฺชติ เหตุปจฺจยา.}

\csromanmarkh{ธมฺมา\-นุโลม\-ปจฺจนีย ทุกทุก\-ปฏฺฐาน\-ปาฬิ}\csromanmarkhii{83. ปิฏฺฐิทุก 1. เหตุทุก}\csromanheadernomark{2}{ธมฺมา\-นุโลม\-ปจฺจนีย ทุกทุก\-ปฏฺฐาน\-ปาฬิ 83. ปิฏฺฐิทุก 1. เหตุทุก}
\proseitem{79}{\csromanfolio{322}ทสฺสเนน ปหาตพฺพํ เหตุํ ธมฺมํ ปฏิจฺจ นทสฺสเนน ปหาตพฺโพ นเหตุ ธมฺโม อุปฺปชฺชติ เหตุปจฺจยา. ตีณิ.}

\prose{นทสฺสเนน ปหาตพฺพํ เหตุํ ธมฺมํ ปฏิจฺจ นทสฺสเนน ปหาตพฺโพ นเหตุ ธมฺโม อุปฺปชฺชติ เหตุปจฺจยา. เอกํ.  เหตุยา จตฺตาริ.}

\prose{ทสฺสเนน ปหาตพฺพํ นเหตุํ ธมฺมํ ปฏิจฺจ นนทสฺสเนน ปหาตพฺโพ นนเหตุ ธมฺโม เหตุปจฺจยา. เอกํ.}

\prose{นทสฺสเนน ปหาตพฺพํ นเหตุํ ธมฺมํ ปฏิจฺจ นทสฺสเนน ปหาตพฺโพ นนเหตุ ธมฺโม อุปฺปชฺชติ เหตุปจฺจยา. เอกํ.  เหตุยา เทฺว.}

\proseitem{80}{ภาวนาย ปหาตพฺพํ เหตุํ ธมฺมํ ปฏิจฺจ นภาวนาย ปหาตพฺโพ นเหตุ ธมฺโม อุปฺปชฺชติ เหตุปจฺจยา. ตีณิ.  นภาวนาย ปหาตพฺพํ เหตุํ ธมฺมํ ปฏิจฺจ นภาวนาย ปหาตพฺโพ นเหตุ ธมฺโม อุปฺปชฺชติ เหตุปจฺจยา. เอกํ.  เหตุยา จตฺตาริ.}

\prose{ภาวนาย ปหาตพฺพํ นเหตุํ ธมฺมํ ปฏิจฺจ นนภาวนาย ปหาตพฺโพ นนเหตุ ธมฺโม อุปฺปชฺชติ เหตุปจฺจยา. เอกํ.}

\prose{นภาวนาย ปหาตพฺพํ นเหตุํ ธมฺมํ ปฏิจฺจ นภาวนาย ปหาตพฺโพ นนเหตุ ธมฺโม อุปฺปชฺชติ เหตุปจฺจยา. เอกํ.  เหตุยา เทฺว.}

\proseitem{81}{ทสฺสเนน ปหาตพฺพ\-เหตุกํ เหตุํ ธมฺมํ ปฏิจฺจ นทสฺสเนน ปหาตพฺพ\-เหตุโก นเหตุ ธมฺโม อุปฺปชฺชติ เหตุปจฺจยา. ตีณิ.}

\prose{นทสฺสเนน ปหาตพฺพ\-เหตุกํ เหตุํ ธมฺมํ ปฏิจฺจ นนทสฺสเนน ปหาตพฺพ\-เหตุโก นเหตุ ธมฺโม อุปฺปชฺชติ เหตุปจฺจยา. ตีณิ.  (สพฺพตฺถ ปเญฺห สํขิตฺตํ.)}

\prose{ทสฺสเนน ปหาตพฺพ\-เหตุกํ นเหตุํ ธมฺมํ ปฏิจฺจ นนทสฺสเนน ปหาตพฺพ\-เหตุโก นนเหตุ ธมฺโม อุปฺปชฺชติ เหตุปจฺจยา.  นทสฺสเนน ปหาตพฺพ\-เหตุกํ นเหตุํ ธมฺมํ ปฏิจฺจ นทสฺสเนน ปหาตพฺพ\-เหตุโก นนเหตุ ธมฺโม อุปฺปชฺชติ เหตุปจฺจยา.  เหตุยา เทฺว.}

\csromanmarkh{100. สรณทุก}\csromanmarkhii{1. เหตุทุกาทิ 100. สรณทุก 1. เหตุทุกาทิ}\csromanheadernomark{2}{100. สรณทุก 1. เหตุทุกาทิ \textbf{100. สรณทุก 1. เหตุทุกาทิ}}
\proseitem{82}{\csromanfolio{323}สรณํ เหตุํ ธมฺมํ ปฏิจฺจ นสรโณ นเหตุ ธมฺโม อุปฺปชฺชติ เหตุปจฺจยา. ตีณิ.}

\prose{อรณํ เหตุํ ธมฺมํ ปฏิจฺจ นสรโณ นเหตุ ธมฺโม อุปฺปชฺชติ เหตุปจฺจยา. เอกํ.  เหตุยา จตฺตาริ.}

\prose{สรณํ นเหตุํ ธมฺมํ ปฏิจฺจ นอรโณ นนเหตุ ธมฺโม อุปฺปชฺชติ เหตุปจฺจยา. เอกํ.}

\prose{อรณํ นเหตุํ ธมฺมํ ปฏิจฺจ นสรโณ นนเหตุ ธมฺโม อุปฺปชฺชติ เหตุปจฺจยา. เอกํ.  เหตุยา เทฺว.}

\proseitem{83}{สรณํ สเหตุกํ ธมฺมํ ปฏิจฺจ นสรโณ นสเหตุโก ธมฺโม อุปฺปชฺชติ เหตุปจฺจยา. เอกํ.}

\prose{อรณํ สเหตุกํ ธมฺมํ ปฏิจฺจ นสรโณ นสเหตุโก ธมฺโม อุปฺปชฺชติ เหตุปจฺจยา. เอกํ.  เหตุยา เทฺว.}

\prose{สรณํ อเหตุกํ ธมฺมํ ปฏิจฺจ นอรโณ นอเหตุโก ธมฺโม อุปฺปชฺชติ เหตุปจฺจยา. เอกํ.}

\prose{อรณํ อเหตุกํ ธมฺมํ ปฏิจฺจ นสรโณ นอเหตุโก ธมฺโม อุปฺปชฺชติ เหตุปจฺจยา. เอกํ.  เหตุยา เทฺว.}

\proseitem{84}{สรณํ เหตุ\-สมฺปยุตฺตํ ธมฺมํ ปฏิจฺจ นสรโณ นเหตุ\-สมฺปยุตฺโต ธมฺโม อุปฺปชฺชติ เหตุปจฺจยา. เอกํ.  อรณํ เหตุ\-สมฺปยุตฺตํ ธมฺมํ ปฏิจฺจ นสรโณ นเหตุ\-สมฺปยุตฺโต ธมฺโม อุปฺปชฺชติ เหตุปจฺจยา. เอกํ.  เหตุยา เทฺว. (สเหตุก\-ทุก\-สทิสํ.)}

\proseitem{85}{สรณํ เหตุญฺเจว สเหตุกญฺจ ธมฺมํ ปฏิจฺจ นอรโณ นเหตุ เจว นอเหตุโก จ ธมฺโม อุปฺปชฺชติ เหตุปจฺจยา. เอกํ.}

\prose{อรณํ เหตุญฺเจว สเหตุกญฺจ ธมฺมํ ปฏิจฺจ นสรโณ นเหตุ เจว นอเหตุโก จ ธมฺโม อุปฺปชฺชติ เหตุปจฺจยา. เอกํ. เหตุยา เทฺว.}

\prose{สรณํ สเหตุกญฺเจว น จ เหตุํ ธมฺมํ ปฏิจฺจ นอรโณ นอเหตุโก เจว นน จ เหตุ ธมฺโม อุปฺปชฺชติ เหตุปจฺจยา. เอกํ.}

\csromanheader{2}{ธมฺมา\-นุโลม\-ปจฺจนีย ทุกทุก\-ปฏฺฐาน\-ปาฬิ}
\prose{\csromanfolio{324}อรณํ สเหตุกญฺเจว น จ เหตุํ ธมฺมํ ปฏิจฺจ นสรโณ นอเหตุโก เจว นน จ เหตุ ธมฺโม อุปฺปชฺชติ เหตุปจฺจยา. เอกํ.  เหตุยา เทฺว.}

\proseitem{86}{สรณํ เหตุญฺเจว เหตุ\-สมฺปยุตฺตญฺจ ธมฺมํ ปฏิจฺจ นอรโณ นเหตุ เจว นเหตุ\-วิปฺปยุตฺโต จ ธมฺโม อุปฺปชฺชติ เหตุปจฺจยา. (เหตุ\-สเหตุก\-ทุก\-สทิสํ.)}

\proseitem{87}{สรณํ นเหตุํ สเหตุกํ ธมฺมํ ปฏิจฺจ นสรโณ นเหตุ นสเหตุโก ธมฺโม อุปฺปชฺชติ เหตุปจฺจยา. เอกํ.}

\prose{อรณํ นเหตุํ สเหตุกํ ธมฺมํ ปฏิจฺจ นสรโณ นเหตุ นสเหตุโก ธมฺโม อุปฺปชฺชติ เหตุปจฺจยา. เอกํ.  เหตุยา เทฺว.}

\prose{อรณํ นเหตุํ อเหตุกํ ธมฺมํ ปฏิจฺจ นสรโณ นเหตุ นอเหตุโก ธมฺโม อุปฺปชฺชติ เหตุปจฺจยา.  เหตุยา เอกํ.}

\csromanmarkh{100. สรณทุก}\csromanmarkhii{7. จูฬนฺตรทุก}\csromanheadernomark{2}{100. \textbf{สรณทุก 7. จูฬนฺตรทุก}}
\proseitem{88}{อรโณ อปฺปจฺจโย ธมฺโม นสรณสฺส นอปฺปจฺจยสฺส ธมฺมสฺส อารมฺมณ\-ปจฺจเยน ปจฺจโย.}

\proseitem{89}{อรโณ อสงฺขโต ธมฺโม นสรณสฺส นอสงฺขตสฺส ธมฺมสฺส อารมฺมณ\-ปจฺจเยน ปจฺจโย.}

\proseitem{90}{อรโณ สนิทสฺสโน ธมฺโม นอรณสฺส นสนิทสฺสนสฺส ธมฺมสฺส อารมฺมณ\-ปจฺจเยน ปจฺจโย.  อรโณ สนิทสฺสโน ธมฺโม นสรณสฺส นสนิทสฺสนสฺส ธมฺมสฺส อารมฺมณ\-ปจฺจเยน ปจฺจโย.}

\csromanmarkh{100. สรณทุก}\csromanmarkhii{14. อาสวโคจฺฉก}\csromanheadernomark{2}{100. \textbf{สรณทุก 1}4. อาสวโคจฺฉก}
\proseitem{91}{สรณํ อาสวํ ธมฺมํ ปฏิจฺจ นสรโณ โนอาสโว ธมฺโม อุปฺปชฺชติ เหตุปจฺจยา.  สรณํ อาสวํ ธมฺมํ ปฏิจฺจ นอรโณ โนอาสโว ธมฺโม อุปฺปชฺชติ เหตุปจฺจยา.  สรณํ อาสวํ ธมฺมํ ปฏิจฺจ นสรโณ โนอาสโว จ นอรโณ โนอาสโว จ ธมฺมา อุปฺปชฺชนฺติ เหตุปจฺจยา.}

\csromanmarkh{100. สรณทุก}\csromanmarkhii{83. ปิฏฺฐิทุก}\csromanheadernomark{2}{100. สรณทุก 83. ปิฏฺฐิทุก}
\prose{\csromanfolio{325}สรณํ โนอาสวํ ธมฺมํ ปฏิจฺจ นอรโณ นโนอาสโว ธมฺโม อุปฺปชฺชติ เหตุปจฺจยา.  เหตุยา เอกํ.}

\proseitem{92}{อรณํ สาสวํ ธมฺมํ ปจฺจยา นสรโณ นสาสโว ธมฺโม อุปฺปชฺชติ เหตุปจฺจยา.}

\prose{อรณํ อนาสวํ ธมฺมํ ปฏิจฺจ นสรโณ นอนาสโว ธมฺโม อุปฺปชฺชติ เหตุปจฺจยา.  เหตุยา เอกํ. (เอเตน อุปาเยน สพฺพตฺถ วิตฺถาเรตพฺพํ.)}

\csromanmarkh{100. สรณทุก}\csromanmarkhii{55. มหนฺตรทุก}\csromanheadernomark{2}{100. \textbf{สรณทุก 5}5. มหนฺตรทุก}
\proseitem{93}{สรณํ สารมฺมณํ ธมฺมํ ปฏิจฺจ นสรโณ นสารมฺมโณ ธมฺโม อุปฺปชฺชติ เหตุปจฺจยา.  อรณํ สารมฺมณํ ธมฺมํ ปฏิจฺจ นสรโณ นสารมฺมโณ ธมฺโม อุปฺปชฺชติ เหตุปจฺจยา. (สํขิตฺตํ.)}

\csromanmarkh{100. สรณทุก}\csromanmarkhii{83. ปิฏฺฐิทุก}\csromanheadernomark{2}{100. สรณทุก 83. ปิฏฺฐิทุก}
\proseitem{94}{สรณํ ทสฺสเนน ปหาตพฺพํ ธมฺมํ ปฏิจฺจ นสรโณ นทสฺสเนน ปหาตพฺโพ ธมฺโม อุปฺปชฺชติ เหตุปจฺจยา.  เหตุยา เอกํ.}

\prose{อรณํ นทสฺสเนน ปหาตพฺพํ ธมฺมํ ปจฺจยา นอรโณ นนทสฺสเนน ปหาตพฺโพ ธมฺโม อุปฺปชฺชติ เหตุปจฺจยา. (สํขิตฺตํ.)}

\proseitem{95}{อรณํ สอุตฺตรํ ธมฺมํ ปจฺจยา นสรโณ นสอุตฺตโร ธมฺโม อุปฺปชฺชติ เหตุปจฺจยา.}

\prose{อรณํ อนุตฺตรํ ธมฺมํ ปฏิจฺจ นสรโณ นอนุตฺตโร ธมฺโม อุปฺปชฺชติ เหตุปจฺจยา.}

\vspace{\tipitakaparskip}
\csromancenter{(สหชาตวารมฺปิ ฯเปฯ ปญฺหาวารมฺปิ วิตฺถาเรตพฺพํ.)}
\nitthitam{ธมฺมา\-นุโลม\-ปจฺจนีเย ทุกทุก\-ปฏฺฐานํ นิฏฺฐิตํ.}
\nitthitam{อนุโลม\-ปจฺจนีย\-ปฏฺฐานํ นิฏฺฐิตํ.}
\pitaka{อภิธมฺม\-ปิฏก}
\csromanheader{2}{ธมฺม\-ปจฺจนียา\-นุโลม}
\namakkaram{นโม ตสฺส ภควโต อรหโต สมฺมา\-สมฺพุทฺธสฺส.}
\csromanmarkh{1. กุสลตฺติก}\csromanmarkhii{1. ปฏิจฺจวาร}\csromanheadernomark{2}{1. \textbf{กุสลตฺติก 1. ปฏิจฺจวาร}}
\prose{\csromanfolio{327}ปจฺจยจตุกฺกเหตุอารมฺมณ}

\proseitem{1}{นกุสลํ ธมฺมํ ปฏิจฺจ อกุสโล ธมฺโม อุปฺปชฺชติ เหตุปจฺจยา. อกุสลํ เอกํ ขนฺธํ ปฏิจฺจ ตโย ขนฺธา ฯเปฯ เทฺว ขนฺเธ ปฏิจฺจ เทฺว ขนฺธา.  นกุสลํ ธมฺมํ ปฏิจฺจ อพฺยากโต ธมฺโม อุปฺปชฺชติ เหตุปจฺจยา. วิปากาพฺยากตํ กิริยาพฺยากตํ เอกํ ขนฺธํ ปฏิจฺจ ตโย ขนฺธา จิตฺต\-สมุฏฺฐานญฺจ รูปํ ฯเปฯ.  นกุสลํ ธมฺมํ ปฏิจฺจ อกุสโล จ อพฺยากโต จ ธมฺมา อุปฺปชฺชนฺติ เหตุปจฺจยา. ตีณิ.}

\prose{นอกุสลํ ธมฺมํ ปฏิจฺจ กุสโล ธมฺโม อุปฺปชฺชติ เหตุปจฺจยา.  นอกุสลํ ธมฺมํ ปฏิจฺจ อพฺยากโต ธมฺโม อุปฺปชฺชติ เหตุปจฺจยา.  นอกุสลํ ธมฺมํ ปฏิจฺจ กุสโล จ อพฺยากโต จ ธมฺมา อุปฺปชฺชนฺติ เหตุปจฺจยา. ตีณิ.}

\prose{นอพฺยากตํ ธมฺมํ ปฏิจฺจ อพฺยากโต ธมฺโม อุปฺปชฺชติ เหตุปจฺจยา.  นอพฺยากตํ ธมฺมํ ปฏิจฺจ กุสโล ธมฺโม อุปฺปชฺชติ เหตุปจฺจยา.  นอพฺยากตํ ธมฺมํ ปฏิจฺจ อกุสโล ธมฺโม อุปฺปชฺชติ เหตุปจฺจยา.  นอพฺยากตํ ธมฺมํ ปฏิจฺจ กุสโล จ อพฺยากโต จ ธมฺมา อุปฺปชฺชนฺติ เหตุปจฺจยา.  นอพฺยากตํ ธมฺมํ ปฏิจฺจ อกุสโล จ อพฺยากโต จ ธมฺมา อุปฺปชฺชนฺติ เหตุปจฺจยา. ปญฺจ.}

\csromanpalirecto
\gambhira{ธมฺม\-ปจฺจนียา\-นุโลม ติกปฏฺฐาน\-ปาฬิ}
\prose{\csromanfolio{328}นกุสลญฺจ นอพฺยากตญฺจ ธมฺมํ ปฏิจฺจ อกุสโล ธมฺโม อุปฺปชฺชติ เหตุปจฺจยา.  นกุสลญฺจ นอพฺยากตญฺจ ธมฺมํ ปฏิจฺจ อพฺยากโต ธมฺโม อุปฺปชฺชติ เหตุปจฺจยา.  นกุสลญฺจ นอพฺยากตญฺจ ธมฺมํ ปฏิจฺจ อกุสโล จ อพฺยากโต จ ธมฺมา อุปฺปชฺชนฺติ เหตุปจฺจยา. ตีณิ.}

\prose{นอกุสลญฺจ นอพฺยากตญฺจ ธมฺมํ ปฏิจฺจ กุสโล ธมฺโม อุปฺปชฺชติ เหตุปจฺจยา.  นอกุสลญฺจ นอพฺยากตญฺจ ธมฺมํ ปฏิจฺจ อพฺยากโต ธมฺโม อุปฺปชฺชติ เหตุปจฺจยา.  นอกุสลญฺจ นอพฺยากตญฺจ ธมฺมํ ปฏิจฺจ กุสโล จ อพฺยากโต จ ธมฺมา อุปฺปชฺชนฺติ เหตุปจฺจยา. ตีณิ.}

\prose{นกุสลญฺจ นอกุสลญฺจ ธมฺมํ ปฏิจฺจ อพฺยากโต ธมฺโม อุปฺปชฺชติ เหตุปจฺจยา. เอกํ.}

\proseitem{2}{นกุสลํ ธมฺมํ ปฏิจฺจ อกุสโล ธมฺโม อุปฺปชฺชติ อารมฺมณ\-ปจฺจยา.  นกุสลํ ธมฺมํ ปฏิจฺจ อพฺยากโต ธมฺโม อุปฺปชฺชติ อารมฺมณ\-ปจฺจยา. เทฺว.}

\prose{นอกุสลํ ธมฺมํ ปฏิจฺจ กุสโล ธมฺโม อุปฺปชฺชติ อารมฺมณ\-ปจฺจยา.  นอกุสลํ ธมฺมํ ปฏิจฺจ อพฺยากโต ธมฺโม อุปฺปชฺชติ อารมฺมณ\-ปจฺจยา. เทฺว.}

\prose{นอพฺยากตํ ธมฺมํ ปฏิจฺจ กุสโล ธมฺโม อุปฺปชฺชติ อารมฺมณ\-ปจฺจยา.  นอพฺยากตํ ธมฺมํ ปฏิจฺจ อกุสโล ธมฺโม อุปฺปชฺชติ อารมฺมณ\-ปจฺจยา. เทฺว.}

\prose{นกุสลญฺจ นอพฺยากตญฺจ ธมฺมํ ปฏิจฺจ อกุสโล ธมฺโม อุปฺปชฺชติ อารมฺมณ\-ปจฺจยา. เอกํ.}

\prose{นอกุสลญฺจ นอพฺยากตญฺจ ธมฺมํ ปฏิจฺจ กุสโล ธมฺโม อุปฺปชฺชติ อารมฺมณ\-ปจฺจยา. เอกํ.}

\prose{นกุสลญฺจ นอกุสลญฺจ ธมฺมํ ปฏิจฺจ อพฺยากโต ธมฺโม อุปฺปชฺชติ อารมฺมณ\-ปจฺจยา. เอกํ.  เหตุยา อฏฺฐารส, อารมฺมเณ นว, อธิปติยา อฏฺฐารส ฯเปฯ อวิคเต อฏฺฐารส.}

\prose{ปจฺจนีย\-น\-เหตุ นอารมฺมณ}

\proseitem{3}{นกุสลํ ธมฺมํ ปฏิจฺจ อกุสโล ธมฺโม อุปฺปชฺชติ นเหตุปจฺจยา.  นกุสลํ ธมฺมํ ปฏิจฺจ อพฺยากโต ธมฺโม อุปฺปชฺชติ เหตุปจฺจยา ฯเปฯ. นกุสลญฺจ นอกุสลญฺจ ธมฺมํ ปฏิจฺจ อพฺยากโต ธมฺโม อุปฺปชฺชติ นเหตุปจฺจยา.}

\proseitem{4}{\csromanfolio{329}นกุสลํ ธมฺมํ ปฏิจฺจ อพฺยากโต ธมฺโม อุปฺปชฺชติ นอารมฺมณ\-ปจฺจยา ฯเปฯ. นกุสลญฺจ นอกุสลญฺจ ธมฺมํ ปฏิจฺจ อพฺยากโต ธมฺโม อุปฺปชฺชติ นอารมฺมณ\-ปจฺจยา. (สํขิตฺตํ.)}

\proseitem{5}{นเหตุยา ฉ, นอารมฺมเณ ฉ, นอธิปติยา อฏฺฐารส ฯเปฯ โนวิคเต ฉ.}

\prose{เหตุปจฺจยา นอารมฺมเณ ฉ. (สํขิตฺตํ.)}

\prose{นเหตุปจฺจยา อารมฺมเณ ฉ. (สํขิตฺตํ.) (สหชาตวาโร ปฏิจฺจ\-วาร\-สทิโส.)}

\csromanmarkh{1. กุสลตฺติก}\csromanmarkhii{3. ปจฺจยวาร}\csromanheadernomark{2}{1. กุสลตฺติก 3. ปจฺจยวาร}
\proseitem{6}{นกุสลํ ธมฺมํ ปจฺจยา กุสโล ธมฺโม อุปฺปชฺชติ เหตุปจฺจยา.  นกุสลํ ธมฺมํ ปจฺจยา อกุสโล ธมฺโม อุปฺปชฺชติ เหตุปจฺจยา.  นกุสลํ ธมฺมํ ปจฺจยา อพฺยากโต ธมฺโม อุปฺปชฺชติ เหตุปจฺจยา.  นกุสลํ ธมฺมํ ปจฺจยา กุสโล จ อพฺยากโต จ ธมฺมา อุปฺปชฺชนฺติ เหตุปจฺจยา.  นกุสลํ ธมฺมํ ปจฺจยา อกุสโล จ อพฺยากโต จ ธมฺมา อุปฺปชฺชนฺติ เหตุปจฺจยา. ปญฺจ.}

\prose{นอกุสลํ ธมฺมํ ปจฺจยา อกุสโล ธมฺโม อุปฺปชฺชติ เหตุปจฺจยา. ปญฺจ.}

\prose{นอพฺยากตํ ธมฺมํ ปจฺจยา อพฺยากโต ธมฺโม อุปฺปชฺชติ เหตุปจฺจยา. ปญฺจ.}

\prose{นกุสลญฺจ นอพฺยากตญฺจ ธมฺมํ ปจฺจยา. ตีณิ.}

\prose{นอกุสลญฺจ นอพฺยากตญฺจ ธมฺมํ ปจฺจยา. ตีณิ. นกุสลญฺจ นอกุสลญฺจ ธมฺมํ ปจฺจยา กุสโล ธมฺโม อุปฺปชฺชติ เหตุปจฺจยา. ปญฺจ. (สํขิตฺตํ.) . เหตุยา ฉพฺพีสติ, อารมฺมเณ เตรส ฯเปฯ อวิคเต ฉพฺพีสติ.}

\prose{สํสฏฺฐวาเร เหตุยา นว ฯเปฯ อวิคเต นว. (สํขิตฺตํ.)}

\csromanmarkh{ธมฺม\-ปจฺจนียา\-นุโลม ติกปฏฺฐาน\-ปาฬิ}\csromanmarkhii{1. กุสลตฺติก 7. ปญฺหาวาร}\csromanheadernomark{2}{ธมฺม\-ปจฺจนียา\-นุโลม ติกปฏฺฐาน\-ปาฬิ 1. กุสลตฺติก 7. ปญฺหาวาร}
\proseitem{7}{\csromanfolio{330}นกุสโล ธมฺโม อกุสลสฺส ธมฺมสฺส เหตุปจฺจเยน ปจฺจโย.  นกุสโล ธมฺโม อพฺยากตสฺส ธมฺมสฺส เหตุปจฺจเยน ปจฺจโย.  นกุสโล ธมฺโม อกุสลสฺส จ อพฺยากตสฺส จ ธมฺมสฺส เหตุปจฺจเยน ปจฺจโย. ตีณิ.}

\prose{นอกุสโล ธมฺโม กุสลสฺส ธมฺมสฺส เหตุปจฺจเยน ปจฺจโย. ตีณิ.}

\prose{นอพฺยากโต ธมฺโม อพฺยากตสฺส ธมฺมสฺส เหตุปจฺจเยน ปจฺจโย. ปญฺจ.}

\prose{นกุสโล จ นอพฺยากโต จ ธมฺมา อกุสลสฺส ธมฺมสฺส เหตุปจฺจเยน ปจฺจโย. ตีณิ.}

\prose{นอกุสโล จ นอพฺยากโต จ ธมฺมา กุสลสฺส ธมฺมสฺส เหตุปจฺจเยน ปจฺจโย. ตีณิ.}

\prose{นกุสโล จ นอกุสโล จ ธมฺมา อพฺยากตสฺส ธมฺมสฺส เหตุปจฺจเยน ปจฺจโย. เอกํ.}

\csromanheader{2}{\textbf{อารมฺมณ}}
\proseitem{8}{นกุสโล ธมฺโม กุสลสฺส ธมฺมสฺส อารมฺมณ\-ปจฺจเยน ปจฺจโย. ตีณิ.}

\prose{นอกุสโล ธมฺโม อกุสลสฺส ธมฺมสฺส อารมฺมณ\-ปจฺจเยน ปจฺจโย. ตีณิ.}

\prose{นอพฺยากโต ธมฺโม อพฺยากตสฺส ธมฺมสฺส อารมฺมณ\-ปจฺจเยน ปจฺจโย. ตีณิ.}

\prose{นกุสโล จ นอพฺยากโต จ ธมฺมา กุสลสฺส ธมฺมสฺส อารมฺมณ\-ปจฺจเยน ปจฺจโย.  นกุสโล จ นอพฺยากโต จ ธมฺมา อกุสลสฺส ธมฺมสฺส อารมฺมณ\-ปจฺจเยน ปจฺจโย.  นกุสโล จ นอพฺยากโต จ ธมฺมา อพฺยากตสฺส ธมฺมสฺส อารมฺมณ\-ปจฺจเยน ปจฺจโย. ตีณิ.}

\prose{นอกุสโล จ นอพฺยากโต จ ธมฺมา กุสลสฺส ธมฺมสฺส อารมฺมณ\-ปจฺจเยน ปจฺจโย.  นอกุสโล จ นอพฺยากโต จ ธมฺมา อกุสลสฺส ธมฺมสฺส อารมฺมณ\-ปจฺจเยน ปจฺจโย.  นอกุสโล จ นอพฺยากโต จ ธมฺมา อพฺยากตสฺส ธมฺมสฺส อารมฺมณ\-ปจฺจเยน ปจฺจโย. ตีณิ.}

\prose{\csromanfolio{331}นกุสโล จ นอกุสโล จ ธมฺมา กุสลสฺส ธมฺมสฺส อารมฺมณ\-ปจฺจเยน ปจฺจโย. ตีณิ. (สํขิตฺตํ.)}

\proseitem{9}{เหตุยา อฏฺฐารส, อารมฺมเณ อฏฺฐารส, อธิปติยา เตวีส ฯเปฯ อวิคเต ทฺวาวีส. (ปญฺหาวารํ วิตฺถาเรตพฺพํ.)}

\csromanmarkh{2. เวทนาตฺติก}\csromanmarkhii{1. ปฏิจฺจวาร}\csromanheadernomark{2}{2. \textbf{เวทนาตฺติก 1. ปฏิจฺจวาร}}
\proseitem{10}{นสุขาย เวทนาย สมฺปยุตฺตํ ธมฺมํ ปฏิจฺจ สุขาย เวทนาย สมฺปยุตฺโต ธมฺโม อุปฺปชฺชติ เหตุปจฺจยา.  นสุขาย เวทนาย สมฺปยุตฺตํ ธมฺมํ ปฏิจฺจ ทุกฺขาย เวทนาย สมฺปยุตฺโต ธมฺโม อุปฺปชฺชติ เหตุปจฺจยา.  นสุขาย เวทนาย สมฺปยุตฺตํ ธมฺมํ ปฏิจฺจ อทุกฺขม\-สุขาย เวทนาย สมฺปยุตฺโต ธมฺโม อุปฺปชฺชติ เหตุปจฺจยา. (3)}

\prose{นทุกฺขาย เวทนาย สมฺปยุตฺตํ ธมฺมํ ปฏิจฺจ ทุกฺขาย เวทนาย สมฺปยุตฺโต ธมฺโม อุปฺปชฺชติ เหตุปจฺจยา.  นทุกฺขาย เวทนาย สมฺปยุตฺตํ ธมฺมํ ปฏิจฺจ สุขาย เวทนาย สมฺปยุตฺโต ธมฺโม อุปฺปชฺชติ เหตุปจฺจยา.  นทุกฺขาย เวทนาย สมฺปยุตฺตํ ธมฺมํ ปฏิจฺจ อทุกฺขม\-สุขาย เวทนาย สมฺปยุตฺโต ธมฺโม อุปฺปชฺชติ เหตุปจฺจยา. (3)}

\prose{นอทุกฺขม\-สุขาย เวทนาย สมฺปยุตฺตํ ธมฺมํ ปฏิจฺจ อทุกฺขม\-สุขาย เวทนาย สมฺปยุตฺโต ธมฺโม อุปฺปชฺชติ เหตุปจฺจยา.  นอทุกฺขม\-สุขาย เวทนาย สมฺปยุตฺตํ ธมฺมํ ปฏิจฺจ สุขาย เวทนาย สมฺปยุตฺโต ธมฺโม อุปฺปชฺชติ เหตุปจฺจยา.  นอทุกฺขม\-สุขาย เวทนาย สมฺปยุตฺตํ ธมฺมํ ปฏิจฺจ ทุกฺขาย เวทนาย สมฺปยุตฺโต ธมฺโม อุปฺปชฺชติ เหตุปจฺจยา. (3) (สํขิตฺตํ.) . เหตุยา เอกวีส, อธิปติยา เอกวีส ฯเปฯ อวิคเต เอกวีส.}

\csromanmarkh{3. วิปากตฺติก}\csromanmarkhii{1. ปฏิจฺจวาร}\csromanheadernomark{2}{3. วิปากตฺติก 1. ปฏิจฺจวาร}
\proseitem{11}{นวิปากํ ธมฺมํ ปฏิจฺจ วิปาโก ธมฺโม อุปฺปชฺชติ เหตุปจฺจยา.  นวิปากํ ธมฺมํ ปฏิจฺจ วิปาก\-ธมฺม\-ธมฺโม อุปฺปชฺชติ เหตุปจฺจยา. ปญฺจ.}

\prose{นวิปาก\-ธมฺม\-ธมฺมํ ปฏิจฺจ วิปาโก ธมฺโม อุปฺปชฺชติ เหตุปจฺจยา.  นวิปาก\-ธมฺม\-ธมฺมํ ปฏิจฺจ เนว\-วิปาก\-น\-วิปาก\-ธมฺม\-ธมฺโม อุปฺปชฺชติ เหตุปจฺจยา. ตีณิ.}

\prose{\csromanfolio{332}นเน\-ว\-วิปาก\-น\-วิปากธมฺม\-ธมฺมํ ปฏิจฺจ เนว\-วิปาก\-น\-วิปาก\-ธมฺม\-ธมฺโม อุปฺปชฺชติ เหตุปจฺจยา.  นเน\-ว\-วิปาก\-น\-วิปากธมฺม\-ธมฺมํ ปฏิจฺจ วิปาโก ธมฺโม อุปฺปชฺชติ เหตุปจฺจยา. ปญฺจ.  เหตุยา ทฺวาวีส.}

\csromanmarkh{4. อุปาทินฺนตฺติก}\csromanmarkhii{1. ปฏิจฺจวาร}\csromanheadernomark{2}{4. \textbf{อุปาทินฺนตฺติก 1. ปฏิจฺจวาร}}
\proseitem{12}{นอุปาทินฺนุปาทานิยํ ธมฺมํ ปฏิจฺจ อนุปาทินฺนุ\-ปาทานิโย ธมฺโม อุปฺปชฺชติ เหตุปจฺจยา. ตีณิ.}

\prose{นอนุปาทินฺนุ\-ปาทานิยํ ธมฺมํ ปฏิจฺจ อนุปาทินฺนุ\-ปาทานิโย ธมฺโม อุปฺปชฺชติ เหตุปจฺจยา. ปญฺจ.}

\prose{นอนุปาทินฺน\-อนุปาทานิยํ ธมฺมํ ปฏิจฺจ อุปาทินฺนุ\-ปาทานิโย ธมฺโม อุปฺปชฺชติ เหตุปจฺจยา. ตีณิ.}

\prose{นอุปาทินฺนุปาทานิยญฺจ นอนุปาทินฺน\-อนุปาทานิยญฺจ ธมฺมํ ปฏิจฺจ อนุปาทินฺนุ\-ปาทานิโย ธมฺโม อุปฺปชฺชติ เหตุปจฺจยา. เอกํ.}

\prose{นอนุปาทินฺนุ\-ปาทานิยญฺจ นอนุปาทินฺน\-อนุปาทานิยญฺจ ธมฺมํ ปฏิจฺจ อุปาทินฺนุ\-ปาทานิโย ธมฺโม อุปฺปชฺชติ เหตุปจฺจยา. ตีณิ.}

\prose{นอุปาทินฺนุปาทานิยญฺจ นอนุปาทินฺนุ\-ปาทานิยญฺจ ธมฺมํ ปฏิจฺจ อนุปาทินฺน\-อนุปาทานิโย ธมฺโม อุปฺปชฺชติ เหตุปจฺจยา. ตีณิ.  เหตุยา อฏฺฐารส.}

\csromanmarkh{5. สํกิลิฏฺฐ\-ตฺติก}\csromanmarkhii{1. ปฏิจฺจวาร}\csromanheadernomark{2}{5. สํกิลิฏฺฐ\-ตฺติก 1. ปฏิจฺจวาร}
\prosecont{นสํกิลิฏฺฐ\-สํกิเลสิกํ ธมฺมํ ปฏิจฺจ อสํกิลิฏฺฐ\-สํกิเลสิโก ธมฺโม อุปฺปชฺชติ เหตุปจฺจยา.  นสํกิลิฏฺฐ\-สํกิเลสิกํ ธมฺมํ ปฏิจฺจ อสํกิลิฏฺฐ\-อสํกิเลสิโก ธมฺโม อุปฺปชฺชติ เหตุปจฺจยา.  นสํกิลิฏฺฐ\-สํกิเลสิกํ ธมฺมํ ปฏิจฺจ อสํกิลิฏฺฐ\-สํกิเลสิโก จ อสํกิลิฏฺฐ\-อสํกิเลสิโก จ ธมฺมา อุปฺปชฺชนฺติ เหตุปจฺจยา. ตีณิ.}

\proseitem{13}{นอสํกิลิฏฺฐ\-สํกิเลสิกํ ธมฺมํ ปฏิจฺจ อสํกิลิฏฺฐ\-สํกิเลสิโก ธมฺโม อุปฺปชฺชติ เหตุปจฺจยา. ปญฺจ.}

\prose{นอสํกิลิฏฺฐ\-อสํกิเลสิกํ ธมฺมํ ปฏิจฺจ สํกิลิฏฺฐ\-สํกิเลสิโก ธมฺโม อุปฺปชฺชติ เหตุปจฺจยา.  นอสํกิลิฏฺฐ\-อสํกิเลสิกํ ธมฺมํ ปฏิจฺจ}

\vspace{\tipitakaparskip}
\csromancenter{ทสฺสนตฺติก อสํกิลิฏฺฐ\-สํกิเลสิโก ธมฺโม อุปฺปชฺชติ เหตุปจฺจยา.  นอสํกิลิฏฺฐ\-อสํกิเลสิกํ ธมฺมํ ปฏิจฺจ สํกิลิฏฺฐ\-สํกิเลสิโก จ อสํกิลิฏฺฐ\-สํกิเลสิโก จ ธมฺมา อุปฺปชฺชนฺติ เหตุปจฺจยา. ตีณิ. เหตุยา อฏฺฐารส. 6. วิตกฺกตฺติก 1. ปฏิจฺจวาร}
\proseitem{14}{\csromanfolio{333}นสวิตกฺก\-สวิจารํ ธมฺมํ ปฏิจฺจ สวิตกฺก\-สวิจาโร ธมฺโม อุปฺปชฺชติ เหตุปจฺจยา.  นสวิตกฺก\-สวิจารํ ธมฺมํ ปฏิจฺจ อวิตกฺก\-วิจารมตฺโต ธมฺโม อุปฺปชฺชติ เหตุปจฺจยา.  นสวิตกฺก\-สวิจารํ ธมฺมํ ปฏิจฺจ อวิตกฺก\-อวิจาโร ธมฺโม อุปฺปชฺชติ เหตุปจฺจยา. สตฺต.}

\prose{นอวิตกฺก\-วิจารมตฺตํ ธมฺมํ ปฏิจฺจ อวิตกฺก\-วิจารมตฺโต ธมฺโม อุปฺปชฺชติ เหตุปจฺจยา. สตฺต.}

\prose{นอวิตกฺก\-อวิจารํ ธมฺมํ ปฏิจฺจ อวิตกฺก\-อวิจาโร ธมฺโม อุปฺปชฺชติ เหตุปจฺจยา. สตฺต.  เหตุยา เอกูนปญฺญาส. 7. ปีติตฺติก 1. ปฏิจฺจวาร}

\proseitem{15}{นปีติสหคตํ ธมฺมํ ปฏิจฺจ ปีติสหคโต ธมฺโม อุปฺปชฺชติ เหตุปจฺจยา. จตฺตาริ.}

\prose{นสุขสหคตํ ธมฺมํ ปฏิจฺจ สุขสหคโต ธมฺโม อุปฺปชฺชติ เหตุปจฺจยา. จตฺตาริ.}

\prose{นอุเปกฺขาสหคตํ ธมฺมํ ปฏิจฺจ อุเปกฺขา\-สหคโต ธมฺโม อุปฺปชฺชติ เหตุปจฺจยา.  นอุเปกฺขาสหคตํ ธมฺมํ ปฏิจฺจ ปีติสหคโต ธมฺโม อุปฺปชฺชติ เหตุปจฺจยา. จตฺตาริ.  เหตุยา อฏฺฐวีส. 8. ทสฺสนตฺติก 1. ปฏิจฺจวาร}

\proseitem{16}{นทสฺสเนน ปหาตพฺพํ ธมฺมํ ปฏิจฺจ ภาวนาย ปหาตพฺโพ ธมฺโม อุปฺปชฺชติ เหตุปจฺจยา.  นทสฺสเนน ปหาตพฺพํ ธมฺมํ ปฏิจฺจ เนวทสฺสเนน นภาวนาย ปหาตพฺโพ ธมฺโม อุปฺปชฺชติ เหตุปจฺจยา.  นทสฺสเนน ปหาตพฺพํ ธมฺมํ ปฏิจฺจ ภาวนาย ปหาตพฺโพ จ เนวทสฺสเนน นภาวนาย ปหาตพฺโพ จ ธมฺมา อุปฺปชฺชนฺติ เหตุปจฺจยา. ตีณิ.}

\prose{\csromanfolio{334}นภาวนาย ปหาตพฺพํ ธมฺมํ ปฏิจฺจ ทสฺสเนน ปหาตพฺโพ ธมฺโม อุปฺปชฺชติ เหตุปจฺจยา.  นภาวนาย ปหาตพฺพํ ธมฺมํ ปฏิจฺจ เนวทสฺสเนน นภาวนาย ปหาตพฺโพ ธมฺโม อุปฺปชฺชติ เหตุปจฺจยา.  นภาวนาย ปหาตพฺพํ ธมฺมํ ปฏิจฺจ ทสฺสเนน ปหาตพฺโพ จ เนวทสฺสเนน นภาวนาย ปหาตพฺโพ จ ธมฺมา อุปฺปชฺชนฺติ เหตุปจฺจยา. ตีณิ.}

\prose{นเน\-ว\-ทสฺสเนน นภาวนาย ปหาตพฺพํ ธมฺมํ ปฏิจฺจ เนวทสฺสเนน นภาวนาย ปหาตพฺโพ ธมฺโม อุปฺปชฺชติ เหตุปจฺจยา.  นเน\-ว\-ทสฺสเนน นภาวนาย ปหาตพฺพํ ธมฺมํ ปฏิจฺจ ทสฺสเนน ปหาตพฺโพ ธมฺโม อุปฺปชฺชติ เหตุปจฺจยา. ปญฺจ. (สํขิตฺตํ.) . เหตุยา อฏฺฐารส.}

\csromanmarkh{9. ทสฺสน\-เหตุ\-ตฺติก}\csromanmarkhii{1. ปฏิจฺจวาร}\csromanheadernomark{2}{9. ทสฺสน\-เหตุ\-ตฺติก 1. ปฏิจฺจวาร}
\proseitem{17}{นทสฺสเนน ปหาตพฺพ\-เหตุกํ ธมฺมํ ปฏิจฺจ ทสฺสเนน ปหาตพฺพ\-เหตุโก ธมฺโม อุปฺปชฺชติ เหตุปจฺจยา.  เหตุยา ฉพฺพีส.}

\csromanmarkh{10. อาจยคามิตฺติก}\csromanmarkhii{1. ปฏิจฺจวาร}\csromanheadernomark{2}{10. \textbf{อาจยคามิตฺติก 1. ปฏิจฺจวาร}}
\proseitem{18}{นอาจยคามิํ ธมฺมํ ปฏิจฺจ อปจยคามี ธมฺโม อุปฺปชฺชติ เหตุปจฺจยา.  เหตุยา อฏฺฐารส.}

\csromanmarkh{11. เสกฺขตฺติก}\csromanmarkhii{1. ปฏิจฺจวาร}\csromanheadernomark{2}{11. \textbf{เสกฺขตฺติก 1. ปฏิจฺจวาร}}
\vspace{\tipitakaparskip}
\csromancenter{นเสกฺขํ ธมฺมํ ปฏิจฺจ อเสกฺโข ธมฺโม อุปฺปชฺชติ เหตุปจฺจยา.  เหตุยา อฏฺฐารส.}
\csromanmarkh{12-16. ปริตฺตตฺติกาทิ}\csromanmarkhii{1. ปฏิจฺจวาร}\csromanheadernomark{2}{12-16. ปริตฺตตฺติกาทิ 1. ปฏิจฺจวาร}
\proseitem{20}{นปริตฺตํ ธมฺมํ ปฏิจฺจ ปริตฺโต ธมฺโม อุปฺปชฺชติ เหตุปจฺจยา. เหตุยา ทฺวาวีส.}

\proseitem{21}{นปริตฺตา\-รมฺมณํ ธมฺมํ ปฏิจฺจ ปริตฺตารมฺมโณ ธมฺโม อุปฺปชฺชติ เหตุปจฺจยา.  เหตุยา เตรส.}

\proseitem{22}{นหีนํ ธมฺมํ ปฏิจฺจ มชฺฌิโม ธมฺโม อุปฺปชฺชติ เหตุปจฺจยา. เหตุยา อฏฺฐารส.}

\csromanheader{2}{20-21. อชฺฌตฺต\-ตฺติก\-ทฺวย}
\proseitem{23}{\csromanfolio{335}นมิจฺฉตฺต\-นิยตํ ธมฺมํ ปฏิจฺจ สมฺมตฺตนิยโต ธมฺโม อุปฺปชฺชติ เหตุปจฺจยา.  เหตุยา อฏฺฐารส.}

\proseitem{24}{นมคฺคา\-รมฺมณํ ธมฺมํ ปฏิจฺจ มคฺคเหตุโก ธมฺโม อุปฺปชฺชติ เหตุปจฺจยา.  เหตุยา ทส\csromansharedfootnote{e0ee4924e42dd2cb}{เหตุยา อฏฺฐ?}.}

\csromanmarkh{17. อุปฺปนฺนตฺติก}\csromanmarkhii{7. ปญฺหาวาร}\csromanheadernomark{2}{17. \textbf{อุปฺปนฺนตฺติก 7. ปญฺหาวาร}}
\proseitem{25}{นอนุปฺปนฺโน ธมฺโม อุปฺปนฺนสฺส ธมฺมสฺส เหตุปจฺจเยน ปจฺจโย.  นอุปฺปาที ธมฺโม อุปฺปนฺนสฺส ธมฺมสฺส เหตุปจฺจเยน ปจฺจโย.  นอนุปฺปนฺโน จ นอุปฺปาที จ ธมฺมา อุปฺปนฺนสฺส ธมฺมสฺส เหตุปจฺจเยน ปจฺจโย. เหตุยา ตีณิ.}

\csromanheader{2}{18-19. อตีต\-ตฺติก\-ทฺวย\csromansharedfootnote{e0ee4924e42dd2cb}{เหตุยา อฏฺฐ?}-7. ปฏิจฺจาทิวาร}
\proseitem{26}{นอตีโต ธมฺโม ปจฺจุปฺปนฺนสฺส ธมฺมสฺส เหตุปจฺจเยน ปจฺจโย.  นอนาคโต ธมฺโม ปจฺจุปฺปนฺนสฺส ธมฺมสฺส เหตุปจฺจเยน ปจฺจโย.  นอตีโต จ นอนาคโต จ ธมฺมา ปจฺจุปฺปนฺนสฺส ธมฺมสฺส เหตุปจฺจเยน ปจฺจโย.  เหตุยา ตีณิ.}

\proseitem{27}{นอตีตา\-รมฺมณํ ธมฺมํ ปฏิจฺจ อตีตารมฺมโณ ธมฺโม อุปฺปชฺชติ เหตุปจฺจยา.  นอตีตา\-รมฺมณํ ธมฺมํ ปฏิจฺจ อนาคตารมฺมโณ ธมฺโม อุปฺปชฺชติ เหตุปจฺจยา.  นอตีตา\-รมฺมณํ ธมฺมํ ปฏิจฺจ ปจฺจุปฺปนฺนา\-รมฺมโณ ธมฺโม อุปฺปชฺชติ เหตุปจฺจยา. ตีณิ.}

\prose{นอนาคตา\-รมฺมณํ ธมฺมํ ปฏิจฺจ อตีตารมฺมโณ ธมฺโม อุปฺปชฺชติ เหตุปจฺจยา.  นอนาคตา\-รมฺมณํ ธมฺมํ ปฏิจฺจ ปจฺจุปฺปนฺนา\-รมฺมโณ ธมฺโม อุปฺปชฺชติ เหตุปจฺจยา. เทฺว.}

\prose{นปจฺจุปฺปนฺนา\-รมฺมณํ ธมฺมํ ปฏิจฺจ ปจฺจุปฺปนฺนา\-รมฺมโณ ธมฺโม อุปฺปชฺชติ เหตุปจฺจยา.  นปจฺจุปฺปนฺนา\-รมฺมณํ ธมฺมํ ปฏิจฺจ อตีตารมฺมโณ ธมฺโม อุปฺปชฺชติ เหตุปจฺจยา.  นปจฺจุปฺปนฺนา\-รมฺมณํ ธมฺมํ ปฏิจฺจ อนาคตารมฺมโณ ธมฺโม อุปฺปชฺชติ เหตุปจฺจยา. ตีณิ. (สํขิตฺตํ.) . เหตุยา สตฺตรส.}

\csromanmarkh{20-21. อชฺฌตฺต\-ตฺติก\-ทฺวย}\csromanmarkhii{1. ปฏิจฺจวาร}\csromanheadernomark{2}{20-21. อชฺฌตฺต\-ตฺติก\-ทฺวย 1. ปฏิจฺจวาร}
\proseitem{28}{นอชฺฌตฺตํ ธมฺมํ ปฏิจฺจ พหิทฺธา ธมฺโม อุปฺปชฺชติ เหตุปจฺจยา. เอกํ.}

\prose{\csromanfolio{336}นพหิทฺธา ธมฺมํ ปฏิจฺจ อชฺฌตฺโต ธมฺโม อุปฺปชฺชติ เหตุปจฺจยา. เอกํ. เหตุยา เทฺว.}

\proseitem{29}{นอชฺฌตฺตา\-รมฺมณํ ธมฺมํ ปฏิจฺจ อชฺฌตฺตา\-รมฺมโณ ธมฺโม อุปฺปชฺชติ เหตุปจฺจยา. เทฺว.}

\prose{นพหิทฺธา\-รมฺมณํ ธมฺมํ ปฏิจฺจ พหิทฺธา\-รมฺมโณ ธมฺโม อุปฺปชฺชติ เหตุปจฺจยา. เทฺว.}

\prose{นอชฺฌตฺตา\-รมฺมณญฺจ นพหิทฺธา\-รมฺมณญฺจ ธมฺมํ ปฏิจฺจ อชฺฌตฺตา\-รมฺมโณ ธมฺโม อุปฺปชฺชติ เหตุปจฺจยา. เทฺว.  เหตุยา ฉ.}

\csromanmarkh{22. สนิทสฺสน\-ตฺติก}\csromanmarkhii{1. ปฏิจฺจวาร}\csromanheadernomark{2}{22. สนิทสฺสน\-ตฺติก 1. ปฏิจฺจวาร}
\prosecont{นสนิทสฺสน\-สปฺปฏิฆํ ธมฺมํ ปฏิจฺจ สนิทสฺสน\-สปฺปฏิโฆ ธมฺโม อุปฺปชฺชติ เหตุปจฺจยา.  นสนิทสฺสน\-สปฺปฏิฆํ ธมฺมํ ปฏิจฺจ อนิทสฺสน\-สปฺปฏิโฆ ธมฺโม อุปฺปชฺชติ เหตุปจฺจยา.  นสนิทสฺสน\-สปฺปฏิฆํ ธมฺมํ ปฏิจฺจ อนิทสฺสน\-อปฺปฏิโฆ ธมฺโม อุปฺปชฺชติ เหตุปจฺจยา.  นสนิทสฺสน\-สปฺปฏิฆํ ธมฺมํ ปฏิจฺจ สนิทสฺสน\-สปฺปฏิโฆ จ อนิทสฺสน\-อปฺปฏิโฆ จ ธมฺมา อุปฺปชฺชนฺติ เหตุปจฺจยา.  นสนิทสฺสน\-สปฺปฏิฆํ ธมฺมํ ปฏิจฺจ อนิทสฺสน\-สปฺปฏิโฆ จ อนิทสฺสน\-อปฺปฏิโฆ จ ธมฺมา อุปฺปชฺชนฺติ เหตุปจฺจยา.  นสนิทสฺสน\-สปฺปฏิฆํ ธมฺมํ ปฏิจฺจ สนิทสฺสน\-สปฺปฏิโฆ จ อนิทสฺสน\-สปฺปฏิโฆ จ ธมฺมา อุปฺปชฺชนฺติ เหตุปจฺจยา.  นสนิทสฺสน\-สปฺปฏิฆํ ธมฺมํ ปฏิจฺจ สนิทสฺสน\-สปฺปฏิโฆ จ อนิทสฺสน\-สปฺปฏิโฆ จ อนิทสฺสน\-อปฺปฏิโฆ จ ธมฺมา อุปฺปชฺชนฺติ เหตุปจฺจยา. สตฺต.}

\proseitem{30}{นอนิทสฺสน\-สปฺปฏิฆํ ธมฺมํ ปฏิจฺจ อนิทสฺสน\-สปฺปฏิโฆ ธมฺโม อุปฺปชฺชติ เหตุปจฺจยา.  นอนิทสฺสน\-สปฺปฏิฆํ ธมฺมํ ปฏิจฺจ สนิทสฺสน\-สปฺปฏิโฆ ธมฺโม อุปฺปชฺชติ เหตุปจฺจยา. สตฺต.}

\prosewithcloser{นอนิทสฺสน\-อปฺปฏิฆํ ธมฺมํ ปฏิจฺจ อนิทสฺสน\-อปฺปฏิโฆ ธมฺโม อุปฺปชฺชติ เหตุปจฺจยา. สตฺต.  เหตุยา ปญฺจติํส. (ปญฺหาวารมฺปิ วิตฺถาเรตพฺพํ.)}{ธมฺม\-ปจฺจนียา\-นุโลเม ติกปฏฺฐานํ นิฏฺฐิตํ.}
\pitaka{อภิธมฺม\-ปิฏก}
\csromanheader{2}{ธมฺม\-ปจฺจนียา\-นุโลม}
\namakkaram{นโม ตสฺส ภควโต อรหโต สมฺมา\-สมฺพุทฺธสฺส.}
\csromanmarkh{1-6. เหตุโคจฺฉก}\csromanmarkhii{1. ปฏิจฺจวาร}\csromanheadernomark{2}{1-6. เหตุโคจฺฉก 1. ปฏิจฺจวาร}
\proseitem{1}{\csromanfolio{337}นเหตุํ ธมฺมํ ปฏิจฺจ เหตุ ธมฺโม อุปฺปชฺชติ เหตุปจฺจยา.  นเหตุํ ธมฺมํ ปฏิจฺจ นเหตุ ธมฺโม อุปฺปชฺชติ เหตุปจฺจยา.  นเหตุํ ธมฺมํ ปฏิจฺจ เหตุ จ นเหตุ จ ธมฺมา อุปฺปชฺชนฺติ เหตุปจฺจยา. (3)}

\prose{นนเหตุํ ธมฺมํ ปฏิจฺจ เหตุ ธมฺโม อุปฺปชฺชติ เหตุปจฺจยา.  นนเหตุํ ธมฺมํ ปฏิจฺจ นเหตุ ธมฺโม อุปฺปชฺชติ เหตุปจฺจยา.  นนเหตุํ ธมฺมํ ปฏิจฺจ เหตุ จ นเหตุ จ ธมฺมา อุปฺปชฺชนฺติ เหตุปจฺจยา. (3)}

\prose{นเหตุญฺจ นนเหตุญฺจ ธมฺมํ ปฏิจฺจ เหตุ ธมฺโม อุปฺปชฺชติ เหตุปจฺจยา.  นเหตุญฺจ นนเหตุญฺจ ธมฺมํ ปฏิจฺจ นเหตุ ธมฺโม อุปฺปชฺชติ เหตุปจฺจยา.  นเหตุญฺจ นนเหตุญฺจ ธมฺมํ ปฏิจฺจ เหตุ จ นเหตุ จ ธมฺมา อุปฺปชฺชนฺติ เหตุปจฺจยา. (3) (สํขิตฺตํ.) . เหตุยา นว, อารมฺมเณ นว ฯเปฯ อวิคเต นว.}

\vspace{\tipitakaparskip}
\csromancenter{(สหชาตวารมฺปิ ฯเปฯ ปญฺหาวารมฺปิ วิตฺถาเรตพฺพํ.)}
\proseitem{2}{นสเหตุกํ ธมฺมํ ปฏิจฺจ สเหตุโก ธมฺโม อุปฺปชฺชติ เหตุปจฺจยา.  นสเหตุกํ ธมฺมํ ปฏิจฺจ อเหตุโก ธมฺโม อุปฺปชฺชติ เหตุปจฺจยา.  นสเหตุกํ ธมฺมํ ปฏิจฺจ สเหตุโก จ อเหตุโก จ ธมฺมา อุปฺปชฺชนฺติ เหตุปจฺจยา. ตีณิ.}

\csromanpalirecto
\gambhira{ธมฺม\-ปจฺจนียา\-นุโลม ทุกปฏฺฐาน\-ปาฬิ}
\prose{\csromanfolio{338}นอเหตุกํ ธมฺมํ ปฏิจฺจ อเหตุโก ธมฺโม อุปฺปชฺชติ เหตุปจฺจยา.  นอเหตุกํ ธมฺมํ ปฏิจฺจ สเหตุโก ธมฺโม อุปฺปชฺชติ เหตุปจฺจยา.  นอเหตุกํ ธมฺมํ ปฏิจฺจ สเหตุโก จ อเหตุโก จ ธมฺมา อุปฺปชฺชนฺติ เหตุปจฺจยา. ตีณิ.}

\prose{นสเหตุกญฺจ นอเหตุกญฺจ ธมฺมํ ปฏิจฺจ สเหตุโก ธมฺโม อุปฺปชฺชติ เหตุปจฺจยา.  นสเหตุกญฺจ นอเหตุกญฺจ ธมฺมํ ปฏิจฺจ อเหตุโก ธมฺโม อุปฺปชฺชติ เหตุปจฺจยา.  นสเหตุกญฺจ นอเหตุกญฺจ ธมฺมํ ปฏิจฺจ สเหตุโก จ อเหตุโก จ ธมฺมา อุปฺปชฺชนฺติ เหตุปจฺจยา. ตีณิ.  เหตุยา นว.}

\proseitem{3}{นเหตุ\-สมฺปยุตฺตํ ธมฺมํ ปฏิจฺจ เหตุสมฺปยุตฺโต ธมฺโม อุปฺปชฺชติ เหตุปจฺจยา.  เหตุยา นว.}

\proseitem{4}{นเหตุญฺเจว นอเหตุกญฺจ ธมฺมํ ปฏิจฺจ เหตุ เจว สเหตุโก ธมฺโม อุปฺปชฺชติ เหตุปจฺจยา.  เหตุยา นว.}

\proseitem{5}{นเหตุญฺเจว นเหตุ\-วิปฺปยุตฺตญฺจ ธมฺมํ ปฏิจฺจ เหตุ เจว เหตุสมฺปยุตฺโต จ ธมฺโม อุปฺปชฺชติ เหตุปจฺจยา.  เหตุยา นว.}

\proseitem{6}{นเหตุํ นสเหตุกํ ธมฺมํ ปฏิจฺจ นเหตุ สเหตุโก ธมฺโม อุปฺปชฺชติ เหตุปจฺจยา.  เหตุยา นว.}

\csromanheader{2}{7-13. จูฬนฺตรทุก}
\proseitem{7}{นอปฺปจฺจยํ ธมฺมํ ปฏิจฺจ สปฺปจฺจโย ธมฺโม อุปฺปชฺชติ เหตุปจฺจยา.  เหตุยา เอกํ.}

\proseitem{8}{นอสงฺขตํ ธมฺมํ ปฏิจฺจ สงฺขโต ธมฺโม อุปฺปชฺชติ เหตุปจฺจยา.  เหตุยา เอกํ.}

\proseitem{9}{นสนิทสฺสนํ ธมฺมํ ปฏิจฺจ สนิทสฺสโน ธมฺโม อุปฺปชฺชติ เหตุปจฺจยา.  นสนิทสฺสนํ ธมฺมํ ปฏิจฺจ อนิทสฺสโน ธมฺโม อุปฺปชฺชติ เหตุปจฺจยา.  นสนิทสฺสนํ ธมฺมํ ปฏิจฺจ สนิทสฺสโน จ อนิทสฺสโน จ ธมฺมา อุปฺปชฺชนฺติ เหตุปจฺจยา.  เหตุยา ตีณิ.}

\csromanheader{2}{14-19. อาสวโคจฺฉก}
\prose{\csromanfolio{339}นสนิทสฺสโน ธมฺโม สนิทสฺสนสฺส ธมฺมสฺส เหตุปจฺจเยน ปจฺจโย.  เหตุยา ตีณิ.}

\prose{นสนิทสฺสโน ธมฺโม อนิทสฺสนสฺส ธมฺมสฺส อารมฺมณ\-ปจฺจเยน ปจฺจโย.  นอนิทสฺสโน ธมฺโม อนิทสฺสนสฺส ธมฺมสฺส อารมฺมณ\-ปจฺจเยน ปจฺจโย.}

\proseitem{10}{นสปฺปฏิฆํ ธมฺมํ ปฏิจฺจ สปฺปฏิโฆ ธมฺโม อุปฺปชฺชติ เหตุปจฺจยา.  นสปฺปฏิฆํ ธมฺมํ ปฏิจฺจ อปฺปฏิโฆ ธมฺโม อุปฺปชฺชติ เหตุปจฺจยา.  นสปฺปฏิฆํ ธมฺมํ ปฏิจฺจ สปฺปฏิโฆ จ อปฺปฏิโฆ จ ธมฺมา อุปฺปชฺชนฺติ เหตุปจฺจยา. ตีณิ.}

\prose{นอปฺปฏิฆํ ธมฺมํ ปฏิจฺจ อปฺปฏิโฆ ธมฺโม อุปฺปชฺชติ เหตุปจฺจยา.  นอปฺปฏิฆํ ธมฺมํ ปฏิจฺจ สปฺปฏิโฆ ธมฺโม อุปฺปชฺชติ เหตุปจฺจยา.  นอปฺปฏิฆํ ธมฺมํ ปฏิจฺจ สปฺปฏิโฆ จ อปฺปฏิโฆ จ ธมฺมา อุปฺปชฺชนฺติ เหตุปจฺจยา. ตีณิ.}

\prose{นสปฺปฏิฆญฺจ นอปฺปฏิฆญฺจ ธมฺมํ ปฏิจฺจ สปฺปฏิโฆ ธมฺโม อุปฺปชฺชติ เหตุปจฺจยา. ตีณิ.  เหตุยา นว.}

\proseitem{11}{นรูปิํ ธมฺมํ ปฏิจฺจ รูปี ธมฺโม อุปฺปชฺชติ เหตุปจฺจยา.  เหตุยา นว.}

\proseitem{12}{นโลกิยํ ธมฺมํ ปฏิจฺจ โลกิโย ธมฺโม อุปฺปชฺชติ เหตุปจฺจยา.  นโลกิยํ ธมฺมํ ปฏิจฺจ โลกุตฺตโร ธมฺโม อุปฺปชฺชติ เหตุปจฺจยา.  นโลกิยํ ธมฺมํ ปฏิจฺจ โลกิโย จ โลกุตฺตโร จ ธมฺมา อุปฺปชฺชนฺติ เหตุปจฺจยา. (3)}

\prose{นโลกุตฺตรํ ธมฺมํ ปฏิจฺจ โลกิโย ธมฺโม อุปฺปชฺชติ เหตุปจฺจยา. (1)}

\prose{นโลกิยญฺจ นโลกุตฺตรญฺจ ธมฺมํ ปฏิจฺจ โลกิโย ธมฺโม อุปฺปชฺชติ เหตุปจฺจยา. (1) เหตุยา ปญฺจ.}

\proseitem{13}{นเกนจิ วิญฺเญยฺยํ ธมฺมํ ปฏิจฺจ เกนจิ วิญฺเญยฺโย ธมฺโม อุปฺปชฺชติ เหตุปจฺจยา.  เหตุยา นว.}

\csromanheader{2}{14-19. อาสวโคจฺฉก}
\proseitem{14}{นอาสวํ ธมฺมํ ปฏิจฺจ อาสโว ธมฺโม อุปฺปชฺชติ เหตุปจฺจยา.  นอาสวํ ธมฺมํ ปฏิจฺจ โนอาสโว ธมฺโม อุปฺปชฺชติ เหตุปจฺจยา.  }

\prose{\csromanfolio{340}นอาสวํ ธมฺมํ ปฏิจฺจ อาสโว จ โนอาสโว จ ธมฺมา อุปฺปชฺชนฺติ เหตุปจฺจยา. ตีณิ.}

\prose{นโนอาสวํ ธมฺมํ ปฏิจฺจ โนอาสโว ธมฺโม อุปฺปชฺชติ เหตุปจฺจยา.  นโนอาสวํ ธมฺมํ ปฏิจฺจ อาสโว ธมฺโม อุปฺปชฺชติ เหตุปจฺจยา.  นโนอาสวํ ธมฺมํ ปฏิจฺจ อาสโว จ โนอาสโว จ ธมฺมา อุปฺปชฺชนฺติ เหตุปจฺจยา. ตีณิ.}

\prose{นอาสวญฺจ นโนอาสวญฺจ ธมฺมํ ปฏิจฺจ อาสโว ธมฺโม อุปฺปชฺชติ เหตุปจฺจยา. ตีณิ.  เหตุยา นว.}

\proseitem{15}{นสาสวํ ธมฺมํ ปฏิจฺจ สาสโว ธมฺโม อุปฺปชฺชติ เหตุปจฺจยา.  นสาสวํ ธมฺมํ ปฏิจฺจ อนาสโว ธมฺโม อุปฺปชฺชติ เหตุปจฺจยา.  นสาสวํ ธมฺมํ ปฏิจฺจ สาสโว จ อนาสโว จ ธมฺมา อุปฺปชฺชนฺติ เหตุปจฺจยา. ตีณิ.}

\prose{นอนาสวํ ธมฺมํ ปฏิจฺจ สาสโว ธมฺโม อุปฺปชฺชติ เหตุปจฺจยา. เอกํ.}

\prose{นสาสวญฺจ นอนาสวญฺจ ธมฺมํ ปฏิจฺจ สาสโว ธมฺโม อุปฺปชฺชติ เหตุปจฺจยา. เอกํ (สํขิตฺตํ.) . เหตุยา ปญฺจ.}

\proseitem{16}{นอาสว\-สมฺปยุตฺตํ ธมฺมํ ปฏิจฺจ อาสว\-สมฺปยุตฺโต ธมฺโม อุปฺปชฺชติ เหตุปจฺจยา.  เหตุยา นว.}

\proseitem{17}{นอาสวญฺเจว อนาสวญฺจ ธมฺมํ ปฏิจฺจ อาสโว เจว สาสโว จ ธมฺโม อุปฺปชฺชติ เหตุปจฺจยา.  เหตุยา นว.}

\proseitem{18}{นอาสวญฺเจว นอาสว\-วิปฺปยุตฺตญฺจ ธมฺมํ ปฏิจฺจ อาสโว เจว อาสว\-สมฺปยุตฺโต จ ธมฺโม อุปฺปชฺชติ เหตุปจฺจยา.  เหตุยา นว.}

\proseitem{19}{อาสว\-วิปฺปยุตฺตํ นสาสวํ ธมฺมํ ปฏิจฺจ อาสว\-วิปฺปยุตฺโต สาสโว ธมฺโม อุปฺปชฺชติ เหตุปจฺจยา.  เหตุยา ปญฺจ.}

\csromanheader{2}{20-54. สญฺโญชนา\-ทิ\-ฉ\-โคจฺฉก}
\proseitem{20}{นสญฺโญชนํ ธมฺมํ ปฏิจฺจ สญฺโญชโน ธมฺโม อุปฺปชฺชติ เหตุปจฺจยา ฯเปฯ.}

\prose{นคนฺถํ ธมฺมํ ปฏิจฺจ คนฺโถ ธมฺโม อุปฺปชฺชติ เหตุปจฺจยา ฯเปฯ.}

\csromanheader{2}{55-68. มหนฺตรทุก}
\prose{\csromanfolio{341}นโอฆํ ธมฺมํ ปฏิจฺจ โอโฆ ธมฺโม อุปฺปชฺชติ เหตุปจฺจยา ฯเปฯ.}

\prose{นโยคํ ธมฺมํ ปฏิจฺจ โยโค ธมฺโม อุปฺปชฺชติ เหตุปจฺจยา ฯเปฯ.}

\prose{นนีวรณํ ธมฺมํ ปฏิจฺจ นีวรโณ ธมฺโม อุปฺปชฺชติ เหตุปจฺจยา.}

\proseitem{21}{นปรามาสํ ธมฺมํ ปฏิจฺจ ปรามาโส ธมฺโม อุปฺปชฺชติ เหตุปจฺจยา.}

\csromanheader{2}{55-68. มหนฺตรทุก}
\proseitem{22}{นสารมฺมณํ ธมฺมํ ปฏิจฺจ สารมฺมโณ ธมฺโม อุปฺปชฺชติ เหตุปจฺจยา.  นสารมฺมณํ ธมฺมํ ปฏิจฺจ อนารมฺมโณ ธมฺโม อุปฺปชฺชติ เหตุปจฺจยา.  นสารมฺมณํ ธมฺมํ ปฏิจฺจ สารมฺมโณ จ อนารมฺมโณ จ ธมฺมา อุปฺปชฺชนฺติ เหตุปจฺจยา. ตีณิ.}

\prose{นอนารมฺมณํ ธมฺมํ ปฏิจฺจ อนารมฺมโณ ธมฺโม อุปฺปชฺชติ เหตุปจฺจยา. ตีณิ.}

\prose{นสารมฺมณญฺจ นอนารมฺมณญฺจ ธมฺมํ ปฏิจฺจ สารมฺมโณ ธมฺโม อุปฺปชฺชติ เหตุปจฺจยา. ตีณิ.  เหตุยา นว.}

\proseitem{23}{นจิตฺตํ ธมฺมํ ปฏิจฺจ จิตฺโต ธมฺโม อุปฺปชฺชติ เหตุปจฺจยา.  นจิตฺตํ ธมฺมํ ปฏิจฺจ โนจิตฺโต ธมฺโม อุปฺปชฺชติ เหตุปจฺจยา.  นจิตฺตํ ธมฺมํ ปฏิจฺจ จิตฺโต จ โนจิตฺโต จ ธมฺมา อุปฺปชฺชนฺติ เหตุปจฺจยา. ตีณิ.}

\prose{นโนจิตฺตํ ธมฺมํ ปฏิจฺจ โนจิตฺโต ธมฺโม อุปฺปชฺชติ เหตุปจฺจยา. เอกํ.}

\prose{นจิตฺตญฺจ นโนจิตฺตญฺจ ธมฺมํ ปฏิจฺจ โนจิตฺโต ธมฺโม อุปฺปชฺชติ เหตุปจฺจยา. เอกํ.  เหตุยา ปญฺจ.}

\proseitem{24}{นเจตสิกํ ธมฺมํ ปฏิจฺจ เจตสิโก ธมฺโม อุปฺปชฺชติ เหตุปจฺจยา.  นเจตสิกํ ธมฺมํ ปฏิจฺจ อเจตสิโก ธมฺโม อุปฺปชฺชติ เหตุปจฺจยา.  นเจตสิกํ ธมฺมํ ปฏิจฺจ เจตสิโก จ อเจตสิโก จ ธมฺมา อุปฺปชฺชนฺติ เหตุปจฺจยา.  เหตุยา นว.}

\proseitem{25}{นจิตฺต\-สมฺปยุตฺตํ ธมฺมํ ปฏิจฺจ จิตฺต\-สมฺปยุตฺโต ธมฺโม อุปฺปชฺชติ เหตุปจฺจยา.  นจิตฺต\-สมฺปยุตฺตํ ธมฺมํ ปฏิจฺจ จิตฺต\-วิปฺปยุตฺโต ธมฺโม อุปฺปชฺชติ เหตุปจฺจยา.  นจิตฺต\-สมฺปยุตฺตํ ธมฺมํ ปฏิจฺจ จิตฺต\-สมฺปยุตฺโต จ จิตฺต\-วิปฺปยุตฺโต จ ธมฺมา อุปฺปชฺชนฺติ เหตุปจฺจยา. (3)}

\prose{\csromanfolio{342}นจิตฺต\-วิปฺปยุตฺตํ ธมฺมํ ปฏิจฺจ จิตฺต\-วิปฺปยุตฺโต ธมฺโม อุปฺปชฺชติ เหตุปจฺจยา.  นจิตฺต\-วิปฺปยุตฺตํ ธมฺมํ ปฏิจฺจ จิตฺต\-สมฺปยุตฺโต ธมฺโม อุปฺปชฺชติ เหตุปจฺจยา.  นจิตฺต\-วิปฺปยุตฺตํ ธมฺมํ ปฏิจฺจ จิตฺต\-สมฺปยุตฺโต จ จิตฺต\-วิปฺปยุตฺโต จ ธมฺมา อุปฺปชฺชนฺติ เหตุปจฺจยา. (3)}

\prose{นจิตฺต\-สมฺปยุตฺตญฺจ นจิตฺต\-วิปฺปยุตฺตญฺจ ธมฺมํ ปฏิจฺจ จิตฺต\-สมฺปยุตฺโต ธมฺโม อุปฺปชฺชติ เหตุปจฺจยา. ตีณิ.  เหตุยา นว.}

\proseitem{26}{นจิตฺต\-สํสฏฺฐํ ธมฺมํ ปฏิจฺจ จิตฺตสํสฏฺโฐ ธมฺโม อุปฺปชฺชติ เหตุปจฺจยา.  นจิตฺต\-สํสฏฺฐํ ธมฺมํ ปฏิจฺจ จิตฺต\-วิสํสฏฺโฐ ธมฺโม อุปฺปชฺชติ เหตุปจฺจยา.  นจิตฺต\-สํสฏฺฐํ ธมฺมํ ปฏิจฺจ จิตฺตสํสฏฺโฐ จ จิตฺต\-วิสํสฏฺโฐ จ ธมฺมา อุปฺปชฺชนฺติ เหตุปจฺจยา.  เหตุยา นว.}

\proseitem{27}{นจิตฺต\-สมุฏฺฐานํ ธมฺมํ ปฏิจฺจ จิตฺต\-สมุฏฺฐาโน ธมฺโม อุปฺปชฺชติ เหตุปจฺจยา.  นจิตฺต\-สมุฏฺฐานํ ธมฺมํ ปฏิจฺจ โนจิตฺต\-สมุฏฺฐาโน ธมฺโม อุปฺปชฺชติ เหตุปจฺจยา.  นจิตฺต\-สมุฏฺฐานํ ธมฺมํ ปฏิจฺจ จิตฺต\-สมุฏฺฐาโน จ โนจิตฺต\-สมุฏฺฐาโน จ ธมฺมา อุปฺปชฺชนฺติ เหตุปจฺจยา.  เหตุยา นว.}

\proseitem{28}{นจิตฺต\-สหภุํ ธมฺมํ ปฏิจฺจ จิตฺตสหภู ธมฺโม อุปฺปชฺชติ เหตุปจฺจยา.  นจิตฺต\-สหภุํ ธมฺมํ ปฏิจฺจ โนจิตฺตสหภู ธมฺโม อุปฺปชฺชติ เหตุปจฺจยา.  นจิตฺต\-สหภุํ ธมฺมํ ปฏิจฺจ จิตฺตสหภู จ โนจิตฺตสหภู จ ธมฺมา อุปฺปชฺชนฺติ เหตุปจฺจยา.  เหตุยา นว.}

\proseitem{29}{นจิตฺตา\-นุปริวตฺติํ ธมฺมํ ปฏิจฺจ จิตฺตา\-นุปริวตฺตี ธมฺโม อุปฺปชฺชติ เหตุปจฺจยา.  นจิตฺตา\-นุปริวตฺติํ ธมฺมํ ปฏิจฺจ โนจิตฺตา\-นุปริวตฺตี ธมฺโม อุปฺปชฺชติ เหตุปจฺจยา.  นจิตฺตา\-นุปริวตฺติํ ธมฺมํ ปฏิจฺจ จิตฺตา\-นุปริวตฺตี จ โนจิตฺตา\-นุปริวตฺตี จ ธมฺมา อุปฺปชฺชนฺติ เหตุปจฺจยา.  เหตุยา นว.}

\proseitem{30}{นจิตฺต\-สํสฏฺฐ\-สมุฏฺฐานํ ธมฺมํ ปฏิจฺจ จิตฺต\-สํสฏฺฐ\-สมุฏฺฐาโน ธมฺโม อุปฺปชฺชติ เหตุปจฺจยา.  เหตุยา นว.}

\proseitem{31}{นจิตฺต\-สํสฏฺฐ\-สมุฏฺฐาน\-สหภุํ ธมฺมํ ปฏิจฺจ จิตฺต\-สํสฏฺฐ\-สมุฏฺฐาน\-สหภู ธมฺโม อุปฺปชฺชติ เหตุปจฺจยา.  เหตุยา นว.}

\proseitem{32}{นจิตฺต\-สํสฏฺฐ\-สมุฏฺฐานา\-นุ\-ปริวตฺติํ ธมฺมํ ปฏิจฺจ จิตฺต\-สํสฏฺฐ\-สมุฏฺฐานา\-นุปริวตฺตี ธมฺโม อุปฺปชฺชติ เหตุปจฺจยา.  เหตุยา นว.}

\csromanheader{2}{83-99. ปิฏฺฐิทุก}
\proseitem{33}{\csromanfolio{343}นอชฺฌตฺติกํ ธมฺมํ ปฏิจฺจ อชฺฌตฺติโก ธมฺโม อุปฺปชฺชติ เหตุปจฺจยา. ตีณิ.}

\prose{นพาหิรํ ธมฺมํ ปฏิจฺจ พาหิโร ธมฺโม อุปฺปชฺชติ เหตุปจฺจยา. ตีณิ.}

\prose{นอชฺฌตฺติกญฺจ นพาหิรญฺจ ธมฺมํ ปฏิจฺจ อชฺฌตฺติโก ธมฺโม อุปฺปชฺชติ เหตุปจฺจยา. ตีณิ.  เหตุยา นว.}

\proseitem{34}{นอุปาทา ธมฺมํ ปฏิจฺจ อุปาทา ธมฺโม อุปฺปชฺชติ เหตุปจฺจยา.  เหตุยา ปญฺจ.}

\proseitem{35}{นอุปาทินฺนํ ธมฺมํ ปฏิจฺจ อนุปาทินฺโน ธมฺโม อุปฺปชฺชติ เหตุปจฺจยา.  เหตุยา ปญฺจ.}

\csromanheader{2}{69-74. อุปาทาน\-โคจฺฉก}
\proseitem{36}{นอุปาทานํ ธมฺมํ ปฏิจฺจ อุปาทาโน ธมฺโม อุปฺปชฺชติ เหตุปจฺจยา. (สํขิตฺตํ.)}

\csromanheader{2}{75-82. กิเลสโคจฺฉก}
\vspace{\tipitakaparskip}
\csromancenter{นกิเลสํ ธมฺมํ ปฏิจฺจ กิเลโส ธมฺโม อุปฺปชฺชติ เหตุปจฺจยา.}
\csromanheader{2}{83-99. ปิฏฺฐิทุก}
\proseitem{38}{นทสฺสเนน ปหาตพฺพํ ธมฺมํ ปฏิจฺจ นทสฺสเนน ปหาตพฺโพ ธมฺโม อุปฺปชฺชติ เหตุปจฺจยา. เอกํ. นนทสฺสเนน ปหาตพฺพํ ธมฺมํ ปฏิจฺจ ทสฺสเนน ปหาตพฺโพ ธมฺโม อุปฺปชฺชติ เหตุปจฺจยา.  นนทสฺสเนน ปหาตพฺพํ ธมฺมํ ปฏิจฺจ นทสฺสเนน ปหาตพฺโพ ธมฺโม อุปฺปชฺชติ เหตุปจฺจยา.  นนทสฺสเนน ปหาตพฺพํ ธมฺมํ ปฏิจฺจ ทสฺสเนน ปหาตพฺโพ จ นทสฺสเนน ปหาตพฺโพ จ ธมฺมา อุปฺปชฺชนฺติ เหตุปจฺจยา. ตีณิ.}

\prose{นทสฺสเนน ปหาตพฺพญฺจ นนทสฺสเนน ปหาตพฺพญฺจ ธมฺมํ ปฏิจฺจ นทสฺสเนน ปหาตพฺโพ ธมฺโม อุปฺปชฺชติ เหตุปจฺจยา. เอกํ.  เหตุยา ปญฺจ.}

\proseitem{39}{นภาวนาย ปหาตพฺพํ ธมฺมํ ปฏิจฺจ นภาวนาย ปหาตพฺโพ ธมฺโม อุปฺปชฺชติ เหตุปจฺจยา.  เหตุยา ปญฺจ.}

\proseitem{40}{\csromanfolio{344}นทสฺสเนน ปหาตพฺพ\-เหตุกํ ธมฺมํ ปฏิจฺจ ทสฺสเนน ปหาตพฺพ\-เหตุโก ธมฺโม อุปฺปชฺชติ เหตุปจฺจยา.  เหตุยา นว.}

\proseitem{41}{นภาวนาย ปหาตพฺพ\-เหตุกํ ธมฺมํ ปฏิจฺจ ภาวนาย ปหาตพฺพ\-เหตุโก ธมฺโม อุปฺปชฺชติ เหตุปจฺจยา.  เหตุยา นว.}

\proseitem{42}{นสวิตกฺกํ ธมฺมํ ปฏิจฺจ สวิตกฺโก ธมฺโม อุปฺปชฺชติ เหตุปจฺจยา.  นสวิตกฺกํ ธมฺมํ ปฏิจฺจ อวิตกฺโก ธมฺโม อุปฺปชฺชติ เหตุปจฺจยา.  นสวิตกฺกํ ธมฺมํ ปฏิจฺจ สวิตกฺโก จ อวิตกฺโก จ ธมฺมา อุปฺปชฺชนฺติ เหตุปจฺจยา. ตีณิ.}

\prose{นอวิตกฺกํ ธมฺมํ ปฏิจฺจ อวิตกฺโก ธมฺโม อุปฺปชฺชติ เหตุปจฺจยา.  นอวิตกฺกํ ธมฺมํ ปฏิจฺจ สวิตกฺโก ธมฺโม อุปฺปชฺชติ เหตุปจฺจยา.  นอวิตกฺกํ ธมฺมํ ปฏิจฺจ สวิตกฺโก จ อวิตกฺโก จ ธมฺมา อุปฺปชฺชนฺติ เหตุปจฺจยา. (3)}

\prose{นสวิตกฺกญฺจ นอวิตกฺกญฺจ ธมฺมํ ปฏิจฺจ สวิตกฺโก ธมฺโม อุปฺปชฺชติ เหตุปจฺจยา.  นสวิตกฺกญฺจ นอวิตกฺกญฺจ ธมฺมํ ปฏิจฺจ อวิตกฺโก ธมฺโม อุปฺปชฺชติ เหตุปจฺจยา.  นสวิตกฺกญฺจ นอวิตกฺกญฺจ ธมฺมํ ปฏิจฺจ สวิตกฺโก จ อวิตกฺโก จ ธมฺมา อุปฺปชฺชนฺติ เหตุปจฺจยา. (3). เหตุยา นว.}

\proseitem{43}{นสวิจารํ ธมฺมํ ปฏิจฺจ สวิจาโร ธมฺโม อุปฺปชฺชติ เหตุปจฺจยา.  เหตุยา นว.}

\proseitem{44}{นสปฺปีติกํ ธมฺมํ ปฏิจฺจ สปฺปีติโก ธมฺโม อุปฺปชฺชติ เหตุปจฺจยา.  เหตุยา นว.}

\proseitem{45}{นปีติสหคตํ ธมฺมํ ปฏิจฺจ ปีติสหคโต ธมฺโม อุปฺปชฺชติ เหตุปจฺจยา. เหตุยา นว.}

\proseitem{46}{นสุขสหคตํ ธมฺมํ ปฏิจฺจ สุขสหคโต ธมฺโม อุปฺปชฺชติ เหตุปจฺจยา.  เหตุยา นว.}

\proseitem{47}{นอุเปกฺขาสหคตํ ธมฺมํ ปฏิจฺจ อุเปกฺขา\-สหคโต ธมฺโม อุปฺปชฺชติ เหตุปจฺจยา.  นอุเปกฺขาสหคตํ ธมฺมํ ปฏิจฺจ นอุเปกฺขา\-สหคโต ธมฺโม อุปฺปชฺชติ เหตุปจฺจยา.  นอุเปกฺขาสหคตํ ธมฺมํ ปฏิจฺจ อุเปกฺขา\-สหคโต จ นอุเปกฺขา\-สหคโต จ ธมฺมา อุปฺปชฺชนฺติ เหตุปจฺจยา. ตีณิ.}

\csromanheader{2}{83-99. ปิฏฺฐิทุก}
\prose{\csromanfolio{345}นนอุเปกฺขาสหคตํ ธมฺมํ ปฏิจฺจ นอุเปกฺขาสหคโต ธมฺโม อุปฺปชฺชติ เหตุปจฺจยา. ตีณิ.}

\prose{นอุเปกฺขาสหคตญฺจ นนอุเปกฺขาสหคตญฺจ ธมฺมํ ปฏิจฺจ อุเปกฺขา\-สหคโต ธมฺโม อุปฺปชฺชติ เหตุปจฺจยา.  นอุเปกฺขาสหคตญฺจ นนอุเปกฺขาสหคตญฺจ ธมฺมํ ปฏิจฺจ นอุเปกฺขา\-สหคโต ธมฺโม อุปฺปชฺชติ เหตุปจฺจยา.  นอุเปกฺขาสหคตญฺจ นนอุเปกฺขาสหคตญฺจ ธมฺมํ ปฏิจฺจ อุเปกฺขา\-สหคโต จ นอุเปกฺขา\-สหคโต จ ธมฺมา อุปฺปชฺชนฺติ เหตุปจฺจยา. ตีณิ.  เหตุยา นว.}

\proseitem{48}{นกามาวจรํ ธมฺมํ ปฏิจฺจ กามาวจโร ธมฺโม อุปฺปชฺชติ เหตุปจฺจยา.  นกามาวจรํ ธมฺมํ ปฏิจฺจ นกามาวจโร ธมฺโม อุปฺปชฺชติ เหตุปจฺจยา.  นกามาวจรํ ธมฺมํ ปฏิจฺจ กามาวจโร จ นกามาวจโร จ ธมฺมา อุปฺปชฺชนฺติ เหตุปจฺจยา. ตีณิ.}

\prose{นนกามา\-วจรํ ธมฺมํ ปฏิจฺจ นกามาวจโร ธมฺโม อุปฺปชฺชติ เหตุปจฺจยา. ตีณิ.}

\prose{นกามาวจรญฺจ นนกามาวจรญฺจ ธมฺมํ ปฏิจฺจ กามาวจโร ธมฺโม อุปฺปชฺชติ เหตุปจฺจยา. ตีณิ.  เหตุยา นว.}

\proseitem{49}{นรูปาวจรํ ธมฺมํ ปฏิจฺจ รูปาวจโร ธมฺโม อุปฺปชฺชติ เหตุปจฺจยา.  เหตุยา นว.}

\proseitem{50}{นอรูปาวจรํ ธมฺมํ ปฏิจฺจ นอรูปาวจโร ธมฺโม อุปฺปชฺชติ เหตุปจฺจยา.  เหตุยา ปญฺจ.}

\proseitem{51}{นปริยาปนฺนํ ธมฺมํ ปฏิจฺจ ปริยาปนฺโน ธมฺโม อุปฺปชฺชติ เหตุปจฺจยา. ตีณิ.}

\prose{นอปริยาปนฺนํ ธมฺมํ ปฏิจฺจ ปริยาปนฺโน ธมฺโม อุปฺปชฺชติ เหตุปจฺจยา. เอกํ.}

\prose{นปริยาปนฺนญฺจ นอปริยาปนฺนญฺจ ธมฺมํ ปฏิจฺจ ปริยาปนฺโน ธมฺโม อุปฺปชฺชติ เหตุปจฺจยา. เอกํ.  เหตุยา ปญฺจ.}

\proseitem{52}{นนิยฺยานิกํ ธมฺมํ ปฏิจฺจ อนิยฺยานิโก ธมฺโม อุปฺปชฺชติ เหตุปจฺจยา. เอกํ.}

\prose{\csromanfolio{346}นอนิยฺยานิกํ ธมฺมํ ปฏิจฺจ อนิยฺยานิโก ธมฺโม อุปฺปชฺชติ เหตุปจฺจยา. ตีณิ.}

\prose{นนิยฺยานิกญฺจ นอนิยฺยานิกญฺจ ธมฺมํ ปฏิจฺจ อนิยฺยานิโก ธมฺโม อุปฺปชฺชติ เหตุปจฺจยา. เอกํ.  เหตุยา ปญฺจ.}

\proseitem{53}{นนิยตํ ธมฺมํ ปฏิจฺจ อนิยโต ธมฺโม อุปฺปชฺชติ เหตุปจฺจยา. เอกํ.}

\prose{นอนิยตํ ธมฺมํ ปฏิจฺจ อนิยโต ธมฺโม อุปฺปชฺชติ เหตุปจฺจยา. ตีณิ.}

\prose{นนิยตญฺจ นอนิยตญฺจ ธมฺมํ ปฏิจฺจ อนิยโต ธมฺโม อุปฺปชฺชติ เหตุปจฺจยา. เอกํ.  เหตุยา ปญฺจ.}

\proseitem{54}{นสอุตฺตรํ ธมฺมํ ปฏิจฺจ สอุตฺตโร ธมฺโม อุปฺปชฺชติ เหตุปจฺจยา.  นสอุตฺตรํ ธมฺมํ ปฏิจฺจ อนุตฺตโร ธมฺโม อุปฺปชฺชติ เหตุปจฺจยา.  นสอุตฺตรํ ธมฺมํ ปฏิจฺจ สอุตฺตโร จ อนุตฺตโร จ ธมฺมา อุปฺปชฺชนฺติ เหตุปจฺจยา. ตีณิ.}

\prose{นอนุตฺตรํ ธมฺมํ ปฏิจฺจ สอุตฺตโร ธมฺโม อุปฺปชฺชติ เหตุปจฺจยา. เอกํ.}

\prose{นสอุตฺตรญฺจ นอนุตฺตรญฺจ ธมฺมํ ปฏิจฺจ สอุตฺตโร ธมฺโม อุปฺปชฺชติ เหตุปจฺจยา.  เหตุยา ปญฺจ.}

\csromanheader{2}{100. \textbf{สรณทุก 1-7. ปฏิจฺจาทิวาร}}
\proseitem{55}{นสรณํ ธมฺมํ ปฏิจฺจ อรโณ ธมฺโม อุปฺปชฺชติ เหตุปจฺจยา. เอกํ.}

\prose{นอรณํ ธมฺมํ ปฏิจฺจ อรโณ ธมฺโม อุปฺปชฺชติ เหตุปจฺจยา. ตีณิ.}

\prose{นสรณญฺจ นอรณญฺจ ธมฺมํ ปฏิจฺจ อรโณ ธมฺโม อุปฺปชฺชติ เหตุปจฺจยา. เอกํ.  เหตุยา ปญฺจ, อารมฺมเณ เทฺว ฯเปฯ อวิคเต ปญฺจ.}

\prose{ปจฺจนีย\-น\-เหตุ}

\proseitem{56}{นสรณํ ธมฺมํ ปฏิจฺจ อรโณ ธมฺโม อุปฺปชฺชติ นเหตุปจฺจยา.  นอรณํ ธมฺมํ ปฏิจฺจ สรโณ ธมฺโม อุปฺปชฺชติ นเหตุปจฺจยา. (สํขิตฺตํ.)}

\prose{นเหตุยา เทฺว, นอารมฺมเณ ตีณิ ฯเปฯ โนวิคเต ตีณิ.}

\prose{(สหชาตวารมฺปิ ฯเปฯ สมฺปยุตฺตวารมฺปิ ปฏิจฺจ\-วาร\-สทิสํ.)}

\csromanheader{2}{100. สรณทุก \textbf{เหตุ อารมฺมณ}}
\proseitem{57}{\csromanfolio{347}นสรโณ ธมฺโม อรณสฺส ธมฺมสฺส เหตุปจฺจเยน ปจฺจโย. (1)}

\prose{นอรโณ ธมฺโม อรณสฺส ธมฺมสฺส เหตุปจฺจเยน ปจฺจโย.  นอรโณ ธมฺโม สรณสฺส ธมฺมสฺส เหตุปจฺจเยน ปจฺจโย.  นอรโณ ธมฺโม สรณสฺส จ อรณสฺส จ ธมฺมสฺส เหตุปจฺจเยน ปจฺจโย. (3)}

\proseitem{58}{นสรโณ ธมฺโม สรณสฺส ธมฺมสฺส อารมฺมณ\-ปจฺจเยน ปจฺจโย.  นสรโณ ธมฺโม อรณสฺส ธมฺมสฺส อารมฺมณ\-ปจฺจเยน ปจฺจโย. เทฺว.}

\prose{นอรโณ ธมฺโม อรณสฺส ธมฺมสฺส อารมฺมณ\-ปจฺจเยน ปจฺจโย.  นอรโณ ธมฺโม สรณสฺส ธมฺมสฺส อารมฺมณ\-ปจฺจเยน ปจฺจโย. เทฺว.}

\proseitem{59}{เหตุยา จตฺตาริ, อารมฺมเณ จตฺตาริ ฯเปฯ อวิคเต สตฺต.}

\csromanheader{2}{ปจฺจนียุ\-ทฺธาร}
\proseitem{60}{นสรโณ ธมฺโม สรณสฺส ธมฺมสฺส อารมฺมณ\-ปจฺจเยน ปจฺจโย. อุปนิสฺสย\-ปจฺจเยน ปจฺจโย. ปุเรชาต\-ปจฺจเยน ปจฺจโย.}

\prose{นเหตุยา สตฺต, นอารมฺมเณ สตฺต ฯเปฯ โนอวิคเต จตฺตาริ.}

\prosewithcloser{(ยถา กุสลตฺติเก ปญฺหาวารสฺส อนุโลมมฺปิ ปจฺจนียมฺปิ อนุโลมปจฺจนียมฺปิ ปจฺจนียานุโลมมฺปิ คณิตํ, เอวํ คเณตพฺพํ.)}{ธมฺม\-ปจฺจนียา\-นุโลเม ทุกปฏฺฐานํ นิฏฺฐิตํ.}
\pitaka{อภิธมฺม\-ปิฏก}
\csromanheader{2}{ธมฺม\-ปจฺจนียา\-นุโลม}
\namakkaram{นโม ตสฺส ภควโต อรหโต สมฺมา\-สมฺพุทฺธสฺส.}
\csromanmarkh{1. เหตุทุก}\csromanmarkhii{1. กุสลตฺติก}\csromanheadernomark{2}{1. \textbf{เหตุทุก 1. กุสลตฺติก}}
\proseitem{1}{\csromanfolio{349}นเหตุํ นกุสลํ ธมฺมํ ปจฺจยา เหตุ กุสโล ธมฺโม อุปฺปชฺชติ เหตุปจฺจยา.  นเหตุํ นกุสลํ ธมฺมํ ปจฺจยา นเหตุ กุสโล ธมฺโม อุปฺปชฺชติ เหตุปจฺจยา.  นเหตุํ นกุสลํ ธมฺมํ ปจฺจยา เหตุ กุสโล จ นเหตุ กุสโล จ ธมฺมา อุปฺปชฺชนฺติ เหตุปจฺจยา.  เหตุยา ตีณิ.}

\prose{นเหตุํ นอกุสลํ ธมฺมํ ปจฺจยา เหตุ อกุสโล ธมฺโม อุปฺปชฺชติ เหตุปจฺจยา.  เหตุยา ตีณิ.}

\proseitem{2}{นเหตุํ นอพฺยากตํ ธมฺมํ ปฏิจฺจ นเหตุ อพฺยากโต ธมฺโม อุปฺปชฺชติ เหตุปจฺจยา.  นนเหตุํ นอพฺยากตํ ธมฺมํ ปฏิจฺจ นเหตุ อพฺยากโต ธมฺโม อุปฺปชฺชติ เหตุปจฺจยา.  นเหตุํ นอพฺยากตญฺจ นนเหตุํ นอพฺยากตญฺจ ธมฺมํ ปฏิจฺจ นเหตุ อพฺยากโต ธมฺโม อุปฺปชฺชติ เหตุปจฺจยา.  เหตุยา ตีณิ.}

\csromanmarkh{2-4. สเหตุกทุกาทิ}\csromanmarkhii{1. กุสลตฺติก}\csromanheadernomark{2}{2-4. สเหตุกทุกาทิ 1. กุสลตฺติก}
\proseitem{3}{นสเหตุกํ นกุสลํ ธมฺมํ ปจฺจยา สเหตุโก กุสโล ธมฺโม อุปฺปชฺชติ เหตุปจฺจยา.  เหตุยา เอกํ.}

\prose{นสเหตุกํ นอกุสลํ ธมฺมํ ปจฺจยา สเหตุโก อกุสโล ธมฺโม อุปฺปชฺชติ เหตุปจฺจยา.  เหตุยา เอกํ.}

\csromanheader{2}{ธมฺม\-ปจฺจนียา\-นุโลม ทุกติก\-ปฏฺฐาน\-ปาฬิ}
\proseitem{4}{\csromanfolio{350}นสเหตุกํ นอพฺยากตํ ธมฺมํ ปฏิจฺจ อเหตุโก อพฺยากโต ธมฺโม อุปฺปชฺชติ เหตุปจฺจยา.  นอเหตุกํ นอพฺยากตํ ธมฺมํ ปฏิจฺจ อเหตุโก อพฺยากโต ธมฺโม อุปฺปชฺชติ เหตุปจฺจยา.  นสเหตุกํ นอพฺยากตญฺจ นอเหตุกํ นอพฺยากตญฺจ ธมฺมํ ปฏิจฺจ อเหตุโก อพฺยากโต ธมฺโม อุปฺปชฺชติ เหตุปจฺจยา.  เหตุยา ตีณิ.}

\proseitem{5}{นเหตุ\-สมฺปยุตฺตํ นกุสลํ ธมฺมํ ปจฺจยา.}

\proseitem{6}{นเหตุ เจว นอเหตุโก จ นกุสโล ธมฺโม เหตุสฺส เจว สเหตุกสฺส จ กุสลสฺส ธมฺมสฺส อารมฺมณ\-ปจฺจเยน ปจฺจโย.  นเหตุ เจว นอเหตุโก จ นกุสโล ธมฺโม สเหตุกสฺส เจว น จ เหตุสฺส กุสลสฺส ธมฺมสฺส อารมฺมณ\-ปจฺจเยน ปจฺจโย.  นเหตุ เจว นอเหตุโก จ นกุสโล ธมฺโม เหตุสฺส เจว สเหตุกสฺส กุสลสฺส จ สเหตุกสฺส เจว น จ เหตุสฺส กุสลสฺส จ ธมฺมสฺส อารมฺมณ\-ปจฺจเยน ปจฺจโย. ตีณิ.}

\prose{นอเหตุโก เจว นน จ เหตุ นกุสโล ธมฺโม เหตุสฺส เจว สเหตุกสฺส จ กุสลสฺส ธมฺมสฺส อารมฺมณ\-ปจฺจเยน ปจฺจโย. ตีณิ.}

\prose{นเหตุ เจว นอเหตุโก นกุสโล จ นอเหตุโก เจว นนเหตุ นกุสโล จ ธมฺมา เหตุสฺส เจว สเหตุกสฺส จ กุสลสฺส ธมฺมสฺส อารมฺมณ\-ปจฺจเยน ปจฺจโย. ตีณิ. (สพฺพตฺถ นว ปญฺหา.)}

\proseitem{7}{นเหตุ เจว นอเหตุโก จ นอกุสโล ธมฺโม เหตุสฺส เจว สเหตุกสฺส จ อกุสลสฺส ธมฺมสฺส อารมฺมณ\-ปจฺจเยน ปจฺจโย. (สํขิตฺตํ. นว ปญฺหา.)}

\prose{นเหตุ เจว นอเหตุโก จ นอพฺยากโต ธมฺโม เหตุสฺส เจว สเหตุกสฺส จ อพฺยากตสฺส ธมฺมสฺส อารมฺมณ\-ปจฺจเยน ปจฺจโย. ตีณิ. (นว ปญฺหา.)}

\csromanmarkh{5. เหตุ\-เหตุ\-สมฺปยุตฺตทุก}\csromanmarkhii{1. กุสลตฺติก}\csromanheadernomark{2}{5. เหตุ\-เหตุ\-สมฺปยุตฺตทุก 1. กุสลตฺติก}
\proseitem{8}{นเหตุ เจว นเหตุ\-วิปฺปยุตฺโต จ นกุสโล ธมฺโม เหตุสฺส เจว เหตุ\-สมฺปยุตฺตสฺส จ กุสลสฺส ธมฺมสฺส อารมฺมณ\-ปจฺจเยน ปจฺจโย. (สํขิตฺตํ. นว ปญฺหา กาตพฺพา.)}

\csromanmarkh{7-13. จูฬนฺตรทุก}\csromanmarkhii{1. กุสลตฺติก}\csromanheadernomark{2}{7-13. จูฬนฺตรทุก 1. กุสลตฺติก}
\prose{\csromanfolio{351}นเหตุ เจว นเหตุ\-วิปฺปยุตฺโต จ นอกุสโล ธมฺโม เหตุสฺส เจว เหตุ\-สมฺปยุตฺตสฺส จ อกุสลสฺส ธมฺมสฺส อารมฺมณ\-ปจฺจเยน ปจฺจโย. (สํขิตฺตํ. นว ปญฺหา กาตพฺพา.)}

\prose{นเหตุ เจว นเหตุ\-วิปฺปยุตฺโต จ นอพฺยากโต ธมฺโม เหตุสฺส เจว เหตุ\-สมฺปยุตฺตสฺส จ อพฺยากตสฺส ธมฺมสฺส อารมฺมณ\-ปจฺจเยน ปจฺจโย. ตีณิ. (นว ปญฺหา กาตพฺพา.)}

\csromanmarkh{6. นเหตุ\-สเหตุกทุก}\csromanmarkhii{1. กุสลตฺติก}\csromanheadernomark{2}{6. นเหตุ\-สเหตุกทุก 1. กุสลตฺติก}
\proseitem{9}{นเหตุํ นสเหตุกํ นกุสลํ ธมฺมํ ปจฺจยา นเหตุ สเหตุโก กุสโล ธมฺโม อุปฺปชฺชติ เหตุปจฺจยา.  เหตุยา เอกํ.}

\prose{นเหตุํ นสเหตุกํ นอกุสลํ ธมฺมํ ปจฺจยา นเหตุ สเหตุโก อกุสโล ธมฺโม อุปฺปชฺชติ เหตุปจฺจยา.  เหตุยา เอกํ.}

\prose{นเหตุํ นอเหตุกํ นอพฺยากตํ ธมฺมํ ปฏิจฺจ นเหตุ อเหตุโก อพฺยากโต ธมฺโม อุปฺปชฺชติ เหตุปจฺจยา.  เหตุยา เอกํ.}

\csromanmarkh{7-13. จูฬนฺตรทุก}\csromanmarkhii{1. กุสลตฺติก}\csromanheadernomark{2}{7-13. จูฬนฺตรทุก 1. กุสลตฺติก}
\proseitem{10}{นอปฺปจฺจยํ นกุสลํ ธมฺมํ ปจฺจยา สปฺปจฺจโย กุสโล ธมฺโม อุปฺปชฺชติ เหตุปจฺจยา.  เหตุยา เอกํ.}

\prose{นอปฺปจฺจยํ นอกุสลํ ธมฺมํ ปจฺจยา สปฺปจฺจโย อกุสโล ธมฺโม อุปฺปชฺชติ เหตุปจฺจยา.  เหตุยา เอกํ.}

\prose{นอปฺปจฺจยํ นอพฺยากตํ ธมฺมํ ปฏิจฺจ สปฺปจฺจโย อพฺยากโต ธมฺโม อุปฺปชฺชติ เหตุปจฺจยา.  เหตุยา เอกํ.}

\prose{11. นอสงฺขตํ ฯเปฯ. (สปฺปจฺจย\-ทุก\-สทิสํ.)}

\proseitem{12}{นสนิทสฺสนํ นกุสลํ ธมฺมํ ปจฺจยา อนิทสฺสโน กุสโล ธมฺโม อุปฺปชฺชติ เหตุปจฺจยา.  เหตุยา เอกํ.}

\prose{นสนิทสฺสนํ นอกุสลํ ธมฺมํ ปจฺจยา อนิทสฺสโน อกุสโล ธมฺโม อุปฺปชฺชติ เหตุปจฺจยา.  เหตุยา เอกํ.}

\csromanheader{2}{ธมฺม\-ปจฺจนียา\-นุโลม ทุกติก\-ปฏฺฐาน\-ปาฬิ}
\prose{\csromanfolio{352}นสนิทสฺสนํ นอพฺยากตํ ธมฺมํ ปฏิจฺจ สนิทสฺสโน อพฺยากโต ธมฺโม อุปฺปชฺชติ เหตุปจฺจยา.  นสนิทสฺสนํ นอพฺยากตํ ธมฺมํ ปฏิจฺจ อนิทสฺสโน อพฺยากโต ธมฺโม อุปฺปชฺชติ เหตุปจฺจยา.  นสนิทสฺสนํ นอพฺยากตํ ธมฺมํ ปฏิจฺจ สนิทสฺสโน อพฺยากโต จ อนิทสฺสโน อพฺยากโต จ ธมฺมา อุปฺปชฺชนฺติ เหตุปจฺจยา. (3). เหตุยา ตีณิ.}

\proseitem{13}{นสปฺปฏิฆํ นกุสลํ ธมฺมํ ปจฺจยา อปฺปฏิโฆ กุสโล ธมฺโม อุปฺปชฺชติ เหตุปจฺจยา.}

\proseitem{14}{นอรูปิํ นกุสลํ ธมฺมํ ปจฺจยา อรูปี กุสโล ธมฺโม อุปฺปชฺชติ เหตุปจฺจยา.  เหตุยา เอกํ.}

\prose{นอรูปิํ นอกุสลํ ธมฺมํ ปจฺจยา อรูปี อกุสโล ธมฺโม อุปฺปชฺชติ เหตุปจฺจยา.  เหตุยา เอกํ.}

\prose{นรูปิํ นอพฺยากตํ ธมฺมํ ปฏิจฺจ รูปี อพฺยากโต ธมฺโม อุปฺปชฺชติ เหตุปจฺจยา.  เหตุยา เอกํ.}

\proseitem{15}{นโลกุตฺตรํ นกุสลํ ธมฺมํ ปจฺจยา โลกิโย กุสโล ธมฺโม อุปฺปชฺชติ เหตุปจฺจยา.  นโลกุตฺตรํ นกุสลํ ธมฺมํ ปจฺจยา โลกุตฺตโร กุสโล ธมฺโม อุปฺปชฺชติ เหตุปจฺจยา.  เหตุยา เทฺว.}

\prose{นโลกุตฺตรํ นอกุสลํ ธมฺมํ ปจฺจยา โลกิโย อกุสโล ธมฺโม อุปฺปชฺชติ เหตุปจฺจยา.  เหตุยา เอกํ.}

\prose{นโลกิยํ นอพฺยากตํ ธมฺมํ ปฏิจฺจ โลกิโย อพฺยากโต ธมฺโม อุปฺปชฺชติ เหตุปจฺจยา.  นโลกุตฺตรํ นอพฺยากตํ ธมฺมํ ปฏิจฺจ โลกิโย อพฺยากโต ธมฺโม อุปฺปชฺชติ เหตุปจฺจยา.  เหตุยา เทฺว.}

\proseitem{16}{นเกนจิ วิญฺเญยฺยํ นกุสลํ ธมฺมํ ปจฺจยา เกนจิ วิญฺเญยฺโย กุสโล ธมฺโม อุปฺปชฺชติ เหตุปจฺจยา. ตีณิ.}

\prose{นเกนจิ นวิญฺเญยฺยํ นกุสลํ ธมฺมํ ปจฺจยา เกนจิ นวิญฺเญยฺโย กุสโล ธมฺโม อุปฺปชฺชติ เหตุปจฺจยา. ตีณิ.}

\csromanmarkh{16. อาสว\-สมฺปยุตฺตทุก}\csromanmarkhii{1. กุสลตฺติก}\csromanheadernomark{2}{16. อาสว\-สมฺปยุตฺตทุก 1. กุสลตฺติก}
\prose{\csromanfolio{353}นเกนจิ วิญฺเญยฺยํ นกุสลญฺจ นเกนจิ นวิญฺเญยฺยํ นกุสลญฺจ ธมฺมํ ปจฺจยา เกนจิ วิญฺเญยฺโย กุสโล ธมฺโม อุปฺปชฺชติ เหตุปจฺจยา. ตีณิ. (สพฺพตฺถ นว.)}

\prose{นเกนจิ วิญฺเญยฺยํ นอกุสลํ ธมฺมํ ปจฺจยา เกนจิ วิญฺเญยฺโย อกุสโล ธมฺโม อุปฺปชฺชติ เหตุปจฺจยา.  นเกนจิ นวิญฺเญยฺยํ นอกุสลํ ธมฺมํ ปจฺจยา เกนจิ นวิญฺเญยฺโย อกุสโล ธมฺโม อุปฺปชฺชติ เหตุปจฺจยา.  เหตุยา นว.}

\prose{นเกนจิ วิญฺเญยฺยํ นอพฺยากตํ ธมฺมํ ปฏิจฺจ เกนจิ วิญฺเญยฺโย อพฺยากโต ธมฺโม อุปฺปชฺชติ เหตุปจฺจยา.  เหตุยา นว.}

\csromanmarkh{14. อาสวทุก}\csromanmarkhii{1. กุสลตฺติก}\csromanheadernomark{2}{14. \textbf{อาสวทุก 1. กุสลตฺติก}}
\proseitem{17}{นอาสวํ นกุสลํ ธมฺมํ ปจฺจยา โนอาสโว กุสโล ธมฺโม อุปฺปชฺชติ เหตุปจฺจยา.  เหตุยา เอกํ.}

\prose{นอาสวํ นอกุสลํ ธมฺมํ ปจฺจยา อาสโว อกุสโล ธมฺโม อุปฺปชฺชติ เหตุปจฺจยา.  เหตุยา ตีณิ.}

\prose{นอาสวํ นอพฺยากตํ ธมฺมํ ปฏิจฺจ โนอาสโว อพฺยากโต ธมฺโม อุปฺปชฺชติ เหตุปจฺจยา.  นโนอาสวํ นอพฺยากตํ ธมฺมํ ปฏิจฺจ โนอาสโว อพฺยากโต ธมฺโม อุปฺปชฺชติ เหตุปจฺจยา. (ทุกมูลกํ เอกํ.) . เหตุยา ตีณิ.}

\csromanmarkh{15. สาสวทุก}\csromanmarkhii{1. กุสลตฺติก}\csromanheadernomark{2}{15. \textbf{สาสวทุก 1. กุสลตฺติก}}
\proseitem{18}{นอนาสวํ นกุสลํ ธมฺมํ ปจฺจยา อนาสโว กุสโล ธมฺโม อุปฺปชฺชติ เหตุปจฺจยา.  นอนาสวํ นกุสลํ ธมฺมํ ปจฺจยา สาสโว กุสโล ธมฺโม อุปฺปชฺชติ เหตุปจฺจยา.  เหตุยา เทฺว.}

\prose{นอนาสวํ นอกุสลํ ธมฺมํ ปจฺจยา สาสโว อกุสโล ธมฺโม อุปฺปชฺชติ เหตุปจฺจยา.  เหตุยา เอกํ.}

\csromanmarkh{16. อาสว\-สมฺปยุตฺตทุก}\csromanmarkhii{1. กุสลตฺติก}\csromanheadernomark{2}{16. อาสว\-สมฺปยุตฺตทุก 1. กุสลตฺติก}
\proseitem{19}{นอาสว\-สมฺปยุตฺตํ นกุสลํ ธมฺมํ ปจฺจยา อาสว\-วิปฺปยุตฺโต กุสโล ธมฺโม อุปฺปชฺชติ เหตุปจฺจยา.  เหตุยา เอกํ.}

\csromanheader{2}{ธมฺม\-ปจฺจนียา\-นุโลม ทุกติก\-ปฏฺฐาน\-ปาฬิ}
\prose{\csromanfolio{354}นอาสว\-สมฺปยุตฺตํ นอกุสลํ ธมฺมํ ปจฺจยา อาสว\-สมฺปยุตฺโต อกุสโล ธมฺโม อุปฺปชฺชติ เหตุปจฺจยา.  เหตุยา ตีณิ.}

\prose{นอาสว\-สมฺปยุตฺตํ นอพฺยากตํ ธมฺมํ ปฏิจฺจ อาสว\-วิปฺปยุตฺโต อพฺยากโต ธมฺโม อุปฺปชฺชติ เหตุปจฺจยา.  เหตุยา ตีณิ.}

\csromanmarkh{17. อาสว\-สาสวทุก}\csromanmarkhii{1. กุสลตฺติก}\csromanheadernomark{2}{17. อาสว\-สาสวทุก 1. กุสลตฺติก}
\proseitem{20}{นอาสวญฺเจว นอนาสวญฺจ นกุสลํ ธมฺมํ ปจฺจยา สาสโว เจว โนอาสโว จ กุสโล ธมฺโม อุปฺปชฺชติ เหตุปจฺจยา.  เหตุยา เอกํ.}

\prose{นอาสวญฺเจว นอนาสวญฺจ นอกุสลํ ธมฺมํ ปจฺจยา อาสโว เจว สาสโว จ อกุสโล ธมฺโม อุปฺปชฺชติ เหตุปจฺจยา.  เหตุยา ตีณิ.}

\prose{นอาสวญฺเจว นอนาสวญฺจ นอพฺยากตํ ธมฺมํ ปฏิจฺจ สาสโว เจว โน จ อาสโว อพฺยากโต ธมฺโม อุปฺปชฺชติ เหตุปจฺจยา.  เหตุยา ตีณิ. (อาสว\-อาสว\-สมฺปยุตฺตทุกํ นตฺถิ.)}

\csromanmarkh{19. อาสว\-วิปฺปยุตฺต\-สาสวทุก}\csromanmarkhii{1. กุสลตฺติก}\csromanheadernomark{2}{19. อาสว\-วิปฺปยุตฺต\-สาสวทุก 1. กุสลตฺติก}
\proseitem{21}{อาสว\-วิปฺปยุตฺตํ นอนาสวํ นกุสลํ ธมฺมํ ปจฺจยา อาสว\-วิปฺปยุตฺโต อนาสโว กุสโล ธมฺโม อุปฺปชฺชติ เหตุปจฺจยา.  อาสว\-วิปฺปยุตฺตํ นอนาสวํ นกุสลํ ธมฺมํ ปจฺจยา อาสว\-วิปฺปยุตฺโต สาสโว กุสโล ธมฺโม อุปฺปชฺชติ เหตุปจฺจยา.  เหตุยา เทฺว.}

\prose{อาสว\-วิปฺปยุตฺตํ นอนาสวํ นอกุสลํ ธมฺมํ ปจฺจยา อาสว\-วิปฺปยุตฺโต สาสโว อกุสโล ธมฺโม อุปฺปชฺชติ เหตุปจฺจยา.  เหตุยา เอกํ.}

\prose{อาสว\-วิปฺปยุตฺตํ นสาสวํ นอพฺยากตํ ธมฺมํ ปฏิจฺจ อาสว\-วิปฺปยุตฺโต สาสโว อพฺยากโต ธมฺโม อุปฺปชฺชติ เหตุปจฺจยา.  อาสว\-วิปฺปยุตฺตํ นอนาสวํ นอพฺยากตํ ธมฺมํ ปฏิจฺจ อาสว\-วิปฺปยุตฺโต สาสโว อพฺยากโต ธมฺโม อุปฺปชฺชติ เหตุปจฺจยา.  เหตุยา เทฺว.}

\vspace{\tipitakaparskip}
\csromancenter{มหนฺตรทุก 1. กุสลตฺติก 20-54. สญฺโญชนา\-ทิ\-ฉ\-โคจฺฉก 1. กุสลตฺติก}
\proseitem{22}{\csromanfolio{355}นสญฺโญชนํ นอกุสลํ ธมฺมํ ปจฺจยา สญฺโญชโน อกุสโล ธมฺโม อุปฺปชฺชติ เหตุปจฺจยา.  เหตุยา ตีณิ.}

\proseitem{23}{นคนฺถํ นอกุสลํ ธมฺมํ ปจฺจยา คนฺโถ อกุสโล ธมฺโม อุปฺปชฺชติ เหตุปจฺจยา.  เหตุยา ตีณิ.}

\proseitem{24}{นโอฆํ นอกุสลํ ธมฺมํ ปจฺจยา โอโฆ อกุสโล ธมฺโม อุปฺปชฺชติ เหตุปจฺจยา.  เหตุยา ตีณิ.}

\proseitem{25}{นโยคํ นอกุสลํ ธมฺมํ ปจฺจยา โยโค อกุสโล ธมฺโม อุปฺปชฺชติ เหตุปจฺจยา.  เหตุยา ตีณิ.}

\proseitem{26}{นนีวรณํ นอกุสลํ ธมฺมํ ปจฺจยา นีวรโณ อกุสโล ธมฺโม อุปฺปชฺชติ เหตุปจฺจยา.  เหตุยา ตีณิ.}

\proseitem{27}{นปรามาสํ นอกุสลํ ธมฺมํ ปจฺจยา ปรามาโส อกุสโล ธมฺโม อุปฺปชฺชติ เหตุปจฺจยา.  เหตุยา ตีณิ.}

\csromanmarkh{55-68. มหนฺตรทุก}\csromanmarkhii{1. กุสลตฺติก}\csromanheadernomark{2}{55-68. มหนฺตรทุก 1. กุสลตฺติก}
\proseitem{28}{นสารมฺมณํ นกุสลํ ธมฺมํ ปจฺจยา สารมฺมโณ กุสโล ธมฺโม อุปฺปชฺชติ เหตุปจฺจยา.  เหตุยา เอกํ.}

\prose{นสารมฺมณํ นอกุสลํ ธมฺมํ ปจฺจยา สารมฺมโณ อกุสโล ธมฺโม อุปฺปชฺชติ เหตุปจฺจยา.  เหตุยา เอกํ.}

\prose{นอนารมฺมณํ นอพฺยากตํ ธมฺมํ ปฏิจฺจ อนารมฺมโณ อพฺยากโต ธมฺโม อุปฺปชฺชติ เหตุปจฺจยา.  เหตุยา เอกํ.}

\proseitem{29}{โนจิตฺตํ นกุสลํ ธมฺมํ ปจฺจยา จิตฺโต กุสโล ธมฺโม อุปฺปชฺชติ เหตุปจฺจยา.  เหตุยา ตีณิ. (นอกุสเล ตีณิ. นอพฺยากเต ตีณิ. สํขิตฺตํ.)}

\proseitem{30}{นเจตสิกํ นกุสลํ ธมฺมํ ปจฺจยา เจตสิโก กุสโล ธมฺโม อุปฺปชฺชติ เหตุปจฺจยา.  เหตุยา ตีณิ.}

\csromanheader{2}{ธมฺม\-ปจฺจนียา\-นุโลม ทุกติก\-ปฏฺฐาน\-ปาฬิ}
\proseitem{31}{\csromanfolio{356}นจิตฺต\-สมฺปยุตฺตํ นกุสลํ ธมฺมํ ปจฺจยา จิตฺต\-สมฺปยุตฺโต กุสโล ธมฺโม อุปฺปชฺชติ เหตุปจฺจยา.  เหตุยา เอกํ.}

\proseitem{32}{นจิตฺต\-สํสฏฺฐํ นกุสลํ ธมฺมํ ปจฺจยา จิตฺตสํสฏฺโฐ กุสโล ธมฺโม อุปฺปชฺชติ เหตุปจฺจยา.  เหตุยา เอกํ.}

\proseitem{33}{นจิตฺต\-สมุฏฺฐานํ นกุสลํ ธมฺมํ ปจฺจยา จิตฺต\-สมุฏฺฐาโน กุสโล ธมฺโม อุปฺปชฺชติ เหตุปจฺจยา.  เหตุยา ตีณิ.}

\proseitem{34}{นจิตฺต\-สหภุํ นกุสลํ ธมฺมํ ปจฺจยา จิตฺตสหภู กุสโล ธมฺโม อุปฺปชฺชติ เหตุปจฺจยา.  เหตุยา ตีณิ.}

\proseitem{35}{นจิตฺตา\-นุปริวตฺติํ นกุสลํ ธมฺมํ ปจฺจยา จิตฺตา\-นุปริวตฺตี กุสโล ธมฺโม อุปฺปชฺชติ เหตุปจฺจยา.  เหตุยา ตีณิ.}

\proseitem{36}{นจิตฺต\-สํสฏฺฐ\-สมุฏฺฐานํ นกุสลํ ธมฺมํ ปจฺจยา จิตฺต\-สํสฏฺฐ\-สมุฏฺฐาโน กุสโล ธมฺโม อุปฺปชฺชติ เหตุปจฺจยา.  เหตุยา ตีณิ.}

\proseitem{37}{นจิตฺต\-สํสฏฺฐ\-สมุฏฺฐาน\-สหภุํ นกุสลํ ธมฺมํ ปจฺจยา จิตฺต\-สํสฏฺฐ\-สมุฏฺฐาน\-สหภู กุสโล ธมฺโม อุปฺปชฺชติ เหตุปจฺจยา.  เหตุยา ตีณิ.}

\proseitem{38}{นจิตฺต\-สํสฏฺฐ\-สมุฏฺฐานา\-นุ\-ปริวตฺติํ นกุสลํ ธมฺมํ ปจฺจยา จิตฺต\-สํสฏฺฐ\-สมุฏฺฐานา\-นุปริวตฺตี กุสโล ธมฺโม อุปฺปชฺชติ เหตุปจฺจยา.  เหตุยา ตีณิ.}

\proseitem{39}{นอชฺฌตฺติกํ นกุสลํ ธมฺมํ ปจฺจยา อชฺฌตฺติโก กุสโล ธมฺโม อุปฺปชฺชติ เหตุปจฺจยา.  เหตุยา ตีณิ.}

\proseitem{40}{นโนอุปาทา นกุสลํ ธมฺมํ ปจฺจยา โนอุปาทา กุสโล ธมฺโม อุปฺปชฺชติ เหตุปจฺจยา.  เหตุยา เอกํ.}

\proseitem{41}{นอนุปาทินฺนํ นกุสลํ ธมฺมํ ปจฺจยา อนุปาทินฺโน กุสโล ธมฺโม อุปฺปชฺชติ เหตุปจฺจยา.  เหตุยา เอกํ.}

\csromanmarkh{83-100. ปิฏฺฐิทุก}\csromanmarkhii{1. กุสลตฺติก 69. อุปาทานทุก 1. กุสลตฺติก}\csromanheadernomark{2}{83-100. ปิฏฺฐิทุก 1. กุสลตฺติก \textbf{69. อุปาทานทุก 1. กุสลตฺติก}}
\proseitem{42}{\csromanfolio{357}นอุปาทานํ นอกุสลํ ธมฺมํ ปจฺจยา อุปาทาโน อกุสโล ธมฺโม อุปฺปชฺชติ เหตุปจฺจยา.  เหตุยา ตีณิ.}

\csromanmarkh{75. กิเลสทุก}\csromanmarkhii{1. กุสลตฺติก}\csromanheadernomark{2}{75. \textbf{กิเลสทุก 1. กุสลตฺติก}}
\proseitem{43}{นกิเลสํ นอกุสลํ ธมฺมํ ปจฺจยา กิเลโส อกุสโล ธมฺโม อุปฺปชฺชติ เหตุปจฺจยา.  เหตุยา ตีณิ.}

\csromanmarkh{83-100. ปิฏฺฐิทุก}\csromanmarkhii{1. กุสลตฺติก}\csromanheadernomark{2}{83-100. ปิฏฺฐิทุก 1. กุสลตฺติก}
\proseitem{44}{นทสฺสเนน ปหาตพฺพํ นกุสลํ ธมฺมํ ปจฺจยา นทสฺสเนน ปหาตพฺโพ กุสโล ธมฺโม อุปฺปชฺชติ เหตุปจฺจยา.  เหตุยา เอกํ.}

\prose{นทสฺสเนน ปหาตพฺพํ นอกุสลํ ธมฺมํ ปจฺจยา ทสฺสเนน ปหาตพฺโพ อกุสโล ธมฺโม อุปฺปชฺชติ เหตุปจฺจยา.  นทสฺสเนน ปหาตพฺพํ นอกุสลํ ธมฺมํ ปจฺจยา นทสฺสเนน ปหาตพฺโพ อกุสโล ธมฺโม อุปฺปชฺชติ เหตุปจฺจยา.  เหตุยา เทฺว.}

\prose{นทสฺสเนน ปหาตพฺพํ นอพฺยากตํ ธมฺมํ ปฏิจฺจ นทสฺสเนน ปหาตพฺโพ อพฺยากโต ธมฺโม อุปฺปชฺชติ เหตุปจฺจยา.  นนทสฺสเนน ปหาตพฺพํ นอพฺยากตํ ธมฺมํ ปฏิจฺจ นทสฺสเนน ปหาตพฺโพ อพฺยากโต ธมฺโม อุปฺปชฺชติ เหตุปจฺจยา.  เหตุยา เทฺว.}

\proseitem{45}{นภาวนาย ปหาตพฺพํ นกุสลํ ธมฺมํ ปจฺจยา นภาวนาย ปหาตพฺโพ กุสโล ธมฺโม อุปฺปชฺชติ เหตุปจฺจยา.  เหตุยา เอกํ.}

\proseitem{46}{นทสฺสเนน ปหาตพฺพ\-เหตุกํ นกุสลํ ธมฺมํ ปจฺจยา นทสฺสเนน ปหาตพฺพ\-เหตุโก กุสโล ธมฺโม อุปฺปชฺชติ เหตุปจฺจยา.  เหตุยา เอกํ.}

\proseitem{47}{นภาวนาย ปหาตพฺพ\-เหตุกํ นกุสลํ ธมฺมํ ปจฺจยา นภาวนาย ปหาตพฺพ\-เหตุโก กุสโล ธมฺโม อุปฺปชฺชติ เหตุปจฺจยา.  เหตุยา เอกํ.}

\proseitem{48}{นสวิตกฺกํ นกุสลํ ธมฺมํ ปจฺจยา สวิตกฺโก กุสโล ธมฺโม อุปฺปชฺชติ เหตุปจฺจยา.  เหตุยา ตีณิ.}

\csromanheader{2}{ธมฺม\-ปจฺจนียา\-นุโลม ทุกติก\-ปฏฺฐาน\-ปาฬิ}
\proseitem{49}{\csromanfolio{358}นสวิจารํ นกุสลํ ธมฺมํ ปจฺจยา สวิจาโร กุสโล ธมฺโม อุปฺปชฺชติ เหตุปจฺจยา.  เหตุยา ตีณิ.}

\proseitem{50}{นสปฺปีติกํ นกุสลํ ธมฺมํ ปจฺจยา สปฺปีติโก กุสโล ธมฺโม อุปฺปชฺชติ เหตุปจฺจยา.  เหตุยา ตีณิ.}

\proseitem{51}{นปีติสหคตํ นกุสลํ ธมฺมํ ปจฺจยา ปีติสหคโต กุสโล ธมฺโม อุปฺปชฺชติ เหตุปจฺจยา.  เหตุยา ตีณิ.}

\proseitem{52}{นสุขสหคตํ นกุสลํ ธมฺมํ ปจฺจยา สุขสหคโต กุสโล ธมฺโม อุปฺปชฺชติ เหตุปจฺจยา.  เหตุยา ตีณิ.}

\proseitem{53}{นอุเปกฺขาสหคตํ นกุสลํ ธมฺมํ ปจฺจยา อุเปกฺขา\-สหคโต กุสโล ธมฺโม อุปฺปชฺชติ เหตุปจฺจยา.  เหตุยา ตีณิ.}

\proseitem{54}{นนกามา\-วจรํ นกุสลํ ธมฺมํ ปจฺจยา นกามาวจโร กุสโล ธมฺโม อุปฺปชฺชติ เหตุปจฺจยา.  นนกามา\-วจรํ นกุสลํ ธมฺมํ ปจฺจยา กามาวจโร กุสโล ธมฺโม อุปฺปชฺชติ เหตุปจฺจยา.  เหตุยา เทฺว.}

\proseitem{55}{นรูปาวจรํ นกุสลํ ธมฺมํ ปจฺจยา รูปาวจโร กุสโล ธมฺโม อุปฺปชฺชติ เหตุปจฺจยา.  นรูปาวจรํ นกุสลํ ธมฺมํ ปจฺจยา นรูปาวจโร กุสโล ธมฺโม อุปฺปชฺชติ เหตุปจฺจยา.  เหตุยา เทฺว.}

\proseitem{56}{นอรูปาวจรํ นกุสลํ ธมฺมํ ปจฺจยา อรูปาวจโร กุสโล ธมฺโม อุปฺปชฺชติ เหตุปจฺจยา.  นอรูปาวจรํ นกุสลํ ธมฺมํ ปจฺจยา นอรูปาวจโร กุสโล ธมฺโม อุปฺปชฺชติ เหตุปจฺจยา.  เหตุยา เทฺว.}

\proseitem{57}{นอปริยาปนฺนํ นกุสลํ ธมฺมํ ปจฺจยา อปริยาปนฺโน กุสโล ธมฺโม อุปฺปชฺชติ เหตุปจฺจยา.  นอปริยาปนฺนํ นกุสลํ ธมฺมํ ปจฺจยา ปริยาปนฺโน กุสโล ธมฺโม อุปฺปชฺชติ เหตุปจฺจยา.  เหตุยา เทฺว.}

\proseitem{58}{นนิยฺยานิกํ นกุสลํ ธมฺมํ ปจฺจยา นิยฺยานิโก กุสโล ธมฺโม อุปฺปชฺชติ เหตุปจฺจยา.  นนิยฺยานิกํ นกุสลํ ธมฺมํ ปจฺจยา นอนิยฺยานิโก กุสโล ธมฺโม อุปฺปชฺชติ เหตุปจฺจยา.  เหตุยา เทฺว.}

\csromanmarkh{1. เหตุทุก}\csromanmarkhii{3. วิปากตฺติก}\csromanheadernomark{2}{1. เหตุทุก 3. วิปากตฺติก}
\proseitem{59}{\csromanfolio{359}นนิยตํ นกุสลํ ธมฺมํ ปจฺจยา นิยโต กุสโล ธมฺโม อุปฺปชฺชติ เหตุปจฺจยา.  นนิยตํ นกุสลํ ธมฺมํ ปจฺจยา อนิยโต กุสโล ธมฺโม อุปฺปชฺชติ เหตุปจฺจยา.  เหตุยา เทฺว.}

\proseitem{60}{นอนุตฺตรํ นกุสลํ ธมฺมํ ปจฺจยา อนุตฺตโร กุสโล ธมฺโม อุปฺปชฺชติ เหตุปจฺจยา.  นอนุตฺตรํ นกุสลํ ธมฺมํ ปจฺจยา สอุตฺตโร กุสโล ธมฺโม อุปฺปชฺชติ เหตุปจฺจยา.  เหตุยา เทฺว.}

\proseitem{61}{นสรณํ นกุสลํ ธมฺมํ ปจฺจยา อรโณ กุสโล ธมฺโม อุปฺปชฺชติ เหตุปจฺจยา.  เหตุยา เอกํ.}

\prose{นสรณํ นอกุสลํ ธมฺมํ ปจฺจยา สรโณ อกุสโล ธมฺโม อุปฺปชฺชติ เหตุปจฺจยา.  เหตุยา เอกํ.}

\prose{นสรณํ นอพฺยากตํ ธมฺมํ ปฏิจฺจ อรโณ อพฺยากโต ธมฺโม อุปฺปชฺชติ เหตุปจฺจยา.  นอรณํ นอพฺยากตํ ธมฺมํ ปฏิจฺจ อรโณ อพฺยากโต ธมฺโม อุปฺปชฺชติ เหตุปจฺจยา.  เหตุยา เทฺว.}

\csromanmarkh{1. เหตุทุก}\csromanmarkhii{2. เวทนาตฺติก}\csromanheadernomark{2}{1. \textbf{เหตุทุก 2. เวทนาตฺติก}}
\proseitem{62}{นเหตุํ นสุขาย เวทนาย สมฺปยุตฺตํ ธมฺมํ ปฏิจฺจ เหตุ สุขาย เวทนาย สมฺปยุตฺโต ธมฺโม อุปฺปชฺชติ เหตุปจฺจยา.  นเหตุํ นสุขาย เวทนาย สมฺปยุตฺตํ ธมฺมํ ปฏิจฺจ นเหตุ สุขาย เวทนาย สมฺปยุตฺโต ธมฺโม อุปฺปชฺชติ เหตุปจฺจยา.  นเหตุํ นสุขาย เวทนาย สมฺปยุตฺตํ ธมฺมํ ปฏิจฺจ เหตุ สุขาย เวทนาย สมฺปยุตฺโต จ นเหตุ สุขาย เวทนาย สมฺปยุตฺโต จ ธมฺมา อุปฺปชฺชนฺติ เหตุปจฺจยา.  เหตุยา ตีณิ.  นเหตุ นทุกฺขาย เวทนาย สมฺปยุตฺตมูลํ ตีณิเยว.}

\prose{นเหตุํ นอทุกฺขามสุขาย เวทนาย สมฺปยุตฺตํ ธมฺมํ ปฏิจฺจ เหตุ อทุกฺขม\-สุขาย เวทนาย สมฺปยุตฺโต ธมฺโม อุปฺปชฺชติ เหตุปจฺจยา.  เหตุยา ตีณิ.}

\csromanmarkh{1. เหตุทุก}\csromanmarkhii{3. วิปากตฺติก}\csromanheadernomark{2}{1. เหตุทุก 3. วิปากตฺติก}
\proseitem{63}{นเหตุํ นวิปากํ ธมฺมํ ปฏิจฺจ เหตุ วิปาโก ธมฺโม อุปฺปชฺชติ เหตุปจฺจยา.  เหตุยา ตีณิ.}

\csromanheader{2}{ธมฺม\-ปจฺจนียา\-นุโลม ทุกติก\-ปฏฺฐาน\-ปาฬิ}
\prose{\csromanfolio{360}นเหตุํ นวิปาก\-ธมฺม\-ธมฺมํ ปจฺจยา เหตุ วิปาก\-ธมฺม\-ธมฺโม อุปฺปชฺชติ เหตุปจฺจยา.  เหตุยา ตีณิ.}

\prose{นเหตุํ นเน\-ว\-วิปาก\-น\-วิปากธมฺม\-ธมฺมํ ปฏิจฺจ นเหตุ เนว\-วิปาก\-น\-วิปาก\-ธมฺม\-ธมฺโม อุปฺปชฺชติ เหตุปจฺจยา.  เหตุยา ตีณิ.}

\csromanmarkh{1. เหตุทุก}\csromanmarkhii{4. อุปาทินฺนตฺติก}\csromanheadernomark{2}{1. \textbf{เหตุทุก 4. อุปาทินฺนตฺติก}}
\proseitem{64}{นเหตุ นอุปาทินฺนุปาทานิโย ธมฺโม เหตุสฺส อุปาทินฺนุ\-ปาทานิยสฺส ธมฺมสฺส อารมฺมณ\-ปจฺจเยน ปจฺจโย.  นเหตุ นอุปาทินฺนุปาทานิโย ธมฺโม นเหตุสฺส อุปาทินฺนุ\-ปาทานิยสฺส ธมฺมสฺส อารมฺมณ\-ปจฺจเยน ปจฺจโย.  นเหตุ นอุปาทินฺนุปาทานิโย ธมฺโม เหตุสฺส อุปาทินฺนุ\-ปาทานิยสฺส จ นเหตุสฺส อุปาทินฺนุ\-ปาทานิยสฺส จ ธมฺมสฺส อารมฺมณ\-ปจฺจเยน ปจฺจโย. (3)}

\prose{นนเหตุ นอุปาทินฺนุปาทานิโย ธมฺโม นเหตุสฺส อุปาทินฺนุ\-ปาทานิยสฺส ธมฺมสฺส อารมฺมณ\-ปจฺจเยน ปจฺจโย. ตีณิ.}

\prose{นเหตุ นอุปาทินฺนุปาทานิโย จ นนเหตุ นอุปาทินฺนุปาทานิโย จ ธมฺมา เหตุสฺส อุปาทินฺนุ\-ปาทานิยสฺส ธมฺมสฺส อารมฺมณ\-ปจฺจเยน ปจฺจโย. ตีณิ. (นว ปญฺหา. สํขิตฺตํ.)}

\prose{นเหตุํ นอนุปาทินฺนุ\-ปาทานิยํ ธมฺมํ ปจฺจยา เหตุ อนุปาทินฺนุ\-ปาทานิโย ธมฺโม อุปฺปชฺชติ เหตุปจฺจยา.  เหตุยา ตีณิ.}

\prose{นเหตุํ นอนุปาทินฺน\-อนุปาทานิยํ ธมฺมํ ปจฺจยา เหตุ อนุปาทินฺน\-อนุปาทานิโย ธมฺโม อุปฺปชฺชติ เหตุปจฺจยา.  เหตุยา ตีณิ.}

\csromanmarkh{1. เหตุทุก}\csromanmarkhii{5. สํกิลิฏฺฐ\-ตฺติก}\csromanheadernomark{2}{1. เหตุทุก 5. สํกิลิฏฺฐ\-ตฺติก}
\proseitem{65}{นเหตุํ นสํกิลิฏฺฐ\-สํกิเลสิกํ ธมฺมํ ปจฺจยา เหตุ สํกิลิฏฺฐ\-สํกิเลสิโก ธมฺโม อุปฺปชฺชติ เหตุปจฺจยา.  เหตุยา ตีณิ.}

\prose{นเหตุํ นอสํกิลิฏฺฐ\-สํกิเลสิกํ ธมฺมํ ปฏิจฺจ นเหตุ อสํกิลิฏฺฐ\-สํกิเลสิโก ธมฺโม อุปฺปชฺชติ เหตุปจฺจยา.  เหตุยา ตีณิ.}

\prose{นเหตุํ นอสํกิลิฏฺฐ\-อสํกิเลสิกํ ธมฺมํ ปจฺจยา เหตุ อสํกิลิฏฺฐ\-อสํกิเลสิโก ธมฺโม อุปฺปชฺชติ เหตุปจฺจยา.  เหตุยา ตีณิ.}

\csromanmarkh{1. เหตุทุก}\csromanmarkhii{9. ทสฺสน\-เหตุ\-ตฺติก 1. เหตุทุก 6. วิตกฺกตฺติก}\csromanheadernomark{2}{1. เหตุทุก 9. ทสฺสน\-เหตุ\-ตฺติก 1. เหตุทุก 6. วิตกฺกตฺติก}
\proseitem{66}{\csromanfolio{361}นเหตุํ นสวิตกฺก\-สวิจารํ ธมฺมํ ปฏิจฺจ เหตุ สวิตกฺก\-สวิจาโร ธมฺโม อุปฺปชฺชติ เหตุปจฺจยา.  เหตุยา ตีณิ.}

\prose{นเหตุํ นอวิตกฺก\-วิจารมตฺตํ ธมฺมํ ปฏิจฺจ เหตุ อวิตกฺก\-วิจารมตฺโต ธมฺโม อุปฺปชฺชติ เหตุปจฺจยา.  เหตุยา ตีณิ.}

\prose{นเหตุํ นอวิตกฺก\-อวิจารํ ธมฺมํ ปฏิจฺจ นเหตุ อวิตกฺก\-อวิจาโร ธมฺโม อุปฺปชฺชติ เหตุปจฺจยา.  เหตุยา ตีณิ.}

\csromanmarkh{1. เหตุทุก}\csromanmarkhii{7. ปีติตฺติก}\csromanheadernomark{2}{1. \textbf{เหตุทุก 7. ปีติตฺติก}}
\proseitem{67}{นเหตุํ นปีติสหคตํ ธมฺมํ ปฏิจฺจ เหตุ ปีติสหคโต ธมฺโม อุปฺปชฺชติ เหตุปจฺจยา.  เหตุยา ตีณิ.}

\prose{นเหตุํ นสุขสหคตํ ธมฺมํ ปฏิจฺจ เหตุ สุขสหคโต ธมฺโม อุปฺปชฺชติ เหตุปจฺจยา.  เหตุยา ตีณิ.}

\prose{นเหตุํ นอุเปกฺขาสหคตํ ธมฺมํ ปฏิจฺจ เหตุ อุเปกฺขา\-สหคโต ธมฺโม อุปฺปชฺชติ เหตุปจฺจยา.  เหตุยา ตีณิ.}

\csromanmarkh{1. เหตุทุก}\csromanmarkhii{8. ทสฺสนตฺติก}\csromanheadernomark{2}{1. \textbf{เหตุทุก 8. ทสฺสนตฺติก}}
\proseitem{68}{นเหตุํ นทสฺสเนน ปหาตพฺพํ ธมฺมํ ปจฺจยา เหตุ ทสฺสเนน ปหาตพฺโพ ธมฺโม อุปฺปชฺชติ เหตุปจฺจยา.  เหตุยา ตีณิ.}

\prose{นเหตุํ นภาวนาย ปหาตพฺพํ ธมฺมํ ปจฺจยา เหตุ ภาวนาย ปหาตพฺโพ ธมฺโม อุปฺปชฺชติ เหตุปจฺจยา.  เหตุยา ตีณิ.}

\proseitem{69}{นเหตุํ นเน\-ว\-ทสฺสเนน นภาวนาย ปหาตพฺพํ ธมฺมํ ปฏิจฺจ นเหตุ เนวทสฺสเนน นภาวนาย ปหาตพฺโพ ธมฺโม อุปฺปชฺชติ เหตุปจฺจยา.  นนเหตุํ นเน\-ว\-ทสฺสเนน นภาวนาย ปหาตพฺพํ ธมฺมํ ปฏิจฺจ นเหตุ เนวทสฺสเนน นภาวนาย ปหาตพฺโพ ธมฺโม อุปฺปชฺชติ เหตุปจฺจยา.  (ทุกมูลํ เอกํ สํขิตฺตํ. สพฺพตฺถ ตีณิ.)}

\csromanmarkh{1. เหตุทุก}\csromanmarkhii{9. ทสฺสน\-เหตุ\-ตฺติก}\csromanheadernomark{2}{1. เหตุทุก 9. ทสฺสน\-เหตุ\-ตฺติก}
\proseitem{70}{นนเหตุํ นทสฺสเนน ปหาตพฺพ\-เหตุกํ ธมฺมํ ปฏิจฺจ นเหตุ ทสฺสเนน ปหาตพฺพ\-เหตุโก ธมฺโม อุปฺปชฺชติ เหตุปจฺจยา.  เหตุยา เอกํ.}

\csromanheader{2}{ธมฺม\-ปจฺจนียา\-นุโลม ทุกติก\-ปฏฺฐาน\-ปาฬิ}
\prose{\csromanfolio{362}นนเหตุํ นภาวนาย ปหาตพฺพ\-เหตุกํ ธมฺมํ ปฏิจฺจ นเหตุ ภาวนาย ปหาตพฺพ\-เหตุโก ธมฺโม อุปฺปชฺชติ เหตุปจฺจยา.  เหตุยา เอกํ.}

\proseitem{71}{นเหตุํ นเน\-ว\-ทสฺสเนน นภาวนาย ปหาตพฺพ\-เหตุกํ ธมฺมํ ปฏิจฺจ นเหตุ เนวทสฺสเนน นภาวนาย ปหาตพฺพ\-เหตุโก ธมฺโม อุปฺปชฺชติ เหตุปจฺจยา.  นนเหตุํ นเน\-ว\-ทสฺสเนน นภาวนาย ปหาตพฺพ\-เหตุกํ ธมฺมํ ปฏิจฺจ นเหตุ เนวทสฺสเนน นภาวนาย ปหาตพฺพ\-เหตุโก ธมฺโม อุปฺปชฺชติ เหตุปจฺจยา. คณิตเกน ตีณิ.}

\csromanmarkh{1. เหตุทุก}\csromanmarkhii{10. อาจยคามิตฺติก}\csromanheadernomark{2}{1. เหตุทุก 10. อาจยคามิตฺติก}
\proseitem{72}{นเหตุํ นอาจยคามิํ ธมฺมํ ปจฺจยา เหตุ อาจยคามี ธมฺโม อุปฺปชฺชติ เหตุปจฺจยา.  เหตุยา ตีณิ.}

\prose{นเหตุํ นอปจยคามิํ ธมฺมํ ปจฺจยา เหตุ อปจยคามี ธมฺโม อุปฺปชฺชติ เหตุปจฺจยา.  เหตุยา ตีณิ.}

\prose{นเหตุํ นเน\-วา\-จยคามิ\-นา\-ปจยคามิํ ธมฺมํ ปฏิจฺจ นเหตุ เนวา\-จยคามิ\-นา\-ปจยคามี ธมฺโม อุปฺปชฺชติ เหตุปจฺจยา.  นนเหตุํ นเน\-วา\-จยคามิ\-นา\-ปจยคามิํ ธมฺมํ ปฏิจฺจ นเหตุ เนวา\-จยคามิ\-นา\-ปจยคามี ธมฺโม อุปฺปชฺชติ เหตุปจฺจยา. คณิตเกน ตีณิ.}

\csromanmarkh{1. เหตุทุก}\csromanmarkhii{11. เสกฺขตฺติก}\csromanheadernomark{2}{1. เหตุทุก 11. เสกฺขตฺติก}
\proseitem{73}{นเหตุํ นเสกฺขํ ธมฺมํ ปจฺจยา เหตุ เสกฺโข ธมฺโม อุปฺปชฺชติ เหตุปจฺจยา.  เหตุยา ตีณิ.}

\prose{นเหตุํ นอเสกฺขํ ธมฺมํ ปจฺจยา เหตุ อเสกฺโข ธมฺโม อุปฺปชฺชติ เหตุปจฺจยา.  เหตุยา ตีณิ.}

\prose{นเหตุํ นเน\-ว\-เสกฺข\-นาเสกฺขํ ธมฺมํ ปฏิจฺจ นเหตุ เนว\-เสกฺข\-นา\-เสกฺโข ธมฺโม อุปฺปชฺชติ เหตุปจฺจยา.  นนเหตุํ นเน\-ว\-เสกฺข\-นาเสกฺขํ ธมฺมํ ปฏิจฺจ นเหตุ เนว\-เสกฺข\-นา\-เสกฺโข ธมฺโม อุปฺปชฺชติ เหตุปจฺจยา. คณิตเกน ตีณิ.}

\vspace{\tipitakaparskip}
\csromancenter{เหตุทุก 15. มิจฺฉตฺตนิยต\-ตฺติก 1. เหตุทุก 12. ปริตฺตตฺติก}
\proseitem{74}{\csromanfolio{363}นเหตุํ นปริตฺตํ ธมฺมํ ปฏิจฺจ นเหตุ ปริตฺโต ธมฺโม อุปฺปชฺชติ เหตุปจฺจยา.  นนเหตุํ นปริตฺตํ ธมฺมํ ปฏิจฺจ นเหตุ ปริตฺโต ธมฺโม อุปฺปชฺชติ เหตุปจฺจยา. คณิตเกน ตีณิ.}

\prose{นเหตุํ นมหคฺคตํ ธมฺมํ ปจฺจยา เหตุ มหคฺคโต ธมฺโม อุปฺปชฺชติ เหตุปจฺจยา.  เหตุยา ตีณิ.}

\prose{นเหตุํ นอปฺปมาณํ ธมฺมํ ปจฺจยา เหตุ อปฺปมาโณ ธมฺโม อุปฺปชฺชติ เหตุปจฺจยา.  เหตุยา ตีณิ.}

\csromanmarkh{1. เหตุทุก}\csromanmarkhii{13. ปริตฺตา\-รมฺมณ\-ตฺติก}\csromanheadernomark{2}{1. เหตุทุก 13. ปริตฺตา\-รมฺมณ\-ตฺติก}
\proseitem{75}{นเหตุํ นปริตฺตา\-รมฺมณํ ธมฺมํ ปฏิจฺจ เหตุ ปริตฺตารมฺมโณ ธมฺโม อุปฺปชฺชติ เหตุปจฺจยา.  เหตุยา ตีณิ.}

\prose{นเหตุํ นมหคฺคตา\-รมฺมณํ ธมฺมํ ปจฺจยา เหตุ มหคฺคตา\-รมฺมโณ ธมฺโม อุปฺปชฺชติ เหตุปจฺจยา.  เหตุยา ตีณิ.}

\prose{นเหตุํ นอปฺปมาณา\-รมฺมณํ ธมฺมํ ปจฺจยา เหตุ อปฺปมาณารมฺมโณ ธมฺโม อุปฺปชฺชติ เหตุปจฺจยา.  เหตุยา ตีณิ.}

\csromanmarkh{1. เหตุทุก}\csromanmarkhii{14. หีนตฺติก}\csromanheadernomark{2}{1. เหตุทุก 14. หีนตฺติก}
\proseitem{76}{นเหตุํ นหีนํ ธมฺมํ ปจฺจยา เหตุ หีโน ธมฺโม อุปฺปชฺชติ เหตุปจฺจยา.  เหตุยา ตีณิ.}

\prose{นเหตุํ นมชฺฌิมํ ธมฺมํ ปฏิจฺจ นเหตุ มชฺฌิโม ธมฺโม อุปฺปชฺชติ เหตุปจฺจยา.  นนเหตุ นมชฺฌิมํ ธมฺมํ ปฏิจฺจ นเหตุ มชฺฌิโม ธมฺโม อุปฺปชฺชติ เหตุปจฺจยา. คณิตเกน ตีณิ.}

\prose{นเหตุํ นปณีตํ ธมฺมํ ปจฺจยา เหตุ ปณีโต ธมฺโม อุปฺปชฺชติ เหตุปจฺจยา.  เหตุยา ตีณิ.}

\csromanmarkh{1. เหตุทุก}\csromanmarkhii{15. มิจฺฉตฺตนิยต\-ตฺติก}\csromanheadernomark{2}{1. เหตุทุก 15. มิจฺฉตฺตนิยต\-ตฺติก}
\proseitem{77}{นเหตุํ นมิจฺฉตฺต\-นิยตํ ธมฺมํ ปจฺจยา เหตุ มิจฺฉตฺต\-นิยโต ธมฺโม อุปฺปชฺชติ เหตุปจฺจยา.  เหตุยา ตีณิ.}

\csromanheader{2}{ธมฺม\-ปจฺจนียา\-นุโลม ทุกติก\-ปฏฺฐาน\-ปาฬิ}
\prose{\csromanfolio{364}นเหตุํ นสมฺมตฺต\-นิยตํ ธมฺมํ ปจฺจยา เหตุ สมฺมตฺตนิยโต ธมฺโม อุปฺปชฺชติ เหตุปจฺจยา.  เหตุยา ตีณิ.}

\prose{นเหตุํ นอนิยตํ ธมฺมํ ปฏิจฺจ นเหตุ อนิยโต ธมฺโม อุปฺปชฺชติ เหตุปจฺจยา.  นนเหตุํ นอนิยตํ ธมฺมํ ปฏิจฺจ นเหตุ อนิยโต ธมฺโม อุปฺปชฺชติ เหตุปจฺจยา. คณิตเกน ตีณิ.}

\csromanmarkh{1. เหตุทุก}\csromanmarkhii{16. มคฺคา\-รมฺมณ\-ตฺติก}\csromanheadernomark{2}{1. เหตุทุก 16. มคฺคา\-รมฺมณ\-ตฺติก}
\proseitem{78}{นเหตุํ นมคฺคา\-รมฺมณํ ธมฺมํ ปจฺจยา เหตุ มคฺคารมฺมโณ ธมฺโม อุปฺปชฺชติ เหตุปจฺจยา.  เหตุยา ตีณิ.}

\prose{นเหตุํ นมคฺคเหตุกํ ธมฺมํ ปจฺจยา เหตุ มคฺคเหตุโก ธมฺโม อุปฺปชฺชติ เหตุปจฺจยา.  เหตุยา ตีณิ.}

\prose{นเหตุํ นมคฺคา\-ธิปติํ ธมฺมํ ปจฺจยา เหตุ มคฺคาธิปติ ธมฺโม อุปฺปชฺชติ เหตุปจฺจยา.  เหตุยา ตีณิ.}

\csromanmarkh{1. เหตุทุก}\csromanmarkhii{17. อุปฺปนฺนตฺติก}\csromanheadernomark{2}{1. เหตุทุก 17. อุปฺปนฺนตฺติก}
\proseitem{79}{นเหตุ นอุปฺปนฺโน ธมฺโม เหตุสฺส อุปฺปนฺนสฺส ธมฺมสฺส อารมฺมณ\-ปจฺจเยน ปจฺจโย.  อารมฺมเณ นว.}

\csromanmarkh{1. เหตุทุก}\csromanmarkhii{18. อตีตตฺติก}\csromanheadernomark{2}{1. เหตุทุก 18. อตีตตฺติก}
\proseitem{80}{นเหตุ นปจฺจุปฺปนฺโน ธมฺโม เหตุสฺส ปจฺจุปฺปนฺนสฺส ธมฺมสฺส อารมฺมณ\-ปจฺจเยน ปจฺจโย.  อารมฺมเณ นว.}

\csromanmarkh{1. เหตุทุก}\csromanmarkhii{19. อตีตา\-รมฺมณ\-ตฺติก}\csromanheadernomark{2}{1. เหตุทุก 19. อตีตา\-รมฺมณ\-ตฺติก}
\proseitem{81}{นเหตุํ นอตีตา\-รมฺมณํ ธมฺมํ ปฏิจฺจ เหตุ อตีตารมฺมโณ ธมฺโม อุปฺปชฺชติ เหตุปจฺจยา.  เหตุยา ตีณิ.}

\prose{นเหตุํ นอนาคตา\-รมฺมณํ ธมฺมํ ปจฺจยา เหตุ อนาคตารมฺมโณ ธมฺโม อุปฺปชฺชติ เหตุปจฺจยา.  เหตุยา ตีณิ.}

\prose{นเหตุํ นปจฺจุปฺปนฺนา\-รมฺมณํ ธมฺมํ ปฏิจฺจ เหตุ ปจฺจุปฺปนฺนา\-รมฺมโณ ธมฺโม อุปฺปชฺชติ เหตุปจฺจยา.  เหตุยา ตีณิ.}

\vspace{\tipitakaparskip}
\csromancenter{สเหตุกทุก 22. สนิทสฺสน\-ตฺติก 1. เหตุทุก 20. อชฺฌตฺตตฺติก}
\proseitem{82}{\csromanfolio{365}นเหตุ นอชฺฌตฺโต ธมฺโม เหตุสฺส อชฺฌตฺตสฺส ธมฺมสฺส อารมฺมณ\-ปจฺจเยน ปจฺจโย. (สํขิตฺตํ.) . อารมฺมเณ นว, อธิปติยา อุปนิสฺสเย นว, ปุเรชาเต อตฺถิยา อวิคเต ตีณิ.}

\prose{นเหตุ นพหิทฺธา ธมฺโม เหตุสฺส พหิทฺธา ธมฺมสฺส อารมฺมณ\-ปจฺจเยน ปจฺจโย. (สํขิตฺตํ.) . อารมฺมเณ นว, อธิปติยา อุปนิสฺสเย นว, ปุเรชาเต อตฺถิยา อวิคเต ตีณิ. (อชฺฌตฺตตฺติโก น ลพฺภติ ปฏิจฺจวาราทีสุ.)}

\vspace{\tipitakaparskip}
\csromancenter{(อชฺฌตฺตตฺติโก น ลพฺภติ ปฏิจฺจ\-วาราทีสุ.)}
\csromanmarkh{1. เหตุทุก}\csromanmarkhii{21. อชฺฌตฺตา\-รมฺมณ\-ตฺติก}\csromanheadernomark{2}{1. เหตุทุก 21. อชฺฌตฺตา\-รมฺมณ\-ตฺติก}
\proseitem{83}{นเหตุํ นอชฺฌตฺตา\-รมฺมณํ ธมฺมํ ปจฺจยา เหตุ อชฺฌตฺตา\-รมฺมโณ ธมฺโม อุปฺปชฺชติ เหตุปจฺจยา.  เหตุยา ตีณิ.}

\prose{นเหตุํ นพหิทฺธา\-รมฺมณํ ธมฺมํ ปจฺจยา เหตุ พหิทฺธา\-รมฺมโณ ธมฺโม อุปฺปชฺชติ เหตุปจฺจยา.  เหตุยา ตีณิ. (อชฺฌตฺต\-พหิทฺธา\-รมฺมณํ นตฺถิ.)}

\csromanmarkh{1. เหตุทุก}\csromanmarkhii{22. สนิทสฺสน\-ตฺติก}\csromanheadernomark{2}{1. เหตุทุก 22. สนิทสฺสน\-ตฺติก}
\proseitem{84}{นเหตุํ นสนิทสฺสน\-สปฺปฏิฆํ ธมฺมํ ปฏิจฺจ นเหตุ สนิทสฺสน\-สปฺปฏิโฆ ธมฺโม อุปฺปชฺชติ เหตุปจฺจยา. เอกํ.}

\prose{นนเหตุํ นสนิทสฺสน\-สปฺปฏิฆํ ธมฺมํ ปฏิจฺจ นเหตุ สนิทสฺสน\-สปฺปฏิโฆ ธมฺโม อุปฺปชฺชติ เหตุปจฺจยา. เอกํ. (นเหตุนอนิทสฺสนสปฺปฏิฆมูเลปิ ตีณิเยว.)}

\vspace{\tipitakaparskip}
\csromancenter{(นเหตุนอนิทสฺสนสปฺปฏิฆมูเลปิ ตีณิเยว.)}
\proseitem{85}{นเหตุํ นอนิทสฺสน\-อปฺปฏิฆํ ธมฺมํ ปฏิจฺจ นเหตุ อนิทสฺสน\-อปฺปฏิโฆ ธมฺโม อุปฺปชฺชติ เหตุปจฺจยา.  เหตุยา เอกํ.}

\csromanmarkh{2. สเหตุกทุก}\csromanmarkhii{22. สนิทสฺสน\-ตฺติก}\csromanheadernomark{2}{2. สเหตุกทุก 22. สนิทสฺสน\-ตฺติก}
\proseitem{86}{นสเหตุกํ นสนิทสฺสน\-สปฺปฏิฆํ ธมฺมํ ปฏิจฺจ อเหตุโก สนิทสฺสน\-สปฺปฏิโฆ ธมฺโม อุปฺปชฺชติ เหตุปจฺจยา. เอกํ.}

\prose{นอเหตุกํ นสนิทสฺสน\-สปฺปฏิฆํ ธมฺมํ ปฏิจฺจ อเหตุโก สนิทสฺสน\-สปฺปฏิโฆ ธมฺโม อุปฺปชฺชติ เหตุปจฺจยา. เอกํ. คณิตเกน ตีณิ ปญฺหา.}

\vspace{\tipitakaparskip}
\csromancenter{(นสเหตุกนอนิทสฺสนสปฺปฏิฆมูเลปิ ตีณิเยว.)}
\csromanheader{2}{ธมฺม\-ปจฺจนียา\-นุโลม ทุกติก\-ปฏฺฐาน\-ปาฬิ}
\prose{\csromanfolio{366}นสเหตุกํ นอนิทสฺสน\-อปฺปฏิฆํ ธมฺมํ ปฏิจฺจ อเหตุโก อนิทสฺสน\-อปฺปฏิโฆ ธมฺโม อุปฺปชฺชติ เหตุปจฺจยา.  เหตุยา เอกํ.}

\prose{(นเหตุ\-สมฺปยุตฺตํ นสนิทสฺสน\-สปฺปฏิฆํ สเหตุก\-ทุก\-สทิสํ. ตีณิ ปญฺหา. เหตุ\-สเหตุก\-ทุเก จ เหตุ\-เหตุ\-สมฺปยุตฺต\-ทุเก จ ปญฺหา โน ลพฺภติ.) 6. นเหตุ\-สเหตุกทุก 22. สนิทสฺสน\-ตฺติก}

\proseitem{87}{นเหตุํ นสเหตุกํ นสนิทสฺสน\-สปฺปฏิฆํ ธมฺมํ ปฏิจฺจ นเหตุ อเหตุโก สนิทสฺสน\-สปฺปฏิโฆ ธมฺโม อุปฺปชฺชติ เหตุปจฺจยา.  เหตุยา เอกํ.}

\prose{นเหตุํ นอเหตุกํ นสนิทสฺสน\-สปฺปฏิฆํ ธมฺมํ ปฏิจฺจ นเหตุ อเหตุโก สนิทสฺสน\-สปฺปฏิโฆ ธมฺโม อุปฺปชฺชติ เหตุปจฺจยา. เอกํ.  คณิตเกน ตีณิ.}

\prose{นเหตุํ นสเหตุกํ นอนิทสฺสน\-สปฺปฏิฆํ ธมฺมํ ปฏิจฺจ. ตีณิเยว.}

\prose{นเหตุํ นสเหตุกํ นอนิทสฺสน\-อปฺปฏิฆํ ธมฺมํ ปฏิจฺจ นเหตุ อเหตุโก อนิทสฺสน\-อปฺปฏิโฆ ธมฺโม อุปฺปชฺชติ เหตุปจฺจยา.  เหตุยา เอกํ. 7. สปฺปจฺจยทุก 22. สนิทสฺสน\-ตฺติก}

\proseitem{88}{นอปฺปจฺจยํ นสนิทสฺสน\-สปฺปฏิฆํ ธมฺมํ ปฏิจฺจ สปฺปจฺจโย สนิทสฺสน\-สปฺปฏิโฆ ธมฺโม อุปฺปชฺชติ เหตุปจฺจยา.  เหตุยา เอกํ.}

\prose{นอปฺปจฺจยํ นอนิทสฺสน\-สปฺปฏิฆํ ธมฺมํ ปฏิจฺจ สปฺปจฺจโย อนิทสฺสน\-สปฺปฏิโฆ ธมฺโม อุปฺปชฺชติ เหตุปจฺจยา.  เหตุยา เอกํ.}

\prose{(นอปฺปจฺจยนอนิทสฺสนอปฺปฏิฆมูเลปิ เอกเมว.) . เหตุยา เอกํ, อธิปติยา เอกํ ฯเปฯ อวิคเต เอกํ.  (นสงฺขตํ นสปฺปจฺจย\-สทิสํ.) 9. สนิทสฺสนทุก 22. สนิทสฺสน\-ตฺติก}

\proseitem{89}{นสนิทสฺสนํ นสนิทสฺสน\-สปฺปฏิฆํ ธมฺมํ ปฏิจฺจ สนิทสฺสโน สนิทสฺสน\-สปฺปฏิโฆ ธมฺโม อุปฺปชฺชติ เหตุปจฺจยา.  เหตุยา เอกํ.}

\csromanmarkh{12. โลกิยทุก}\csromanmarkhii{22. สนิทสฺสน\-ตฺติก}\csromanheadernomark{2}{12. โลกิยทุก 22. สนิทสฺสน\-ตฺติก}
\prose{\csromanfolio{367}นสนิทสฺสนํ นอนิทสฺสน\-สปฺปฏิฆํ ธมฺมํ ปฏิจฺจ อนิทสฺสโน อนิทสฺสน\-สปฺปฏิโฆ ธมฺโม อุปฺปชฺชติ เหตุปจฺจยา.  เหตุยา เอกํ.}

\prose{นสนิทสฺสนํ นอนิทสฺสน\-อปฺปฏิฆํ ธมฺมํ ปฏิจฺจ อนิทสฺสโน อนิทสฺสน\-อปฺปฏิโฆ ธมฺโม อุปฺปชฺชติ เหตุปจฺจยา.  เหตุยา เอกํ.}

\csromanmarkh{10. สปฺปฏิฆทุกฺอ}\csromanmarkhii{22. สนิทสฺสน\-ตฺติก}\csromanheadernomark{2}{10. สปฺปฏิฆทุกฺอ 22. สนิทสฺสน\-ตฺติก}
\proseitem{90}{นสปฺปฏิฆํ นสนิทสฺสน\-สปฺปฏิฆํ ธมฺมํ ปฏิจฺจ สปฺปฏิโฆ สนิทสฺสน\-สปฺปฏิโฆ ธมฺโม อุปฺปชฺชติ เหตุปจฺจยา. เอกํ.}

\prose{นอปฺปฏิฆํ นสนิทสฺสน\-สปฺปฏิฆํ ธมฺมํ ปฏิจฺจ สปฺปฏิโฆ สนิทสฺสน\-สปฺปฏิโฆ ธมฺโม อุปฺปชฺชติ เหตุปจฺจยา. เอกํ.}

\prose{นสปฺปฏิฆํ นสนิทสฺสนสปฺปฏิฆญฺจ นอปฺปฏิฆํ นสนิทสฺสนสปฺปฏิฆญฺจ ธมฺมํ ปฏิจฺจ สปฺปฏิโฆ สนิทสฺสน\-สปฺปฏิโฆ ธมฺโม อุปฺปชฺชติ เหตุปจฺจยา. เอกํ. (สพฺพตฺถ ตีณิ.)}

\prose{นสปฺปฏิฆํ นอนิทสฺสน\-อปฺปฏิฆํ ธมฺมํ ปฏิจฺจ สปฺปฏิโฆ อนิทสฺสน\-สปฺปฏิโฆ ธมฺโม อุปฺปชฺชติ เหตุปจฺจยา.  เหตุยา เอกํ.}

\prose{นอปฺปฏิฆํ นอนิทสฺสน\-อปฺปฏิฆํ ธมฺมํ ปฏิจฺจ อปฺปฏิโฆ อนิทสฺสน\-อปฺปฏิโฆ ธมฺโม อุปฺปชฺชติ เหตุปจฺจยา.  เหตุยา เอกํ.}

\csromanmarkh{11. รูปีทุก}\csromanmarkhii{22. สนิทสฺสน\-ตฺติก}\csromanheadernomark{2}{11. รูปีทุก 22. สนิทสฺสน\-ตฺติก}
\proseitem{91}{นรูปิํ นสนิทสฺสน\-สปฺปฏิฆํ ธมฺมํ ปฏิจฺจ รูปี สนิทสฺสน\-สปฺปฏิโฆ ธมฺโม อุปฺปชฺชติ เหตุปจฺจยา.  เหตุยา ตีณิ.}

\prose{นรูปิํ นอนิทสฺสน\-สปฺปฏิฆํ ธมฺมํ ปฏิจฺจ รูปี อนิทสฺสน\-สปฺปฏิโฆ ธมฺโม อุปฺปชฺชติ เหตุปจฺจยา.  เหตุยา ตีณิ.}

\prose{นอรูปิํ นอนิทสฺสน\-อปฺปฏิฆํ ธมฺมํ ปฏิจฺจ รูปี อนิทสฺสน\-อปฺปฏิโฆ ธมฺโม อุปฺปชฺชติ เหตุปจฺจยา.  เหตุยา เอกํ.}

\csromanmarkh{12. โลกิยทุก}\csromanmarkhii{22. สนิทสฺสน\-ตฺติก}\csromanheadernomark{2}{12. โลกิยทุก 22. สนิทสฺสน\-ตฺติก}
\proseitem{92}{นโลกิยํ นสนิทสฺสน\-สปฺปฏิฆํ ธมฺมํ ปฏิจฺจ โลกิโย สนิทสฺสน\-สปฺปฏิโฆ ธมฺโม อุปฺปชฺชติ เหตุปจฺจยา. เอกํ.}

\csromanheader{2}{ธมฺม\-ปจฺจนียา\-นุโลม ทุกติก\-ปฏฺฐาน\-ปาฬิ}
\prose{\csromanfolio{368}นโลกุตฺตรํ นสนิทสฺสน\-สปฺปฏิฆํ ธมฺมํ ปฏิจฺจ โลกิโย สนิทสฺสน\-สปฺปฏิโฆ ธมฺโม อุปฺปชฺชติ เหตุปจฺจยา. เอกํ.}

\prose{(คณิตเกน คเณตพฺพา ตีณิ ปญฺหา.)}

\prose{นโลกิยํ นอนิทสฺสน\-สปฺปฏิฆํ ธมฺมํ ปฏิจฺจ โลกิโย อนิทสฺสน\-สปฺปฏิโฆ. (ตีณิเยว ปญฺหา กาตพฺพา.)}

\prose{นโลกุตฺตรํ นอนิทสฺสน\-อปฺปฏิฆํ ธมฺมํ ปฏิจฺจ โลกิโย อนิทสฺสน\-อปฺปฏิโฆ ธมฺโม อุปฺปชฺชติ เหตุปจฺจยา.  เหตุยา เอกํ.}

\csromanmarkh{13. เกนจิวิญฺเญยฺยทฺอุก}\csromanmarkhii{22. สนิทสฺสน\-ตฺติก}\csromanheadernomark{2}{13. เกนจิวิญฺเญยฺยทฺอุก 22. สนิทสฺสน\-ตฺติก}
\prosecont{นเกนจิ วิญฺเญยฺยํ นสนิทสฺสน\-สปฺปฏิฆํ ธมฺมํ ปฏิจฺจ เกนจิ วิญฺเญยฺโย สนิทสฺสน\-สปฺปฏิโฆ ธมฺโม อุปฺปชฺชติ เหตุปจฺจยา.  นเกนจิ วิญฺเญยฺยํ นสนิทสฺสน\-สปฺปฏิฆํ ธมฺมํ ปฏิจฺจ เกนจิ นวิญฺเญยฺโย สนิทสฺสน\-สปฺปฏิโฆ ธมฺโม อุปฺปชฺชติ เหตุปจฺจยา.  นเกนจิ วิญฺเญยฺยํ นสนิทสฺสน\-สปฺปฏิฆํ ธมฺมํ ปฏิจฺจ เกนจิ วิญฺเญยฺโย สนิทสฺสน\-สปฺปฏิโฆ จ เกนจิ นวิญฺเญยฺโย สนิทสฺสน\-สปฺปฏิโฆ จ ธมฺมา อุปฺปชฺชนฺติ เหตุปจฺจยา. (3)}

\proseitem{93}{นเกนจิ นวิญฺเญยฺยํ นสนิทสฺสน\-สปฺปฏิฆํ ธมฺมํ ปฏิจฺจ. ตีณิเยว.}

\prose{นเกนจิ วิญฺเญยฺยํ นสนิทสฺสนสปฺปฏิฆญฺจ นเกนจิ นวิญฺเญยฺยํ นสนิทสฺสนสปฺปฏิฆญฺจ ธมฺมํ ปฏิจฺจ. ตีณิเยว. (สพฺพตฺถ นว.) นเกนจิ วิญฺเญยฺยํ นอนิทสฺสน\-สปฺปฏิฆสฺส ปุริมสทิสํ นวปญฺหํ กาตพฺพํ. นเกนจิ วิญฺเญยฺยํ นอนิทสฺสนอปฺปฏิฆ\-มูลสฺส นวปญฺหเมว.  เหตุยา นว, อธิปติยา นว ฯเปฯ อวิคเต นว.}

\csromanmarkh{14. อาสวทุก}\csromanmarkhii{22. สนิทสฺสน\-ตฺติก}\csromanheadernomark{2}{14. อาสวทุก 22. สนิทสฺสน\-ตฺติก}
\proseitem{94}{นอาสวํ นสนิทสฺสน\-สปฺปฏิฆํ ธมฺมํ ปฏิจฺจ โนอาสโว สนิทสฺสน\-สปฺปฏิโฆ ธมฺโม อุปฺปชฺชติ เหตุปจฺจยา. เอกํ.}

\prose{นโนอาสวํ นสนิทสฺสน\-สปฺปฏิฆํ ธมฺมํ ปฏิจฺจ โนอาสโว สนิทสฺสน\-สปฺปฏิโฆ ธมฺโม อุปฺปชฺชติ เหตุปจฺจยา. เอกํ. คณิตเกน ตีณิ.}

\csromanmarkh{17. อาสว\-สาสวทุก}\csromanmarkhii{22. สนิทสฺสน\-ตฺติก}\csromanheadernomark{2}{17. อาสว\-สาสวทุก 22. สนิทสฺสน\-ตฺติก}
\prose{\csromanfolio{369}โนอาสวํ นอนิทสฺสน\-สปฺปฏิฆํ ธมฺมํ ปฏิจฺจ. (ปุริมนเยน ตีณิ ปญฺหา.)}

\prose{โนอาสวํ นอนิทสฺสน\-อปฺปฏิฆํ ธมฺมํ ปฏิจฺจ โนอาสโว อนิทสฺสน\-อปฺปฏิโฆ ธมฺโม อุปฺปชฺชติ เหตุปจฺจยา.  เหตุยา เอกํ.}

\csromanmarkh{15. สาสวทุก}\csromanmarkhii{22. สนิทสฺสน\-ตฺติก}\csromanheadernomark{2}{15. สาสวทุก 22. สนิทสฺสน\-ตฺติก}
\proseitem{95}{นสาสวํ นสนิทสฺสน\-สปฺปฏิฆํ ธมฺมํ ปฏิจฺจ สาสโว สนิทสฺสน\-สปฺปฏิโฆ ธมฺโม อุปฺปชฺชติ เหตุปจฺจยา. เอกํ.}

\prose{นอนาสวํ นสนิทสฺสน\-สปฺปฏิฆํ ธมฺมํ ปฏิจฺจ สาสโว สนิทสฺสน\-สปฺปฏิโฆ ธมฺโม อุปฺปชฺชติ เหตุปจฺจยา. เอกํ. คณิตเกน ตีณิ.}

\prose{(นสาสวํ นอนิทสฺสนสปฺปฏิฆ\-มูลสฺส ปุริมนเยน ตีณิ ปญฺหา.)}

\prose{นอนาสวํ นอนิทสฺสน\-อปฺปฏิฆํ ธมฺมํ ปฏิจฺจ สาสโว อนิทสฺสน\-อปฺปฏิโฆ ธมฺโม อุปฺปชฺชติ เหตุปจฺจยา.  เหตุยา เอกํ.}

\csromanmarkh{16. อาสว\-สมฺปยุตฺตทุก}\csromanmarkhii{22. สนิทสฺสน\-ตฺติก}\csromanheadernomark{2}{16. อาสว\-สมฺปยุตฺตทุก 22. สนิทสฺสน\-ตฺติก}
\proseitem{96}{นอาสว\-สมฺปยุตฺตํ นสนิทสฺสน\-สปฺปฏิฆํ ธมฺมํ ปฏิจฺจ อาสว\-วิปฺปยุตฺโต สนิทสฺสน\-สปฺปฏิโฆ ธมฺโม อุปฺปชฺชติ เหตุปจฺจยา. เอกํ.}

\prose{นอาสว\-วิปฺปยุตฺตํ นสนิทสฺสน\-สปฺปฏิฆํ ธมฺมํ ปฏิจฺจ อาสว\-วิปฺปยุตฺโต สนิทสฺสน\-สปฺปฏิโฆ ธมฺโม อุปฺปชฺชติ เหตุปจฺจยา. เอกํ. คณิตเกน ตีณิ.}

\prose{(นอาสวสมฺปยุตฺตนอนิทสฺสนสปฺปฏิฆมูเล ตีณิเยว.)}

\prose{นอาสว\-สมฺปยุตฺตํ นอนิทสฺสน\-อปฺปฏิฆํ ธมฺมํ ปฏิจฺจ อาสว\-วิปฺปยุตฺโต อนิทสฺสน\-อปฺปฏิโฆ ธมฺโม อุปฺปชฺชติ เหตุปจฺจยา.  เหตุยา เอกํ.}

\csromanmarkh{17. อาสว\-สาสวทุก}\csromanmarkhii{22. สนิทสฺสน\-ตฺติก}\csromanheadernomark{2}{17. อาสว\-สาสวทุก 22. สนิทสฺสน\-ตฺติก}
\proseitem{97}{นอาสวญฺเจว นอนาสวญฺจ นสนิทสฺสน\-สปฺปฏิฆํ ธมฺมํ ปฏิจฺจ สาสโว เจว โน จ อาสโว สนิทสฺสน\-สปฺปฏิโฆ ธมฺโม อุปฺปชฺชติ เหตุปจฺจยา.}

\csromanheader{2}{ธมฺม\-ปจฺจนียา\-นุโลม ทุกติก\-ปฏฺฐาน\-ปาฬิ}
\prose{\csromanfolio{370}นอนาสวญฺเจว นโน จ อาสวํ นสนิทสฺสน\-สปฺปฏิฆํ ธมฺมํ ปฏิจฺจ สาสโว เจว โน จ อาสโว สนิทสฺสน\-สปฺปฏิโฆ ธมฺโม อุปฺปชฺชติ เหตุปจฺจยา. คณิตเกน ตีณิ.}

\prose{(นอาสวญฺเจว นอนาสวํ นอนิทสฺสนสปฺปฏิฆมูเลปิ ปุริมนเยน ตีณิเยว.)}

\prose{นอาสวญฺเจว นอนาสวญฺจ นอนิทสฺสน\-อปฺปฏิฆํ ธมฺมํ ปฏิจฺจ สาสโว เจว โน จ อาสโว อนิทสฺสน\-อปฺปฏิโฆ ธมฺโม อุปฺปชฺชติ เหตุปจฺจยา.  เหตุยา เอกํ. (อาสวอาสวสมฺปยุตฺตทุเก ปญฺหา นลพฺภติ.)}

\vspace{\tipitakaparskip}
\csromancenter{(อาสว\-อาสว\-สมฺปยุตฺต\-ทุเก ปญฺหา นลพฺภติ.)}
\prose{19. อาสววิปฺปยุตฺตสาสฺอวทุก 22. สนิทสฺสน\-ตฺติก}

\proseitem{98}{อาสว\-วิปฺปยุตฺตํ นสาสวํ นสนิทสฺสน\-สปฺปฏิฆํ ธมฺมํ ปฏิจฺจ อาสว\-วิปฺปยุตฺโต สาสโว สนิทสฺสน\-สปฺปฏิโฆ ธมฺโม อุปฺปชฺชติ เหตุปจฺจยา. เอกํ.}

\prose{อาสว\-วิปฺปยุตฺตํ นอนาสวํ นสนิทสฺสน\-สปฺปฏิฆํ ธมฺมํ ปฏิจฺจ อาสว\-วิปฺปยุตฺโต สาสโว สนิทสฺสน\-สปฺปฏิโฆ ธมฺโม อุปฺปชฺชติ เหตุปจฺจยา. เอกํ. คณิตเกน ตีณิ.}

\prose{(อาสววิปฺปยุตฺตนสาสวนอนิทสฺสนสปฺปฏิฆมูเลปิ ปุริมนเยเนว ตีณิ ปญฺหา.)}

\prose{อาสว\-วิปฺปยุตฺตํ นอนาสวํ นอนิทสฺสน\-อปฺปฏิฆํ ธมฺมํ ปฏิจฺจ อาสว\-วิปฺปยุตฺโต สาสโว อนิทสฺสน\-อปฺปฏิโฆ ธมฺโม อุปฺปชฺชติ เหตุปจฺจยา.  เหตุยา เอกํ.}

\prose{20-54. สญฺโญชนา\-ทิ\-ฉ\-โคจฺฉก 22. สนิทสฺสน\-ตฺติก}

\proseitem{99}{โนสญฺโญชนํ นสนิทสฺสน\-สปฺปฏิฆํ ธมฺมํ ปฏิจฺจ โนสญฺโญชโน สนิทสฺสน\-สปฺปฏิโฆ ธมฺโม อุปฺปชฺชติ เหตุปจฺจยา. (ตีณิ ปญฺหา.)}

\proseitem{100}{โนคนฺถํ นสนิทสฺสน\-สปฺปฏิฆํ ธมฺมํ ปฏิจฺจ โนคนฺโถ สนิทสฺสน\-สปฺปฏิโฆ ธมฺโม อุปฺปชฺชติ เหตุปจฺจยา.  เหตุยา ตีณิ.}

\csromanmarkh{55-68. มหนฺตรทุก}\csromanmarkhii{22. สนิทสฺสน\-ตฺติก}\csromanheadernomark{2}{55-68. มหนฺตรทุก 22. สนิทสฺสน\-ตฺติก}
\proseitem{101}{\csromanfolio{371}โนโอฆํ นสนิทสฺสน\-สปฺปฏิฆํ ธมฺมํ ปฏิจฺจ โนโอโฆ สนิทสฺสน\-สปฺปฏิโฆ ธมฺโม อุปฺปชฺชติ เหตุปจฺจยา.  เหตุยา ตีณิ.}

\proseitem{102}{โนโยคํ นสนิทสฺสน\-สปฺปฏิฆํ ธมฺมํ ปฏิจฺจ โนโยโค สนิทสฺสน\-สปฺปฏิโฆ ธมฺโม อุปฺปชฺชติ เหตุปจฺจยา.  เหตุยา ตีณิ.}

\proseitem{103}{โนนีวรณํ นสนิทสฺสน\-สปฺปฏิฆํ ธมฺมํ ปฏิจฺจ โนนีวรโณ สนิทสฺสน\-สปฺปฏิโฆ ธมฺโม อุปฺปชฺชติ เหตุปจฺจยา.  เหตุยา ตีณิ.}

\proseitem{104}{โนปรามาสํ นสนิทสฺสน\-สปฺปฏิฆํ ธมฺมํ ปฏิจฺจ โนปรามาโส สนิทสฺสน\-สปฺปฏิโฆ ธมฺโม อุปฺปชฺชติ เหตุปจฺจยา.  เหตุยา ตีณิ.}

\csromanmarkh{55-68. มหนฺตรทุก}\csromanmarkhii{22. สนิทสฺสน\-ตฺติก}\csromanheadernomark{2}{55-68. มหนฺตรทุก 22. สนิทสฺสน\-ตฺติก}
\proseitem{105}{นสารมฺมณํ นสนิทสฺสน\-สปฺปฏิฆํ ธมฺมํ ปฏิจฺจ อนารมฺมโณ สนิทสฺสน\-สปฺปฏิโฆ ธมฺโม อุปฺปชฺชติ เหตุปจฺจยา. เอกํ.}

\prose{นอนารมฺมณํ นสนิทสฺสน\-สปฺปฏิฆํ ธมฺมํ ปฏิจฺจ อนารมฺมโณ สนิทสฺสน\-สปฺปฏิโฆ ธมฺโม อุปฺปชฺชติ เหตุปจฺจยา. เอกํ. คณิตเกน ตีณิ.}

\prose{(อนิทสฺสน\-สปฺปฏิเฆ ตีณิเยว.)}

\prose{นสารมฺมณํ นอนิทสฺสน\-อปฺปฏิฆํ ธมฺมํ ปฏิจฺจ อนารมฺมโณ อนิทสฺสน\-อปฺปฏิโฆ ธมฺโม อุปฺปชฺชติ เหตุปจฺจยา.  เหตุยา เอกํ.}

\proseitem{106}{โนจิตฺตํ นสนิทสฺสน\-สปฺปฏิฆํ ธมฺมํ ปฏิจฺจ โนจิตฺโต สนิทสฺสน\-สปฺปฏิโฆ ธมฺโม อุปฺปชฺชติ เหตุปจฺจยา. เอกํ.}

\prose{นโนจิตฺตํ นสนิทสฺสน\-สปฺปฏิฆํ ธมฺมํ ปฏิจฺจ โนจิตฺโต สนิทสฺสน\-สปฺปฏิโฆ ธมฺโม อุปฺปชฺชติ เหตุปจฺจยา. เอกํ. คณิตเกน ตีณิ.}

\proseitem{107}{นเจตสิกํ นสนิทสฺสน\-สปฺปฏิฆํ ธมฺมํ ปฏิจฺจ นเจตสิโก สนิทสฺสน\-สปฺปฏิโฆ ธมฺโม อุปฺปชฺชติ เหตุปจฺจยา.  เหตุยา ตีณิ.}

\proseitem{108}{นจิตฺต\-สมฺปยุตฺตํ นสนิทสฺสน\-สปฺปฏิฆํ ธมฺมํ ปฏิจฺจ จิตฺต\-วิปฺปยุตฺโต สนิทสฺสน\-สปฺปฏิโฆ ธมฺโม อุปฺปชฺชติ เหตุปจฺจยา.  เหตุยา ตีณิ.}

\csromanheader{2}{ธมฺม\-ปจฺจนียา\-นุโลม ทุกติก\-ปฏฺฐาน\-ปาฬิ}
\proseitem{109}{\csromanfolio{372}นจิตฺต\-สํสฏฺฐํ นสนิทสฺสน\-สปฺปฏิฆํ ธมฺมํ ปฏิจฺจ จิตฺต\-วิสํสฏฺโฐ สนิทสฺสน\-สปฺปฏิโฆ ธมฺโม อุปฺปชฺชติ เหตุปจฺจยา.  เหตุยา ตีณิ.}

\proseitem{110}{โนจิตฺต\-สมุฏฺฐานํ นสนิทสฺสน\-สปฺปฏิฆํ ธมฺมํ ปฏิจฺจ จิตฺต\-สมุฏฺฐาโน สนิทสฺสน\-สปฺปฏิโฆ ธมฺโม อุปฺปชฺชติ เหตุปจฺจยา. (จิตฺตสมุฏฺฐาน\-รูเปเนว ตีณิ. สํขิตฺตํ.)}

\proseitem{111}{โนจิตฺต\-สหภุํ นสนิทสฺสน\-สปฺปฏิฆํ ธมฺมํ ปฏิจฺจ โนจิตฺตสหภู สนิทสฺสน\-สปฺปฏิโฆ ธมฺโม อุปฺปชฺชติ เหตุปจฺจยา.  เหตุยา ตีณิ.}

\proseitem{112}{โนจิตฺตา\-นุปริวตฺติํ นสนิทสฺสน\-สปฺปฏิฆํ ธมฺมํ ปฏิจฺจ โนจิตฺตา\-นุปริวตฺตี สนิทสฺสน\-สปฺปฏิโฆ ธมฺโม อุปฺปชฺชติ เหตุปจฺจยา.  เหตุยา ตีณิ.}

\proseitem{113}{โนจิตฺต\-สํสฏฺฐ\-สมุฏฺฐานํ นสนิทสฺสน\-สปฺปฏิฆํ ธมฺมํ ปฏิจฺจ โนจิตฺต\-สํสฏฺฐ\-สมุฏฺฐาโน สนิทสฺสน\-สปฺปฏิโฆ ธมฺโม อุปฺปชฺชติ เหตุปจฺจยา.  เหตุยา ตีณิ.}

\proseitem{114}{โนจิตฺต\-สํสฏฺฐ\-สมุฏฺฐาน\-สหภุํ นสนิทสฺสน\-สปฺปฏิฆํ ธมฺมํ ปฏิจฺจ โนจิตฺต\-สํสฏฺฐ\-สมุฏฺฐาน\-สหภู สนิทสฺสน\-สปฺปฏิโฆ ธมฺโม อุปฺปชฺชติ เหตุปจฺจยา.  เหตุยา ตีณิ.}

\proseitem{115}{โนจิตฺต\-สํสฏฺฐ\-สมุฏฺฐานา\-นุปริวตฺติํ นสนิทสฺสน\-สปฺปฏิฆํ ธมฺมํ ปฏิจฺจ โนจิตฺต\-สํสฏฺฐ\-สมุฏฺฐานา\-นุปริวตฺตี สนิทสฺสน\-สปฺปฏิโฆ ธมฺโม อุปฺปชฺชติ เหตุปจฺจยา.  เหตุยา ตีณิ.}

\proseitem{116}{นอชฺฌตฺติกํ นสนิทสฺสน\-สปฺปฏิฆํ ธมฺมํ ปฏิจฺจ พาหิโร สนิทสฺสน\-สปฺปฏิโฆ ธมฺโม อุปฺปชฺชติ เหตุปจฺจยา. เอกํ.}

\prose{นพาหิรํ นสนิทสฺสน\-สปฺปฏิฆํ ธมฺมํ ปฏิจฺจ พาหิโร สนิทสฺสน\-สปฺปฏิโฆ ธมฺโม อุปฺปชฺชติ เหตุปจฺจยา. เอกํ. คณิตเกน ตีณิ.}

\proseitem{117}{โนอุปาทา นสนิทสฺสน\-สปฺปฏิฆํ ธมฺมํ ปฏิจฺจ อุปาทา สนิทสฺสน\-สปฺปฏิโฆ ธมฺโม อุปฺปชฺชติ เหตุปจฺจยา.  เหตุยา เอกํ.}

\csromanmarkh{83-100. ปิฏฺฐิทุก}\csromanmarkhii{22. สนิทสฺสน\-ตฺติก}\csromanheadernomark{2}{83-100. ปิฏฺฐิทุก 22. สนิทสฺสน\-ตฺติก}
\proseitem{118}{\csromanfolio{373}นอุปาทินฺนํ นสนิทสฺสน\-สปฺปฏิฆํ ธมฺมํ ปฏิจฺจ อนุปาทินฺโน สนิทสฺสน\-สปฺปฏิโฆ ธมฺโม อุปฺปชฺชติ เหตุปจฺจยา.  เหตุยา ตีณิ.}

\csromanmarkh{69. อุปาทานทุก}\csromanmarkhii{22. สนิทสฺสน\-ตฺติก}\csromanheadernomark{2}{69. อุปาทานทุก 22. สนิทสฺสน\-ตฺติก}
\proseitem{119}{โนอุปาทานํ นสนิทสฺสน\-สปฺปฏิฆํ ธมฺมํ ปฏิจฺจ โนอุปาทาโน สนิทสฺสน\-สปฺปฏิโฆ ธมฺโม อุปฺปชฺชติ เหตุปจฺจยา.  เหตุยา ตีณิ.}

\csromanmarkh{75. กิเลสทุก}\csromanmarkhii{22. สนิทสฺสน\-ตฺติก}\csromanheadernomark{2}{75. กิเลสทุก 22. สนิทสฺสน\-ตฺติก}
\proseitem{120}{โนกิเลสํ นสนิทสฺสน\-สปฺปฏิฆํ ธมฺมํ ปฏิจฺจ โนกิเลโส สนิทสฺสน\-สปฺปฏิโฆ ธมฺโม อุปฺปชฺชติ เหตุปจฺจยา.  เหตุยา ตีณิ.}

\csromanmarkh{83-100. ปิฏฺฐิทุก}\csromanmarkhii{22. สนิทสฺสน\-ตฺติก}\csromanheadernomark{2}{83-100. ปิฏฺฐิทุก 22. สนิทสฺสน\-ตฺติก}
\proseitem{121}{นทสฺสเนน ปหาตพฺพํ นสนิทสฺสน\-สปฺปฏิฆํ ธมฺมํ ปฏิจฺจ นทสฺสเนน ปหาตพฺโพ สนิทสฺสน\-สปฺปฏิโฆ ธมฺโม อุปฺปชฺชติ เหตุปจฺจยา.  เหตุยา ตีณิ.}

\proseitem{122}{นภาวนาย ปหาตพฺพํ นสนิทสฺสน\-สปฺปฏิฆํ ธมฺมํ ปฏิจฺจ นภาวนาย ปหาตพฺโพ สนิทสฺสน\-สปฺปฏิโฆ ธมฺโม อุปฺปชฺชติ เหตุปจฺจยา.  เหตุยา ตีณิ.}

\proseitem{123}{นทสฺสเนน ปหาตพฺพ\-เหตุกํ นสนิทสฺสน\-สปฺปฏิฆํ ธมฺมํ ปฏิจฺจ นทสฺสเนน ปหาตพฺพ\-เหตุโก สนิทสฺสน\-สปฺปฏิโฆ ธมฺโม อุปฺปชฺชติ เหตุปจฺจยา.  เหตุยา ตีณิ.}

\proseitem{124}{นภาวนาย ปหาตพฺพ\-เหตุกํ นสนิทสฺสน\-สปฺปฏิฆํ ธมฺมํ ปฏิจฺจ นภาวนาย ปหาตพฺพ\-เหตุโก สนิทสฺสน\-สปฺปฏิโฆ ธมฺโม อุปฺปชฺชติ เหตุปจฺจยา.  เหตุยา ตีณิ.}

\proseitem{125}{นสวิตกฺกํ นสนิทสฺสน\-สปฺปฏิฆํ ธมฺมํ ปฏิจฺจ อวิตกฺโก สนิทสฺสน\-สปฺปฏิโฆ ธมฺโม อุปฺปชฺชติ เหตุปจฺจยา.  เหตุยา ตีณิ.}

\proseitem{126}{นสวิจารํ นสนิทสฺสน\-สปฺปฏิฆํ ธมฺมํ ปฏิจฺจ อวิจาโร สนิทสฺสน\-สปฺปฏิโฆ ธมฺโม อุปฺปชฺชติ เหตุปจฺจยา.  เหตุยา ตีณิ.}

\csromanheader{2}{ธมฺม\-ปจฺจนียา\-นุโลม ทุกติก\-ปฏฺฐาน\-ปาฬิ}
\proseitem{127}{\csromanfolio{374}นสปฺปีติกํ นสนิทสฺสน\-สปฺปฏิฆํ ธมฺมํ ปฏิจฺจ นสปฺปีติโก สนิทสฺสน\-สปฺปฏิโฆ ธมฺโม อุปฺปชฺชติ เหตุปจฺจยา.  เหตุยา ตีณิ.}

\proseitem{128}{นปีติสหคตํ นสนิทสฺสน\-สปฺปฏิฆํ ธมฺมํ ปฏิจฺจ นปีติสหคโต สนิทสฺสน\-สปฺปฏิโฆ ธมฺโม อุปฺปชฺชติ เหตุปจฺจยา.  เหตุยา ตีณิ.}

\proseitem{129}{นสุขสหคตํ นสนิทสฺสน\-สปฺปฏิฆํ ธมฺมํ ปฏิจฺจ นสุขสหคโต สนิทสฺสน\-สปฺปฏิโฆ ธมฺโม อุปฺปชฺชติ เหตุปจฺจยา.  เหตุยา ตีณิ.}

\proseitem{130}{นอุเปกฺขาสหคตํ นสนิทสฺสน\-สปฺปฏิฆํ ธมฺมํ ปฏิจฺจ นอุเปกฺขาสหคโต สนิทสฺสน\-สปฺปฏิโฆ ธมฺโม อุปฺปชฺชติ เหตุปจฺจยา.  เหตุยา ตีณิ.}

\proseitem{131}{นกามาวจรํ นสนิทสฺสน\-สปฺปฏิฆํ ธมฺมํ ปฏิจฺจ กามาวจโร สนิทสฺสน\-สปฺปฏิโฆ ธมฺโม อุปฺปชฺชติ เหตุปจฺจยา.  เหตุยา ตีณิ.}

\proseitem{132}{นรูปาวจรํ นสนิทสฺสน\-สปฺปฏิฆํ ธมฺมํ ปฏิจฺจ นรูปาวจโร สนิทสฺสน\-สปฺปฏิโฆ ธมฺโม อุปฺปชฺชติ เหตุปจฺจยา.  เหตุยา ตีณิ.}

\proseitem{133}{นอรูปาวจรํ นสนิทสฺสน\-สปฺปฏิฆํ ธมฺมํ ปฏิจฺจ นอรูปาวจโร สนิทสฺสน\-สปฺปฏิโฆ ธมฺโม อุปฺปชฺชติ เหตุปจฺจยา.  เหตุยา ตีณิ.}

\proseitem{134}{นปริยาปนฺนํ นสนิทสฺสน\-สปฺปฏิฆํ ธมฺมํ ปฏิจฺจ ปริยาปนฺโน สนิทสฺสน\-สปฺปฏิโฆ ธมฺโม อุปฺปชฺชติ เหตุปจฺจยา.  เหตุยา ตีณิ.}

\proseitem{135}{นนิยฺยานิกํ นสนิทสฺสน\-สปฺปฏิฆํ ธมฺมํ ปฏิจฺจ อนิยฺยานิโก สนิทสฺสน\-สปฺปฏิโฆ ธมฺโม อุปฺปชฺชติ เหตุปจฺจยา.  เหตุยา ตีณิ.}

\proseitem{136}{นนิยตํ นสนิทสฺสน\-สปฺปฏิฆํ ธมฺมํ ปฏิจฺจ อนิยโต สนิทสฺสน\-สปฺปฏิโฆ ธมฺโม อุปฺปชฺชติ เหตุปจฺจยา.  เหตุยา ตีณิ.}

\proseitem{137}{นสอุตฺตรํ นสนิทสฺสน\-สปฺปฏิฆํ ธมฺมํ ปฏิจฺจ สอุตฺตโร สนิทสฺสน\-สปฺปฏิโฆ ธมฺโม อุปฺปชฺชติ เหตุปจฺจยา. เหตุยา ตีณิ.}

\proseitem{138}{นสรณํ นสนิทสฺสน\-สปฺปฏิฆํ ธมฺมํ ปฏิจฺจ อรโณ สนิทสฺสน\-สปฺปฏิโฆ ธมฺโม อุปฺปชฺชติ เหตุปจฺจยา. เอกํ.}

\csromanmarkh{83-100. ปิฏฺฐิทุก}\csromanmarkhii{22. สนิทสฺสน\-ตฺติก}\csromanheadernomark{2}{83-100. ปิฏฺฐิทุก 22. สนิทสฺสน\-ตฺติก}
\prose{\csromanfolio{375}นอรณํ นสนิทสฺสน\-สปฺปฏิฆํ ธมฺมํ ปฏิจฺจ อรโณ สนิทสฺสน\-สปฺปฏิโฆ ธมฺโม อุปฺปชฺชติ เหตุปจฺจยา. เอกํ.}

\prose{นสรณํ นสนิทสฺสนสปฺปฏิฆญฺจ นอรณํ นสนิทสฺสนสปฺปฏิฆญฺจ ธมฺมํ ปฏิจฺจ อรโณ สนิทสฺสน\-สปฺปฏิโฆ ธมฺโม อุปฺปชฺชติ เหตุปจฺจยา.  เหตุยา ตีณิ.  (อนิทสฺสน\-สปฺปฏิเฆ ตีณิเยว.)}

\prose{นสรณํ นอนิทสฺสน\-อปฺปฏิฆํ ธมฺมํ ปฏิจฺจ อรโณ อนิทสฺสน\-อปฺปฏิโฆ ธมฺโม อุปฺปชฺชติ เหตุปจฺจยา.  เหตุยา เอกํ.}

\prose{ปจฺจนีย\-น\-เหตุ นอารมฺมณ}

\proseitem{139}{นสรณํ นสนิทสฺสน\-สปฺปฏิฆํ ธมฺมํ ปฏิจฺจ อรโณ สนิทสฺสน\-สปฺปฏิโฆ ธมฺโม อุปฺปชฺชติ นเหตุปจฺจยา. เอกํ.}

\prose{นสรณํ นสนิทสฺสน\-สปฺปฏิฆํ ธมฺมํ ปฏิจฺจ อรโณ สนิทสฺสน\-สปฺปฏิโฆ ธมฺโม อุปฺปชฺชติ นอารมฺมณ\-ปจฺจยา. ตีณิ.  นเหตุยา เอกํ, นอารมฺมเณ ตีณิ, นอธิปติยา ตีณิ ฯเปฯ โนวิคเต ตีณิ.}

\prose{เหตุปจฺจยา นอารมฺมเณ ตีณิ.  นอารมฺมณ\-ปจฺจยา เหตุยา ตีณิ.}

\prose{(สหชาตวารมฺปิ ปจฺจยวารมฺปิ นิสฺสยวารมฺปิ สํสฏฺฐ\-วารมฺปิ สมฺปยุตฺตวารมฺปิ วิตฺถาเรตพฺพํ.)}

\prose{ปญฺหาวารเหตุ}

\proseitem{140}{นสรโณ นสนิทสฺสน\-สปฺปฏิโฆ ธมฺโม อรณสฺส สนิทสฺสน\-สปฺปฏิฆสฺส ธมฺมสฺส เหตุปจฺจเยน ปจฺจโย.  นอรโณ นสนิทสฺสน\-สปฺปฏิโฆ ธมฺโม อรณสฺส สนิทสฺสน\-สปฺปฏิฆสฺส ธมฺมสฺส เหตุปจฺจเยน ปจฺจโย.}

\prose{เหตุยา เทฺว, อธิปติยา เทฺว ฯเปฯ อวิคเต ตีณิ.}

\prosewithcloser{(ยถา กุสลตฺติเก ปญฺหาวารสฺส อนุโลมมฺปิ ปจฺจนียมฺปิ อนุโลมปจฺจนียมฺปิ ปจฺจนียานุโลมมฺปิ คณิตํ, เอวํ คเณตพฺพํ.)}{ธมฺม\-ปจฺจนียา\-นุโลเม ทุกติก\-ปฏฺฐานํ นิฏฺฐิตํ.}
\pitaka{อภิธมฺม\-ปิฏก}
\csromanheader{2}{ธมฺม\-ปจฺจนียา\-นุโลม}
\namakkaram{นโม ตสฺส ภควโต อรหโต สมฺมา\-สมฺพุทฺธสฺส.}
\csromanmarkh{1. กุสลตฺติก}\csromanmarkhii{1. เหตุทุก}\csromanheadernomark{2}{1. \textbf{กุสลตฺติก 1. เหตุทุก}}
\proseitem{1}{\csromanfolio{377}นกุสลํ นเหตุํ ธมฺมํ ปฏิจฺจ อกุสโล เหตุ ธมฺโม อุปฺปชฺชติ เหตุปจฺจยา.  นกุสลํ นเหตุํ ธมฺมํ ปฏิจฺจ อพฺยากโต เหตุ ธมฺโม อุปฺปชฺชติ เหตุปจฺจยา. (2)}

\prose{นอกุสลํ นเหตุํ ธมฺมํ ปฏิจฺจ กุสโล เหตุ ธมฺโม อุปฺปชฺชติ เหตุปจฺจยา.  นอกุสลํ นเหตุํ ธมฺมํ ปฏิจฺจ อพฺยากโต เหตุ ธมฺโม อุปฺปชฺชติ เหตุปจฺจยา. (2)}

\prose{นอพฺยากตํ นเหตุํ ธมฺมํ ปฏิจฺจ กุสโล เหตุ ธมฺโม อุปฺปชฺชติ เหตุปจฺจยา.  นอพฺยากตํ นเหตุํ ธมฺมํ ปฏิจฺจ อกุสโล เหตุ ธมฺโม อุปฺปชฺชติ เหตุปจฺจยา. (2)}

\prose{นกุสลํ นเหตุญฺจ นอพฺยากตํ นเหตุญฺจ ธมฺมํ ปฏิจฺจ อกุสโล เหตุ ธมฺโม อุปฺปชฺชติ เหตุปจฺจยา.  นอกุสลํ นเหตุญฺจ นอพฺยากตํ นเหตุญฺจ ธมฺมํ ปฏิจฺจ กุสโล เหตุ ธมฺโม อุปฺปชฺชติ เหตุปจฺจยา.  นกุสลํ นเหตุญฺจ นอกุสลํ นเหตุญฺจ ธมฺมํ ปฏิจฺจ อพฺยากโต เหตุ ธมฺโม อุปฺปชฺชติ เหตุปจฺจยา. ตีณิ. (สํขิตฺตํ.) . เหตุยา นว.}

\proseitem{2}{นกุสลํ นนเหตุํ ธมฺมํ ปฏิจฺจ อกุสโล นเหตุ ธมฺโม อุปฺปชฺชติ เหตุปจฺจยา.  นกุสลํ นนเหตุํ ธมฺมํ ปฏิจฺจ อพฺยากโต นเหตุ}

\vspace{\tipitakaparskip}
\csromancenter{ธมฺม\-ปจฺจนียา\-นุโลม ติกทุก\-ปฏฺฐาน\-ปาฬิ ธมฺโม อุปฺปชฺชติ เหตุปจฺจยา.  นกุสลํ นนเหตุํ ธมฺมํ ปฏิจฺจ อกุสโล นเหตุ จ อพฺยากโต นเหตุ จ ธมฺมา อุปฺปชฺชนฺติ เหตุปจฺจยา. ตีณิ.}
\prose{\csromanfolio{378}นอกุสลํ นนเหตุํ ธมฺมํ ปฏิจฺจ กุสโล นเหตุ ธมฺโม อุปฺปชฺชติ เหตุปจฺจยา.  นอกุสลํ นนเหตุํ ธมฺมํ ปฏิจฺจ อพฺยากโต นเหตุ ธมฺโม อุปฺปชฺชติ เหตุปจฺจยา.  นอกุสลํ นนเหตุํ ธมฺมํ ปฏิจฺจ กุสโล นเหตุ จ อพฺยากโต นเหตุ จ ธมฺมา อุปฺปชฺชนฺติ เหตุปจฺจยา. ตีณิ.}

\prose{นอพฺยากตํ นนเหตุํ ธมฺมํ ปฏิจฺจ อพฺยากโต นเหตุ ธมฺโม อุปฺปชฺชติ เหตุปจฺจยา. ปญฺจ.}

\prose{นกุสลํ นนเหตุญฺจ นอพฺยากตํ นนเหตุญฺจ ธมฺมํ ปฏิจฺจ อกุสโล นเหตุ ธมฺโม อุปฺปชฺชติ เหตุปจฺจยา. ตีณิ.}

\prose{นอกุสลํ นนเหตุญฺจ นอพฺยากตํ นนเหตุญฺจ ธมฺมํ ปฏิจฺจ กุสโล นเหตุ ธมฺโม อุปฺปชฺชติ เหตุปจฺจยา. ตีณิ.}

\prose{นกุสลํ นนเหตุญฺจ นอกุสลํ นนเหตุญฺจ ธมฺมํ ปฏิจฺจ อพฺยากโต นเหตุ ธมฺโม อุปฺปชฺชติ เหตุปจฺจยา. เอกํ.  เหตุยา อฏฺฐารส.}

\csromanmarkh{1. กุสลตฺติก}\csromanmarkhii{2. สเหตุกทุก}\csromanheadernomark{2}{1. \textbf{กุสลตฺติก 2. สเหตุกทุก}}
\proseitem{3}{นกุสลํ นสเหตุกํ ธมฺมํ ปฏิจฺจ อกุสโล สเหตุโก ธมฺโม อุปฺปชฺชติ เหตุปจฺจยา.  นกุสลํ นสเหตุกํ ธมฺมํ ปฏิจฺจ อพฺยากโต สเหตุโก ธมฺโม อุปฺปชฺชติ เหตุปจฺจยา.  นอกุสลํ นสเหตุกํ ธมฺมํ ปฏิจฺจ อพฺยากโต สเหตุโก ธมฺโม อุปฺปชฺชติ เหตุปจฺจยา.  นอพฺยากตํ นสเหตุกํ ธมฺมํ ปฏิจฺจ อกุสโล สเหตุโก ธมฺโม อุปฺปชฺชติ เหตุปจฺจยา.  นกุสลํ นสเหตุกญฺจ นอพฺยากตํ นสเหตุกญฺจ ธมฺมํ ปฏิจฺจ อกุสโล สเหตุโก ธมฺโม อุปฺปชฺชติ เหตุปจฺจยา.  นกุสลํ นสเหตุกญฺจ นอกุสลํ นสเหตุกญฺจ ธมฺมํ ปฏิจฺจ อพฺยากโต สเหตุโก ธมฺโม อุปฺปชฺชติ เหตุปจฺจยา.  เหตุยา ฉ.}

\proseitem{4}{นกุสลํ นอเหตุกํ ธมฺมํ ปฏิจฺจ อพฺยากโต อเหตุโก ธมฺโม อุปฺปชฺชติ เหตุปจฺจยา.  นอกุสลํ นอเหตุกํ ธมฺมํ ปฏิจฺจ อพฺยากโต อเหตุโก ธมฺโม อุปฺปชฺชติ เหตุปจฺจยา.}

\vspace{\tipitakaparskip}
\csromancenter{4-6. เหตุ\-สเหตุก\-ทุกาทิ นอพฺยากตํ นอเหตุกํ ธมฺมํ ปฏิจฺจ อพฺยากโต อเหตุโก ธมฺโม อุปฺปชฺชติ เหตุปจฺจยา.  นกุสลํ นอเหตุกญฺจ นอพฺยากตํ นอเหตุกญฺจ ธมฺมํ ปฏิจฺจ อพฺยากโต อเหตุโก ธมฺโม อุปฺปชฺชติ เหตุปจฺจยา.  นอกุสลํ นอเหตุกญฺจ นอพฺยากตํ นอเหตุกญฺจ ธมฺมํ ปฏิจฺจ อพฺยากโต อเหตุโก ธมฺโม อุปฺปชฺชติ เหตุปจฺจยา.  นกุสลํ นอเหตุกญฺจ นอกุสลํ นอเหตุกญฺจ ธมฺมํ ปฏิจฺจ อพฺยากโต อเหตุโก ธมฺโม อุปฺปชฺชติ เหตุปจฺจยา.  เหตุยา ฉ.}
\csromanmarkh{1. กุสลตฺติก}\csromanmarkhii{3. เหตุ\-สมฺปยุตฺตทุก}\csromanheadernomark{2}{1. กุสลตฺติก 3. เหตุ\-สมฺปยุตฺตทุก}
\proseitem{5}{\csromanfolio{379}นกุสลํ นเหตุ\-สมฺปยุตฺตํ ธมฺมํ ปฏิจฺจ อกุสโล เหตุสมฺปยุตฺโต ธมฺโม อุปฺปชฺชติ เหตุปจฺจยา.  นกุสลํ นเหตุ\-สมฺปยุตฺตํ ธมฺมํ ปฏิจฺจ อพฺยากโต เหตุสมฺปยุตฺโต ธมฺโม อุปฺปชฺชติ เหตุปจฺจยา.  นอกุสลํ นเหตุ\-สมฺปยุตฺตํ ธมฺมํ ปฏิจฺจ อพฺยากโต เหตุสมฺปยุตฺโต ธมฺโม อุปฺปชฺชติ เหตุปจฺจยา. ฉ. (สเหตุก\-สทิสํ.)}

\prose{นกุสลํ นเหตุ\-วิปฺปยุตฺตํ ธมฺมํ ปฏิจฺจ อพฺยากโต เหตุวิปฺปยุตฺโต ธมฺโม อุปฺปชฺชติ เหตุปจฺจยา. ฉ. (อเหตุกสทิสํ. สํขิตฺตํ.)}

\csromanheader{2}{1. กุสลตฺติก 4-6. เหตุ\-สเหตุก\-ทุกาทิ}
\proseitem{6}{นกุสลํ นเหตุญฺเจว นอเหตุกญฺจ ธมฺมํ ปฏิจฺจ อกุสโล เหตุ เจว สเหตุโก จ ธมฺโม อุปฺปชฺชติ เหตุปจฺจยา.  นกุสลํ นเหตุญฺเจว นอเหตุกญฺจ ธมฺมํ ปฏิจฺจ อพฺยากโต เหตุ เจว สเหตุโก จ ธมฺโม อุปฺปชฺชติ เหตุปจฺจยา. เทฺว.}

\prose{(นอกุสเล เทฺว. นอพฺยากเต เทฺว. ปฐมํ คณิตเกน เอกํ, ทุติยํ คณิตเกน เอกํ, ตติยํ คณิตเกน เอกํ, สพฺพตฺถ นว ปญฺหา.)}

\prose{นกุสลํ นอเหตุกญฺเจว น จ เหตุํ ธมฺมํ ปฏิจฺจ อกุสโล สเหตุโก เจว น จ เหตุ ธมฺโม อุปฺปชฺชติ เหตุปจฺจยา.  เหตุยา นว.}

\proseitem{7}{นกุสลํ นเหตุญฺเจว นเหตุ\-วิปฺปยุตฺตญฺจ ธมฺมํ ปฏิจฺจ อกุสโล เหตุ เจว เหตุสมฺปยุตฺโต จ ธมฺโม อุปฺปชฺชติ เหตุปจฺจยา.  เหตุยา นว.}

\csromanheader{2}{ธมฺม\-ปจฺจนียา\-นุโลม ติกทุก\-ปฏฺฐาน\-ปาฬิ}
\prose{\csromanfolio{380}นกุสลํ นเหตุ\-วิปฺปยุตฺต\-ญฺเจว นนเหตุญฺจ ธมฺมํ ปฏิจฺจ อกุสโล เหตุสมฺปยุตฺโต เจว น จ เหตุ ธมฺโม อุปฺปชฺชติ เหตุปจฺจยา.  เหตุยา นว.}

\proseitem{8}{นกุสลํ นเหตุํ นสเหตุกํ ธมฺมํ ปฏิจฺจ อพฺยากโต นเหตุ สเหตุโก ธมฺโม อุปฺปชฺชติ เหตุปจฺจยา.  เหตุยา ตีณิ.}

\proseitem{9}{นกุสลํ นเหตุํ นอเหตุกํ ธมฺมํ ปฏิจฺจ อกุสโล นเหตุ อเหตุโก ธมฺโม อุปฺปชฺชติ เหตุปจฺจยา. (1)}

\prose{นอกุสลํ นเหตุํ นอเหตุกํ ธมฺมํ ปฏิจฺจ อพฺยากโต นเหตุ อเหตุโก ธมฺโม อุปฺปชฺชติ เหตุปจฺจยา. (1)}

\prose{นอพฺยากตํ นเหตุํ นอเหตุกํ ธมฺมํ ปฏิจฺจ อพฺยากโต นเหตุ อเหตุโก ธมฺโม อุปฺปชฺชติ เหตุปจฺจยา. (1) (คณิตเก ตีณิ ปญฺหา.) . เหตุยา ฉ.}

\csromanheader{2}{1. \textbf{กุสลตฺติก} 7-13. จูฬนฺตรทุก}
\proseitem{10}{นกุสโล นสปฺปจฺจโย ธมฺโม กุสลสฺส สปฺปจฺจยสฺส ธมฺมสฺส อารมฺมณ\-ปจฺจเยน ปจฺจโย.  อารมฺมเณ ฉ. (สงฺขตํ สปฺปจฺจย\-สทิสํ.)}

\proseitem{11}{นกุสลํ นสนิทสฺสนํ ธมฺมํ ปฏิจฺจ อพฺยากโต สนิทสฺสโน ธมฺโม อุปฺปชฺชติ เหตุปจฺจยา.  เหตุยา ฉ.}

\prose{นกุสโล นอนิทสฺสโน ธมฺโม กุสลสฺส อนิทสฺสนสฺส ธมฺมสฺส อารมฺมณ\-ปจฺจเยน ปจฺจโย. (นว ปญฺหา กาตพฺพา.)}

\proseitem{12}{นกุสลํ นสปฺปฏิฆํ ธมฺมํ ปฏิจฺจ อพฺยากโต สปฺปฏิโฆ ธมฺโม อุปฺปชฺชติ เหตุปจฺจยา.  เหตุยา ฉ.}

\prose{นกุสลํ นอปฺปฏิฆํ ธมฺมํ ปฏิจฺจ อพฺยากโต อปฺปฏิโฆ ธมฺโม อุปฺปชฺชติ เหตุปจฺจยา.  เหตุยา ตีณิ.}

\proseitem{13}{นกุสลํ นรูปิํ ธมฺมํ ปฏิจฺจ อพฺยากโต รูปี ธมฺโม อุปฺปชฺชติ เหตุปจฺจยา.  เหตุยา ฉ.}

\csromanmarkh{1. กุสลตฺติก}\csromanmarkhii{14. อาสวโคจฺฉก}\csromanheadernomark{2}{1. กุสลตฺติก 14. อาสวโคจฺฉก}
\prose{\csromanfolio{381}นกุสลํ นอรูปิํ ธมฺมํ ปฏิจฺจ อพฺยากโต อรูปี ธมฺโม อุปฺปชฺชติ เหตุปจฺจยา.  เหตุยา ตีณิ.}

\proseitem{14}{นกุสลํ นโลกิยํ ธมฺมํ ปฏิจฺจ อพฺยากโต โลกิโย ธมฺโม อุปฺปชฺชติ เหตุปจฺจยา.  เหตุยา ปญฺจ.}

\prose{นกุสลํ นโลกุตฺตรํ ธมฺมํ ปจฺจยา กุสโล โลกุตฺตโร ธมฺโม อุปฺปชฺชติ เหตุปจฺจยา.  นกุสลํ นโลกุตฺตรํ ธมฺมํ ปจฺจยา อพฺยากโต โลกุตฺตโร ธมฺโม อุปฺปชฺชติ เหตุปจฺจยา. (ฉ ปญฺหา กาตพฺพา.)}

\proseitem{15}{นกุสลํ นเกนจิ วิญฺเญยฺยํ ธมฺมํ ปฏิจฺจ อกุสโล เกนจิ วิญฺเญยฺโย ธมฺโม อุปฺปชฺชติ เหตุปจฺจยา.  เหตุยา อฏฺฐารส.}

\prose{นกุสลํ นเกนจิ นวิญฺเญยฺยํ ธมฺมํ ปฏิจฺจ อกุสโล เกนจิ วิญฺเญยฺโย ธมฺโม อุปฺปชฺชติ เหตุปจฺจยา.  เหตุยา อฏฺฐารส.}

\csromanmarkh{1. กุสลตฺติก}\csromanmarkhii{14. อาสวโคจฺฉก}\csromanheadernomark{2}{1. กุสลตฺติก 14. อาสวโคจฺฉก}
\proseitem{16}{นกุสลํ โนอาสวํ ธมฺมํ ปฏิจฺจ อกุสโล อาสโว ธมฺโม อุปฺปชฺชติ เหตุปจฺจยา.  เหตุยา ตีณิ.}

\prose{นกุสลํ นโนอาสวํ ธมฺมํ ปฏิจฺจ อกุสโล โนอาสโว ธมฺโม อุปฺปชฺชติ เหตุปจฺจยา. เหตุยา นว.}

\proseitem{17}{นกุสลํ นสาสวํ ธมฺมํ ปฏิจฺจ อพฺยากโต สาสโว ธมฺโม อุปฺปชฺชติ เหตุปจฺจยา.  เหตุยา ปญฺจ.}

\prose{นกุสลํ นอนาสวํ ธมฺมํ ปจฺจยา กุสโล อนาสโว ธมฺโม อุปฺปชฺชติ เหตุปจฺจยา. (สํขิตฺตํ.) . เหตุยา ฉ ฯเปฯ วิปาเก ตีณิ ฯเปฯ อวิคเต ฉ.}

\proseitem{18}{นกุสลํ นอาสว\-สมฺปยุตฺตํ ธมฺมํ ปฏิจฺจ อกุสโล อาสว\-สมฺปยุตฺโต ธมฺโม อุปฺปชฺชติ เหตุปจฺจยา.  เหตุยา ตีณิ.}

\prose{นกุสลํ นอาสว\-วิปฺปยุตฺตํ ธมฺมํ ปฏิจฺจ อกุสโล อาสว\-วิปฺปยุตฺโต ธมฺโม อุปฺปชฺชติ เหตุปจฺจยา.  เหตุยา นว.}

\csromanheader{2}{ธมฺม\-ปจฺจนียา\-นุโลม ติกทุก\-ปฏฺฐาน\-ปาฬิ}
\proseitem{19}{\csromanfolio{382}นกุสลํ นอาสวญฺเจว นอนาสวญฺจ ธมฺมํ ปฏิจฺจ อกุสโล อาสโว เจว สาสโว จ ธมฺโม อุปฺปชฺชติ เหตุปจฺจยา.  เหตุยา ติณิ.}

\prose{นกุสลํ นอนาสวญฺเจว นโน จ อาสวํ ธมฺมํ ปฏิจฺจ อกุสโล สาสโว เจว โน จ อาสโว ธมฺโม อุปฺปชฺชติ เหตุปจฺจยา.  เหตุยา นว.}

\proseitem{20}{นกุสลํ นอาสวญฺเจว นอาสว\-วิปฺปยุตฺตญฺจ ธมฺมํ ปฏิจฺจ อกุสโล อาสโว เจว อาสว\-สมฺปยุตฺโต จ ธมฺโม อุปฺปชฺชติ เหตุปจฺจยา.  เหตุยา ตีณิ.}

\prose{นกุสลํ นอาสว\-วิปฺปยุตฺต\-ญฺเจว นโน จ อาสวํ ธมฺมํ ปฏิจฺจ อกุสโล อาสว\-สมฺปยุตฺโต เจว โน อาสโว ธมฺโม อุปฺปชฺชติ เหตุปจฺจยา.  เหตุยา ตีณิ.}

\proseitem{21}{นกุสลํ อาสว\-วิปฺปยุตฺตํ นสาสวํ ธมฺมํ ปฏิจฺจ อพฺยากโต อาสว\-วิปฺปยุตฺโต สาสโว ธมฺโม อุปฺปชฺชติ เหตุปจฺจยา.  เหตุยา ปญฺจ.}

\csromanmarkh{1. กุสลตฺติก}\csromanmarkhii{20. สญฺโญชนาทิโคจฺฉก}\csromanheadernomark{2}{1. \textbf{กุสลตฺติก 2}0. สญฺโญชนาทิโคจฺฉก}
\proseitem{22}{นกุสลํ โนสญฺโญชนํ ธมฺมํ ปฏิจฺจ อกุสโล สญฺโญชโน ธมฺโม อุปฺปชฺชติ เหตุปจฺจยา.}

\proseitem{23}{นกุสลํ โนคนฺถํ ธมฺมํ ปฏิจฺจ อกุสโล คนฺโถ ธมฺโม อุปฺปชฺชติ เหตุปจฺจยา.}

\proseitem{24}{นกุสลํ โนโอฆํ ธมฺมํ ปฏิจฺจ อกุสโล โอโฆ ธมฺโม อุปฺปชฺชติ เหตุปจฺจยา.}

\proseitem{25}{นกุสลํ โนโยคํ ธมฺมํ ปฏิจฺจ อกุสโล โยโค ธมฺโม อุปฺปชฺชติ เหตุปจฺจยา.}

\proseitem{26}{นกุสลํ โนนีวรณํ ธมฺมํ ปฏิจฺจ อกุสโล นีวรโณ ธมฺโม อุปฺปชฺชติ เหตุปจฺจยา.}

\csromanmarkh{1. กุสลตฺติก}\csromanmarkhii{96. ปริยาปนฺนทุก}\csromanheadernomark{2}{1. กุสลตฺติก 96. ปริยาปนฺนทุก}
\proseitem{27}{\csromanfolio{383}นกุสลํ โนปรามาสํ ธมฺมํ ปฏิจฺจ อกุสโล ปรามาโส ธมฺโม อุปฺปชฺชติ เหตุปจฺจยา.}

\csromanmarkh{1. กุสลตฺติก}\csromanmarkhii{55. มหนฺตรทุก}\csromanheadernomark{2}{1. \textbf{กุสลตฺติก} 55. มหนฺตรทุก}
\proseitem{28}{นกุสลํ นสารมฺมณํ ธมฺมํ ปฏิจฺจ อพฺยากโต สารมฺมโณ ธมฺโม อุปฺปชฺชติ เหตุปจฺจยา ฯเปฯ.}

\csromanmarkh{1. กุสลตฺติก}\csromanmarkhii{93. กามาวจรทุก}\csromanheadernomark{2}{1. \textbf{กุสลตฺติก} 93. กามาวจรทุก}
\proseitem{29}{นกุสลํ นกามาวจรํ ธมฺมํ ปฏิจฺจ อพฺยากโต กามาวจโร ธมฺโม อุปฺปชฺชติ เหตุปจฺจยา.  เหตุยา ปญฺจ.}

\prose{นกุสลํ นนกามา\-วจรํ ธมฺมํ ปฏิจฺจ อพฺยากโต นกามาวจโร ธมฺโม อุปฺปชฺชติ เหตุปจฺจยา.  เหตุยา ตีณิ.}

\csromanmarkh{1. กุสลตฺติก}\csromanmarkhii{94. รูปาวจรทุก}\csromanheadernomark{2}{1. \textbf{กุสลตฺติก} 94. รูปาวจรทุก}
\proseitem{30}{นกุสลํ นรูปาวจรํ ธมฺมํ ปฏิจฺจ อพฺยากโต รูปาวจโร ธมฺโม อุปฺปชฺชติ เหตุปจฺจยา.  เหตุยา ตีณิ.}

\prose{นกุสลํ นนรูปาวจรํ ธมฺมํ ปฏิจฺจ อพฺยากโต นรูปาวจโร ธมฺโม อุปฺปชฺชติ เหตุปจฺจยา.  เหตุยา ปญฺจ.}

\csromanmarkh{1. กุสลตฺติก}\csromanmarkhii{95. อรูปาวจรทุก}\csromanheadernomark{2}{1. \textbf{กุสลตฺติก} 95. อรูปาวจรทุก}
\proseitem{31}{นกุสลํ นอรูปาวจรํ ธมฺมํ ปจฺจยา กุสโล อรูปาวจโร ธมฺโม อุปฺปชฺชติ เหตุปจฺจยา.  เหตุยา ฉ.}

\prose{นกุสลํ นนอรูปาวจรํ ธมฺมํ ปฏิจฺจ อพฺยากโต น อรูปาวจโร ธมฺโม อุปฺปชฺชติ เหตุปจฺจยา.  เหตุยา ปญฺจ.}

\csromanmarkh{1. กุสลตฺติก}\csromanmarkhii{96. ปริยาปนฺนทุก}\csromanheadernomark{2}{1. กุสลตฺติก 96. ปริยาปนฺนทุก}
\proseitem{32}{นกุสลํ นปริยาปนฺนํ ธมฺมํ ปฏิจฺจ อพฺยากโต ปริยาปนฺโน ธมฺโม อุปฺปชฺชติ เหตุปจฺจยา.  นอกุสลํ นปริยาปนฺนํ ธมฺมํ ปฏิจฺจ อพฺยากโต ปริยาปนฺโน ธมฺโม อุปฺปชฺชติ เหตุปจฺจยา.  นอพฺยากตํ นปริยาปนฺนํ ธมฺมํ ปฏิจฺจ อพฺยากโต ปริยาปนฺโน ธมฺโม อุปฺปชฺชติ เหตุปจฺจยา. เหตุยา ปญฺจ.}

\csromanheader{2}{ธมฺม\-ปจฺจนียา\-นุโลม ติกทุก\-ปฏฺฐาน\-ปาฬิ}
\prose{\csromanfolio{384}นกุสลํ นอปริยาปนฺนํ ธมฺมํ ปจฺจยา กุสโล อปริยาปนฺโน ธมฺโม อุปฺปชฺชติ เหตุปจฺจยา.  นกุสลํ นอปริยาปนฺนํ ธมฺมํ ปจฺจยา อพฺยากโต อปริยาปนฺโน ธมฺโม อุปฺปชฺชติ เหตุปจฺจยา. เทฺว. (อกุสลมูเล เทฺว. ทุกมูเล เทฺว.) . เหตุยา ฉ.}

\csromanmarkh{1. กุสลตฺติก}\csromanmarkhii{97. นิยฺยานิกทุก}\csromanheadernomark{2}{1. \textbf{กุสลตฺติก} 97. นิยฺยานิกทุก}
\proseitem{33}{นกุสลํ นนิยฺยานิกํ ธมฺมํ ปจฺจยา กุสโล นิยฺยานิโก ธมฺโม อุปฺปชฺชติ เหตุปจฺจยา.  เหตุยา ตีณิ.}

\prose{นอกุสลํ นอนิยฺยานิกํ ธมฺมํ ปฏิจฺจ อพฺยากโต อนิยฺยานิโก ธมฺโม อุปฺปชฺชติ เหตุปจฺจยา.  นอพฺยากตํ นอนิยฺยานิกํ ธมฺมํ ปฏิจฺจ อพฺยากโต อนิยฺยานิโก ธมฺโม อุปฺปชฺชติ เหตุปจฺจยา.  นอกุสลํ นอนิยฺยานิกญฺจ นอพฺยากตํ นอนิยฺยานิกญฺจ ธมฺมํ ปฏิจฺจ อพฺยากโต อนิยฺยานิโก ธมฺโม อุปฺปชฺชติ เหตุปจฺจยา.  เหตุยา ตีณิ.}

\csromanmarkh{1. กุสลตฺติก}\csromanmarkhii{98. นิยตทุก}\csromanheadernomark{2}{1. \textbf{กุสลตฺติก} 98. นิยตทุก}
\proseitem{34}{นกุสลํ นนิยตํ ธมฺมํ ปจฺจยา กุสโล นิยโต ธมฺโม อุปฺปชฺชติ เหตุปจฺจยา. (อปริยาปนฺน\-สทิสํ. ฉ ปญฺหา.)}

\prose{นกุสลํ นอนิยตํ ธมฺมํ ปฏิจฺจ อพฺยากโต อนิยโต ธมฺโม อุปฺปชฺชติ เหตุปจฺจยา.  นอกุสลํ นอนิยตํ ธมฺมํ ปฏิจฺจ อพฺยากโต อนิยโต ธมฺโม อุปฺปชฺชติ เหตุปจฺจยา.  นอพฺยากตํ นอนิยตํ ธมฺมํ ปฏิจฺจ อพฺยากโต อนิยโต ธมฺโม อุปฺปชฺชติ เหตุปจฺจยา.  นกุสลํ นอนิยตญฺจ นอพฺยากตํ นอนิยตญฺจ ธมฺมํ ปฏิจฺจ อพฺยากโต อนิยโต ธมฺโม อุปฺปชฺชติ เหตุปจฺจยา.  นกุสลํ นอนิยตญฺจ นอพฺยากตํ นอนิยตญฺจ ธมฺมํ ปฏิจฺจ อพฺยากโต อนิยโต ธมฺโม อุปฺปชฺชติ เหตุปจฺจยา.  เหตุยา ปญฺจ.}

\prose{1. กุสลตฺติก 99. สอุตฺตรทุก}

\proseitem{35}{นกุสลํ นสอุตฺตรํ ธมฺมํ ปฏิจฺจ อพฺยากโต สอุตฺตโร ธมฺโม อุปฺปชฺชติ เหตุปจฺจยา.  นอกุสลํ นสอุตฺตรํ ธมฺมํ ปฏิจฺจ อพฺยากโต สอุตฺตโร ธมฺโม อุปฺปชฺชติ เหตุปจฺจยา.  นอพฺยากตํ นสอุตฺตรํ ธมฺมํ}

\vspace{\tipitakaparskip}
\csromancenter{เวทนาตฺติก 1. เหตุทุก ปฏิจฺจ อพฺยากโต สอุตฺตโร ธมฺโม อุปฺปชฺชติ เหตุปจฺจยา.  นอกุสลํ นสอุตฺตรญฺจ นอพฺยากตํ นสอุตฺตรญฺจ ธมฺมํ ปฏิจฺจ อพฺยากโต สอุตฺตโร ธมฺโม อุปฺปชฺชติ เหตุปจฺจยา.  นกุสลํ นสอุตฺตรญฺจ นอกุสลํ นสอุตฺตรญฺจ ธมฺมํ ปฏิจฺจ อพฺยากโต สอุตฺตโร ธมฺโม อุปฺปชฺชติ เหตุปจฺจยา.  เหตุยา ปญฺจ.}
\prose{\csromanfolio{385}นกุสลํ นอนุตฺตรํ ธมฺมํ ปจฺจยา กุสโล อนุตฺตโร ธมฺโม อุปฺปชฺชติ เหตุปจฺจยา. (สํขิตฺตํ.) . เหตุยา ฉ ฯเปฯ วิปาเก ตีณิ ฯเปฯ อวิคเต ฉ.}

\csromanmarkh{1. กุสลตฺติก}\csromanmarkhii{100. สรณทุก}\csromanheadernomark{2}{1. \textbf{กุสลตฺติก} 100. สรณทุก}
\proseitem{36}{นกุสลํ นสรณํ ธมฺมํ ปจฺจยา อกุสโล สรโณ ธมฺโม อุปฺปชฺชติ เหตุปจฺจยา.  เหตุยา ตีณิ.}

\prose{นกุสลํ นอรณํ ธมฺมํ ปฏิจฺจ อพฺยากโต อรโณ ธมฺโม อุปฺปชฺชติ เหตุปจฺจยา.  นอพฺยากตํ นอรณํ ธมฺมํ ปฏิจฺจ อพฺยากโต อรโณ ธมฺโม อุปฺปชฺชติ เหตุปจฺจยา.  นกุสลํ นอรณญฺจ นอพฺยากตํ นอรณญฺจ ธมฺมํ ปฏิจฺจ อพฺยากโต อรโณ ธมฺโม อุปฺปชฺชติ เหตุปจฺจยา.  เหตุยา ตีณิ.}

\csromanmarkh{2. เวทนาตฺติก}\csromanmarkhii{1. เหตุทุก}\csromanheadernomark{2}{2. เวทนาตฺติก 1. เหตุทุก}
\prosecont{นสุขาย เวทนาย สมฺปยุตฺตํ นเหตุํ ธมฺมํ ปฏิจฺจ สุขาย เวทนาย สมฺปยุตฺโต เหตุ ธมฺโม อุปฺปชฺชติ เหตุปจฺจยา.  นสุขาย เวทนาย สมฺปยุตฺตํ นเหตุํ ธมฺมํ ปฏิจฺจ ทุกฺขาย เวทนาย สมฺปยุตฺโต เหตุ ธมฺโม อุปฺปชฺชติ เหตุปจฺจยา.  นสุขาย เวทนาย สมฺปยุตฺตํ นเหตุํ ธมฺมํ ปฏิจฺจ อทุกฺขม\-สุขาย เวทนาย สมฺปยุตฺโต เหตุ ธมฺโม อุปฺปชฺชติ เหตุปจฺจยา. ตีณิ.}

\proseitem{37}{นทุกฺขาย เวทนาย สมฺปยุตฺตํ นเหตุํ ธมฺมํ ปฏิจฺจ ทุกฺขาย เวทนาย สมฺปยุตฺโต เหตุ ธมฺโม อุปฺปชฺชติ เหตุปจฺจยา. ตีณิ.}

\prose{นอทุกฺขม\-สุขาย เวทนาย สมฺปยุตฺตํ นเหตุํ ธมฺมํ ปฏิจฺจ อทุกฺขม\-สุขาย เวทนาย สมฺปยุตฺโต เหตุ ธมฺโม อุปฺปชฺชติ เหตุปจฺจยา. ตีณิ.}

\csromanheader{2}{ธมฺม\-ปจฺจนียา\-นุโลม ติกทุก\-ปฏฺฐาน\-ปาฬิ}
\prose{\csromanfolio{386}นสุขาย เวทนาย สมฺปยุตฺตํ นเหตุญฺจ นอทุกฺขม\-สุขาย เวทนาย สมฺปยุตฺตํ นเหตุญฺจ ธมฺมํ ปฏิจฺจ สุขาย เวทนาย สมฺปยุตฺโต เหตุ ธมฺโม อุปฺปชฺชติ เหตุปจฺจยา. ตีณิ. (ทุติยํ คณิตเกน ตีณิ.) . เหตุยา เอกวีส.}

\prose{นสุขาย เวทนาย สมฺปยุตฺตํ นนเหตุํ ธมฺมํ ปฏิจฺจ ทุกฺขาย เวทนาย สมฺปยุตฺโต นเหตุ ธมฺโม อุปฺปชฺชติ เหตุปจฺจยา. เทฺว.}

\prose{นทุกฺขาย เวทนาย สมฺปยุตฺตํ นนเหตุํ ธมฺมํ ปฏิจฺจ. (เทฺว ปญฺหาเยว.)}

\prose{นอทุกฺขม\-สุขาย เวทนาย สมฺปยุตฺตํ นนเหตุํ ธมฺมํ ปฏิจฺจ. (เทฺวเยว. ปฐมํ คณิตเกน เอกํ, ทุติยํ คณิตเกน เอกํ, ตติยํ คณิตเกน เอกํ กาตพฺพํ.) . เหตุยา นว.}

\csromanmarkh{3. วิปากตฺติก}\csromanmarkhii{1. เหตุทุก}\csromanheadernomark{2}{3. \textbf{วิปากตฺติก 1. เหตุทุก}}
\proseitem{38}{นวิปากํ นเหตุํ ธมฺมํ ปฏิจฺจ วิปาโก เหตุ ธมฺโม อุปฺปชฺชติ เหตุปจฺจยา.  นวิปากํ นเหตุํ ธมฺมํ ปฏิจฺจ วิปาก\-ธมฺม\-ธมฺโม เหตุ ธมฺโม อุปฺปชฺชติ เหตุปจฺจยา.  นวิปากํ นเหตุํ ธมฺมํ ปฏิจฺจ เนว\-วิปาก\-น\-วิปาก\-ธมฺม\-ธมฺโม เหตุ ธมฺโม อุปฺปชฺชติ เหตุปจฺจยา. ตีณิ.}

\prose{นวิปาก\-ธมฺม\-ธมฺมํ นเหตุํ ธมฺมํ ปฏิจฺจ วิปาโก เหตุ ธมฺโม อุปฺปชฺชติ เหตุปจฺจยา.  นวิปาก\-ธมฺม\-ธมฺมํ นเหตุํ ธมฺมํ ปฏิจฺจ เนว\-วิปาก\-น\-วิปาก\-ธมฺม\-ธมฺโม เหตุ ธมฺโม อุปฺปชฺชติ เหตุปจฺจยา. เทฺว.}

\prose{นเน\-ว\-วิปาก\-น\-วิปากธมฺม\-ธมฺมํ นเหตุํ ธมฺมํ ปฏิจฺจ วิปาโก เหตุ ธมฺโม อุปฺปชฺชติ เหตุปจฺจยา.  นเน\-ว\-วิปาก\-น\-วิปากธมฺม\-ธมฺมํ นเหตุํ ธมฺมํ ปฏิจฺจ วิปาก\-ธมฺม\-ธมฺโม เหตุ ธมฺโม อุปฺปชฺชติ เหตุปจฺจยา. เทฺว.}

\prose{นวิปากํ นเหตุญฺจ นเน\-ว\-วิปาก\-น\-วิปากธมฺม\-ธมฺมํ นเหตุญฺจ ธมฺมํ ปฏิจฺจ วิปาก\-ธมฺม\-ธมฺโม เหตุ ธมฺโม อุปฺปชฺชติ เหตุปจฺจยา. เอกํ.}

\prose{นวิปาก\-ธมฺม\-ธมฺมํ นเหตุญฺจ นเน\-ว\-วิปาก\-น\-วิปากธมฺม\-ธมฺมํ นเหตุญฺจ ธมฺมํ ปฏิจฺจ วิปาโก เหตุ ธมฺโม อุปฺปชฺชติ เหตุปจฺจยา. เอกํ.}

\prose{นวิปากํ นเหตุญฺจ นวิปาก\-ธมฺม\-ธมฺมํ นเหตุญฺจ ธมฺมํ ปฏิจฺจ วิปาโก เหตุ ธมฺโม อุปฺปชฺชติ เหตุปจฺจยา.  นวิปากํ นเหตุญฺจ นวิปาก\-ธมฺม\-ธมฺมํ}

\vspace{\tipitakaparskip}
\csromancenter{วิตกฺกตฺติกาทิ 1. เหตุทุก นเหตุญฺจ ธมฺมํ ปฏิจฺจ เนว\-วิปาก\-น\-วิปาก\-ธมฺม\-ธมฺโม เหตุ ธมฺโม อุปฺปชฺชติ เหตุปจฺจยา. เทฺว. (สํขิตฺตํ.) . เหตุยา เอกาทส.}
\prose{\csromanfolio{387}นวิปากํ นนเหตุํ ธมฺมํ ปฏิจฺจ วิปาก\-ธมฺม\-ธมฺโม นเหตุ ธมฺโม อุปฺปชฺชติ เหตุปจฺจยา.  เหตุยา อฏฺฐารส. (สํขิตฺตํ.)}

\csromanmarkh{4. อุปาทินฺนตฺติก}\csromanmarkhii{1. เหตุทุก}\csromanheadernomark{2}{4. \textbf{อุปาทินฺนตฺติก 1. เหตุทุก}}
\proseitem{39}{นอุปาทินฺนุปาทานิยํ นเหตุํ ธมฺมํ ปฏิจฺจ อนุปาทินฺนุ\-ปาทานิโย เหตุ ธมฺโม อุปฺปชฺชติ เหตุปจฺจยา.  เหตุยา นว.}

\prose{นอุปาทินฺนุปาทานิยํ นนเหตุํ ธมฺมํ ปฏิจฺจ อนุปาทินฺนุ\-ปาทานิโย นเหตุ ธมฺโม อุปฺปชฺชติ เหตุปจฺจยา.  เหตุยา อฏฺฐารส.}

\csromanmarkh{5. สํกิลิฏฺฐ\-ตฺติก}\csromanmarkhii{1. เหตุทุก}\csromanheadernomark{2}{5. สํกิลิฏฺฐ\-ตฺติก 1. เหตุทุก}
\proseitem{40}{นสํกิลิฏฺฐ\-สํกิเลสิกํ นเหตุํ ธมฺมํ ปฏิจฺจ อสํกิลิฏฺฐ\-สํกิเลสิโก เหตุ ธมฺโม อุปฺปชฺชติ เหตุปจฺจยา.  เหตุยา นว.}

\csromanmarkh{6-11. วิตกฺกตฺติกาทิ}\csromanmarkhii{1. เหตุทุก}\csromanheadernomark{2}{6-11. วิตกฺกตฺติกาทิ 1. เหตุทุก}
\proseitem{41}{นสวิตกฺก\-สวิจารํ นเหตุํ ธมฺมํ ปฏิจฺจ สวิตกฺก\-สวิจาโร เหตุ ธมฺโม อุปฺปชฺชติ เหตุปจฺจยา.  เหตุยา ปนฺนรส.}

\proseitem{42}{นปีติสหคตํ นเหตุํ ธมฺมํ ปฏิจฺจ ปีติสหคโต เหตุ ธมฺโม อุปฺปชฺชติ เหตุปจฺจยา.  เหตุยา อฏฺฐวีส.}

\proseitem{43}{นทสฺสเนน ปหาตพฺพํ นเหตุํ ธมฺมํ ปฏิจฺจ ภาวนาย ปหาตพฺโพ เหตุ ธมฺโม อุปฺปชฺชติ เหตุปจฺจยา.  เหตุยา นว.}

\proseitem{44}{นทสฺสเนน ปหาตพฺพ\-เหตุกํ นเหตุํ ธมฺมํ ปฏิจฺจ ภาวนาย ปหาตพฺพ\-เหตุโก เหตุ ธมฺโม อุปฺปชฺชติ เหตุปจฺจยา.  เหตุยา นว.}

\proseitem{45}{นอาจยคามิํ นเหตุํ ธมฺมํ ปฏิจฺจ อปจยคามี เหตุ ธมฺโม อุปฺปชฺชติ เหตุปจฺจยา.  เหตุยา นว.}

\csromanheader{2}{ธมฺม\-ปจฺจนียา\-นุโลม ติกทุก\-ปฏฺฐาน\-ปาฬิ}
\proseitem{46}{\csromanfolio{388}นเสกฺขํ นเหตุํ ธมฺมํ ปฏิจฺจ อเสกฺโข เหตุ ธมฺโม อุปฺปชฺชติ เหตุปจฺจยา.  เหตุยา นว.}

\csromanmarkh{12-13. ปริตฺต\-ตฺติก\-ทฺวย}\csromanmarkhii{1. เหตุทุก}\csromanheadernomark{2}{12-13. ปริตฺต\-ตฺติก\-ทฺวย 1. เหตุทุก}
\proseitem{47}{นปริตฺตํ นเหตุํ ธมฺมํ ปฏิจฺจ มหคฺคโต เหตุ ธมฺโม อุปฺปชฺชติ เหตุปจฺจยา.  เหตุยา เอกาทส.}

\proseitem{48}{นปริตฺตา\-รมฺมณํ นเหตุํ ธมฺมํ ปฏิจฺจ ปริตฺตารมฺมโณ เหตุ ธมฺโม อุปฺปชฺชติ เหตุปจฺจยา.  เหตุยา เตรส.}

\csromanmarkh{14-17. หีนตฺติกาทิ}\csromanmarkhii{1. เหตุทุก}\csromanheadernomark{2}{14-17. หีนตฺติกาทิ 1. เหตุทุก}
\proseitem{49}{นหีนํ นเหตุํ ธมฺมํ ปฏิจฺจ มชฺฌิโม เหตุ ธมฺโม อุปฺปชฺชติ เหตุปจฺจยา.  เหตุยา นว.}

\proseitem{50}{นมิจฺฉตฺต\-นิยตํ นเหตุํ ธมฺมํ ปฏิจฺจ สมฺมตฺตนิยโต เหตุ ธมฺโม อุปฺปชฺชติ เหตุปจฺจยา.  เหตุยา นว.}

\proseitem{51}{นมคฺคา\-รมฺมณํ นเหตุํ ธมฺมํ ปฏิจฺจ มคฺคเหตุโก เหตุ ธมฺโม อุปฺปชฺชติ เหตุปจฺจยา. (สํขิตฺตํ.) . เหตุยา ทส.}

\proseitem{52}{นอนุปฺปนฺโน นนเหตุ ธมฺโม อุปฺปนฺนสฺส นเหตุสฺส ธมฺมสฺส เหตุปจฺจเยน ปจฺจโย.  นอุปฺปาที นนเหตุ ธมฺโม อุปฺปนฺนสฺส นเหตุสฺส ธมฺมสฺส เหตุปจฺจเยน ปจฺจโย.  นอนุปฺปนฺโน นนเหตุ จ นอุปฺปาที นนเหตุ จ ธมฺมา อุปฺปนฺนสฺส นเหตุสฺส ธมฺมสฺส เหตุปจฺจเยน ปจฺจโย.  เหตุยา ตีณิ.}

\csromanmarkh{18-19. อตีต\-ตฺติก\-ทฺวย}\csromanmarkhii{1. เหตุทุก}\csromanheadernomark{2}{18-19. อตีต\-ตฺติก\-ทฺวย 1. เหตุทุก}
\proseitem{53}{นอตีโต นนเหตุ ธมฺโม ปจฺจุปฺปนฺนสฺส นเหตุสฺส ธมฺมสฺส เหตุปจฺจเยน ปจฺจโย. นอนาคโต นนเหตุ ธมฺโม ปจฺจุปฺปนฺนสฺส นเหตุสฺส ธมฺมสฺส เหตุปจฺจเยน ปจฺจโย.  นอตีโต นนเหตุ จ นอนาคโต นนเหตุ จ ธมฺมา ปจฺจุปฺปนฺนสฺส นเหตุสฺส ธมฺมสฺส เหตุปจฺจเยน ปจฺจโย.  เหตุยา ตีณิ.}

\csromanmarkh{22. สนิทสฺสน\-ตฺติก}\csromanmarkhii{2. สเหตุกทุก}\csromanheadernomark{2}{22. สนิทสฺสน\-ตฺติก 2. สเหตุกทุก}
\proseitem{54}{\csromanfolio{389}นอตีตา\-รมฺมณํ นเหตุํ ธมฺมํ ปฏิจฺจ อตีตารมฺมโณ เหตุ ธมฺโม อุปฺปชฺชติ เหตุปจฺจยา.  เหตุยา สตฺตรส.}

\csromanmarkh{20-21. อชฺฌตฺต\-ตฺติก\-ทฺวย}\csromanmarkhii{1. เหตุทุก}\csromanheadernomark{2}{20-21. อชฺฌตฺต\-ตฺติก\-ทฺวย 1. เหตุทุก}
\proseitem{55}{นอชฺฌตฺตํ นเหตุํ ธมฺมํ ปฏิจฺจ พหิทฺธา เหตุ ธมฺโม อุปฺปชฺชติ เหตุปจฺจยา.}

\prose{นพหิทฺธา นเหตุํ ธมฺมํ ปฏิจฺจ อชฺฌตฺโต เหตุ ธมฺโม อุปฺปชฺชติ เหตุปจฺจยา.  เหตุยา เทฺว.}

\proseitem{56}{นอชฺฌตฺตา\-รมฺมณํ นเหตุํ ธมฺมํ ปฏิจฺจ อชฺฌตฺตา\-รมฺมโณ เหตุ ธมฺโม อุปฺปชฺชติ เหตุปจฺจยา. เทฺว.}

\prose{นพหิทฺธา\-รมฺมณํ นเหตุํ ธมฺมํ ปฏิจฺจ พหิทฺธา\-รมฺมโณ เหตุ ธมฺโม อุปฺปชฺชติ เหตุปจฺจยา.  เหตุยา ฉ.}

\csromanmarkh{22. สนิทสฺสน\-ตฺติก}\csromanmarkhii{1. เหตุทุก}\csromanheadernomark{2}{22. สนิทสฺสน\-ตฺติก 1. เหตุทุก}
\proseitem{57}{นสนิทสฺสน\-สปฺปฏิฆํ นเหตุํ ธมฺมํ ปฏิจฺจ อนิทสฺสน\-อปฺปฏิโฆ เหตุ ธมฺโม อุปฺปชฺชติ เหตุปจฺจยา.  นอนิทสฺสน\-สปฺปฏิฆํ นเหตุํ ธมฺมํ ปฏิจฺจ อนิทสฺสน\-อปฺปฏิโฆ เหตุ ธมฺโม อุปฺปชฺชติ เหตุปจฺจยา.  นสนิทสฺสน\-สปฺปฏิฆํ นเหตุญฺจ นอนิทสฺสน\-สปฺปฏิฆํ นเหตุญฺจ ธมฺมํ ปฏิจฺจ อนิทสฺสน\-อปฺปฏิโฆ เหตุ ธมฺโม อุปฺปชฺชติ เหตุปจฺจยา.  เหตุยา ตีณิ.}

\prose{นสนิทสฺสน\-สปฺปฏิฆํ นนเหตุํ ธมฺมํ ปฏิจฺจ สนิทสฺสน\-สปฺปฏิโฆ นเหตุ ธมฺโม อุปฺปชฺชติ เหตุปจฺจยา.  นสนิทสฺสน\-สปฺปฏิฆํ นนเหตุํ ธมฺมํ ปฏิจฺจ อนิทสฺสน\-สปฺปฏิโฆ นเหตุ ธมฺโม อุปฺปชฺชติ เหตุปจฺจยา. (นเหตุ นนเหตุ โอวตฺตา. ตีณิ มูลานิ. เอกวีสติ ปญฺหา กาตพฺพา.)}

\csromanmarkh{22. สนิทสฺสน\-ตฺติก}\csromanmarkhii{2. สเหตุกทุก}\csromanheadernomark{2}{22. สนิทสฺสน\-ตฺติก 2. สเหตุกทุก}
\proseitem{58}{นสนิทสฺสน\-สปฺปฏิฆํ นสเหตุกํ ธมฺมํ ปฏิจฺจ อนิทสฺสน\-อปฺปฏิโฆ สเหตุโก ธมฺโม อุปฺปชฺชติ เหตุปจฺจยา.  นอนิทสฺสน\-สปฺปฏิฆํ นสเหตุกํ ธมฺมํ ปฏิจฺจ อนิทสฺสน\-อปฺปฏิโฆ สเหตุโก ธมฺโม อุปฺปชฺชติ}

\vspace{\tipitakaparskip}
\csromancenter{ธมฺม\-ปจฺจนียา\-นุโลม ติกทุก\-ปฏฺฐาน\-ปาฬิ เหตุปจฺจยา.  นสนิทสฺสน\-สปฺปฏิฆํ นสเหตุกญฺจ นอนิทสฺสน\-สปฺปฏิฆํ นสเหตุกญฺจ ธมฺมํ ปฏิจฺจ อนิทสฺสน\-อปฺปฏิโฆ สเหตุโก ธมฺโม อุปฺปชฺชติ เหตุปจฺจยา.  เหตุยา ตีณิ.}
\prose{\csromanfolio{390}นสนิทสฺสน\-สปฺปฏิฆํ นอเหตุกํ ธมฺมํ ปฏิจฺจ สนิทสฺสน\-สปฺปฏิโฆ อเหตุโก ธมฺโม อุปฺปชฺชติ เหตุปจฺจยา. สตฺต.}

\prose{นอนิทสฺสน\-สปฺปฏิฆํ นอเหตุกํ ธมฺมํ ปฏิจฺจ อนิทสฺสน\-สปฺปฏิโฆ อเหตุโก ธมฺโม อุปฺปชฺชติ เหตุปจฺจยา. สตฺต.}

\prose{นสนิทสฺสน\-สปฺปฏิฆํ นอเหตุกญฺจ นอนิทสฺสน\-สปฺปฏิฆํ นอเหตุกญฺจ ธมฺมํ ปฏิจฺจ อนิทสฺสน\-สปฺปฏิโฆ อเหตุโก ธมฺโม อุปฺปชฺชติ เหตุปจฺจยา. สตฺต.  เหตุยา เอกวีส.}

\csromanheader{2}{22. สนิทสฺสน\-ตฺติก 3-6. เหตุ\-สมฺปยุตฺต\-ทุกาทิ}
\proseitem{59}{นสนิทสฺสน\-สปฺปฏิฆํ นเหตุ\-สมฺปยุตฺตํ ธมฺมํ ปฏิจฺจ อนิทสฺสน\-อปฺปฏิโฆ เหตุสมฺปยุตฺโต ธมฺโม อุปฺปชฺชติ เหตุปจฺจยา.  เหตุยา ตีณิ.}

\prose{นสนิทสฺสน\-สปฺปฏิฆํ นเหตุ\-วิปฺปยุตฺตํ ธมฺมํ ปฏิจฺจ สนิทสฺสน\-สปฺปฏิโฆ เหตุวิปฺปยุตฺโต ธมฺโม อุปฺปชฺชติ เหตุปจฺจยา.  เหตุยา เอกวีส.}

\proseitem{60}{นสนิทสฺสน\-สปฺปฏิฆํ นเหตุญฺเจว นอเหตุกญฺจ ธมฺมํ ปฏิจฺจ อนิทสฺสน\-อปฺปฏิโฆ เหตุ เจว สเหตุโก จ ธมฺโม อุปฺปชฺชติ เหตุปจฺจยา.  เหตุยา ตีณิ.}

\prose{นสนิทสฺสน\-สปฺปฏิฆํ นอเหตุกญฺเจว นน จ เหตุํ ธมฺมํ ปฏิจฺจ อนิทสฺสน\-อปฺปฏิโฆ สเหตุโก เจว น จ เหตุ ธมฺโม อุปฺปชฺชติ เหตุปจฺจยา.  เหตุยา ตีณิ.}

\proseitem{61}{นสนิทสฺสน\-สปฺปฏิฆํ นเหตุญฺเจว นเหตุ\-วิปฺปยุตฺตญฺจ ธมฺมํ ปฏิจฺจ อนิทสฺสน\-อปฺปฏิโฆ เหตุ เจว เหตุสมฺปยุตฺโต จ ธมฺโม อุปฺปชฺชติ เหตุปจฺจยา.  เหตุยา ตีณิ.}

\prose{นสนิทสฺสน\-สปฺปฏิฆํ นเหตุ\-วิปฺปยุตฺต\-ญฺเจว นน จ เหตุํ ธมฺมํ ปฏิจฺจ อนิทสฺสน\-อปฺปฏิโฆ เหตุสมฺปยุตฺโต เจว น จ เหตุ ธมฺโม อุปฺปชฺชติ เหตุปจฺจยา.  เหตุยา ตีณิ.}

\csromanmarkh{22. สนิทสฺสน\-ตฺติก}\csromanmarkhii{7. จูฬนฺตรทุก}\csromanheadernomark{2}{22. สนิทสฺสน\-ตฺติก 7. จูฬนฺตรทุก}
\proseitem{62}{\csromanfolio{391}นสนิทสฺสน\-สปฺปฏิฆํ นเหตุํ นสเหตุกํ ธมฺมํ ปฏิจฺจ อนิทสฺสน\-อปฺปฏิโฆ นเหตุ สเหตุโก ธมฺโม อุปฺปชฺชติ เหตุปจฺจยา.  เหตุยา ตีณิ.}

\prose{นสนิทสฺสน\-สปฺปฏิฆํ นเหตุํ นอเหตุกํ ธมฺมํ ปฏิจฺจ สนิทสฺสน\-สปฺปฏิโฆ นเหตุ อเหตุโก ธมฺโม อุปฺปชฺชติ เหตุปจฺจยา.  เหตุยา เอกวีส.}

\csromanmarkh{22. สนิทสฺสน\-ตฺติก}\csromanmarkhii{7. จูฬนฺตรทุก}\csromanheadernomark{2}{22. สนิทสฺสน\-ตฺติก 7. จูฬนฺตรทุก}
\proseitem{63}{นสนิทสฺสน\-สปฺปฏิโฆ นสปฺปจฺจโย ธมฺโม อนิทสฺสน\-อปฺปฏิฆสฺส สปฺปจฺจยสฺส ธมฺมสฺส อารมฺมณ\-ปจฺจเยน ปจฺจโย. (สํขิตฺตํ.) . อารมฺมเณ ตีณิ, อธิปติยา อุปนิสฺสเย ตีณิ.}

\proseitem{64}{นสนิทสฺสน\-สปฺปฏิฆํ นสนิทสฺสนํ ธมฺมํ ปฏิจฺจ สนิทสฺสน\-สปฺปฏิโฆ สนิทสฺสโน ธมฺโม อุปฺปชฺชติ เหตุปจฺจยา.  เหตุยา ปญฺจ.}

\prose{นอนิทสฺสน\-อปฺปฏิโฆ นนส\-นิทสฺสโน ธมฺโม อนิทสฺสน\-อปฺปฏิฆสฺส นสนิทสฺสนสฺส ธมฺมสฺส อารมฺมณ\-ปจฺจเยน ปจฺจโย. (สํขิตฺตํ.) . อารมฺมเณ ตีณิ, อธิปติยา อุปนิสฺสเย ปุเรชาเต อตฺถิยา อวิปเต ตีณิ.}

\proseitem{65}{นสนิทสฺสน\-สปฺปฏิฆํ นสปฺปฏิฆํ ธมฺมํ ปฏิจฺจ สนิทสฺสน\-สปฺปฏิโฆ สปฺปฏิโฆ ธมฺโม อุปฺปชฺชติ เหตุปจฺจยา.  เหตุยา นว.}

\prose{นสนิทสฺสน\-สปฺปฏิฆํ นอปฺปฏิฆํ ธมฺมํ ปฏิจฺจ อนิทสฺสน\-อปฺปฏิโฆ อปฺปฏิโฆ ธมฺโม อุปฺปชฺชติ เหตุปจฺจยา.  เหตุยา ตีณิ.}

\proseitem{66}{นสนิทสฺสน\-สปฺปฏิฆํ นรูปิํ ธมฺมํ ปฏิจฺจ สนิทสฺสน\-สปฺปฏิโฆ รูปี ธมฺโม อุปฺปชฺชติ เหตุปจฺจยา.  เหตุยา เอกวีส.}

\prose{นสนิทสฺสน\-สปฺปฏิฆํ นอรูปิํ ธมฺมํ ปฏิจฺจ อนิทสฺสน\-อปฺปฏิโฆ อรูปี ธมฺโม อุปฺปชฺชติ เหตุปจฺจยา.  เหตุยา ตีณิ.}

\proseitem{67}{นสนิทสฺสน\-สปฺปฏิฆํ นโลกิยํ ธมฺมํ ปฏิจฺจ สนิทสฺสน\-สปฺปฏิโฆ โลกิโย ธมฺโม อุปฺปชฺชติ เหตุปจฺจยา. สตฺต.}

\csromanheader{2}{ธมฺม\-ปจฺจนียา\-นุโลม ติกทุก\-ปฏฺฐาน\-ปาฬิ}
\prose{\csromanfolio{392}นอนิทสฺสน\-สปฺปฏิฆํ นโลกิยํ ธมฺมํ ปฏิจฺจ อนิทสฺสน\-สปฺปฏิโฆ โลกิโย ธมฺโม อุปฺปชฺชติ เหตุปจฺจยา. สตฺต.  เหตุยา เอกวีส.}

\prose{นสนิทสฺสน\-สปฺปฏิฆํ นโลกุตฺตรํ ธมฺมํ ปจฺจยา อนิทสฺสน\-อปฺปฏิโฆ โลกุตฺตโร ธมฺโม อุปฺปชฺชติ เหตุปจฺจยา. (สํขิตฺตํ.) . เหตุยา ตีณิ ฯเปฯ อวิคเต ตีณิ.}

\prose{นสนิทสฺสน\-สปฺปฏิฆํ นโลกุตฺตรํ ธมฺมํ ปจฺจยา อนิทสฺสน\-อปฺปฏิโฆ โลกุตฺตโร ธมฺโม อุปฺปชฺชติ เหตุปจฺจยา. (สํขิตฺตํ.) . เหตุยา ตีณิ ฯเปฯ อวิคเต ตีณิ.}

\proseitem{68}{นสนิทสฺสน\-สปฺปฏิฆํ นเกนจิ วิญฺเญยฺยํ ธมฺมํ ปฏิจฺจ สนิทสฺสน\-สปฺปฏิโฆ เกนจิ วิญฺเญยฺโย ธมฺโม อุปฺปชฺชติ เหตุปจฺจยา.  เหตุยา ปญฺจติํส.}

\prose{นสนิทสฺสน\-สปฺปฏิฆํ นนเกนจิ วิญฺเญยฺยํ ธมฺมํ ปฏิจฺจ สนิทสฺสน\-สปฺปฏิโฆ นเกนจิ วิญฺเญยฺโย ธมฺโม อุปฺปชฺชติ เหตุปจฺจยา.  เหตุยา ปญฺจติํส.}

\csromanmarkh{22. สนิทสฺสนตฺตฺอิก}\csromanmarkhii{14. อาสวโคจฺฉก}\csromanheadernomark{2}{22. \textbf{สนิทสฺสนตฺตฺ}อิก 14. อาสวโคจฺฉก}
\proseitem{69}{นสนิทสฺสน\-สปฺปฏิฆํ โนอาสวํ ธมฺมํ ปฏิจฺจ อนิทสฺสน\-อปฺปฏิโฆ อาสโว ธมฺโม อุปฺปชฺชติ เหตุปจฺจยา.  เหตุยา ตีณิ.}

\prose{นสนิทสฺสน\-สปฺปฏิฆํ นโนอาสวํ ธมฺมํ ปฏิจฺจ สนิทสฺสน\-สปฺปฏิโฆ โนอาสโว ธมฺโม อุปฺปชฺชติ เหตุปจฺจยา.  เหตุยา เอกวีส.}

\proseitem{70}{นสนิทสฺสน\-สปฺปฏิฆํ นสาสวํ ธมฺมํ ปฏิจฺจ สนิทสฺสน\-สปฺปฏิโฆ สาสโว ธมฺโม อุปฺปชฺชติ เหตุปจฺจยา.  เหตุยา เอกวีส.}

\prose{นสนิทสฺสน\-สปฺปฏิฆํ นอนาสวํ ธมฺมํ ปจฺจยา อนิทสฺสน\-อปฺปฏิโฆ อนาสโว ธมฺโม อุปฺปชฺชติ เหตุปจฺจยา. (สํขิตฺตํ.) . เหตุยา ตีณิ ฯเปฯ อวิคเต ตีณิ.}

\proseitem{71}{นสนิทสฺสน\-สปฺปฏิฆํ นอาสว\-สมฺปยุตฺตํ ธมฺมํ ปฏิจฺจ อนิทสฺสน\-อปฺปฏิโฆ อาสว\-สมฺปยุตฺโต ธมฺโม อุปฺปชฺชติ เหตุปจฺจยา.  เหตุยา ตีณิ.}

\csromanmarkh{22. สนิทสฺสน\-ตฺติก}\csromanmarkhii{20. สญฺโญชนาทิโคจฺฉก}\csromanheadernomark{2}{22. สนิทสฺสน\-ตฺติก 20. สญฺโญชนาทิโคจฺฉก}
\prose{\csromanfolio{393}นสนิทสฺสน\-สปฺปฏิฆํ นอาสว\-วิปฺปยุตฺตํ ธมฺมํ ปฏิจฺจ สนิทสฺสน\-สปฺปฏิโฆ อาสว\-วิปฺปยุตฺโต ธมฺโม อุปฺปชฺชติ เหตุปจฺจยา.  เหตุยา เอกวีส.}

\proseitem{72}{นสนิทสฺสน\-สปฺปฏิฆํ นอาสวญฺเจว นอนาสวญฺจ ธมฺมํ ปฏิจฺจ อนิทสฺสน\-อปฺปฏิโฆ อาสโว เจว สาสโว จ ธมฺโม อุปฺปชฺชติ เหตุปจฺจยา.  เหตุยา ตีณิ.}

\prose{นสนิทสฺสน\-สปฺปฏิฆํ นอนาสวญฺเจว นโน จ อาสวํ ธมฺมํ ปฏิจฺจ สนิทสฺสน\-สปฺปฏิโฆ สาสโว เจว โน จ อาสโว ธมฺโม อุปฺปชฺชติ เหตุปจฺจยา.  เหตุยา เอกวีส.}

\proseitem{73}{นสนิทสฺสน\-สปฺปฏิฆํ นอาสวญฺเจว นอาสว\-วิปฺปยุตฺตญฺจ ธมฺมํ ปฏิจฺจ อนิทสฺสน\-อปฺปฏิโฆ อาสโว เจว อาสว\-สมฺปยุตฺโต จ ธมฺโม อุปฺปชฺชติ เหตุปจฺจยา.  เหตุยา ตีณิ.}

\prose{นสนิทสฺสน\-สปฺปฏิฆํ นอาสว\-วิปฺปยุตฺต\-ญฺเจว นโน จ อาสวํ ธมฺมํ ปฏิจฺจ อนิทสฺสน\-อปฺปฏิโฆ อาสว\-สมฺปยุตฺโต เจว โน จ อาสโว ธมฺโม อุปฺปชฺชติ เหตุปจฺจยา.  เหตุยา ตีณิ.}

\proseitem{74}{นสนิทสฺสน\-สปฺปฏิฆํ อาสว\-วิปฺปยุตฺตํ นสาสวํ ธมฺมํ ปฏิจฺจ สนิทสฺสน\-สปฺปฏิโฆ อาสว\-วิปฺปยุตฺโต สาสโว ธมฺโม อุปฺปชฺชติ เหตุปจฺจยา.  เหตุยา เอกวีส.}

\prose{นสนิทสฺสน\-สปฺปฏิฆํ อาสว\-วิปฺปยุตฺตํ นอนาสวํ ธมฺมํ ปจฺจยา อนิทสฺสน\-อปฺปฏิโฆ อาสว\-วิปฺปยุตฺโต อนาสโว ธมฺโม อุปฺปชฺชติ เหตุปจฺจยา.  เหตุยา ตีณิ.}

\csromanmarkh{22. สนิทสฺสน\-ตฺติก}\csromanmarkhii{20. สญฺโญชนาทิโคจฺฉก}\csromanheadernomark{2}{22. สนิทสฺสน\-ตฺติก 20. สญฺโญชนาทิโคจฺฉก}
\proseitem{75}{นสนิทสฺสน\-สปฺปฏิฆํ โนสญฺโญชนํ ธมฺมํ ปฏิจฺจ อนิทสฺสน\-อปฺปฏิโฆ สญฺโญชโน ธมฺโม อุปฺปชฺชติ เหตุปจฺจยา.  เหตุยา ตีณิ.}

\proseitem{76}{นสนิทสฺสน\-สปฺปฏิฆํ โนคนฺถํ ธมฺมํ ปฏิจฺจ อนิทสฺสน\-อปฺปฏิโฆ คนฺโถ ธมฺโม อุปฺปชฺชติ เหตุปจฺจยา.  เหตุยา ตีณิ.}

\csromanheader{2}{ธมฺม\-ปจฺจนียา\-นุโลม ติกทุก\-ปฏฺฐาน\-ปาฬิ}
\proseitem{77}{\csromanfolio{394}นสนิทสฺสน\-สปฺปฏิฆํ โนโอฆํ ธมฺมํ ปฏิจฺจ อนิทสฺสน\-อปฺปฏิโฆ โอโฆ ธมฺโม อุปฺปชฺชติ เหตุปจฺจยา.  เหตุยา ตีณิ.}

\proseitem{78}{นสนิทสฺสน\-สปฺปฏิฆํ โนโยคํ ธมฺมํ ปฏิจฺจ อนิทสฺสน\-อปฺปฏิโฆ โยโค ธมฺโม อุปฺปชฺชติ เหตุปจฺจยา.  เหตุยา ตีณิ.}

\proseitem{79}{นสนิทสฺสน\-สปฺปฏิฆํ โนนีวรณํ ธมฺมํ ปฏิจฺจ อนิทสฺสน\-อปฺปฏิโฆ นีวรโณ ธมฺโม อุปฺปชฺชติ เหตุปจฺจยา.  เหตุยา ตีณิ.}

\proseitem{80}{นสนิทสฺสน\-สปฺปฏิฆํ โนปรามาสํ ธมฺมํ ปฏิจฺจ อนิทสฺสน\-อปฺปฏิโฆ ปรามาโส ธมฺโม อุปฺปชฺชติ เหตุปจฺจยา.  เหตุยา ตีณิ. 22. สนิทสฺสน\-ตฺติก 55. มหนฺตรทุก}

\proseitem{81}{นสนิทสฺสน\-สปฺปฏิฆํ นสารมฺมณํ ธมฺมํ ปฏิจฺจ อนิทสฺสน\-อปฺปฏิโฆ สารมฺมโณ ธมฺโม อุปฺปชฺชติ เหตุปจฺจยา.  เหตุยา ตีณิ.}

\prose{นสนิทสฺสน\-สปฺปฏิฆํ นอนารมฺมณํ ธมฺมํ ปฏิจฺจ สนิทสฺสน\-สปฺปฏิโฆ อนารมฺมโณ ธมฺโม อุปฺปชฺชติ เหตุปจฺจยา.  เหตุยา เอกวีส.}

\proseitem{82}{นสนิทสฺสน\-สปฺปฏิฆํ โนจิตฺตํ ธมฺมํ ปฏิจฺจ อนิทสฺสน\-อปฺปฏิโฆ จิตฺโต ธมฺโม อุปฺปชฺชติ เหตุปจฺจยา.  เหตุยา ตีณิ.}

\prose{นสนิทสฺสน\-สปฺปฏิฆํ นโนจิตฺตํ ธมฺมํ ปฏิจฺจ สนิทสฺสน\-สปฺปฏิโฆ โนจิตฺโต ธมฺโม อุปฺปชฺชติ เหตุปจฺจยา.  เหตุยา เอกวีส.}

\proseitem{83}{นสนิทสฺสน\-สปฺปฏิฆํ นเจตสิกํ ธมฺมํ ปฏิจฺจ อนิทสฺสน\-อปฺปฏิโฆ เจตสิโก ธมฺโม อุปฺปชฺชติ เหตุปจฺจยา.  เหตุยา ตีณิ.}

\prose{นสนิทสฺสน\-สปฺปฏิฆํ นอเจตสิกํ ธมฺมํ ปฏิจฺจ สนิทสฺสน\-สปฺปฏิโฆ อเจตสิโก ธมฺโม อุปฺปชฺชติ เหตุปจฺจยา.  เหตุยา เอกวีส.}

\proseitem{84}{นสนิทสฺสน\-สปฺปฏิฆํ นจิตฺต\-สมฺปยุตฺตํ ธมฺมํ ปฏิจฺจ อนิทสฺสน\-อปฺปฏิโฆ จิตฺต\-สมฺปยุตฺโต ธมฺโม อุปฺปชฺชติ เหตุปจฺจยา.  เหตุยา ตีณิ.}

\csromanmarkh{22. สนิทสฺสน\-ตฺติก}\csromanmarkhii{55. มหนฺตรทุก}\csromanheadernomark{2}{22. สนิทสฺสน\-ตฺติก 55. มหนฺตรทุก}
\prose{\csromanfolio{395}นสนิทสฺสน\-สปฺปฏิฆํ นจิตฺต\-วิปฺปยุตฺตํ ธมฺมํ ปฏิจฺจ สนิทสฺสน\-สปฺปฏิโฆ จิตฺต\-วิปฺปยุตฺโต ธมฺโม อุปฺปชฺชติ เหตุปจฺจยา.  เหตุยา เอกวีส.}

\proseitem{85}{นสนิทสฺสน\-สปฺปฏิฆํ นจิตฺต\-สํสฏฺฐํ ธมฺมํ ปฏิจฺจ อนิทสฺสน\-อปฺปฏิโฆ จิตฺตสํสฏฺโฐ ธมฺโม อุปฺปชฺชติ เหตุปจฺจยา.  เหตุยา ตีณิ.}

\prose{นสนิทสฺสน\-สปฺปฏิฆํ นจิตฺต\-วิสํสฏฺฐํ ธมฺมํ ปฏิจฺจ สนิทสฺสน\-สปฺปฏิโฆ จิตฺต\-วิสํสฏฺโฐ ธมฺโม อุปฺปชฺชติ เหตุปจฺจยา.  เหตุยา เอกวีส.}

\proseitem{86}{นสนิทสฺสน\-สปฺปฏิฆํ โนจิตฺต\-สมุฏฺฐานํ ธมฺมํ ปฏิจฺจ สนิทสฺสน\-สปฺปฏิโฆ จิตฺต\-สมุฏฺฐาโน ธมฺโม อุปฺปชฺชติ เหตุปจฺจยา.  เหตุยา เอกวีส.}

\prose{นสนิทสฺสน\-สปฺปฏิฆํ นโนจิตฺตสมุฏฺฐานํ ธมฺมํ ปฏิจฺจ อนิทสฺสน\-อปฺปฏิโฆ โนจิตฺต\-สมุฏฺฐาโน ธมฺโม อุปฺปชฺชติ เหตุปจฺจยา.  เหตุยา ตีณิ.}

\proseitem{87}{นสนิทสฺสน\-สปฺปฏิฆํ โนจิตฺต\-สหภุํ ธมฺมํ ปฏิจฺจ อนิทสฺสน\-อปฺปฏิโฆ จิตฺตสหภู ธมฺโม อุปฺปชฺชติ เหตุปจฺจยา.  เหตุยา ปญฺจ.}

\prose{นสนิทสฺสน\-สปฺปฏิฆํ นโนจิตฺตสหภุํ ธมฺมํ ปฏิจฺจ สนิทสฺสน\-สปฺปฏิโฆ โนจิตฺตสหภู ธมฺโม อุปฺปชฺชติ เหตุปจฺจยา.  เหตุยา เอกวีส.}

\proseitem{88}{นสนิทสฺสน\-สปฺปฏิฆํ โนจิตฺตา\-นุปริวตฺติํ ธมฺมํ ปฏิจฺจ อนิทสฺสน\-อปฺปฏิโฆ จิตฺตา\-นุปริวตฺตี ธมฺโม อุปฺปชฺชติ เหตุปจฺจยา.  เหตุยา ปญฺจ.}

\prose{นสนิทสฺสน\-สปฺปฏิฆํ นโนจิตฺตนุปริวตฺติํ ธมฺมํ ปฏิจฺจ สนิทสฺสน\-สปฺปฏิโฆ โนจิตฺตา\-นุปริวตฺตี ธมฺโม อุปฺปชฺชติ เหตุปจฺจยา.  เหตุยา เอกวีส.}

\proseitem{89}{นสนิทสฺสน\-สปฺปฏิฆํ โนจิตฺต\-สํสฏฺฐ\-สมุฏฺฐานํ ธมฺมํ ปฏิจฺจ อนิทสฺสน\-อปฺปฏิโฆ จิตฺต\-สํสฏฺฐ\-สมุฏฺฐาโน ธมฺโม อุปฺปชฺชติ เหตุปจฺจยา.  เหตุยา ตีณิ.}

\prose{นสนิทสฺสน\-สปฺปฏิฆํ นโนจิตฺตสํสฏฺฐสมุฏฺฐานํ ธมฺมํ ปฏิจฺจ สนิทสฺสน\-สปฺปฏิโฆ โนจิตฺต\-สํสฏฺฐ\-สมุฏฺฐาโน ธมฺโม อุปฺปชฺชติ เหตุปจฺจยา.  เหตุยา เอกวีส.}

\csromanheader{2}{ธมฺม\-ปจฺจนียา\-นุโลม ติกทุก\-ปฏฺฐาน\-ปาฬิ}
\proseitem{90}{\csromanfolio{396}นสนิทสฺสน\-สปฺปฏิฆํ โนจิตฺต\-สํสฏฺฐ\-สมุฏฺฐาน\-สหภุํ ธมฺมํ ปฏิจฺจ อนิทสฺสน\-อปฺปฏิโฆ จิตฺต\-สํสฏฺฐ\-สมุฏฺฐาน\-สหภู ธมฺโม อุปฺปชฺชติ เหตุปจฺจยา.  เหตุยา ตีณิ.}

\prose{นสนิทสฺสน\-สปฺปฏิฆํ นโนจิตฺตสํสฏฺฐสมุฏฺฐานสหภุํ ธมฺมํ ปฏิจฺจ สนิทสฺสน\-สปฺปฏิโฆ โนจิตฺต\-สํสฏฺฐ\-สมุฏฺฐาน\-สหภู ธมฺโม อุปฺปชฺชติ เหตุปจฺจยา.  เหตุยา เอกวีส.}

\proseitem{91}{นสนิทสฺสน\-สปฺปฏิฆํ โนจิตฺต\-สํสฏฺฐ\-สมุฏฺฐานา\-นุปริวตฺติํ ธมฺมํ ปฏิจฺจ อนิทสฺสน\-อปฺปฏิโฆ จิตฺต\-สํสฏฺฐ\-สมุฏฺฐานา\-นุปริวตฺตี ธมฺโม อุปฺปชฺชติ เหตุปจฺจยา.  เหตุยา ตีณิ.}

\prose{นสนิทสฺสน\-สปฺปฏิฆํ นโนจิตฺตสํสฏฺฐสมุฏฺฐานานุปริวตฺติํ ธมฺมํ ปฏิจฺจ สนิทสฺสน\-สปฺปฏิโฆ โนจิตฺต\-สํสฏฺฐ\-สมุฏฺฐานา\-นุปริวตฺตี ธมฺโม อุปฺปชฺชติ เหตุปจฺจยา.  เหตุยา เอกวีส.}

\proseitem{92}{นสนิทสฺสน\-สปฺปฏิฆํ นอชฺฌตฺติกํ ธมฺมํ ปฏิจฺจ อนิทสฺสน\-สปฺปฏิโฆ อชฺฌตฺติโก ธมฺโม อุปฺปชฺชติ เหตุปจฺจยา.  เหตุยา เอกาทส.}

\prose{นสนิทสฺสน\-สปฺปฏิฆํ นพาหิรํ ธมฺมํ ปฏิจฺจ สนิทสฺสน\-สปฺปฏิโฆ พาหิโร ธมฺโม อุปฺปชฺชติ เหตุปจฺจยา.  เหตุยา เอกวีส.}

\proseitem{93}{นสนิทสฺสน\-สปฺปฏิฆํ โนอุปาทา ธมฺมํ ปฏิจฺจ สนิทสฺสน\-สปฺปฏิโฆ อุปาทา ธมฺโม อุปฺปชฺชติ เหตุปจฺจยา.  เหตุยา ปญฺจติํส.}

\prose{นสนิทสฺสน\-สปฺปฏิฆํ นโนอุปาทา ธมฺมํ ปฏิจฺจ อนิทสฺสน\-อปฺปฏิโฆ โนอุปาทา ธมฺโม อุปฺปชฺชติ เหตุปจฺจยา.  เหตุยา ตีณิ.}

\proseitem{94}{นสนิทสฺสน\-สปฺปฏิฆํ นอนุปาทินฺนํ ธมฺมํ ปฏิจฺจ สนิทสฺสน\-สปฺปฏิโฆ อนุปาทินฺโน ธมฺโม อุปฺปชฺชติ เหตุปจฺจยา.  เหตุยา เอกวีส.}

\csromanmarkh{22. สนิทสฺสน\-ตฺติก}\csromanmarkhii{69. อุปาทานทุก}\csromanheadernomark{2}{22. สนิทสฺสน\-ตฺติก 69. อุปาทานทุก}
\proseitem{95}{นสนิทสฺสน\-สปฺปฏิฆํ โนอุปาทานํ ธมฺมํ ปฏิจฺจ อนิทสฺสน\-อปฺปฏิโฆ อุปาทาโน ธมฺโม อุปฺปชฺชติ เหตุปจฺจยา.  เหตุยา ตีณิ.}

\vspace{\tipitakaparskip}
\csromancenter{สนิทสฺสน\-ตฺติก 83-99. ปิฏฺฐิทุก 22. สนิทสฺสน\-ตฺติก 75. กิเลสทุก}
\proseitem{96}{\csromanfolio{397}นสนิทสฺสน\-สปฺปฏิฆํ โนกิเลสํ ธมฺมํ ปฏิจฺจ อนิทสฺสน\-อปฺปฏิโฆ กิเลโส ธมฺโม อุปฺปชฺชติ เหตุปจฺจยา.  เหตุยา ตีณิ.}

\csromanheader{2}{22. สนิทสฺสน\-ตฺติก 83-99. ปิฏฺฐิทุก}
\proseitem{97}{นสนิทสฺสน\-สปฺปฏิฆํ นทสฺสเนน ปหาตพฺพํ ธมฺมํ ปจฺจยา อนิทสฺสน\-อปฺปฏิโฆ ทสฺสเนน ปหาตพฺโพ ธมฺโม อุปฺปชฺชติ เหตุปจฺจยา.  เหตุยา ตีณิ.}

\prose{นสนิทสฺสน\-สปฺปฏิฆํ นนทสฺสเนน ปหาตพฺพํ ธมฺมํ ปฏิจฺจ สนิทสฺสน\-สปฺปฏิโฆ นทสฺสเนน ปหาตพฺโพ ธมฺโม อุปฺปชฺชติ เหตุปจฺจยา.  เหตุยา เอกวีส.}

\proseitem{98}{นสนิทสฺสน\-สปฺปฏิฆํ นภาวนาย ปหาตพฺพํ ธมฺมํ ปจฺจยา อนิทสฺสน\-อปฺปฏิโฆ ภาวนาย ปหาตพฺโพ ธมฺโม อุปฺปชฺชติ เหตุปจฺจยา.  เหตุยา ตีณิ.}

\prose{นสนิทสฺสน\-สปฺปฏิฆํ นนภาวนาย ปหาตพฺพํ ธมฺมํ ปฏิจฺจ สนิทสฺสน\-สปฺปฏิโฆ นภาวนาย ปหาตพฺโพ ธมฺโม อุปฺปชฺชติ เหตุปจฺจยา.  เหตุยา เอกวีส.}

\proseitem{99}{นสนิทสฺสน\-สปฺปฏิฆํ นทสฺสเนน ปหาตพฺพ\-เหตุกํ ธมฺมํ ปฏิจฺจ อนิทสฺสน\-อปฺปฏิโฆ ทสฺสเนน ปหาตพฺพ\-เหตุโก ธมฺโม อุปฺปชฺชติ เหตุปจฺจยา. เหตุยา ตีณิ.}

\prose{นสนิทสฺสน\-สปฺปฏิฆํ นนทสฺสเนน ปหาตพฺพ\-เหตุกํ ธมฺมํ ปฏิจฺจ สนิทสฺสน\-สปฺปฏิโฆ นทสฺสเนน ปหาตพฺพ\-เหตุโก ธมฺโม อุปฺปชฺชติ เหตุปจฺจยา.  เหตุยา เอกวีส.}

\proseitem{100}{นสนิทสฺสน\-สปฺปฏิฆํ นภาวนาย ปหาตพฺพ\-เหตุกํ ธมฺมํ ปจฺจยา อนิทสฺสน\-อปฺปฏิโฆ ภาวนาย ปหาตพฺพ\-เหตุโก ธมฺโม อุปฺปชฺชติ เหตุปจฺจยา.  เหตุยา ตีณิ.}

\prose{นสนิทสฺสน\-สปฺปฏิฆํ นนภาวนาย ปหาตพฺพ\-เหตุกํ ธมฺมํ ปฏิจฺจ สนิทสฺสน\-สปฺปฏิโฆ นภาวนาย ปหาตพฺพ\-เหตุโก ธมฺโม อุปฺปชฺชติ เหตุปจฺจยา.  เหตุยา เอกวีส.}

\csromanheader{2}{ธมฺม\-ปจฺจนียา\-นุโลม ติกทุก\-ปฏฺฐาน\-ปาฬิ}
\proseitem{101}{\csromanfolio{398}นสนิทสฺสน\-สปฺปฏิฆํ นสวิตกฺกํ ธมฺมํ ปฏิจฺจ อนิทสฺสน\-อปฺปฏิโฆ สวิตกฺโก ธมฺโม อุปฺปชฺชติ เหตุปจฺจยา.  เหตุยา ตีณิ.}

\prose{นสนิทสฺสน\-สปฺปฏิฆํ นอวิตกฺกํ ธมฺมํ ปฏิจฺจ สนิทสฺสน\-สปฺปฏิโฆ อวิตกฺโก ธมฺโม อุปฺปชฺชติ เหตุปจฺจยา.  เหตุยา เอกวีส.}

\proseitem{102}{นสนิทสฺสน\-สปฺปฏิฆํ นสวิจารํ ธมฺมํ ปฏิจฺจ อนิทสฺสน\-อปฺปฏิโฆ สวิจาโร ธมฺโม อุปฺปชฺชติ เหตุปจฺจยา.  เหตุยา ตีณิ.}

\prose{นสนิทสฺสน\-สปฺปฏิฆํ นอวิจารํ ธมฺมํ ปฏิจฺจ สนิทสฺสน\-สปฺปฏิโฆ อวิจาโร ธมฺโม อุปฺปชฺชติ เหตุปจฺจยา.  เหตุยา เอกวีส.}

\proseitem{103}{นสนิทสฺสน\-สปฺปฏิฆํ นสปฺปีติกํ ธมฺมํ ปฏิจฺจ อนิทสฺสน\-อปฺปฏิโฆ สปฺปีติโก ธมฺโม อุปฺปชฺชติ เหตุปจฺจยา.  เหตุยา ตีณิ.}

\prose{นสนิทสฺสน\-สปฺปฏิฆํ นอปฺปีติกํ ธมฺมํ ปฏิจฺจ สนิทสฺสน\-สปฺปฏิโฆ อปฺปีติโก ธมฺโม อุปฺปชฺชติ เหตุปจฺจยา.  เหตุยา เอกวีส.}

\proseitem{104}{นสนิทสฺสน\-สปฺปฏิฆํ นปีติสหคตํ ธมฺมํ ปฏิจฺจ อนิทสฺสน\-อปฺปฏิโฆ ปีติสหคโต ธมฺโม อุปฺปชฺชติ เหตุปจฺจยา.  เหตุยา ตีณิ.}

\prose{นสนิทสฺสน\-สปฺปฏิฆํ นนปี\-ติ\-สหคตํ ธมฺมํ ปฏิจฺจ สนิทสฺสน\-สปฺปฏิโฆ นปีติสหคโต ธมฺโม อุปฺปชฺชติ เหตุปจฺจยา.  เหตุยา เอกวีส.}

\proseitem{105}{นสนิทสฺสน\-สปฺปฏิฆํ นสุขสหคตํ ธมฺมํ ปฏิจฺจ อนิทสฺสน\-อปฺปฏิโฆ สุขสหคโต ธมฺโม อุปฺปชฺชติ เหตุปจฺจยา.  เหตุยา ตีณิ.}

\prose{นสนิทสฺสน\-สปฺปฏิฆํ นนสุข\-สหคตํ ธมฺมํ ปฏิจฺจ สนิทสฺสน\-สปฺปฏิโฆ นสุขสหคโต ธมฺโม อุปฺปชฺชติ เหตุปจฺจยา.  เหตุยา เอกวีส.}

\proseitem{106}{นสนิทสฺสน\-สปฺปฏิฆํ นอุเปกฺขาสหคตํ ธมฺมํ ปฏิจฺจ อนิทสฺสน\-อปฺปฏิโฆ อุเปกฺขา\-สหคโต ธมฺโม อุปฺปชฺชติ เหตุปจฺจยา.  เหตุยา ตีณิ.}

\csromanheader{2}{22. สนิทสฺสน\-ตฺติก 83-99. ปิฏฺฐิทุก}
\prose{\csromanfolio{399}นสนิทสฺสน\-สปฺปฏิฆํ นนอุเปกฺขาสหคตํ ธมฺมํ ปฏิจฺจ สนิทสฺสน\-สปฺปฏิโฆ นอุเปกฺขาสหคโต ธมฺโม อุปฺปชฺชติ เหตุปจฺจยา.  เหตุยา เอกวีส.}

\proseitem{107}{นสนิทสฺสน\-สปฺปฏิฆํ นกามาวจรํ ธมฺมํ ปฏิจฺจ สนิทสฺสน\-สปฺปฏิโฆ กามาวจโร ธมฺโม อุปฺปชฺชติ เหตุปจฺจยา.  เหตุยา เอกวีส.}

\prose{นสนิทสฺสน\-สปฺปฏิฆํ นนกามา\-วจรํ ธมฺมํ ปฏิจฺจ อนิทสฺสน\-อปฺปฏิโฆ นกามาวจโร ธมฺโม อุปฺปชฺชติ เหตุปจฺจยา.  เหตุยา ตีณิ.}

\proseitem{108}{นสนิทสฺสน\-สปฺปฏิฆํ นรูปาวจรํ ธมฺมํ ปฏิจฺจ อนิทสฺสน\-อปฺปฏิโฆ รูปาวจโร ธมฺโม อุปฺปชฺชติ เหตุปจฺจยา.  เหตุยา ติณิ.}

\prose{นสนิทสฺสน\-สปฺปฏิฆํ นนรูปาวจรํ ธมฺมํ ปฏิจฺจ สนิทสฺสน\-สปฺปฏิโฆ นรูปาวจโร ธมฺโม อุปฺปชฺชติ เหตุปจฺจยา.  เหตุยา เอกวีส.}

\proseitem{109}{นสนิทสฺสน\-สปฺปฏิฆํ นอรูปาวจรํ ธมฺมํ ปจฺจยา อนิทสฺสน\-อปฺปฏิโฆ อรูปาวจโร ธมฺโม อุปฺปชฺชติ เหตุปจฺจยา.  เหตุยา ตีณิ.}

\prose{นสนิทสฺสน\-สปฺปฏิฆํ นนอรูปาวจรํ ธมฺมํ ปฏิจฺจ สนิทสฺสน\-สปฺปฏิโฆ นอรูปาวจโร ธมฺโม อุปฺปชฺชติ เหตุปจฺจยา.  เหตุยา เอกวีส.}

\proseitem{110}{นสนิทสฺสน\-สปฺปฏิฆํ นปริยาปนฺนํ ธมฺมํ ปฏิจฺจ สนิทสฺสน\-สปฺปฏิโฆ ปริยาปนฺโน ธมฺโม อุปฺปชฺชติ เหตุปจฺจยา.  เหตุยา เอกวีส.}

\prose{นสนิทสฺสน\-สปฺปฏิฆํ นอปริยาปนฺนํ ธมฺมํ ปจฺจยา อนิทสฺสน\-อปฺปฏิโฆ อปริยาปนฺโน ธมฺโม อุปฺปชฺชติ เหตุปจฺจยา.  เหตุยา ตีณิ.}

\proseitem{111}{นสนิทสฺสน\-สปฺปฏิฆํ นนิยฺยานิกํ ธมฺมํ ปจฺจยา อนิทสฺสน\-อปฺปฏิโฆ นิยฺยานิโก ธมฺโม อุปฺปชฺชติ เหตุปจฺจยา.  เหตุยา ตีณิ.}

\prose{นสนิทสฺสน\-สปฺปฏิฆํ นอนิยฺยานิกํ ธมฺมํ ปฏิจฺจ สนิทสฺสน\-สปฺปฏิโฆ อนิยฺยานิโก ธมฺโม อุปฺปชฺชติ เหตุปจฺจยา.  เหตุยา เอกวีส.}

\csromanheader{2}{ธมฺม\-ปจฺจนียา\-นุโลม ติกทุก\-ปฏฺฐาน\-ปาฬิ}
\proseitem{112}{\csromanfolio{400}นสนิทสฺสน\-สปฺปฏิฆํ นนิยตํ ธมฺมํ ปจฺจยา อนิทสฺสน\-อปฺปฏิโฆ นิยโต ธมฺโม อุปฺปชฺชติ เหตุปจฺจยา.  เหตุยา ตีณิ.}

\prose{นสนิทสฺสน\-สปฺปฏิฆํ นอนิยตํ ธมฺมํ ปฏิจฺจ สนิทสฺสน\-สปฺปฏิโฆ อนิยโต ธมฺโม อุปฺปชฺชติ เหตุปจฺจยา.  เหตุยา เอกวีส.}

\proseitem{113}{นสนิทสฺสน\-สปฺปฏิฆํ นสอุตฺตรํ ธมฺมํ ปฏิจฺจ สนิทสฺสน\-สปฺปฏิโฆ สอุตฺตโร ธมฺโม อุปฺปชฺชติ เหตุปจฺจยา.  เหตุยา เอกวีส.}

\prose{นสนิทสฺสน\-สปฺปฏิฆํ นอนุตฺตรํ ธมฺมํ ปจฺจยา อนิทสฺสน\-อปฺปฏิโฆ อนุตฺตโร ธมฺโม อุปฺปชฺชติ เหตุปจฺจยา.  เหตุยา ตีณิ.}

\csromanmarkh{22. สนิทสฺสน\-ตฺติก}\csromanmarkhii{100. สรณทุก}\csromanheadernomark{2}{22. สนิทสฺสน\-ตฺติก 100. สรณทุก}
\proseitem{114}{นสนิทสฺสน\-สปฺปฏิฆํ นสรณํ ธมฺมํ ปจฺจยา อนิทสฺสน\-อปฺปฏิโฆ สรโณ ธมฺโม อุปฺปชฺชติ เหตุปจฺจยา.  เหตุยา ตีณิ.}

\prose{นสนิทสฺสน\-สปฺปฏิฆํ นอรณํ ธมฺมํ ปฏิจฺจ สนิทสฺสน\-สปฺปฏิโฆ อรโณ ธมฺโม อุปฺปชฺชติ เหตุปจฺจยา. สตฺต.}

\prose{นอนิทสฺสน\-สปฺปฏิฆํ นอรณํ ธมฺมํ ปฏิจฺจ สนิทสฺสน\-สปฺปฏิโฆ อรโณ ธมฺโม อุปฺปชฺชติ เหตุปจฺจยา. สตฺต.}

\prose{นสนิทสฺสน\-สปฺปฏิฆํ นอรณญฺจ นอนิทสฺสน\-สปฺปฏิฆํ นอรณญฺจ ธมฺมํ ปฏิจฺจ สนิทสฺสน\-สปฺปฏิโฆ อรโณ ธมฺโม อุปฺปชฺชติ เหตุปจฺจยา. สตฺต.  เหตุยา เอกวีส.}

\vspace{\tipitakaparskip}
\csromancenter{(สหชาตวารมฺปิ ปฏิจฺจ\-วาร\-สทิสํ วิตฺถาเรตพฺพํ.)}
\csromanheader{2}{\textbf{เหตุ อารมฺมณ}}
\proseitem{115}{นสนิทสฺสน\-สปฺปฏิโฆ นอรโณ ธมฺโม สนิทสฺสน\-สปฺปฏิฆสฺส อรณสฺส ธมฺมสฺส เหตุปจฺจเยน ปจฺจโย. สตฺต.}

\prose{นอนิทสฺสน\-สปฺปฏิโฆ นอรโณ ธมฺโม สนิทสฺสน\-สปฺปฏิฆสฺส อรณสฺส ธมฺมสฺส เหตุปจฺจเยน ปจฺจโย. สตฺต.}

\csromanmarkh{22. สนิทสฺสน\-ตฺติก}\csromanmarkhii{100. สรณทุก}\csromanheadernomark{2}{22. สนิทสฺสน\-ตฺติก 100. สรณทุก}
\prose{\csromanfolio{401}นสนิทสฺสน\-สปฺปฏิโฆ นอรโณ จ นอนิทสฺสน\-สปฺปฏิโฆ นอรโณ จ ธมฺมา สนิทสฺสน\-สปฺปฏิฆสฺส อรณสฺส ธมฺมสฺส เหตุปจฺจเยน ปจฺจโย. สตฺต.}

\proseitem{116}{นสนิทสฺสน\-สปฺปฏิโฆ นอรโณ ธมฺโม อนิทสฺสน\-อปฺปฏิฆสฺส อรณสฺส ธมฺมสฺส อารมฺมณ\-ปจฺจเยน ปจฺจโย.  นอนิทสฺสน\-สปฺปฏิโฆ นอรโณ ธมฺโม อนิทสฺสน\-อปฺปฏิฆสฺส อรณสฺส ธมฺมสฺส อารมฺมณ\-ปจฺจเยน ปจฺจโย.  นสนิทสฺสน\-สปฺปฏิโฆ นอรโณ จ นอนิทสฺสน\-สปฺปฏิโฆ นอรโณ จ ธมฺมา อนิทสฺสน\-อปฺปฏิฆสฺส อรณสฺส ธมฺมสฺส อารมฺมณ\-ปจฺจเยน ปจฺจโย.}

\proseitem{117}{เหตุยา เอกวีส, อารมฺมเณ ตีณิ ฯเปฯ อวิคเต เอกวีส.}

\prosewithcloser{(ยถา กุสลตฺติเก ปญฺหาวารสฺส อนุโลมมฺปิ ปจฺจนียมฺปิ อนุโลมปจฺจนียมฺปิ ปจฺจนียานุโลมมฺปิ คณิตํ, เอวํ คเณตพฺพํ.)}{ธมฺม\-ปจฺจนียา\-นุโลเม ติกทุก\-ปฏฺฐานํ นิฏฺฐิตํ.}
\pitaka{อภิธมฺม\-ปิฏก}
\csromanheader{2}{ธมฺม\-ปจฺจนียา\-นุโลม}
\namakkaram{นโม ตสฺส ภควโต อรหโต สมฺมา\-สมฺพุทฺธสฺส.}
\csromanmarkh{1. กุสลตฺติก}\csromanmarkhii{1. เวทนาตฺติก}\csromanheadernomark{2}{1. \textbf{กุสลตฺติก 1. เวทนาตฺติก}}
\prosecont{\csromanfolio{403}นกุสลํ นสุขาย เวทนาย สมฺปยุตฺตํ ธมฺมํ ปฏิจฺจ อกุสโล สุขาย เวทนาย สมฺปยุตฺโต ธมฺโม อุปฺปชฺชติ เหตุปจฺจยา.  นกุสลํ นสุขาย เวทนาย สมฺปยุตฺตํ ธมฺมํ ปฏิจฺจ อพฺยากโต สุขาย เวทนาย สมฺปยุตฺโต ธมฺโม อุปฺปชฺชติ เหตุปจฺจยา. เทฺว.}

\proseitem{1}{นอกุสลํ นสุขาย เวทนาย สมฺปยุตฺตํ ธมฺมํ ปฏิจฺจ กุสโล สุขาย เวทนาย สมฺปยุตฺโต ธมฺโม อุปฺปชฺชติ เหตุปจฺจยา.  นอกุสลํ นสุขาย เวทนาย สมฺปยุตฺตํ ธมฺมํ ปฏิจฺจ อพฺยากโต สุขาย เวทนาย สมฺปยุตฺโต ธมฺโม อุปฺปชฺชติ เหตุปจฺจยา. เทฺว.}

\prose{นอพฺยากตํ นสุขาย เวทนาย สมฺปยุตฺตํ ธมฺมํ ปฏิจฺจ กุสโล สุขาย เวทนาย สมฺปยุตฺโต ธมฺโม อุปฺปชฺชติ เหตุปจฺจยา.  นอพฺยากตํ นสุขาย เวทนาย สมฺปยุตฺตํ ธมฺมํ ปฏิจฺจ อกุสโล สุขาย เวทนาย สมฺปยุตฺโต ธมฺโม อุปฺปชฺชติ เหตุปจฺจยา. เทฺว.}

\prose{นกุสลํ นสุขาย เวทนาย สมฺปยุตฺตญฺจ นอพฺยากตํ นสุขาย เวทนาย สมฺปยุตฺตญฺจ ธมฺมํ ปฏิจฺจ อกุสโล สุขาย เวทนาย สมฺปยุตฺโต ธมฺโม อุปฺปชฺชติ เหตุปจฺจยา. เอกํ.}

\csromanheader{2}{ธมฺม\-ปจฺจนียา\-นุโลม ติกติก\-ปฏฺฐาน\-ปาฬิ}
\prose{\csromanfolio{404}นอกุสลํ นสุขาย เวทนาย สมฺปยุตฺตญฺจ นอพฺยากตํ นสุขาย เวทนาย สมฺปยุตฺตญฺจ ธมฺมํ ปฏิจฺจ กุสโล สุขาย เวทนาย สมฺปยุตฺโต ธมฺโม อุปฺปชฺชติ เหตุปจฺจยา. เอกํ.}

\prose{นกุสลํ นสุขาย เวทนาย สมฺปยุตฺตญฺจ นอกุสลํ นสุขาย เวทนาย สมฺปยุตฺตญฺจ ธมฺมํ ปฏิจฺจ อพฺยากโต สุขาย เวทนาย สมฺปยุตฺโต ธมฺโม อุปฺปชฺชติ เหตุปจฺจยา. เอกํ.  เหตุยา นว.}

\proseitem{2}{นกุสลํ นทุกฺขาย เวทนาย สมฺปยุตฺตํ ธมฺมํ ปฏิจฺจ อกุสโล ทุกฺขาย เวทนาย สมฺปยุตฺโต ธมฺโม อุปฺปชฺชติ เหตุปจฺจยา. เอกํ.}

\prose{นอพฺยากตํ นทุกฺขาย เวทนาย สมฺปยุตฺตํ ธมฺมํ ปฏิจฺจ อกุสโล ทุกฺขาย เวทนาย สมฺปยุตฺโต ธมฺโม อุปฺปชฺชติ เหตุปจฺจยา. เอกํ.}

\prose{นกุสลํ นทุกฺขาย เวทนาย สมฺปยุตฺตญฺจ นอพฺยากตํ นทุกฺขาย เวทนาย สมฺปยุตฺตญฺจ ธมฺมํ ปฏิจฺจ อกุสโล ทุกฺขาย เวทนาย สมฺปยุตฺโต ธมฺโม อุปฺปชฺชติ เหตุปจฺจยา. เอกํ. เหตุยา ตีณิ.}

\proseitem{3}{นกุสลํ นอทุกฺขม\-สุขาย เวทนาย สมฺปยุตฺตํ ธมฺมํ ปฏิจฺจ อกุสโล อทุกฺขม\-สุขาย เวทนาย สมฺปยุตฺโต ธมฺโม อุปฺปชฺชติ เหตุปจฺจยา.  นกุสลํ นอทุกฺขม\-สุขาย เวทนาย สมฺปยุตฺตํ ธมฺมํ ปฏิจฺจ อพฺยากโต อทุกฺขม\-สุขาย เวทนาย สมฺปยุตฺโต ธมฺโม อุปฺปชฺชติ เหตุปจฺจยา. เทฺว.}

\prose{นอกุสลํ นอทุกฺขม\-สุขาย เวทนาย สมฺปยุตฺตํ ธมฺมํ ปฏิจฺจ กุสโล อทุกฺขม\-สุขาย เวทนาย สมฺปยุตฺโต ธมฺโม อุปฺปชฺชติ เหตุปจฺจยา.  นอกุสลํ นอทุกฺขม\-สุขาย เวทนาย สมฺปยุตฺตํ ธมฺมํ ปฏิจฺจ อพฺยากโต อทุกฺขม\-สุขาย เวทนาย สมฺปยุตฺโต ธมฺโม อุปฺปชฺชติ เหตุปจฺจยา. เทฺว.}

\prose{นอพฺยากตํ นอทุกฺขม\-สุขาย เวทนาย สมฺปยุตฺตํ ธมฺมํ ปฏิจฺจ กุสโล อทุกฺขม\-สุขาย เวทนาย สมฺปยุตฺโต ธมฺโม อุปฺปชฺชติ เหตุปจฺจยา.  นอพฺยากตํ นอทุกฺขม\-สุขาย เวทนาย สมฺปยุตฺตํ ธมฺมํ ปฏิจฺจ อกุสโล อทุกฺขม\-สุขาย เวทนาย สมฺปยุตฺโต ธมฺโม อุปฺปชฺชติ เหตุปจฺจยา. เทฺว.}

\prose{นกุสลํ นอทุกฺขม\-สุขาย เวทนาย สมฺปยุตฺตญฺจ นอพฺยากตํ นอทุกฺขม\-สุขาย เวทนาย สมฺปยุตฺตญฺจ ธมฺมํ ปฏิจฺจ อกุสโล อทุกฺขม\-สุขาย เวทนาย สมฺปยุตฺโต ธมฺโม อุปฺปชฺชติ เหตุปจฺจยา. เอกํ.}

\csromanmarkh{1. กุสลตฺติก}\csromanmarkhii{3. อุปาทินฺนตฺติก}\csromanheadernomark{2}{1. กุสลตฺติก 3. อุปาทินฺนตฺติก}
\prose{\csromanfolio{405}นอกุสลํ นอทุกฺขม\-สุขาย เวทนาย สมฺปยุตฺตญฺจ นอพฺยากตํ นอทุกฺขม\-สุขาย เวทนาย สมฺปยุตฺตญฺจ ธมฺมํ ปฏิจฺจ กุสโล อทุกฺขม\-สุขาย เวทนาย สมฺปยุตฺโต ธมฺโม อุปฺปชฺชติ เหตุปจฺจยา. เอกํ.}

\prose{นกุสลํ นอทุกฺขม\-สุขาย เวทนาย สมฺปยุตฺตญฺจ นอกุสลํ นอทุกฺขม\-สุขาย เวทนาย สมฺปยุตฺตญฺจ ธมฺมํ ปฏิจฺจ อพฺยากโต อทุกฺขม\-สุขาย เวทนาย สมฺปยุตฺโต ธมฺโม อุปฺปชฺชติ เหตุปจฺจยา. เอกํ.  เหตุยา นว.}

\csromanmarkh{1. กุสลตฺติก}\csromanmarkhii{2. วิปากตฺติก}\csromanheadernomark{2}{1. \textbf{กุสลตฺติก 2. วิปากตฺติก}}
\proseitem{4}{นกุสลํ นวิปากํ ธมฺมํ ปฏิจฺจ อพฺยากโต วิปาโก ธมฺโม อุปฺปชฺชติ เหตุปจฺจยา. เอกํ.}

\prose{นอกุสลํ นวิปากํ ธมฺมํ ปฏิจฺจ อพฺยากโต วิปาโก ธมฺโม อุปฺปชฺชติ เหตุปจฺจยา. เอกํ.}

\prose{นกุสลํ นวิปากญฺจ นอกุสลํ นวิปากญฺจ ธมฺมํ ปฏิจฺจ อพฺยากโต วิปาโก ธมฺโม อุปฺปชฺชติ เหตุปจฺจยา. เอกํ.  เหตุยา ตีณิ.}

\proseitem{5}{นกุสลํ นวิปาก\-ธมฺม\-ธมฺมํ ปจฺจยา กุสโล วิปาก\-ธมฺม\-ธมฺโม อุปฺปชฺชติ เหตุปจฺจยา.  นกุสลํ นวิปาก\-ธมฺม\-ธมฺมํ ปจฺจยา อกุสโล วิปาก\-ธมฺม\-ธมฺโม อุปฺปชฺชติ เหตุปจฺจยา. เทฺว.}

\prose{นอกุสลํ นวิปาก\-ธมฺม\-ธมฺมํ ปจฺจยา อกุสโล วิปาก\-ธมฺม\-ธมฺโม อุปฺปชฺชติ เหตุปจฺจยา. เทฺว.}

\prose{นกุสลํ นวิปาก\-ธมฺม\-ธมฺมญฺจ นอกุสลํ นวิปาก\-ธมฺม\-ธมฺมญฺจ ธมฺมํ ปจฺจยา กุสโล วิปาก\-ธมฺม\-ธมฺโม อุปฺปชฺชติ เหตุปจฺจยา. เทฺว. (สํขิตฺตํ.) . เหตุยา ฉ ปญฺหา.}

\proseitem{6}{นกุสลํ นเน\-ว\-วิปาก\-น\-วิปากธมฺม\-ธมฺมํ ปฏิจฺจ อพฺยากโต เนว\-วิปาก\-น\-วิปาก\-ธมฺม\-ธมฺโม อุปฺปชฺชติ เหตุปจฺจยา.  เหตุยา ฉ.}

\csromanmarkh{1. กุสลตฺติก}\csromanmarkhii{3. อุปาทินฺนตฺติก}\csromanheadernomark{2}{1. กุสลตฺติก 3. อุปาทินฺนตฺติก}
\proseitem{7}{นกุสโล นอุปาทินฺนุปาทานิโย ธมฺโม อพฺยากตสฺส อุปาทินฺนุ\-ปาทานิยสฺส ธมฺมสฺส อารมฺมณ\-ปจฺจเยน ปจฺจโย. ฉ ปญฺหา.}

\csromanheader{2}{ธมฺม\-ปจฺจนียา\-นุโลม ติกติก\-ปฏฺฐาน\-ปาฬิ}
\prose{\csromanfolio{406}นกุสลํ นอนุปาทินฺนุ\-ปาทานิยํ ธมฺมํ ปฏิจฺจ อพฺยากโต อนุปาทินฺนุ\-ปาทานิโย ธมฺโม อุปฺปชฺชติ เหตุปจฺจยา.  เหตุยา ปญฺจ.}

\prose{นกุสลํ นอนุปาทินฺน\-อนุปาทานิยํ ธมฺมํ ปจฺจยา กุสโล อนุปาทินฺน\-อนุปาทานิโย ธมฺโม อุปฺปชฺชติ เหตุปจฺจยา.  เหตุยา ฉ.}

\csromanmarkh{1. กุสลตฺติก}\csromanmarkhii{4. สํกิลิฏฺฐ\-ตฺติก}\csromanheadernomark{2}{1. กุสลตฺติก 4. สํกิลิฏฺฐ\-ตฺติก}
\proseitem{8}{นกุสลํ นสํกิลิฏฺฐ\-สํกิเลสิกํ ธมฺมํ ปจฺจยา อกุสโล สํกิลิฏฺฐ\-สํกิเลสิโก ธมฺโม อุปฺปชฺชติ เหตุปจฺจยา. (อกุสลาเนว ตีณิ.)}

\prose{นกุสลํ นอสํกิลิฏฺฐ\-สํกิเลสิกํ ธมฺมํ ปฏิจฺจ อพฺยากโต อสํกิลิฏฺฐ\-อสํกิเลสิโก ธมฺโม อุปฺปชฺชติ เหตุปจฺจยา.  เหตุยา ฉ.}

\prose{นกุสลํ นอสํกิลิฏฺฐ\-อสํกิเลสิกํ ธมฺมํ ปจฺจยา กุสโล อสํกิลิฏฺฐ\-อสํกิเลสิโก ธมฺโม อุปฺปชฺชติ เหตุปจฺจยา.  เหตุยา ฉ.}

\csromanmarkh{1. กุสลตฺติก}\csromanmarkhii{5. วิตกฺกตฺติก}\csromanheadernomark{2}{1. \textbf{กุสลตฺติก 5. วิตกฺกตฺติก}}
\proseitem{9}{นกุสลํ นสวิตกฺก\-สวิจารํ ธมฺมํ ปฏิจฺจ อกุสโล สวิตกฺก\-สวิจาโร ธมฺโม อุปฺปชฺชติ เหตุปจฺจยา.  เหตุยา นว.}

\prose{นกุสลํ นอวิตกฺก\-วิจารมตฺตํ ธมฺมํ ปฏิจฺจ อกุสโล อวิตกฺก\-วิจารมตฺโต ธมฺโม อุปฺปชฺชติ เหตุปจฺจยา.  เหตุยา นว.}

\prose{นกุสลํ นอวิตกฺก\-อวิจารํ ธมฺมํ ปฏิจฺจ อพฺยากโต อวิตกฺก\-อวิจาโร ธมฺโม อุปฺปชฺชติ เหตุปจฺจยา.  เหตุยา ทฺวาทส.}

\csromanmarkh{1. กุสลตฺติก}\csromanmarkhii{6. ปีติตฺติก}\csromanheadernomark{2}{1. \textbf{กุสลตฺติก 6. ปีติตฺติก}}
\proseitem{10}{นกุสลํ นปีติสหคตํ ธมฺมํ ปฏิจฺจ อกุสโล ปีติสหคโต ธมฺโม อุปฺปชฺชติ เหตุปจฺจยา.  เหตุยา นว.}

\prose{นกุสลํ นสุขสหคตํ ธมฺมํ ปฏิจฺจ อกุสโล สุขสหคโต ธมฺโม อุปฺปชฺชติ เหตุปจฺจยา.  เหตุยา นว.}

\prose{นกุสลํ นอุเปกฺขาสหคตํ ธมฺมํ ปฏิจฺจ อกุสโล อุเปกฺขา\-สหคโต ธมฺโม อุปฺปชฺชติ เหตุปจฺจยา.  เหตุยา นว.}

\csromanmarkh{1. กุสลตฺติก}\csromanmarkhii{9. อาจยคามิตฺติก 1. กุสลตฺติก 7. ทสฺสนตฺติก}\csromanheadernomark{2}{1. กุสลตฺติก 9. อาจยคามิตฺติก \textbf{1. กุสลตฺติก 7. ทสฺสนตฺติก}}
\proseitem{11}{\csromanfolio{407}นกุสลํ นทสฺสเนน ปหาตพฺพํ ธมฺมํ ปจฺจยา อกุสโล ทสฺสเนน ปหาตพฺโพ ธมฺโม อุปฺปชฺชติ เหตุปจฺจยา.  เหตุยา ตีณิ.}

\prose{นกุสลํ นภาวนาย ปหาตพฺพํ ธมฺมํ ปจฺจยา อกุสโล ภาวนาย ปหาตพฺโพ ธมฺโม อุปฺปชฺชติ เหตุปจฺจยา.  เหตุยา ตีณิ.}

\prose{นกุสลํ นเน\-ว\-ทสฺสเนน นภาวนาย ปหาตพฺพํ ธมฺมํ ปฏิจฺจ อพฺยากโต เนวทสฺสเนน นภาวนาย ปหาตพฺโพ ธมฺโม อุปฺปชฺชติ เหตุปจฺจยา.  เหตุยา ตีณิ.}

\csromanmarkh{1. กุสลตฺติก}\csromanmarkhii{8. ทสฺสน\-เหตุ\-ตฺติก}\csromanheadernomark{2}{1. กุสลตฺติก 8. ทสฺสน\-เหตุ\-ตฺติก}
\proseitem{12}{นกุสลํ นทสฺสเนน ปหาตพฺพ\-เหตุกํ ธมฺมํ ปฏิจฺจ อกุสโล ทสฺสเนน ปหาตพฺพ\-เหตุโก ธมฺโม อุปฺปชฺชติ เหตุปจฺจยา.  เหตุยา ตีณิ.}

\prose{นกุสลํ นภาวนาย ปหาตพฺพ\-เหตุกํ ธมฺมํ ปฏิจฺจ อกุสโล ภาวนาย ปหาตพฺพ\-เหตุโก ธมฺโม อุปฺปชฺชติ เหตุปจฺจยา.  เหตุยา ตีณิ.}

\prose{นกุสลํ นเน\-ว\-ทสฺสเนน นภาวนาย ปหาตพฺพ\-เหตุกํ ธมฺมํ ปฏิจฺจ อพฺยากโต เนวทสฺสเนน นภาวนาย ปหาตพฺพ\-เหตุโก ธมฺโม อุปฺปชฺชติ เหตุปจฺจยา.  เหตุยา ตีณิ.}

\csromanmarkh{1. กุสลตฺติก}\csromanmarkhii{9. อาจยคามิตฺติก}\csromanheadernomark{2}{1. กุสลตฺติก 9. อาจยคามิตฺติก}
\proseitem{13}{นกุสลํ นอาจยคามิํ ธมฺมํ ปจฺจยา กุสโล อาจยคามี ธมฺโม อุปฺปชฺชติ เหตุปจฺจยา.  เหตุยา ฉ.}

\prose{นกุสลํ นอปจยคามิํ ธมฺมํ ปจฺจยา กุสโล อปจยคามี ธมฺโม อุปฺปชฺชติ เหตุปจฺจยา.  เหตุยา ตีณิ.}

\prose{นกุสลํ นเน\-วา\-จยคามิ\-นา\-ปจยคามิํ ธมฺมํ ปฏิจฺจ อพฺยากโต เนวา\-จยคามิ\-นา\-ปจยคามี ธมฺโม อุปฺปชฺชติ เหตุปจฺจยา.  เหตุยา ปญฺจ.}

\csromanmarkh{ธมฺม\-ปจฺจนียา\-นุโลม ติกติก\-ปฏฺฐาน\-ปาฬิ}\csromanmarkhii{1. กุสลตฺติก 10. เสกฺขตฺติก}\csromanheadernomark{2}{ธมฺม\-ปจฺจนียา\-นุโลม ติกติก\-ปฏฺฐาน\-ปาฬิ 1. กุสลตฺติก 10. เสกฺขตฺติก}
\proseitem{14}{\csromanfolio{408}นกุสลํ นเสกฺขํ ธมฺมํ ปจฺจยา กุสโล เสกฺโข ธมฺโม อุปฺปชฺชติ เหตุปจฺจยา.  เหตุยา ฉ.}

\prose{นกุสลํ นอเสกฺขํ ธมฺมํ ปจฺจยา อพฺยากโต อเสกฺโข ธมฺโม อุปฺปชฺชติ เหตุปจฺจยา.  เหตุยา ตีณิ.}

\prose{นกุสลํ นเน\-ว\-เสกฺข\-นาเสกฺขํ ธมฺมํ ปฏิจฺจ อพฺยากโต เนว\-เสกฺข\-นา\-เสกฺโข ธมฺโม อุปฺปชฺชติ เหตุปจฺจยา.  เหตุยา ปญฺจ.}

\csromanmarkh{1. กุสลตฺติก}\csromanmarkhii{11. ปริตฺตตฺติก}\csromanheadernomark{2}{1. \textbf{กุสลตฺติก} 11. ปริตฺตตฺติก}
\proseitem{15}{นกุสลํ นปริตฺตํ ธมฺมํ ปฏิจฺจ อพฺยากโต ปริตฺโต ธมฺโม อุปฺปชฺชติ เหตุปจฺจยา.  เหตุยา ปญฺจ.}

\prose{นกุสลํ นมหคฺคตํ ธมฺมํ ปฏิจฺจ อพฺยากโต มหคฺคโต ธมฺโม อุปฺปชฺชติ เหตุปจฺจยา.  เหตุยา ตีณิ.}

\prose{นกุสลํ นอปฺปมาณํ ธมฺมํ ปจฺจยา กุสโล อปฺปมาโณ ธมฺโม อุปฺปชฺชติ เหตุปจฺจยา.  เหตุยา ฉ.}

\csromanmarkh{1. กุสลตฺติก}\csromanmarkhii{12. ปริตฺตา\-รมฺมณ\-ตฺติก}\csromanheadernomark{2}{1. กุสลตฺติก 12. ปริตฺตา\-รมฺมณ\-ตฺติก}
\proseitem{16}{นกุสลํ นปริตฺตา\-รมฺมณํ ธมฺมํ ปฏิจฺจ อพฺยากโต ปริตฺตารมฺมโณ ธมฺโม อุปฺปชฺชติ เหตุปจฺจยา.  เหตุยา ตีณิ.}

\prose{นกุสลํ นมหคฺคตา\-รมฺมณํ ธมฺมํ ปจฺจยา กุสโล มหคฺคตา\-รมฺมโณ ธมฺโม อุปฺปชฺชติ เหตุปจฺจยา.  เหตุยา นว.}

\prose{นกุสลํ นอปฺปมาณา\-รมฺมณํ ธมฺมํ ปจฺจยา กุสโล อปฺปมาณารมฺมโณ ธมฺโม อุปฺปชฺชติ เหตุปจฺจยา.  เหตุยา ฉ.}

\csromanmarkh{1. กุสลตฺติก}\csromanmarkhii{13. หีนตฺติก}\csromanheadernomark{2}{1. \textbf{กุสลตฺติก} 13. หีนตฺติก}
\proseitem{17}{นกุสลํ นหีนํ ธมฺมํ ปจฺจยา อกุสโล หีโน ธมฺโม อุปฺปชฺชติ เหตุปจฺจยา.  เหตุยา ตีณิ.}

\prose{นกุสลํ นมชฺฌิมํ ธมฺมํ ปฏิจฺจ อพฺยากโต มชฺฌิโม ธมฺโม อุปฺปชฺชติ เหตุปจฺจยา.  เหตุยา ฉ.}

\csromanheader{2}{1. กุสลตฺติก 17-18. อตีต\-ตฺติก\-ทฺวย}
\prose{\csromanfolio{409}นกุสลํ นปณีตํ ธมฺมํ ปจฺจยา กุสโล ปณีโต ธมฺโม อุปฺปชฺชติ เหตุปจฺจยา.  เหตุยา ฉ.}

\csromanmarkh{1. กุสลตฺติก}\csromanmarkhii{14. มิจฺฉตฺตนิยต\-ตฺติก}\csromanheadernomark{2}{1. กุสลตฺติก 14. มิจฺฉตฺตนิยต\-ตฺติก}
\proseitem{18}{นกุสลํ นมิจฺฉตฺต\-นิยตํ ธมฺมํ ปจฺจยา อกุสโล มิจฺฉตฺต\-นิยโต ธมฺโม อุปฺปชฺชติ เหตุปจฺจยา.  เหตุยา ตีณิ.}

\prose{นกุสลํ นสมฺมตฺต\-นิยตํ ธมฺมํ ปจฺจยา กุสโล สมฺมตฺตนิยโต ธมฺโม อุปฺปชฺชติ เหตุปจฺจยา.  เหตุยา ตีณิ.}

\prose{นกุสลํ นอนิยตํ ธมฺมํ ปฏิจฺจ อพฺยากโต อนิยโต ธมฺโม อุปฺปชฺชติ เหตุปจฺจยา.  เหตุยา ปญฺจ.}

\csromanmarkh{1. กุสลตฺติก}\csromanmarkhii{15. มคฺคา\-รมฺมณ\-ตฺติก}\csromanheadernomark{2}{1. กุสลตฺติก 15. มคฺคา\-รมฺมณ\-ตฺติก}
\proseitem{19}{นกุสลํ นมคฺคา\-รมฺมณํ ธมฺมํ ปจฺจยา กุสโล มคฺคารมฺมโณ ธมฺโม อุปฺปชฺชติ เหตุปจฺจยา.  เหตุยา ฉ.}

\prose{นกุสลํ นมคฺคเหตุกํ ธมฺมํ ปจฺจยา กุสโล มคฺคเหตุโก ธมฺโม อุปฺปชฺชติ เหตุปจฺจยา.  เหตุยา ตีณิ.}

\prose{นกุสลํ นมคฺคา\-ธิปติํ ธมฺมํ ปจฺจยา กุสโล มคฺคาธิปติ ธมฺโม อุปฺปชฺชติ เหตุปจฺจยา.  เหตุยา ฉ.}

\csromanmarkh{1. กุสลตฺติก}\csromanmarkhii{16. อุปฺปนฺนตฺติก}\csromanheadernomark{2}{1. \textbf{กุสลตฺติก} 16. อุปฺปนฺนตฺติก}
\proseitem{20}{นกุสโล นอุปฺปนฺโน ธมฺโม กุสลสฺส อุปฺปนฺนสฺส ธมฺมสฺส อารมฺมณ\-ปจฺจเยน ปจฺจโย.  อารมฺมเณ อฏฺฐารส.}

\csromanheader{2}{1. กุสลตฺติก 17-18. อตีต\-ตฺติก\-ทฺวย}
\proseitem{21}{นกุสโล นปจฺจุปฺปนฺโน ธมฺโม กุสลสฺส ปจฺจุปฺปนฺนสฺส ธมฺมสฺส อารมฺมณ\-ปจฺจเยน ปจฺจโย.  อารมฺมเณ อฏฺฐารส.}

\proseitem{22}{นกุสลํ นอตีตา\-รมฺมณํ ธมฺมํ ปฏิจฺจ อพฺยากโต อตีตารมฺมโณ ธมฺโม อุปฺปชฺชติ เหตุปจฺจยา.  เหตุยา ตีณิ.}

\csromanheader{2}{ธมฺม\-ปจฺจนียา\-นุโลม ติกติก\-ปฏฺฐาน\-ปาฬิ}
\prose{\csromanfolio{410}นกุสลํ นอนาคตา\-รมฺมณํ ธมฺมํ ปจฺจยา กุสโล อนาคตารมฺมโณ ธมฺโม อุปฺปชฺชติ เหตุปจฺจยา.  เหตุยา นว.}

\prose{นกุสลํ นปจฺจุปฺปนฺนา\-รมฺมณํ ธมฺมํ ปฏิจฺจ อพฺยากโต ปจฺจุปฺปนฺนา\-รมฺมโณ ธมฺโม อุปฺปชฺชติ เหตุปจฺจยา.  เหตุยา ตีณิ.}

\csromanheader{2}{1. กุสลตฺติก 19-20. อชฺฌตฺต\-ตฺติก\-ทฺวย}
\proseitem{23}{นกุสโล นอชฺฌตฺโต ธมฺโม อชฺฌตฺตสฺส กุสลสฺส ธมฺมสฺส อารมฺมณ\-ปจฺจเยน ปจฺจโย. (สํขิตฺตํ.) . อารมฺมเณ อฏฺฐารส, อธิปติยา โสฬส, อุปนิสฺสเย อฏฺฐารส, ปุเรชาเต อตฺถิยา อวิคเต นว.}

\prose{นกุสโล นพหิทฺธา ธมฺโม พหิทฺธา กุสลสฺส ธมฺมสฺส อารมฺมณ\-ปจฺจเยน ปจฺจโย. (สํขิตฺตํ.) . อารมฺมเณ อฏฺฐารส, อธิปติยา ฉ, อุปนิสฺสเย อฏฺฐารส, ปุเรชาเต อตฺถิยา อวิคเต นว.}

\proseitem{24}{นกุสลํ นอชฺฌตฺตา\-รมฺมณํ ธมฺมํ ปฏิจฺจ อพฺยากโต อชฺฌตฺตา\-รมฺมโณ ธมฺโม อุปฺปชฺชติ เหตุปจฺจยา.  เหตุยา ตีณิ.}

\prose{นกุสลํ นพหิทฺธา\-รมฺมณํ ธมฺมํ ปฏิจฺจ อพฺยากโต พหิทฺธา\-รมฺมโณ ธมฺโม อุปฺปชฺชติ เหตุปจฺจยา.  เหตุยา ตีณิ.}

\csromanmarkh{1. กุสลตฺติก}\csromanmarkhii{21. สนิทสฺสน\-ตฺติก}\csromanheadernomark{2}{1. กุสลตฺติก 21. สนิทสฺสน\-ตฺติก}
\proseitem{25}{นกุสลํ นสนิทสฺสน\-สปฺปฏิฆํ ธมฺมํ ปฏิจฺจ อพฺยากโต สนิทสฺสน\-สปฺปฏิโฆ ธมฺโม อุปฺปชฺชติ เหตุปจฺจยา. เอกํ.}

\prose{นอกุสลํ นสนิทสฺสน\-สปฺปฏิฆํ ธมฺมํ ปฏิจฺจ อพฺยากโต สนิทสฺสน\-สปฺปฏิโฆ ธมฺโม อุปฺปชฺชติ เหตุปจฺจยา. เอกํ.}

\prose{นอพฺยากตํ นสนิทสฺสน\-สปฺปฏิฆํ ธมฺมํ ปฏิจฺจ อพฺยากโต สนิทสฺสน\-สปฺปฏิโฆ ธมฺโม อุปฺปชฺชติ เหตุปจฺจยา. เอกํ.}

\prose{นกุสลํ นสนิทสฺสนสปฺปฏิฆญฺจ นอพฺยากตํ นสนิทสฺสนสปฺปฏิฆญฺจ ธมฺมํ ปฏิจฺจ อพฺยากโต สนิทสฺสน\-สปฺปฏิโฆ ธมฺโม อุปฺปชฺชติ เหตุปจฺจยา. เอกํ.}

\csromanmarkh{2. เวทนาตฺติก}\csromanmarkhii{1. กุสลตฺติก}\csromanheadernomark{2}{2. เวทนาตฺติก 1. กุสลตฺติก}
\prose{\csromanfolio{411}นอกุสลํ นสนิทสฺสนสปฺปฏิฆญฺจ นอพฺยากตํ นสนิทสฺสนสปฺปฏิฆญฺจ ธมฺมํ ปฏิจฺจ อพฺยากโต สนิทสฺสน\-สปฺปฏิโฆ ธมฺโม อุปฺปชฺชติ เหตุปจฺจยา. เอกํ.}

\prose{นกุสลํ นสนิทสฺสนสปฺปฏิฆญฺจ นอกุสลํ นสนิทสฺสนสปฺปฏิฆญฺจ ธมฺมํ ปฏิจฺจ อพฺยากโต สนิทสฺสน\-สปฺปฏิโฆ ธมฺโม อุปฺปชฺชติ เหตุปจฺจยา. เอกํ.  เหตุยา ฉ.}

\proseitem{26}{นกุสลํ นอนิทสฺสน\-สปฺปฏิฆํ ธมฺมํ ปฏิจฺจ อพฺยากโต อนิทสฺสน\-สปฺปฏิโฆ ธมฺโม อุปฺปชฺชติ เหตุปจฺจยา.  เหตุยา ฉ.}

\prose{นกุสลํ นอนิทสฺสน\-อปฺปฏิฆํ ธมฺมํ ปฏิจฺจ อพฺยากโต อนิทสฺสน\-สปฺปฏิโฆ ธมฺโม อุปฺปชฺชติ เหตุปจฺจยา.  เหตุยา ตีณิ.}

\csromanmarkh{2. เวทนาตฺติก}\csromanmarkhii{1. กุสลตฺติก}\csromanheadernomark{2}{2. เวทนาตฺติก 1. กุสลตฺติก}
\proseitem{27}{นสุขาย เวทนาย สมฺปยุตฺตํ นกุสลํ ธมฺมํ ปจฺจยา สุขาย เวทนาย สมฺปยุตฺโต กุสโล ธมฺโม อุปฺปชฺชติ เหตุปจฺจยา.  นสุขาย เวทนาย สมฺปยุตฺตํ นกุสลํ ธมฺมํ ปจฺจยา อทุกฺขม\-สุขาย เวทนาย สมฺปยุตฺโต กุสโล ธมฺโม อุปฺปชฺชติ เหตุปจฺจยา. เทฺว.}

\prose{นทุกฺขาย เวทนาย สมฺปยุตฺตํ นกุสลํ ธมฺมํ ปจฺจยา สุขาย เวทนาย สมฺปยุตฺโต กุสโล ธมฺโม อุปฺปชฺชติ เหตุปจฺจยา.  นทุกฺขาย เวทนาย สมฺปยุตฺตํ นกุสลํ ธมฺมํ ปจฺจยา อทุกฺขม\-สุขาย เวทนาย สมฺปยุตฺโต กุสโล ธมฺโม อุปฺปชฺชติ เหตุปจฺจยา. เทฺว.}

\prose{นอทุกฺขม\-สุขาย เวทนาย สมฺปยุตฺตํ นกุสลํ ธมฺมํ ปจฺจยา อทุกฺขม\-สุขาย เวทนาย สมฺปยุตฺโต กุสโล ธมฺโม อุปฺปชฺชติ เหตุปจฺจยา.  นอทุกฺขม\-สุขาย เวทนาย สมฺปยุตฺตํ นกุสลํ ธมฺมํ ปจฺจยา สุขาย เวทนาย สมฺปยุตฺโต กุสโล ธมฺโม อุปฺปชฺชติ เหตุปจฺจยา. เทฺว. (จตฺตาริ คณิตเกน เทฺว เทฺว ปญฺหา กาตพฺพา.) . เหตุยา จุทฺทส.}

\proseitem{28}{นสุขาย เวทนาย สมฺปยุตฺตํ นอกุสลํ ธมฺมํ ปจฺจยา สุขาย เวทนาย สมฺปยุตฺโต อกุสโล ธมฺโม อุปฺปชฺชติ เหตุปจฺจยา.  นสุขาย เวทนาย สมฺปยุตฺตํ นอกุสลํ ธมฺมํ ปจฺจยา ทุกฺขาย เวทนาย สมฺปยุตฺโต อกุสโล ธมฺโม อุปฺปชฺชติ เหตุปจฺจยา.  เหตุยา เอกวีส.}

\csromanheader{2}{ธมฺม\-ปจฺจนียา\-นุโลม ติกติก\-ปฏฺฐาน\-ปาฬิ}
\proseitem{29}{\csromanfolio{412}นสุขาย เวทนาย สมฺปยุตฺโต นอพฺยากโต ธมฺโม สุขาย เวทนาย สมฺปยุตฺตสฺส อพฺยากตสฺส ธมฺมสฺส อารมฺมณ\-ปจฺจเยน ปจฺจโย. อารมฺมเณ จุทฺทส.}

\csromanmarkh{3. วิปากตฺติก}\csromanmarkhii{1. กุสลตฺติก}\csromanheadernomark{2}{3. \textbf{วิปากตฺติก 1. กุสลตฺติก}}
\proseitem{30}{นวิปากํ นกุสลํ ธมฺมํ ปจฺจยา วิปาก\-ธมฺม\-ธมฺโม กุสโล ธมฺโม อุปฺปชฺชติ เหตุปจฺจยา.  เหตุยา ตีณิ.}

\prose{นวิปากํ นอกุสลํ ธมฺมํ ปจฺจยา วิปาก\-ธมฺม\-ธมฺโม อกุสโล ธมฺโม อุปฺปชฺชติ เหตุปจฺจยา.  เหตุยา ตีณิ.}

\prose{นวิปากํ นอพฺยากตํ ธมฺมํ ปฏิจฺจ เนว\-วิปาก\-น\-วิปาก\-ธมฺม\-ธมฺโม อพฺยากโต ธมฺโม อุปฺปชฺชติ เหตุปจฺจยา.  เหตุยา ตีณิ.}

\csromanmarkh{4. อุปาทินฺนตฺติก}\csromanmarkhii{1. กุสลตฺติก}\csromanheadernomark{2}{4. \textbf{อุปาทินฺนตฺติก 1. กุสลตฺติก}}
\proseitem{31}{นอนุปาทินฺนุ\-ปาทานิยํ นกุสลํ ธมฺมํ ปจฺจยา อนุปาทินฺนุ\-ปาทานิโย กุสโล ธมฺโม อุปฺปชฺชติ เหตุปจฺจยา.  เหตุยา ฉ.}

\prose{นอนุปาทินฺนุ\-ปาทานิยํ นอกุสลํ ธมฺมํ ปจฺจยา อนุปาทินฺนุ\-ปาทานิโย อกุสโล ธมฺโม อุปฺปชฺชติ เหตุปจฺจยา.  เหตุยา ตีณิ.}

\prose{นอุปาทินฺนุปาทานิยํ นอพฺยากตํ ธมฺมํ ปฏิจฺจ อนุปาทินฺนุ\-ปาทานิโย อพฺยากโต ธมฺโม อุปฺปชฺชติ เหตุปจฺจยา.  เหตุยา ปญฺจ.}

\csromanmarkh{5. สํกิลิฏฺฐ\-ตฺติก}\csromanmarkhii{1. กุสลตฺติก}\csromanheadernomark{2}{5. สํกิลิฏฺฐ\-ตฺติก 1. กุสลตฺติก}
\proseitem{32}{นสํกิลิฏฺฐ\-สํกิเลสิกํ นกุสลํ ธมฺมํ ปจฺจยา อสํกิลิฏฺฐ\-สํกิเลสิโก กุสโล ธมฺโม อุปฺปชฺชติ เหตุปจฺจยา.  เหตุยา ฉ.}

\prose{นสํกิลิฏฺฐ\-สํกิเลสิกํ นอกุสลํ ธมฺมํ ปจฺจยา สํกิลิฏฺฐ\-สํกิเลสิโก อกุสโล ธมฺโม อุปฺปชฺชติ เหตุปจฺจยา.  เหตุยา ตีณิ.}

\prose{นสํกิลิฏฺฐ\-สํกิเลสิกํ นอพฺยากตํ ธมฺมํ ปฏิจฺจ อสํกิลิฏฺฐ\-สํกิเลสิโก อพฺยากโต ธมฺโม อุปฺปชฺชติ เหตุปจฺจยา.  เหตุยา ฉ.}

\csromanmarkh{8. ทสฺสนตฺติก}\csromanmarkhii{1. กุสลตฺติก 6. วิตกฺกตฺติก 1. กุสลตฺติก}\csromanheadernomark{2}{8. ทสฺสนตฺติก 1. กุสลตฺติก \textbf{6. วิตกฺกตฺติก 1. กุสลตฺติก}}
\proseitem{33}{\csromanfolio{413}นสวิตกฺก\-สวิจารํ นกุสลํ ธมฺมํ ปจฺจยา สวิตกฺก\-สวิจาโร กุสโล ธมฺโม อุปฺปชฺชติ เหตุปจฺจยา.  เหตุยา ปนฺนรส.}

\prose{นสวิตกฺก\-สวิจารํ นอกุสลํ ธมฺมํ ปจฺจยา สวิตกฺก\-สวิจาโร อกุสโล ธมฺโม อุปฺปชฺชติ เหตุปจฺจยา.  เหตุยา นว.}

\prose{นสวิตกฺก\-สวิจารํ นอพฺยากตํ ธมฺมํ ปฏิจฺจ อวิตกฺก\-อวิจาโร อพฺยากโต ธมฺโม อุปฺปชฺชติ เหตุปจฺจยา.  เหตุยา สตฺต.}

\csromanmarkh{7. ปีติตฺติก}\csromanmarkhii{1. กุสลตฺติก}\csromanheadernomark{2}{7. \textbf{ปีติตฺติก 1. กุสลตฺติก}}
\proseitem{34}{นปีติสหคตํ นกุสลํ ธมฺมํ ปจฺจยา ปีติสหคโต กุสโล ธมฺโม อุปฺปชฺชติ เหตุปจฺจยา.  เหตุยา อฏฺฐวีส.}

\prose{นปีติสหคตํ นอกุสลํ ธมฺมํ ปจฺจยา ปีติสหคโต อกุสโล ธมฺโม อุปฺปชฺชติ เหตุปจฺจยา.  เหตุยา อฏฺฐวีส.}

\prose{นปีติสหคโต นอพฺยากโต ธมฺโม ปีติสหคตสฺส อพฺยากตสฺส ธมฺมสฺส อารมฺมณ\-ปจฺจเยน ปจฺจโย. (สํขิตฺตํ.) . อารมฺมเณ อฏฺฐวีส, อธิปติยา อนนฺตเร อฏฺฐวีส ฯเปฯ อุปนิสฺสเย อฏฺฐวีส, กมฺเม จตุวีส, นตฺถิยา วิคเต อฏฺฐวีส.}

\csromanmarkh{8. ทสฺสนตฺติก}\csromanmarkhii{1. กุสลตฺติก}\csromanheadernomark{2}{8. ทสฺสนตฺติก 1. กุสลตฺติก}
\proseitem{35}{นทสฺสเนน ปหาตพฺพํ นกุสลํ ธมฺมํ ปจฺจยา เนวทสฺสเนน นภาวนาย ปหาตพฺโพ กุสโล ธมฺโม อุปฺปชฺชติ เหตุปจฺจยา.  เหตุยา ตีณิ.}

\prose{นทสฺสเนน ปหาตพฺพํ นอกุสลํ ธมฺมํ ปจฺจยา ทสฺสเนน ปหาตพฺโพ อกุสโล ธมฺโม อุปฺปชฺชติ เหตุปจฺจยา.  เหตุยา ฉ.}

\prose{นทสฺสเนน ปหาตพฺพํ นอพฺยากตํ ธมฺมํ ปฏิจฺจ เนวทสฺสเนน นภาวนาย ปหาตพฺโพ อพฺยากโต ธมฺโม อุปฺปชฺชติ เหตุปจฺจยา.  เหตุยา ฉ.}

\csromanmarkh{ธมฺม\-ปจฺจนียา\-นุโลม ติกติก\-ปฏฺฐาน\-ปาฬิ}\csromanmarkhii{9. ทสฺสน\-เหตุ\-ตฺติก 1. กุสลตฺติก}\csromanheadernomark{2}{ธมฺม\-ปจฺจนียา\-นุโลม ติกติก\-ปฏฺฐาน\-ปาฬิ 9. ทสฺสน\-เหตุ\-ตฺติก 1. กุสลตฺติก}
\proseitem{36}{\csromanfolio{414}นทสฺสเนน ปหาตพฺพ\-เหตุกํ นกุสลํ ธมฺมํ ปจฺจยา เนวทสฺสเนน นภาวนาย ปหาตพฺพ\-เหตุโก กุสโล ธมฺโม อุปฺปชฺชติ เหตุปจฺจยา.  เหตุยา ตีณิ.}

\prose{นทสฺสเนน ปหาตพฺพ\-เหตุกํ นอกุสลํ ธมฺมํ ปจฺจยา ทสฺสเนน ปหาตพฺพ\-เหตุโก อกุสโล ธมฺโม อุปฺปชฺชติ เหตุปจฺจยา.  เหตุยา ฉ.}

\prose{นเน\-ว\-ทสฺสเนน นภาวนาย ปหาตพฺพ\-เหตุกํ นอพฺยากตํ ธมฺมํ ปฏิจฺจ เนวทสฺสเนน นภาวนาย ปหาตพฺพ\-เหตุโก อพฺยากโต ธมฺโม อุปฺปชฺชติ เหตุปจฺจยา.  เหตุยา ฉ.}

\csromanmarkh{10. อาจยคามิตฺติก}\csromanmarkhii{1. กุสลตฺติก}\csromanheadernomark{2}{10. \textbf{อาจยคามิตฺติก 1. กุสลตฺติก}}
\proseitem{37}{นอาจยคามิํ นกุสลํ ธมฺมํ ปจฺจยา อาจยคามี กุสโล ธมฺโม อุปฺปชฺชติ เหตุปจฺจยา.  เหตุยา ฉ.}

\prose{นอาจยคามิํ นอกุสลํ ธมฺมํ ปจฺจยา อาจยคามี อกุสโล ธมฺโม อุปฺปชฺชติ เหตุปจฺจยา.  เหตุยา ตีณิ.}

\prose{นอาจยคามิํ นอพฺยากตํ ธมฺมํ ปฏิจฺจ เนวา\-จยคามิ\-นา\-ปจยคามี อพฺยากโต ธมฺโม อุปฺปชฺชติ เหตุปจฺจยา.  เหตุยา ปญฺจ.}

\csromanmarkh{11. เสกฺขตฺติก}\csromanmarkhii{1. กุสลตฺติก}\csromanheadernomark{2}{11. \textbf{เสกฺขตฺติก 1. กุสลตฺติก}}
\proseitem{38}{นเสกฺขํ นกุสลํ ธมฺมํ ปจฺจยา เสกฺโข กุสโล ธมฺโม อุปฺปชฺชติ เหตุปจฺจยา.  เหตุยา ฉ.}

\prose{นเสกฺขํ นอกุสลํ ธมฺมํ ปจฺจยา เนว\-เสกฺข\-นา\-เสกฺโข อกุสโล ธมฺโม อุปฺปชฺชติ เหตุปจฺจยา.  เหตุยา ตีณิ.}

\prose{นเสกฺขํ นอพฺยากตํ ธมฺมํ ปฏิจฺจ เนว\-เสกฺข\-นา\-เสกฺโข อพฺยากโต ธมฺโม อุปฺปชฺชติ เหตุปจฺจยา.  เหตุยา ปญฺจ.}

\csromanmarkh{12. ปริตฺตตฺติก}\csromanmarkhii{1. กุสลตฺติก}\csromanheadernomark{2}{12. \textbf{ปริตฺตตฺติก 1. กุสลตฺติก}}
\proseitem{39}{นมหคฺคตํ นกุสลํ ธมฺมํ ปจฺจยา มหคฺคโต กุสโล ธมฺโม อุปฺปชฺชติ เหตุปจฺจยา.  เหตุยา นว.}

\csromanmarkh{15. มิจฺฉตฺตนิยต\-ตฺติก}\csromanmarkhii{1. กุสลตฺติก}\csromanheadernomark{2}{15. มิจฺฉตฺตนิยต\-ตฺติก 1. กุสลตฺติก}
\prose{\csromanfolio{415}นมหคฺคตํ นอกุสลํ ธมฺมํ ปจฺจยา ปริตฺโต อกุสโล ธมฺโม อุปฺปชฺชติ เหตุปจฺจยา.  เหตุยา ตีณิ.}

\prose{นปริตฺตํ นอพฺยากตํ ธมฺมํ ปฏิจฺจ ปริตฺโต อพฺยากโต ธมฺโม อุปฺปชฺชติ เหตุปจฺจยา.  เหตุยา ฉ.}

\csromanmarkh{13. ปริตฺตา\-รมฺมณ\-ตฺติก}\csromanmarkhii{1. กุสลตฺติก}\csromanheadernomark{2}{13. ปริตฺตา\-รมฺมณ\-ตฺติก 1. กุสลตฺติก}
\proseitem{40}{นปริตฺตา\-รมฺมณํ นกุสลํ ธมฺมํ ปจฺจยา ปริตฺตารมฺมโณ กุสโล ธมฺโม อุปฺปชฺชติ เหตุปจฺจยา.  เหตุยา เอกวีส.}

\prose{นปริตฺตา\-รมฺมณํ นอกุสลํ ธมฺมํ ปจฺจยา ปริตฺตารมฺมโณ อกุสโล ธมฺโม อุปฺปชฺชติ เหตุปจฺจยา.  เหตุยา จุทฺทส.}

\prose{นปริตฺตา\-รมฺมโณ นอพฺยากโต ธมฺโม ปริตฺตา\-รมฺมณสฺส อพฺยากตสฺส ธมฺมสฺส อารมฺมณ\-ปจฺจเยน ปจฺจโย. (สํขิตฺตํ.) . อารมฺมเณ เอกวีส.}

\csromanmarkh{14. หีนตฺติก}\csromanmarkhii{1. กุสลตฺติก}\csromanheadernomark{2}{14. \textbf{หีนตฺติก 1. กุสลตฺติก}}
\proseitem{41}{นหีนํ นกุสลํ ธมฺมํ ปจฺจยา มชฺฌิโม กุสโล ธมฺโม อุปฺปชฺชติ เหตุปจฺจยา.  เหตุยา ฉ.}

\prose{นหีนํ นอกุสลํ ธมฺมํ ปจฺจยา หีโน อกุสโล ธมฺโม อุปฺปชฺชติ เหตุปจฺจยา.  เหตุยา ตีณิ.}

\prose{นหีนํ นอพฺยากตํ ธมฺมํ ปฏิจฺจ มชฺฌิโม อพฺยากโต ธมฺโม อุปฺปชฺชติ เหตุปจฺจยา.  เหตุยา ฉ.}

\csromanmarkh{15. มิจฺฉตฺตนิยต\-ตฺติก}\csromanmarkhii{1. กุสลตฺติก}\csromanheadernomark{2}{15. มิจฺฉตฺตนิยต\-ตฺติก 1. กุสลตฺติก}
\proseitem{42}{นมิจฺฉตฺต\-นิยตํ นกุสลํ ธมฺมํ ปจฺจยา สมฺมตฺตนิยโต กุสโล ธมฺโม อุปฺปชฺชติ เหตุปจฺจยา.  เหตุยา ฉ.}

\prose{นมิจฺฉตฺต\-นิยตํ นอกุสลํ ธมฺมํ ปจฺจยา มิจฺฉตฺต\-นิยโต อกุสโล ธมฺโม อุปฺปชฺชติ เหตุปจฺจยา.  เหตุยา ฉ.}

\prose{นมิจฺฉตฺต\-นิยตํ นอพฺยากตํ ธมฺมํ ปฏิจฺจ อนิยโต อพฺยากโต ธมฺโม อุปฺปชฺชติ เหตุปจฺจยา.  เหตุยา ฉ.}

\csromanmarkh{ธมฺม\-ปจฺจนียา\-นุโลม ติกติก\-ปฏฺฐาน\-ปาฬิ}\csromanmarkhii{16. มคฺคา\-รมฺมณ\-ตฺติก 1. กุสลตฺติก}\csromanheadernomark{2}{ธมฺม\-ปจฺจนียา\-นุโลม ติกติก\-ปฏฺฐาน\-ปาฬิ 16. มคฺคา\-รมฺมณ\-ตฺติก 1. กุสลตฺติก}
\proseitem{43}{\csromanfolio{416}นมคฺคา\-รมฺมณํ นกุสลํ ธมฺมํ ปจฺจยา มคฺคารมฺมโณ กุสโล ธมฺโม อุปฺปชฺชติ เหตุปจฺจยา.  เหตุยา ปญฺจติํส. (นอกุสลํ นอพฺยากตํ นตฺถิ.)}

\csromanmarkh{17. อุปฺปนฺนตฺติก}\csromanmarkhii{1. กุสลตฺติก}\csromanheadernomark{2}{17. \textbf{อุปฺปนฺนตฺติก 1. กุสลตฺติก}}
\proseitem{44}{นอุปฺปนฺโน นกุสโล ธมฺโม อุปฺปนฺนสฺส กุสลสฺส ธมฺมสฺส อารมฺมณ\-ปจฺจเยน ปจฺจโย.  อารมฺมเณ สตฺต.}

\prose{นอุปฺปนฺโน นอกุสโล ธมฺโม อุปฺปนฺนสฺส อกุสลสฺส ธมฺมสฺส อารมฺมณ\-ปจฺจเยน ปจฺจโย.  อารมฺมเณ ฉ.}

\prose{นอุปฺปนฺโน นอพฺยากโต ธมฺโม อุปฺปนฺนสฺส อพฺยากตสฺส ธมฺมสฺส อารมฺมณ\-ปจฺจเยน ปจฺจโย.  อารมฺมเณ สตฺต. (อตีตตฺติกํ อุปฺปนฺน\-ตฺติก\-สทิสํ.)}

\csromanmarkh{19. อตีตา\-รมฺมณ\-ตฺติก}\csromanmarkhii{1. กุสลตฺติก}\csromanheadernomark{2}{19. อตีตา\-รมฺมณ\-ตฺติก 1. กุสลตฺติก}
\proseitem{45}{นอตีตา\-รมฺมณํ นกุสลํ ธมฺมํ ปจฺจยา อตีตารมฺมโณ กุสโล ธมฺโม อุปฺปชฺชติ เหตุปจฺจยา.  เหตุยา เอกวีส.}

\prose{นอตีตา\-รมฺมณํ นอกุสลํ ธมฺมํ ปจฺจยา อตีตารมฺมโณ อกุสโล ธมฺโม อุปฺปชฺชติ เหตุปจฺจยา.  เหตุยา เอกวีส.}

\csromanmarkh{20. อชฺฌตฺตตฺติก}\csromanmarkhii{1. กุสลตฺติก}\csromanheadernomark{2}{20. \textbf{อชฺฌตฺตตฺติก 1. กุสลตฺติก}}
\proseitem{46}{นอชฺฌตฺตํ นกุสลํ ธมฺมํ ปจฺจยา พหิทฺธา กุสโล ธมฺโม อุปฺปชฺชติ เหตุปจฺจยา.  นพหิทฺธา นกุสลํ ธมฺมํ ปจฺจยา อชฺฌตฺโต กุสโล ธมฺโม อุปฺปชฺชติ เหตุปจฺจยา.  เหตุยา เทฺว.}

\prose{นอชฺฌตฺตํ นอกุสลํ ธมฺมํ ปจฺจยา พหิทฺธา อกุสโล ธมฺโม อุปฺปชฺชติ เหตุปจฺจยา. นพหิทฺธา นอกุสลํ ธมฺมํ ปจฺจยา อชฺฌตฺโต อกุสโล ธมฺโม อุปฺปชฺชติ เหตุปจฺจยา.  เหตุยา เทฺว.}

\vspace{\tipitakaparskip}
\csromancenter{สนิทสฺสน\-ตฺติก 2. เวทนาตฺติก 21. อชฺฌตฺตา\-รมฺมณ\-ตฺติก 1. กุสลตฺติก}
\proseitem{47}{\csromanfolio{417}นอชฺฌตฺตา\-รมฺมณํ นกุสลํ ธมฺมํ ปจฺจยา อชฺฌตฺตา\-รมฺมโณ กุสโล ธมฺโม อุปฺปชฺชติ เหตุปจฺจยา.  เหตุยา ฉ.}

\prose{นอชฺฌตฺตา\-รมฺมณํ นอกุสลํ ธมฺมํ ปจฺจยา อชฺฌตฺตา\-รมฺมโณ อกุสโล ธมฺโม อุปฺปชฺชติ เหตุปจฺจยา.  เหตุยา ฉ.}

\csromanmarkh{22. สนิทสฺสน\-ตฺติก}\csromanmarkhii{1. กุสลตฺติก}\csromanheadernomark{2}{22. สนิทสฺสน\-ตฺติก 1. กุสลตฺติก}
\proseitem{48}{นสนิทสฺสน\-สปฺปฏิฆํ นกุสลํ ธมฺมํ ปจฺจยา อนิทสฺสน\-อปฺปฏิโฆ กุสโล ธมฺโม อุปฺปชฺชติ เหตุปจฺจยา.  นอนิทสฺสน\-สปฺปฏิฆํ นกุสลํ ธมฺมํ ปจฺจยา อนิทสฺสน\-อปฺปฏิโฆ กุสโล ธมฺโม อุปฺปชฺชติ เหตุปจฺจยา.  นสนิทสฺสน\-สปฺปฏิฆํ นกุสลญฺจ นอนิทสฺสน\-สปฺปฏิฆํ นกุสลญฺจ ธมฺมํ ปจฺจยา อนิทสฺสน\-อปฺปฏิโฆ กุสโล ธมฺโม อุปฺปชฺชติ เหตุปจฺจยา.  เหตุยา ตีณิ.}

\prose{นสนิทสฺสน\-สปฺปฏิฆํ นอกุสลํ ธมฺมํ ปจฺจยา อนิทสฺสน\-อปฺปฏิโฆ อกุสโล ธมฺโม อุปฺปชฺชติ เหตุปจฺจยา.  เหตุยา ตีณิ.}

\prose{นสนิทสฺสน\-สปฺปฏิฆํ นอพฺยากตํ ธมฺมํ ปฏิจฺจ สนิทสฺสน\-สปฺปฏิโฆ อพฺยากโต ธมฺโม อุปฺปชฺชติ เหตุปจฺจยา. สตฺต.}

\prose{(นอนิทสฺสนสปฺปฏิฆนอพฺยากตมูลานิ สตฺตเมว. ทุกมูลานิ สตฺตเมว. สพฺพํ เอกวีสติเมว.)}

\csromanmarkh{22. สนิทสฺสน\-ตฺติก}\csromanmarkhii{2. เวทนาตฺติก}\csromanheadernomark{2}{22. สนิทสฺสน\-ตฺติก 2. เวทนาตฺติก}
\proseitem{49}{นสนิทสฺสน\-สปฺปฏิฆํ นสุขาย เวทนาย สมฺปยุตฺตํ ธมฺมํ ปฏิจฺจ อนิทสฺสน\-อปฺปฏิโฆ สุขาย เวทนาย สมฺปยุตฺโต ธมฺโม อุปฺปชฺชติ เหตุปจฺจยา.  เหตุยา ตีณิ.}

\prose{นสนิทสฺสน\-สปฺปฏิฆํ นทุกฺขาย เวทนาย สมฺปยุตฺตํ ธมฺมํ ปฏิจฺจ อนิทสฺสน\-อปฺปฏิโฆ ทุกฺขาย เวทนาย สมฺปยุตฺโต ธมฺโม อุปฺปชฺชติ เหตุปจฺจยา.  เหตุยา ตีณิ.}

\prose{นสนิทสฺสน\-สปฺปฏิฆํ นอทุกฺขม\-สุขาย เวทนาย สมฺปยุตฺตํ ธมฺมํ ปฏิจฺจ อนิทสฺสน\-อปฺปฏิโฆ อทุกฺขม\-สุขาย เวทนาย สมฺปยุตฺโต ธมฺโม อุปฺปชฺชติ เหตุปจฺจยา.  เหตุยา ตีณิ.}

\csromanmarkh{ธมฺม\-ปจฺจนียา\-นุโลม ติกติก\-ปฏฺฐาน\-ปาฬิ}\csromanmarkhii{22. สนิทสฺสน\-ตฺติก 3. วิปากตฺติก}\csromanheadernomark{2}{ธมฺม\-ปจฺจนียา\-นุโลม ติกติก\-ปฏฺฐาน\-ปาฬิ 22. สนิทสฺสน\-ตฺติก 3. วิปากตฺติก}
\proseitem{50}{\csromanfolio{418}นสนิทสฺสน\-สปฺปฏิฆํ นวิปากํ ธมฺมํ ปฏิจฺจ อนิทสฺสน\-อปฺปฏิโฆ วิปาโก ธมฺโม อุปฺปชฺชติ เหตุปจฺจยา.  เหตุยา ตีณิ.}

\prose{นสนิทสฺสน\-สปฺปฏิฆํ นวิปาก\-ธมฺม\-ธมฺมํ ปจฺจยา อนิทสฺสน\-อปฺปฏิโฆ วิปาก\-ธมฺม\-ธมฺโม อุปฺปชฺชติ เหตุปจฺจยา.  เหตุยา ตีณิ.}

\prose{นสนิทสฺสน\-สปฺปฏิฆํ นเน\-ว\-วิปาก\-น\-วิปากธมฺม\-ธมฺมํ ปฏิจฺจ สนิทสฺสน\-สปฺปฏิโฆ เนว\-วิปาก\-น\-วิปาก\-ธมฺม\-ธมฺโม อุปฺปชฺชติ เหตุปจฺจยา.  เหตุยา เอกวีส.}

\csromanmarkh{22. สนิทสฺสน\-ตฺติก}\csromanmarkhii{4. อุปาทินฺนตฺติก}\csromanheadernomark{2}{22. สนิทสฺสน\-ตฺติก 4. อุปาทินฺนตฺติก}
\proseitem{51}{นสนิทสฺสน\-สปฺปฏิโฆ นอุปาทินฺนุปาทานิโย ธมฺโม อนิทสฺสน\-อปฺปฏิฆสฺส อุปาทินฺนุ\-ปาทานิยสฺส ธมฺมสฺส อารมฺมณ\-ปจฺจเยน ปจฺจโย.  อารมฺมเณ ฉ.}

\prose{นสนิทสฺสน\-สปฺปฏิฆํ นอนุปาทินฺนุ\-ปาทานิยํ ธมฺมํ ปฏิจฺจ สนิทสฺสน\-สปฺปฏิโฆ อนุปาทินฺนุ\-ปาทานิโย ธมฺโม อุปฺปชฺชติ เหตุปจฺจยา.  เหตุยา เอกวีส.}

\prose{นสนิทสฺสน\-สปฺปฏิฆํ นอนุปาทินฺน\-อนุปาทานิยํ ธมฺมํ ปจฺจยา อนิทสฺสน\-อปฺปฏิโฆ อนุปาทินฺน\-อนุปาทานิโย ธมฺโม อุปฺปชฺชติ เหตุปจฺจยา.  เหตุยา ตีณิ.}

\csromanmarkh{22. สนิทสฺสน\-ตฺติก}\csromanmarkhii{5. สํกิลิฏฺฐ\-ตฺติก}\csromanheadernomark{2}{22. สนิทสฺสน\-ตฺติก 5. สํกิลิฏฺฐ\-ตฺติก}
\proseitem{52}{นสนิทสฺสน\-สปฺปฏิฆํ นสํกิลิฏฺฐ\-สํกิเลสิกํ ธมฺมํ ปจฺจยา อนิทสฺสน\-อปฺปฏิโฆ สํกิลิฏฺฐ\-สํกิเลสิโก ธมฺโม อุปฺปชฺชติ เหตุปจฺจยา.  เหตุยา ตีณิ.}

\prose{นสนิทสฺสน\-สปฺปฏิฆํ นอสํกิลิฏฺฐ\-สํกิเลสิกํ ธมฺมํ ปฏิจฺจ สนิทสฺสน\-สปฺปฏิโฆ อสํกิลิฏฺฐ\-สํกิเลสิโก ธมฺโม อุปฺปชฺชติ เหตุปจฺจยา.  เหตุยา เอกวีส.}

\prose{นสนิทสฺสน\-สปฺปฏิฆํ นอสํกิลิฏฺฐ\-อสํกิเลสิกํ ธมฺมํ ปจฺจยา อนิทสฺสน\-อปฺปฏิโฆ อสํกิลิฏฺฐ\-อสํกิเลสิโก ธมฺโม อุปฺปชฺชติ เหตุปจฺจยา.  เหตุยา ตีณิ.}

\vspace{\tipitakaparskip}
\csromancenter{สนิทสฺสน\-ตฺติก 8. ทสฺสนตฺติก 22. สนิทสฺสน\-ตฺติก 6. วิตกฺกตฺติก}
\proseitem{53}{\csromanfolio{419}นสนิทสฺสน\-สปฺปฏิฆํ นสวิตกฺก\-สวิจารํ ธมฺมํ ปฏิจฺจ อนิทสฺสน\-อปฺปฏิโฆ สวิตกฺก\-สวิจาโร ธมฺโม อุปฺปชฺชติ เหตุปจฺจยา.  เหตุยา ตีณิ.}

\prose{นสนิทสฺสน\-สปฺปฏิฆํ นอวิตกฺก\-วิจารมตฺตํ ธมฺมํ ปฏิจฺจ อนิทสฺสน\-อปฺปฏิโฆ อวิตกฺก\-วิจารมตฺโต ธมฺโม อุปฺปชฺชติ เหตุปจฺจยา.  เหตุยา ตีณิ.}

\prose{นสนิทสฺสน\-สปฺปฏิฆํ นอวิตกฺก\-อวิจารํ ธมฺมํ ปฏิจฺจ สนิทสฺสน\-สปฺปฏิโฆ อวิตกฺก\-อวิจาโร ธมฺโม อุปฺปชฺชติ เหตุปจฺจยา.  เหตุยา เอกวีส.}

\csromanmarkh{22. สนิทสฺสน\-ตฺติก}\csromanmarkhii{7. ปีติตฺติก}\csromanheadernomark{2}{22. สนิทสฺสน\-ตฺติก 7. ปีติตฺติก}
\proseitem{54}{นสนิทสฺสน\-สปฺปฏิฆํ นปีติสหคตํ ธมฺมํ ปฏิจฺจ อนิทสฺสน\-อปฺปฏิโฆ ปีติสหคโต ธมฺโม อุปฺปชฺชติ เหตุปจฺจยา.  เหตุยา ตีณิ.}

\prose{นสนิทสฺสน\-สปฺปฏิฆํ นสุขสหคตํ ธมฺมํ ปฏิจฺจ อนิทสฺสน\-อปฺปฏิโฆ สุขสหคโต ธมฺโม อุปฺปชฺชติ เหตุปจฺจยา.  เหตุยา ตีณิ.}

\prose{นสนิทสฺสน\-สปฺปฏิฆํ นอุเปกฺขาสหคตํ ธมฺมํ ปฏิจฺจ อนิทสฺสน\-อปฺปฏิโฆ อุเปกฺขา\-สหคโต ธมฺโม อุปฺปชฺชติ เหตุปจฺจยา.  เหตุยา ตีณิ.}

\csromanmarkh{22. สนิทสฺสน\-ตฺติก}\csromanmarkhii{8. ทสฺสนตฺติก}\csromanheadernomark{2}{22. สนิทสฺสน\-ตฺติก 8. ทสฺสนตฺติก}
\proseitem{55}{นสนิทสฺสน\-สปฺปฏิฆํ นทสฺสเนน ปหาตพฺพํ ธมฺมํ ปจฺจยา อนิทสฺสน\-อปฺปฏิโฆ ทสฺสเนน ปหาตพฺโพ ธมฺโม อุปฺปชฺชติ เหตุปจฺจยา.  เหตุยา ตีณิ.}

\prose{นสนิทสฺสน\-สปฺปฏิฆํ นภาวนาย ปหาตพฺพํ ธมฺมํ ปจฺจยา อนิทสฺสน\-อปฺปฏิโฆ ภาวนาย ปหาตพฺโพ ธมฺโม อุปฺปชฺชติ เหตุปจฺจยา.  เหตุยา ตีณิ.}

\prose{นสนิทสฺสน\-สปฺปฏิฆํ นเน\-ว\-ทสฺสเนน นภาวนาย ปหาตพฺพํ ธมฺมํ ปฏิจฺจ สนิทสฺสน\-สปฺปฏิโฆ เนวทสฺสเนน นภาวนาย ปหาตพฺโพ ธมฺโม อุปฺปชฺชติ เหตุปจฺจยา.  เหตุยา เอกวีส.}

\vspace{\tipitakaparskip}
\csromancenter{ธมฺม\-ปจฺจนียา\-นุโลม ติกติก\-ปฏฺฐาน\-ปาฬิ 22. สนิทสฺสน\-ตฺติก 9. ทสฺสน\-เหตุ\-ตฺติก}
\proseitem{56}{\csromanfolio{420}นสนิทสฺสน\-สปฺปฏิฆํ นทสฺสเนน ปหาตพฺพ\-เหตุกํ ธมฺมํ ปฏิจฺจ อนิทสฺสน\-อปฺปฏิโฆ ทสฺสเนน ปหาตพฺพ\-เหตุโก ธมฺโม อุปฺปชฺชติ เหตุปจฺจยา.  เหตุยา ตีณิ.}

\prose{นสนิทสฺสน\-สปฺปฏิฆํ นภาวนาย ปหาตพฺพ\-เหตุกํ ธมฺมํ ปฏิจฺจ อนิทสฺสน\-อปฺปฏิโฆ ภาวนาย ปหาตพฺพ\-เหตุโก ธมฺโม อุปฺปชฺชติ เหตุปจฺจยา.  เหตุยา ตีณิ.}

\prose{นสนิทสฺสน\-สปฺปฏิฆํ นเน\-ว\-ทสฺสเนน นภาวนาย ปหาตพฺพ\-เหตุกํ ธมฺมํ ปฏิจฺจ สนิทสฺสน\-สปฺปฏิโฆ เนวทสฺสเนน นภาวนาย ปหาตพฺพ\-เหตุโก ธมฺโม อุปฺปชฺชติ เหตุปจฺจยา.  เหตุยา เอกวีส. 22. สนิทสฺสน\-ตฺติก 10. อาจยคามิตฺติก}

\proseitem{57}{นสนิทสฺสน\-สปฺปฏิฆํ นอาจยคามิํ ธมฺมํ ปจฺจยา อนิทสฺสน\-อปฺปฏิโฆ อาจยคามี ธมฺโม อุปฺปชฺชติ เหตุปจฺจยา.  เหตุยา ตีณิ.}

\prose{นสนิทสฺสน\-สปฺปฏิฆํ นอปจยคามิํ ธมฺมํ ปจฺจยา อนิทสฺสน\-อปฺปฏิโฆ อปจยคามี ธมฺโม อุปฺปชฺชติ เหตุปจฺจยา.  เหตุยา ตีณิ.}

\prose{นสนิทสฺสน\-สปฺปฏิฆํ นเน\-วา\-จยคามิ\-นา\-ปจยคามิํ ธมฺมํ ปฏิจฺจ สนิทสฺสน\-สปฺปฏิโฆ เนวา\-จยคามิ\-นา\-ปจยคามี ธมฺโม อุปฺปชฺชติ เหตุปจฺจยา.  เหตุยา เอกวีส. 22. สนิทสฺสน\-ตฺติก 11. เสกฺขตฺติก}

\proseitem{58}{นสนิทสฺสน\-สปฺปฏิฆํ นเสกฺขํ ธมฺมํ ปจฺจยา อนิทสฺสน\-อปฺปฏิโฆ เสกฺโข ธมฺโม อุปฺปชฺชติ เหตุปจฺจยา.  เหตุยา ตีณิ.}

\prose{นสนิทสฺสน\-สปฺปฏิฆํ นอเสกฺขํ ธมฺมํ ปจฺจยา อนิทสฺสน\-อปฺปฏิโฆ อเสกฺโข ธมฺโม อุปฺปชฺชติ เหตุปจฺจยา.  เหตุยา ตีณิ.}

\prose{นสนิทสฺสน\-สปฺปฏิฆํ นเน\-ว\-เสกฺข\-นาเสกฺขํ ธมฺมํ ปฏิจฺจ สนิทสฺสน\-สปฺปฏิโฆ เนว\-เสกฺข\-นา\-เสกฺโข ธมฺโม อุปฺปชฺชติ เหตุปจฺจยา.  เหตุยา เอกวีส.}

\vspace{\tipitakaparskip}
\csromancenter{สนิทสฺสน\-ตฺติก 15. มิจฺฉตฺตนิยต\-ตฺติก 22. สนิทสฺสน\-ตฺติก 12. ปริตฺตตฺติก}
\proseitem{59}{\csromanfolio{421}นสนิทสฺสน\-สปฺปฏิฆํ นปริตฺตํ ธมฺมํ ปฏิจฺจ สนิทสฺสน\-สปฺปฏิโฆ ปริตฺโต ธมฺโม อุปฺปชฺชติ เหตุปจฺจยา.  เหตุยา เอกวีส.}

\prose{นสนิทสฺสน\-สปฺปฏิฆํ นมหคฺคตํ ธมฺมํ ปฏิจฺจ อนิทสฺสน\-อปฺปฏิโฆ มหคฺคโต ธมฺโม อุปฺปชฺชติ เหตุปจฺจยา.  เหตุยา ตีณิ.}

\prose{นสนิทสฺสน\-สปฺปฏิฆํ นอปฺปมาณํ ธมฺมํ ปจฺจยา อนิทสฺสน\-อปฺปฏิโฆ อปฺปมาโณ ธมฺโม อุปฺปชฺชติ เหตุปจฺจยา.  เหตุยา ตีณิ.}

\csromanmarkh{22. สนิทสฺสน\-ตฺติกฺอ}\csromanmarkhii{13. ปริตฺตา\-รมฺมณ\-ตฺติก}\csromanheadernomark{2}{22. สนิทสฺสน\-ตฺติกฺอ 13. ปริตฺตา\-รมฺมณ\-ตฺติก}
\proseitem{60}{นสนิทสฺสน\-สปฺปฏิฆํ นปริตฺตา\-รมฺมณํ ธมฺมํ ปฏิจฺจ อนิทสฺสน\-อปฺปฏิโฆ ปริตฺตารมฺมโณ ธมฺโม อุปฺปชฺชติ เหตุปจฺจยา.  เหตุยา ตีณิ.}

\prose{นสนิทสฺสน\-สปฺปฏิฆํ นมหคฺคตา\-รมฺมณํ ธมฺมํ ปจฺจยา อนิทสฺสน\-อปฺปฏิโฆ มหคฺคตา\-รมฺมโณ ธมฺโม อุปฺปชฺชติ เหตุปจฺจยา.  เหตุยา ตีณิ.}

\prose{นสนิทสฺสน\-สปฺปฏิฆํ นอปฺปมาณา\-รมฺมณํ ธมฺมํ ปจฺจยา อนิทสฺสน\-อปฺปฏิโฆ อปฺปมาณารมฺมโณ ธมฺโม อุปฺปชฺชติ เหตุปจฺจยา.  เหตุยา ตีณิ.}

\csromanmarkh{22. สนิทสฺสน\-ตฺติก}\csromanmarkhii{14. หีนตฺติก}\csromanheadernomark{2}{22. สนิทสฺสน\-ตฺติก 14. หีนตฺติก}
\proseitem{61}{นสนิทสฺสน\-สปฺปฏิฆํ นหีนํ ธมฺมํ ปจฺจยา อนิทสฺสน\-อปฺปฏิโฆ หีโน ธมฺโม อุปฺปชฺชติ เหตุปจฺจยา.  เหตุยา ตีณิ.}

\prose{นสนิทสฺสน\-สปฺปฏิฆํ นมชฺฌิมํ ธมฺมํ ปฏิจฺจ สนิทสฺสน\-สปฺปฏิโฆ มชฺฌิโม ธมฺโม อุปฺปชฺชติ เหตุปจฺจยา.  เหตุยา เอกวีส.}

\prose{นสนิทสฺสน\-สปฺปฏิฆํ นปณีตํ ธมฺมํ ปจฺจยา อนิทสฺสน\-อปฺปฏิโฆ ปณีโต ธมฺโม อุปฺปชฺชติ เหตุปจฺจยา.  เหตุยา ตีณิ.}

\csromanmarkh{22. สนิทสฺสน\-ตฺติก}\csromanmarkhii{15. มิจฺฉตฺตนิยต\-ตฺติก}\csromanheadernomark{2}{22. สนิทสฺสน\-ตฺติก 15. มิจฺฉตฺตนิยต\-ตฺติก}
\proseitem{62}{นสนิทสฺสน\-สปฺปฏิฆํ นมิจฺฉตฺต\-นิยตํ ธมฺมํ ปจฺจยา อนิทสฺสน\-อปฺปฏิโฆ มิจฺฉตฺต\-นิยโต ธมฺโม อุปฺปชฺชติ เหตุปจฺจยา.  เหตุยา ตีณิ.}

\csromanheader{2}{ธมฺม\-ปจฺจนียา\-นุโลม ติกติก\-ปฏฺฐาน\-ปาฬิ}
\prose{\csromanfolio{422}นสนิทสฺสน\-สปฺปฏิฆํ นสมฺมตฺต\-นิยตํ ธมฺมํ ปจฺจยา อนิทสฺสน\-อปฺปฏิโฆ สมฺมตฺตนิยโต ธมฺโม อุปฺปชฺชติ เหตุปจฺจยา.  เหตุยา ตีณิ.}

\prose{นสนิทสฺสน\-สปฺปฏิฆํ นอนิยตํ ธมฺมํ ปฏิจฺจ สนิทสฺสน\-สปฺปฏิโฆ อนิยโต ธมฺโม อุปฺปชฺชติ เหตุปจฺจยา.  เหตุยา เอกวีส.}

\csromanmarkh{22. สนิทสฺสน\-ตฺติก}\csromanmarkhii{16. มคฺคา\-รมฺมณ\-ตฺติก}\csromanheadernomark{2}{22. สนิทสฺสน\-ตฺติก 16. มคฺคา\-รมฺมณ\-ตฺติก}
\proseitem{63}{นสนิทสฺสน\-สปฺปฏิฆํ นมคฺคา\-รมฺมณํ ธมฺมํ ปจฺจยา อนิทสฺสน\-อปฺปฏิโฆ มคฺคารมฺมโณ ธมฺโม อุปฺปชฺชติ เหตุปจฺจยา.  เหตุยา ตีณิ.}

\prose{นสนิทสฺสน\-สปฺปฏิฆํ นมคฺคเหตุกํ ธมฺมํ ปจฺจยา อนิทสฺสน\-อปฺปฏิโฆ มคฺคเหตุโก ธมฺโม อุปฺปชฺชติ เหตุปจฺจยา.  เหตุยา ตีณิ.}

\prose{นสนิทสฺสน\-สปฺปฏิฆํ นมคฺคา\-ธิปติํ ธมฺมํ ปจฺจยา อนิทสฺสน\-อปฺปฏิโฆ มคฺคาธิปติ ธมฺโม อุปฺปชฺชติ เหตุปจฺจยา.  เหตุยา ตีณิ.}

\csromanmarkh{22. สนิทสฺสน\-ตฺติก}\csromanmarkhii{17. อุปฺปนฺนตฺติก}\csromanheadernomark{2}{22. สนิทสฺสน\-ตฺติก 17. อุปฺปนฺนตฺติก}
\proseitem{64}{นสนิทสฺสน\-สปฺปฏิโฆ นอุปฺปนฺโน ธมฺโม อนิทสฺสน\-อปฺปฏิฆสฺส อุปฺปนฺนสฺส ธมฺมสฺส อารมฺมณ\-ปจฺจเยน ปจฺจโย.  อารมฺมเณ ฉ.}

\csromanmarkh{22. สนิทสฺสน\-ตฺติก}\csromanmarkhii{18. อตีตตฺติก}\csromanheadernomark{2}{22. สนิทสฺสน\-ตฺติก 18. อตีตตฺติก}
\proseitem{65}{นสนิทสฺสน\-สปฺปฏิโฆ นปจฺจุปฺปนฺโน ธมฺโม อนิทสฺสน\-อปฺปฏิฆสฺส ปจฺจุปฺปนฺนสฺส ธมฺมสฺส อารมฺมณ\-ปจฺจเยน ปจฺจโย.  อารมฺมเณ ฉ.}

\csromanmarkh{22. สนิทสฺสน\-ตฺติก}\csromanmarkhii{19. อตีตา\-รมฺมณ\-ตฺติก}\csromanheadernomark{2}{22. สนิทสฺสน\-ตฺติก 19. อตีตา\-รมฺมณ\-ตฺติก}
\proseitem{66}{นสนิทสฺสน\-สปฺปฏิฆํ นอตีตา\-รมฺมณํ ธมฺมํ ปฏิจฺจ อนิทสฺสน\-อปฺปฏิโฆ อตีตารมฺมโณ ธมฺโม อุปฺปชฺชติ เหตุปจฺจยา.  เหตุยา ตีณิ.}

\prose{นสนิทสฺสน\-สปฺปฏิฆํ นอนาคตา\-รมฺมณํ ธมฺมํ ปจฺจยา อนิทสฺสน\-อปฺปฏิโฆ อนาคตารมฺมโณ ธมฺโม อุปฺปชฺชติ เหตุปจฺจยา.  เหตุยา ตีณิ.}

\prose{นสนิทสฺสน\-สปฺปฏิฆํ นปจฺจุปฺปนฺนา\-รมฺมณํ ธมฺมํ ปฏิจฺจ อนิทสฺสน\-อปฺปฏิโฆ ปจฺจุปฺปนฺนา\-รมฺมโณ ธมฺโม อุปฺปชฺชติ เหตุปจฺจยา.  เหตุยา ตีณิ.}

\vspace{\tipitakaparskip}
\csromancenter{สนิทสฺสน\-ตฺติก 20-21. อชฺฌตฺต\-ตฺติก\-ทฺวย 22. สนิทสฺสน\-ตฺติก 20-21. อชฺฌตฺต\-ตฺติก\-ทฺวย}
\proseitem{67}{\csromanfolio{423}นสนิทสฺสน\-สปฺปฏิโฆ นอชฺฌตฺโต ธมฺโม อนิทสฺสน\-อปฺปฏิฆสฺส อชฺฌตฺตสฺส ธมฺมสฺส อารมฺมณ\-ปจฺจเยน ปจฺจโย. (สํขิตฺตํ.) . อารมฺมเณ ฉ, อธิปติยา ฉ, ปุเรชาเต อตฺถิยา อวิคเต ฉ.}

\prose{นสนิทสฺสน\-สปฺปฏิโฆ นพหิทฺธา ธมฺโม อนิทสฺสน\-อปฺปฏิฆสฺส พหิทฺธา ธมฺมสฺส อารมฺมณ\-ปจฺจเยน ปจฺจโย. (สํขิตฺตํ.) . อารมฺมเณ ฉ, อธิปติยา ฉ, ปุเรชาเต อตฺถิยา อวิคเต ฉ.}

\proseitem{68}{นสนิทสฺสน\-สปฺปฏิฆํ นอชฺฌตฺตา\-รมฺมณํ ธมฺมํ ปฏิจฺจ อนิทสฺสน\-อปฺปฏิโฆ อชฺฌตฺตา\-รมฺมโณ ธมฺโม อุปฺปชฺชติ เหตุปจฺจยา.  เหตุยา ตีณิ.}

\prose{นสนิทสฺสน\-สปฺปฏิฆํ นพหิทฺธา\-รมฺมณํ ธมฺมํ ปฏิจฺจ อนิทสฺสน\-อปฺปฏิโฆ พหิทฺธา\-รมฺมโณ ธมฺโม อุปฺปชฺชติ เหตุปจฺจยา.  เหตุยา ตีณิ ฯเปฯ อวิคเต ตีณิ.}

\prosewithcloser{(สหชาตวารมฺปิ ปจฺจยวารมฺปิ นิสฺสยวารมฺปิ สํสฏฺฐ\-วารมฺปิ สมฺปยุตฺตวารมฺปิ ปญฺหาวารมฺปิ วิตฺถาเรตพฺพํ.)}{ธมฺม\-ปจฺจนียา\-นุโลเม ติกติก\-ปฏฺฐานํ นิฏฺฐิตํ.}
\pitaka{อภิธมฺม\-ปิฏก}
\csromanheader{2}{ธมฺม\-ปจฺจนียา\-นุโลม}
\namakkaram{นโม ตสฺส ภควโต อรหโต สมฺมา\-สมฺพุทฺธสฺส.}
\csromanmarkh{1. เหตุทุก}\csromanmarkhii{1. สเหตุกทุก}\csromanheadernomark{2}{1. \textbf{เหตุทุก 1. สเหตุกทุก}}
\proseitem{1}{\csromanfolio{425}นเหตุํ นสเหตุกํ ธมฺมํ ปฏิจฺจ เหตุ สเหตุโก ธมฺโม อุปฺปชฺชติ เหตุปจฺจยา.  นเหตุํ นสเหตุกํ ธมฺมํ ปฏิจฺจ นเหตุ สเหตุโก ธมฺโม อุปฺปชฺชติ เหตุปจฺจยา.  นเหตุํ นสเหตุกํ ธมฺมํ ปฏิจฺจ เหตุ สเหตุโก จ นเหตุ สเหตุโก จ ธมฺมา อุปฺปชฺชนฺติ เหตุปจฺจยา. ตีณิ.}

\prose{นนเหตุํ นสเหตุกํ ธมฺมํ ปฏิจฺจ นเหตุ สเหตุโก ธมฺโม อุปฺปชฺชติ เหตุปจฺจยา. เอกํ.}

\prose{เหตุยา จตฺตาริ, อารมฺมเณ จตฺตาริ ฯเปฯ อวิคเต จตฺตาริ.}

\proseitem{2}{นนเหตุ นอเหตุโก ธมฺโม นเหตุสฺส อเหตุกสฺส ธมฺมสฺส เหตุปจฺจเยน ปจฺจโย.}

\prose{เหตุยา เอกํ, อารมฺมเณ ฉ ฯเปฯ อวิคเต ปญฺจ. (ปญฺหาวารํ วิตฺถาเรตพฺพํ.)}

\proseitem{3}{นเหตุํ นอเหตุกํ ธมฺมํ ปฏิจฺจ นเหตุ อเหตุโก ธมฺโม อุปฺปชฺชติ เหตุปจฺจยา. เอกํ.}

\csromanheader{2}{ธมฺม\-ปจฺจนียา\-นุโลม ทุกทุก\-ปฏฺฐาน\-ปาฬิ}
\prose{\csromanfolio{426}นนเหตุํ นอเหตุกํ ธมฺมํ ปฏิจฺจ นเหตุ อเหตุโก ธมฺโม อุปฺปชฺชติ เหตุปจฺจยา. เอกํ.}

\prose{นเหตุํ นอเหตุกญฺจ นนเหตุํ นอเหตุกญฺจ ธมฺมํ ปฏิจฺจ นเหตุ อเหตุโก ธมฺโม อุปฺปชฺชติ เหตุปจฺจยา. เอกํ.  เหตุยา ตีณิ.}

\csromanheader{2}{1. เหตุทุก 2-5. เหตุ\-สมฺปยุตฺต\-ทุกาทิ}
\proseitem{4}{นเหตุํ นเหตุ\-สมฺปยุตฺตํ ธมฺมํ ปฏิจฺจ เหตุ เหตุสมฺปยุตฺโต ธมฺโม อุปฺปชฺชติ เหตุปจฺจยา.  เหตุยา จตฺตาริ.}

\prose{นเหตุํ นเหตุ\-วิปฺปยุตฺตํ ธมฺมํ ปฏิจฺจ นเหตุ เหตุวิปฺปยุตฺโต ธมฺโม อุปฺปชฺชติ เหตุปจฺจยา.  เหตุยา ตีณิ.}

\proseitem{5}{นเหตุํ นเหตุญฺเจว นอเหตุกญฺจ ธมฺมํ ปฏิจฺจ เหตุ เหตุ เจว สเหตุโก จ ธมฺโม อุปฺปชฺชติ เหตุปจฺจยา.  เหตุยา เอกํ.}

\prose{นนเหตุํ นอเหตุกญฺเจว นนเหตุญฺจ ธมฺมํ ปฏิจฺจ นเหตุ สเหตุโก เจว น จ เหตุ ธมฺโม อุปฺปชฺชติ เหตุปจฺจยา.  เหตุยา เอกํ.}

\proseitem{6}{นเหตุํ นเหตุญฺเจว นเหตุ\-วิปฺปยุตฺตญฺจ ธมฺมํ ปฏิจฺจ เหตุ เหตุ เจว เหตุสมฺปยุตฺโต จ ธมฺโม อุปฺปชฺชติ เหตุปจฺจยา.  เหตุยา เอกํ.}

\prose{นนเหตุํ นเหตุ\-วิปฺปยุตฺต\-ญฺเจว นนเหตุญฺจ ธมฺมํ ปฏิจฺจ นเหตุ เหตุสมฺปยุตฺโต เจว น จ เหตุ ธมฺโม อุปฺปชฺชติ เหตุปจฺจยา.  เหตุยา เอกํ.}

\proseitem{7}{นเหตุํ นเหตุํ นสเหตุกํ ธมฺมํ ปฏิจฺจ นเหตุ นเหตุ สเหตุโก ธมฺโม อุปฺปชฺชติ เหตุปจฺจยา.  เหตุยา เอกํ.}

\prose{นเหตุํ นเหตุํ นอเหตุกํ ธมฺมํ ปฏิจฺจ นเหตุ นเหตุ อเหตุโก ธมฺโม อุปฺปชฺชติ เหตุปจฺจยา.  เหตุยา เอกํ.}

\csromanheader{2}{1. \textbf{เหตุทุก 6-1}2. จูฬนฺตรทุก}
\proseitem{8}{นเหตุ นสปฺปจฺจโย ธมฺโม เหตุสฺส สปฺปจฺจยสฺส ธมฺมสฺส อารมฺมณ\-ปจฺจเยน ปจฺจโย.  อารมฺมเณ ตีณิ. (สงฺขตํ สปฺปจฺจย\-สทิสํ.)}

\csromanheader{2}{1. เหตุทุก 13-18. อาสวโคจฺฉก}
\proseitem{9}{\csromanfolio{427}นเหตุํ นสนิทสฺสนํ ธมฺมํ ปฏิจฺจ นเหตุ สนิทสฺสโน ธมฺโม อุปฺปชฺชติ เหตุปจฺจยา.  เหตุยา ตีณิ.}

\prose{นเหตุ นอนิทสฺสโน ธมฺโม เหตุสฺส อนิทสฺสนสฺส ธมฺมสฺส อารมฺมณ\-ปจฺจเยน ปจฺจโย.  อารมฺมเณ ตีณิ.}

\proseitem{10}{นเหตุํ นสปฺปฏิฆํ ธมฺมํ ปฏิจฺจ นเหตุ สปฺปฏิโฆ ธมฺโม อุปฺปชฺชติ เหตุปจฺจยา.  เหตุยา ตีณิ.}

\prose{นเหตุํ นอปฺปฏิฆํ ธมฺมํ ปฏิจฺจ นเหตุ อปฺปฏิโฆ ธมฺโม อุปฺปชฺชติ เหตุปจฺจยา.  เหตุยา เอกํ.}

\proseitem{11}{นเหตุํ นรูปิํ ธมฺมํ ปฏิจฺจ เหตุ รูปี ธมฺโม อุปฺปชฺชติ เหตุปจฺจยา.  เหตุยา ตีณิ.}

\prose{นเหตุํ นอรูปิํ ธมฺมํ ปฏิจฺจ เหตุ อรูปี ธมฺโม อุปฺปชฺชติ เหตุปจฺจยา.  เหตุยา ตีณิ.}

\proseitem{12}{นเหตุํ นโลกิยํ ธมฺมํ ปฏิจฺจ นเหตุ โลกิโย ธมฺโม อุปฺปชฺชติ เหตุปจฺจยา.  นนเหตุํ นโลกิยํ ธมฺมํ ปฏิจฺจ นเหตุ โลกิโย ธมฺโม อุปฺปชฺชติ เหตุปจฺจยา. (คณิตเกน ตีณิ.)}

\prose{นเหตุํ นโลกุตฺตรํ ธมฺมํ ปจฺจยา เหตุ โลกุตฺตโร ธมฺโม อุปฺปชฺชติ เหตุปจฺจยา.  เหตุยา ตีณิ.}

\proseitem{13}{นเหตุํ นเกนจิ วิญฺเญยฺยํ ธมฺมํ ปฏิจฺจ เหตุ เกนจิ วิญฺเญยฺโย ธมฺโม อุปฺปชฺชติ เหตุปจฺจยา.  เหตุยา นว.}

\prose{นเหตุํ นเกนจิ นวิญฺเญยฺยํ ธมฺมํ ปฏิจฺจ เหตุ เกนจิ นวิญฺเญยฺโย ธมฺโม อุปฺปชฺชติ เหตุปจฺจยา.  เหตุยา นว.}

\csromanheader{2}{1. เหตุทุก 13-18. อาสวโคจฺฉก}
\proseitem{14}{นเหตุํ โนอาสวํ ธมฺมํ ปฏิจฺจ เหตุ อาสโว ธมฺโม อุปฺปชฺชติ เหตุปจฺจยา.  นเหตุํ โนอาสวํ ธมฺมํ ปฏิจฺจ นเหตุ อาสโว ธมฺโม อุปฺปชฺชติ เหตุปจฺจยา.  นเหตุํ โนอาสวํ ธมฺมํ ปฏิจฺจ เหตุ อาสโว จ นเหตุ อาสโว จ ธมฺมา อุปฺปชฺชนฺติ เหตุปจฺจยา. ตีณิ.}

\csromanheader{2}{ธมฺม\-ปจฺจนียา\-นุโลม ทุกทุก\-ปฏฺฐาน\-ปาฬิ}
\prose{\csromanfolio{428}นนเหตุํ โนอาสวํ ธมฺมํ ปฏิจฺจ เหตุ อาสโว ธมฺโม อุปฺปชฺชติ เหตุปจฺจยา. เอกํ.}

\prose{นเหตุํ โนอาสวญฺจ นนเหตุํ โนอาสวญฺจ ธมฺมํ ปฏิจฺจ เหตุ อาสโว ธมฺโม อุปฺปชฺชติ เหตุปจฺจยา. เอกํ.  เหตุยา ปญฺจ.}

\prose{นเหตุํ นโนอาสวํ ธมฺมํ ปฏิจฺจ นเหตุ โนอาสโว ธมฺโม อุปฺปชฺชติ เหตุปจฺจยา. เอกํ.}

\prose{นนเหตุํ นโนอาสวํ ธมฺมํ ปฏิจฺจ นเหตุ โนอาสโว ธมฺโม อุปฺปชฺชติ เหตุปจฺจยา.  นนเหตุํ นโนอาสวํ ธมฺมํ ปฏิจฺจ เหตุ โนอาสโว ธมฺโม อุปฺปชฺชติ เหตุปจฺจยา.  นนเหตุํ นโนอาสวํ ธมฺมํ ปฏิจฺจ เหตุ โนอาสโว จ นเหตุ โนอาสโว จ ธมฺมา อุปฺปชฺชนฺติ เหตุปจฺจยา. ตีณิ.}

\prose{นเหตุํ นโนอาสวญฺจ นนเหตุํ นโนอาสวญฺจ ธมฺมํ ปฏิจฺจ นเหตุ โนอาสโว ธมฺโม อุปฺปชฺชติ เหตุปจฺจยา. เอกํ.  เหตุยา ปญฺจ.}

\proseitem{15}{นเหตุํ นสาสวํ ธมฺมํ ปฏิจฺจ นเหตุ สาสโว ธมฺโม อุปฺปชฺชติ เหตุปจฺจยา.  เหตุยา ตีณิ.}

\prose{นเหตุํ นอนาสวํ ธมฺมํ ปจฺจยา เหตุ อนาสโว ธมฺโม อุปฺปชฺชติ เหตุปจฺจยา.  เหตุยา ตีณิ.}

\proseitem{16}{นนเหตุํ นอาสว\-สมฺปยุตฺตํ ธมฺมํ ปฏิจฺจ นเหตุ อาสว\-สมฺปยุตฺโต ธมฺโม อุปฺปชฺชติ เหตุปจฺจยา.  เหตุยา ตีณิ.}

\prose{นเหตุํ นอาสว\-วิปฺปยุตฺตํ ธมฺมํ ปฏิจฺจ เหตุ อาสว\-วิปฺปยุตฺโต ธมฺโม อุปฺปชฺชติ เหตุปจฺจยา.  เหตุยา นว.}

\proseitem{17}{นเหตุํ นอาสวญฺเจว นอนาสวญฺจ ธมฺมํ ปฏิจฺจ เหตุ อาสโว เจว สาสโว จ ธมฺโม อุปฺปชฺชติ เหตุปจฺจยา.  เหตุยา ปญฺจ.}

\prose{นเหตุํ นอนาสวญฺเจว นโน จ อาสวํ ธมฺมํ ปฏิจฺจ นเหตุ สาสโว เจว โน จ อาสโว ธมฺโม อุปฺปชฺชติ เหตุปจฺจยา.  เหตุยา ปญฺจ.}

\csromanmarkh{1. เหตุทุก}\csromanmarkhii{19. สญฺโญชนาทิโคจฺฉก}\csromanheadernomark{2}{1. เหตุทุก 19. สญฺโญชนาทิโคจฺฉก}
\proseitem{18}{\csromanfolio{429}นเหตุํ นอาสวญฺเจว นอาสว\-วิปฺปยุตฺตญฺจ ธมฺมํ ปฏิจฺจ เหตุ อาสโว เจว อาสว\-สมฺปยุตฺโต จ ธมฺโม อุปฺปชฺชติ เหตุปจฺจยา.  เหตุยา ตีณิ.}

\prose{นเหตุํ นอาสว\-วิปฺปยุตฺต\-ญฺเจว นโน จ อาสวํ ธมฺมํ ปฏิจฺจ นเหตุ อาสว\-สมฺปยุตฺโต เจว โน จ อาสโว ธมฺโม อุปฺปชฺชติ เหตุปจฺจยา.  เหตุยา ตีณิ.}

\proseitem{19}{นเหตุํ อาสว\-วิปฺปยุตฺตํ นสาสวํ ธมฺมํ ปฏิจฺจ นเหตุ อาสว\-วิปฺปยุตฺโต สาสโว ธมฺโม อุปฺปชฺชติ เหตุปจฺจยา.  เหตุยา ตีณิ.}

\prose{นเหตุํ อาสว\-วิปฺปยุตฺตํ นอนาสวํ ธมฺมํ ปจฺจยา เหตุ อาสว\-วิปฺปยุตฺโต อนาสโว ธมฺโม อุปฺปชฺชติ เหตุปจฺจยา.  เหตุยา ตีณิ.}

\csromanmarkh{1. เหตุทุก}\csromanmarkhii{19. สญฺโญชนาทิโคจฺฉก}\csromanheadernomark{2}{1. เหตุทุก 19. สญฺโญชนาทิโคจฺฉก}
\proseitem{20}{นเหตุํ โนสญฺโญชนํ ธมฺมํ ปฏิจฺจ เหตุ สญฺโญชโน ธมฺโม อุปฺปชฺชติ เหตุปจฺจยา.  เหตุยา ตีณิ.}

\proseitem{21}{นเหตุํ โนคนฺถํ ธมฺมํ ปฏิจฺจ เหตุ คนฺโถ ธมฺโม อุปฺปชฺชติ เหตุปจฺจยา.  เหตุยา นว.}

\proseitem{22}{นเหตุํ โนโอฆํ ธมฺมํ ปฏิจฺจ เหตุ โอโฆ ธมฺโม อุปฺปชฺชติ เหตุปจฺจยา.  เหตุยา ปญฺจ.}

\proseitem{23}{นเหตุํ โนโยคํ ธมฺมํ ปฏิจฺจ เหตุ โยโค ธมฺโม อุปฺปชฺชติ เหตุปจฺจยา.  เหตุยา ปญฺจ.}

\proseitem{24}{นเหตุํ โนนีวรณํ ธมฺมํ ปฏิจฺจ เหตุ นีวรโณ ธมฺโม อุปฺปชฺชติ เหตุปจฺจยา.  เหตุยา ตีณิ.}

\proseitem{25}{นเหตุํ โนปรามาสํ ธมฺมํ ปฏิจฺจ นเหตุ ปรามาโส ธมฺโม อุปฺปชฺชติ เหตุปจฺจยา.  เหตุยา ตีณิ.}

\csromanmarkh{ธมฺม\-ปจฺจนียา\-นุโลม ทุกทุก\-ปฏฺฐาน\-ปาฬิ}\csromanmarkhii{1. เหตุทุก 54. มหนฺตรทุก}\csromanheadernomark{2}{ธมฺม\-ปจฺจนียา\-นุโลม ทุกทุก\-ปฏฺฐาน\-ปาฬิ 1. เหตุทุก 54. มหนฺตรทุก}
\proseitem{26}{\csromanfolio{430}นเหตุํ นสารมฺมณํ ธมฺมํ ปฏิจฺจ เหตุ สารมฺมโณ ธมฺโม อุปฺปชฺชติ เหตุปจฺจยา.  เหตุยา ตีณิ.}

\prose{นเหตุํ นอนารมฺมณํ ธมฺมํ ปฏิจฺจ นเหตุ อนารมฺมโณ ธมฺโม อุปฺปชฺชติ เหตุปจฺจยา.  เหตุยา ตีณิ.}

\proseitem{27}{นเหตุํ โนจิตฺตํ ธมฺมํ ปฏิจฺจ นเหตุ จิตฺโต ธมฺโม อุปฺปชฺชติ เหตุปจฺจยา.  เหตุยา ตีณิ.}

\prose{นเหตุํ นโนจิตฺตํ ธมฺมํ ปฏิจฺจ เหตุ โนจิตฺโต ธมฺโม อุปฺปชฺชติ เหตุปจฺจยา.  เหตุยา ตีณิ.}

\proseitem{28}{นเหตุํ นเจตสิกํ ธมฺมํ ปฏิจฺจ เหตุ เจตสิโก ธมฺโม อุปฺปชฺชติ เหตุปจฺจยา.  เหตุยา ตีณิ.}

\prose{นเหตุํ นอเจตสิกํ ธมฺมํ ปฏิจฺจ นเหตุ อเจตสิโก ธมฺโม อุปฺปชฺชติ เหตุปจฺจยา.  เหตุยา ตีณิ.}

\proseitem{29}{นเหตุํ นจิตฺต\-สมฺปยุตฺตํ ธมฺมํ ปฏิจฺจ เหตุ จิตฺต\-สมฺปยุตฺโต ธมฺโม อุปฺปชฺชติ เหตุปจฺจยา.  เหตุยา ตีณิ.}

\prose{นเหตุํ นจิตฺต\-วิปฺปยุตฺตํ ธมฺมํ ปฏิจฺจ นเหตุ จิตฺต\-วิปฺปยุตฺโต ธมฺโม อุปฺปชฺชติ เหตุปจฺจยา.  เหตุยา ตีณิ.}

\proseitem{30}{นเหตุํ นจิตฺต\-สํสฏฺฐํ ธมฺมํ ปฏิจฺจ เหตุ จิตฺตสํสฏฺโฐ ธมฺโม อุปฺปชฺชติ เหตุปจฺจยา.  เหตุยา ตีณิ.}

\prose{นเหตุํ นจิตฺต\-วิสํสฏฺฐํ ธมฺมํ ปฏิจฺจ นเหตุ จิตฺต\-วิสํสฏฺโฐ ธมฺโม อุปฺปชฺชติ เหตุปจฺจยา.  เหตุยา ตีณิ.}

\proseitem{31}{นเหตุํ โนจิตฺต\-สมุฏฺฐานํ ธมฺมํ ปฏิจฺจ เหตุ จิตฺต\-สมุฏฺฐาโน ธมฺโม อุปฺปชฺชติ เหตุปจฺจยา.  เหตุยา ตีณิ.}

\prose{นเหตุํ นโนจิตฺตสมุฏฺฐานํ ธมฺมํ ปฏิจฺจ เหตุ โนจิตฺต\-สมุฏฺฐาโน ธมฺโม อุปฺปชฺชติ เหตุปจฺจยา.  เหตุยา ตีณิ.}

\csromanmarkh{1. เหตุทุก}\csromanmarkhii{54. มหนฺตรทุก}\csromanheadernomark{2}{1. เหตุทุก 54. มหนฺตรทุก}
\proseitem{32}{\csromanfolio{431}นเหตุํ โนจิตฺต\-สหภุํ ธมฺมํ ปฏิจฺจ เหตุ จิตฺตสหภู ธมฺโม อุปฺปชฺชติ เหตุปจฺจยา.  เหตุยา ตีณิ.}

\prose{นเหตุํ นโนจิตฺตสหภุํ ธมฺมํ ปฏิจฺจ นเหตุ โนจิตฺตสหภู ธมฺโม อุปฺปชฺชติ เหตุปจฺจยา.  เหตุยา ตีณิ.}

\proseitem{33}{นเหตุํ นจิตฺตา\-นุปริวตฺติํ ธมฺมํ ปฏิจฺจ เหตุ จิตฺตา\-นุปริวตฺตี ธมฺโม อุปฺปชฺชติ เหตุปจฺจยา.  เหตุยา ตีณิ.}

\prose{นเหตุํ นโนจิตฺตานุปริวตฺติํ ธมฺมํ ปฏิจฺจ นเหตุ โนจิตฺตา\-นุปริวตฺตี ธมฺโม อุปฺปชฺชติ เหตุปจฺจยา.  เหตุยา ตีณิ.}

\proseitem{34}{นเหตุํ นจิตฺต\-สํสฏฺฐ\-สมุฏฺฐานํ ธมฺมํ ปฏิจฺจ เหตุ จิตฺต\-สํสฏฺฐ\-สมุฏฺฐาโน ธมฺโม อุปฺปชฺชติ เหตุปจฺจยา.  เหตุยา ตีณิ.}

\prose{นเหตุํ นโนจิตฺตสํสฏฺฐสมุฏฺฐานํ ธมฺมํ ปฏิจฺจ นเหตุ โนจิตฺต\-สํสฏฺฐ\-สมุฏฺฐาโน ธมฺโม อุปฺปชฺชติ เหตุปจฺจยา.  เหตุยา ตีณิ.}

\proseitem{35}{นเหตุํ โนจิตฺต\-สํสฏฺฐ\-สมุฏฺฐาน\-สหภุํ ธมฺมํ ปฏิจฺจ เหตุ จิตฺต\-สํสฏฺฐ\-สมุฏฺฐาน\-สหภู ธมฺโม อุปฺปชฺชติ เหตุปจฺจยา.  เหตุยา ตีณิ.}

\prose{นเหตุํ นโนจิตฺตสํสฏฺฐสมุฏฺฐานสหภุํ ธมฺมํ ปฏิจฺจ นเหตุ โนจิตฺต\-สํสฏฺฐ\-สมุฏฺฐาน\-สหภู ธมฺโม อุปฺปชฺชติ เหตุปจฺจยา.  เหตุยา ตีณิ.}

\proseitem{36}{นเหตุํ โนจิตฺต\-สํสฏฺฐ\-สมุฏฺฐานา\-นุปริวตฺติํ ธมฺมํ ปฏิจฺจ เหตุ จิตฺต\-สํสฏฺฐ\-สมุฏฺฐานา\-นุปริวตฺตี ธมฺโม อุปฺปชฺชติ เหตุปจฺจยา.  เหตุยา ตีณิ.}

\prose{นเหตุํ นโนจิตฺตสํสฏฺฐสมุฏฺฐานานุปริวตฺติํ ธมฺมํ ปฏิจฺจ นเหตุ โนจิตฺต\-สํสฏฺฐ\-สมุฏฺฐานา\-นุปริวตฺตี ธมฺโม อุปฺปชฺชติ เหตุปจฺจยา.  เหตุยา ตีณิ.}

\proseitem{37}{นเหตุํ นอชฺฌตฺติกํ ธมฺมํ ปฏิจฺจ นเหตุ อชฺฌตฺติโก ธมฺโม อุปฺปชฺชติ เหตุปจฺจยา.  เหตุยา ตีณิ.}

\csromanheader{2}{ธมฺม\-ปจฺจนียา\-นุโลม ทุกทุก\-ปฏฺฐาน\-ปาฬิ}
\prose{\csromanfolio{432}นเหตุํ นพาหิรํ ธมฺมํ ปฏิจฺจ เหตุ พาหิโร ธมฺโม อุปฺปชฺชติ เหตุปจฺจยา.  เหตุยา ตีณิ.}

\proseitem{38}{นเหตุํ นอุปาทา ธมฺมํ ปฏิจฺจ นเหตุ อุปาทา ธมฺโม อุปฺปชฺชติ เหตุปจฺจยา.  เหตุยา ตีณิ.}

\prose{นเหตุํ นโนอุปาทา ธมฺมํ ปฏิจฺจ เหตุ โนอุปาทา ธมฺโม อุปฺปชฺชติ เหตุปจฺจยา.  เหตุยา ตีณิ.}

\proseitem{39}{นเหตุํ นอนุปาทินฺนํ ธมฺมํ ปฏิจฺจ นเหตุ อนุปาทินฺโน ธมฺโม อุปฺปชฺชติ เหตุปจฺจยา.  เหตุยา ตีณิ.}

\csromanmarkh{1. เหตุทุก}\csromanmarkhii{68. อุปาทาน\-โคจฺฉก}\csromanheadernomark{2}{1. เหตุทุก 68. อุปาทาน\-โคจฺฉก}
\proseitem{40}{นเหตุํ นอุปาทานํ ธมฺมํ ปฏิจฺจ เหตุ อุปาทาโน ธมฺโม อุปฺปชฺชติ เหตุปจฺจยา.  เหตุยา นว.}

\csromanmarkh{1. เหตุทุก}\csromanmarkhii{74. กิเลสโคจฺฉก}\csromanheadernomark{2}{1. เหตุทุก 74. กิเลสโคจฺฉก}
\proseitem{41}{นเหตุํ นกิเลสํ ธมฺมํ ปฏิจฺจ เหตุ กิเลโส ธมฺโม อุปฺปชฺชติ เหตุปจฺจยา.  เหตุยา ตีณิ.}

\csromanmarkh{1. เหตุทุก}\csromanmarkhii{82. ปิฏฺฐิทุก}\csromanheadernomark{2}{1. เหตุทุก 82. ปิฏฺฐิทุก}
\proseitem{42}{นเหตุํ นทสฺสเนน ปหาตพฺพํ ธมฺมํ ปจฺจยา เหตุ ทสฺสเนน ปหาตพฺโพ ธมฺโม อุปฺปชฺชติ เหตุปจฺจยา.  เหตุยา ตีณิ.}

\prose{นเหตุํ นนทสฺสเนน ปหาตพฺพํ ธมฺมํ ปฏิจฺจ นเหตุ นทสฺสเนน ปหาตพฺโพ ธมฺโม อุปฺปชฺชติ เหตุปจฺจยา.  เหตุยา ตีณิ.}

\proseitem{43}{นเหตุํ นภาวนาย ปหาตพฺพํ ธมฺมํ ปจฺจยา เหตุ ภาวนาย ปหาตพฺโพ ธมฺโม อุปฺปชฺชติ เหตุปจฺจยา.  เหตุยา ตีณิ.}

\prose{นเหตุํ นนภาวนาย ปหาตพฺพํ ธมฺมํ ปฏิจฺจ นเหตุ นภาวนาย ปหาตพฺโพ ธมฺโม อุปฺปชฺชติ เหตุปจฺจยา.  เหตุยา ตีณิ.}

\proseitem{44}{นนเหตุํ นทสฺสเนน ปหาตพฺพ\-เหตุกํ ธมฺมํ ปฏิจฺจ นเหตุ ทสฺสเนน ปหาตพฺพ\-เหตุโก ธมฺโม อุปฺปชฺชติ เหตุปจฺจยา.  เหตุยา เอกํ.}

\csromanmarkh{2. สเหตุกาทิทุก}\csromanmarkhii{1. เหตุทุก}\csromanheadernomark{2}{2. สเหตุกาทิทุก 1. เหตุทุก}
\prose{\csromanfolio{433}นเหตุํ นนทสฺสเนน ปหาตพฺพ\-เหตุกํ ธมฺมํ ปฏิจฺจ นเหตุ นทสฺสเนน ปหาตพฺพ\-เหตุโก ธมฺโม อุปฺปชฺชติ เหตุปจฺจยา.  เหตุยา ตีณิ.}

\proseitem{45}{นนเหตุํ นภาวนาย ปหาตพฺพ\-เหตุกํ ธมฺมํ ปฏิจฺจ นเหตุ ภาวนาย ปหาตพฺพ\-เหตุโก ธมฺโม อุปฺปชฺชติ เหตุปจฺจยา.  เหตุยา เอกํ.}

\prose{นเหตุํ นนภาวนาย ปหาตพฺพ\-เหตุกํ ธมฺมํ ปฏิจฺจ นเหตุ นภาวนาย ปหาตพฺพ\-เหตุโก ธมฺโม อุปฺปชฺชติ เหตุปจฺจยา.  เหตุยา ตีณิ.}

\proseitem{46}{นเหตุํ นสวิตกฺกํ ธมฺมํ ปฏิจฺจ เหตุ สวิตกฺโก ธมฺโม อุปฺปชฺชติ เหตุปจฺจยา.  เหตุยา ตีณิ. (สพฺพตฺถ สํขิตฺตํ.)}

\proseitem{47}{นเหตุํ นสรณํ ธมฺมํ ปจฺจยา เหตุ สรโณ ธมฺโม อุปฺปชฺชติ เหตุปจฺจยา.  เหตุยา ตีณิ.}

\prose{นเหตุํ นอรณํ ธมฺมํ ปฏิจฺจ นเหตุ อรโณ ธมฺโม อุปฺปชฺชติ เหตุปจฺจยา.  นนเหตุํ นอรณํ ธมฺมํ ปฏิจฺจ นเหตุ อรโณ ธมฺโม อุปฺปชฺชติ เหตุปจฺจยา.  นเหตุํ นอรณญฺจ นนเหตุํ นอรณญฺจ ธมฺมํ ปฏิจฺจ นเหตุ อรโณ ธมฺโม อุปฺปชฺชติ เหตุปจฺจยา.  เหตุยา ตีณิ. 2. สเหตุกาทิทุก 1. เหตุทุก}

\proseitem{48}{นสเหตุกํ นเหตุํ ธมฺมํ ปฏิจฺจ สเหตุโก เหตุ ธมฺโม อุปฺปชฺชติ เหตุปจฺจยา.  นอเหตุกํ นเหตุํ ธมฺมํ ปฏิจฺจ สเหตุโก เหตุ ธมฺโม อุปฺปชฺชติ เหตุปจฺจยา.  นสเหตุกํ นเหตุญฺจ นอเหตุกํ นเหตุญฺจ ธมฺมํ ปฏิจฺจ สเหตุโก เหตุ ธมฺโม อุปฺปชฺชติ เหตุปจฺจยา.  เหตุยา ตีณิ.}

\prose{นสเหตุกํ นนเหตุํ ธมฺมํ ปฏิจฺจ สเหตุโก นเหตุ ธมฺโม อุปฺปชฺชติ เหตุปจฺจยา. ตีณิ.}

\prose{นอเหตุกํ นนเหตุํ ธมฺมํ ปฏิจฺจ อเหตุโก นเหตุ ธมฺโม อุปฺปชฺชติ เหตุปจฺจยา. ตีณิ.  เหตุยา ฉ.}

\proseitem{49}{นเหตุ\-สมฺปยุตฺตํ นเหตุํ ธมฺมํ ปฏิจฺจ เหตุสมฺปยุตฺโต เหตุ ธมฺโม อุปฺปชฺชติ เหตุปจฺจยา. ตีณิ. (สเหตุก\-ทุก\-สทิสํ.)}

\csromanheader{2}{ธมฺม\-ปจฺจนียา\-นุโลม ทุกทุก\-ปฏฺฐาน\-ปาฬิ}
\proseitem{50}{\csromanfolio{434}นเหตุญฺเจว นอเหตุกญฺจ นเหตุํ ธมฺมํ ปฏิจฺจ เหตุ เจว สเหตุโก จ เหตุ ธมฺโม อุปฺปชฺชติ เหตุปจฺจยา.  เหตุยา เอกํ.}

\prose{นอเหตุกญฺเจว นน จ เหตุํ นนเหตุํ ธมฺมํ ปฏิจฺจ สเหตุโก เจว น จ เหตุ นเหตุ ธมฺโม อุปฺปชฺชติ เหตุปจฺจยา.  เหตุยา เอกํ.}

\proseitem{51}{นเหตุญฺเจว นเหตุ\-วิปฺปยุตฺตญฺจ นเหตุํ ธมฺมํ ปฏิจฺจ เหตุ เจว เหตุสมฺปยุตฺโต จ เหตุ ธมฺโม อุปฺปชฺชติ เหตุปจฺจยา.  เหตุยา เอกํ.}

\prose{นเหตุ\-วิปฺปยุตฺต\-ญฺเจว นนเหตุญฺจ นนเหตุํ ธมฺมํ ปฏิจฺจ เหตุสมฺปยุตฺโต เจว น จ เหตุ นเหตุ ธมฺโม อุปฺปชฺชติ เหตุปจฺจยา.  เหตุยา เอกํ. (อนฺติมทุกํ น ลพฺภติ.)}

\csromanmarkh{7-13. จูฬนฺตรทุก}\csromanmarkhii{1. เหตุทุก}\csromanheadernomark{2}{7-13. จูฬนฺตรทุก 1. เหตุทุก}
\proseitem{52}{นอปฺปจฺจยํ นเหตุํ ธมฺมํ ปฏิจฺจ สปฺปจฺจโย เหตุ ธมฺโม อุปฺปชฺชติ เหตุปจฺจยา.  เหตุยา เอกํ.}

\prose{นอปฺปจฺจยํ นนเหตุํ ธมฺมํ ปฏิจฺจ สปฺปจฺจโย นเหตุ ธมฺโม อุปฺปชฺชติ เหตุปจฺจยา.  เหตุยา เอกํ. (สงฺขตํ สปฺปจฺจย\-สทิสํ.)}

\proseitem{53}{นสนิทสฺสนํ นเหตุํ ธมฺมํ ปฏิจฺจ อนิทสฺสโน เหตุ ธมฺโม อุปฺปชฺชติ เหตุปจฺจยา.  เหตุยา เอกํ.}

\prose{นสนิทสฺสนํ นนเหตุํ ธมฺมํ ปฏิจฺจ สนิทสฺสโน นเหตุ ธมฺโม อุปฺปชฺชติ เหตุปจฺจยา.  เหตุยา ตีณิ.}

\proseitem{54}{นสปฺปฏิฆํ นเหตุํ ธมฺมํ ปฏิจฺจ อปฺปฏิโฆ เหตุ ธมฺโม อุปฺปชฺชติ เหตุปจฺจยา.  เหตุยา เอกํ.}

\prose{นสปฺปฏิฆํ นนเหตุํ ธมฺมํ ปฏิจฺจ สปฺปฏิโฆ นเหตุ ธมฺโม อุปฺปชฺชติ เหตุปจฺจยา.  เหตุยา ตีณิ.}

\proseitem{55}{นรูปิํ นเหตุํ ธมฺมํ ปฏิจฺจ อรูปี เหตุ ธมฺโม อุปฺปชฺชติ เหตุปจฺจยา.}

\csromanmarkh{14. อาสวโคจฺฉก}\csromanmarkhii{1. เหตุทุก}\csromanheadernomark{2}{14. อาสวโคจฺฉก 1. เหตุทุก}
\prose{\csromanfolio{435}นรูปิํ นนเหตุํ ธมฺมํ ปฏิจฺจ อรูปี นเหตุ ธมฺโม อุปฺปชฺชติ เหตุปจฺจยา. คณิตเกน ตีณิ.}

\proseitem{56}{นโลกิยํ นเหตุํ ธมฺมํ ปฏิจฺจ โลกุตฺตโร เหตุ ธมฺโม อุปฺปชฺชติ เหตุปจฺจยา. เอกํ.}

\prose{นโลกุตฺตรํ นเหตุํ ธมฺมํ ปฏิจฺจ โลกิโย เหตุ ธมฺโม อุปฺปชฺชติ เหตุปจฺจยา. เอกํ.  เหตุยา เทฺว.}

\prose{นโลกิยํ นนเหตุํ ธมฺมํ ปฏิจฺจ โลกิโย นเหตุ ธมฺโม อุปฺปชฺชติ เหตุปจฺจยา. ตีณิ.}

\prose{นโลกุตฺตรํ นนเหตุํ ธมฺมํ ปฏิจฺจ โลกิโย นเหตุ ธมฺโม อุปฺปชฺชติ เหตุปจฺจยา. เอกํ.  เหตุยา จตฺตาริ.}

\proseitem{57}{นเกนจิ วิญฺเญยฺยํ นเหตุํ ธมฺมํ ปฏิจฺจ เกนจิ วิญฺเญยฺโย เหตุ ธมฺโม อุปฺปชฺชติ เหตุปจฺจยา.  เหตุยา นว.}

\prose{นเกนจิ นวิญฺเญยฺยํ นนเหตุํ ธมฺมํ ปฏิจฺจ เกนจิ นวิญฺเญยฺโย นเหตุ ธมฺโม อุปฺปชฺชติ เหตุปจฺจยา.  เหตุยา นว.}

\csromanmarkh{14. อาสวโคจฺฉก}\csromanmarkhii{1. เหตุทุก}\csromanheadernomark{2}{14. อาสวโคจฺฉก 1. เหตุทุก}
\proseitem{58}{โนอาสวํ นเหตุํ ธมฺมํ ปฏิจฺจ อาสโว เหตุ ธมฺโม อุปฺปชฺชติ เหตุปจฺจยา. ตีณิ.}

\prose{นโนอาสวํ นเหตุํ ธมฺมํ ปฏิจฺจ อาสโว เหตุ ธมฺโม อุปฺปชฺชติ เหตุปจฺจยา. เอกํ.}

\prose{โนอาสวํ นเหตุญฺจ นโนอาสวํ นเหตุญฺจ ธมฺมํ ปฏิจฺจ อาสโว เหตุ ธมฺโม อุปฺปชฺชติ เหตุปจฺจยา. เอกํ.  เหตุยา ปญฺจ.}

\prose{โนอาสวํ นนเหตุํ ธมฺมํ ปฏิจฺจ โนอาสโว นเหตุ ธมฺโม อุปฺปชฺชติ เหตุปจฺจยา. เอกํ.}

\prose{นโนอาสวํ นนเหตุํ ธมฺมํ ปฏิจฺจ โนอาสโว นเหตุ ธมฺโม อุปฺปชฺชติ เหตุปจฺจยา. ตีณิ.}

\csromanheader{2}{ธมฺม\-ปจฺจนียา\-นุโลม ทุกทุก\-ปฏฺฐาน\-ปาฬิ}
\prose{\csromanfolio{436}โนอาสวํ นนเหตุญฺจ นโนอาสวํ นนเหตุญฺจ ธมฺมํ ปฏิจฺจ โนอาสโว นเหตุ ธมฺโม อุปฺปชฺชติ เหตุปจฺจยา. เอกํ.  เหตุยา ปญฺจ.}

\csromanmarkh{55. มหนฺตรทุก}\csromanmarkhii{1. เหตุทุก}\csromanheadernomark{2}{55. \textbf{มหนฺตรทุก 1. เหตุทุก}}
\proseitem{59}{นสารมฺมณํ นเหตุํ ธมฺมํ ปฏิจฺจ สารมฺมโณ เหตุ ธมฺโม อุปฺปชฺชติ เหตุปจฺจยา.  เหตุยา ตีณิ.}

\prose{นอนารมฺมณํ นนเหตุํ ธมฺมํ ปฏิจฺจ สารมฺมโณ นเหตุ ธมฺโม อุปฺปชฺชติ เหตุปจฺจยา.  เหตุยา ตีณิ. (สํขิตฺตํ.)}

\csromanmarkh{100. สรณทุก}\csromanmarkhii{1. เหตุทุก}\csromanheadernomark{2}{100. \textbf{สรณทุก 1. เหตุทุก}}
\proseitem{60}{นสรณํ นเหตุํ ธมฺมํ ปฏิจฺจ อรโณ เหตุ ธมฺโม อุปฺปชฺชติ เหตุปจฺจยา.  นอรณํ นเหตุํ ธมฺมํ ปฏิจฺจ สรโณ เหตุ ธมฺโม อุปฺปชฺชติ เหตุปจฺจยา.  เหตุยา เทฺว.}

\prose{นสรณํ นนเหตุํ ธมฺมํ ปฏิจฺจ อรโณ นเหตุ ธมฺโม อุปฺปชฺชติ เหตุปจฺจยา. (1)}

\prose{นอรณํ นนเหตุํ ธมฺมํ ปฏิจฺจ อรโณ นเหตุ ธมฺโม อุปฺปชฺชติ เหตุปจฺจยา.  นอรณํ นนเหตุํ ธมฺมํ ปฏิจฺจ สรโณ นเหตุ ธมฺโม อุปฺปชฺชติ เหตุปจฺจยา.  นอรณํ นนเหตุํ ธมฺมํ ปฏิจฺจ สรโณ นเหตุ จ อรโณ นเหตุ จ ธมฺมา อุปฺปชฺชนฺติ เหตุปจฺจยา. ตีณิ. (สํขิตฺตํ.)}

\csromanmarkh{100. สรณทุก}\csromanmarkhii{2. สเหตุกทุก}\csromanheadernomark{2}{100. \textbf{สรณทุก 2. สเหตุกทุก}}
\proseitem{61}{นสรณํ นสเหตุกํ ธมฺมํ ปฏิจฺจ อรโณ สเหตุโก ธมฺโม อุปฺปชฺชติ เหตุปจฺจยา.  นอรณํ นสเหตุกํ ธมฺมํ ปฏิจฺจ สรโณ สเหตุโก ธมฺโม อุปฺปชฺชติ เหตุปจฺจยา.  เหตุยา เทฺว.}

\proseitem{62}{นสรณํ นอเหตุกํ ธมฺมํ ปฏิจฺจ อรโณ อเหตุโก ธมฺโม อุปฺปชฺชติ เหตุปจฺจยา.  นอรณํ นอเหตุกํ ธมฺมํ ปฏิจฺจ อรโณ อเหตุโก ธมฺโม อุปฺปชฺชติ เหตุปจฺจยา.  เหตุยา เทฺว.}

\vspace{\tipitakaparskip}
\csromancenter{สรณทุก 6. นเหตุ\-สเหตุกทุก 100. สรณทุก 3. เหตุ\-สมฺปยุตฺตทุก}
\proseitem{63}{\csromanfolio{437}นสรณํ นเหตุ\-สมฺปยุตฺตํ ธมฺมํ ปฏิจฺจ อรโณ เหตุสมฺปยุตฺโต ธมฺโม อุปฺปชฺชติ เหตุปจฺจยา.  นอรณํ นเหตุ\-สมฺปยุตฺตํ ธมฺมํ ปฏิจฺจ สรโณ เหตุสมฺปยุตฺโต ธมฺโม อุปฺปชฺชติ เหตุปจฺจยา.  เหตุยา เทฺว.}

\proseitem{64}{นสรณํ นเหตุ\-วิปฺปยุตฺตํ ธมฺมํ ปฏิจฺจ อรโณ เหตุวิปฺปยุตฺโต ธมฺโม อุปฺปชฺชติ เหตุปจฺจยา.  นอรณํ นเหตุ\-วิปฺปยุตฺตํ ธมฺมํ ปฏิจฺจ อรโณ เหตุวิปฺปยุตฺโต ธมฺโม อุปฺปชฺชติ เหตุปจฺจยา.  เหตุยา เทฺว.}

\csromanmarkh{100. สรณทุก}\csromanmarkhii{4. เหตุ\-สเหตุกทุก}\csromanheadernomark{2}{100. สรณทุก 4. เหตุ\-สเหตุกทุก}
\proseitem{65}{นสรณํ นเหตุญฺเจว นอเหตุกญฺจ ธมฺมํ ปฏิจฺจ อรโณ เหตุ เจว สเหตุโก จ ธมฺโม อุปฺปชฺชติ เหตุปจฺจยา.  นอรณํ นเหตุญฺเจว นอเหตุกญฺจ ธมฺมํ ปฏิจฺจ สรโณ เหตุ เจว สเหตุโก จ ธมฺโม อุปฺปชฺชติ เหตุปจฺจยา.  เหตุยา เทฺว.}

\prose{นสรณํ นอเหตุกญฺเจว นนเหตุญฺจ ธมฺมํ ปฏิจฺจ อรโณ สเหตุโก เจว น จ เหตุ ธมฺโม อุปฺปชฺชติ เหตุปจฺจยา.  นอรณํ นอเหตุกญฺเจว นนเหตุญฺจ ธมฺมํ ปฏิจฺจ สรโณ สเหตุโก เจว น จ เหตุ ธมฺโม อุปฺปชฺชติ เหตุปจฺจยา.  เหตุยา เทฺว.  (เหตุ\-เหตุ\-สมฺปยุตฺตทุกํ สํขิตฺตํ.)}

\csromanmarkh{100. สรณทุก}\csromanmarkhii{6. นเหตุ\-สเหตุกทุก}\csromanheadernomark{2}{100. สรณทุก 6. นเหตุ\-สเหตุกทุก}
\proseitem{66}{นสรณํ นเหตุํ นสเหตุกํ ธมฺมํ ปฏิจฺจ อรโณ นเหตุ สเหตุโก ธมฺโม อุปฺปชฺชติ เหตุปจฺจยา.  เหตุยา เอกํ.}

\prose{นสรณํ นเหตุํ นอเหตุกํ ธมฺมํ ปฏิจฺจ อรโณ นเหตุ อเหตุโก ธมฺโม อุปฺปชฺชติ เหตุปจฺจยา. เอกํ.}

\prose{นอรณํ นเหตุํ นอเหตุกํ ธมฺมํ ปฏิจฺจ อรโณ นเหตุ อเหตุโก ธมฺโม อุปฺปชฺชติ เหตุปจฺจยา. เอกํ.  เหตุยา เทฺว.}

\csromanmarkh{ธมฺม\-ปจฺจนียา\-นุโลม ทุกทุก\-ปฏฺฐาน\-ปาฬิ}\csromanmarkhii{100. สรณทุก 7. จูฬนฺตรทุก}\csromanheadernomark{2}{ธมฺม\-ปจฺจนียา\-นุโลม ทุกทุก\-ปฏฺฐาน\-ปาฬิ 100. สรณทุก 7. จูฬนฺตรทุก}
\proseitem{67}{\csromanfolio{438}นสรโณ นสปฺปจฺจโย ธมฺโม อรณสฺส สปฺปจฺจยสฺส ธมฺมสฺส อารมฺมณ\-ปจฺจเยน ปจฺจโย.  อารมฺมเณ เอกํ.}

\prose{68. นสรโณ นสงฺขโต ธมฺโม. (สํขิตฺตํ.)}

\proseitem{69}{นสรณํ นสนิทสฺสนํ ธมฺมํ ปฏิจฺจ อรโณ สนิทสฺสโน ธมฺโม อุปฺปชฺชติ เหตุปจฺจยา.  เหตุยา ตีณิ.}

\prose{นสรโณ นอนิทสฺสโน ธมฺโม สรณสฺส อนิทสฺสนสฺส ธมฺมสฺส อารมฺมณ\-ปจฺจเยน ปจฺจโย.  นสรโณ นอนิทสฺสโน ธมฺโม อรณสฺส อนิทสฺสนสฺส ธมฺมสฺส อารมฺมณ\-ปจฺจเยน ปจฺจโย.  อารมฺมเณ เทฺว.}

\proseitem{70}{นสรณํ นสปฺปฏิฆํ ธมฺมํ ปฏิจฺจ อรโณ สปฺปฏิโฆ ธมฺโม อุปฺปชฺชติ เหตุปจฺจยา.  เหตุยา ตีณิ.}

\proseitem{71}{นสรณํ นรูปิํ ธมฺมํ ปฏิจฺจ อรโณ รูปี ธมฺโม อุปฺปชฺชติ เหตุปจฺจยา.  นอรณํ นรูปิํ ธมฺมํ ปฏิจฺจ อรโณ รูปี ธมฺโม อุปฺปชฺชติ เหตุปจฺจยา.  เหตุยา เทฺว.}

\prose{นสรณํ นอรูปิํ ธมฺมํ ปจฺจยา สรโณ อรูปี ธมฺโม อุปฺปชฺชติ เหตุปจฺจยา.  นสรณํ นอรูปิํ ธมฺมํ ปจฺจยา อรโณ อรูปี ธมฺโม อุปฺปชฺชติ เหตุปจฺจยา.  เหตุยา เทฺว.}

\proseitem{72}{นสรณํ นโลกิยํ ธมฺมํ ปฏิจฺจ อรโณ โลกิโย ธมฺโม อุปฺปชฺชติ เหตุปจฺจยา.  เหตุยา เอกํ.}

\prose{นสรณํ นโลกุตฺตรํ ธมฺมํ ปจฺจยา อรโณ โลกุตฺตโร ธมฺโม อุปฺปชฺชติ เหตุปจฺจยา.  เหตุยา เอกํ.}

\proseitem{73}{นสรณํ นเกนจิ วิญฺเญยฺยํ ธมฺมํ ปฏิจฺจ อรโณ เกนจิ วิญฺเญยฺโย ธมฺโม อุปฺปชฺชติ เหตุปจฺจยา.  เหตุยา ปญฺจ.}

\prose{นสรณํ นเกนจิ นวิญฺเญยฺยํ ธมฺมํ ปฏิจฺจ อรโณ เกนจิ นวิญฺเญยฺโย ธมฺโม อุปฺปชฺชติ เหตุปจฺจยา.  เหตุยา ปญฺจ.}

\vspace{\tipitakaparskip}
\csromancenter{สรณทุก 55. มหนฺตรทุก 100. สรณทุก 14. อาสว\-โคจฺฉกาทิ}
\proseitem{74}{\csromanfolio{439}นอรณํ นอาสวํ ธมฺมํ ปฏิจฺจ สรโณ อาสโว ธมฺโม อุปฺปชฺชติ เหตุปจฺจยา.  เหตุยา เอกํ.}

\proseitem{75}{นอรณํ นสญฺโญชนํ ธมฺมํ ปฏิจฺจ สรโณ สญฺโญชโน ธมฺโม อุปฺปชฺชติ เหตุปจฺจยา.  เหตุยา เอกํ.}

\proseitem{76}{นอรณํ นคนฺถํ ธมฺมํ ปฏิจฺจ สรโณ คนฺโถ ธมฺโม อุปฺปชฺชติ เหตุปจฺจยา.  เหตุยา เอกํ.}

\proseitem{77}{นอรณํ นโอฆํ ธมฺมํ ปฏิจฺจ สรโณ โอโฆ ธมฺโม อุปฺปชฺชติ เหตุปจฺจยา.  เหตุยา เอกํ.}

\proseitem{78}{นอรณํ โนโยคํ ธมฺมํ ปฏิจฺจ สรโณ โยโค ธมฺโม อุปฺปชฺชติ เหตุปจฺจยา.  เหตุยา เอกํ.}

\proseitem{79}{นอรณํ นนีวรณํ ธมฺมํ ปฏิจฺจ สรโณ นีวรโณ ธมฺโม อุปฺปชฺชติ เหตุปจฺจยา.  เหตุยา เอกํ.}

\proseitem{80}{นอรณํ นปรามาสํ ธมฺมํ ปฏิจฺจ สรโณ ปรามาโส ธมฺโม อุปฺปชฺชติ เหตุปจฺจยา.  เหตุยา เอกํ.}

\csromanmarkh{100. สรณทุก}\csromanmarkhii{55. มหนฺตรทุก}\csromanheadernomark{2}{100. สรณทุก 55. มหนฺตรทุก}
\proseitem{81}{นสรณํ นสารมฺมณํ ธมฺมํ ปฏิจฺจ อรโณ สารมฺมโณ ธมฺโม อุปฺปชฺชติ เหตุปจฺจยา.  เหตุยา เอกํ.}

\prose{นสรณํ นอนารมฺมณํ ธมฺมํ ปฏิจฺจ อรโณ อนารมฺมโณ ธมฺโม อุปฺปชฺชติ เหตุปจฺจยา.  นอรณํ นอนารมฺมณํ ธมฺมํ ปฏิจฺจ อรโณ อนารมฺมโณ ธมฺโม อุปฺปชฺชติ เหตุปจฺจยา.  เหตุยา เทฺว.}

\proseitem{82}{นสรณํ นจิตฺตํ ธมฺมํ ปฏิจฺจ อรโณ จิตฺโต ธมฺโม อุปฺปชฺชติ เหตุปจฺจยา.  นอรณํ นจิตฺตํ ธมฺมํ ปฏิจฺจ สรโณ จิตฺโต ธมฺโม อุปฺปชฺชติ เหตุปจฺจยา.  เหตุยา เทฺว. (สํขิตฺตํ.)}

\csromanheader{2}{ธมฺม\-ปจฺจนียา\-นุโลม ทุกทุก\-ปฏฺฐาน\-ปาฬิ}
\proseitem{83}{\csromanfolio{440}นสรณํ นเจตสิกํ ธมฺมํ ปฏิจฺจ อรโณ เจตสิโก ธมฺโม อุปฺปชฺชติ เหตุปจฺจยา.  นอรณํ นเจตสิกํ ธมฺมํ ปฏิจฺจ สรโณ เจตสิโก ธมฺโม อุปฺปชฺชติ เหตุปจฺจยา.  เหตุยา เทฺว.}

\proseitem{84}{นสรณํ นจิตฺต\-สมฺปยุตฺตํ ธมฺมํ ปฏิจฺจ อรโณ จิตฺต\-สมฺปยุตฺโต ธมฺโม อุปฺปชฺชติ เหตุปจฺจยา.  นอรณํ นจิตฺต\-สมฺปยุตฺตํ ธมฺมํ ปฏิจฺจ สรโณ จิตฺต\-สมฺปยุตฺโต ธมฺโม อุปฺปชฺชติ เหตุปจฺจยา.  เหตุยา เทฺว.}

\proseitem{85}{นสรณํ นจิตฺต\-สํสฏฺฐํ ธมฺมํ ปฏิจฺจ อรโณ จิตฺตสํสฏฺโฐ ธมฺโม อุปฺปชฺชติ เหตุปจฺจยา.  นอรณํ นจิตฺต\-สํสฏฺฐํ ธมฺมํ ปฏิจฺจ สรโณ จิตฺตสํสฏฺโฐ ธมฺโม อุปฺปชฺชติ เหตุปจฺจยา.  เหตุยา เทฺว.}

\csromanmarkh{100. สรณทุก}\csromanmarkhii{83. ปิฏฺฐิทุก}\csromanheadernomark{2}{100. \textbf{สรณทุก 8}3. ปิฏฺฐิทุก}
\proseitem{86}{นสรณํ นทสฺสเนน ปหาตพฺพํ ธมฺมํ ปฏิจฺจ สรโณ ทสฺสเนน ปหาตพฺโพ ธมฺโม อุปฺปชฺชติ เหตุปจฺจยา.  เหตุยา เอกํ.}

\prose{นอรณํ นนทสฺสเนน ปหาตพฺพํ ธมฺมํ ปฏิจฺจ อรโณ นทสฺสเนน ปหาตพฺโพ ธมฺโม อุปฺปชฺชติ เหตุปจฺจยา.  เหตุยา เอกํ.}

\proseitem{87}{นสรณํ นสอุตฺตรํ ธมฺมํ ปฏิจฺจ อรโณ สอุตฺตโร ธมฺโม อุปฺปชฺชติ เหตุปจฺจยา.  เหตุยา เอกํ.}

\prose{(สหชาตวารมฺปิ ปจฺจยวารมฺปิ นิสฺสยวารมฺปิ สํสฏฺฐ\-วารมฺปิ สมฺปยุตฺตวารมฺปิ ปฏิจฺจ\-วาร\-สทิสํ.)}

\prose{ปญฺหาวาร เหตุอารมฺมณ}

\proseitem{88}{นสรโณ นสอุตฺตโร ธมฺโม อรณสฺส สอุตฺตรสฺส ธมฺมสฺส เหตุปจฺจเยน ปจฺจโย. เอกํ.}

\prose{นสรโณ นสอุตฺตโร ธมฺโม อรณสฺส สอุตฺตรสฺส ธมฺมสฺส อารมฺมณ\-ปจฺจเยน ปจฺจโย. เอกํ.  เหตุยา เอกํ.}

\csromanheader{2}{ปจฺจนียุ\-ทฺธาร ปจฺจนียุ\-ทฺธาร}
\proseitem{89}{\csromanfolio{441}นสรโณ นสอุตฺตโร ธมฺโม อรณสฺส สอุตฺตรสฺส ธมฺมสฺส อารมฺมณ\-ปจฺจเยน ปจฺจโย. สหชาต\-ปจฺจเยน ปจฺจโย. อุปนิสฺสย\-ปจฺจเยน ปจฺจโย. ปจฺฉาชาต\-ปจฺจเยน ปจฺจโย. (สํขิตฺตํ.) . นเหตุยา เอกํ, นอารมฺมเณ เอกํ.}

\prose{เหตุปจฺจยา นอารมฺมเณ เอกํ. (สํขิตฺตํ.)}

\prose{นเหตุปจฺจยา อารมฺมเณ เอกํ. (สํขิตฺตํ.) (ยถา กุสลตฺติเก ปญฺหาวารํ เอวํ วิตฺถาเรตพฺพํ.)}

\prose{อนุตฺตรปทเหตุอนนฺตร}

\proseitem{90}{นสรณํ นอนุตฺตรํ ธมฺมํ ปจฺจยา อรโณ อนุตฺตโร ธมฺโม อุปฺปชฺชติ เหตุปจฺจยา.  เหตุยา เอกํ ฯเปฯ อวิคเต เอกํ.}

\proseitem{91}{นสรโณ นอนุตฺตโร ธมฺโม อรณสฺส อนุตฺตรสฺส ธมฺมสฺส อนนฺตร\-ปจฺจเยน ปจฺจโย.  อนนฺตเร เอกํ, สมนนฺตเร เอกํ, อุปนิสฺสเย เทฺว, ปุเรชาเต เอกํ อาเสวเน เอกํ, วิปฺปยุตฺเต เอกํ, อตฺถิยา เอกํ, นตฺถิยา เอกํ, วิคเต เอกํ, อวิคเต เอกํ.}

\csromanheader{2}{ปจฺจนียุ\-ทฺธาร}
\proseitem{92}{นสรโณ นอนุตฺตโร ธมฺโม อรณสฺส อนุตฺตรสฺส ธมฺมสฺส อุปนิสฺสย\-ปจฺจเยน ปจฺจโย. ปุเรชาต\-ปจฺจเยน ปจฺจโย.}

\prose{นอรโณ นอนุตฺตโร ธมฺโม อรณสฺส อนุตฺตรสฺส ธมฺมสฺส อุปนิสฺสย\-ปจฺจเยน ปจฺจโย.  นเหตุยา เทฺว, นอารมฺมเณ เทฺว ฯเปฯ นอุปนิสฺสเย เอกํ, นปุเรชาเต เทฺว ฯเปฯ โนอวิคเต เทฺว.}

\prose{อุปนิสฺสย\-ปจฺจยา นเหตุยา เทฺว. (สํขิตฺตํ.)}

\csromanheader{2}{ธมฺม\-ปจฺจนียา\-นุโลม ทุกทุก\-ปฏฺฐาน\-ปาฬิ}
\prose{\csromanfolio{442}นเหตุปจฺจยา อุปนิสฺสเย เทฺว, ปุเรชาเต เอกํ ฯเปฯ อตฺถิยา เอกํ ฯเปฯ อวิคเต เอกํ. (สํขิตฺตํ.) (ยถา กุสลตฺติเก ปญฺหาวารํ เอวํ วิตฺถาเรตพฺพํ.)}

\prosewithcloser{อนุโลม\-ทุ\-กติก\-ปฏฺฐานโต ปฏฺฐาย ยาว ปริโยสานา ติํสมตฺเตหิ ภาณวาเรหิ ปฏฺฐานํ.}{ธมฺม\-ปจฺจนียา\-นุโลเม ทุกทุก\-ปฏฺฐานํ นิฏฺฐิตํ.}
\nitthitam{ปจฺจนียา\-นุโลม\-ปฏฺฐานํ นิฏฺฐิตํ.}
\nitthitam{ปฏฺฐาน\-ปกรณํ นิฏฺฐิตํ.}
